% =====================================================================
%  Eletricidade Basica, Circuitos Eletricos e Eletromagnetismo
%  Prof. Francisco Viana - 3a Edicao, Volume 1
%  Arquivo principal
%

%  Compilar com: pdflatex Livro.tex  (duas vezes, para o sumario)
% =====================================================================
\documentclass[11pt,a4paper,oneside]{article}
\usepackage[utf8]{inputenc}
\usepackage{textgreek}
\usepackage{etoolbox}
\usepackage[left=2.2cm,right=2.2cm,top=2.4cm,bottom=2.2cm,headsep=12pt]{geometry}
\usepackage{estilo-eletricidade}
\DeclareSIUnit{\gauss}{G}

\begin{document}

% ------------------------------ CAPA ---------------------------------
\begin{titlepage}
\thispagestyle{empty}
\begin{tikzpicture}[remember picture,overlay]
  % fundo
  \fill[azulprof] (current page.north west) rectangle
       ($(current page.south east)+(0,0)$);
  \fill[white] ($(current page.north west)+(0,-7.2cm)$) rectangle
       ($(current page.south east)+(0,4.6cm)$);
  % faixa decorativa: senoide + circuito estilizado
  \begin{scope}[shift={($(current page.north west)+(2.2cm,-4.4cm)$)}]
     \draw[white,line width=1.2pt,opacity=0.55]
        plot[domain=0:16,samples=200] (\x,{0.75*sin(\x*40)});
     \foreach \x in {0,4,8,12,16}
        \fill[laranja,opacity=0.9] (\x,{0.75*sin(\x*40)}) circle (3.5pt);
  \end{scope}
  % titulo
  \node[anchor=north west,text=white,align=left,inner sep=0]
       at ($(current page.north west)+(2.2cm,-1.6cm)$)
       {\sffamily\bfseries\Huge Eletricidade Básica,\\[2mm]
        Circuitos Elétricos\\[2mm] e Eletromagnetismo};
  % bloco central
  \node[anchor=north,align=center,inner sep=0]
       at ($(current page.north)+(0,-9.5cm)$)
       {\sffamily\LARGE\color{azulprof}\bfseries Coletânea de Exercícios
        Resolvidos\\[4mm]
        \sffamily\large\color{cinza} Volume 1 --- Corrente Contínua};
  % icones de circuito
  \node[anchor=north] at ($(current page.north)+(0,-13.2cm)$) {
    \begin{circuitikz}[scale=0.95,transform shape,line width=0.9pt]
      \draw (0,0) to[battery1,l_=$\EMF$] (0,2.4)
                  to[R,l=$R_1$,color=Rcol] (3.2,2.4)
                  to[R,l=$R_2$,color=Rcol] (3.2,0) -- (0,0);
      \draw (3.2,2.4) -- (5.4,2.4) to[C,l=$C$] (5.4,0) -- (3.2,0);
      \draw[->,Icol,line width=1pt] (0.8,2.75) -- (2.0,2.75)
            node[midway,above,font=\footnotesize\sffamily]{$I$};
    \end{circuitikz}};
  % rodape
  \node[anchor=south west,text=white,align=left,inner sep=0]
       at ($(current page.south west)+(2.2cm,2.0cm)$)
       {\sffamily\LARGE\bfseries Prof. Francisco Viana\\[2mm]
        \sffamily\normalsize 3\textsuperscript{a} Edição};
\end{tikzpicture}
\end{titlepage}

% -------------------------- APRESENTACAO -----------------------------
\newpage
\section*{Apresentação}
\addcontentsline{toc}{section}{Apresentação}
\thispagestyle{fancy}

Este material foi desenvolvido como ferramenta prática de apoio ao ensino e à
formação técnica de estudantes e profissionais da área elétrica, especialmente
daqueles envolvidos com \textbf{Eletricidade Básica}, \textbf{Análise de
Circuitos Elétricos} e \textbf{Eletromagnetismo}.

A proposta desta coletânea é oferecer uma base sólida e progressiva de
conhecimentos por meio da resolução de problemas contextualizados, que auxiliam
na fixação dos conceitos teóricos e no desenvolvimento do raciocínio aplicado à
eletricidade e ao magnetismo.

\begin{conceito}[Como usar este material]
\textbf{1. Níveis de dificuldade.} Cada questão traz uma marcação de nível:
\nivelmarca{1}~\emph{fixação} (aplicação direta de uma definição),
\nivelmarca{2}~\emph{aplicação} (exige interpretação e uma ou duas etapas de
cálculo), \nivelmarca{3}~\emph{aprofundamento} (várias etapas ou combinação de
conceitos) e \nivelmarca{4}~\emph{desafio} (síntese, demonstração, otimização ou
modelagem não rotineira). O Nível 4 é usado apenas quando a complexidade realmente
justifica essa classificação. Dentro de cada seção as questões estão
\textbf{ordenadas em ordem crescente de dificuldade}.

\smallskip
\textbf{2. Gabarito comentado.} Ao final de cada seção há a chave de respostas
seguida das \emph{resoluções comentadas}. Resolva primeiro; consulte depois.

\smallskip
\textbf{3. Lembretes teóricos.} As caixas azuis no início de cada seção reúnem as
fórmulas e os princípios necessários para resolver as questões daquele bloco.
\end{conceito}

\vspace{1em}
\begin{flushright}
\textit{Desejo a todos uma excelente jornada de estudos.}\\[2mm]
\textit{Prof. Francisco}
\end{flushright}

% ---------------------------- SUMARIO --------------------------------
\newpage
\renewcommand{\contentsname}{\sffamily\color{azulprof}Sumário}
\tableofcontents

% --------------------------- CAPITULOS -------------------------------
% A numeracao dos capitulos segue a estrutura final da obra (12 capitulos).
% Cada arquivo de capitulo NAO deve conter \setcounter: o controle fica aqui,
% para que capitulos ainda nao migrados nao quebrem a sequencia.
\newpage

% Cap. 1 -------------------------------------------------------------
% =====================================================================
%  CAPITULO 1 - ELETROSTATICA
%  Subsecao 1.1 - Cargas eletricas  (versao revisada)
% =====================================================================
\section{Eletrostática}

\begin{conceito}[Antes de começar]
\textbf{Quantização da carga:} $Q = n\,e$, com $e = \SI{1.6e-19}{\coulomb}$.\par
\textbf{Conservação da carga:} em um sistema isolado, a soma algébrica das cargas é constante.\par
\textbf{Contato entre esferas idênticas:} a carga total se divide igualmente,
$Q'_1 = Q'_2 = \dfrac{Q_1+Q_2}{2}$. Para esferas de raios diferentes, a divisão é
proporcional ao raio: $\dfrac{Q'_1}{R_1} = \dfrac{Q'_2}{R_2}$.\par
\textbf{Indução:} a aproximação de um corpo eletrizado separa cargas no condutor
neutro, que continua com carga total nula até ser aterrado.
\end{conceito}

\legendaniveis

\subsection{Cargas elétricas}

% ---------------------------------------------------------------
% NIVEL 1 - FIXACAO
% ---------------------------------------------------------------

\blocofixacao

\begin{questao}[1]{}
O princípio da conservação da carga elétrica estabelece que:
\begin{alternativas}
\item as cargas elétricas de mesmo sinal se repelem.
\item cargas elétricas de sinais opostos se atraem.
\item a soma algébrica das cargas elétricas é constante em um sistema eletricamente isolado.
\item a soma das cargas positivas e negativas é diferente de zero em um sistema neutro.
\item os elétrons livres se atraem.
\end{alternativas}
\resposta{(c)}
\resolucao{A conservação da carga afirma que, em um sistema isolado, cargas podem
ser transferidas de um corpo para outro, mas a soma algébrica permanece a mesma.
As alternativas (a) e (b) descrevem a \emph{lei de atração e repulsão}, e não a
conservação; (d) e (e) são falsas.}
\end{questao}

\begin{questao}[1]{}
Dois corpos de materiais diferentes, inicialmente neutros, são atritados entre si,
de forma isolada do meio. É correto afirmar que:
\begin{alternativas}
\item um fica eletrizado positivamente e o outro, negativamente.
\item um fica eletrizado negativamente e o outro permanece neutro.
\item um fica eletrizado positivamente e o outro permanece neutro.
\item ambos ficam eletrizados negativamente.
\item ambos ficam eletrizados positivamente.
\end{alternativas}
\resposta{(a)}
\resolucao{Na eletrização por atrito há apenas \emph{transferência} de elétrons de
um corpo para o outro. Pela conservação da carga, o que um perde o outro ganha:
as cargas finais têm o mesmo módulo e sinais opostos.}
\end{questao}

\begin{questao}[1]{}
Atrita-se um bastão de vidro com um pano de lã, ambos inicialmente neutros.
Pode-se afirmar que:
\begin{alternativas}
\item só a lã fica eletrizada.
\item só o bastão fica eletrizado.
\item o bastão e a lã se eletrizam com cargas de mesmo sinal.
\item o bastão e a lã se eletrizam com cargas de mesmo valor absoluto e sinais opostos.
\item nenhuma das anteriores.
\end{alternativas}
\resposta{(d)}
\resolucao{Mesma justificativa da questão anterior. Na série triboelétrica o vidro
cede elétrons à lã: o vidro fica positivo e a lã, negativa.}
\end{questao}

\begin{questao}[1]{}
A tabela apresenta parte de uma \textbf{série triboelétrica}: ao serem atritados,
o material que aparece \emph{mais acima} tende a ceder elétrons ao que aparece
mais abaixo.
\begin{table}[H]
\centering
\small\sffamily
\setlength{\arrayrulewidth}{0.5pt}
\begin{tabular}{@{}c l@{}}
\toprule
\textbf{Ordem} & \textbf{Material}\\
\midrule
1 & Pele de coelho\\
2 & Vidro\\
3 & Lã\\
4 & Seda\\
5 & Plástico (PVC)\\
6 & Teflon\\
\bottomrule
\end{tabular}
\caption{Série triboelétrica simplificada.}
\end{table}
Atritando-se um bastão de \textbf{PVC} com um pedaço de \textbf{lã}, os sinais das
cargas adquiridas pelo PVC e pela lã são, respectivamente:
\begin{alternativas}[2]
\item positivo e positivo.
\item positivo e negativo.
\item negativo e positivo.
\item negativo e negativo.
\item nulo e nulo.
\end{alternativas}
\resposta{(c)}
\resolucao{A lã (posição 3) está acima do PVC (posição 5), logo cede elétrons ao PVC.
O PVC \emph{ganha} elétrons (fica negativo) e a lã \emph{perde} elétrons (fica positiva).}
\end{questao}

\begin{questao}[1]{}
Duas cargas elétricas $Q_1$ e $Q_2$ atraem-se quando colocadas próximas uma da outra.
\begin{itensq}
\item O que se pode afirmar sobre os sinais de $Q_1$ e $Q_2$?
\item Sabe-se que $Q_1$ é repelida por uma terceira carga $Q_3$, positiva.
      Qual é o sinal de $Q_2$?
\end{itensq}
\resposta{(a) sinais opostos; (b) negativo}
\resolucao{(a) A atração ocorre entre cargas de sinais contrários (ou entre corpo
eletrizado e corpo neutro; como o enunciado fala em \emph{cargas}, os sinais são
opostos).\\
(b) $Q_1$ é repelida por $Q_3>0$, logo $Q_1>0$. Como $Q_1$ e $Q_2$ se atraem,
$Q_2<0$.}
\end{questao}

\begin{questao}[1]{}
Com relação à eletrização de um corpo, assinale \textbf{todas} as afirmativas corretas:
\begin{alternativas}
\item Um corpo eletricamente neutro que perde elétrons fica eletrizado positivamente.
\item Um corpo eletricamente neutro não possui cargas elétricas.
\item Um dos processos de eletrização consiste em retirar prótons do corpo.
\item Um corpo eletricamente neutro não pode ser atraído por um corpo eletrizado.
\item Friccionando-se dois corpos constituídos do mesmo material, um se eletriza
      positivamente e o outro, negativamente.
\end{alternativas}
\resposta{Apenas (a)}
\resolucao{(a) Correta. (b) Errada: o corpo neutro possui igual número de prótons
e elétrons, não ausência de cargas. (c) Errada: prótons estão presos ao núcleo;
nos sólidos só os elétrons se movem. (d) Errada: por indução, um corpo neutro é
atraído por um corpo eletrizado. (e) Errada: materiais iguais têm a mesma
tendência de ceder ou receber elétrons, e o atrito não produz eletrização
apreciável.}
\end{questao}

\begin{questao}[1]{}
Em uma cabine de \textbf{pintura eletrostática}, as gotículas de tinta são
eletrizadas negativamente e a peça metálica a ser pintada é ligada ao terminal
positivo da fonte. Esse arranjo é usado porque:
\begin{alternativas}
\item as gotículas se repelem entre si e são atraídas pela peça, reduzindo o desperdício.
\item as gotículas se atraem entre si, formando gotas maiores.
\item a tinta se torna condutora e adere melhor.
\item o processo elimina a necessidade de secagem.
\item a peça precisa estar eletrizada com o mesmo sinal da tinta.
\end{alternativas}
\resposta{(a)}
\resolucao{Gotículas de mesmo sinal se repelem, o que espalha o jato e evita
acúmulos; simultaneamente, todas são atraídas pela peça de sinal oposto,
inclusive contornando-a (efeito \emph{wrap-around}). O aproveitamento da tinta
passa de cerca de \SI{50}{\percent} para mais de \SI{90}{\percent}.}
\end{questao}

\begin{questao}[1]{}
Um corpo, inicialmente neutro, é eletrizado com carga $Q = \SI{48}{\micro\coulomb}$.
\dados{carga elementar $e = \SI{1.6e-19}{\coulomb}$.}
\begin{itensq}
\item O corpo ficou com falta ou com excesso de elétrons?
\item Qual é o número de elétrons retirado do corpo ou a ele fornecido?
\end{itensq}
\resposta{(a) falta de elétrons; (b) $n = \pot{3.0}{14}$ elétrons}
\resolucao{(a) A carga final é positiva, logo o corpo \emph{perdeu} elétrons
(ficou com falta).\\
(b) Da quantização da carga, $Q = n\,e$:
\[
n=\frac{Q}{e}=\frac{\SI{48e-6}{\coulomb}}{\SI{1.6e-19}{\coulomb}}
 = \pot{3.0}{14}\ \text{elétrons}.
\]}
\end{questao}

\begin{questao}[1]{}
Um bastão carregado positivamente atrai um objeto isolado, suspenso por um fio.
Sobre esse objeto, é correto afirmar que:
\begin{alternativas}
\item necessariamente possui elétrons em excesso.
\item é obrigatoriamente condutor.
\item trata-se de um isolante.
\item está carregado positivamente.
\item pode estar neutro.
\end{alternativas}
\resposta{(e)}
\resolucao{A atração é compatível com duas situações: objeto negativo, ou objeto
\emph{neutro} polarizado por indução. Como não há informação adicional, a única
afirmação segura é que ele \emph{pode} estar neutro --- por isso (a) e (d) estão
erradas. A polarização ocorre tanto em condutores quanto em isolantes, o que
elimina (b) e (c).}
\end{questao}

\blocoaplicacao

\begin{questao}[2]{}
Uma barra eletrizada negativamente é aproximada de um corpo metálico AB,
inicialmente neutro e apoiado sobre suportes isolantes, conforme a figura.

\begin{figure}[H]
\centering
\begin{tikzpicture}[scale=1,>=Stealth]
  % barra eletrizada
  \draw[fill=laranja!15,draw=laranja,thick,rounded corners=2pt]
        (-4.3,-0.3) rectangle (-2.6,0.3);
  \foreach \x in {-4.1,-3.7,-3.3,-2.9}
     \node[laranja,font=\bfseries\small] at (\x,0) {$-$};
  \node[below,font=\footnotesize\sffamily] at (-3.45,-0.45) {barra eletrizada};
  % corpo metalico
  \shade[left color=cinza!35,right color=cinza!15,draw=cinza,thick,rounded corners=6pt]
        (-1.9,-0.45) rectangle (2.4,0.45);
  \node[font=\bfseries] at (-1.55,0) {A};
  \node[font=\bfseries] at ( 2.05,0) {B};
  % suportes isolantes
  \foreach \x in {-1.2,1.7}{
    \draw[fill=cinza!25,draw=cinza] (\x-0.18,-1.1) rectangle (\x+0.18,-0.45);
    \draw[fill=cinza!45,draw=cinza] (\x-0.45,-1.25) rectangle (\x+0.45,-1.1);}
  \node[font=\footnotesize\sffamily,cinza] at (0.25,-1.5) {suportes isolantes};
  % linha de base
  \draw[cinza!60] (-2.6,-1.25) -- (2.9,-1.25);
\end{tikzpicture}
\caption{Barra negativa aproximada do condutor neutro AB.}
\label{fig:inducao-ab}
\end{figure}

Podemos afirmar que:
\begin{alternativas}
\item não haverá movimento de elétrons livres no corpo AB.
\item os elétrons livres do corpo AB deslocam-se para a extremidade A.
\item o sinal da carga que aparece em B é positivo.
\item ocorre no corpo metálico o fenômeno da indução eletrostática.
\item após a separação de cargas, a carga total do corpo é não nula.
\end{alternativas}
\resposta{(d)}
\resolucao{A barra negativa \emph{repele} os elétrons livres do metal, que migram
para a extremidade \emph{mais distante}, B. Logo B fica negativa e A, positiva
--- o que elimina (a), (b) e (c). Como não houve aterramento nem contato, a carga
total de AB continua nula, eliminando (e). O fenômeno descrito é exatamente a
indução eletrostática.}
\end{questao}

\begin{questao}[2]{}
Um corpo eletrizado \textbf{positivamente} é aproximado de um corpo metálico
neutro, cujas extremidades são X (mais próxima) e Y (mais afastada).

\begin{figure}[H]
\centering
\begin{tikzpicture}[scale=1,>=Stealth]
  \draw[fill=Vcol!12,draw=Vcol,thick,rounded corners=2pt] (-4.2,-0.3) rectangle (-2.7,0.3);
  \foreach \x in {-4.0,-3.6,-3.2,-2.9} \node[Vcol,font=\bfseries\small] at (\x,0) {$+$};
  \node[below,font=\footnotesize\sffamily] at (-3.45,-0.45) {corpo eletrizado};
  \shade[left color=cinza!35,right color=cinza!15,draw=cinza,thick,rounded corners=6pt]
        (-2.0,-0.45) rectangle (2.3,0.45);
  \node[font=\bfseries] at (-1.65,0) {X};
  \node[font=\bfseries] at ( 1.95,0) {Y};
  \foreach \x in {-1.3,1.6}{
    \draw[fill=cinza!25,draw=cinza] (\x-0.18,-1.1) rectangle (\x+0.18,-0.45);
    \draw[fill=cinza!45,draw=cinza] (\x-0.45,-1.25) rectangle (\x+0.45,-1.1);}
  \draw[cinza!60] (-2.7,-1.25) -- (2.8,-1.25);
\end{tikzpicture}
\caption{Indução eletrostática em um condutor neutro.}
\label{fig:inducao-xy}
\end{figure}

Podemos afirmar que:
\begin{alternativas}
\item não haverá movimentação de cargas negativas no corpo neutro.
\item a carga que aparece em X é positiva.
\item a carga que aparece em Y é negativa.
\item haverá força de interação elétrica entre os dois corpos.
\item todas as afirmativas acima estão erradas.
\end{alternativas}
\resposta{(d)}
\resolucao{O corpo positivo \emph{atrai} os elétrons livres: X fica negativa e Y,
positiva --- o que torna (a), (b) e (c) falsas. Como a extremidade negativa X está
\emph{mais próxima} da barra do que a positiva Y, a atração supera a repulsão e
existe força resultante de atração entre os corpos.}
\end{questao}

\begin{questao}[2]{}
Duas esferas condutoras descarregadas, $x$ e $y$, colocadas sobre suportes
isolantes, estão inicialmente em contato. Um bastão carregado positivamente é
aproximado da esfera $x$, como mostra a figura.

\begin{figure}[H]
\centering
\begin{tikzpicture}[scale=1]
  \draw[fill=Vcol!12,draw=Vcol,thick,rounded corners=2pt] (-4.6,-0.25) rectangle (-3.2,0.25);
  \foreach \x in {-4.4,-4.0,-3.6,-3.35} \node[Vcol,font=\bfseries\small] at (\x,0) {$+$};
  % esferas
  \shade[ball color=cinza!30] (-1.9,0) circle (0.75);
  \shade[ball color=cinza!30] ( 0.1,0) circle (0.75);
  \draw[cinza,thick] (-1.9,0) circle (0.75);
  \draw[cinza,thick] ( 0.1,0) circle (0.75);
  \node at (-1.9,0) {\bfseries $x$};
  \node at ( 0.1,0) {\bfseries $y$};
  % contato
  \draw[line width=2pt,cinza] (-1.15,0) -- (-0.65,0);
  % suportes
  \foreach \x in {-1.9,0.1}{
    \draw[fill=cinza!25,draw=cinza] (\x-0.12,-1.55) rectangle (\x+0.12,-0.75);
    \draw[fill=cinza!45,draw=cinza] (\x-0.42,-1.7) rectangle (\x+0.42,-1.55);}
  \draw[cinza!60] (-3.0,-1.7) -- (1.1,-1.7);
\end{tikzpicture}
\caption{Eletrização por indução com separação das esferas.}
\label{fig:inducao-xy2}
\end{figure}

Em seguida, a esfera $y$ é afastada de $x$, \textbf{mantendo-se o bastão em sua
posição}. Ao final, as cargas de $x$ e $y$ são, respectivamente:
\begin{alternativas}[2]
\item nula e positiva.
\item negativa e positiva.
\item nula e nula.
\item negativa e nula.
\item positiva e negativa.
\end{alternativas}
\resposta{(b)}
\resolucao{Com as esferas em contato, elas formam um único condutor: o bastão
positivo atrai os elétrons para $x$ e deixa $y$ com falta de elétrons.
\textbf{Separando as esferas com o bastão ainda presente}, essa separação de
cargas fica "aprisionada": $x$ permanece negativa e $y$, positiva.
(Se o bastão fosse afastado \emph{antes} da separação, ambas voltariam a ser
neutras.)}
\end{questao}

\begin{questao}[2]{}
A figura mostra um eletroscópio de folhas, inicialmente neutro. Um bastão
eletrizado negativamente é aproximado do disco, \textbf{sem tocá-lo}.

\begin{figure}[H]
\centering
\begin{tikzpicture}[scale=1]
  % bastao
  \draw[fill=laranja!15,draw=laranja,thick,rounded corners=2pt] (-2.9,1.9) rectangle (-1.3,2.3);
  \foreach \x in {-2.7,-2.35,-2.0,-1.6} \node[laranja,font=\bfseries\small] at (\x,2.1) {$-$};
  % disco
  \draw[fill=cinza!30,draw=cinza,thick] (-0.9,2.0) rectangle (0.9,2.2);
  % haste
  \draw[line width=1.6pt,cinza] (0,2.0) -- (0,0.55);
  % frasco de vidro
  \draw[azulprof!35,thick,fill=azulclaro!35]
        (-1.05,-1.5) -- (-1.05,1.35) .. controls (-1.05,1.65) and (-0.6,1.7) ..
        (-0.45,1.7) -- (0.45,1.7) .. controls (0.6,1.7) and (1.05,1.65) ..
        (1.05,1.35) -- (1.05,-1.5) -- cycle;
  \draw[fill=cinza!45,draw=cinza] (-1.25,-1.65) rectangle (1.25,-1.5);
  % folhas
  \draw[line width=1.1pt,laranja] (0,0.55) -- (-0.42,-0.75);
  \draw[line width=1.1pt,laranja] (0,0.55) -- ( 0.42,-0.75);
  \node[font=\footnotesize\sffamily,cinza,right] at (1.35,-0.2) {folhas metálicas};
  \node[font=\footnotesize\sffamily,cinza,right] at (1.35,2.1) {disco};
\end{tikzpicture}
\caption{Eletroscópio de folhas sob indução.}
\label{fig:eletroscopio}
\end{figure}

Sobre a situação descrita, é correto afirmar que:
\begin{alternativas}
\item as folhas permanecem fechadas, pois o eletroscópio continua neutro.
\item as folhas se abrem e o disco fica com carga positiva.
\item as folhas se abrem e o disco fica com carga negativa.
\item as folhas se abrem porque o eletroscópio adquiriu carga negativa total.
\item as folhas se fecham ainda mais.
\end{alternativas}
\resposta{(b)}
\resolucao{O bastão negativo repele os elétrons livres do disco, que descem para
as folhas. O disco fica \emph{positivo} e as folhas ficam \emph{negativas}; como
ambas têm o mesmo sinal, elas se repelem e se abrem. A carga total do
eletroscópio, porém, continua nula --- o que elimina (a) e (d).}
\end{questao}

\begin{questao}[2]{}
Três bolas metálicas podem ser carregadas eletricamente. Observa-se que cada uma
das três bolas \textbf{atrai uma e repele outra}. Considere as hipóteses:

\smallskip
\begin{tabular}{@{}ll@{}}
I   & Apenas uma das bolas está carregada.\\
II  & Duas das bolas estão carregadas.\\
III & As três bolas estão carregadas.\\
\end{tabular}
\smallskip

O fenômeno pode ser explicado:
\begin{alternativas}
\item somente pelas hipóteses II ou III.
\item somente pela hipótese I.
\item somente pela hipótese III.
\item somente pela hipótese II.
\item somente pelas hipóteses I ou II.
\end{alternativas}
\resposta{(c)}
\resolucao{A \emph{repulsão} só ocorre entre corpos eletrizados de mesmo sinal.
Como cada bola repele alguma outra, cada uma das três precisa estar eletrizada:
apenas a hipótese III explica o fenômeno. Um arranjo possível é $(+,+,-)$.}
\end{questao}

\begin{questao}[2]{}
Três esferas metálicas podem ser carregadas eletricamente. Aproximando-as duas a
duas, observa-se que \textbf{em todos os casos ocorre atração}. Considere:

\smallskip
\begin{tabular}{@{}ll@{}}
I   & Somente uma das esferas está carregada.\\
II  & Duas esferas estão carregadas.\\
III & As três esferas estão carregadas.\\
\end{tabular}
\smallskip

Quais hipóteses explicam o fenômeno descrito?
\begin{alternativas}
\item Apenas a hipótese I.
\item Apenas a hipótese II.
\item Apenas a hipótese III.
\item Apenas as hipóteses I e II.
\item Nenhuma das três hipóteses.
\end{alternativas}
\resposta{(d)}
\resolucao{Se as três estivessem carregadas, haveria pelo menos um par de mesmo
sinal (princípio da casa dos pombos com dois sinais e três esferas), gerando
repulsão --- III fica descartada. Com apenas uma carregada (I), ela atrai as duas
neutras por indução e as neutras se atraem mutuamente por polarização mútua.
Com duas carregadas de sinais opostos (II), elas se atraem entre si e cada uma
atrai a neutra. Logo I e II são possíveis.}
\end{questao}

\begin{questao}[2]{}
Uma estudante colocou sobre uma mesa horizontal três pêndulos eletrostáticos
idênticos, equidistantes entre si, como se cada um ocupasse o vértice de um
triângulo equilátero. As esferas metálicas dos pêndulos atraíram-se mutuamente.
A estudante pode concluir corretamente que:
\begin{alternativas}
\item as três esferas estavam eletrizadas com cargas de mesmo sinal.
\item duas esferas estavam eletrizadas com cargas de mesmo sinal e uma com sinal oposto.
\item duas esferas estavam eletrizadas com cargas de mesmo sinal e uma neutra.
\item duas esferas estavam eletrizadas com cargas de sinais opostos e uma neutra.
\item uma esfera estava eletrizada e duas, neutras.
\end{alternativas}
\resposta{(d)}
\resolucao{Cargas de mesmo sinal se repeliriam, o que elimina (a) e (c).
Em (b), as duas de mesmo sinal se repeliriam. Em (e), as duas esferas neutras
não interagiriam apreciavelmente entre si. Resta (d): as duas eletrizadas se
atraem por terem sinais opostos, e cada uma atrai a neutra por indução.}
\end{questao}

\begin{questao}[2]{}
Eletriza-se por atrito um bastão de plástico com um pedaço de papel. Em seguida,
o bastão é aproximado de um pêndulo eletrostático já eletrizado e observa-se
\textbf{repulsão}. Em qual alternativa a carga de cada elemento corresponde a
essa descrição?

\begin{table}[H]
\centering
\small\sffamily
\setlength{\arraystretch}{1.25}
\begin{tabular}{@{}c c c c@{}}
\toprule
 & \textbf{Papel} & \textbf{Bastão} & \textbf{Pêndulo}\\
\midrule
a) & positiva & positiva & positiva\\
b) & negativa & positiva & negativa\\
c) & negativa & negativa & positiva\\
d) & positiva & positiva & negativa\\
e) & positiva & negativa & negativa\\
\bottomrule
\end{tabular}
\caption{Alternativas da questão.}
\end{table}
\resposta{(e)}
\resolucao{Duas condições devem ser satisfeitas simultaneamente:
\textbf{(i)} bastão e papel, eletrizados por atrito, têm sinais \emph{opostos};
\textbf{(ii)} bastão e pêndulo se repelem, logo têm sinais \emph{iguais}.
Somente a alternativa (e) satisfaz ambas (papel $+$, bastão $-$, pêndulo $-$).
Em (a) o par bastão/papel tem sinais iguais; em (b) e (d) o bastão e o pêndulo
têm sinais opostos; em (c) o par bastão/papel tem sinais iguais.}
\end{questao}

\begin{questao}[2]{}
Duas esferas condutoras de dimensões iguais, A e B, estão eletrizadas com
$Q_A = \SI{+5}{\coulomb}$ e $Q_B = \SI{-1}{\coulomb}$, e são ligadas por um fio
condutor com uma chave C, inicialmente aberta.

\begin{figure}[H]
\centering
\begin{tikzpicture}[scale=1]
  \shade[ball color=Vcol!25] (-2.6,0) circle (0.8);
  \shade[ball color=Icol!25] ( 2.6,0) circle (0.8);
  \draw[Vcol,thick] (-2.6,0) circle (0.8);
  \draw[Icol,thick] ( 2.6,0) circle (0.8);
  \node[font=\bfseries] at (-2.6,0.18) {A};
  \node[font=\small] at (-2.6,-0.28) {$+\SI{5}{\coulomb}$};
  \node[font=\bfseries] at ( 2.6,0.18) {B};
  \node[font=\small] at ( 2.6,-0.28) {$-\SI{1}{\coulomb}$};
  % fio com chave
  \draw[line width=1.1pt,cinza] (-1.8,0) -- (-0.55,0);
  \draw[line width=1.1pt,cinza] ( 0.55,0) -- ( 1.8,0);
  \fill[cinza] (-0.55,0) circle (0.06);
  \fill[cinza] ( 0.55,0) circle (0.06);
  \draw[line width=1.1pt,cinza] (-0.55,0) -- (0.42,0.52);
  \node[font=\footnotesize\sffamily,cinza,above] at (0,0.62) {C (aberta)};
  % suportes
  \foreach \x in {-2.6,2.6}{
    \draw[fill=cinza!25,draw=cinza] (\x-0.12,-1.6) rectangle (\x+0.12,-0.8);
    \draw[fill=cinza!45,draw=cinza] (\x-0.42,-1.75) rectangle (\x+0.42,-1.6);}
  \draw[cinza!60] (-3.6,-1.75) -- (3.6,-1.75);
\end{tikzpicture}
\caption{Esferas idênticas ligadas por fio condutor com chave.}
\label{fig:esferas-chave-ab}
\end{figure}

Quando a chave é fechada, passam elétrons:
\begin{alternativas}
\item de A para B, e a nova carga de A é $\SI{+2}{\coulomb}$.
\item de A para B, e a nova carga de B é $\SI{-1}{\coulomb}$.
\item de B para A, e a nova carga de A é $\SI{+1}{\coulomb}$.
\item de B para A, e a nova carga de B é $\SI{-1}{\coulomb}$.
\item de B para A, e a nova carga de A é $\SI{+2}{\coulomb}$.
\end{alternativas}
\resposta{(e)}
\resolucao{As esferas são idênticas, então a carga total se divide igualmente:
\[
Q'_A=Q'_B=\frac{Q_A+Q_B}{2}=\frac{(+5)+(-1)}{2}=\SI{+2}{\coulomb}.
\]
A esfera A passou de $+5$ para $+2$: ela \emph{ganhou} \SI{3}{\coulomb} em
elétrons. Como os elétrons são negativos, eles se deslocam \emph{para} o corpo de
maior potencial, ou seja, de B para A.}
\end{questao}

\begin{questao}[2]{}
Duas esferas condutoras idênticas, K e L, com cargas
$Q_K = \SI{+3}{\micro\coulomb}$ e $Q_L = \SI{+1}{\micro\coulomb}$, são ligadas por
um fio condutor com uma chave C, inicialmente aberta.

\begin{figure}[H]
\centering
\begin{tikzpicture}[scale=1]
  \shade[ball color=Vcol!25] (-2.6,0) circle (0.8);
  \shade[ball color=Vcol!12] ( 2.6,0) circle (0.8);
  \draw[Vcol,thick] (-2.6,0) circle (0.8);
  \draw[Vcol,thick] ( 2.6,0) circle (0.8);
  \node[font=\bfseries] at (-2.6,0.18) {K};
  \node[font=\small] at (-2.6,-0.30) {$+\SI{3}{\micro\coulomb}$};
  \node[font=\bfseries] at ( 2.6,0.18) {L};
  \node[font=\small] at ( 2.6,-0.30) {$+\SI{1}{\micro\coulomb}$};
  \draw[line width=1.1pt,cinza] (-1.8,0) -- (-0.55,0);
  \draw[line width=1.1pt,cinza] ( 0.55,0) -- ( 1.8,0);
  \fill[cinza] (-0.55,0) circle (0.06);
  \fill[cinza] ( 0.55,0) circle (0.06);
  \draw[line width=1.1pt,cinza] (-0.55,0) -- (0.42,0.52);
  \node[font=\footnotesize\sffamily,cinza,above] at (0,0.62) {C (aberta)};
  \foreach \x in {-2.6,2.6}{
    \draw[fill=cinza!25,draw=cinza] (\x-0.12,-1.6) rectangle (\x+0.12,-0.8);
    \draw[fill=cinza!45,draw=cinza] (\x-0.42,-1.75) rectangle (\x+0.42,-1.6);}
  \draw[cinza!60] (-3.6,-1.75) -- (3.6,-1.75);
\end{tikzpicture}
\caption{Esferas idênticas com cargas de mesmo sinal.}
\label{fig:esferas-chave-kl}
\end{figure}

Quando a chave é fechada:
\begin{alternativas}
\item não ocorre fluxo de elétrons, pois ambas estão com cargas positivas.
\item passam elétrons de L para K, transportando carga de $\SI{+2}{\micro\coulomb}$.
\item passam elétrons de L para K, transportando carga de $\SI{-1}{\micro\coulomb}$.
\item passam elétrons de K para L, transportando carga de $\SI{+1}{\micro\coulomb}$.
\item passam elétrons de K para L, transportando carga de $\SI{-1}{\micro\coulomb}$.
\end{alternativas}
\resposta{(c)}
\resolucao{Carga final de cada esfera:
$Q' = \dfrac{3+1}{2} = \SI{+2}{\micro\coulomb}$.
A esfera K passou de $+3$ para $+2\,\si{\micro\coulomb}$: recebeu
$\SI{-1}{\micro\coulomb}$ (isto é, elétrons). Portanto os elétrons vão de L para K
transportando $\SI{-1}{\micro\coulomb}$. O que impulsiona o fluxo é a diferença de
\emph{potencial}, e não o sinal das cargas --- por isso (a) é falsa.}
\end{questao}

\begin{questao}[2]{}
Duas esferas metálicas idênticas estão carregadas com $5Q$ e $-3Q$.
Encostando-se uma na outra e separando-as em seguida, cada esfera ficará com
carga elétrica igual a:
\begin{alternativas}[3]
\item $8Q$
\item $5Q$
\item $3Q$
\item $2Q$
\item $Q$
\end{alternativas}
\resposta{(e)}
\resolucao{$Q' = \dfrac{5Q + (-3Q)}{2} = \dfrac{2Q}{2} = Q$ em cada esfera.}
\end{questao}

\begin{questao}[2]{}
Uma pequena esfera metálica está eletrizada com carga $\pot{8.0}{-8}\,\si{\coulomb}$.
Colocando-a em contato com outra idêntica, porém neutra, o número de elétrons que
passa de uma esfera para a outra é:
\dados{$e = \SI{1.6e-19}{\coulomb}$.}
\begin{alternativas}[3]
\item $\pot{4.0}{12}$
\item $\pot{4.0}{11}$
\item $\pot{4.0}{10}$
\item $\pot{2.5}{12}$
\item $\pot{2.5}{11}$
\end{alternativas}
\resposta{(e)}
\resolucao{Após o contato, cada esfera fica com
$Q' = \dfrac{\pot{8.0}{-8}+0}{2} = \pot{4.0}{-8}\,\si{\coulomb}$.
A carga transferida tem módulo $\Delta Q = \pot{4.0}{-8}\,\si{\coulomb}$, logo
\[
n=\frac{\Delta Q}{e}=\frac{\pot{4.0}{-8}}{\pot{1.6}{-19}}=\pot{2.5}{11}
\ \text{elétrons}.
\]
Note que a esfera eletrizada é positiva: os elétrons vão da esfera neutra para ela.}
\end{questao}

\begin{questao}[2]{}
Um estudante afirma ter medido, em um pequeno corpo eletrizado, a carga
$Q = \pot{2.4}{-19}\,\si{\coulomb}$. Sabendo que $e = \SI{1.6e-19}{\coulomb}$,
esse resultado é possível? Justifique.
\resposta{Não: a carga não é múltipla inteira de $e$.}
\resolucao{Pela quantização, toda carga deve satisfazer $Q=n\,e$ com $n$ inteiro.
Aqui,
\[
n=\frac{\pot{2.4}{-19}}{\pot{1.6}{-19}}=1{,}5,
\]
que não é inteiro. Não existe carga livre igual a $1{,}5\,e$, logo a medida está
incorreta (erro experimental ou de cálculo).}
\end{questao}

\blocoaprofundamento

\begin{questao}[3]{}
Três esferas metálicas A, B e C, idênticas e no vácuo, encontram-se inicialmente
com A carregada com $+Q$ e B e C neutras. A esfera A é colocada em contato com B
e, posteriormente, com C. Os valores finais das cargas de A, B e C são,
respectivamente:
\begin{alternativas}[2]
\item $\dfrac{Q}{3},\ \dfrac{Q}{3},\ \dfrac{Q}{3}$
\item $\dfrac{Q}{2},\ \dfrac{Q}{2},\ \dfrac{Q}{2}$
\item $\dfrac{Q}{4},\ \dfrac{Q}{4},\ \dfrac{Q}{4}$
\item $\dfrac{Q}{4},\ \dfrac{Q}{2},\ \dfrac{Q}{4}$
\item $\dfrac{Q}{2},\ \dfrac{Q}{2},\ \dfrac{Q}{4}$
\end{alternativas}
\resposta{(d)}
\resolucao{\textbf{1\textsuperscript{o} contato (A--B):}
$Q_A = Q_B = \dfrac{Q+0}{2} = \dfrac{Q}{2}$.\\
\textbf{2\textsuperscript{o} contato (A--C):}
$Q_A = Q_C = \dfrac{Q/2+0}{2} = \dfrac{Q}{4}$.\\
Situação final: $A=\dfrac{Q}{4}$, $B=\dfrac{Q}{2}$, $C=\dfrac{Q}{4}$.
O erro mais comum é esquecer que B \emph{conserva} a carga obtida no primeiro
contato.}
\end{questao}

\begin{questao}[3]{}
Quatro esferas metálicas idênticas estão isoladas entre si. As esferas A, B e C
estão neutras e a esfera D está eletrizada com carga $Q$. A esfera D é colocada
em contato sucessivamente com A, depois com B e, por fim, com C.
Ao final do processo, as cargas de C e D são, respectivamente:
\begin{alternativas}[2]
\item $\dfrac{Q}{8}$ e $\dfrac{Q}{8}$
\item $\dfrac{Q}{8}$ e $\dfrac{Q}{4}$
\item $\dfrac{Q}{4}$ e $\dfrac{Q}{8}$
\item $\dfrac{Q}{2}$ e $\dfrac{Q}{2}$
\item $Q$ e $-Q$
\end{alternativas}
\resposta{(a)}
\resolucao{A cada contato com uma esfera \emph{neutra} e idêntica, a carga de D
é dividida por 2:
\[
Q \xrightarrow{\text{A}} \frac{Q}{2}
  \xrightarrow{\text{B}} \frac{Q}{4}
  \xrightarrow{\text{C}} \frac{Q}{8}.
\]
No último contato, C fica com o mesmo valor de D: ambas com $\dfrac{Q}{8}$.
De modo geral, após $n$ contatos com esferas neutras idênticas,
$Q_D = \dfrac{Q}{2^{\,n}}$.}
\end{questao}

\begin{questao}[3]{}
Considere três esferas metálicas X, Y e Z, de diâmetros iguais. Y e Z estão fixas
e suficientemente afastadas para que efeitos de indução sejam desprezíveis. A
situação inicial é: X neutra, Y com carga $+Q$ e Z com carga $-Q$. As esferas não
trocam cargas com o ambiente. Fazendo-se X tocar \textbf{primeiro} Y e
\textbf{depois} Z, a carga final de X será:
\begin{alternativas}[2]
\item zero
\item $\dfrac{2Q}{3}$
\item $-\dfrac{Q}{2}$
\item $\dfrac{Q}{8}$
\item $-\dfrac{Q}{4}$
\end{alternativas}
\resposta{(e)}
\resolucao{\textbf{Contato X--Y:} $Q_X = Q_Y = \dfrac{0+Q}{2} = \dfrac{Q}{2}$.\\
\textbf{Contato X--Z:} agora X tem $+\dfrac{Q}{2}$ e Z tem $-Q$:
\[
Q_X=\frac{\dfrac{Q}{2}+(-Q)}{2}=\frac{-\dfrac{Q}{2}}{2}=-\frac{Q}{4}.
\]
Observe a inversão de sinal: X termina \emph{negativa}. A alternativa (a) é a
armadilha típica de quem soma as cargas iniciais de Y e Z.}
\end{questao}

\begin{questao}[3]{}
Duas esferas metálicas idênticas, A e B, estão carregadas com
$\SI{16}{\micro\coulomb}$ e $\SI{4}{\micro\coulomb}$, respectivamente. Uma terceira
esfera C, idêntica às anteriores, está inicialmente descarregada. Coloca-se C em
contato com A; desfeito o contato, C é colocada em contato com B. Supondo que não
haja troca de cargas com o exterior, a carga final de C é:
\begin{alternativas}[3]
\item $\SI{8}{\micro\coulomb}$
\item $\SI{6}{\micro\coulomb}$
\item $\SI{4}{\micro\coulomb}$
\item $\SI{3}{\micro\coulomb}$
\item nula
\end{alternativas}
\resposta{(b)}
\resolucao{\textbf{Contato C--A:}
$Q_C = Q_A = \dfrac{0+16}{2} = \SI{8}{\micro\coulomb}$.\\
\textbf{Contato C--B:}
$Q_C = Q_B = \dfrac{8+4}{2} = \SI{6}{\micro\coulomb}$.\\
Situação final: $A=\SI{8}{\micro\coulomb}$, $B=Q_C=\SI{6}{\micro\coulomb}$.
Confira a conservação: $16+4+0 = 8+6+6 = \SI{20}{\micro\coulomb}$.}
\end{questao}

\begin{questao}[3]{}
Três esferas condutoras A, B e C têm o mesmo diâmetro. A esfera A está
inicialmente neutra, e as outras duas estão carregadas com
$q_B = \SI{6}{\micro\coulomb}$ e $q_C = \SI{7}{\micro\coulomb}$. Com a esfera A,
toca-se primeiro B e depois C. As cargas de A, B e C após os contatos são,
respectivamente:
\begin{alternativas}
\item zero, zero e $\SI{13}{\micro\coulomb}$
\item $\SI{7}{\micro\coulomb}$, $\SI{3}{\micro\coulomb}$ e $\SI{5}{\micro\coulomb}$
\item $\SI{5}{\micro\coulomb}$, $\SI{3}{\micro\coulomb}$ e $\SI{5}{\micro\coulomb}$
\item $\SI{6}{\micro\coulomb}$, $\SI{7}{\micro\coulomb}$ e zero
\item todas iguais a $\SI{4.3}{\micro\coulomb}$
\end{alternativas}
\resposta{(c)}
\resolucao{\textbf{Contato A--B:}
$q_A = q_B = \dfrac{0+6}{2} = \SI{3}{\micro\coulomb}$.\\
\textbf{Contato A--C:}
$q_A = q_C = \dfrac{3+7}{2} = \SI{5}{\micro\coulomb}$.\\
Final: $A=\SI{5}{\micro\coulomb}$, $B=\SI{3}{\micro\coulomb}$,
$C=\SI{5}{\micro\coulomb}$.
Verificação: $0+6+7 = 5+3+5 = \SI{13}{\micro\coulomb}$.}
\end{questao}

\begin{questao}[3]{}
Na eletrização por contato entre duas esferas condutoras de raios $R$ e $2R$,
sabe-se que, após o contato, a razão entre a carga adquirida por cada esfera e o
seu respectivo raio é constante. Se, inicialmente, a esfera de raio $R$ possui
carga $3Q$ e a de raio $2R$ possui $2Q$, calcule a carga final de cada esfera.
\resposta{$Q_1 = \dfrac{5Q}{3}$ (raio $R$) e $Q_2 = \dfrac{10Q}{3}$ (raio $2R$)}
\resolucao{Sejam $Q_1$ e $Q_2$ as cargas finais. Aplicando as duas condições:
\[
\textbf{(i) conservação: } Q_1+Q_2 = 3Q+2Q = 5Q,
\qquad
\textbf{(ii) proporcionalidade: } \frac{Q_1}{R}=\frac{Q_2}{2R}
\ \Rightarrow\ Q_2 = 2Q_1 .
\]
Substituindo (ii) em (i): $Q_1 + 2Q_1 = 5Q \Rightarrow Q_1 = \dfrac{5Q}{3}$ e
$Q_2 = \dfrac{10Q}{3}$.
Fisicamente, a proporcionalidade decorre de as esferas ficarem com o \emph{mesmo
potencial} após o contato: $V = \dfrac{\ke Q}{R}$ igual para ambas.}
\end{questao}

\begin{questao}[3]{}
Uma esfera metálica isolada possui carga inicial $Q_0$. Ela é colocada em contato,
uma única vez e sucessivamente, com $n$ esferas idênticas a ela e inicialmente
neutras.
\begin{itensq}
\item Obtenha a expressão da carga final da esfera original em função de $Q_0$ e $n$.
\item A partir de quantos contatos a carga da esfera original cai abaixo de
      \SI{1}{\percent} de $Q_0$?
\end{itensq}
\resposta{(a) $Q_n = Q_0/2^{\,n}$; \quad (b) $n = 7$}
\resolucao{(a) Cada contato com uma esfera idêntica e neutra reparte a carga em
duas partes iguais, logo $Q_1 = Q_0/2$, $Q_2 = Q_0/4$ e, por indução,
\[
\boxed{\,Q_n = \frac{Q_0}{2^{\,n}}\,}.
\]
(b) Queremos $\dfrac{1}{2^{\,n}} < 0{,}01 \Rightarrow 2^{\,n} > 100$. Como
$2^6 = 64$ e $2^7 = 128$, são necessários $n = 7$ contatos.
Repare que a carga nunca se anula: ela apenas tende a zero, o que ilustra por
que o aterramento (contato com um condutor "infinito") é a forma prática de
descarregar um corpo.}
\end{questao}

\begin{questao}[3]{}
Um gerador de Van de Graaff transporta carga para sua cúpula a uma taxa
constante de $\SI{1}{\micro\ampere}$.
\begin{itensq}
\item Que carga é acumulada em $\SI{10}{\second}$?
\item Quantos elétrons isso representa?
\item Se a cúpula é uma esfera de raio $\SI{25}{\centi\meter}$, qual o potencial
      atingido? O ar rompe a cerca de $\SI{3}{\mega\volt\per\meter}$ --- haverá
      descarga?
\end{itensq}
\dados{$e=\SI{1.6e-19}{\coulomb}$; $\ke=\SI{9e9}{\newton\meter\squared\per\coulomb\squared}$.}
\resposta{(a) \SI{10}{\micro\coulomb}; (b) $\pot{6.25}{13}$;
(c) $V=\SI{360}{\kilo\volt}$, $E=\SI{1.44}{\mega\volt\per\meter}$ --- sem descarga}
\resolucao{(a) $Q = I\,t = \pot{1}{-6}\cdot10 = \SI{10}{\micro\coulomb}$.\\
(b) $n = \dfrac{Q}{e} = \dfrac{\pot{1}{-5}}{\pot{1.6}{-19}} = \pot{6.25}{13}$
elétrons --- \emph{removidos} da cúpula (que fica positiva).\\
(c) Para uma esfera condutora,
\[
V=\frac{\ke Q}{R}=\frac{9\cdot10^9\cdot\pot{1}{-5}}{0{,}25}=\SI{360}{\kilo\volt},
\qquad
E_{\text{sup}}=\frac{\ke Q}{R^{2}}=\frac{V}{R}=\frac{\pot{3.6}{5}}{0{,}25}
=\SI{1.44}{\mega\volt\per\meter}.
\]
Como $\SI{1.44}{} < \SI{3}{\mega\volt\per\meter}$, \textbf{não há ruptura} --- a
carga continua se acumulando.\\
\textbf{Limite do gerador:} a descarga ocorreria quando
$E_{\text{sup}} = \SI{3}{\mega\volt\per\meter}$, ou seja
$V_{\max}=E\,R = 3\cdot10^6\cdot0{,}25=\SI{750}{\kilo\volt}$.
\textbf{Consequência de projeto:} como $V_{\max} \propto R$, cúpulas
\emph{maiores} sustentam tensões maiores --- por isso os geradores de demonstração
têm esferas polidas e volumosas, sem pontas (onde $E$ dispararia).}
\end{questao}

\begin{questao}[3]{}
Um fio de cobre tem massa $\SI{1}{\gram}$. Sabendo que cada átomo de cobre
contribui com \emph{um} elétron livre:
\begin{itensq}
\item Calcule o número de elétrons livres no fio.
\item Determine a carga total desses elétrons.
\item Compare com a carga típica obtida por atrito ($\sim\SI{1}{\micro\coulomb}$)
      e interprete o resultado.
\end{itensq}
\dados{$N_A=\SI{6.02e23}{\per\mole}$; $M_{\text{Cu}}=\SI{63.5}{\gram\per\mole}$;
$e=\SI{1.6e-19}{\coulomb}$.}
\resposta{(a) $\approx\pot{9.5}{21}$; (b) $\approx\SI{1517}{\coulomb}$;
(c) razão $\sim10^9$}
\resolucao{(a) O número de mols é $n=\dfrac{1}{63{,}5}$ mol, logo
\[
N=\frac{1}{63{,}5}\cdot\pot{6.02}{23}\approx\pot{9.5}{21}\ \text{elétrons livres}.
\]
(b) $|Q| = N\,e = \pot{9.5}{21}\cdot\pot{1.6}{-19}\approx\SI{1517}{\coulomb}$.\\
(c) A razão entre essa carga e a obtida por atrito é
\[
\frac{1517}{\pot{1}{-6}}\approx\pot{1.5}{9},
\]
ou seja, cerca de \textbf{um bilhão e meio de vezes} maior.

\textbf{Interpretação --- por que isso importa.} Eletrizar um corpo por atrito
remove ou acrescenta uma fração absolutamente ínfima dos elétrons disponíveis
(da ordem de $10^{-9}$). O material permanece, para todos os efeitos,
\emph{eletricamente neutro em sua estrutura}; a "carga" que percebemos é apenas
um desequilíbrio superficial minúsculo.\\
Isso explica dois fatos:
\begin{itemize}[leftmargin=1.6em,itemsep=1pt,topsep=2pt]
\item por que um condutor pode fornecer corrente indefinidamente sem "acabar" os
      elétrons;
\item por que forças elétricas macroscópicas são tão intensas --- se
      conseguíssemos separar \emph{toda} a carga de $\SI{1}{\gram}$ de cobre e
      colocá-la a $\SI{1}{\meter}$ de distância de outra igual, a força seria da
      ordem de $10^{16}$\,N, milhões de vezes o peso de um transatlântico.
\end{itemize}}
\end{questao}

\begin{questao}[3]{}
Três esferas condutoras isoladas têm raios $R$, $2R$ e $3R$. A primeira está
carregada com $Q$ e as demais estão neutras. Sabe-se que, ao se tocarem, duas
esferas de raios diferentes repartem a carga de modo que o \emph{potencial} de
ambas fique igual, o que implica $\dfrac{Q_1}{R_1}=\dfrac{Q_2}{R_2}$.

Encosta-se a esfera $A$ (raio $R$) na esfera $B$ (raio $2R$); separa-se; em
seguida, encosta-se $A$ na esfera $C$ (raio $3R$).
\begin{itensq}
\item Determine a carga final de cada esfera.
\item Verifique a conservação da carga total.
\item Generalize: qual a carga final de $A$ após tocar sucessivamente esferas de
      raios $2R, 3R, \dots, nR$?
\end{itensq}
\resposta{(a) $Q_A=\dfrac{Q}{12}$, $Q_B=\dfrac{2Q}{3}$, $Q_C=\dfrac{Q}{4}$;
(b) soma $=Q$; (c) $Q_A^{(n)}=Q\displaystyle\prod_{j=2}^{n}\frac{1}{j+1}$}
\resolucao{\textbf{Contato A--B.} A carga total $Q$ se reparte proporcionalmente
aos raios ($R$ e $2R$), ou seja, na razão $1:2$:
\[
Q_A=\frac{R}{R+2R}\,Q=\frac{Q}{3},
\qquad
Q_B=\frac{2R}{3R}\,Q=\frac{2Q}{3}.
\]
\textbf{Contato A--C.} Agora $A$ traz $\dfrac{Q}{3}$ e encontra $C$ (raio $3R$,
neutra). A repartição é na razão $R:3R = 1:3$:
\[
Q_A=\frac{1}{4}\cdot\frac{Q}{3}=\frac{Q}{12},
\qquad
Q_C=\frac{3}{4}\cdot\frac{Q}{3}=\frac{Q}{4}.
\]
\textbf{(b) Conservação:}
$\dfrac{Q}{12}+\dfrac{2Q}{3}+\dfrac{Q}{4}
=\dfrac{Q+8Q+3Q}{12}=\dfrac{12Q}{12}=Q$. \checkmark

\textbf{(c) Generalização.} Ao tocar uma esfera de raio $jR$, a esfera $A$ retém a
fração $\dfrac{R}{R+jR}=\dfrac{1}{1+j}$ da carga que possuía. Logo, após a
sequência $2R, 3R, \dots, nR$:
\[
Q_A^{(n)} = Q\prod_{j=2}^{n}\frac{1}{1+j}
= Q\cdot\frac{1}{3}\cdot\frac{1}{4}\cdots\frac{1}{n+1}
= \frac{2\,Q}{(n+1)!}.
\]
Confira para $n=3$: $\dfrac{2Q}{4!}=\dfrac{2Q}{24}=\dfrac{Q}{12}$. \checkmark\\
\textbf{Por que isso importa:} a carga de $A$ decai \emph{fatorialmente} --- muito
mais rápido que o decaimento $2^{-n}$ das esferas idênticas. É o princípio do
\emph{gerador de Van de Graaff}: transferir carga para um condutor de raio muito
maior esvazia quase completamente o menor.}
\end{questao}

% BEGIN BANCO COMPLEMENTAR Cargas elétricas
% =====================================================================
%  Banco complementar --- Cargas Eletricas  (REVISADO)
%  Substitui a versao gerada por template (2 clones em N3 e 8 questoes
%  conceituais marcadas como N4).
%  RESTRICAO SEVERA: o cap01 e extenso (31 questoes) e ja cobre
%  praticamente todos os temas das N4 antigas: contato entre esferas
%  identicas (varias vezes), contatos sucessivos, esferas de raios
%  diferentes, eletroscopio, inducao com aterramento, verificacao de
%  quantizacao (Q = 2,4e-19 C impossivel), serie triboeletrica,
%  gerador de Van de Graaff e contagem de eletrons em fio de cobre.
%  Este banco explora o que restou: limite de eletrizacao pela rigidez
%  dieletrica, poder das pontas, ordem de grandeza do atrito, forca
%  eletrica x peso, conservacao de carga em processos nucleares,
%  gaiola/copo de Faraday e a serie geometrica dos contatos sucessivos.
%  Distribuicao mantida: 2 N3 + 8 N4 = 10
% =====================================================================

\subsubsection*{Banco complementar --- Cargas elétricas}

\begin{atencao}[Organização do banco]
Dois princípios sustentam o tópico: a carga é \textbf{conservada} (em sistema
isolado, a soma algébrica não muda) e é \textbf{quantizada} (todo valor é
múltiplo inteiro de $e=\pot{1.6}{-19}\,\si{\coulomb}$). O N4 explora as
consequências quantitativas menos óbvias --- há um teto para a carga que um
corpo pode reter, o atrito envolve uma fração ínfima dos átomos, e a
distribuição em condutores é governada pela curvatura.
\end{atencao}

\blocoaprofundamento

\begin{questao}[3]{}
Um sistema isolado contém $5$ prótons e $3$ elétrons livres. Após uma
reorganização interna, dois elétrons passam a orbitar dois dos prótons,
formando átomos neutros.
\begin{itensq}
\item Determine a carga total antes e depois.
\item Justifique o resultado.
\end{itensq}
\dados{$e=\pot{1.6}{-19}\,\si{\coulomb}$}
\resposta{$Q=+2e=\pot{3.2}{-19}\,\si{\coulomb}$ nos dois instantes}
\resolucao{(a) \textbf{Antes:} $Q=5(+e)+3(-e)=+2e=\pot{3.2}{-19}\,\si{\coulomb}$.\\
\textbf{Depois:} dois pares próton--elétron formam átomos neutros, restando
$3$ prótons livres e $1$ elétron livre:
\[
Q=3(+e)+1(-e)=+2e=\pot{3.2}{-19}\,\si{\coulomb}.
\]
(b) \textbf{Nada mudou} --- e não poderia mudar. A reorganização apenas
\emph{redistribuiu} as cargas dentro do sistema; nenhuma foi criada nem
destruída, e o sistema é isolado.\\
\textbf{Ponto conceitual:} um corpo \emph{neutro} não é um corpo \emph{sem
carga} --- ele tem tantas cargas positivas quanto negativas. Agrupar cargas em
átomos neutros esconde-as, mas não as elimina do balanço.}
\end{questao}

\begin{questao}[3]{}
Um bastão isolante adquire, por atrito, carga de $\SI{-5}{\nano\coulomb}$.
\begin{itensq}
\item Determine o número de elétrons transferidos.
\item Determine a massa correspondente e compare com a massa do bastão
      ($\SI{20}{\gram}$).
\end{itensq}
\dados{$e=\pot{1.6}{-19}\,\si{\coulomb}$;
$m_e=\pot{9.11}{-31}\,\si{\kilogram}$}
\resposta{$n=\pot{3.1}{10}$ elétrons; $\Delta m=\pot{2.8}{-20}\,\si{\kilogram}$}
\resolucao{(a)
\[
n=\frac{|Q|}{e}=\frac{\pot{5}{-9}}{\pot{1.6}{-19}}=\pot{3.125}{10}.
\]
Trinta bilhões de elétrons --- número enorme em termos absolutos.\\
(b)
\[
\Delta m=n\,m_e=\pot{3.125}{10}\times\pot{9.11}{-31}
=\pot{2.85}{-20}\,\si{\kilogram}.
\]
Comparando com os $\SI{20}{\gram}$ do bastão:
\[
\frac{\Delta m}{m}=\frac{\pot{2.85}{-20}}{\pot{2}{-2}}=\pot{1.4}{-18}.
\]
\textbf{Conclusão:} a eletrização por atrito \emph{não} altera a massa de forma
mensurável --- o desvio está $15$ ordens de grandeza abaixo da melhor balança
analítica. É por isso que a massa não serve como indicador de eletrização, e
por que o fenômeno passou séculos sendo descrito apenas por seus efeitos de
força.}
\end{questao}

\blocodesafio

\begin{questao}[4]{}
\textbf{(Série geométrica dos contatos.)} Uma esfera $A$ com carga $Q$ é
tocada sucessivamente por esferas idênticas \emph{neutras}
$B_1,B_2,\dots,B_n$, uma de cada vez, sendo cada uma retirada após o contato.
\begin{itensq}
\item Determine a carga de $A$ após $n$ contatos.
\item Determine a carga de cada esfera $B_k$.
\item Verifique a conservação da carga total.
\end{itensq}
\resposta{$Q_A=Q/2^{n}$; $Q_{B_k}=Q/2^{k}$; a soma reconstitui $Q$}
\resolucao{(a) Cada contato entre esferas \emph{idênticas} reparte igualmente a
carga presente. Como cada $B_k$ chega neutra, $A$ perde metade do que tem:
\[
Q\to\frac{Q}{2}\to\frac{Q}{4}\to\cdots
\;\Longrightarrow\;
Q_A=\frac{Q}{2^{n}}.
\]
(b) A esfera $B_k$ leva metade do que $A$ possuía imediatamente antes, ou seja
$\dfrac{1}{2}\cdot\dfrac{Q}{2^{k-1}}=\dfrac{Q}{2^{k}}$.\\
(c) \textbf{Conservação:}
\[
Q_A+\sum_{k=1}^{n}Q_{B_k}
=\frac{Q}{2^{n}}+Q\sum_{k=1}^{n}\frac{1}{2^{k}}
=\frac{Q}{2^{n}}+Q\left(1-\frac{1}{2^{n}}\right)=Q.\ \checkmark
\]
A soma geométrica $\sum 1/2^{k}$ tende a $1$, e o termo residual de $A$ fecha
exatamente a diferença.\\
\textbf{Leitura:} nenhuma carga se perde --- ela apenas se dilui. Com
$n=10$, $A$ retém menos de $0{,}1\%$ de $Q$, mas os $99{,}9\%$ restantes estão
distribuídos pelas dez esferas retiradas.\\
\textbf{Uso prático:} é assim que se \emph{descarrega} progressivamente um corpo
sem aterrá-lo, e também como se obtêm frações controladas de uma carga de
referência.}
\end{questao}

\begin{questao}[4]{}
\textbf{(Limite de eletrização.)} Uma esfera condutora de raio
$R=\SI{10}{\centi\meter}$ é carregada no ar, cuja rigidez dielétrica é
$E_{rup}=\SI{3e6}{\volt\per\meter}$.
\begin{itensq}
\item Determine a carga máxima que ela pode reter.
\item Determine o potencial correspondente.
\item Explique o que ocorre se tentarmos ultrapassar esse valor.
\end{itensq}
\dados{$\ke=\SI{9e9}{\newton\meter\squared\per\coulomb\squared}$}
\resposta{$Q_{\max}=\SI{3.33}{\micro\coulomb}$; $V=\SI{300}{\kilo\volt}$}
\resolucao{(a) O campo na superfície de uma esfera condutora vale
$E=\ke Q/R^{2}$. Impondo $E=E_{rup}$:
\[
Q_{\max}=\frac{E_{rup}R^{2}}{\ke}=\frac{\pot{3}{6}\times0{,}01}{\pot{9}{9}}
=\pot{3.33}{-6}\,\si{\coulomb}=\SI{3.33}{\micro\coulomb}.
\]
(b)
\[
V=\frac{\ke Q_{\max}}{R}=E_{rup}\,R=\pot{3}{6}(0{,}10)=\SI{300}{\kilo\volt}.
\]
\textbf{Relação notável:} $V_{\max}=E_{rup}R$ --- o potencial máximo é
\emph{proporcional ao raio}. Dobrar o raio dobra a tensão suportável, e é por
isso que geradores de Van de Graaff têm cúpulas grandes e lisas.\\
(c) \textbf{Acima do limite}, o ar ioniza e passa a conduzir: a carga escapa por
descarga corona (o brilho azulado e o chiado característicos) ou por faísca.
A esfera \emph{não consegue} reter mais carga --- ela se autolimita.\\
\textbf{Consequência de projeto:} para armazenar mais carga é preciso aumentar
$R$, usar meio de maior rigidez (óleo, $\mathrm{SF_6}$, vácuo) ou eliminar
irregularidades da superfície, onde o campo local é muito maior --- assunto da
questão seguinte.}
\end{questao}

\begin{questao}[4]{}
\textbf{(Poder das pontas.)} Duas esferas condutoras de raios
$R=\SI{10}{\centi\meter}$ e $r=\SI{1}{\centi\meter}$ são ligadas por um fio
longo e carregadas até $V=\SI{30}{\kilo\volt}$.
\begin{itensq}
\item Determine a razão entre as cargas e entre as densidades superficiais.
\item Determine o campo na superfície de cada uma.
\item Explique o efeito e cite duas aplicações.
\end{itensq}
\resposta{$Q_R/Q_r=10$; $\sigma_r/\sigma_R=10$;
$E_R=\SI{300}{\kilo\volt\per\meter}$ e $E_r=\SI{3}{\mega\volt\per\meter}$}
\resolucao{(a) Ligadas por fio, as duas ficam no \emph{mesmo potencial}:
\[
V=\frac{\ke Q}{R}\Rightarrow Q\propto R
\;\Longrightarrow\;
\frac{Q_R}{Q_r}=\frac{R}{r}=10.
\]
A densidade superficial é $\sigma=Q/(4\pi R^{2})\propto R/R^{2}=1/R$:
\[
\frac{\sigma_r}{\sigma_R}=\frac{R}{r}=10.
\]
\textbf{A esfera maior tem mais carga, mas a menor tem carga mais
concentrada.}\\
(b) $E=V/R$ na superfície de cada uma:
\[
E_R=\frac{\pot{30}{3}}{0{,}10}=\SI{300}{\kilo\volt\per\meter};
\qquad
E_r=\frac{\pot{30}{3}}{0{,}01}=\SI{3}{\mega\volt\per\meter}.
\]
(c) \textbf{O campo na esfera pequena atingiu exatamente a rigidez do ar.} A
ionização começa \emph{ali}, embora as duas estejam ao mesmo potencial ---
regiões de menor raio de curvatura concentram campo. É o \textbf{poder das
pontas}.\\
\textbf{Aplicações:}
\begin{itemize}
\item \textbf{Para-raios:} a ponta afiada ioniza o ar e estabelece o caminho
preferencial da descarga, protegendo a estrutura.
\item \textbf{Impressão e pintura eletrostática:} eletrodos pontiagudos ionizam
o ar para carregar tinta ou papel.
\end{itemize}
\textbf{E a contrapartida:} em equipamentos de alta tensão evita-se qualquer
ponta ou aresta viva --- daí os terminais arredondados e as blindagens
toroidais.}
\end{questao}

\begin{questao}[4]{}
\textbf{(Ordem de grandeza do atrito.)} Um bastão de vidro de $\SI{20}{\gram}$
adquire $\SI{5}{\nano\coulomb}$ por atrito. A massa molar média do vidro é
cerca de $\SI{60}{\gram\per\mole}$.
\begin{itensq}
\item Estime o número de átomos do bastão.
\item Determine a fração de átomos que perdeu um elétron.
\item Interprete o resultado.
\end{itensq}
\dados{$N_A=\pot{6.02}{23}\,\si{\per\mole}$}
\resposta{$\approx\pot{2}{23}$ átomos; fração de $\pot{1.6}{-13}$}
\resolucao{(a)
\[
N=\frac{m}{M}N_A=\frac{20}{60}\times\pot{6.02}{23}\approx\pot{2.0}{23}.
\]
(b) Já se sabe que $\pot{3.1}{10}$ elétrons foram transferidos:
\[
f=\frac{\pot{3.1}{10}}{\pot{2.0}{23}}=\pot{1.6}{-13},
\]
ou cerca de \textbf{um átomo em seis trilhões}.\\
(c) \textbf{Duas leituras que se complementam.}\\
\emph{Do ponto de vista microscópico}, a eletrização por atrito é um fenômeno
de superfície absolutamente marginal --- envolve uma fração desprezível da
matéria, e apenas a camada atômica mais externa.\\
\emph{Do ponto de vista macroscópico}, esses mesmos
$\pot{3.1}{10}$ elétrons produzem forças facilmente visíveis, capazes de
levantar pedaços de papel contra o peso.\\
\textbf{O contraste explica a intensidade da força elétrica:} ela é tão maior
que a gravitacional que um desequilíbrio de uma parte em $10^{13}$ já produz
efeitos observáveis. E explica também por que a matéria comum é neutra com
altíssima precisão --- qualquer desequilíbrio maior seria violentamente
desfeito.}
\end{questao}

\begin{questao}[4]{}
\textbf{(Força elétrica contra peso.)} Duas esferas idênticas de
$\SI{1}{\gram}$ estão a $\SI{5}{\centi\meter}$ uma da outra, cada uma com
carga $q$.
\begin{itensq}
\item Determine a força para $q=\SI{5}{\nano\coulomb}$ e para
      $q=\SI{50}{\nano\coulomb}$.
\item Determine a carga necessária para que a força iguale o peso.
\item Comente a viabilidade.
\end{itensq}
\dados{$g=\SI{9.8}{\meter\per\second\squared}$}
\resposta{$\SI{90}{\micro\newton}$ e $\SI{9}{\milli\newton}$;
$q\approx\SI{52}{\nano\coulomb}$}
\resolucao{(a)
\[
F=\frac{\ke q^{2}}{d^{2}}:\quad
q=\SI{5}{\nano\coulomb}\Rightarrow\SI{9.0e-5}{\newton};
\qquad
q=\SI{50}{\nano\coulomb}\Rightarrow\SI{9.0e-3}{\newton}.
\]
O peso vale $P=mg=\pot{1}{-3}(9{,}8)=\SI{9.8}{\milli\newton}$.\\
(b) Impondo $F=P$:
\[
q=d\sqrt{\frac{P}{\ke}}=0{,}05\sqrt{\frac{\pot{9.8}{-3}}{\pot{9}{9}}}
=0{,}05\times\pot{1.04}{-6}=\SI{52}{\nano\coulomb}.
\]
(c) \textbf{Perfeitamente viável.} Cinquenta nanocoulombs correspondem a
$\pot{3.3}{11}$ elétrons --- ordem de grandeza obtida rotineiramente por
atrito, como mostrou a questão anterior.\\
\textbf{É por isso que a eletrostática é observável a olho nu:} um simples
canudo atritado sustenta pedaços de papel contra a gravidade de todo o planeta.
\\
\textbf{Verificação do limite:} essas esferas suportariam essa carga? Para
$R=\SI{5}{\milli\meter}$, o campo na superfície seria
$\ke q/R^{2}=\SI{1.9}{\mega\volt\per\meter}$ --- abaixo da rigidez do ar, mas
já na mesma ordem. Cargas muito maiores escapariam por descarga.}
\end{questao}

\begin{questao}[4]{}
\textbf{(Conservação em processos nucleares.)} Verifique a conservação da carga
nos processos abaixo.
\begin{itensq}
\item Decaimento beta: $n\to p+e^{-}+\bar\nu_e$.
\item Aniquilação: $e^{-}+e^{+}\to2\gamma$.
\item Discuta o alcance do princípio.
\end{itensq}
\resposta{Conservada em ambos: $0=+e-e+0$ e $-e+e=0$}
\resolucao{(a) \textbf{Decaimento beta.} Antes: nêutron, carga $0$. Depois:
próton ($+e$), elétron ($-e$) e antineutrino ($0$):
\[
0=(+e)+(-e)+0=0.\ \checkmark
\]
Note que \emph{o número de partículas mudou} e que partículas foram
\emph{criadas} --- o elétron não estava dentro do nêutron. Ainda assim a carga
total se conserva.\\
(b) \textbf{Aniquilação.} Antes: $-e+(+e)=0$. Depois: dois fótons, ambos de
carga nula. \checkmark\\
Aqui a \emph{matéria} desaparece por completo, convertida em radiação --- e a
carga continua conservada, porque já era nula.\\
(c) \textbf{Alcance do princípio.} A conservação da carga é uma das leis mais
robustas da física:
\begin{itemize}
\item Vale em processos \textbf{químicos}, \textbf{nucleares} e de
\textbf{física de partículas}, sem exceção conhecida.
\item Vale mesmo quando partículas são criadas ou destruídas --- o que a
distingue da conservação do número de partículas, que \emph{não} é uma lei.
\item Experimentalmente, o decaimento do elétron (que violaria a conservação) foi
buscado e não observado; o limite atual para sua vida média é superior a
$10^{26}$ anos.
\end{itemize}
\textbf{Origem teórica:} a conservação da carga decorre de uma simetria
fundamental do eletromagnetismo, do mesmo modo que a conservação da energia
decorre da invariância no tempo. Ela não é um fato empírico isolado, mas
consequência da estrutura da teoria.}
\end{questao}

\begin{questao}[4]{}
\textbf{(Copo de Faraday.)} Um corpo carregado é introduzido, sem tocar as
paredes, no interior de um recipiente metálico fechado ligado a um eletrômetro.
\begin{itensq}
\item Explique o que o instrumento indica e por quê.
\item Explique por que a leitura não depende da posição do corpo dentro do
      recipiente.
\item Descreva o que ocorre se o corpo tocar a parede interna.
\end{itensq}
\resposta{O eletrômetro indica exatamente a carga do corpo, independentemente da
posição}
\resolucao{(a) \textbf{Indução total.} O corpo de carga $+Q$ induz $-Q$ na
superfície \emph{interna} do recipiente --- exatamente $-Q$, porque todas as
linhas de campo que saem do corpo terminam nela. Por conservação, sobra
$+Q$ na superfície \emph{externa}, e é essa carga que o eletrômetro mede.\\
\textbf{Resultado:} a leitura é $+Q$, o valor exato da carga introduzida.\\
(b) \textbf{Independência da posição.} O campo elétrico dentro do metal é nulo
em equilíbrio. Aplicando a lei de Gauss a uma superfície inteiramente contida no
metal, o fluxo é zero --- logo a carga interna total (corpo mais superfície
interna) é nula, \emph{qualquer que seja a posição do corpo}. A carga externa
fica sempre $+Q$, distribuída de acordo apenas com a geometria \emph{externa}
do recipiente.\\
\textbf{Consequência prática:} não é preciso centralizar o corpo nem controlar
sua orientação --- o que torna o copo de Faraday o método padrão de medição
absoluta de carga.\\
(c) \textbf{Se o corpo tocar a parede:} a carga $+Q$ neutraliza a
$-Q$ induzida, o corpo sai neutro e o recipiente fica com $+Q$ na superfície
externa. \textbf{A leitura do eletrômetro não muda} --- continua $+Q$.\\
Isso permite \emph{transferir} integralmente a carga de um corpo para o
recipiente, e é a base do método clássico de acumulação sucessiva de carga em
geradores eletrostáticos.}
\end{questao}

\begin{questao}[4]{}
\textbf{(Blindagem eletrostática.)} Uma malha metálica aterrada envolve um
instrumento sensível.
\begin{itensq}
\item Explique por que o campo externo não penetra.
\item Determine se a blindagem funciona também no sentido inverso.
\item Discuta a diferença entre malha aterrada e não aterrada.
\end{itensq}
\resposta{Blinda o interior contra campos externos; o aterramento é necessário
para blindar o exterior contra fontes internas}
\resolucao{(a) \textbf{Campo externo não penetra.} As cargas livres do condutor
se redistribuem até anular o campo em seu interior --- é a condição de
equilíbrio eletrostático. Qualquer campo externo induz uma distribuição
superficial cujo campo cancela exatamente o campo aplicado na região interna.\\
Isso vale \emph{mesmo sem aterramento}: a gaiola de Faraday isolada blinda o
interior. É o efeito que protege passageiros de um carro atingido por raio.\\
(b) \textbf{Sentido inverso: só com aterramento.} Uma carga $+Q$ colocada
\emph{dentro} da gaiola induz $-Q$ na face interna e $+Q$ na externa --- e essa
carga externa produz campo do lado de fora. \textbf{A blindagem falha.}\\
Com a malha \textbf{aterrada}, a carga $+Q$ da face externa escoa para a terra,
e o campo externo desaparece. Só então a blindagem funciona nos dois sentidos.\\
(c) \textbf{Resumo da diferença.}
\begin{center}
\begin{tabular}{lcc}
\hline
& não aterrada & aterrada\\
\hline
Protege o interior de campos externos & sim & sim\\
Impede que fontes internas afetem o exterior & não & sim\\
\hline
\end{tabular}
\end{center}
\textbf{Limitações práticas:} a malha só é eficaz para aberturas muito menores
que o comprimento de onda do sinal a blindar --- em eletrostática (campo
constante), qualquer malha fina basta; em radiofrequência, a exigência é
severa. E a impedância da conexão de terra importa: um fio longo de aterramento
tem indutância e degrada a blindagem em altas frequências.}
\end{questao}

% END BANCO COMPLEMENTAR

\gabarito[1.1 Cargas elétricas]

% ---------------------------------------------------------------
% Subsecoes seguintes do Capitulo 1
% =====================================================================
%  CAPITULO 1 - continuacao
%  1.2 Lei de Coulomb  |  1.3 Campo eletrico
% =====================================================================

% =====================================================================
\subsection{Lei de Coulomb}
% =====================================================================

\begin{conceito}[Antes de começar]
\textbf{Lei de Coulomb} --- o módulo da força entre duas cargas puntiformes é
\[
F = \ke\,\frac{|q_1||q_2|}{d^2},\qquad
\ke = \SI{9e9}{\newton\meter\squared\per\coulomb\squared}\ \text{(no vácuo)}.
\]
A força é \emph{atrativa} para cargas de sinais opostos e \emph{repulsiva} para
sinais iguais, sempre na direção da reta que une as cargas.\par
\textbf{Princípio da superposição:} a força resultante sobre uma carga é a
\emph{soma vetorial} das forças que cada uma das demais exerce sobre ela.\par
\textbf{Dependências úteis:} $F \propto \dfrac{1}{d^2}$ (dobrar a distância divide
a força por 4) e $F \propto q_1 q_2$.
\end{conceito}

\legendaniveis

\blocofixacao

\begin{questao}[1]{}
Duas cargas $q_1 = \SI{+2}{\micro\coulomb}$ e $q_2 = \SI{-3}{\micro\coulomb}$
estão separadas por $\SI{4}{\centi\meter}$ no vácuo. Determine o módulo e o tipo
(atração ou repulsão) da força entre elas.
\dados{$\ke = \SI{9e9}{\newton\meter\squared\per\coulomb\squared}$.}
\resposta{$F = \SI{33.75}{\newton}$, atrativa}
\resolucao{Com $d = \SI{4}{\centi\meter} = \SI{0.04}{\meter}$:
\[
F=\ke\frac{|q_1||q_2|}{d^2}
=9\cdot10^9\cdot\frac{(\pot{2}{-6})(\pot{3}{-6})}{(0{,}04)^2}
=\SI{33.75}{\newton}.
\]
Como as cargas têm sinais opostos, a força é \textbf{atrativa}.}
\end{questao}

\begin{questao}[1]{}
Duas cargas $q_1 = q_2 = \SI{+5}{\micro\coulomb}$ estão a
$\SI{2}{\centi\meter}$ no vácuo. Determine o módulo e o tipo da força.
\dados{$\ke = \SI{9e9}{\newton\meter\squared\per\coulomb\squared}$.}
\resposta{$F = \SI{562.5}{\newton}$, repulsiva}
\resolucao{$F = 9\cdot10^9\cdot\dfrac{(\pot{5}{-6})^2}{(0{,}02)^2}
=\SI{562.5}{\newton}$, \textbf{repulsiva} (sinais iguais).}
\end{questao}

\begin{questao}[1]{}
Duas cargas puntiformes se repelem com força de intensidade
$\pot{42}{-5}\,\si{\newton}$. Reduzindo-se a distância entre elas à
\textbf{metade}, a nova força de repulsão será:
\begin{alternativas}[2]
\item $\pot{10.5}{-5}\,\si{\newton}$
\item $\pot{21}{-5}\,\si{\newton}$
\item $\pot{84}{-5}\,\si{\newton}$
\item $\pot{168}{-5}\,\si{\newton}$
\item inalterada
\end{alternativas}
\resposta{(d)}
\resolucao{Como $F \propto \dfrac{1}{d^2}$, reduzir $d$ à metade \emph{multiplica}
a força por 4:
\[
F' = 4F = 4\cdot\pot{42}{-5} = \pot{168}{-5}\,\si{\newton}.
\]}
\end{questao}

\blocoaplicacao

\begin{questao}[2]{}
A força elétrica entre duas partículas de cargas $q$ e $q/2$, separadas por uma
distância $d$ no vácuo, é $F$. A força entre duas partículas de cargas $q$ e $2q$,
separadas por $d/2$, também no vácuo, é:
\begin{alternativas}[3]
\item $F$
\item $2F$
\item $4F$
\item $8F$
\item $16F$
\end{alternativas}
\resposta{(e)}
\resolucao{Situação inicial: $F = \ke\dfrac{q\cdot(q/2)}{d^2} = \dfrac{\ke q^2}{2d^2}$.\\
Situação nova: $F' = \ke\dfrac{q\cdot 2q}{(d/2)^2} = \dfrac{2\ke q^2}{d^2/4}
= \dfrac{8\ke q^2}{d^2}$.\\
Dividindo: $\dfrac{F'}{F} = \dfrac{8\ke q^2/d^2}{\ke q^2/(2d^2)} = 16$, logo
$F' = 16F$. O produto das cargas quadruplicou e a distância caiu à metade
(fator 4 na força): $4\times4 = 16$.}
\end{questao}

\begin{questao}[2]{}
Se a força entre dois objetos carregados se mantém \emph{constante} mesmo quando a
carga de cada um é reduzida à metade, então a distância entre eles:
\begin{alternativas}
\item foi quadruplicada.
\item foi duplicada.
\item foi reduzida à quarta parte.
\item foi reduzida à metade.
\item permaneceu constante.
\end{alternativas}
\resposta{(d)}
\resolucao{Reduzir cada carga à metade divide o produto $q_1 q_2$ por 4, o que
sozinho reduziria $F$ a $1/4$. Para $F$ permanecer constante, o fator $1/d^2$
precisa \emph{quadruplicar}, ou seja, $d^2$ deve cair a $1/4$, o que significa $d$
reduzido à \textbf{metade}:
\[
F=\ke\frac{(q_1/2)(q_2/2)}{d'^2}=\ke\frac{q_1q_2}{4d'^2}
\stackrel{!}{=}\ke\frac{q_1q_2}{d^2}
\ \Rightarrow\ d'^2=\frac{d^2}{4}\ \Rightarrow\ d'=\frac{d}{2}.
\]}
\end{questao}

\begin{questao}[2]{}
Três cargas iguais de $\SI{+1}{\micro\coulomb}$ ocupam os vértices de um triângulo
equilátero de lado $\SI{5}{\centi\meter}$. Determine o módulo da força resultante
sobre uma delas.

\circuitoq{%
\begin{tikzpicture}[scale=1,line width=0.8pt,>=Stealth]
 \coordinate (A) at (0,0);
 \coordinate (B) at (3,0);
 \coordinate (C) at (1.5,2.6);
 \foreach \p/\n in {A/A,B/B,C/C}{
   \shade[ball color=Vcol!25] (\p) circle (0.28);
   \draw[Vcol] (\p) circle (0.28);
   \node at (\p) {\footnotesize$+q$};}
 \draw[dashed,cinza] (A)--(B)--(C)--cycle;
 \node[below] at (1.5,-0.35) {$\SI{5}{\centi\meter}$};
\end{tikzpicture}}{Três cargas iguais nos vértices de um triângulo equilátero.}

\resposta{$F_R \approx \SI{6.24}{\newton}$}
\resolucao{Cada par de cargas exerce uma força de módulo
\[
F=\ke\frac{q^2}{d^2}=9\cdot10^9\cdot\frac{(\pot{1}{-6})^2}{(0{,}05)^2}
=\SI{3.6}{\newton}.
\]
Sobre a carga escolhida atuam \emph{duas} forças iguais de \SI{3.6}{\newton}, com
$\SI{60}{\degree}$ entre elas (geometria do equilátero). A resultante é
\[
F_R = \sqrt{F^2+F^2+2F^2\cos 60^\circ}=F\sqrt{3}
=3{,}6\sqrt{3}\approx\SI{6.24}{\newton},
\]
dirigida para fora do triângulo, ao longo da bissetriz do vértice.}
\end{questao}

\begin{questao}[2]{}
Quatro cargas positivas iguais de $\SI{1}{\micro\coulomb}$ ocupam os vértices de
um quadrado de lado $\SI{10}{\centi\meter}$. Determine o módulo da força
resultante sobre uma delas.
\dados{$\ke = \SI{9e9}{\newton\meter\squared\per\coulomb\squared}$; $\sqrt2\approx1{,}41$.}

\circuitoq{%
\begin{tikzpicture}[scale=1,line width=0.8pt,>=Stealth]
 \foreach \p in {(0,0),(2.4,0),(2.4,2.4),(0,2.4)}{
   \shade[ball color=Vcol!25] \p circle (0.22);}
 \node at (0,0){\footnotesize$+q$};\node at (2.4,0){\footnotesize$+q$};
 \node at (2.4,2.4){\footnotesize$+q$};\node at (0,2.4){\footnotesize$+q$};
 \draw[dashed,cinza] (0,0)--(2.4,0)--(2.4,2.4)--(0,2.4)--cycle;
 \draw[dashed,cinza] (0,0)--(2.4,2.4);
 \node[below] at (1.2,-0.3){$\SI{10}{\centi\meter}$};
\end{tikzpicture}}{Quatro cargas iguais nos vértices de um quadrado.}

\resposta{$F_R \approx \SI{1.72}{\newton}$}
\resolucao{Sobre a carga de um vértice atuam três forças repulsivas:
duas dos vértices \emph{adjacentes} (ao longo dos lados, perpendiculares entre si)
e uma do vértice \emph{oposto} (ao longo da diagonal).\\
\textbf{Força de um lado} ($d = \SI{10}{\centi\meter}$):
$F_L = \ke\dfrac{q^2}{d^2} = 9\cdot10^9\dfrac{(\pot{1}{-6})^2}{(0{,}1)^2}
= \SI{0.9}{\newton}$.\\
\textbf{Força da diagonal} ($d' = d\sqrt2$):
$F_D = \ke\dfrac{q^2}{(d\sqrt2)^2} = \dfrac{F_L}{2} = \SI{0.45}{\newton}$.\\
As duas forças de lado, perpendiculares, têm resultante
$F_L\sqrt2 \approx \SI{1.27}{\newton}$, na direção da diagonal --- \emph{mesma}
direção de $F_D$. Somando:
\[
F_R = F_L\sqrt2 + F_D \approx 1{,}27 + 0{,}45 = \SI{1.72}{\newton}.
\]}
\end{questao}

\begin{questao}[2]{}
Uma pequena esfera de massa $\SI{2}{\gram}$, eletrizada com carga $q$, permanece
em equilíbrio, suspensa, sob a ação combinada do peso e de um campo elétrico
uniforme vertical de intensidade $\SI{e5}{\newton\per\coulomb}$, dirigido para
cima. Determine a carga $q$.
\dados{$g = \SI{10}{\meter\per\second\squared}$.}
\resposta{$q = \SI{0.2}{\micro\coulomb}$}
\resolucao{No equilíbrio, a força elétrica (para cima) equilibra o peso (para
baixo): $qE = mg$, logo
\[
q=\frac{mg}{E}=\frac{\pot{2}{-3}\cdot10}{\pot{1}{5}}
=\pot{2}{-7}\,\si{\coulomb}=\SI{0.2}{\micro\coulomb}.
\]
Para a força ser para cima (mesmo sentido de $E$), a carga deve ser
\emph{positiva}. É o princípio do experimento de Millikan, que mediu a carga do
elétron equilibrando gotas de óleo.}
\end{questao}

\blocoaprofundamento

\begin{questao}[3]{}
Duas cargas $Q_1$ e $Q_2 = 4Q_1$ (ambas positivas) estão fixas nos pontos A e B,
distantes $\SI{30}{\centi\meter}$. Em que posição $x$, medida a partir de A sobre
a reta AB, deve-se colocar uma carga $Q_3$ para que ela fique em equilíbrio sob
ação apenas de forças elétricas?

\circuitoq{%
\begin{tikzpicture}[scale=1,line width=0.8pt]
 \draw[cinza] (0,0)--(6,0);
 \shade[ball color=Vcol!25] (0,0) circle (0.25);
 \shade[ball color=Vcol!25] (6,0) circle (0.25);
 \node at (0,0) {\footnotesize$Q_1$}; \node at (6,0) {\footnotesize$Q_2$};
 \node[below] at (0,-0.35) {A}; \node[below] at (6,-0.35) {B};
 \fill[laranja] (2,0) circle (0.09);
 \node[above,laranja] at (2,0.15) {$Q_3$};
 \draw[<->,cinza] (0,-0.9)--(2,-0.9) node[midway,below]{$x$};
 \draw[<->,cinza] (0,-1.5)--(6,-1.5) node[midway,below]{$\SI{30}{\centi\meter}$};
\end{tikzpicture}}{Posição de equilíbrio entre duas cargas.}

\resposta{$x = \SI{10}{\centi\meter}$ a partir de A}
\resolucao{No equilíbrio, as forças que $Q_1$ e $Q_2$ exercem sobre $Q_3$ têm
mesmo módulo e sentidos opostos (o que só é possível \emph{entre} as cargas, já
que ambas são positivas). Igualando os módulos, com $d = \SI{30}{\centi\meter}$:
\[
\ke\frac{Q_1 Q_3}{x^2}=\ke\frac{4Q_1\,Q_3}{(d-x)^2}
\ \Rightarrow\
\frac{1}{x^2}=\frac{4}{(d-x)^2}
\ \Rightarrow\
\frac{d-x}{x}=2.
\]
Logo $d - x = 2x \Rightarrow x = \dfrac{d}{3} = \SI{10}{\centi\meter}$. O ponto de
equilíbrio fica \emph{mais próximo} da carga menor ($Q_1$), como esperado.
Note que o valor de $Q_3$ e seu sinal não influenciam a posição.}
\end{questao}

\begin{questao}[3]{}
Quatro cargas ocupam os vértices de um quadrado: duas delas são $+Q$ e $-Q$ (mesmo
módulo), e as outras duas, $q_1$ e $q_2$, são desconhecidas. Coloca-se uma carga
de prova positiva no centro e observa-se que a força resultante sobre ela aponta
na direção da diagonal, \emph{do vértice $-Q$ para o vértice $+Q$}. Sobre $q_1$ e
$q_2$ (nos outros dois vértices), pode-se afirmar que, para que essa força
resultante fique exatamente sobre a diagonal $-Q\,$--$\,+Q$:
\begin{alternativas}
\item devem ser iguais em módulo e de mesmo sinal entre si.
\item devem ser iguais em módulo e de sinais opostos.
\item devem ser ambas nulas.
\item devem ser ambas positivas e diferentes.
\item nada se pode afirmar.
\end{alternativas}
\resposta{(a)}
\resolucao{A força devida ao par $+Q$/$-Q$ já está sobre a diagonal
$-Q\,$--$\,+Q$. Para que a resultante \emph{total} continue sobre essa mesma
diagonal, a contribuição do par $q_1$/$q_2$ (na outra diagonal) não pode ter
componente perpendicular a ela: as forças de $q_1$ e $q_2$ sobre a carga central
devem se \emph{cancelar} nessa direção transversal. Isso ocorre quando $q_1$ e
$q_2$ têm o mesmo módulo e o mesmo sinal (por simetria em relação à diagonal
principal), produzindo forças opostas sobre a carga de prova que se anulam.}
\end{questao}

\begin{questao}[3]{}
Duas cargas fixas, $q_1 = \SI{+9}{\micro\coulomb}$ em $x=0$ e
$q_2 = \SI{+4}{\micro\coulomb}$ em $x = \SI{10}{\centi\meter}$, estão sobre o eixo
$x$. Uma terceira carga $q_3$ é colocada sobre esse eixo.
\begin{itensq}
\item Determine a posição em que $q_3$ fica em equilíbrio.
\item Mostre que existe uma \emph{segunda} raiz matemática e explique por que ela
      deve ser descartada.
\item O equilíbrio encontrado é \emph{estável} ou \emph{instável}? Analise para
      $q_3$ positiva e negativa.
\end{itensq}
\resposta{(a) $x = \SI{6}{\centi\meter}$; (b) $x=\SI{30}{\centi\meter}$ (fora do
segmento); (c) estável para $q_3<0$, instável para $q_3>0$}
\resolucao{(a) Igualando os módulos das forças (a carga $q_3$ cancela-se):
\[
\frac{\ke\,9}{x^{2}}=\frac{\ke\,4}{(10-x)^{2}}
\;\Longrightarrow\;
\frac{3}{x}=\frac{2}{10-x}
\;\Longrightarrow\;
30-3x=2x
\;\Longrightarrow\;
x=\SI{6}{\centi\meter}.
\]
(b) Ao elevar ao quadrado, a equação admite também $\dfrac{3}{x}=-\dfrac{2}{10-x}$,
que fornece $x = \SI{30}{\centi\meter}$. Nesse ponto, \emph{fora} do segmento,
as duas forças apontam no \textbf{mesmo sentido} (ambas as cargas são positivas e
empurram $q_3$ para longe) --- elas se somam e nunca se cancelam. A raiz é
espúria, introduzida pela elevação ao quadrado.

(c) \textbf{Análise de estabilidade.} Suponha $q_3$ deslocada ligeiramente para a
direita (aproximando-se de $q_2$):
\begin{itemize}[leftmargin=1.6em,itemsep=1pt,topsep=2pt]
\item Se $q_3 > 0$: a repulsão de $q_2$ (agora mais próxima) supera a de $q_1$, e
      $q_3$ é empurrada \emph{ainda mais} para longe do ponto de equilíbrio ---
      \textbf{instável}.
\item Se $q_3 < 0$: a atração por $q_2$ cresce, puxando $q_3$ \emph{de volta}?
      Não --- a atração por $q_2$ a puxa para longe de $x=6$. Nesse eixo,
      portanto, também é instável.
\end{itemize}
\textbf{Análise completa:} para $q_3<0$, o equilíbrio é estável na direção
\emph{transversal} ao eixo (qualquer deslocamento lateral gera força restauradora),
mas instável ao longo do eixo. Para $q_3>0$, ocorre o oposto.\\
\textbf{Conclusão profunda --- teorema de Earnshaw:} \emph{nenhuma} configuração
de cargas estáticas produz equilíbrio estável em \textbf{todas} as direções. Sempre
há ao menos uma direção instável. É por isso que não se consegue levitar um objeto
apenas com cargas fixas --- são necessários campos variáveis (armadilhas de Paul)
ou diamagnetismo.}
\end{questao}

\begin{questao}[3]{}
Compare a força elétrica e a força gravitacional entre dois elétrons.
\begin{itensq}
\item Mostre que a razão $F_e/F_g$ \emph{independe} da distância.
\item Calcule seu valor numérico.
\item Que conclusão isso permite sobre a estrutura da matéria e sobre a
      astronomia?
\end{itensq}
\dados{$\ke=\SI{9e9}{}$; $G=\SI{6.67e-11}{}$; $e=\SI{1.6e-19}{\coulomb}$;
$m_e=\SI{9.11e-31}{\kilogram}$ (SI).}
\resposta{(a) demonstração; (b) $\approx\pot{4.2}{42}$; (c) ver comentário}
\resolucao{(a) Ambas as forças seguem a lei do inverso do quadrado:
\[
F_e=\ke\frac{e^{2}}{d^{2}},
\qquad
F_g=G\frac{m_e^{2}}{d^{2}}
\;\Longrightarrow\;
\frac{F_e}{F_g}=\frac{\ke\,e^{2}}{G\,m_e^{2}}.
\]
O fator $d^{2}$ \textbf{cancela-se}: a razão é uma constante universal.\\
(b) Numericamente:
\[
\frac{F_e}{F_g}=\frac{(9\cdot10^{9})(\pot{1.6}{-19})^{2}}
{(\pot{6.67}{-11})(\pot{9.11}{-31})^{2}}\approx\pot{4.2}{42}.
\]
(c) \textbf{Duas conclusões profundas:}
\begin{itemize}[leftmargin=1.6em,itemsep=2pt,topsep=3pt]
\item \textbf{Na escala atômica, a gravidade é irrelevante.} A força elétrica é
      $10^{42}$ vezes maior --- por isso a estrutura de átomos, moléculas,
      cristais e de toda a química é governada \emph{exclusivamente} por forças
      eletromagnéticas. Nenhum modelo atômico precisa incluir gravidade.
\item \textbf{Na escala astronômica, ocorre o oposto.} Corpos celestes são
      eletricamente \emph{neutros} (as cargas positivas e negativas se cancelam
      quase perfeitamente), então a força elétrica resultante é praticamente
      nula. Já a massa só se \emph{acumula} --- não existe "massa negativa". Por
      isso a gravidade, apesar de absurdamente mais fraca, domina o cosmos.
\end{itemize}
\emph{A natureza usa a força fraca para construir galáxias e a força forte para
construir átomos --- precisamente porque uma se cancela e a outra não.}}
\end{questao}

\begin{questao}[3]{}
Duas pequenas esferas idênticas, de massa $\SI{0.5}{\gram}$ cada, são suspensas
do mesmo ponto por fios isolantes de $\SI{20}{\centi\meter}$. Ao receberem cargas
iguais $q$, elas se afastam até que cada fio forme $\SI{10}{\degree}$ com a
vertical.

\circuitoq{%
\begin{tikzpicture}[scale=1,line width=0.8pt,>=Stealth]
 \draw[cinza!60,line width=2pt] (-1.6,3)--(1.6,3);
 \coordinate (O) at (0,3);
 \coordinate (P1) at ({-2.4*sin(10)},{3-2.4*cos(10)});
 \coordinate (P2) at ({ 2.4*sin(10)},{3-2.4*cos(10)});
 \draw[cinza,thick] (O)--(P1); \draw[cinza,thick] (O)--(P2);
 \draw[dashed,cinza!70] (O)--(0,0.3);
 \shade[ball color=Vcol!30] (P1) circle (0.2);
 \shade[ball color=Vcol!30] (P2) circle (0.2);
 \node[font=\footnotesize] at (0.34,2.55) {$10^\circ$};
 \node[left,font=\footnotesize] at ($(P1)+(-0.25,0)$) {$q$};
 \node[right,font=\footnotesize] at ($(P2)+(0.25,0)$) {$q$};
 \draw[<->,cinza] ($(P1)+(0,-0.45)$)--($(P2)+(0,-0.45)$)
       node[midway,below,font=\footnotesize]{$d$};
 \node[right,font=\footnotesize,cinza] at (0.9,1.9) {$L=\SI{20}{\centi\meter}$};
\end{tikzpicture}}{Pêndulos eletrostáticos idênticos.}

\begin{itensq}
\item Determine a separação $d$ entre as esferas.
\item Calcule a força elétrica entre elas.
\item Determine a carga $q$ de cada esfera.
\end{itensq}
\dados{$g=\SI{10}{\meter\per\second\squared}$; $\sin10^\circ=0{,}174$;
$\tan10^\circ=0{,}176$; $\ke=\SI{9e9}{}$.}
\resposta{(a) $d\approx\SI{6.9}{\centi\meter}$; (b) $F\approx\SI{0.88}{\milli\newton}$;
(c) $q\approx\SI{22}{\nano\coulomb}$}
\resolucao{(a) Cada esfera se afasta $L\sin\theta$ da vertical, logo
\[
d=2L\sin\theta=2(0{,}20)(0{,}174)\approx\SI{0.069}{\meter}=\SI{6.9}{\centi\meter}.
\]
(b) No equilíbrio de cada esfera (tração, peso e força elétrica),
a decomposição fornece $\tan\theta = \dfrac{F_e}{mg}$:
\[
F_e = mg\tan\theta=(\pot{0.5}{-3})(10)(0{,}176)
\approx\pot{8.8}{-4}\,\si{\newton}=\SI{0.88}{\milli\newton}.
\]
(c) Da lei de Coulomb, com as duas cargas iguais:
\[
F_e=\ke\frac{q^{2}}{d^{2}}
\;\Longrightarrow\;
q=d\sqrt{\frac{F_e}{\ke}}
=0{,}069\sqrt{\frac{\pot{8.8}{-4}}{9\cdot10^{9}}}
\approx\pot{2.2}{-8}\,\si{\coulomb}=\SI{22}{\nano\coulomb}.
\]
\textbf{Ordem de grandeza:} apenas $\sim\pot{1.4}{11}$ elétrons em excesso ---
uma quantidade ínfima --- já produz um afastamento visível. \emph{Isso mostra
quão intensa é a força elétrica em escala macroscópica.}}
\end{questao}

\begin{questao}[3]{}
Quatro cargas de mesmo módulo $q=\SI{1}{\micro\coulomb}$ ocupam os vértices de um
quadrado de lado $a=\SI{10}{\centi\meter}$, \textbf{alternando os sinais}
($+,-,+,-$ ao percorrer o perímetro). Determine o módulo e a direção da força
resultante sobre uma das cargas positivas.
\dados{$\ke=\SI{9e9}{}$; $\sqrt2\approx1{,}41$.}
\resposta{$F_R \approx \SI{0.82}{\newton}$, dirigida ao centro do quadrado}
\resolucao{\textbf{Forças envolvidas.} Sobre a carga positiva escolhida atuam:
\begin{itemize}[leftmargin=1.6em,itemsep=1pt,topsep=2pt]
\item duas forças dos vizinhos \emph{adjacentes} (que são \textbf{negativos}):
      são de \textbf{atração}, ao longo dos lados;
\item uma força da carga da \emph{diagonal} (que é \textbf{positiva}):
      \textbf{repulsão}, ao longo da diagonal.
\end{itemize}
\textbf{Módulos:}
\[
F_L=\ke\frac{q^{2}}{a^{2}}=9\cdot10^{9}\frac{(\pot{1}{-6})^{2}}{(0{,}1)^{2}}
=\SI{0.9}{\newton},
\qquad
F_D=\ke\frac{q^{2}}{(a\sqrt2)^{2}}=\frac{F_L}{2}=\SI{0.45}{\newton}.
\]
\textbf{Composição.} As duas atrações, perpendiculares entre si, têm resultante
$F_L\sqrt2 \approx \SI{1.27}{\newton}$ apontando \emph{para o centro} do quadrado
(ao longo da diagonal). A repulsão da diagonal aponta \emph{para fora}, na mesma
reta. Logo elas se \textbf{subtraem}:
\[
F_R=F_L\sqrt2-F_D = 1{,}27-0{,}45\approx\SI{0.82}{\newton},
\]
dirigida \textbf{para o centro} do quadrado.\\
\textbf{Contraste com o caso de cargas iguais} (visto anteriormente): lá as três
forças eram repulsivas e se \emph{somavam} ($F_R=F_L\sqrt2+F_D=\SI{1.72}{\newton}$),
apontando para fora. Aqui, alternar os sinais inverte o sentido e reduz o módulo
para menos da metade. \emph{A geometria é a mesma; a configuração de sinais muda
tudo.}}
\end{questao}

\begin{questao}[3]{}
Duas cargas fixas, $q_1=+q$ e $q_2=+4q$, estão separadas por $d$. Deseja-se
colocar uma terceira carga $q_3$ de modo que \textbf{as três} fiquem em
equilíbrio simultaneamente.
\begin{itensq}
\item Determine a posição de $q_3$.
\item Determine o valor e o sinal de $q_3$.
\item Esse equilíbrio é estável?
\end{itensq}
\resposta{(a) a $d/3$ de $q_1$; (b) $q_3=-\dfrac{4q}{9}$; (c) instável}
\resolucao{(a) Para $q_3$ estar em equilíbrio, ela deve ficar onde os campos de
$q_1$ e $q_2$ se cancelam (calculado anteriormente):
\[
\frac{q}{x^{2}}=\frac{4q}{(d-x)^{2}}
\;\Rightarrow\;
\frac{d-x}{x}=2
\;\Rightarrow\;
x=\frac{d}{3}.
\]
(b) Mas isso garante apenas o equilíbrio de $q_3$. Para que \emph{$q_1$} também
esteja equilibrada, a repulsão de $q_2$ sobre ela deve ser cancelada pela atração
de $q_3$:
\[
\underbrace{\ke\frac{q\cdot 4q}{d^{2}}}_{\text{repulsão de }q_2}
=\underbrace{\ke\frac{q\,|q_3|}{(d/3)^{2}}}_{\text{atração de }q_3}
\;\Longrightarrow\;
\frac{4q}{d^{2}}=\frac{9|q_3|}{d^{2}}
\;\Longrightarrow\;
|q_3|=\frac{4q}{9}.
\]
O sinal deve ser \textbf{negativo} (para atrair), logo
$q_3 = -\dfrac{4q}{9}$.\\
\textbf{Verificação em $q_2$:} a repulsão de $q_1$ sobre $q_2$ é
$\ke\dfrac{4q^2}{d^2}$; a atração de $q_3$ (à distância $2d/3$) é
$\ke\dfrac{4q\cdot(4q/9)}{(2d/3)^{2}}=\ke\dfrac{(16q^2/9)}{(4d^2/9)}
=\ke\dfrac{4q^{2}}{d^{2}}$. \checkmark\ São iguais --- as três estão em
equilíbrio.\\
(c) \textbf{Instável}, pelo teorema de Earnshaw: qualquer deslocamento infinitesimal
de qualquer das cargas destrói o equilíbrio, e as forças resultantes \emph{afastam}
ainda mais o sistema da configuração original. Não existe arranjo eletrostático
puro que seja estável.}
\end{questao}

\begin{questao}[3]{}
Uma pequena esfera de massa $m = \SI{10}{\gram}$, eletrizada com carga $q$, é
suspensa por um fio isolante de massa desprezível em uma região de campo elétrico
\emph{horizontal} e uniforme, de intensidade $E = \SI{e4}{\newton\per\coulomb}$.
No equilíbrio, o fio forma $\SI{30}{\degree}$ com a vertical.

\circuitoq{%
\begin{tikzpicture}[scale=1,line width=0.8pt,>=Stealth]
 \draw[cinza!60,line width=2pt] (-2,2.6)--(2,2.6);
 \coordinate (O) at (0,2.6);
 \coordinate (P) at ({2.2*sin(30)},{2.6-2.2*cos(30)});
 \draw[cinza,thick] (O)--(P);
 \draw[dashed,cinza!70] (O)--(0,0.2);
 \shade[ball color=Vcol!30] (P) circle (0.22);
 \node[right,font=\footnotesize] at ($(P)+(0.28,0)$) {$m,\,q$};
 \draw[->,Vcol,thick] (P)--($(P)+(1.0,0)$) node[right,font=\footnotesize]{$\vec F_e$};
 \draw[->,cinza,thick] (P)--($(P)+(0,-0.9)$) node[below,font=\footnotesize]{$\vec P$};
 \node[font=\footnotesize] at (0.42,1.9) {$30^\circ$};
 \foreach \y in {0.4,1.0,1.6} \draw[->,Vcol!45] (-1.9,\y)--(-1.2,\y);
 \node[Vcol,font=\footnotesize] at (-1.55,2.05) {$\vec E$};
\end{tikzpicture}}{Pêndulo eletrostático em campo horizontal.}

\begin{itensq}
\item Determine a carga $q$ da esfera.
\item Calcule a tração no fio.
\item Se o campo dobrasse, o ângulo dobraria? Justifique quantitativamente.
\end{itensq}
\dados{$g = \SI{10}{\meter\per\second\squared}$; $\tan30^\circ\approx0{,}577$.}
\resposta{(a) $q \approx \SI{5.8}{\micro\coulomb}$; (b) $T\approx\SI{0.115}{\newton}$;
(c) não --- o ângulo cresce menos que proporcionalmente}
\resolucao{\textbf{(a)} No equilíbrio, decompondo a tração:
\[
T\sin\theta = F_e = qE,
\qquad
T\cos\theta = P = mg
\quad\Longrightarrow\quad
\tan\theta=\frac{qE}{mg}.
\]
Isolando $q$:
\[
q=\frac{mg\tan\theta}{E}
=\frac{(\pot{10}{-3})(10)(0{,}577)}{\pot{1}{4}}
\approx\pot{5.8}{-6}\,\si{\coulomb}=\SI{5.8}{\micro\coulomb}.
\]
\textbf{(b)} $T = \dfrac{mg}{\cos\theta} = \dfrac{0{,}1}{0{,}866}
\approx\SI{0.115}{\newton}$.\\
\textbf{(c)} \textbf{Não.} A relação é $\tan\theta \propto E$, e \emph{não}
$\theta \propto E$. Dobrando o campo:
\[
\tan\theta' = 2\tan30^\circ = 1{,}155
\;\Rightarrow\;
\theta' = \arctan(1{,}155)\approx 49{,}1^\circ,
\]
que é \emph{menos} que $2\times30^\circ = 60^\circ$. A função tangente cresce mais
rápido que o ângulo, então o ângulo responde de forma \emph{sublinear}. No limite
$E\to\infty$, $\theta\to90^\circ$ mas nunca o atinge --- o fio jamais fica
horizontal.}
\end{questao}

\begin{questao}[3]{}
Três cargas puntiformes estão sobre uma reta: $+Q$ na origem, $-2Q$ em $x=a$ e
$+Q$ em $x=2a$ (um \emph{quadrupolo linear}).
\begin{itensq}
\item Determine a força resultante sobre a carga central.
\item Determine a força resultante sobre a carga da extremidade em $x=0$.
\item Explique por que o sistema, embora tenha carga total nula, ainda exerce
      força sobre uma carga externa distante.
\end{itensq}
\resposta{(a) nula, por simetria; (b) $F=\dfrac{7\ke Q^{2}}{4a^{2}}$ dirigida para
$+x$; (c) ver comentário}
\resolucao{(a) A carga central $-2Q$ está \emph{equidistante} das duas cargas
$+Q$, que são iguais. As duas atrações têm mesmo módulo e sentidos opostos: a
resultante é \textbf{nula} por simetria.\\
(b) Sobre a carga em $x=0$ atuam:
\begin{itemize}[leftmargin=1.6em,itemsep=1pt,topsep=2pt]
\item atração de $-2Q$ (distância $a$):
      $F_1=\ke\dfrac{Q\cdot2Q}{a^{2}}=\dfrac{2\ke Q^{2}}{a^{2}}$, para $+x$;
\item repulsão de $+Q$ (distância $2a$):
      $F_2=\ke\dfrac{Q^{2}}{4a^{2}}$, para $-x$.
\end{itemize}
Resultante:
\[
F=F_1-F_2=\frac{2\ke Q^{2}}{a^{2}}-\frac{\ke Q^{2}}{4a^{2}}
=\frac{\ke Q^{2}}{a^{2}}\left(2-\frac14\right)
=\frac{7\ke Q^{2}}{4a^{2}},
\]
dirigida para $+x$ (para dentro do conjunto).\\
(c) \textbf{Por que a carga total nula não implica campo nulo.} A carga total é
$Q-2Q+Q=0$, então \emph{a longas distâncias} o campo decai muito mais rápido que
o de uma carga isolada. Mas ele não é nulo: as cargas ocupam \emph{posições
diferentes}, e essa separação espacial cria momentos de ordem superior.\\
\textbf{Hierarquia dos decaimentos:}
\begin{center}\small\sffamily
\begin{tabular}{@{}l c c@{}}
\toprule
\textbf{Configuração} & \textbf{Carga total} & \textbf{Campo distante}\\
\midrule
Monopolo (carga isolada) & $\neq0$ & $\propto 1/r^{2}$\\
Dipolo ($+q,-q$)         & $0$     & $\propto 1/r^{3}$\\
Quadrupolo (este caso)   & $0$     & $\propto 1/r^{4}$\\
\bottomrule
\end{tabular}
\end{center}
\emph{Cada cancelamento sucessivo acelera o decaimento --- é o princípio da
expansão multipolar, base da análise de antenas e de moléculas polares.}}
\end{questao}

\blocodesafio

\begin{questao}[4]{}
Uma carga $Q$ é dividida em duas partes, $q$ e $(Q-q)$, que são então colocadas a
uma distância fixa $d$ uma da outra.
\begin{itensq}
\item Obtenha a expressão da força entre elas em função de $q$.
\item Determine o valor de $q$ que \textbf{maximiza} essa força.
\item Qual é a força máxima?
\end{itensq}
\resposta{(a) $F=\dfrac{\ke\,q(Q-q)}{d^{2}}$; (b) $q=Q/2$;
(c) $F_{\max}=\dfrac{\ke Q^{2}}{4d^{2}}$}
\resolucao{(a) Pela lei de Coulomb,
\[
F(q)=\ke\,\frac{q\,(Q-q)}{d^{2}}=\frac{\ke}{d^{2}}\left(Qq-q^{2}\right).
\]
(b) É uma \textbf{parábola} com concavidade para baixo em função de $q$. O máximo
ocorre no vértice. Derivando:
\[
\frac{\mathrm{d}F}{\mathrm{d}q}=\frac{\ke}{d^{2}}\left(Q-2q\right)=0
\;\Longrightarrow\;
\boxed{\,q=\frac{Q}{2}\,}
\]
(Sem cálculo: o vértice de $Qq-q^2$ está em $q = -\dfrac{b}{2a}=\dfrac{Q}{2}$.)\\
(c) Substituindo:
\[
F_{\max}=\frac{\ke}{d^{2}}\cdot\frac{Q}{2}\cdot\frac{Q}{2}
=\frac{\ke\,Q^{2}}{4d^{2}}.
\]
\textbf{Interpretação.} A força é máxima quando a carga é dividida
\emph{igualmente}. Isso decorre da desigualdade entre médias: para soma fixa
($q+(Q-q)=Q$), o produto é máximo quando as parcelas são iguais.\\
\textbf{Casos-limite:} se $q\to0$ ou $q\to Q$, uma das partes fica sem carga e
$F\to0$ --- coerente. \emph{Este é um problema de otimização disfarçado de
eletrostática: reconhecer a estrutura matemática vale mais que decorar a
fórmula.}}
\end{questao}

% BEGIN BANCO COMPLEMENTAR Lei de Coulomb
% =====================================================================
%  Banco complementar --- Lei de Coulomb  (REVISADO)
%  Substitui a versao gerada por template: 32 questoes, todas com
%  enunciado distinto. Nivel 4 reclassificado conforme o criterio da
%  obra (sintese, demonstracao, otimizacao ou modelagem nao rotineira).
%  Distribuicao mantida: 7 N1 + 10 N2 + 6 N3 + 9 N4 = 32
% =====================================================================

\subsubsection*{Banco complementar --- Lei de Coulomb}

\begin{atencao}[Organização do banco]
As questões abaixo complementam o tópico sem repetir enunciados da seção
principal. Cada nível explora uma \emph{estrutura de raciocínio} diferente:
N1 aplica a fórmula em suas quatro formas isoladas ($F$, $q$, $r$, meio);
N2 exige composição ou reinterpretação; N3 combina vetores, equilíbrio e
redistribuição de carga; N4 pede demonstração, integração, otimização ou
modelagem.
\end{atencao}

\blocofixacao

\begin{questao}[1]{}
Duas cargas puntiformes $q_1=\SI{20}{\nano\coulomb}$ e
$q_2=\SI{30}{\nano\coulomb}$ estão separadas por $\SI{3}{\centi\meter}$ no
vácuo. Determine o módulo da força elétrica e classifique-a.
\dados{$\ke=\SI{9e9}{\newton\meter\squared\per\coulomb\squared}$}
\resposta{$F=\SI{6}{\milli\newton}$, repulsiva}
\resolucao{Convertendo antes de substituir --- este é o erro mais comum:
$q_1=\pot{2}{-8}\,\si{\coulomb}$, $q_2=\pot{3}{-8}\,\si{\coulomb}$,
$r=\pot{3}{-2}\,\si{\meter}$.
\[
F=\ke\frac{q_1q_2}{r^{2}}
 =\frac{\pot{9}{9}\cdot\pot{2}{-8}\cdot\pot{3}{-8}}{(\pot{3}{-2})^{2}}
 =\frac{\pot{5.4}{-6}}{\pot{9}{-4}}=\pot{6}{-3}\,\si{\newton}.
\]
Ambas positivas $\Rightarrow$ repulsão.}
\end{questao}

\begin{questao}[1]{}
Duas esferas idênticas recebem cargas iguais $q$ e se repelem com força de
$\SI{0.90}{\newton}$ quando separadas por $\SI{20}{\centi\meter}$.
\begin{itensq}
\item Determine $q$.
\item Quantos elétrons foram \textbf{retirados} de cada esfera?
\end{itensq}
\dados{$e=\pot{1.6}{-19}\,\si{\coulomb}$}
\resposta{(a) $q=\SI{2}{\micro\coulomb}$; (b) $n=\pot{1.25}{13}$ elétrons}
\resolucao{(a) De $F=\ke q^{2}/r^{2}$,
\[
q=\sqrt{\frac{Fr^{2}}{\ke}}=\sqrt{\frac{0{,}90\cdot(0{,}20)^{2}}{\pot{9}{9}}}
 =\sqrt{\pot{4}{-12}}=\pot{2}{-6}\,\si{\coulomb}.
\]
(b) Repulsão com cargas iguais e sinal positivo indica \emph{falta} de
elétrons: $n=q/e=\pot{2}{-6}/\pot{1.6}{-19}=\pot{1.25}{13}$.}
\end{questao}

\begin{questao}[1]{}
Duas cargas de $\SI{5}{\micro\coulomb}$ cada devem exercer uma sobre a outra
força de $\SI{2.5}{\newton}$ no vácuo. A que distância devem ser colocadas?
\resposta{$r=\SI{0.30}{\meter}$}
\resolucao{Isolando $r$ na lei de Coulomb:
\[
r=\sqrt{\frac{\ke q_1q_2}{F}}
 =\sqrt{\frac{\pot{9}{9}\cdot(\pot{5}{-6})^{2}}{2{,}5}}
 =\sqrt{\frac{0{,}225}{2{,}5}}=\sqrt{0{,}09}=\SI{0.30}{\meter}.
\]}
\end{questao}

\begin{questao}[1]{}
Duas cargas puntiformes se atraem com força $F$. Mantendo as cargas e
\textbf{dobrando} a distância entre elas, a nova força será:
\begin{alternativas}[3]
\item $F/4$
\item $F/2$
\item $F$
\item $2F$
\item $4F$
\end{alternativas}
\resposta{(a)}
\resolucao{$F\propto 1/r^{2}$. Dobrar $r$ divide a força por $2^{2}=4$:
$F'=F/4$. Note que a natureza (atrativa) não muda --- ela depende apenas dos
sinais, nunca da distância.}
\end{questao}

\begin{questao}[1]{}
Duas cargas separadas por uma distância fixa exercem força de
$\SI{0.80}{\newton}$ no vácuo. Todo o conjunto é então imerso em um líquido
isolante de permissividade relativa $\varepsilon_r=4$. Determine a nova força.
\resposta{$F'=\SI{0.20}{\newton}$}
\resolucao{Em um meio dielétrico, $F_{\text{meio}}=F_{\text{vácuo}}/\varepsilon_r$,
pois $\ke$ é substituída por $\ke/\varepsilon_r$:
\[
F'=\frac{0{,}80}{4}=\SI{0.20}{\newton}.
\]
O dielétrico se polariza e blinda parcialmente as cargas --- por isso a força
\emph{sempre} diminui ($\varepsilon_r\ge 1$).}
\end{questao}

\begin{questao}[1]{}
Um bastão carregado positivamente é aproximado --- sem tocar --- de uma esfera
metálica \textbf{neutra} e isolada. A força entre eles é:
\begin{alternativas}[2]
\item nula, pois a esfera é neutra
\item atrativa
\item repulsiva
\item atrativa ou repulsiva, dependendo da distância
\end{alternativas}
\resposta{(b)}
\resolucao{A carga \emph{líquida} da esfera continua nula, mas o bastão
positivo atrai elétrons livres para a face próxima e repele carga positiva
para a face distante (\textbf{indução eletrostática}). A face próxima, negativa,
está \emph{mais perto} do bastão que a face distante, positiva. Como
$F\propto 1/r^{2}$, a atração vence: resultante \textbf{atrativa}. É o mesmo
princípio pelo qual um canudo atritado atrai pedaços de papel neutros.}
\end{questao}

\begin{questao}[1]{}
Estime a força entre duas cargas de $\SI{1}{\coulomb}$ separadas por
$\SI{1}{\kilo\meter}$ e comente o resultado.
\resposta{$F=\SI{9}{\kilo\newton}$ (o peso de cerca de \SI{900}{\kilogram})}
\resolucao{
\[
F=\frac{\pot{9}{9}\cdot 1\cdot 1}{(10^{3})^{2}}=\pot{9}{3}\,\si{\newton}.
\]
Nove mil newtons --- o peso de quase uma tonelada --- a \emph{um quilômetro}
de distância. Isso mostra por que o coulomb é uma unidade enorme para a
eletrostática: cargas realmente separáveis por atrito ficam na faixa de
$\si{\nano\coulomb}$ a $\si{\micro\coulomb}$.}
\end{questao}

\blocoaplicacao

\begin{questao}[2]{}
Três cargas estão fixas sobre o eixo $x$: $q_1=\SI{+3}{\micro\coulomb}$ em
$x=0$, $q_2=\SI{+2}{\micro\coulomb}$ em $x=\SI{20}{\centi\meter}$ e
$q_3=\SI{-5}{\micro\coulomb}$ em $x=\SI{50}{\centi\meter}$. Determine a força
resultante sobre $q_2$ (módulo e sentido).
\resposta{$F=\SI{2.35}{\newton}$, no sentido $+x$}
\resolucao{Trate cada par separadamente e some \emph{com sinal}.
\[
F_{12}=\frac{\pot{9}{9}\cdot\pot{3}{-6}\cdot\pot{2}{-6}}{(0{,}20)^{2}}
=\SI{1.35}{\newton}\quad(\text{repulsão}\Rightarrow +x)
\]
\[
F_{32}=\frac{\pot{9}{9}\cdot\pot{2}{-6}\cdot\pot{5}{-6}}{(0{,}30)^{2}}
=\SI{1.00}{\newton}\quad(\text{atração por } q_3 \Rightarrow +x)
\]
As duas apontam no mesmo sentido: $F=1{,}35+1{,}00=\SI{2.35}{\newton}$ em $+x$.
\textbf{Cuidado:} não some as cargas antes; a lei de Coulomb vale par a par
(princípio da superposição).}
\end{questao}

\begin{questao}[2]{}
Duas esferas metálicas idênticas com $q_1=\SI{+8}{\micro\coulomb}$ e
$q_2=\SI{-2}{\micro\coulomb}$ estão a $\SI{30}{\centi\meter}$. Elas são postas
em contato e recolocadas na mesma posição.
\begin{itensq}
\item Calcule a força antes do contato.
\item Calcule a força depois.
\end{itensq}
\resposta{(a) $\SI{1.6}{\newton}$, atrativa; (b) $\SI{0.9}{\newton}$, repulsiva}
\resolucao{(a) $F=\dfrac{\pot{9}{9}\cdot\pot{8}{-6}\cdot\pot{2}{-6}}{(0{,}30)^{2}}
=\dfrac{0{,}144}{0{,}09}=\SI{1.6}{\newton}$, atrativa (sinais opostos).\\
(b) Esferas \emph{idênticas} em contato dividem a carga total igualmente:
\[
q'=\frac{q_1+q_2}{2}=\frac{+8-2}{2}=\SI{+3}{\micro\coulomb}\ \text{cada}.
\]
\[
F'=\frac{\pot{9}{9}\cdot(\pot{3}{-6})^{2}}{0{,}09}=\SI{0.9}{\newton},
\ \text{agora repulsiva}.
\]
A força \emph{e} a natureza mudaram: o contato conserva a carga total, não a
individual.}
\end{questao}

\begin{questao}[2]{}
Duas cargas fixas exercem entre si $\SI{0.60}{\newton}$ no vácuo. O conjunto é
mergulhado em óleo com $\varepsilon_r=2{,}5$.
\begin{itensq}
\item Qual é a nova força, mantida a distância?
\item Que fração da distância original restauraria a força de
      $\SI{0.60}{\newton}$?
\end{itensq}
\resposta{(a) $\SI{0.24}{\newton}$; (b) $r'=r/\sqrt{2{,}5}\approx 0{,}63\,r$}
\resolucao{(a) $F'=F/\varepsilon_r=0{,}60/2{,}5=\SI{0.24}{\newton}$.\\
(b) Queremos $\dfrac{\ke q_1q_2}{\varepsilon_r r'^{2}}=\dfrac{\ke q_1q_2}{r^{2}}$,
ou seja $\varepsilon_r r'^{2}=r^{2}$:
\[
r'=\frac{r}{\sqrt{\varepsilon_r}}=\frac{r}{\sqrt{2{,}5}}=0{,}632\,r.
\]
Aproximar as cargas em cerca de $37\%$ compensa exatamente o efeito do meio.}
\end{questao}

\begin{questao}[2]{}
A força entre duas cargas $q_1$ e $q_2$ separadas por $d$ vale $F$. Se $q_1$ for
\textbf{triplicada} e a distância for reduzida à \textbf{metade}, a nova força
será:
\begin{alternativas}[3]
\item $1{,}5F$
\item $3F$
\item $6F$
\item $12F$
\item $24F$
\end{alternativas}
\resposta{(d)}
\resolucao{Os efeitos são multiplicativos e independentes:
\[
F'=\ke\frac{(3q_1)q_2}{(d/2)^{2}}=3\cdot 4\cdot \ke\frac{q_1q_2}{d^{2}}=12F.
\]
Fator $3$ da carga $\times$ fator $4$ da distância $=12$.}
\end{questao}

\begin{questao}[2]{}
Em uma impressora a jato eletrostático, cada gota de tinta tem massa
$\pot{1.0}{-11}\,\si{\kilogram}$ e é carregada pela remoção de
$\pot{2.0}{5}$ elétrons. Duas gotas vizinhas ficam a $\SI{1.0}{\milli\meter}$.
Compare a força elétrica entre elas com o peso de uma gota.
\dados{$e=\pot{1.6}{-19}\,\si{\coulomb}$; $g=\SI{9.8}{\meter\per\second\squared}$}
\resposta{$F_e=\pot{9.2}{-12}\,\si{\newton}$; $P=\pot{9.8}{-11}\,\si{\newton}$;
$F_e/P\approx 0{,}094$}
\resolucao{Carga da gota: $q=ne=\pot{2.0}{5}\cdot\pot{1.6}{-19}
=\pot{3.2}{-14}\,\si{\coulomb}$.
\[
F_e=\frac{\pot{9}{9}\cdot(\pot{3.2}{-14})^{2}}{(\pot{1.0}{-3})^{2}}
=\pot{9.2}{-12}\,\si{\newton}.
\]
Peso: $P=mg=\pot{1.0}{-11}\cdot 9{,}8=\pot{9.8}{-11}\,\si{\newton}$.
Razão $\approx 9{,}4\%$: a repulsão entre gotas é pequena diante do peso, mas
não desprezível --- é ela que degrada a nitidez quando as gotas são emitidas
muito próximas.}
\end{questao}

\begin{questao}[2]{}
Duas cargas positivas, $q_1=\SI{4}{\micro\coulomb}$ em $x=0$ e
$q_2=\SI{9}{\micro\coulomb}$ em $x=\SI{50}{\centi\meter}$, estão fixas.
Determine o ponto onde uma terceira carga ficaria em equilíbrio.
\resposta{$x=\SI{20}{\centi\meter}$}
\resolucao{Como as duas são positivas, o ponto de força nula está
\emph{entre} elas. Seja $x$ a distância a $q_1$:
\[
\frac{\ke q_1 q_3}{x^{2}}=\frac{\ke q_2 q_3}{(0{,}50-x)^{2}}
\;\Longrightarrow\;
\frac{4}{x^{2}}=\frac{9}{(0{,}50-x)^{2}}.
\]
Extraindo a raiz (ambos os lados positivos):
$\dfrac{2}{x}=\dfrac{3}{0{,}50-x}\Rightarrow 2(0{,}50-x)=3x
\Rightarrow x=\SI{0.20}{\meter}$.
O ponto fica \emph{mais próximo da carga menor}, como esperado. Note que $q_3$
se cancela: a posição não depende dela.}
\end{questao}

\begin{questao}[2]{}
Uma carga $q$ e outra $100q$ interagem. Sobre os módulos das forças que uma
exerce sobre a outra, é correto afirmar:
\begin{alternativas}[2]
\item a força sobre $q$ é 100 vezes maior
\item a força sobre $100q$ é 100 vezes maior
\item as forças têm o mesmo módulo
\item depende das massas
\end{alternativas}
\resposta{(c)}
\resolucao{A lei de Coulomb é simétrica no produto $q_1q_2$: ambas valem
$\ke\,q\cdot 100q/r^{2}$. Formam um par ação--reação (3\textsuperscript{a} lei
de Newton): mesmo módulo, mesma direção, sentidos opostos. O que \emph{difere}
é a \textbf{aceleração}, se as massas forem diferentes --- confundir força com
aceleração é o erro clássico aqui.}
\end{questao}

\begin{questao}[2]{}
No modelo de Bohr para o hidrogênio, o elétron orbita o próton a
$\pot{5.3}{-11}\,\si{\meter}$. Determine a força elétrica e a aceleração
centrípeta do elétron.
\dados{$e=\pot{1.6}{-19}\,\si{\coulomb}$; $m_e=\pot{9.11}{-31}\,\si{\kilogram}$}
\resposta{$F=\pot{8.2}{-8}\,\si{\newton}$;
$a=\pot{9.0}{22}\,\si{\meter\per\second\squared}$}
\resolucao{
\[
F=\frac{\pot{9}{9}\cdot(\pot{1.6}{-19})^{2}}{(\pot{5.3}{-11})^{2}}
=\frac{\pot{2.304}{-28}}{\pot{2.81}{-21}}=\pot{8.2}{-8}\,\si{\newton}.
\]
\[
a=\frac{F}{m_e}=\frac{\pot{8.2}{-8}}{\pot{9.11}{-31}}
=\pot{9.0}{22}\,\si{\meter\per\second\squared}.
\]
Uma força minúscula produz aceleração de $10^{22}\,\si{\meter\per\second\squared}$
porque a massa do elétron é ainda mais minúscula.}
\end{questao}

\begin{questao}[2]{}
Duas esferas metálicas idênticas com cargas $+Q$ e $+3Q$, separadas por $d$,
repelem-se com força $F$. Elas são postas em contato e depois afastadas até
$2d$. Determine a nova força em função de $F$.
\resposta{$F'=F/3$}
\resolucao{Antes: $F=\ke\dfrac{Q\cdot 3Q}{d^{2}}=\dfrac{3\ke Q^{2}}{d^{2}}$.\\
Após o contato, cada esfera fica com $\dfrac{Q+3Q}{2}=2Q$. A $2d$:
\[
F'=\ke\frac{(2Q)(2Q)}{(2d)^{2}}=\ke\frac{4Q^{2}}{4d^{2}}=\frac{\ke Q^{2}}{d^{2}}
=\frac{F}{3}.
\]
O produto das cargas subiu (de $3Q^2$ para $4Q^2$), mas o fator $4$ da
distância dominou.}
\end{questao}

\begin{questao}[2]{}
Compare a força entre duas cargas de $\SI{1.0}{\micro\coulomb}$ separadas por
$\SI{1.0}{\centi\meter}$ com o peso de um corpo de $\SI{1.0}{\kilogram}$, e
interprete.
\dados{$g=\SI{9.8}{\meter\per\second\squared}$}
\resposta{$F=\SI{90}{\newton}$ contra $P=\SI{9.8}{\newton}$; razão $\approx 9{,}2$}
\resolucao{
\[
F=\frac{\pot{9}{9}\cdot(\pot{1}{-6})^{2}}{(\pot{1}{-2})^{2}}
=\frac{\pot{9}{-3}}{\pot{1}{-4}}=\SI{90}{\newton};
\qquad P=1{,}0\cdot 9{,}8=\SI{9.8}{\newton}.
\]
Um milionésimo de coulomb, a um centímetro, produz nove vezes o peso de um
quilograma. Essa desproporção entre a força elétrica e a gravitacional é a razão
pela qual é a matéria macroscópica é praticamente neutra: qualquer desequilíbrio
apreciável de carga seria violentamente desfeito.}
\end{questao}

\blocoaprofundamento

\begin{questao}[3]{}
Duas cargas fixas sobre o eixo $x$: $q_1=\SI{+4}{\micro\coulomb}$ em $x=0$ e
$q_2=\SI{-1}{\micro\coulomb}$ em $x=\SI{30}{\centi\meter}$. Determine todos os
pontos do eixo onde a força sobre uma carga de prova seria nula.
\resposta{Um único ponto: $x=\SI{60}{\centi\meter}$}
\resolucao{Com \textbf{sinais opostos}, o ponto não pode estar entre as cargas
(ali as duas forças apontam para o mesmo lado). Deve estar \emph{fora} do
segmento e, mais especificamente, do lado da carga de \emph{menor} módulo ---
só assim a maior proximidade compensa o menor valor. Testando $x>0{,}30$:
\[
\frac{\ke\cdot 4\mu}{x^{2}}=\frac{\ke\cdot 1\mu}{(x-0{,}30)^{2}}
\;\Longrightarrow\;
\frac{2}{x}=\frac{1}{x-0{,}30}
\;\Longrightarrow\;
2x-0{,}60=x \Rightarrow x=\SI{0.60}{\meter}.
\]
\textbf{Verificação:}
$\ke\cdot4\mu/(0{,}60)^{2}=\ke\cdot1\mu/(0{,}30)^{2}$ --- confere. A raiz
$x=0{,}20$ da equação quadrática completa é \emph{espúria}: corresponde ao
ponto interno, onde as forças se somam em vez de se cancelar.}
\end{questao}

\begin{questao}[3]{}
Seis cargas iguais $Q=\SI{2}{\micro\coulomb}$ ocupam os vértices de um hexágono
regular de lado $a=\SI{10}{\centi\meter}$, com uma carga
$q=\SI{1}{\micro\coulomb}$ no centro.
\begin{itensq}
\item Qual é a força resultante sobre $q$?
\item Uma das seis cargas é removida. Qual passa a ser a força sobre $q$
      (módulo e direção)?
\end{itensq}
\resposta{(a) nula; (b) $\SI{1.8}{\newton}$, apontando \emph{para} o vértice vazio}
\resolucao{(a) Em um hexágono regular, o centro dista $a$ de todos os vértices
e as cargas ocupam posições diametralmente opostas duas a duas. Cada par
contribui com forças de mesmo módulo e sentidos opostos: resultante nula, por
simetria.\\
(b) Use o argumento da superposição \emph{ao contrário}: 
$\vec{F}_{6}= \vec{F}_{5}+\vec{F}_{\text{removida}}=\vec{0}$, logo
$\vec{F}_{5}=-\vec{F}_{\text{removida}}$.
\[
|\vec{F}_{5}|=\frac{\ke Qq}{a^{2}}
=\frac{\pot{9}{9}\cdot\pot{2}{-6}\cdot\pot{1}{-6}}{(0{,}10)^{2}}
=\SI{1.8}{\newton}.
\]
A carga removida \emph{repelia} $q$ para longe do vértice; sem ela, a resultante
aponta \textbf{para} o vértice vazio. Este truque de simetria evita somar cinco
vetores.}
\end{questao}

\begin{questao}[3]{}
Três cargas ocupam os vértices de um triângulo retângulo:
$q_A=\SI{+2}{\micro\coulomb}$ na origem, $q_B=\SI{+3}{\micro\coulomb}$ em
$(\SI{30}{\centi\meter},0)$ e $q_C=\SI{+4}{\micro\coulomb}$ em
$(0,\SI{40}{\centi\meter})$. Determine a força resultante sobre $q_A$: módulo e
ângulo com o eixo $x$.
\resposta{$F=\SI{0.75}{\newton}$, a $\SI{216.9}{\degree}$ (isto é, $\SI{36.9}{\degree}$ abaixo
do semieixo $-x$)}
\resolucao{As duas forças sobre $q_A$ são perpendiculares --- o que torna a
composição imediata.
\[
F_{AB}=\frac{\pot{9}{9}\cdot\pot{2}{-6}\cdot\pot{3}{-6}}{(0{,}30)^{2}}
=\SI{0.60}{\newton}\ \ (\text{repulsão}\Rightarrow -x)
\]
\[
F_{AC}=\frac{\pot{9}{9}\cdot\pot{2}{-6}\cdot\pot{4}{-6}}{(0{,}40)^{2}}
=\SI{0.45}{\newton}\ \ (\text{repulsão}\Rightarrow -y)
\]
\[
F=\sqrt{0{,}60^{2}+0{,}45^{2}}=\sqrt{0{,}5625}=\SI{0.75}{\newton},
\qquad
\tan\theta=\frac{0{,}45}{0{,}60}=0{,}75 \Rightarrow \theta=\SI{36.87}{\degree}.
\]
Um clássico triângulo $3$--$4$--$5$ escondido. \textbf{Atenção:} a hipotenusa
$BC=\SI{50}{\centi\meter}$ é irrelevante --- ela governa a força \emph{entre}
$q_B$ e $q_C$, não a força sobre $q_A$.}
\end{questao}

\begin{questao}[3]{}
Duas pequenas esferas metálicas idênticas, separadas por $\SI{10}{\centi\meter}$,
\textbf{atraem-se} com força de $\SI{0.108}{\newton}$. Postas em contato e
recolocadas na mesma distância, passam a \textbf{repelir-se} com
$\SI{0.036}{\newton}$. Determine as cargas iniciais.
\resposta{$q_1=\SI{+0.6}{\micro\coulomb}$ e $q_2=\SI{-0.2}{\micro\coulomb}$
(ou a permuta)}
\resolucao{\textbf{Etapa 1 --- produto.} A atração indica sinais opostos:
\[
|q_1q_2|=\frac{F_1r^{2}}{\ke}=\frac{0{,}108\cdot 0{,}01}{\pot{9}{9}}
=\pot{1.2}{-13}\ \Rightarrow\ q_1q_2=-\pot{1.2}{-13}\,\si{\coulomb\squared}.
\]
\textbf{Etapa 2 --- soma.} Após o contato cada esfera fica com
$q'=(q_1+q_2)/2$:
\[
q'=\sqrt{\frac{F_2r^{2}}{\ke}}=\sqrt{\frac{0{,}036\cdot0{,}01}{\pot{9}{9}}}
=\pot{2}{-7}\,\si{\coulomb}
\ \Rightarrow\ q_1+q_2=\pot{4}{-7}\,\si{\coulomb}.
\]
\textbf{Etapa 3 --- soma e produto.} As cargas são raízes de
$x^{2}-Sx+P=0$ com $S=\pot{4}{-7}$ e $P=-\pot{1.2}{-13}$:
\[
x=\frac{\pot{4}{-7}\pm\sqrt{\pot{1.6}{-13}+\pot{4.8}{-13}}}{2}
 =\frac{\pot{4}{-7}\pm\pot{8}{-7}}{2}
 \Rightarrow x=\pot{6}{-7}\ \text{ou}\ -\pot{2}{-7}.
\]
A repulsão final confirma que a soma é não nula. Se as cargas fossem simétricas
($q$ e $-q$), o contato as neutralizaria e a força final seria zero.}
\end{questao}

\begin{questao}[3]{}
Um pequeno bloco de massa $m=\SI{20}{\gram}$, com carga $q=\SI{1}{\micro\coulomb}$,
repousa sobre um plano inclinado \textbf{liso} de $\SI{30}{\degree}$. Uma carga fixa
$Q=\SI{1}{\micro\coulomb}$ está presa na base do plano, e a força elétrica atua
ao longo da superfície. Determine a distância de equilíbrio $d$ medida sobre o
plano.
\dados{$g=\SI{9.8}{\meter\per\second\squared}$}
\resposta{$d\approx\SI{30.3}{\centi\meter}$}
\resolucao{Sobre o plano liso, o equilíbrio na direção da superfície envolve
apenas a componente do peso e a repulsão elétrica (a normal é perpendicular):
\[
\frac{\ke Qq}{d^{2}}=mg\sin\theta.
\]
\[
mg\sin\SI{30}{\degree}=0{,}020\cdot 9{,}8\cdot 0{,}5=\SI{0.098}{\newton};
\qquad
d=\sqrt{\frac{\pot{9}{9}\cdot(\pot{1}{-6})^{2}}{0{,}098}}
=\sqrt{0{,}0918}=\SI{0.303}{\meter}.
\]
\textbf{Estabilidade:} se o bloco descer, $d$ diminui e a repulsão cresce como
$1/d^{2}$, empurrando-o de volta; se subir, a repulsão cai e o peso o traz de
volta. O equilíbrio é \emph{estável}.}
\end{questao}

\begin{questao}[3]{}
Um dipolo é formado por $+q$ e $-q$, com $q=\SI{1}{\micro\coulomb}$, fixos em
$(\pm\SI{5}{\centi\meter},0)$. Uma carga $Q=\SI{2}{\micro\coulomb}$ é colocada
em $(0,\SI{12}{\centi\meter})$, sobre a mediatriz.
\begin{itensq}
\item Mostre que a resultante sobre $Q$ é paralela ao eixo do dipolo.
\item Calcule seu módulo.
\end{itensq}
\resposta{(a) as componentes verticais se cancelam; (b) $F=\SI{0.82}{\newton}$,
no sentido $-x$ (de $+q$ para $-q$)}
\resolucao{(a) Cada carga dista
$r=\sqrt{0{,}05^{2}+0{,}12^{2}}=\SI{0.13}{\meter}$ de $Q$ (triângulo
$5$--$12$--$13$). A carga $+q$ \emph{repele} $Q$ e $-q$ a \emph{atrai}: as duas
forças têm o mesmo módulo e as componentes em $y$ apontam em sentidos opostos,
cancelando-se. Restam as componentes em $x$, que se \textbf{somam}.\\
(b) Cada força vale $F_1=\ke qQ/r^{2}$ e sua componente horizontal é
$F_1\cdot(a/r)$, com $a=\SI{0.05}{\meter}$:
\[
F=2\,\frac{\ke qQ}{r^{2}}\cdot\frac{a}{r}=\frac{2\ke qQ\,a}{r^{3}}
=\frac{2\cdot\pot{9}{9}\cdot\pot{1}{-6}\cdot\pot{2}{-6}\cdot 0{,}05}
       {(0{,}13)^{3}}=\SI{0.82}{\newton}.
\]
Note a dependência $1/r^{3}$: longe do dipolo a força cai \emph{mais rápido}
que a de uma carga isolada, porque as duas cargas quase se cancelam.}
\end{questao}

\blocodesafio

\begin{questao}[4]{}
\textbf{(Modelagem experimental.)} Em uma reprodução da balança de torção, duas
esferas com cargas iguais $q=\SI{0.10}{\micro\coulomb}$ foram mantidas a várias
distâncias e a força foi medida:

\begin{center}
\begin{tabular}{cccccc}
\hline
$r$ (m) & 0,020 & 0,030 & 0,040 & 0,050 & 0,060\\
$F$ (N) & 0,223 & 0,101 & 0,0558 & 0,0362 & 0,0249\\
\hline
\end{tabular}
\end{center}

\begin{itensq}
\item Proponha uma linearização que permita testar a hipótese $F=C\,r^{n}$.
\item Estime $n$ e discuta se os dados sustentam a lei do inverso do quadrado.
\item Obtenha o valor experimental de $\ke$ e compare com o tabelado.
\end{itensq}
\resposta{(a) $\log F=\log C+n\log r$; (b) $n\approx-2{,}00$; (c)
$\ke\approx\pot{9.1}{9}\,\si{\newton\meter\squared\per\coulomb\squared}$
(erro $\approx 1\%$)}
\resolucao{(a) Tomando logaritmos, uma lei de potência vira uma reta:
$\log F = \log C + n\log r$. O coeficiente angular do gráfico
$\log F \times \log r$ é o expoente procurado.\\
(b) Ajustando os cinco pontos por mínimos quadrados obtém-se
$n=-1{,}997$ e $C=\pot{9.09}{-5}$. O expoente reproduz $-2$ com desvio de
$0{,}15\%$ --- dentro do que se espera de medidas de força dessa ordem. Um teste
rápido sem regressão: dobrar $r$ de $0{,}020$ para $0{,}040$ deveria dividir $F$
por $4$; de fato $0{,}223/0{,}0558=4{,}00$.\\
(c) Como $C=\ke q^{2}$ e $q^{2}=\pot{1}{-14}\,\si{\coulomb\squared}$:
\[
\ke=\frac{C}{q^{2}}=\frac{\pot{9.09}{-5}}{\pot{1}{-14}}
=\pot{9.09}{9}\,\si{\newton\meter\squared\per\coulomb\squared},
\]
contra $\pot{8.99}{9}$ tabelado --- erro relativo de $1{,}1\%$.\\
\textbf{Observação de método:} a linearização logarítmica é preferível a ajustar
$F\times 1/r^{2}$ diretamente, porque \emph{não pressupõe} o expoente --- ela o
mede. Esse é exatamente o argumento que Coulomb precisou construir em 1785.}
\end{questao}

\begin{questao}[4]{}
\textbf{(Otimização com cálculo.)} Um anel fino de raio $R$, com carga total $Q$
uniformemente distribuída, está no plano $yz$. Uma carga puntiforme $q$ é
colocada sobre o eixo do anel, a uma distância $x$ do centro.
\begin{itensq}
\item Mostre que $F(x)=\dfrac{\ke Q q\,x}{(x^{2}+R^{2})^{3/2}}$.
\item Determine o valor de $x$ que \textbf{maximiza} a força.
\item Calcule $F_{\max}$ para $Q=\SI{10}{\micro\coulomb}$,
      $q=\SI{2}{\micro\coulomb}$ e $R=\SI{10}{\centi\meter}$.
\end{itensq}
\resposta{(b) $x=R/\sqrt{2}\approx\SI{7.07}{\centi\meter}$;
(c) $F_{\max}=\dfrac{2\ke Qq}{3\sqrt{3}\,R^{2}}=\SI{6.93}{\newton}$}
\resolucao{(a) Divida o anel em elementos $\mathrm{d} Q$. Cada um dista
$s=\sqrt{x^{2}+R^{2}}$ de $q$ e contribui com $\ke q\,\mathrm{d} Q/s^{2}$. Por
simetria, as componentes perpendiculares ao eixo se cancelam aos pares; sobra a
projeção axial, fator $\cos\alpha=x/s$:
\[
F=\int \frac{\ke q\,\mathrm{d} Q}{s^{2}}\cdot\frac{x}{s}
 =\frac{\ke q x}{(x^{2}+R^{2})^{3/2}}\int \mathrm{d} Q
 =\frac{\ke Q q\,x}{(x^{2}+R^{2})^{3/2}}.
\]
(b) Derivando e igualando a zero:
\[
\frac{\mathrm{d} F}{\mathrm{d} x}\propto (x^{2}+R^{2})^{3/2}-x\cdot 3x(x^{2}+R^{2})^{1/2}=0
\;\Rightarrow\; x^{2}+R^{2}-3x^{2}=0 \;\Rightarrow\; x=\frac{R}{\sqrt{2}}.
\]
(c) Substituindo $x=R/\sqrt2$:
\[
F_{\max}=\frac{2\,\ke Qq}{3\sqrt{3}\,R^{2}}
=\frac{2\cdot\pot{9}{9}\cdot\pot{1}{-5}\cdot\pot{2}{-6}}{3\sqrt3\cdot 0{,}01}
=\frac{0{,}36}{0{,}0520}=\SI{6.93}{\newton}.
\]
\textbf{Casos-limite:} em $x=0$ a força é nula (simetria perfeita); para
$x\gg R$, $F\to\ke Qq/x^{2}$ --- o anel se comporta como carga puntiforme.
Entre esses dois zeros deve existir um máximo, e ele está em $R/\sqrt2$.}
\end{questao}

\begin{questao}[4]{}
\textbf{(Distribuição contínua.)} Um bastão isolante retilíneo de comprimento
$L$ possui carga $Q$ uniformemente distribuída. Uma carga puntiforme $q$ está
sobre o prolongamento do bastão, a uma distância $d$ da extremidade mais
próxima.
\begin{itensq}
\item Deduza $F=\dfrac{\ke Qq}{d\,(d+L)}$.
\item Verifique o comportamento para $L\ll d$.
\item Calcule $F$ para $Q=\SI{4}{\micro\coulomb}$, $L=\SI{20}{\centi\meter}$,
      $q=\SI{1}{\micro\coulomb}$ e $d=\SI{10}{\centi\meter}$.
\end{itensq}
\resposta{(b) $F\to\ke Qq/d^{2}$; (c) $F=\SI{1.2}{\newton}$}
\resolucao{(a) Densidade linear $\lambda=Q/L$. Um elemento $\mathrm{d} x$ a uma
distância $x$ da carga carrega $\mathrm{d} Q=\lambda\,\mathrm{d} x$ e todas as contribuições
são \emph{colineares} --- por isso a integral é escalar:
\[
F=\int_{d}^{d+L}\frac{\ke q\lambda\,\mathrm{d} x}{x^{2}}
 =\ke q\lambda\left[-\frac{1}{x}\right]_{d}^{d+L}
 =\ke q\lambda\left(\frac{1}{d}-\frac{1}{d+L}\right)
 =\frac{\ke q\lambda L}{d(d+L)}=\frac{\ke Qq}{d(d+L)}.
\]
(b) Se $L\ll d$, então $d+L\approx d$ e $F\to\ke Qq/d^{2}$: recuperamos a lei de
Coulomb puntiforme, como deve ser.\\
(c) $F=\dfrac{\pot{9}{9}\cdot\pot{4}{-6}\cdot\pot{1}{-6}}{0{,}10\cdot 0{,}30}
=\dfrac{0{,}036}{0{,}03}=\SI{1.2}{\newton}$.\\
\textbf{Compare:} se toda a carga estivesse concentrada no centro do bastão
($\SI{20}{\centi\meter}$), teríamos $\SI{0.9}{\newton}$. A força real é maior
porque as porções mais próximas pesam com $1/x^{2}$ --- concentrar carga no
centro geométrico \emph{subestima} a força.}
\end{questao}

\begin{questao}[4]{}
\textbf{(Demonstração --- estabilidade.)} Duas cargas $+Q$ estão fixas em
$x=-d$ e $x=+d$. Uma carga $q$ é colocada na origem.
\begin{itensq}
\item Mostre que a origem é posição de equilíbrio para qualquer $q$.
\item Para $q>0$, analise a estabilidade a um pequeno deslocamento
      \emph{longitudinal} ($\pm x$).
\item Repita para um deslocamento \emph{transversal} ($\pm y$) e conclua.
\end{itensq}
\resposta{(a) por simetria; (b) estável, com $F\approx-4\ke Qq\,x/d^{3}$;
(c) instável, com $F\approx+2\ke Qq\,y/d^{3}$ --- não há equilíbrio estável em
todas as direções}
\resolucao{(a) As duas forças têm o mesmo módulo e sentidos opostos: resultante
nula, independentemente do sinal e do valor de $q$.\\
(b) Deslocando $q>0$ para $x$ pequeno:
\[
F_x=\ke Qq\left[\frac{1}{(d+x)^{2}}-\frac{1}{(d-x)^{2}}\right].
\]
Com $(d\pm x)^{-2}\approx d^{-2}(1\mp 2x/d)$:
\[
F_x\approx\frac{\ke Qq}{d^{2}}\left(-\frac{4x}{d}\right)=-\frac{4\ke Qq}{d^{3}}\,x.
\]
Sinal negativo $\Rightarrow$ força \textbf{restauradora}: estável nessa direção
(e o movimento é harmônico simples).\\
(c) Deslocando para $y$ pequeno, cada carga dista $\sqrt{d^{2}+y^{2}}\approx d$ e
repele $q$ com componente transversal $\ke Qq\,y/d^{3}$; as duas se somam:
\[
F_y\approx +\frac{2\ke Qq}{d^{3}}\,y,
\]
sinal positivo $\Rightarrow$ força que \textbf{afasta}: instável.\\
\textbf{Conclusão.} O equilíbrio é do tipo \emph{ponto de sela}: estável em uma
direção, instável na outra. Este é um caso particular do \textbf{teorema de
Earnshaw}: nenhuma configuração de cargas fixas pode manter uma carga puntiforme
em equilíbrio estável apenas por forças eletrostáticas. É por isso que modelos
puramente eletrostáticos do átomo não se sustentam.}
\end{questao}

\begin{questao}[4]{}
\textbf{(Série infinita.)} Cargas puntiformes de mesmo módulo $Q$ e
\textbf{sinais alternados} ocupam as posições $x=a,\,2a,\,3a,\dots$ sobre um
eixo: $-Q$ em $a$, $+Q$ em $2a$, $-Q$ em $3a$ e assim por diante,
indefinidamente. Uma carga $+q$ está na origem.
\begin{itensq}
\item Escreva a força resultante sobre $q$ como uma série.
\item Some a série e obtenha a expressão fechada.
\end{itensq}
\dados{$\displaystyle\sum_{n=1}^{\infty}\frac{(-1)^{n+1}}{n^{2}}=\frac{\pi^{2}}{12}$}
\resposta{$F=\dfrac{\pi^{2}}{12}\cdot\dfrac{\ke Qq}{a^{2}}
\approx 0{,}822\,\dfrac{\ke Qq}{a^{2}}$, atrativa (sentido $+x$)}
\resolucao{(a) A $n$-ésima carga está em $x=na$ e tem sinal $(-1)^{n}Q$. Uma
carga negativa \emph{atrai} $q$ (puxa no sentido $+x$); uma positiva
\emph{repele} (sentido $-x$). Tomando $+x$ como positivo:
\[
F=\frac{\ke Qq}{a^{2}}\left(\frac{1}{1^{2}}-\frac{1}{2^{2}}+\frac{1}{3^{2}}
-\frac{1}{4^{2}}+\cdots\right)
=\frac{\ke Qq}{a^{2}}\sum_{n=1}^{\infty}\frac{(-1)^{n+1}}{n^{2}}.
\]
(b) Usando o resultado dado:
\[
F=\frac{\pi^{2}}{12}\cdot\frac{\ke Qq}{a^{2}}\approx 0{,}822\,\frac{\ke Qq}{a^{2}}.
\]
\textbf{Interpretação.} Apesar de haver infinitas cargas, a força é
\emph{finita} e vale menos que a da primeira carga sozinha: a alternância de
sinais produz cancelamento progressivo, e o decaimento $1/n^{2}$ garante
convergência absoluta. Se todas as cargas fossem negativas, a soma seria
$\pi^2/6\approx1{,}645$ --- exatamente o dobro.}
\end{questao}

\begin{questao}[4]{}
\textbf{(Oscilações.)} Uma partícula de massa $m=\SI{1.0}{\gram}$ e carga
$-q=\SI{-2}{\micro\coulomb}$ é liberada em repouso sobre o eixo de um anel de
raio $R=\SI{10}{\centi\meter}$ e carga $Q=\SI{5}{\micro\coulomb}$, a uma
distância $x_0\ll R$ do centro.
\begin{itensq}
\item Mostre que, para $x\ll R$, o movimento é harmônico simples.
\item Obtenha $\omega$ e o período $T$.
\end{itensq}
\resposta{(b) $\omega=\sqrt{\ke Qq/(mR^{3})}=\SI{300}{\radian\per\second}$;
$T\approx\SI{20.9}{\milli\second}$}
\resolucao{(a) A força axial sobre a carga negativa é atrativa, dirigida ao
centro, de módulo $F=\ke Qq\,x/(x^{2}+R^{2})^{3/2}$. Para $x\ll R$, o
denominador tende a $R^{3}$:
\[
F_x\approx-\frac{\ke Qq}{R^{3}}\,x.
\]
Força proporcional ao deslocamento e de sinal oposto --- exatamente a forma
$F=-kx$ da lei de Hooke, com constante efetiva $k_{\text{ef}}=\ke Qq/R^{3}$.
Logo o MHS.\\
(b)
\[
\omega=\sqrt{\frac{k_{\text{ef}}}{m}}=\sqrt{\frac{\ke Qq}{mR^{3}}}
=\sqrt{\frac{\pot{9}{9}\cdot\pot{5}{-6}\cdot\pot{2}{-6}}
{\pot{1}{-3}\cdot(0{,}10)^{3}}}
=\sqrt{\frac{0{,}09}{\pot{1}{-6}}}=\SI{300}{\radian\per\second},
\]
\[
T=\frac{2\pi}{\omega}=\SI{20.9}{\milli\second}.
\]
\textbf{Observe:} o período não depende de $x_0$ --- é a assinatura do MHS. Fora
da aproximação $x\ll R$, a força deixa de ser linear e o movimento continua
periódico, mas com período dependente da amplitude.}
\end{questao}

\begin{questao}[4]{}
\textbf{(Modelagem com taxas.)} Duas esferas idênticas de massa
$m=\SI{2.0}{\gram}$, cada uma com carga $q$, pendem de fios isolantes de
comprimento $L=\SI{60}{\centi\meter}$ presos ao mesmo ponto. Elas se afastam até
uma separação $x\ll L$. Por fuga de carga para o ar úmido, as esferas se
aproximam lentamente.
\begin{itensq}
\item Mostre que $x^{3}=\dfrac{2\ke q^{2}L}{mg}$.
\item Demonstre que
      $\dfrac{\mathrm{d} x}{\mathrm{d} t}=\dfrac{2x}{3q}\dfrac{\mathrm{d} q}{\mathrm{d} t}$.
\item Se em certo instante $x=\SI{6.0}{\centi\meter}$ e as esferas se aproximam
      a $\SI{1.0}{\milli\meter\per\second}$, calcule $\mathrm{d} q/\mathrm{d} t$.
\end{itensq}
\dados{$g=\SI{9.8}{\meter\per\second\squared}$}
\resposta{(c) $q\approx\SI{19.8}{\nano\coulomb}$ e
$\mathrm{d} q/\mathrm{d} t\approx\SI{-0.49}{\nano\coulomb\per\second}$}
\resolucao{(a) Para cada esfera: tensão, peso e repulsão. Do triângulo de
forças, $\tan\theta=F/(mg)$. Com $x\ll L$,
$\tan\theta\approx\sin\theta=(x/2)/L$. Igualando:
\[
\frac{x}{2L}=\frac{\ke q^{2}}{x^{2}mg}
\;\Longrightarrow\;
x^{3}=\frac{2\ke q^{2}L}{mg}.
\]
(b) Derivando implicitamente em relação ao tempo:
\[
3x^{2}\frac{\mathrm{d} x}{\mathrm{d} t}=\frac{4\ke L}{mg}\,q\,\frac{\mathrm{d} q}{\mathrm{d} t}.
\]
De (a), $\dfrac{2\ke L}{mg}=\dfrac{x^{3}}{q^{2}}$; substituindo:
\[
3x^{2}\frac{\mathrm{d} x}{\mathrm{d} t}=\frac{2x^{3}}{q}\frac{\mathrm{d} q}{\mathrm{d} t}
\;\Longrightarrow\;
\frac{\mathrm{d} x}{\mathrm{d} t}=\frac{2x}{3q}\frac{\mathrm{d} q}{\mathrm{d} t}.
\]
(c) Primeiro $q$, por (a):
\[
q=\sqrt{\frac{x^{3}mg}{2\ke L}}
=\sqrt{\frac{(0{,}060)^{3}\cdot 0{,}0020\cdot 9{,}8}{2\cdot\pot{9}{9}\cdot0{,}60}}
=\pot{1.98}{-8}\,\si{\coulomb}.
\]
Com $\mathrm{d} x/\mathrm{d} t=-\pot{1.0}{-3}\,\si{\meter\per\second}$ (aproximação):
\[
\frac{\mathrm{d} q}{\mathrm{d} t}=\frac{3q}{2x}\frac{\mathrm{d} x}{\mathrm{d} t}
=\frac{3\cdot\pot{1.98}{-8}}{2\cdot 0{,}060}\cdot(-\pot{1}{-3})
=-\pot{4.9}{-10}\,\si{\coulomb\per\second}.
\]
\textbf{Leitura física:} medir uma \emph{velocidade} converte-se em medir uma
\emph{corrente de fuga} de meio nanoampère --- pequena demais para muitos
amperímetros, mas visível a olho nu pelo movimento das esferas.}
\end{questao}

\begin{questao}[4]{}
\textbf{(Simetria e integração.)} Um fio semicircular de raio
$R=\SI{5.0}{\centi\meter}$ possui carga $Q=\SI{3}{\micro\coulomb}$ distribuída
uniformemente. Uma carga $q=\SI{1}{\micro\coulomb}$ ocupa o centro do
semicírculo.
\begin{itensq}
\item Justifique por que a resultante é dirigida ao longo do eixo de simetria.
\item Mostre que $F=\dfrac{2\ke Qq}{\pi R^{2}}$ e calcule seu valor.
\end{itensq}
\resposta{$F=\SI{6.9}{\newton}$, ao longo do eixo de simetria}
\resolucao{(a) Para cada elemento do fio existe outro simétrico em relação ao
eixo vertical. Suas contribuições têm componentes perpendiculares ao eixo iguais
e opostas --- que se cancelam --- e componentes ao longo do eixo que se somam.\\
(b) Densidade linear $\lambda=Q/(\pi R)$. O elemento em $\theta$ tem
$\mathrm{d} Q=\lambda R\,\mathrm{d}\theta$, dista $R$ do centro e contribui com
$\ke q\,\mathrm{d} Q/R^{2}$; sua projeção no eixo vale $\sin\theta$:
\[
F=\int_{0}^{\pi}\frac{\ke q\lambda R\,\mathrm{d}\theta}{R^{2}}\sin\theta
 =\frac{\ke q\lambda}{R}\underbrace{\int_{0}^{\pi}\sin\theta\,\mathrm{d}\theta}_{=2}
 =\frac{2\ke q\lambda}{R}=\frac{2\ke Qq}{\pi R^{2}}.
\]
\[
F=\frac{2\cdot\pot{9}{9}\cdot\pot{3}{-6}\cdot\pot{1}{-6}}
{\pi\,(0{,}050)^{2}}=\frac{0{,}054}{\pot{7.85}{-3}}=\SI{6.9}{\newton}.
\]
\textbf{Compare:} se a mesma carga $Q$ estivesse concentrada em um ponto a $R$,
teríamos $\ke Qq/R^{2}=\SI{10.8}{\newton}$. O fator $2/\pi\approx 0{,}64$ mede
exatamente a perda por cancelamento vetorial ao longo do arco. Um anel
\emph{completo} levaria esse fator a zero.}
\end{questao}

\begin{questao}[4]{}
\textbf{(Síntese --- equilíbrio de um sistema.)} Quatro cargas iguais
$Q=\SI{2}{\micro\coulomb}$ ocupam os vértices de um quadrado de lado $a$. Uma
quinta carga $q$ é colocada no centro.
\begin{itensq}
\item Determine $q$ (sinal e módulo, em função de $Q$) para que \textbf{cada}
      carga dos vértices fique em equilíbrio.
\item Calcule $q$ numericamente.
\item O sistema completo está em equilíbrio estável? Justifique.
\end{itensq}
\resposta{(a) $q=-\dfrac{Q(2\sqrt{2}+1)}{4}$; (b) $q\approx\SI{-1.91}{\micro\coulomb}$;
(c) não --- o equilíbrio é instável}
\resolucao{(a) Por simetria, basta impor o equilíbrio de \emph{um} vértice; a
resultante das outras três cargas aponta ao longo da diagonal, para fora.
\begin{itemize}
\item Duas vizinhas, a distância $a$: cada uma contribui $\ke Q^{2}/a^{2}$, e
suas projeções sobre a diagonal valem $\cos\SI{45}{\degree}=\sqrt2/2$ cada, somando
$\sqrt{2}\,\ke Q^{2}/a^{2}$.
\item A oposta, a distância $a\sqrt{2}$: contribui
$\ke Q^{2}/(2a^{2})$, já ao longo da diagonal.
\end{itemize}
\[
F_{\text{vértices}}=\frac{\ke Q^{2}}{a^{2}}\left(\sqrt{2}+\frac{1}{2}\right).
\]
A carga central dista $a/\sqrt{2}$ do vértice, logo atua com
$F_{c}=\ke|q|Q/(a^{2}/2)=2\ke|q|Q/a^{2}$ e deve ser \textbf{atrativa} ---
portanto $q<0$. Igualando:
\[
2|q|=Q\left(\sqrt2+\frac12\right)
\;\Longrightarrow\;
|q|=\frac{Q\,(2\sqrt{2}+1)}{4}\approx 0{,}957\,Q.
\]
(b) $q\approx-0{,}957\cdot 2=\SI{-1.91}{\micro\coulomb}$.\\
(c) \textbf{Não.} A condição imposta zera a resultante sobre os vértices, mas
não sobre a carga central --- que já está em equilíbrio por simetria, embora
instável. Mais fundamentalmente, o teorema de Earnshaw proíbe equilíbrio
eletrostático estável: qualquer perturbação de $q$ ao longo de uma diagonal a
faz cair sobre um vértice. O equilíbrio existe, mas não sobrevive a uma
perturbação infinitesimal.}
\end{questao}

% END BANCO COMPLEMENTAR

\gabarito[1.2 Lei de Coulomb]

% =====================================================================
\subsection{Campo elétrico}
% =====================================================================

\begin{conceito}[Antes de começar]
\textbf{Campo de uma carga puntiforme:}
\[
E = \ke\,\frac{|q|}{d^2}\quad(\si{\newton\per\coulomb}\ \text{ou}\ \si{\volt\per\meter}).
\]
O campo \emph{aponta para fora} de cargas positivas e \emph{para dentro} de cargas
negativas.\par
\textbf{Força sobre uma carga no campo:} $F = qE$ (mesmo sentido de $E$ se $q>0$;
oposto se $q<0$).\par
\textbf{Superposição:} o campo resultante é a \emph{soma vetorial} dos campos
individuais.\par
\textbf{Ponto de campo nulo} entre duas cargas de mesmo sinal: onde os módulos se
igualam, $\dfrac{|q_1|}{d_1^2} = \dfrac{|q_2|}{d_2^2}$.
\end{conceito}

\legendaniveis

\blocofixacao

\begin{questao}[1]{}
Qual a intensidade do campo elétrico gerado por uma carga de
$\SI{-15}{\micro\coulomb}$, a $\SI{10}{\milli\meter}$ dela, em um meio de constante
$\ke = \SI{8e9}{\newton\meter\squared\per\coulomb\squared}$?
\resposta{$E = \pot{1.2}{9}\,\si{\newton\per\coulomb}$}
\resolucao{Com $d = \SI{10}{\milli\meter} = \SI{0.01}{\meter}$:
\[
E=\ke\frac{|q|}{d^2}
=8\cdot10^9\cdot\frac{\pot{15}{-6}}{(0{,}01)^2}=\pot{1.2}{9}\,\si{\newton\per\coulomb}.
\]
O campo aponta \emph{para} a carga, por ela ser negativa.}
\end{questao}

\begin{questao}[1]{}
O módulo do campo elétrico de uma carga puntiforme $q$, em um ponto P a uma dada
distância, vale $E$. Afastando-se a carga de modo que a distância a P
\textbf{dobre}, o campo em P passa a ser:
\begin{alternativas}[3]
\item $E$
\item $E/2$
\item $E/4$
\item $2E$
\item $4E$
\end{alternativas}
\resposta{(c)}
\resolucao{Como $E \propto \dfrac{1}{d^2}$, dobrar a distância divide o campo por
$2^2 = 4$: o novo valor é $E/4$.}
\end{questao}

\blocoaplicacao

\begin{questao}[2]{}
Uma carga de prova de $\SI{e-9}{\coulomb}$, colocada em um ponto P de um campo
elétrico, fica sujeita a uma força de $\SI{e-2}{\newton}$, vertical e para baixo.
Determine:
\begin{itensq}
\item a intensidade, direção e sentido do campo em P;
\item a força que atuaria sobre uma carga de $\SI{3}{\milli\coulomb}$ colocada em P.
\end{itensq}
\resposta{(a) $E=\pot{1.0}{7}\,\si{\newton\per\coulomb}$, vertical para baixo;
(b) $F = \pot{3}{4}\,\si{\newton}$}
\resolucao{(a) $E = \dfrac{F}{q} = \dfrac{\pot{1}{-2}}{\pot{1}{-9}}
= \pot{1.0}{7}\,\si{\newton\per\coulomb}$. Como a carga de prova é positiva, o
campo tem o \emph{mesmo} sentido da força: vertical, para baixo.\\
(b) $F = qE = \pot{3}{-3}\cdot\pot{1}{7} = \pot{3}{4}\,\si{\newton}$,
também vertical para baixo (carga positiva).}
\end{questao}

\begin{questao}[2]{}
Uma carga positiva $Q = \SI{4.0}{\micro\coulomb}$ está no ar. Determine a
intensidade, a direção e o sentido do campo elétrico em um ponto M, a
$\SI{20}{\centi\meter}$ da carga.
\dados{$\ke = \SI{9e9}{\newton\meter\squared\per\coulomb\squared}$.}
\resposta{$E = \pot{9.0}{5}\,\si{\newton\per\coulomb}$, radial, afastando-se de $Q$}
\resolucao{$E = \ke\dfrac{Q}{d^2}
= 9\cdot10^9\cdot\dfrac{\pot{4}{-6}}{(0{,}2)^2} = \pot{9.0}{5}\,\si{\newton\per\coulomb}$.
A direção é a da reta que une $Q$ a M, e o sentido é \emph{para fora} de $Q$
(carga positiva).}
\end{questao}

\begin{questao}[2]{}
Duas cargas puntiformes, $Q$ e $4Q$ (de mesmo sinal), estão fixas no vácuo. O
módulo do campo elétrico resultante é \textbf{nulo} em um ponto sobre a reta que
as une. Esse ponto está tal que sua distância à carga menor, $d_1$, e à carga
maior, $d_2$, satisfazem:
\begin{alternativas}[2]
\item $d_2 = d_1$
\item $d_2 = 2d_1$
\item $d_2 = 4d_1$
\item $d_1 = 2d_2$
\item $d_1 = 4d_2$
\end{alternativas}
\resposta{(b)}
\resolucao{O campo se anula \emph{entre} as cargas, onde os módulos se igualam:
\[
\ke\frac{Q}{d_1^2}=\ke\frac{4Q}{d_2^2}
\ \Rightarrow\
\frac{d_2^2}{d_1^2}=4
\ \Rightarrow\
d_2 = 2d_1.
\]
O ponto de campo nulo fica \emph{mais próximo} da carga menor.}
\end{questao}

\begin{questao}[2]{}
Duas cargas de mesmo módulo $Q$ estão fixas, separadas por $\SI{20}{\centi\meter}$.
Analise o campo elétrico resultante no ponto médio M do segmento que as une, em
dois casos:

\circuitoq{%
\begin{tikzpicture}[scale=0.9,line width=0.8pt,>=Stealth]
 % caso 1: +Q e -Q
 \begin{scope}
   \shade[ball color=Vcol!25] (-1.6,0) circle (0.25); \node at (-1.6,0){\footnotesize$+$};
   \shade[ball color=Icol!25] (1.6,0) circle (0.25);   \node at (1.6,0){\footnotesize$-$};
   \fill[laranja] (0,0) circle (0.06); \node[above,font=\footnotesize] at (0,0.12){M};
   \draw[Vcol,->,thick] (0,-0.5)--(1.1,-0.5); \node[Vcol,font=\footnotesize] at (0.55,-0.8){$\vec E_1$};
   \draw[Vcol,->,thick] (0,-1.1)--(1.1,-1.1); \node[Vcol,font=\footnotesize] at (0.55,-1.4){$\vec E_2$};
   \node[font=\footnotesize] at (0,1.0){Caso I: $+Q$ e $-Q$};
 \end{scope}
\end{tikzpicture}}{Campo no ponto médio entre duas cargas.}

No \textbf{Caso I} ($+Q$ e $-Q$) e no \textbf{Caso II} (ambas $+Q$), o campo
resultante em M é, respectivamente:
\begin{alternativas}
\item nulo em I e nulo em II.
\item não-nulo em I e nulo em II.
\item nulo em I e não-nulo em II.
\item não-nulo em ambos.
\item nulo em ambos.
\end{alternativas}
\resposta{(b)}
\resolucao{Em M, cada carga produz um campo de mesmo módulo (mesma distância,
mesmo $|Q|$).\\
\textbf{Caso I ($+Q$, $-Q$):} os dois campos apontam no \emph{mesmo} sentido (para
longe da positiva e em direção à negativa), então se \textbf{somam}: campo
resultante \emph{não-nulo}, $E = 2\cdot\dfrac{\ke Q}{d^2}$.\\
\textbf{Caso II ($+Q$, $+Q$):} os campos apontam em sentidos \emph{opostos} e se
\textbf{cancelam}: campo resultante \emph{nulo}.\\
Cuidado: para o \emph{potencial} a conclusão se inverte (nulo no caso I, não-nulo
no caso II), porque potencial é escalar e campo é vetorial.}
\end{questao}

\blocoaprofundamento

\begin{questao}[3]{}
As cargas $q_1 = \SI{20}{\coulomb}$ e $q_2 = \SI{64}{\coulomb}$ estão fixas no
vácuo nos pontos A e B, separados por $\SI{9}{\meter}$. Determine a posição, sobre
o segmento AB, onde o campo elétrico resultante é nulo.
\dados{$\ke = \SI{9e9}{\newton\meter\squared\per\coulomb\squared}$.}
\resposta{A $\SI{4}{\meter}$ de A (e $\SI{5}{\meter}$ de B)}
\resolucao{Como as cargas têm o mesmo sinal, o campo se anula entre elas. Seja $x$
a distância a partir de A; a distância a B é $9 - x$. Igualando os módulos:
\[
\ke\frac{q_1}{x^2}=\ke\frac{q_2}{(9-x)^2}
\ \Rightarrow\
\frac{20}{x^2}=\frac{64}{(9-x)^2}.
\]
Extraindo a raiz (grandezas positivas):
$\dfrac{9-x}{x}=\sqrt{\dfrac{64}{20}}=\dfrac{8}{\sqrt{20}}=\dfrac{8}{2\sqrt5}
=\dfrac{4}{\sqrt5}$. Resolvendo,
$9 - x = \dfrac{4}{\sqrt5}x \Rightarrow x = \dfrac{9}{1+4/\sqrt5}\approx\SI{4}{\meter}$.
O ponto neutro fica a cerca de \SI{4}{\meter} de A e \SI{5}{\meter} de B --- mais
perto da carga menor, como sempre ocorre.}
\end{questao}

\begin{questao}[3]{}
A intensidade do campo elétrico de uma carga puntiforme positiva, em função da
distância $d$, obedece ao gráfico $E \times d$ abaixo (curva do tipo $E = c/d^2$).
Sabendo que a $\SI{3}{\meter}$ o campo vale $\SI{80}{\newton\per\coulomb}$,
determine o campo a $\SI{6}{\meter}$.

\circuitoqq{%
\begin{tikzpicture}
 \begin{axis}[width=9cm,height=5cm,xlabel={$d$ (m)},ylabel={$E$ (N/C)},
   xmin=0,xmax=8,ymin=0,ymax=90,xtick={0,2,3,4,6,8},ytick={0,20,40,60,80},
   grid=both,axis lines=left,domain=1.6:8,samples=100]
   \addplot[Vcol,very thick] {720/(x^2)};
   \addplot[only marks,mark=*,Vcol] coordinates {(3,80)(6,20)};
   \node[Vcol,anchor=south west,font=\footnotesize] at (axis cs:3,80) {$(3,\,80)$};
   \node[Vcol,anchor=south west,font=\footnotesize] at (axis cs:6,20) {$(6,\,20)$};
 \end{axis}
\end{tikzpicture}}{Campo elétrico em função da distância ($E \propto 1/d^2$).}

\resposta{$E = \SI{20}{\newton\per\coulomb}$}
\resolucao{Como $E \propto \dfrac{1}{d^2}$, dobrar a distância (de \SI{3}{} para
\SI{6}{\meter}) divide o campo por 4:
\[
E(6)=\frac{E(3)}{4}=\frac{80}{4}=\SI{20}{\newton\per\coulomb}.
\]
Formalmente, $E(3)\cdot3^2 = E(6)\cdot6^2$, ou seja,
$80\cdot9 = E(6)\cdot36 \Rightarrow E(6) = \SI{20}{\newton\per\coulomb}$.}
\end{questao}

\begin{questao}[3]{}
Duas cargas $q_1 = \SI{+9}{\micro\coulomb}$ e $q_2 = \SI{+4}{\micro\coulomb}$
estão fixas e separadas por $\SI{10}{\centi\meter}$. Determine a que distância da
carga $q_1$, sobre a reta que as une, o campo elétrico resultante é nulo.
\resposta{A $\SI{6}{\centi\meter}$ de $q_1$}
\resolucao{Cargas de mesmo sinal: o campo se anula \emph{entre} elas. Seja $x$ a
distância a $q_1$; a distância a $q_2$ é $(10 - x)$. Igualando os módulos dos
campos:
\[
\frac{\ke q_1}{x^2}=\frac{\ke q_2}{(10-x)^2}
\ \Rightarrow\
\frac{9}{x^2}=\frac{4}{(10-x)^2}.
\]
Extraindo a raiz: $\dfrac{3}{x}=\dfrac{2}{10-x} \Rightarrow 3(10-x)=2x
\Rightarrow 30 = 5x \Rightarrow x = \SI{6}{\centi\meter}$.
O ponto neutro fica mais perto da carga \emph{menor} ($q_2$), a
\SI{4}{\centi\meter} dela.}
\end{questao}

\begin{questao}[3]{}
Duas cargas puntiformes, $+4Q$ e $-Q$, estão fixas no vácuo, separadas por uma
distância $d$.
\begin{itensq}
\item Mostre que existe um ponto sobre a reta que as une onde o campo elétrico
      resultante é \textbf{nulo}, e determine sua posição.
\item Explique por que esse ponto está \emph{fora} do segmento que une as cargas,
      e do lado da carga de \emph{menor} módulo.
\item Nesse mesmo ponto, o potencial elétrico é nulo? Justifique.
\end{itensq}
\resposta{(a) a $2d$ da carga $+4Q$ (isto é, a $d$ além de $-Q$);
(b) e (c) ver comentário}
\resolucao{\textbf{(a)} Seja $x$ a distância medida a partir de $+4Q$, com o ponto
além de $-Q$ (logo a distância a $-Q$ é $x-d$). Igualando os módulos:
\[
\frac{\ke\,4Q}{x^{2}}=\frac{\ke\,Q}{(x-d)^{2}}
\;\Longrightarrow\;
4(x-d)^2 = x^2
\;\Longrightarrow\;
2(x-d)=\pm x.
\]
A raiz fisicamente válida é $2x-2d = x \Rightarrow \boxed{x = 2d}$, ou seja, a
$2d$ da carga maior e a $d$ além da carga menor.
(A outra raiz, $x = \tfrac{2d}{3}$, cai \emph{entre} as cargas, onde os campos
apontam no \emph{mesmo} sentido e portanto não podem se cancelar --- deve ser
descartada.)

\textbf{(b)} Entre cargas de sinais \emph{opostos}, os dois vetores campo apontam
no mesmo sentido (saindo da positiva, entrando na negativa) e se \emph{somam};
nunca se anulam ali. Fora do segmento, eles têm sentidos contrários, e o
cancelamento só é possível onde o campo da carga \emph{menor} --- que está mais
próxima --- consegue igualar o da maior. Por isso o ponto neutro fica sempre do
lado da carga de menor módulo.

\textbf{(c) Não.} Potencial é grandeza \emph{escalar}:
\[
V = \ke\frac{4Q}{2d}+\ke\frac{(-Q)}{d}
= \frac{2\ke Q}{d}-\frac{\ke Q}{d}
= \frac{\ke Q}{d}\neq 0.
\]
\textbf{Conclusão conceitual:} campo nulo e potencial nulo são condições
\emph{independentes}. Neste caso, $\vec E = 0$ mas $V \neq 0$. No ponto médio de um
dipolo simétrico ocorre o oposto: $V = 0$ mas $\vec E \neq 0$.}
\end{questao}

\begin{questao}[3]{}
Um elétron entra com velocidade horizontal $v_0=\SI{2e7}{\meter\per\second}$ na
região entre duas placas paralelas de comprimento $L=\SI{5}{\centi\meter}$, onde
existe campo elétrico uniforme $E=\SI{e4}{\volt\per\meter}$, perpendicular à
velocidade.

\circuitoq{%
\begin{tikzpicture}[scale=1,line width=0.8pt,>=Stealth]
 \draw[cinza!70,line width=2.2pt] (0,1.3)--(4.2,1.3);
 \draw[cinza!70,line width=2.2pt] (0,-1.3)--(4.2,-1.3);
 \node[above,font=\footnotesize] at (2.1,1.35) {$+$\ placa};
 \node[below,font=\footnotesize] at (2.1,-1.35) {$-$\ placa};
 \foreach \x in {0.6,1.5,2.4,3.3}
   \draw[->,Vcol!70] (\x,1.2)--(\x,-1.2);
 \node[Vcol,font=\footnotesize] at (3.85,0.55) {$\vec E$};
 \draw[->,laranja,line width=1.2pt] (-1.4,0)--(0,0)
       node[midway,above,font=\footnotesize]{$v_0$};
 \draw[laranja,line width=1.2pt,domain=0:4.2,samples=40,smooth]
       plot (\x,{0.055*\x*\x});
 \draw[->,laranja,line width=1.2pt] (4.2,{0.055*4.2*4.2})--(5.3,{0.055*4.2*4.2+0.55});
 \draw[<->,cinza] (0,-1.75)--(4.2,-1.75) node[midway,below,font=\footnotesize]{$L$};
 \draw[<->,cinza] (4.5,0)--(4.5,{0.055*4.2*4.2}) node[midway,right,font=\footnotesize]{$y$};
\end{tikzpicture}}{Deflexão de um elétron em campo elétrico uniforme.}

\begin{itensq}
\item Calcule a aceleração do elétron.
\item Determine o tempo de permanência entre as placas e o desvio vertical $y$.
\item Calcule o ângulo de saída em relação à direção inicial.
\end{itensq}
\dados{$e=\SI{1.6e-19}{\coulomb}$; $m_e=\SI{9.11e-31}{\kilogram}$.}
\resposta{(a) $a\approx\SI{1.76e15}{\meter\per\second\squared}$;
(b) $t=\SI{2.5}{\nano\second}$, $y\approx\SI{5.5}{\milli\meter}$;
(c) $\theta\approx\SI{12.4}{\degree}$}
\resolucao{Este é um \textbf{lançamento de projétil} eletrostático: movimento
uniforme na horizontal e uniformemente acelerado na vertical.\\
(a) $a=\dfrac{eE}{m_e}=\dfrac{(\pot{1.6}{-19})(\pot{1}{4})}{\pot{9.11}{-31}}
\approx\pot{1.76}{15}\,\si{\meter\per\second\squared}$
--- cerca de $10^{14}$ vezes a gravidade, o que justifica desprezar o peso.\\
(b) Na horizontal (MU):
$t=\dfrac{L}{v_0}=\dfrac{0{,}05}{\pot{2}{7}}=\pot{2.5}{-9}\,\si{\second}
=\SI{2.5}{\nano\second}$.\\
Na vertical (MUV, partindo do repouso):
\[
y=\tfrac12 a t^{2}=\tfrac12(\pot{1.76}{15})(\pot{2.5}{-9})^{2}
\approx\pot{5.5}{-3}\,\si{\meter}=\SI{5.5}{\milli\meter}.
\]
(c) A velocidade vertical na saída é $v_y = a\,t \approx
\pot{4.4}{6}\,\si{\meter\per\second}$, logo
\[
\tan\theta=\frac{v_y}{v_0}=\frac{\pot{4.4}{6}}{\pot{2}{7}}=0{,}22
\;\Longrightarrow\;
\theta\approx\SI{12.4}{\degree}.
\]
\textbf{Aplicação:} é exatamente esse o mecanismo de deflexão dos antigos tubos de
raios catódicos (TV e osciloscópios) e das impressoras a jato de tinta, que
defletem gotículas carregadas para formar caracteres.}
\end{questao}

\begin{questao}[3]{}
Quatro cargas de módulo $q=\SI{1}{\micro\coulomb}$ ocupam os vértices de um
quadrado de lado $a=\SI{10}{\centi\meter}$. Compare o \textbf{campo elétrico no
centro} em três configurações:
\begin{itensq}
\item todas positivas;
\item duas positivas e duas negativas, alternadas ($+,-,+,-$);
\item duas positivas e duas negativas, adjacentes ($+,+,-,-$).
\end{itensq}
\dados{$\ke=\SI{9e9}{}$; $\sqrt2\approx1{,}41$.}
\resposta{(a) $\vec E=0$; (b) $\vec E=0$; (c) $E\approx\pot{5.1}{6}\,\si{\newton\per\coulomb}$}
\resolucao{Cada carga dista $\dfrac{a\sqrt2}{2}=\SI{0.0707}{\meter}$ do centro,
produzindo individualmente
\[
E_1=\frac{\ke q}{(a\sqrt2/2)^{2}}=\frac{9\cdot10^{9}\cdot\pot{1}{-6}}{0{,}005}
=\pot{1.8}{6}\,\si{\newton\per\coulomb}.
\]
\textbf{(a) Todas positivas.} Os quatro vetores apontam radialmente para fora,
aos pares diametralmente opostos: cancelam-se dois a dois. $\vec E = 0$.\\
\textbf{(b) Alternadas ($+,-,+,-$).} As cargas diametralmente opostas têm o
\emph{mesmo} sinal, e seus campos continuam se cancelando aos pares.
$\vec E = 0$ novamente.\\
\textbf{(c) Adjacentes ($+,+,-,-$).} Agora cada par diametralmente oposto tem
sinais \emph{contrários}: os campos se \textbf{somam} em vez de cancelar. Os dois
pares contribuem na mesma direção (perpendicular à linha que separa as positivas
das negativas):
\[
E=2\,(E_1\sqrt2)=2\sqrt2\,E_1\approx 2{,}83\cdot\pot{1.8}{6}
\approx\pot{5.1}{6}\,\si{\newton\per\coulomb}.
\]
\textbf{Lição central.} Em (a) e (b), configurações \emph{muito} diferentes
produzem o mesmo resultado nulo. Em (c), com as mesmas quatro cargas, apenas
reposicionadas, o campo é máximo. \emph{O campo resultante depende da geometria
tanto quanto dos valores das cargas} --- e a simetria é a ferramenta que revela
isso sem nenhum cálculo pesado.}
\end{questao}

\begin{questao}[3]{}
Uma carga $-Q$ está na origem e uma carga $+4Q$ está em $x=d$.
\begin{itensq}
\item Determine o ponto sobre o eixo $x$ onde o campo resultante é nulo.
\item Explique por que ele fica \emph{fora} do segmento e do lado da carga de
      menor módulo.
\item Nesse ponto, quanto vale o potencial elétrico?
\end{itensq}
\resposta{(a) $x=-d$ (a uma distância $d$ à esquerda de $-Q$);
(b) e (c) ver comentário; $V=\dfrac{\ke Q}{d}$}
\resolucao{(a) Como as cargas têm sinais \emph{opostos}, entre elas os campos
apontam no mesmo sentido e se somam --- o cancelamento só pode ocorrer fora do
segmento. Seja $x<0$ a posição procurada; a distância a $-Q$ é $|x|$ e a $+4Q$ é
$|x|+d$:
\[
\frac{\ke Q}{|x|^{2}}=\frac{\ke\,4Q}{(|x|+d)^{2}}
\;\Longrightarrow\;
\frac{|x|+d}{|x|}=2
\;\Longrightarrow\;
|x|=d.
\]
O ponto está em $x=-d$.\\
(b) Do lado da carga \emph{maior}, seu campo sempre domina (está mais perto
\emph{e} é mais intensa) --- nunca poderia ser equilibrado. Do lado da carga
menor, a proximidade compensa a menor intensidade, e existe exatamente um ponto
de equilíbrio.\\
(c) O potencial é escalar:
\[
V=\ke\frac{(-Q)}{d}+\ke\frac{4Q}{2d}
=-\frac{\ke Q}{d}+\frac{2\ke Q}{d}
=\frac{\ke Q}{d}\neq0.
\]
\textbf{Mais um exemplo de que $\vec E=0$ não implica $V=0$} --- e vice-versa. O
campo mede a \emph{taxa de variação} do potencial, não o seu valor.}
\end{questao}

\begin{questao}[3]{}
Um anel isolante de raio $R$ está uniformemente carregado com carga total $Q$.
Considere um ponto P sobre o eixo do anel, a uma distância $x$ do centro.
\begin{itensq}
\item Explique, por simetria, por que o campo em P é paralelo ao eixo.
\item Sabendo que $E(x)=\dfrac{\ke\,Q\,x}{\left(x^{2}+R^{2}\right)^{3/2}}$,
      determine $E$ no centro do anel e no limite $x\gg R$.
\item Mostre que o campo é \emph{máximo} em $x = \dfrac{R}{\sqrt2}$.
\end{itensq}
\resposta{(b) $E(0)=0$ e $E\to\ke Q/x^{2}$; (c) demonstração}
\resolucao{(a) \textbf{Argumento de simetria.} Para cada elemento de carga do
anel existe outro diametralmente oposto. As componentes de campo
\emph{perpendiculares} ao eixo desses dois elementos são iguais e opostas,
cancelando-se. Restam apenas as componentes \emph{paralelas} ao eixo, que se
somam. Logo $\vec E$ é axial.

(b) \textbf{No centro} ($x=0$): substituindo,
\[
E(0)=\frac{\ke Q\cdot 0}{R^{3}}=0.
\]
Faz sentido: no centro, as contribuições de elementos opostos se cancelam
\emph{integralmente}.\\
\textbf{Longe} ($x\gg R$): então $\left(x^{2}+R^{2}\right)^{3/2}\approx x^{3}$ e
\[
E\approx\frac{\ke Q x}{x^{3}}=\frac{\ke Q}{x^{2}},
\]
o campo de uma carga puntiforme --- como esperado, pois de longe o anel "vira"
um ponto.

(c) \textbf{Máximo.} Derivando $E(x)$ e igualando a zero:
\[
\frac{\mathrm{d}E}{\mathrm{d}x}
=\ke Q\,\frac{\left(x^{2}+R^{2}\right)^{3/2}-x\cdot\tfrac32\left(x^{2}+R^{2}\right)^{1/2}(2x)}
{\left(x^{2}+R^{2}\right)^{3}}=0.
\]
O numerador, dividido por $\left(x^{2}+R^{2}\right)^{1/2}$, dá
\[
\left(x^{2}+R^{2}\right)-3x^{2}=0
\;\Longrightarrow\;
R^{2}=2x^{2}
\;\Longrightarrow\;
\boxed{\,x=\frac{R}{\sqrt2}\,}
\]
\textbf{Interpretação física do máximo.} Perto do centro, as contribuições
axiais são pequenas (os vetores são quase perpendiculares ao eixo). Muito longe,
a distância cresce e o campo decai. Entre os dois regimes existe um ponto ótimo,
a $x=R/\sqrt2\approx0{,}71R$.\\
\textbf{Onde isso é usado:} essa geometria é a base das \emph{bobinas de
Helmholtz} (versão magnética do mesmo cálculo) e dos anéis focalizadores em
aceleradores de partículas e microscópios eletrônicos.}
\end{questao}

\blocodesafio

\begin{questao}[4]{}
Uma haste isolante fina de comprimento $L$ está uniformemente carregada com carga
total $Q$ (densidade linear $\lambda=Q/L$). Considere um ponto P sobre o
prolongamento da haste, a uma distância $a$ da extremidade mais próxima.
\begin{itensq}
\item Monte a integral que fornece o campo elétrico em P.
\item Mostre que $E=\dfrac{\ke\,Q}{a\,(a+L)}$.
\item Verifique o comportamento no limite $a \gg L$ e interprete.
\end{itensq}
\resposta{(b) demonstração; (c) $E\to \ke Q/a^{2}$ (carga puntiforme)}
\resolucao{(a) Tome um elemento de comprimento $\mathrm{d}x$ à distância $x$ de P
(com $x$ variando de $a$ até $a+L$). Sua carga é
$\mathrm{d}q=\lambda\,\mathrm{d}x$ e contribui com
\[
\mathrm{d}E=\frac{\ke\,\mathrm{d}q}{x^{2}}=\frac{\ke\lambda\,\mathrm{d}x}{x^{2}}.
\]
Como todos os elementos estão alinhados com P, as contribuições têm a
\emph{mesma direção} e somam-se escalarmente:
\[
E=\int_{a}^{a+L}\frac{\ke\lambda}{x^{2}}\,\mathrm{d}x.
\]
(b) Integrando:
\[
E=\ke\lambda\left[-\frac{1}{x}\right]_{a}^{a+L}
=\ke\lambda\left(\frac{1}{a}-\frac{1}{a+L}\right)
=\ke\lambda\,\frac{L}{a(a+L)}.
\]
Como $\lambda L = Q$:
\[
\boxed{\;E=\frac{\ke\,Q}{a\,(a+L)}\;}
\]
(c) \textbf{Limite $a\gg L$:} então $a+L\approx a$ e
\[
E\approx\frac{\ke Q}{a^{2}},
\]
que é exatamente o campo de uma \textbf{carga puntiforme} $Q$.\\
\textbf{Interpretação:} de longe, qualquer distribuição finita de carga
"parece" puntiforme. Esse teste de consistência --- verificar se a expressão
obtida reproduz o caso conhecido em algum limite --- é uma das melhores formas de
validar um resultado. \emph{Se o limite não fechasse, haveria erro na
integração.}\\
\textbf{Limite oposto ($a\ll L$):} $E\approx\dfrac{\ke Q}{aL}=\dfrac{\ke\lambda}{a}$
--- campo de um fio "infinito" visto de perto, decaindo com $1/a$ em vez de
$1/a^{2}$.}
\end{questao}

\begin{questao}[4]{}
Um dipolo elétrico é formado por cargas $+q$ e $-q$ separadas por uma pequena
distância $2\ell$. Define-se o \textbf{momento de dipolo} $p=q\,(2\ell)$.
\begin{itensq}
\item Mostre que, sobre a mediatriz do dipolo, a uma distância $r\gg\ell$ do
      centro, o campo vale $E\approx\dfrac{\ke\,p}{r^{3}}$.
\item Por que o campo do dipolo decai mais rápido que o de uma carga isolada?
\item Um dipolo colocado em campo uniforme sofre força resultante? E torque?
\end{itensq}
\resposta{(a) demonstração; (b) e (c) ver comentário}
\resolucao{(a) Sobre a mediatriz, cada carga dista
$d=\sqrt{r^{2}+\ell^{2}}$ do ponto e produz campo de módulo
$E_1=\dfrac{\ke q}{r^{2}+\ell^{2}}$. As componentes \emph{paralelas} à mediatriz
se cancelam (simetria); restam as componentes ao longo do eixo do dipolo, cada uma
valendo $E_1\cos\alpha$, com $\cos\alpha=\dfrac{\ell}{\sqrt{r^{2}+\ell^{2}}}$:
\[
E=2E_1\cos\alpha
=\frac{2\ke q\,\ell}{\left(r^{2}+\ell^{2}\right)^{3/2}}.
\]
Para $r\gg\ell$, $\left(r^{2}+\ell^{2}\right)^{3/2}\approx r^{3}$, e como
$p=2q\ell$:
\[
\boxed{\;E\approx\frac{\ke\,p}{r^{3}}\;}
\]
(b) Porque a carga \emph{total} do dipolo é nula. De longe, os campos das duas
cargas quase se cancelam --- sobra apenas o efeito da pequena separação entre
elas, que é de ordem superior. Daí o decaimento $1/r^{3}$ em vez de $1/r^{2}$.\\
(c) Em campo \textbf{uniforme}:
\begin{itemize}[leftmargin=1.6em,itemsep=1pt,topsep=2pt]
\item \textbf{Força resultante nula}: as forças $+qE$ e $-qE$ sobre as duas cargas
      têm mesmo módulo e sentidos opostos.
\item \textbf{Torque não nulo}: como elas atuam em \emph{pontos diferentes},
      formam um binário de momento
      \[
      \tau = p\,E\sin\theta,
      \]
      que tende a \emph{alinhar} o dipolo com o campo.
\end{itemize}
\textbf{Aplicações:} é esse torque que alinha moléculas polares (como a água) em
um campo --- princípio do forno de micro-ondas, onde o campo oscilante faz as
moléculas girarem e o atrito molecular gera calor. Em campo \emph{não} uniforme,
surge também força resultante --- razão pela qual um objeto neutro é atraído por
um corpo eletrizado.}
\end{questao}

% BEGIN BANCO COMPLEMENTAR Campo elétrico
% =====================================================================
%  Banco complementar --- Campo Eletrico  (REVISADO)
%  Substitui a versao gerada por template (8 clones em N1, 11 em N2,
%  7 em N3 e 8 questoes conceituais marcadas como N4).
%  O cap01b e denso e ja cobre: E de carga puntiforme, proporcionalidade
%  com a distancia, carga de prova com forca dada, direcao e sentido do
%  campo, PONTO DE CAMPO NULO (tres vezes: Q e 4Q; 9 e 4 uC; -Q e +4Q),
%  duas cargas de mesmo modulo a 20 cm, grafico E x d, ELETRON ENTRE
%  PLACAS PARALELAS com deflexao, quadrado de quatro cargas, anel
%  isolante no eixo, haste finita por integracao (N4) e momento dipolar
%  (N4).
%  Este banco evita esses casos e explora: E=0 com V nao nulo, linhas de
%  campo e por que nao se cruzam, PROJETO de placas paralelas, cavidade
%  em condutor, tracado numerico de linhas, esfera isolante macica por
%  Gauss, par de placas com +-sigma e o limite imposto pela rigidez
%  dieletrica.
%  Distribuicao mantida: 8 N1 + 11 N2 + 7 N3 + 8 N4 = 34
% =====================================================================

\subsubsection*{Banco complementar --- Campo elétrico}

\begin{atencao}[Organização do banco]
O campo é \textbf{vetorial}: contribuições somam-se por componentes, e podem
cancelar-se. Ele é definido por $\vec E=\vec F/q_0$, mas \emph{não depende} da
carga de prova --- é propriedade do ponto. Duas expressões dominam o tópico:
$E=\ke Q/r^{2}$ para carga puntiforme e $E=\sigma/\varepsilon_0$ para plano
carregado, esta última \emph{independente da distância}. O N4 trata de projeto,
condutores, distribuições contínuas e limites físicos.
\end{atencao}

\blocofixacao

\begin{questao}[1]{}
Determine o módulo do campo elétrico produzido por $Q=+\SI{3}{\micro\coulomb}$
a $\SI{30}{\centi\meter}$ de distância.
\dados{$\ke=\SI{9e9}{\newton\meter\squared\per\coulomb\squared}$}
\resposta{$E=\SI{300}{\kilo\newton\per\coulomb}$}
\resolucao{
\[
E=\frac{\ke Q}{r^{2}}=\frac{\pot{9}{9}\times\pot{3}{-6}}{(0{,}30)^{2}}
=\frac{\pot{2.7}{4}}{0{,}09}=\pot{3}{5}\,\si{\newton\per\coulomb}.
\]
O campo aponta \textbf{para longe} da carga, por ela ser positiva.\\
\textbf{Atenção à potência:} o campo cai com $1/r^{2}$, o potencial com
$1/r$. Dobrar a distância divide $E$ por quatro e $V$ por dois.}
\end{questao}

\begin{questao}[1]{}
Determine o campo produzido por $Q=-\SI{5}{\micro\coulomb}$ a
$\SI{50}{\centi\meter}$, indicando o sentido.
\resposta{$E=\SI{180}{\kilo\newton\per\coulomb}$, apontando \emph{para} a carga}
\resolucao{
\[
E=\frac{\pot{9}{9}\times\pot{5}{-6}}{0{,}25}
=\pot{1.8}{5}\,\si{\newton\per\coulomb}.
\]
\textbf{Sentido:} cargas negativas \emph{atraem} cargas de prova positivas, logo
o campo aponta \textbf{para} a carga fonte.\\
\textbf{Regra mnemônica:} o campo \emph{sai} do positivo e \emph{entra} no
negativo --- é a mesma regra que orienta as linhas de campo.}
\end{questao}

\begin{questao}[1]{}
Uma carga $q=+\SI{2}{\micro\coulomb}$ é colocada em um ponto onde
$E=\SI{5000}{\newton\per\coulomb}$. Determine a força sobre ela.
\resposta{$F=\SI{10}{\milli\newton}$, no sentido do campo}
\resolucao{
\[
F=qE=\pot{2}{-6}\times5000=\pot{1}{-2}\,\si{\newton}=\SI{10}{\milli\newton}.
\]
Como a carga é positiva, a força tem o \emph{mesmo} sentido de $\vec E$. Se
fosse negativa, teria sentido oposto --- e apenas o sentido mudaria, não o
módulo.}
\end{questao}

\begin{questao}[1]{}
Determine a carga que produz $E=\SI{2000}{\newton\per\coulomb}$ a
$\SI{30}{\centi\meter}$.
\resposta{$Q=\SI{20}{\nano\coulomb}$}
\resolucao{
\[
Q=\frac{Er^{2}}{\ke}=\frac{2000\times0{,}09}{\pot{9}{9}}
=\pot{2}{-8}\,\si{\coulomb}=\SI{20}{\nano\coulomb}.
\]
Uma carga de nanocoulombs --- ordem de grandeza típica da eletrização por
atrito --- já produz milhares de newtons por coulomb a trinta centímetros.}
\end{questao}

\begin{questao}[1]{}
A que distância de $Q=+\SI{8}{\micro\coulomb}$ o campo vale
$\SI{200}{\kilo\newton\per\coulomb}$?
\resposta{$r=\SI{0.6}{\meter}$}
\resolucao{
\[
r=\sqrt{\frac{\ke Q}{E}}=\sqrt{\frac{\pot{9}{9}\times\pot{8}{-6}}{\pot{2}{5}}}
=\sqrt{0{,}36}=\SI{0.6}{\meter}.
\]
\textbf{Superfície de campo constante:} todos os pontos a
$\SI{0.6}{\meter}$ da carga têm o mesmo \emph{módulo} de campo --- mas
\emph{direções} diferentes, pois o campo é radial. Diferentemente das
equipotenciais, essas superfícies não têm significado especial.}
\end{questao}

\begin{questao}[1]{}
O campo elétrico em um ponto, medido com carga de prova $q_0$, vale $E$. Se a
carga de prova for \textbf{dobrada}, o valor medido:
\begin{alternativas}[2]
\item dobra
\item cai à metade
\item não se altera
\item depende do sinal de $q_0$
\end{alternativas}
\resposta{(c)}
\resolucao{Por definição, $\vec E=\vec F/q_0$. Dobrar $q_0$ dobra a força
\emph{e} o denominador --- a razão permanece.\\
\textbf{Ponto conceitual:} o campo é uma propriedade \textbf{do ponto do
espaço}, criada pelas cargas-fonte. A carga de prova apenas o \emph{revela};
ela não o cria nem o modifica.\\
\textbf{Ressalva real:} se $q_0$ for grande demais, ela perturba as
cargas-fonte (por indução, em condutores) e a medida deixa de valer. Daí a
definição rigorosa usar o limite $q_0\to0$.}
\end{questao}

\begin{questao}[1]{}
Mostre que $\si{\newton\per\coulomb}$ e $\si{\volt\per\meter}$ são a mesma
unidade.
\resposta{$\si{\volt}=\si{\joule\per\coulomb}$ e
$\si{\joule}=\si{\newton\meter}$}
\resolucao{
\[
\frac{\si{\volt}}{\si{\meter}}
=\frac{\si{\joule}}{\si{\coulomb}\cdot\si{\meter}}
=\frac{\si{\newton}\cdot\si{\meter}}{\si{\coulomb}\cdot\si{\meter}}
=\frac{\si{\newton}}{\si{\coulomb}}.\ \checkmark
\]
\textbf{As duas leituras do campo.} $\si{\newton\per\coulomb}$ enfatiza a
\emph{força} por unidade de carga; $\si{\volt\per\meter}$ enfatiza a
\emph{taxa de queda do potencial} por unidade de comprimento. São descrições
equivalentes de $\vec E=-\nabla V$.\\
Na prática, usa-se $\si{\volt\per\meter}$ quando há tensões envolvidas (placas,
cabos) e $\si{\newton\per\coulomb}$ em problemas de dinâmica de partículas.}
\end{questao}

\begin{questao}[1]{}
Uma carga negativa é colocada em uma região de campo elétrico uniforme. Sobre
seu movimento inicial, é correto afirmar que ela:
\begin{alternativas}[2]
\item acelera no sentido do campo
\item acelera no sentido oposto ao campo
\item permanece em repouso
\item move-se perpendicularmente ao campo
\end{alternativas}
\resposta{(b)}
\resolucao{A força vale $\vec F=q\vec E$ com $q<0$, logo $\vec F$ é
\textbf{antiparalela} a $\vec E$.\\
\textbf{Consequência energética:} a carga negativa move-se espontaneamente no
sentido de $V$ \emph{crescente} --- o oposto de uma carga positiva. Em ambos os
casos, a energia potencial $U=qV$ \emph{diminui}, como deve ocorrer em qualquer
movimento espontâneo.}
\end{questao}

\blocoaplicacao

\begin{questao}[2]{}
Uma carga $q_0=-\SI{1}{\micro\coulomb}$ é colocada em um campo uniforme de
$E=\SI{2000}{\newton\per\coulomb}$.
\begin{itensq}
\item Determine a força.
\item Determine sua aceleração, se a massa for $\SI{4}{\gram}$.
\end{itensq}
\resposta{$F=\SI{2}{\milli\newton}$, oposta a $\vec E$;
$a=\SI{0.5}{\meter\per\second\squared}$}
\resolucao{(a) $F=|q|E=\pot{1}{-6}(2000)=\pot{2}{-3}\,\si{\newton}$, com sentido
\emph{oposto} ao campo.\\
(b)
\[
a=\frac{F}{m}=\frac{\pot{2}{-3}}{\pot{4}{-3}}=\SI{0.5}{\meter\per\second\squared}.
\]
\textbf{Comparação com a gravidade:} $a$ é apenas $5\%$ de $g$. Para objetos
macroscópicos, a força elétrica só compete com o peso se a carga for
relativamente grande ou a massa muito pequena --- daí os experimentos
eletrostáticos clássicos usarem pedaços de papel e gotas de óleo.}
\end{questao}

\begin{questao}[2]{}
Um próton é abandonado em um campo uniforme de
$E=\SI{1e4}{\newton\per\coulomb}$. Determine sua aceleração.
\dados{$e=\pot{1.6}{-19}\,\si{\coulomb}$; $m_p=\pot{1.67}{-27}\,\si{\kilogram}$}
\resposta{$a=\pot{9.6}{11}\,\si{\meter\per\second\squared}$}
\resolucao{
\[
a=\frac{eE}{m_p}=\frac{(\pot{1.6}{-19})(\pot{1}{4})}{\pot{1.67}{-27}}
=\pot{9.58}{11}\,\si{\meter\per\second\squared}.
\]
\textbf{Cerca de $10^{11}$ vezes $g$} --- a gravidade é completamente
desprezível na dinâmica de partículas carregadas.\\
\textbf{Para o elétron}, a aceleração seria $1836$ vezes maior
($\pot{1.76}{15}\,\si{\meter\per\second\squared}$), pela razão das massas. É por
isso que, em tubos e aceleradores, são os elétrons que respondem primeiro.}
\end{questao}

\begin{questao}[2]{}
Duas cargas estão sobre o eixo $x$: $+\SI{4}{\micro\coulomb}$ em $x=0$ e
$-\SI{2}{\micro\coulomb}$ em $x=\SI{0.30}{\meter}$. Determine o campo em
$x=\SI{0.60}{\meter}$.
\resposta{$E=\SI{100}{\kilo\newton\per\coulomb}$, no sentido $-x$}
\resolucao{
\[
E_1=\frac{\pot{9}{9}(\pot{4}{-6})}{(0{,}60)^{2}}
=\pot{1}{5}\,\si{\newton\per\coulomb}\ (\text{sentido }+x);
\]
\[
E_2=\frac{\pot{9}{9}(\pot{2}{-6})}{(0{,}30)^{2}}
=\pot{2}{5}\,\si{\newton\per\coulomb}\ (\text{sentido }-x).
\]
\[
E=2\times10^{5}-1\times10^{5}=\pot{1}{5}\,\si{\newton\per\coulomb}\ (-x).
\]
\textbf{Observe:} a carga \emph{menor} produziu o campo \emph{maior}, por estar
duas vezes mais próxima --- e $1/r^{2}$ vence o fator $2$ de carga.
Compare os fatores: carga $\times\frac12$, distância$^{2}$
$\times\frac14$.}
\end{questao}

\begin{questao}[2]{}
Duas cargas iguais $Q=+\SI{2}{\micro\coulomb}$ estão em
$(\pm\SI{0.30}{\meter},\,0)$. Determine o campo no ponto
$(0,\ \SI{0.40}{\meter})$.
\resposta{$E=\SI{115.2}{\kilo\newton\per\coulomb}$, ao longo de $+y$}
\resolucao{Cada carga dista $r=\sqrt{0{,}30^{2}+0{,}40^{2}}=\SI{0.5}{\meter}$
(triângulo $3$--$4$--$5$):
\[
E_1=E_2=\frac{\pot{9}{9}(\pot{2}{-6})}{0{,}25}
=\pot{7.2}{4}\,\si{\newton\per\coulomb}.
\]
\textbf{Por simetria}, as componentes horizontais se cancelam e as verticais se
somam, cada uma multiplicada por $\cos\theta=0{,}40/0{,}50=0{,}8$:
\[
E=2(\pot{7.2}{4})(0{,}8)=\pot{1.152}{5}\,\si{\newton\per\coulomb}.
\]
\textbf{Uso da simetria:} reconhecê-la elimina metade do trabalho. Sempre que a
configuração for simétrica em relação ao ponto de interesse, identifique quais
componentes se anulam \emph{antes} de calcular.}
\end{questao}

\begin{questao}[2]{}
Duas placas paralelas separadas por $\SI{2}{\milli\meter}$ estão submetidas a
$\SI{600}{\volt}$. Determine o campo entre elas.
\resposta{$E=\SI{300}{\kilo\volt\per\meter}$}
\resolucao{
\[
E=\frac{U}{d}=\frac{600}{\pot{2}{-3}}=\pot{3}{5}\,\si{\volt\per\meter}.
\]
\textbf{Campo uniforme:} entre placas extensas e próximas, $E$ é praticamente o
mesmo em todos os pontos --- e independe da posição. É a configuração de
referência para acelerar e defletir partículas.\\
\textbf{Verificação de segurança:} $\SI{300}{\kilo\volt\per\meter}$ é um décimo
da rigidez do ar ($\SI{3}{\mega\volt\per\meter}$) --- há folga confortável.
Reduzir a separação para $\SI{0.2}{\milli\meter}$ com a mesma tensão provocaria
ruptura.}
\end{questao}

\begin{questao}[2]{}
Uma placa condutora extensa tem densidade superficial de carga
$\sigma=\SI{2}{\micro\coulomb\per\meter\squared}$. Determine o campo próximo à
sua superfície.
\dados{$\varepsilon_0=\pot{8.85}{-12}\,\si{\farad\per\meter}$}
\resposta{$E=\SI{226}{\kilo\newton\per\coulomb}$}
\resolucao{Para um condutor em equilíbrio,
\[
E=\frac{\sigma}{\varepsilon_0}=\frac{\pot{2}{-6}}{\pot{8.85}{-12}}
=\pot{2.26}{5}\,\si{\newton\per\coulomb}.
\]
\textbf{Fato notável:} o campo \textbf{não depende da distância} à placa (desde
que ela seja extensa e o ponto próximo). Diferentemente da carga puntiforme, não
há queda com $1/r^{2}$.\\
\textbf{Cuidado com a fórmula:} para um \emph{plano isolante} com carga nas duas
faces, o campo vale $\sigma/2\varepsilon_0$ --- metade. A diferença é que no
condutor toda a carga fica em uma face, e o campo interno é nulo.}
\end{questao}

\begin{questao}[2]{}
Compare o campo produzido por $Q_1=+\SI{9}{\micro\coulomb}$ a
$\SI{30}{\centi\meter}$ com o de $Q_2=+\SI{4}{\micro\coulomb}$ a
$\SI{20}{\centi\meter}$.
\resposta{$E_1=E_2=\SI{900}{\kilo\newton\per\coulomb}$ --- iguais}
\resolucao{
\[
E_1=\frac{\pot{9}{9}(\pot{9}{-6})}{0{,}09}=\pot{9}{5};
\qquad
E_2=\frac{\pot{9}{9}(\pot{4}{-6})}{0{,}04}=\pot{9}{5}\,\si{\newton\per\coulomb}.
\]
\textbf{Coincidência não é acaso:} as razões $Q/r^{2}$ são iguais, pois
$9/9=4/4$. Sempre que $Q\propto r^{2}$, o campo se mantém.\\
\textbf{Consequência:} se essas duas cargas estivessem em lados opostos de um
ponto, o campo resultante ali seria \emph{nulo} --- e essa é exatamente a
condição do "ponto de campo nulo" entre cargas de mesmo sinal.}
\end{questao}

\begin{questao}[2]{}
Uma carga em um ponto sofre campos de $\SI{3}{\kilo\newton\per\coulomb}$ na
direção $x$ e $\SI{4}{\kilo\newton\per\coulomb}$ na direção $y$, produzidos por
duas fontes distintas. Determine o campo resultante.
\resposta{$E=\SI{5}{\kilo\newton\per\coulomb}$, a $\ang{53.1}$ do eixo $x$}
\resolucao{Sendo perpendiculares, compõem-se por Pitágoras:
\[
E=\sqrt{3^{2}+4^{2}}=\SI{5}{\kilo\newton\per\coulomb};
\qquad
\theta=\arctan\frac{4}{3}=\ang{53.1}.
\]
\textbf{Contraste com o potencial:} se esses mesmos pontos tivessem potenciais
de $3$ e $\SI{4}{\kilo\volt}$, o total seria simplesmente
$\SI{7}{\kilo\volt}$ --- soma algébrica, sem geometria.\\
É essa diferença que torna o potencial a ferramenta preferida para cálculos, e o
campo a ferramenta preferida para entender forças e movimento.}
\end{questao}

\begin{questao}[2]{}
Em uma representação por linhas de campo, o que indica uma região onde as
linhas estão mais próximas entre si?
\begin{alternativas}[2]
\item potencial maior
\item campo mais intenso
\item carga maior naquele ponto
\item ausência de campo
\end{alternativas}
\resposta{(b)}
\resolucao{A \textbf{densidade de linhas} é proporcional ao módulo do campo ---
essa é a convenção que dá utilidade quantitativa ao desenho.\\
\textbf{Três regras das linhas de campo:}
\begin{enumerate}
\item Nascem em cargas positivas e morrem em negativas (ou no infinito).
\item São tangentes a $\vec E$ em cada ponto.
\item \textbf{Nunca se cruzam} --- em um cruzamento, o campo teria duas
direções, o que é impossível.
\end{enumerate}
\textbf{Cuidado com (a):} potencial e campo são grandezas distintas. Linhas
próximas indicam campo intenso, o que corresponde a equipotenciais próximas ---
mas não diz nada sobre o \emph{valor} de $V$.}
\end{questao}

\begin{questao}[2]{}
Uma carga produz campo de $\SI{8}{\kilo\newton\per\coulomb}$ no vácuo. O
conjunto é imerso em um dielétrico de $\varepsilon_r=4$. Determine o novo campo.
\resposta{$E=\SI{2}{\kilo\newton\per\coulomb}$}
\resolucao{
\[
E_{meio}=\frac{E_{vácuo}}{\varepsilon_r}=\frac{8}{4}
=\SI{2}{\kilo\newton\per\coulomb}.
\]
\textbf{Mecanismo:} o dielétrico se \emph{polariza} --- suas moléculas se
orientam e criam um campo interno oposto, que cancela parcialmente o campo
original.\\
\textbf{Consequência tecnológica:} o dielétrico reduz o campo para a mesma
carga, o que permite armazenar \emph{mais} carga para a mesma tensão --- é
exatamente por isso que $C=\varepsilon_r C_0$ em capacitores.}
\end{questao}

\begin{questao}[2]{}
Uma gota de massa $\SI{2}{\milli\gram}$ e carga $\SI{1}{\nano\coulomb}$ deve
ficar suspensa em repouso. Determine o campo elétrico necessário.
\dados{$g=\SI{9.8}{\meter\per\second\squared}$}
\resposta{$E=\SI{19.6}{\kilo\newton\per\coulomb}$}
\resolucao{Equilíbrio entre peso e força elétrica:
\[
qE=mg
\;\Longrightarrow\;
E=\frac{mg}{q}=\frac{(\pot{2}{-6})(9{,}8)}{\pot{1}{-9}}
=\pot{1.96}{4}\,\si{\newton\per\coulomb}.
\]
\textbf{Orientação:} o campo deve produzir força \emph{para cima}. Se a gota for
positiva, $\vec E$ aponta para cima; se negativa, para baixo.\\
\textbf{Este é o princípio do experimento de Millikan}, que determinou a carga
elementar. Medindo $E$ e $m$ de gotas suspensas, ele encontrou que todas as
cargas eram múltiplas inteiras de $\pot{1.6}{-19}\,\si{\coulomb}$ --- a
primeira evidência direta da quantização.}
\end{questao}

\blocoaprofundamento

\begin{questao}[3]{}
Duas cargas iguais $Q=+\SI{1}{\micro\coulomb}$ ocupam
$(\pm\SI{0.20}{\meter},\,0)$.
\begin{itensq}
\item Determine o campo em $(0,\ \SI{0.20}{\meter})$.
\item Determine o campo no ponto médio e justifique.
\end{itensq}
\resposta{(a) $E=\SI{159}{\kilo\newton\per\coulomb}$ ao longo de $+y$;
(b) $\vec E=\vec0$}
\resolucao{(a) Cada carga dista
$r=\sqrt{2}\,(0{,}20)=\SI{0.283}{\meter}$:
\[
E_1=E_2=\frac{\pot{9}{9}(\pot{1}{-6})}{0{,}08}
=\pot{1.125}{5}\,\si{\newton\per\coulomb}.
\]
As componentes em $x$ cancelam; as em $y$ somam-se, cada uma multiplicada por
$\cos\ang{45}=0{,}707$:
\[
E=2(\pot{1.125}{5})(0{,}707)=\pot{1.59}{5}\,\si{\newton\per\coulomb}.
\]
(b) \textbf{No ponto médio}, as duas contribuições têm módulos iguais e sentidos
\emph{opostos} (cada carga empurra para longe de si):
\[
\vec E=\vec 0.
\]
\textbf{Mas o potencial ali não é nulo:}
$V=2\ke Q/0{,}20=\SI{90}{\kilo\volt}$. É o caso complementar ao do dipolo, em
que $V=0$ e $\vec E\ne\vec0$ --- e mostra novamente que as duas condições são
independentes.\\
\textbf{Estabilidade:} uma carga positiva no ponto médio está em equilíbrio, mas
\emph{instável} na direção perpendicular ao eixo --- conforme o teorema de
Earnshaw.}
\end{questao}

\begin{questao}[3]{}
Um dipolo é formado por $+q$ e $-q$ separadas por $2a$. Considere um ponto do
eixo perpendicular à linha que as une, a distância $y$ do centro.
\begin{itensq}
\item Deduza a expressão exata do campo.
\item Obtenha a forma assintótica para $y\gg a$.
\item Compare a queda com a de uma carga isolada.
\end{itensq}
\resposta{$E=\dfrac{2\ke qa}{(y^{2}+a^{2})^{3/2}}\to\dfrac{2\ke qa}{y^{3}}$}
\resolucao{(a) Cada carga dista $r=\sqrt{y^{2}+a^{2}}$ e produz
$\ke q/r^{2}$. Por simetria, as componentes ao longo de $y$ se cancelam (uma
aponta para cima, a outra para baixo) e restam as componentes \emph{paralelas ao
eixo do dipolo}, cada uma multiplicada por $a/r$:
\[
E=2\cdot\frac{\ke q}{y^{2}+a^{2}}\cdot\frac{a}{\sqrt{y^{2}+a^{2}}}
=\frac{2\ke qa}{(y^{2}+a^{2})^{3/2}}.
\]
(b) Para $y\gg a$, $(y^{2}+a^{2})^{3/2}\to y^{3}$:
\[
E\approx\frac{2\ke qa}{y^{3}}=\frac{\ke p}{y^{3}},
\]
com $p=2qa$ o momento dipolar.\\
(c) \textbf{O campo do dipolo cai com $1/y^{3}$}, muito mais rápido que o
$1/r^{2}$ de uma carga isolada.\\
\textbf{Razão física:} de longe, as duas cargas quase se cancelam --- o sistema
é \emph{globalmente neutro}. O que sobra é o efeito de sua \emph{separação},
uma correção de ordem superior.\\
\textbf{Padrão geral:} carga isolada $\sim1/r^{2}$; dipolo
$\sim1/r^{3}$; quadrupolo $\sim1/r^{4}$. Quanto mais neutra a distribuição, mais
rápido o campo se extingue --- e é por isso que a matéria comum, neutra, não
exerce força elétrica a distância.}
\end{questao}

\begin{questao}[3]{}
Duas placas condutoras extensas e paralelas têm densidades superficiais
$+\sigma$ e $-\sigma$, com
$\sigma=\SI{1}{\micro\coulomb\per\meter\squared}$.
\begin{itensq}
\item Determine o campo entre as placas.
\item Determine o campo fora delas.
\item Justifique fisicamente.
\end{itensq}
\resposta{Entre: $E=\SI{113}{\kilo\newton\per\coulomb}$; fora: $E=0$}
\resolucao{(a) Cada plano contribui com $\sigma/2\varepsilon_0$, e
\textbf{entre} as placas as duas contribuições têm o \emph{mesmo} sentido (do
positivo para o negativo):
\[
E=2\cdot\frac{\sigma}{2\varepsilon_0}=\frac{\sigma}{\varepsilon_0}
=\frac{\pot{1}{-6}}{\pot{8.85}{-12}}=\pot{1.13}{5}\,\si{\newton\per\coulomb}.
\]
(b) \textbf{Fora}, as duas contribuições têm sentidos \emph{opostos} e módulos
iguais:
\[
E=\frac{\sigma}{2\varepsilon_0}-\frac{\sigma}{2\varepsilon_0}=0.
\]
(c) \textbf{Justificativa:} o campo de um plano infinito é uniforme e aponta
para longe dele (se positivo) ou para ele (se negativo). Na região externa,
qualquer ponto está do mesmo lado das duas placas --- e os campos se opõem.
Entre elas, o ponto está "depois" da positiva e "antes" da negativa, e os campos
se somam.\\
\textbf{Consequência: blindagem natural.} O par de placas confina completamente
seu campo na região interna --- não há campo escapando para fora. É por isso que
um capacitor de placas paralelas não interfere no circuito vizinho, e por que a
energia armazenada está inteiramente entre as armaduras.}
\end{questao}

\begin{questao}[3]{}
Uma esfera condutora de raio $R=\SI{10}{\centi\meter}$ tem carga
$Q=\SI{5}{\micro\coulomb}$.
\begin{itensq}
\item Determine o campo no interior, na superfície e a $\SI{30}{\centi\meter}$
      do centro.
\item Avalie se a configuração é fisicamente sustentável no ar.
\end{itensq}
\resposta{$0$; $\SI{4.5}{\mega\volt\per\meter}$;
$\SI{500}{\kilo\volt\per\meter}$ --- \textbf{insustentável}}
\resolucao{(a) \textbf{Interior:} $E=0$, condição de equilíbrio de qualquer
condutor.\\
\textbf{Superfície:} toda a carga se comporta como concentrada no centro, para
pontos externos:
\[
E=\frac{\ke Q}{R^{2}}=\frac{\pot{9}{9}(\pot{5}{-6})}{0{,}01}
=\pot{4.5}{6}\,\si{\volt\per\meter}.
\]
\textbf{A $\SI{30}{\centi\meter}$:}
$E=\dfrac{\pot{4.5}{4}}{0{,}09}=\pot{5}{5}\,\si{\volt\per\meter}$.\\
(b) \textbf{Insustentável.} O campo na superfície,
$\SI{4.5}{\mega\volt\per\meter}$, excede a rigidez dielétrica do ar
($\SI{3}{\mega\volt\per\meter}$) em $50\%$. O ar ioniza e a carga escapa por
descarga corona.\\
\textbf{Carga máxima real:}
$Q_{\max}=E_{rup}R^{2}/\ke=\SI{3.33}{\micro\coulomb}$. A esfera simplesmente
\emph{não consegue} reter os $\SI{5}{\micro\coulomb}$ enunciados.\\
\textbf{Lição de método:} sempre verifique se o resultado é fisicamente
possível. Um problema pode ser aritmeticamente consistente e descrever uma
situação inexistente.}
\end{questao}

\begin{questao}[3]{}
Um elétron parte do repouso em um campo uniforme de
$E=\SI{1000}{\volt\per\meter}$ e percorre $\SI{2}{\centi\meter}$.
\begin{itensq}
\item Determine a aceleração e o tempo de percurso.
\item Determine a velocidade final e a energia adquirida.
\end{itensq}
\dados{$e=\pot{1.6}{-19}\,\si{\coulomb}$; $m_e=\pot{9.11}{-31}\,\si{\kilogram}$}
\resposta{$a=\pot{1.76}{14}\,\si{\meter\per\second\squared}$;
$t=\SI{15.1}{\nano\second}$; $v=\pot{2.65}{6}\,\si{\meter\per\second}$;
$\SI{20}{\electronvolt}$}
\resolucao{(a)
\[
a=\frac{eE}{m_e}=\frac{(\pot{1.6}{-19})(1000)}{\pot{9.11}{-31}}
=\pot{1.756}{14}\,\si{\meter\per\second\squared};
\]
\[
t=\sqrt{\frac{2d}{a}}=\sqrt{\frac{2(0{,}02)}{\pot{1.756}{14}}}
=\pot{1.51}{-8}\,\si{\second}=\SI{15.1}{\nano\second}.
\]
(b) $v=at=(\pot{1.756}{14})(\pot{1.51}{-8})
=\pot{2.65}{6}\,\si{\meter\per\second}$.\\
\textbf{Energia:} o mais rápido é usar a tensão percorrida,
$U=Ed=1000(0{,}02)=\SI{20}{\volt}$:
\[
W=eU=\SI{20}{\electronvolt}=\pot{3.2}{-18}\,\si{\joule}.
\]
\emph{(Conferindo: $\frac12m_ev^{2}=\frac12(\pot{9.11}{-31})(\pot{2.65}{6})^{2}
=\pot{3.2}{-18}\,\si{\joule}$ \checkmark.)}\\
\textbf{Verificação relativística:} $v/c=0{,}9\%$ --- plenamente no regime
clássico.}
\end{questao}

\begin{questao}[3]{}
Compare a queda do campo com a distância para três configurações: carga
puntiforme, fio infinito carregado e plano infinito carregado.
\begin{itensq}
\item Escreva a dependência de cada uma.
\item Explique a origem geométrica das diferenças.
\end{itensq}
\resposta{$1/r^{2}$, $1/r$ e constante}
\resolucao{(a)
\begin{center}
\begin{tabular}{lcc}
\hline
Distribuição & Campo & Queda\\
\hline
Carga puntiforme & $\dfrac{\ke Q}{r^{2}}$ & $1/r^{2}$\\[6pt]
Fio infinito ($\lambda$) & $\dfrac{\lambda}{2\pi\varepsilon_0 r}$ & $1/r$\\[6pt]
Plano infinito ($\sigma$) & $\dfrac{\sigma}{2\varepsilon_0}$ & constante\\
\hline
\end{tabular}
\end{center}
(b) \textbf{Origem geométrica.} Pela lei de Gauss, o fluxo total é fixo pela
carga envolvida, e o campo é o fluxo dividido pela área da superfície:
\begin{itemize}
\item \textbf{Ponto:} a superfície é uma \emph{esfera}, de área
$4\pi r^{2}$ --- o campo cai com $1/r^{2}$.
\item \textbf{Fio:} a superfície é um \emph{cilindro}, de área lateral
$2\pi r\ell$ --- cresce apenas com $r$, e o campo cai com $1/r$.
\item \textbf{Plano:} a superfície é um \emph{prisma} cuja área \emph{não
depende} de $r$ --- e o campo não cai.
\end{itemize}
\textbf{Regra unificadora:} o campo cai como $1/r^{n}$, onde $n$ é o número de
dimensões em que a fonte é \emph{limitada}. Ponto: limitado em três
($n=2$). Fio: em duas ($n=1$). Plano: em uma ($n=0$).\\
\textbf{Consequência prática:} perto de uma placa carregada extensa, afastar-se
não ajuda --- o campo permanece. É o oposto da intuição formada com cargas
puntiformes.}
\end{questao}

\begin{questao}[3]{}
Uma haste isolante fina é carregada de modo que produz, em certo ponto, campo
de $\SI{600}{\newton\per\coulomb}$. A haste é então partida ao meio, e apenas
uma das metades permanece.
\begin{itensq}
\item O campo cai à metade? Justifique.
\item Determine em que condição isso valeria.
\end{itensq}
\resposta{Não necessariamente --- só vale se a metade removida contribuísse com
exatamente metade do campo}
\resolucao{(a) \textbf{Não, em geral.} O campo é a soma \emph{vetorial} das
contribuições de todos os elementos da haste, e cada elemento contribui com
módulo e direção diferentes, dependendo de sua distância e posição angular em
relação ao ponto.\\
Se o ponto estiver \emph{sobre a mediatriz} da haste, as duas metades
contribuem igualmente e o campo de fato cai à metade --- mas a \emph{direção}
resultante também muda, pois as componentes que antes se cancelavam agora não
se cancelam mais.\\
Se o ponto estiver \emph{sobre o prolongamento} da haste, a metade mais próxima
contribui muito mais que a distante, e remover a distante reduz o campo bem
menos que $50\%$.\\
(b) \textbf{A queda é exatamente à metade quando:}
\begin{itemize}
\item as duas metades contribuem com vetores \emph{idênticos} --- o que exige
simetria completa em relação ao ponto; \emph{e}
\item essa condição só se realiza em pontos muito distantes, onde a haste inteira
se comporta como carga puntiforme e a divisão é simplesmente $Q\to Q/2$.
\end{itemize}
\textbf{Lição:} a superposição vale \emph{sempre}, mas ela é vetorial. Reduzir
a fonte pela metade não reduz o campo pela metade, exceto em casos de simetria
particular.}
\end{questao}

\blocodesafio

\begin{questao}[4]{}
\textbf{(Campo nulo com potencial não nulo.)} Construa uma configuração em que
$\vec E=\vec0$ em um ponto onde $V\ne0$.
\begin{itensq}
\item Proponha o arranjo mais simples e determine os valores.
\item Explique por que isso é possível.
\item Determine se o inverso ($V=0$ com $\vec E\ne\vec0$) também ocorre.
\end{itensq}
\resposta{Duas cargas iguais $+Q$; no ponto médio $\vec E=\vec0$ e
$V=2\ke Q/a$}
\resolucao{(a) \textbf{Arranjo:} duas cargas iguais $Q=+\SI{2}{\micro\coulomb}$
em $(\pm\SI{0.20}{\meter},\,0)$. No ponto médio:
\[
\vec E=\vec 0
\quad\text{(contribuições opostas)};
\]
\[
V=2\cdot\frac{\ke Q}{0{,}20}
=2\cdot\frac{\pot{9}{9}(\pot{2}{-6})}{0{,}20}
=\pot{1.8}{5}\,\si{\volt}=\SI{180}{\kilo\volt}.
\]
Campo nulo, potencial de cento e oitenta mil volts.\\
(b) \textbf{Por que é possível.} As duas grandezas se relacionam por
$\vec E=-\nabla V$: o campo é a \emph{taxa de variação} de $V$, não seu valor.
Uma função pode ter valor grande e derivada nula --- é exatamente o que ocorre
em um ponto \emph{estacionário}.\\
No ponto médio, $V$ tem um mínimo ao longo do eixo que une as cargas: o valor é
alto, mas a inclinação é zero.\\
(c) \textbf{Sim, o inverso também ocorre.} Substituindo uma das cargas por
$-Q$, o ponto médio passa a ter $V=0$ (contribuições opostas se cancelam
algebricamente) e $\vec E\ne\vec0$ (os vetores apontam para o mesmo lado e se
somam).\\
\textbf{Síntese:} $V=0$ e $\vec E=\vec0$ são condições \textbf{logicamente
independentes}. Nenhuma implica a outra, em nenhum dos dois sentidos.}
\end{questao}

\begin{questao}[4]{}
\textbf{(Linhas de campo.)} Analise as linhas de campo de um dipolo elétrico.
\begin{itensq}
\item Descreva seu traçado.
\item Demonstre que linhas de campo nunca se cruzam.
\item Explique o que ocorre em pontos de campo nulo.
\end{itensq}
\resposta{Saem de $+q$ e entram em $-q$; não se cruzam porque $\vec E$ tem
direção única em cada ponto}
\resolucao{(a) \textbf{Traçado.} As linhas saem radialmente de $+q$, curvam-se e
entram em $-q$. A linha que segue o eixo do dipolo é reta; as demais formam
arcos progressivamente mais abertos. Algumas partem para o infinito e retornam.\\
A densidade de linhas é máxima entre as cargas --- onde o campo é mais intenso
--- e decresce com $1/r^{3}$ ao se afastar, conforme demonstrado antes.\\
(b) \textbf{Demonstração.} Suponha, por absurdo, que duas linhas se cruzem em
um ponto $P$. Por definição, a linha de campo é \emph{tangente} a $\vec E$ em
cada ponto. Se duas linhas distintas passam por $P$, elas têm tangentes
distintas ali --- e portanto $\vec E(P)$ teria \textbf{duas direções}
simultaneamente.\\
Mas $\vec E$ é uma função vetorial \emph{bem definida}: a cada ponto do espaço
corresponde \emph{um} vetor. Contradição. Logo, as linhas não se cruzam.\\
\textbf{Corolário:} o mesmo argumento vale para linhas de corrente, linhas de
campo magnético e qualquer outro campo vetorial.\\
(c) \textbf{Pontos de campo nulo são a exceção.} Onde $\vec E=\vec0$, não há
direção definida --- e ali as linhas \emph{podem} se encontrar, sem contradição.
São os \textbf{pontos de sela} do potencial.\\
Em um dipolo ($+q$ e $-q$) não existe tal ponto no espaço finito. Mas em duas
cargas \emph{iguais}, o ponto médio é exatamente um ponto de campo nulo --- e
ali as linhas se aproximam de todos os lados sem se tocarem, definindo uma
separatriz.}
\end{questao}

\begin{questao}[4]{}
\textbf{(Projeto de acelerador de placas.)} Projete um arranjo de placas
paralelas que imprima aceleração de
$a=\pot{2}{12}\,\si{\meter\per\second\squared}$ a um elétron.
\begin{itensq}
\item Determine o campo necessário.
\item Determine a tensão para separação de $\SI{1}{\centi\meter}$.
\item Determine o tempo de trânsito e a velocidade de saída.
\end{itensq}
\dados{$e=\pot{1.6}{-19}\,\si{\coulomb}$; $m_e=\pot{9.11}{-31}\,\si{\kilogram}$}
\resposta{$E=\SI{11.4}{\volt\per\meter}$; $U=\SI{0.114}{\volt}$;
$t=\SI{100}{\nano\second}$; $v=\pot{2}{5}\,\si{\meter\per\second}$}
\resolucao{(a)
\[
E=\frac{m_ea}{e}=\frac{(\pot{9.11}{-31})(\pot{2}{12})}{\pot{1.6}{-19}}
=\SI{11.4}{\volt\per\meter}.
\]
(b) $U=Ed=11{,}4(0{,}01)=\SI{0.114}{\volt}$.\\
(c)
\[
t=\sqrt{\frac{2d}{a}}=\sqrt{\frac{0{,}02}{\pot{2}{12}}}
=\pot{1}{-7}\,\si{\second}=\SI{100}{\nano\second};
\]
\[
v=at=(\pot{2}{12})(\pot{1}{-7})=\pot{2}{5}\,\si{\meter\per\second}.
\]
\textbf{O resultado é surpreendente e instrutivo.} Uma aceleração de
$\pot{2}{12}\,\si{\meter\per\second\squared}$ --- duzentos bilhões de vezes
$g$ --- exige apenas $\SI{0.114}{\volt}$, tensão menor que a de uma pilha
descarregada.\\
\textbf{Razão:} a razão carga/massa do elétron é enorme
($\pot{1.76}{11}\,\si{\coulomb\per\kilogram}$). Campos elétricos modestos
produzem acelerações colossais em partículas leves.\\
\textbf{Ressalva prática:} com tensão tão baixa, o projeto seria dominado por
efeitos parasitas --- diferenças de função-trabalho entre os metais das placas
(décimos de volt), cargas superficiais e ruído térmico. Um projeto real usaria
tensões de dezenas de volts e aceleração proporcionalmente maior.}
\end{questao}

\begin{questao}[4]{}
\textbf{(Cavidade em condutor.)} Uma carga puntiforme $+q$ é colocada no
interior de uma cavidade vazia dentro de um condutor \emph{neutro} e isolado.
\begin{itensq}
\item Determine a carga induzida em cada superfície.
\item Determine o campo dentro do metal e fora do condutor.
\item Analise o que muda se o condutor for aterrado.
\end{itensq}
\resposta{$-q$ na parede da cavidade, $+q$ na superfície externa; campo nulo no
metal e não nulo fora}
\resolucao{(a) \textbf{Parede da cavidade:} aplicando a lei de Gauss a uma
superfície inteiramente dentro do metal (onde $\vec E=\vec0$), o fluxo é nulo e
a carga envolvida também. Como a superfície envolve $+q$ e a parede interna:
\[
q+q_{interna}=0
\;\Longrightarrow\;
q_{interna}=-q.
\]
\textbf{Superfície externa:} o condutor era neutro, então a carga total deve
permanecer nula:
\[
q_{externa}=+q.
\]
(b) \textbf{Dentro do metal:} $\vec E=\vec0$, por equilíbrio.\\
\textbf{Fora:} a carga $+q$ da superfície externa produz campo --- e ele é
\textbf{esfericamente simétrico} se a superfície externa for esférica,
\emph{independentemente} de onde a carga esteja dentro da cavidade.\\
\textbf{Consequência notável:} de fora, é impossível saber a posição da carga
interna, nem a forma da cavidade. O condutor "apaga" essa informação.\\
(c) \textbf{Se o condutor for aterrado:} a carga $+q$ da superfície externa
escoa para a terra. Resta apenas $-q$ na parede da cavidade, e
\[
\vec E=\vec 0\quad\text{em todo o exterior}.
\]
\textbf{A blindagem torna-se completa nos dois sentidos.} É exatamente esse o
princípio da gaiola de Faraday aterrada --- e a diferença entre blindar contra
campos externos (que dispensa aterramento) e impedir que fontes internas
irradiem (que o exige).}
\end{questao}

\begin{questao}[4]{}
\textbf{(Esfera isolante uniformemente carregada.)} Uma esfera isolante maciça
de raio $R$ tem carga $Q$ distribuída uniformemente em seu \emph{volume}.
\begin{itensq}
\item Determine o campo para $r<R$ e para $r>R$.
\item Identifique onde o campo é máximo.
\item Compare com a esfera \emph{condutora} de mesma carga.
\end{itensq}
\resposta{$E=\dfrac{\ke Qr}{R^{3}}$ dentro e $\dfrac{\ke Q}{r^{2}}$ fora;
máximo na superfície}
\resolucao{(a) \textbf{Fora ($r>R$):} pela lei de Gauss com superfície esférica,
toda a carga está envolvida:
\[
E=\frac{\ke Q}{r^{2}},
\]
idêntico ao de uma carga puntiforme no centro.\\
\textbf{Dentro ($r<R$):} apenas a fração de carga interna à superfície gaussiana
conta. Como a densidade é uniforme, essa fração é $(r/R)^{3}$:
\[
E=\frac{\ke\,Q(r/R)^{3}}{r^{2}}=\frac{\ke Q\,r}{R^{3}},
\]
\textbf{proporcional a $r$} --- cresce linearmente do centro à superfície.\\
(b) \textbf{Máximo na superfície} ($r=R$), onde as duas expressões coincidem em
$E=\ke Q/R^{2}$. O campo é contínuo ali, embora com derivadas diferentes de
cada lado.\\
\emph{No centro, $E=0$ por simetria.}\\
(c) \textbf{Comparação com a esfera condutora.}
\begin{center}
\begin{tabular}{lcc}
\hline
Região & Isolante & Condutora\\
\hline
$r<R$ & $\ke Qr/R^{3}$ & $0$\\
$r=R$ & $\ke Q/R^{2}$ & $\ke Q/R^{2}$\\
$r>R$ & $\ke Q/r^{2}$ & $\ke Q/r^{2}$\\
\hline
\end{tabular}
\end{center}
\textbf{São indistinguíveis do lado de fora} --- a lei de Gauss só "enxerga" a
carga total envolvida, não sua distribuição. A diferença está inteiramente no
interior: no condutor a carga migra para a superfície e o campo interno é nulo;
no isolante ela permanece onde foi colocada.}
\end{questao}

\begin{questao}[4]{}
\textbf{(Traçado numérico de linhas de campo.)} Descreva um procedimento
computacional para traçar linhas de campo a partir de uma expressão conhecida de
$\vec E(x,y)$.
\begin{itensq}
\item Formule o algoritmo.
\item Indique como escolher os pontos de partida.
\item Aponte as dificuldades numéricas e como contorná-las.
\end{itensq}
\resposta{Integrar $\dfrac{\mathrm{d}\vec r}{\mathrm{d}s}
=\dfrac{\vec E}{|\vec E|}$ a partir de pontos distribuídos em torno das cargas}
\resolucao{(a) \textbf{Algoritmo.} A linha de campo é, por definição, tangente a
$\vec E$. Parametrizando pelo comprimento de arco $s$:
\[
\frac{\mathrm{d}\vec r}{\mathrm{d}s}=\frac{\vec E(\vec r)}{|\vec E(\vec r)|}.
\]
Numericamente, com passo $h$:
\begin{enumerate}
\item Escolha um ponto inicial $\vec r_0$.
\item Calcule $\vec E(\vec r_n)$ e normalize.
\item Avance: $\vec r_{n+1}=\vec r_n+h\,\hat E$ (Euler) ou, melhor, por
Runge--Kutta de quarta ordem.
\item Repita até sair do domínio ou atingir uma carga negativa.
\end{enumerate}
\textbf{Normalizar é essencial:} sem isso, o passo seria proporcional a
$|\vec E|$ e a linha "voaria" perto das cargas e travaria longe delas.\\
(b) \textbf{Pontos de partida.} Distribua-os em pequenos círculos ao redor de
cada carga \emph{positiva}, com espaçamento angular uniforme. Para que a
densidade de linhas represente corretamente a intensidade, o número de linhas
por carga deve ser \emph{proporcional a $|q|$} --- oito linhas para $+q$,
dezesseis para $+2q$.\\
(c) \textbf{Dificuldades e soluções.}
\begin{itemize}
\item \textbf{Singularidades:} $|\vec E|\to\infty$ nas cargas. Solução: iniciar
a uma distância pequena mas finita, e interromper a integração ao entrar em um
raio mínimo em torno de uma carga negativa.
\item \textbf{Pontos de campo nulo:} ali $\hat E$ é indefinido e o algoritmo
trava ou oscila. Solução: detectar $|\vec E|<\epsilon$ e encerrar a linha.
\item \textbf{Acúmulo de erro:} o método de Euler desvia sistematicamente em
regiões de curvatura alta. Solução: Runge--Kutta e passo adaptativo, menor onde
o campo varia rápido.
\item \textbf{Linhas que escapam:} defina um domínio e encerre ao ultrapassá-lo.
\end{itemize}
\textbf{Verificação do resultado:} as linhas traçadas devem ser perpendiculares
às equipotenciais calculadas independentemente --- teste que detecta quase todo
erro de implementação.}
\end{questao}

\begin{questao}[4]{}
\textbf{(Limite de campo em capacitor.)} Um capacitor de placas paralelas com
$d=\SI{0.1}{\milli\meter}$ usa dielétrico de rigidez
$\SI{20}{\kilo\volt\per\milli\meter}$ e $\varepsilon_r=3$.
\begin{itensq}
\item Determine a tensão máxima e o campo correspondente.
\item Determine a densidade superficial de carga no limite.
\item Discuta o compromisso entre capacitância e tensão.
\end{itensq}
\dados{$\varepsilon_0=\pot{8.85}{-12}\,\si{\farad\per\meter}$}
\resposta{$U_{\max}=\SI{2}{\kilo\volt}$;
$E=\SI{2e7}{\volt\per\meter}$;
$\sigma=\SI{531}{\micro\coulomb\per\meter\squared}$}
\resolucao{(a)
\[
U_{\max}=E_{rup}\,d=20\times0{,}1=\SI{2}{\kilo\volt};
\qquad
E=\pot{2}{7}\,\si{\volt\per\meter}.
\]
(b) No dielétrico, $E=\dfrac{\sigma}{\varepsilon_r\varepsilon_0}$, logo
\[
\sigma=\varepsilon_r\varepsilon_0E
=3(\pot{8.85}{-12})(\pot{2}{7})
=\pot{5.31}{-4}\,\si{\coulomb\per\meter\squared}.
\]
(c) \textbf{O compromisso.} A capacitância por unidade de área é
$C/A=\varepsilon_r\varepsilon_0/d$: reduzir $d$ \emph{aumenta} $C$. Mas
$U_{\max}=E_{rup}d$: reduzir $d$ \emph{diminui} a tensão suportável.\\
O produto revela o que realmente importa:
\[
\frac{C\,U_{\max}}{A}=\varepsilon_r\varepsilon_0E_{rup},
\]
que \textbf{não depende de $d$}. E a energia por unidade de volume,
\[
u=\frac12\varepsilon_r\varepsilon_0E_{rup}^{2},
\]
também não. \textbf{A geometria não decide nada} --- só o material.\\
\textbf{Consequência:} não existe truque de projeto que aumente a energia
armazenada por volume. Todo avanço em capacitores vem de materiais com
$\varepsilon_r$ ou $E_{rup}$ maiores --- e esses dois parâmetros costumam ser
antagônicos, o que torna o problema difícil.}
\end{questao}

\begin{questao}[4]{}
\textbf{(Superposição vetorial em anel.)} Um anel de raio $R$ tem carga
$Q$ uniformemente distribuída.
\begin{itensq}
\item Explique por que o campo é nulo no centro.
\item Determine, sem integrar, o campo em um ponto muito distante no eixo.
\item Discuta o que ocorre se metade do anel for removida.
\end{itensq}
\resposta{Centro: $\vec E=\vec0$ por simetria; longe: $E\to\ke Q/x^{2}$;
com meio anel, campo não nulo no centro}
\resolucao{(a) \textbf{No centro,} para cada elemento de carga existe outro
diametralmente oposto, à mesma distância $R$. Suas contribuições têm módulos
iguais e sentidos opostos, cancelando-se aos pares:
\[
\vec E_{centro}=\vec 0.
\]
\textbf{Note:} o \emph{potencial} no centro não é nulo --- vale
$V=\ke Q/R$, pois potenciais somam-se sem cancelamento entre contribuições de
mesmo sinal. Mais um caso de $\vec E=\vec0$ com $V\ne0$.\\
(b) \textbf{Muito longe} ($x\gg R$), todos os elementos estão praticamente à
mesma distância $x$ e praticamente na mesma direção. O anel se comporta como uma
carga puntiforme:
\[
E\to\frac{\ke Q}{x^{2}}.
\]
\textbf{Este é o teste que toda expressão deve passar:} qualquer distribuição
finita, vista de longe, deve reproduzir a lei de Coulomb.\\
(c) \textbf{Com meio anel}, a simetria de cancelamento se perde. Cada elemento
não tem mais seu par oposto, e o campo no centro \textbf{deixa de ser nulo}.\\
Pelo argumento da superposição ao contrário: como o anel inteiro dá campo nulo,
o campo do semianel restante é o \emph{negativo} do campo do semianel removido
--- ou seja, os dois têm o mesmo módulo. Integrando (como no caso do fio
semicircular), obtém-se
\[
E=\frac{2\ke Q'}{\pi R^{2}},
\]
com $Q'=Q/2$ a carga de cada metade, dirigido ao longo do eixo de simetria.}
\end{questao}

% END BANCO COMPLEMENTAR

\gabarito[1.3 Campo elétrico]

% =====================================================================
%  CAPITULO 1 - continuacao
%  1.4 Potencial eletrico e superficies equipotenciais
% =====================================================================
\subsection{Potencial elétrico e superfícies equipotenciais}

\begin{conceito}[Antes de começar]
\textbf{Potencial de uma carga puntiforme:}
\[
V = \ke\,\frac{q}{d}\quad(\si{\volt}),
\]
com o \emph{sinal} da carga (positivo para $q>0$, negativo para $q<0$).\par
\textbf{Diferença de potencial (ddp)} e trabalho: o trabalho para levar uma carga
$q$ entre dois pontos é $W = q\,U_{AB} = q\,(V_A - V_B)$.\par
\textbf{Campo uniforme:} $U = E\,d$, onde $d$ é a distância medida \emph{ao longo}
da direção do campo.\par
\textbf{Superfície equipotencial:} conjunto de pontos com o mesmo potencial. O
campo elétrico é sempre \textbf{perpendicular} às equipotenciais e aponta no
sentido dos potenciais \emph{decrescentes}. Não há trabalho ao mover uma carga
sobre uma mesma equipotencial.
\end{conceito}

\legendaniveis

\blocofixacao

\begin{questao}[1]{}
Sobre superfícies equipotenciais, é correto afirmar que:
\begin{alternativas}
\item o campo elétrico é tangente a elas.
\item o campo elétrico é perpendicular a elas.
\item o potencial varia ao longo de uma mesma superfície.
\item é necessário realizar trabalho para mover uma carga sobre a superfície.
\item elas sempre coincidem com as linhas de campo.
\end{alternativas}
\resposta{(b)}
\resolucao{Por definição, o potencial é constante sobre uma equipotencial (isso
elimina c). Como não há variação de potencial ao longo dela, o trabalho para mover
uma carga sobre a superfície é nulo (elimina d). O campo, que aponta na direção de
maior variação do potencial, é sempre \emph{perpendicular} às equipotenciais ---
nunca tangente, o que elimina (a) e (e).}
\end{questao}

\begin{questao}[1]{}
Considere dois pontos A e B de um campo elétrico uniforme de intensidade
$E = \SI{e4}{\newton\per\coulomb}$, separados por $\SI{5}{\centi\meter}$ ao longo
da direção do campo. Calcule a ddp entre A e B.
\resposta{$U_{AB} = \SI{500}{\volt}$}
\resolucao{Em campo uniforme, $U = E\,d = \pot{1}{4}\cdot0{,}05 = \SI{500}{\volt}$
(com A no potencial maior, a favor do campo).}
\end{questao}

\blocoaplicacao

\begin{questao}[2]{}
A figura representa três superfícies equipotenciais de um campo elétrico
uniforme, com $V_A = \SI{60}{\volt}$, $V_B = \SI{40}{\volt}$ e
$V_C = \SI{20}{\volt}$, igualmente espaçadas de $\SI{2}{\centi\meter}$.

\circuitoq{%
\begin{tikzpicture}[scale=1,line width=0.8pt,>=Stealth]
 \foreach \x/\v in {0/A,2/B,4/C}{
   \draw[azulprof,thick] (\x,-1.3) -- (\x,1.3);
   \node[above,azulprof] at (\x,1.3) {$V_{\v}$};}
 \node[below] at (0,-1.3) {\footnotesize$\SI{60}{\volt}$};
 \node[below] at (2,-1.3) {\footnotesize$\SI{40}{\volt}$};
 \node[below] at (4,-1.3) {\footnotesize$\SI{20}{\volt}$};
 % campo (do maior para o menor potencial)
 \foreach \y in {-0.7,0,0.7}
   \draw[->,Vcol,thick] (0.3,\y) -- (3.7,\y);
 \node[Vcol] at (2,1.0) {$\vec{E}$};
 \draw[<->,cinza] (0,-1.9) -- (2,-1.9) node[midway,below]{\footnotesize$\SI{2}{\centi\meter}$};
\end{tikzpicture}}{Superfícies equipotenciais em campo uniforme.}

\begin{itensq}
\item Determine a intensidade do campo elétrico.
\item Qual é o trabalho para levar uma carga de $\SI{5}{\micro\coulomb}$ de A até C?
\end{itensq}
\resposta{(a) $E = \SI{1000}{\volt\per\meter}$; (b) $W = \SI{2e-4}{\joule}$}
\resolucao{(a) Entre superfícies adjacentes, $U = \SI{20}{\volt}$ em
$d = \SI{2}{\centi\meter}$:
\[
E=\frac{U}{d}=\frac{20}{0{,}02}=\SI{1000}{\volt\per\meter}.
\]
(b) $U_{AC} = V_A - V_C = 60 - 20 = \SI{40}{\volt}$:
\[
W = q\,U_{AC} = \pot{5}{-6}\cdot40 = \pot{2}{-4}\,\si{\joule}=\SI{2e-4}{\joule}.
\]
O trabalho é positivo porque a carga positiva se move no sentido do campo
(de maior para menor potencial).}
\end{questao}

\begin{questao}[2]{}
A figura mostra linhas equipotenciais no campo gerado por duas cargas de mesmo
módulo e sinais opostos. São dados $V_A = \SI{30}{\volt}$, $V_B = \SI{10}{\volt}$
e $V_C = \SI{-10}{\volt}$. Calcule a ddp entre A e B, e entre B e C.

\circuitoq{%
\begin{tikzpicture}[scale=0.85,line width=0.8pt]
 % duas cargas
 \shade[ball color=Vcol!25] (-3,0) circle (0.3); \node at (-3,0){\footnotesize$+$};
 \shade[ball color=Icol!25] (3,0) circle (0.3);  \node at (3,0){\footnotesize$-$};
 % equipotenciais (arcos estilizados)
 \foreach \r/\lab/\x in {1.2/A/-1.6, 2.0/B/0, 1.2/C/1.6}{}
 \draw[azulprof,thick] (-1.6,-1.6) .. controls (-2.2,0) .. (-1.6,1.6);
 \node[azulprof] at (-1.6,1.85) {$A$};
 \draw[azulprof,thick] (0,-1.9) -- (0,1.9);
 \node[azulprof] at (0,2.15) {$B$};
 \draw[azulprof,thick] (1.6,-1.6) .. controls (2.2,0) .. (1.6,1.6);
 \node[azulprof] at (1.6,1.85) {$C$};
 \node[below,font=\footnotesize] at (-1.6,-1.7) {$\SI{30}{\volt}$};
 \node[below,font=\footnotesize] at (0,-2.0) {$\SI{10}{\volt}$};
 \node[below,font=\footnotesize] at (1.6,-1.7) {$\SI{-10}{\volt}$};
\end{tikzpicture}}{Linhas equipotenciais de um dipolo elétrico.}

\resposta{$U_{AB} = \SI{20}{\volt}$; $U_{BC} = \SI{20}{\volt}$}
\resolucao{A ddp é sempre a diferença dos potenciais:
\[
U_{AB}=V_A-V_B=30-10=\SI{20}{\volt},
\qquad
U_{BC}=V_B-V_C=10-(-10)=\SI{20}{\volt}.
\]
Embora $U_{AB}=U_{BC}$, note que $V_C$ é \emph{negativo}: perto da carga negativa
o potencial fica abaixo de zero. O ponto médio B, equidistante das duas cargas de
módulos iguais, tem potencial intermediário.}
\end{questao}

\begin{questao}[2]{}
As linhas cheias da figura representam linhas de força de um campo eletrostático
e as tracejadas, linhas equipotenciais, com $V_A > V_B > V_C$. Uma partícula de
carga $q = \SI{2e-6}{\coulomb}$ é levada de A até C.

\circuitoq{%
\begin{tikzpicture}[scale=1,line width=0.8pt,>=Stealth]
 % linhas equipotenciais (verticais tracejadas)
 \foreach \x/\lab/\v in {0/A/{V_A},2.2/B/{V_B},4.4/C/{V_C}}{
   \draw[azulprof,dashed,thick] (\x,-1.4) -- (\x,1.4);
   \node[azulprof,above] at (\x,1.4) {$\v$};
   \fill[laranja] (\x,0) circle (0.07);
   \node[below,font=\footnotesize] at (\x,-0.15) {\lab};}
 % linhas de forca (horizontais com seta)
 \foreach \y in {-1,0,1}
   \draw[Vcol,->,thick] (-0.6,\y) -- (5.0,\y);
 \node[Vcol] at (5.2,0.7) {$\vec{E}$};
\end{tikzpicture}}{Linhas de força (cheias) e equipotenciais (tracejadas).}

Sobre o trabalho da força elétrica nesse deslocamento, é correto afirmar que:
\begin{alternativas}
\item é nulo, pois A e C estão na mesma linha de força.
\item é positivo, pois a carga se move no sentido do campo.
\item é negativo, pois a carga se move contra o campo.
\item depende do caminho seguido entre A e C.
\item é máximo se a carga for levada de A até B.
\end{alternativas}
\resposta{(b)}
\resolucao{A carga positiva se desloca de A (maior potencial) para C (menor
potencial), ou seja, \emph{a favor} do campo elétrico. O trabalho da força
elétrica é $W = q(V_A - V_C) > 0$. O trabalho \emph{não} depende do caminho (força
conservativa), o que elimina (d).}
\end{questao}

\begin{questao}[2]{}
A figura representa as linhas equipotenciais e as linhas de força do campo gerado
por uma carga puntiforme $Q$ fixa no vácuo. As equipotenciais são círculos
concêntricos.

\circuitoq{%
\begin{tikzpicture}[scale=0.9,line width=0.8pt,>=Stealth]
 \shade[ball color=Vcol!30] (0,0) circle (0.28);
 \node at (0,0) {\footnotesize$Q$};
 \foreach \r in {1,1.7,2.4}
   \draw[azulprof,dashed] (0,0) circle (\r);
 \foreach \a in {0,45,...,315}
   \draw[Vcol,->,thick] (\a:0.35) -- (\a:2.7);
 \fill[laranja] (0:1) circle (0.06) node[above right,font=\footnotesize]{A};
 \fill[laranja] (0:2.4) circle (0.06) node[above right,font=\footnotesize]{C};
\end{tikzpicture}}{Campo e equipotenciais de uma carga puntiforme positiva.}

Como as linhas de força \emph{apontam para fora} de $Q$, pode-se concluir que:
\begin{alternativas}
\item $Q$ é negativa e o potencial aumenta com a distância.
\item $Q$ é positiva e o potencial diminui com a distância.
\item $Q$ é positiva e o potencial aumenta com a distância.
\item o potencial é o mesmo em A e C.
\item o campo é uniforme.
\end{alternativas}
\resposta{(b)}
\resolucao{Linhas de força saindo indicam carga \textbf{positiva}. O potencial de
uma carga puntiforme, $V = \ke Q/d$, \emph{diminui} à medida que a distância
cresce: $V_A > V_C$. As equipotenciais, mais afastadas, têm menor potencial. O
campo não é uniforme --- é radial e enfraquece com $1/d^2$.}
\end{questao}

\begin{questao}[2]{}
A figura mostra a configuração das linhas de força do campo gerado por duas
partículas carregadas A e B. As linhas \emph{saem} de A e \emph{entram} em B.

\circuitoq{%
\begin{tikzpicture}[scale=0.85,line width=0.7pt,>=Stealth]
 \shade[ball color=Vcol!25] (-2,0) circle (0.3); \node at (-2,0){\footnotesize A};
 \shade[ball color=Icol!25] ( 2,0) circle (0.3); \node at (2,0){\footnotesize B};
 \foreach \a in {30,60,90,120,150,-30,-60,-90,-120,-150}
   \draw[Vcol,->] (-2,0) .. controls (\a:1.3) and ({180-\a}:1.3) .. (2,0);
 \draw[Vcol,->] (-2,0)--(-3.3,0);
 \draw[Vcol,->] (3.3,0)--(2,0);
\end{tikzpicture}}{Linhas de força de duas cargas A e B.}

Os sinais das cargas A e B são, respectivamente:
\begin{alternativas}[2]
\item positivo e positivo.
\item positivo e negativo.
\item negativo e positivo.
\item negativo e negativo.
\item ambos nulos.
\end{alternativas}
\resposta{(b)}
\resolucao{As linhas de força \emph{nascem} nas cargas positivas e \emph{morrem}
nas negativas. Como saem de A, esta é \textbf{positiva}; como entram em B, esta é
\textbf{negativa}. A configuração é a de um dipolo elétrico.}
\end{questao}

\begin{questao}[2]{}
Uma carga elétrica cria, em um ponto P situado a $\SI{20}{\centi\meter}$ dela, um
campo de intensidade $\SI{900}{\newton\per\coulomb}$. O módulo do potencial
elétrico em P é:
\begin{alternativas}[3]
\item \SI{45}{\volt}
\item \SI{90}{\volt}
\item \SI{180}{\volt}
\item \SI{450}{\volt}
\item \SI{900}{\volt}
\end{alternativas}
\resposta{(c)}
\resolucao{Para uma carga puntiforme, $E = \dfrac{\ke Q}{d^2}$ e
$V = \dfrac{\ke Q}{d}$. Dividindo, $\dfrac{V}{E} = d$, logo
\[
V = E\cdot d = 900\cdot0{,}20 = \SI{180}{\volt}.
\]
Essa relação $V = E\,d$ vale para carga puntiforme; não confundir com o campo
uniforme, onde também $U = E\,d$, mas por outro motivo.}
\end{questao}

\begin{questao}[2]{}
Próximo ao solo, em um dia limpo, a atmosfera apresenta um campo elétrico
aproximadamente uniforme, vertical, apontando para baixo, de intensidade
$\SI{120}{\volt\per\meter}$. Considere dois pontos M (mais alto) e N (mais baixo),
separados verticalmente por $\SI{10}{\meter}$.

\circuitoq{%
\begin{tikzpicture}[scale=0.9,line width=0.8pt,>=Stealth]
 \fill[cinza!25] (-2.5,-1.6) rectangle (2.5,-1.2);
 \node[cinza,font=\footnotesize] at (0,-1.4) {solo};
 \foreach \x in {-2,-1,0,1,2}
   \draw[Vcol,->,thick] (\x,1.4) -- (\x,-1.1);
 \node[Vcol] at (2.5,0.4) {$\vec{E}$};
 \fill[laranja] (-1,1.1) circle (0.08) node[left,font=\footnotesize]{M};
 \fill[laranja] (-1,-0.9) circle (0.08) node[left,font=\footnotesize]{N};
 \draw[<->,cinza] (0.6,1.1)--(0.6,-0.9) node[midway,right,font=\footnotesize]{$\SI{10}{\meter}$};
\end{tikzpicture}}{Campo elétrico atmosférico próximo ao solo.}

A diferença de potencial entre M e N vale:
\begin{alternativas}[2]
\item \SI{12}{\volt}
\item \SI{120}{\volt}
\item \SI{1200}{\volt}
\item \SI{12000}{\volt}
\item zero
\end{alternativas}
\resposta{(c)}
\resolucao{Em campo uniforme, $U = E\,d = 120\cdot10 = \SI{1200}{\volt}$. Como o
campo aponta para baixo (de M para N), M está a maior potencial:
$V_M - V_N = \SI{1200}{\volt}$.}
\end{questao}

\blocoaprofundamento

\begin{questao}[3]{}
Duas cargas $Q_1 = \SI{+2}{\micro\coulomb}$ e $Q_2 = \SI{-2}{\micro\coulomb}$
estão fixas, separadas por $\SI{20}{\centi\meter}$. Determine o potencial elétrico
no ponto médio M do segmento que as une.
\dados{$\ke = \SI{9e9}{\newton\meter\squared\per\coulomb\squared}$.}
\resposta{$V_M = 0$}
\resolucao{O potencial é uma grandeza \emph{escalar}: soma-se com sinal. No ponto
médio, cada carga está a $d = \SI{10}{\centi\meter}$:
\[
V_M = \ke\frac{Q_1}{d}+\ke\frac{Q_2}{d}
= \ke\frac{(+2)+(-2)}{d}\,\si{\micro\coulomb}=0.
\]
O potencial é nulo, mas o \emph{campo} em M \emph{não} é zero --- os vetores campo
das duas cargas apontam no mesmo sentido e se somam. Este é o ponto central: campo
nulo e potencial nulo são condições independentes.}
\end{questao}

\begin{questao}[3]{}
Uma carga puntiforme $Q = \SI{+5}{\micro\coulomb}$ está fixa no vácuo. Uma segunda
carga $q = \SI{+2}{\micro\coulomb}$ é trazida do infinito até um ponto a
$\SI{30}{\centi\meter}$ de $Q$.
\begin{itensq}
\item Qual é o potencial criado por $Q$ nesse ponto?
\item Qual é o trabalho realizado pela força externa para trazer $q$ (supondo-a
      trazida lentamente)?
\end{itensq}
\dados{$\ke = \SI{9e9}{\newton\meter\squared\per\coulomb\squared}$;
$V_\infty = 0$.}
\resposta{(a) $V = \SI{1.5e5}{\volt}$; (b) $W = \SI{0.3}{\joule}$}
\resolucao{(a) $V = \ke\dfrac{Q}{d}
= 9\cdot10^9\cdot\dfrac{\pot{5}{-6}}{0{,}3}=\pot{1.5}{5}\,\si{\volt}$.\\
(b) O trabalho da força externa iguala a energia potencial adquirida (partindo do
infinito, onde $V=0$):
\[
W = q\,(V - V_\infty) = \pot{2}{-6}\cdot\pot{1.5}{5} = \SI{0.3}{\joule}.
\]
Como ambas as cargas são positivas, elas se repelem, e o trabalho externo é
\emph{positivo} (precisamos empurrar $q$ contra a repulsão). Se as cargas tivessem
sinais opostos, o trabalho externo seria negativo.}
\end{questao}

\begin{questao}[3]{}
Um elétron ($q = \SI{-1.6e-19}{\coulomb}$, $m = \SI{9.1e-31}{\kilogram}$) parte do
repouso e é acelerado por uma ddp de $\SI{200}{\volt}$. Determine a energia
cinética adquirida e a velocidade final do elétron.
\resposta{$E_c = \SI{3.2e-17}{\joule}$; $v \approx \SI{8.4e6}{\meter\per\second}$}
\resolucao{Toda a energia potencial elétrica se converte em energia cinética:
\[
E_c = |q|\,U = \pot{1.6}{-19}\cdot200 = \pot{3.2}{-17}\,\si{\joule}.
\]
Da definição de energia cinética, $E_c = \tfrac12 m v^2$:
\[
v=\sqrt{\frac{2E_c}{m}}=\sqrt{\frac{2\cdot\pot{3.2}{-17}}{\pot{9.1}{-31}}}
\approx\SI{8.4e6}{\meter\per\second}.
\]
É o princípio de funcionamento dos antigos tubos de raios catódicos e dos
aceleradores de partículas: usar uma ddp para acelerar cargas.}
\end{questao}

\begin{questao}[3]{}
O diagrama representa o potencial elétrico $V$ criado por uma carga puntiforme em
função da distância $d$ até ela.

\circuitoqq{%
\begin{tikzpicture}
 \begin{axis}[width=9cm,height=5cm,xlabel={$d$ (m)},ylabel={$V$ (V)},
   xmin=0,xmax=6,ymin=0,ymax=90,xtick={0,1,2,3,4},ytick={0,20,40,80},
   grid=both,axis lines=left,domain=0.9:6,samples=100]
   \addplot[Vcol,very thick] {80/x};
   \addplot[only marks,mark=*,Vcol] coordinates {(1,80)(2,40)(4,20)};
   \node[Vcol,anchor=south west,font=\footnotesize] at (axis cs:2,40){$(2,\,40)$};
 \end{axis}
\end{tikzpicture}}{Potencial de uma carga puntiforme em função da distância.}

Sabendo que, a $\SI{2}{\meter}$, o potencial vale $\SI{40}{\volt}$, o potencial a
$\SI{4}{\meter}$ e o campo elétrico a $\SI{2}{\meter}$ valem, respectivamente:
\begin{alternativas}[2]
\item \SI{20}{\volt} e \SI{20}{\volt\per\meter}
\item \SI{20}{\volt} e \SI{40}{\volt\per\meter}
\item \SI{10}{\volt} e \SI{10}{\volt\per\meter}
\item \SI{40}{\volt} e \SI{40}{\volt\per\meter}
\item \SI{10}{\volt} e \SI{20}{\volt\per\meter}
\end{alternativas}
\resposta{(a)}
\resolucao{O potencial varia com $1/d$: dobrando a distância (de \SI{2}{} para
\SI{4}{\meter}), ele cai à metade, $V(4) = \SI{20}{\volt}$.\\
Para o campo, usamos $E = \dfrac{V}{d}$ (válido para carga puntiforme):
$E(2) = \dfrac{40}{2} = \SI{20}{\volt\per\meter}$.}
\end{questao}

\begin{questao}[3]{}
Na figura, uma partícula de carga $Q > 0$ está fixa em P. As superfícies A, B e C
são equipotenciais, e um ponto D situa-se sobre uma linha de força entre B e C.

\circuitoq{%
\begin{tikzpicture}[scale=0.9,line width=0.8pt,>=Stealth]
 \shade[ball color=Vcol!30] (0,0) circle (0.25); \node at (0,0){\footnotesize P};
 \foreach \r/\lab in {1.1/A,1.8/B,2.5/C}{
   \draw[azulprof,dashed] (0,0) circle (\r);
   \node[azulprof,font=\footnotesize] at (\r,0.28) {\lab};}
 \draw[Vcol,->,thick] (20:0.3)--(20:2.7);
 \fill[laranja] (20:2.15) circle (0.07) node[right,font=\footnotesize]{D};
\end{tikzpicture}}{Equipotenciais de uma carga puntiforme positiva.}

Sobre os potenciais, é correto afirmar que:
\begin{alternativas}
\item $V_A < V_B < V_C$.
\item $V_A > V_B > V_C$.
\item $V_A = V_B = V_C$.
\item $V_C$ é o maior, por estar mais longe.
\item nada se pode afirmar sem o valor de $Q$.
\end{alternativas}
\resposta{(b)}
\resolucao{Para carga positiva, $V = \ke Q/d$ \emph{decresce} com a distância. Como
A é a superfície mais interna (menor $d$) e C a mais externa (maior $d$):
$V_A > V_B > V_C$. O sinal de $Q$ determina a tendência; o \emph{valor} de $Q$ não
é necessário para ordenar os potenciais.}
\end{questao}

\begin{questao}[3]{}
Uma esfera condutora de raio $R = \SI{10}{\centi\meter}$ está carregada com
$Q = \SI{6}{\micro\coulomb}$, isolada no vácuo. Os gráficos mostram o campo
elétrico e o potencial em função da distância $d$ ao centro.

\circuitoqq{%
\begin{tikzpicture}
 \begin{axis}[width=6.2cm,height=4.2cm,xlabel={$d$ (m)},ylabel={$E$ (kN/C)},
   xmin=0,xmax=0.45,ymin=0,ymax=6200,xtick={0,0.1,0.2,0.3,0.4},
   ytick={0,2000,4000,5400},yticklabels={0,2000,4000,5400},
   grid=both,axis lines=left,domain=0.1:0.45,samples=80,
   scaled y ticks=false,y tick label style={/pgf/number format/fixed}]
   \addplot[Vcol,very thick] {9e9*6e-6/(x^2)/1000};
   \addplot[Vcol,very thick] coordinates {(0,0)(0.1,0)};
   \addplot[Vcol,dashed] coordinates {(0.1,0)(0.1,5400)};
 \end{axis}
\end{tikzpicture}\hspace{4mm}
\begin{tikzpicture}
 \begin{axis}[width=6.2cm,height=4.2cm,xlabel={$d$ (m)},ylabel={$V$ (kV)},
   xmin=0,xmax=0.45,ymin=0,ymax=620,xtick={0,0.1,0.2,0.3,0.4},
   ytick={0,180,540},grid=both,axis lines=left,domain=0.1:0.45,samples=80]
   \addplot[Icol,very thick] {9e9*6e-6/x/1000};
   \addplot[Icol,very thick] coordinates {(0,540)(0.1,540)};
 \end{axis}
\end{tikzpicture}}{Campo e potencial de uma esfera condutora carregada.}

\begin{itensq}
\item Calcule $E$ e $V$ na superfície da esfera.
\item Determine $E$ e $V$ a $\SI{30}{\centi\meter}$ do centro.
\item Explique por que, \emph{dentro} da esfera, $E = 0$ mas $V \neq 0$.
\end{itensq}
\dados{$\ke = \SI{9e9}{\newton\meter\squared\per\coulomb\squared}$.}
\resposta{(a) $E=\SI{5.4e6}{\newton\per\coulomb}$, $V=\SI{540}{\kilo\volt}$;
(b) $E=\SI{6e5}{\newton\per\coulomb}$, $V=\SI{180}{\kilo\volt}$; (c) ver comentário}
\resolucao{(a) Fora da esfera (inclusive na superfície), ela se comporta como uma
carga puntiforme concentrada no centro:
\[
E(R)=\frac{\ke Q}{R^{2}}=\frac{9\cdot10^9\cdot\pot{6}{-6}}{(0{,}1)^2}
=\pot{5.4}{6}\,\si{\newton\per\coulomb},
\]
\[
V(R)=\frac{\ke Q}{R}=\frac{9\cdot10^9\cdot\pot{6}{-6}}{0{,}1}
=\pot{5.4}{5}\,\si{\volt}=\SI{540}{\kilo\volt}.
\]
(b) A $\SI{30}{\centi\meter}$ (três vezes o raio): $E$ cai com $1/d^2$ (fator 9) e
$V$ com $1/d$ (fator 3):
\[
E=\frac{5{,}4\cdot10^6}{9}=\pot{6}{5}\,\si{\newton\per\coulomb},
\qquad
V=\frac{540}{3}=\SI{180}{\kilo\volt}.
\]
(c) \textbf{A distinção conceitual central.} Em um condutor em equilíbrio
eletrostático, toda a carga se distribui na \emph{superfície} e o campo interno é
nulo --- se não fosse, as cargas livres se moveriam e não haveria equilíbrio.\\
Mas o potencial \emph{não} é nulo: ele é a \emph{integral} do campo desde o
infinito até o ponto. Ao entrar na esfera, o campo passa a ser zero, então o
potencial \textbf{para de variar} --- ele permanece constante e igual ao valor da
superfície:
\[
V_{\text{interno}} = V(R) = \SI{540}{\kilo\volt}.
\]
Isso é coerente com $E = -\dfrac{\mathrm{d}V}{\mathrm{d}d}$: campo nulo significa
potencial \emph{constante}, não necessariamente nulo. Todo o volume do condutor é
uma única região equipotencial.\\
\textbf{Aplicação:} é o princípio da \textbf{gaiola de Faraday} --- dentro de um
condutor fechado o campo é nulo, protegendo o conteúdo, ainda que o conjunto esteja
a um potencial elevadíssimo. É por isso que se está seguro dentro de um carro
atingido por um raio.}
\end{questao}

\begin{questao}[3]{}
Três cargas puntiformes iguais, $q = \SI{1}{\micro\coulomb}$, são trazidas
\emph{uma a uma} desde o infinito até ocuparem os vértices de um triângulo
equilátero de lado $L = \SI{10}{\centi\meter}$.
\begin{itensq}
\item Calcule o trabalho externo necessário para trazer cada carga.
\item Determine a energia potencial total da configuração.
\item Se as cargas forem liberadas simultaneamente, qual é a energia cinética total
      quando estiverem infinitamente afastadas?
\end{itensq}
\dados{$\ke = \SI{9e9}{\newton\meter\squared\per\coulomb\squared}$.}
\resposta{(a) $0$, \SI{0.09}{\joule}, \SI{0.18}{\joule};
(b) $U = \SI{0.27}{\joule}$; (c) \SI{0.27}{\joule}}
\resolucao{\textbf{(a) Trabalho carga a carga.}
\begin{itemize}[leftmargin=1.6em,itemsep=1pt,topsep=2pt]
\item \textbf{1\textsuperscript{a} carga:} não há campo algum ainda, logo
      $W_1 = 0$.
\item \textbf{2\textsuperscript{a} carga:} vem contra o campo da primeira, a uma
      distância final $L$:
      \[
      W_2=\ke\frac{q^2}{L}=9\cdot10^9\frac{(\pot{1}{-6})^2}{0{,}1}=\SI{0.09}{\joule}.
      \]
\item \textbf{3\textsuperscript{a} carga:} sofre a ação das \emph{duas} anteriores,
      ambas à distância $L$:
      \[
      W_3=2\,\ke\frac{q^2}{L}=\SI{0.18}{\joule}.
      \]
\end{itemize}
\textbf{(b) Energia total da configuração.} É a soma dos trabalhos --- ou,
equivalentemente, a soma sobre os \emph{três pares} distintos:
\[
U = W_1+W_2+W_3 = 0+0{,}09+0{,}18 = \SI{0.27}{\joule}
= 3\,\ke\frac{q^2}{L}. \checkmark
\]
Note que a resposta \emph{não depende} da ordem em que as cargas foram trazidas
--- consequência de a força elétrica ser conservativa.

\textbf{(c)} Ao serem liberadas, toda a energia potencial armazenada se converte
em cinética (conservação de energia, sem dissipação):
\[
E_{c,\text{total}} = U = \SI{0.27}{\joule}.
\]
\textbf{Cuidado com o erro clássico:} a energia total \emph{não} é
$3\times W_3$ nem a soma de "cada carga vezes o potencial das outras" sem dividir
por 2 --- isso contaria cada par duas vezes. A fórmula geral é
$U=\tfrac12\sum_i q_i V_i$, onde $V_i$ é o potencial criado \emph{pelas demais}
na posição de $i$.}
\end{questao}

\begin{questao}[3]{}
Duas esferas condutoras, de raios $R$ e $3R$, estão distantes entre si e são
ligadas por um fio condutor longo e fino. A carga total do sistema é $Q$.
\begin{itensq}
\item Determine a carga final de cada esfera.
\item Compare os campos elétricos nas superfícies das duas esferas.
\item Explique o "poder das pontas" a partir desse resultado.
\end{itensq}
\resposta{(a) $Q_1=Q/4$, $Q_2=3Q/4$; (b) $E_1/E_2=3$; (c) ver comentário}
\resolucao{(a) Ligadas por um condutor, as esferas atingem o \emph{mesmo
potencial}:
\[
\frac{\ke Q_1}{R}=\frac{\ke Q_2}{3R}
\;\Longrightarrow\;
Q_2 = 3Q_1.
\]
Com $Q_1+Q_2=Q$: $Q_1=\dfrac{Q}{4}$ e $Q_2=\dfrac{3Q}{4}$.\\
(b) Campos nas superfícies:
\[
E_1=\frac{\ke Q_1}{R^{2}}=\frac{\ke Q}{4R^{2}},
\qquad
E_2=\frac{\ke Q_2}{(3R)^{2}}=\frac{3\ke Q}{4\cdot9R^{2}}=\frac{\ke Q}{12R^{2}}.
\]
\[
\frac{E_1}{E_2}=3.
\]
A esfera \textbf{menor}, embora tenha \emph{menos} carga, apresenta campo
\textbf{três vezes mais intenso} em sua superfície.\\
\textbf{Relação geral:} como $V=\ke Q/R$ e $E=\ke Q/R^{2}$, tem-se
$E = V/R$. Com $V$ comum às duas, \emph{o campo é inversamente proporcional ao
raio}.\\
(c) \textbf{Poder das pontas.} Uma "ponta" é uma região de raio de curvatura
muito pequeno. Sendo $E \propto 1/R$, o campo ali fica extremamente intenso ---
suficiente para ionizar o ar e provocar descarga, mesmo com o objeto a um
potencial moderado.\\
\textbf{Consequências práticas:}
\begin{itemize}[leftmargin=1.6em,itemsep=1pt,topsep=2pt]
\item \textbf{Para-raios} são pontiagudos \emph{de propósito}: concentram o campo
      e induzem a descarga em local controlado.
\item Equipamentos de alta tensão têm cantos \emph{arredondados} e cúpulas
      polidas, justamente para \emph{evitar} campos intensos e o efeito corona.
\end{itemize}}
\end{questao}

\blocodesafio

\begin{questao}[4]{}
Quatro cargas iguais $q=\SI{1}{\micro\coulomb}$ ocupam os vértices de um quadrado
de lado $a=\SI{10}{\centi\meter}$.
\begin{itensq}
\item Calcule a energia potencial elétrica total da configuração.
\item Quanto trabalho externo seria necessário para trazer uma quinta carga $q$
      desde o infinito até o centro do quadrado?
\item Se todas fossem liberadas, qual seria a energia cinética total no infinito?
\end{itensq}
\dados{$\ke=\SI{9e9}{}$; $\sqrt2\approx1{,}41$.}
\resposta{(a) $U\approx\SI{0.49}{\joule}$; (b) $W\approx\SI{0.51}{\joule}$;
(c) $\approx\SI{0.49}{\joule}$}
\resolucao{(a) Há $\binom{4}{2}=6$ pares: \textbf{4 lados} (distância $a$) e
\textbf{2 diagonais} (distância $a\sqrt2$):
\[
U=4\cdot\frac{\ke q^{2}}{a}+2\cdot\frac{\ke q^{2}}{a\sqrt2}
=\frac{\ke q^{2}}{a}\left(4+\frac{2}{\sqrt2}\right)
=\frac{\ke q^{2}}{a}\left(4+\sqrt2\right).
\]
Com $\dfrac{\ke q^{2}}{a}=\dfrac{9\cdot10^{9}(\pot{1}{-6})^{2}}{0{,}1}
=\SI{0.09}{\joule}$:
\[
U=0{,}09\,(4+1{,}41)\approx\SI{0.49}{\joule}.
\]
(b) A quinta carga, no centro, dista $\dfrac{a\sqrt2}{2}=\SI{0.0707}{\meter}$ de
\emph{cada} uma das quatro. O potencial no centro é
\[
V_c=4\cdot\frac{\ke q}{a\sqrt2/2}
=4\cdot\frac{9\cdot10^{9}\cdot\pot{1}{-6}}{0{,}0707}
\approx\pot{5.09}{5}\,\si{\volt},
\]
e o trabalho externo é
\[
W=q\,V_c=\pot{1}{-6}\cdot\pot{5.09}{5}\approx\SI{0.51}{\joule}.
\]
(c) Sem a quinta carga, ao liberar as quatro, toda a energia potencial da
configuração converte-se em cinética:
\[
E_{c}=U\approx\SI{0.49}{\joule}.
\]
\textbf{Atenção ao método.} Em (a) somam-se os \emph{pares} (cada par uma vez
só); em (b) usa-se $W=qV$ porque as quatro cargas já estavam posicionadas e
imóveis. Confundir os dois procedimentos --- contando pares duas vezes --- é o
erro mais comum neste tipo de problema.}
\end{questao}

\begin{questao}[4]{}
Uma esfera condutora oca de raio interno $a$ e raio externo $b$ envolve, no
centro, uma carga puntiforme $+Q$. A casca está inicialmente \emph{neutra} e
isolada.
\begin{itensq}
\item Determine as cargas induzidas nas superfícies interna e externa da casca.
\item Esboce o comportamento de $E(r)$ nas quatro regiões ($r<a$, $a<r<b$,
      $r>b$).
\item Se a casca for \textbf{aterrada}, o que muda?
\end{itensq}
\resposta{(a) $-Q$ na interna, $+Q$ na externa; (b) e (c) ver comentário}
\resolucao{(a) Dentro do \emph{material} do condutor ($a<r<b$) o campo deve ser
\textbf{nulo} em equilíbrio. Aplicando a lei de Gauss a uma superfície esférica
nessa região, a carga interna total deve ser zero:
\[
Q + q_{\text{int}} = 0 \;\Longrightarrow\; q_{\text{int}}=-Q.
\]
Como a casca é neutra no total, a superfície externa fica com
$q_{\text{ext}}=+Q$.\\
(b) \textbf{Comportamento do campo:}
\begin{center}\small\sffamily
\begin{tabular}{@{}l l l@{}}
\toprule
\textbf{Região} & \textbf{Carga envolvida} & \textbf{Campo}\\
\midrule
$r<a$      & $+Q$        & $E=\ke Q/r^{2}$ (radial, para fora)\\
$a<r<b$    & $Q-Q=0$     & $E=0$ (dentro do condutor)\\
$r>b$      & $Q-Q+Q=+Q$  & $E=\ke Q/r^{2}$\\
\bottomrule
\end{tabular}
\end{center}
Fora da casca, o campo é \emph{idêntico} ao que existiria sem ela --- a casca
neutra não blinda o exterior.\\
(c) \textbf{Aterrando a casca}, a superfície externa perde sua carga $+Q$ para a
terra (os elétrons sobem para neutralizá-la). Resultado:
\[
q_{\text{int}}=-Q,\qquad q_{\text{ext}}=0
\;\Longrightarrow\;
E = 0 \ \text{ para } r>b.
\]
\textbf{Agora sim há blindagem completa}: nenhum campo escapa para fora.\\
\textbf{Aplicação --- blindagem eletrostática.} É por isso que cabos coaxiais e
gabinetes de equipamentos sensíveis têm sua malha metálica \textbf{aterrada}:
uma blindagem flutuante (não aterrada) não impede o campo de vazar. O aterramento
é o que transforma o condutor em uma verdadeira gaiola de Faraday.}
\end{questao}

% BEGIN BANCO COMPLEMENTAR Potencial elétrico e superfícies equipotenciais
% =====================================================================
%  Banco complementar --- Potencial Eletrico e Superficies
%  Equipotenciais  (REVISADO)
%  Substitui a versao gerada por template: as 6 questoes N1 eram o
%  mesmo enunciado ("determine V a r de uma carga Q") com numeros
%  trocados, as 5 N2 eram "carga q atravessa uma ddp" repetido, e as
%  4 N3 eram "no ponto P as distancias as cargas q1 e q2 sao..."
%  igualmente clonado. As 6 N4 eram frases conceituais de uma linha.
%  O cap01c ja cobre: equipotenciais (MC), campo uniforme entre A e B,
%  tres equipotenciais, dipolo com linhas, linhas de forca com
%  equipotenciais, carga puntiforme, duas particulas A e B, campo a
%  20 cm, campo atmosferico, cargas +-2 uC, trabalho, eletron
%  acelerado, grafico V x d, superficies A/B/C, esfera condutora de
%  R=10 cm, tres cargas trazidas do infinito, duas esferas ligadas por
%  fio, quadrado de quatro cargas e esfera oca com carga central.
%  Distribuicao mantida: 6 N1 + 5 N2 + 4 N3 + 6 N4 = 21
% =====================================================================

\subsubsection*{Banco complementar --- Potencial elétrico e superfícies equipotenciais}

\begin{atencao}[Organização do banco]
O potencial é \textbf{escalar}: contribuições de várias cargas somam-se
algebricamente, com sinal, sem decomposição vetorial --- o que torna $V$ muito
mais fácil de calcular que $\vec E$. Em compensação, $V$ carrega menos
informação: pontos de $V=0$ podem ter campo intenso, e é o \emph{gradiente},
não o valor, que determina a força. O N4 explora essa distinção, a
independência do caminho e a estabilidade.
\end{atencao}

\blocofixacao

\begin{questao}[1]{}
Determine o potencial a $\SI{30}{\centi\meter}$ de uma carga
$Q=+\SI{2}{\micro\coulomb}$, adotando $V(\infty)=0$.
\dados{$\ke=\SI{9e9}{\newton\meter\squared\per\coulomb\squared}$}
\resposta{$V=\SI{60}{\kilo\volt}$}
\resolucao{
\[
V=\frac{\ke Q}{r}=\frac{\pot{9}{9}\times\pot{2}{-6}}{0{,}30}
=\frac{\pot{1.8}{4}}{0{,}30}=\pot{6}{4}\,\si{\volt}=\SI{60}{\kilo\volt}.
\]
\textbf{Atenção à potência de $r$:} o potencial cai com $1/r$, enquanto o campo
cai com $1/r^{2}$. Confundir as duas é o erro mais comum do tópico --- e o
resultado sai com a dimensão errada.}
\end{questao}

\begin{questao}[1]{}
Determine o potencial a $\SI{60}{\centi\meter}$ de uma carga
$Q=-\SI{4}{\micro\coulomb}$.
\resposta{$V=\SI{-60}{\kilo\volt}$}
\resolucao{
\[
V=\frac{\pot{9}{9}\times(-\pot{4}{-6})}{0{,}60}=-\pot{6}{4}\,\si{\volt}
=\SI{-60}{\kilo\volt}.
\]
\textbf{O sinal é do potencial, não do módulo.} Diferentemente do campo --- cuja
\emph{orientação} é dada por um vetor --- o potencial é um escalar com sinal
próprio. Cargas negativas produzem potencial negativo em toda parte.}
\end{questao}

\begin{questao}[1]{}
A que distância de $Q=+\SI{3}{\micro\coulomb}$ o potencial vale
$\SI{45}{\kilo\volt}$?
\resposta{$r=\SI{0.6}{\meter}$}
\resolucao{
\[
r=\frac{\ke Q}{V}=\frac{\pot{9}{9}\times\pot{3}{-6}}{\pot{4.5}{4}}
=\frac{\pot{2.7}{4}}{\pot{4.5}{4}}=\SI{0.6}{\meter}.
\]
\textbf{Superfície equipotencial:} \emph{todos} os pontos a
$\SI{0.6}{\meter}$ da carga estão a $\SI{45}{\kilo\volt}$. Em torno de uma
carga puntiforme, as equipotenciais são esferas concêntricas.}
\end{questao}

\begin{questao}[1]{}
Determine a carga que produz $\SI{-18}{\kilo\volt}$ a
$\SI{0.5}{\meter}$ de distância.
\resposta{$Q=-\SI{1}{\micro\coulomb}$}
\resolucao{
\[
Q=\frac{Vr}{\ke}=\frac{(-\pot{1.8}{4})(0{,}5)}{\pot{9}{9}}
=-\pot{1}{-6}\,\si{\coulomb}=-\SI{1}{\micro\coulomb}.
\]
O sinal negativo do potencial revelou diretamente o sinal da carga --- vantagem
prática de trabalhar com uma grandeza escalar.}
\end{questao}

\begin{questao}[1]{}
Sobre o potencial criado por várias cargas em um ponto, é correto afirmar:
\begin{alternativas}[2]
\item soma-se vetorialmente
\item soma-se algebricamente, com sinal
\item soma-se em módulo
\item depende do caminho até o ponto
\end{alternativas}
\resposta{(b)}
\resolucao{O potencial é \textbf{escalar}: basta somar as contribuições
$\ke q_k/r_k$ com seus sinais, sem ângulos nem componentes.\\
\textbf{Vantagem prática enorme:} calcular $V$ de dez cargas exige dez divisões
e uma soma; calcular $\vec E$ exigiria dez vetores decompostos em componentes.\\
A alternativa (d) descreve a propriedade oposta --- justamente por o campo
eletrostático ser \emph{conservativo}, o potencial de um ponto independe do
caminho, e é isso que permite defini-lo.}
\end{questao}

\begin{questao}[1]{}
Uma carga $q=+\SI{3}{\micro\coulomb}$ é colocada em um ponto onde
$V=\SI{200}{\volt}$. Determine sua energia potencial elétrica.
\resposta{$U=\SI{0.6}{\milli\joule}$}
\resolucao{
\[
U=qV=\pot{3}{-6}\times200=\pot{6}{-4}\,\si{\joule}=\SI{0.6}{\milli\joule}.
\]
\textbf{Distinção essencial:} $V$ é uma propriedade \emph{do ponto} --- existe
mesmo sem carga alguma ali. Já $U$ é uma propriedade \emph{do sistema}
carga-campo, e só faz sentido quando a carga está presente.\\
A relação $U=qV$ é o análogo elétrico de $E_p=mgh$: potencial faz o papel de
$gh$, e carga faz o papel de massa.}
\end{questao}

\blocoaplicacao

\begin{questao}[2]{}
Uma carga $q=+\SI{2}{\micro\coulomb}$ desloca-se de um ponto a
$\SI{50}{\volt}$ para outro a $\SI{20}{\volt}$.
\begin{itensq}
\item Determine o trabalho realizado pela força elétrica.
\item Determine a variação da energia potencial.
\end{itensq}
\resposta{$W=\SI{60}{\micro\joule}$; $\Delta U=\SI{-60}{\micro\joule}$}
\resolucao{(a)
\[
W=q\left(V_A-V_B\right)=\pot{2}{-6}(50-20)=\pot{6}{-5}\,\si{\joule}
=\SI{60}{\micro\joule}.
\]
(b) $\Delta U=-W=\SI{-60}{\micro\joule}$.\\
\textbf{Coerência dos sinais:} a carga é positiva e foi de maior para menor
potencial --- movimento espontâneo. A força elétrica realiza trabalho
\emph{positivo} e a energia potencial \emph{diminui}, convertendo-se em energia
cinética.\\
\textbf{Regra:} cargas positivas "descem" o potencial espontaneamente; cargas
negativas "sobem".}
\end{questao}

\begin{questao}[2]{}
Duas cargas estão fixas: $q_1=+\SI{3}{\micro\coulomb}$ a
$\SI{30}{\centi\meter}$ do ponto $P$, e $q_2=-\SI{2}{\micro\coulomb}$ a
$\SI{20}{\centi\meter}$ de $P$. Determine o potencial em $P$.
\resposta{$V=0$}
\resolucao{Somando algebricamente:
\[
V=\frac{\pot{9}{9}(\pot{3}{-6})}{0{,}30}
+\frac{\pot{9}{9}(-\pot{2}{-6})}{0{,}20}
=\pot{9}{4}-\pot{9}{4}=\SI{0}{\volt}.
\]
\textbf{Cuidado:} $V=0$ \textbf{não} significa campo nulo nem ausência de força.
As duas cargas produzem campos que \emph{não} se cancelam --- elas têm sinais
opostos, e seus campos apontam para o mesmo lado ao longo da linha que as une.\\
Uma carga colocada em $P$ teria energia potencial nula, mas sofreria força.
Essa distinção é aprofundada no bloco seguinte.}
\end{questao}

\begin{questao}[2]{}
Em um campo uniforme de $E=\SI{500}{\volt\per\meter}$, dois pontos estão
separados por $\SI{4}{\centi\meter}$ ao longo da direção do campo. Determine a
diferença de potencial.
\resposta{$U=\SI{20}{\volt}$}
\resolucao{
\[
U=E\,d=500\times0{,}04=\SI{20}{\volt},
\]
com o ponto \emph{a montante} do campo em potencial maior --- o campo aponta
sempre no sentido de $V$ decrescente.\\
\textbf{Se o deslocamento fosse perpendicular} ao campo, a diferença seria
\emph{nula}: os dois pontos estariam na mesma equipotencial. Só a componente do
deslocamento \emph{ao longo} do campo contribui.\\
\textbf{Unidades:} $\si{\volt\per\meter}$ e $\si{\newton\per\coulomb}$ são a
mesma coisa --- verifique multiplicando por metro e por coulomb.}
\end{questao}

\begin{questao}[2]{}
Um elétron é acelerado através de uma diferença de potencial de
$\SI{500}{\volt}$.
\begin{itensq}
\item Determine a energia adquirida, em joules e em elétron-volts.
\item Determine sua velocidade final, partindo do repouso.
\end{itensq}
\dados{$e=\pot{1.6}{-19}\,\si{\coulomb}$;
$m_e=\pot{9.11}{-31}\,\si{\kilogram}$}
\resposta{$\SI{500}{\electronvolt}=\pot{8}{-17}\,\si{\joule}$;
$v=\pot{1.33}{7}\,\si{\meter\per\second}$}
\resolucao{(a)
\[
W=eU=\pot{1.6}{-19}(500)=\pot{8}{-17}\,\si{\joule}.
\]
Em elétron-volts, a leitura é imediata: $\SI{500}{\electronvolt}$ --- \emph{por
definição}, $\SI{1}{\electronvolt}$ é a energia que a carga elementar adquire
sob $\SI{1}{\volt}$.\\
(b) Toda a energia vira cinética:
\[
v=\sqrt{\frac{2W}{m}}=\sqrt{\frac{2(\pot{8}{-17})}{\pot{9.11}{-31}}}
=\sqrt{\pot{1.756}{14}}=\pot{1.33}{7}\,\si{\meter\per\second}.
\]
\textbf{Verificação de consistência:} isso é $4{,}4\%$ da velocidade da luz ---
ainda no regime não relativístico, onde $\frac12mv^{2}$ é válido. Acima de cerca
de $\SI{50}{\kilo\electronvolt}$, a correção relativística passa a ser
necessária.}
\end{questao}

\begin{questao}[2]{}
Duas cargas, $q_1=+\SI{4}{\micro\coulomb}$ em $x=0$ e
$q_2=-\SI{1}{\micro\coulomb}$ em $x=\SI{0.5}{\meter}$, estão fixas. Determine o
ponto do segmento entre elas onde o potencial é nulo.
\resposta{$x=\SI{0.4}{\meter}$}
\resolucao{Impondo $V=0$ com $x$ medido a partir de $q_1$:
\[
\frac{\ke(4\mu)}{x}+\frac{\ke(-1\mu)}{0{,}5-x}=0
\;\Longrightarrow\;
\frac{4}{x}=\frac{1}{0{,}5-x}.
\]
\[
4(0{,}5-x)=x
\;\Longrightarrow\;
2=5x
\;\Longrightarrow\;
x=\SI{0.4}{\meter}.
\]
\textbf{Contraste importante com o campo:} o ponto de \emph{campo} nulo dessas
mesmas cargas fica \emph{fora} do segmento (a $\SI{1}{\meter}$ de $q_1$),
porque o campo é vetorial e só pode se cancelar onde os dois vetores forem
opostos. Já o potencial é escalar e se anula sempre que as contribuições tiverem
módulos iguais e sinais opostos --- o que ocorre dentro do segmento.\\
\textbf{Existe também um segundo ponto de $V=0$} fora do segmento, do lado da
carga menor, em $x=\SI{0.667}{\meter}$. A superfície completa de $V=0$ é uma
esfera.}
\end{questao}

\blocoaprofundamento

\begin{questao}[3]{}
Duas cargas $+Q$ e $-Q$, com $Q=\SI{2}{\micro\coulomb}$, estão separadas por
$\SI{40}{\centi\meter}$.
\begin{itensq}
\item Determine o potencial no ponto médio.
\item Determine o campo no ponto médio.
\item Explique a aparente contradição.
\end{itensq}
\resposta{$V=0$; $E=\SI{900}{\kilo\volt\per\meter}$}
\resolucao{(a) As duas contribuições têm módulos iguais e sinais opostos:
\[
V=\frac{\ke Q}{0{,}20}-\frac{\ke Q}{0{,}20}=0.
\]
(b) Os \emph{campos}, ao contrário, apontam ambos da carga positiva para a
negativa --- eles se \textbf{somam}:
\[
E=2\times\frac{\ke Q}{(0{,}20)^{2}}
=2\times\frac{\pot{9}{9}(\pot{2}{-6})}{0{,}04}
=\pot{9}{5}\,\si{\volt\per\meter}.
\]
(c) \textbf{Não há contradição.} O potencial não determina o campo --- quem o
determina é o \emph{gradiente} do potencial:
\[
\vec E=-\nabla V.
\]
Uma função pode valer zero em um ponto e ter derivada grande ali. É exatamente o
caso: caminhando ao longo do eixo, $V$ passa de positivo a negativo cruzando o
zero com inclinação acentuada.\\
\textbf{Analogia:} a altitude do nível do mar é zero, mas isso nada diz sobre a
inclinação do terreno naquele ponto. É a inclinação --- e não a altitude --- que
faz uma bola rolar.\\
\textbf{Consequência:} uma carga solta no ponto médio \emph{não} fica parada,
apesar de $V=0$ e $U=0$.}
\end{questao}

\begin{questao}[3]{}
O potencial ao longo de um eixo é dado por
$V(x)=\SI{200}{}-\SI{50}{}\,x$ (com $V$ em volts e $x$ em metros), no trecho
$0\le x\le\SI{4}{\meter}$.
\begin{itensq}
\item Determine o campo elétrico.
\item Determine o trabalho para levar $q=+\SI{2}{\micro\coulomb}$ de
      $x=1$ a $x=\SI{3}{\meter}$.
\item Identifique as superfícies equipotenciais.
\end{itensq}
\resposta{$E=\SI{50}{\volt\per\meter}$ no sentido $+x$;
$W=\SI{200}{\micro\joule}$}
\resolucao{(a)
\[
E_x=-\frac{\mathrm{d}V}{\mathrm{d}x}=-(-50)=\SI{+50}{\volt\per\meter}.
\]
Campo uniforme, apontando no sentido \emph{crescente} de $x$ --- ou seja, no
sentido em que $V$ \emph{decresce}, como sempre.\\
(b) $V(1)=\SI{150}{\volt}$ e $V(3)=\SI{50}{\volt}$:
\[
W=q\left(V_1-V_3\right)=\pot{2}{-6}(150-50)=\SI{200}{\micro\joule}.
\]
\emph{(Conferindo por força: $W=qE\,d=\pot{2}{-6}(50)(2)
=\SI{200}{\micro\joule}$ \checkmark.)}\\
(c) Como $V$ depende \emph{apenas} de $x$, as equipotenciais são
\textbf{planos perpendiculares ao eixo} $x$. O plano $x=\SI{1}{\meter}$ é toda
a superfície com $V=\SI{150}{\volt}$.\\
\textbf{Regra geral:} equipotenciais são sempre \textbf{perpendiculares} às
linhas de campo. Se não fossem, haveria componente de campo ao longo da
superfície, e mover uma carga sobre ela exigiria trabalho --- contradizendo a
definição de equipotencial.}
\end{questao}

\begin{questao}[3]{}
Uma carga $q=+\SI{1}{\micro\coulomb}$ é trazida do infinito até
$\SI{20}{\centi\meter}$ de uma carga fixa $Q=+\SI{5}{\micro\coulomb}$.
\begin{itensq}
\item Determine o trabalho realizado por um agente externo.
\item Determine a energia potencial do par.
\item Interprete o sinal.
\end{itensq}
\resposta{$W_{ext}=U=\SI{0.225}{\joule}$}
\resolucao{(a) e (b) Como a carga parte do repouso no infinito e chega em
repouso, todo o trabalho externo vira energia potencial:
\[
U=\frac{\ke Qq}{r}=\frac{\pot{9}{9}(\pot{5}{-6})(\pot{1}{-6})}{0{,}20}
=\frac{\pot{4.5}{-2}}{0{,}20}=\SI{0.225}{\joule}.
\]
\emph{(Equivalentemente, $U=qV$ com $V=\ke Q/r=\SI{225}{\kilo\volt}$.)}\\
(c) \textbf{Sinal positivo:} as cargas se repelem, então aproximá-las exige
trabalho \emph{contra} a força elétrica. O agente externo gasta energia, que
fica armazenada no sistema.\\
\textbf{Se as cargas tivessem sinais opostos}, $U$ seria negativa: elas se
atraem, e seria preciso \emph{segurá-las} durante a aproximação --- o agente
externo receberia energia.\\
\textbf{Verificação de coerência:} soltando a carga, ela é repelida e ganha
$\SI{0.225}{\joule}$ de energia cinética ao voltar ao infinito. A energia
potencial é integralmente recuperável --- marca de um campo conservativo.}
\end{questao}

\begin{questao}[3]{}
Compare, para duas cargas positivas iguais, o lugar geométrico dos pontos de
potencial nulo e o dos pontos de campo nulo.
\begin{itensq}
\item Determine onde $V=0$.
\item Determine onde $\vec E=\vec 0$.
\item Explique a diferença.
\end{itensq}
\resposta{$V$ nunca se anula (exceto no infinito); $\vec E$ se anula no ponto
médio}
\resolucao{(a) Ambas as cargas são positivas, logo cada uma contribui com
potencial \textbf{positivo} em todo ponto finito. A soma de dois positivos nunca
é zero:
\[
V=\frac{\ke Q}{r_1}+\frac{\ke Q}{r_2}>0\quad\text{sempre}.
\]
$V\to0$ apenas no limite $r\to\infty$.\\
(b) O campo é vetorial e \emph{pode} se cancelar. No ponto médio do segmento,
os dois vetores têm módulos iguais e sentidos opostos:
\[
\vec E=\vec 0\ \text{no ponto médio}.
\]
(c) \textbf{A diferença é a natureza das grandezas.}
\begin{itemize}
\item $V$ é escalar e, para cargas de mesmo sinal, as contribuições não podem se
cancelar --- só se somam.
\item $\vec E$ é vetorial e se cancela sempre que os vetores forem opostos, o que
não impõe restrição alguma sobre os sinais das cargas.
\end{itemize}
\textbf{Contraste com o caso de cargas opostas} ($+Q$ e $-Q$): ali ocorre
exatamente o inverso --- $V=0$ no ponto médio, mas $\vec E\neq\vec 0$.\\
\textbf{Síntese:} $V=0$ e $\vec E=\vec 0$ são condições \emph{independentes}.
Nenhuma implica a outra, em nenhum sentido.}
\end{questao}

\blocodesafio

\begin{questao}[4]{}
\textbf{(Demonstração --- trabalho e independência do caminho.)} Prove que o
trabalho da força elétrica entre dois pontos independe do caminho, e deduza a
relação com a diferença de potencial.
\begin{itensq}
\item Calcule o trabalho para o campo de uma carga puntiforme.
\item Generalize para uma distribuição qualquer.
\item Explique por que isso permite \emph{definir} o potencial.
\end{itensq}
\resposta{$W_{AB}=q(V_A-V_B)$, com
$V=\ke Q/r$}
\resolucao{(a) Para uma carga fixa $Q$ e uma carga de prova $q$, a força é
radial: $\vec F=\dfrac{\ke Qq}{r^{2}}\hat r$. Em um deslocamento
$\mathrm{d}\vec\ell$, apenas a componente radial contribui, pois
$\vec F\cdot\mathrm{d}\vec\ell=F\,\mathrm{d}r$:
\[
W_{AB}=\int_{r_A}^{r_B}\frac{\ke Qq}{r^{2}}\,\mathrm{d}r
=\ke Qq\left[-\frac{1}{r}\right]_{r_A}^{r_B}
=\ke Qq\left(\frac{1}{r_A}-\frac{1}{r_B}\right).
\]
\textbf{O resultado depende apenas de $r_A$ e $r_B$} --- as componentes
transversais do caminho não contribuem, por serem perpendiculares à força.\\
(b) \textbf{Generalização:} para $n$ cargas, a força total é a soma vetorial das
forças individuais, e o trabalho total é a soma dos trabalhos. Como cada parcela
independe do caminho, a soma também independe.\\
Definindo $V=\sum\ke Q_k/r_k$, obtém-se
\[
\boxed{W_{AB}=q\left(V_A-V_B\right)}
\]
(c) \textbf{Por que isso permite definir $V$.} Se o trabalho dependesse do
caminho, não haveria como atribuir um \emph{número} a cada ponto --- o "potencial"
seria ambíguo. A independência do caminho garante que
$V(P)=W_{\infty\to P}/q$ é bem definido.\\
\textbf{Formulação equivalente:} $\oint\vec E\cdot\mathrm{d}\vec\ell=0$ em
qualquer curva fechada. O campo eletrostático é \textbf{conservativo}, e essa é
exatamente a lei das malhas de Kirchhoff em sua forma de campo.\\
\textbf{Ressalva importante:} isso vale para campos \emph{eletrostáticos}.
Campos elétricos induzidos por variação de fluxo magnético \emph{não} são
conservativos, e para eles o potencial escalar deixa de existir --- assunto da
lei de Faraday.}
\end{questao}

\begin{questao}[4]{}
\textbf{(Projeto --- $V=0$ com $E\neq0$.)} Projete um par de cargas de sinais
opostos e módulos \emph{diferentes} tal que exista um ponto do segmento onde
$V=0$ mas $\vec E\neq\vec0$.
\begin{itensq}
\item Determine a condição geral.
\item Escolha valores e localize o ponto.
\item Verifique que o campo não é nulo ali.
\end{itensq}
\resposta{Condição: $|q_1|/r_1=|q_2|/r_2$; com $+4\,\si{\micro\coulomb}$ e
$-1\,\si{\micro\coulomb}$ a $\SI{0.5}{\meter}$, o ponto é
$x=\SI{0.4}{\meter}$}
\resolucao{(a) \textbf{Condição para $V=0$:}
\[
\frac{\ke q_1}{r_1}+\frac{\ke q_2}{r_2}=0
\;\Longrightarrow\;
\frac{|q_1|}{r_1}=\frac{|q_2|}{r_2},
\]
com $q_1$ e $q_2$ de sinais \emph{opostos}.\\
\textbf{Condição para $\vec E=\vec0$ no mesmo ponto:}
$\dfrac{|q_1|}{r_1^{2}}=\dfrac{|q_2|}{r_2^{2}}$.\\
As duas só coincidem se $r_1=r_2$ --- e então $|q_1|=|q_2|$, que é o caso
excluído pelo enunciado. \textbf{Com módulos diferentes, é impossível anular
ambos no mesmo ponto.}\\
(b) Tome $q_1=+\SI{4}{\micro\coulomb}$ em $x=0$ e
$q_2=-\SI{1}{\micro\coulomb}$ em $x=\SI{0.5}{\meter}$. De
$\dfrac{4}{x}=\dfrac{1}{0{,}5-x}$ vem $x=\SI{0.4}{\meter}$.\\
\emph{(Confira: $r_1=0{,}4$ e $r_2=0{,}1$, com
$4/0{,}4=1/0{,}1=10$ \checkmark.)}\\
(c) \textbf{Campo no ponto.} Ambos apontam no sentido $+x$ (afastando-se de
$q_1$, aproximando-se de $q_2$) e portanto se somam:
\[
E=\frac{\pot{9}{9}(\pot{4}{-6})}{0{,}16}
+\frac{\pot{9}{9}(\pot{1}{-6})}{0{,}01}
=\pot{2.25}{5}+\pot{9}{5}=\pot{1.125}{6}\,\si{\volt\per\meter}.
\]
Mais de um megavolt por metro --- longe de nulo.\\
\textbf{Interpretação:} uma carga colocada ali teria energia potencial nula, mas
sofreria força intensa. \textbf{$V$ e $\vec E$ carregam informações
diferentes:} $V$ é o valor, $\vec E$ é a taxa de variação.}
\end{questao}

\begin{questao}[4]{}
\textbf{(Estabilidade em um mínimo de potencial.)} Uma curva
$V(x)$ apresenta mínimo local em $x_0$.
\begin{itensq}
\item Analise a estabilidade de uma carga positiva colocada em $x_0$.
\item Repita para uma carga negativa.
\item Relacione com o teorema de Earnshaw.
\end{itensq}
\resposta{Positiva: \textbf{instável}; negativa: estável em $x$ --- mas nenhuma
é estável em três dimensões}
\resolucao{(a) A energia potencial de uma carga positiva é $U=qV$, com
$q>0$ --- ela tem a \emph{mesma} forma de $V$ e, portanto, também um mínimo em
$x_0$. Mínimo de energia significa \textbf{equilíbrio estável}... \emph{ao longo
de $x$}.\\
\emph{(Atenção ao raciocínio: é o mínimo de $U$, e não o de $V$, que define
estabilidade. Para carga positiva os dois coincidem; para negativa, invertem-se.)}\\
(b) Para carga negativa, $U=qV$ com $q<0$ \textbf{inverte} a curva: onde $V$
tem mínimo, $U$ tem \emph{máximo}. Equilíbrio \textbf{instável} em $x$.\\
(c) \textbf{Teorema de Earnshaw.} Em uma região sem cargas, o potencial
eletrostático satisfaz a equação de Laplace, $\nabla^{2}V=0$, ou seja
\[
\frac{\partial^{2}V}{\partial x^{2}}
+\frac{\partial^{2}V}{\partial y^{2}}
+\frac{\partial^{2}V}{\partial z^{2}}=0.
\]
Se $V$ tivesse um mínimo \emph{tridimensional} em $x_0$, as três segundas
derivadas seriam positivas e a soma não poderia ser zero. \textbf{Contradição.}\\
Logo, um mínimo em uma direção obriga a haver \emph{máximo} em outra --- todo
ponto crítico é um \textbf{ponto de sela}.\\
\textbf{Conclusão:} nenhuma carga pode ser mantida em equilíbrio estável apenas
por forças eletrostáticas, seja ela positiva ou negativa. A análise
unidimensional do item (a) é enganosa justamente por ignorar as outras duas
direções.\\
\textbf{Como se contorna na prática:} campos variáveis no tempo (armadilha de
Paul), campos magnéticos, ou realimentação ativa --- todas soluções que escapam
da hipótese eletrostática.}
\end{questao}

\begin{questao}[4]{}
\textbf{(Condutor em equilíbrio.)} Demonstre que toda a superfície de um
condutor em equilíbrio eletrostático é equipotencial.
\begin{itensq}
\item Prove a partir da condição de equilíbrio.
\item Mostre que o interior também está no mesmo potencial.
\item Deduza a orientação do campo na superfície.
\end{itensq}
\resposta{Campo interno nulo $\Rightarrow$ $V$ constante em todo o condutor; o
campo externo é \textbf{perpendicular} à superfície}
\resolucao{(a) \textbf{Condição de equilíbrio:} não há movimento de cargas
livres, logo o campo dentro do metal é \textbf{nulo}. Se houvesse campo, as
cargas se moveriam e o equilíbrio não existiria.\\
Tomando dois pontos $A$ e $B$ quaisquer do condutor e um caminho
\emph{interno} entre eles:
\[
V_A-V_B=\int_A^B\vec E\cdot\mathrm{d}\vec\ell=\int_A^B\vec 0\cdot\mathrm{d}\vec\ell=0.
\]
Portanto $V_A=V_B$ --- \textbf{todo o condutor é equipotencial}.\\
(b) O argumento não distinguiu superfície de interior: ele vale para
\emph{qualquer} par de pontos ligados por caminho dentro do metal. Superfície e
interior estão no mesmo potencial, e o condutor inteiro é um volume
equipotencial.\\
(c) \textbf{Orientação do campo externo.} Se houvesse componente
\emph{tangencial} à superfície, ela empurraria as cargas livres ao longo dela --
contradizendo o equilíbrio. Logo o campo externo é \textbf{perpendicular} à
superfície em todo ponto.\\
\textbf{Consequências, todas verificáveis:}
\begin{itemize}
\item A carga em excesso reside \emph{integralmente} na superfície externa.
\item A densidade superficial é maior onde a curvatura é maior (poder das
pontas).
\item Uma cavidade vazia dentro do condutor tem campo nulo --- é a blindagem
eletrostática.
\item Ligar dois condutores por fio iguala seus potenciais, e a carga se
redistribui proporcionalmente às capacitâncias.
\end{itemize}}
\end{questao}

\begin{questao}[4]{}
\textbf{(Do potencial ao campo.)} O potencial em uma região é
$V(x,y)=\SI{100}{}-\SI{20}{}\,x+\SI{5}{}\,y^{2}$ (volts, com $x,y$ em metros).
\begin{itensq}
\item Determine $\vec E(x,y)$.
\item Avalie em $(x,y)=(2;\,1)$.
\item Determine a forma das equipotenciais.
\end{itensq}
\resposta{$\vec E=(20;\,-10y)\,\si{\volt\per\meter}$; em $(2;1)$,
$\vec E=(20;\,-10)$ com $|E|=\SI{22.4}{\volt\per\meter}$}
\resolucao{(a) $\vec E=-\nabla V$, componente a componente:
\[
E_x=-\frac{\partial V}{\partial x}=-(-20)=\SI{+20}{\volt\per\meter};
\qquad
E_y=-\frac{\partial V}{\partial y}=-(10y)=-10y.
\]
Note que $E_x$ é \emph{uniforme}, enquanto $E_y$ cresce linearmente com $y$.\\
(b) Em $(2;\,1)$:
\[
\vec E=(20;\,-10)\,\si{\volt\per\meter};
\qquad
|E|=\sqrt{400+100}=\SI{22.4}{\volt\per\meter},
\]
formando $\ang{-26.6}$ com o eixo $x$.\\
\textbf{Observe:} o valor de $V$ no ponto ($V=100-40+5=\SI{65}{\volt}$) não
entrou no cálculo do campo. Só as \emph{derivadas} importam --- somar uma
constante a $V$ não altera $\vec E$.\\
(c) \textbf{Equipotenciais:} $100-20x+5y^{2}=\text{const}$, ou seja
\[
x=\frac{5y^{2}+100-C}{20},
\]
uma família de \textbf{parábolas} com eixo horizontal.\\
\textbf{Verificação de coerência:} o vetor $\vec E$ calculado deve ser
perpendicular à parábola que passa por $(2;1)$. A tangente à curva ali tem
inclinação $\mathrm{d}y/\mathrm{d}x=2$, e o campo tem inclinação
$-10/20=-1/2$ --- produto $-1$, portanto perpendiculares \checkmark.}
\end{questao}

\begin{questao}[4]{}
\textbf{(Equipotenciais e linhas de campo.)} Compare os dois conjuntos de
curvas para duas cargas positivas iguais.
\begin{itensq}
\item Descreva as equipotenciais perto de cada carga e muito longe do par.
\item Descreva as linhas de campo, incluindo o ponto médio.
\item Estabeleça as regras gerais que relacionam os dois conjuntos.
\end{itensq}
\resposta{Perto: esferas em torno de cada carga; longe: esferas em torno do par;
o ponto médio é ponto de sela do potencial}
\resolucao{(a) \textbf{Equipotenciais.}
\begin{itemize}
\item \emph{Muito perto de uma carga:} a contribuição dela domina, e as
superfícies são praticamente \textbf{esferas} centradas nessa carga.
\item \emph{Muito longe do par:} o sistema parece uma única carga $2Q$, e as
superfícies voltam a ser esferas, agora centradas no \textbf{centro do par}.
\item \emph{Na região intermediária:} as superfícies deixam de ser esferas e
assumem forma de "amendoim", chegando a se fundir em uma única superfície que
envolve as duas cargas.
\end{itemize}
(b) \textbf{Linhas de campo.} Saem radialmente de cada carga, curvam-se
afastando-se uma da outra e seguem para o infinito. \textbf{No ponto médio o
campo é nulo} --- as linhas não passam por ali. Não há linha ligando as duas
cargas, pois ambas são fontes.\\
(c) \textbf{Regras gerais.}
\begin{enumerate}
\item As linhas de campo são \textbf{perpendiculares} às equipotenciais em todo
ponto.
\item O campo aponta no sentido de $V$ \textbf{decrescente}.
\item Onde as equipotenciais estão \textbf{mais próximas}, o campo é mais
\textbf{intenso} --- pois $|E|=|\mathrm{d}V/\mathrm{d}n|$.
\item Equipotenciais nunca se cruzam (um ponto teria dois potenciais); linhas de
campo também não (haveria duas direções de força).
\item Onde $\vec E=\vec0$, as duas regras de perpendicularidade perdem sentido:
é um \textbf{ponto de sela} do potencial, e ali as equipotenciais podem se
cruzar.
\end{enumerate}
\textbf{O ponto médio é justamente esse caso excepcional:} $V$ tem mínimo ao
longo da linha que une as cargas e máximo na direção perpendicular --- exatamente
a estrutura de sela exigida pelo teorema de Earnshaw.}
\end{questao}

% END BANCO COMPLEMENTAR

\gabarito[1.4 Potencial elétrico e superfícies equipotenciais]

          % (parcial: 1.1 Cargas eletricas migrada)

% Cap. 2 -------------------------------------------------------------
\setcounter{section}{1}
% =====================================================================
%  CAPITULO 2 - CONCEITOS INICIAIS
%  Leis de Ohm, Potencia, Fontes e Leis de Kirchhoff (versao revisada)
% =====================================================================
\section{Conceitos Iniciais: Leis de Ohm, Potência, Fontes e Kirchhoff}

\begin{conceito}[Antes de começar]
\textbf{Corrente elétrica:} $I = \dfrac{Q}{t}$, medida em ampères
($\SI{1}{\ampere}=\SI{1}{\coulomb\per\second}$). Como $Q = n\,e$, o número de
elétrons é $n = Q/e$, com $e=\SI{1.6e-19}{\coulomb}$.\par
\textbf{1\textsuperscript{a} lei de Ohm:} $U = R\,I$.\par
\textbf{2\textsuperscript{a} lei de Ohm (resistividade):}
$R = \rho\,\dfrac{L}{A}$, onde $\rho$ é a resistividade do material, $L$ o
comprimento e $A$ a área da seção.\par
\textbf{Potência elétrica:}
$P = U\,I = R\,I^2 = \dfrac{U^2}{R}$.\par
\textbf{Fonte real:} $U = \EMF - r\,I$, onde $r$ é a resistência interna.\par
\textbf{Leis de Kirchhoff:} \emph{dos nós} --- a soma das correntes que entram é
igual à soma das que saem; \emph{das malhas} --- a soma algébrica das tensões ao
longo de uma malha fechada é nula.
\end{conceito}

\legendaniveis

% =====================================================================
\subsection{Corrente elétrica e carga}
% =====================================================================

\blocofixacao

\begin{questao}[1]{}
Se uma corrente de $\SI{3}{\ampere}$ percorre um fio durante $\SI{2}{\minute}$, a
carga que atravessa sua seção reta, em coulombs, vale:
\begin{alternativas}[3]
\item 60
\item 110
\item 360
\item 220
\item 180
\end{alternativas}
\resposta{(c)}
\resolucao{Convertendo o tempo para segundos, $t = 2\cdot 60 = \SI{120}{\second}$:
\[
Q = I\,t = 3\cdot 120 = \SI{360}{\coulomb}.
\]
O erro típico (alternativa e) é usar $t = \SI{2}{\minute}$ diretamente, sem
converter para segundos.}
\end{questao}

\begin{questao}[1]{}
Uma carga de $\SI{120}{\coulomb}$ passa uniformemente pela seção de um fio durante
$\SI{1}{\minute}$. A intensidade da corrente é:
\begin{alternativas}[3]
\item $\tfrac{1}{30}\,\si{\ampere}$
\item $\tfrac{1}{2}\,\si{\ampere}$
\item \SI{2}{\ampere}
\item \SI{30}{\ampere}
\item \SI{120}{\ampere}
\end{alternativas}
\resposta{(c)}
\resolucao{$I = \dfrac{Q}{t} = \dfrac{120}{60} = \SI{2}{\ampere}$.}
\end{questao}

\begin{questao}[1]{}
Uma lâmpada permanece acesa durante $\SI{5}{\minute}$ sob corrente de
$\SI{2}{\ampere}$. A carga total que circula, em coulombs, é:
\begin{alternativas}[3]
\item 0,4
\item 2,5
\item 10
\item 150
\item 600
\end{alternativas}
\resposta{(e)}
\resolucao{$Q = I\,t = 2\cdot(5\cdot60) = 2\cdot 300 = \SI{600}{\coulomb}$.}
\end{questao}

\blocoaplicacao

\begin{questao}[2]{}
Um condutor metálico é percorrido por uma corrente contínua e constante de
$\SI{32}{\milli\ampere}$. Determine:
\begin{itensq}
\item a carga que atravessa uma seção reta por segundo;
\item o número de elétrons correspondente.
\end{itensq}
\dados{$e = \SI{1.6e-19}{\coulomb}$.}
\resposta{(a) \SI{32}{\milli\coulomb}; (b) $n = \pot{2.0}{17}$ elétrons}
\resolucao{(a) Em $t = \SI{1}{\second}$:
$Q = I\,t = \pot{32}{-3}\cdot 1 = \SI{32}{\milli\coulomb}$.\\
(b) $n = \dfrac{Q}{e} = \dfrac{\pot{32}{-3}}{\pot{1.6}{-19}} = \pot{2.0}{17}$
elétrons por segundo.}
\end{questao}

\begin{questao}[2]{}
O gráfico mostra a intensidade da corrente $i$ em um condutor em função do tempo.
Determine a carga que atravessa uma seção do condutor nos cinco primeiros
segundos.

\circuitoqq{%
\begin{tikzpicture}
 \begin{axis}[width=9cm,height=5cm,xlabel={$t$ (s)},ylabel={$i$ (A)},
   xmin=0,xmax=5.5,ymin=0,ymax=5,xtick={0,1,2,3,4,5},ytick={0,1,2,3,4},
   grid=both,axis lines=left]
   \addplot[Icol,very thick] coordinates {(0,4)(5,4)};
   \addplot[Icol!30] coordinates {(0,0)(0,4)};
 \end{axis}
\end{tikzpicture}}{Corrente constante em função do tempo.}

\resposta{$Q = \SI{20}{\coulomb}$}
\resolucao{A carga é a \emph{área} sob o gráfico $i \times t$. Como a corrente é
constante e igual a $\SI{4}{\ampere}$ durante $\SI{5}{\second}$, a área é um
retângulo:
\[
Q = i\,t = 4\cdot 5 = \SI{20}{\coulomb}.
\]
Essa interpretação por área vale mesmo quando a corrente varia --- basta calcular
a área da figura correspondente.}
\end{questao}

\begin{questao}[2]{}
O gráfico representa a corrente elétrica $i$ que atravessa um condutor em função
do tempo. Determine a carga total que passa por uma seção reta entre $t=0$ e
$t=\SI{6}{\second}$.

\circuitoqq{%
\begin{tikzpicture}
 \begin{axis}[width=9.5cm,height=5cm,xlabel={$t$ (s)},ylabel={$i$ (A)},
   xmin=0,xmax=6.5,ymin=0,ymax=5,xtick={0,1,2,3,4,5,6},ytick={0,2,4},
   grid=both,axis lines=left]
   \addplot[Icol,very thick] coordinates {(0,4)(2,4)(6,0)};
   \addplot[Icol!12] coordinates {(0,0)(0,4)};
 \end{axis}
\end{tikzpicture}}{Corrente variável em função do tempo.}

\resposta{$Q = \SI{16}{\coulomb}$}
\resolucao{A carga é a \emph{área} sob o gráfico $i \times t$. A figura é composta
de um retângulo (de $0$ a \SI{2}{\second}) e um triângulo (de $2$ a \SI{6}{\second}):
\[
Q = \underbrace{4\cdot2}_{\text{retângulo}}
  + \underbrace{\tfrac12\cdot4\cdot4}_{\text{triângulo}}
  = 8 + 8 = \SI{16}{\coulomb}.
\]
Mesmo com corrente variável, a área continua valendo como carga total.}
\end{questao}

\begin{questao}[2]{}
Qual é a carga elétrica que atravessa, durante $\SI{10}{\hour}$, a seção de um
condutor percorrido por corrente constante de $\SI{5}{\ampere}$?
\begin{alternativas}[3]
\item \SI{50}{\coulomb}
\item \SI{1000}{\coulomb}
\item \SI{90000}{\coulomb}
\item \SI{180000}{\coulomb}
\item \SI{360000}{\coulomb}
\end{alternativas}
\resposta{(d)}
\resolucao{Convertendo o tempo: $\SI{10}{\hour} = 10\cdot3600 = \SI{36000}{\second}$.
\[
Q = I\,t = 5\cdot36000 = \SI{180000}{\coulomb}.
\]
O erro comum é esquecer de converter horas em segundos.}
\end{questao}

\blocoaprofundamento

\begin{questao}[3]{}
O gráfico mostra a corrente em um condutor variando linearmente de
$\SI{6}{\ampere}$ a $\SI{2}{\ampere}$ entre $t=0$ e $t=\SI{8}{\second}$.

\circuitoqq{%
\begin{tikzpicture}
 \begin{axis}[width=9cm,height=4.6cm,xlabel={$t$ (s)},ylabel={$i$ (A)},
   xmin=0,xmax=9,ymin=0,ymax=7,xtick={0,2,4,6,8},ytick={0,2,4,6},
   grid=both,axis lines=left]
   \addplot[Icol,very thick] coordinates {(0,6)(8,2)};
 \end{axis}
\end{tikzpicture}}{Corrente decrescente em função do tempo.}

\begin{itensq}
\item Determine a carga total que atravessa o condutor nesse intervalo.
\item Calcule a corrente \emph{média} e verifique a coerência com o item (a).
\item Em que instante a corrente instantânea iguala a média?
\end{itensq}
\resposta{(a) $Q = \SI{32}{\coulomb}$; (b) $i_m = \SI{4}{\ampere}$;
(c) $t = \SI{4}{\second}$}
\resolucao{(a) A carga é a área sob o gráfico --- aqui, um trapézio de bases
\SI{6}{} e \SI{2}{\ampere} e altura \SI{8}{\second}:
\[
Q=\frac{(6+2)}{2}\cdot 8 = \SI{32}{\coulomb}.
\]
(b) $i_m = \dfrac{Q}{\Delta t} = \dfrac{32}{8} = \SI{4}{\ampere}$, que coincide
com a média aritmética $\dfrac{6+2}{2}$ --- válido \emph{porque} a variação é
linear.\\
(c) Como a reta decresce uniformemente, ela cruza o valor médio exatamente no
\emph{meio} do intervalo: $t = \SI{4}{\second}$.
Conferindo pela equação da reta $i(t) = 6 - 0{,}5t$:
$i(4) = 6-2 = \SI{4}{\ampere}$. \checkmark\\
\textbf{Cuidado:} se a corrente variasse de forma \emph{não} linear (por exemplo,
exponencialmente), a corrente média \textbf{não} seria a média dos extremos ---
seria preciso calcular a área real sob a curva.}
\end{questao}

\begin{questao}[3]{}
Um fio de cobre de seção $A = \SI{1}{\milli\meter\squared}$ conduz corrente
constante de $\SI{10}{\ampere}$. O cobre possui aproximadamente
$n = \SI{8.5e28}{}$ elétrons livres por metro cúbico.
\begin{itensq}
\item Determine a \emph{velocidade de deriva} dos elétrons.
\item Quanto tempo um elétron leva para percorrer $\SI{1}{\meter}$ de fio?
\item Como conciliar esse resultado com o fato de uma lâmpada acender
      \emph{instantaneamente} ao acionarmos o interruptor?
\end{itensq}
\dados{$e = \SI{1.6e-19}{\coulomb}$.}
\resposta{(a) $v_d \approx \SI{0.74}{\milli\meter\per\second}$;
(b) $\approx\SI{1360}{\second}$ ($\approx\SI{23}{\minute}$); (c) ver comentário}
\resolucao{\textbf{(a)} A corrente relaciona-se com a velocidade de deriva por
$I = n\,e\,A\,v_d$. Isolando:
\[
v_d=\frac{I}{n\,e\,A}
=\frac{10}{(\pot{8.5}{28})(\pot{1.6}{-19})(\pot{1}{-6})}
\approx\pot{7.4}{-4}\,\si{\meter\per\second}=\SI{0.74}{\milli\meter\per\second}.
\]
\textbf{(b)} $t = \dfrac{1}{\pot{7.4}{-4}} \approx \SI{1360}{\second}
\approx \SI{23}{\minute}$ --- um elétron leva aproximadamente 23 minutos para percorrer um metro.

\textbf{(c) A resolução do aparente paradoxo.} A lâmpada acende
instantaneamente porque \textbf{não é preciso que um elétron viaje da chave até a
lâmpada}. O que se propaga rapidamente é o \emph{campo elétrico} ao longo do
condutor, a uma velocidade próxima à da luz ($\sim\SI{2e8}{\meter\per\second}$
em cabos reais). Esse campo põe em movimento, quase simultaneamente,
\emph{todos} os elétrons livres ao longo de todo o fio --- inclusive os que já
estavam dentro do filamento da lâmpada.\\
\textbf{Analogia:} é como um tubo já cheio de água. Ao empurrar na entrada, sai
água na saída imediatamente, embora cada molécula individual se desloque
lentamente. A informação (pressão) viaja rápido; a matéria, devagar.\\
\textbf{Ordens de grandeza envolvidas:} deriva $\sim\SI{e-4}{\meter\per\second}$;
agitação térmica dos elétrons $\sim\SI{e6}{\meter\per\second}$; sinal elétrico
$\sim\SI{e8}{\meter\per\second}$ --- doze ordens de grandeza separando a mais
lenta da mais rápida.}
\end{questao}

% BEGIN BANCO COMPLEMENTAR Corrente elétrica e carga
% =====================================================================
%  Banco complementar --- Corrente Eletrica e Carga  (REVISADO)
%  Substitui a versao gerada por template (8 clones em N2, 10 em N3 e
%  8 questoes conceituais marcadas como N4).
%  O cap02 ja cobre: Q=It em tres questoes N1, condutor com 32 mA,
%  DOIS graficos i x t, carga em 10 h, rampa linear de 6 A a 2 A e a
%  velocidade de deriva em fio de cobre de 1 mm2 com 10 A.
%  Este banco explora: capacidade em A.h e autonomia sob corrente
%  variavel, formas de onda (triangular, trapezoidal, pulsada),
%  carga liquida x carga total em correntes bipolares, densidade de
%  corrente e limite termico, eletrolise, corrente de fuga, o erro de
%  estimar carga pelo pico e a distincao entre velocidade de deriva e
%  velocidade de propagacao do sinal.
%  Distribuicao mantida: 5 N1 + 8 N2 + 10 N3 + 8 N4 = 31
% =====================================================================

\subsubsection*{Banco complementar --- Corrente elétrica e carga}

\begin{atencao}[Organização do banco]
A relação $Q=It$ vale apenas para corrente \textbf{constante}. No caso geral,
$Q=\int i\,\mathrm{d}t$ --- ou seja, a carga é a \textbf{área} sob a curva
$i\times t$, com sinal. O N3 explora formas de onda, eletrólise e limites
térmicos; o N4 trata de autonomia de baterias, erros de estimativa por pico,
correntes bipolares e a diferença entre velocidade de deriva e velocidade do
sinal.
\end{atencao}

\blocofixacao

\begin{questao}[1]{}
Uma corrente constante de $\SI{1.5}{\ampere}$ circula durante
$\SI{40}{\second}$. Determine a carga transportada.
\resposta{$Q=\SI{60}{\coulomb}$}
\resolucao{
\[
Q=It=1{,}5\times40=\SI{60}{\coulomb}.
\]
\textbf{Fundamento:} a corrente é a taxa de passagem de carga,
$i=\mathrm{d}q/\mathrm{d}t$. Sendo constante, a integral vira produto simples.
Um ampère é, por definição, um coulomb por segundo.}
\end{questao}

\begin{questao}[1]{}
Uma carga de $\SI{90}{\coulomb}$ atravessa a seção de um condutor em
$\SI{45}{\second}$. Determine a corrente média.
\resposta{$I=\SI{2}{\ampere}$}
\resolucao{
\[
I=\frac{Q}{t}=\frac{90}{45}=\SI{2}{\ampere}.
\]
\textbf{Atenção ao termo "média":} o resultado é a corrente \emph{média} no
intervalo. Ela pode ter variado --- inclusive muito --- sem que isso altere a
carga total. Duas formas de onda completamente diferentes podem ter a mesma
média.}
\end{questao}

\begin{questao}[1]{}
Determine o tempo necessário para transportar $\SI{60}{\coulomb}$ com corrente
de $\SI{250}{\milli\ampere}$.
\resposta{$t=\SI{240}{\second}=\SI{4}{\minute}$}
\resolucao{
\[
t=\frac{Q}{I}=\frac{60}{0{,}250}=\SI{240}{\second}.
\]
\textbf{Erro comum:} usar $I=\SI{250}{\ampere}$ e obter
$\SI{0.24}{\second}$. Converta o prefixo \emph{antes} de dividir --- o
resultado errado é mil vezes menor e ainda assim parece plausível.}
\end{questao}

\begin{questao}[1]{}
Determine o número de elétrons que atravessa uma seção quando circula
$\SI{16}{\coulomb}$.
\dados{$e=\pot{1.6}{-19}\,\si{\coulomb}$}
\resposta{$n=\pot{1}{20}$ elétrons}
\resolucao{
\[
n=\frac{Q}{e}=\frac{16}{\pot{1.6}{-19}}=\pot{1}{20}.
\]
Cem quintilhões de elétrons --- e essa carga passa em apenas
$\SI{8}{\second}$ com corrente de $\SI{2}{\ampere}$.\\
\textbf{Ordem de grandeza:} o coulomb é uma unidade \emph{enorme} em termos de
partículas. É por isso que a quantização da carga é imperceptível em circuitos:
um erro de um elétron em $10^{20}$ não se mede.}
\end{questao}

\begin{questao}[1]{}
Sobre o sentido convencional da corrente, é correto afirmar:
\begin{alternativas}[2]
\item coincide com o movimento dos elétrons em um metal
\item é oposto ao movimento dos elétrons em um metal
\item existe apenas em semicondutores
\item depende da intensidade da corrente
\end{alternativas}
\resposta{(b)}
\resolucao{O sentido \textbf{convencional} é o do movimento de cargas
\emph{positivas} --- escolha feita antes de se conhecer o elétron. Em metais,
os portadores são elétrons (negativos), que se movem no sentido
\textbf{oposto}.\\
\textbf{Por que a convenção não atrapalha:} uma carga $-q$ movendo-se para a
esquerda é eletricamente equivalente a $+q$ movendo-se para a direita. Todas as
leis de circuitos funcionam identicamente.\\
\textbf{Onde a distinção importa:} em efeitos que dependem do \emph{portador
real}, como o efeito Hall --- cujo sinal revela se os portadores são elétrons ou
lacunas, e é assim que se distingue um semicondutor tipo N de um tipo P.}
\end{questao}

\blocoaplicacao

\begin{questao}[2]{}
Por um condutor passam $\SI{240}{\coulomb}$ em $\SI{120}{\second}$.
\begin{itensq}
\item Determine a corrente média.
\item Determine o número de elétrons.
\end{itensq}
\resposta{$I=\SI{2}{\ampere}$; $n=\pot{1.5}{21}$}
\resolucao{
\[
I=\frac{240}{120}=\SI{2}{\ampere};
\qquad
n=\frac{240}{\pot{1.6}{-19}}=\pot{1.5}{21}\ \text{elétrons}.
\]
\textbf{Taxa por segundo:} $\pot{1.25}{19}$ elétrons atravessam a seção a cada
segundo. Se fosse possível contá-los um por vez, a $\SI{1}{\per\second}$,
levaria cerca de $400$ bilhões de anos --- trinta vezes a idade do universo.}
\end{questao}

\begin{questao}[2]{}
Uma bateria de celular tem capacidade de $\SI{2000}{\milli\ampere\hour}$.
\begin{itensq}
\item Converta para coulombs.
\item Determine por quanto tempo ela alimenta uma carga de
      $\SI{250}{\milli\ampere}$.
\end{itensq}
\resposta{$Q=\SI{7200}{\coulomb}$; $t=\SI{8}{\hour}$}
\resolucao{(a)
\[
Q=2000\,\si{\milli\ampere\hour}=2\,\si{\ampere}\times3600\,\si{\second}
=\SI{7200}{\coulomb}.
\]
(b)
\[
t=\frac{2000\,\si{\milli\ampere\hour}}{250\,\si{\milli\ampere}}=\SI{8}{\hour}.
\]
\textbf{Por que se usa $\si{\ampere\hour}$:} a conta acima sai em uma linha,
enquanto com coulombs seria preciso converter duas vezes. A unidade é prática,
mas é preciso lembrar que ela mede \emph{carga}, não energia.\\
\textbf{Ressalva:} a capacidade nominal pressupõe uma corrente de descarga
específica. Descargas muito rápidas entregam menos carga total --- efeito
Peukert.}
\end{questao}

\begin{questao}[2]{}
Uma corrente pulsada vale $\SI{5}{\ampere}$ durante
$\SI{20}{\milli\second}$ e zero durante os $\SI{80}{\milli\second}$
seguintes, repetindo-se.
\begin{itensq}
\item Determine a carga por ciclo.
\item Determine a corrente média.
\end{itensq}
\resposta{$Q=\SI{100}{\milli\coulomb}$; $I_{méd}=\SI{1}{\ampere}$}
\resolucao{(a) $Q=5\times0{,}020=\SI{0.1}{\coulomb}$.\\
(b) O período é $\SI{100}{\milli\second}$:
\[
I_{méd}=\frac{Q}{T}=\frac{0{,}1}{0{,}1}=\SI{1}{\ampere},
\]
ou, equivalentemente, $I_{pico}\times D=5(0{,}20)$.\\
\textbf{Advertência que vale repetir:} a corrente média de
$\SI{1}{\ampere}$ descreve corretamente a \emph{carga} transportada, mas
\emph{não} o aquecimento. Para efeito Joule, o que vale é o valor eficaz,
$I_{rms}=5\sqrt{0{,}2}=\SI{2.24}{\ampere}$ --- que dissipa cinco vezes mais que
$\SI{1}{\ampere}$ contínuo.}
\end{questao}

\begin{questao}[2]{}
Um condutor de seção $\SI{1.5}{\milli\meter\squared}$ conduz
$\SI{15}{\ampere}$. Determine a densidade de corrente.
\resposta{$J=\SI{10}{\ampere\per\milli\meter\squared}=\pot{1}{7}\,
\si{\ampere\per\meter\squared}$}
\resolucao{
\[
J=\frac{I}{A}=\frac{15}{1{,}5}=\SI{10}{\ampere\per\milli\meter\squared}
=\frac{15}{\pot{1.5}{-6}}=\pot{1}{7}\,\si{\ampere\per\meter\squared}.
\]
\textbf{Por que $J$ importa mais que $I$:} o aquecimento e a durabilidade de um
condutor dependem da \emph{concentração} de corrente, não do total. Valores
usuais em instalações prediais ficam entre $4$ e
$\SI{6}{\ampere\per\milli\meter\squared}$ --- este condutor está bem acima, e
aqueceria demais.\\
\textbf{Comparação:} trilhas de circuito impresso admitem
$\SI{20}{\ampere\per\milli\meter\squared}$ ou mais, porque dissipam calor
diretamente para o ar e para a placa.}
\end{questao}

\begin{questao}[2]{}
Um capacitor de $\SI{100}{\micro\farad}$, inicialmente descarregado, recebe
corrente constante de $\SI{2}{\milli\ampere}$ durante $\SI{5}{\second}$.
\begin{itensq}
\item Determine a carga acumulada.
\item Determine a tensão final.
\end{itensq}
\resposta{$Q=\SI{10}{\milli\coulomb}$; $U=\SI{100}{\volt}$}
\resolucao{(a) $Q=It=\pot{2}{-3}\times5=\SI{10}{\milli\coulomb}$.\\
(b) $U=\dfrac{Q}{C}=\dfrac{\pot{1}{-2}}{\pot{1}{-4}}=\SI{100}{\volt}$.\\
\textbf{Observação:} com corrente \emph{constante}, a tensão cresce
\textbf{linearmente} no tempo --- uma rampa de
$\SI{20}{\volt\per\second}$. É exatamente o oposto do carregamento por fonte de
tensão através de resistor, que produz crescimento exponencial.\\
\textbf{Aplicação:} é assim que se geram rampas precisas em geradores de sinal e
em conversores analógico-digitais de dupla rampa.}
\end{questao}

\begin{questao}[2]{}
Em uma solução de $\mathrm{NaCl}$, íons $\mathrm{Na^{+}}$ movem-se para a
direita transportando $\SI{3}{\coulomb}$ por segundo, enquanto
$\mathrm{Cl^{-}}$ movem-se para a esquerda transportando
$\SI{2}{\coulomb}$ por segundo em módulo. Determine a corrente resultante.
\resposta{$I=\SI{5}{\ampere}$ para a direita}
\resolucao{Cargas positivas indo para a direita e cargas negativas indo para a
esquerda produzem corrente convencional \textbf{no mesmo sentido} --- elas se
\emph{somam}:
\[
I=3+2=\SI{5}{\ampere}\ \text{para a direita}.
\]
\textbf{Erro comum:} subtrair, obtendo $\SI{1}{\ampere}$. O sinal negativo do
íon \emph{e} o sentido oposto do movimento se cancelam, resultando em
contribuição positiva.\\
\textbf{Regra:} corrente convencional $=$ (cargas positivas no sentido adotado)
$+$ (cargas negativas no sentido contrário). Em eletrólitos e em plasmas, os
dois tipos de portador contribuem simultaneamente.}
\end{questao}

\begin{questao}[2]{}
Uma bateria descarregada de $\SI{1200}{\milli\ampere\hour}$ é recarregada com
corrente constante de $\SI{400}{\milli\ampere}$.
\begin{itensq}
\item Determine o tempo teórico de recarga.
\item Explique por que o tempo real é maior.
\end{itensq}
\resposta{$t=\SI{3}{\hour}$ teórico; na prática, $20$ a $40\%$ mais}
\resolucao{(a)
\[
t=\frac{1200}{400}=\SI{3}{\hour}.
\]
(b) \textbf{Três razões para o tempo real ser maior.}
\begin{itemize}
\item \textbf{Rendimento de carga:} parte da carga injetada não fica armazenada
--- ela se perde em reações secundárias e em aquecimento. O rendimento
coulômbico típico fica entre $85$ e $95\%$.
\item \textbf{Fase de tensão constante:} carregadores modernos aplicam corrente
constante até certa tensão e depois \emph{reduzem} a corrente
progressivamente. Essa fase final é lenta e pode responder por metade do tempo
total.
\item \textbf{Limitação térmica:} se a bateria aquecer, o carregador reduz a
corrente para protegê-la.
\end{itemize}
\textbf{Estimativa prática:} multiplique o tempo teórico por $1{,}3$ a
$1{,}4$.}
\end{questao}

\begin{questao}[2]{}
Uma corrente de $\SI{8}{\ampere}$ atravessa um fusível cuja especificação é
$I^{2}t=\SI{50}{\ampere\squared\second}$. Determine o tempo até a fusão.
\resposta{$t\approx\SI{0.78}{\second}$}
\resolucao{
\[
t=\frac{I^{2}t}{I^{2}}=\frac{50}{64}=\SI{0.78}{\second}.
\]
\textbf{Por que a especificação usa $I^{2}t$:} a energia dissipada no elemento
fusível é $RI^{2}t$, e a fusão ocorre quando essa energia atinge o valor
necessário para fundir o metal. Como $R$ é fixo, o parâmetro característico é
$I^{2}t$.\\
\textbf{Consequência:} dobrar a corrente reduz o tempo de atuação por
\textbf{quatro}. Com $\SI{16}{\ampere}$, o mesmo fusível abriria em
$\SI{0.20}{\second}$; com $\SI{80}{\ampere}$, em
$\SI{7.8}{\milli\second}$.\\
\textbf{Note a diferença conceitual:} aqui a grandeza relevante é
$I^{2}t$ (energia), e não $It$ (carga). São perguntas distintas sobre a mesma
corrente.}
\end{questao}

\blocoaprofundamento

\begin{questao}[3]{}
A corrente em um circuito varia segundo $i(t)=1+2t$ (ampères, com $t$ em
segundos), no intervalo $0\le t\le\SI{2}{\second}$.
\begin{itensq}
\item Determine a carga total transferida.
\item Determine a corrente média e compare com a média dos extremos.
\end{itensq}
\resposta{$Q=\SI{6}{\coulomb}$; $I_{méd}=\SI{3}{\ampere}$}
\resolucao{(a)
\[
Q=\int_0^{2}(1+2t)\,\mathrm{d}t=\Big[t+t^{2}\Big]_0^{2}=2+4=\SI{6}{\coulomb}.
\]
(b) $I_{méd}=Q/T=6/2=\SI{3}{\ampere}$.\\
\textbf{Média dos extremos:} $i(0)=1$ e $i(2)=\SI{5}{\ampere}$, cuja média é
$\SI{3}{\ampere}$ --- \emph{coincide}.\\
\textbf{Mas cuidado:} a coincidência ocorre porque a função é \textbf{linear}.
Para $i(t)=t^{2}$ no mesmo intervalo, a média verdadeira seria
$\frac{1}{2}\int_0^{2}t^{2}\mathrm{d}t=\SI{1.33}{\ampere}$, enquanto a média
dos extremos daria $\SI{2}{\ampere}$ --- erro de $50\%$.\\
\textbf{Regra:} a média dos extremos só vale para variação linear. No caso
geral, é preciso integrar.}
\end{questao}

\begin{questao}[3]{}
Um pulso de corrente cresce linearmente de zero a $\SI{6}{\ampere}$ em
$\SI{2}{\milli\second}$ e retorna a zero em mais $\SI{3}{\milli\second}$.
\begin{itensq}
\item Determine a carga transferida.
\item Determine a corrente média no pulso.
\end{itensq}
\resposta{$Q=\SI{15}{\milli\coulomb}$; $I_{méd}=\SI{3}{\ampere}$}
\resolucao{(a) A carga é a \textbf{área} do triângulo:
\[
Q=\frac{1}{2}\,I_{pico}\,T=\frac{1}{2}(6)(\pot{5}{-3})
=\pot{1.5}{-2}\,\si{\coulomb}=\SI{15}{\milli\coulomb}.
\]
\textbf{Observe:} a repartição $2+3$ do tempo \emph{não} importa para a área ---
qualquer triângulo de base $\SI{5}{\milli\second}$ e altura
$\SI{6}{\ampere}$ tem a mesma área. A assimetria afetaria o valor eficaz, não a
carga.\\
(b) $I_{méd}=Q/T=0{,}015/0{,}005=\SI{3}{\ampere}$, exatamente metade do pico
--- característica de qualquer forma triangular.}
\end{questao}

\begin{questao}[3]{}
Uma corrente trapezoidal sobe de $0$ a $\SI{4}{\ampere}$ em
$\SI{1}{\second}$, permanece constante por $\SI{3}{\second}$ e cai a zero em
$\SI{2}{\second}$.
\begin{itensq}
\item Determine a carga total.
\item Determine a corrente média.
\end{itensq}
\resposta{$Q=\SI{18}{\coulomb}$; $I_{méd}=\SI{3}{\ampere}$}
\resolucao{(a) Some as três áreas:
\[
Q=\underbrace{\tfrac12(1)(4)}_{2}+\underbrace{(3)(4)}_{12}
+\underbrace{\tfrac12(2)(4)}_{4}=\SI{18}{\coulomb}.
\]
(b) O intervalo total é $\SI{6}{\second}$:
$I_{méd}=18/6=\SI{3}{\ampere}$.\\
\textbf{Atalho para trapézios:} a área vale
$I_{pico}\times\dfrac{B+b}{2}$, com $B=\SI{6}{\second}$ (base maior) e
$b=\SI{3}{\second}$ (base menor):
$4\times4{,}5=\SI{18}{\coulomb}$ \checkmark.}
\end{questao}

\begin{questao}[3]{}
Uma corrente alterna entre $+\SI{4}{\ampere}$ por $\SI{3}{\second}$ e
$-\SI{6}{\ampere}$ por $\SI{2}{\second}$, repetindo-se.
\begin{itensq}
\item Determine a carga líquida por ciclo.
\item Determine a carga total \emph{transportada} em módulo.
\item Interprete a diferença.
\end{itensq}
\resposta{Líquida: $\SI{0}{\coulomb}$; total: $\SI{24}{\coulomb}$}
\resolucao{(a)
\[
Q_{líq}=4(3)+(-6)(2)=12-12=\SI{0}{\coulomb}.
\]
(b) Em módulo, passaram $12+12=\SI{24}{\coulomb}$.\\
(c) \textbf{As duas respostas são corretas para perguntas diferentes.}
\begin{itemize}
\item A \textbf{carga líquida} nula significa que nada se acumula: um capacitor
nesse circuito voltaria à tensão inicial a cada ciclo, e uma bateria não se
carregaria nem descarregaria.
\item A \textbf{carga transportada} de $\SI{24}{\coulomb}$ é o que
efetivamente atravessou a seção. Em eletrólise, ela é que produziria depósito
--- se o processo fosse irreversível.
\item O \textbf{aquecimento} depende de outra grandeza ainda:
$I_{rms}=\sqrt{\frac{4^{2}(3)+6^{2}(2)}{5}}=\SI{4.9}{\ampere}$.
\end{itemize}
\textbf{Moral:} "quanta corrente passou" é uma pergunta ambígua. Especifique se
o interesse é acúmulo, transporte ou dissipação.}
\end{questao}

\begin{questao}[3]{}
Uma corrente de descarga vale $i(t)=I_0e^{-t/\tau}$, com
$I_0=\SI{2}{\ampere}$ e $\tau=\SI{50}{\milli\second}$.
\begin{itensq}
\item Determine a carga total transferida.
\item Determine a fração transferida no primeiro $\tau$.
\end{itensq}
\resposta{$Q=\SI{0.1}{\coulomb}$; $63{,}2\%$}
\resolucao{(a)
\[
Q=\int_0^{\infty}I_0e^{-t/\tau}\,\mathrm{d}t
=I_0\tau=2(0{,}050)=\SI{0.1}{\coulomb}.
\]
\textbf{Interpretação geométrica:} a carga equivale à corrente inicial mantida
por exatamente $\tau$ --- o retângulo $I_0\times\tau$ tem a mesma área da
exponencial completa.\\
(b)
\[
\frac{Q(\tau)}{Q}=1-e^{-1}=63{,}2\%.
\]
\textbf{Compare com a energia:} na mesma descarga, a energia decai com
constante $\tau/2$ (pois $p\propto i^{2}$), e $86{,}5\%$ dela é dissipada em um
$\tau$. \textbf{Carga e energia têm ritmos diferentes} --- a energia se esgota
mais depressa.}
\end{questao}

\begin{questao}[3]{}
Uma corrente de $\SI{2}{\ampere}$ atravessa uma célula eletrolítica de cobre
durante $\SI{30}{\minute}$.
\begin{itensq}
\item Determine a carga transferida.
\item Determine a massa de cobre depositada.
\end{itensq}
\dados{$M_{Cu}=\SI{63.5}{\gram\per\mole}$; valência $z=2$;
$F=\SI{96500}{\coulomb\per\mole}$}
\resposta{$Q=\SI{3600}{\coulomb}$; $m=\SI{1.18}{\gram}$}
\resolucao{(a) $Q=2\times1800=\SI{3600}{\coulomb}$.\\
(b) Pela lei de Faraday da eletrólise:
\[
m=\frac{QM}{zF}=\frac{3600(63{,}5)}{2(96500)}
=\frac{228600}{193000}=\SI{1.18}{\gram}.
\]
\textbf{Significado de $z=2$:} cada íon $\mathrm{Cu^{2+}}$ precisa de
\emph{dois} elétrons para se neutralizar. São necessários, portanto,
$2F$ coulombs por mol depositado.\\
\textbf{Conexão conceitual:} esta é a medida mais direta de que a corrente
transporta carga em quantidades proporcionais à matéria depositada --- foi
justamente por experimentos assim que Faraday inferiu a existência de uma
unidade elementar de carga, décadas antes da descoberta do elétron.}
\end{questao}

\begin{questao}[3]{}
Um condutor de $\SI{2.5}{\milli\meter\squared}$ deve conduzir
$\SI{20}{\ampere}$.
\begin{itensq}
\item Determine a densidade de corrente.
\item Compare com o limite usual de $\SI{5}{\ampere\per\milli\meter\squared}$ e
      determine a seção adequada.
\end{itensq}
\resposta{$J=\SI{8}{\ampere\per\milli\meter\squared}$ --- excessiva; adotar
$\SI{4}{\milli\meter\squared}$ ou mais}
\resolucao{(a) $J=20/2{,}5=\SI{8}{\ampere\per\milli\meter\squared}$.\\
(b) O valor está $60\%$ acima do limite usual. A seção mínima seria
\[
A=\frac{20}{5}=\SI{4}{\milli\meter\squared},
\]
e a bitola comercial adequada é $\SI{4}{\milli\meter\squared}$ (ou
$\SI{6}{\milli\meter\squared}$, com folga).\\
\textbf{O que acontece se mantivermos $\SI{2.5}{\milli\meter\squared}$:} a
potência dissipada por metro é $\rho J^{2}A=\rho I^{2}/A$, que com a seção
menor é $1{,}6$ vezes maior. A temperatura do condutor sobe, a isolação
envelhece mais depressa e a resistência aumenta --- realimentando o
aquecimento.\\
\textbf{Ressalva:} o limite de $\SI{5}{\ampere\per\milli\meter\squared}$ é
apenas uma regra de bolso. O valor normativo depende do tipo de isolação, do
método de instalação (eletroduto, bandeja, ao ar livre), do agrupamento de
circuitos e da temperatura ambiente.}
\end{questao}

\begin{questao}[3]{}
Um capacitor de $\SI{470}{\micro\farad}$ carregado a $\SI{100}{\volt}$
apresenta corrente de fuga aproximadamente constante de
$\SI{20}{\micro\ampere}$.
\begin{itensq}
\item Determine a carga armazenada.
\item Estime o tempo para descarga completa.
\item Comente a validade da hipótese de corrente constante.
\end{itensq}
\resposta{$Q=\SI{47}{\milli\coulomb}$; $t\approx\SI{39}{\minute}$}
\resolucao{(a) $Q=CU=\pot{470}{-6}(100)=\SI{47}{\milli\coulomb}$.\\
(b) Com $I$ constante:
\[
t=\frac{Q}{I}=\frac{\pot{4.7}{-2}}{\pot{2}{-5}}=\SI{2350}{\second}
\approx\SI{39}{\minute}.
\]
(c) \textbf{A hipótese é otimista.} A corrente de fuga é aproximadamente
proporcional à tensão ($I=U/R_{fuga}$), de modo que ela \emph{diminui} durante
a descarga --- e o processo é exponencial, não linear.\\
Com $R_{fuga}=100/\pot{2}{-5}=\SI{5}{\mega\ohm}$, tem-se
$\tau=R_{fuga}C=\SI{2350}{\second}$ --- o \emph{mesmo} número, mas com
significado diferente: após $\SI{39}{\minute}$ resta $37\%$ da tensão, não
zero. A descarga prática ($<1\%$) leva $5\tau\approx\SI{3.3}{\hour}$.\\
\textbf{Implicação de segurança:} um banco de capacitores pode reter tensão
letal por horas. Resistores de descarga permanentes não são opcionais.}
\end{questao}

\begin{questao}[3]{}
Compare a carga transportada e o aquecimento produzido por duas correntes de
mesma média: (A) $\SI{2}{\ampere}$ constantes; (B) $\SI{10}{\ampere}$ durante
$20\%$ do tempo e zero no restante.
\begin{itensq}
\item Determine a carga em $\SI{10}{\second}$ para cada uma.
\item Determine o valor eficaz de cada uma.
\item Compare o aquecimento em um resistor de $\SI{1}{\ohm}$.
\end{itensq}
\resposta{Mesma carga ($\SI{20}{\coulomb}$); $I_{rms}=2$ e
$\SI{4.47}{\ampere}$; aquecimento $5$ vezes maior em (B)}
\resolucao{(a) Ambas transportam $Q=I_{méd}t=2(10)=\SI{20}{\coulomb}$ ---
por construção.\\
(b) \textbf{(A)} $I_{rms}=\SI{2}{\ampere}$ (constante).\\
\textbf{(B)} $I_{rms}=\sqrt{0{,}2(10)^{2}}=\sqrt{20}=\SI{4.47}{\ampere}$.\\
(c)
\[
P_A=1(2)^{2}=\SI{4}{\watt};
\qquad
P_B=1(4{,}47)^{2}=\SI{20}{\watt}.
\]
\textbf{Cinco vezes mais}, exatamente o inverso do ciclo de trabalho.\\
\textbf{Síntese fundamental:} \emph{mesma carga, aquecimento completamente
diferente}. A carga depende da média; o aquecimento, do valor eficaz. Duas
grandezas independentes de uma mesma forma de onda.\\
\textbf{Consequência prática:} circuitos pulsados exigem condutores
dimensionados pelo valor \emph{eficaz}, ainda que a carga transportada seja
modesta.}
\end{questao}

\begin{questao}[3]{}
Um circuito recebe pulsos de $\SI{50}{\ampere}$ com duração de
$\SI{100}{\micro\second}$, a uma taxa de $\SI{200}{\per\second}$.
\begin{itensq}
\item Determine a carga por pulso e a corrente média.
\item Determine o ciclo de trabalho.
\item Determine o valor eficaz.
\end{itensq}
\resposta{$\SI{5}{\milli\coulomb}$ por pulso; $I_{méd}=\SI{1}{\ampere}$;
$D=2\%$; $I_{rms}=\SI{7.07}{\ampere}$}
\resolucao{(a) $Q=50\times\pot{100}{-6}=\SI{5}{\milli\coulomb}$.
Com $200$ pulsos por segundo,
\[
I_{méd}=200\times\pot{5}{-3}=\SI{1}{\ampere}.
\]
(b) $D=\dfrac{\pot{100}{-6}}{1/200}=\dfrac{\pot{100}{-6}}{\pot{5}{-3}}
=0{,}02=2\%$.\\
(c) $I_{rms}=I_{pico}\sqrt{D}=50\sqrt{0{,}02}=\SI{7.07}{\ampere}$.\\
\textbf{Três números, três finalidades.}
\begin{itemize}
\item $I_{méd}=\SI{1}{\ampere}$ dimensiona a \textbf{fonte} e a autonomia da
bateria.
\item $I_{rms}=\SI{7.07}{\ampere}$ dimensiona o \textbf{aquecimento} de
condutores e componentes.
\item $I_{pico}=\SI{50}{\ampere}$ dimensiona a \textbf{corrente máxima}
suportável pelos semicondutores e a queda instantânea na fonte.
\end{itemize}
Usar o número errado leva a projetos ou superdimensionados ou destruídos.}
\end{questao}

\blocodesafio

\begin{questao}[4]{}
\textbf{(Autonomia sob corrente variável.)} Uma bateria de
$\SI{2.4}{\ampere\hour}$ alimenta uma carga cuja corrente alterna entre
$\SI{0.4}{\ampere}$ por $\SI{20}{\minute}$ e $\SI{1.2}{\ampere}$ por
$\SI{5}{\minute}$, repetidamente.
\begin{itensq}
\item Determine a carga consumida por ciclo e a corrente média.
\item Determine a autonomia.
\item Discuta por que a autonomia real será menor.
\end{itensq}
\resposta{$\SI{0.233}{\ampere\hour}$ por ciclo; $I_{méd}=\SI{0.56}{\ampere}$;
autonomia $\approx\SI{4.3}{\hour}$}
\resolucao{(a) Convertendo os tempos para horas:
\[
Q_{ciclo}=0{,}4\left(\tfrac{20}{60}\right)+1{,}2\left(\tfrac{5}{60}\right)
=0{,}1333+0{,}1000=\SI{0.2333}{\ampere\hour}.
\]
O ciclo dura $\SI{25}{\minute}$, logo
\[
I_{méd}=\frac{0{,}2333}{25/60}=\SI{0.56}{\ampere}.
\]
(b)
\[
n=\frac{2{,}4}{0{,}2333}=10{,}3\ \text{ciclos}
\;\Longrightarrow\;
t=10{,}3(25)=\SI{257}{\minute}=\SI{4.3}{\hour}.
\]
\emph{(Conferindo: $2{,}4/0{,}56=\SI{4.3}{\hour}$ \checkmark.)}\\
\textbf{Observe:} o patamar de $\SI{1.2}{\ampere}$ dura apenas um quinto do
tempo do outro, mas consome $43\%$ da carga.\\
(c) \textbf{Por que a autonomia real é menor.}
\begin{itemize}
\item \textbf{Efeito Peukert:} a capacidade útil diminui com correntes mais
altas. Os picos de $\SI{1.2}{\ampere}$ extraem menos carga útil do que a
proporção sugere.
\item \textbf{Tensão de corte:} o equipamento para de funcionar antes de a
bateria se esgotar, ao atingir sua tensão mínima --- e isso ocorre \emph{mais
cedo} durante um pico de corrente, por causa da queda em $r$.
\item \textbf{Resistência interna crescente} ao longo da descarga, o que agrava
o item anterior.
\item \textbf{Temperatura:} capacidade menor no frio.
\end{itemize}
\textbf{Estimativa realista:} $\SI{4.3}{\hour}$ é um teto teórico; espere de
$3{,}5$ a $\SI{4}{\hour}$.}
\end{questao}

\begin{questao}[4]{}
\textbf{(Erro de estimativa pelo pico.)} Um técnico estima a carga de um pulso
triangular ($0\to\SI{6}{\ampere}\to0$ em $\SI{5}{\milli\second}$) como
$Q\approx I_{pico}T$.
\begin{itensq}
\item Determine o erro cometido.
\item Generalize para formas de onda quaisquer.
\item Proponha uma regra prática.
\end{itensq}
\resposta{Erro de $+100\%$ --- o dobro do valor real}
\resolucao{(a) \textbf{Valor real:}
$Q=\frac12(6)(0{,}005)=\SI{15}{\milli\coulomb}$.\\
\textbf{Estimativa:} $6(0{,}005)=\SI{30}{\milli\coulomb}$.\\
\[
\text{erro}=\frac{30-15}{15}=+100\%.
\]
(b) \textbf{Generalização.} Definindo o \emph{fator de forma}
$k=I_{méd}/I_{pico}$, a estimativa por pico erra sempre por um fator $1/k$:
\begin{center}
\begin{tabular}{lcc}
\hline
Forma de onda & $k$ & erro de $I_{pico}T$\\
\hline
Retangular (contínua) & $1{,}00$ & $0\%$\\
Trapezoidal (exemplo anterior) & $0{,}75$ & $+33\%$\\
Triangular & $0{,}50$ & $+100\%$\\
Exponencial ($T=5\tau$) & $0{,}20$ & $+400\%$\\
Pulsada ($D=2\%$) & $0{,}02$ & $+4900\%$\\
\hline
\end{tabular}
\end{center}
\textbf{O erro é sempre para mais} --- a estimativa por pico é
\emph{conservadora}.\\
(c) \textbf{Regra prática.}
\begin{itemize}
\item Se a intenção é \textbf{garantir margem} (dimensionar uma bateria, por
exemplo), a estimativa por pico é segura, mas pode superdimensionar
absurdamente.
\item Se a intenção é \textbf{prever} a carga, integre --- ou ao menos use o
fator de forma da onda em questão.
\item \textbf{Nunca} use $I_{pico}T$ para estimar autonomia: o
superdimensionamento chega a ordens de grandeza em sinais pulsados.
\end{itemize}}
\end{questao}

\begin{questao}[4]{}
\textbf{(Projeto de forma de onda.)} Projete uma corrente em \textbf{dois
patamares} que transfira exatamente $\SI{10}{\coulomb}$ em
$\SI{4}{\second}$, com valores comerciais simples.
\begin{itensq}
\item Estabeleça a condição geral.
\item Proponha duas soluções e compare o valor eficaz.
\item Indique qual é preferível e por quê.
\end{itensq}
\resposta{Condição: $I_1t_1+I_2t_2=10$ com $t_1+t_2=4$; a solução mais
uniforme minimiza o aquecimento}
\resolucao{(a) \textbf{Condição:}
\[
I_1t_1+I_2t_2=\SI{10}{\coulomb};
\qquad
t_1+t_2=\SI{4}{\second}.
\]
Duas equações e quatro incógnitas --- o problema tem \textbf{dois graus de
liberdade}, e admite infinitas soluções. É preciso um critério adicional.\\
(b) \textbf{Solução A:} $\SI{4}{\ampere}$ por $\SI{1}{\second}$ e
$\SI{2}{\ampere}$ por $\SI{3}{\second}$.
\[
Q=4+6=\SI{10}{\coulomb}\ \checkmark;
\qquad
I_{rms}=\sqrt{\frac{16(1)+4(3)}{4}}=\sqrt{7}=\SI{2.65}{\ampere}.
\]
\textbf{Solução B:} $\SI{8}{\ampere}$ por $\SI{1}{\second}$ e
$\SI{0.667}{\ampere}$ por $\SI{3}{\second}$.
\[
Q=8+2=\SI{10}{\coulomb}\ \checkmark;
\qquad
I_{rms}=\sqrt{\frac{64(1)+0{,}444(3)}{4}}=\sqrt{16{,}3}=\SI{4.04}{\ampere}.
\]
(c) \textbf{A solução A é preferível.} Ela tem valor eficaz $34\%$ menor e,
portanto, aquecimento $2{,}3$ vezes menor --- entregando exatamente a mesma
carga.\\
\textbf{Princípio geral:} para uma dada carga em um dado tempo, o valor eficaz é
\textbf{mínimo} quando a corrente é \emph{constante}
($I=10/4=\SI{2.5}{\ampere}$, $I_{rms}=\SI{2.5}{\ampere}$). Qualquer
variação aumenta $I_{rms}$ acima da média --- é a mesma desigualdade
$\langle i^{2}\rangle\ge\langle i\rangle^{2}$ vista antes.\\
\textbf{Conclusão de projeto:} se a aplicação permitir, use corrente constante.
Se exigir patamares, mantenha-os o mais próximos possível entre si.}
\end{questao}

\begin{questao}[4]{}
\textbf{(Corrente bipolar.)} Uma corrente periódica é positiva em parte do
ciclo e negativa no restante.
\begin{itensq}
\item Formule a condição para carga líquida nula.
\item Determine se isso implica ausência de efeitos.
\item Cite duas aplicações em que essa condição é imposta deliberadamente.
\end{itensq}
\resposta{$\int_0^{T}i\,\mathrm{d}t=0$; os efeitos térmicos e mecânicos
permanecem}
\resolucao{(a) \textbf{Condição:}
\[
\int_0^{T}i(t)\,\mathrm{d}t=0,
\]
ou seja, as áreas positiva e negativa devem ser iguais em módulo. Em termos de
patamares, $I_{+}t_{+}=|I_{-}|t_{-}$.\\
(b) \textbf{Não implica ausência de efeitos.} O que se anula é apenas o
\emph{acúmulo} de carga. Permanecem:
\begin{itemize}
\item \textbf{Aquecimento:} $P=RI_{rms}^{2}$, com
$I_{rms}=\sqrt{\frac1T\int i^{2}\mathrm{d}t}>0$ sempre.
\item \textbf{Forças magnéticas} sobre condutores, que dependem de $i$
instantâneo.
\item \textbf{Desgaste eletroquímico}, se as reações direta e inversa não forem
perfeitamente simétricas.
\end{itemize}
(c) \textbf{Duas aplicações da condição imposta deliberadamente.}
\begin{itemize}
\item \textbf{Estimulação elétrica médica} (marca-passos, eletroestimuladores):
pulsos bifásicos com carga líquida rigorosamente nula evitam acúmulo de
subprodutos eletroquímicos no tecido e corrosão dos eletrodos. É requisito
normativo, não preferência.
\item \textbf{Acionamento de transformadores:} carga líquida não nula
introduziria componente contínua, que satura o núcleo magnético e provoca
corrente de magnetização excessiva. Conversores em ponte são projetados para
garantir simetria.
\end{itemize}
\textbf{E há um terceiro caso instrutivo:} em galvanoplastia usa-se
\emph{deliberadamente} carga líquida não nula --- ali o acúmulo é o objetivo.}
\end{questao}

\begin{questao}[4]{}
\textbf{(Medição --- média contra instantâneo.)} Dois instrumentos medem a
mesma corrente pulsada: um indica a média temporal, outro o valor instantâneo.
\begin{itensq}
\item Explique por que os resultados diferem.
\item Determine qual usar para estimar carga, aquecimento e pico.
\item Projete uma medição que forneça os três.
\end{itensq}
\resposta{Média para carga, eficaz para aquecimento, instantâneo para pico}
\resolucao{(a) \textbf{Os instrumentos medem grandezas diferentes de uma mesma
forma de onda.} Um amperímetro de bobina móvel responde à média (por inércia
mecânica); um osciloscópio mostra o valor instantâneo; um multímetro "true RMS"
calcula o valor eficaz.\\
Para a corrente pulsada de $\SI{50}{\ampere}$ com $D=2\%$, os três leriam
$\SI{1}{\ampere}$, $\SI{50}{\ampere}$ e $\SI{7.07}{\ampere}$ --- todos
corretos, todos respondendo perguntas distintas.\\
(b)
\begin{center}
\begin{tabular}{ll}
\hline
Objetivo & Grandeza\\
\hline
Carga transferida, autonomia de bateria & \textbf{média}\\
Aquecimento, dimensionamento de condutor & \textbf{eficaz}\\
Corrente máxima em semicondutor, saturação & \textbf{pico}\\
\hline
\end{tabular}
\end{center}
(c) \textbf{Projeto de medição completa.} Use um \emph{resistor shunt} de baixo
valor e observe a tensão sobre ele com osciloscópio digital:
\begin{enumerate}
\item O \textbf{traço instantâneo} fornece o pico diretamente.
\item A função \textbf{média} do osciloscópio integra a forma de onda e entrega
$I_{méd}$ --- e portanto a carga, multiplicando pelo tempo.
\item A função \textbf{RMS} entrega $I_{rms}$.
\end{enumerate}
\textbf{Cuidados:} a largura de banda do osciloscópio e do shunt deve ser muito
maior que a taxa de subida dos pulsos, sob pena de arredondá-los e subestimar o
pico. E a indutância parasita do shunt introduz erro proporcional a
$\mathrm{d}i/\mathrm{d}t$ --- use shunts coaxiais em pulsos rápidos.}
\end{questao}

\begin{questao}[4]{}
\textbf{(Integração de sinal ruidoso.)} Uma corrente medida contém ruído
aditivo de média zero.
\begin{itensq}
\item Explique por que integrar melhora a estimativa da carga.
\item Determine como o erro decresce com o tempo de integração.
\item Indique em que situação integrar \emph{piora} o resultado.
\end{itensq}
\resposta{O erro relativo decresce como $1/\sqrt{N}$; integrar piora se houver
\emph{deriva} (ruído de média não nula)}
\resolucao{(a) \textbf{Por que integrar ajuda.} A carga é
$Q=\int i\,\mathrm{d}t$. Se o ruído tem média zero, suas contribuições
positivas e negativas tendem a se cancelar na integral, enquanto o sinal
verdadeiro se acumula. A relação sinal-ruído \emph{melhora} com o tempo.\\
(b) \textbf{Decrescimento do erro.} Para $N$ amostras independentes de ruído de
desvio $\sigma$, o desvio da \emph{soma} é $\sigma\sqrt{N}$, enquanto a soma do
sinal cresce com $N$. Logo o erro relativo cai como
\[
\frac{\sigma\sqrt{N}}{SN}=\frac{\sigma}{S\sqrt{N}}\propto\frac{1}{\sqrt{T}}.
\]
\textbf{Quantificando:} integrar por $\SI{100}{\second}$ em vez de
$\SI{1}{\second}$ reduz o erro relativo por um fator $10$. Para ganhar outro
fator $10$, seriam necessários $\SI{10000}{\second}$ --- rendimento
decrescente.\\
(c) \textbf{Quando integrar piora.} Se o ruído \textbf{não} tiver média zero ---
isto é, se houver \emph{offset} ou \emph{deriva} no instrumento --- o erro
\textbf{também se acumula}, e cresce \emph{linearmente} com $T$:
\[
\text{erro}=\varepsilon_{offset}\times T.
\]
Enquanto o ruído aleatório cai com $\sqrt{T}$, o erro sistemático cresce com
$T$. Existe, portanto, um \textbf{tempo ótimo} de integração, além do qual a
deriva domina.\\
\textbf{Prática de laboratório:} meça a leitura com \emph{entrada em curto}
antes do ensaio, para quantificar o \emph{offset}, e subtraia-o. Sem essa
correção, integrações longas são pior que inúteis --- elas produzem números
precisos e errados.}
\end{questao}

\begin{questao}[4]{}
\textbf{(Por que a LKC é tão precisa.)} A lei dos nós afirma que a carga não se
acumula em um nó. Quantifique o quanto ela poderia, em princípio, ser violada.
\begin{itensq}
\item Estime a capacitância parasita de um nó típico e relacione carga acumulada
      com potencial.
\item Determine o desequilíbrio de corrente necessário para elevar o nó em
      $\SI{1}{\volt}$.
\item Conclua sobre o caráter da lei.
\end{itensq}
\resposta{Com $C\approx\SI{1}{\pico\farad}$, basta $\SI{1}{\pico\coulomb}$ para
$\SI{1}{\volt}$ --- um desequilíbrio de $\SI{1}{\micro\ampere}$ por
$\SI{1}{\micro\second}$}
\resolucao{(a) Um nó real (uma solda, um trecho de trilha) tem capacitância
parasita para o ambiente da ordem de $C\approx\SI{1}{\pico\farad}$. Se uma
carga $\Delta Q$ se acumulasse ali, o potencial subiria:
\[
\Delta V=\frac{\Delta Q}{C}.
\]
Para $\Delta V=\SI{1}{\volt}$:
\[
\Delta Q=\pot{1}{-12}\,\si{\coulomb}
=\frac{\pot{1}{-12}}{\pot{1.6}{-19}}=\pot{6.25}{6}\ \text{elétrons}.
\]
Apenas seis milhões de elétrons --- número minúsculo diante dos
$10^{19}$ que atravessam a seção a cada segundo com $\SI{2}{\ampere}$.\\
(b) O desequilíbrio necessário é irrisório:
\[
\Delta I\,\Delta t=\SI{1}{\pico\coulomb}
\;\Longrightarrow\;
\SI{1}{\micro\ampere}\ \text{durante}\ \SI{1}{\micro\second},
\]
ou $\SI{1}{\nano\ampere}$ durante $\SI{1}{\milli\second}$. Em um circuito de
$\SI{2}{\ampere}$, isso corresponde a um desequilíbrio relativo de
$\pot{5}{-7}$ --- meio milionésimo.\\
(c) \textbf{Conclusão: a LKC não é uma aproximação, é uma consequência
autorregulada.} Qualquer tentativa de acumular carga em um nó eleva
imediatamente seu potencial, e esse potencial \emph{expulsa} o excesso --- as
correntes se reajustam em nanossegundos.\\
\textbf{Onde a lei realmente falha:} em frequências altíssimas, quando o tempo
de propagação pelo circuito deixa de ser desprezível, correntes de deslocamento
por capacitâncias parasitas passam a importar e a LKC "de nó concentrado" perde
sentido. É preciso então tratar o circuito por linhas de transmissão ou por
campos.\\
\textbf{Critério prático:} a LKC vale enquanto as dimensões do circuito forem
muito menores que o comprimento de onda do sinal --- a chamada aproximação de
\emph{parâmetros concentrados}.}
\end{questao}

\begin{questao}[4]{}
\textbf{(Deriva contra propagação.)} Um fio de cobre de
$\SI{2.5}{\milli\meter\squared}$ conduz $\SI{10}{\ampere}$, com densidade de
portadores $n=\pot{8.5}{28}\,\si{\per\meter\cubed}$.
\begin{itensq}
\item Determine a velocidade de deriva dos elétrons.
\item Determine o tempo para um elétron percorrer $\SI{1}{\meter}$.
\item Compare com o tempo de propagação do sinal e explique o paradoxo
      aparente.
\end{itensq}
\resposta{$v_d=\SI{0.29}{\milli\meter\per\second}$; $\SI{57}{\minute}$ para
$\SI{1}{\meter}$; o sinal leva $\SI{5}{\nano\second}$}
\resolucao{(a)
\[
v_d=\frac{I}{nAe}
=\frac{10}{(\pot{8.5}{28})(\pot{2.5}{-6})(\pot{1.6}{-19})}
=\frac{10}{\pot{3.4}{4}}=\pot{2.94}{-4}\,\si{\meter\per\second}.
\]
Menos de um terço de milímetro por segundo.\\
(b)
\[
t=\frac{1}{\pot{2.94}{-4}}=\SI{3400}{\second}\approx\SI{57}{\minute}.
\]
(c) \textbf{O sinal viaja a cerca de $\pot{2}{8}\,\si{\meter\per\second}$}
(dois terços de $c$, em cabo típico), percorrendo o metro em
$\SI{5}{\nano\second}$ --- \textbf{$\pot{6.8}{11}$ vezes mais rápido} que os
elétrons.\\
\textbf{Resolução do paradoxo aparente.} O que se propaga rapidamente não são os
elétrons, mas o \textbf{campo eletromagnético} que os põe em movimento. Ao
fechar o interruptor, o campo se estabelece ao longo de todo o condutor quase
instantaneamente, e \emph{todos} os elétrons começam a se mover praticamente ao
mesmo tempo --- cada um localmente, e devagar.\\
\textbf{Analogia:} um cano cheio de água. Ao abrir a torneira, a água sai
imediatamente na outra ponta, embora nenhuma molécula tenha percorrido o cano.
O que se propaga é a \emph{pressão}, não a matéria.\\
\textbf{Consequência prática:} a lâmpada acende no instante em que se aciona o
interruptor, mas os elétrons que estavam junto dele levariam quase uma hora para
chegar até ela.}
\end{questao}

% END BANCO COMPLEMENTAR

\gabarito[2.1 Corrente elétrica e carga]

% =====================================================================
\subsection{Lei de Ohm e resistividade}
% =====================================================================

\blocofixacao

\begin{questao}[1]{}
De acordo com a segunda lei de Ohm, a resistência de um fio condutor:
\begin{alternativas}
\item aumenta com a área da seção e diminui com o comprimento.
\item aumenta com o comprimento e diminui com a área da seção.
\item depende apenas do material.
\item independe da temperatura.
\item é sempre igual para todos os materiais.
\end{alternativas}
\resposta{(b)}
\resolucao{De $R = \rho\dfrac{L}{A}$: a resistência é \emph{diretamente}
proporcional ao comprimento $L$ e \emph{inversamente} proporcional à área $A$.
O material entra pela resistividade $\rho$, que por sua vez varia com a
temperatura --- o que torna (c) e (d) incorretas.}
\end{questao}

\begin{questao}[1]{}
O gráfico mostra a resistência de um fio em função de seu comprimento, mantida a
área constante.

\circuitoqq{%
\begin{tikzpicture}
 \begin{axis}[width=9cm,height=5cm,xlabel={Comprimento (m)},
   ylabel={$R$ ($\Omega$)},xmin=0,xmax=4,ymin=0,ymax=7,
   xtick={0,1,2,3},ytick={0,2,4,6},grid=major,axis lines=left]
   \addplot[Rcol,very thick,domain=0:3] {2*x};
   \addplot[only marks,mark=*,Rcol] coordinates {(1,2)(2,4)(3,6)};
 \end{axis}
\end{tikzpicture}}{Resistência proporcional ao comprimento.}

\begin{itensq}
\item Qual é a constante de proporcionalidade $k$ da reta $R = kL$?
\item Qual seria a resistência de um fio de $\SI{5}{\meter}$ do mesmo tipo?
\item Que lei essa proporcionalidade confirma?
\end{itensq}
\resposta{(a) $k=\SI{2}{\ohm\per\meter}$; (b) \SI{10}{\ohm}; (c) 2ª lei de Ohm}
\resolucao{(a) A reta passa por $(3, 6)$, logo
$k = \dfrac{R}{L} = \dfrac{6}{3} = \SI{2}{\ohm\per\meter}$.\\
(b) $R = kL = 2\cdot 5 = \SI{10}{\ohm}$.\\
(c) A proporcionalidade direta entre $R$ e $L$ (com $A$ e material fixos) confirma
a segunda lei de Ohm, $R = \rho\dfrac{L}{A}$, em que $k = \dfrac{\rho}{A}$.}
\end{questao}

\blocoaplicacao

\begin{questao}[2]{}
Um estudante quer construir um resistor de $\SI{10}{\ohm}$ com fio de
níquel-cromo, de resistividade $\rho = \SI{1.1e-6}{\ohm\meter}$ e seção
$A = \SI{0.5}{\milli\meter\squared}$. Qual é o comprimento necessário?
\resposta{$L \approx \SI{4.55}{\meter}$}
\resolucao{Isolando $L$ na segunda lei de Ohm (com
$A = \SI{0.5}{\milli\meter\squared} = \pot{0.5}{-6}\,\si{\meter\squared}$):
\[
L=\frac{R\,A}{\rho}=\frac{10\cdot\pot{0.5}{-6}}{\pot{1.1}{-6}}
=\frac{\pot{5}{-6}}{\pot{1.1}{-6}}\approx\SI{4.55}{\meter}.
\]}
\end{questao}

\begin{questao}[2]{}
Um fio condutor de comprimento $L$, diâmetro $D$ e resistência $R$ tem seu
comprimento \textbf{e} seu diâmetro \emph{duplicados}. A nova resistência será:
\begin{alternativas}[3]
\item $R$
\item $2R$
\item $R/2$
\item $4R$
\item $R/4$
\end{alternativas}
\resposta{(c)}
\resolucao{A área da seção é proporcional ao quadrado do diâmetro, $A \propto D^2$.
Ao dobrar $D$, a área quadruplica ($A' = 4A$); ao dobrar $L$, o comprimento
duplica ($L' = 2L$):
\[
R'=\rho\frac{L'}{A'}=\rho\frac{2L}{4A}=\frac{1}{2}\,\rho\frac{L}{A}=\frac{R}{2}.
\]
A armadilha (alternativa b) é considerar só o efeito do comprimento e esquecer que
a área cresce com o \emph{quadrado} do diâmetro.}
\end{questao}

\begin{questao}[2]{}
O gráfico representa a diferença de potencial $U$ entre dois pontos de um fio em
função da corrente $i$. A resistência do trecho, em ohms, vale:

\circuitoqq{%
\begin{tikzpicture}
 \begin{axis}[width=9cm,height=5cm,xlabel={$i$ (A)},ylabel={$U$ (V)},
   xmin=0,xmax=5,ymin=0,ymax=22,xtick={0,1,2,3,4},ytick={0,4,8,12,16,20},
   grid=major,axis lines=left]
   \addplot[Vcol,very thick,domain=0:4] {4*x};
   \addplot[only marks,mark=*,Vcol] coordinates {(1,4)(2,8)(3,12)(4,16)};
 \end{axis}
\end{tikzpicture}}{Curva característica $U \times i$ de um resistor ôhmico.}

\begin{alternativas}[3]
\item 0,05
\item 4
\item 20
\item 80
\item 160
\end{alternativas}
\resposta{(b)}
\resolucao{Para um condutor ôhmico, $U = R\,i$: a resistência é a \emph{inclinação}
da reta. Tomando o ponto $(4, 16)$:
\[
R=\frac{U}{i}=\frac{16}{4}=\SI{4}{\ohm}.
\]
A reta passar pela origem confirma que o resistor obedece à lei de Ohm.}
\end{questao}

\begin{questao}[2]{}
Dois fios do mesmo material têm resistências $R_1 = \SI{6}{\ohm}$ e
$R_2 = \SI{3}{\ohm}$; o fio 1 tem o dobro do comprimento do fio 2. Qual é a relação
entre as áreas das seções transversais, $A_1/A_2$?
\resposta{$A_1/A_2 = 1$ (áreas iguais)}
\resolucao{Como o material é o mesmo, $\rho$ é comum. De $R = \rho L/A$:
\[
\frac{R_1}{R_2}=\frac{L_1/A_1}{L_2/A_2}=\frac{L_1}{L_2}\cdot\frac{A_2}{A_1}.
\]
Com $\dfrac{R_1}{R_2}=\dfrac{6}{3}=2$ e $\dfrac{L_1}{L_2}=2$:
\[
2 = 2\cdot\frac{A_2}{A_1}\ \Rightarrow\ \frac{A_2}{A_1}=1
\ \Rightarrow\ A_1 = A_2.
\]
Ou seja, dobrar o comprimento já explica a resistência dobrada; as seções são
iguais.}
\end{questao}

\begin{questao}[2]{}
A tabela apresenta a resistividade de alguns materiais a $\SI{20}{\celsius}$.

\begin{table}[H]
\centering\small\sffamily
\begin{tabular}{@{}l c@{}}
\toprule
\textbf{Material} & \textbf{$\rho$ ($\times10^{-8}\,\si{\ohm\meter}$)}\\
\midrule
Prata      & 1,6\\
Cobre      & 1,7\\
Alumínio   & 2,8\\
Ferro      & 10\\
Níquel-cromo & 110\\
\bottomrule
\end{tabular}
\caption{Resistividade de materiais condutores.}
\end{table}

Dois fios de mesmo comprimento e mesma seção, um de cobre e outro de alumínio,
têm resistências $R_{\text{Cu}}$ e $R_{\text{Al}}$. A razão
$R_{\text{Al}}/R_{\text{Cu}}$ vale aproximadamente:
\begin{alternativas}[3]
\item 0,6
\item 1,0
\item 1,6
\item 2,8
\item 4,4
\end{alternativas}
\resposta{(c)}
\resolucao{Com $L$ e $A$ iguais, $R = \rho\dfrac{L}{A}$ é proporcional a $\rho$:
\[
\frac{R_{\text{Al}}}{R_{\text{Cu}}}=\frac{\rho_{\text{Al}}}{\rho_{\text{Cu}}}
=\frac{2{,}8}{1{,}7}\approx1{,}6.
\]
Por isso, para a mesma resistência, um fio de alumínio precisa ser mais grosso
que um de cobre --- o que se compensa por ele ser mais leve e barato, motivo do
seu uso em linhas de transmissão.}
\end{questao}

\begin{questao}[2]{}
Um fio de cobre tem comprimento $L = \SI{2}{\meter}$ e seção transversal
$A = \SI{1}{\milli\meter\squared}$. Calcule sua resistência.
\dados{$\rho_{\text{Cu}} = \SI{1.7e-8}{\ohm\meter}$.}
\resposta{$R = \SI{34}{\milli\ohm}$}
\resolucao{Com $A = \SI{1}{\milli\meter\squared} = \pot{1}{-6}\,\si{\meter\squared}$:
\[
R=\rho\frac{L}{A}=\pot{1.7}{-8}\cdot\frac{2}{\pot{1}{-6}}
=\pot{3.4}{-2}\,\si{\ohm}=\SI{34}{\milli\ohm}.
\]
Resistência bem baixa, como esperado de um bom condutor em curto comprimento.}
\end{questao}

\begin{questao}[2]{Apostila original revisada}
Um fio possui resistência elétrica $R$. Mantendo o mesmo material, seu comprimento é
duplicado e sua área de seção transversal é reduzida à metade. Determine a nova
resistência em função de $R$.
\resposta{$R'=4R$}
\resolucao{Pela segunda lei de Ohm, $R=\rho L/A$. Logo,
\[
R'=\rho\frac{2L}{A/2}=4\rho\frac{L}{A}=4R.
\]}
\end{questao}

\blocoaprofundamento

\begin{questao}[3]{}
Um fio de aço e um de cobre, de mesmo comprimento e mesma área, são ligados
individualmente à mesma tensão. Sabendo que
$\rho_{\text{aço}} = 10\,\rho_{\text{cobre}}$, compare:
\begin{itensq}
\item as correntes em cada fio;
\item as potências dissipadas em cada um.
\end{itensq}
\resposta{(a) $I_{\text{cobre}} = 10\,I_{\text{aço}}$;
(b) $P_{\text{cobre}} = 10\,P_{\text{aço}}$}
\resolucao{Como $L$ e $A$ são iguais, $R = \rho L/A \Rightarrow
R_{\text{aço}} = 10\,R_{\text{cobre}}$.\\
(a) Sob a mesma tensão, $I = U/R$, logo a corrente é inversamente proporcional a
$R$:
\[
\frac{I_{\text{cobre}}}{I_{\text{aço}}}
=\frac{R_{\text{aço}}}{R_{\text{cobre}}}=10
\ \Rightarrow\ I_{\text{cobre}}=10\,I_{\text{aço}}.
\]
(b) $P = U^2/R$; com $U$ igual, a potência também é inversamente proporcional a
$R$:
\[
\frac{P_{\text{cobre}}}{P_{\text{aço}}}=\frac{R_{\text{aço}}}{R_{\text{cobre}}}=10.
\]
O cobre, menos resistivo, conduz mais corrente e dissipa mais potência sob a mesma
tensão --- por isso é preferido em condutores.}
\end{questao}

\begin{questao}[3]{}
Dois fios de \emph{mesmo comprimento} e \emph{mesma seção}, um de cobre
($\rho_{\text{Cu}}=\SI{1.7e-8}{\ohm\meter}$) e outro de alumínio
($\rho_{\text{Al}}=\SI{2.8e-8}{\ohm\meter}$), são emendados em série e ligados a
$\SI{12}{\volt}$.
\begin{itensq}
\item Determine a razão entre as tensões em cada trecho.
\item Calcule a tensão em cada fio.
\item Em qual dos dois há maior dissipação de calor? Justifique.
\end{itensq}
\resposta{(a) $U_{\text{Al}}/U_{\text{Cu}}\approx1{,}65$;
(b) $\SI{4.53}{\volt}$ e $\SI{7.47}{\volt}$; (c) no alumínio}
\resolucao{(a) Com $L$ e $A$ iguais, $R \propto \rho$; e como estão em série
(mesma corrente), $U = RI \propto \rho$:
\[
\frac{U_{\text{Al}}}{U_{\text{Cu}}}=\frac{\rho_{\text{Al}}}{\rho_{\text{Cu}}}
=\frac{2{,}8}{1{,}7}\approx1{,}65.
\]
(b) Pelo divisor de tensão, a repartição segue as resistividades:
\[
U_{\text{Cu}}=12\cdot\frac{1{,}7}{4{,}5}\approx\SI{4.53}{\volt},
\qquad
U_{\text{Al}}=12\cdot\frac{2{,}8}{4{,}5}\approx\SI{7.47}{\volt}.
\]
(c) A potência é $P = UI$ com $I$ comum, logo \emph{maior tensão implica maior
potência}: o \textbf{alumínio} aquece mais. Fisicamente, sendo mais resistivo,
ele opõe mais resistência à mesma corrente.\\
\textbf{Implicação prática:} emendas entre metais diferentes são pontos críticos
em instalações --- além do aquecimento desigual, os coeficientes de dilatação
distintos afrouxam o contato com o tempo, aumentando a resistência local e
realimentando o aquecimento. É a causa clássica de incêndios em conexões
Cu--Al, razão pela qual normas exigem conectores bimetálicos específicos.}
\end{questao}

\begin{questao}[3]{}
A resistência de um condutor metálico varia com a temperatura segundo
$R = R_0\big[1+\alpha(T-T_0)\big]$, onde $\alpha$ é o coeficiente de temperatura.
Para o cobre, $\alpha \approx \SI{4e-3}{\per\celsius}$.

Um enrolamento de cobre de um motor mede $\SI{10}{\ohm}$ a $\SI{20}{\celsius}$
(motor frio). Após operar em regime, sua resistência sobe para
$\SI{12.4}{\ohm}$.
\begin{itensq}
\item Determine a temperatura do enrolamento em operação.
\item Se a tensão aplicada for mantida em $\SI{220}{\volt}$, qual é a variação
      percentual da potência dissipada entre o motor frio e o quente?
\item Explique como esse efeito é usado, na prática, para medir a temperatura
      interna de máquinas elétricas sem sensores.
\end{itensq}
\resposta{(a) $T = \SI{80}{\celsius}$; (b) a potência cai $\approx\SI{19}{\percent}$;
(c) ver comentário}
\resolucao{\textbf{(a)} Isolando a temperatura:
\[
12{,}4 = 10\big[1+\pot{4}{-3}(T-20)\big]
\;\Rightarrow\;
1{,}24 = 1+\pot{4}{-3}(T-20)
\;\Rightarrow\;
T-20 = \frac{0{,}24}{\pot{4}{-3}} = 60,
\]
logo $T = \SI{80}{\celsius}$.

\textbf{(b)} Com tensão fixa, $P = \dfrac{U^2}{R}$ é inversamente proporcional a
$R$:
\[
\frac{P_{\text{quente}}}{P_{\text{frio}}}=\frac{R_{\text{frio}}}{R_{\text{quente}}}
=\frac{10}{12{,}4}=0{,}806,
\]
ou seja, a potência dissipada \emph{cai} cerca de \SI{19.4}{\percent}.

\textbf{(c) Método da variação de resistência.} Como $\alpha$ do cobre é bem
conhecido e estável, mede-se a resistência do enrolamento a frio (com o motor
parado e em temperatura ambiente conhecida) e novamente logo após o desligamento.
Invertendo a fórmula:
\[
T = T_0 + \frac{1}{\alpha}\left(\frac{R}{R_0}-1\right),
\]
obtém-se a temperatura \emph{média} de todo o enrolamento --- justamente o que
importa para a vida útil do isolamento --- sem instalar nenhum sensor interno.
Esse é o procedimento normalizado em ensaios de elevação de temperatura de
motores e transformadores.\\
\textbf{Observação crítica:} o método fornece a média, não o \emph{ponto quente}.
Como a degradação do isolante segue uma lei exponencial com a temperatura, o ponto
mais quente pode falhar antes do que a média sugeriria --- por isso as normas
aplicam margens de segurança sobre esse valor.}
\end{questao}

\begin{questao}[3]{Apostila original revisada}
A resistividade de um material varia aproximadamente de forma linear com a
temperatura. Foram medidos os valores
$\rho(\SI{0}{\celsius})=\SI{1.5e-8}{\ohm\metre}$ e
$\rho(\SI{100}{\celsius})=\SI{2.5e-8}{\ohm\metre}$.
\begin{itensq}
\item Determine a taxa de variação da resistividade por grau Celsius.
\item Estime a resistividade a $\SI{150}{\celsius}$, mantendo o modelo linear.
\end{itensq}
\resposta{(a) $\SI{1.0e-10}{\ohm\metre\per\celsius}$; (b) $\SI{3.0e-8}{\ohm\metre}$}
\resolucao{A inclinação é
\[
\frac{\Delta\rho}{\Delta T}=\frac{(2.5-1.5)\times10^{-8}}{100}
=\SI{1.0e-10}{\ohm\metre\per\celsius}.
\]
Extrapolando até $\SI{150}{\celsius}$:
$\rho=1.5\times10^{-8}+150\times10^{-10}=\SI{3.0e-8}{\ohm\metre}$.}
\end{questao}

% BEGIN BANCO COMPLEMENTAR Lei de Ohm e resistividade
% =====================================================================
%  Banco complementar --- Lei de Ohm e Resistividade  (REVISADO)
%  Substitui a versao gerada por template (8 clones em N2, 11 em N3 e
%  10 questoes conceituais marcadas como N4).
%  A secao correspondente do cap02 e densa e ja cobre: 2a lei de Ohm
%  (MC), grafico R x L, resistor de 10 ohm com niquel-cromo, fio com
%  L e D dobrados, grafico U x i, dois fios de mesmo material, tabela
%  de resistividades, fio de cobre 2 m / 1 mm2, L dobrado com A
%  dobrada, aco x cobre, cobre x aluminio, R=R0[1+alfa(T-T0)] e
%  resistividade linear com T a partir de dados.
%  Este banco evita esses casos e explora: projeto de resistor de fio,
%  faixa total de R (tolerancia + deriva + autoaquecimento), R estatica
%  x dinamica, secao de cabo por queda de tensao, fio ESTICADO a volume
%  constante, secao variavel por integracao, sensor a corrente
%  constante, medicao a quatro fios e estabilidade termica.
%  Distribuicao mantida: 8 N1 + 8 N2 + 11 N3 + 10 N4 = 37
% =====================================================================

\subsubsection*{Banco complementar --- Lei de Ohm e resistividade}

\begin{atencao}[Organização do banco]
Duas leis distintas convivem aqui. A \textbf{primeira} ($U=RI$) é uma relação
\emph{de circuito} e vale apenas para componentes ôhmicos. A \textbf{segunda}
($R=\rho L/A$) é uma relação \emph{de geometria e material}, e vale sempre. O N3
trata de projeto de resistores, seção de cabos, sensores e medição; o N4 exige
integração, demonstrações e análise de estabilidade térmica.
\end{atencao}

\blocofixacao

\begin{questao}[1]{}
Um resistor de $\SI{470}{\ohm}$ é percorrido por $\SI{15}{\milli\ampere}$.
Determine a tensão em seus terminais.
\resposta{$U=\SI{7.05}{\volt}$}
\resolucao{
\[
U=RI=470\times\pot{15}{-3}=\SI{7.05}{\volt}.
\]
\textbf{Cuidado com o prefixo:} usar $I=\SI{15}{\ampere}$ daria
$\SI{7050}{\volt}$ --- resultado absurdo, mas que passa despercebido se a
conversão não for feita explicitamente.}
\end{questao}

\begin{questao}[1]{}
Uma tensão de $\SI{9}{\volt}$ produz corrente de
$\SI{30}{\milli\ampere}$ em um componente ôhmico. Determine sua resistência.
\resposta{$R=\SI{300}{\ohm}$}
\resolucao{
\[
R=\frac{U}{I}=\frac{9}{0{,}030}=\SI{300}{\ohm}.
\]
\textbf{Atenção à hipótese:} a razão $U/I$ sempre pode ser calculada, mas só se
chama \emph{resistência} no sentido pleno se o componente for ôhmico --- isto é,
se essa razão for a mesma para qualquer tensão aplicada. Em componentes não
lineares ela varia, e recebe o nome de \emph{resistência estática}.}
\end{questao}

\begin{questao}[1]{}
Determine a corrente em um resistor de $\SI{2.2}{\kilo\ohm}$ submetido a
$\SI{12}{\volt}$.
\resposta{$I=\SI{5.45}{\milli\ampere}$}
\resolucao{
\[
I=\frac{U}{R}=\frac{12}{2200}=\pot{5.45}{-3}\,\si{\ampere}
=\SI{5.45}{\milli\ampere}.
\]
\textbf{Atalho útil:} volts divididos por quilohms dão \emph{miliampères}
diretamente --- $12/2{,}2=5{,}45$. Essa combinação de prefixos aparece o tempo
todo em eletrônica e poupa conversões.}
\end{questao}

\begin{questao}[1]{}
Um condutor de cobre ($\rho=\pot{1.7}{-8}\,\si{\ohm\meter}$) tem
$\SI{50}{\meter}$ e seção de $\SI{2.5}{\milli\meter\squared}$. Determine sua
resistência.
\resposta{$R=\SI{0.34}{\ohm}$}
\resolucao{Convertendo a área para $\si{\meter\squared}$:
$A=\SI{2.5}{\milli\meter\squared}=\pot{2.5}{-6}\,\si{\meter\squared}$.
\[
R=\frac{\rho L}{A}=\frac{\pot{1.7}{-8}\times50}{\pot{2.5}{-6}}
=\SI{0.34}{\ohm}.
\]
\textbf{Erro clássico:} esquecer que
$\SI{1}{\milli\meter\squared}=\pot{1}{-6}\,\si{\meter\squared}$, e não
$\pot{1}{-3}$. A área é o \emph{quadrado} de um comprimento, logo o fator de
conversão também é quadrático.}
\end{questao}

\begin{questao}[1]{}
Determine a condutância de um resistor de $\SI{25}{\ohm}$.
\resposta{$G=\SI{0.04}{\ensuremath{\mathrm{S}}}$}
\resolucao{
\[
G=\frac{1}{R}=\frac{1}{25}=\SI{0.04}{\ensuremath{\mathrm{S}}}.
\]
A condutância mede a \emph{facilidade} de conduzir. Ela é a grandeza natural
quando os elementos estão em paralelo (somam-se diretamente) e na análise nodal.
Nas relações de material, define-se analogamente a \textbf{condutividade}
$\sigma=1/\rho$.}
\end{questao}

\begin{questao}[1]{}
Segundo a segunda lei de Ohm, dobrar o \textbf{comprimento} de um fio, mantida
a seção, faz a resistência:
\begin{alternativas}[2]
\item dobrar
\item cair à metade
\item quadruplicar
\item não se alterar
\end{alternativas}
\resposta{(a)}
\resolucao{Como $R=\rho L/A$, tem-se $R\propto L$: dobrar $L$ dobra $R$.\\
\textbf{Interpretação física:} o dobro do comprimento equivale a dois fios
idênticos em \emph{série} --- e resistências em série se somam. A analogia
hidráulica ajuda: um cano duas vezes mais longo oferece o dobro de oposição ao
escoamento.}
\end{questao}

\begin{questao}[1]{}
Dobrar a \textbf{área} da seção de um fio, mantido o comprimento, faz a
resistência:
\begin{alternativas}[2]
\item dobrar
\item cair à metade
\item quadruplicar
\item não se alterar
\end{alternativas}
\resposta{(b)}
\resolucao{$R\propto1/A$: dobrar $A$ divide $R$ por dois.\\
\textbf{Interpretação:} o dobro da área equivale a dois fios idênticos em
\emph{paralelo}, e o paralelo de dois iguais vale metade.\\
\textbf{Cuidado com o diâmetro:} dobrar o \emph{diâmetro} quadruplica a área
($A=\pi d^{2}/4$) e divide $R$ por \textbf{quatro} --- confundir diâmetro com
área é o deslize mais comum aqui.}
\end{questao}

\begin{questao}[1]{}
Dois fios de mesmo material e mesmo comprimento têm diâmetros $d$ e $2d$.
Determine a razão entre suas resistências.
\resposta{$R_1/R_2=4$}
\resolucao{A área é proporcional ao quadrado do diâmetro:
\[
\frac{A_2}{A_1}=\left(\frac{2d}{d}\right)^{2}=4
\;\Longrightarrow\;
\frac{R_1}{R_2}=\frac{A_2}{A_1}=4.
\]
O fio mais grosso tem um quarto da resistência. É por isso que a bitola dos
condutores é especificada em $\si{\milli\meter\squared}$ (área) e não em
milímetros de diâmetro --- a área é o que entra linearmente na conta.}
\end{questao}

\blocoaplicacao

\begin{questao}[2]{}
Determine o comprimento de um fio de níquel-cromo
($\rho=\pot{1.1}{-6}\,\si{\ohm\meter}$) de seção
$\SI{0.2}{\milli\meter\squared}$ necessário para obter
$\SI{22}{\ohm}$.
\resposta{$L=\SI{4}{\meter}$}
\resolucao{Isolando $L$:
\[
L=\frac{RA}{\rho}=\frac{22\times\pot{0.2}{-6}}{\pot{1.1}{-6}}
=\frac{\pot{4.4}{-6}}{\pot{1.1}{-6}}=\SI{4}{\meter}.
\]
\textbf{Por que níquel-cromo:} sua resistividade é cerca de $65$ vezes a do
cobre, o que permite resistências úteis com comprimentos manejáveis. Com cobre,
os mesmos $\SI{22}{\ohm}$ exigiriam $\SI{259}{\meter}$ de fio --- inviável.}
\end{questao}

\begin{questao}[2]{}
Um cabo de alumínio ($\rho=\pot{2.8}{-8}\,\si{\ohm\meter}$) de
$\SI{200}{\meter}$ deve ter resistência de $\SI{0.5}{\ohm}$. Determine a seção
e o diâmetro.
\resposta{$A=\SI{11.2}{\milli\meter\squared}$; $d=\SI{3.78}{\milli\meter}$}
\resolucao{
\[
A=\frac{\rho L}{R}=\frac{\pot{2.8}{-8}\times200}{0{,}5}
=\pot{1.12}{-5}\,\si{\meter\squared}=\SI{11.2}{\milli\meter\squared}.
\]
\[
d=\sqrt{\frac{4A}{\pi}}=\sqrt{\frac{4\times\pot{1.12}{-5}}{\pi}}
=\SI{3.78}{\milli\meter}.
\]
\textbf{Valor comercial:} a bitola imediatamente superior é
$\SI{16}{\milli\meter\squared}$, que daria $R=\SI{0.35}{\ohm}$ --- com folga.
Adotar $\SI{10}{\milli\meter\squared}$ resultaria em $\SI{0.56}{\ohm}$,
acima do especificado.}
\end{questao}

\begin{questao}[2]{}
Mediu-se $R=\SI{2.5}{\ohm}$ em um fio de $\SI{30}{\meter}$ e seção
$\SI{0.5}{\milli\meter\squared}$. Determine a resistividade e identifique o
material provável.
\resposta{$\rho=\pot{4.2}{-8}\,\si{\ohm\meter}$ --- compatível com latão ou
zinco}
\resolucao{
\[
\rho=\frac{RA}{L}=\frac{2{,}5\times\pot{0.5}{-6}}{30}
=\pot{4.17}{-8}\,\si{\ohm\meter}.
\]
\textbf{Comparação:} cobre vale $\pot{1.7}{-8}$, alumínio
$\pot{2.8}{-8}$, zinco $\pot{5.9}{-8}$ e latão em torno de
$\pot{7}{-8}\,\si{\ohm\meter}$. O valor obtido está entre alumínio e zinco.\\
\textbf{Ressalva honesta:} a resistividade sozinha \emph{não} identifica o
material com segurança --- ligas de composições distintas podem coincidir. Ela
serve para descartar hipóteses (não é cobre puro, certamente) e como um entre
vários indícios.}
\end{questao}

\begin{questao}[2]{}
Um fio cilíndrico é \textbf{esticado} até dobrar seu comprimento, conservando o
volume e o material. Determine a nova resistência.
\resposta{$R'=4R$}
\resolucao{\textbf{Volume constante:} $V=AL=A'L'$. Com $L'=2L$:
\[
A'=\frac{AL}{2L}=\frac{A}{2}.
\]
\textbf{Nova resistência:}
\[
R'=\frac{\rho L'}{A'}=\frac{\rho(2L)}{A/2}=4\,\frac{\rho L}{A}=4R.
\]
\textbf{Dois efeitos que se multiplicam:} o comprimento dobra (fator $2$) e a
área cai à metade (outro fator $2$).\\
\textbf{Não confunda} com o caso em que $L$ e $A$ são alterados
\emph{independentemente} --- ali os fatores podem se cancelar. No estiramento,
eles se reforçam, porque o volume amarra as duas dimensões.}
\end{questao}

\begin{questao}[2]{}
Um resistor de fio tem $\SI{50}{\ohm}$ a $\SI{20}{\celsius}$, com
$\alpha=\SI{0.0039}{\per\celsius}$. Determine sua resistência a
$\SI{80}{\celsius}$.
\resposta{$R=\SI{61.7}{\ohm}$}
\resolucao{
\[
R=R_0\left[1+\alpha(T-T_0)\right]=50\left[1+0{,}0039(60)\right]
=50(1{,}234)=\SI{61.7}{\ohm}.
\]
Aumento de $23{,}4\%$ --- nada desprezível.\\
\textbf{Consequência para medição:} um ohmímetro que aplique corrente suficiente
para aquecer o componente lerá um valor \emph{maior} que o real a temperatura
ambiente. É por isso que medições de precisão especificam a corrente de teste e
a temperatura.}
\end{questao}

\begin{questao}[2]{}
Um estudante mediu, em um componente, os pares
$(\SI{2}{\volt};\SI{20}{\milli\ampere})$,
$(\SI{4}{\volt};\SI{40}{\milli\ampere})$ e
$(\SI{6}{\volt};\SI{57}{\milli\ampere})$. O componente é ôhmico?
\resposta{Não --- a razão $U/I$ cresce de $100$ para $\SI{105}{\ohm}$}
\resolucao{Calculando $R=U/I$ em cada ponto:
\[
\frac{2}{0{,}020}=100;
\qquad
\frac{4}{0{,}040}=100;
\qquad
\frac{6}{0{,}057}=\SI{105}{\ohm}.
\]
Os dois primeiros pontos são consistentes; o terceiro destoa em $5\%$.\\
\textbf{Diagnóstico provável: autoaquecimento.} Em $\SI{6}{\volt}$ a potência é
$\SI{0.34}{\watt}$, contra $\SI{0.04}{\watt}$ no primeiro ponto --- oito vezes
mais. Se o material tiver coeficiente positivo, a resistência sobe com a
temperatura.\\
\textbf{Como confirmar:} refaça o terceiro ponto \emph{rapidamente}, antes que o
componente aqueça. Se a leitura voltar a $\SI{100}{\ohm}$, o material é ôhmico
--- e o desvio era térmico, não uma propriedade elétrica intrínseca.}
\end{questao}

\begin{questao}[2]{}
Compare a seção necessária para obter a mesma resistência com cobre
($\rho=\pot{1.7}{-8}$) e com alumínio ($\pot{2.8}{-8}\,\si{\ohm\meter}$),
mantido o comprimento.
\resposta{$A_{Al}/A_{Cu}=1{,}65$}
\resolucao{Para $R$ e $L$ iguais, $A\propto\rho$:
\[
\frac{A_{Al}}{A_{Cu}}=\frac{\rho_{Al}}{\rho_{Cu}}=\frac{2{,}8}{1{,}7}=1{,}65.
\]
O alumínio exige $65\%$ mais seção.\\
\textbf{Por que ainda assim se usa alumínio:} sua densidade é
$\SI{2.7}{\gram\per\centi\meter\cubed}$ contra $\SI{8.9}{}$ do cobre. Para a
mesma resistência, o cabo de alumínio pesa
$1{,}65\times(2{,}7/8{,}9)=0{,}50$ --- \textbf{metade} do peso, e custa bem
menos. É por isso que linhas de transmissão aéreas são de alumínio, onde o peso
é crítico, enquanto instalações prediais são de cobre, onde o espaço nos
eletrodutos é que limita.}
\end{questao}

\begin{questao}[2]{}
Um fio de cobre de $\SI{100}{\meter}$ e $\SI{1.5}{\milli\meter\squared}$
conduz $\SI{10}{\ampere}$. Determine a resistência, a queda de tensão e a
potência dissipada.
\resposta{$R=\SI{1.13}{\ohm}$; $\Delta U=\SI{11.3}{\volt}$;
$P=\SI{113}{\watt}$}
\resolucao{
\[
R=\frac{\pot{1.7}{-8}\times100}{\pot{1.5}{-6}}=\SI{1.13}{\ohm};
\qquad
\Delta U=RI=\SI{11.3}{\volt};
\qquad
P=RI^{2}=\SI{113}{\watt}.
\]
\textbf{Avaliação:} em uma rede de $\SI{220}{\volt}$, essa queda representa
$5{,}2\%$ --- acima do limite usual de $4\%$ para circuitos terminais. E
$\SI{113}{\watt}$ dissipados ao longo do cabo é aquecimento real, que exige
verificação de capacidade de condução.\\
\textbf{Conclusão:} a bitola está subdimensionada para esse comprimento. O
critério de queda de tensão, e não o de corrente admissível, é que domina em
circuitos longos.}
\end{questao}

\blocoaprofundamento

\begin{questao}[3]{}
Projete um resistor de fio de $\SI{10}{\ohm}$ usando material de
$\rho=\pot{5}{-7}\,\si{\ohm\meter}$ e fio de $\SI{0.5}{\milli\meter}$ de
diâmetro.
\begin{itensq}
\item Determine o comprimento necessário.
\item Determine a corrente máxima para $\SI{5}{\watt}$ de dissipação.
\item Comente a construção.
\end{itensq}
\resposta{(a) $L=\SI{3.93}{\meter}$; (b) $I=\SI{0.707}{\ampere}$}
\resolucao{(a) Seção:
$A=\dfrac{\pi d^{2}}{4}=\dfrac{\pi(\pot{5}{-4})^{2}}{4}
=\pot{1.963}{-7}\,\si{\meter\squared}$.
\[
L=\frac{RA}{\rho}=\frac{10\times\pot{1.963}{-7}}{\pot{5}{-7}}
=\SI{3.93}{\meter}.
\]
(b) $I=\sqrt{P/R}=\sqrt{5/10}=\SI{0.707}{\ampere}$, com
$U=\SI{7.07}{\volt}$.\\
(c) \textbf{Construção.} Quase quatro metros de fio precisam ser enrolados
sobre um corpo cerâmico. Três cuidados:
\begin{itemize}
\item \textbf{Espaçamento entre espiras}, para evitar curtos e permitir
circulação de ar.
\item \textbf{Área de dissipação suficiente} --- $\SI{5}{\watt}$ exigem corpo de
alguns centímetros, e a temperatura de superfície pode passar de
$\SI{100}{\celsius}$.
\item \textbf{Indutância parasita}: um enrolamento simples de $\SI{3.93}{\meter}$
comporta-se como bobina. Em corrente contínua isso é irrelevante, mas em sinais
rápidos é preciso enrolamento \emph{bifilar}, em que a ida e a volta se
cancelam.
\end{itemize}}
\end{questao}

\begin{questao}[3]{}
Um resistor nominal de $\SI{100}{\ohm}$ tem tolerância de $\SI{5}{\percent}$ e
coeficiente térmico de $400\,\mathrm{ppm/^{\circ}C}$, especificado a
$\SI{20}{\celsius}$. Determine a faixa total de resistência na operação entre
$\SI{0}{\celsius}$ e $\SI{70}{\celsius}$.
\resposta{de $\SI{94.2}{\ohm}$ a $\SI{107.1}{\ohm}$ ($-5{,}8\%$ a
$+7{,}1\%$)}
\resolucao{\textbf{Deriva térmica:} $400\,\mathrm{ppm}=\pot{4}{-4}$ por
$\si{\celsius}$.
\[
\SI{70}{\celsius}:\ \Delta T=+50 \Rightarrow +2{,}0\%;
\qquad
\SI{0}{\celsius}:\ \Delta T=-20 \Rightarrow -0{,}8\%.
\]
\textbf{Combinação de pior caso} (os dois efeitos no mesmo sentido):
\[
R_{\max}=100(1{,}05)(1{,}020)=\SI{107.1}{\ohm};
\qquad
R_{\min}=100(0{,}95)(0{,}992)=\SI{94.2}{\ohm}.
\]
\textbf{Observações de projeto.}
\begin{itemize}
\item A tolerância de fabricação domina --- ela responde por $5$ dos
$7{,}1$ pontos.
\item A faixa é \emph{assimétrica}, porque a temperatura de operação se afasta
mais para cima que para baixo do valor de referência.
\item Se $7\%$ for demais, apertar a tolerância para $1\%$ traria a faixa a
$-1{,}8\%$/$+3{,}0\%$; trocar o coeficiente para
$25\,\mathrm{ppm}$ traria a $-5{,}0\%$/$+5{,}1\%$. \textbf{Ataque primeiro a
maior parcela.}
\end{itemize}}
\end{questao}

\begin{questao}[3]{}
Uma lâmpada incandescente apresenta $\SI{2.00}{\ampere}$ a
$\SI{12.0}{\volt}$ e $\SI{2.02}{\ampere}$ a $\SI{12.5}{\volt}$.
\begin{itensq}
\item Determine a resistência estática no primeiro ponto.
\item Determine a resistência dinâmica.
\item Explique a diferença.
\end{itensq}
\resposta{$R_{est}=\SI{6}{\ohm}$; $r_{din}=\SI{25}{\ohm}$}
\resolucao{(a) $R_{est}=\dfrac{U}{I}=\dfrac{12{,}0}{2{,}00}=\SI{6}{\ohm}$.\\
(b) $r_{din}=\dfrac{\Delta U}{\Delta I}=\dfrac{0{,}5}{0{,}02}=\SI{25}{\ohm}$.\\
(c) \textbf{A dinâmica é mais de quatro vezes maior.} A razão é o
\textbf{autoaquecimento}: aumentar a tensão eleva a potência, o filamento
esquenta e sua resistência sobe --- de modo que a corrente cresce muito menos
que proporcionalmente.\\
\textbf{Cada grandeza responde a uma pergunta diferente:}
\begin{itemize}
\item $R_{est}$ governa o \emph{balanço de potência} no ponto de operação:
$P=U^{2}/R_{est}=\SI{24}{\watt}$ \checkmark.
\item $r_{din}$ governa a \emph{resposta a pequenas variações} --- é ela que
entra em qualquer modelo incremental.
\end{itemize}
\textbf{A lâmpada não é ôhmica em regime}, mas é aproximadamente ôhmica a
temperatura constante. É o acoplamento térmico, e não a física da condução, que
quebra a linearidade.}
\end{questao}

\begin{questao}[3]{}
Um cabo de cobre deve alimentar uma carga de $\SI{20}{\ampere}$ a
$\SI{220}{\volt}$, a $\SI{40}{\meter}$ de distância, com queda de tensão
inferior a $\SI{3}{\percent}$.
\begin{itensq}
\item Determine a resistência máxima admissível do cabo.
\item Determine a seção mínima.
\item Escolha a bitola comercial.
\end{itensq}
\resposta{$R\le\SI{0.33}{\ohm}$; $A\ge\SI{4.12}{\milli\meter\squared}$;
adotar $\SI{6}{\milli\meter\squared}$}
\resolucao{(a) Queda admissível: $0{,}03\times220=\SI{6.6}{\volt}$.
\[
R_{\max}=\frac{6{,}6}{20}=\SI{0.33}{\ohm}.
\]
(b) \textbf{Atenção ao comprimento:} a corrente percorre \emph{ida e volta},
logo $L=2\times40=\SI{80}{\meter}$.
\[
A\ge\frac{\rho L}{R_{\max}}=\frac{\pot{1.7}{-8}\times80}{0{,}33}
=\pot{4.12}{-6}\,\si{\meter\squared}=\SI{4.12}{\milli\meter\squared}.
\]
(c) A bitola comercial imediatamente superior é
$\SI{6}{\milli\meter\squared}$, que dá $R=\SI{0.227}{\ohm}$ e queda de
$\SI{4.5}{\volt}$ ($2{,}1\%$) --- dentro da especificação com folga.\\
\textbf{Verificação obrigatória do outro critério:} $\SI{6}{\milli\meter\squared}$
suporta cerca de $\SI{41}{\ampere}$ em eletroduto, bem acima dos
$\SI{20}{\ampere}$ da carga. \textbf{Aqui a queda de tensão é o critério
dominante} --- e é o que costuma ocorrer em circuitos longos. Em circuitos
curtos, a corrente admissível é que decide.}
\end{questao}

\begin{questao}[3]{}
Um fio é esticado até que seu comprimento fique $n$ vezes maior, conservando o
volume.
\begin{itensq}
\item Deduza a nova resistência em função de $n$.
\item Avalie para $n=3$ e $n=5$.
\item Determine a redução do diâmetro em cada caso.
\end{itensq}
\resposta{$R'=n^{2}R$; $9R$ e $25R$; $d'/d=1/\sqrt{n}$}
\resolucao{(a) Volume constante: $AL=A'L'$ com $L'=nL$, logo
$A'=A/n$.
\[
R'=\frac{\rho(nL)}{A/n}=n^{2}\frac{\rho L}{A}=n^{2}R.
\]
(b) $n=3\Rightarrow9R$; $n=5\Rightarrow25R$.\\
(c) Como $A\propto d^{2}$ e $A'=A/n$:
\[
\frac{d'}{d}=\frac{1}{\sqrt{n}}
\Rightarrow
n=3:\ 0{,}577;
\qquad
n=5:\ 0{,}447.
\]
\textbf{Aplicação prática --- o extensômetro.} Um fio colado a uma peça que se
deforma sofre exatamente esse efeito. Para pequenas deformações
$\varepsilon=\Delta L/L$, a expansão de $n^{2}$ dá
\[
\frac{\Delta R}{R}\approx2\varepsilon,
\]
ou seja, o fio funciona como sensor de deformação com \emph{fator de
sensibilidade} próximo de $2$. É o princípio das células de carga e das balanças
eletrônicas.}
\end{questao}

\begin{questao}[3]{}
Uma barra condutora de comprimento $L=\SI{1}{\meter}$ tem seção que varia
linearmente segundo $A(x)=A_0\left(1+x/L\right)$, com
$A_0=\SI{1}{\milli\meter\squared}$ e $\rho=\pot{1.7}{-8}\,\si{\ohm\meter}$.
\begin{itensq}
\item Determine a resistência total.
\item Compare com uma barra uniforme de seção $A_0$ e com uma de seção média.
\end{itensq}
\resposta{$R=\dfrac{\rho L\ln2}{A_0}=\SI{11.8}{\milli\ohm}$}
\resolucao{(a) Somando as resistências de fatias infinitesimais em série:
\[
R=\int_0^{L}\frac{\rho\,\mathrm{d}x}{A(x)}
=\frac{\rho}{A_0}\int_0^{L}\frac{\mathrm{d}x}{1+x/L}
=\frac{\rho L}{A_0}\Big[\ln(1+x/L)\Big]_0^{L}
=\frac{\rho L\ln2}{A_0}.
\]
\[
R=\frac{\pot{1.7}{-8}\times1\times0{,}693}{\pot{1}{-6}}
=\SI{11.8}{\milli\ohm}.
\]
(b)
\begin{center}
\begin{tabular}{lc}
\hline
Modelo & $R$\\
\hline
uniforme com $A_0$ (menor seção) & $\SI{17.0}{\milli\ohm}$\\
\textbf{seção variável (exato)} & $\SI{11.8}{\milli\ohm}$\\
uniforme com $\bar A=1{,}5A_0$ & $\SI{11.3}{\milli\ohm}$\\
\hline
\end{tabular}
\end{center}
\textbf{O valor exato é maior que o da seção média} --- porque a resistência
depende de $1/A$, e a média de $1/A$ é maior que $1/\bar A$ (desigualdade das
médias). O erro de usar a seção média é de $4{,}3\%$, para menos.\\
\textbf{Regra:} em seções variáveis, integrar $\rho\,\mathrm{d}x/A(x)$. Usar a
seção média subestima; usar a seção mínima superestima. As duas aproximações
servem como limites.}
\end{questao}

\begin{questao}[3]{}
Um sensor Pt100 tem $R_0=\SI{100}{\ohm}$ a $\SI{0}{\celsius}$ e
$\alpha=\SI{0.00385}{\per\celsius}$, alimentado por corrente constante de
$\SI{1}{\milli\ampere}$.
\begin{itensq}
\item Determine a tensão a $\SI{0}{\celsius}$ e a sensibilidade em
      $\si{\milli\volt\per\celsius}$.
\item Explique a vantagem da alimentação por corrente constante.
\item Verifique o autoaquecimento.
\end{itensq}
\resposta{$U_0=\SI{100}{\milli\volt}$;
$\SI{0.385}{\milli\volt\per\celsius}$}
\resolucao{(a) $U_0=IR_0=\pot{1}{-3}(100)=\SI{100}{\milli\volt}$.
\[
\frac{\mathrm{d}U}{\mathrm{d}T}=I\,R_0\,\alpha
=\pot{1}{-3}(100)(0{,}00385)=\SI{0.385}{\milli\volt\per\celsius}.
\]
(b) \textbf{Vantagem decisiva: linearidade exata.} Com $I$ constante,
$U=IR$ é \emph{proporcional} a $R$, e como $R$ é linear em $T$, a tensão também
é. Não há distorção alguma.\\
Com \emph{tensão} constante em um divisor, a saída seria
$U\,R/(R+R_{fixo})$ --- uma função \textbf{não linear} de $R$, exigindo
linearização posterior.\\
(c) \textbf{Autoaquecimento:} $P=RI^{2}=100(\pot{1}{-6})=\SI{0.1}{\milli\watt}$.
Com resistência térmica típica de $\SI{100}{\celsius\per\watt}$ ao ar, a
elevação é de $\SI{0.01}{\celsius}$ --- desprezível.\\
\textbf{Mas cuidado:} elevar a corrente para $\SI{10}{\milli\ampere}$ (para
ganhar sinal) multiplicaria a potência por $100$, levando o erro a
$\SI{1}{\celsius}$. \textbf{A corrente de excitação é um compromisso entre
relação sinal-ruído e erro de autoaquecimento} --- e $\SI{1}{\milli\ampere}$ é
o valor de consenso para Pt100.}
\end{questao}

\begin{questao}[3]{}
Mede-se um resistor de $\SI{0.1}{\ohm}$ com um ohmímetro cujos cabos têm
$\SI{0.2}{\ohm}$ cada.
\begin{itensq}
\item Determine a leitura pelo método de \textbf{dois fios}.
\item Determine o erro.
\item Explique como o método de \textbf{quatro fios} o elimina.
\end{itensq}
\resposta{Leitura de $\SI{0.5}{\ohm}$ --- erro de $400\%$}
\resolucao{(a) O instrumento mede tudo o que está em série:
\[
R_{lido}=0{,}1+2(0{,}2)=\SI{0.5}{\ohm}.
\]
(b) Erro de $\SI{0.4}{\ohm}$, ou $400\%$ --- a leitura é \emph{cinco vezes} o
valor real. Para resistências pequenas, o método de dois fios é simplesmente
inutilizável.\\
(c) \textbf{Método de quatro fios (Kelvin).} Usam-se dois pares independentes:
\begin{itemize}
\item Um par \textbf{injeta} a corrente conhecida $I$. A queda nesses cabos
existe, mas é irrelevante --- ela apenas exige um pouco mais de tensão da fonte.
\item O outro par \textbf{mede} a tensão diretamente sobre o resistor, com um
voltímetro de altíssima impedância. Como praticamente não circula corrente
nesses cabos, sua resistência \emph{não produz queda}.
\end{itemize}
Assim $R=U_{medido}/I$ devolve exatamente $\SI{0.1}{\ohm}$.\\
\textbf{Quando usar:} sempre que $R$ for comparável ou menor que a resistência
dos cabos --- tipicamente abaixo de alguns ohms. É o método obrigatório para
shunts, enrolamentos de motor e contatos.}
\end{questao}

\begin{questao}[3]{}
Dois resistores têm o mesmo valor nominal de $\SI{100}{\ohm}$, mas potências de
$\frac14\,\si{\watt}$ e $\SI{5}{\watt}$.
\begin{itensq}
\item Determine a tensão e a corrente máximas de cada um.
\item Ambos obedecem igualmente à lei de Ohm?
\item Explique o que a especificação de potência realmente limita.
\end{itensq}
\resposta{$\SI{5}{\volt}$/$\SI{50}{\milli\ampere}$ e
$\SI{22.4}{\volt}$/$\SI{224}{\milli\ampere}$; ambos são ôhmicos}
\resolucao{(a)
\[
U_{\max}=\sqrt{PR}:\quad \sqrt{0{,}25(100)}=\SI{5}{\volt};
\qquad \sqrt{5(100)}=\SI{22.4}{\volt};
\]
\[
I_{\max}=\sqrt{P/R}:\quad \SI{50}{\milli\ampere}
\quad\text{e}\quad \SI{224}{\milli\ampere}.
\]
(b) \textbf{Sim, ambos obedecem igualmente.} A lei de Ohm relaciona $U$ e $I$
através de $R$ --- e $R$ é o mesmo. A potência nominal \emph{não} aparece na
lei.\\
(c) \textbf{A potência nominal limita a faixa de validade}, não a relação em
si. Ela informa até onde o componente pode operar \emph{antes de se
autodestruir ou sair da linearidade por aquecimento}.\\
\textbf{Analogia útil:} a lei de Ohm é como a relação entre força e deformação
em uma mola; a potência nominal é o limite de escoamento do material. Duas
molas de mesma constante e materiais diferentes obedecem à mesma lei --- até que
uma delas se rompa.\\
\textbf{Diferença física:} o resistor de $\SI{5}{\watt}$ é maior, com mais área
de troca térmica. Mesma resistência, mesma lei, capacidade de dissipação
$20$ vezes maior.}
\end{questao}

\begin{questao}[3]{}
Um resistor de $\SI{100}{\ohm}$ com $\alpha=\SI{+0.004}{\per\celsius}$ dissipa
$\SI{1}{\watt}$ e tem resistência térmica de
$\SI{100}{\celsius\per\watt}$.
\begin{itensq}
\item Determine a elevação de temperatura e a nova resistência.
\item Analise o efeito sob tensão constante e sob corrente constante.
\end{itensq}
\resposta{$\Delta T=\SI{100}{\celsius}$; $R=\SI{140}{\ohm}$}
\resolucao{(a)
\[
\Delta T=R_{\text{térm}}P=100(1)=\SI{100}{\celsius};
\qquad
R=100\left[1+0{,}004(100)\right]=\SI{140}{\ohm}.
\]
Aumento de $40\%$ apenas por autoaquecimento --- muito acima da tolerância
típica de fabricação.\\
(b) \textbf{Sob tensão constante:} $P=U^{2}/R$ \emph{diminui} quando $R$ sobe.
Menos potência, menos aquecimento, $R$ estabiliza. \textbf{Realimentação
negativa --- estável.}\\
\textbf{Sob corrente constante:} $P=RI^{2}$ \emph{aumenta} quando $R$ sobe.
Mais potência, mais aquecimento, $R$ cresce ainda mais.
\textbf{Realimentação positiva.}\\
\textbf{Condição de estabilidade} no caso de corrente constante: o incremento de
potência por grau, $\alpha R I^{2}$, deve ser menor que a capacidade de
escoar calor, $1/R_{\text{térm}}$:
\[
\alpha R I^{2}R_{\text{térm}}<1.
\]
Aqui, $0{,}004(100)(0{,}01)(100)=0{,}4<1$ \checkmark --- estável, mas com apenas
$2{,}5$ vezes de margem. Dobrar a corrente levaria o produto a $1{,}6$ e
provocaria fuga térmica.\\
\textbf{Materiais com $\alpha<0$} (termistores NTC, carbono) invertem os dois
casos --- e são instáveis justamente sob \emph{tensão} constante.}
\end{questao}

\begin{questao}[3]{}
Um resistor de $\SI{100}{\ohm}$ e $\frac14\,\si{\watt}$ deve ser testado quanto
à ohmicidade. Projete o ensaio.
\begin{itensq}
\item Determine a faixa segura de tensão.
\item Proponha os pontos de medição e o critério de decisão.
\item Indique os cuidados experimentais.
\end{itensq}
\resposta{Até $\SI{5}{\volt}$; usar $20\%$ da potência nominal como teto
prático}
\resolucao{(a) $U_{\max}=\sqrt{PR}=\SI{5}{\volt}$ no limite absoluto. Para
evitar autoaquecimento --- que é justamente o que falsearia o ensaio --- adote
$20\%$ da potência nominal:
\[
P_{ensaio}=\SI{50}{\milli\watt}
\;\Longrightarrow\;
U_{\max}=\sqrt{0{,}05(100)}=\SI{2.24}{\volt}.
\]
(b) \textbf{Pontos:} cinco valores igualmente espaçados entre $0{,}5$ e
$\SI{2.2}{\volt}$, medindo $U$ e $I$ em cada um.\\
\textbf{Critério:} calcular $R=U/I$ em todos os pontos. O componente é ôhmico se
a dispersão dos valores ficar dentro da incerteza combinada dos instrumentos
(tipicamente $1$ a $2\%$). Alternativamente, ajustar uma reta a $U\times I$ e
verificar se o intercepto é nulo dentro da incerteza --- intercepto não nulo
indica tensão de limiar, como em diodos.\\
(c) \textbf{Cuidados.}
\begin{itemize}
\item \textbf{Medir rápido} e desligar entre pontos, para o componente não
aquecer cumulativamente.
\item \textbf{Ordem aleatória} dos pontos, e não crescente --- assim um eventual
aquecimento progressivo aparece como dispersão, e não como tendência que
poderia ser confundida com não linearidade real.
\item \textbf{Método de quatro fios} se $R$ for pequeno.
\item \textbf{Repetir o primeiro ponto ao final:} se o valor mudou, houve
aquecimento e o ensaio precisa ser refeito com potência menor.
\end{itemize}}
\end{questao}

\blocodesafio

\begin{questao}[4]{}
\textbf{(Demonstração --- deformação e resistência.)} Um fio de comprimento $L$
e seção $A$ sofre pequena deformação axial $\varepsilon=\Delta L/L$, a volume
constante.
\begin{itensq}
\item Deduza $\Delta R/R$ em primeira ordem.
\item Determine o fator de sensibilidade.
\item Discuta o efeito de a resistividade também variar com a deformação.
\end{itensq}
\resposta{$\Delta R/R\approx2\varepsilon$; fator $\approx2$}
\resolucao{(a) Com volume constante, $L'=L(1+\varepsilon)$ e
$A'=A/(1+\varepsilon)$:
\[
R'=\frac{\rho L(1+\varepsilon)}{A/(1+\varepsilon)}
=R(1+\varepsilon)^{2}\approx R(1+2\varepsilon),
\]
para $\varepsilon\ll1$. Logo
\[
\frac{\Delta R}{R}\approx2\varepsilon.
\]
(b) O \textbf{fator de sensibilidade} (\emph{gauge factor}) é
$k=\dfrac{\Delta R/R}{\varepsilon}\approx2$.\\
\textbf{Ordem de grandeza prática:} uma deformação de
$\SI{1000}{\micro\meter\per\meter}$ ($\varepsilon=\pot{1}{-3}$) produz
$\Delta R/R=0{,}2\%$ --- em um extensômetro de $\SI{350}{\ohm}$, apenas
$\SI{0.7}{\ohm}$. Daí a necessidade de ponte de Wheatstone e amplificação de
instrumentação.\\
(c) \textbf{Efeito piezorresistivo.} Em metais, $\rho$ varia pouco com a
deformação e $k\approx2$ decorre quase só da geometria. Em
\textbf{semicondutores}, a deformação altera fortemente a mobilidade dos
portadores, e $k$ chega a $100$ ou mais --- ganho enorme de sensibilidade.\\
\textbf{Contrapartida:} extensômetros semicondutores são muito mais sensíveis à
temperatura e menos lineares. \textbf{A escolha entre metal e semicondutor é um
compromisso entre sensibilidade e estabilidade} --- e por isso células de carga
de precisão ainda usam extensômetros metálicos.}
\end{questao}

\begin{questao}[4]{}
\textbf{(Integração --- seção variável.)} Uma barra tem seção
$A(x)=A_0\left(1+x/L\right)$.
\begin{itensq}
\item Deduza $R$ por integração.
\item Generalize para $A(x)=A_0(1+kx/L)$.
\item Analise os limites $k\to0$ e $k\to\infty$.
\end{itensq}
\resposta{$R=\dfrac{\rho L}{A_0}\cdot\dfrac{\ln(1+k)}{k}$; para $k=1$,
$\ln2$}
\resolucao{(a) Fatias de espessura $\mathrm{d}x$ estão em série:
\[
R=\int_0^{L}\frac{\rho\,\mathrm{d}x}{A_0(1+x/L)}
=\frac{\rho L}{A_0}\ln 2.
\]
(b) Com $A(x)=A_0(1+kx/L)$:
\[
R=\frac{\rho}{A_0}\int_0^{L}\frac{\mathrm{d}x}{1+kx/L}
=\frac{\rho L}{A_0k}\Big[\ln(1+kx/L)\Big]_0^{L}
=\frac{\rho L}{A_0}\cdot\frac{\ln(1+k)}{k}.
\]
(c) \textbf{Limite $k\to0$} (seção uniforme). Usando
$\ln(1+k)\approx k-k^{2}/2$:
\[
\frac{\ln(1+k)}{k}\to1
\;\Longrightarrow\;
R\to\frac{\rho L}{A_0}.\ \checkmark
\]
Recupera-se a fórmula elementar --- teste obrigatório de qualquer generalização.\\
\textbf{Limite $k\to\infty$} (seção crescendo muito rápido):
\[
\frac{\ln(1+k)}{k}\to0
\;\Longrightarrow\;
R\to0.
\]
A barra fica tão larga na maior parte de seu comprimento que sua resistência se
concentra apenas na região inicial estreita --- e o total tende a zero, embora
\emph{logaritmicamente devagar}. Mesmo com $k=1000$, o fator ainda vale
$0{,}0069$, e não zero.\\
\textbf{Leitura geral:} a resistência é dominada pela região de \emph{menor}
seção. É por isso que um estrangulamento local em um condutor concentra
praticamente toda a queda de tensão --- e todo o aquecimento.}
\end{questao}

\begin{questao}[4]{}
\textbf{(Ponto de operação com fonte não ideal.)} Um resistor de
$\SI{100}{\ohm}$ é ligado a uma fonte cuja característica é
$U=24-8I$ (com $U$ em volts e $I$ em ampères).
\begin{itensq}
\item Determine o ponto de operação.
\item Determine a potência no resistor e o rendimento.
\item Repita para um resistor de $\SI{8}{\ohm}$ e compare.
\end{itensq}
\resposta{$I=\SI{0.222}{\ampere}$, $U=\SI{22.2}{\volt}$;
$P=\SI{4.94}{\watt}$, $\eta=92{,}6\%$}
\resolucao{(a) A fonte impõe $U=24-8I$ e o resistor impõe $U=100I$. Igualando:
\[
100I=24-8I
\;\Longrightarrow\;
108I=24
\;\Longrightarrow\;
I=\SI{0.222}{\ampere};
\qquad
U=\SI{22.2}{\volt}.
\]
(b) $P_L=UI=22{,}2(0{,}222)=\SI{4.94}{\watt}$.\\
Rendimento: $\eta=\dfrac{R_L}{R_L+R_{Th}}=\dfrac{100}{108}=92{,}6\%$.\\
(c) \textbf{Com $\SI{8}{\ohm}$:} $8I=24-8I\Rightarrow I=\SI{1.5}{\ampere}$ e
$U=\SI{12}{\volt}$.
\[
P_L=12(1{,}5)=\SI{18}{\watt};
\qquad
\eta=\frac{8}{16}=50\%.
\]
\textbf{Comparação reveladora.} O resistor de $\SI{8}{\ohm}$ está \emph{casado}
com a fonte e recebe a potência máxima ($\SI{18}{\watt}$, quase quatro vezes
mais), mas desperdiça metade da energia internamente. O de
$\SI{100}{\ohm}$ recebe pouco, porém com rendimento de $92{,}6\%$.\\
\textbf{Ponto de método:} o problema é a interseção de duas características
$U\times I$ --- uma reta de fonte e uma reta de carga. Se a carga fosse não
linear (um diodo, por exemplo), o procedimento seria idêntico: basta trocar a
segunda reta pela curva do dispositivo.}
\end{questao}

\begin{questao}[4]{}
\textbf{(Estabilidade térmica.)} Um resistor de resistência $R$, coeficiente
$\alpha$ e resistência térmica $R_{\text{térm}}$ é alimentado por corrente constante
$I$.
\begin{itensq}
\item Deduza a condição de estabilidade térmica.
\item Avalie para $R=\SI{100}{\ohm}$, $\alpha=\SI{0.004}{\per\celsius}$,
      $R_{\text{térm}}=\SI{100}{\celsius\per\watt}$ e $I=\SI{100}{\milli\ampere}$.
\item Determine a corrente crítica.
\end{itensq}
\resposta{Condição: $\alpha R I^{2}R_{\text{térm}}<1$; aqui $0{,}4$ --- estável;
$I_c=\SI{158}{\milli\ampere}$}
\resolucao{(a) Em equilíbrio, $\Delta T=R_{\text{térm}}P=R_{\text{térm}}RI^{2}$ e
$R=R_0(1+\alpha\Delta T)$. Substituindo,
\[
R=R_0\left(1+\alpha R_{\text{térm}}RI^{2}\right)
\;\Longrightarrow\;
R\left(1-\alpha R_{\text{térm}}R_0I^{2}\right)=R_0.
\]
Existe solução finita e estável se, e somente se, o parênteses for positivo:
\[
\boxed{\alpha\,R_0\,I^{2}\,R_{\text{térm}}<1}
\]
Fisicamente: o acréscimo de potência por grau de aquecimento deve ser
\emph{menor} que a capacidade de escoar calor. Caso contrário, cada grau
adicional gera mais calor do que consegue dissipar --- \textbf{fuga térmica}.\\
(b)
\[
0{,}004\times100\times(0{,}1)^{2}\times100=0{,}4<1\ \checkmark
\]
Estável, com margem de $2{,}5$ vezes. A resistência de equilíbrio é
$R=100/(1-0{,}4)=\SI{167}{\ohm}$, e a elevação, $\SI{167}{\celsius}$ ---
elevada, mas estável.\\
(c) Corrente crítica:
\[
I_c=\sqrt{\frac{1}{\alpha R_0R_{\text{térm}}}}=\sqrt{\frac{1}{0{,}004(100)(100)}}
=\sqrt{0{,}025}=\SI{158}{\milli\ampere}.
\]
\textbf{Margem estreita:} apenas $58\%$ acima da corrente de operação. E note
que a estabilidade \emph{matemática} não garante a \emph{sobrevivência} --- a
$\SI{167}{\celsius}$ o componente já pode estar além de sua temperatura máxima.\\
\textbf{Sob tensão constante} o problema não existe: $P=U^{2}/R$ decresce com
$R$, e a realimentação é sempre negativa.}
\end{questao}

\begin{questao}[4]{}
\textbf{(Dois critérios de dimensionamento.)} Um cabo deve alimentar
$\SI{30}{\ampere}$ a $\SI{220}{\volt}$, com queda máxima de
$\SI{4}{\percent}$.
\begin{itensq}
\item Determine a seção exigida pelo critério de queda de tensão, para
      $\SI{20}{\meter}$ e para $\SI{80}{\meter}$.
\item Sabendo que $\SI{6}{\milli\meter\squared}$ suporta cerca de
      $\SI{41}{\ampere}$, determine qual critério domina em cada caso.
\item Formule a regra geral.
\end{itensq}
\resposta{$\SI{3.87}{\milli\meter\squared}$ e
$\SI{15.5}{\milli\meter\squared}$; a corrente domina no cabo curto, a queda no
longo}
\resolucao{(a) Queda admissível: $0{,}04(220)=\SI{8.8}{\volt}$, logo
$R_{\max}=8{,}8/30=\SI{0.293}{\ohm}$.
\[
\SI{20}{\meter}:\ A\ge\frac{\pot{1.7}{-8}(40)}{0{,}293}
=\SI{2.32}{\milli\meter\squared};
\]
\[
\SI{80}{\meter}:\ A\ge\frac{\pot{1.7}{-8}(160)}{0{,}293}
=\SI{9.28}{\milli\meter\squared}.
\]
(lembrando o fator $2$ de ida e volta).\\
(b) \textbf{Cabo de $\SI{20}{\meter}$:} a queda exige
$\SI{2.32}{\milli\meter\squared}$, mas a corrente de $\SI{30}{\ampere}$ exige
pelo menos $\SI{6}{\milli\meter\squared}$. \textbf{O critério de corrente
domina.}\\
\textbf{Cabo de $\SI{80}{\meter}$:} a queda exige
$\SI{9.28}{\milli\meter\squared}$, acima do que a corrente pediria.
\textbf{O critério de queda domina}, e adota-se
$\SI{10}{\milli\meter\squared}$.\\
(c) \textbf{Regra geral.} Os dois critérios escalam de formas diferentes:
\begin{itemize}
\item \textbf{Corrente admissível:} depende só de $I$ --- \emph{independe do
comprimento}. É um critério de segurança térmica do próprio cabo.
\item \textbf{Queda de tensão:} $A\propto IL/\Delta U$ --- cresce
\emph{linearmente} com o comprimento. É um critério de qualidade da energia
entregue.
\end{itemize}
Existe um \textbf{comprimento de cruzamento} em que os dois coincidem; abaixo
dele manda a corrente, acima manda a queda. Aqui, com
$\SI{6}{\milli\meter\squared}$, o cruzamento ocorre por volta de
$\SI{52}{\meter}$.\\
\textbf{Procedimento correto:} calcular os \emph{dois} e adotar a maior seção.
Verificar apenas um deles é a origem mais comum de instalações que "funcionam"
mas apresentam tensão baixa na ponta.}
\end{questao}

\begin{questao}[4]{}
\textbf{(Sensor a corrente constante.)} Compare duas formas de converter a
variação de um sensor resistivo em tensão.
\begin{itensq}
\item Obtenha $U(R)$ para excitação por corrente constante e para divisor de
      tensão.
\item Compare a linearidade das duas.
\item Determine a sensibilidade de cada uma no ponto de operação
      $R=R_{fixo}$.
\end{itensq}
\resposta{Corrente constante: $U=IR$, exatamente linear; divisor:
$U=U_0R/(R+R_f)$, não linear}
\resolucao{(a) \textbf{Corrente constante:} $U=IR$ --- proporcional a $R$.\\
\textbf{Divisor:} $U=U_0\dfrac{R}{R+R_f}$ --- função racional de $R$.\\
(b) \textbf{Linearidade.} A primeira é \emph{exatamente} linear, sempre. A
segunda só é aproximadamente linear para pequenas variações; expandindo em torno
de $R=R_f$:
\[
U\approx\frac{U_0}{2}\left[1+\frac{\Delta R}{2R_f}
-\frac{1}{4}\left(\frac{\Delta R}{R_f}\right)^{2}+\cdots\right].
\]
O termo quadrático produz erro de linearidade. Para
$\Delta R/R_f=10\%$, ele vale $0{,}25\%$ da leitura --- pequeno, mas suficiente
para limitar a exatidão em medições finas.\\
(c) \textbf{Sensibilidades.}
\[
\text{corrente constante:}\quad \frac{\mathrm{d}U}{\mathrm{d}R}=I;
\qquad
\text{divisor:}\quad \frac{\mathrm{d}U}{\mathrm{d}R}
=\frac{U_0R_f}{(R+R_f)^{2}}\bigg|_{R=R_f}=\frac{U_0}{4R_f}.
\]
Ambas podem ser ajustadas ao mesmo valor escolhendo $I=U_0/(4R_f)$ --- a
sensibilidade não distingue os métodos no ponto de operação.\\
\textbf{O que decide é o resto.}
\begin{itemize}
\item \textbf{Corrente constante:} linearidade exata, mas exige fonte de
corrente estável e o sinal parte de zero apenas se $R\to0$.
\item \textbf{Divisor:} circuito trivial e barato, e a saída pode ser centrada
em meia alimentação --- conveniente para conversores A/D. Em contrapartida,
carrega o erro de linearidade e a deriva de $R_f$ e de $U_0$.
\end{itemize}
\textbf{Solução de compromisso muito usada:} a \emph{ponte} de Wheatstone, que
cancela a componente comum e entrega diretamente o desequilíbrio --- combinando
boa linearidade com circuito passivo.}
\end{questao}

\begin{questao}[4]{}
\textbf{(Resistividade e temperatura.)} A resistividade de um metal varia como
$\rho(T)=\rho_0\left[1+\alpha(T-T_0)\right]$.
\begin{itensq}
\item Mostre que a mesma lei vale para $R$, e identifique a hipótese oculta.
\item Determine o erro cometido ao ignorar a dilatação térmica das dimensões.
\item Avalie para cobre entre $\SI{20}{\celsius}$ e
      $\SI{120}{\celsius}$.
\end{itensq}
\resposta{A hipótese oculta é dimensões constantes; o erro é da ordem de
$\SI{0.05}{\percent}$ --- desprezível}
\resolucao{(a) De $R=\rho L/A$, se $L$ e $A$ forem constantes, então
$R\propto\rho$ e a lei se transfere diretamente. \textbf{A hipótese oculta é
justamente essa}: que a geometria não muda com a temperatura --- o que é falso,
pois todo metal se dilata.\\
(b) Com coeficiente de dilatação linear $\lambda$:
$L'=L(1+\lambda\Delta T)$ e $A'=A(1+\lambda\Delta T)^{2}$, logo
\[
\frac{L'}{A'}=\frac{L}{A}\cdot\frac{1}{1+\lambda\Delta T}.
\]
Assim
\[
R=\frac{\rho_0(1+\alpha\Delta T)}{1+\lambda\Delta T}\cdot\frac{L}{A}
\approx R_0\left[1+(\alpha-\lambda)\Delta T\right].
\]
\textbf{O efeito da dilatação é subtrair $\lambda$ de $\alpha$} --- a peça fica
mais longa (aumentando $R$) mas também mais grossa (reduzindo $R$), e o segundo
efeito predomina por ser quadrático na área.\\
(c) Para o cobre, $\alpha\approx\SI{3.9e-3}{\per\celsius}$ e
$\lambda\approx\SI{1.7e-5}{\per\celsius}$:
\[
\frac{\lambda}{\alpha}=\frac{1{,}7\times10^{-5}}{3{,}9\times10^{-3}}
=0{,}0044=0{,}44\%.
\]
Com $\Delta T=\SI{100}{\celsius}$, $R$ aumenta $39\%$ pela resistividade e
diminui $0{,}17\%$ pela dilatação --- efeito líquido de $38{,}8\%$.\\
\textbf{Conclusão:} ignorar a dilatação introduz erro de cerca de
$0{,}4\%$ \emph{do efeito térmico}, ou $0{,}17\%$ do valor de $R$. Desprezível
em quase toda aplicação --- mas não em padrões de resistência de laboratório,
onde se buscam partes por milhão.}
\end{questao}

\begin{questao}[4]{}
\textbf{(Medição a quatro fios.)} Quantifique a vantagem do método de Kelvin.
\begin{itensq}
\item Deduza o erro relativo do método de dois fios em função de
      $R_{cabo}/R_x$.
\item Determine a partir de qual valor de $R_x$ o método de dois fios dá erro
      inferior a $\SI{1}{\percent}$, com cabos de $\SI{0.2}{\ohm}$ cada.
\item Explique por que o método de quatro fios não elimina \emph{todos} os
      erros.
\end{itensq}
\resposta{Erro $=2R_{cabo}/R_x$; é preciso $R_x\ge\SI{40}{\ohm}$}
\resolucao{(a) O instrumento lê $R_x+2R_{cabo}$, logo
\[
\varepsilon=\frac{2R_{cabo}}{R_x}.
\]
(b) Impondo $\varepsilon\le0{,}01$ com $R_{cabo}=\SI{0.2}{\ohm}$:
\[
R_x\ge\frac{2(0{,}2)}{0{,}01}=\SI{40}{\ohm}.
\]
\textbf{Regra prática:} o método de dois fios só é aceitável para resistências
de dezenas de ohms ou mais. Abaixo disso, o erro cresce rapidamente --- em
$\SI{1}{\ohm}$ ele já é de $40\%$.\\
(c) \textbf{O que o método de quatro fios \emph{não} elimina.}
\begin{itemize}
\item \textbf{Tensões termoelétricas.} Junções entre metais diferentes nos
contatos geram forças eletromotrizes de dezenas de microvolts --- comparáveis ao
sinal quando se mede miliohms. Combate-se invertendo o sentido da corrente e
tomando a média das duas leituras.
\item \textbf{Autoaquecimento} pela corrente de excitação, que altera o próprio
$R_x$.
\item \textbf{Localização dos contatos de tensão.} Eles devem estar
\emph{entre} os de corrente, e a distribuição de corrente na peça influencia o
que se mede. Em amostras curtas, a definição do que é "o resistor" torna-se
ambígua.
\item \textbf{Ruído e deriva} do amplificador, agravados porque o sinal é de
poucos milivolts.
\end{itemize}
\textbf{Síntese:} o método de quatro fios elimina o erro \emph{sistemático e
dominante} --- a resistência dos cabos. Os erros remanescentes são de outra
natureza e exigem outras técnicas.}
\end{questao}

\begin{questao}[4]{}
\textbf{(Condutor composto.)} Um cabo de alta tensão tem alma de aço
($\rho_{\text{aço}}=\pot{1.5}{-7}\,\si{\ohm\meter}$, $A=\SI{20}{\milli\meter\squared}$)
envolvida por alumínio ($\rho_{Al}=\pot{2.8}{-8}\,\si{\ohm\meter}$,
$A=\SI{100}{\milli\meter\squared}$), com $\SI{1}{\kilo\meter}$ de comprimento.
\begin{itensq}
\item Determine a resistência de cada material e a do conjunto.
\item Determine a fração de corrente em cada um.
\item Explique a função da alma de aço.
\end{itensq}
\resposta{$R_{\text{aço}}=\SI{7.5}{\ohm}$, $R_{Al}=\SI{0.28}{\ohm}$,
$R_{eq}=\SI{0.270}{\ohm}$; o aço conduz $3{,}6\%$}
\resolucao{(a) Os dois materiais estão em \textbf{paralelo} (mesmo comprimento,
mesmos terminais):
\[
R_{\text{aço}}=\frac{\pot{1.5}{-7}(1000)}{\pot{20}{-6}}=\SI{7.5}{\ohm};
\qquad
R_{Al}=\frac{\pot{2.8}{-8}(1000)}{\pot{100}{-6}}=\SI{0.28}{\ohm};
\]
\[
R_{eq}=\frac{7{,}5(0{,}28)}{7{,}78}=\SI{0.270}{\ohm}.
\]
(b) Pelo divisor de corrente em condutâncias:
\[
\frac{I_{Al}}{I_T}=\frac{1/0{,}28}{1/0{,}28+1/7{,}5}=96{,}4\%;
\qquad
\frac{I_{\text{aço}}}{I_T}=3{,}6\%.
\]
(c) \textbf{A alma de aço praticamente não conduz --- e não é essa a sua
função.} Ela existe por razões \textbf{mecânicas}:
\begin{itemize}
\item O alumínio tem resistência à tração baixa e escoa sob carga permanente. O
aço sustenta o peso próprio do cabo entre torres, que podem distar centenas de
metros.
\item Reduz a flecha (o "barrigamento") do vão e, com ela, a altura necessária
das torres.
\item Suporta os esforços de vento e de acúmulo de gelo.
\end{itemize}
\textbf{Observe a divisão de papéis:} o alumínio conduz $96\%$ da corrente e o
aço sustenta praticamente toda a carga mecânica. É o princípio do cabo
\textbf{ACSR} (\emph{aluminium conductor steel-reinforced}), padrão em linhas de
transmissão aéreas --- um caso em que a análise elétrica revela que o componente
"elétrico" está ali por outro motivo.}
\end{questao}

\begin{questao}[4]{}
\textbf{(Limites da lei de Ohm.)} Discuta situações em que a relação
$U=RI$ deixa de descrever o comportamento de um condutor.
\begin{itensq}
\item Enumere ao menos quatro mecanismos, com exemplos.
\item Explique por que a lei de Ohm não é uma lei fundamental.
\item Indique como se procede quando ela falha.
\end{itensq}
\resposta{Autoaquecimento, não linearidade intrínseca, efeitos de alta
frequência e ruptura do meio}
\resolucao{(a) \textbf{Quatro mecanismos.}
\begin{enumerate}
\item \textbf{Autoaquecimento.} A resistência depende da temperatura, que
depende da potência dissipada. Exemplo: lâmpada incandescente, cuja resistência
a frio é cerca de dez vezes menor que a quente.
\item \textbf{Não linearidade intrínseca.} Junções semicondutoras seguem
relação exponencial, não linear. Exemplo: diodo, cuja "resistência" varia
ordens de grandeza em uma década de tensão.
\item \textbf{Efeitos de frequência.} O efeito pelicular concentra a corrente na
superfície do condutor, aumentando a resistência efetiva; somam-se indutância e
capacitância parasitas. Exemplo: um fio que tem $\SI{1}{\milli\ohm}$ em CC pode
apresentar dezenas de miliohms em megahertz.
\item \textbf{Ruptura do meio.} Acima da rigidez dielétrica, um isolante
torna-se condutor abruptamente. Exemplo: arco elétrico, cuja característica
$U\times I$ tem inclinação \emph{negativa}.
\end{enumerate}
\emph{(Acrescente-se ainda a saturação de velocidade dos portadores em campos
muito intensos e a supercondutividade, em que $R$ colapsa a zero.)}\\
(b) \textbf{A lei de Ohm é uma lei empírica e aproximada}, não fundamental. Ela
descreve o comportamento de \emph{certos} materiais (metais, em faixa limitada
de temperatura e frequência) e emerge de um modelo estatístico de colisões dos
portadores --- não de um princípio de conservação.\\
Compare com as leis de Kirchhoff, que decorrem da conservação de carga e de
energia e valem \emph{sempre}, para qualquer componente, linear ou não.\\
(c) \textbf{Quando a lei falha, procede-se assim:}
\begin{itemize}
\item Substituir $R$ pela \textbf{característica $U\times I$} completa do
dispositivo, obtida por medição ou por modelo.
\item Resolver o circuito por \textbf{interseção de características} (reta de
carga), como se faz com diodos.
\item Para pequenos sinais, linearizar em torno do ponto de operação usando a
\textbf{resistência dinâmica} --- e aí toda a maquinaria linear volta a valer.
\end{itemize}
\textbf{As leis de Kirchhoff continuam válidas em todos esses casos} --- é
sempre por elas que se começa.}
\end{questao}

% END BANCO COMPLEMENTAR

\gabarito[2.2 Lei de Ohm e resistividade]

% =====================================================================
\subsection{Potência elétrica e energia}
% =====================================================================

\blocofixacao

\begin{questao}[1]{}
A tensão nos terminais de um condutor é $\SI{12}{\volt}$ e ele dissipa
$\SI{24}{\watt}$.
\begin{itensq}
\item Qual é a corrente que o percorre?
\item Qual é o valor de sua resistência?
\end{itensq}
\resposta{(a) \SI{2}{\ampere}; (b) \SI{6}{\ohm}}
\resolucao{(a) De $P = U\,I$: $I = \dfrac{P}{U} = \dfrac{24}{12} = \SI{2}{\ampere}$.\\
(b) $R = \dfrac{U}{I} = \dfrac{12}{2} = \SI{6}{\ohm}$
(ou $R = \dfrac{U^2}{P} = \dfrac{144}{24} = \SI{6}{\ohm}$).}
\end{questao}

\begin{questao}[1]{}
Um aparelho opera em $\SI{220}{\volt}$ e consome $\SI{440}{\watt}$.
\begin{itensq}
\item Qual é a corrente que ele consome?
\item Qual sua resistência equivalente?
\end{itensq}
\resposta{(a) \SI{2}{\ampere}; (b) \SI{110}{\ohm}}
\resolucao{(a) $I = \dfrac{P}{U} = \dfrac{440}{220} = \SI{2}{\ampere}$.\\
(b) $R = \dfrac{U^2}{P} = \dfrac{220^2}{440} = \SI{110}{\ohm}$.}
\end{questao}

\begin{questao}[1]{Apostila original revisada}
Uma fonte fornece $\SI{100}{\watt}$ a um circuito, mas somente
$\SI{95}{\watt}$ são aproveitados pela carga. Determine o rendimento do circuito
e a potência dissipada em perdas.
\resposta{$\eta=95\%$; perdas de $\SI{5}{\watt}$}
\resolucao{$\eta=P_{\text{útil}}/P_{\mathrm{entrada}}=95/100=0{,}95$.
A diferença, $100-95=\SI{5}{\watt}$, é dissipada em fios, contatos e outros elementos.}
\end{questao}

\blocoaplicacao

\begin{questao}[2]{}
Uma lâmpada de especificação $\SI{60}{\watt}$/$\SI{220}{\volt}$ é ligada, por
engano, a uma rede de $\SI{110}{\volt}$. Sua potência dissipada passa a ser:
\begin{alternativas}
\item inalterada.
\item reduzida à metade.
\item duplicada.
\item reduzida à quarta parte.
\item aumentada quatro vezes.
\end{alternativas}
\resposta{(d)}
\resolucao{A resistência da lâmpada é fixa, determinada por seus valores nominais:
$R = \dfrac{U_{\text{nom}}^2}{P_{\text{nom}}} = \dfrac{220^2}{60} \approx
\SI{807}{\ohm}$. Sob a nova tensão, $P' = \dfrac{U'^2}{R}$. Como a tensão caiu à
metade e $P \propto U^2$:
\[
P' = \frac{(110)^2}{R} = \frac{60}{4} = \SI{15}{\watt},
\]
ou seja, a quarta parte. Meia tensão dissipa um quarto da potência.}
\end{questao}

\begin{questao}[2]{}
Um chuveiro elétrico tem especificações $\SI{220}{\volt}$/$\SI{4400}{\watt}$. Se
for ligado a uma rede de $\SI{110}{\volt}$, a potência dissipada será:
\begin{alternativas}[2]
\item \SI{550}{\watt}
\item \SI{1100}{\watt}
\item \SI{2200}{\watt}
\item \SI{4400}{\watt}
\item \SI{8800}{\watt}
\end{alternativas}
\resposta{(b)}
\resolucao{Mesma ideia da questão anterior: metade da tensão dissipa um quarto da
potência:
\[
P' = \frac{4400}{4} = \SI{1100}{\watt}.
\]
Por isso um chuveiro projetado para \SI{220}{\volt} "esquenta pouco" em rede de
\SI{110}{\volt}.}
\end{questao}

\begin{questao}[2]{}
Dois resistores $R_1 = \SI{10}{\ohm}$ e $R_2 = \SI{20}{\ohm}$ estão em série sob
uma tensão total de $\SI{60}{\volt}$.
\begin{itensq}
\item Qual é a corrente do circuito?
\item Qual é a potência dissipada em cada resistor?
\end{itensq}
\resposta{(a) \SI{2}{\ampere}; (b) $P_1=\SI{40}{\watt}$, $P_2=\SI{80}{\watt}$}
\resolucao{(a) $I = \dfrac{60}{10+20} = \SI{2}{\ampere}$.\\
(b) Com a mesma corrente, $P = R\,I^2$:
\[
P_1 = 10\cdot 2^2 = \SI{40}{\watt},\qquad
P_2 = 20\cdot 2^2 = \SI{80}{\watt}.
\]
\textbf{Verificação:} $P_{\text{total}} = U\,I = 60\cdot2 = \SI{120}{\watt}
= 40+80$. Na série, o maior resistor dissipa a maior potência.}
\end{questao}

\begin{questao}[2]{}
Um chuveiro elétrico de $\SI{5500}{\watt}$ opera em $\SI{220}{\volt}$.
\begin{itensq}
\item Qual é a corrente que ele exige da rede?
\item Qual é a resistência da sua resistência de aquecimento?
\item O disjuntor do circuito é de $\SI{20}{\ampere}$. Ele suporta esse chuveiro?
\end{itensq}
\resposta{(a) \SI{25}{\ampere}; (b) \SI{8.8}{\ohm}; (c) não}
\resolucao{(a) $I = \dfrac{P}{U} = \dfrac{5500}{220} = \SI{25}{\ampere}$.\\
(b) $R = \dfrac{U^2}{P} = \dfrac{220^2}{5500} = \SI{8.8}{\ohm}$.\\
(c) A corrente exigida (\SI{25}{\ampere}) \emph{excede} os \SI{20}{\ampere} do
disjuntor, que desarmaria. Seria necessário um disjuntor de pelo menos
\SI{25}{\ampere} (na prática, \SI{32}{\ampere}) e fiação compatível --- por isso
chuveiros exigem circuito exclusivo.}
\end{questao}

\begin{questao}[2]{Apostila original revisada}
Em uma residência, a tarifa média foi obtida de uma conta de
R\$~80,00 para um consumo de $\SI{200}{\kilo\watt\hour}$. Um chuveiro de
$\SI{3}{\kilo\watt}$ é utilizado durante $\SI{30}{\minute}$ por dia, ao longo de
30 dias. Calcule a energia mensal consumida pelo chuveiro e seu custo.
\resposta{$\SI{45}{\kilo\watt\hour}$; R\$~18,00}
\resolucao{A tarifa é $80/200=\text{R\$ }0{,}40/\si{\kilo\watt\hour}$.
O tempo mensal é $0{,}5\times30=\SI{15}{\hour}$, então
$E=3\times15=\SI{45}{\kilo\watt\hour}$ e o custo é $45\times0{,}40=\text{R\$ }18{,}00$.}
\end{questao}

\blocoaprofundamento

\begin{questao}[3]{}
Em uma residência permanecem ligados, durante $\SI{2}{\hour}$, um ferro elétrico
de $\SI{1440}{\watt}$/$\SI{120}{\volt}$ e duas lâmpadas de
$\SI{60}{\watt}$/$\SI{120}{\volt}$. Mantida a tensão em $\SI{120}{\volt}$,
determine a corrente total no fusível e a energia elétrica consumida.
\resposta{$I = \SI{13}{\ampere}$; $E = \SI{3.12}{\kilo\watt\hour}$}
\resolucao{Potência total: $P = 1440 + 2\cdot60 = \SI{1560}{\watt}$.\\
Corrente no fusível (todos em paralelo, sob \SI{120}{\volt}):
\[
I = \frac{P}{U} = \frac{1560}{120} = \SI{13}{\ampere}.
\]
Energia em $\SI{2}{\hour}$:
\[
E = P\,t = 1{,}560\,\si{\kilo\watt}\cdot2\,\si{\hour} = \SI{3.12}{\kilo\watt\hour}.
\]
O kWh é a unidade da conta de luz: potência (em kW) vezes tempo (em h).}
\end{questao}

\begin{questao}[3]{}
Uma residência utiliza, por dia: uma lâmpada de $\SI{100}{\watt}$ acesa
$\SI{5}{\hour}$, um chuveiro de $\SI{4400}{\watt}$ ligado $\SI{0.5}{\hour}$ e uma
geladeira de $\SI{200}{\watt}$ que funciona, em média, $\SI{8}{\hour}$ por dia.
Calcule o consumo mensal de energia (30 dias) e o custo, considerando a tarifa de
$\SI{0.80}{\text{R\$}\per\kilo\watt\hour}$.
\resposta{$E \approx \SI{127.5}{\kilo\watt\hour}$; custo $\approx$ R\$\,102,00}
\resolucao{Energia diária de cada aparelho ($E = P\,t$, em kWh):
\[
\text{lâmpada: } 0{,}1\cdot5 = \SI{0.5}{\kilo\watt\hour};\quad
\text{chuveiro: } 4{,}4\cdot0{,}5 = \SI{2.2}{\kilo\watt\hour};\quad
\text{geladeira: } 0{,}2\cdot8 = \SI{1.6}{\kilo\watt\hour}.
\]
Total diário: $0{,}5 + 2{,}2 + 1{,}6 = \SI{4.3}{\kilo\watt\hour}$. Em 30 dias:
\[
E = 4{,}3\cdot30 = \SI{129}{\kilo\watt\hour}
\quad\Rightarrow\quad
\text{custo} = 129\cdot0{,}80 \approx \text{R\$}\,103{,}00.
\]
(Valores próximos de \SI{127.5}{\kilo\watt\hour} conforme arredondamentos.) É
assim que se estima a conta de luz --- e fica claro que o chuveiro, apesar de
usado poucos minutos, é um dos maiores vilões pelo seu alto \emph{potência}.}
\end{questao}

\begin{questao}[3]{}
Duas lâmpadas incandescentes de especificações $\SI{60}{\watt}$/$\SI{220}{\volt}$
e $\SI{100}{\watt}$/$\SI{220}{\volt}$ são ligadas \textbf{em série} à rede de
$\SI{220}{\volt}$.
\begin{itensq}
\item Calcule a resistência de cada lâmpada (suposta constante).
\item Determine a potência dissipada por cada uma nessa ligação.
\item \textbf{Qual delas brilha mais?} O resultado contraria a intuição?
      Explique.
\end{itensq}
\resposta{(a) $\SI{806.7}{\ohm}$ e $\SI{484}{\ohm}$;
(b) $\SI{23.4}{\watt}$ e $\SI{14.1}{\watt}$; (c) a de \SI{60}{\watt}}
\resolucao{\textbf{(a)} Das especificações nominais, $R = \dfrac{U^2}{P}$:
\[
R_{60}=\frac{220^2}{60}\approx\SI{806.7}{\ohm},
\qquad
R_{100}=\frac{220^2}{100}=\SI{484}{\ohm}.
\]
Note: \emph{a lâmpada de menor potência tem MAIOR resistência.}

\textbf{(b)} Em série, a corrente é comum:
\[
I=\frac{220}{806{,}7+484}=\frac{220}{1290{,}7}\approx\SI{0.170}{\ampere}.
\]
Com $P = R I^2$:
\[
P_{60}=806{,}7\,(0{,}170)^2\approx\SI{23.4}{\watt},
\qquad
P_{100}=484\,(0{,}170)^2\approx\SI{14.1}{\watt}.
\]
\textbf{Verificação:} $23{,}4+14{,}1 = \SI{37.5}{\watt} = 220\cdot0{,}170$.
\checkmark

\textbf{(c) A lâmpada de \SI{60}{\watt} brilha MAIS} --- resultado que contraria
a intuição de quem associa "100 W" a "mais brilhante".\\
\textbf{Por quê:} as especificações valem \emph{apenas} na tensão nominal, com
cada lâmpada recebendo os \SI{220}{\volt} completos. Em série, elas
\textbf{compartilham} a mesma corrente, e quem dita a potência passa a ser a
resistência: $P = RI^2$. A de \SI{60}{\watt}, tendo maior $R$, fica com maior
parcela da tensão e dissipa mais.\\
\textbf{Regra geral:}
\begin{center}\small\sffamily
\begin{tabular}{@{}l l l@{}}
\toprule
\textbf{Ligação} & \textbf{Grandeza comum} & \textbf{Brilha mais}\\
\midrule
Série    & corrente $I$ & a de \emph{maior} $R$ (menor potência nominal)\\
Paralelo & tensão $U$   & a de \emph{menor} $R$ (maior potência nominal)\\
\bottomrule
\end{tabular}
\end{center}
Observe ainda que \emph{ambas} ficam muito abaixo do nominal (\SI{37.5}{\watt} no
total contra \SI{160}{\watt} se ligadas corretamente em paralelo).}
\end{questao}

\begin{questao}[3]{}
Um chuveiro elétrico de $\SI{220}{\volt}$ possui uma resistência de aquecimento
com uma \emph{derivação central}, permitindo duas posições: na posição "verão",
usa-se a resistência inteira ($\SI{8.8}{\ohm}$); na posição "inverno", apenas
\emph{metade} dela.
\begin{itensq}
\item Calcule a potência em cada posição.
\item Determine a corrente em cada caso e verifique se um disjuntor de
      $\SI{40}{\ampere}$ é adequado.
\item Sabendo que a água entra a $\SI{20}{\celsius}$ com vazão de
      $\SI{3}{\liter\per\minute}$, estime a temperatura de saída no inverno.
\end{itensq}
\dados{$c_{\text{água}}=\SI{4180}{\joule\per\kilogram\per\celsius}$;
$\rho_{\text{água}}=\SI{1}{\kilogram\per\liter}$.}
\resposta{(a) \SI{5500}{\watt} e \SI{11000}{\watt}; (b) \SI{25}{\ampere} e
\SI{50}{\ampere} --- disjuntor insuficiente; (c) $\approx\SI{72}{\celsius}$}
\resolucao{\textbf{(a)} Com $P = \dfrac{U^2}{R}$:
\[
P_{\text{verão}}=\frac{220^2}{8{,}8}=\SI{5500}{\watt},
\qquad
P_{\text{inverno}}=\frac{220^2}{4{,}4}=\SI{11000}{\watt}.
\]
\emph{Metade da resistência dobra a potência} --- e não a reduz, como a intuição
poderia sugerir.

\textbf{(b)} $I = \dfrac{P}{U}$:
$I_{\text{verão}} = \SI{25}{\ampere}$ e $I_{\text{inverno}} = \SI{50}{\ampere}$.
O disjuntor de \SI{40}{\ampere} \textbf{não é adequado} --- ele desarmaria na
posição inverno. Seria necessário um de \SI{63}{\ampere} com fiação
de $\SI{10}{\milli\meter\squared}$.

\textbf{(c) Balanço térmico.} A vazão mássica é
\[
\dot m = \SI{3}{\liter\per\minute} = \frac{3}{60}=\SI{0.05}{\kilogram\per\second}.
\]
Toda a potência elétrica vira calor na água:
\[
P = \dot m\,c\,\Delta T
\;\Rightarrow\;
\Delta T=\frac{P}{\dot m\,c}=\frac{11000}{0{,}05\cdot4180}\approx\SI{52.6}{\celsius}.
\]
Logo a saída fica a $20+52{,}6 \approx \SI{72}{\celsius}$ --- alta demais para
banho, o que na prática é ajustado aumentando a vazão.\\
\textbf{Síntese de engenharia:} este exercício conecta \emph{três} domínios ---
circuito elétrico (lei de Ohm), proteção (dimensionamento de disjuntor) e
termodinâmica (balanço de energia). É exatamente esse tipo de integração que
caracteriza o trabalho profissional real.}
\end{questao}

% BEGIN BANCO COMPLEMENTAR Potência elétrica e energia
% =====================================================================
%  Banco complementar --- Potencia Eletrica e Energia  (REVISADO)
%  Substitui a versao gerada por template (7 clones em N2, 8 em N3 e
%  8 questoes conceituais marcadas como N4).
%  O conteudo de ENERGIA esta no cap02, que ja cobre: rendimento
%  100 W/95 W, tarifa media a partir de conta de R\$ 80 / 200 kWh,
%  ferro eletrico e outros aparelhos por 2 h, e o consumo diario
%  residencial (lampada, chuveiro etc.). O cap05 cobre P=RI^2, P=VI
%  e maxima transferencia. O banco de circuitos serie revisado ja usa
%  um calculo de custo mensal (300 W, 6 h/dia, 30 dias).
%  Este banco explora: potencia MEDIA x valor eficaz, regime pulsado
%  e constante termica, rendimento em cascata, tensao errada de
%  alimentacao, energia em transitorio exponencial, payback e o
%  dimensionamento economico de condutor.
%  Distribuicao mantida: 5 N1 + 7 N2 + 8 N3 + 8 N4 = 28
% =====================================================================

\subsubsection*{Banco complementar --- Potência elétrica e energia}

\begin{atencao}[Organização do banco]
As três formas de $P$ ($UI$, $RI^{2}$ e $U^{2}/R$) são equivalentes apenas
quando vale a lei de Ohm; escolha a que usa os dados disponíveis sem passos
intermediários. Para \emph{energia}, lembre que $\SI{1}{\kilo\watt\hour}
=\SI{3.6}{\mega\joule}$. O N3 trata de regime pulsado, rendimento em cascata e
tensão incorreta de alimentação; o N4 exige integração, demonstração de que a
potência média usa o valor \emph{eficaz}, e análises de retorno econômico.
\end{atencao}

\blocofixacao

\begin{questao}[1]{}
Um resistor de $\SI{8}{\ohm}$ é percorrido por $\SI{2.5}{\ampere}$. Determine a
potência dissipada.
\resposta{$P=\SI{50}{\watt}$}
\resolucao{
\[
P=RI^{2}=8(2{,}5)^{2}=8(6{,}25)=\SI{50}{\watt}.
\]
\textbf{Escolha da fórmula:} conhecendo $R$ e $I$, use $RI^{2}$ diretamente.
Calcular $U=RI=\SI{20}{\volt}$ e depois $P=UI$ dá o mesmo resultado, mas
acrescenta um passo e uma oportunidade de erro.}
\end{questao}

\begin{questao}[1]{}
Um resistor de $\SI{44}{\ohm}$ é submetido a $\SI{220}{\volt}$. Determine a
potência.
\resposta{$P=\SI{1100}{\watt}$}
\resolucao{
\[
P=\frac{U^{2}}{R}=\frac{220^{2}}{44}=\frac{48400}{44}=\SI{1100}{\watt}.
\]
Note a dependência \emph{quadrática} na tensão: uma queda de $10\%$ na rede
($\SI{198}{\volt}$) reduziria a potência para
$\SI{891}{\watt}$ --- $19\%$ a menos. É por isso que chuveiros "esquentam
menos" em horários de pico.}
\end{questao}

\begin{questao}[1]{}
Converta $\SI{250}{\watt\hour}$ para joules.
\resposta{$\SI{900}{\kilo\joule}$}
\resolucao{
\[
E=Pt=250\,\si{\watt}\times3600\,\si{\second}=\SI{900000}{\joule}
=\SI{900}{\kilo\joule}.
\]
\textbf{Fator de conversão útil:}
$\SI{1}{\kilo\watt\hour}=\SI{3.6}{\mega\joule}$. O joule é pequeno demais para
contas de energia elétrica --- daí o uso do quilowatt-hora, que é uma unidade
de energia, não de potência. Confundir os dois é o erro mais comum do tópico.}
\end{questao}

\begin{questao}[1]{}
Conhecendo apenas a tensão sobre um resistor e o valor de sua resistência, a
expressão adequada para a potência é:
\begin{alternativas}[2]
\item $P=UI$
\item $P=RI^{2}$
\item $P=U^{2}/R$
\item $P=U/R$
\end{alternativas}
\resposta{(c)}
\resolucao{As alternativas (a) e (b) exigiriam calcular $I$ antes. A (d) está
dimensionalmente errada: $\si{\volt}/\si{\ohm}$ é ampère, não watt.\\
\textbf{Verificação dimensional rápida:} $U^{2}/R$ dá
$\si{\volt\squared}/\si{\ohm}=\si{\volt}\cdot\si{\ampere}=\si{\watt}$
\checkmark. Sempre que houver dúvida entre fórmulas parecidas, confira a
dimensão --- ela elimina a maioria das alternativas erradas.}
\end{questao}

\begin{questao}[1]{}
Um dispositivo de $\SI{60}{\watt}$ opera por $\SI{5}{\minute}$. Determine a
energia consumida em joules.
\resposta{$E=\SI{18}{\kilo\joule}$}
\resolucao{Convertendo o tempo para segundos:
\[
E=Pt=60\times(5\times60)=60(300)=\SI{18000}{\joule}=\SI{18}{\kilo\joule}.
\]
\textbf{Cuidado com as unidades:} $P$ em watts exige $t$ em \emph{segundos}
para obter joules. Se o tempo ficar em minutos, o resultado sai
$60\times$ menor --- erro silencioso, porque o número parece plausível.}
\end{questao}

\blocoaplicacao

\begin{questao}[2]{}
Um chuveiro de $\SI{5500}{\watt}$ opera em $\SI{220}{\volt}$. Determine a
corrente e a resistência.
\resposta{$I=\SI{25}{\ampere}$; $R=\SI{8.8}{\ohm}$}
\resolucao{
\[
I=\frac{P}{U}=\frac{5500}{220}=\SI{25}{\ampere};
\qquad
R=\frac{U^{2}}{P}=\frac{48400}{5500}=\SI{8.8}{\ohm}.
\]
\textbf{Conferência:} $R=U/I=220/25=\SI{8.8}{\ohm}$ \checkmark.\\
\textbf{Implicação de instalação:} $\SI{25}{\ampere}$ é uma corrente elevada.
Ela exige circuito exclusivo, condutor de seção adequada e disjuntor de
$\SI{32}{\ampere}$ --- e é por isso que chuveiro nunca compartilha circuito com
tomadas.}
\end{questao}

\begin{questao}[2]{}
Projete a resistência de um aquecedor de $\SI{1200}{\watt}$ para
$\SI{220}{\volt}$ e determine a corrente nominal e o disjuntor adequado.
\resposta{$R=\SI{40.3}{\ohm}$; $I=\SI{5.45}{\ampere}$; disjuntor de
$\SI{10}{\ampere}$}
\resolucao{
\[
R=\frac{U^{2}}{P}=\frac{48400}{1200}=\SI{40.3}{\ohm};
\qquad
I=\frac{1200}{220}=\SI{5.45}{\ampere}.
\]
\textbf{Disjuntor:} o valor comercial imediatamente acima é
$\SI{10}{\ampere}$, que oferece folga de $83\%$ --- adequado, pois aquecedores
resistivos não têm corrente de partida elevada.\\
\textbf{Nota sobre o material:} $\SI{40.3}{\ohm}$ deve ser obtido com fio de
liga resistiva (níquel-cromo) de comprimento e seção calculados. E como a
resistência desses materiais varia pouco com a temperatura, a potência
permanece próxima do nominal em regime.}
\end{questao}

\begin{questao}[2]{}
Duas lâmpadas incandescentes de $\SI{60}{\watt}$ são fabricadas para
$\SI{127}{\volt}$ e para $\SI{220}{\volt}$. Determine a resistência de cada
uma.
\resposta{$\SI{268.8}{\ohm}$ e $\SI{806.7}{\ohm}$}
\resolucao{
\[
R_{127}=\frac{127^{2}}{60}=\SI{268.8}{\ohm};
\qquad
R_{220}=\frac{220^{2}}{60}=\SI{806.7}{\ohm}.
\]
A lâmpada de tensão maior tem resistência \textbf{três vezes} maior --- razão
$(220/127)^{2}=3{,}0$.\\
\textbf{Consequência construtiva:} para a mesma potência, o filamento de
$\SI{220}{\volt}$ é mais longo e mais fino. Ele é, portanto, mais frágil
mecanicamente --- fato que se observa na prática, com lâmpadas de tensão maior
tendendo a durar menos sob vibração.}
\end{questao}

\begin{questao}[2]{}
Um aparelho opera em dois modos: $\SI{1500}{\watt}$ durante
$\SI{20}{\minute}$ e $\SI{600}{\watt}$ durante $\SI{40}{\minute}$, todos os
dias. Determine o consumo mensal ($\SI{30}{\day}$) e o custo a
$\mathrm{R\$}\,0{,}95$ por $\si{\kilo\watt\hour}$.
\resposta{$\SI{27}{\kilo\watt\hour}$; $\mathrm{R\$}\,25{,}65$}
\resolucao{Energia diária, com os tempos em horas:
\[
E=1{,}5\left(\tfrac13\right)+0{,}6\left(\tfrac23\right)
=0{,}50+0{,}40=\SI{0.90}{\kilo\watt\hour}.
\]
\[
E_{\text{mês}}=0{,}90(30)=\SI{27}{\kilo\watt\hour};
\qquad
\text{custo}=27(0{,}95)=\mathrm{R\$}\,25{,}65.
\]
\textbf{Observação:} o modo de $\SI{600}{\watt}$, embora tenha menos da metade
da potência, responde por $44\%$ do consumo --- porque funciona o dobro do
tempo. Em energia, \textbf{potência sozinha não decide}; o produto pelo tempo é
que conta.}
\end{questao}

\begin{questao}[2]{}
Um resistor é submetido a pulsos de $\SI{50}{\watt}$ com ciclo de trabalho de
$\SI{10}{\percent}$. Determine a potência média.
\resposta{$P_{\text{média}}=\SI{5}{\watt}$}
\resolucao{
\[
P_{\text{média}}=P_{pico}\times D=50(0{,}10)=\SI{5}{\watt}.
\]
\textbf{Mas cuidado:} isso \emph{não} significa que um resistor de
$\SI{5}{\watt}$ sirva. Se a duração do pulso for comparável ou maior que a
constante térmica do componente, ele atinge temperaturas próximas às de
$\SI{50}{\watt}$ contínuos e pode falhar.\\
A questão de quando a potência média basta é retomada quantitativamente no bloco
de aprofundamento.}
\end{questao}

\begin{questao}[2]{}
Uma fonte com rendimento de $\SI{90}{\percent}$ alimenta um conversor de
$\SI{85}{\percent}$. Determine o rendimento global e a potência de entrada
necessária para entregar $\SI{100}{\watt}$ úteis.
\resposta{$\eta=76{,}5\%$; $P_{entrada}=\SI{130.7}{\watt}$}
\resolucao{Rendimentos em cascata \textbf{multiplicam-se}:
\[
\eta=0{,}90\times0{,}85=0{,}765=76{,}5\%.
\]
\[
P_{entrada}=\frac{100}{0{,}765}=\SI{130.7}{\watt}.
\]
\textbf{Erro comum:} \emph{somar} ou fazer a média dos rendimentos, o que daria
$87{,}5\%$ --- valor otimista demais. Cada estágio desperdiça uma fração do que
recebe, e as perdas se compõem multiplicativamente.\\
\textbf{Perdas:} $\SI{13.1}{\watt}$ na fonte e $\SI{17.6}{\watt}$ no
conversor, totalizando $\SI{30.7}{\watt}$ dissipados.}
\end{questao}

\begin{questao}[2]{}
Uma carga alterna entre $\SI{500}{\watt}$ por $\SI{5}{\minute}$ e
$\SI{100}{\watt}$ por $\SI{15}{\minute}$, em ciclos de
$\SI{20}{\minute}$. Determine a potência média e a energia por ciclo.
\resposta{$P_{\text{média}}=\SI{200}{\watt}$; $E=\SI{66.7}{\watt\hour}$}
\resolucao{A potência média é a \textbf{média ponderada pelo tempo}:
\[
P_{\text{média}}=\frac{500(5)+100(15)}{20}=\frac{2500+1500}{20}=\SI{200}{\watt}.
\]
\[
E_{ciclo}=P_{\text{média}}\times t_{ciclo}=200\times\frac{20}{60}
=\SI{66.7}{\watt\hour}.
\]
\textbf{Atenção:} a média \emph{aritmética} das potências ($300\,\si{\watt}$)
estaria errada --- ela ignoraria que o patamar baixo dura três vezes mais.
Média de potência é sempre ponderada pelo tempo.}
\end{questao}

\blocoaprofundamento

\begin{questao}[3]{}
Um resistor de $\SI{10}{\ohm}$ conduz $\SI{4}{\ampere}$ durante
$\SI{25}{\percent}$ do tempo e nada nos $\SI{75}{\percent}$ restantes.
\begin{itensq}
\item Determine a corrente média e a potência calculada a partir dela.
\item Determine a potência média verdadeira.
\item Explique a discrepância.
\end{itensq}
\resposta{(a) $\SI{1}{\ampere}$ e $\SI{10}{\watt}$ (errado); (b)
$\SI{40}{\watt}$; (c) a potência depende de $\langle i^{2}\rangle$}
\resolucao{(a) $I_{\text{média}}=0{,}25(4)+0{,}75(0)=\SI{1}{\ampere}$, e
$RI_{\text{média}}^{2}=10(1)=\SI{10}{\watt}$.\\
(b) A potência instantânea vale $\SI{160}{\watt}$ durante $25\%$ do tempo e
zero no resto:
\[
P_{\text{média}}=0{,}25(160)+0{,}75(0)=\SI{40}{\watt}.
\]
\textbf{Quatro vezes} o valor obtido em (a).\\
(c) A potência é \emph{quadrática} na corrente, e a média de um quadrado não é
o quadrado da média:
\[
P_{\text{média}}=R\left\langle i^{2}\right\rangle
\ne R\left\langle i\right\rangle^{2}.
\]
Define-se então o \textbf{valor eficaz}
$I_{rms}=\sqrt{\langle i^{2}\rangle}=\sqrt{0{,}25(16)}=\SI{2}{\ampere}$, com o
qual $P=RI_{rms}^{2}=\SI{40}{\watt}$ \checkmark.\\
\textbf{Definição operacional do valor eficaz:} é a corrente contínua que
produziria o \emph{mesmo aquecimento}. Aqui, $\SI{2}{\ampere}$ contínuos
aqueceriam tanto quanto o regime pulsado --- e é por isso que multímetros de
qualidade medem "true RMS".}
\end{questao}

\begin{questao}[3]{}
Um resistor com potência nominal de $\SI{10}{\watt}$ recebe pulsos de
$\SI{50}{\watt}$ com ciclo de trabalho de $\SI{10}{\percent}$.
\begin{itensq}
\item Determine a potência média e compare com o nominal.
\item Estabeleça o critério para decidir se o componente sobrevive.
\item Analise dois casos: pulsos de $\SI{1}{\milli\second}$ e de
      $\SI{10}{\second}$.
\end{itensq}
\resposta{$P_{\text{média}}=\SI{5}{\watt}$; a decisão depende da constante térmica}
\resolucao{(a) $P_{\text{média}}=\SI{5}{\watt}$, metade do nominal --- aparentemente
seguro.\\
(b) \textbf{O critério é a comparação entre a duração do pulso e a constante
térmica $\tau_t$ do componente} (tipicamente de segundos a dezenas de segundos
para resistores de fio; menos para os de filme).
\begin{itemize}
\item Se $t_{pulso}\ll\tau_t$, a massa térmica \emph{integra} os pulsos: a
temperatura acompanha a potência \textbf{média}, e o critério de
$\SI{5}{\watt}$ vale.
\item Se $t_{pulso}\gg\tau_t$, o resistor atinge o equilíbrio térmico
\emph{durante} cada pulso: a temperatura corresponde a $\SI{50}{\watt}$
contínuos, e o componente queima.
\end{itemize}
(c) \textbf{Pulsos de $\SI{1}{\milli\second}$:} muito menores que qualquer
$\tau_t$ realista --- o resistor "enxerga" $\SI{5}{\watt}$. \textbf{Seguro.}\\
\textbf{Pulsos de $\SI{10}{\second}$:} comparáveis ou maiores que $\tau_t$ --- o
resistor esquenta como se recebesse $\SI{50}{\watt}$ por dez segundos.
\textbf{Falha provável.}\\
\textbf{Consequência prática:} fabricantes publicam curvas de \emph{potência de
pulso} justamente porque o valor nominal contínuo não responde a essa pergunta.
Nunca dimensione regime pulsado apenas pela média.}
\end{questao}

\begin{questao}[3]{}
Um sistema tem uma fonte de $\SI{92}{\percent}$ em cascata com um conversor de
$\SI{85}{\percent}$.
\begin{itensq}
\item Determine o rendimento global.
\item Compare o ganho de melhorar o estágio pior (de $85$ para
      $\SI{92}{\percent}$) com o de melhorar o melhor (de $92$ para
      $\SI{97}{\percent}$).
\item Formule a regra geral.
\end{itensq}
\resposta{(a) $78{,}2\%$; (b) $84{,}6\%$ contra $82{,}5\%$ --- melhorar o pior
rende mais}
\resolucao{(a) $\eta=0{,}92(0{,}85)=0{,}782=78{,}2\%$.\\
(b) \textbf{Melhorando o pior:} $0{,}92(0{,}92)=84{,}6\%$ --- ganho de
$6{,}4$ pontos.\\
\textbf{Melhorando o melhor:} $0{,}97(0{,}85)=82{,}5\%$ --- ganho de
$4{,}3$ pontos.\\
E note que o segundo caso exigiu uma melhoria maior em pontos absolutos
($5$ contra $7$) para render \emph{menos}.\\
(c) \textbf{Regra geral.} Como $\eta=\prod\eta_k$, o ganho relativo global é a
soma dos ganhos relativos:
\[
\frac{\Delta\eta}{\eta}=\sum_k\frac{\Delta\eta_k}{\eta_k}.
\]
Um mesmo acréscimo \emph{absoluto} rende mais no estágio de menor $\eta_k$,
pois ali representa maior acréscimo relativo. \textbf{Ataque sempre o elo mais
fraco} --- e é também nele que a margem de melhoria costuma ser maior e mais
barata.}
\end{questao}

\begin{questao}[3]{}
Um aparelho consome $\SI{4}{\watt}$ em espera, permanentemente.
\begin{itensq}
\item Determine o consumo anual e o custo a
      $\mathrm{R\$}\,0{,}95/\si{\kilo\watt\hour}$.
\item Estime o impacto de dez aparelhos semelhantes.
\item Proponha uma decisão de engenharia.
\end{itensq}
\resposta{$\SI{35}{\kilo\watt\hour}$ e $\mathrm{R\$}\,33{,}29$ por aparelho;
$\mathrm{R\$}\,333$ para dez}
\resolucao{(a) Um ano tem $\SI{8760}{\hour}$:
\[
E=4\times8760=\SI{35040}{\watt\hour}=\SI{35.04}{\kilo\watt\hour};
\]
\[
\text{custo}=35{,}04(0{,}95)=\mathrm{R\$}\,33{,}29.
\]
(b) Dez aparelhos: $\SI{350}{\kilo\watt\hour}$ e cerca de
$\mathrm{R\$}\,333$ por ano --- comparável ao consumo anual de uma geladeira
eficiente.\\
(c) \textbf{Decisões possíveis, com seus custos.}
\begin{itemize}
\item \textbf{Filtro com chave} para desligar grupos de aparelhos: custo de
dezenas de reais, retorno em poucos meses. Melhor relação custo-benefício.
\item \textbf{Redesenho do produto:} normas modernas exigem espera abaixo de
$\SI{0.5}{\watt}$, o que reduziria o custo a menos de
$\mathrm{R\$}\,4{,}20$ por ano.
\item \textbf{Não fazer nada:} defensável para um aparelho isolado, insustentável
em escala --- o consumo agregado de espera responde por vários por cento do
consumo residencial nacional.
\end{itemize}
\textbf{O ponto de engenharia:} $\SI{4}{\watt}$ parecem irrelevantes porque a
potência é pequena. Mas energia é potência \emph{vezes tempo}, e
$\SI{8760}{\hour}$ é um multiplicador enorme.}
\end{questao}

\begin{questao}[3]{}
Um aparelho projetado para $\SI{220}{\volt}$ e $\SI{1000}{\watt}$ é ligado por
engano em $\SI{127}{\volt}$.
\begin{itensq}
\item Determine a potência resultante, supondo resistência constante.
\item Analise o caso inverso: aparelho de $\SI{127}{\volt}$ em
      $\SI{220}{\volt}$.
\end{itensq}
\resposta{(a) $\SI{333}{\watt}$ ($33\%$); (b) potência \emph{triplicada} ---
destruição}
\resolucao{(a) $R=\dfrac{220^{2}}{1000}=\SI{48.4}{\ohm}$, e em
$\SI{127}{\volt}$:
\[
P=\frac{127^{2}}{48{,}4}=\SI{333}{\watt}
=\left(\frac{127}{220}\right)^{2}\times1000=33{,}3\%.
\]
\textbf{Consequência:} o aparelho funciona, mas com um terço da potência --- um
chuveiro esquenta pouco, um aquecedor não aquece. Situação incômoda, porém
\emph{segura}.\\
(b) $P'=\left(\dfrac{220}{127}\right)^{2}P=3{,}0\,P$ --- a potência
\textbf{triplica}, e a corrente quase dobra.\\
\textbf{Resultado:} superaquecimento imediato, fumaça e, muitas vezes, incêndio.
O disjuntor pode nem atuar, pois a corrente ($\SI{13.6}{\ampere}$ para um
aparelho de $\SI{1000}{\watt}$/$\SI{127}{\volt}$) fica abaixo do valor de
disparo --- a proteção existe para o \emph{circuito}, não para o aparelho.\\
\textbf{Assimetria fundamental:} subtensão é inconveniente; sobretensão é
destrutiva. Por isso a etiqueta de tensão merece conferência antes de ligar
qualquer aparelho resistivo.}
\end{questao}

\begin{questao}[3]{}
Um resistor de $\SI{4}{\ohm}$ é percorrido por
$i(t)=\SI{5}{\ampere}\cdot e^{-t/\tau}$, com $\tau=\SI{20}{\milli\second}$.
\begin{itensq}
\item Determine a potência inicial.
\item Determine a energia total dissipada.
\item Determine o instante em que metade da energia já foi dissipada.
\end{itensq}
\resposta{(a) $\SI{100}{\watt}$; (b) $\SI{1}{\joule}$; (c)
$t=\SI{6.93}{\milli\second}$}
\resolucao{(a) $P_0=RI_0^{2}=4(25)=\SI{100}{\watt}$.\\
(b) A potência decai como $e^{-2t/\tau}$:
\[
E=\int_0^{\infty}RI_0^{2}e^{-2t/\tau}\,\mathrm{d}t
=RI_0^{2}\cdot\frac{\tau}{2}
=100\times\frac{0{,}020}{2}=\SI{1}{\joule}.
\]
\textbf{Regra útil:} a energia equivale à potência inicial mantida por
$\tau/2$. A "constante de tempo da potência" é metade da constante de tempo da
corrente, porque a potência vai com o quadrado.\\
(c) Metade da energia:
\[
\int_0^{t}e^{-2s/\tau}\mathrm{d}s=\frac12\int_0^{\infty}
\;\Longrightarrow\;
1-e^{-2t/\tau}=\frac12
\;\Longrightarrow\;
t=\frac{\tau\ln2}{2}=\SI{6.93}{\milli\second}.
\]
\textbf{Interpretação:} metade de toda a energia é liberada em apenas
$0{,}35\,\tau$ --- o transitório é fortemente concentrado no início. É o pico
inicial, e não a duração, que dimensiona o componente.}
\end{questao}

\begin{questao}[3]{}
Compare uma lâmpada incandescente de $\SI{60}{\watt}$ com uma de LED de
$\SI{9}{\watt}$ e mesmo fluxo luminoso, usadas $\SI{5}{\hour}$ por dia.
\begin{itensq}
\item Determine o consumo e o custo anual de cada uma, a
      $\mathrm{R\$}\,0{,}95/\si{\kilo\watt\hour}$.
\item Determine o tempo de retorno se a lâmpada de LED custar
      $\mathrm{R\$}\,15{,}00$ a mais.
\end{itensq}
\resposta{(a) $\mathrm{R\$}\,104{,}02$ contra $\mathrm{R\$}\,15{,}60$;
(b) cerca de $2 meses$}
\resolucao{(a) Horas anuais: $5\times365=\SI{1825}{\hour}$.
\[
E_{inc}=0{,}060(1825)=\SI{109.5}{\kilo\watt\hour}
\Rightarrow \mathrm{R\$}\,104{,}02;
\]
\[
E_{LED}=0{,}009(1825)=\SI{16.4}{\kilo\watt\hour}
\Rightarrow \mathrm{R\$}\,15{,}60.
\]
Economia anual: $\mathrm{R\$}\,88{,}42$.\\
(b)
\[
t=\frac{15{,}00}{88{,}42}=0{,}17\ \text{ano}\approx2 meses.
\]
\textbf{Retorno excepcional} --- poucos investimentos em engenharia se pagam em
dois meses.\\
\textbf{E há mais:} a lâmpada de LED dura tipicamente $\SI{25000}{\hour}$
contra $\SI{1000}{\hour}$ da incandescente, evitando cerca de $24$ trocas ao
longo da vida. Somando o custo das lâmpadas substituídas e da mão de obra, a
vantagem cresce ainda mais.\\
\textbf{Nota física:} a incandescente converte cerca de $95\%$ da energia em
calor e apenas $5\%$ em luz --- ela é, essencialmente, um aquecedor que
ilumina de lado.}
\end{questao}

\begin{questao}[3]{}
Uma instalação alimenta uma carga de $\SI{10}{\kilo\watt}$ por um cabo de
resistência total $\SI{0.2}{\ohm}$, em $\SI{220}{\volt}$.
\begin{itensq}
\item Determine a corrente e a perda no cabo.
\item Determine a fração da energia perdida.
\item Repita para $\SI{440}{\volt}$ e comente.
\end{itensq}
\resposta{(a) $\SI{45.5}{\ampere}$, $\SI{414}{\watt}$; (b) $4{,}1\%$;
(c) $\SI{103}{\watt}$, $1{,}0\%$}
\resolucao{(a) $I=\dfrac{10000}{220}=\SI{45.5}{\ampere}$ e
\[
P_{cabo}=RI^{2}=0{,}2(45{,}5)^{2}=\SI{414}{\watt}.
\]
(b) $414/10000=4{,}1\%$ da potência útil vira calor no cabo.\\
(c) Em $\SI{440}{\volt}$: $I=\SI{22.7}{\ampere}$ e
$P_{cabo}=0{,}2(22{,}7)^{2}=\SI{103}{\watt}$, ou $1{,}0\%$.\\
\textbf{Dobrar a tensão reduziu a perda a um quarto.} A razão é direta: para a
mesma potência, $I\propto1/U$, e a perda vai com $I^{2}$ --- logo
$P_{perda}\propto1/U^{2}$.\\
\textbf{É esta a razão de existir a alta tensão na transmissão.} Linhas de
$\SI{500}{\kilo\volt}$ transportam gigawatts com perdas de poucos por cento; nas
mesmas linhas em baixa tensão, a energia não chegaria ao destino. O preço é o
isolamento e os transformadores nas duas pontas.}
\end{questao}

\blocodesafio

\begin{questao}[4]{}
\textbf{(Demonstração --- potência média e valor eficaz.)} Prove que a potência
média em um resistor depende do valor \emph{eficaz} da corrente.
\begin{itensq}
\item Deduza $P_{\text{média}}=R\,I_{rms}^{2}$ a partir da definição.
\item Mostre que $I_{rms}\ge I_{\text{média}}$, com igualdade apenas para corrente
      constante.
\item Aplique a uma onda quadrada de $\SI{0}{\ampere}$ e
      $\SI{4}{\ampere}$ com ciclo de $\SI{25}{\percent}$.
\end{itensq}
\resposta{$I_{rms}=\SI{2}{\ampere}$ contra $I_{\text{média}}=\SI{1}{\ampere}$;
$P=R I_{rms}^{2}$}
\resolucao{(a) Por definição, a potência média em um período $T$ é
\[
P_{\text{média}}=\frac{1}{T}\int_0^{T}R\,i^{2}(t)\,\mathrm{d}t
=R\left[\frac{1}{T}\int_0^{T}i^{2}\,\mathrm{d}t\right]
=R\,I_{rms}^{2},
\]
onde a definição de valor eficaz é justamente
$I_{rms}=\sqrt{\frac1T\int_0^{T}i^{2}\mathrm{d}t}$ --- a
\emph{raiz} da \emph{média} do \emph{quadrado}.\\
(b) Pela desigualdade de Cauchy--Schwarz (ou pelo fato de a variância ser não
negativa):
\[
\left\langle i^{2}\right\rangle-\left\langle i\right\rangle^{2}
=\operatorname{Var}(i)\ge0
\;\Longrightarrow\;
I_{rms}^{2}\ge I_{\text{média}}^{2}.
\]
A igualdade vale se e somente se a variância é nula, isto é, se $i(t)$ é
\textbf{constante}.\\
\textbf{Leitura:} qualquer flutuação da corrente aumenta o aquecimento em
relação ao que a média sugeriria. Corrente ondulada aquece mais que corrente
contínua de mesmo valor médio --- resultado com consequências diretas em
retificadores e conversores.\\
(c) $\langle i\rangle=0{,}25(4)=\SI{1}{\ampere}$;
$\langle i^{2}\rangle=0{,}25(16)=\SI{4}{\ampere\squared}$, logo
$I_{rms}=\SI{2}{\ampere}$.\\
Com $R=\SI{10}{\ohm}$: $P=10(4)=\SI{40}{\watt}$, e não os
$\SI{10}{\watt}$ que a corrente média sugeriria. Fator $4$ --- exatamente
$(I_{rms}/I_{\text{média}})^{2}$.}
\end{questao}

\begin{questao}[4]{}
\textbf{(Integração --- energia em transitório.)} Um resistor $R$ conduz
$i(t)=I_0e^{-t/\tau}$.
\begin{itensq}
\item Deduza a energia total dissipada.
\item Determine a "potência média equivalente" ao longo de $5\tau$.
\item Compare com a energia de um pulso retangular de mesma amplitude e
      duração $\tau$.
\end{itensq}
\resposta{$E=\dfrac{RI_0^{2}\tau}{2}$; equivale à potência inicial mantida por
$\tau/2$}
\resolucao{(a)
\[
E=\int_0^{\infty}R\,I_0^{2}e^{-2t/\tau}\,\mathrm{d}t
=RI_0^{2}\left[-\frac{\tau}{2}e^{-2t/\tau}\right]_0^{\infty}
=\frac{RI_0^{2}\tau}{2}.
\]
\textbf{Observe o fator $\tau/2$, e não $\tau$:} a potência decai com constante
$\tau/2$ porque é proporcional a $i^{2}$, e o quadrado de uma exponencial tem
metade da constante de tempo.\\
(b) Distribuindo $E$ por $5\tau$:
\[
P_{\text{média}}=\frac{RI_0^{2}\tau/2}{5\tau}=\frac{RI_0^{2}}{10}=\frac{P_0}{10}.
\]
Apenas um décimo da potência inicial --- confirmando que o transitório é curto
diante da janela considerada.\\
(c) \textbf{Pulso retangular} de amplitude $I_0$ e duração $\tau$:
$E_{ret}=RI_0^{2}\tau$, exatamente o \textbf{dobro}.\\
\textbf{Consequência de dimensionamento:} substituir o decaimento real por um
pulso retangular equivalente de duração $\tau$ é uma aproximação
\emph{conservadora} --- ela superestima a energia por um fator $2$. Útil como
estimativa rápida e segura, desde que se saiba que há folga embutida.}
\end{questao}

\begin{questao}[4]{}
\textbf{(Projeto de proteção.)} Um resistor de $\SI{100}{\ohm}$ e
$\frac14\,\si{\watt}$ deve ser protegido contra sobrecarga.
\begin{itensq}
\item Determine a tensão e a corrente máximas admissíveis.
\item Projete um resistor limitador em série, para uma fonte de
      $\SI{24}{\volt}$.
\item Verifique a potência do limitador.
\end{itensq}
\resposta{$U_{\max}=\SI{5}{\volt}$, $I_{\max}=\SI{50}{\milli\ampere}$;
limitador de $\SI{380}{\ohm}$}
\resolucao{(a)
\[
U_{\max}=\sqrt{PR}=\sqrt{0{,}25(100)}=\SI{5}{\volt};
\qquad
I_{\max}=\sqrt{\frac{P}{R}}=\sqrt{0{,}0025}=\SI{50}{\milli\ampere}.
\]
(b) Com $\SI{24}{\volt}$ e a exigência $I\le\SI{50}{\milli\ampere}$:
\[
R_{total}\ge\frac{24}{0{,}050}=\SI{480}{\ohm}
\;\Longrightarrow\;
R_{lim}\ge480-100=\SI{380}{\ohm}.
\]
O valor E24 imediatamente superior é $\SI{390}{\ohm}$, que dá
$I=24/490=\SI{49}{\milli\ampere}$ \checkmark.\\
(c) $P_{lim}=390(0{,}049)^{2}=\SI{0.94}{\watt}$ --- especifique
$\SI{2}{\watt}$.\\
\textbf{Observação incômoda:} o limitador dissipa \emph{quatro vezes} mais que o
resistor protegido, e o conjunto consome $\SI{1.18}{\watt}$ para entregar
$\SI{0.24}{\watt}$ ao componente de interesse.\\
\textbf{Quando isso se justifica:} apenas se o resistor de
$\SI{100}{\ohm}$ tiver função específica (sensor, elemento de calibração) que o
torne insubstituível. Caso contrário, redimensionar o próprio resistor para uma
potência maior é mais simples e mais eficiente.}
\end{questao}

\begin{questao}[4]{}
\textbf{(Análise de retorno.)} Dois equipamentos entregam a mesma energia útil
de $\SI{1000}{\kilo\watt\hour}$ por ano: o modelo A tem rendimento de
$\SI{80}{\percent}$ e o B, de $\SI{92}{\percent}$.
\begin{itensq}
\item Determine a energia de entrada de cada um e a diferença.
\item Determine o tempo de retorno se B custar $\mathrm{R\$}\,800$ a mais, com
      tarifa de $\mathrm{R\$}\,0{,}95/\si{\kilo\watt\hour}$.
\item Discuta os fatores que a análise simples ignora.
\end{itensq}
\resposta{(a) $1250$ e $\SI{1087}{\kilo\watt\hour}$, diferença de
$\SI{163}{\kilo\watt\hour}$; (b) cerca de $\SI{5.2}{year}$}
\resolucao{(a)
\[
E_A=\frac{1000}{0{,}80}=\SI{1250}{\kilo\watt\hour};
\qquad
E_B=\frac{1000}{0{,}92}=\SI{1087}{\kilo\watt\hour};
\]
diferença de $\SI{163}{\kilo\watt\hour}$ por ano.\\
(b) Economia anual: $163(0{,}95)=\mathrm{R\$}\,154{,}89$.
\[
t=\frac{800}{154{,}89}=\SI{5.2}{year}.
\]
(c) \textbf{O que a conta simples ignora.}
\begin{itemize}
\item \textbf{Reajuste da tarifa:} se a energia subir acima da inflação, o
retorno acelera. Historicamente é o que ocorre.
\item \textbf{Custo do capital:} $\mathrm{R\$}\,800$ aplicados renderiam algo
--- o retorno real é mais longo que o nominal.
\item \textbf{Vida útil:} se o equipamento durar $\SI{10}{year}$, B gera lucro
nos últimos cinco. Se durar $\SI{4}{year}$, a troca não se paga.
\item \textbf{Custo do calor rejeitado:} A dissipa $\SI{250}{\kilo\watt\hour}$
contra $\SI{87}{\kilo\watt\hour}$ de B. Em ambiente climatizado, esse calor
extra precisa ser \emph{removido} --- o que pode dobrar a economia efetiva.
\item \textbf{Regime de uso:} rendimentos variam com a carga, e o valor de
catálogo costuma ser o de plena carga.
\end{itemize}
\textbf{Conclusão:} $\SI{5.2}{year}$ é aceitável para equipamento industrial de
vida longa e uso contínuo; é ruim para uso intermitente ou equipamento de
substituição rápida. \textbf{A resposta depende do horizonte, não apenas do
número.}}
\end{questao}

\begin{questao}[4]{}
\textbf{(Restrição orçamentária.)} Deseja-se limitar a
$\mathrm{R\$}\,90{,}00$ o custo mensal de um aquecedor de
$\SI{1.5}{\kilo\watt}$, com tarifa de
$\mathrm{R\$}\,1{,}00/\si{\kilo\watt\hour}$.
\begin{itensq}
\item Determine o tempo de uso admissível.
\item Determine a potência que permitiria uso de $\SI{4}{\hour}$ diárias dentro
      do mesmo orçamento.
\item Discuta qual das duas soluções é preferível.
\end{itensq}
\resposta{(a) $\SI{60}{\hour}$ por mês, ou $\SI{2}{\hour}$ por dia;
(b) $\SI{750}{\watt}$}
\resolucao{(a) O orçamento permite $\SI{90}{\kilo\watt\hour}$ por mês:
\[
t=\frac{90}{1{,}5}=\SI{60}{\hour\per month}
=\SI{2}{\hour\per day}.
\]
(b) Para $4\times30=\SI{120}{\hour}$ mensais dentro dos mesmos
$\SI{90}{\kilo\watt\hour}$:
\[
P=\frac{90}{120}=\SI{0.75}{\kilo\watt}=\SI{750}{\watt}.
\]
(c) \textbf{Depende do objetivo térmico.} As duas soluções entregam a
\emph{mesma energia} --- e portanto o mesmo calor total no mês.
\begin{itemize}
\item \textbf{$\SI{1.5}{\kilo\watt}$ por $\SI{2}{\hour}$:} aquece rápido e
atinge temperatura mais alta, mas o ambiente esfria nas outras
$\SI{22}{\hour}$.
\item \textbf{$\SI{750}{\watt}$ por $\SI{4}{\hour}$:} aquecimento mais suave e
prolongado, com temperatura média mais estável.
\end{itemize}
\textbf{Para conforto térmico, a segunda costuma ser superior} --- o ambiente
tem inércia térmica, e potência menor por mais tempo reduz as oscilações.\\
\textbf{E há uma terceira solução, melhor que ambas:} melhorar o
\emph{isolamento} do ambiente. Isso reduz a energia necessária, e não apenas a
redistribui --- atacando a causa em vez do sintoma.}
\end{questao}

\begin{questao}[4]{}
\textbf{(Otimização econômica de condutor.)} Um cabo alimenta uma carga de
$\SI{40}{\ampere}$ por $\SI{4000}{\hour}$ ao ano. O custo do cabo é
proporcional à seção, e a resistência é inversamente proporcional a ela.
\begin{itensq}
\item Escreva o custo total como função da seção.
\item Determine a condição de mínimo.
\item Interprete o resultado.
\end{itensq}
\resposta{No ótimo, o custo anualizado do cabo iguala o custo anual da energia
perdida}
\resolucao{(a) Sejam $S$ a seção, $k_1$ o custo anualizado do cabo por unidade
de seção e $R=\rho\ell/S$. O custo total anual é
\[
C(S)=\underbrace{k_1S}_{\text{investimento}}
+\underbrace{\frac{\rho\ell}{S}I^{2}\,t\,c_E}_{\text{energia perdida}}
=k_1S+\frac{k_2}{S},
\]
com $k_2=\rho\ell I^{2}t\,c_E$.\\
(b) Derivando e anulando:
\[
\frac{\mathrm{d}C}{\mathrm{d}S}=k_1-\frac{k_2}{S^{2}}=0
\;\Longrightarrow\;
S^{*}=\sqrt{\frac{k_2}{k_1}}.
\]
A segunda derivada, $2k_2/S^{3}>0$, confirma o mínimo.\\
(c) \textbf{Substituindo $S^{*}$ nas duas parcelas:}
\[
k_1S^{*}=\sqrt{k_1k_2}
\qquad\text{e}\qquad
\frac{k_2}{S^{*}}=\sqrt{k_1k_2}.
\]
As duas são \textbf{iguais} no ótimo --- resultado conhecido como \emph{lei de
Kelvin}: \textbf{a seção econômica é aquela em que o custo anualizado do
condutor iguala o custo anual das perdas}.\\
\textbf{Leituras práticas.}
\begin{itemize}
\item Cabos usados \emph{muitas horas} por ano ($t$ grande) justificam seções
maiores --- $S^{*}\propto\sqrt{t}$.
\item Correntes maiores também: $S^{*}\propto I$.
\item Energia mais cara empurra a seção para cima: $S^{*}\propto\sqrt{c_E}$.
\end{itemize}
\textbf{Importante:} a seção econômica é frequentemente \emph{maior} que a
mínima exigida pela norma (que garante apenas segurança térmica e queda de
tensão). Dimensionar pelo mínimo normativo é seguro, mas costuma ser
economicamente subótimo em instalações de uso intenso.}
\end{questao}

\begin{questao}[4]{}
\textbf{(Balanço energético de um chuveiro.)} Um chuveiro de
$\SI{5500}{\watt}$ aquece água de $\SI{20}{\celsius}$ a
$\SI{40}{\celsius}$.
\begin{itensq}
\item Determine a vazão máxima, supondo rendimento de $\SI{100}{\percent}$.
\item Determine o custo de um banho de $\SI{8}{\minute}$ a
      $\mathrm{R\$}\,0{,}95/\si{\kilo\watt\hour}$.
\item Discuta por que o rendimento é tão alto neste caso.
\end{itensq}
\dados{$c_{\text{água}}=\SI{4180}{\joule\per\kilogram\per\kelvin}$}
\resposta{(a) $\SI{0.0658}{\kilogram\per\second}\approx
\SI{3.9}{\liter\per\minute}$; (b) $\mathrm{R\$}\,0{,}70$}
\resolucao{(a) A potência necessária para aquecer uma vazão mássica $\dot m$
por $\Delta T=\SI{20}{\kelvin}$ é $P=\dot m\,c\,\Delta T$:
\[
\dot m=\frac{P}{c\,\Delta T}=\frac{5500}{4180(20)}
=\SI{0.0658}{\kilogram\per\second}\approx\SI{3.9}{\liter\per\minute}.
\]
Vazão modesta --- e é exatamente por isso que abrir demais o registro faz a água
esfriar: a potência é fixa, e mais vazão significa menor $\Delta T$.\\
(b)
\[
E=5{,}5\,\si{\kilo\watt}\times\frac{8}{60}\,\si{\hour}
=\SI{0.733}{\kilo\watt\hour}
\;\Longrightarrow\;
\text{custo}=0{,}733(0{,}95)=\mathrm{R\$}\,0{,}70.
\]
Um banho diário custa cerca de $\mathrm{R\$}\,21$ por mês --- e para uma
família de quatro pessoas, mais de $\mathrm{R\$}\,80$.\\
(c) \textbf{Por que o rendimento é quase $100\%$.} A resistência está
\emph{imersa} na água, e \textbf{toda} a energia dissipada por efeito Joule vira
calor exatamente onde ele é desejado. Não há conversão intermediária, nem calor
"desperdiçado" --- o que seria perda em outro contexto é aqui o produto.\\
\textbf{Contraste instrutivo:} uma lâmpada incandescente tem os mesmos
$\approx100\%$ de conversão em calor, mas rendimento \emph{luminoso} de $5\%$.
O rendimento só se define em relação ao \emph{objetivo}: a mesma física é
excelente para aquecer e péssima para iluminar.}
\end{questao}

\begin{questao}[4]{}
\textbf{(Diagnóstico por consumo.)} A conta de energia de uma residência subiu
de $\SI{180}{\kilo\watt\hour}$ para $\SI{320}{\kilo\watt\hour}$ mensais, sem
mudança de hábitos.
\begin{itensq}
\item Determine a potência contínua equivalente ao acréscimo.
\item Levante hipóteses compatíveis com esse valor.
\item Proponha um procedimento de investigação.
\end{itensq}
\resposta{Acréscimo de $\SI{140}{\kilo\watt\hour}$, equivalente a
$\SI{194}{\watt}$ contínuos}
\resolucao{(a) O mês tem cerca de $\SI{720}{\hour}$:
\[
P=\frac{140000}{720}=\SI{194}{\watt}\ \text{contínuos}.
\]
(b) \textbf{Hipóteses compatíveis com $\approx\SI{200}{\watt}$ permanentes.}
\begin{itemize}
\item \textbf{Fuga para o terra} por isolamento danificado --- especialmente em
resistência de chuveiro ou aquecedor de água. Uma fuga de
$\SI{0.88}{\ampere}$ em $\SI{220}{\volt}$ dá exatamente
$\SI{194}{\watt}$.
\item \textbf{Geladeira com vedação ruim ou gás insuficiente}, operando em
regime quase contínuo em vez de ciclado.
\item \textbf{Bomba ou motor} funcionando sem parar por falha de boia ou
pressostato.
\item \textbf{Ligação clandestina} a jusante do medidor.
\end{itemize}
(c) \textbf{Procedimento de investigação, em ordem.}
\begin{enumerate}
\item \textbf{Desligar o disjuntor geral} e verificar se o medidor continua
girando. Se sim, o problema está \emph{antes} do quadro --- ligação irregular ou
medidor defeituoso.
\item \textbf{Religar circuito por circuito}, observando o medidor a cada
etapa, para isolar o ramo responsável.
\item \textbf{Medir a corrente} de cada circuito suspeito com alicate
amperímetro, com todos os aparelhos desligados. Corrente não nula indica fuga.
\item \textbf{Testar o dispositivo DR}, se houver --- fuga acima de
$\SI{30}{\milli\ampere}$ deveria ter provocado desarme. A ausência de desarme
com fuga confirmada indica DR defeituoso, e é problema de segurança, não apenas
de conta.
\end{enumerate}
\textbf{Prioridade:} uma fuga de $\SI{0.88}{\ampere}$ é risco de choque e de
incêndio. O custo na conta é o sintoma menos grave.}
\end{questao}

% END BANCO COMPLEMENTAR

\gabarito[2.3 Potência elétrica e energia]

% =====================================================================
\subsection{Fontes reais e Leis de Kirchhoff}
% =====================================================================

\blocofixacao

\begin{questao}[1]{}
Um circuito tem duas fontes ideais de $\SI{6}{\volt}$ e $\SI{4}{\volt}$ ligadas em
série e em \textbf{oposição}.
\begin{itensq}
\item Qual é a tensão total fornecida ao circuito?
\item Se a resistência total for $\SI{5}{\ohm}$, qual é a corrente?
\end{itensq}
\resposta{(a) \SI{2}{\volt}; (b) \SI{0.4}{\ampere}}
\resolucao{(a) Em oposição, as fem se subtraem:
$U = 6 - 4 = \SI{2}{\volt}$.\\
(b) $I = \dfrac{U}{R} = \dfrac{2}{5} = \SI{0.4}{\ampere}$.
Se estivessem em série \emph{concordante}, teríamos $6+4 = \SI{10}{\volt}$ e
$I = \SI{2}{\ampere}$.}
\end{questao}

\begin{questao}[1]{}
No nó da figura, as correntes $I_1 = \SI{5}{\ampere}$ e $I_2 = \SI{3}{\ampere}$
entram, e $I_3$ sai. Pela lei dos nós de Kirchhoff, $I_3$ vale:

\circuitoq{%
\begin{circuitikz}[scale=1,transform shape,line width=0.8pt]
 \coordinate (o) at (0,0);
 \draw[->,Icol,line width=1.1pt] (-2,0.6) -- (o) node[midway,above]{$I_1$};
 \draw[->,Icol,line width=1.1pt] (-2,-0.6) -- (o) node[midway,below]{$I_2$};
 \draw[->,Icol,line width=1.1pt] (o) -- (2,0) node[midway,above]{$I_3$};
 \fill (o) circle (0.06);
\end{circuitikz}}

\resposta{$I_3 = \SI{8}{\ampere}$}
\resolucao{A lei dos nós afirma que a soma das correntes que entram é igual à das
que saem:
\[
I_3 = I_1 + I_2 = 5 + 3 = \SI{8}{\ampere}.
\]
Ela expressa a conservação da carga: o nó não acumula carga.}
\end{questao}

\blocoaplicacao

\begin{questao}[2]{}
Uma bateria de fem $\EMF = \SI{12}{\volt}$ e resistência interna
$r = \SI{0.5}{\ohm}$ alimenta um resistor externo $R = \SI{5.5}{\ohm}$. Determine
a corrente e a tensão nos terminais da bateria.
\resposta{$I = \SI{2}{\ampere}$; $U = \SI{11}{\volt}$}
\resolucao{A resistência interna está em série com a carga:
\[
I=\frac{\EMF}{R+r}=\frac{12}{5{,}5+0{,}5}=\SI{2}{\ampere}.
\]
A tensão nos terminais é a fem menos a queda interna:
\[
U=\EMF - r\,I = 12 - 0{,}5\cdot2 = \SI{11}{\volt}.
\]
A diferença de \SI{1}{\volt} "some" dentro da bateria --- é a queda em $r$. Quanto
maior a corrente, menor a tensão entregue: por isso baterias velhas (com $r$
elevado) fornecem menos tensão sob carga.}
\end{questao}

\begin{questao}[2]{}
Uma bateria de fem $\EMF$ e resistência interna $r = \SI{0.4}{\ohm}$ apresenta
tensão de $\SI{12}{\volt}$ em seus terminais quando fornece
$\SI{5}{\ampere}$. Determine a fem e a tensão em vazio (corrente nula).
\resposta{$\EMF = \SI{14}{\volt}$; em vazio, $U = \SI{14}{\volt}$}
\resolucao{Da equação da fonte real, $U = \EMF - r\,I$:
\[
\EMF = U + r\,I = 12 + 0{,}4\cdot5 = \SI{14}{\volt}.
\]
Em vazio ($I = 0$), não há queda interna: $U = \EMF = \SI{14}{\volt}$. A tensão de
circuito aberto de uma fonte é sempre igual à sua fem.}
\end{questao}

\begin{questao}[2]{}
No nó da figura, são conhecidas quatro das cinco correntes. Determine $I_5$,
indicando se entra ou sai do nó.

\circuitoq{%
\begin{tikzpicture}[scale=1,transform shape,line width=0.9pt,>=Stealth]
 \coordinate (o) at (0,0);
 \draw[->,Icol] (-2,0.8) -- (o) node[midway,above]{$\SI{6}{\ampere}$};
 \draw[->,Icol] (-2,-0.8) -- (o) node[midway,below]{$\SI{4}{\ampere}$};
 \draw[->,Icol] (o) -- (2,0.8) node[midway,above]{$\SI{3}{\ampere}$};
 \draw[->,Icol] (o) -- (2,-0.8) node[midway,below]{$\SI{5}{\ampere}$};
 \draw[->,Vcol,thick] (o) -- (0,-1.8) node[midway,right]{$I_5$};
 \fill (o) circle (0.07);
\end{tikzpicture}}{Nó com cinco ramos.}

\resposta{$I_5 = \SI{2}{\ampere}$, saindo do nó}
\resolucao{Pela lei dos nós, a soma das que entram é igual à soma das que saem.
Entram $6 + 4 = \SI{10}{\ampere}$; saem $3 + 5 = \SI{8}{\ampere}$ pelos ramos
conhecidos. Para equilibrar, $I_5$ deve \emph{sair} com:
\[
I_5 = 10 - 8 = \SI{2}{\ampere}.
\]
Como o sentido adotado na figura já era de saída, o valor positivo confirma essa
suposição.}
\end{questao}

\blocoaprofundamento

\begin{questao}[3]{}
No circuito de duas malhas da figura, aplique as leis de Kirchhoff para determinar
as correntes $I_1$, $I_2$ e a corrente $I_3$ no resistor central.

\circuitoq{%
\begin{circuitikz}[scale=1,transform shape,line width=0.8pt]
 % malha esquerda
 \draw (0,0) to[battery1,l_={$E_1=\SI{18}{\volt}$}] (0,3)
       to[R,l={$R_1=\SI{2}{\ohm}$}] (3,3) coordinate(top);
 \draw (top) to[R,l={$R_2=\SI{6}{\ohm}$},i>_={$I_3$}] (3,0) coordinate(bot) -- (0,0);
 % malha direita
 \draw (top) to[R,l={$R_3=\SI{6}{\ohm}$}] (6,3)
       to[battery1,l={$E_2=\SI{6}{\volt}$},invert] (6,0) -- (bot);
 \draw[->,Icol,line width=1pt] (0.6,3.4)--(1.5,3.4) node[midway,above,font=\footnotesize]{$I_1$};
 \draw[->,Icol,line width=1pt] (5.4,3.4)--(4.5,3.4) node[midway,above,font=\footnotesize]{$I_2$};
\end{circuitikz}}

\resposta{$I_1 = \SI{3}{\ampere}$, $I_2 = \SI{1}{\ampere}$, $I_3 = \SI{2}{\ampere}$}
\resolucao{\textbf{Lei dos nós} (no topo): $I_1 = I_2 + I_3$, logo $I_3 = I_1 - I_2$.\\
\textbf{Malha esquerda} ($E_1$, $R_1$, $R_2$):
\[
E_1 = R_1 I_1 + R_2 I_3 \;\Rightarrow\; 18 = 2I_1 + 6(I_1 - I_2).
\]
\textbf{Malha direita} ($E_2$, $R_3$, $R_2$):
\[
E_2 = R_3 I_2 + R_2 I_3 \;\Rightarrow\; 6 = 6I_2 + 6(I_1 - I_2)\cdot(-1)?
\]
Organizando o sistema (com $I_3 = I_1 - I_2$):
\[
\begin{cases}
8I_1 - 6I_2 = 18\\[2pt]
-6I_1 + 12I_2 = -6
\end{cases}
\;\Rightarrow\;
I_1 = \SI{3}{\ampere},\ I_2 = \SI{1}{\ampere}.
\]
Logo $I_3 = 3 - 1 = \SI{2}{\ampere}$.\\
\textbf{Verificação} pela malha esquerda: $2\cdot3 + 6\cdot2 = 6+12 =
\SI{18}{\volt} = E_1$. Resultado consistente.}
\end{questao}

\begin{questao}[3]{}
No circuito de duas malhas da figura, determine as três correntes de ramo
($I_1$, $I_2$, $I_3$) aplicando as leis de Kirchhoff.

\circuitoq{%
\begin{circuitikz}[scale=1,transform shape,line width=0.8pt]
 \draw (0,0) to[battery1,l_={$E_1=\SI{12}{\volt}$}] (0,3)
       to[R,l={$R_1=\SI{2}{\ohm}$}] (3,3) coordinate(top);
 \draw (top) to[R,l={$R_2=\SI{4}{\ohm}$},i>_={$I_3$}] (3,0) coordinate(bot) -- (0,0);
 \draw (top) to[R,l={$R_3=\SI{2}{\ohm}$}] (6,3)
       to[battery1,l={$E_2=\SI{6}{\volt}$}] (6,0) -- (bot);
 \draw[->,Icol,line width=1pt] (0.6,3.4)--(1.5,3.4) node[midway,above,font=\footnotesize]{$I_1$};
 \draw[->,Icol,line width=1pt] (5.4,3.4)--(4.5,3.4) node[midway,above,font=\footnotesize]{$I_2$};
\end{circuitikz}}

\resposta{$I_1 = \SI{2.4}{\ampere}$, $I_2 = \SI{-0.6}{\ampere}$, $I_3 = \SI{1.8}{\ampere}$}
\resolucao{\textbf{Nó superior:} $I_1 + I_2 = I_3$.\\
\textbf{Malha esquerda:} $E_1 = R_1 I_1 + R_2 I_3
\Rightarrow 12 = 2I_1 + 4I_3$.\\
\textbf{Malha direita:} $E_2 = R_3 I_2 + R_2 I_3
\Rightarrow 6 = 2I_2 + 4I_3$.\\
Resolvendo o sistema (com $I_3 = I_1 + I_2$):
\[
I_1 = \SI{2.4}{\ampere},\quad I_2 = \SI{-0.6}{\ampere},\quad
I_3 = \SI{1.8}{\ampere}.
\]
O sinal \emph{negativo} de $I_2$ significa que a corrente real nesse ramo é
\emph{oposta} ao sentido adotado: a fonte $E_2$, mais fraca, na verdade
\textbf{recebe} corrente (está sendo carregada) em vez de fornecer.\\
\textbf{Verificação} (malha esquerda): $2\cdot2{,}4 + 4\cdot1{,}8 = 4{,}8+7{,}2
= \SI{12}{\volt} = E_1$. Resultado consistente. Este é um bom lembrete: adotar um sentido
errado não invalida a solução --- o sinal negativo se encarrega de corrigi-lo.}
\end{questao}

\begin{questao}[3]{}
No circuito da figura, quatro fontes ideais de tensão estão interligadas, e os
terminais $a$, $b$ e $c$ estão \textbf{abertos} (nenhuma carga conectada).
Observe atentamente a \emph{polaridade} indicada dentro de cada fonte.

\circuitoq{%
\begin{circuitikz}[scale=1.05,transform shape,line width=0.9pt]
 % barramento vertical
 \coordinate (N1) at (0,4.2);      % no superior
 \coordinate (N2) at (0,2.2);      % no inferior (referencia)
 \draw (N1) -- (0,4.2);
 % fonte 5V vertical (+ em cima, - embaixo)
 \draw (N2) to[\fonteV,l_={$\SI{5}{\volt}$},invert] (N1);
 % ramo superior -> a  (4V com '-' a esquerda, '+' a direita)
 \draw (N1) -- (0.9,4.2) to[\fonteV,l={$\SI{4}{\volt}$}] (2.9,4.2)
       to[short,-o] (4.0,4.2) node[right,font=\Large]{$a$};
 % barramento inferior
 \draw (N2) -- (0,0.2);
 % ramo do meio -> b  (5V com '+' a esquerda, '-' a direita)
 \draw (0,1.4) -- (0.9,1.4) to[\fonteV,l_={$\SI{5}{\volt}$},invert] (2.9,1.4)
       to[short,-o] (4.0,1.4) node[right,font=\Large]{$b$};
 \fill (0,1.4) circle (0.07);
 % ramo inferior -> c  (3V com '-' a esquerda, '+' a direita)
 \draw (0,0.2) -- (0.9,0.2) to[\fonteV,l_={$\SI{3}{\volt}$}] (2.9,0.2)
       to[short,-o] (4.0,0.2) node[right,font=\Large]{$c$};
\end{circuitikz}}

\begin{itensq}
\item Determine $V_{ac}$.
\item Determine também $V_{ab}$ e $V_{bc}$, e verifique a consistência
      ($V_{ab}+V_{bc}=V_{ac}$).
\item Por que os valores \emph{não} dependem das resistências dos fios?
\end{itensq}
\resposta{(a) $V_{ac}=\SI{6}{\volt}$; (b) $V_{ab}=\SI{14}{\volt}$,
$V_{bc}=\SI{-8}{\volt}$; (c) ver comentário}
\resolucao{\textbf{Estratégia.} Como os terminais estão abertos, \emph{não circula
corrente} em nenhum ramo. Sem corrente, não há queda em resistência alguma: os
potenciais são determinados exclusivamente pelas fem, somadas com o \emph{sinal}
correto. Basta adotar uma referência e "caminhar" pelo circuito.

\medskip
\textbf{Passo 1 --- adotar a referência.} Seja o nó inferior do barramento
(ao qual se ligam os ramos de $b$ e $c$) o nó de referência: $V_{N_2}=0$.

\textbf{Passo 2 --- subir pela fonte de \SI{5}{\volt} vertical.} Ela tem o
terminal $+$ \emph{em cima} e $-$ embaixo. Indo de $N_2$ para $N_1$, entramos
pelo $-$ e saímos pelo $+$: \emph{ganhamos} \SI{5}{\volt}.
\[
V_{N_1}=V_{N_2}+5=\SI{5}{\volt}.
\]

\textbf{Passo 3 --- percorrer cada ramo até seu terminal.}
\begin{itemize}[leftmargin=1.7em,itemsep=2pt,topsep=3pt]
\item \textbf{Ramo de $a$} (fonte \SI{4}{\volt}, com $-$ à esquerda e $+$ à
      direita): indo de $N_1$ para $a$, entramos pelo $-$ e saímos pelo $+$,
      ganhando \SI{4}{\volt}:
      \[V_a = V_{N_1}+4 = 5+4=\SI{9}{\volt}.\]
\item \textbf{Ramo de $b$} (fonte \SI{5}{\volt}, com $+$ à esquerda e $-$ à
      direita): indo de $N_2$ para $b$, entramos pelo $+$ e saímos pelo $-$,
      \emph{perdendo} \SI{5}{\volt}:
      \[V_b = V_{N_2}-5 = \SI{-5}{\volt}.\]
\item \textbf{Ramo de $c$} (fonte \SI{3}{\volt}, com $-$ à esquerda e $+$ à
      direita): ganhamos \SI{3}{\volt}:
      \[V_c = V_{N_2}+3 = \SI{3}{\volt}.\]
\end{itemize}

\textbf{Passo 4 --- calcular as diferenças.}
\[
\boxed{\;V_{ac}=V_a-V_c = 9-3 = \SI{6}{\volt}\;}
\]
\[
V_{ab}=V_a-V_b = 9-(-5)=\SI{14}{\volt},
\qquad
V_{bc}=V_b-V_c = -5-3=\SI{-8}{\volt}.
\]
\textbf{Consistência:} $V_{ab}+V_{bc} = 14+(-8) = \SI{6}{\volt} = V_{ac}$
\checkmark --- como deve ser, pois a diferença de potencial é uma função de
estado (independe do caminho).

\medskip
\textbf{(c) Por que os fios não importam?} A queda em um resistor é $U=RI$. Com
$I=0$ em todos os ramos (terminais abertos), \emph{qualquer} resistência produz
queda nula. Os potenciais dos terminais dependem só das fem e de suas
polaridades.\\
\textbf{Atenção ao erro mais comum:} somar as fontes ignorando a polaridade
(obtendo, por exemplo, $4+5+3=12$). O sinal de cada fonte é determinado pelo
\emph{sentido em que a percorremos}: entrando pelo $-$ e saindo pelo $+$, soma-se;
o contrário, subtrai-se. \emph{Sempre marque a polaridade antes de escrever a
primeira equação.}}
\end{questao}

\blocodesafio

\begin{questao}[4]{}
Duas baterias de fem $\EMF_1 = \SI{12}{\volt}$ ($r_1=\SI{0.2}{\ohm}$) e
$\EMF_2 = \SI{12.6}{\volt}$ ($r_2=\SI{0.1}{\ohm}$) são ligadas em
\textbf{paralelo} para alimentar uma carga $R_L = \SI{2}{\ohm}$ --- situação
comum em bancos de baterias e em sistemas com gerador e bateria em paralelo.

\circuitoq{%
\begin{circuitikz}[scale=0.95,transform shape,line width=0.8pt]
 \draw (0,0) to[battery1,l_={$\EMF_1$}] (0,1.6) to[R,l_={$r_1$},color=Vcol] (0,3);
 \draw (2.4,0) to[battery1,l_={$\EMF_2$}] (2.4,1.6) to[R,l_={$r_2$},color=Vcol] (2.4,3);
 \draw (0,3) -- (2.4,3) -- (5,3) to[R,l={$R_L=\SI{2}{\ohm}$}] (5,0) -- (2.4,0) -- (0,0);
 \node[above,Vcol] at (3.6,3.1) {$V$};
 \fill (2.4,3) circle (0.07); \fill (2.4,0) circle (0.07);
\end{circuitikz}}

\begin{itensq}
\item Determine a tensão $V$ sobre a carga (use análise nodal).
\item Calcule a corrente fornecida por \emph{cada} bateria.
\item Interprete o sinal obtido e explique o risco prático dessa associação.
\end{itensq}
\resposta{(a) $V = \SI{12}{\volt}$; (b) $I_1 = 0$, $I_2 = \SI{6}{\ampere}$;
(c) ver comentário}
\resolucao{\textbf{(a)} Tomando o nó superior com potencial $V$ e o inferior como
terra, a lei dos nós exige que a soma das correntes que \emph{chegam} ao nó seja
igual à que sai pela carga:
\[
\frac{\EMF_1-V}{r_1}+\frac{\EMF_2-V}{r_2}=\frac{V}{R_L}.
\]
Substituindo:
\[
\frac{12-V}{0{,}2}+\frac{12{,}6-V}{0{,}1}=\frac{V}{2}.
\]
Multiplicando por 2: $10(12-V)+20(12{,}6-V)=V$, ou seja,
\[
120-10V+252-20V=V
\;\Longrightarrow\;
372 = 31V
\;\Longrightarrow\;
V = \SI{12}{\volt}.
\]

\textbf{(b)} As correntes de cada fonte:
\[
I_1=\frac{12-12}{0{,}2}=\SI{0}{\ampere},
\qquad
I_2=\frac{12{,}6-12}{0{,}1}=\SI{6}{\ampere},
\qquad
I_L=\frac{12}{2}=\SI{6}{\ampere}.\ \checkmark
\]
\textbf{(c) Interpretação e risco.} A bateria de \emph{maior} fem assume
praticamente toda a carga, enquanto a de menor fem fica ociosa. Se a diferença de
fem fosse maior --- por exemplo, uma bateria descarregada em paralelo com uma
carregada --- a corrente na bateria mais fraca ficaria \textbf{negativa}: ela
seria \emph{carregada} pela outra, com corrente limitada apenas pelas
resistências internas, tipicamente muito pequenas.\\
\textbf{Consequência prática:} correntes de circulação elevadíssimas, aquecimento
e risco de dano ou explosão. Por isso o paralelismo de baterias exige unidades
\emph{idênticas em fem, tecnologia e estado de carga}, ou o uso de diodos de
isolamento e controladores de balanceamento --- prática obrigatória em bancos de
baterias e sistemas fotovoltaicos.}
\end{questao}

% BEGIN BANCO COMPLEMENTAR Fontes reais e Leis de Kirchhoff
% =====================================================================
%  Banco complementar --- Fontes Reais e Leis de Kirchhoff  (REVISADO)
%  Substitui a versao gerada por template (9 clones em N2, 9 em N3 e
%  7 questoes conceituais marcadas como N4).
%  O cap02 ja cobre: duas fontes ideais em oposicao, LKC em no com
%  I1 e I2 entrando, bateria 12 V / 0,5 ohm, bateria com r=0,4 ohm,
%  no com cinco correntes, DOIS circuitos de duas malhas resolvidos
%  por Kirchhoff, quatro fontes ideais com terminais abertos e --- como
%  N4 --- DUAS BATERIAS DE FEM DIFERENTES EM PARALELO. Este ultimo tema
%  motivou a remocao de uma N4 do banco antigo, que o duplicava.
%  O banco de circuitos serie revisado ja traz regulacao de tensao e a
%  determinacao de fem e r por dois pontos de carga.
%  Este banco explora: fonte ABSORVENDO potencia, banco de baterias
%  serie x paralelo, LKC generalizada, contagem de equacoes
%  independentes, propagacao de incerteza na verificacao da LKC e o
%  teorema de Tellegen em forma elementar.
%  Distribuicao mantida: 6 N1 + 9 N2 + 9 N3 + 7 N4 = 31
% =====================================================================

\subsubsection*{Banco complementar --- Fontes reais e leis de Kirchhoff}

\begin{atencao}[Organização do banco]
A fonte real é o modelo $\EMF$ em série com $r$, e dele decorre
$V=\EMF-rI$ quando ela \emph{fornece} --- mas $V=\EMF+rI$ quando ela
\textbf{recebe} corrente, situação de carregamento que aparece com frequência e
costuma ser esquecida. As leis de Kirchhoff, por sua vez, decorrem da
conservação de carga e de energia: elas valem para \emph{qualquer} elemento,
linear ou não. O N3 explora fontes absorvendo potência, bancos de baterias e
contagem de equações; o N4 exige demonstração, ensaio e propagação de
incertezas.
\end{atencao}

\blocofixacao

\begin{questao}[1]{}
Uma fonte de $\EMF=\SI{6}{\volt}$ e $r=\SI{1}{\ohm}$ fornece
$\SI{1}{\ampere}$. Determine a tensão em seus terminais.
\resposta{$V=\SI{5}{\volt}$}
\resolucao{
\[
V=\EMF-rI=6-1(1)=\SI{5}{\volt}.
\]
A diferença de $\SI{1}{\volt}$ é a queda \emph{interna}, que não chega à carga.
Ela cresce com a corrente --- é por isso que a tensão de uma pilha cai quando se
liga um aparelho que exige mais.}
\end{questao}

\begin{questao}[1]{}
Uma fonte de $\EMF=\SI{9}{\volt}$ apresenta $\SI{8.4}{\volt}$ nos terminais
quando fornece $\SI{2}{\ampere}$. Determine sua resistência interna.
\resposta{$r=\SI{0.3}{\ohm}$}
\resolucao{
\[
r=\frac{\EMF-V}{I}=\frac{9-8{,}4}{2}=\SI{0.3}{\ohm}.
\]
\textbf{Método geral:} a queda interna dividida pela corrente. Uma única medição
sob carga, somada ao valor de $\EMF$, determina $r$ --- e $\EMF$ pode ser lida
em vazio com um voltímetro de alta impedância.}
\end{questao}

\begin{questao}[1]{}
Em um nó entram $\SI{2.5}{\ampere}$ e $\SI{1.5}{\ampere}$, e sai uma única
corrente. Determine seu valor.
\resposta{$I=\SI{4}{\ampere}$}
\resolucao{Pela lei dos nós,
\[
\sum I_{entra}=\sum I_{sai}
\;\Longrightarrow\;
2{,}5+1{,}5=\SI{4}{\ampere}.
\]
\textbf{Fundamento:} a LKC expressa a \textbf{conservação da carga}. Um nó não
acumula carga --- tudo o que entra, sai. Isso vale independentemente dos
elementos ligados a ele, e mesmo que sejam não lineares.}
\end{questao}

\begin{questao}[1]{}
Uma malha simples contém uma fonte ideal de $\SI{12}{\volt}$ e resistores de
$\SI{4}{\ohm}$ e $\SI{8}{\ohm}$. Determine a corrente pela lei das malhas.
\resposta{$I=\SI{1}{\ampere}$}
\resolucao{Pela LKT, a soma das quedas iguala a f.e.m.:
\[
12=4I+8I=12I
\;\Longrightarrow\;
I=\SI{1}{\ampere}.
\]
\textbf{Fundamento:} a LKT expressa a \textbf{conservação da energia}. Uma carga
que percorre a malha e retorna ao ponto de partida deve ter ganho e perdido a
mesma energia --- o potencial é uma função de estado.}
\end{questao}

\begin{questao}[1]{}
Sobre a diferença entre fonte ideal e fonte real, é correto afirmar:
\begin{alternativas}[2]
\item a ideal mantém a tensão constante para qualquer corrente
\item a real mantém a tensão constante para qualquer corrente
\item ambas mantêm a corrente constante
\item a ideal tem resistência interna maior
\end{alternativas}
\resposta{(a)}
\resolucao{A fonte \textbf{ideal} tem $r=0$: sua tensão terminal é
$\EMF$ independentemente da corrente, inclusive em curto-circuito --- o que
exigiria corrente infinita, revelando que se trata de uma idealização.\\
A fonte \textbf{real} tem $r>0$, e sua tensão cai como $\EMF-rI$.\\
\textbf{Quando a idealização é aceitável:} sempre que $r\ll R_{carga}$. Uma
fonte de bancada com $r=\SI{10}{\milli\ohm}$ alimentando
$\SI{100}{\ohm}$ é, para todos os efeitos, ideal.}
\end{questao}

\begin{questao}[1]{}
Uma pilha de $\EMF=\SI{1.5}{\volt}$ tem $r=\SI{0.25}{\ohm}$. Determine sua
corrente de curto-circuito.
\resposta{$I_{cc}=\SI{6}{\ampere}$}
\resolucao{
\[
I_{cc}=\frac{\EMF}{r}=\frac{1{,}5}{0{,}25}=\SI{6}{\ampere}.
\]
\textbf{Ressalva importante:} o valor pressupõe $r$ constante --- hipótese
otimista. Em uma pilha real, $r$ cresce com a corrente e com o aquecimento, e a
corrente efetiva é menor. Ainda assim, $\SI{6}{\ampere}$ em uma pilha pequena
significa $\SI{9}{\watt}$ dissipados \emph{dentro} dela: aquecimento rápido,
vazamento de eletrólito e risco de ruptura.}
\end{questao}

\blocoaplicacao

\begin{questao}[2]{}
Uma bateria de $\EMF=\SI{12}{\volt}$ e $r=\SI{0.6}{\ohm}$ alimenta
$R=\SI{5.4}{\ohm}$.
\begin{itensq}
\item Determine a corrente e a tensão terminal.
\item Determine as potências e o rendimento.
\end{itensq}
\resposta{$I=\SI{2}{\ampere}$, $V=\SI{10.8}{\volt}$;
$P_R=\SI{21.6}{\watt}$, $P_r=\SI{2.4}{\watt}$, $\eta=90\%$}
\resolucao{(a)
\[
I=\frac{\EMF}{r+R}=\frac{12}{6}=\SI{2}{\ampere};
\qquad
V=12-0{,}6(2)=\SI{10.8}{\volt}.
\]
(b)
\[
P_{total}=\EMF\,I=\SI{24}{\watt};
\quad
P_R=5{,}4(4)=\SI{21.6}{\watt};
\quad
P_r=0{,}6(4)=\SI{2.4}{\watt};
\]
\[
\eta=\frac{R}{R+r}=\frac{5{,}4}{6}=90\%.
\]
\textbf{Fórmula útil:} o rendimento de uma fonte real é sempre
$R/(R+r)$ --- ele depende só da razão entre as resistências, e não da tensão
nem da corrente.}
\end{questao}

\begin{questao}[2]{}
Um carregador de $\SI{14}{\volt}$ é ligado a uma bateria de
$\SI{12}{\volt}$, com resistência total de $\SI{0.4}{\ohm}$ no circuito.
\begin{itensq}
\item Determine a corrente de carga.
\item Determine a tensão nos terminais da bateria.
\end{itensq}
\resposta{$I=\SI{5}{\ampere}$; $V=\SI{14}{\volt}$}
\resolucao{(a) As duas f.e.m. estão em \textbf{oposição}, e a diferença é que
impulsiona a corrente:
\[
I=\frac{14-12}{0{,}4}=\SI{5}{\ampere}.
\]
(b) A bateria está \textbf{recebendo} corrente, então a queda interna
\emph{soma-se} à f.e.m.:
\[
V=\EMF+rI=12+0{,}4(5)=\SI{14}{\volt}.
\]
\textbf{Inversão do sinal:} quando a fonte fornece, $V=\EMF-rI<\EMF$; quando
recebe, $V=\EMF+rI>\EMF$. É por isso que a tensão de uma bateria de automóvel
sobe de $\SI{12.6}{\volt}$ para cerca de $\SI{14}{\volt}$ com o motor ligado ---
sinal de que o alternador está carregando.}
\end{questao}

\begin{questao}[2]{}
Em um nó, entram $\SI{3}{\ampere}$, $\SI{4}{\ampere}$ e $\SI{2}{\ampere}$;
sai $\SI{6}{\ampere}$ por um ramo e $I_5$ por outro. Determine $I_5$.
\resposta{$I_5=\SI{3}{\ampere}$, saindo}
\resolucao{
\[
3+4+2=6+I_5
\;\Longrightarrow\;
I_5=9-6=\SI{3}{\ampere}.
\]
O sinal positivo confirma o sentido admitido (saindo).\\
\textbf{Convenção alternativa:} adotar todas as correntes como \emph{entrando} e
impor soma nula. As que realmente saem aparecem com sinal negativo. As duas
convenções são equivalentes --- escolha uma e mantenha-a.}
\end{questao}

\begin{questao}[2]{}
Uma malha contém fontes ideais de $\SI{24}{\volt}$ e $\SI{6}{\volt}$ em
oposição, com resistores somando $\SI{6}{\ohm}$. Determine a corrente e
identifique qual fonte absorve energia.
\resposta{$I=\SI{3}{\ampere}$; a de $\SI{6}{\volt}$ absorve}
\resolucao{
\[
I=\frac{24-6}{6}=\SI{3}{\ampere},
\]
no sentido imposto pela fonte maior.\\
\textbf{Balanço:} a de $\SI{24}{\volt}$ fornece $24(3)=\SI{72}{\watt}$; a de
$\SI{6}{\volt}$ \emph{absorve} $6(3)=\SI{18}{\watt}$; os resistores dissipam
$6(9)=\SI{54}{\watt}$. Soma: $18+54=\SI{72}{\watt}$ \checkmark.\\
\textbf{Critério para identificar:} a fonte \emph{absorve} quando a corrente
entra pelo seu terminal positivo. Se ela for recarregável, está sendo
carregada; se não for (uma pilha comum), a energia vira calor e o dispositivo
pode ser danificado.}
\end{questao}

\begin{questao}[2]{}
Duas pilhas de $\SI{1.5}{\volt}$ e $r=\SI{0.2}{\ohm}$ cada são associadas.
Determine a f.e.m. e a resistência interna equivalentes em série e em paralelo.
\resposta{Série: $\SI{3}{\volt}$ e $\SI{0.4}{\ohm}$; paralelo:
$\SI{1.5}{\volt}$ e $\SI{0.1}{\ohm}$}
\resolucao{\textbf{Série:} as f.e.m. se somam e as resistências internas também:
$\EMF=\SI{3}{\volt}$, $r=\SI{0.4}{\ohm}$.\\
\textbf{Paralelo} (pilhas \emph{idênticas}): a f.e.m. permanece
$\SI{1.5}{\volt}$ e as resistências internas ficam em paralelo:
$r=\SI{0.1}{\ohm}$.\\
\textbf{Correntes de curto:} série $3/0{,}4=\SI{7.5}{\ampere}$; paralelo
$1{,}5/0{,}1=\SI{15}{\ampere}$.\\
\textbf{Advertência:} a associação em paralelo só é segura com pilhas
\emph{idênticas e no mesmo estado de carga}. Com f.e.m. diferentes, circula
corrente entre elas mesmo sem carga externa.}
\end{questao}

\begin{questao}[2]{}
Determine o número de equações independentes de LKC e de LKT para um circuito
com $n=5$ nós e $b=8$ ramos.
\resposta{$4$ de LKC e $4$ de LKT, totalizando $8$}
\resolucao{\textbf{LKC:} $n-1=4$ equações. A do quinto nó é combinação linear
das anteriores --- redundante.\\
\textbf{LKT:} $b-n+1=8-5+1=4$ equações, uma por malha independente.\\
\textbf{Total:} $4+4=8=b$ --- exatamente o número de correntes de ramo a
determinar. O sistema é quadrado e bem posto.\\
\textbf{Não é coincidência:} $(n-1)+(b-n+1)=b$ sempre. As leis de Kirchhoff
fornecem \emph{precisamente} as equações necessárias, nem uma a mais nem a
menos.}
\end{questao}

\begin{questao}[2]{}
Uma fonte de $\EMF=\SI{12}{\volt}$ e $r=\SI{0.5}{\ohm}$ alimenta duas cargas
em paralelo, de $\SI{12}{\ohm}$ e $\SI{6}{\ohm}$. Determine a tensão terminal e
a corrente de cada carga.
\resposta{$V=\SI{10.67}{\volt}$; $\SI{0.89}{\ampere}$ e
$\SI{1.78}{\ampere}$}
\resolucao{
\[
R_{par}=\frac{12(6)}{18}=\SI{4}{\ohm};
\qquad
I=\frac{12}{4{,}5}=\SI{2.67}{\ampere};
\qquad
V=12-0{,}5(2{,}67)=\SI{10.67}{\volt}.
\]
\emph{(Conferindo: $V=R_{par}I=4(2{,}67)=\SI{10.67}{\volt}$ \checkmark.)}
\[
I_{12}=\frac{10{,}67}{12}=\SI{0.89}{\ampere};
\qquad
I_{6}=\frac{10{,}67}{6}=\SI{1.78}{\ampere}.
\]
\textbf{Acoplamento pelas cargas:} se a carga de $\SI{6}{\ohm}$ for desligada,
a corrente total cai para $12/12{,}5=\SI{0.96}{\ampere}$ e a tensão sobe para
$\SI{11.52}{\volt}$ --- o que \emph{aumenta} a corrente da carga remanescente.
As cargas se influenciam através de $r$.}
\end{questao}

\begin{questao}[2]{}
Uma bateria em uso apresenta $\EMF=\SI{1.5}{\volt}$ e $r=\SI{0.3}{\ohm}$
quando nova, e $\EMF=\SI{1.3}{\volt}$ com $r=\SI{1.2}{\ohm}$ ao final da vida.
Compare a tensão entregue a uma carga de $\SI{3}{\ohm}$.
\resposta{$\SI{1.36}{\volt}$ contra $\SI{0.93}{\volt}$ --- queda de $32\%$}
\resolucao{\textbf{Nova:} $I=1{,}5/3{,}3=\SI{0.455}{\ampere}$ e
$V=3(0{,}455)=\SI{1.36}{\volt}$.\\
\textbf{Gasta:} $I=1{,}3/4{,}2=\SI{0.310}{\ampere}$ e
$V=3(0{,}310)=\SI{0.93}{\volt}$.\\
\textbf{Observe o que domina:} a f.e.m. caiu apenas $13\%$, mas a resistência
interna \textbf{quadruplicou} --- e é ela a responsável pela maior parte da
perda de desempenho.\\
\textbf{Consequência para diagnóstico:} medir a tensão de uma pilha \emph{em
vazio} é enganoso, pois ela mostra quase $\EMF$ mesmo já esgotada. O teste
correto é sob carga, ou a medição direta de $r$.}
\end{questao}

\begin{questao}[2]{}
Verifique o balanço de potência de um circuito com fonte de
$\EMF=\SI{20}{\volt}$, $r=\SI{1}{\ohm}$ e carga de $\SI{9}{\ohm}$.
\resposta{Fornecido $\SI{40}{\watt}$; dissipado $4+\SI{36}{\watt}$}
\resolucao{
\[
I=\frac{20}{10}=\SI{2}{\ampere};
\qquad
P_{fonte}=\EMF\,I=\SI{40}{\watt};
\]
\[
P_r=1(4)=\SI{4}{\watt};
\qquad
P_R=9(4)=\SI{36}{\watt};
\qquad 4+36=40\ \checkmark
\]
\textbf{Cuidado com a distinção:} a potência \emph{gerada} usa $\EMF$, não a
tensão terminal. Usando $V=\SI{18}{\volt}$ obter-se-ia $\SI{36}{\watt}$ ---
que é a potência \emph{entregue aos terminais}, e não a produzida internamente.
A diferença de $\SI{4}{\watt}$ é exatamente o que $r$ consome.}
\end{questao}

\blocoaprofundamento

\begin{questao}[3]{}
Duas fontes reais alimentam a mesma carga: $\EMF_1=\SI{12}{\volt}$ com
$r_1=\SI{1}{\ohm}$ e $\EMF_2=\SI{6}{\volt}$ com $r_2=\SI{1}{\ohm}$, ambas
ligadas a $R=\SI{4}{\ohm}$.
\begin{itensq}
\item Determine a tensão comum e as três correntes.
\item Identifique o comportamento de cada fonte.
\end{itensq}
\resposta{$V=\SI{8}{\volt}$; $I_1=\SI{4}{\ampere}$,
$I_2=\SI{-2}{\ampere}$, $I_R=\SI{2}{\ampere}$ --- a fonte 2 é
\textbf{carregada}}
\resolucao{(a) LKC no nó comum:
\[
\frac{V-12}{1}+\frac{V-6}{1}+\frac{V}{4}=0
\;\Longrightarrow\;
4V-48+4V-24+V=0
\;\Longrightarrow\;
V=\SI{8}{\volt}.
\]
\[
I_1=\frac{12-8}{1}=\SI{4}{\ampere};
\qquad
I_2=\frac{6-8}{1}=\SI{-2}{\ampere};
\qquad
I_R=\frac{8}{4}=\SI{2}{\ampere}.
\]
\textbf{LKC:} $4+(-2)=2$ \checkmark.\\
(b) \textbf{A fonte 2 tem corrente negativa} --- ela não fornece, \emph{recebe}
$\SI{2}{\ampere}$. A fonte 1, mais forte, alimenta a carga \emph{e} carrega a
fonte 2.\\
\textbf{Balanço completo:}
\begin{center}
\begin{tabular}{lr}
\hline
Fonte 1 fornece & $\SI{48}{\watt}$\\
$r_1$ dissipa & $\SI{16}{\watt}$\\
$R$ dissipa & $\SI{16}{\watt}$\\
Fonte 2 absorve (f.e.m.) & $\SI{12}{\watt}$\\
$r_2$ dissipa & $\SI{4}{\watt}$\\
\hline
\end{tabular}
\end{center}
Total consumido: $16+16+12+4=\SI{48}{\watt}$ \checkmark.\\
\textbf{Se a fonte 2 for recarregável}, os $\SI{12}{\watt}$ são armazenados. Se
não for, viram calor --- e o dispositivo aquece perigosamente.}
\end{questao}

\begin{questao}[3]{}
Um banco é formado por quatro pilhas de $\SI{1.5}{\volt}$ e
$r=\SI{0.2}{\ohm}$.
\begin{itensq}
\item Compare série e paralelo quanto a $\EMF$, $r$ e corrente de
      curto.
\item Determine a potência máxima disponível em cada arranjo.
\item Estabeleça o critério de escolha.
\end{itensq}
\resposta{Série: $\SI{6}{\volt}$/$\SI{0.8}{\ohm}$; paralelo:
$\SI{1.5}{\volt}$/$\SI{0.05}{\ohm}$; $P_{\max}=\SI{11.25}{\watt}$ nos
\textbf{dois}}
\resolucao{(a)
\begin{center}
\begin{tabular}{lccc}
\hline
Arranjo & $\EMF$ & $r$ & $I_{cc}$\\
\hline
Série & $\SI{6}{\volt}$ & $\SI{0.8}{\ohm}$ & $\SI{7.5}{\ampere}$\\
Paralelo & $\SI{1.5}{\volt}$ & $\SI{0.05}{\ohm}$ & $\SI{30}{\ampere}$\\
\hline
\end{tabular}
\end{center}
(b)
\[
\text{Série: } \frac{6^{2}}{4(0{,}8)}=\SI{11.25}{\watt};
\qquad
\text{Paralelo: } \frac{1{,}5^{2}}{4(0{,}05)}=\SI{11.25}{\watt}.
\]
\textbf{Idênticas} --- e não por acaso. Ambos os arranjos usam as mesmas quatro
pilhas, e a energia total disponível é a mesma. A topologia decide \emph{como}
ela é entregue, não \emph{quanto}.\\
(c) \textbf{Critério de escolha.}
\begin{itemize}
\item \textbf{Série} para cargas de \emph{alta} resistência, que precisam de
tensão. A carga ótima é $\SI{0.8}{\ohm}$.
\item \textbf{Paralelo} para cargas de \emph{baixa} resistência, que precisam de
corrente. A carga ótima é $\SI{0.05}{\ohm}$.
\item \textbf{Série-paralelo} (dois ramos de duas) daria $\SI{3}{\volt}$ e
$\SI{0.2}{\ohm}$ --- e, de novo, $\SI{11.25}{\watt}$.
\end{itemize}
\textbf{Regra:} escolha o arranjo cuja resistência interna se aproxime da carga
--- ou, mais realisticamente, cuja \emph{tensão} atenda ao que a carga exige.}
\end{questao}

\begin{questao}[3]{}
Aplique a LKC a uma \textbf{superfície fechada} que envolva dois nós de um
circuito, e explique por que o resultado é válido.
\begin{itensq}
\item Enuncie a forma generalizada.
\item Justifique a partir da conservação de carga.
\item Dê um exemplo de uso.
\end{itensq}
\resposta{A soma das correntes que atravessam qualquer superfície fechada é
nula}
\resolucao{(a) \textbf{LKC generalizada:} para qualquer superfície fechada que
corte o circuito, a soma algébrica das correntes que a atravessam é zero.\\
(b) \textbf{Justificativa:} a carga total dentro da superfície não pode variar
(em regime permanente). Como corrente é fluxo de carga, o que entra deve
igualar o que sai --- exatamente o mesmo argumento que vale para um nó, agora
aplicado a uma região que \emph{contém} vários nós.\\
Formalmente, escrevendo a LKC para cada nó interno e somando, as correntes dos
ramos \emph{internos} aparecem duas vezes com sinais opostos e se cancelam;
sobram apenas as que cruzam a fronteira.\\
(c) \textbf{Exemplo de uso.} Considere um bloco de circuito ligado ao restante
por três fios apenas. Se dois deles conduzem $\SI{5}{\ampere}$ e
$\SI{3}{\ampere}$ para dentro, o terceiro conduz necessariamente
$\SI{8}{\ampere}$ para fora --- \textbf{sem que seja preciso conhecer nada do
interior do bloco}, que pode conter dezenas de componentes.\\
\textbf{Aplicação prática:} é o que permite verificar a corrente de um circuito
inteiro medindo apenas os condutores de entrada, e é também o princípio do
\emph{dispositivo diferencial residual}, que compara fase e neutro --- qualquer
diferença indica corrente escapando por um terceiro caminho.}
\end{questao}

\begin{questao}[3]{}
Uma fonte foi ensaiada sob diferentes cargas, obtendo-se:
\begin{center}
\begin{tabular}{cccc}
\hline
$I$ ($\si{\ampere}$) & $0{,}50$ & $1{,}00$ & $2{,}00$\\
$V$ ($\si{\volt}$) & $11{,}75$ & $11{,}50$ & $11{,}00$\\
\hline
\end{tabular}
\end{center}
\begin{itensq}
\item Determine $\EMF$ e $r$.
\item Verifique a linearidade do modelo.
\item Estime a corrente de curto e discuta a extrapolação.
\end{itensq}
\resposta{$\EMF=\SI{12}{\volt}$, $r=\SI{0.5}{\ohm}$;
$I_{cc}=\SI{24}{\ampere}$}
\resolucao{(a) O modelo $V=\EMF-rI$ é uma reta. De dois pontos quaisquer:
\[
r=-\frac{\Delta V}{\Delta I}=-\frac{11{,}00-11{,}75}{2{,}00-0{,}50}
=\frac{0{,}75}{1{,}50}=\SI{0.5}{\ohm};
\]
\[
\EMF=11{,}75+0{,}5(0{,}50)=\SI{12}{\volt}.
\]
(b) \textbf{Verificação com o terceiro ponto:}
$12-0{,}5(1{,}00)=\SI{11.50}{\volt}$ \checkmark. Os três pontos são
\emph{exatamente} colineares --- o modelo linear descreve bem a fonte nessa
faixa.\\
\textbf{Nota de método:} com apenas dois pontos, sempre passa uma reta e a
linearidade não pode ser testada. O terceiro ponto é o que transforma o ajuste
em \emph{verificação}.\\
(c) $I_{cc}=\EMF/r=12/0{,}5=\SI{24}{\ampere}$.\\
\textbf{Cautela com a extrapolação:} o ensaio cobriu de $0{,}5$ a
$\SI{2}{\ampere}$, e o valor de curto está \emph{doze vezes} além do maior
ponto medido. Fontes reais aquecem, $r$ cresce e proteções atuam --- os
$\SI{24}{\ampere}$ são um \textbf{limite superior}, útil para dimensionar
proteção, não uma previsão.}
\end{questao}

\begin{questao}[3]{}
Um circuito tem $n=6$ nós, $b=9$ ramos e $2$ fontes de corrente ideais.
\begin{itensq}
\item Determine o número de equações de Kirchhoff necessárias.
\item Explique como as fontes de corrente reduzem o trabalho.
\item Compare com a contagem para análise nodal.
\end{itensq}
\resposta{$5$ LKC $+$ $4$ LKT $=9$; as fontes de corrente fixam duas correntes
de ramo, restando $7$ incógnitas}
\resolucao{(a) LKC: $n-1=5$. LKT: $b-n+1=9-6+1=4$. Total: $9$ --- igual ao
número de correntes de ramo.\\
(b) \textbf{Cada fonte de corrente ideal fixa diretamente a corrente de seu
ramo.} Duas delas eliminam duas incógnitas, restando $7$. Em compensação, a
tensão sobre cada fonte de corrente é desconhecida, o que impede escrever a LKT
de malhas que a atravessem --- é preciso escolher malhas que as contornem, ou
tratar essas tensões como novas incógnitas.\\
\textbf{Resultado líquido:} sistema $7\times7$ em vez de $9\times9$.\\
(c) \textbf{Análise nodal:} $n-1=5$ equações, com as fontes de corrente entrando
\emph{diretamente} no vetor independente, sem nenhuma manipulação. Sistema
$5\times5$.\\
\textbf{Conclusão:} a análise nodal é claramente superior aqui. Ela é, na
verdade, uma forma organizada de aplicar a LKC com a LKT já embutida na escolha
dos potenciais --- e por isso produz o sistema menor sempre que há fontes de
corrente.}
\end{questao}

\begin{questao}[3]{}
Uma fonte tem resistência interna que \emph{cresce} com a corrente, segundo
$r=0{,}2+0{,}1\,I$ (com $r$ em $\si{\ohm}$ e $I$ em $\si{\ampere}$), com
$\EMF=\SI{6}{\volt}$.
\begin{itensq}
\item Determine a corrente para uma carga de $\SI{1}{\ohm}$.
\item Compare com a previsão do modelo linear com $r=\SI{0.2}{\ohm}$.
\end{itensq}
\resposta{$I=\SI{3.80}{\ampere}$ contra $\SI{5}{\ampere}$ previstos}
\resolucao{(a) A LKT dá
\[
6=(0{,}2+0{,}1I)I+1\cdot I
\;\Longrightarrow\;
0{,}1I^{2}+1{,}2I-6=0.
\]
\[
I=\frac{-1{,}2+\sqrt{1{,}44+2{,}4}}{0{,}2}
=\frac{-1{,}2+1{,}960}{0{,}2}=\SI{3.80}{\ampere}.
\]
\emph{(Verificação: $r=0{,}2+0{,}380=\SI{0.58}{\ohm}$ e
$6/(0{,}58+1)=\SI{3.80}{\ampere}$ \checkmark.)}\\
(b) \textbf{Modelo linear:} $I=6/1{,}2=\SI{5}{\ampere}$ --- \textbf{$32\%$
acima} do valor real.\\
\textbf{Duas lições.}\\
\emph{(i)} As \textbf{leis de Kirchhoff continuam válidas} --- elas não exigem
linearidade. O que muda é que a equação resultante deixa de ser linear e passa a
exigir solução de segundo grau.\\
\emph{(ii)} O modelo de $r$ constante é uma \emph{aproximação de pequenos
sinais}. Ele descreve bem a fonte perto do ponto onde $r$ foi medido, e falha
progressivamente longe dele. Baterias reais comportam-se assim, e é por isso
que fabricantes especificam $r$ para uma corrente de ensaio definida.}
\end{questao}

\begin{questao}[3]{}
Uma bateria de $\EMF=\SI{12}{\volt}$ e $r=\SI{0.5}{\ohm}$ deve entregar tensão
terminal de pelo menos $\SI{11.4}{\volt}$.
\begin{itensq}
\item Determine a corrente máxima admissível.
\item Determine a menor carga que pode ser conectada.
\item Determine a potência entregue nessa condição.
\end{itensq}
\resposta{$I\le\SI{1.2}{\ampere}$; $R\ge\SI{9.5}{\ohm}$;
$P=\SI{13.68}{\watt}$}
\resolucao{(a) De $V=\EMF-rI\ge11{,}4$:
\[
0{,}5I\le0{,}6
\;\Longrightarrow\;
I\le\SI{1.2}{\ampere}.
\]
(b) $R=V/I=11{,}4/1{,}2=\SI{9.5}{\ohm}$ --- cargas menores exigiriam mais
corrente e derrubariam a tensão.\\
(c) $P=VI=11{,}4(1{,}2)=\SI{13.68}{\watt}$, com rendimento
$9{,}5/10=95\%$.\\
\textbf{Contraste com a máxima transferência:} a potência máxima ocorreria em
$R=\SI{0.5}{\ohm}$, valendo $\SI{72}{\watt}$ --- mas com tensão terminal de
apenas $\SI{6}{\volt}$, muito abaixo do requisito.\\
\textbf{A restrição de tensão é quase sempre a que manda} em eletrônica, porque
circuitos alimentados simplesmente não funcionam abaixo de sua tensão mínima. A
saída de engenharia é reduzir $r$, não aceitar tensão menor.}
\end{questao}

\begin{questao}[3]{}
Um circuito contém uma fonte \textbf{dependente} de tensão de valor
$3\,I_x$ (em volts, com $I_x$ em ampères) em série com $\SI{2}{\ohm}$, além de
uma fonte independente de $\SI{20}{\volt}$ e um resistor de $\SI{8}{\ohm}$, em
malha única. A corrente de controle $I_x$ é a própria corrente da malha.
\begin{itensq}
\item Escreva a LKT e resolva.
\item Explique por que Kirchhoff continua válido.
\end{itensq}
\resposta{$I=\SI{1.54}{\ampere}$ (fonte dependente em oposição)}
\resolucao{(a) Percorrendo a malha, com a fonte dependente \emph{opondo-se}:
\[
20=2I+8I+3I=13I
\;\Longrightarrow\;
I=\frac{20}{13}=\SI{1.54}{\ampere}.
\]
\emph{(Se a fonte dependente auxiliasse, viria $20=10I-3I=7I$ e
$I=\SI{2.86}{\ampere}$ --- o sentido da fonte dependente altera
completamente o resultado, e por isso a convenção de polaridade precisa estar
explícita no esquema.)}\\
(b) \textbf{As leis de Kirchhoff decorrem de conservação de carga e de
energia}, e nenhuma dessas conservações depende de os elementos serem lineares
ou independentes. A LKT continua sendo "a soma das quedas em uma malha é nula";
a LKC continua sendo "o que entra em um nó, sai".\\
\textbf{O que muda:} a fonte dependente introduz uma \textbf{equação de
restrição} adicional, ligando seu valor a uma variável do circuito. Essa
equação é resolvida \emph{junto} com as de Kirchhoff.\\
\textbf{Consequência estrutural:} o termo migra para o lado das incógnitas, o
que pode tornar a matriz assimétrica --- como se observa na análise nodal e na
de malhas. A reciprocidade se perde; as leis, não.}
\end{questao}

\begin{questao}[3]{}
Uma fonte alimenta uma carga através de um cabo de $\SI{0.3}{\ohm}$ (ida e
volta). A fonte tem $\EMF=\SI{24}{\volt}$ e $r=\SI{0.2}{\ohm}$; a carga vale
$\SI{9.5}{\ohm}$.
\begin{itensq}
\item Determine a corrente, a tensão na carga e a tensão nos terminais da
      fonte.
\item Determine onde a energia é perdida.
\end{itensq}
\resposta{$I=\SI{2.4}{\ampere}$; $V_{carga}=\SI{22.8}{\volt}$;
$V_{fonte}=\SI{23.52}{\volt}$}
\resolucao{(a)
\[
I=\frac{24}{0{,}2+0{,}3+9{,}5}=\frac{24}{10}=\SI{2.4}{\ampere};
\]
\[
V_{carga}=9{,}5(2{,}4)=\SI{22.8}{\volt};
\qquad
V_{fonte}=24-0{,}2(2{,}4)=\SI{23.52}{\volt}.
\]
A diferença de $\SI{0.72}{\volt}$ entre os dois é exatamente a queda no cabo.\\
(b) \textbf{Distribuição das perdas:}
\begin{center}
\begin{tabular}{lrr}
\hline
Elemento & Potência & Fração\\
\hline
Carga ($\SI{9.5}{\ohm}$) & $\SI{54.72}{\watt}$ & $95{,}0\%$\\
Cabo ($\SI{0.3}{\ohm}$) & $\SI{1.73}{\watt}$ & $3{,}0\%$\\
Interna ($\SI{0.2}{\ohm}$) & $\SI{1.15}{\watt}$ & $2{,}0\%$\\
\hline
Total & $\SI{57.6}{\watt}$ & $100\%$\\
\hline
\end{tabular}
\end{center}
\textbf{Observação de diagnóstico:} medindo apenas nos terminais da fonte
($\SI{23.52}{\volt}$), a instalação pareceria melhor do que é. A tensão que
importa é a que chega à \emph{carga} --- e ela está $5\%$ abaixo da f.e.m.\\
\textbf{Regra de medição:} sempre meça no ponto de uso, não na fonte. As duas
leituras diferem exatamente pela queda no caminho.}
\end{questao}

\blocodesafio

\begin{questao}[4]{}
\textbf{(Ensaio de caracterização.)} Descreva como obter $\EMF$ e $r$ de uma
fonte por ensaios em vazio e em curto-circuito.
\begin{itensq}
\item Deduza as relações.
\item Analise os riscos do ensaio em curto.
\item Proponha uma alternativa segura e compare.
\end{itensq}
\resposta{$\EMF=V_{oc}$ e $r=V_{oc}/I_{cc}$; o curto é destrutivo em fontes de
baixa impedância}
\resolucao{(a) \textbf{Em vazio} ($I=0$): $V_{oc}=\EMF$ --- a leitura direta,
desde que o voltímetro tenha impedância muito maior que $r$.\\
\textbf{Em curto} ($V=0$): $\EMF=rI_{cc}$, logo
\[
r=\frac{V_{oc}}{I_{cc}}.
\]
(b) \textbf{Riscos do curto.}
\begin{itemize}
\item Uma pilha AA de $\SI{1.5}{\volt}$ com $r=\SI{0.25}{\ohm}$ entrega
$\SI{6}{\ampere}$ e dissipa $\SI{9}{\watt}$ internamente --- aquecimento
rápido e vazamento.
\item Uma bateria de automóvel ($r\approx\SI{5}{\milli\ohm}$) chega a
$\SI{2500}{\ampere}$, com centelhamento, fusão de terminais e risco de
explosão por ignição do hidrogênio liberado.
\item Baterias de lítio podem entrar em fuga térmica e incendiar.
\item Além disso, $r$ \emph{aumenta} durante o curto, de modo que a leitura
subestima a corrente inicial e superestima $r$.
\end{itemize}
(c) \textbf{Alternativa segura: dois pontos de carga.} Meça $(I_1,V_1)$ e
$(I_2,V_2)$ com resistências conhecidas e moderadas, e ajuste a reta:
\[
r=-\frac{V_2-V_1}{I_2-I_1};
\qquad
\EMF=V_1+rI_1.
\]
\textbf{Comparação.} O método de dois pontos é seguro, não perturba a fonte, e
com três ou mais pontos ainda \emph{testa} a linearidade do modelo --- coisa que
o ensaio de curto não faz. O ensaio em curto só se justifica em fontes
projetadas para suportá-lo, com limitação eletrônica de corrente.}
\end{questao}

\begin{questao}[4]{}
\textbf{(Demonstração --- balanço de potência.)} Prove que a soma algébrica das
potências de todos os elementos de um circuito é nula.
\begin{itensq}
\item Estabeleça a convenção de sinais.
\item Demonstre a partir das leis de Kirchhoff.
\item Explique por que o resultado independe da natureza dos elementos.
\end{itensq}
\resposta{$\sum_k v_ki_k=0$ --- consequência direta de LKC e LKT}
\resolucao{(a) \textbf{Convenção do receptor:} para cada ramo $k$, adote a
corrente $i_k$ entrando pelo terminal de maior potencial. A potência
$p_k=v_ki_k$ é então \emph{positiva} se o elemento absorve e \emph{negativa} se
fornece.\\
(b) \textbf{Demonstração.} Escreva cada tensão de ramo como diferença de
potenciais nodais: $v_k=V_a-V_b$, onde $a$ e $b$ são os nós do ramo $k$. Então
\[
\sum_k p_k=\sum_k (V_a-V_b)\,i_k.
\]
Reagrupando por \emph{nó} em vez de por ramo, cada potencial $V_n$ aparece
multiplicado pela soma algébrica das correntes que chegam ao nó $n$:
\[
\sum_k p_k=\sum_n V_n\left(\sum_{k\in n}\pm i_k\right)=\sum_n V_n\cdot 0=0,
\]
pois o parêntese é exatamente a \textbf{LKC} do nó $n$.\\
\emph{(A escrita $v_k=V_a-V_b$ já embute a \textbf{LKT}: ela pressupõe que o
potencial é bem definido em cada nó.)}\\
(c) \textbf{Independência da natureza dos elementos.} A demonstração usou
\emph{apenas} as duas leis de Kirchhoff e a definição $p=vi$. Em nenhum momento
foi preciso saber a relação entre $v_k$ e $i_k$ de cada elemento.\\
Portanto o resultado vale para resistores, fontes, diodos, capacitores,
indutores, elementos ativos --- lineares ou não, com ou sem memória. É o
\textbf{teorema de Tellegen}, e sua generalidade é notável: ele é uma
propriedade da \emph{topologia} do circuito, não dos componentes.\\
\textbf{Uso prático:} verificar o balanço de potência é o teste de consistência
mais completo de qualquer solução --- se ele não fecha, há erro em algum lugar.}
\end{questao}

\begin{questao}[4]{}
\textbf{(Formulação com fonte ideal entre dois nós.)} Um circuito tem dois nós
não referenciais, $A$ e $B$, com uma fonte ideal de tensão $\EMF$ ligada
diretamente entre eles.
\begin{itensq}
\item Explique por que a LKC não pode ser escrita isoladamente em $A$ nem em
      $B$.
\item Formule o sistema corretamente.
\item Determine a corrente na fonte após resolver.
\end{itensq}
\resposta{Usar supernó: uma LKC envolvendo $A$ e $B$, mais a restrição
$V_A-V_B=\EMF$}
\resolucao{(a) A corrente que atravessa uma fonte de tensão \textbf{ideal} é
determinada pelo circuito, não pela fonte --- não existe relação
$i=f(v)$ para ela. Ao escrever a LKC em $A$, essa corrente desconhecida aparece
como uma incógnita a mais, e o mesmo ocorre em $B$. Duas equações, três
incógnitas.\\
(b) \textbf{Formulação por supernó.} Englobe $A$ e $B$ em uma única superfície
fechada. Pela LKC generalizada, a corrente interna da fonte \emph{se cancela}
--- ela entra em um nó e sai do outro:
\[
\sum_{\text{ramos que cruzam a fronteira}} i=0.
\]
Essa é \emph{uma} equação. A segunda vem da própria fonte:
\[
V_A-V_B=\EMF.
\]
Duas equações, duas incógnitas ($V_A$ e $V_B$). \textbf{Sistema fechado.}\\
(c) \textbf{Corrente na fonte:} obtida \emph{depois}, aplicando a LKC a
\emph{um} dos dois nós isoladamente. Com $V_A$ e $V_B$ já conhecidos, todas as
demais correntes daquele nó são calculáveis, e a da fonte é o que falta para
fechar o balanço.\\
\textbf{Observação:} se um dos terminais da fonte for o próprio nó de
referência, o problema desaparece --- o outro potencial fica \emph{conhecido} e
deixa de ser incógnita. \textbf{Escolher o terra em um terminal de fonte é a
simplificação mais barata que existe.}}
\end{questao}

\begin{questao}[4]{}
\textbf{(Projeto de verificação experimental.)} Elabore um ensaio de laboratório
para verificar a LKC em um nó com três ramos, considerando incertezas dos
instrumentos.
\begin{itensq}
\item Descreva o procedimento.
\item Propague as incertezas e estabeleça o critério de aceitação.
\item Interprete um resíduo de $\SI{0.35}{\ampere}$ em correntes de alguns
      ampères.
\end{itensq}
\resposta{Critério: resíduo compatível com $u_c\approx\sqrt{\sum u_k^{2}}$;
$\SI{0.35}{\ampere}$ indica erro real}
\resolucao{(a) \textbf{Procedimento.} Meça as três correntes
\emph{simultaneamente}, com três amperímetros --- ou, se houver apenas um,
faça as três medições sem alterar o circuito entre elas. Registre sentidos e
adote uma convenção única (por exemplo, positivo entrando).\\
\textbf{Cuidado essencial:} inserir o amperímetro \emph{altera} o circuito, pois
acrescenta sua resistência interna ao ramo. Com três inserções sucessivas, o
circuito muda três vezes e a verificação perde sentido. Use amperímetro de baixa
resistência ou, melhor, alicate amperímetro, que não abre o ramo.\\
(b) \textbf{Propagação.} Suponha
$I_1=(5{,}02\pm0{,}05)\,\si{\ampere}$,
$I_2=(3{,}01\pm0{,}03)\,\si{\ampere}$ entrando, e $I_3$ saindo. A previsão é
$I_1+I_2=\SI{8.03}{\ampere}$, com incerteza combinada
\[
u_c=\sqrt{0{,}05^{2}+0{,}03^{2}}=\SI{0.058}{\ampere}.
\]
Incluindo a incerteza da própria medição de $I_3$ (digamos
$\SI{0.08}{\ampere}$), o resíduo esperado tem incerteza
$\sqrt{0{,}058^{2}+0{,}08^{2}}=\SI{0.099}{\ampere}$.\\
\textbf{Critério de aceitação:} $|I_1+I_2-I_3|\le2u\approx
\SI{0.20}{\ampere}$ (dois desvios, cerca de $95\%$ de confiança).\\
(c) \textbf{Resíduo de $\SI{0.35}{\ampere}$} é $3{,}5$ vezes a incerteza
combinada --- \textbf{não é compatível com erro aleatório}. Causas prováveis, em
ordem:
\begin{itemize}
\item \textbf{Erro de sinal ou de sentido} em uma das leituras --- o mais comum.
\item \textbf{Um quarto ramo não contabilizado}, inclusive uma fuga para o terra
ou pela bancada.
\item \textbf{Instrumento descalibrado} ou operando fora de escala.
\item \textbf{Circuito alterado entre medições} (aquecimento, contato
intermitente).
\end{itemize}
\textbf{Conclusão importante:} a LKC \emph{não} pode ser violada --- ela é
conservação de carga. Um resíduo significativo é sempre erro de medição ou de
modelo, nunca da lei. O ensaio, portanto, verifica o \emph{procedimento}, e não
a física.}
\end{questao}

\begin{questao}[4]{}
\textbf{(Baterias em série desequilibradas.)} Quatro células de
$\SI{1.5}{\volt}$ estão em série alimentando uma carga, mas uma delas está
gasta, com $\EMF=\SI{0.9}{\volt}$ e $r=\SI{1.5}{\ohm}$ (as demais mantêm
$r=\SI{0.2}{\ohm}$). A carga vale $\SI{5}{\ohm}$.
\begin{itensq}
\item Determine a corrente e a tensão em cada célula.
\item Identifique o que ocorre com a célula gasta.
\item Discuta a implicação prática.
\end{itensq}
\resposta{$I=\SI{0.76}{\ampere}$; a célula gasta apresenta
$\SI{-0.24}{\volt}$ --- está sendo \textbf{invertida}}
\resolucao{(a)
\[
\EMF_{total}=3(1{,}5)+0{,}9=\SI{5.4}{\volt};
\qquad
r_{total}=3(0{,}2)+1{,}5=\SI{2.1}{\ohm};
\]
\[
I=\frac{5{,}4}{2{,}1+5}=\frac{5{,}4}{7{,}1}=\SI{0.76}{\ampere}.
\]
\textbf{Tensão de cada célula boa:} $1{,}5-0{,}2(0{,}76)=\SI{1.35}{\volt}$.\\
\textbf{Tensão da célula gasta:}
$0{,}9-1{,}5(0{,}76)=0{,}9-1{,}14=\SI{-0.24}{\volt}$.\\
(b) \textbf{A tensão da célula gasta é negativa} --- ou seja, ela está sendo
percorrida por corrente no sentido de \emph{carga}, forçada pelas três células
boas. Ela deixou de ser fonte e passou a ser carga, dissipando
$1{,}5(0{,}76)^{2}=\SI{0.87}{\watt}$ internamente.\\
\textbf{Em células não recarregáveis}, isso provoca a chamada \textbf{inversão
de polaridade}: reações internas indesejadas, geração de gás, aquecimento e
vazamento de eletrólito --- exatamente o que se observa quando uma pilha
"vaza" em um controle remoto esquecido ligado.\\
(c) \textbf{Implicações práticas.}
\begin{itemize}
\item \textbf{Nunca misture pilhas novas e usadas}, nem de marcas ou tecnologias
diferentes, em série. A mais fraca será invertida pelas demais.
\item \textbf{Substitua sempre o conjunto completo.}
\item Em baterias de lítio em série, o problema é grave a ponto de exigir
\textbf{circuito de balanceamento} (BMS), que monitora cada célula
individualmente e impede que qualquer uma seja levada à inversão.
\end{itemize}
\textbf{Diagnóstico:} a célula gasta domina o conjunto por sua \emph{resistência
interna}, não por sua f.e.m. --- $\SI{1.5}{\ohm}$ contra
$\SI{0.6}{\ohm}$ das outras três somadas.}
\end{questao}

\begin{questao}[4]{}
\textbf{(Modelagem --- fonte com limitação de corrente.)} Uma fonte de bancada
tem $\EMF=\SI{20}{\volt}$ e $r=\SI{0.1}{\ohm}$, com limitação eletrônica em
$\SI{2}{\ampere}$.
\begin{itensq}
\item Descreva a característica $V\times I$ completa.
\item Determine o ponto de operação para cargas de $\SI{50}{\ohm}$,
      $\SI{10}{\ohm}$ e $\SI{2}{\ohm}$.
\item Explique por que o modelo de fonte real deixa de valer.
\end{itensq}
\resposta{Duas regiões: fonte de tensão até $\SI{2}{\ampere}$, fonte de corrente
depois}
\resolucao{(a) \textbf{Característica em dois trechos.}
\begin{itemize}
\item \textbf{Região de tensão constante} ($I<\SI{2}{\ampere}$):
$V=20-0{,}1I$ --- reta quase horizontal.
\item \textbf{Região de corrente constante} ($V<\SI{19.8}{\volt}$):
$I=\SI{2}{\ampere}$ --- reta vertical.
\end{itemize}
O "joelho" fica em $(\SI{2}{\ampere};\ \SI{19.8}{\volt})$, com
$R_{joelho}=\SI{9.9}{\ohm}$.\\
(b)
\begin{center}
\begin{tabular}{lccl}
\hline
$R_L$ & $I$ & $V$ & Região\\
\hline
$\SI{50}{\ohm}$ & $\SI{0.399}{\ampere}$ & $\SI{19.96}{\volt}$ & tensão\\
$\SI{10}{\ohm}$ & $\SI{1.98}{\ampere}$ & $\SI{19.80}{\volt}$ & tensão (limite)\\
$\SI{2}{\ohm}$ & $\SI{2.00}{\ampere}$ & $\SI{4.00}{\volt}$ & corrente\\
\hline
\end{tabular}
\end{center}
Note o terceiro caso: o modelo linear preveria
$20/2{,}1=\SI{9.5}{\ampere}$ --- quase cinco vezes o real.\\
(c) \textbf{Por que o modelo falha.} A fonte real é um modelo \emph{linear} de
dois parâmetros, e sua característica é uma \emph{única} reta. A limitação
eletrônica introduz um segundo trecho, tornando a característica
\textbf{quebrada} --- e nenhuma reta descreve as duas regiões.\\
\textbf{Como proceder corretamente:}
\begin{itemize}
\item Determinar em qual região a carga opera, comparando $R_L$ com
$R_{joelho}=\SI{9.9}{\ohm}$.
\item $R_L>R_{joelho}$: use o modelo de \emph{fonte de tensão} $(\EMF,r)$.
\item $R_L<R_{joelho}$: use o modelo de \emph{fonte de corrente} ideal de
$\SI{2}{\ampere}$, e a tensão sai de $V=R_LI$.
\end{itemize}
\textbf{As leis de Kirchhoff seguem válidas nas duas regiões} --- o que muda é
apenas a relação constitutiva da fonte. É o mesmo procedimento usado com
diodos: identificar o trecho e aplicar o modelo correspondente.}
\end{questao}

\begin{questao}[4]{}
\textbf{(Sistema completo por Kirchhoff.)} Um circuito de duas malhas tem
$\EMF_1=\SI{18}{\volt}$ com $r_1=\SI{1}{\ohm}$ no primeiro ramo,
$\EMF_2=\SI{12}{\volt}$ com $r_2=\SI{2}{\ohm}$ no segundo, e
$R=\SI{6}{\ohm}$ no ramo comum.
\begin{itensq}
\item Escreva as equações de Kirchhoff e resolva.
\item Verifique o balanço de potência.
\item Identifique o papel de cada fonte.
\end{itensq}
\resposta{$I_1=\SI{3.6}{\ampere}$, $I_2=\SI{-1.2}{\ampere}$,
$I_R=\SI{2.4}{\ampere}$}
\resolucao{(a) Chamando $V$ a tensão do nó comum:
\[
\frac{V-18}{1}+\frac{V-12}{2}+\frac{V}{6}=0.
\]
Multiplicando por $6$: $6V-108+3V-36+V=0\Rightarrow10V=144
\Rightarrow V=\SI{14.4}{\volt}$.
\[
I_1=\frac{18-14{,}4}{1}=\SI{3.6}{\ampere};
\quad
I_2=\frac{12-14{,}4}{2}=\SI{-1.2}{\ampere};
\quad
I_R=\frac{14{,}4}{6}=\SI{2.4}{\ampere}.
\]
\textbf{LKC:} $3{,}6-1{,}2=2{,}4$ \checkmark.\\
(b) \textbf{Balanço.}
\begin{center}
\begin{tabular}{lr}
\hline
$\EMF_1$ fornece $18(3{,}6)$ & $\SI{64.8}{\watt}$\\
$r_1$ dissipa $1(3{,}6)^{2}$ & $\SI{12.96}{\watt}$\\
$R$ dissipa $6(2{,}4)^{2}$ & $\SI{34.56}{\watt}$\\
$\EMF_2$ absorve $12(1{,}2)$ & $\SI{14.4}{\watt}$\\
$r_2$ dissipa $2(1{,}2)^{2}$ & $\SI{2.88}{\watt}$\\
\hline
Total consumido & $\SI{64.8}{\watt}$\\
\hline
\end{tabular}
\end{center}
Fecha exatamente \checkmark.\\
(c) \textbf{Papéis.} A fonte 1 é a única \emph{geradora}: ela alimenta a carga
$R$ \emph{e} carrega a fonte 2. A fonte 2 comporta-se como \textbf{receptor},
porque sua f.e.m. ($\SI{12}{\volt}$) é menor que a tensão do nó
($\SI{14.4}{\volt}$) --- e corrente sempre entra por onde o potencial externo é
maior.\\
\textbf{Critério geral:} compare a f.e.m. de cada fonte com a tensão do nó a
que está ligada. Maior $\Rightarrow$ fornece; menor $\Rightarrow$ recebe. O
sinal da corrente calculada confirma automaticamente.}
\end{questao}

% END BANCO COMPLEMENTAR

\gabarito[2.4 Fontes reais e Leis de Kirchhoff]
     % Ohm, Potencia, Fontes, Kirchhoff

% Cap. 3 -------------------------------------------------------------
\setcounter{section}{2}
% =====================================================================
%  CAPITULO 3 - ASSOCIACAO DE RESISTORES E TRANSFORMACOES
%  (versao revisada e ampliada: 26 -> 48 questoes)
% =====================================================================
\section{Associação de Resistores e Transformações}

\begin{conceito}[Antes de começar]
\textbf{Série} --- a corrente é a \emph{mesma} em todos os elementos e as tensões
se somam:
\[
\Req = R_1+R_2+\dots+R_n ,\qquad U = U_1+U_2+\dots+U_n .
\]
\textbf{Paralelo} --- a tensão é a \emph{mesma} em todos os ramos e as correntes
se somam:
\[
\frac{1}{\Req}=\frac{1}{R_1}+\dots+\frac{1}{R_n},\qquad
\text{dois resistores: }\ \Req=\frac{R_1R_2}{R_1+R_2},\qquad
\text{$n$ iguais: }\ \Req=\frac{R}{n}.
\]
\textbf{Transformação $\Delta \to Y$} (resistores do triângulo no denominador comum
$S=R_{AB}+R_{BC}+R_{CA}$):
\[
R_A=\frac{R_{AB}\,R_{CA}}{S},\qquad
R_B=\frac{R_{AB}\,R_{BC}}{S},\qquad
R_C=\frac{R_{BC}\,R_{CA}}{S}.
\]
\textbf{Transformação $Y \to \Delta$} (com $P=R_AR_B+R_BR_C+R_CR_A$):
\[
R_{AB}=\frac{P}{R_C},\qquad R_{BC}=\frac{P}{R_A},\qquad R_{CA}=\frac{P}{R_B}.
\]
\end{conceito}

\begin{atencao}[Dois atalhos que economizam muito tempo]
\textbf{Curto-circuito:} um fio ideal em paralelo com um resistor \emph{anula} esse
resistor ($\Req = 0$). \quad
\textbf{Circuito aberto:} um ramo interrompido não conduz e pode ser apagado do
desenho. Antes de calcular, procure por esses dois casos.
\end{atencao}

\legendaniveis

% =====================================================================
\subsection{Circuitos em série}
% =====================================================================

\blocofixacao

\begin{questao}[1]{}
Em uma associação de resistores \textbf{em série}, é correto afirmar que:
\begin{alternativas}
\item a tensão é a mesma em todos os resistores.
\item a corrente é a mesma em todos os resistores.
\item a resistência equivalente é menor que a menor das resistências.
\item a corrente se divide proporcionalmente às resistências.
\item a potência dissipada é a mesma em todos os resistores.
\end{alternativas}
\resposta{(b)}
\resolucao{Só existe um caminho para a corrente, logo ela é a mesma em todos os
elementos --- (a) e (d) descrevem o circuito paralelo. A resistência equivalente
é a \emph{soma}, portanto maior que a maior delas, o que elimina (c).
A potência $P=R\,I^2$ depende de $R$: só é igual se os resistores forem iguais,
o que elimina (e).}
\end{questao}

\begin{questao}[1]{}
Três resistores de $\SI{4}{\ohm}$, $\SI{6}{\ohm}$ e $\SI{10}{\ohm}$ são ligados em
série a uma fonte de $\SI{40}{\volt}$ de resistência interna desprezível.
Determine:
\begin{itensq}
\item a resistência equivalente;
\item a corrente do circuito;
\item a tensão em cada resistor.
\end{itensq}
\resposta{(a) \SI{20}{\ohm}; (b) \SI{2}{\ampere}; (c) \SI{8}{\volt}, \SI{12}{\volt} e \SI{20}{\volt}}
\resolucao{(a) $\Req = 4+6+10 = \SI{20}{\ohm}$.\\
(b) $I = \dfrac{U}{\Req} = \dfrac{40}{20} = \SI{2}{\ampere}$.\\
(c) $U_1 = 4\cdot 2 = \SI{8}{\volt}$, $U_2 = 6\cdot 2 = \SI{12}{\volt}$,
$U_3 = 10\cdot 2 = \SI{20}{\volt}$.\\
\textbf{Verificação:} $8+12+20 = \SI{40}{\volt}$ --- confere com a tensão da fonte.
Repare que a maior tensão cai sobre o maior resistor.}
\end{questao}

\begin{questao}[1]{}
Para cada circuito, calcule a resistência equivalente e a corrente $I$.

\circuitoq{%
\subcirc{a}{%
\begin{circuitikz}[scale=0.78,transform shape,line width=0.8pt]
 \draw (0,0) to[battery1,l_={$E=\SI{10}{\volt}$}] (0,3)
       to[R,l={$R_1=\SI{3}{\ohm}$}] (3,3)
       to[R,l={$R_2=\SI{2}{\ohm}$}] (3,0) -- (0,0);
 \draw[->,Icol,line width=1pt] (0.7,3.4) -- (1.7,3.4)
       node[midway,above,font=\footnotesize]{$I$};
\end{circuitikz}}
\hfill
\subcirc{b}{%
\begin{circuitikz}[scale=0.78,transform shape,line width=0.8pt]
 \draw (0,0) to[battery1,l_={$E=\SI{15}{\volt}$}] (0,3)
       to[R,l={$R_1=\SI{4}{\ohm}$}] (3,3)
       to[R,l={$R_2=\SI{6}{\ohm}$}] (3,0) -- (0,0);
 \draw[->,Icol,line width=1pt] (0.7,3.4) -- (1.7,3.4)
       node[midway,above,font=\footnotesize]{$I$};
\end{circuitikz}}

\vspace{1.1em}

\subcirc{c}{%
\begin{circuitikz}[scale=0.78,transform shape,line width=0.8pt]
 \draw (0,0) to[battery1,l_={$E=\SI{25}{\volt}$}] (0,3)
       to[R,l={$R_1=\SI{4.5}{\ohm}$}] (3.4,3)
       to[R,l={$R_2=\SI{3}{\ohm}$}] (3.4,0)
       to[R,l={$R_3=\SI{2.5}{\ohm}$}] (0,0);
 \draw[->,Icol,line width=1pt] (0.7,3.4) -- (1.7,3.4)
       node[midway,above,font=\footnotesize]{$I$};
\end{circuitikz}}
\hfill
\subcirc{d}{%
\begin{circuitikz}[scale=0.78,transform shape,line width=0.8pt]
 \draw (0,0) to[battery1,l_={$E=\SI{30}{\volt}$}] (0,4.4)
       to[R,l={$R_1=\SI{2.5}{\ohm}$}] (3.6,4.4)
       to[R,l={$R_2=\SI{4}{\ohm}$}] (3.6,2.2)
       to[R,l={$R_3=\SI{1}{\ohm}$}] (3.6,0)
       to[R,l={$R_4=\SI{2.5}{\ohm}$}] (0,0);
 \draw[->,Icol,line width=1pt] (0.7,4.8) -- (1.7,4.8)
       node[midway,above,font=\footnotesize]{$I$};
\end{circuitikz}}}

\resposta{(a) \SI{5}{\ohm}, \SI{2}{\ampere}; (b) \SI{10}{\ohm}, \SI{1.5}{\ampere};
(c) \SI{10}{\ohm}, \SI{2.5}{\ampere}; (d) \SI{10}{\ohm}, \SI{3}{\ampere}}
\resolucao{Em todos os casos, $\Req$ é a soma das resistências e
$I = E/\Req$:\\
(a) $\Req = 3+2 = \SI{5}{\ohm}$, $I = 10/5 = \SI{2}{\ampere}$.\\
(b) $\Req = 4+6 = \SI{10}{\ohm}$, $I = 15/10 = \SI{1.5}{\ampere}$.\\
(c) $\Req = 4{,}5+3+2{,}5 = \SI{10}{\ohm}$, $I = 25/10 = \SI{2.5}{\ampere}$.\\
(d) $\Req = 2{,}5+4+1+2{,}5 = \SI{10}{\ohm}$, $I = 30/10 = \SI{3}{\ampere}$.}
\end{questao}

\begin{questao}[1]{}
Uma associação em série de $n$ resistores iguais de $\SI{15}{\ohm}$ apresenta
resistência equivalente de $\SI{105}{\ohm}$. O valor de $n$ é:
\begin{alternativas}[3]
\item 3
\item 5
\item 6
\item 7
\item 9
\end{alternativas}
\resposta{(d)}
\resolucao{Para $n$ resistores iguais em série, $\Req = n\,R$, logo
$n = \dfrac{105}{15} = 7$.}
\end{questao}

\blocoaplicacao

\begin{questao}[2]{}
No circuito da figura, duas fontes ideais estão ligadas em \textbf{oposição}.

\circuitoq{%
\begin{circuitikz}[scale=0.9,transform shape,line width=0.8pt]
 \draw (0,0) to[battery1,l_={$E_1=\SI{30}{\volt}$}] (0,3.6)
       to[R,l={$R_1=\SI{2.5}{\ohm}$}] (4.5,3.6)
       to[battery1,l={$E_2=\SI{10}{\volt}$},invert] (4.5,1.8)
       to[R,l={$R_2=\SI{1.5}{\ohm}$}] (4.5,0)
       to[R,l={$R_3=\SI{1}{\ohm}$}] (0,0);
 \draw[->,Icol,line width=1pt] (1.0,4.0) -- (2.2,4.0)
       node[midway,above,font=\footnotesize]{$I$};
\end{circuitikz}}

Determine a corrente do circuito e a tensão sobre $R_1$.
\resposta{$I = \SI{4}{\ampere}$; $U_{R_1} = \SI{10}{\volt}$}
\resolucao{Percorrendo a malha, $E_1$ impulsiona a corrente no sentido adotado e
$E_2$ se opõe. Pela lei das malhas:
\[
E_1 - E_2 = I\,(R_1+R_2+R_3)
\quad\Rightarrow\quad
I = \frac{30-10}{2{,}5+1{,}5+1} = \frac{20}{5} = \SI{4}{\ampere}.
\]
$U_{R_1} = R_1 I = 2{,}5\cdot 4 = \SI{10}{\volt}$.\\
\textbf{Observação:} se $E_2$ fosse maior que $E_1$, a corrente resultaria
negativa --- o que significa apenas que ela circula no sentido oposto ao adotado.
Nesse caso $E_2$ estaria \emph{sendo carregada} por $E_1$.}
\end{questao}

\begin{questao}[2]{}
Uma bateria de força eletromotriz $\EMF = \SI{12}{\volt}$ e resistência interna
$r = \SI{0.5}{\ohm}$ alimenta um resistor $R = \SI{5.5}{\ohm}$.

\circuitoq{%
\begin{circuitikz}[scale=0.9,transform shape,line width=0.8pt]
 \draw (0,0) to[battery1,l_={$\EMF=\SI{12}{\volt}$}] (0,2)
       to[R,l_={$r=\SI{0.5}{\ohm}$},color=Vcol] (0,3.4) -- (3.6,3.4)
       to[R,l={$R=\SI{5.5}{\ohm}$}] (3.6,0) -- (0,0);
 \draw[dashed,cinza] (-1.15,-0.5) rectangle (0.95,3.9);
 \node[cinza,font=\footnotesize\sffamily,above] at (-0.1,3.95) {bateria real};
 \draw[->,Icol,line width=1pt] (1.2,3.75) -- (2.3,3.75)
       node[midway,above,font=\footnotesize]{$I$};
\end{circuitikz}}

Determine:
\begin{itensq}
\item a corrente do circuito;
\item a tensão nos terminais da bateria;
\item a potência dissipada internamente.
\end{itensq}
\resposta{(a) \SI{2}{\ampere}; (b) \SI{11}{\volt}; (c) \SI{2}{\watt}}
\resolucao{(a) A resistência interna está \emph{em série} com a carga:
\[
I=\frac{\EMF}{R+r}=\frac{12}{5{,}5+0{,}5}=\frac{12}{6}=\SI{2}{\ampere}.
\]
(b) $U = \EMF - r\,I = 12 - 0{,}5\cdot 2 = \SI{11}{\volt}$
(ou, equivalentemente, $U = R\,I = 5{,}5\cdot 2 = \SI{11}{\volt}$).\\
(c) $P_r = r\,I^2 = 0{,}5\cdot 2^2 = \SI{2}{\watt}$ --- energia perdida dentro da
própria bateria, que é o que a faz aquecer em curtos-circuitos.}
\end{questao}

\begin{questao}[2]{}
Seis lâmpadas idênticas, de resistência $\SI{40}{\ohm}$ cada, são ligadas
\textbf{em série} a uma fonte de $\SI{120}{\volt}$.
\begin{itensq}
\item Calcule a corrente e a tensão em cada lâmpada.
\item Uma das lâmpadas queima (filamento aberto). O que acontece com as demais?
      Justifique.
\end{itensq}
\resposta{(a) $I=\SI{0.5}{\ampere}$ e $U=\SI{20}{\volt}$ por lâmpada;
(b) todas apagam}
\resolucao{(a) $\Req = 6\cdot 40 = \SI{240}{\ohm}$;
$I = \dfrac{120}{240} = \SI{0.5}{\ampere}$;
$U = 40\cdot 0{,}5 = \SI{20}{\volt}$ em cada lâmpada
(confira: $6\cdot 20 = \SI{120}{\volt}$).\\
(b) O filamento aberto interrompe o \emph{único} caminho da corrente:
$I = 0$ e todas as lâmpadas apagam. É exatamente esse o inconveniente dos antigos
cordões de luzes de Natal ligados em série.}
\end{questao}

\begin{questao}[2]{}
Três resistores estão em série sob $\SI{120}{\volt}$: $R_1=\SI{20}{\ohm}$,
$R_2$ desconhecido e $R_3=\SI{30}{\ohm}$. Deseja-se que a tensão sobre $R_2$
seja de $\SI{20}{\volt}$. O valor de $R_2$ deve ser:
\begin{alternativas}[3]
\item \SI{5}{\ohm}
\item \SI{10}{\ohm}
\item \SI{15}{\ohm}
\item \SI{20}{\ohm}
\item \SI{25}{\ohm}
\end{alternativas}
\resposta{(b)}
\resolucao{Se $U_2 = \SI{20}{\volt}$, sobram $120-20 = \SI{100}{\volt}$ para
$R_1$ e $R_3$, que somam \SI{50}{\ohm}. Como a corrente é a mesma em toda a série:
\[
I=\frac{100}{50}=\SI{2}{\ampere}
\qquad\Rightarrow\qquad
R_2=\frac{U_2}{I}=\frac{20}{2}=\SI{10}{\ohm}.
\]}
\end{questao}

\blocoaprofundamento

\begin{questao}[3]{}
Uma carga de $\SI{9.6}{\ohm}$ é alimentada por uma fonte de $\SI{24}{\volt}$ por
meio de um cabo de dois condutores, cada um com resistência de
$\SI{0.2}{\ohm}$.

\circuitoq{%
\begin{circuitikz}[scale=0.9,transform shape,line width=0.8pt]
 \draw (0,0) to[battery1,l_={$E=\SI{24}{\volt}$}] (0,3)
       to[R,l=$\SI{0.2}{\ohm}$,color=Vcol] (4.4,3)
       to[R,l={$R_c=\SI{9.6}{\ohm}$}] (4.4,0)
       to[R,l=$\SI{0.2}{\ohm}$,color=Vcol] (0,0);
 \node[Vcol,font=\footnotesize\sffamily] at (2.2,3.85) {condutores do cabo};
 \draw[->,Icol,line width=1pt] (0.5,2.55) -- (1.5,2.55)
       node[midway,below,font=\footnotesize]{$I$};
\end{circuitikz}}

Determine:
\begin{itensq}
\item a corrente na linha;
\item a queda de tensão total no cabo;
\item a tensão efetivamente entregue à carga;
\item o rendimento da transmissão, $\eta = P_{\text{carga}}/P_{\text{total}}$.
\end{itensq}
\resposta{(a) \SI{2.4}{\ampere}; (b) \SI{0.96}{\volt}; (c) \SI{23.04}{\volt};
(d) \SI{96}{\percent}}
\resolucao{Os dois condutores (ida e volta) ficam \emph{em série} com a carga:
\[
\Req = 0{,}2+9{,}6+0{,}2=\SI{10}{\ohm}
\quad\Rightarrow\quad
I=\frac{24}{10}=\SI{2.4}{\ampere}.
\]
(b) $U_{\text{cabo}} = 0{,}4\cdot 2{,}4 = \SI{0.96}{\volt}$.\\
(c) $U_{\text{carga}} = 24-0{,}96 = \SI{23.04}{\volt}$ (ou $9{,}6\cdot2{,}4$).\\
(d) Como a corrente é a mesma em tudo,
$\eta = \dfrac{R_c}{\Req} = \dfrac{9{,}6}{10} = 0{,}96 = \SI{96}{\percent}$.\\
\textbf{Por que isso importa:} a perda cresce com $I^2$. Se a mesma potência fosse
entregue com o dobro da tensão (e metade da corrente), a perda no cabo cairia à
quarta parte --- é o argumento a favor da transmissão em alta tensão.}
\end{questao}

\begin{questao}[3]{}
Dispõe-se de resistores de $\SI{100}{\ohm}$, todos idênticos, e de uma fonte
ideal de $\SI{50}{\volt}$. Deseja-se obter, a partir dessa fonte, uma saída de
exatamente $\SI{20}{\volt}$ em vazio (sem carga conectada à saída), usando apenas
uma associação série desses resistores como divisor de tensão.
\begin{itensq}
\item Qual é o menor número de resistores capaz de fornecer essa tensão?
\item Nessa configuração, quanto vale a corrente que circula pelo divisor?
\end{itensq}
\resposta{(a) 5 resistores (saída sobre 2 deles); (b) \SI{0.1}{\ampere}}
\resolucao{(a) Com $n$ resistores iguais em série, a tensão sobre $k$ deles é
\[
U_k = E\,\frac{k}{n} \quad\Rightarrow\quad 20 = 50\cdot\frac{k}{n}
\quad\Rightarrow\quad \frac{k}{n}=\frac{2}{5}.
\]
A menor solução inteira é $k=2$ e $n=5$: cinco resistores em série, com a saída
tomada sobre dois deles.\\
(b) $\Req = 5\cdot 100 = \SI{500}{\ohm}$ e
$I = \dfrac{50}{500} = \SI{0.1}{\ampere}$.\\
\textbf{Cuidado prático:} essa tensão de \SI{20}{\volt} só se mantém \emph{em
vazio}. Ao ligar uma carga em paralelo com os dois resistores, a tensão de saída
cai --- por isso divisores resistivos não servem como fonte de alimentação.}
\end{questao}

\begin{questao}[3]{}
Uma linha de transmissão de $\SI{2}{\kilo\meter}$ alimenta uma carga de
$\SI{10}{\ohm}$ a partir de uma fonte de $\SI{240}{\volt}$. O cabo é de cobre
($\rho=\SI{1.7e-8}{\ohm\meter}$) e deve-se escolher sua seção $A$ de modo que a
queda de tensão na linha não ultrapasse $\SI{4}{\percent}$ da tensão da fonte.
\begin{itensq}
\item Determine a resistência máxima admissível para o par de condutores.
\item Calcule a seção mínima do cabo.
\item Determine o rendimento da transmissão nessa condição e a potência perdida.
\end{itensq}
\resposta{(a) $R_{\text{linha}} \le \SI{0.417}{\ohm}$;
(b) $A \ge \SI{163}{\milli\meter\squared}$; (c) $\eta = \SI{96}{\percent}$,
$P_{\text{perdida}}\approx\SI{221}{\watt}$}
\resolucao{\textbf{(a)} A linha e a carga estão em série, formando um divisor de
tensão. Para a queda na linha ser no máximo \SI{4}{\percent}:
\[
\frac{R_{\text{linha}}}{R_{\text{linha}}+R_L}\le 0{,}04
\;\Longrightarrow\;
R_{\text{linha}} \le \frac{0{,}04}{0{,}96}\,R_L = \frac{10}{24}
\approx\SI{0.417}{\ohm}.
\]
\textbf{(b)} O comprimento \emph{total} de condutor é o dobro da distância
(ida e volta): $L = \SI{4000}{\meter}$. De $R=\rho L/A$:
\[
A \ge \frac{\rho L}{R_{\text{linha}}}
=\frac{(\pot{1.7}{-8})(4000)}{0{,}417}
\approx\pot{1.63}{-4}\,\si{\meter\squared}=\SI{163}{\milli\meter\squared}.
\]
Comercialmente, adotaria-se $\SI{185}{\milli\meter\squared}$ (valor normalizado
imediatamente superior).

\textbf{(c)} Com $R_{\text{linha}} = \SI{0.417}{\ohm}$:
\[
I=\frac{240}{10{,}417}\approx\SI{23}{\ampere},
\qquad
\eta=\frac{R_L}{R_L+R_{\text{linha}}}=\frac{10}{10{,}417}\approx\SI{96}{\percent},
\]
\[
P_{\text{perdida}}=R_{\text{linha}}I^{2}\approx 0{,}417\cdot23^{2}
\approx\SI{221}{\watt}.
\]
\textbf{Reflexão de projeto:} note o custo do critério. Exigir apenas
\SI{4}{\percent} de queda impôs um cabo de $\SI{163}{\milli\meter\squared}$ ---
muito caro. Se a mesma potência fosse transmitida em $\SI{2400}{\volt}$
(dez vezes mais), a corrente cairia a $1/10$ e a perda ($\propto I^2$) a
$1/100$, permitindo um cabo cem vezes menor. \textbf{É por isso que se transmite
energia em alta tensão} --- a economia está no cobre, não na energia.}
\end{questao}

% BEGIN BANCO COMPLEMENTAR Circuitos em série
% =====================================================================
%  Banco complementar --- Circuitos Serie  (REVISADO)
%  Substitui a versao gerada por template (11 clones em N2, 12 em N3,
%  10 questoes conceituais indevidamente marcadas como N4).
%  Sem sobreposicao com o cap03 (fonte em oposicao, seis lampadas,
%  queda em cabo e linha de transmissao ja estao la).
%  Distribuicao mantida: 6 N1 + 11 N2 + 12 N3 + 10 N4 = 39
% =====================================================================

\subsubsection*{Banco complementar --- Circuitos em série}

\begin{atencao}[Organização do banco]
Toda associação série repousa em dois fatos: a \textbf{corrente é a mesma} em
todos os elementos e as \textbf{quedas se somam}. Deles decorre que a tensão e
a potência se distribuem \emph{proporcionalmente} a $R$. O N3 explora falhas,
temperatura e tolerância; o N4 trata de projeto, sensibilidade, propagação de
incerteza e otimização.
\end{atencao}

\blocofixacao

\begin{questao}[1]{}
Determine a resistência equivalente de $R_1=\SI{1.2}{\kilo\ohm}$,
$R_2=\SI{470}{\ohm}$ e $R_3=\SI{2.2}{\kilo\ohm}$ em série.
\resposta{$R_{eq}=\SI{3.87}{\kilo\ohm}$}
\resolucao{Converta tudo para a mesma unidade \emph{antes} de somar:
\[
R_{eq}=1200+470+2200=\SI{3870}{\ohm}=\SI{3.87}{\kilo\ohm}.
\]
Em série a soma é direta --- não há inversos envolvidos. Misturar
$\si{\kilo\ohm}$ com $\si{\ohm}$ sem converter é o erro mais frequente aqui.}
\end{questao}

\begin{questao}[1]{}
Uma associação série de três resistores tem $R_{eq}=\SI{75}{\ohm}$. Sabendo que
$R_1=\SI{22}{\ohm}$ e $R_2=\SI{33}{\ohm}$, determine $R_3$.
\resposta{$R_3=\SI{20}{\ohm}$}
\resolucao{
\[
R_3=R_{eq}-R_1-R_2=75-22-33=\SI{20}{\ohm}.
\]
A soma série é reversível: conhecidos o total e todas as parcelas menos uma,
a que falta sai por subtração --- diferente do paralelo, em que seria preciso
subtrair \emph{condutâncias}.}
\end{questao}

\begin{questao}[1]{}
Uma fonte ideal de $\SI{24}{\volt}$ alimenta $R_1=\SI{4}{\ohm}$ em série com
$R_2=\SI{8}{\ohm}$. Determine a corrente do circuito.
\resposta{$I=\SI{2}{\ampere}$}
\resolucao{
\[
R_{eq}=4+8=\SI{12}{\ohm}
\;\Longrightarrow\;
I=\frac{V}{R_{eq}}=\frac{24}{12}=\SI{2}{\ampere}.
\]
Essa corrente é a \emph{mesma} nos dois resistores --- não existe "corrente de
$R_1$" diferente da "corrente de $R_2$" em uma série.}
\end{questao}

\begin{questao}[1]{}
Em uma cadeia série percorrida por $\SI{0.5}{\ampere}$, determine a queda de
tensão em um resistor de $\SI{68}{\ohm}$.
\resposta{$U=\SI{34}{\volt}$}
\resolucao{Lei de Ohm aplicada ao elemento:
\[
U=R\,I=68\cdot0{,}5=\SI{34}{\volt}.
\]
Como $I$ é comum, a queda de cada resistor é \emph{diretamente proporcional} ao
seu valor: o maior resistor sempre fica com a maior parcela da tensão.}
\end{questao}

\begin{questao}[1]{}
Acrescentando-se mais um resistor \textbf{em série} a um circuito alimentado
por fonte ideal, a corrente total:
\begin{alternativas}[2]
\item aumenta
\item diminui
\item não se altera
\item pode aumentar ou diminuir, conforme o valor
\end{alternativas}
\resposta{(b)}
\resolucao{$R_{eq}$ só pode crescer ao se acrescentar um resistor em série
($R_{eq}=\sum R_k$, com todas as parcelas positivas). Com $V$ fixo,
$I=V/R_{eq}$ necessariamente diminui. Consequência prática: em uma cadeia de
lâmpadas, acrescentar mais uma faz \emph{todas} brilharem menos.}
\end{questao}

\begin{questao}[1]{}
Dois resistores de $\SI{15}{\ohm}$ e $\SI{45}{\ohm}$ são ligados em série a uma
fonte. Trocando-se a \textbf{ordem} deles no circuito, o que muda?
\begin{alternativas}[2]
\item a resistência equivalente
\item a corrente do circuito
\item a queda em cada resistor
\item nada
\end{alternativas}
\resposta{(d)}
\resolucao{A soma $R_{eq}=15+45=\SI{60}{\ohm}$ é comutativa, a corrente
$I=V/R_{eq}$ depende só de $R_{eq}$, e cada queda $U_k=R_kI$ depende só do
próprio resistor e de $I$. Nada muda.\\
O que muda é o \emph{potencial} de cada ponto intermediário em relação ao terra
--- relevante apenas se algum ponto do circuito estiver aterrado ou se houver
medição referida a um nó específico.}
\end{questao}

\blocoaplicacao

\begin{questao}[2]{}
$R_1=\SI{5}{\ohm}$, $R_2=\SI{10}{\ohm}$ e $R_3=\SI{15}{\ohm}$ estão em série
sob $\SI{60}{\volt}$. Determine a corrente e as três quedas de tensão.
\resposta{$I=\SI{2}{\ampere}$; $U_1=\SI{10}{\volt}$, $U_2=\SI{20}{\volt}$,
$U_3=\SI{30}{\volt}$}
\resolucao{
\[
R_{eq}=5+10+15=\SI{30}{\ohm}
\;\Longrightarrow\;
I=\frac{60}{30}=\SI{2}{\ampere}.
\]
\[
U_1=5\cdot2=10,\quad U_2=10\cdot2=20,\quad U_3=15\cdot2=\SI{30}{\volt}.
\]
\textbf{Verificação (LKT):} $10+20+30=\SI{60}{\volt}$ --- fecha com a fonte.
Note a proporcionalidade: as quedas estão entre si como $5:10:15=1:2:3$.}
\end{questao}

\begin{questao}[2]{}
No circuito da questão anterior, calcule a potência dissipada em cada resistor
e confronte com a potência fornecida pela fonte.
\resposta{$\SI{20}{\watt}$, $\SI{40}{\watt}$, $\SI{60}{\watt}$; total
$\SI{120}{\watt}$}
\resolucao{Com $I=\SI{2}{\ampere}$ comum, use $P=RI^{2}$:
\[
P_1=5\cdot4=20,\quad P_2=10\cdot4=40,\quad P_3=15\cdot4=\SI{60}{\watt}.
\]
Fonte: $P=VI=60\cdot2=\SI{120}{\watt}=20+40+60$ \checkmark.\\
\textbf{Regra da série:} como $I$ é comum, $P\propto R$. O maior resistor
dissipa mais --- exatamente o oposto do que ocorre no paralelo, onde $V$ é
comum e $P\propto 1/R$.}
\end{questao}

\begin{questao}[2]{}
Em uma cadeia série de três resistores, mediu-se $\SI{18}{\volt}$ sobre o
resistor de $\SI{90}{\ohm}$. Os outros dois valem $\SI{30}{\ohm}$ e
$\SI{60}{\ohm}$. Determine a tensão da fonte.
\resposta{$V=\SI{36}{\volt}$}
\resolucao{A medida sobre um único elemento revela a corrente da cadeia
inteira:
\[
I=\frac{18}{90}=\SI{0.2}{\ampere}.
\]
\[
R_{eq}=30+60+90=\SI{180}{\ohm}
\;\Longrightarrow\;
V=R_{eq}I=180\cdot0{,}2=\SI{36}{\volt}.
\]
Este é o procedimento padrão de bancada: medir a queda em um resistor conhecido
transforma-o em um \emph{sensor de corrente} (resistor \emph{shunt}), sem
abrir o circuito para inserir o amperímetro.}
\end{questao}

\begin{questao}[2]{}
Duas lâmpadas incandescentes, de resistências $\SI{240}{\ohm}$ e
$\SI{960}{\ohm}$, são ligadas \textbf{em série} a $\SI{120}{\volt}$. Qual
brilha mais e por quê?
\resposta{A de $\SI{960}{\ohm}$; $P=\SI{9.6}{\watt}$ contra $\SI{2.4}{\watt}$}
\resolucao{
\[
I=\frac{120}{240+960}=\frac{120}{1200}=\SI{0.1}{\ampere}
\]
\[
P_{240}=240\cdot(0{,}1)^{2}=\SI{2.4}{\watt},\qquad
P_{960}=960\cdot(0{,}1)^{2}=\SI{9.6}{\watt}.
\]
Como a corrente é comum, $P\propto R$: a lâmpada de \emph{maior} resistência
brilha mais. Isso contraria a intuição de quem associa "mais potência" a "menos
resistência" --- válido apenas quando a \emph{tensão} é comum, isto é, no
paralelo. Uma lâmpada de $\SI{100}{\watt}$ e outra de $\SI{25}{\watt}$ ligadas
em série: a de $\SI{25}{\watt}$ é que brilha mais.}
\end{questao}

\begin{questao}[2]{}
Deseja-se limitar a $\SI{150}{\milli\ampere}$ a corrente de uma carga de
$\SI{60}{\ohm}$ alimentada por $\SI{24}{\volt}$. Determine o resistor em série
necessário e sua potência.
\resposta{$R_s=\SI{100}{\ohm}$; $P_s=\SI{2.25}{\watt}$}
\resolucao{
\[
R_{eq}=\frac{V}{I}=\frac{24}{0{,}150}=\SI{160}{\ohm}
\;\Longrightarrow\;
R_s=160-60=\SI{100}{\ohm}.
\]
\[
P_s=R_sI^{2}=100\cdot(0{,}15)^{2}=\SI{2.25}{\watt}.
\]
Especifique pelo menos $\SI{5}{\watt}$ (fator $2\times$ de folga). Observe que o
resistor limitador dissipa $\SI{2.25}{\watt}$ dos $\SI{3.6}{\watt}$ totais ---
$62\%$ da energia vira calor. Limitar corrente com resistor é simples, mas
nunca eficiente.}
\end{questao}

\begin{questao}[2]{}
Duas montagens têm a mesma $R_{eq}=\SI{100}{\ohm}$ sob $\SI{50}{\volt}$:
(A) dois resistores de $\SI{50}{\ohm}$; (B) um de $\SI{10}{\ohm}$ e outro de
$\SI{90}{\ohm}$. Compare corrente total, potência total e a potência do
resistor mais solicitado.
\resposta{Corrente e potência totais iguais ($\SI{0.5}{\ampere}$,
$\SI{25}{\watt}$); resistor mais solicitado: $\SI{12.5}{\watt}$ em (A) e
$\SI{22.5}{\watt}$ em (B)}
\resolucao{Em ambas: $I=50/100=\SI{0.5}{\ampere}$ e
$P_{tot}=50\cdot0{,}5=\SI{25}{\watt}$.\\
(A) $P=50\cdot0{,}25=\SI{12.5}{\watt}$ em cada.\\
(B) $P_{10}=\SI{2.5}{\watt}$ e $P_{90}=\SI{22.5}{\watt}$.\\
\textbf{Conclusão de projeto:} $R_{eq}$ igual \emph{não} significa
especificação igual. A montagem (B) exige um resistor de pelo menos
$\SI{25}{\watt}$, enquanto (A) se resolve com dois de $\SI{15}{\watt}$.
Distribuir a resistência em parcelas iguais distribui também o aquecimento.}
\end{questao}

\begin{questao}[2]{}
Uma cadeia série de $\SI{48}{\ohm}$ é alimentada por $\SI{12}{\volt}$ e
protegida por fusível. Determine a corrente normal de operação e justifique a
escolha de um fusível de $\SI{500}{\milli\ampere}$.
\resposta{$I=\SI{250}{\milli\ampere}$; o fusível deve ficar acima da corrente
normal e abaixo da corrente de falha}
\resolucao{
\[
I=\frac{12}{48}=\SI{0.25}{\ampere}=\SI{250}{\milli\ampere}.
\]
Um fusível de $\SI{500}{\milli\ampere}$ dá folga de $2\times$: não atua em
transitórios normais de energização, mas interrompe bem antes das correntes de
curto. \textbf{Estimativa da falha:} se metade da cadeia entrar em curto,
$R=\SI{24}{\ohm}$ e $I=\SI{500}{\milli\ampere}$ --- exatamente no limiar; um
curto franco ($R\to0$) levaria a corrente a valores muito maiores, e o fusível
atua rapidamente. Um fusível \emph{muito} próximo da corrente nominal produz
desarmes espúrios; muito acima, não protege.}
\end{questao}

\begin{questao}[2]{}
Um reostato de $\SI{0}{\ohm}$ a $\SI{100}{\ohm}$ está em série com uma carga de
$\SI{50}{\ohm}$, sob $\SI{30}{\volt}$. Determine a faixa de variação da
corrente.
\resposta{de $\SI{0.2}{\ampere}$ a $\SI{0.6}{\ampere}$}
\resolucao{
\[
I_{\max}=\frac{30}{50+0}=\SI{0.6}{\ampere},
\qquad
I_{\min}=\frac{30}{50+100}=\SI{0.2}{\ampere}.
\]
A faixa é $3{:}1$ --- \textbf{não} $\infty{:}1$, porque a carga fixa de
$\SI{50}{\ohm}$ impõe um teto à corrente por mais que o reostato seja reduzido.
Para ampliar a faixa seria preciso um reostato de valor maior em relação à
carga. Note também que a relação $I(R)$ é hiperbólica, não linear: o controle é
muito mais "sensível" perto de $R=0$.}
\end{questao}

\begin{questao}[2]{}
Em uma cadeia série sob tensão fixa, se um dos resistores tiver seu valor
\textbf{dobrado}, a queda sobre ele:
\begin{alternativas}[2]
\item dobra
\item aumenta, mas menos que o dobro
\item não se altera
\item diminui
\end{alternativas}
\resposta{(b)}
\resolucao{A queda é $U_k=V\dfrac{R_k}{R_{eq}}$. Dobrar $R_k$ dobra o
numerador, mas \emph{também} aumenta o denominador --- o crescimento é menos
que proporcional.\\
\textbf{Exemplo:} $R_1=R_2=\SI{10}{\ohm}$ sob $\SI{12}{\volt}$ dá
$U_1=\SI{6}{\volt}$. Com $R_1=\SI{20}{\ohm}$: $U_1=12\cdot20/30=\SI{8}{\volt}$
--- aumento de $33\%$, não de $100\%$. No limite $R_1\to\infty$, $U_1\to V$: a
queda satura no valor da fonte, nunca o ultrapassa.}
\end{questao}

\begin{questao}[2]{}
Uma cadeia série dissipa $\SI{300}{\watt}$ e opera $\SI{6}{\hour}$ por dia
durante $\SI{30}{\day}$. Calcule a energia consumida e o custo, a
$\mathrm{R\$}\,0{,}80$ por $\si{\kilo\watt\hour}$.
\resposta{$E=\SI{54}{\kilo\watt\hour}$; custo $=\mathrm{R\$}\,43{,}20$}
\resolucao{
\[
E=P\,t=0{,}300\,\si{\kilo\watt}\times6\times30
=\SI{54}{\kilo\watt\hour}.
\]
\[
\text{Custo}=54\times0{,}80=\mathrm{R\$}\,43{,}20.
\]
A conta não depende de \emph{como} a resistência está distribuída na cadeia ---
só da potência total. A distribuição importa para o dimensionamento térmico de
cada componente, não para a fatura.}
\end{questao}

\begin{questao}[2]{}
Em uma cadeia série sob $\SI{100}{\volt}$ mediram-se as quedas
$\SI{22}{\volt}$, $\SI{35}{\volt}$ e $\SI{41}{\volt}$. A medição é consistente?
Se a corrente for $\SI{50}{\milli\ampere}$, determine os três resistores.
\resposta{Sim ($22+35+41=98\approx100$, dentro de $2\%$);
$\SI{440}{\ohm}$, $\SI{700}{\ohm}$ e $\SI{820}{\ohm}$}
\resolucao{\textbf{Consistência (LKT):} a soma das quedas deve igualar a tensão
da fonte. Aqui $22+35+41=\SI{98}{\volt}$, desvio de $\SI{2}{\volt}$ ($2\%$) ---
compatível com a exatidão típica de um multímetro de bancada. Um desvio de,
digamos, $\SI{20}{\volt}$ indicaria erro de medida ou um quarto elemento não
contabilizado (resistência de contato, por exemplo).\\
\textbf{Resistores:} com $I=\SI{0.050}{\ampere}$ comum a todos,
\[
R_1=\frac{22}{0{,}05}=440,\quad
R_2=\frac{35}{0{,}05}=700,\quad
R_3=\frac{41}{0{,}05}=\SI{820}{\ohm}.
\]}
\end{questao}

\blocoaprofundamento

\begin{questao}[3]{}
$R_1=\SI{10}{\ohm}$ e $R_2=\SI{20}{\ohm}$ estão em série sob
$\SI{30}{\volt}$.
\begin{itensq}
\item Calcule a potência de cada um.
\item Especifique a potência nominal comercial de cada resistor, com fator de
      segurança $2$.
\item Explique por que especificar ambos igualmente seria desperdício ou risco.
\end{itensq}
\resposta{(a) $\SI{10}{\watt}$ e $\SI{20}{\watt}$; (b) $\SI{25}{\watt}$ e
$\SI{50}{\watt}$}
\resolucao{(a) $I=30/30=\SI{1}{\ampere}$; $P_1=10\cdot1=\SI{10}{\watt}$,
$P_2=20\cdot1=\SI{20}{\watt}$.\\
(b) Com fator $2$: $\SI{20}{\watt}$ e $\SI{40}{\watt}$; os valores comerciais
imediatamente superiores são $\SI{25}{\watt}$ e $\SI{50}{\watt}$.\\
(c) Especificar \emph{ambos} como $\SI{25}{\watt}$ deixaria $R_2$ operando a
$80\%$ do nominal --- sem margem para elevação de temperatura ambiente ou
ventilação deficiente. Especificar \emph{ambos} como $\SI{50}{\watt}$ dobraria
volume e custo de $R_1$ sem ganho. \textbf{Regra:} em série, $P\propto R$;
dimensione resistor a resistor.}
\end{questao}

\begin{questao}[3]{}
Um resistor de fio tem $\SI{50}{\ohm}$ a $\SI{20}{\celsius}$ e coeficiente
térmico $\alpha=\SI{0.004}{\per\celsius}$. Ele está em série com um resistor de
filme de $\SI{50}{\ohm}$ e $\alpha\approx0$, sob $\SI{124}{\volt}$.
\begin{itensq}
\item Determine a resistência do resistor de fio a $\SI{80}{\celsius}$.
\item Calcule a corrente a $\SI{20}{\celsius}$ e a $\SI{80}{\celsius}$.
\item Comente a variação percentual.
\end{itensq}
\resposta{(a) $\SI{62}{\ohm}$; (b) $\SI{1.24}{\ampere}$ e
$\SI{1.107}{\ampere}$; (c) queda de $10{,}7\%$}
\resolucao{(a)
$R=R_0[1+\alpha(T-T_0)]=50[1+0{,}004(60)]=50\cdot1{,}24=\SI{62}{\ohm}$.\\
(b) A $\SI{20}{\celsius}$: $R_{eq}=\SI{100}{\ohm}$ e
$I=124/100=\SI{1.24}{\ampere}$.\\
A $\SI{80}{\celsius}$: $R_{eq}=62+50=\SI{112}{\ohm}$ e
$I=124/112=\SI{1.107}{\ampere}$.\\
(c) Redução de $10{,}7\%$. Note que o resistor de fio variou $24\%$, mas a
corrente variou apenas $10{,}7\%$ --- o resistor de filme, estável, "dilui" o
efeito. Essa é a base do item de projeto explorado adiante: combinar
coeficientes para estabilizar a cadeia inteira.}
\end{questao}

\begin{questao}[3]{}
$R_1=\SI{100}{\ohm}\pm5\%$ e $R_2=\SI{200}{\ohm}\pm5\%$ estão em série.
\begin{itensq}
\item Determine $R_{eq}$ e seus limites de pior caso.
\item Expresse a tolerância resultante em porcentagem.
\item Repita supondo $R_2$ com tolerância de $1\%$.
\end{itensq}
\resposta{(a) $\SI{300}{\ohm}$, entre $\SI{285}{\ohm}$ e $\SI{315}{\ohm}$;
(b) $\pm5\%$; (c) $\pm2{,}33\%$}
\resolucao{(a) Em série as incertezas \emph{absolutas} se somam:
$\Delta R=5+10=\SI{15}{\ohm}$, logo $300\pm15$, isto é, de
$\SI{285}{\ohm}$ a $\SI{315}{\ohm}$.\\
(b) $15/300=5\%$ --- \emph{idêntica} às parcelas.\\
(c) Agora $\Delta R=5+2=\SI{7}{\ohm}$, e $7/300=2{,}33\%$.\\
\textbf{Padrão geral:} a tolerância relativa da série é a \emph{média
ponderada} das tolerâncias individuais, com pesos $R_k/R_{eq}$:
\[
t_{eq}=\frac{\sum R_k t_k}{\sum R_k}.
\]
Em (b), $t_1=t_2$ e a média devolve o mesmo valor. Em (c), o resistor
\emph{maior} domina o resultado --- por isso, ao apertar tolerâncias, comece
sempre pela maior parcela da cadeia.}
\end{questao}

\begin{questao}[3]{}
Na cadeia $R_1=\SI{10}{\ohm}$, $R_2=\SI{20}{\ohm}$, $R_3=\SI{30}{\ohm}$ sob
$\SI{60}{\volt}$, o resistor $R_2$ sofre \textbf{curto-circuito}. Compare
corrente e potências antes e depois.
\resposta{$I$: $\SI{1}{\ampere}\to\SI{1.5}{\ampere}$;
$P_1$: $10\to\SI{22.5}{\watt}$; $P_3$: $30\to\SI{67.5}{\watt}$}
\resolucao{\textbf{Antes:} $R_{eq}=\SI{60}{\ohm}$, $I=\SI{1}{\ampere}$,
$P=(10;20;30)\,\si{\watt}$, total $\SI{60}{\watt}$.\\
\textbf{Depois:} $R_{eq}=10+30=\SI{40}{\ohm}$, $I=60/40=\SI{1.5}{\ampere}$.
\[
P_1=10(1{,}5)^{2}=\SI{22.5}{\watt},\qquad P_3=30(1{,}5)^{2}=\SI{67.5}{\watt},
\]
total $\SI{90}{\watt}$ --- $50\%$ a mais consumido da fonte.\\
\textbf{Ponto crítico:} $P_1$ e $P_3$ subiram $125\%$ (fator $1{,}5^{2}=2{,}25$).
Um curto em \emph{um} componente pode destruir os \emph{outros}, que passam a
operar muito acima do nominal. Essa é a razão de as falhas em cadeia serem
frequentemente sequenciais, e não simultâneas.}
\end{questao}

\begin{questao}[3]{}
Na mesma cadeia ($\SI{10}{\ohm}$, $\SI{20}{\ohm}$, $\SI{30}{\ohm}$ sob
$\SI{60}{\volt}$), agora $R_2$ \textbf{abre}. Determine a corrente e a leitura
de um voltímetro ideal colocado sobre cada resistor.
\resposta{$I=0$; $\SI{0}{\volt}$ em $R_1$ e $R_3$; $\SI{60}{\volt}$ em $R_2$}
\resolucao{Com o circuito interrompido, $I=0$ em toda a série. Pela lei de Ohm,
$U=RI=0$ nos resistores íntegros --- mesmo estando eles perfeitamente
conectados à fonte.\\
Pela LKT, a soma das quedas deve continuar valendo $\SI{60}{\volt}$; como as
outras são nulas, \textbf{toda} a tensão da fonte aparece sobre a
descontinuidade: o voltímetro sobre $R_2$ lê $\SI{60}{\volt}$.\\
\textbf{Uso prático:} é assim que se localiza o componente aberto em uma cadeia
--- percorre-se a série com o voltímetro e o defeito é o único ponto onde
aparece tensão. Note que o voltímetro precisa ser de alta impedância; um
instrumento de baixa impedância fecharia o circuito e falsearia a leitura.}
\end{questao}

\begin{questao}[3]{}
Um LED com queda direta de $\SI{2.0}{\volt}$ deve operar com
$\SI{15}{\milli\ampere}$ a partir de uma fonte de $\SI{5.0}{\volt}$.
\begin{itensq}
\item Determine o resistor em série ideal.
\item Escolha o valor comercial (série E12) e recalcule a corrente real.
\item Determine a potência do resistor.
\end{itensq}
\resposta{(a) $\SI{200}{\ohm}$; (b) $\SI{220}{\ohm}\Rightarrow
\SI{13.6}{\milli\ampere}$; (c) $\SI{41}{\milli\watt}$}
\resolucao{(a) O resistor absorve a diferença entre a fonte e a queda do LED:
\[
R=\frac{V_{fonte}-V_{LED}}{I}=\frac{5{,}0-2{,}0}{0{,}015}=\SI{200}{\ohm}.
\]
(b) A E12 oferece $\SI{180}{\ohm}$ e $\SI{220}{\ohm}$. Escolha o
\textbf{maior} --- errar para menos aumenta a corrente e reduz a vida do LED:
\[
I=\frac{3{,}0}{220}=\SI{13.6}{\milli\ampere}.
\]
(c) $P=R I^{2}=220(0{,}0136)^{2}=\SI{41}{\milli\watt}$; um resistor de
$\frac{1}{8}\,\si{\watt}$ já basta com folga generosa.\\
\textbf{Observação:} o LED \emph{não} obedece à lei de Ohm --- sua queda é
aproximadamente constante. Por isso ele entra na conta como uma "fonte de
tensão" a ser subtraída, e não como uma resistência a ser somada.}
\end{questao}

\begin{questao}[3]{}
Uma bateria de $\EMF=\SI{12}{\volt}$ e resistência interna $r=\SI{0.5}{\ohm}$
alimenta uma carga de $\SI{9.5}{\ohm}$.
\begin{itensq}
\item Determine a corrente e a tensão nos terminais.
\item Calcule a regulação de tensão percentual.
\item Determine a fração da potência perdida internamente.
\end{itensq}
\resposta{(a) $\SI{1.2}{\ampere}$, $\SI{11.4}{\volt}$; (b) $5{,}26\%$;
(c) $5\%$}
\resolucao{(a) $I=\dfrac{12}{9{,}5+0{,}5}=\SI{1.2}{\ampere}$;
$V_{term}=12-0{,}5(1{,}2)=\SI{11.4}{\volt}$.\\
(b) Tomando a tensão em vazio como referência de partida:
\[
\text{Reg}=\frac{V_{vazio}-V_{carga}}{V_{carga}}
=\frac{12-11{,}4}{11{,}4}=5{,}26\%.
\]
(c) $P_r=rI^{2}=0{,}5(1{,}44)=\SI{0.72}{\watt}$ contra
$P_{tot}=12(1{,}2)=\SI{14.4}{\watt}$: exatamente $5\%$, que é a razão
$r/(r+R_L)$.\\
\textbf{Compromisso:} reduzir $r$ melhora regulação \emph{e} rendimento
simultaneamente --- por isso $r$ baixo é sempre desejável em fontes. É diferente
da máxima transferência de potência, que exigiria $R_L=r$ e um rendimento de
apenas $50\%$.}
\end{questao}

\begin{questao}[3]{}
Duas cadeias têm $R_{eq}=\SI{120}{\ohm}$ sob $\SI{120}{\volt}$: (A) quatro
resistores de $\SI{30}{\ohm}$; (B) um de $\SI{100}{\ohm}$ e dois de
$\SI{10}{\ohm}$. Compare a potência do componente mais solicitado e a potência
total, e escolha a montagem preferível.
\resposta{Total $\SI{120}{\watt}$ em ambas; pior componente:
$\SI{30}{\watt}$ em (A) e $\SI{100}{\watt}$ em (B) --- (A) é preferível}
\resolucao{$I=120/120=\SI{1}{\ampere}$ nas duas.\\
(A) $P=30\cdot1=\SI{30}{\watt}$ em cada um dos quatro.\\
(B) $P_{100}=\SI{100}{\watt}$; $P_{10}=\SI{10}{\watt}$ (dois).\\
Total: $\SI{120}{\watt}$ nos dois casos --- a fonte não distingue as montagens.\\
\textbf{Decisão:} (A) exige quatro resistores de $\SI{50}{\watt}$ nominais;
(B) exige um de $\SI{200}{\watt}$, componente caro, volumoso e com ponto quente
concentrado. Sempre que possível, \textbf{divida a dissipação} em parcelas
iguais: a área de troca térmica cresce e a temperatura máxima cai.}
\end{questao}

\begin{questao}[3]{}
Um termistor NTC de $\SI{120}{\ohm}$ (frio) e $\SI{5}{\ohm}$ (quente) é
colocado em série com uma carga de $\SI{20}{\ohm}$, sob $\SI{100}{\volt}$.
Determine a corrente no instante da energização e em regime, e explique a
função do componente.
\resposta{$\SI{0.71}{\ampere}$ na partida e $\SI{4.0}{\ampere}$ em regime}
\resolucao{\textbf{Partida (frio):}
$I=\dfrac{100}{120+20}=\SI{0.71}{\ampere}$.\\
\textbf{Regime (aquecido):}
$I=\dfrac{100}{5+20}=\SI{4.0}{\ampere}$.\\
\textbf{Função:} sem o NTC, a corrente de energização seria imediatamente
$100/20=\SI{5}{\ampere}$. O termistor a limita a $\SI{0.71}{\ampere}$ e vai
"saindo do caminho" à medida que aquece --- um \textbf{limitador de corrente de
partida}. É o que protege capacitores de entrada e contatos de relé em fontes
chaveadas.\\
\textbf{Limitação:} em regime o NTC ainda dissipa $5(4)^{2}=\SI{80}{\watt}$, e
precisa de tempo para esfriar antes de um novo religamento --- por isso ele não
protege contra reenergizações rápidas e sucessivas.}
\end{questao}

\begin{questao}[3]{}
Em uma cadeia de três resistores sob $\SI{90}{\volt}$, sabe-se que
$R_2=2R_1$, $R_3=3R_1$ e que a corrente é $\SI{0.5}{\ampere}$. Determine os
três valores e as quedas.
\resposta{$R_1=\SI{30}{\ohm}$, $R_2=\SI{60}{\ohm}$, $R_3=\SI{90}{\ohm}$;
quedas $\SI{15}{\volt}$, $\SI{30}{\volt}$, $\SI{45}{\volt}$}
\resolucao{
\[
R_{eq}=\frac{V}{I}=\frac{90}{0{,}5}=\SI{180}{\ohm}.
\]
Como $R_{eq}=R_1+2R_1+3R_1=6R_1$:
\[
R_1=\frac{180}{6}=\SI{30}{\ohm},\quad R_2=\SI{60}{\ohm},\quad R_3=\SI{90}{\ohm}.
\]
Quedas: $U_k=R_kI$, ou seja $15$, $30$ e $\SI{45}{\volt}$; soma
$=\SI{90}{\volt}$ \checkmark.\\
Como as quedas guardam a mesma proporção $1{:}2{:}3$ dos resistores, poderiam
ter sido obtidas direto: $90$ dividido em $6$ partes de $\SI{15}{\volt}$.}
\end{questao}

\begin{questao}[3]{}
Uma cadeia de $n$ resistores iguais de $\SI{47}{\ohm}$ deve limitar a corrente
a no máximo $\SI{40}{\milli\ampere}$ sob $\SI{24}{\volt}$. Determine o menor
$n$ e a corrente resultante.
\resposta{$n=13$; $I=\SI{39.3}{\milli\ampere}$}
\resolucao{Exige-se $R_{eq}\ge V/I_{\max}=24/0{,}040=\SI{600}{\ohm}$:
\[
47n\ge600 \;\Longrightarrow\; n\ge12{,}77 \;\Longrightarrow\; n=13.
\]
\[
R_{eq}=13\cdot47=\SI{611}{\ohm}
\;\Longrightarrow\;
I=\frac{24}{611}=\SI{39.3}{\milli\ampere}\ \le\ \SI{40}{\milli\ampere}\ \checkmark
\]
\textbf{Cuidado com o arredondamento:} $n=12$ daria
$R_{eq}=\SI{564}{\ohm}$ e $I=\SI{42.6}{\milli\ampere}$ --- \emph{acima} do
limite. Em restrições do tipo "no máximo", arredonde sempre \textbf{para cima},
mesmo que a fração seja pequena.}
\end{questao}

\begin{questao}[3]{}
Em uma cadeia série de dois resistores sob tensão fixa $V$, demonstre que a
soma das potências é $P=\dfrac{V^{2}}{R_1+R_2}$ e explique por que aumentar
\emph{qualquer} resistor reduz a potência total, embora possa \emph{aumentar} a
potência do outro.
\resposta{Aumentar $R_k$ reduz $P_{tot}$; a potência do outro sempre diminui
--- é a do próprio $R_k$ que pode subir}
\resolucao{Com $I=V/(R_1+R_2)$:
\[
P=VI=\frac{V^{2}}{R_1+R_2}.
\]
Como $R_1+R_2$ cresce, $P$ diminui --- \emph{sempre}.\\
Para o outro resistor, $P_2=R_2I^{2}=\dfrac{V^{2}R_2}{(R_1+R_2)^{2}}$: com
$R_2$ fixo e $R_1$ crescendo, $P_2$ só pode \textbf{diminuir}.\\
Já a potência do \emph{próprio} $R_1$,
$P_1=\dfrac{V^{2}R_1}{(R_1+R_2)^{2}}$, cresce enquanto $R_1<R_2$, atinge um
máximo em $R_1=R_2$ e depois decresce. \textbf{Conclusão:} o enunciado
intuitivo "mais resistência, mais aquecimento" é falso em série --- o
aquecimento de um resistor é máximo quando ele iguala a soma dos demais, e cai
para os dois lados. Este resultado é retomado quantitativamente no bloco de
desafio.}
\end{questao}

\blocodesafio

\begin{questao}[4]{}
\textbf{(Projeto.)} Uma cadeia série de três resistores deve dividir
$\SI{24}{\volt}$ em quedas de $\SI{6}{\volt}$, $\SI{8}{\volt}$ e
$\SI{10}{\volt}$, com corrente de $\SI{20}{\milli\ampere}$.
\begin{itensq}
\item Determine os três resistores e suas potências.
\item Escolha valores comerciais E24 e recalcule as quedas reais.
\item Avalie se o erro é aceitável para alimentar cargas com tolerância de
      $\pm5\%$.
\end{itensq}
\resposta{(a) $\SI{300}{\ohm}$, $\SI{400}{\ohm}$, $\SI{500}{\ohm}$;
$\SI{0.12}{\watt}$, $\SI{0.16}{\watt}$, $\SI{0.20}{\watt}$;
(b) $300/390/510\ \si{\ohm}$ $\Rightarrow$ $5{,}98/7{,}77/10{,}16\ \si{\volt}$;
(c) sim, erro máximo de $2{,}9\%$}
\resolucao{(a) Como a corrente é comum, $R_k=U_k/I$:
\[
R_1=\frac{6}{0{,}02}=300,\quad R_2=\frac{8}{0{,}02}=400,\quad
R_3=\frac{10}{0{,}02}=\SI{500}{\ohm}.
\]
Potências: $P_k=U_kI$ $\Rightarrow$ $0{,}12$; $0{,}16$;
$\SI{0.20}{\watt}$. Um resistor de $\frac{1}{2}\,\si{\watt}$ em cada posição
oferece folga confortável.\\
(b) A E24 tem $300$ e $510$, mas não $400$ --- o vizinho é $\SI{390}{\ohm}$.
Com $R_{eq}=300+390+510=\SI{1200}{\ohm}$, a corrente real é
$24/1200=\SI{20.0}{\milli\ampere}$ (coincidência favorável), e as quedas ficam
\[
U_1=\SI{6.00}{\volt},\quad U_2=\SI{7.80}{\volt},\quad U_3=\SI{10.20}{\volt}.
\]
(c) Desvios: $0\%$, $-2{,}5\%$ e $+2{,}0\%$ --- todos dentro de $\pm5\%$.
\textbf{Ressalva essencial:} este projeto só vale \emph{em vazio}. Qualquer
carga real conectada em paralelo com uma das quedas altera a corrente da cadeia
e derruba todas as tensões. Um divisor resistivo puro serve para referências de
altíssima impedância, nunca para alimentar cargas.}
\end{questao}

\begin{questao}[4]{}
\textbf{(Propagação de incerteza.)} Uma cadeia é formada por $N$ resistores
nominalmente iguais, cada um com tolerância relativa $t$.
\begin{itensq}
\item Obtenha a tolerância de $R_{eq}$ no critério de \textbf{pior caso}.
\item Obtenha-a no critério \textbf{estatístico} (erros independentes).
\item Compare para $N=5$ e $t=5\%$, e discuta qual usar.
\end{itensq}
\resposta{(a) $t$; (b) $t/\sqrt{N}$; (c) $5\%$ contra $2{,}24\%$}
\resolucao{(a) Pior caso: todos desviam para o mesmo lado.
$\Delta R_{eq}=N(tR)$ e $R_{eq}=NR$, logo
\[
\frac{\Delta R_{eq}}{R_{eq}}=\frac{NtR}{NR}=t.
\]
A tolerância relativa \textbf{não melhora nem piora} --- resultado que
surpreende, mas decorre de a série somar tanto os valores quanto os desvios.\\
(b) Desvios independentes combinam-se em quadratura:
\[
u(R_{eq})=\sqrt{N}\,(tR)
\;\Longrightarrow\;
\frac{u(R_{eq})}{R_{eq}}=\frac{\sqrt{N}\,tR}{NR}=\frac{t}{\sqrt{N}}.
\]
(c) Para $N=5$, $t=5\%$: pior caso $5\%$; estatístico
$5/\sqrt5=2{,}24\%$.\\
\textbf{Qual usar:} o \emph{pior caso} para projetos de segurança e para lotes
pequenos, em que os desvios podem ser correlacionados (resistores do mesmo
rolo, mesma batelada, mesma temperatura). O \emph{estatístico} para produção em
série com componentes de origens independentes, onde o pior caso seria
excessivamente conservador. Note que a hipótese de independência é quase sempre
otimista na prática.}
\end{questao}

\begin{questao}[4]{}
\textbf{(Projeto com compensação térmica.)} Deseja-se uma associação série cuja
resistência total seja \emph{independente} da temperatura, combinando um
resistor de $R_1=\SI{100}{\ohm}$ com $\alpha_1=+4000\,\mathrm{ppm/^{\circ}C}$ e
outro com $\alpha_2=-1000\,\mathrm{ppm/^{\circ}C}$.
\begin{itensq}
\item Deduza a condição de compensação.
\item Determine $R_2$ e a resistência total.
\item Discuta a limitação prática dessa compensação.
\end{itensq}
\resposta{(a) $R_1\alpha_1+R_2\alpha_2=0$; (b) $R_2=\SI{400}{\ohm}$,
$R_{eq}=\SI{500}{\ohm}$; (c) só é exata em primeira ordem e numa temperatura}
\resolucao{(a) Para pequenas variações,
$R_k(T)=R_k[1+\alpha_k\Delta T]$, logo
\[
R_{eq}(T)=(R_1+R_2)+(R_1\alpha_1+R_2\alpha_2)\Delta T.
\]
Anular o termo em $\Delta T$ exige
\[
R_1\alpha_1+R_2\alpha_2=0.
\]
(b) $R_2=-\dfrac{R_1\alpha_1}{\alpha_2}
=\dfrac{100\cdot4000}{1000}=\SI{400}{\ohm}$, e
$R_{eq}=100+400=\SI{500}{\ohm}$.\\
\textbf{Verificação a $\Delta T=\SI{50}{\celsius}$:}
$R_1\to100(1{,}20)=\SI{120}{\ohm}$ e $R_2\to400(0{,}95)=\SI{380}{\ohm}$; soma
$=\SI{500}{\ohm}$ \checkmark.\\
(c) Três limitações. \textbf{(i)} A compensação é de \emph{primeira ordem};
$\alpha$ real varia com $T$, e o cancelamento se degrada longe da temperatura
de projeto. \textbf{(ii)} Os dois resistores precisam estar \emph{à mesma
temperatura} --- exige acoplamento térmico, o que é difícil quando as
dissipações diferem (aqui, $P_2=4P_1$). \textbf{(iii)} O valor de $R_2$ fica
\emph{amarrado} pela razão $\alpha_1/\alpha_2$: não se pode escolher
livremente $R_{eq}$ \emph{e} exigir compensação. Materiais dedicados
(manganina, constantan, com $\alpha$ naturalmente próximo de zero) resolvem o
problema de forma mais direta.}
\end{questao}

\begin{questao}[4]{}
\textbf{(Projeto sob incerteza.)} Um LED deve ser alimentado por uma fonte
entre $\SI{9}{\volt}$ e $\SI{12}{\volt}$. Sua queda direta varia entre
$\SI{1.8}{\volt}$ e $\SI{2.2}{\volt}$, e a corrente não pode exceder
$\SI{20}{\milli\ampere}$.
\begin{itensq}
\item Identifique a combinação de pior caso e determine o menor $R$ admissível.
\item Determine a corrente mínima resultante e avalie o impacto no brilho.
\item Especifique a potência do resistor.
\end{itensq}
\resposta{(a) $\SI{12}{\volt}$ com $\SI{1.8}{\volt}$ $\Rightarrow$
$R\ge\SI{510}{\ohm}$; (b) $\SI{13.3}{\milli\ampere}$;
(c) $\SI{0.204}{\watt}\Rightarrow\frac{1}{2}\,\si{\watt}$}
\resolucao{(a) A corrente é máxima quando a fonte está no \textbf{máximo} e a
queda do LED no \textbf{mínimo} --- é preciso combinar os extremos que atuam no
mesmo sentido, e não os extremos "nominais":
\[
R\ge\frac{12-1{,}8}{0{,}020}=\SI{510}{\ohm}.
\]
O valor $\SI{510}{\ohm}$ existe na série E24.\\
(b) O outro extremo (fonte mínima, queda máxima):
\[
I_{\min}=\frac{9-2{,}2}{510}=\SI{13.3}{\milli\ampere}.
\]
A corrente varia $13{,}3$--$\SI{20}{\milli\ampere}$, razão de $1{,}5{:}1$. Como
o fluxo luminoso é aproximadamente proporcional à corrente, o brilho varia
cerca de $33\%$ --- perceptível, porém aceitável para sinalização. Para brilho
constante seria preciso uma fonte de \emph{corrente}, não um resistor.\\
(c) $P=RI_{\max}^{2}=510(0{,}020)^{2}=\SI{0.204}{\watt}$; especifique
$\frac{1}{2}\,\si{\watt}$.\\
\textbf{Método:} em projeto com faixas, nunca use valores nominais --- enumere
os extremos e identifique qual combinação é crítica para \emph{cada} requisito.}
\end{questao}

\begin{questao}[4]{}
\textbf{(Sensibilidade.)} Em uma cadeia $I=\dfrac{V}{R_1+R_2}$, com
$V=\SI{12}{\volt}$, $R_1=\SI{100}{\ohm}$ e $R_2=\SI{200}{\ohm}$.
\begin{itensq}
\item Obtenha $\partial I/\partial R_1$ e avalie-a numericamente.
\item Obtenha a sensibilidade \emph{relativa}
      $S=\dfrac{\partial I/I}{\partial R_1/R_1}$.
\item Interprete o resultado para $R_1\ll R_2$ e $R_1\gg R_2$.
\end{itensq}
\resposta{(a) $-\pot{1.33}{-4}\,\si{\ampere\per\ohm}$;
(b) $S=-R_1/(R_1+R_2)=-1/3$; (c) $S\to0$ e $S\to-1$}
\resolucao{(a)
\[
\frac{\partial I}{\partial R_1}=-\frac{V}{(R_1+R_2)^{2}}
=-\frac{12}{300^{2}}=-\pot{1.33}{-4}\,\si{\ampere\per\ohm}.
\]
Um aumento de $\SI{1}{\ohm}$ em $R_1$ reduz a corrente em
$\SI{0.133}{\milli\ampere}$ (de $\SI{40}{\milli\ampere}$).\\
(b)
\[
S=\frac{\partial I}{\partial R_1}\cdot\frac{R_1}{I}
=-\frac{V}{(R_1+R_2)^{2}}\cdot\frac{R_1(R_1+R_2)}{V}
=-\frac{R_1}{R_1+R_2}=-\frac13.
\]
Ou seja: $1\%$ de erro em $R_1$ produz apenas $0{,}33\%$ de erro em $I$.\\
(c) Se $R_1\ll R_2$, $S\to0$: a corrente é governada por $R_2$ e $R_1$ torna-se
irrelevante --- é o caso do resistor \emph{shunt} de medição, deliberadamente
pequeno para não perturbar o circuito. Se $R_1\gg R_2$, $S\to-1$: $R_1$ passa a
determinar sozinho a corrente, e toda a sua imprecisão é transferida
integralmente. \textbf{Regra de projeto:} coloque a tolerância apertada no
elemento cuja sensibilidade relativa se aproxima de $1$; nos demais, use
componentes baratos.}
\end{questao}

\begin{questao}[4]{}
\textbf{(Análise de falhas.)} Dez resistores de $\SI{10}{\ohm}$ formam uma
cadeia sob $\SI{100}{\volt}$.
\begin{itensq}
\item Determine corrente, queda e potência por resistor em operação normal.
\item Repita supondo que \emph{um} resistor entre em curto.
\item Compare com o caso de um resistor \emph{aberto} e discuta qual falha é
      mais perigosa.
\end{itensq}
\resposta{(a) $\SI{1}{\ampere}$, $\SI{10}{\volt}$, $\SI{10}{\watt}$;
(b) $\SI{1.11}{\ampere}$, $\SI{11.1}{\volt}$, $\SI{12.3}{\watt}$;
(c) o curto é mais perigoso}
\resolucao{(a) $R_{eq}=\SI{100}{\ohm}$, $I=\SI{1}{\ampere}$,
$U=\SI{10}{\volt}$, $P=\SI{10}{\watt}$ cada; total $\SI{100}{\watt}$.\\
(b) $R_{eq}=\SI{90}{\ohm}$, $I=100/90=\SI{1.111}{\ampere}$ ($+11\%$),
$U=\SI{11.1}{\volt}$, $P=10(1{,}111)^{2}=\SI{12.35}{\watt}$ ($+23{,}5\%$);
total $\SI{111}{\watt}$.\\
(c) O \textbf{aberto} zera a corrente: o circuito para de funcionar, mas nada
sofre --- é uma falha \emph{segura} (\emph{fail-safe}), fácil de localizar com
voltímetro.\\
O \textbf{curto} é insidioso: o circuito continua operando, aparentemente
normal, mas cada resistor remanescente passa a dissipar $23{,}5\%$ a mais. Se
já operavam próximos do nominal, entram em sobrecarga --- e cada novo curto
agrava o quadro de forma \emph{acelerada}. Com dois curtos,
$P=\SI{15.6}{\watt}$ ($+56\%$); com três, $\SI{20.4}{\watt}$ ($+104\%$).
\textbf{Conclusão:} falhas que \emph{reduzem} a resistência de um circuito em
série tendem a ser progressivas e devem ser detectadas por monitoração de
corrente, não por observação funcional.}
\end{questao}

\begin{questao}[4]{}
\textbf{(Método experimental.)} Para determinar $\EMF$ e $r$ de uma bateria,
mediu-se a tensão de terminal em duas condições de carga:
$(\SI{1}{\ampere};\ \SI{11.5}{\volt})$ e
$(\SI{3}{\ampere};\ \SI{10.5}{\volt})$.
\begin{itensq}
\item Obtenha $r$ e $\EMF$.
\item Explique a vantagem de usar dois pontos em vez de medir a tensão em
      vazio.
\item Estime a corrente de curto-circuito e comente sua validade.
\end{itensq}
\resposta{(a) $r=\SI{0.5}{\ohm}$, $\EMF=\SI{12}{\volt}$;
(c) $I_{cc}=\SI{24}{\ampere}$, valor apenas indicativo}
\resolucao{(a) O modelo é $V=\EMF-rI$ --- uma reta. Dois pontos a determinam:
\[
r=-\frac{\Delta V}{\Delta I}=-\frac{10{,}5-11{,}5}{3-1}=\SI{0.5}{\ohm},
\]
\[
\EMF=V+rI=11{,}5+0{,}5(1)=\SI{12}{\volt}.
\]
(b) A medida em vazio exigiria um voltímetro de impedância \emph{infinita}; com
impedância finita, circula uma pequena corrente e a leitura já não é $\EMF$.
Além disso, baterias reais apresentam recuperação de tensão em repouso, o que
falseia a leitura em vazio. O ajuste por dois (ou mais) pontos de carga
\textbf{extrapola} para $I=0$ sem precisar chegar lá, e ainda entrega $r$ de
brinde. Com três ou mais pontos, uma regressão linear também revela se o modelo
é adequado.\\
(c) $I_{cc}=\EMF/r=12/0{,}5=\SI{24}{\ampere}$. \textbf{Ressalva:} o valor
pressupõe $r$ constante, hipótese que falha em correntes altas --- $r$ cresce
com a corrente e com a descarga, e há aquecimento. O número serve como
\emph{ordem de grandeza} para dimensionar proteção, jamais como previsão
precisa.}
\end{questao}

\begin{questao}[4]{}
\textbf{(Projeto com componentes disponíveis.)} Deseja-se uma resistência de
$\SI{1}{\kilo\ohm}$ capaz de dissipar $\SI{2}{\watt}$, dispondo apenas de
resistores de $\SI{200}{\ohm}$ e $\frac{1}{2}\,\si{\watt}$.
\begin{itensq}
\item Proponha a associação e verifique o valor.
\item Determine a potência de cada resistor e a margem de segurança.
\item Calcule a tensão e a corrente nominais da associação.
\end{itensq}
\resposta{(a) cinco em série; (b) $\SI{0.4}{\watt}$ cada, margem de $25\%$;
(c) $\SI{44.7}{\volt}$ e $\SI{44.7}{\milli\ampere}$}
\resolucao{(a) $5\times\SI{200}{\ohm}=\SI{1000}{\ohm}$ \checkmark.\\
(b) Em série a corrente é comum e os resistores são iguais, logo a potência se
reparte igualmente:
\[
P_k=\frac{2}{5}=\SI{0.4}{\watt}\ <\ \SI{0.5}{\watt}\ \checkmark
\]
Margem: $(0{,}5-0{,}4)/0{,}4=25\%$. É uma folga modesta --- em ambiente quente
ou sem ventilação, considere sete unidades de $\SI{150}{\ohm}$
($\SI{1050}{\ohm}$) ou resistores de $\SI{1}{\watt}$.\\
(c)
\[
V=\sqrt{PR}=\sqrt{2\cdot1000}=\SI{44.7}{\volt},
\qquad
I=\sqrt{P/R}=\sqrt{0{,}002}=\SI{44.7}{\milli\ampere}.
\]
\textbf{Por que a série resolve:} ela multiplica \emph{simultaneamente} o valor
e a capacidade de dissipação por $N$. Uma associação em paralelo de cinco
unidades também daria $\SI{2.5}{\watt}$, mas $\SI{40}{\ohm}$ --- resolveria a
potência e destruiria o valor.}
\end{questao}

\begin{questao}[4]{}
\textbf{(Instrumentação.)} Cinco sensores resistivos de $\SI{1}{\kilo\ohm}$
estão em série, alimentados por uma fonte de \textbf{corrente constante} de
$\SI{1}{\milli\ampere}$. Mede-se a tensão de cada nó em relação ao terra.
\begin{itensq}
\item Determine as cinco tensões nodais em condição nominal.
\item Um sensor passa a $\SI{1.2}{\kilo\ohm}$. Mostre como identificar
      \emph{qual} deles mudou.
\item Explique por que a fonte de corrente é preferível a uma fonte de tensão.
\end{itensq}
\resposta{(a) $1$, $2$, $3$, $4$ e $\SI{5}{\volt}$; (b) pela quebra do degrau
uniforme de $\SI{1}{\volt}$; (c) desacopla os sensores entre si}
\resolucao{(a) Com $I=\SI{1}{\milli\ampere}$, cada sensor cai
$\SI{1}{\volt}$. Acumulando a partir do terra: $1$, $2$, $3$, $4$,
$\SI{5}{\volt}$.\\
(b) Se o \emph{terceiro} sensor for a $\SI{1.2}{\kilo\ohm}$, as tensões passam
a $1$, $2$, $3{,}2$, $4{,}2$, $\SI{5.2}{\volt}$. A \textbf{diferença entre nós
consecutivos} vale $1;\,1;\,1{,}2;\,1;\,1$ --- o degrau anômalo aponta
diretamente o sensor defeituoso, e sua magnitude dá o novo valor:
$\Delta R=\Delta U/I=0{,}2/0{,}001=\SI{200}{\ohm}$.\\
(c) Com \textbf{corrente constante}, a queda de cada sensor é $U_k=IR_k$ ---
depende \emph{só} do próprio sensor. A variação de um não afeta a leitura dos
demais, e a relação é rigorosamente linear.\\
Com fonte de \textbf{tensão}, a corrente seria $V/\sum R_k$: a mudança de um
sensor alteraria a corrente e, portanto, \emph{todas} as leituras
simultaneamente --- o sistema ficaria acoplado, com resposta não linear e sem
identificação direta do sensor. É exatamente por isso que sistemas de RTD e
extensômetros são excitados por corrente.}
\end{questao}

\begin{questao}[4]{}
\textbf{(Otimização.)} Uma cadeia série sob tensão fixa $V$ contém um resistor
ajustável $R$ e um bloco fixo de resistência total $R_f$.
\begin{itensq}
\item Obtenha $P_R(R)$ e mostre que ela é máxima em $R=R_f$.
\item Calcule $P_{R,\max}$ e o rendimento nesse ponto.
\item Compare com a condição de máxima potência \emph{total} e explique a
      diferença.
\end{itensq}
\resposta{(a) $R=R_f$; (b) $P_{\max}=V^{2}/(4R_f)$, rendimento $50\%$;
(c) a potência total é máxima em $R\to0$}
\resolucao{(a)
\[
P_R=RI^{2}=\frac{V^{2}R}{(R+R_f)^{2}}.
\]
\[
\frac{\mathrm{d}P_R}{\mathrm{d}R}
=V^{2}\frac{(R+R_f)^{2}-R\cdot2(R+R_f)}{(R+R_f)^{4}}
=V^{2}\frac{R_f-R}{(R+R_f)^{3}}=0
\;\Longrightarrow\;
R=R_f.
\]
A derivada é positiva para $R<R_f$ e negativa para $R>R_f$: máximo confirmado.\\
(b) $P_{\max}=\dfrac{V^{2}R_f}{(2R_f)^{2}}=\dfrac{V^{2}}{4R_f}$. A potência
total nesse ponto é $\dfrac{V^{2}}{2R_f}$, de modo que $R$ recebe exatamente
\textbf{metade} --- rendimento de $50\%$.\\
(c) A potência \emph{total} $P=\dfrac{V^{2}}{R+R_f}$ é monotonicamente
decrescente: ela é máxima quando $R\to0$, e aí $P\to V^{2}/R_f$ --- o dobro do
valor anterior, porém integralmente entregue ao bloco fixo.\\
\textbf{Distinção que importa:} "máxima potência \emph{no elemento}" e "máxima
potência \emph{do circuito}" são objetivos diferentes e incompatíveis. O
primeiro é o critério de casamento (relevante quando o bloco fixo é a
resistência interna de uma fonte); o segundo é o critério de eficiência. Este
mesmo antagonismo reaparece no teorema da máxima transferência de potência.}
\end{questao}

% END BANCO COMPLEMENTAR

\gabarito[3.1 Circuitos em série]

% =====================================================================
\subsection{Circuitos em paralelo}
% =====================================================================

\blocofixacao

\begin{questao}[1]{}
Sobre uma associação de resistores \textbf{em paralelo}, é correto afirmar que:
\begin{alternativas}
\item a corrente é a mesma em todos os ramos.
\item a resistência equivalente é sempre maior que a maior das resistências.
\item a tensão é a mesma em todos os ramos.
\item a tensão se divide proporcionalmente às resistências.
\item a resistência equivalente é a soma das resistências.
\end{alternativas}
\resposta{(c)}
\resolucao{Todos os ramos estão ligados ao mesmo par de nós, logo compartilham a
mesma tensão. A corrente, essa sim, se divide --- e de forma \emph{inversamente}
proporcional às resistências. A resistência equivalente do paralelo é sempre
\textbf{menor} que a menor das resistências associadas, o que elimina (b) e (e).}
\end{questao}

\begin{questao}[1]{}
Para cada circuito, calcule a resistência equivalente, as correntes de ramo e a
corrente total $I_t$.

\circuitoq{%
\subcirc{a}{%
\begin{circuitikz}[scale=0.72,transform shape,line width=0.8pt]
 \draw (0,0) to[battery1,l_={$E=\SI{30}{\volt}$}] (0,3.4);
 \draw (0,3.4) -- (9,3.4);  \draw (0,0) -- (9,0);
 \draw (2.2,3.4) to[R,l=$\SI{4}{\ohm}$,*-*] (2.2,0);
 \draw (4.5,3.4) to[R,l=$\SI{3}{\ohm}$,*-*] (4.5,0);
 \draw (6.8,3.4) to[R,l=$\SI{6}{\ohm}$,*-*] (6.8,0);
 \draw (9,3.4)   to[R,l=$\SI{12}{\ohm}$] (9,0);
 \draw[->,Icol,line width=1pt] (0.5,3.75) -- (1.6,3.75)
       node[midway,above,font=\footnotesize]{$I_t$};
\end{circuitikz}}
\hfill
\subcirc{b}{%
\begin{circuitikz}[scale=0.72,transform shape,line width=0.8pt]
 \draw (0,0) to[battery1,l_={$E=\SI{24}{\volt}$}] (0,3.4);
 \draw (0,3.4) -- (6.6,3.4);  \draw (0,0) -- (6.6,0);
 \draw (2.2,3.4) to[R,l=$\SI{12}{\ohm}$,*-*] (2.2,0);
 \draw (4.4,3.4) to[R,l=$\SI{12}{\ohm}$,*-*] (4.4,0);
 \draw (6.6,3.4) to[R,l=$\SI{12}{\ohm}$] (6.6,0);
 \draw[->,Icol,line width=1pt] (0.5,3.75) -- (1.6,3.75)
       node[midway,above,font=\footnotesize]{$I_t$};
\end{circuitikz}}

\vspace{1.1em}

\subcirc{c}{%
\begin{circuitikz}[scale=0.72,transform shape,line width=0.8pt]
 \draw (0,0) to[battery1,l_={$E=\SI{15}{\volt}$}] (0,3.4);
 \draw (0,3.4) -- (6.6,3.4);  \draw (0,0) -- (6.6,0);
 \draw (2.2,3.4) to[R,l=$\SI{12}{\ohm}$,*-*] (2.2,0);
 \draw (4.4,3.4) to[R,l=$\SI{12}{\ohm}$,*-*] (4.4,0);
 \draw (6.6,3.4) to[R,l=$\SI{3}{\ohm}$] (6.6,0);
 \draw[->,Icol,line width=1pt] (0.5,3.75) -- (1.6,3.75)
       node[midway,above,font=\footnotesize]{$I_t$};
\end{circuitikz}}
\hfill
\subcirc{d}{%
\begin{circuitikz}[scale=0.72,transform shape,line width=0.8pt]
 \draw (0,0) to[battery1,l_={$E=\SI{9}{\volt}$}] (0,3.4);
 \draw (0,3.4) -- (4.4,3.4);  \draw (0,0) -- (4.4,0);
 \draw (2.2,3.4) to[R,l=$\SI{30}{\ohm}$,*-*] (2.2,0);
 \draw (4.4,3.4) to[R,l=$\SI{60}{\ohm}$] (4.4,0);
 \draw[->,Icol,line width=1pt] (0.5,3.75) -- (1.6,3.75)
       node[midway,above,font=\footnotesize]{$I_t$};
\end{circuitikz}}}

\resposta{(a) \SI{1.2}{\ohm}, $I_t=\SI{25}{\ampere}$;
(b) \SI{4}{\ohm}, $I_t=\SI{6}{\ampere}$;
(c) \SI{2}{\ohm}, $I_t=\SI{7.5}{\ampere}$;
(d) \SI{20}{\ohm}, $I_t=\SI{0.45}{\ampere}$}
\resolucao{Em todos os casos a tensão da fonte aparece \emph{inteira} sobre cada
ramo, então cada corrente sai direto da lei de Ohm.\\
\textbf{(a)} $I_1=30/4=\SI{7.5}{\ampere}$, $I_2=30/3=\SI{10}{\ampere}$,
$I_3=30/6=\SI{5}{\ampere}$, $I_4=30/12=\SI{2.5}{\ampere}$;
$I_t=\SI{25}{\ampere}$ e $\Req=30/25=\SI{1.2}{\ohm}$.\\
\textbf{(b)} Três resistores iguais: $\Req = 12/3 = \SI{4}{\ohm}$;
cada ramo com $24/12=\SI{2}{\ampere}$ e $I_t = \SI{6}{\ampere}$.\\
\textbf{(c)} $\dfrac{1}{\Req}=\dfrac{1}{12}+\dfrac{1}{12}+\dfrac{1}{3}
=\dfrac{1+1+4}{12}=\dfrac{6}{12}$, logo $\Req=\SI{2}{\ohm}$.
$I_1=I_2=\SI{1.25}{\ampere}$, $I_3=\SI{5}{\ampere}$, $I_t=\SI{7.5}{\ampere}$.\\
\textbf{(d)} $\Req=\dfrac{30\cdot 60}{30+60}=\SI{20}{\ohm}$;
$I_1=\SI{0.3}{\ampere}$, $I_2=\SI{0.15}{\ampere}$, $I_t=\SI{0.45}{\ampere}$.\\
Note em (a) que $\Req$ ficou menor que o menor resistor (\SI{3}{\ohm}) --- sempre
é assim no paralelo.}
\end{questao}

\begin{questao}[1]{}
$n$ resistores iguais de $\SI{60}{\ohm}$ são associados em paralelo, resultando
em $\SI{10}{\ohm}$. O valor de $n$ é:
\begin{alternativas}[3]
\item 2
\item 3
\item 4
\item 6
\item 10
\end{alternativas}
\resposta{(d)}
\resolucao{Para $n$ resistores iguais em paralelo, $\Req = \dfrac{R}{n}$, logo
$n = \dfrac{60}{10} = 6$.}
\end{questao}

\blocoaplicacao

\begin{questao}[2]{}
Um resistor de $\SI{20}{\ohm}$ é associado em paralelo com um resistor $R$
desconhecido, resultando em uma resistência equivalente de $\SI{12}{\ohm}$.
Determine $R$.
\resposta{$R = \SI{30}{\ohm}$}
\resolucao{Da definição de paralelo,
\[
\frac{1}{12}=\frac{1}{20}+\frac{1}{R}
\quad\Rightarrow\quad
\frac{1}{R}=\frac{1}{12}-\frac{1}{20}=\frac{5-3}{60}=\frac{2}{60}
\quad\Rightarrow\quad R=\SI{30}{\ohm}.
\]
\textbf{Atalho útil:} com dois resistores,
$R = \dfrac{R_1\,\Req}{R_1-\Req} = \dfrac{20\cdot 12}{20-12} = \SI{30}{\ohm}$.
Note que $\Req$ precisa necessariamente ser menor que \SI{20}{\ohm}; se o problema
pedisse $\Req = \SI{25}{\ohm}$, não haveria solução com resistor positivo.}
\end{questao}

\begin{questao}[2]{}
Uma corrente total de $\SI{12}{\ampere}$ chega a um nó e se divide entre dois
resistores em paralelo, de $\SI{4}{\ohm}$ e $\SI{12}{\ohm}$. As correntes em cada
resistor valem, respectivamente:
\begin{alternativas}[2]
\item \SI{3}{\ampere} e \SI{9}{\ampere}
\item \SI{9}{\ampere} e \SI{3}{\ampere}
\item \SI{6}{\ampere} e \SI{6}{\ampere}
\item \SI{4}{\ampere} e \SI{8}{\ampere}
\item \SI{8}{\ampere} e \SI{4}{\ampere}
\end{alternativas}
\resposta{(b)}
\resolucao{A corrente se divide de forma \emph{inversamente} proporcional às
resistências: o menor resistor recebe a maior corrente. Pelo divisor de corrente,
\[
I_1 = I_t\,\frac{R_2}{R_1+R_2}=12\cdot\frac{12}{16}=\SI{9}{\ampere},
\qquad
I_2 = 12-9 = \SI{3}{\ampere}.
\]
\textbf{Verificação alternativa:} $\Req = \dfrac{4\cdot12}{16}=\SI{3}{\ohm}$ e
$U = 3\cdot 12 = \SI{36}{\volt}$; daí $I_1 = 36/4 = \SI{9}{\ampere}$ e
$I_2 = 36/12 = \SI{3}{\ampere}$. A armadilha da alternativa (a) é aplicar a
proporção direta, como se fosse divisor de tensão.}
\end{questao}

\begin{questao}[2]{}
Três lâmpadas incandescentes de $\SI{60}{\watt}$/$\SI{120}{\volt}$ são ligadas
\textbf{em paralelo} à rede de $\SI{120}{\volt}$.
\begin{itensq}
\item Calcule a resistência de cada lâmpada e a resistência equivalente.
\item Determine a corrente total fornecida pela rede.
\item Se uma das lâmpadas queimar, o que acontece com o brilho das outras duas?
\end{itensq}
\resposta{(a) \SI{240}{\ohm} cada, $\Req=\SI{80}{\ohm}$;
(b) \SI{1.5}{\ampere}; (c) nada muda}
\resolucao{(a) De $P = \dfrac{U^2}{R}$:
$R = \dfrac{120^2}{60} = \SI{240}{\ohm}$ por lâmpada.
Três iguais em paralelo: $\Req = \dfrac{240}{3} = \SI{80}{\ohm}$.\\
(b) $I_t = \dfrac{120}{80} = \SI{1.5}{\ampere}$
(ou $3\times\SI{0.5}{\ampere}$ por lâmpada).\\
(c) As demais continuam com os mesmos \SI{120}{\volt} nos seus terminais, logo o
brilho \emph{não muda}: cada ramo é independente. É por isso que as instalações
residenciais são feitas em paralelo, e não em série.}
\end{questao}

\begin{questao}[2]{}
No trecho da figura, um fio ideal (sem resistência) é ligado em paralelo com o
resistor $R_2$.

\circuitoq{%
\begin{circuitikz}[scale=0.95,transform shape,line width=0.8pt]
 \draw (0,0) to[battery1,l_={$E=\SI{20}{\volt}$}] (0,3)
       to[R,l={$R_1=\SI{10}{\ohm}$}] (3,3)
       to[R,l={$R_2=\SI{20}{\ohm}$}] (6,3) -- (6,0) -- (0,0);
 \draw[Vcol,line width=1.1pt] (3,3) -- (3,4.2) -- (6,4.2) -- (6,3);
 \node[Vcol,font=\footnotesize\sffamily,above] at (4.5,4.25) {fio ideal};
 \fill (3,3) circle (0.06); \fill (6,3) circle (0.06);
\end{circuitikz}}

A corrente fornecida pela fonte vale:
\begin{alternativas}[3]
\item \SI{0.67}{\ampere}
\item \SI{1}{\ampere}
\item \SI{1.5}{\ampere}
\item \SI{2}{\ampere}
\item \SI{3}{\ampere}
\end{alternativas}
\resposta{(d)}
\resolucao{O fio ideal curto-circuita $R_2$: os dois terminais de $R_2$ ficam no
\emph{mesmo potencial}, logo $U_{R_2} = 0$ e nenhuma corrente passa por ele.
Formalmente, $R_2 \parallel 0 = 0$. Resta apenas $R_1$:
\[
I=\frac{20}{10}=\SI{2}{\ampere}.
\]
\textbf{Erro comum:} calcular $10 + 20 = \SI{30}{\ohm}$ e responder
\SI{0.67}{\ampere}, ignorando o curto.}
\end{questao}

\blocoaprofundamento

\begin{questao}[3]{}
Um resistor de $\SI{100}{\ohm}$ suporta no máximo $\SI{2}{\watt}$ de dissipação.
Deseja-se construir, com resistores desse mesmo tipo, uma associação
\textbf{em paralelo} de $\SI{25}{\ohm}$.
\begin{itensq}
\item Quantos resistores são necessários?
\item Qual é a máxima tensão que pode ser aplicada à associação?
\item Qual é a potência máxima que a associação suporta?
\end{itensq}
\resposta{(a) 4 resistores; (b) $\approx\SI{14.1}{\volt}$;
(c) \SI{8}{\watt}}
\resolucao{(a) $\Req = R/n \Rightarrow n = 100/25 = 4$ resistores.\\
(b) Todos ficam sob a \emph{mesma} tensão, então basta respeitar o limite de um
deles:
\[
P_{\max}=\frac{U^2}{R}\ \Rightarrow\ U_{\max}=\sqrt{P_{\max}R}=\sqrt{2\cdot100}
=\sqrt{200}\approx\SI{14.1}{\volt}.
\]
(c) Como os 4 resistores trabalham no limite simultaneamente,
$P_{\text{total}} = 4\cdot 2 = \SI{8}{\watt}$ --- que também pode ser conferido
por $P = U^2/\Req = 200/25 = \SI{8}{\watt}$.\\
\textbf{Ideia geral:} associar resistores iguais em paralelo (ou em série)
multiplica por $n$ a capacidade de dissipação do conjunto.}
\end{questao}

\begin{questao}[3]{}
Determine a resistência equivalente entre os terminais A e B do circuito da
figura, no qual dois trechos são interligados por fios ideais.

\circuitoq{%
\begin{circuitikz}[scale=0.85,transform shape,line width=0.8pt]
 \draw (0,0) node[left,font=\bfseries] {A}
       to[R,l=$\SI{15}{\ohm}$,o-*] (3,0)
       to[R,l=$\SI{30}{\ohm}$,-*] (6,0)
       to[R,l=$\SI{30}{\ohm}$,-o] (9,0)
       node[right,font=\bfseries] {B};
 \draw[Vcol,line width=1.1pt] (3,0) -- (3,1.8) -- (6,1.8) -- (6,0);
 \node[Vcol,font=\footnotesize\sffamily,above] at (4.5,1.85) {fio ideal};
\end{circuitikz}}

\resposta{$R_{AB} = \SI{45}{\ohm}$}
\resolucao{O fio ideal liga os dois terminais do resistor de \SI{30}{\ohm}
central, curto-circuitando-o. Ele deixa de conduzir e pode ser desconsiderado no circuito equivalente. Sobram
apenas os resistores de \SI{15}{\ohm} e \SI{30}{\ohm} em série:
\[
R_{AB}=15+30=\SI{45}{\ohm}.
\]
\textbf{Método seguro:} antes de calcular, redesenhe o circuito nomeando os nós.
Dois pontos ligados por fio ideal são \emph{o mesmo nó} --- qualquer componente
entre eles está curto-circuitado.}
\end{questao}

\blocodesafio

\begin{questao}[4]{}
$n$ resistores \emph{diferentes} $R_1, R_2, \dots, R_n$ são associados em
paralelo.
\begin{itensq}
\item Demonstre que a resistência equivalente é sempre \textbf{menor} que a menor
      das resistências.
\item Mostre que, se um dos resistores tende a zero, $\Req \to 0$, e interprete
      fisicamente.
\item Dois resistores em paralelo dão $\SI{6}{\ohm}$; em série, dariam
      $\SI{25}{\ohm}$. Determine seus valores.
\end{itensq}
\resposta{(a) e (b) demonstrações; (c) $\SI{10}{\ohm}$ e $\SI{15}{\ohm}$}
\resolucao{\textbf{(a)} Seja $R_{\min}$ a menor das resistências. Da definição,
\[
\frac{1}{\Req}=\frac{1}{R_1}+\dots+\frac{1}{R_n}
> \frac{1}{R_{\min}},
\]
pois a soma inclui o termo $\dfrac{1}{R_{\min}}$ \emph{mais} outras parcelas
positivas. Invertendo a desigualdade (ambos os lados positivos):
$\Req < R_{\min}$. \checkmark\\
Interpretação: acrescentar um caminho \emph{sempre} facilita a passagem de
corrente.

\textbf{(b)} Se $R_k \to 0$, então $\dfrac{1}{R_k}\to\infty$, e a soma das
condutâncias diverge, levando $\Req \to 0$. Fisicamente, esse ramo é um
\textbf{curto-circuito}: ele desvia toda a corrente e anula a tensão sobre o
conjunto, tornando irrelevantes os demais resistores.

\textbf{(c)} Temos o sistema
\[
\begin{cases}
R_1+R_2 = 25\\[2pt]
\dfrac{R_1R_2}{R_1+R_2}=6 \;\Rightarrow\; R_1R_2 = 6\cdot25 = 150.
\end{cases}
\]
Soma e produto conhecidos: $R_1$ e $R_2$ são as raízes de
$x^2-25x+150=0$, isto é,
\[
x=\frac{25\pm\sqrt{625-600}}{2}=\frac{25\pm5}{2}
\;\Longrightarrow\;
R_1=\SI{15}{\ohm},\quad R_2=\SI{10}{\ohm}.
\]
\textbf{Verificação:} $15+10=25$ \checkmark\ e
$\dfrac{15\cdot10}{25}=\SI{6}{\ohm}$ \checkmark\\
\textbf{Observação:} o problema só tem solução real se
$(R_s)^2 \ge 4R_sR_p$, ou seja, $R_s \ge 4R_p$. Aqui, $25 \ge 24$ --- a condição é satisfeita por pequena margem. Se o enunciado pedisse $\Req=\SI{7}{\ohm}$ em paralelo com a mesma soma
$\SI{25}{\ohm}$, não existiriam resistores reais que satisfizessem.}
\end{questao}

% BEGIN BANCO COMPLEMENTAR Circuitos em paralelo
% =====================================================================
%  Banco complementar --- Circuitos Paralelo  (REVISADO)
%  Substitui a versao gerada por template (11 clones em N2, 13 em N3,
%  9 questoes conceituais indevidamente marcadas como N4).
%  Sem sobreposicao com o cap03 (n resistores iguais, R desconhecido a
%  partir de Req, divisor 12 A entre 4 e 12 ohm, tres lampadas de 60 W,
%  fio ideal em paralelo, banco de 100 ohm/2 W e a demonstracao de que
%  Req < menor resistor ja estao la).
%  Distribuicao mantida: 7 N1 + 11 N2 + 13 N3 + 9 N4 = 40
% =====================================================================

\subsubsection*{Banco complementar --- Circuitos paralelo}

\begin{atencao}[Organização do banco]
No paralelo, a \textbf{tensão é comum} e as \textbf{correntes se somam}. Daí
decorre que corrente e potência se distribuem \emph{inversamente} a $R$ ---
exatamente o oposto da série. Trabalhe com \textbf{condutâncias}: elas somam
diretamente e tornam quase toda conta trivial. O N3 explora falhas de ramo,
fuga, deriva térmica e instrumentação; o N4 trata de projeto de \emph{shunt},
sensibilidade, propagação de incerteza e seletividade de proteção.
\end{atencao}

\blocofixacao

\begin{questao}[1]{}
Determine a resistência equivalente de $R_1=\SI{30}{\ohm}$ e
$R_2=\SI{60}{\ohm}$ em paralelo.
\resposta{$R_{eq}=\SI{20}{\ohm}$}
\resolucao{Para dois resistores, a regra do produto pela soma é a mais rápida:
\[
R_{eq}=\frac{R_1R_2}{R_1+R_2}=\frac{30\cdot60}{90}=\SI{20}{\ohm}.
\]
\textbf{Conferência instantânea:} o resultado ($\SI{20}{\ohm}$) é menor que o
menor dos dois ($\SI{30}{\ohm}$) --- se não for, há erro. Cuidado: a regra
produto/soma vale \emph{apenas} para dois resistores.}
\end{questao}

\begin{questao}[1]{}
Três resistores de $\SI{4}{\ohm}$, $\SI{8}{\ohm}$ e $\SI{8}{\ohm}$ estão em
paralelo. Calcule a condutância total e $R_{eq}$.
\resposta{$G=\SI{0.5}{\siemens}$; $R_{eq}=\SI{2}{\ohm}$}
\resolucao{
\[
G=\frac14+\frac18+\frac18=0{,}25+0{,}125+0{,}125=\SI{0.5}{\siemens}
\;\Longrightarrow\;
R_{eq}=\frac{1}{0{,}5}=\SI{2}{\ohm}.
\]
Com três ou mais ramos, some as condutâncias e inverta \emph{uma única vez} no
final. Inverter parcialmente a cada passo é a origem mais comum de erro.}
\end{questao}

\begin{questao}[1]{}
Três resistores iguais de $\SI{90}{\ohm}$ são associados em paralelo. Determine
$R_{eq}$.
\resposta{$R_{eq}=\SI{30}{\ohm}$}
\resolucao{Para $n$ resistores \emph{iguais},
\[
R_{eq}=\frac{R}{n}=\frac{90}{3}=\SI{30}{\ohm}.
\]
O atalho vale só quando todos são iguais; com valores diferentes é obrigatório
somar condutâncias.}
\end{questao}

\begin{questao}[1]{}
Acrescentando-se mais um ramo \textbf{em paralelo} a um circuito alimentado por
fonte ideal de tensão, a corrente total:
\begin{alternativas}[2]
\item aumenta
\item diminui
\item não se altera
\item pode aumentar ou diminuir, conforme o valor
\end{alternativas}
\resposta{(a)}
\resolucao{Cada novo ramo acrescenta uma condutância positiva:
$G_{eq}=\sum G_k$ só pode crescer, logo $R_{eq}$ só pode diminuir. Com $V$
fixo, $I=V/R_{eq}$ aumenta.\\
Note o contraste com a série, em que acrescentar elemento \emph{reduz} a
corrente. E note também que os ramos já existentes continuam com \emph{a mesma}
corrente --- quem cresce é o total.}
\end{questao}

\begin{questao}[1]{}
Um ramo de $\SI{48}{\ohm}$ pertence a uma associação paralela ligada a
$\SI{24}{\volt}$. Determine a corrente desse ramo.
\resposta{$I=\SI{0.5}{\ampere}$}
\resolucao{
\[
I=\frac{V}{R}=\frac{24}{48}=\SI{0.5}{\ampere}.
\]
No paralelo, a tensão é \emph{comum} a todos os ramos: a corrente de cada um
depende só do próprio resistor. Não é preciso conhecer os demais ramos nem
$R_{eq}$.}
\end{questao}

\begin{questao}[1]{}
Uma associação paralela sob $\SI{40}{\volt}$ tem ramos que conduzem
$\SI{2}{\ampere}$, $\SI{3}{\ampere}$ e $\SI{5}{\ampere}$. Determine a corrente
total e $R_{eq}$.
\resposta{$I_t=\SI{10}{\ampere}$; $R_{eq}=\SI{4}{\ohm}$}
\resolucao{Pela LKC, as correntes de ramo se somam:
$I_t=2+3+5=\SI{10}{\ampere}$.
\[
R_{eq}=\frac{V}{I_t}=\frac{40}{10}=\SI{4}{\ohm}.
\]
Observe que não foi preciso conhecer os resistores individualmente --- embora
eles saiam de imediato: $20$, $13{,}3$ e $\SI{8}{\ohm}$.}
\end{questao}

\begin{questao}[1]{}
Em uma associação paralela, o ramo que conduz a \textbf{maior} corrente é:
\begin{alternativas}[2]
\item o de maior resistência
\item o de menor resistência
\item o mais próximo da fonte
\item todos conduzem igualmente
\end{alternativas}
\resposta{(b)}
\resolucao{Com $V$ comum, $I_k=V/R_k$: a corrente é \emph{inversamente}
proporcional à resistência. O ramo de menor $R$ é o caminho de menor oposição e
leva a maior corrente.\\
A alternativa (c) é sedutora, mas falsa: em circuitos de parâmetros
concentrados, a posição física dos ramos é irrelevante --- só importa a
topologia.}
\end{questao}

\blocoaplicacao

\begin{questao}[2]{}
$R_1=\SI{10}{\ohm}$, $R_2=\SI{20}{\ohm}$ e $R_3=\SI{20}{\ohm}$ estão em
paralelo sob $\SI{60}{\volt}$. Determine as correntes de ramo, a corrente total
e $R_{eq}$.
\resposta{$\SI{6}{\ampere}$, $\SI{3}{\ampere}$, $\SI{3}{\ampere}$;
$I_t=\SI{12}{\ampere}$; $R_{eq}=\SI{5}{\ohm}$}
\resolucao{Tensão comum, então cada ramo é independente:
\[
I_1=\frac{60}{10}=6,\quad I_2=I_3=\frac{60}{20}=\SI{3}{\ampere}.
\]
\[
I_t=6+3+3=\SI{12}{\ampere}
\;\Longrightarrow\;
R_{eq}=\frac{60}{12}=\SI{5}{\ohm}.
\]
\textbf{Conferência:} $G=0{,}1+0{,}05+0{,}05=\SI{0.2}{\siemens}$ e
$1/0{,}2=\SI{5}{\ohm}$ \checkmark.}
\end{questao}

\begin{questao}[2]{}
No circuito da questão anterior, calcule a potência de cada ramo e confronte
com a potência total.
\resposta{$\SI{360}{\watt}$, $\SI{180}{\watt}$, $\SI{180}{\watt}$; total
$\SI{720}{\watt}$}
\resolucao{Com $V$ comum, use $P=V^{2}/R$:
\[
P_1=\frac{3600}{10}=360,\qquad P_2=P_3=\frac{3600}{20}=\SI{180}{\watt}.
\]
Fonte: $P=VI_t=60\cdot12=\SI{720}{\watt}=360+180+180$ \checkmark.\\
\textbf{Regra do paralelo:} $P\propto1/R$ --- o \emph{menor} resistor dissipa
mais. Compare com a série, em que $P\propto R$. Trocar essas duas regras é o
erro conceitual mais frequente em associações.}
\end{questao}

\begin{questao}[2]{}
Uma corrente de $\SI{15}{\ampere}$ chega a um nó e se divide entre
$R_1=\SI{8}{\ohm}$ e $R_2=\SI{12}{\ohm}$. Determine as correntes de ramo pela
fórmula do divisor.
\resposta{$I_1=\SI{9}{\ampere}$, $I_2=\SI{6}{\ampere}$}
\resolucao{Para \emph{dois} ramos, a corrente de um deles é proporcional à
resistência do \textbf{outro}:
\[
I_1=I_t\frac{R_2}{R_1+R_2}=15\cdot\frac{12}{20}=\SI{9}{\ampere},
\qquad
I_2=15\cdot\frac{8}{20}=\SI{6}{\ampere}.
\]
\textbf{Verificação:} $9+6=\SI{15}{\ampere}$ \checkmark, e a maior corrente ficou
no \emph{menor} resistor \checkmark. O "cruzamento" dos índices é o ponto que
mais confunde --- ele existe porque $I_k\propto G_k$, e escrever a fórmula em
condutâncias ($I_1=I_tG_1/(G_1+G_2)$) elimina a inversão.}
\end{questao}

\begin{questao}[2]{}
Três resistores em paralelo sob $\SI{100}{\volt}$ drenam $\SI{5}{\ampere}$ no
total. Dois deles valem $\SI{50}{\ohm}$ e $\SI{100}{\ohm}$. Determine o
terceiro.
\resposta{$R_3=\SI{50}{\ohm}$}
\resolucao{Trabalhe em condutâncias, onde a subtração é direta:
\[
G_{eq}=\frac{I_t}{V}=\frac{5}{100}=\SI{0.05}{\siemens},
\qquad
G_1+G_2=0{,}02+0{,}01=\SI{0.03}{\siemens}.
\]
\[
G_3=0{,}05-0{,}03=\SI{0.02}{\siemens}
\;\Longrightarrow\;
R_3=\SI{50}{\ohm}.
\]
Tentar subtrair \emph{resistências} aqui levaria a um resultado sem sentido:
no paralelo, o que é aditivo é $G$, não $R$.}
\end{questao}

\begin{questao}[2]{}
Duas lâmpadas de resistências $\SI{240}{\ohm}$ e $\SI{960}{\ohm}$ são ligadas
\textbf{em paralelo} a $\SI{120}{\volt}$. Qual brilha mais? Compare com o
resultado da mesma dupla em série.
\resposta{Em paralelo, a de $\SI{240}{\ohm}$ ($\SI{60}{\watt}$ contra
$\SI{15}{\watt}$) --- inverte-se em relação à série}
\resolucao{
\[
P_{240}=\frac{120^{2}}{240}=\SI{60}{\watt},
\qquad
P_{960}=\frac{120^{2}}{960}=\SI{15}{\watt}.
\]
Em paralelo, $V$ é comum e $P\propto1/R$: brilha mais a de \emph{menor}
resistência. Em série (mesma dupla), $I$ é comum e $P\propto R$: brilharia mais
a de $\SI{960}{\ohm}$.\\
\textbf{Conclusão que vale guardar:} não existe "a lâmpada mais potente" em
termos absolutos --- a potência depende da \emph{ligação}. É por isso que
lâmpadas residenciais, sempre em paralelo, obedecem à intuição usual
($\SI{100}{\watt}$ brilha mais que $\SI{25}{\watt}$), mas guirlandas em série
se comportam ao contrário.}
\end{questao}

\begin{questao}[2]{}
Um galvanômetro tem resistência interna $\SI{100}{\ohm}$ e deflexão total com
$\SI{1}{\milli\ampere}$. Determine o resistor \emph{shunt} que amplia a escala
para $\SI{11}{\milli\ampere}$.
\resposta{$R_{sh}=\SI{10}{\ohm}$}
\resolucao{O shunt fica em paralelo com o galvanômetro e desvia o excedente:
$I_{sh}=11-1=\SI{10}{\milli\ampere}$. Como a tensão é comum,
\[
R_{sh}I_{sh}=R_gI_g
\;\Longrightarrow\;
R_{sh}=\frac{R_gI_g}{I_{sh}}=\frac{100\cdot1}{10}=\SI{10}{\ohm}.
\]
A regra geral é $R_{sh}=\dfrac{R_g}{n-1}$, com $n=I_{total}/I_g$ o fator de
ampliação --- aqui $n=11$ e $R_{sh}=100/10$ \checkmark.}
\end{questao}

\begin{questao}[2]{}
Um circuito residencial de $\SI{220}{\volt}$ alimenta, em paralelo, cargas de
$\SI{4400}{\watt}$, $\SI{1100}{\watt}$ e $\SI{550}{\watt}$. Determine a
corrente total e escolha o disjuntor.
\resposta{$I_t=\SI{27.5}{\ampere}$; disjuntor de $\SI{32}{\ampere}$}
\resolucao{Cada carga drena $I=P/V$:
\[
I_1=\frac{4400}{220}=20,\quad
I_2=\frac{1100}{220}=5,\quad
I_3=\frac{550}{220}=\SI{2.5}{\ampere}.
\]
Total: $\SI{27.5}{\ampere}$ (as correntes se somam porque as cargas estão em
paralelo).\\
O disjuntor comercial imediatamente superior é o de $\SI{32}{\ampere}$ ---
$\SI{25}{\ampere}$ seria insuficiente. \textbf{Atenção:} escolhido o disjuntor,
a seção do condutor deve suportar \emph{a corrente do disjuntor}, não a da
carga; caso contrário, o cabo vira o elo fraco da proteção.}
\end{questao}

\begin{questao}[2]{}
Dois condutores idênticos de $\SI{0.8}{\ohm}$ cada são ligados em paralelo para
alimentar uma carga. Determine a resistência resultante e explique a prática.
\resposta{$\SI{0.4}{\ohm}$ --- metade}
\resolucao{$R_{eq}=0{,}8/2=\SI{0.4}{\ohm}$.\\
Dobrar o número de condutores idênticos equivale a dobrar a seção: a
resistência cai à metade, e com ela a queda de tensão e as perdas
$RI^{2}$ na linha.\\
\textbf{Condição prática:} os condutores devem ter mesmo comprimento, seção e
material. Se um for mais curto, ele terá menor resistência e assumirá corrente
desproporcional --- podendo sobreaquecer enquanto o outro fica ocioso. Por isso
condutores em paralelo são sempre especificados com o mesmo lance.}
\end{questao}

\begin{questao}[2]{}
Um resistor de $\SI{1}{\mega\ohm}$ é ligado em paralelo com um de
$\SI{100}{\ohm}$. A resistência equivalente:
\begin{alternativas}[2]
\item fica praticamente igual a $\SI{100}{\ohm}$
\item fica praticamente igual a $\SI{1}{\mega\ohm}$
\item vale aproximadamente $\SI{500}{\kilo\ohm}$
\item vale exatamente $\SI{50}{\ohm}$
\end{alternativas}
\resposta{(a)}
\resolucao{$G=0{,}01+10^{-6}\approx\SI{0.010001}{\siemens}$, logo
$R_{eq}\approx\SI{99.99}{\ohm}$ --- diferença de $0{,}01\%$.\\
\textbf{Princípio geral:} no paralelo, quem manda é o \emph{menor} resistor;
ramos muito maiores são praticamente circuitos abertos. É exatamente por isso
que um voltímetro de alta impedância "não carrega" o circuito sob medição.}
\end{questao}

\begin{questao}[2]{}
Um chuveiro possui duas resistências de $\SI{22}{\ohm}$ que podem ser ligadas
isoladamente ou em paralelo, sob $\SI{220}{\volt}$. Determine a potência em
cada caso.
\resposta{$\SI{2200}{\watt}$ com uma; $\SI{4400}{\watt}$ com as duas}
\resolucao{Uma resistência:
$P=\dfrac{220^{2}}{22}=\SI{2200}{\watt}$.\\
Duas em paralelo: $R_{eq}=\SI{11}{\ohm}$ e
$P=\dfrac{220^{2}}{11}=\SI{4400}{\watt}$.\\
A potência \emph{dobra} porque a resistência caiu à metade sob a mesma tensão.
É assim que funcionam as chaves inverno/verão. \textbf{Consequência para a
instalação:} a corrente também dobra ($10\to\SI{20}{\ampere}$), e o
dimensionamento de disjuntor e condutor deve ser feito para a posição de maior
potência.}
\end{questao}

\begin{questao}[2]{}
Em uma associação paralela mediram-se as correntes de ramo $\SI{4}{\ampere}$,
$\SI{1}{\ampere}$ e $\SI{1}{\ampere}$, com corrente total de
$\SI{6}{\ampere}$ e tensão de $\SI{20}{\volt}$. Verifique a consistência e
determine as três resistências e a potência total.
\resposta{Consistente; $\SI{5}{\ohm}$, $\SI{20}{\ohm}$, $\SI{20}{\ohm}$;
$P=\SI{120}{\watt}$}
\resolucao{\textbf{LKC:} $4+1+1=\SI{6}{\ampere}$ \checkmark.\\
\textbf{Resistências:} $R_k=V/I_k$ $\Rightarrow$ $20/4=5$ e
$20/1=\SI{20}{\ohm}$ (duas vezes).\\
\textbf{Potência:} $P=VI_t=20\cdot6=\SI{120}{\watt}$; por ramo,
$80+20+20=\SI{120}{\watt}$ \checkmark.\\
Dois testes independentes --- LKC e balanço de potência --- fecham. Se apenas o
primeiro fechasse, o erro estaria na leitura da tensão.}
\end{questao}

\blocoaprofundamento

\begin{questao}[3]{}
Uma corrente total de $\SI{6}{\ampere}$ divide-se entre $\SI{5}{\ohm}$,
$\SI{20}{\ohm}$ e $\SI{20}{\ohm}$ em paralelo.
\begin{itensq}
\item Determine $R_{eq}$ e a tensão da associação.
\item Obtenha as três correntes pelo divisor em condutâncias.
\item Verifique por dois caminhos independentes.
\end{itensq}
\resposta{$R_{eq}=\SI{3.33}{\ohm}$, $V=\SI{20}{\volt}$;
$\SI{4}{\ampere}$, $\SI{1}{\ampere}$, $\SI{1}{\ampere}$}
\resolucao{(a) $G=0{,}2+0{,}05+0{,}05=\SI{0.3}{\siemens}$,
$R_{eq}=1/0{,}3=\SI{3.33}{\ohm}$ e $V=I_tR_{eq}=6\cdot3{,}33=\SI{20}{\volt}$.\\
(b) Com três ou mais ramos, a forma correta do divisor usa condutâncias:
\[
I_k=I_t\frac{G_k}{G_{eq}}
\;\Rightarrow\;
I_1=6\cdot\frac{0{,}2}{0{,}3}=4,\quad
I_2=I_3=6\cdot\frac{0{,}05}{0{,}3}=\SI{1}{\ampere}.
\]
\textbf{Alerta:} a fórmula $I_1=I_tR_2/(R_1+R_2)$ \emph{não} se generaliza para
três ramos --- ela é um caso particular que só funciona com dois.\\
(c) \emph{LKC:} $4+1+1=6$ \checkmark. \emph{Lei de Ohm por ramo:}
$20/5=4$ e $20/20=1$ \checkmark.}
\end{questao}

\begin{questao}[3]{}
Três ramos de $\SI{10}{\ohm}$, $\SI{20}{\ohm}$ e $\SI{20}{\ohm}$ operam sob
fonte \textbf{ideal} de $\SI{60}{\volt}$. O ramo de $\SI{10}{\ohm}$ se
\textbf{abre}. Determine o que muda nas correntes de ramo, na corrente total e
em $R_{eq}$.
\resposta{Ramos restantes inalterados ($\SI{3}{\ampere}$ cada);
$I_t$: $12\to\SI{6}{\ampere}$; $R_{eq}$: $5\to\SI{10}{\ohm}$}
\resolucao{Com fonte ideal, $V=\SI{60}{\volt}$ permanece fixo. Como cada ramo
depende \emph{só} da tensão comum e do próprio resistor, os dois ramos íntegros
continuam com $60/20=\SI{3}{\ampere}$ --- \textbf{nada muda para eles}.\\
O que muda é o agregado: $I_t=3+3=\SI{6}{\ampere}$ (queda de $50\%$) e
$R_{eq}=\SI{10}{\ohm}$.\\
\textbf{Por isso a ligação residencial é em paralelo:} a queima de uma lâmpada
não afeta as demais. Em série, uma falha derruba o circuito inteiro.\\
\emph{Ressalva importante:} isso depende de a fonte ser ideal. Com resistência
interna, $V$ sobe quando um ramo abre --- caso tratado no bloco de desafio.}
\end{questao}

\begin{questao}[3]{}
Um dos ramos de uma associação paralela sofre \textbf{curto-circuito}. Analise
$R_{eq}$, a corrente total e o efeito sobre os demais ramos, e justifique a
necessidade de proteção.
\resposta{$R_{eq}\to0$; $I_t\to\infty$ (limitada só pela fonte e pela fiação);
os demais ramos ficam sem tensão}
\resolucao{Um ramo em curto tem $R=0$, portanto $G\to\infty$. Como
$G_{eq}=\sum G_k$, a condutância total explode e $R_{eq}\to0$:
\[
R_{eq}=\frac{1}{\infty+G_2+G_3}\to0.
\]
Com fonte ideal, $I_t=V/R_{eq}\to\infty$. Na prática, a corrente é limitada
pela resistência interna da fonte e pela impedância da fiação --- mas ainda
assim atinge dezenas ou centenas de ampères.\\
\textbf{Efeito sobre os demais:} o curto derruba a tensão do barramento a
praticamente zero. Os outros ramos, embora íntegros, \emph{param de funcionar}
--- eles são vítimas, não causa.\\
\textbf{Proteção:} é este o cenário que o disjuntor combate. Diferentemente do
ramo aberto (falha benigna e localizada), o curto em um ramo compromete o
sistema inteiro e libera energia suficiente para incendiar a instalação. Daí a
assimetria de tratamento: monitoram-se \emph{sobrecorrentes}, não
subcorrentes.}
\end{questao}

\begin{questao}[3]{}
Dispõe-se de resistores de $\SI{220}{\ohm}$ com dissipação máxima de
$\SI{1}{\watt}$. Projete uma associação paralela de $\SI{55}{\ohm}$ capaz de
dissipar pelo menos $\SI{4}{\watt}$, e determine a tensão máxima aplicável.
\resposta{Quatro em paralelo; $P_{\max}=\SI{4}{\watt}$;
$V_{\max}=\SI{14.8}{\volt}$}
\resolucao{\textbf{Valor:} $220/n=55\Rightarrow n=4$.\\
\textbf{Potência:} sendo iguais e sob a mesma tensão, os quatro dividem a carga
térmica igualmente, então $P_{tot}=4\times1=\SI{4}{\watt}$ \checkmark.\\
\textbf{Tensão máxima:} limita-se pelo componente individual:
\[
V_{\max}=\sqrt{P_1R_1}=\sqrt{1\cdot220}=\SI{14.8}{\volt}.
\]
\textbf{Conferência:} $P_{tot}=V^{2}/R_{eq}=220/55=\SI{4}{\watt}$ \checkmark.\\
\textbf{Observação de projeto:} o paralelo multiplica a capacidade térmica
\emph{e} divide o valor por $n$; a série multiplica ambos. Se o objetivo exige
valor \emph{maior} com mais potência, use série; valor \emph{menor}, paralelo.
Para manter o valor e aumentar a potência, é preciso uma matriz
série--paralelo de $n\times n$ unidades.}
\end{questao}

\begin{questao}[3]{}
Uma rede tem cinco ramos em paralelo: $\SI{10}{\ohm}$, $\SI{20}{\ohm}$,
$\SI{20}{\ohm}$, $\SI{40}{\ohm}$ e $\SI{40}{\ohm}$. Calcule $R_{eq}$ e explique
por que o cálculo em condutâncias é preferível.
\resposta{$G=\SI{0.25}{\siemens}$; $R_{eq}=\SI{4}{\ohm}$}
\resolucao{
\[
G=0{,}100+0{,}050+0{,}050+0{,}025+0{,}025=\SI{0.25}{\siemens}
\;\Longrightarrow\;
R_{eq}=\SI{4}{\ohm}.
\]
\textbf{Três vantagens da condutância.} \emph{(i)} A operação é uma soma
simples, com um único inverso no fim --- contra quatro aplicações sucessivas da
regra produto/soma. \emph{(ii)} A distribuição de corrente sai de graça:
$I_k/I_t=G_k/G_{eq}$, aqui $40\%$, $20\%$, $20\%$, $10\%$ e $10\%$.
\emph{(iii)} Para incerteza, a propagação em $G$ é \emph{linear}, ao passo que
em $R$ envolve derivadas de uma função racional.\\
É também a razão pela qual a análise nodal --- construída sobre a LKC --- adota
condutâncias como grandeza natural.}
\end{questao}

\begin{questao}[3]{}
Um galvanômetro de $\SI{99}{\ohm}$ e $\SI{1}{\milli\ampere}$ de fundo de escala
deve medir até $\SI{100}{\milli\ampere}$.
\begin{itensq}
\item Determine o shunt.
\item Calcule a resistência do amperímetro resultante.
\item Explique por que essa resistência precisa ser pequena.
\end{itensq}
\resposta{(a) $R_{sh}=\SI{1}{\ohm}$; (b) $\SI{0.99}{\ohm}$}
\resolucao{(a) O shunt conduz $100-1=\SI{99}{\milli\ampere}$:
\[
R_{sh}=\frac{R_gI_g}{I_{sh}}=\frac{99\cdot1}{99}=\SI{1}{\ohm}.
\]
(b) $R_A=R_g\parallel R_{sh}=\dfrac{99\cdot1}{100}=\SI{0.99}{\ohm}$.\\
(c) O amperímetro é inserido \textbf{em série} no ramo sob medição, somando-se
à resistência do circuito. Se $R_A$ não for desprezível diante do circuito, a
própria medição altera a corrente que se pretende medir --- o chamado
\emph{erro de inserção}.\\
\textbf{Exemplo:} medindo um ramo de $\SI{100}{\ohm}$, o erro é
$0{,}99/100\approx1\%$; medindo um de $\SI{1}{\ohm}$, seria de $50\%$.\\
Note a dualidade: o \emph{amperímetro} deve ter resistência quase nula e vai em
série; o \emph{voltímetro}, resistência quase infinita e vai em paralelo.}
\end{questao}

\begin{questao}[3]{}
Uma carga de $\SI{100}{\ohm}$ opera a $\SI{200}{\volt}$. A isolação
degradada cria um caminho de fuga de $\SI{10}{\kilo\ohm}$ em paralelo.
\begin{itensq}
\item Determine a corrente de fuga e a corrente total.
\item Calcule $R_{eq}$ e o erro relativo cometido ao ignorar a fuga.
\item Comente o risco.
\end{itensq}
\resposta{(a) $\SI{20}{\milli\ampere}$ e $\SI{2.02}{\ampere}$;
(b) $\SI{99.01}{\ohm}$, erro de $1\%$; (c) risco de choque}
\resolucao{(a) $I_{carga}=200/100=\SI{2}{\ampere}$;
$I_{fuga}=200/10000=\SI{20}{\milli\ampere}$; total $\SI{2.02}{\ampere}$.\\
(b) $R_{eq}=\dfrac{100\cdot10000}{10100}=\SI{99.01}{\ohm}$ --- desvio de
$\approx1\%$ em relação a $\SI{100}{\ohm}$.\\
(c) Do ponto de vista \emph{elétrico}, $1\%$ é irrelevante: nenhum instrumento
de painel acusaria a anomalia. Do ponto de vista da \textbf{segurança}, é
gravíssimo: $\SI{20}{\milli\ampere}$ atravessando um corpo humano já provoca
tetanização muscular, e a partir de $\SI{30}{\milli\ampere}$ o risco de
fibrilação é real.\\
\textbf{Por isso existe o DR (dispositivo diferencial residual):} ele não mede
a corrente total, e sim o \emph{desequilíbrio} entre fase e neutro --- que é
exatamente a corrente de fuga. Um DR de $\SI{30}{\milli\ampere}$ detecta o que
os $\SI{2}{\ampere}$ da carga mascaram completamente.}
\end{questao}

\begin{questao}[3]{}
$R_1=\SI{100}{\ohm}\pm5\%$ e $R_2=\SI{400}{\ohm}\pm5\%$ estão em paralelo.
Determine $R_{eq}$ nominal e seus extremos de pior caso.
\resposta{$R_{eq}=\SI{80}{\ohm}$, entre $\SI{76}{\ohm}$ e $\SI{84}{\ohm}$
($\pm5\%$)}
\resolucao{Nominal: $R_{eq}=\dfrac{100\cdot400}{500}=\SI{80}{\ohm}$.\\
Como $R_{eq}$ cresce monotonicamente com $R_1$ e com $R_2$, os extremos ocorrem
com \emph{ambos} nos extremos:
\[
R_{eq}^{\min}=\frac{95\cdot380}{475}=\SI{76}{\ohm},
\qquad
R_{eq}^{\max}=\frac{105\cdot420}{525}=\SI{84}{\ohm}.
\]
Ou seja, $80\pm4$, isto é, $\pm5\%$ --- exatamente a tolerância das parcelas.\\
\textbf{Resultado notável:} o paralelo \emph{também} preserva a tolerância
relativa, assim como a série. Isso não é coincidência, e a razão é demonstrada
no bloco de desafio.}
\end{questao}

\begin{questao}[3]{}
$R_1=\SI{100}{\ohm}$ com $\alpha=\SI{0.004}{\per\celsius}$ está em paralelo com
$R_2=\SI{100}{\ohm}$ de coeficiente desprezível. Determine $R_{eq}$ a
$\SI{20}{\celsius}$ e a $\SI{80}{\celsius}$, e calcule a variação percentual.
\resposta{$\SI{50}{\ohm}$ e $\SI{55.36}{\ohm}$; variação de $+10{,}7\%$}
\resolucao{A $\SI{20}{\celsius}$: $R_{eq}=100/2=\SI{50}{\ohm}$.\\
A $\SI{80}{\celsius}$: $R_1=100[1+0{,}004(60)]=\SI{124}{\ohm}$ e
\[
R_{eq}=\frac{124\cdot100}{224}=\SI{55.36}{\ohm}.
\]
Variação: $(55{,}36-50)/50=+10{,}7\%$.\\
\textbf{Comparação instrutiva:} $R_1$ variou $24\%$, mas $R_{eq}$ variou
$10{,}7\%$ --- o ramo estável "amortece" a deriva. Curiosamente, o valor é o
mesmo obtido para a \emph{corrente} da associação série equivalente: em ambos
os casos o elemento estável divide o efeito aproximadamente pela metade, porque
as duas resistências são comparáveis.}
\end{questao}

\begin{questao}[3]{}
Um resistor de $\SI{1000}{\ohm}$ está em paralelo com um de $\SI{10}{\ohm}$.
\begin{itensq}
\item Calcule $R_{eq}$.
\item Recalcule supondo que $R_1$ varie $+10\%$.
\item Interprete o resultado em termos de sensibilidade.
\end{itensq}
\resposta{(a) $\SI{9.901}{\ohm}$; (b) $\SI{9.910}{\ohm}$; (c) variação de
apenas $0{,}09\%$}
\resolucao{(a) $R_{eq}=\dfrac{1000\cdot10}{1010}=\SI{9.901}{\ohm}$.\\
(b) Com $R_1=\SI{1100}{\ohm}$:
$R_{eq}=\dfrac{1100\cdot10}{1110}=\SI{9.910}{\ohm}$.\\
(c) Uma variação de $10\%$ em $R_1$ produziu $0{,}09\%$ em $R_{eq}$ ---
atenuação de mais de cem vezes. O ramo dominante ($\SI{10}{\ohm}$) fixa
praticamente sozinho o resultado.\\
\textbf{Uso prático:} é assim que se faz um ajuste \emph{fino}. Colocando um
resistor grande em paralelo com um pequeno, obtém-se um deslocamento pequeno e
controlado --- técnica usada em calibração de shunts e de divisores de
precisão. A quantificação exata dessa sensibilidade aparece no bloco de
desafio.}
\end{questao}

\begin{questao}[3]{}
As cargas da questão residencial ($\SI{27.5}{\ampere}$ totais) são alimentadas
por um circuito cuja fiação tem $\SI{0.2}{\ohm}$ de ida e volta, a partir de um
quadro de $\SI{220}{\volt}$.
\begin{itensq}
\item Determine a queda de tensão na fiação e a tensão nas cargas.
\item Calcule a perda de potência na fiação.
\item Avalie se a queda atende ao limite usual de $4\%$.
\end{itensq}
\resposta{(a) $\SI{5.5}{\volt}$; $\SI{214.5}{\volt}$;
(b) $\SI{151.25}{\watt}$; (c) $2{,}5\%$ --- atende}
\resolucao{(a) $\Delta V=R_fI=0{,}2\cdot27{,}5=\SI{5.5}{\volt}$, logo as cargas
recebem $220-5{,}5=\SI{214.5}{\volt}$.\\
(b) $P_f=R_fI^{2}=0{,}2(27{,}5)^{2}=\SI{151.25}{\watt}$ --- dissipados em puro
aquecimento do cabo.\\
(c) $5{,}5/220=2{,}5\%$, abaixo do limite típico de $4\%$ para circuitos
terminais. \checkmark\\
\textbf{Observação sutil:} a fiação está em \emph{série} com o conjunto, mesmo
que as cargas estejam em paralelo entre si. É por isso que ligar mais uma carga
--- que não afeta as demais em uma rede ideal --- na prática \emph{reduz} a
tensão de todas: a corrente total sobe e a queda na fiação também. O efeito é
visível quando as luzes piscam ao ligar um chuveiro.}
\end{questao}

\begin{questao}[3]{}
Deseja-se obter $\SI{75}{\ohm}$ dispondo de um resistor de $\SI{100}{\ohm}$.
Determine o resistor a ser associado em paralelo e discuta a viabilidade.
\resposta{$R=\SI{300}{\ohm}$}
\resolucao{
\[
\frac{1}{75}=\frac{1}{100}+\frac{1}{R}
\;\Longrightarrow\;
\frac{1}{R}=\frac{1}{75}-\frac{1}{100}=\frac{4-3}{300}
\;\Longrightarrow\;
R=\SI{300}{\ohm}.
\]
\textbf{Viabilidade:} $\SI{300}{\ohm}$ existe na série E24, e o alvo
($\SI{75}{\ohm}$) está a $25\%$ do resistor de partida --- folga confortável.
A precisão do resultado depende principalmente do resistor de
$\SI{100}{\ohm}$, que domina o paralelo.\\
\textbf{Limite do método:} só se pode \emph{reduzir} o valor de partida; alvos
acima de $\SI{100}{\ohm}$ exigem série. E, como se verá no desafio, alvos muito
próximos de $\SI{100}{\ohm}$ tornam a solução impraticável.}
\end{questao}

\begin{questao}[3]{}
Compare a mesma dupla $R_1=\SI{30}{\ohm}$ e $R_2=\SI{60}{\ohm}$ ligada em série
e em paralelo, sob $\SI{90}{\volt}$: resistência, corrente total, distribuição
de corrente e de potência.
\resposta{Série: $\SI{90}{\ohm}$, $\SI{1}{\ampere}$, $P=30$ e
$\SI{60}{\watt}$. Paralelo: $\SI{20}{\ohm}$, $\SI{4.5}{\ampere}$,
$P=270$ e $\SI{135}{\watt}$}
\resolucao{\textbf{Série:} $R=\SI{90}{\ohm}$, $I=\SI{1}{\ampere}$ (comum),
$P_1=30$, $P_2=\SI{60}{\watt}$; total $\SI{90}{\watt}$. O \emph{maior}
resistor dissipa mais.\\
\textbf{Paralelo:} $R=\SI{20}{\ohm}$, $I_1=3$, $I_2=\SI{1.5}{\ampere}$,
$I_t=\SI{4.5}{\ampere}$, $P_1=\SI{270}{\watt}$, $P_2=\SI{135}{\watt}$; total
$\SI{405}{\watt}$. O \emph{menor} resistor dissipa mais.\\
\textbf{Razões:} $R_{s}/R_{p}=90/20=4{,}5$, e a potência total é $4{,}5$ vezes
maior no paralelo --- o mesmo fator, pois $P=V^{2}/R$ com $V$ fixo.\\
\textbf{Síntese:} de dois resistores obtêm-se quatro valores
($30$, $60$, $20$ e $\SI{90}{\ohm}$) e duas distribuições opostas de potência.
A ligação é uma variável de projeto tão importante quanto os componentes.}
\end{questao}

\blocodesafio

\begin{questao}[4]{}
\textbf{(Projeto de instrumentação.)} Um galvanômetro tem $R_g=\SI{99}{\ohm}$ e
fundo de escala $I_g=\SI{1}{\milli\ampere}$. Deseja-se um amperímetro de três
escalas: $\SI{10}{\milli\ampere}$, $\SI{100}{\milli\ampere}$ e
$\SI{1}{\ampere}$.
\begin{itensq}
\item Determine os três shunts na configuração de shunts comutados.
\item Aponte o risco dessa configuração e descreva a alternativa (shunt de
      Ayrton).
\item Calcule a resistência do instrumento em cada escala.
\end{itensq}
\resposta{(a) $\SI{11}{\ohm}$, $\SI{1}{\ohm}$ e $\SI{0.0991}{\ohm}$;
(c) $\SI{9.9}{\ohm}$, $\SI{0.99}{\ohm}$ e $\SI{0.099}{\ohm}$}
\resolucao{(a) Com $n=I_{\max}/I_g$, vale
$R_{sh}=\dfrac{R_g}{n-1}$:
\[
n=10\Rightarrow\frac{99}{9}=\SI{11}{\ohm};\quad
n=100\Rightarrow\frac{99}{99}=\SI{1}{\ohm};\quad
n=1000\Rightarrow\frac{99}{999}=\SI{0.0991}{\ohm}.
\]
(b) \textbf{Risco:} durante a comutação, a chave passa por um instante em que
\emph{nenhum} shunt está conectado. Toda a corrente de linha --- até
$\SI{1}{\ampere}$ --- atravessaria o galvanômetro, projetado para
$\SI{1}{\milli\ampere}$: destruição imediata.\\
\textbf{Shunt de Ayrton:} os resistores ficam permanentemente ligados em série
entre si e o conjunto em paralelo com o galvanômetro; a chave apenas escolhe
\emph{em que ponto} da cadeia a corrente entra. Assim o galvanômetro nunca
fica desprotegido, mesmo durante a comutação.\\
(c) $R_A=R_g\parallel R_{sh}$: $99\parallel11=\SI{9.9}{\ohm}$;
$99\parallel1=\SI{0.99}{\ohm}$; $99\parallel0{,}0991=\SI{0.099}{\ohm}$.\\
Note que $R_A=R_g/n$ exatamente --- escalas maiores exigem instrumentos de
resistência menor, o que é justamente o desejável.}
\end{questao}

\begin{questao}[4]{}
\textbf{(Demonstração --- escalonamento de bancos.)} Considere $N$ resistores
iguais de valor $R$ e potência nominal $P_1$.
\begin{itensq}
\item Mostre que a associação paralela vale $R/N$ e suporta $NP_1$.
\item Mostre que a série vale $NR$ e também suporta $NP_1$.
\item Deduza a matriz $m\times n$ necessária para manter o valor $R$ e obter
      potência $mnP_1$.
\end{itensq}
\resposta{(c) $m$ colunas em paralelo de $n$ unidades em série, com $m=n$}
\resolucao{(a) $G_{eq}=NG\Rightarrow R_{eq}=R/N$. Sob tensão $V$ comum, cada
unidade dissipa $V^{2}/R$; como todas são iguais e submetidas à mesma tensão, o
total é $N$ vezes o individual: $P_{\max}=NP_1$.\\
(b) $R_{eq}=NR$. Com corrente $I$ comum, cada unidade dissipa $RI^{2}$ e o
total é novamente $NP_1$.\\
\textbf{Fato central:} \emph{qualquer} associação de $N$ resistores iguais
suporta $NP_1$, desde que a dissipação se reparta igualmente. O que a topologia
escolhe é o \emph{valor}, não a capacidade térmica.\\
(c) Uma matriz de $n$ unidades em série, repetida em $m$ colunas paralelas:
\[
R_{eq}=\frac{nR}{m},\qquad P_{\max}=mnP_1.
\]
Para preservar o valor, $n=m$; então $R_{eq}=R$ e
$P_{\max}=n^{2}P_1$.\\
\textbf{Exemplo:} para quadruplicar a potência sem alterar o valor, use
$2\times2=4$ unidades. Para $\times9$, use $3\times3=9$. A potência cresce com
o \emph{quadrado} do lado da matriz --- e o custo, também.}
\end{questao}

\begin{questao}[4]{}
\textbf{(Sensibilidade.)} Para $R_{eq}=\dfrac{R_1R_2}{R_1+R_2}$:
\begin{itensq}
\item Obtenha $\dfrac{\partial R_{eq}}{\partial R_1}$ e mostre que vale
      $\left(\dfrac{R_2}{R_1+R_2}\right)^{2}$.
\item Obtenha a sensibilidade relativa e mostre que
      $S_1+S_2=1$.
\item Avalie para $R_1=\SI{100}{\ohm}$ e $R_2=\SI{400}{\ohm}$, e interprete.
\end{itensq}
\resposta{(a) $0{,}64$; (b) $S_1=R_2/(R_1+R_2)$, $S_2=R_1/(R_1+R_2)$;
(c) $S_1=0{,}8$ e $S_2=0{,}2$}
\resolucao{(a) Derivando o quociente em relação a $R_1$:
\[
\frac{\partial R_{eq}}{\partial R_1}
=\frac{R_2(R_1+R_2)-R_1R_2}{(R_1+R_2)^{2}}
=\frac{R_2^{2}}{(R_1+R_2)^{2}}
=\left(\frac{R_2}{R_1+R_2}\right)^{2}.
\]
Com $R_1=100$ e $R_2=400$: $(400/500)^{2}=0{,}64$ --- adimensional, pois é
$\si{\ohm}$ por $\si{\ohm}$.\\
(b)
\[
S_1=\frac{\partial R_{eq}}{\partial R_1}\cdot\frac{R_1}{R_{eq}}
=\frac{R_2^{2}}{(R_1+R_2)^{2}}\cdot\frac{R_1(R_1+R_2)}{R_1R_2}
=\frac{R_2}{R_1+R_2}.
\]
Por simetria, $S_2=\dfrac{R_1}{R_1+R_2}$, e a soma é $1$ --- consequência de
$R_{eq}$ ser \textbf{homogênea de grau 1}.\\
(c) $S_1=0{,}8$ e $S_2=0{,}2$: um erro de $1\%$ em $R_1$ produz $0{,}8\%$ em
$R_{eq}$, e o mesmo erro em $R_2$ produz apenas $0{,}2\%$.\\
\textbf{Regra de projeto:} o resistor \emph{menor} domina o paralelo e deve
receber a tolerância apertada --- exatamente o inverso da série, onde manda o
maior. E se $R_2\gg R_1$, então $S_1\to1$ e $S_2\to0$: o ramo grande vira
ajuste fino, o que confirma quantitativamente o observado na questão do
$\SI{1000}{\ohm}$ com $\SI{10}{\ohm}$.}
\end{questao}

\begin{questao}[4]{}
\textbf{(Divisão de corrente sob tolerância.)} Três shunts iguais de
$\SI{0.01}{\ohm}$ devem dividir igualmente $\SI{30}{\ampere}$.
\begin{itensq}
\item Determine a corrente ideal por shunt e a tensão sobre eles.
\item Se um deles for $2\%$ menor e os outros dois nominais, calcule o
      desequilíbrio de corrente.
\item Discuta a implicação térmica.
\end{itensq}
\resposta{(a) $\SI{10}{\ampere}$ e $\SI{0.1}{\volt}$;
(b) $\SI{10.13}{\ampere}$ contra $\SI{9.93}{\ampere}$ --- desvio de $1{,}3\%$;
(c) o mais "fraco" aquece mais}
\resolucao{(a) $I=30/3=\SI{10}{\ampere}$;
$V=0{,}01\cdot10=\SI{0.1}{\volt}$.\\
(b) Com $R_a=\SI{0.0098}{\ohm}$ e $R_b=R_c=\SI{0.0100}{\ohm}$:
\[
G=\frac{1}{0{,}0098}+\frac{2}{0{,}0100}=102{,}04+200{,}00=\SI{302.04}{\siemens},
\]
\[
I_a=30\cdot\frac{102{,}04}{302{,}04}=\SI{10.13}{\ampere},
\qquad
I_b=I_c=30\cdot\frac{100{,}00}{302{,}04}=\SI{9.93}{\ampere}.
\]
\textbf{Verificação:} $10{,}13+9{,}93+9{,}93=\SI{29.99}{\ampere}$ \checkmark.\\
\textbf{Fato tranquilizador:} um desvio de $2\%$ na resistência produziu apenas
$1{,}3\%$ de desequilíbrio --- a divisão é \emph{mais estável} que os
componentes, porque os outros dois ramos "puxam" a média.\\
(c) $P_a=0{,}0098(10{,}13)^{2}=\SI{1.006}{\watt}$ contra
$P_b=0{,}0100(9{,}93)^{2}=\SI{0.986}{\watt}$: o shunt de menor resistência
dissipa $2\%$ a mais. Se tiver coeficiente térmico positivo, aquecerá,
aumentará sua resistência e \emph{devolverá} corrente aos vizinhos --- uma
realimentação \textbf{estabilizadora}. Com coeficiente negativo ocorreria o
oposto: fuga térmica. É por isso que resistores para divisão de corrente são
feitos em ligas de $\alpha$ positivo e pequeno, como a manganina.}
\end{questao}

\begin{questao}[4]{}
\textbf{(Compensação térmica em paralelo.)} Deseja-se uma associação paralela
com resistência independente da temperatura, usando
$R_1=\SI{100}{\ohm}$ com $\alpha_1=+4000\,\mathrm{ppm/^{\circ}C}$ e um segundo
resistor com $\alpha_2=-1000\,\mathrm{ppm/^{\circ}C}$.
\begin{itensq}
\item Deduza a condição de compensação em termos de condutâncias.
\item Determine $R_2$ e $R_{eq}$.
\item Verifique a compensação e comente sua exatidão.
\end{itensq}
\resposta{(a) $G_1\alpha_1+G_2\alpha_2=0$; (b) $R_2=\SI{25}{\ohm}$,
$R_{eq}=\SI{20}{\ohm}$; (c) exata só em primeira ordem}
\resolucao{(a) Partindo de $G_{eq}=G_1+G_2$ e de
$R_k(T)=R_k(1+\alpha_k\Delta T)$, tem-se
$G_k\approx G_k(1-\alpha_k\Delta T)$ para $\alpha_k\Delta T\ll1$. Então
\[
G_{eq}(T)\approx(G_1+G_2)-(G_1\alpha_1+G_2\alpha_2)\Delta T,
\]
e a condição de estabilidade é
\[
G_1\alpha_1+G_2\alpha_2=0
\quad\Longleftrightarrow\quad
\frac{\alpha_1}{R_1}+\frac{\alpha_2}{R_2}=0.
\]
(b) $R_2=-\dfrac{\alpha_2R_1}{\alpha_1}=\dfrac{1000\cdot100}{4000}
=\SI{25}{\ohm}$, e $R_{eq}=100\parallel25=\SI{20}{\ohm}$.\\
(c) \emph{Para $\Delta T=\SI{1}{\celsius}$:}
$R_1=\SI{100.4}{\ohm}$, $R_2=\SI{24.975}{\ohm}$, $R_{eq}=\SI{20.000}{\ohm}$ ---
compensação praticamente perfeita.\\
\emph{Para $\Delta T=\SI{50}{\celsius}$:} $R_1=\SI{120}{\ohm}$,
$R_2=\SI{23.75}{\ohm}$, $R_{eq}=\SI{19.83}{\ohm}$ --- desvio de $0{,}9\%$.\\
A compensação é de \textbf{primeira ordem}: excelente perto da temperatura de
projeto, degradando-se quadraticamente conforme se afasta.\\
\textbf{Compare com a série}, cuja condição era $R_1\alpha_1+R_2\alpha_2=0$
(ponderada por resistências). Aqui a ponderação é por \emph{condutâncias} ---
mais uma manifestação da dualidade série/paralelo.}
\end{questao}

\begin{questao}[4]{}
\textbf{(Falha de ramo com fonte real.)} Uma fonte de $\SI{24}{\volt}$ com
$r=\SI{2}{\ohm}$ alimenta três ramos de $\SI{6}{\ohm}$ em paralelo.
\begin{itensq}
\item Determine a tensão do barramento, a corrente total e a de cada ramo.
\item Um ramo se abre. Recalcule tudo.
\item Explique por que os ramos remanescentes \emph{mudam}, ao contrário do
      caso com fonte ideal.
\end{itensq}
\resposta{(a) $\SI{12}{\volt}$, $\SI{6}{\ampere}$, $\SI{2}{\ampere}$ por ramo;
(b) $\SI{14.4}{\volt}$, $\SI{4.8}{\ampere}$, $\SI{2.4}{\ampere}$ por ramo}
\resolucao{(a) $R_{eq}=6/3=\SI{2}{\ohm}$;
\[
I_t=\frac{24}{2+2}=\SI{6}{\ampere},\qquad
V_{barra}=24-2(6)=\SI{12}{\volt},\qquad
I_{ramo}=\frac{12}{6}=\SI{2}{\ampere}.
\]
(b) Com dois ramos, $R_{eq}=\SI{3}{\ohm}$:
\[
I_t=\frac{24}{2+3}=\SI{4.8}{\ampere},\qquad
V_{barra}=24-2(4{,}8)=\SI{14.4}{\volt},\qquad
I_{ramo}=\frac{14{,}4}{6}=\SI{2.4}{\ampere}.
\]
(c) Os ramos remanescentes \emph{não mudaram de valor}, mas a tensão que os
alimenta subiu de $12$ para $\SI{14.4}{\volt}$ ($+20\%$), e com ela a corrente
de cada um ($+20\%$).\\
\textbf{Mecanismo:} $r$ está em série com o conjunto. Menos ramos $\Rightarrow$
menor corrente total $\Rightarrow$ menor queda interna $\Rightarrow$ maior
tensão disponível.\\
\textbf{Implicação de projeto:} o isolamento entre ramos paralelos é uma
\emph{idealização}. Com fonte real, ramos interagem via $r$ --- e a perda de
uma carga pode sobrecarregar as demais em $20\%$ ou mais. Fontes reguladas
existem justamente para levar $r$ efetiva a valores próximos de zero e
restaurar o desacoplamento.}
\end{questao}

\begin{questao}[4]{}
\textbf{(Limite do método de ajuste.)} Dispõe-se de um resistor fixo de
$\SI{100}{\ohm}$ e deseja-se obter valores-alvo por associação paralela.
\begin{itensq}
\item Deduza $R(alvo)$ e calcule para alvos de $\SI{90}{\ohm}$,
      $\SI{95}{\ohm}$ e $\SI{99}{\ohm}$.
\item Obtenha a sensibilidade do alvo em relação a $R$ e avalie no caso de
      $\SI{99}{\ohm}$.
\item Conclua sobre a viabilidade prática.
\end{itensq}
\resposta{(a) $\SI{900}{\ohm}$, $\SI{1900}{\ohm}$ e $\SI{9900}{\ohm}$;
(c) inviável para alvos próximos do resistor fixo}
\resolucao{(a) De $\dfrac{1}{R_{alvo}}=\dfrac{1}{R_f}+\dfrac{1}{R}$:
\[
R=\frac{R_fR_{alvo}}{R_f-R_{alvo}}.
\]
\[
90\to\frac{100\cdot90}{10}=\SI{900}{\ohm};\quad
95\to\frac{100\cdot95}{5}=\SI{1900}{\ohm};\quad
99\to\frac{100\cdot99}{1}=\SI{9900}{\ohm}.
\]
O denominador $R_f-R_{alvo}$ tende a zero e $R$ \textbf{diverge}.\\
(b) Pela sensibilidade relativa do paralelo,
$S=\dfrac{R}{R_f+R}$. Para $R=\SI{9900}{\ohm}$:
$S=9900/10000=0{,}99$.\\
Ou seja: $1\%$ de erro em $R$ produz $0{,}99\%$ de erro no alvo. Como o ajuste
pretendido é de apenas $1\%$ (de $100$ para $\SI{99}{\ohm}$), \emph{o erro tem
a mesma ordem de grandeza do ajuste} --- ele o anula.\\
(c) \textbf{Inviável.} Com resistor de $5\%$ de tolerância, o "ajuste" de
$\SI{1}{\ohm}$ vem com incerteza de $\SI{0.99}{\ohm}$: o alvo pode ficar em
qualquer ponto entre $98$ e $\SI{100}{\ohm}$. Além disso,
$\SI{9900}{\ohm}$ não é valor comercial, e a resistência de contato já é
comparável ao efeito buscado.\\
\textbf{Regra prática:} ajuste por paralelo só é confiável para correções de
\emph{dezenas} de por cento. Para correções finas, use potenciômetro de ajuste
ou selecione o componente por medição direta.}
\end{questao}

\begin{questao}[4]{}
\textbf{(Demonstração --- propagação de incerteza.)} Mostre que, se
\emph{todos} os resistores de uma rede puramente resistiva tiverem a mesma
tolerância relativa $t$, então $R_{eq}$ terá exatamente a tolerância $t$,
qualquer que seja a topologia.
\begin{itensq}
\item Demonstre a propriedade de homogeneidade.
\item Verifique em um exemplo série e em um paralelo.
\item Discuta o que muda quando as tolerâncias são diferentes.
\end{itensq}
\resposta{$R_{eq}$ é homogênea de grau 1: $R_{eq}(\lambda R_k)=\lambda
R_{eq}(R_k)$; logo a tolerância relativa se conserva}
\resolucao{(a) Qualquer $R_{eq}$ de rede resistiva é construída por somas
($R_i+R_j$) e por combinações do tipo $\dfrac{R_iR_j}{R_i+R_j}$ --- ambas
\textbf{homogêneas de grau 1}. Multiplicando \emph{todos} os resistores por
$\lambda$:
\[
\lambda R_i+\lambda R_j=\lambda(R_i+R_j),
\qquad
\frac{(\lambda R_i)(\lambda R_j)}{\lambda R_i+\lambda R_j}
=\lambda\frac{R_iR_j}{R_i+R_j}.
\]
Por indução sobre a construção da rede, $R_{eq}(\lambda R_k)=\lambda R_{eq}$.
Fazendo $\lambda=1\pm t$, obtém-se $R_{eq}(1\pm t)$: tolerância relativa
exatamente $t$.\\
(b) \emph{Série:} $100\pm5\%$ e $200\pm5\%$ dão $300\pm\SI{15}{\ohm}=\pm5\%$
\checkmark. \emph{Paralelo:} $100\pm5\%$ e $400\pm5\%$ dão
$80\pm\SI{4}{\ohm}=\pm5\%$ \checkmark. A demonstração explica por que os dois
casos, aparentemente distintos, deram o mesmo resultado.\\
(c) Com tolerâncias \emph{diferentes}, a homogeneidade não se aplica --- os
resistores não escalam pelo mesmo $\lambda$. Vale então a combinação ponderada
pelas sensibilidades:
\[
t_{eq}=\sum_k S_k\,t_k,
\qquad \text{com } \sum_k S_k=1.
\]
Em série, $S_k=R_k/\sum R_j$ (manda o maior); em paralelo,
$S_k=G_k/\sum G_j$ (manda o menor). A restrição $\sum S_k=1$ é justamente a
identidade de Euler para funções homogêneas de grau 1 --- e garante que
$t_{eq}$ jamais ultrapasse a maior das tolerâncias individuais.}
\end{questao}

\begin{questao}[4]{}
\textbf{(Seletividade da proteção.)} Um barramento de $\SI{220}{\volt}$
alimenta três circuitos em paralelo, com correntes nominais de
$\SI{20}{\ampere}$, $\SI{5}{\ampere}$ e $\SI{2.5}{\ampere}$, protegidos
individualmente e por um disjuntor geral.
\begin{itensq}
\item Especifique os disjuntores de ramo e o geral.
\item Explique o requisito de seletividade e o que ocorre se ele falhar.
\item Discuta por que a soma dos disjuntores de ramo pode exceder o geral.
\end{itensq}
\resposta{(a) ramos de $25$, $10$ e $\SI{10}{\ampere}$; geral de
$\SI{32}{\ampere}$; (b) o ramo deve atuar antes do geral}
\resolucao{(a) Cada disjuntor de ramo fica acima da corrente nominal do próprio
circuito: $\SI{25}{\ampere}$, $\SI{10}{\ampere}$ e $\SI{10}{\ampere}$
(valores comerciais). O geral cobre a soma das cargas,
$\SI{27.5}{\ampere}$: disjuntor de $\SI{32}{\ampere}$.\\
(b) \textbf{Seletividade} significa que, diante de uma falta em um ramo,
\emph{apenas} o disjuntor daquele ramo abre --- os demais circuitos continuam
alimentados. Isso exige que a curva tempo--corrente do dispositivo de ramo
esteja inteiramente \emph{abaixo} da do geral.\\
Se a seletividade falhar, um curto em uma tomada derruba a instalação inteira:
além do transtorno, perde-se a informação de \emph{onde} está o defeito, e a
localização passa a exigir religamentos sucessivos.\\
(c) Porque as cargas dificilmente operam \emph{todas} no máximo ao mesmo tempo
--- é o \textbf{fator de demanda}. Aqui $25+10+10=\SI{45}{\ampere}$ contra um
geral de $\SI{32}{\ampere}$: perfeitamente admissível, desde que a demanda
simultânea real permaneça abaixo de $\SI{32}{\ampere}$.\\
\textbf{Consequência do paralelo:} como as correntes de ramo se somam no
alimentador, é o \emph{comportamento estatístico} do conjunto --- e não a soma
aritmética dos máximos --- que dimensiona a entrada. Superdimensionar o geral
para $\SI{45}{\ampere}$ exigiria condutores muito mais caros sem ganho real de
disponibilidade.}
\end{questao}

% END BANCO COMPLEMENTAR

\gabarito[3.2 Circuitos em paralelo]

% =====================================================================
\subsection{Circuitos mistos}
% =====================================================================

\blocofixacao

\begin{questao}[1]{}
A estratégia geral para reduzir um circuito misto é:
\begin{alternativas}
\item somar todas as resistências, independentemente da ligação.
\item resolver primeiro as associações série mais internas, depois os paralelos.
\item identificar os trechos em série e em paralelo e reduzi-los, dos mais internos
      para os mais externos, um de cada vez.
\item aplicar sempre a transformação estrela-triângulo.
\item calcular a corrente antes de reduzir o circuito.
\end{alternativas}
\resposta{(c)}
\resolucao{Reduz-se o circuito por etapas, sempre pelo trecho onde a ligação
(série ou paralelo) é inequívoca --- em geral o mais interno ---, substituindo-o
por um resistor equivalente e redesenhando. Repete-se até restar um único
resistor. A transformação estrela-triângulo (alternativa d) só é necessária
quando \emph{nenhum} trecho é puramente série ou paralelo, como nas pontes.}
\end{questao}

\begin{questao}[1]{}
Para cada circuito, calcule a resistência equivalente vista pela fonte e a
corrente total.

\circuitoq{%
\subcirc{a}{%
\begin{circuitikz}[scale=0.8,transform shape,line width=0.8pt]
 \draw (0,0) to[battery1,l_=$\SI{24}{\volt}$] (0,3)
       to[R,l=$\SI{6}{\ohm}$] (2.6,3) coordinate(n);
 \draw (n) to[R,l=$\SI{6}{\ohm}$] (2.6,0) -- (0,0);
 \draw (n) -- (4.6,3) to[R,l=$\SI{3}{\ohm}$] (4.6,0) -- (2.6,0);
 \draw[->,Icol,line width=1pt] (0.6,3.4)--(1.5,3.4) node[midway,above,font=\footnotesize]{$I$};
\end{circuitikz}}
\hfill
\subcirc{b}{%
\begin{circuitikz}[scale=0.8,transform shape,line width=0.8pt]
 \draw (0,0) to[battery1,l_=$\SI{20}{\volt}$] (0,3) coordinate(t);
 \draw (t) -- (2,3) to[R,l=$\SI{4}{\ohm}$] (2,0) coordinate(b);
 \draw (2,3) -- (4,3) to[R,l=$\SI{4}{\ohm}$] (4,0) -- (2,0);
 \draw (4,3) to[R,l=$\SI{8}{\ohm}$] (6,3) -- (6,0) -- (b) -- (0,0);
 \draw[->,Icol,line width=1pt] (0.5,3.4)--(1.4,3.4) node[midway,above,font=\footnotesize]{$I$};
\end{circuitikz}}}

\resposta{(a) \SI{8}{\ohm}, \SI{3}{\ampere}; (b) \SI{10}{\ohm}, \SI{2}{\ampere}}
\resolucao{\textbf{(a)} Os resistores de \SI{6}{\ohm} e \SI{3}{\ohm} da direita
estão em paralelo: $6\parallel 3 = \SI{2}{\ohm}$. Esse resultado fica em série com
o \SI{6}{\ohm} de entrada: $\Req = 6+2 = \SI{8}{\ohm}$;
$I = 24/8 = \SI{3}{\ampere}$.\\
\textbf{(b)} Os dois \SI{4}{\ohm} em paralelo dão \SI{2}{\ohm}, em série com o
\SI{8}{\ohm}: $\Req = 2+8 = \SI{10}{\ohm}$; $I = 20/10 = \SI{2}{\ampere}$.}
\end{questao}

\blocoaplicacao

\begin{questao}[2]{}
No circuito da figura, determine a resistência equivalente vista pela fonte, a
corrente total e a corrente em cada resistor de \SI{20}{\ohm}.

\circuitoq{%
\begin{circuitikz}[scale=0.9,transform shape,line width=0.8pt]
 \draw (0,0) to[battery1,l_={$E=\SI{40}{\volt}$}] (0,3)
       to[R,l=$\SI{10}{\ohm}$] (3,3) coordinate(n);
 \draw (n) to[R,l=$\SI{20}{\ohm}$] (3,0);
 \draw (n) -- (5.5,3) to[R,l=$\SI{20}{\ohm}$] (5.5,0) -- (3,0) -- (0,0);
 \draw[->,Icol,line width=1pt] (0.6,3.4)--(1.6,3.4) node[midway,above,font=\footnotesize]{$I_t$};
\end{circuitikz}}

\resposta{$\Req=\SI{20}{\ohm}$; $I_t=\SI{2}{\ampere}$; \SI{1}{\ampere} em cada \SI{20}{\ohm}}
\resolucao{Os dois \SI{20}{\ohm} em paralelo equivalem a $20/2 = \SI{10}{\ohm}$,
em série com o \SI{10}{\ohm} de entrada:
$\Req = 10+10 = \SI{20}{\ohm}$; $I_t = 40/20 = \SI{2}{\ampere}$.
Como os ramos são iguais, a corrente total se divide igualmente:
\SI{1}{\ampere} em cada resistor de \SI{20}{\ohm}.\\
\textbf{Confira as tensões:} sobre o \SI{10}{\ohm} caem $10\cdot 2 = \SI{20}{\volt}$;
sobram \SI{20}{\volt} no bloco paralelo, e $20/20 = \SI{1}{\ampere}$ em cada ramo.}
\end{questao}

\begin{questao}[2]{}
Calcule a resistência equivalente entre A e B.

\circuitoq{%
\begin{circuitikz}[scale=0.9,transform shape,line width=0.8pt]
 \draw (0,0) node[left,font=\bfseries]{A} to[short,o-*] (1.2,0) coordinate(a);
 % bloco 1: R=2 em serie
 \draw (a) to[R,l=$\SI{2}{\ohm}$] (3.2,0) coordinate(b);
 % bloco 2: 6 || 3
 \draw (b) to[R,l=$\SI{6}{\ohm}$] (5.4,0) coordinate(c);
 \draw (b) ++(0,1.4) coordinate(b2) to[R,l=$\SI{3}{\ohm}$] ($(c)+(0,1.4)$) coordinate(c2);
 \draw (b) -- (b2);  \draw (c) -- (c2);
 % bloco 3: 10 || 15
 \draw (c) to[R,l=$\SI{10}{\ohm}$] (7.6,0) coordinate(d);
 \draw (c) ++(0,1.4) coordinate(c3) to[R,l=$\SI{15}{\ohm}$] ($(d)+(0,1.4)$) coordinate(d2);
 \draw (c) -- (c3);  \draw (d) -- (d2);
 \draw (d) to[short,-o] (8.8,0) node[right,font=\bfseries]{B};
\end{circuitikz}}

\resposta{$R_{AB} = \SI{10}{\ohm}$}
\resolucao{Há dois blocos em paralelo, separados por um resistor em série:
\[
6\parallel 3=\frac{6\cdot3}{9}=\SI{2}{\ohm},
\qquad
10\parallel 15=\frac{10\cdot15}{25}=\SI{6}{\ohm}.
\]
Os três trechos estão em série (o \SI{2}{\ohm} de entrada, o primeiro bloco e o
segundo bloco):
\[
R_{AB}=2+2+6=\SI{10}{\ohm}.
\]}
\end{questao}

\blocoaprofundamento

\begin{questao}[3]{}
Determine a resistência equivalente entre A e B da rede em escada da figura.

\circuitoq{%
\begin{circuitikz}[scale=0.9,transform shape,line width=0.8pt]
 \draw (0,2) node[left,font=\bfseries]{A} to[short,o-*] (1,2) coordinate(a);
 \draw (a) to[R,l=$\SI{12}{\ohm}$] (1,0) coordinate(b);
 \draw (a) to[R,l=$\SI{4}{\ohm}$] (4,2) coordinate(c);
 \draw (c) to[R,l=$\SI{6}{\ohm}$] (4,0) coordinate(d);
 \draw (c) to[R,l=$\SI{3}{\ohm}$] (6.5,2) coordinate(e);
 \draw (e) -- (6.5,0) -- (d) -- (b);
 \draw (b) to[short,*-o] (0,0) node[left,font=\bfseries]{B};
 \draw (e) -- (7,2) node[right]{};
\end{circuitikz}}

\resposta{$R_{AB} = \SI{4}{\ohm}$}
\resolucao{Resolve-se do trecho mais distante para o mais próximo da entrada:
\begin{align*}
6 \parallel 3 &= \SI{2}{\ohm} && \text{(último ramo)}\\
4 + 2 &= \SI{6}{\ohm} && \text{(em série com o } \SI{4}{\ohm}\text{)}\\
12 \parallel 6 &= \SI{4}{\ohm} && \text{(em paralelo com o } \SI{12}{\ohm}\text{ de entrada)}.
\end{align*}
Logo $R_{AB} = \SI{4}{\ohm}$. O nome "escada" vem justamente desse padrão de
reduzir degrau a degrau.}
\end{questao}

\begin{questao}[3]{}
No circuito da figura, um fio ideal curto-circuita parte da rede. Determine a
resistência equivalente entre A e B.

\circuitoq{%
\begin{circuitikz}[scale=0.9,transform shape,line width=0.8pt]
 \draw (0,1.5) node[left,font=\bfseries]{A} to[short,o-*] (1.5,1.5) coordinate(a);
 \draw (a) to[R,l=$\SI{10}{\ohm}$] (1.5,3.2) -- (5,3.2) coordinate(top);
 \draw (a) to[R,l=$\SI{20}{\ohm}$] (5,1.5) coordinate(m);
 \draw (top) -- (m);
 \draw[Vcol,line width=1.1pt] (a) -- (1.5,0) -- (5,0) -- (m);
 \node[Vcol,font=\footnotesize\sffamily,below] at (3.25,0) {fio ideal};
 \draw (m) to[R,l=$\SI{10}{\ohm}$] (7,1.5) to[short,-o] (7.8,1.5)
       node[right,font=\bfseries]{B};
\end{circuitikz}}

\resposta{$R_{AB} = \SI{10}{\ohm}$}
\resolucao{O fio ideal liga diretamente o nó A ao nó central: eles passam a ser
\emph{o mesmo nó}. Com isso, os resistores de \SI{10}{\ohm} e \SI{20}{\ohm} do topo
ficam ambos ligados entre A e o nó central --- que agora coincidem ---, ou seja,
têm os dois terminais no mesmo ponto e ficam curto-circuitados. Sobra apenas o
resistor de \SI{10}{\ohm} até B:
\[
R_{AB} = \SI{10}{\ohm}.
\]
\textbf{Lição:} identificar nós em curto \emph{antes} de somar resistências evita
contas longas e erradas.}
\end{questao}

\begin{questao}[3]{}
Doze resistores iguais de $\SI{6}{\ohm}$ são soldados formando as arestas de um
\textbf{cubo}. Determine a resistência equivalente entre dois vértices ligados por
uma \emph{aresta} (A e B na figura).

\circuitoq{%
\begin{tikzpicture}[scale=1,line width=0.8pt]
 \def\a{2}
 % vertices do cubo em perspectiva
 \coordinate (A) at (0,0);
 \coordinate (B) at (\a,0);
 \coordinate (C) at (\a,\a);
 \coordinate (D) at (0,\a);
 \coordinate (E) at (0.7,0.6);
 \coordinate (F) at (\a+0.7,0.6);
 \coordinate (G) at (\a+0.7,\a+0.6);
 \coordinate (H) at (0.7,\a+0.6);
 \draw[Rcol] (A)--(B)--(C)--(D)--cycle;
 \draw[Rcol] (E)--(F)--(G)--(H)--cycle;
 \draw[Rcol] (A)--(E); \draw[Rcol] (B)--(F);
 \draw[Rcol] (C)--(G); \draw[Rcol] (D)--(H);
 \fill[Vcol] (A) circle (0.09) node[below left,font=\bfseries]{A};
 \fill[Vcol] (B) circle (0.09) node[below right,font=\bfseries]{B};
 \node[Rcol,font=\footnotesize] at (\a/2,-0.3){$\SI{6}{\ohm}$ (cada aresta)};
\end{tikzpicture}}{Cubo de resistores.}

\resposta{$R_{AB} = \SI{5}{\ohm}$}
\resolucao{Explorando a \emph{simetria} do cubo, injeta-se corrente $I$ em A e
retira-se em B. Por simetria, a corrente se distribui de forma previsível:
o resultado clássico para a resistência entre vértices de uma \emph{aresta} é
\[
R_{AB} = \frac{7}{12}\,R.
\]
Com $R = \SI{6}{\ohm}$: $R_{AB} = \dfrac{7}{12}\cdot6 = \SI{3.5}{\ohm}$.\\
(Observação: os três casos clássicos do cubo são $\tfrac{7}{12}R$ pela aresta,
$\tfrac{3}{4}R$ pela diagonal de face e $\tfrac{5}{6}R$ pela diagonal principal.
Com $R=6$: \SI{3.5}{}, \SI{4.5}{} e \SI{5}{\ohm}, respectivamente.) A técnica-chave
é usar a simetria para agrupar nós de mesmo potencial antes de calcular.}
\end{questao}

\begin{questao}[3]{}
Doze resistores iguais de $\SI{12}{\ohm}$ formam as arestas de um cubo.
Determine a resistência equivalente entre:
\begin{itensq}
\item dois vértices de uma mesma \emph{aresta};
\item dois vértices opostos de uma mesma \emph{face} (diagonal de face);
\item dois vértices \emph{diametralmente opostos} (diagonal principal).
\end{itensq}
\resposta{(a) \SI{7}{\ohm}; (b) \SI{9}{\ohm}; (c) \SI{10}{\ohm}}
\resolucao{A chave é explorar a \textbf{simetria} para identificar nós de mesmo
potencial, que podem ser unidos sem alterar o circuito.

\textbf{(c) Diagonal principal --- o caso mais simétrico.} Injetando corrente $I$
em A, ela se divide igualmente pelas 3 arestas que saem de A ($I/3$ em cada). Em
seguida, cada uma se divide em duas ($I/6$ nas 6 arestas intermediárias) e,
finalmente, recombina-se em 3 arestas ($I/3$) chegando em B. Somando as quedas ao
longo de um caminho:
\[
U = \frac{I}{3}R+\frac{I}{6}R+\frac{I}{3}R
= IR\left(\frac13+\frac16+\frac13\right)=\frac{5}{6}IR,
\]
logo $R_{AB}=\dfrac{5R}{6}=\dfrac{5\cdot12}{6}=\SI{10}{\ohm}$.

\textbf{(b) Diagonal de face:} por raciocínio análogo,
$R = \dfrac{3R}{4}=\dfrac{3\cdot12}{4}=\SI{9}{\ohm}$.

\textbf{(a) Aresta:} $R = \dfrac{7R}{12}=\dfrac{7\cdot12}{12}=\SI{7}{\ohm}$.

\medskip
\begin{center}\small\sffamily
\begin{tabular}{@{}l c c@{}}
\toprule
\textbf{Par de vértices} & \textbf{Fórmula} & \textbf{Valor ($R=\SI{12}{\ohm}$)}\\
\midrule
Aresta            & $\tfrac{7}{12}R$ & \SI{7}{\ohm}\\
Diagonal de face  & $\tfrac{3}{4}R$  & \SI{9}{\ohm}\\
Diagonal principal& $\tfrac{5}{6}R$  & \SI{10}{\ohm}\\
\bottomrule
\end{tabular}
\end{center}
\textbf{Note a ordenação:} quanto mais "distantes" os vértices, maior a
resistência --- mas a diferença é modesta (7 a 10 $\Omega$), porque há sempre
muitos caminhos paralelos. \emph{A simetria é a ferramenta central: sem ela,
seria preciso resolver um sistema de 7 equações nodais.}}
\end{questao}

\begin{questao}[3]{}
No circuito da figura, todos os resistores valem $\SI{6}{\ohm}$, exceto o
resistor central da ponte, de $\SI{5}{\ohm}$. Determine $R_{AB}$ e justifique por
que o valor do resistor central é irrelevante.

\circuitoq{%
\begin{circuitikz}[scale=1,transform shape,line width=0.8pt]
 \draw (0,0) node[left,font=\bfseries]{A} coordinate (A);
 \draw (5,0) node[right,font=\bfseries]{B} coordinate (B);
 \coordinate (C) at (2.5,1.6); \coordinate (D) at (2.5,-1.6);
 \draw (A) to[R,l={$\SI{4}{\ohm}$}] (C);
 \draw (C) to[R,l={$\SI{8}{\ohm}$}] (B);
 \draw (A) to[R,l_={$\SI{4}{\ohm}$}] (D);
 \draw (D) to[R,l_={$\SI{8}{\ohm}$}] (B);
 \draw (C) to[R,l={$\SI{5}{\ohm}$},color=Vcol] (D);
\end{circuitikz}}

\resposta{$R_{AB}=\SI{6}{\ohm}$; a ponte está equilibrada}
\resolucao{\textbf{Teste de equilíbrio:} os braços opostos satisfazem
\[
\frac{R_{AC}}{R_{CB}}=\frac{4}{8}=\frac12,
\qquad
\frac{R_{AD}}{R_{DB}}=\frac{4}{8}=\frac12.
\]
As razões são \emph{iguais}, logo a ponte está \textbf{equilibrada}: os nós C e D
estão no \emph{mesmo potencial}. Sem diferença de potencial, não circula corrente
pelo resistor central --- ele pode ser \emph{removido} (ou substituído por
qualquer valor, inclusive um curto ou um circuito aberto) sem alterar nada.

Restam dois ramos em série, associados em paralelo:
\[
R_{AB}=(4+8)\parallel(4+8)=\frac{12}{2}=\SI{6}{\ohm}.
\]
\textbf{Consequência prática --- a ponte como instrumento:} é exatamente essa
insensibilidade que torna a ponte de Wheatstone tão precisa. No equilíbrio, o
detector central indica zero \emph{independentemente} de sua própria resistência
e da tensão da fonte. A medição depende apenas das \emph{razões} entre
resistores conhecidos.}
\end{questao}

\begin{questao}[3]{}
Dispõe-se de um número ilimitado de resistores de $\SI{10}{\ohm}$. Projete
associações que resultem em (a) $\SI{15}{\ohm}$; (b) $\SI{25}{\ohm}$;
(c) $\SI{4}{\ohm}$, usando o \textbf{menor número possível} de resistores em cada
caso.
\resposta{(a) 3 resistores; (b) 4 resistores; (c) 4 resistores}
\resolucao{A estratégia é decompor o valor-alvo em somas de múltiplos
($nR$: série) e submúltiplos ($R/n$: paralelo) de $R = \SI{10}{\ohm}$.

\medskip
\textbf{(a) \SI{15}{\ohm} --- 3 resistores.}
\[
15 = 10 + 5 = R + \frac{R}{2}.
\]
Um resistor em série com dois em paralelo:
$10 + \dfrac{10\cdot10}{20} = 10+5 = \SI{15}{\ohm}$.

\medskip
\textbf{(b) \SI{25}{\ohm} --- 4 resistores.}
\[
25 = 20 + 5 = 2R + \frac{R}{2}.
\]
Dois em série ($\SI{20}{\ohm}$) somados a dois em paralelo ($\SI{5}{\ohm}$).

\medskip
\textbf{(c) \SI{4}{\ohm} --- 4 resistores.} Aqui a soma direta não funciona; é
preciso um \emph{paralelo de dois grupos}:
\[
\SI{4}{\ohm} = 20 \parallel 5 = \frac{20\cdot5}{25},
\]
ou seja, \textbf{dois em série} ($\SI{20}{\ohm}$) em paralelo com \textbf{dois em
paralelo} ($\SI{5}{\ohm}$). Verificando:
$\dfrac{20\cdot5}{20+5} = \dfrac{100}{25} = \SI{4}{\ohm}$.

\medskip
\textbf{Método geral.} Uma busca exaustiva sobre todas as combinações série e
paralelo confirma que 3, 4 e 4 são de fato os mínimos. O raciocínio de projeto é:
\begin{enumerate}[leftmargin=1.7em,itemsep=1pt,topsep=2pt]
\item se o alvo for \emph{maior} que $R$, comece somando (série);
\item se for \emph{menor} que $R$, comece dividindo (paralelo);
\item se não fechar, combine \emph{grupos} em paralelo --- como no item (c), onde
      $4 < 10$ exigiu paralelo de blocos, e não de resistores individuais.
\end{enumerate}
Essa habilidade é usada na prática quando um valor comercial não existe e precisa
ser sintetizado com os valores disponíveis na série E12/E24.}
\end{questao}

\blocodesafio

\begin{questao}[4]{}
A figura mostra uma rede em escada \textbf{infinita}, em que todos os resistores
valem $R$. Determine a resistência equivalente $\Req$ entre os terminais A e B.

\circuitoq{%
\begin{circuitikz}[scale=0.9,transform shape,line width=0.8pt]
 \draw (0,2) node[left,font=\bfseries]{A} to[short,o-] (0.8,2);
 \draw (0.8,2) to[R,l=$R$] (2.8,2) coordinate(a1);
 \draw (a1) to[R,l=$R$] (2.8,0) coordinate(b1);
 \draw (a1) to[R,l=$R$] (4.8,2) coordinate(a2);
 \draw (a2) to[R,l=$R$] (4.8,0) coordinate(b2);
 \draw (a2) to[R,l=$R$] (6.8,2) coordinate(a3);
 \draw (a3) to[R,l=$R$] (6.8,0) coordinate(b3);
 \draw[dotted,line width=1.2pt] (6.9,2) -- (8.0,2);
 \draw[dotted,line width=1.2pt] (6.9,0) -- (8.0,0);
 \node[font=\footnotesize] at (8.4,1) {$\cdots\infty$};
 \draw (0.8,0) to[short,o-] (0,0) node[left,font=\bfseries]{B};
 \draw (0.8,2)--(0.8,2); \draw (0.8,0)--(2.8,0);
 \draw (2.8,0)--(4.8,0); \draw (4.8,0)--(6.8,0);
\end{circuitikz}}

\resposta{$\Req = \dfrac{R(1+\sqrt5)}{2} \approx 1{,}618\,R$}
\resolucao{\textbf{Passo 1 --- explorar a autossimilaridade.} Se a rede é
infinita, retirar a primeira célula (um $R$ em série e um $R$ em paralelo) deixa
\emph{exatamente a mesma rede}. Logo, o bloco à direita do primeiro par também
vale $\Req$:
\[
\Req = R + \big(R \parallel \Req\big) = R + \frac{R\,\Req}{R+\Req}.
\]
\textbf{Passo 2 --- resolver a equação.} Multiplicando por $(R+\Req)$:
\[
\Req(R+\Req) = R(R+\Req) + R\,\Req
\;\Rightarrow\;
\Req^{\,2} - R\,\Req - R^{2} = 0.
\]
\textbf{Passo 3 --- raiz positiva.} Pela fórmula quadrática,
\[
\Req = \frac{R \pm \sqrt{R^2+4R^2}}{2} = \frac{R(1\pm\sqrt5)}{2}.
\]
Descarta-se a raiz negativa (resistência é positiva):
\[
\boxed{\ \Req = \frac{R(1+\sqrt5)}{2} \approx 1{,}618\,R\ }
\]
\textbf{Observação notável:} o fator $\dfrac{1+\sqrt5}{2}$ é a \textbf{razão
áurea} $\varphi$ --- a mesma constante da sequência de Fibonacci e da geometria
clássica. Ela aparece aqui porque a recorrência do circuito tem a mesma forma
algébrica $x = 1 + 1/x$.}
\end{questao}

\begin{questao}[4]{}
Quantas seções da escada da questão anterior são necessárias para que a
resistência de entrada difira do valor infinito em menos de
$\SI{0.1}{\percent}$? Considere $R = \SI{1}{\ohm}$ e use a recorrência
$\Req^{(n)} = R + \big(R \parallel \Req^{(n-1)}\big)$, com
$\Req^{(1)} = R$.

\begin{table}[H]
\centering\small\sffamily
\begin{tabular}{@{}c c c@{}}
\toprule
\textbf{$n$ (seções)} & \textbf{$\Req$ ($\Omega$)} & \textbf{erro relativo}\\
\midrule
1  & 1,000000 & \SI{38}{\percent}\\
2  & 1,500000 & \SI{7.3}{\percent}\\
3  & 1,600000 & \SI{1.1}{\percent}\\
5  & 1,617647 & \SI{0.024}{\percent}\\
10 & 1,618034 & $<\SI{0.001}{\percent}$\\
\bottomrule
\end{tabular}
\caption{Convergência da escada finita para o valor infinito.}
\end{table}

\resposta{$n = 5$ seções já bastam}
\resolucao{Aplicando a recorrência iterativamente com $R=\SI{1}{\ohm}$:
\[
\Req^{(1)}=1,\quad
\Req^{(2)}=1+\frac{1\cdot1}{2}=1{,}5,\quad
\Req^{(3)}=1+\frac{1{,}5}{2{,}5}=1{,}6,\ \dots
\]
Comparando com o limite $\varphi \approx 1{,}618034$, o erro cai
\emph{geometricamente}: já em $n=5$ o desvio é de \SI{0.024}{\percent}, abaixo do
critério de \SI{0.1}{\percent}.\\
\textbf{Lição de engenharia:} redes "infinitas" são uma \emph{idealização útil}.
Na prática, meia dúzia de seções já reproduz o comportamento assintótico dentro
da tolerância dos componentes reais (resistores comerciais têm
\SI{1}{\percent} a \SI{5}{\percent} de tolerância). É por isso que atenuadores e
conversores R-2R usam poucos estágios.}
\end{questao}

\begin{questao}[4]{}
Considere agora uma escada infinita em que os resistores \emph{em série} valem $R$
e os \emph{em paralelo} valem $2R$ --- a chamada rede \textbf{R-2R}, usada em
conversores digital-analógico (DAC).

\circuitoq{%
\begin{circuitikz}[scale=0.9,transform shape,line width=0.8pt]
 \draw (0,2) node[left,font=\bfseries]{A} to[short,o-] (0.8,2);
 \draw (0.8,2) to[R,l=$R$] (2.8,2) coordinate(a1);
 \draw (a1) to[R,l=$2R$,color=Rcol] (2.8,0);
 \draw (a1) to[R,l=$R$] (4.8,2) coordinate(a2);
 \draw (a2) to[R,l=$2R$,color=Rcol] (4.8,0);
 \draw (a2) to[R,l=$R$] (6.8,2) coordinate(a3);
 \draw (a3) to[R,l=$2R$,color=Rcol] (6.8,0);
 \draw[dotted,line width=1.2pt] (6.9,2) -- (8.0,2);
 \draw[dotted,line width=1.2pt] (6.9,0) -- (8.0,0);
 \node[font=\footnotesize] at (8.4,1) {$\cdots\infty$};
 \draw (0.8,0) to[short,o-] (0,0) node[left,font=\bfseries]{B};
 \draw (0.8,0)--(6.8,0);
\end{circuitikz}}

\begin{itensq}
\item Determine $\Req$ entre A e B.
\item Explique por que esse resultado torna a rede R-2R especialmente útil.
\end{itensq}
\resposta{(a) $\Req = 2R$ (exato); (b) ver comentário}
\resolucao{(a) Pela autossimilaridade,
\[
\Req = R + \big(2R \parallel \Req\big) = R + \frac{2R\,\Req}{2R+\Req}.
\]
Multiplicando por $(2R+\Req)$:
\[
\Req(2R+\Req) = R(2R+\Req) + 2R\,\Req
\;\Rightarrow\;
\Req^{\,2} - R\,\Req - 2R^{2}=0.
\]
Fatorando: $(\Req - 2R)(\Req + R) = 0$, cuja raiz positiva é
\[
\boxed{\ \Req = 2R\ }
\]
--- um valor \textbf{exato e inteiro}, sem irracionais.\\
(b) Essa é a propriedade notável da rede R-2R: a impedância vista de qualquer
ponto da escada é sempre $2R$, o que a torna \emph{perfeitamente casada} em todos
os estágios. Consequência: a corrente se divide exatamente ao meio a cada nó,
gerando os pesos binários $\tfrac12, \tfrac14, \tfrac18,\dots$ --- exatamente o
que um \textbf{conversor D/A} precisa. Além disso, o circuito usa apenas dois
valores de resistor, o que é decisivo para a precisão em circuitos integrados
(é muito mais fácil casar razões de resistores iguais do que valores absolutos).}
\end{questao}

\begin{questao}[4]{}
Generalize: para uma escada infinita com resistores $R_s$ em série e $R_p$ em
paralelo, obtenha a expressão de $\Req$ e determine a condição sobre
$R_p$ para que $\Req$ seja um múltiplo \emph{racional} de $R_s$.
\resposta{$\Req = \dfrac{R_s + \sqrt{R_s^2+4R_sR_p}}{2}$; racional se
$R_s^2+4R_sR_p$ for quadrado perfeito}
\resolucao{Da recorrência $\Req = R_s + \dfrac{R_p\Req}{R_p+\Req}$, multiplicando
e reorganizando:
\[
\Req^{\,2} - R_s\,\Req - R_s R_p = 0
\quad\Rightarrow\quad
\Req = \frac{R_s + \sqrt{R_s^{2}+4R_sR_p}}{2}.
\]
O resultado é racional quando o discriminante $\Delta = R_s^2+4R_sR_p$ é um
\emph{quadrado perfeito}. Casos particulares:
\begin{itemize}[leftmargin=1.6em,itemsep=1pt,topsep=2pt]
\item $R_p = R_s$: $\Delta = 5R_s^2 \Rightarrow \Req = \varphi R_s$ (irracional);
\item $R_p = 2R_s$: $\Delta = 9R_s^2 \Rightarrow \Req = 2R_s$ (racional);
\item $R_p = 6R_s$: $\Delta = 25R_s^2 \Rightarrow \Req = 3R_s$ (racional).
\end{itemize}
De modo geral, $R_p = \dfrac{k(k-1)}{1}\,R_s$ com $k$ inteiro fornece
$\Req = k\,R_s$. Verifique para $k=2$: $R_p = 2R_s \Rightarrow \Req = 2R_s$.
\textbf{Este é o tipo de raciocínio que separa o técnico do projetista}: em vez de
calcular um caso, obtém-se a \emph{família} de soluções e escolhe-se a que atende
ao requisito.}
\end{questao}

\begin{questao}[4]{}
A figura mostra uma rede \textbf{infinita} em que cada estágio $n$ é formado por
$(n+1)$ resistores \emph{em paralelo}, ligados entre dois barramentos
consecutivos. Os valores seguem o padrão do \textbf{triângulo de Pascal}: no
estágio $n$, os resistores valem $\dfrac{1}{\binom{n}{k}}\,\si{\ohm}$, com
$k = 0,1,\dots,n$.

\circuitoq{%
\begin{circuitikz}[scale=0.82,transform shape,line width=0.75pt,font=\footnotesize]
 % barramentos verticais
 \foreach \x/\h in {0/0.6, 2.4/1.2, 4.8/1.8, 7.2/2.4}
   \draw[cinza] (\x,-\h) -- (\x,\h);
 \draw[cinza] (9.6,-3.0) -- (9.6,3.0);
 \node[left,font=\bfseries\normalsize] at (-0.35,0) {A};
 \draw (-0.35,0)--(0,0);
 % estagio 1: 2 resistores (1,1)
 \foreach \y in {-0.55,0.55}
   \draw (0,\y) to[R,l={$1$}] (2.4,\y);
 % estagio 2: 3 resistores (1, 1/2, 1)
 \foreach \y/\lab in {-1.1/{$1$}, 0/{$\frac12$}, 1.1/{$1$}}
   \draw (2.4,\y) to[R,l={\lab}] (4.8,\y);
 % estagio 3: 4 resistores (1, 1/3, 1/3, 1)
 \foreach \y/\lab in {-1.65/{$1$}, -0.55/{$\frac13$}, 0.55/{$\frac13$}, 1.65/{$1$}}
   \draw (4.8,\y) to[R,l={\lab}] (7.2,\y);
 % estagio 4: 5 resistores (1, 1/4, 1/6, 1/4, 1)
 \foreach \y/\lab in {-2.2/{$1$}, -1.1/{$\frac14$}, 0/{$\frac16$}, 1.1/{$\frac14$}, 2.2/{$1$}}
   \draw (7.2,\y) to[R,l={\lab}] (9.6,\y);
 % reticencias
 \foreach \y in {-2.6,-1.3,0,1.3,2.6}
   \draw[dotted,line width=1pt] (9.7,\y) -- (10.9,\y);
 \node[font=\normalsize] at (11.4,0) {$\infty$};
 \node[cinza,font=\scriptsize] at (1.2,-1.7) {estágio 1};
 \node[cinza,font=\scriptsize] at (3.6,-2.2) {2};
 \node[cinza,font=\scriptsize] at (6.0,-2.6) {3};
 \node[cinza,font=\scriptsize] at (8.4,-3.0) {4};
\end{circuitikz}}

Determine a resistência equivalente $R_{A\infty}$ vista do terminal A.
\resposta{$R_{A\infty} = \SI{1}{\ohm}$ (exatamente)}
\resolucao{\textbf{Passo 1 --- reconhecer a estrutura.} Cada estágio é um
conjunto de resistores \emph{em paralelo} entre dois barramentos. Estágios
consecutivos ficam \emph{em série}. Logo:
\[
R_{A\infty}=\sum_{n=1}^{\infty} R_n,
\qquad\text{onde } R_n \text{ é o paralelo do estágio } n.
\]

\textbf{Passo 2 --- calcular cada estágio pela condutância.} Em paralelo, somam-se
as \emph{condutâncias}. Como o resistor $k$ do estágio $n$ vale
$\dfrac{1}{\binom{n}{k}}$, sua condutância é exatamente $\binom{n}{k}$:
\[
G_n=\sum_{k=0}^{n}\binom{n}{k}.
\]

\textbf{Passo 3 --- a identidade decisiva.} A soma de uma linha inteira do
triângulo de Pascal é uma potência de dois:
\[
\sum_{k=0}^{n}\binom{n}{k}=2^{\,n}
\]
(consequência do binômio de Newton com $x=y=1$: $(1+1)^n = 2^n$). Portanto
\[
G_n = 2^{\,n}
\qquad\Longrightarrow\qquad
R_n=\frac{1}{2^{\,n}}\ \si{\ohm}.
\]
Conferindo os primeiros estágios:
\begin{center}\small\sffamily
\begin{tabular}{@{}c l c c@{}}
\toprule
$n$ & \textbf{Condutâncias} & $G_n$ & $R_n$ ($\Omega$)\\
\midrule
1 & $1,1$              & $2$  & $1/2$\\
2 & $1,2,1$            & $4$  & $1/4$\\
3 & $1,3,3,1$          & $8$  & $1/8$\\
4 & $1,4,6,4,1$        & $16$ & $1/16$\\
5 & $1,5,10,10,5,1$    & $32$ & $1/32$\\
\bottomrule
\end{tabular}
\end{center}

\textbf{Passo 4 --- somar a série.} Os estágios estão em série, então basta somar:
\[
R_{A\infty}=\sum_{n=1}^{\infty}\frac{1}{2^{\,n}}
=\frac12+\frac14+\frac18+\frac{1}{16}+\cdots
\]
É uma \textbf{série geométrica} de razão $q=\tfrac12$ e primeiro termo
$a_1=\tfrac12$:
\[
R_{A\infty}=\frac{a_1}{1-q}=\frac{1/2}{1-1/2}=1.
\]
\[
\boxed{\;R_{A\infty}=\SI{1}{\ohm}\;}
\]

\textbf{Convergência.} Somando apenas os primeiros estágios:
\[
0{,}5 \to 0{,}75 \to 0{,}875 \to 0{,}9375 \to 0{,}96875 \to \cdots \to 1.
\]
Com 8 estágios já se atinge $\SI{99.6}{\percent}$ do valor final --- o erro cai
pela metade a cada estágio acrescentado.

\medskip
\textbf{Por que esta questão é notável.} Ela reúne, em um único problema:
\begin{itemize}[leftmargin=1.7em,itemsep=1pt,topsep=2pt]
\item associação \emph{paralelo} (soma de condutâncias) e \emph{série}
      (soma de resistências), nas duas direções do circuito;
\item uma identidade combinatória (soma da linha do triângulo de Pascal $=2^n$);
\item a convergência de uma série geométrica infinita.
\end{itemize}
Um circuito de complexidade infinita colapsa em uma resposta \textbf{exata e
unitária} --- resultado que só se enxerga reconhecendo o padrão, jamais
calculando estágio a estágio.}
\end{questao}

% BEGIN BANCO COMPLEMENTAR Circuitos mistos
% =====================================================================
%  Banco complementar --- Circuitos Mistos  (REVISADO)
%  Substitui a versao gerada por template (13 clones em N2, 9 em N3,
%  5 questoes conceituais indevidamente marcadas como N4).
%  Sem sobreposicao com o cap03, que ja traz: redes em escada finita e
%  infinita, R-2R, escada generalizada, cubo de resistores (duas
%  versoes), ponte com resistor central e simetria, projeto de valores
%  a partir de resistores de 10 ohm e reducao com fio ideal.
%  Sem transformacao estrela-triangulo (banco proprio).
%  Distribuicao mantida: 8 N1 + 13 N2 + 9 N3 + 5 N4 = 35
% =====================================================================

\subsubsection*{Banco complementar --- Circuitos mistos}

\begin{atencao}[Organização do banco]
Toda rede mista se resolve com o mesmo procedimento: identificar os blocos
puramente série ou puramente paralelo, reduzir \textbf{dos mais internos para
os mais externos} até chegar a $R_{eq}$, e depois \emph{expandir} de volta para
recuperar tensões e correntes de cada elemento. O N3 trata de carregamento,
falhas e saturação térmica; o N4 aborda projeto sob restrição, enumeração de
topologias, diagnóstico e redução por simetria.
\end{atencao}

\blocofixacao

\begin{questao}[1]{}
$R_1=\SI{10}{\ohm}$ está em série com o paralelo de $R_2=\SI{20}{\ohm}$ e
$R_3=\SI{20}{\ohm}$. Determine $R_{eq}$.
\resposta{$R_{eq}=\SI{20}{\ohm}$}
\resolucao{Reduza primeiro o bloco \emph{interno} (o paralelo):
\[
R_{23}=\frac{20\cdot20}{40}=\SI{10}{\ohm}.
\]
Depois o bloco externo (a série):
\[
R_{eq}=10+10=\SI{20}{\ohm}.
\]
A ordem importa: somar os três de uma vez daria $\SI{50}{\ohm}$, resultado sem
sentido físico.}
\end{questao}

\begin{questao}[1]{}
Calcule $R_{eq}$ de $(\SI{10}{\ohm}\parallel\SI{10}{\ohm})$ em série com
$(\SI{30}{\ohm}\parallel\SI{60}{\ohm})$.
\resposta{$R_{eq}=\SI{25}{\ohm}$}
\resolucao{Dois blocos paralelos independentes, depois somados:
\[
R_{a}=\frac{10}{2}=\SI{5}{\ohm},
\qquad
R_{b}=\frac{30\cdot60}{90}=\SI{20}{\ohm},
\]
\[
R_{eq}=5+20=\SI{25}{\ohm}.
\]
Blocos que não compartilham nós podem ser reduzidos em qualquer ordem entre si
--- só a hierarquia interna de cada um é obrigatória.}
\end{questao}

\begin{questao}[1]{}
Dois resistores estão ligados \emph{aos mesmos dois nós} do circuito. Eles
estão:
\begin{alternativas}[2]
\item em série
\item em paralelo
\item em série ou paralelo, conforme o desenho
\item nem em série nem em paralelo
\end{alternativas}
\resposta{(b)}
\resolucao{O critério é \textbf{topológico}, não gráfico: dois elementos ligados
ao mesmo par de nós estão submetidos à mesma tensão --- definição de paralelo.
Já a série exige que compartilhem \emph{um} nó \emph{e} que nenhum outro
elemento se conecte a ele, garantindo a mesma corrente.\\
Não se guie pelo desenho: resistores que "parecem" em série no esquema podem
estar em paralelo se um fio os ligar pelos dois lados.}
\end{questao}

\begin{questao}[1]{}
A estratégia correta para reduzir uma rede mista é:
\begin{alternativas}[2]
\item começar sempre pelos resistores mais próximos da fonte
\item reduzir dos blocos mais internos para os mais externos
\item somar todas as resistências e depois corrigir
\item reduzir os paralelos antes de qualquer série, sempre
\end{alternativas}
\resposta{(b)}
\resolucao{Reduz-se sempre \emph{de dentro para fora}, isto é, partindo do
bloco mais distante da fonte (ou mais aninhado) em direção aos terminais. Não
existe regra fixa de "paralelo antes de série": a ordem é ditada pela
hierarquia da própria rede.\\
A alternativa (d) falha, por exemplo, em $R_1+(R_2+R_3)\parallel R_4$, onde a
\emph{série} $R_2+R_3$ precisa ser feita antes do paralelo.}
\end{questao}

\begin{questao}[1]{}
Determine $R_{eq}$ de $(\SI{10}{\ohm}+\SI{20}{\ohm})$ em paralelo com
$(\SI{15}{\ohm}+\SI{15}{\ohm})$.
\resposta{$R_{eq}=\SI{15}{\ohm}$}
\resolucao{Cada ramo é reduzido primeiro:
$R_a=10+20=\SI{30}{\ohm}$ e $R_b=15+15=\SI{30}{\ohm}$. Depois o paralelo:
\[
R_{eq}=\frac{30}{2}=\SI{15}{\ohm}.
\]
Aqui a série \emph{obrigatoriamente} vem antes: só depois de reduzidos os dois
ramos é que eles ficam ligados ao mesmo par de nós.}
\end{questao}

\begin{questao}[1]{}
$R_1=\SI{12}{\ohm}$ está em série com o paralelo entre $R_2=\SI{6}{\ohm}$ e a
série $R_3=\SI{4}{\ohm}$ com $R_4=\SI{8}{\ohm}$. Determine $R_{eq}$.
\resposta{$R_{eq}=\SI{16}{\ohm}$}
\resolucao{Três níveis, do mais interno ao mais externo:
\[
R_{34}=4+8=\SI{12}{\ohm};
\qquad
R_{234}=\frac{6\cdot12}{18}=\SI{4}{\ohm};
\qquad
R_{eq}=12+4=\SI{16}{\ohm}.
\]}
\end{questao}

\begin{questao}[1]{}
Uma fonte de $\SI{24}{\volt}$ alimenta a rede
$\SI{10}{\ohm}+(\SI{20}{\ohm}\parallel\SI{20}{\ohm})$. Determine a corrente
total e a tensão sobre o bloco paralelo.
\resposta{$I=\SI{1.2}{\ampere}$; $U_{par}=\SI{12}{\volt}$}
\resolucao{Já se sabe que $R_{eq}=\SI{20}{\ohm}$, logo
\[
I=\frac{24}{20}=\SI{1.2}{\ampere}.
\]
Toda a corrente atravessa $R_1$, e depois o bloco paralelo:
\[
U_{par}=R_{23}\,I=10\cdot1{,}2=\SI{12}{\volt}.
\]
\textbf{Conferência (LKT):} $U_1=10\cdot1{,}2=\SI{12}{\volt}$, e
$12+12=\SI{24}{\volt}$ \checkmark.}
\end{questao}

\begin{questao}[1]{}
Em uma rede mista, a corrente que atravessa o resistor imediatamente ligado ao
terminal da fonte é:
\begin{alternativas}[2]
\item igual à corrente total
\item metade da corrente total
\item dividida entre os ramos internos
\item indeterminada sem mais dados
\end{alternativas}
\resposta{(a)}
\resolucao{Se esse resistor está \emph{em série} com a fonte --- isto é, se não
há bifurcação antes dele --- toda a corrente que sai da fonte passa por ele.
Por isso ele é o primeiro a ser recuperado na etapa de expansão, e costuma ser
também o mais solicitado termicamente.\\
Note a ressalva: se houver um nó entre a fonte e o resistor, a alternativa (c)
passaria a valer. O critério é sempre topológico.}
\end{questao}

\blocoaplicacao

\begin{questao}[2]{}
Uma fonte de $\SI{24}{\volt}$ alimenta $R_1=\SI{10}{\ohm}$ em série com
$R_2=\SI{20}{\ohm}$ e $R_3=\SI{20}{\ohm}$ em paralelo. Determine todas as
correntes e tensões.
\resposta{$I_t=\SI{1.2}{\ampere}$; $U_1=\SI{12}{\volt}$;
$U_{23}=\SI{12}{\volt}$; $I_2=I_3=\SI{0.6}{\ampere}$}
\resolucao{\textbf{Redução:} $R_{23}=\SI{10}{\ohm}$, $R_{eq}=\SI{20}{\ohm}$,
$I_t=24/20=\SI{1.2}{\ampere}$.\\
\textbf{Expansão:} $U_1=10(1{,}2)=\SI{12}{\volt}$;
$U_{23}=24-12=\SI{12}{\volt}$;
\[
I_2=I_3=\frac{12}{20}=\SI{0.6}{\ampere}.
\]
\textbf{Verificação (LKC):} $0{,}6+0{,}6=\SI{1.2}{\ampere}$ \checkmark.
Como $R_2=R_3$, a corrente se divide igualmente --- resultado que a simetria
antecipa sem nenhum cálculo.}
\end{questao}

\begin{questao}[2]{}
No circuito anterior, $R_3$ é trocado por $\SI{60}{\ohm}$. Recalcule
$R_{eq}$, a corrente total e as correntes de ramo.
\resposta{$R_{eq}=\SI{25}{\ohm}$; $I_t=\SI{0.96}{\ampere}$;
$I_2=\SI{0.72}{\ampere}$, $I_3=\SI{0.24}{\ampere}$}
\resolucao{$R_{23}=\dfrac{20\cdot60}{80}=\SI{15}{\ohm}$ e
$R_{eq}=10+15=\SI{25}{\ohm}$, logo $I_t=24/25=\SI{0.96}{\ampere}$.\\
$U_{23}=15(0{,}96)=\SI{14.4}{\volt}$, e daí
\[
I_2=\frac{14{,}4}{20}=\SI{0.72}{\ampere},
\qquad
I_3=\frac{14{,}4}{60}=\SI{0.24}{\ampere}.
\]
\textbf{Divisor de corrente (alternativa):}
$I_2=0{,}96\cdot\frac{60}{80}=\SI{0.72}{\ampere}$ \checkmark. Note que a razão
$I_2/I_3=3$ é exatamente $R_3/R_2$ --- inversa das resistências.}
\end{questao}

\begin{questao}[2]{}
Na rede $\SI{10}{\ohm}+(\SI{20}{\ohm}\parallel\SI{20}{\ohm})$ sob
$\SI{24}{\volt}$, calcule a potência de cada resistor e confronte com a
potência da fonte.
\resposta{$\SI{14.4}{\watt}$; $\SI{7.2}{\watt}$; $\SI{7.2}{\watt}$; total
$\SI{28.8}{\watt}$}
\resolucao{
\[
P_1=R_1I_t^{2}=10(1{,}2)^{2}=\SI{14.4}{\watt};
\qquad
P_2=P_3=\frac{U_{23}^{2}}{20}=\frac{144}{20}=\SI{7.2}{\watt}.
\]
Fonte: $P=24(1{,}2)=\SI{28.8}{\watt}=14{,}4+7{,}2+7{,}2$ \checkmark.\\
\textbf{Observação:} $R_1$ dissipa exatamente metade de toda a potência, embora
seja o menor resistor da rede --- porque é o único que conduz a corrente
inteira. Em redes mistas, "quem aquece mais" depende da posição, não apenas do
valor.}
\end{questao}

\begin{questao}[2]{}
Determine $R_{eq}$ de $(\SI{5}{\ohm}+\SI{15}{\ohm})$ em paralelo com
$(\SI{10}{\ohm}+\SI{10}{\ohm})$, e a corrente de cada ramo sob
$\SI{30}{\volt}$.
\resposta{$R_{eq}=\SI{10}{\ohm}$; $I_t=\SI{3}{\ampere}$; $\SI{1.5}{\ampere}$
em cada ramo}
\resolucao{Ramos: $R_a=\SI{20}{\ohm}$ e $R_b=\SI{20}{\ohm}$, logo
$R_{eq}=\SI{10}{\ohm}$ e $I_t=30/10=\SI{3}{\ampere}$.\\
Ambos os ramos estão sob os mesmos $\SI{30}{\volt}$:
$I_a=I_b=30/20=\SI{1.5}{\ampere}$.\\
\textbf{Cuidado:} apesar de $R_a=R_b$, as \emph{quedas internas} diferem. No
ramo $a$: $\SI{7.5}{\volt}$ e $\SI{22.5}{\volt}$; no ramo $b$:
$\SI{15}{\volt}$ e $\SI{15}{\volt}$. Ramos equivalentes nos terminais não são
equivalentes internamente.}
\end{questao}

\begin{questao}[2]{}
Uma rede é formada por $R_1=\SI{8}{\ohm}$ em série com o paralelo de
$\SI{24}{\ohm}$ e $R$. Sabendo que $R_{eq}=\SI{24}{\ohm}$, determine $R$.
\resposta{$R=\SI{48}{\ohm}$}
\resolucao{O bloco paralelo deve valer $24-8=\SI{16}{\ohm}$. Em condutâncias:
\[
\frac{1}{R}=\frac{1}{16}-\frac{1}{24}=\frac{3-2}{48}
\;\Longrightarrow\;
R=\SI{48}{\ohm}.
\]
\textbf{Conferência:} $24\parallel48=\dfrac{24\cdot48}{72}=\SI{16}{\ohm}$ e
$8+16=\SI{24}{\ohm}$ \checkmark.\\
Note que a subtração foi feita em \emph{condutâncias}: tentar subtrair
resistências ($16-24$) daria valor negativo, sinal claro de método errado.}
\end{questao}

\begin{questao}[2]{}
Determine $R_{eq}$ de $R_1=\SI{10}{\ohm}$ em série com o paralelo entre
$(\SI{20}{\ohm}+\SI{20}{\ohm})$ e $\SI{10}{\ohm}$.
\resposta{$R_{eq}=\SI{18}{\ohm}$}
\resolucao{
\[
R_{23}=20+20=\SI{40}{\ohm};
\qquad
R_{234}=\frac{40\cdot10}{50}=\SI{8}{\ohm};
\qquad
R_{eq}=10+8=\SI{18}{\ohm}.
\]
O bloco paralelo resultou em $\SI{8}{\ohm}$ --- menor que o menor dos dois
($\SI{10}{\ohm}$), como sempre deve ocorrer. Esse teste de sanidade a cada
etapa evita que um erro se propague por toda a redução.}
\end{questao}

\begin{questao}[2]{}
Na rede $\SI{10}{\ohm}+(\SI{20}{\ohm}\parallel\SI{20}{\ohm})$ sob
$\SI{24}{\volt}$, o nó entre $R_1$ e o bloco paralelo é chamado $A$, e o
terminal negativo da fonte é o terra. Determine $V_A$.
\resposta{$V_A=\SI{12}{\volt}$}
\resolucao{O potencial de $A$ em relação ao terra é a tensão sobre o bloco
paralelo, pois é ele que separa $A$ do terra:
\[
V_A=U_{23}=\SI{12}{\volt}.
\]
\textbf{Alternativamente,} descendo da fonte: $V_A=24-U_1=24-12=\SI{12}{\volt}$
\checkmark.\\
Os dois caminhos precisam concordar --- é a LKT em ação. Se discordassem,
haveria erro na redução.}
\end{questao}

\begin{questao}[2]{}
Na mesma rede, um dos resistores de $\SI{20}{\ohm}$ sofre
\textbf{curto-circuito}. Recalcule $R_{eq}$, a corrente total e a potência de
$R_1$.
\resposta{$R_{eq}=\SI{10}{\ohm}$; $I_t=\SI{2.4}{\ampere}$;
$P_1=\SI{57.6}{\watt}$}
\resolucao{Um curto no bloco paralelo anula o bloco inteiro: o outro resistor
de $\SI{20}{\ohm}$ fica em paralelo com $\SI{0}{\ohm}$ e deixa de conduzir.
\[
R_{eq}=10+0=\SI{10}{\ohm},
\qquad
I_t=\frac{24}{10}=\SI{2.4}{\ampere},
\qquad
P_1=10(2{,}4)^{2}=\SI{57.6}{\watt}.
\]
$P_1$ saltou de $14{,}4$ para $\SI{57.6}{\watt}$ --- fator \textbf{quatro},
pois a corrente dobrou. Se $R_1$ fosse especificado para
$\SI{25}{\watt}$, queimaria em seguida.\\
\textbf{Padrão de falha:} o curto em um ramo interno transfere toda a
solicitação para o elemento em série --- e o defeito "migra" da causa para a
vítima, dificultando o diagnóstico.}
\end{questao}

\begin{questao}[2]{}
Na mesma rede, agora um dos resistores de $\SI{20}{\ohm}$ \textbf{abre}.
Recalcule $R_{eq}$, a corrente total e a potência do resistor remanescente.
\resposta{$R_{eq}=\SI{30}{\ohm}$; $I_t=\SI{0.8}{\ampere}$;
$P=\SI{12.8}{\watt}$}
\resolucao{Sem um dos ramos, o bloco paralelo passa a valer $\SI{20}{\ohm}$:
\[
R_{eq}=10+20=\SI{30}{\ohm},
\qquad
I_t=\frac{24}{30}=\SI{0.8}{\ampere},
\qquad
U_{par}=20(0{,}8)=\SI{16}{\volt},
\]
\[
P=\frac{16^{2}}{20}=\SI{12.8}{\watt}.
\]
\textbf{Contraste com o curto:} aqui a corrente total \emph{caiu} (de $1{,}2$
para $\SI{0.8}{\ampere}$), mas o resistor sobrevivente passou a dissipar
$12{,}8$ em vez de $\SI{7.2}{\watt}$ --- $78\%$ mais, porque agora assume
sozinho a corrente antes dividida. Falhas de abertura em blocos paralelos
\emph{sobrecarregam} os ramos irmãos.}
\end{questao}

\begin{questao}[2]{}
Duas redes têm $R_{eq}=\SI{20}{\ohm}$: (A) $\SI{10}{\ohm}$ em série com
$\SI{20}{\ohm}\parallel\SI{20}{\ohm}$; (B) $\SI{30}{\ohm}\parallel\SI{60}{\ohm}$.
Sob $\SI{60}{\volt}$, compare a potência total e a do componente mais
solicitado.
\resposta{Total $\SI{180}{\watt}$ em ambas; pior componente:
$\SI{90}{\watt}$ em (A) e $\SI{120}{\watt}$ em (B)}
\resolucao{$I_t=60/20=\SI{3}{\ampere}$ e $P_{tot}=60(3)=\SI{180}{\watt}$ nas
duas.\\
\textbf{(A)} $P_{10}=10(3)^{2}=\SI{90}{\watt}$; bloco sob
$\SI{30}{\volt}$, cada $\SI{20}{\ohm}$ com $\SI{45}{\watt}$. Soma:
$90+45+45=\SI{180}{\watt}$ \checkmark.\\
\textbf{(B)} $P_{30}=\dfrac{3600}{30}=\SI{120}{\watt}$;
$P_{60}=\dfrac{3600}{60}=\SI{60}{\watt}$. Soma $\SI{180}{\watt}$ \checkmark.\\
\textbf{Conclusão:} resistência equivalente idêntica, potência total idêntica,
mas o ponto quente difere em $33\%$. A escolha da topologia é decisão térmica,
não elétrica.}
\end{questao}

\begin{questao}[2]{}
Uma fonte de $\EMF=\SI{24}{\volt}$ e $r=\SI{2}{\ohm}$ alimenta a rede mista de
$\SI{10}{\ohm}$. Determine a corrente, a tensão nos terminais e a potência
entregue à rede.
\resposta{$I=\SI{2}{\ampere}$; $V=\SI{20}{\volt}$; $P=\SI{40}{\watt}$}
\resolucao{
\[
I=\frac{24}{2+10}=\SI{2}{\ampere};
\qquad
V=24-2(2)=\SI{20}{\volt};
\qquad
P_{rede}=20(2)=\SI{40}{\watt}.
\]
A fonte gera $24(2)=\SI{48}{\watt}$, dos quais $2(2)^{2}=\SI{8}{\watt}$ se
perdem internamente. Rendimento: $40/48=83{,}3\%$.\\
\textbf{Regra útil:} o rendimento é $R_{rede}/(R_{rede}+r)$ --- aqui
$10/12$. Ele depende apenas dessa razão, não da topologia interna da rede.}
\end{questao}

\begin{questao}[2]{}
Em uma rede mista de $R_{eq}=\SI{16}{\ohm}$ sob $\SI{32}{\volt}$, verificou-se
$P_{total}=\SI{64}{\watt}$. A medição é consistente? Justifique de duas
maneiras.
\resposta{Sim: $P=V^{2}/R_{eq}=VI=\SI{64}{\watt}$}
\resolucao{\textbf{Primeiro caminho:}
$P=\dfrac{V^{2}}{R_{eq}}=\dfrac{32^{2}}{16}=\SI{64}{\watt}$ \checkmark.\\
\textbf{Segundo caminho:} $I=32/16=\SI{2}{\ampere}$ e
$P=VI=32(2)=\SI{64}{\watt}$ \checkmark.\\
\textbf{Terceiro (se os elementos forem conhecidos):} somar $R_kI_k^{2}$ de
cada resistor deve devolver os mesmos $\SI{64}{\watt}$. Essa terceira
verificação é a mais forte, porque testa a \emph{expansão} da rede --- as duas
primeiras só testam $R_{eq}$ e passariam mesmo se a distribuição interna
estivesse errada.}
\end{questao}

\begin{questao}[2]{}
Na rede $R_1=\SI{12}{\ohm}$ em série com $R_2=\SI{6}{\ohm}\parallel
(R_3=\SI{4}{\ohm}+R_4=\SI{8}{\ohm})$, sob $\SI{32}{\volt}$, determine a
corrente em $R_3$.
\resposta{$I_3=\SI{0.667}{\ampere}$}
\resolucao{\textbf{Redução:} $R_{34}=\SI{12}{\ohm}$;
$R_{234}=6\parallel12=\SI{4}{\ohm}$; $R_{eq}=12+4=\SI{16}{\ohm}$.\\
\textbf{Expansão:} $I_t=32/16=\SI{2}{\ampere}$;
$U_{234}=4(2)=\SI{8}{\volt}$.\\
Esse é o mesmo $\SI{8}{\volt}$ aplicado ao ramo $R_3+R_4$:
\[
I_3=\frac{8}{12}=\SI{0.667}{\ampere}.
\]
\textbf{Conferência pelo divisor:}
$I_3=I_t\cdot\dfrac{6}{6+12}=2\cdot\dfrac13=\SI{0.667}{\ampere}$ \checkmark. E
$I_2=8/6=\SI{1.333}{\ampere}$, com $0{,}667+1{,}333=\SI{2}{\ampere}$
\checkmark.}
\end{questao}

\blocoaprofundamento

\begin{questao}[3]{}
Uma rede consiste em $(R_1+R_2)$ em paralelo com $(R_3+R_4)$.
\begin{itensq}
\item Obtenha a expressão de $R_{eq}$.
\item Avalie para $R_1=\SI{8}{\ohm}$, $R_2=\SI{16}{\ohm}$,
      $R_3=R_4=\SI{12}{\ohm}$.
\item Analise o caso-limite $R_3+R_4\gg R_1+R_2$.
\end{itensq}
\resposta{(a) $R_{eq}=\dfrac{(R_1+R_2)(R_3+R_4)}{R_1+R_2+R_3+R_4}$;
(b) $\SI{12}{\ohm}$; (c) $R_{eq}\to R_1+R_2$}
\resolucao{(a) Reduzindo cada ramo e aplicando produto pela soma:
\[
R_{eq}=\frac{(R_1+R_2)(R_3+R_4)}{(R_1+R_2)+(R_3+R_4)}.
\]
(b) $R_a=\SI{24}{\ohm}$, $R_b=\SI{24}{\ohm}$, logo
$R_{eq}=\dfrac{24\cdot24}{48}=\SI{12}{\ohm}$.\\
(c) Se $R_b\gg R_a$, divida numerador e denominador por $R_b$:
\[
R_{eq}=\frac{R_a}{\dfrac{R_a}{R_b}+1}\longrightarrow R_a=R_1+R_2.
\]
O ramo de alta resistência comporta-se como circuito aberto e não participa.
\textbf{Uso prático:} é essa aproximação que justifica desprezar a resistência
de um voltímetro ligado em paralelo a um trecho de circuito --- desde que ela
seja, digamos, cem vezes maior que o trecho medido.}
\end{questao}

\begin{questao}[3]{}
Um divisor formado por $R_1=\SI{140}{\ohm}$ e $R_2=\SI{100}{\ohm}$ é alimentado
por $\SI{12}{\volt}$.
\begin{itensq}
\item Determine a tensão de saída em vazio.
\item Recalcule com uma carga de $\SI{1}{\kilo\ohm}$ conectada à saída.
\item Determine o erro percentual introduzido pela carga.
\end{itensq}
\resposta{(a) $\SI{5.00}{\volt}$; (b) $\SI{4.72}{\volt}$; (c) $-5{,}5\%$}
\resolucao{(a) $U_{2}=12\cdot\dfrac{100}{240}=\SI{5.00}{\volt}$.\\
(b) A carga entra em \textbf{paralelo} com $R_2$:
\[
R_2'=\frac{100\cdot1000}{1100}=\SI{90.91}{\ohm},
\qquad
U_2=12\cdot\frac{90{,}91}{140+90{,}91}=\SI{4.72}{\volt}.
\]
(c) Erro $=(4{,}72-5{,}00)/5{,}00=-5{,}5\%$.\\
\textbf{Diagnóstico:} a carga de $\SI{1}{\kilo\ohm}$ é apenas $10\times$ maior
que $R_2$ --- insuficiente. A regra de bolso exige
$R_L\ge10R_2$ para erro em torno de $5\%$, e $R_L\ge100R_2$ para erro em torno
de $0{,}5\%$. Um divisor resistivo não é fonte de tensão: ele tem resistência
interna $R_1\parallel R_2=\SI{58.3}{\ohm}$, e é ela que produz a queda ao
ser carregada. O projeto para uma carga \emph{conhecida} é feito no bloco de
desafio.}
\end{questao}

\begin{questao}[3]{}
Uma rede tem $R_1=\SI{6}{\ohm}$ em série com $(\SI{12}{\ohm}\parallel
\SI{12}{\ohm})$, e este em série com $(\SI{8}{\ohm}\parallel\SI{8}{\ohm})$, sob
$\SI{32}{\volt}$. Determine a potência de cada um dos cinco resistores.
\resposta{$\SI{24}{\watt}$ em $R_1$; $\SI{12}{\watt}$ em cada
$\SI{12}{\ohm}$; $\SI{8}{\watt}$ em cada $\SI{8}{\ohm}$}
\resolucao{\textbf{Redução:} $12\parallel12=\SI{6}{\ohm}$;
$8\parallel8=\SI{4}{\ohm}$; $R_{eq}=6+6+4=\SI{16}{\ohm}$ e
$I_t=32/16=\SI{2}{\ampere}$.\\
\textbf{Expansão:}
\begin{itemize}
\item $R_1$: conduz $\SI{2}{\ampere}$ $\Rightarrow$
      $P=6(4)=\SI{24}{\watt}$.
\item Bloco de $\SI{12}{\ohm}$: $U=6(2)=\SI{12}{\volt}$; cada ramo com
      $\SI{1}{\ampere}$ $\Rightarrow$ $P=12(1)=\SI{12}{\watt}$.
\item Bloco de $\SI{8}{\ohm}$: $U=4(2)=\SI{8}{\volt}$; cada ramo com
      $\SI{1}{\ampere}$ $\Rightarrow$ $P=8(1)=\SI{8}{\watt}$.
\end{itemize}
\textbf{Balanço:} $24+12+12+8+8=\SI{64}{\watt}=32(2)$ \checkmark.\\
Observe que $R_1$, sendo o menor resistor da rede, é o que mais dissipa ---
porque conduz o dobro da corrente dos demais e a potência vai com $I^{2}$.}
\end{questao}

\begin{questao}[3]{}
Um reostato de $\SI{0}{\ohm}$ a $\SI{60}{\ohm}$ está em paralelo com
$\SI{30}{\ohm}$, e o conjunto em série com $\SI{20}{\ohm}$, sob
$\SI{60}{\volt}$.
\begin{itensq}
\item Determine a faixa de $R_{eq}$.
\item Determine a faixa da corrente total.
\item Explique por que a faixa é limitada.
\end{itensq}
\resposta{(a) de $\SI{20}{\ohm}$ a $\SI{40}{\ohm}$; (b) de $\SI{1.5}{\ampere}$
a $\SI{3}{\ampere}$}
\resolucao{(a) Com o reostato em $\SI{0}{\ohm}$, o bloco paralelo é
curto-circuitado e vale $\SI{0}{\ohm}$: $R_{eq}=\SI{20}{\ohm}$.\\
Com o reostato em $\SI{60}{\ohm}$: $30\parallel60=\SI{20}{\ohm}$ e
$R_{eq}=\SI{40}{\ohm}$.\\
(b) $I_{\max}=60/20=\SI{3}{\ampere}$ e
$I_{\min}=60/40=\SI{1.5}{\ampere}$: faixa de $2{:}1$.\\
(c) Dois limites atuam. O resistor \emph{série} de $\SI{20}{\ohm}$ impõe um
teto à corrente ($\SI{3}{\ampere}$), por mais que o bloco paralelo seja
anulado. E o resistor de $\SI{30}{\ohm}$ \emph{em paralelo} com o reostato
impõe um piso a $R_{eq}$: mesmo com o reostato em $\infty$, o bloco não
passaria de $\SI{30}{\ohm}$, e $R_{eq}$ de $\SI{50}{\ohm}$.\\
\textbf{Consequência de projeto:} para ampliar a faixa de controle, é preciso
reduzir o resistor série \emph{e} aumentar o resistor em paralelo --- alterar
apenas o reostato tem efeito saturante.}
\end{questao}

\begin{questao}[3]{}
Na rede $\SI{10}{\ohm}+(\SI{20}{\ohm}\parallel\SI{20}{\ohm})$, todos os
resistores têm potência nominal de $\SI{10}{\watt}$. Determine a tensão máxima
aplicável e identifique o componente crítico.
\resposta{$V_{\max}=\SI{20}{\volt}$; o crítico é $R_1$}
\resolucao{Expresse cada potência em função da corrente total $I$:
\[
P_1=10I^{2};
\qquad
P_2=P_3=20\left(\frac{I}{2}\right)^{2}=5I^{2}.
\]
$R_1$ dissipa \textbf{o dobro} de cada ramo paralelo --- é ele que satura
primeiro:
\[
10I^{2}=10 \;\Longrightarrow\; I=\SI{1}{\ampere}
\;\Longrightarrow\;
V_{\max}=R_{eq}I=20(1)=\SI{20}{\volt}.
\]
Nessa condição os ramos paralelos operam a $\SI{5}{\watt}$ --- metade do
nominal, ou seja, subutilizados.\\
\textbf{Melhoria possível:} trocar $R_1$ por dois de $\SI{20}{\ohm}$ em
paralelo (mesmos $\SI{10}{\ohm}$, mas $\SI{20}{\watt}$ de capacidade) elevaria
$V_{\max}$ para o limite imposto pelos ramos:
$5I^{2}=10\Rightarrow I=\SI{1.41}{\ampere}$ e
$V_{\max}=\SI{28.3}{\volt}$. Identificar o componente crítico é o primeiro
passo para melhorar qualquer rede.}
\end{questao}

\begin{questao}[3]{}
Obtenha $\SI{24}{\ohm}$ dispondo apenas de resistores de $\SI{36}{\ohm}$.
Determine a associação de \textbf{menor} número de unidades e prove que não há
solução com duas.
\resposta{Três unidades: $(\SI{36}{\ohm}+\SI{36}{\ohm})\parallel\SI{36}{\ohm}$}
\resolucao{\textbf{Com duas unidades} só existem dois valores possíveis:
$36+36=\SI{72}{\ohm}$ e $36\parallel36=\SI{18}{\ohm}$. Nenhum é
$\SI{24}{\ohm}$ --- logo duas são insuficientes.\\
\textbf{Com três:}
\[
(36+36)\parallel36=\frac{72\cdot36}{108}=\SI{24}{\ohm}.\ \checkmark
\]
\textbf{Verificação por outro caminho:} três em paralelo dariam
$\SI{12}{\ohm}$; três em série, $\SI{108}{\ohm}$; e
$36+(36\parallel36)=\SI{54}{\ohm}$. Das quatro topologias possíveis com três
unidades, apenas a indicada resolve.\\
\textbf{Observação:} $24=\frac23\cdot36$, e a razão $2/3$ é característica da
combinação "dois em série contra um em paralelo" --- padrão útil de memorizar,
pois $\dfrac{2R\cdot R}{3R}=\dfrac{2R}{3}$ para qualquer $R$.}
\end{questao}

\begin{questao}[3]{}
Em uma rede mista alimentada por fonte ideal, um resistor de um ramo paralelo
abre. Descreva um procedimento sistemático para localizar a falha usando apenas
um voltímetro, e aplique-o à rede
$\SI{10}{\ohm}+(\SI{20}{\ohm}\parallel\SI{20}{\ohm})$ sob $\SI{24}{\volt}$.
\resposta{Comparar as tensões medidas com as previstas; o bloco defeituoso
apresenta tensão \emph{maior} que a esperada}
\resolucao{\textbf{Procedimento.} \emph{(i)} Calcule as tensões esperadas em
operação normal. \emph{(ii)} Meça, bloco a bloco, a partir da fonte.
\emph{(iii)} Compare: um aumento de tensão indica bloco com resistência
\emph{maior} que a prevista (abertura); uma queda indica resistência
\emph{menor} (curto).\\
\textbf{Aplicação.} Previsto: $U_1=\SI{12}{\volt}$ e
$U_{par}=\SI{12}{\volt}$. Com um ramo aberto,
$R_{eq}=\SI{30}{\ohm}$, $I=\SI{0.8}{\ampere}$ e:
\[
U_1=10(0{,}8)=\SI{8}{\volt}\ (\downarrow),
\qquad
U_{par}=20(0{,}8)=\SI{16}{\volt}\ (\uparrow).
\]
O voltímetro acusa $\SI{16}{\volt}$ onde se esperavam $\SI{12}{\volt}$ ---
o bloco paralelo é o suspeito.\\
\textbf{Confirmação:} isolando o bloco, um ohmímetro leria
$\SI{20}{\ohm}$ em vez de $\SI{10}{\ohm}$. E \emph{qual} dos dois ramos
abriu se identifica medindo cada um separadamente --- a tensão sozinha não
distingue ramos idênticos em paralelo.\\
\textbf{Por que $U_1$ cai:} a corrente total diminuiu. Note a assinatura
característica: em uma abertura, as tensões se \emph{redistribuem} com soma
constante ($8+16=24$), enquanto em um curto elas colapsam localmente.}
\end{questao}

\begin{questao}[3]{}
Uma rede mista de $R_{eq}=\SI{20}{\ohm}$ é alimentada por fonte de
$\EMF=\SI{60}{\volt}$ com $r=\SI{5}{\ohm}$.
\begin{itensq}
\item Determine a corrente, a tensão na rede e o rendimento.
\item Determine o valor de $R_{eq}$ que maximizaria a potência na rede.
\item Compare o rendimento nas duas situações.
\end{itensq}
\resposta{(a) $\SI{2.4}{\ampere}$, $\SI{48}{\volt}$, $80\%$;
(b) $R_{eq}=\SI{5}{\ohm}$; (c) $80\%$ contra $50\%$}
\resolucao{(a) $I=\dfrac{60}{25}=\SI{2.4}{\ampere}$;
$V=20(2{,}4)=\SI{48}{\volt}$;
$P_{rede}=48(2{,}4)=\SI{115.2}{\watt}$; rendimento
$=20/25=80\%$.\\
(b) A potência na rede é máxima quando $R_{eq}=r=\SI{5}{\ohm}$. Então
$I=60/10=\SI{6}{\ampere}$, $V=\SI{30}{\volt}$ e
$P=\SI{180}{\watt}$.\\
(c) Rendimento na condição de casamento: $5/10=50\%$. A fonte passa a
gerar $\SI{360}{\watt}$, dos quais $\SI{180}{\watt}$ se perdem internamente.\\
\textbf{Conflito de objetivos:} entregar \emph{mais} potência
($180$ contra $\SI{115.2}{\watt}$) custa despachar $\SI{360}{\watt}$ em vez de
$\SI{144}{\watt}$ --- e aquecer a fonte com $\SI{180}{\watt}$. Em eletrônica de
sinal, o casamento vale a pena; em eletrotécnica de potência, jamais.}
\end{questao}

\begin{questao}[3]{}
Uma fonte de \textbf{corrente} de $\SI{3}{\ampere}$ alimenta $R_1=\SI{20}{\ohm}$
em paralelo com o ramo formado por $R_2=\SI{10}{\ohm}$ em série com
$(\SI{20}{\ohm}\parallel\SI{20}{\ohm})$.
\begin{itensq}
\item Determine $R_{eq}$ e a tensão da associação.
\item Obtenha a corrente em cada um dos cinco resistores.
\item Verifique o balanço de potência.
\end{itensq}
\resposta{(a) $\SI{10}{\ohm}$ e $\SI{30}{\volt}$;
(b) $\SI{1.5}{\ampere}$ em $R_1$ e em $R_2$, $\SI{0.75}{\ampere}$ em cada
$\SI{20}{\ohm}$ do bloco; (c) $\SI{90}{\watt}$}
\resolucao{(a) O ramo vale $10+(20\parallel20)=10+10=\SI{20}{\ohm}$, e
\[
R_{eq}=20\parallel20=\SI{10}{\ohm}
\;\Longrightarrow\;
U=I_tR_{eq}=3(10)=\SI{30}{\volt}.
\]
\textbf{Diferença essencial em relação à fonte de tensão:} aqui $U$ é
\emph{resultado}, não dado. Se a rede mudar, a tensão muda --- a corrente é que
permanece fixa.\\
(b) Os dois ramos estão sob $\SI{30}{\volt}$ e têm $\SI{20}{\ohm}$ cada:
$I_{R_1}=I_{ramo}=\SI{1.5}{\ampere}$. Dentro do ramo,
$U_{R_2}=10(1{,}5)=\SI{15}{\volt}$ e restam $\SI{15}{\volt}$ para o bloco
paralelo:
\[
I=\frac{15}{20}=\SI{0.75}{\ampere}\ \text{em cada}.
\]
\textbf{LKC:} $0{,}75+0{,}75=\SI{1.5}{\ampere}$ \checkmark; e
$1{,}5+1{,}5=\SI{3}{\ampere}$ \checkmark.\\
(c) $P_1=\dfrac{30^{2}}{20}=45$; $P_2=10(1{,}5)^{2}=22{,}5$; e
$20(0{,}75)^{2}=\SI{11.25}{\watt}$ em cada resistor do bloco:
\[
45+22{,}5+11{,}25+11{,}25=\SI{90}{\watt}=30(3)\ \checkmark
\]
\textbf{Observação:} a potência da fonte de corrente só pôde ser calculada
\emph{depois} de obtida a tensão. Com fonte de tensão ocorre o inverso --- é a
corrente que precisa ser encontrada primeiro.}
\end{questao}

\blocodesafio

\begin{questao}[4]{}
\textbf{(Projeto de divisor carregado.)} Deseja-se obter $\SI{5}{\volt}$ a
partir de $\SI{12}{\volt}$ para alimentar uma carga fixa de
$\SI{1}{\kilo\ohm}$.
\begin{itensq}
\item Deduza o critério de "corrente de sangria" e determine $R_1,R_2$ para
      erro inferior a $2\%$.
\item Projete o divisor \emph{compensado}, que entregue exatamente
      $\SI{5}{\volt}$ com a carga presente.
\item Compare as duas soluções em potência dissipada e em robustez.
\end{itensq}
\resposta{(a) $R_2\approx\SI{25}{\ohm}$, $R_1=\SI{35}{\ohm}$ (erro $1{,}4\%$);
(b) $R_1=\SI{140}{\ohm}$, $R_2=\SI{111}{\ohm}$;
(c) a compensada dissipa muito menos, mas só é exata para aquela carga}
\resolucao{(a) A \textbf{corrente de sangria} é a que circula por $R_2$ sem
passar pela carga. Se ela for muito maior que a da carga, conectar a carga
perturba pouco o divisor. Com $I_L=5/1000=\SI{5}{\milli\ampere}$, exigir
sangria $\ge20I_L=\SI{100}{\milli\ampere}$ dá
$R_2\le5/0{,}1=\SI{50}{\ohm}$. Testando valores (com
$R_1=1{,}4R_2$ para manter a razão $5/12$ em vazio):
\begin{center}
\begin{tabular}{cccc}
\hline
$R_2$ & $R_1$ & $U_{saída}$ & erro\\
\hline
$\SI{100}{\ohm}$ & $\SI{140}{\ohm}$ & $\SI{4.72}{\volt}$ & $-5{,}5\%$\\
$\SI{50}{\ohm}$  & $\SI{70}{\ohm}$  & $\SI{4.86}{\volt}$ & $-2{,}8\%$\\
$\SI{25}{\ohm}$  & $\SI{35}{\ohm}$  & $\SI{4.93}{\volt}$ & $-1{,}4\%$\\
\hline
\end{tabular}
\end{center}
Escolha $R_2=\SI{25}{\ohm}$, $R_1=\SI{35}{\ohm}$.\\
(b) \textbf{Compensação:} projete com a carga \emph{já incluída}. Exigindo
$12\cdot\dfrac{x}{R_1+x}=5$ com $x=R_2\parallel1000$, vem $x=\dfrac{5R_1}{7}$.
Tomando $R_1=\SI{140}{\ohm}$: $x=\SI{100}{\ohm}$ e
\[
\frac{1}{R_2}=\frac{1}{100}-\frac{1}{1000}
\;\Longrightarrow\;
R_2=\SI{111}{\ohm}.
\]
Verificação: $111\parallel1000=\SI{100}{\ohm}$ e
$12\cdot100/240=\SI{5.00}{\volt}$ \checkmark.\\
(c) \emph{Potência:} na solução (a), a corrente do divisor é
$\approx\SI{200}{\milli\ampere}$ e a dissipação total
$\approx\SI{2.4}{\watt}$; na (b), $\approx\SI{50}{\milli\ampere}$ e
$\approx\SI{0.6}{\watt}$ --- quatro vezes menos.\\
\emph{Robustez:} a solução (b) é exata \textbf{apenas} para
$R_L=\SI{1}{\kilo\ohm}$; se a carga for retirada, a saída sobe para
$12\cdot111/251=\SI{5.31}{\volt}$ ($+6\%$). A solução (a) é menos precisa,
porém quase indiferente à carga. \textbf{A escolha depende de a carga ser fixa
e conhecida ou variável} --- e, se for variável, nenhuma das duas serve:
o caso pede um regulador.}
\end{questao}

\begin{questao}[4]{}
\textbf{(Enumeração de topologias.)} Dispõe-se de quatro resistores iguais de
$\SI{100}{\ohm}$.
\begin{itensq}
\item Determine todos os valores distintos de $R_{eq}$ obteníveis usando os
      quatro.
\item Identifique os dois arranjos que resultam em $\SI{100}{\ohm}$ e
      compare a potência do componente mais solicitado sob
      $\SI{100}{\volt}$.
\item Justifique por que os extremos são $\SI{25}{\ohm}$ e
      $\SI{400}{\ohm}$.
\end{itensq}
\resposta{(a) $25$, $40$, $60$, $75$, $100$, $133{,}3$, $166{,}7$, $250$ e
$\SI{400}{\ohm}$; (b) $\SI{25}{\watt}$ em ambos}
\resolucao{(a) Enumerando as topologias possíveis:
\begin{center}
\begin{tabular}{ll}
\hline
Arranjo & $R_{eq}$\\
\hline
4 em paralelo & $\SI{25}{\ohm}$\\
$(2\text{ série})\parallel(2\text{ paralelo})$ & $\SI{40}{\ohm}$\\
$[(2\parallel)+1]\parallel 1$ & $\SI{60}{\ohm}$\\
$(3\text{ série})\parallel 1$ & $\SI{75}{\ohm}$\\
$(2\text{ série})\parallel(2\text{ série})$ & $\SI{100}{\ohm}$\\
$(2\parallel)+(2\parallel)$ & $\SI{100}{\ohm}$\\
$(3\parallel)+1$ & $\SI{133.3}{\ohm}$\\
$1+[(2\text{ série})\parallel 1]$ & $\SI{166.7}{\ohm}$\\
$(2\parallel)+2\text{ série}$ & $\SI{250}{\ohm}$\\
4 em série & $\SI{400}{\ohm}$\\
\hline
\end{tabular}
\end{center}
Nove valores distintos.\\
(b) Sob $\SI{100}{\volt}$, ambos os arranjos de $\SI{100}{\ohm}$ dão
$I_t=\SI{1}{\ampere}$ e $P_{tot}=\SI{100}{\watt}$.\\
\emph{Dois ramos de dois em série:} cada ramo conduz $\SI{0.5}{\ampere}$
$\Rightarrow$ $P=100(0{,}25)=\SI{25}{\watt}$ por resistor.\\
\emph{Dois blocos paralelos em série:} cada bloco sob $\SI{50}{\volt}$, cada
resistor com $\SI{0.5}{\ampere}$ $\Rightarrow$ também
$\SI{25}{\watt}$.\\
Os dois arranjos são \textbf{topologicamente distintos, mas termicamente
equivalentes} --- ambos distribuem a dissipação em quatro partes iguais. É o
melhor aproveitamento possível dos componentes.\\
(c) O mínimo é o paralelo total, pois cada resistor acrescentado em paralelo só
pode \emph{reduzir} $R_{eq}$; o máximo é a série total, pela razão simétrica.
Qualquer topologia mista fica necessariamente entre $R/n$ e $nR$.}
\end{questao}

\begin{questao}[4]{}
\textbf{(Projeto sob duas restrições.)} Uma rede deve apresentar
$R_{eq}=\SI{20}{\ohm}$ e dissipar $\SI{45}{\watt}$.
\begin{itensq}
\item Determine a tensão e a corrente nominais.
\item Compare três topologias quanto à potência do componente mais solicitado.
\item Escolha a melhor e justifique.
\end{itensq}
\resposta{(a) $\SI{30}{\volt}$ e $\SI{1.5}{\ampere}$; (b) pior componente:
$\SI{30}{\watt}$, $\SI{22.5}{\watt}$ e $\SI{11.25}{\watt}$;
(c) a matriz $2\times2$ de $\SI{20}{\ohm}$}
\resolucao{(a) De $P=V^{2}/R_{eq}$ e $P=RI^{2}$:
\[
V=\sqrt{PR_{eq}}=\sqrt{45\cdot20}=\SI{30}{\volt},
\qquad
I=\sqrt{P/R_{eq}}=\sqrt{2{,}25}=\SI{1.5}{\ampere}.
\]
(b) Três topologias que dão $\SI{20}{\ohm}$:
\begin{center}
\begin{tabular}{lll}
\hline
Topologia & Potências & Pior\\
\hline
$\SI{30}{\ohm}\parallel\SI{60}{\ohm}$ & $30$; $\SI{15}{\watt}$
  & $\SI{30}{\watt}$\\
$\SI{10}{\ohm}+(\SI{20}{\ohm}\parallel\SI{20}{\ohm})$
  & $22{,}5$; $11{,}25$; $\SI{11.25}{\watt}$ & $\SI{22.5}{\watt}$\\
$(\SI{20}{\ohm}+\SI{20}{\ohm})\parallel(\SI{20}{\ohm}+\SI{20}{\ohm})$
  & $11{,}25\times4$ & $\SI{11.25}{\watt}$\\
\hline
\end{tabular}
\end{center}
Em todas, o total é $\SI{45}{\watt}$ --- como tem de ser, pois $R_{eq}$ e $V$
são os mesmos.\\
(c) \textbf{A matriz $2\times2$} de resistores de $\SI{20}{\ohm}$. Ela reparte
a dissipação em quatro parcelas \emph{iguais} de $\SI{11.25}{\watt}$,
permitindo componentes de $\SI{25}{\watt}$ (fator $2$ de folga) e um único
valor de estoque. As alternativas exigiriam um componente de $\SI{50}{\watt}$
ou de $\SI{50}{\watt}$/$\SI{25}{\watt}$ mistos, com ponto quente concentrado.\\
\textbf{Princípio geral:} para uma dada $R_{eq}$ e potência total, a topologia
ótima é a que \textbf{equaliza} a dissipação --- matrizes $n\times n$ de
resistores idênticos são o caso limite dessa ideia.}
\end{questao}

\begin{questao}[4]{}
\textbf{(Diagnóstico quantitativo.)} Uma rede
$R_1=\SI{10}{\ohm}+(R_2=\SI{20}{\ohm}\parallel R_3=\SI{20}{\ohm})$ opera sob
$\SI{24}{\volt}$. Mediu-se $I_t=\SI{2.4}{\ampere}$.
\begin{itensq}
\item Mostre que a rede está defeituosa e identifique o tipo de falha.
\item Localize o componente e prove que a localização é única.
\item Proponha uma medição de confirmação.
\end{itensq}
\resposta{$R_{eq}$ medida $=\SI{10}{\ohm}$ contra $\SI{20}{\ohm}$ prevista;
o bloco paralelo está em curto}
\resolucao{(a) $R_{eq}^{med}=24/2{,}4=\SI{10}{\ohm}$, contra
$\SI{20}{\ohm}$ previstos. Resistência \emph{menor} que a nominal
$\Rightarrow$ falha do tipo \textbf{curto}. (Se fosse maior, seria abertura.)\\
(b) A diferença é de exatamente $\SI{10}{\ohm}$, que é o valor do bloco
paralelo --- ele está anulado. \emph{Unicidade:} $R_1$ não pode estar em curto,
pois isso daria $R_{eq}=\SI{10}{\ohm}$ também... \textbf{e aqui está o ponto:}
as duas hipóteses são numericamente idênticas! Um curto em $R_1$ deixaria
$R_{eq}=0+10=\SI{10}{\ohm}$, exatamente o medido.\\
Logo a medição de corrente \textbf{não} distingue as duas falhas --- e afirmar
localização única com esse único dado seria erro de método. O item (c) é
obrigatório, não opcional.\\
(c) \textbf{Medição decisiva:} um voltímetro sobre $R_1$.
\begin{itemize}
\item Se o bloco paralelo está em curto: $U_1=10(2{,}4)=\SI{24}{\volt}$
      (toda a fonte).
\item Se $R_1$ está em curto: $U_1=\SI{0}{\volt}$, e os
      $\SI{24}{\volt}$ aparecem sobre o bloco paralelo.
\end{itemize}
\textbf{Lição de diagnóstico:} uma única medição escalar raramente identifica
uma falha em rede mista, porque topologias diferentes podem produzir a mesma
$R_{eq}$. É preciso pelo menos uma medida de \emph{distribuição} --- uma
tensão interna --- para quebrar a ambiguidade.}
\end{questao}

\begin{questao}[4]{}
\textbf{(Redução por simetria.)} Seis resistores de $\SI{6}{\ohm}$ formam as
arestas de um \textbf{tetraedro} regular (quatro vértices, todos ligados dois a
dois).
\begin{itensq}
\item Determine $R_{AB}$ entre dois vértices quaisquer, usando simetria.
\item Determine a corrente em cada aresta para $I_t=\SI{4}{\ampere}$.
\item Recalcule $R_{AB}$ se a aresta $AB$ for removida.
\end{itensq}
\resposta{(a) $R_{AB}=\SI{3}{\ohm}$; (c) $\SI{6}{\ohm}$}
\resolucao{(a) Sejam $A$ e $B$ os terminais e $C$, $D$ os outros dois vértices.
Como $C$ e $D$ ocupam posições \textbf{simétricas} em relação ao par $A$--$B$
(cada um ligado a $A$, a $B$ e ao outro, com resistências iguais), tem-se
$V_C=V_D$. A aresta $CD$ não conduz e pode ser \textbf{removida}.\\
Restam três caminhos em paralelo: a aresta direta $AB$ ($\SI{6}{\ohm}$), o
percurso $A$--$C$--$B$ ($\SI{12}{\ohm}$) e o percurso $A$--$D$--$B$
($\SI{12}{\ohm}$):
\[
R_{AB}=6\parallel12\parallel12=6\parallel6=\SI{3}{\ohm}=\frac{R}{2}.
\]
(b) Com $I_t=\SI{4}{\ampere}$, $U_{AB}=3(4)=\SI{12}{\volt}$:
\begin{itemize}
\item aresta $AB$: $12/6=\SI{2}{\ampere}$;
\item cada percurso lateral: $12/12=\SI{1}{\ampere}$;
\item aresta $CD$: $\SI{0}{\ampere}$.
\end{itemize}
Soma: $2+1+1=\SI{4}{\ampere}$ \checkmark. A aresta direta leva metade da
corrente total.\\
(c) Sem a aresta $AB$, a simetria $C\leftrightarrow D$ \emph{persiste} --- ela
não dependia daquela aresta. Logo $V_C=V_D$ ainda, e $CD$ continua sem
conduzir:
\[
R_{AB}=12\parallel12=\SI{6}{\ohm}=R.
\]
\textbf{Resultado elegante:} remover uma aresta dobra a resistência
($3\to\SI{6}{\ohm}$), porque justamente essa aresta carregava metade da
corrente.\\
\textbf{Por que a simetria funciona:} a solução de uma rede resistiva é
\emph{única}; se permutar $C$ com $D$ devolve o mesmo circuito, a solução
permutada é a mesma solução --- e isso exige $V_C=V_D$. O argumento é o mesmo
que resolve a ponte equilibrada e o cubo de resistores, e dispensa qualquer
sistema de equações.}
\end{questao}

% END BANCO COMPLEMENTAR

\gabarito[3.3 Circuitos mistos]

% =====================================================================
\subsection{Transformação estrela-triângulo}
% =====================================================================

\begin{conceito}[Quando usar]
Use a transformação $\Delta\leftrightarrow Y$ quando o circuito \textbf{não}
puder ser reduzido apenas por série e paralelo --- o caso típico é a
\emph{ponte de Wheatstone} desequilibrada. Converta um dos triângulos (ou uma das
estrelas) para "abrir" o circuito e revelar novas associações série/paralelo.
\end{conceito}

\legendaniveis

\blocofixacao

\begin{questao}[1]{}
Considere uma configuração \textbf{estrela} com $R_1=\SI{3}{\ohm}$,
$R_2=\SI{6}{\ohm}$ e $R_3=\SI{9}{\ohm}$. Determine os resistores da configuração
\textbf{triângulo} equivalente.
\resposta{$R_{AB}=\SI{11}{\ohm}$, $R_{BC}=\SI{33}{\ohm}$, $R_{CA}=\SI{16.5}{\ohm}$}
\resolucao{Com $P = R_1R_2+R_2R_3+R_3R_1 = 18+54+27 = \SI{99}{\ohm\squared}$, cada
resistor do triângulo é $P$ dividido pelo resistor \emph{oposto} da estrela:
\[
R_{AB}=\frac{P}{R_3}=\frac{99}{9}=\SI{11}{\ohm},\quad
R_{BC}=\frac{P}{R_1}=\frac{99}{3}=\SI{33}{\ohm},\quad
R_{CA}=\frac{P}{R_2}=\frac{99}{6}=\SI{16.5}{\ohm}.
\]}
\end{questao}

\begin{questao}[1]{}
Um circuito \textbf{triângulo} possui $R_{AB}=\SI{10}{\ohm}$,
$R_{BC}=\SI{5}{\ohm}$ e $R_{CA}=\SI{15}{\ohm}$. Calcule os resistores da
\textbf{estrela} equivalente.
\resposta{$R_A=\SI{5}{\ohm}$, $R_B=\SI{1.67}{\ohm}$, $R_C=\SI{2.5}{\ohm}$}
\resolucao{Com $S = R_{AB}+R_{BC}+R_{CA} = 10+5+15 = \SI{30}{\ohm}$, cada resistor
da estrela é o produto dos dois lados \emph{adjacentes} ao seu vértice, dividido
por $S$:
\[
R_A=\frac{R_{AB}R_{CA}}{S}=\frac{10\cdot15}{30}=\SI{5}{\ohm},\quad
R_B=\frac{R_{AB}R_{BC}}{S}=\frac{10\cdot5}{30}\approx\SI{1.67}{\ohm},\quad
R_C=\frac{R_{BC}R_{CA}}{S}=\frac{5\cdot15}{30}=\SI{2.5}{\ohm}.
\]}
\end{questao}

\begin{questao}[1]{}
Explique, com suas palavras, por que a transformação estrela-triângulo é útil e
cite um exemplo prático de circuito em que ela é necessária.
\resposta{Resposta dissertativa --- ver comentário}
\resolucao{Ela permite substituir um agrupamento de três resistores em triângulo
por três em estrela (ou vice-versa) \emph{sem alterar} o comportamento elétrico
visto pelos três terminais externos. Isso "desmancha" configurações que não são
nem série nem paralelo, revelando novas associações simples. O exemplo clássico é
a \textbf{ponte de Wheatstone desequilibrada}: convertendo um dos triângulos em
estrela, os ramos passam a ficar em série/paralelo e o circuito se resolve sem
precisar montar um sistema de equações. Outra aplicação é a análise de redes
trifásicas, onde cargas ligadas em $\Delta$ e em $Y$ são convertidas entre si.}
\end{questao}

\blocoaplicacao

\begin{questao}[2]{}
No triângulo da figura, com $R_{AB}=\SI{9}{\ohm}$ (lado direto entre A e B),
$R_{AC}=\SI{6}{\ohm}$ e $R_{CB}=\SI{3}{\ohm}$, determine a resistência
equivalente entre A e B.

\circuitoq{%
\begin{circuitikz}[scale=1,transform shape,line width=0.8pt]
 \draw (0,0) node[below left,font=\bfseries]{A} coordinate(A)
       to[R,l=$\SI{6}{\ohm}$] (2,2.4) node[above,font=\bfseries]{C} coordinate(C)
       to[R,l=$\SI{3}{\ohm}$] (4,0) node[below right,font=\bfseries]{B} coordinate(B);
 \draw (A) to[R,l_=$\SI{9}{\ohm}$] (B);
\end{circuitikz}}

\resposta{$R_{AB} = \SI{4.5}{\ohm}$}
\resolucao{O caminho superior A--C--B está em \emph{série}:
$6+3 = \SI{9}{\ohm}$. Esse ramo está em \emph{paralelo} com o lado direto
$R_{AB} = \SI{9}{\ohm}$:
\[
R_{AB}=9\parallel 9=\frac{9}{2}=\SI{4.5}{\ohm}.
\]
Aqui a redução série/paralelo bastou. Guarde a transformação $\Delta$-$Y$ para
quando há um resistor \emph{atravessando} o triângulo (a ponte), como nas
questões seguintes.}
\end{questao}

\begin{questao}[2]{}
Três resistores iguais de $\SI{2}{\ohm}$ estão ligados em \textbf{estrela}.
Convertendo-os em \textbf{triângulo} e aplicando uma fonte de $\SI{12}{\volt}$
entre dois dos terminais (A e B), determine a corrente total fornecida pela fonte.
\resposta{$I = \SI{3}{\ampere}$}
\resolucao{A estrela de três resistores iguais $R$ vira um triângulo de
resistores $3R = \SI{6}{\ohm}$ cada. Entre A e B temos o lado direto de
\SI{6}{\ohm} em paralelo com o caminho A--C--B ($6+6 = \SI{12}{\ohm}$):
\[
R_{AB}=6\parallel 12=\frac{6\cdot12}{18}=\SI{4}{\ohm}
\quad\Rightarrow\quad
I=\frac{12}{4}=\SI{3}{\ampere}.
\]
\textbf{Confira:} entre A e B, a estrela original também dá o mesmo resultado ---
$R_A + (R_B \parallel$ com o caminho por C$)$ ---, como esperado, já que as duas
representações são equivalentes vistas dos terminais.}
\end{questao}

\begin{questao}[2]{}
Um triângulo entre A, B e C tem $R_{AB}=\SI{12}{\ohm}$, $R_{BC}=\SI{6}{\ohm}$ e
$R_{CA}=\SI{6}{\ohm}$.
\begin{itensq}
\item Converta para estrela.
\item Calcule a resistência equivalente entre A e B após a conversão.
\end{itensq}
\resposta{(a) $R_A=\SI{3}{\ohm}$, $R_B=\SI{3}{\ohm}$, $R_C=\SI{1.5}{\ohm}$;
(b) $R_{AB}=\SI{6}{\ohm}$}
\resolucao{(a) $S = 12+6+6 = \SI{24}{\ohm}$:
\[
R_A=\frac{R_{AB}R_{CA}}{S}=\frac{12\cdot6}{24}=\SI{3}{\ohm},\quad
R_B=\frac{R_{AB}R_{BC}}{S}=\frac{12\cdot6}{24}=\SI{3}{\ohm},\quad
R_C=\frac{R_{BC}R_{CA}}{S}=\frac{6\cdot6}{24}=\SI{1.5}{\ohm}.
\]
(b) Na estrela, entre A e B a corrente percorre $R_A$ e $R_B$ em série (o braço
$R_C$ fica em aberto, pois C não é terminal):
$R_{AB} = R_A + R_B = 3+3 = \SI{6}{\ohm}$.}
\end{questao}

\blocoaprofundamento

\begin{questao}[3]{}
A figura mostra uma \textbf{ponte de Wheatstone}. Verifique se ela está
equilibrada e, caso esteja, determine a resistência equivalente entre A e B.

\circuitoq{%
\begin{circuitikz}[scale=1,transform shape,line width=0.8pt]
 \draw (0,0) node[left,font=\bfseries]{A} coordinate(A);
 \draw (5,0) node[right,font=\bfseries]{B} coordinate(B);
 \coordinate (C) at (2.5,1.7);
 \coordinate (D) at (2.5,-1.7);
 \draw (A) to[R,l=$\SI{10}{\ohm}$] (C);
 \draw (C) to[R,l=$\SI{20}{\ohm}$] (B);
 \draw (A) to[R,l_=$\SI{15}{\ohm}$] (D);
 \draw (D) to[R,l_=$\SI{30}{\ohm}$] (B);
 \draw (C) to[R,l=$\SI{8}{\ohm}$] (D);
 \node[font=\footnotesize] at (3.1,-0.27) {galv.};
\end{circuitikz}}

\resposta{Equilibrada; $R_{AB} = \SI{18}{\ohm}$}
\resolucao{A ponte está \textbf{equilibrada} quando os produtos dos braços opostos
são iguais:
\[
\frac{R_{AC}}{R_{CB}}=\frac{10}{20}=\frac12,
\qquad
\frac{R_{AD}}{R_{DB}}=\frac{15}{30}=\frac12.
\]
As razões são iguais, logo C e D estão no \emph{mesmo potencial}: o resistor
central de \SI{8}{\ohm} não conduz e pode ser removido. Restam dois ramos em série,
associados em paralelo:
\[
R_{AB}=(10+20)\parallel(15+30)=30\parallel45=\frac{30\cdot45}{75}=\SI{18}{\ohm}.
\]}
\end{questao}

\begin{questao}[3]{}
Na ponte da figura, $R_{AC}=\SI{6}{\ohm}$, $R_{CB}=\SI{18}{\ohm}$,
$R_{AD}=\SI{18}{\ohm}$, $R_{DB}=\SI{6}{\ohm}$ e o resistor central (ponte) vale
$R_{CD}=\SI{6}{\ohm}$. Determine a resistência equivalente entre A e B.

\circuitoq{%
\begin{circuitikz}[scale=1,transform shape,line width=0.8pt]
 \draw (0,0) node[left,font=\bfseries]{A} coordinate(A);
 \draw (5,0) node[right,font=\bfseries]{B} coordinate(B);
 \coordinate (C) at (2.5,1.7);
 \coordinate (D) at (2.5,-1.7);
 \draw (A) to[R,l=$\SI{6}{\ohm}$] (C);
 \draw (C) to[R,l=$\SI{18}{\ohm}$] (B);
 \draw (A) to[R,l_=$\SI{18}{\ohm}$] (D);
 \draw (D) to[R,l_=$\SI{6}{\ohm}$] (B);
 \draw (C) to[R,l=$\SI{6}{\ohm}$] (D);
\end{circuitikz}}

\resposta{$R_{AB} = \SI{10}{\ohm}$}
\resolucao{Como $\dfrac{6}{18}\neq\dfrac{18}{6}$, a ponte está
\textbf{desequilibrada} e o resistor central conduz: não dá para reduzir só por
série/paralelo. Convertemos o triângulo superior A--C--D
($R_{AC}=6$, $R_{CD}=6$, $R_{AD}=18$) em estrela, com $S = 6+6+18 = \SI{30}{\ohm}$:
\[
R_{A}=\frac{6\cdot18}{30}=\SI{3.6}{\ohm},\quad
R_{C}=\frac{6\cdot6}{30}=\SI{1.2}{\ohm},\quad
R_{D}=\frac{6\cdot18}{30}=\SI{3.6}{\ohm}.
\]
Chamando de N o centro da estrela, os caminhos até B ficam:
$N$--$C$--$B = 1{,}2+18 = \SI{19.2}{\ohm}$ e
$N$--$D$--$B = 3{,}6+6 = \SI{9.6}{\ohm}$, em paralelo:
\[
19{,}2\parallel 9{,}6 = \frac{19{,}2\cdot9{,}6}{28{,}8}=\SI{6.4}{\ohm}.
\]
Somando o braço $R_A$ que liga A ao centro:
\[
R_{AB}=3{,}6+6{,}4=\SI{10}{\ohm}.
\]}
\end{questao}

\begin{questao}[3]{}
Três resistores formam um triângulo entre os nós X, Y e Z, com
$R_{XY}=\SI{9}{\ohm}$, $R_{YZ}=\SI{12}{\ohm}$ e $R_{ZX}=\SI{15}{\ohm}$. Além
disso, cada nó é ligado ao terra por um resistor de $\SI{6}{\ohm}$. Descreva o
procedimento para obter a resistência entre o nó X e o terra e monte a expressão
final (sem exigir o valor numérico fechado).
\resposta{Ver roteiro comentado}
\resolucao{\textbf{Passo 1.} Converta o triângulo XYZ em estrela, com centro N e
$S = 9+12+15 = \SI{36}{\ohm}$:
\[
R_X=\frac{R_{XY}R_{ZX}}{S}=\frac{9\cdot15}{36}=\SI{3.75}{\ohm},\quad
R_Y=\frac{R_{XY}R_{YZ}}{S}=\frac{9\cdot12}{36}=\SI{3}{\ohm},\quad
R_Z=\frac{R_{YZ}R_{ZX}}{S}=\frac{12\cdot15}{36}=\SI{5}{\ohm}.
\]
\textbf{Passo 2.} Agora cada terminal externo Y e Z leva ao terra por
$R_Y+6$ e $R_Z+6$, ambos em paralelo (pois se encontram no terra):
$(3+6)\parallel(5+6) = 9\parallel 11$. \textbf{Passo 3.} Esse resultado fica em
série com $R_X$, e todo o conjunto em paralelo com o resistor de \SI{6}{\ohm} que
liga X diretamente ao terra:
\[
R_{X\to\text{terra}}
=\Big[\,R_X+\big(9\parallel11\big)\Big]\ \big\|\ 6
=\Big[\,3{,}75+\tfrac{99}{20}\Big]\ \big\|\ 6 .
\]
O que se pretende com esta questão é o \emph{método}: transformar o triângulo em
estrela desmonta o emaranhado e transforma tudo em série/paralelo.}
\end{questao}

\begin{questao}[3]{}
Uma ponte de Wheatstone é usada para medir uma resistência desconhecida $R_x$.
Os braços conhecidos são $R_1 = \SI{100}{\ohm}$, $R_2 = \SI{200}{\ohm}$ e
$R_3$ é uma década ajustável. O equilíbrio é obtido com
$R_3 = \SI{347}{\ohm}$.
\begin{itensq}
\item Determine $R_x$.
\item Se $R_1$ e $R_2$ têm tolerância de $\SI{0.1}{\percent}$ cada e $R_3$ de
      $\SI{0.05}{\percent}$, estime a incerteza relativa de $R_x$.
\item Por que a ponte é mais precisa que medir $R_x$ com um ohmímetro comum?
\end{itensq}
\resposta{(a) $R_x = \SI{694}{\ohm}$; (b) $\approx\SI{0.25}{\percent}$;
(c) ver comentário}
\resolucao{(a) No equilíbrio, os produtos cruzados se igualam:
\[
\frac{R_1}{R_2}=\frac{R_3}{R_x}
\;\Rightarrow\;
R_x = R_3\,\frac{R_2}{R_1} = 347\cdot\frac{200}{100} = \SI{694}{\ohm}.
\]
(b) Para um produto/quociente, as incertezas \emph{relativas} se somam (no pior
caso):
\[
\frac{\delta R_x}{R_x} = \frac{\delta R_3}{R_3}+\frac{\delta R_2}{R_2}
+\frac{\delta R_1}{R_1} = 0{,}05\% + 0{,}1\% + 0{,}1\% = \SI{0.25}{\percent},
\]
ou seja, $R_x = \SI{694}{\ohm} \pm \SI{1.7}{\ohm}$.
(Somando em quadratura, para erros independentes, obtém-se
$\sqrt{0{,}05^2+0{,}1^2+0{,}1^2}\approx\SI{0.15}{\percent}$ --- estimativa mais
realista.)\\
(c) Porque a ponte é um método de \textbf{zero} (ou de detecção nula): no
equilíbrio, o galvanômetro indica \emph{zero}, e a medida não depende da precisão
do próprio galvanômetro nem da tensão da fonte --- depende apenas das
\emph{razões} entre resistores conhecidos. Um ohmímetro comum, ao contrário,
depende da estabilidade da sua fonte interna e da calibração do seu indicador.
\textbf{Este é um princípio geral da metrologia:} comparar contra um padrão é
sempre mais preciso do que medir em escala absoluta.}
\end{questao}

% BEGIN BANCO COMPLEMENTAR Transformação estrela-triângulo
% =====================================================================
%  Banco complementar --- Transformacao Estrela-Triangulo  (REVISADO)
%  Substitui a versao gerada por template (9 clones em N2, 8 em N3 e
%  8 questoes conceituais marcadas como N4).
%  O cap03 ja cobre: estrela 3/6/9 -> triangulo, triangulo 10/5/15 ->
%  estrela, "explique por que a transformacao e util", triangulo
%  9/6/3 com figura, tres de 2 ohm em estrela com fonte, triangulo
%  12/6/6 com itens, e o triangulo XYZ 9/12/15. As conversoes basicas,
%  portanto, ja estao exaustivamente cobertas la.
%  Este banco explora: reversibilidade e verificacao, casos-limite
%  (lado nulo, lado infinito), ponte desequilibrada por dois caminhos,
%  condicionamento numerico, valores comerciais, e as DEDUCOES das
%  duas formulas a partir das resistencias entre terminais.
%  Distribuicao mantida: 5 N1 + 9 N2 + 8 N3 + 8 N4 = 30
% =====================================================================

\subsubsection*{Banco complementar --- Transformação estrela-triângulo}

\begin{atencao}[Organização do banco]
As duas fórmulas têm estruturas espelhadas. \textbf{Triângulo para estrela:}
cada braço é o \emph{produto} dos dois lados que tocam aquele terminal, dividido
pela \emph{soma} dos três lados. \textbf{Estrela para triângulo:} cada lado é a
soma dos três produtos de braços, dividida pelo braço \emph{oposto}. O N3 trata
de pontes, casos-limite e condicionamento; o N4 exige as deduções, provas de
positividade e a análise de valores comerciais.
\end{atencao}

\blocofixacao

\begin{questao}[1]{}
Um triângulo equilátero tem três lados de $\SI{30}{\ohm}$. Determine a estrela
equivalente.
\resposta{Três braços de $\SI{10}{\ohm}$}
\resolucao{Para o caso equilibrado, a fórmula geral colapsa:
\[
R_Y=\frac{R_\Delta\cdot R_\Delta}{3R_\Delta}=\frac{R_\Delta}{3}
=\frac{30}{3}=\SI{10}{\ohm}.
\]
\textbf{Regra rápida:} $R_Y=R_\Delta/3$. Guarde o sentido --- a estrela é
sempre \emph{menor}, porque seus dois braços em série ($2R_Y=\SI{20}{\ohm}$)
devem igualar o paralelo do triângulo
($30\parallel60=\SI{20}{\ohm}$) \checkmark.}
\end{questao}

\begin{questao}[1]{}
Uma estrela equilibrada tem três braços de $\SI{5}{\ohm}$. Determine o
triângulo equivalente.
\resposta{Três lados de $\SI{15}{\ohm}$}
\resolucao{
\[
R_\Delta=\frac{3R_Y^{2}}{R_Y}=3R_Y=3(5)=\SI{15}{\ohm}.
\]
\textbf{Regra rápida:} $R_\Delta=3R_Y$ --- inversa da anterior. Se
$R_Y=R_\Delta/3$ e $R_\Delta=3R_Y$ não fossem consistentes, uma das fórmulas
estaria trocada. É a conferência mais barata que existe para o caso
equilibrado.}
\end{questao}

\begin{questao}[1]{}
Escreva, \textbf{sem calcular}, a expressão do braço $R_B$ de uma estrela
equivalente a um triângulo de lados $R_{AB}$, $R_{BC}$ e $R_{CA}$.
\resposta{$R_B=\dfrac{R_{AB}\,R_{BC}}{R_{AB}+R_{BC}+R_{CA}}$}
\resolucao{A regra é posicional: o braço ligado ao terminal $B$ é o
\textbf{produto dos dois lados que tocam $B$} --- que são $R_{AB}$ e
$R_{BC}$ --- dividido pela \textbf{soma dos três lados}.\\
\textbf{Como não errar:} identifique o terminal, olhe quais dois lados chegam
nele e multiplique. O lado \emph{oposto} ($R_{CA}$) aparece apenas no
denominador, junto com os outros dois.}
\end{questao}

\begin{questao}[1]{}
Compare estrela e triângulo quanto ao número de elementos e de nós.
\begin{alternativas}[2]
\item ambos têm 3 resistores; a estrela tem um nó interno adicional
\item a estrela tem 3 e o triângulo tem 4 resistores
\item ambos têm o mesmo número de nós
\item o triângulo tem um nó interno adicional
\end{alternativas}
\resposta{(a)}
\resolucao{Ambas as configurações usam \textbf{três} resistores e possuem os
mesmos três terminais externos ($A$, $B$, $C$). A diferença é o \textbf{nó
central} da estrela, que não tem correspondente no triângulo.\\
\textbf{Consequência importante:} a equivalência vale apenas para o que se
observa nos \emph{três terminais externos}. O potencial do nó central da estrela
não corresponde a ponto algum do triângulo --- e nenhuma corrente interna é
preservada pela transformação.}
\end{questao}

\begin{questao}[1]{}
Identifique qual das expressões é a conversão triângulo-estrela:
\begin{alternativas}[2]
\item $R=\dfrac{\text{produto de dois}}{\text{soma dos três}}$
\item $R=\dfrac{\text{soma dos produtos}}{\text{oposto}}$
\item $R=\dfrac{\text{soma dos três}}{3}$
\item $R=\dfrac{\text{produto dos três}}{\text{soma dos três}}$
\end{alternativas}
\resposta{(a)}
\resolucao{\textbf{Triângulo $\to$ estrela:} produto de dois sobre a soma dos
três --- alternativa (a). O resultado tem dimensão de $\si{\ohm}$
($\si{\ohm^2}/\si{\ohm}$) \checkmark.\\
\textbf{Estrela $\to$ triângulo:} soma dos produtos sobre o braço oposto ---
alternativa (b). Também $\si{\ohm}$ \checkmark.\\
\textbf{Teste de sanidade:} a alternativa (d) daria $\si{\ohm^{2}}$ --- errada
por análise dimensional. E, no caso equilibrado, (a) devolve $R/3$ e (b)
devolve $3R$, o que confirma qual é qual.}
\end{questao}

\blocoaplicacao

\begin{questao}[2]{}
Um triângulo tem $R_{AB}=\SI{20}{\ohm}$, $R_{BC}=\SI{30}{\ohm}$ e
$R_{CA}=\SI{50}{\ohm}$. Determine a estrela equivalente.
\resposta{$R_A=\SI{10}{\ohm}$, $R_B=\SI{6}{\ohm}$, $R_C=\SI{15}{\ohm}$}
\resolucao{A soma dos lados é $20+30+50=\SI{100}{\ohm}$.
\[
R_A=\frac{R_{AB}R_{CA}}{100}=\frac{20(50)}{100}=\SI{10}{\ohm};
\]
\[
R_B=\frac{R_{AB}R_{BC}}{100}=\frac{20(30)}{100}=\SI{6}{\ohm};
\qquad
R_C=\frac{R_{BC}R_{CA}}{100}=\frac{30(50)}{100}=\SI{15}{\ohm}.
\]
\textbf{Conferência de plausibilidade:} cada braço é menor que os lados que o
originaram, e o maior braço ($R_C$) está no terminal tocado pelos dois maiores
lados. Coerente.}
\end{questao}

\begin{questao}[2]{}
Converta de volta a estrela obtida ($10$, $6$ e $\SI{15}{\ohm}$) para
triângulo e verifique a reversibilidade.
\resposta{$20$, $30$ e $\SI{50}{\ohm}$ --- os valores originais}
\resolucao{Soma dos produtos:
\[
\Sigma=R_AR_B+R_BR_C+R_CR_A=60+90+150=\SI{300}{\ohm\squared}.
\]
Cada lado é $\Sigma$ dividido pelo braço \emph{oposto}:
\[
R_{AB}=\frac{300}{R_C}=\frac{300}{15}=\SI{20}{\ohm};
\quad
R_{BC}=\frac{300}{R_A}=\SI{30}{\ohm};
\quad
R_{CA}=\frac{300}{R_B}=\SI{50}{\ohm}.
\]
\textbf{Reversibilidade confirmada.} As duas transformações são inversas exatas
--- fato garantido pela construção, já que ambas derivam do mesmo sistema de
três equações. Fazer a ida e a volta é a verificação mais completa possível de
uma conversão.}
\end{questao}

\begin{questao}[2]{}
Para o triângulo $20/30/\SI{50}{\ohm}$ e sua estrela equivalente, calcule a
resistência entre $A$ e $B$ (com $C$ aberto) nas duas configurações.
\resposta{$\SI{16}{\ohm}$ em ambas}
\resolucao{\textbf{No triângulo:} o lado direto $R_{AB}$ em paralelo com o
caminho $A$--$C$--$B$:
\[
R_{AB}^{eq}=20\parallel(50+30)=\frac{20(80)}{100}=\SI{16}{\ohm}.
\]
\textbf{Na estrela:} com $C$ aberto, o braço $R_C$ não conduz; sobram
$R_A$ e $R_B$ em série:
\[
R_{AB}^{eq}=10+6=\SI{16}{\ohm}.\ \checkmark
\]
\textbf{É exatamente esta a definição de equivalência:} a resistência vista
entre cada par de terminais, com o terceiro aberto, deve coincidir. As três
igualdades desse tipo constituem o sistema de onde as fórmulas são deduzidas ---
como se mostra no bloco de desafio.}
\end{questao}

\begin{questao}[2]{}
Uma estrela tem $R_A=\SI{4}{\ohm}$, $R_B=\SI{8}{\ohm}$ e
$R_C=\SI{12}{\ohm}$. Determine o triângulo equivalente e identifique o maior
lado.
\resposta{$R_{AB}=\SI{14.67}{\ohm}$, $R_{BC}=\SI{44}{\ohm}$,
$R_{CA}=\SI{22}{\ohm}$; o maior é $R_{BC}$}
\resolucao{
\[
\Sigma=4(8)+8(12)+12(4)=32+96+48=\SI{176}{\ohm\squared}.
\]
\[
R_{AB}=\frac{176}{12}=\SI{14.67}{\ohm};
\qquad
R_{BC}=\frac{176}{4}=\SI{44}{\ohm};
\qquad
R_{CA}=\frac{176}{8}=\SI{22}{\ohm}.
\]
\textbf{Padrão:} o maior lado do triângulo é o \emph{oposto} ao menor braço da
estrela. Como $R_A=\SI{4}{\ohm}$ é o menor braço, $R_{BC}$ é o maior lado ---
a divisão por um número pequeno produz um resultado grande.\\
\textbf{Interpretação:} um braço pequeno "aproxima" seu terminal do centro,
o que equivale, no triângulo, a afastar os outros dois entre si.}
\end{questao}

\begin{questao}[2]{}
Em um triângulo, o lado $R_{BC}$ é substituído por um \textbf{circuito aberto}.
Determine a estrela equivalente e interprete.
\resposta{$R_A=\dfrac{R_{AB}R_{CA}}{R_{AB}+R_{CA}}$, $R_B=R_C=0$ --- ou seja,
apenas dois resistores em série}
\resolucao{Com $R_{BC}\to\infty$, a soma dos lados também diverge. Nos numeradores:
\[
R_A=\frac{R_{AB}R_{CA}}{R_{AB}+R_{BC}+R_{CA}}\to0
\]
--- mas os braços $R_B$ e $R_C$ contêm $R_{BC}$ no numerador e o limite exige
cuidado:
\[
R_B=\frac{R_{AB}R_{BC}}{R_{AB}+R_{BC}+R_{CA}}\to R_{AB};
\qquad
R_C=\frac{R_{BC}R_{CA}}{R_{AB}+R_{BC}+R_{CA}}\to R_{CA}.
\]
\textbf{Resultado:} $R_A=0$, $R_B=R_{AB}$ e $R_C=R_{CA}$. A estrela degenera em
$R_{AB}$ e $R_{CA}$ ligados diretamente em série pelo terminal $A$ --- que é
exatamente o circuito restante. \checkmark\\
\textbf{Lição:} os casos-limite não quebram as fórmulas; eles as confirmam. Um
lado ausente do triângulo apenas devolve a associação série trivial.}
\end{questao}

\begin{questao}[2]{}
Em um triângulo, o lado $R_{AB}$ é substituído por um \textbf{curto-circuito}.
Determine a estrela equivalente.
\resposta{$R_A=R_B=0$ e $R_C=\dfrac{R_{BC}R_{CA}}{R_{BC}+R_{CA}}$}
\resolucao{Com $R_{AB}=0$, a soma vale $R_{BC}+R_{CA}$ e:
\[
R_A=\frac{0\cdot R_{CA}}{\Sigma}=0;
\qquad
R_B=\frac{0\cdot R_{BC}}{\Sigma}=0;
\qquad
R_C=\frac{R_{BC}R_{CA}}{R_{BC}+R_{CA}}.
\]
\textbf{Interpretação:} os terminais $A$ e $B$ estão em curto, logo ambos
coincidem com o nó central --- daí $R_A=R_B=0$. E $R_C$ é o paralelo de
$R_{BC}$ com $R_{CA}$, o que é correto: com $A$ e $B$ unidos, os dois resistores
ficam mesmo em paralelo entre $C$ e esse nó comum. \checkmark\\
\textbf{Observação:} a conversão \emph{inversa} não existe aqui --- uma estrela
com dois braços nulos daria $\Sigma=0$ e divisão por zero. A transformação
$Y\to\Delta$ exige todos os braços estritamente positivos.}
\end{questao}

\begin{questao}[2]{}
Verifique numericamente, para o triângulo $20/30/\SI{50}{\ohm}$, a identidade
$R_{AB}R_{BC}+R_{BC}R_{CA}+R_{CA}R_{AB}
=\left(\sum R_\Delta\right)\left(\sum R_Y\right)$.
\resposta{$3100=100\times31$ \checkmark}
\resolucao{\textbf{Lado esquerdo:}
\[
20(30)+30(50)+50(20)=600+1500+1000=3100.
\]
\textbf{Lado direito:} $\sum R_\Delta=100$ e
$\sum R_Y=10+6+15=31$, logo $100(31)=3100$ \checkmark.\\
\textbf{Utilidade:} a identidade permite obter a soma dos braços da estrela
\emph{sem calcular cada um}:
\[
\sum R_Y=\frac{\sum_{\text{pares}}R_\Delta R_\Delta}{\sum R_\Delta}
=\frac{3100}{100}=\SI{31}{\ohm}.
\]
É uma conferência rápida do conjunto das três conversões --- se a soma não
bater, ao menos um braço está errado.}
\end{questao}

\begin{questao}[2]{}
Em uma rede, três resistores formam um triângulo entre os nós $P$, $Q$ e $R$,
e há mais resistores ligados a esses três nós. Explique por que convém
converter o triângulo em estrela, e não o contrário.
\resposta{A estrela cria um nó interno isolado, liberando associações
série/paralelo}
\resolucao{\textbf{O problema do triângulo:} cada um de seus lados liga dois nós
que \emph{também} recebem outros elementos. Nenhum lado fica em série ou em
paralelo com nada --- a redução trava.\\
\textbf{O que a estrela resolve:} ela substitui os três lados por três braços
que se encontram em um nó \textbf{novo}, ao qual nada mais está ligado. Cada
braço passa a estar em série com o que já existia em seu terminal, e essas
séries podem então ser combinadas em paralelo.\\
\textbf{Quando fazer o contrário:} se a rede contiver uma \emph{estrela} cujo
nó central recebe outros elementos, converta-a em triângulo --- isso elimina o
nó problemático. \textbf{A regra é sempre a mesma:} transforme na direção que
\emph{elimina} o nó que está impedindo a redução.}
\end{questao}

\begin{questao}[2]{}
Uma estrela tem braços $R_A=R_B=\SI{6}{\ohm}$ e $R_C=\SI{3}{\ohm}$. Determine o
triângulo equivalente e verifique pela resistência entre $A$ e $B$.
\resposta{$R_{AB}=\SI{24}{\ohm}$, $R_{BC}=R_{CA}=\SI{12}{\ohm}$}
\resolucao{
\[
\Sigma=6(6)+6(3)+3(6)=36+18+18=\SI{72}{\ohm\squared};
\]
\[
R_{AB}=\frac{72}{3}=\SI{24}{\ohm};
\qquad
R_{BC}=R_{CA}=\frac{72}{6}=\SI{12}{\ohm}.
\]
\textbf{Verificação (entre $A$ e $B$, com $C$ aberto):}\\
\emph{Estrela:} $R_A+R_B=\SI{12}{\ohm}$.\\
\emph{Triângulo:} $24\parallel(12+12)=\dfrac{24(24)}{48}=\SI{12}{\ohm}$
\checkmark.\\
\textbf{Simetria preservada:} como $R_A=R_B$, o triângulo também sai simétrico
em relação ao par $A$--$B$. Qualquer simetria da estrela reaparece no
triângulo --- propriedade útil para conferir resultados rapidamente.}
\end{questao}

\blocoaprofundamento

\begin{questao}[3]{}
Uma ponte é alimentada por $\SI{84}{\volt}$ entre o nó $S$ e o terra. Do nó
$S$ partem $\SI{6}{\ohm}$ até $A$ e $\SI{12}{\ohm}$ até $B$; de $A$ sai
$\SI{12}{\ohm}$ ao terra e de $B$ sai $\SI{6}{\ohm}$ ao terra; entre $A$ e
$B$ há $\SI{6}{\ohm}$.
\begin{itensq}
\item Verifique que a ponte está desequilibrada.
\item Converta o triângulo $S$--$A$--$B$ em estrela e determine $R_{eq}$.
\item Determine a corrente total.
\end{itensq}
\resposta{(b) $R_{eq}=\SI{8.4}{\ohm}$; (c) $I=\SI{10}{\ampere}$}
\resolucao{(a) Equilíbrio exigiria $\dfrac{6}{12}=\dfrac{12}{6}$, ou seja
$0{,}5=2$ --- falso. \textbf{Desequilibrada}, e o resistor central conduz.
Nenhuma associação série/paralelo é possível.\\
(b) O triângulo formado por $S$, $A$ e $B$ tem lados
$R_{SA}=6$, $R_{SB}=12$ e $R_{AB}=\SI{6}{\ohm}$, com soma
$\SI{24}{\ohm}$. Convertendo em estrela de centro $N$:
\[
R_S=\frac{6(12)}{24}=\SI{3}{\ohm};
\quad
R_A=\frac{6(6)}{24}=\SI{1.5}{\ohm};
\quad
R_B=\frac{12(6)}{24}=\SI{3}{\ohm}.
\]
Agora a rede é redutível: de $N$ partem dois caminhos ao terra,
\[
(1{,}5+12)=\SI{13.5}{\ohm}
\quad\text{e}\quad
(3+6)=\SI{9}{\ohm},
\]
em paralelo: $\dfrac{13{,}5(9)}{22{,}5}=\SI{5.4}{\ohm}$. Somando $R_S$:
\[
R_{eq}=3+5{,}4=\SI{8.4}{\ohm}.
\]
(c) $I=\dfrac{84}{8{,}4}=\SI{10}{\ampere}$.}
\end{questao}

\begin{questao}[3]{}
Resolva a mesma ponte por \textbf{análise nodal} e confronte com o resultado da
conversão.
\resposta{$V_A=\SI{48}{\volt}$, $V_B=\SI{36}{\volt}$; $R_{eq}=\SI{8.4}{\ohm}$
--- confere}
\resolucao{Com $V_S=\SI{84}{\volt}$ conhecido:
\[
\begin{cases}
V_A\left(\tfrac16+\tfrac1{12}+\tfrac16\right)-\tfrac16V_B=\tfrac{84}{6}=14\\[4pt]
-\tfrac16V_A+V_B\left(\tfrac1{12}+\tfrac16+\tfrac16\right)=\tfrac{84}{12}=7
\end{cases}
\]
Resolvendo: $V_A=\SI{48}{\volt}$ e $V_B=\SI{36}{\volt}$.\\
\textbf{Corrente total} (a que sai de $S$):
\[
I=\frac{84-48}{6}+\frac{84-36}{12}=6+4=\SI{10}{\ampere}
\;\Longrightarrow\;
R_{eq}=\frac{84}{10}=\SI{8.4}{\ohm}.\ \checkmark
\]
\textbf{Corrente no ramo central:} $(48-36)/6=\SI{2}{\ampere}$, de $A$ para
$B$ --- confirmando o desequilíbrio.\\
\textbf{Vantagem da conversão:} ela entrega $R_{eq}$ sem resolver sistema
algum. \textbf{Vantagem da nodal:} ela entrega \emph{também} os potenciais
internos e a corrente do ramo central --- que a conversão destrói. Escolha
conforme o que se precisa saber.}
\end{questao}

\begin{questao}[3]{}
Um triângulo tem um lado muito menor que os outros dois:
$R_{AB}=\SI{0.1}{\ohm}$, $R_{BC}=R_{CA}=\SI{100}{\ohm}$.
\begin{itensq}
\item Determine a estrela equivalente.
\item Compare $R_A+R_B$ com $R_{AB}$ e interprete.
\end{itensq}
\resposta{$R_A=R_B=\SI{49.98}{\milli\ohm}$, $R_C=\SI{49.98}{\ohm}$;
$R_A+R_B\approx R_{AB}$}
\resolucao{(a) $\Sigma=0{,}1+100+100=\SI{200.1}{\ohm}$:
\[
R_A=R_B=\frac{0{,}1(100)}{200{,}1}=\SI{49.98}{\milli\ohm};
\qquad
R_C=\frac{100(100)}{200{,}1}=\SI{49.98}{\ohm}.
\]
(b) $R_A+R_B=\SI{99.95}{\milli\ohm}\approx R_{AB}=\SI{0.1}{\ohm}$.\\
\textbf{Por quê:} a resistência entre $A$ e $B$ com $C$ aberto vale
$R_{AB}\parallel(R_{BC}+R_{CA})=0{,}1\parallel200\approx\SI{0.1}{\ohm}$, e na
estrela ela é exatamente $R_A+R_B$. O lado pequeno domina o paralelo, e a soma
dos dois braços apenas reproduz esse domínio.\\
\textbf{Consequência prática:} quando um lado do triângulo é muito menor que os
outros, a estrela resultante tem \emph{dois braços minúsculos e um braço
médio}. Os dois braços pequenos costumam ser desprezíveis diante do restante da
rede --- mas verifique antes de descartá-los, porque em redes de baixa
resistência eles podem ser exatamente o que importa.}
\end{questao}

\begin{questao}[3]{}
Uma estrela tem um braço muito pequeno: $R_A=R_B=\SI{10}{\ohm}$ e
$R_C=\SI{0.1}{\ohm}$.
\begin{itensq}
\item Determine o triângulo equivalente.
\item Analise o comportamento quando $R_C\to0$.
\end{itensq}
\resposta{$R_{AB}=\SI{1020}{\ohm}$, $R_{BC}=R_{CA}=\SI{10.2}{\ohm}$;
$R_{AB}\to\infty$}
\resolucao{(a) $\Sigma=10(10)+10(0{,}1)+0{,}1(10)=100+1+1
=\SI{102}{\ohm\squared}$:
\[
R_{AB}=\frac{102}{0{,}1}=\SI{1020}{\ohm};
\qquad
R_{BC}=R_{CA}=\frac{102}{10}=\SI{10.2}{\ohm}.
\]
(b) Com $R_C\to0$, o numerador tende a $R_AR_B=100$ e
\[
R_{AB}=\frac{100}{R_C}\to\infty,
\]
enquanto $R_{BC}\to R_B$ e $R_{CA}\to R_A$.\\
\textbf{Interpretação:} $R_C=0$ significa que o terminal $C$ coincide com o nó
central. No triângulo, isso equivale a $A$ e $B$ se ligarem a $C$ diretamente
por $R_A$ e $R_B$, sem lado direto entre eles --- e $R_{AB}=\infty$ é
exatamente a ausência desse lado. \checkmark\\
\textbf{Nota numérica:} $\SI{1020}{\ohm}$ em paralelo com $\SI{10.2}{\ohm}$
contribui com menos de $1\%$. Nesses casos, o lado enorme pode ser descartado
sem prejuízo --- e vale conferir se descartá-lo simplifica a rede antes de
seguir com o valor exato.}
\end{questao}

\begin{questao}[3]{}
Discuta o condicionamento das fórmulas de conversão para resistores positivos.
\begin{itensq}
\item Identifique se há risco de cancelamento catastrófico.
\item Determine onde, na prática, a precisão realmente se perde.
\end{itensq}
\resposta{As fórmulas são bem condicionadas; a perda de precisão ocorre na
\emph{redução posterior} da rede}
\resolucao{(a) \textbf{Não há risco.} Ambas as fórmulas envolvem apenas
\emph{produtos} e \emph{somas} de quantidades positivas, seguidos de uma
divisão por um número positivo. Não ocorre subtração de valores próximos ---
que é a única fonte de cancelamento catastrófico. Erros relativos de entrada se
propagam com fator próximo de $1$.\\
\textbf{Verificação:} arredondando a estrela do exemplo anterior
($49{,}98\,\mathrm{m}$; $49{,}98\,\mathrm{m}$; $\SI{49.98}{\ohm}$) para três
algarismos e convertendo de volta, recupera-se
$R_{AB}=\SI{0.1000}{\ohm}$ --- erro inferior a $0{,}1\%$.\\
(b) \textbf{A precisão se perde depois.} Situações típicas:
\begin{itemize}
\item Um braço convertido, muito grande, fica em paralelo com um pequeno: sua
contribuição é desprezível, e carregar seus algarismos é inútil --- mas
\emph{descartá-lo cedo demais} pode mascarar um efeito real.
\item O resultado final é uma \emph{diferença} entre valores próximos --- por
exemplo, ao calcular a corrente de uma ponte quase equilibrada, em que
$V_A-V_B$ é pequeno diante de $V_A$.
\end{itemize}
\textbf{Regra prática:} conduza a conversão com todos os algarismos e arredonde
\emph{apenas no resultado final}. E desconfie sempre que a resposta for uma
pequena diferença de números grandes --- ali, e não na transformação, é onde a
precisão evapora.}
\end{questao}

\begin{questao}[3]{}
Uma conversão produz um triângulo de $36{,}67$, $110$ e $\SI{55}{\ohm}$.
Escolha valores comerciais E24 e avalie o erro.
\begin{itensq}
\item Indique os valores adotados.
\item Converta de volta e compare com a estrela original de $10$, $20$ e
      $\SI{30}{\ohm}$.
\end{itensq}
\resposta{$36$, $110$ e $\SI{56}{\ohm}$; erros de $-0{,}2\%$, $-2{,}0\%$ e
$+1{,}7\%$ nos braços}
\resolucao{(a) Valores E24 mais próximos: $\SI{36}{\ohm}$, $\SI{110}{\ohm}$ e
$\SI{56}{\ohm}$ (desvios de $-1{,}8\%$, $0\%$ e $+1{,}8\%$ nos lados).\\
(b) Convertendo o triângulo comercial de volta a estrela
($\Sigma=36+110+56=\SI{202}{\ohm}$):
\[
R_A=\frac{36(56)}{202}=\SI{9.98}{\ohm};
\quad
R_B=\frac{36(110)}{202}=\SI{19.60}{\ohm};
\quad
R_C=\frac{110(56)}{202}=\SI{30.50}{\ohm}.
\]
Erros: $-0{,}2\%$, $-2{,}0\%$ e $+1{,}7\%$ em relação a $10$, $20$ e
$\SI{30}{\ohm}$.\\
\textbf{Observação importante:} os erros dos braços \emph{não} reproduzem os
erros dos lados --- eles se recombinam. Um lado com erro nulo
($\SI{110}{\ohm}$) contribuiu para dois braços que erraram $2\%$.\\
\textbf{Método correto de avaliação:} nunca julgue o erro pelos componentes
isolados. Converta de volta e compare com o alvo, ou --- melhor ainda --- calcule
a grandeza que realmente interessa (a resistência entre os terminais usados) nos
dois casos.}
\end{questao}

\begin{questao}[3]{}
Compare o esforço de resolver a ponte desequilibrada por conversão
estrela-triângulo e por análise nodal.
\begin{itensq}
\item Enumere as etapas de cada método.
\item Indique o critério de escolha.
\end{itensq}
\resposta{Conversão: 3 fórmulas e reduções; nodal: sistema $2\times2$. A
conversão dá só $R_{eq}$; a nodal dá tudo}
\resolucao{(a) \textbf{Conversão:} identificar o triângulo (1); aplicar três
fórmulas (2); somar duas séries (3); um paralelo (4); uma soma final (5). Sem
sistemas de equações.\\
\textbf{Nodal:} montar a matriz $2\times2$ (1); resolver (2); calcular a
corrente total a partir dos potenciais (3).\\
(b) \textbf{Critério.}
\begin{itemize}
\item Se a pergunta é \textbf{só} $R_{eq}$ --- ou a corrente total --- a
conversão é mais rápida e não exige álgebra linear.
\item Se são necessários potenciais internos, correntes de ramo ou potências
individuais, use \textbf{nodal}. A conversão destrói a estrutura interna: o nó
central da estrela não existe no circuito real, e nenhuma corrente calculada
nele tem significado físico.
\item Para pontes com \emph{mais} de um triângulo irredutível, a conversão
precisa ser aplicada repetidamente e perde a vantagem; a nodal escala melhor.
\end{itemize}
\textbf{Estratégia mista frequente:} converter para obter $R_{eq}$ e a corrente
total, e depois voltar ao circuito \emph{original} para distribuir essa corrente
--- aproveitando o melhor dos dois.}
\end{questao}

\begin{questao}[3]{}
Uma rede de três terminais é formada por resistores positivos em uma topologia
qualquer. Discuta se ela sempre admite equivalente em estrela.
\begin{itensq}
\item Estabeleça a condição.
\item Dê um exemplo em que a estrela equivalente teria braço negativo.
\end{itensq}
\resposta{A condição é a desigualdade triangular sobre as resistências entre
pares}
\resolucao{(a) Sejam $R_{ab}$, $R_{bc}$ e $R_{ca}$ as resistências medidas
entre cada par de terminais (com o terceiro aberto). Da estrela, valem
$R_{ab}=R_A+R_B$, $R_{bc}=R_B+R_C$ e $R_{ca}=R_C+R_A$, de onde
\[
R_A=\frac{R_{ab}+R_{ca}-R_{bc}}{2},
\]
e analogamente para $R_B$ e $R_C$.\\
\textbf{Condição:} os três braços são não negativos se e somente se cada soma
de duas resistências medidas for maior ou igual à terceira --- exatamente a
\textbf{desigualdade triangular}.\\
(b) \textbf{Contraexemplo.} Suponha
$R_{ab}=\SI{1}{\ohm}$, $R_{bc}=\SI{10}{\ohm}$ e
$R_{ca}=\SI{1}{\ohm}$. Então
\[
R_B=\frac{1+10-1}{2}=\SI{5}{\ohm};
\qquad
R_A=\frac{1+1-10}{2}=\SI{-4}{\ohm}.
\]
Braço negativo --- fisicamente irrealizável com resistores passivos.\\
\textbf{Mas atenção:} essa combinação de medidas também não é realizável por
uma rede passiva de três terminais. A desigualdade triangular é uma
\emph{propriedade} de tais redes, não uma restrição adicional. Braços negativos
só surgem como artifício de cálculo --- em transformações intermediárias ou em
redes com elementos ativos.}
\end{questao}

\blocodesafio

\begin{questao}[4]{}
\textbf{(Dedução --- triângulo para estrela.)} Deduza as fórmulas de conversão a
partir das resistências entre terminais.
\begin{itensq}
\item Escreva as três equações de equivalência.
\item Resolva o sistema para $R_A$, $R_B$ e $R_C$.
\item Obtenha a forma final em função dos lados do triângulo.
\end{itensq}
\resposta{$R_A=\dfrac{R_{AB}R_{CA}}{R_{AB}+R_{BC}+R_{CA}}$ e análogas}
\resolucao{(a) A equivalência exige que a resistência entre cada par de
terminais, com o terceiro \emph{aberto}, seja a mesma nas duas configurações:
\[
\begin{cases}
R_A+R_B=R_{AB}\parallel(R_{BC}+R_{CA})\\[3pt]
R_B+R_C=R_{BC}\parallel(R_{CA}+R_{AB})\\[3pt]
R_C+R_A=R_{CA}\parallel(R_{AB}+R_{BC})
\end{cases}
\]
(b) Somando as três e dividindo por $2$, obtém-se $R_A+R_B+R_C$. Subtraindo
dessa soma a segunda equação, resta
\[
R_A=\frac{(R_A+R_B)+(R_C+R_A)-(R_B+R_C)}{2}.
\]
(c) Substituindo os paralelos, com $\Sigma=R_{AB}+R_{BC}+R_{CA}$:
\[
R_A+R_B=\frac{R_{AB}(R_{BC}+R_{CA})}{\Sigma};
\qquad
R_C+R_A=\frac{R_{CA}(R_{AB}+R_{BC})}{\Sigma};
\]
\[
R_B+R_C=\frac{R_{BC}(R_{CA}+R_{AB})}{\Sigma}.
\]
Levando à expressão de $R_A$ e simplificando o numerador:
\[
R_{AB}R_{BC}+R_{AB}R_{CA}+R_{CA}R_{AB}+R_{CA}R_{BC}
-R_{BC}R_{CA}-R_{BC}R_{AB}=2R_{AB}R_{CA},
\]
donde
\[
R_A=\frac{2R_{AB}R_{CA}}{2\Sigma}=\frac{R_{AB}R_{CA}}{\Sigma}.
\]
\textbf{A regra mnemônica agora tem origem:} o produto dos dois lados que tocam
o terminal aparece porque os demais termos se cancelam aos pares na combinação
$(R_{ab}+R_{ca}-R_{bc})/2$.}
\end{questao}

\begin{questao}[4]{}
\textbf{(Dedução --- estrela para triângulo.)} Obtenha a transformação inversa.
\begin{itensq}
\item Deduza $R_{AB}$ em função dos braços.
\item Verifique que a composição das duas transformações é a identidade.
\item Explique por que o denominador é o braço \emph{oposto}.
\end{itensq}
\resposta{$R_{AB}=\dfrac{R_AR_B+R_BR_C+R_CR_A}{R_C}$}
\resolucao{(a) Tomando as três fórmulas diretas e formando a soma dos produtos
dois a dois:
\[
R_AR_B+R_BR_C+R_CR_A
=\frac{R_{AB}R_{BC}R_{CA}\left(R_{AB}+R_{BC}+R_{CA}\right)}{\Sigma^{2}}
=\frac{R_{AB}R_{BC}R_{CA}}{\Sigma}.
\]
Chamando esse valor de $P$ e observando que
$R_C=\dfrac{R_{BC}R_{CA}}{\Sigma}$, tem-se
\[
\frac{P}{R_C}
=\frac{R_{AB}R_{BC}R_{CA}/\Sigma}{R_{BC}R_{CA}/\Sigma}=R_{AB}.
\]
Logo
\[
R_{AB}=\frac{R_AR_B+R_BR_C+R_CR_A}{R_C}.
\]
(b) A dedução \emph{é} a verificação: partiu-se das fórmulas diretas e
chegou-se de volta aos lados originais. Numericamente, o exemplo
$20/30/\SI{50}{\ohm}\to10/6/\SI{15}{\ohm}\to20/30/\SI{50}{\ohm}$ confirma.\\
(c) \textbf{Por que o oposto.} O braço $R_C$ é o único que \emph{não} contém
$R_{AB}$ em seu numerador --- ele é construído a partir de $R_{BC}$ e
$R_{CA}$. Ao dividir $P$ (que contém o produto dos três lados) por $R_C$,
cancelam-se exatamente $R_{BC}$ e $R_{CA}$, isolando $R_{AB}$. É por isso que
o denominador tem de ser o braço oposto, e não outro.}
\end{questao}

\begin{questao}[4]{}
\textbf{(Positividade.)} Prove que as transformações preservam a positividade
dos elementos.
\begin{itensq}
\item Mostre que $\Delta\to Y$ com lados positivos gera braços positivos.
\item Mostre que $Y\to\Delta$ com braços positivos gera lados positivos.
\item Determine o que ocorre quando um elemento tende a zero ou a infinito.
\end{itensq}
\resposta{Ambas as fórmulas são quocientes de quantidades positivas ---
positividade garantida}
\resolucao{(a) $R_A=\dfrac{R_{AB}R_{CA}}{R_{AB}+R_{BC}+R_{CA}}$. Com todos os
lados positivos, numerador e denominador são positivos, logo $R_A>0$. O mesmo
para $R_B$ e $R_C$.\\
\textbf{Além disso, há um limite superior:} como
$\Sigma>R_{CA}$, segue $R_A<R_{AB}$; e como $\Sigma>R_{AB}$, segue
$R_A<R_{CA}$. Ou seja, \textbf{cada braço é menor que os dois lados que o
originaram} --- teste de sanidade imediato para qualquer conversão.\\
(b) $R_{AB}=\dfrac{R_AR_B+R_BR_C+R_CR_A}{R_C}$: soma de três produtos positivos
dividida por um positivo, logo $R_{AB}>0$.\\
\textbf{E há um limite inferior:} $R_{AB}>\dfrac{R_AR_C+R_BR_C}{R_C}=R_A+R_B$
--- \textbf{cada lado é maior que a soma dos dois braços correspondentes}.
Segundo teste de sanidade, dual do anterior.\\
(c) \textbf{Casos-limite.}
\begin{itemize}
\item $R_{AB}\to0$ (lado em curto) $\Rightarrow$ $R_A,R_B\to0$: os terminais
$A$ e $B$ colapsam sobre o nó central.
\item $R_{AB}\to\infty$ (lado aberto) $\Rightarrow$ a estrela degenera em uma
série simples.
\item $R_C\to0$ na estrela $\Rightarrow$ $R_{AB}\to\infty$: o lado direto
desaparece.
\item $R_C\to\infty$ $\Rightarrow$ $R_{AB}\to R_A+R_B$, e os outros dois lados
divergem.
\end{itemize}
Em todos os casos, o limite reproduz a interpretação física esperada --- as
fórmulas nunca produzem resultados sem sentido para entradas positivas.}
\end{questao}

\begin{questao}[4]{}
\textbf{(Alcance da equivalência.)} Analise \emph{o que} a transformação
preserva e o que ela destrói.
\begin{itensq}
\item Prove que as resistências entre os três terminais são preservadas.
\item Determine se correntes e potências internas são preservadas.
\item Estabeleça o procedimento correto quando grandezas internas são
      necessárias.
\end{itensq}
\resposta{Preserva o comportamento nos três terminais; \emph{não} preserva nada
do interior}
\resolucao{(a) Por construção, as três equações de equivalência impõem
igualdade das resistências entre pares. Como a rede de três terminais é linear e
tem apenas três acessos, seu comportamento externo é completamente descrito por
esses três valores --- qualquer excitação aplicada aos terminais produz a mesma
resposta nas duas configurações.\\
\emph{(Formalmente: a matriz de condutâncias de três terminais tem três
parâmetros independentes, e as três equações os fixam todos.)}\\
(b) \textbf{Não.} A estrela possui um nó central inexistente no triângulo, e o
triângulo possui um lado direto inexistente na estrela. Consequências:
\begin{itemize}
\item As correntes dos elementos são \emph{completamente diferentes}.
\item As potências individuais também --- embora a potência \emph{total} seja a
mesma, pois o comportamento externo é idêntico.
\item O potencial do nó central não corresponde a ponto algum do triângulo.
\end{itemize}
\textbf{Exemplo:} na ponte resolvida anteriormente, o resistor central de
$\SI{6}{\ohm}$ conduz $\SI{2}{\ampere}$. Após a conversão, esse resistor
\emph{não existe} --- e nenhum braço da estrela conduz $\SI{2}{\ampere}$.\\
(c) \textbf{Procedimento correto:} use a conversão para obter as grandezas
\emph{externas} (resistência equivalente, corrente total, tensões nos
terminais). Depois \textbf{volte ao circuito original} e, com essas grandezas já
conhecidas, distribua correntes e tensões pelos elementos reais.\\
Se muitas grandezas internas forem necessárias, a análise nodal costuma ser
mais direta desde o início.}
\end{questao}

\begin{questao}[4]{}
\textbf{(Projeto com valores comerciais.)} Uma estrela de $10$, $20$ e
$\SI{30}{\ohm}$ deve ser substituída por um triângulo montado com resistores
E24.
\begin{itensq}
\item Determine o triângulo exato.
\item Escolha valores comerciais e avalie o erro na resistência entre $A$ e
      $B$.
\item Proponha uma estratégia para reduzir o erro.
\end{itensq}
\resposta{Exato: $36{,}67$, $110$ e $\SI{55}{\ohm}$; com $36/110/56$, o erro em
$R_{AB}$ é de cerca de $1\%$}
\resolucao{(a) $\Sigma=10(20)+20(30)+30(10)=200+600+300=\SI{1100}{\ohm\squared}$:
\[
R_{AB}=\frac{1100}{30}=\SI{36.67}{\ohm};
\quad
R_{BC}=\frac{1100}{10}=\SI{110}{\ohm};
\quad
R_{CA}=\frac{1100}{20}=\SI{55}{\ohm}.
\]
(b) E24: $\SI{36}{\ohm}$, $\SI{110}{\ohm}$, $\SI{56}{\ohm}$.\\
\textbf{Grandeza que interessa} --- resistência entre $A$ e $B$ com $C$ aberto:
\[
\text{alvo (estrela)}: R_A+R_B=\SI{30}{\ohm};
\]
\[
\text{obtido}: 36\parallel(110+56)=\frac{36(166)}{202}=\SI{29.58}{\ohm}.
\]
Erro de $-1{,}4\%$ --- aceitável para a maioria das aplicações, e menor que a
tolerância de $5\%$ dos próprios resistores.\\
(c) \textbf{Estratégias, em ordem de custo.}
\begin{itemize}
\item \textbf{Testar combinações:} $39/110/56$ daria
$39\parallel166=\SI{31.6}{\ohm}$ ($+5{,}3\%$) --- pior. O par
$36/110/56$ já é bom.
\item \textbf{Série E96} ($1\%$): contém $36{,}5$ e $54{,}9$, reduzindo o erro
a frações de por cento.
\item \textbf{Associação:} $\SI{36.67}{\ohm}$ sai de $\SI{33}{\ohm}$ em série
com $\SI{3.6}{\ohm}$, ou de $\SI{110}{\ohm}$ em paralelo com
$\SI{55}{\ohm}$ --- e note que esses dois valores já estão na lista!
\end{itemize}
\textbf{Critério de avaliação:} sempre meça o erro na \emph{grandeza de
interesse}, não nos componentes. Erros de lados individuais podem se compensar
ou se somar, e só o cálculo da resistência terminal revela o efeito real.}
\end{questao}

\begin{questao}[4]{}
\textbf{(Ponte quase equilibrada.)} Uma ponte tem $R_1=\SI{100}{\ohm}$,
$R_2=\SI{100}{\ohm}$, $R_3=\SI{100}{\ohm}$ e
$R_4=\SI{101}{\ohm}$, com $R_g=\SI{100}{\ohm}$ e alimentação de
$\SI{10}{\volt}$.
\begin{itensq}
\item Estime a corrente do galvanômetro por análise nodal.
\item Discuta a viabilidade de obter esse valor por conversão
      estrela-triângulo.
\end{itensq}
\resposta{$I_g\approx\SI{-124}{\micro\ampere}$ (de $B$ para $A$); a conversão é
inadequada para esta grandeza}
\resolucao{(a) Com $V_S=\SI{10}{\volt}$ e o terra na base, o sistema nodal
\[
\begin{cases}
V_A\left(\tfrac{1}{100}+\tfrac{1}{100}+\tfrac{1}{100}\right)
-\tfrac{V_B}{100}=\tfrac{10}{100}\\[4pt]
-\tfrac{V_A}{100}
+V_B\left(\tfrac{1}{100}+\tfrac{1}{101}+\tfrac{1}{100}\right)=\tfrac{10}{100}
\end{cases}
\]
dá $V_A=\SI{5.00621}{\volt}$ e $V_B=\SI{5.01863}{\volt}$, de onde
\[
I_g=\frac{V_A-V_B}{100}=\frac{-0{,}01242}{100}=\SI{-124}{\micro\ampere}.
\]
O sinal indica corrente de $B$ para $A$ --- coerente, pois $R_4>R_3$ eleva
$V_B$.\\
(b) \textbf{A conversão é inadequada, por duas razões independentes.}\\
\emph{(i) Ela não fornece a grandeza pedida.} Após converter o triângulo, o
resistor do galvanômetro deixa de existir como elemento --- e é justamente sua
corrente que se quer. Seria necessário converter, obter $R_{eq}$ e a corrente
total, voltar ao circuito original e só então distribuir --- caminho mais longo
que a nodal direta.\\
\emph{(ii) O problema é mal condicionado.} A resposta é a diferença
$V_A-V_B=\SI{-12.4}{\milli\volt}$ entre dois valores próximos de
$\SI{5}{\volt}$ --- uma diferença relativa de apenas $0{,}25\%$.
\begin{center}
\begin{tabular}{ccc}
\hline
Algarismos usados & $V_A-V_B$ obtido & erro em $I_g$\\
\hline
$6$ ($5{,}00621$; $5{,}01863$) & $\SI{-12.42}{\milli\volt}$ & $<0{,}1\%$\\
$4$ ($5{,}006$; $5{,}019$) & $\SI{-13}{\milli\volt}$ & $\approx5\%$\\
$3$ ($5{,}01$; $5{,}02$) & $\SI{-10}{\milli\volt}$ & $\approx20\%$\\
\hline
\end{tabular}
\end{center}
Arredondar cedo destrói o resultado.\\
\textbf{Método adequado:} tratar o desequilíbrio como \emph{perturbação}. Para
$\delta=\Delta R/R$ pequeno, $I_g$ é aproximadamente proporcional a $\delta$, e
a expressão linearizada evita completamente a subtração de valores próximos.\\
\textbf{Moral:} escolher o método é também escolher \emph{como} o erro numérico
se propaga. A conversão estrela-triângulo é excelente para $R_{eq}$ e péssima
para desequilíbrios pequenos.}
\end{questao}

\begin{questao}[4]{}
\textbf{(Síntese.)} Uma rede de três terminais deve apresentar
$R_{AB}=\SI{20}{\ohm}$, $R_{BC}=\SI{30}{\ohm}$ e $R_{CA}=\SI{40}{\ohm}$
(medidas com o terceiro terminal aberto).
\begin{itensq}
\item Verifique a realizabilidade e determine a estrela que atende.
\item Determine o triângulo equivalente.
\item Escolha valores comerciais e avalie qual das duas topologias erra menos.
\end{itensq}
\resposta{$R_A=\SI{15}{\ohm}$, $R_B=\SI{5}{\ohm}$, $R_C=\SI{25}{\ohm}$;
triângulo $23$, $38{,}33$ e $\SI{115}{\ohm}$; a estrela é preferível}
\resolucao{(a) \textbf{Realizabilidade.} A desigualdade triangular exige que
cada resistência medida seja menor que a soma das outras duas:
$20<70$, $30<60$, $40<50$ \checkmark. A rede é realizável com resistores
positivos.\\
Invertendo o sistema $R_A+R_B=R_{AB}$ e análogos:
\[
R_A=\frac{20+40-30}{2}=\SI{15}{\ohm};
\quad
R_B=\frac{20+30-40}{2}=\SI{5}{\ohm};
\quad
R_C=\frac{30+40-20}{2}=\SI{25}{\ohm}.
\]
\textbf{Conferência:} $15+5=20$; $5+25=30$; $25+15=40$ \checkmark.\\
(b) $\Sigma=15(5)+5(25)+25(15)=75+125+375=\SI{575}{\ohm\squared}$:
\[
R_{AB}=\frac{575}{25}=\SI{23}{\ohm};
\quad
R_{BC}=\frac{575}{15}=\SI{38.33}{\ohm};
\quad
R_{CA}=\frac{575}{5}=\SI{115}{\ohm}.
\]
(c) \textbf{Estrela com E24:} $\SI{15}{\ohm}$ e $\SI{24}{\ohm}$ existem, e o
mais próximo de $\SI{5}{\ohm}$ é $\SI{5.1}{\ohm}$. Com
$(15;\,5{,}1;\,24)$:
\[
R_{AB}=\SI{20.1}{\ohm}\ (+0{,}5\%);
\quad
R_{BC}=\SI{29.1}{\ohm}\ (-3{,}0\%);
\quad
R_{CA}=\SI{39.0}{\ohm}\ (-2{,}5\%).
\]
\textbf{Triângulo com E24:} $(24;\,39;\,110)$:
\[
R_{AB}=\SI{20.67}{\ohm}\ (+3{,}4\%);
\quad
R_{BC}=\SI{30.21}{\ohm}\ (+0{,}7\%);
\quad
R_{CA}=\SI{40.06}{\ohm}\ (+0{,}1\%).
\]
\textbf{Conclusão.} Os dois erram no mesmo patamar (pior caso de $3{,}0\%$
contra $3{,}4\%$), mas a \textbf{estrela é preferível} por três razões: seus
valores são menores (componentes mais baratos), o erro é dominado por um único
braço difícil ($\SI{5}{\ohm}$), e cada resistência medida depende de
\emph{apenas dois} braços --- o que torna o ajuste previsível.\\
No triângulo, cada medida envolve os \emph{três} lados através de um paralelo, e
os erros se recombinam de forma difícil de controlar.\\
\textbf{Se a precisão for crítica:} o braço de $\SI{5}{\ohm}$ pode ser obtido
associando $\SI{10}{\ohm}$ em paralelo com $\SI{10}{\ohm}$, eliminando o
único desvio significativo.}
\end{questao}

\begin{questao}[4]{}
\textbf{(Generalização.)} A transformação estrela-triângulo é o caso $n=3$ de
uma família mais ampla.
\begin{itensq}
\item Descreva a transformação estrela-polígono para $n$ terminais.
\item Determine quantos elementos há em cada configuração.
\item Explique por que a transformação inversa nem sempre existe para $n>3$.
\end{itensq}
\resposta{Estrela: $n$ braços; polígono: $n(n-1)/2$ lados. A inversa só existe
para $n=3$}
\resolucao{(a) Uma estrela de $n$ braços $R_1,\dots,R_n$ ligados a um nó comum
equivale a um polígono completo em que o lado entre os terminais $i$ e $j$ vale
\[
R_{ij}=R_iR_j\sum_{k=1}^{n}\frac{1}{R_k}.
\]
Para $n=3$ isso reproduz a fórmula conhecida:
$R_{12}=R_1R_2\left(\frac1{R_1}+\frac1{R_2}+\frac1{R_3}\right)
=\frac{R_1R_2+R_2R_3+R_3R_1}{R_3}$ \checkmark.\\
(b) \textbf{Estrela:} $n$ elementos. \textbf{Polígono completo:}
$\dbinom{n}{2}=\dfrac{n(n-1)}{2}$ elementos.
\begin{center}
\begin{tabular}{ccc}
\hline
$n$ & estrela & polígono\\
\hline
$3$ & $3$ & $3$\\
$4$ & $4$ & $6$\\
$5$ & $5$ & $10$\\
\hline
\end{tabular}
\end{center}
(c) \textbf{Contagem de graus de liberdade.} A transformação
$Y\to$ polígono sempre existe. A inversa exigiria determinar $n$ braços a
partir de $n(n-1)/2$ lados --- um sistema \textbf{sobredeterminado} para
$n>3$.\\
Só quando $n=3$ os dois números coincidem ($3=3$), e a correspondência é
biunívoca. Para $n=4$, seria preciso obter $4$ incógnitas de $6$ equações: um
polígono \emph{genérico} de quatro terminais simplesmente não admite estrela
equivalente.\\
\textbf{Consequência prática:} a conversão $\Delta\to Y$ é uma ferramenta
peculiar ao caso de três terminais. Redes de quatro ou mais acessos exigem
descrição matricial completa --- e é por isso que a análise nodal, que não
depende de topologias especiais, é o método geral.}
\end{questao}

% END BANCO COMPLEMENTAR

\gabarito[3.4 Transformação estrela-triângulo]
  % completo

% Cap. 4 -------------------------------------------------------------
\setcounter{section}{3}
% =====================================================================
%  CAPITULO 4 - DIVISOR DE TENSAO E CORRENTE
%  (versao revisada e ampliada: 2 -> 16 questoes)
% =====================================================================
\section{Divisor de Tensão e de Corrente}

\begin{conceito}[Antes de começar]
\textbf{Divisor de tensão} (resistores em \emph{série}): a tensão total se
reparte na proporção direta das resistências.
\[
U_k = U_{\text{total}}\;\frac{R_k}{R_1+R_2+\dots+R_n}.
\]
\textbf{Divisor de corrente} (resistores em \emph{paralelo}): a corrente total se
reparte na proporção \emph{inversa} das resistências.
\[
I_k = I_{\text{total}}\;\frac{G_k}{G_1+\dots+G_n}
     = I_{\text{total}}\;\frac{1/R_k}{1/R_1+\dots+1/R_n}.
\]
\textbf{Caso de dois resistores} (o mais usado):
\[
U_1 = U\,\frac{R_1}{R_1+R_2}, \qquad
I_1 = I\,\frac{R_2}{R_1+R_2}\quad(\text{note o índice \emph{trocado}}).
\]
\end{conceito}

\begin{atencao}[O erro clássico]
No divisor de \textbf{corrente} de dois resistores, o ramo $R_1$ recebe a
corrente proporcional a $R_2$ (o \emph{outro} resistor). Trocar os índices faz a
maior corrente ir para o maior resistor --- o oposto do correto. Regra de ouro:
\emph{a maior corrente sempre passa pelo menor resistor}.
\end{atencao}

\legendaniveis

% =====================================================================
\subsection{Divisor de tensão}
% =====================================================================

\blocofixacao

\begin{questao}[1]{}
Em um divisor de tensão formado por dois resistores em série, a maior parcela de
tensão aparece:
\begin{alternativas}
\item sobre o menor resistor.
\item sobre o maior resistor.
\item igualmente sobre os dois, sempre.
\item sobre o resistor mais próximo do terminal positivo.
\item sobre nenhum: a tensão fica toda na fonte.
\end{alternativas}
\resposta{(b)}
\resolucao{Como a corrente é a mesma nos dois ($U = R\,I$), a tensão é
proporcional à resistência: o maior resistor fica com a maior tensão. A posição
em relação ao terminal positivo (alternativa d) é irrelevante.}
\end{questao}

\begin{questao}[1]{}
No circuito, três resistores em série de $\SI{3}{\ohm}$, $\SI{4}{\ohm}$ e
$\SI{5}{\ohm}$ estão ligados a uma fonte de $\SI{120}{\volt}$. Calcule a tensão
sobre cada resistor.

\circuitoq{%
\begin{circuitikz}[scale=0.95,transform shape,line width=0.8pt]
 \draw (0,0) to[battery1,l_={$E=\SI{120}{\volt}$}] (0,3)
       to[R,l={$R_1=\SI{3}{\ohm}$}] (3.2,3)
       to[R,l={$R_2=\SI{4}{\ohm}$}] (3.2,0)
       to[R,l={$R_3=\SI{5}{\ohm}$}] (0,0);
 \draw[->,Icol,line width=1pt] (0.7,3.4)--(1.7,3.4)
       node[midway,above,font=\footnotesize]{$I$};
\end{circuitikz}}

\resposta{$U_1=\SI{30}{\volt}$, $U_2=\SI{40}{\volt}$, $U_3=\SI{50}{\volt}$}
\resolucao{A corrente é $I = \dfrac{120}{3+4+5} = \dfrac{120}{12} = \SI{10}{\ampere}$.
Aplicando o divisor (ou simplesmente $U_k = R_k I$):
\[
U_1 = 3\cdot10 = \SI{30}{\volt},\quad
U_2 = 4\cdot10 = \SI{40}{\volt},\quad
U_3 = 5\cdot10 = \SI{50}{\volt}.
\]
\textbf{Verificação:} $30+40+50 = \SI{120}{\volt}$, igual à fonte.}
\end{questao}

\begin{questao}[1]{}
Uma fonte de $\SI{100}{\volt}$ alimenta dois resistores em série,
$R_1=\SI{60}{\ohm}$ e $R_2=\SI{40}{\ohm}$. A tensão de saída, tomada sobre $R_2$,
vale:
\begin{alternativas}[3]
\item \SI{20}{\volt}
\item \SI{40}{\volt}
\item \SI{50}{\volt}
\item \SI{60}{\volt}
\item \SI{100}{\volt}
\end{alternativas}
\resposta{(b)}
\resolucao{Aplicando diretamente a fórmula do divisor:
\[
U_2 = E\,\frac{R_2}{R_1+R_2} = 100\cdot\frac{40}{100} = \SI{40}{\volt}.
\]}
\end{questao}

\begin{questao}[1]{}
Três resistores iguais estão em série sob $\SI{60}{\volt}$. A tensão sobre cada um
vale:
\begin{alternativas}[3]
\item \SI{10}{\volt}
\item \SI{15}{\volt}
\item \SI{20}{\volt}
\item \SI{30}{\volt}
\item \SI{60}{\volt}
\end{alternativas}
\resposta{(c)}
\resolucao{Resistores iguais dividem a tensão igualmente:
$U = \dfrac{60}{3} = \SI{20}{\volt}$ em cada um. É o divisor de tensão mais
simples --- a base dos potenciômetros de precisão.}
\end{questao}

\blocoaplicacao

\begin{questao}[2]{}
Deseja-se obter $\SI{5}{\volt}$ a partir de uma fonte de $\SI{15}{\volt}$ usando
um divisor de dois resistores. Sabendo que $R_1 = \SI{2}{\kilo\ohm}$ (ligado ao
terminal positivo), determine $R_2$ (a saída é tomada sobre $R_2$).
\resposta{$R_2 = \SI{1}{\kilo\ohm}$}
\resolucao{Da relação do divisor,
\[
\frac{U_{\text{out}}}{E}=\frac{R_2}{R_1+R_2}
\ \Rightarrow\
\frac{5}{15}=\frac{R_2}{2000+R_2}
\ \Rightarrow\
2000+R_2 = 3R_2
\ \Rightarrow\
R_2 = \SI{1}{\kilo\ohm}.
\]
Repare que só importa a \emph{razão} $R_1{:}R_2 = 2{:}1$; qualquer par nessa
proporção daria \SI{5}{\volt} em vazio.}
\end{questao}

\begin{questao}[2]{}
No circuito da figura, calcule a tensão de saída $U_{\text{out}}$ (a) em vazio e
(b) quando uma carga $R_L = \SI{1}{\kilo\ohm}$ é ligada aos terminais de saída.
Comente o resultado.

\circuitoq{%
\begin{circuitikz}[scale=0.95,transform shape,line width=0.8pt]
 \draw (0,0) to[battery1,l_={$E=\SI{12}{\volt}$}] (0,3.4)
       to[R,l={$R_1=\SI{1}{\kilo\ohm}$}] (3.4,3.4) coordinate(a);
 \draw (a) to[R,l={$R_2=\SI{1}{\kilo\ohm}$}] (3.4,0) coordinate(b) -- (0,0);
 \draw (a) to[short,-o] (5,3.4) node[right]{$+$};
 \draw (b) to[short,-o] (5,0) node[right]{$-$};
 \draw[->] (5.3,3.0) -- (5.3,0.4) node[midway,right,font=\footnotesize]{$U_{\text{out}}$};
\end{circuitikz}}

\resposta{(a) \SI{6}{\volt}; (b) \SI{4}{\volt}}
\resolucao{(a) \textbf{Em vazio}, nenhuma corrente sai pelos terminais, e os dois
resistores iguais dividem a tensão ao meio:
$U_{\text{out}} = 12\cdot\dfrac{1}{2} = \SI{6}{\volt}$.\\
(b) \textbf{Com carga}, $R_L$ fica em paralelo com $R_2$:
\[
R_2 \parallel R_L = 1\text{k}\parallel 1\text{k} = \SI{0.5}{\kilo\ohm}
\ \Rightarrow\
U_{\text{out}} = 12\cdot\frac{0{,}5}{1+0{,}5} = \SI{4}{\volt}.
\]
\textbf{Conclusão:} a conexão da carga reduz a tensão de saída. Divisores resistivos só
mantêm a saída se $R_L \gg R_2$; por isso não servem como fonte de alimentação
para cargas variáveis --- para isso usa-se um regulador de tensão.}
\end{questao}

\begin{questao}[2]{}
Um potenciômetro de $\SI{10}{\kilo\ohm}$ é ligado como divisor de tensão a uma
fonte de $\SI{9}{\volt}$. Quando o cursor está a $\SI{30}{\percent}$ do curso (a
partir do terminal de referência), qual é a tensão de saída em vazio?
\resposta{\SI{2.7}{\volt}}
\resolucao{O cursor divide a resistência total na proporção do curso. A \SI{30}{\percent},
a saída corresponde a $\SI{30}{\percent}$ da tensão total:
\[
U_{\text{out}} = 9\cdot 0{,}30 = \SI{2.7}{\volt}.
\]
O valor da resistência total (\SI{10}{\kilo\ohm}) não entra na conta em vazio ---
apenas a \emph{fração} do curso importa. É exatamente assim que um controle de
volume ajusta o sinal.}
\end{questao}

\begin{questao}[2]{}
Para o circuito da figura (malha única com cinco resistores), calcule a corrente
e a tensão sobre o resistor de $\SI{8}{\ohm}$.

\circuitoq{%
\begin{circuitikz}[scale=0.9,transform shape,line width=0.8pt]
 \draw (0,0) to[battery1,l_={$E=\SI{100}{\volt}$},invert] (0,3)
       to[R,l={$\SI{10}{\ohm}$}] (3,3)
       to[R,l={$\SI{6}{\ohm}$}] (6,3)
       to[R,l={$\SI{8}{\ohm}$}] (6,0)
       to[R,l={$\SI{12}{\ohm}$}] (3,0)
       to[R,l={$\SI{6}{\ohm}$}] (0,0);
 \draw[->,Icol,line width=1pt] (0.7,3.4)--(1.7,3.4)
       node[midway,above,font=\footnotesize]{$I$};
\end{circuitikz}}

\resposta{$I = \SI{2.38}{\ampere}$ (aprox.); $U_{8} \approx \SI{19}{\volt}$}
\resolucao{Todos os resistores estão em série:
$\Req = 10+6+8+12+6 = \SI{42}{\ohm}$.
\[
I=\frac{100}{42}\approx\SI{2.38}{\ampere},
\qquad
U_8 = 8\cdot I = \frac{800}{42}\approx\SI{19.0}{\volt}.
\]
Pelo divisor: $U_8 = 100\cdot\dfrac{8}{42}\approx\SI{19.0}{\volt}$, o mesmo
resultado.}
\end{questao}

\blocoaprofundamento

\begin{questao}[3]{}
Um divisor de tensão deve fornecer $\SI{5}{\volt}$ a uma carga que consome
$\SI{10}{\milli\ampere}$, a partir de uma fonte de $\SI{15}{\volt}$. Como regra
de projeto, adota-se a \emph{corrente de sangria} (a que circula pelo divisor sem
passar pela carga) igual a $\SI{10}{\milli\ampere}$. Determine $R_1$ e $R_2$.

\circuitoq{%
\begin{circuitikz}[scale=0.95,transform shape,line width=0.8pt]
 \draw (0,0) to[battery1,l_={$E=\SI{15}{\volt}$}] (0,3.4)
       to[R,l={$R_1$},i={$I_1$}] (3.4,3.4) coordinate(a);
 \draw (a) to[R,l={$R_2$},i={$I_2$}] (3.4,0) coordinate(b) -- (0,0);
 \draw (a) to[short,-o] (5.2,3.4) coordinate(o1);
 \draw (b) to[short,-o] (5.2,0) coordinate(o2);
 \draw (o1) to[R,l={carga},i={$\SI{10}{\milli\ampere}$}] (o2);
\end{circuitikz}}

\resposta{$R_2 = \SI{500}{\ohm}$; $R_1 = \SI{500}{\ohm}$}
\resolucao{\textbf{Ramo de $R_2$ (sangria):} $R_2$ está em paralelo com a carga e
sob \SI{5}{\volt}. Com corrente de sangria $I_2 = \SI{10}{\milli\ampere}$:
\[
R_2 = \frac{U_2}{I_2}=\frac{5}{\pot{10}{-3}}=\SI{500}{\ohm}.
\]
\textbf{Ramo de $R_1$:} por ele passa a soma da corrente de sangria com a da carga,
$I_1 = 10+10 = \SI{20}{\milli\ampere}$, sob a tensão restante
$15-5 = \SI{10}{\volt}$:
\[
R_1=\frac{U_1}{I_1}=\frac{10}{\pot{20}{-3}}=\SI{500}{\ohm}.
\]
\textbf{Por que a sangria?} Se $R_2$ conduzisse muito pouco, a tensão de saída
variaria fortemente com a carga. Uma sangria comparável à corrente da carga
estabiliza a saída, ao custo de mais potência dissipada --- o compromisso típico
de projeto.}
\end{questao}

\begin{questao}[3]{}
Um divisor de tensão com $R_1 = \SI{4}{\kilo\ohm}$ e $R_2 = \SI{2}{\kilo\ohm}$ é
alimentado por $\SI{12}{\volt}$; a saída é tomada sobre $R_2$. Determine a faixa
de valores da carga $R_L$ para que a tensão de saída não caia mais que
$\SI{10}{\percent}$ em relação ao valor em vazio.
\resposta{$R_L \gtrsim \SI{12}{\kilo\ohm}$}
\resolucao{Em vazio, $U_0 = 12\cdot\dfrac{2}{6} = \SI{4}{\volt}$. Queremos
$U \geq 0{,}9\cdot U_0 = \SI{3.6}{\volt}$. Com carga, $R_2$ é substituído por
$R_2' = R_2\parallel R_L$:
\[
U = 12\cdot\frac{R_2'}{R_1+R_2'}\geq 3{,}6.
\]
Isolando, $\dfrac{R_2'}{4000+R_2'}\geq 0{,}3 \Rightarrow R_2'\geq \SI{1714}{\ohm}$
(aprox.). De $R_2' = \dfrac{2000\,R_L}{2000+R_L}\geq 1714$ resulta
$R_L \geq \SI{12}{\kilo\ohm}$ (aproximadamente). Ou seja, a carga precisa ser cerca
de \textbf{6 vezes maior} que $R_2$ para "não pesar" no divisor. Confirma a regra
prática: divisores resistivos só servem para cargas de alta impedância, como a
entrada de um amplificador.}
\end{questao}

\begin{questao}[3]{}
Um sistema de alimentação deve entregar $\SI{5}{\volt}$ a uma carga de
$\SI{100}{\ohm}$, a partir de uma fonte de $\SI{15}{\volt}$, usando um divisor
resistivo. Especifica-se que a tensão na carga não deve cair mais de
$\SI{2}{\percent}$ se a carga variar de $\SI{100}{\ohm}$ a $\SI{200}{\ohm}$.
\begin{itensq}
\item Explique por que um divisor puramente resistivo tem dificuldade em atender
      a essa especificação.
\item Determine a ordem de grandeza que $R_2$ (resistor em paralelo com a carga)
      deveria ter e discuta o custo dessa escolha.
\end{itensq}
\resposta{Ver análise; $R_2 \ll R_L$, com alto consumo}
\resolucao{(a) A tensão de saída de um divisor \emph{depende da carga}, pois
$R_L$ entra em paralelo com $R_2$. Quando $R_L$ varia, o paralelo
$R_2 \parallel R_L$ varia, e a saída varia junto. A regulação só é boa quando
$R_2 \parallel R_L \approx R_2$, isto é, quando $R_2 \ll R_L$.\\
(b) Quantificando: para que a variação de $R_L$ (de \SI{100}{} para
\SI{200}{\ohm}) altere o paralelo em menos de \SI{2}{\percent}, precisamos de
\[
R_2 \parallel 100 \approx R_2 \parallel 200
\quad\Longrightarrow\quad
R_2 \lesssim \SI{4}{\ohm} \ (\text{ordem de } R_L/25).
\]
Com $R_2 = \SI{4}{\ohm}$ e a razão do divisor exigindo
$R_1 = 2R_2 = \SI{8}{\ohm}$, a corrente do divisor seria
$I \approx \dfrac{15}{12} = \SI{1.25}{\ampere}$, dissipando cerca de
$\SI{19}{\watt}$ --- para entregar apenas $\dfrac{5^2}{100} = \SI{0.25}{\watt}$
à carga!\\
\textbf{Conclusão de projeto:} o rendimento seria de aproximadamente
\SI{1}{\percent}. Isso demonstra por que divisores resistivos servem apenas como
\emph{referência de tensão} para entradas de alta impedância, e por que fontes
reais usam \textbf{reguladores} (série ou chaveados), que mantêm a saída
independentemente da carga sem desperdiçar potência.}
\end{questao}

% BEGIN BANCO COMPLEMENTAR Divisor de tensão
% =====================================================================
%  Banco complementar --- Divisor de Tensao  (REVISADO)
%  Substitui a versao gerada por template (8 clones em N2, 9 em N3 e
%  8 questoes conceituais marcadas como N4).
%  O cap04 ja cobre: maior parcela (MC), tres resistores em serie,
%  100 V com 60/40, tres iguais sob 60 V, projeto de 5 V a partir de
%  15 V, saida em vazio e com carga R_L, potenciometro de 10 kohm como
%  divisor, malha unica com cinco resistores, divisor para carga de
%  10 mA, divisor 4k/2k carregado e sistema de 5 V para carga de
%  100 ohm. O EFEITO DE CARGA basico, portanto, ja esta no capitulo --
%  e tambem no banco de circuitos mistos. Aqui ele so reaparece na
%  forma da CONDICAO DE PROJETO (N4) e da nao linearidade do
%  potenciometro carregado, que sao angulos novos.
%  Distribuicao mantida: 4 N1 + 8 N2 + 9 N3 + 8 N4 = 29
% =====================================================================

\subsubsection*{Banco complementar --- Divisor de tensão}

\begin{atencao}[Organização do banco]
A relação $U_o=U_{in}\dfrac{R_2}{R_1+R_2}$ pressupõe \textbf{corrente comum} ---
ou seja, saída em vazio. Todo o resto do tópico decorre de levar a sério essa
hipótese: o divisor tem resistência interna $R_1\parallel R_2$, e qualquer carga
comparável a ela derruba a saída. O N3 explora tomadas múltiplas, sensores e
fontes não ideais; o N4 trata de projeto sob restrição de potência, tolerância
de razão, sonda compensada e a comparação com reguladores.
\end{atencao}

\blocofixacao

\begin{questao}[1]{}
Três resistores de $\SI{1}{\kilo\ohm}$, $\SI{2}{\kilo\ohm}$ e
$\SI{3}{\kilo\ohm}$ estão em série sob $\SI{12}{\volt}$. Determine a tensão
sobre o resistor \emph{intermediário}.
\resposta{$U=\SI{4}{\volt}$}
\resolucao{A regra do divisor vale para qualquer número de elementos em série:
cada um fica com a fração de sua própria resistência sobre o total.
\[
U_2=12\cdot\frac{2}{1+2+3}=12\cdot\frac13=\SI{4}{\volt}.
\]
\textbf{Conferência:} as três tensões são $2$, $4$ e $\SI{6}{\volt}$, somando
$\SI{12}{\volt}$ \checkmark, e guardam a proporção $1{:}2{:}3$ dos
resistores.}
\end{questao}

\begin{questao}[1]{}
Um divisor deve entregar $25\%$ da tensão de entrada. Determine a razão
$R_1/R_2$.
\resposta{$R_1/R_2=3$}
\resolucao{
\[
\frac{R_2}{R_1+R_2}=0{,}25
\;\Longrightarrow\;
R_1+R_2=4R_2
\;\Longrightarrow\;
R_1=3R_2.
\]
\textbf{Observação importante:} apenas a \emph{razão} é fixada --- $3\,$k$/1\,$k
e $30\,$k$/10\,$k produzem a mesma saída em vazio. O que os distingue é a
corrente de repouso e a resistência interna do divisor, e é isso que decide o
projeto quando há carga.}
\end{questao}

\begin{questao}[1]{}
Em um divisor em vazio, ambos os resistores são \textbf{dobrados}. A tensão de
saída:
\begin{alternativas}[2]
\item dobra
\item cai à metade
\item não se altera
\item depende dos valores originais
\end{alternativas}
\resposta{(c)}
\resolucao{A razão $\dfrac{2R_2}{2R_1+2R_2}=\dfrac{R_2}{R_1+R_2}$ é
inalterada --- a saída em vazio depende só da proporção.\\
\textbf{O que muda:} a corrente de repouso cai à metade (menos consumo) e a
resistência interna $R_1\parallel R_2$ \emph{dobra} (mais sensível à carga).
Escalar um divisor troca consumo por robustez, e essa é sempre a decisão
central do projeto.}
\end{questao}

\begin{questao}[1]{}
Analise os casos-limite de um divisor: (a) $R_2\to0$; (b) $R_2\to\infty$.
\resposta{(a) $U_o\to0$; (b) $U_o\to U_{in}$}
\resolucao{(a) Com $R_2=0$, a saída está em \textbf{curto} com o terra:
$U_o=0$. A corrente vale $U_{in}/R_1$ --- toda a potência vai para $R_1$.\\
(b) Com $R_2\to\infty$ (circuito aberto), não circula corrente, não há queda em
$R_1$ e $U_o\to U_{in}$.\\
\textbf{Conclusão:} a saída de um divisor está sempre entre $0$ e $U_{in}$ --- ele
só \emph{reduz} tensão, nunca a eleva. Obter tensão maior que a entrada exige
elemento armazenador (indutor ou capacitor) e comutação.}
\end{questao}

\blocoaplicacao

\begin{questao}[2]{}
Um divisor de $\SI{12}{\volt}$ usa três resistores em série de
$\SI{1}{\kilo\ohm}$, $\SI{1}{\kilo\ohm}$ e $\SI{2}{\kilo\ohm}$, com duas
tomadas intermediárias.
\begin{itensq}
\item Determine as duas tensões de saída em relação ao terra.
\item Determine a corrente de repouso.
\end{itensq}
\resposta{(a) $\SI{9}{\volt}$ e $\SI{6}{\volt}$; (b) $\SI{3}{\milli\ampere}$}
\resolucao{Total: $R=\SI{4}{\kilo\ohm}$ e
$I=12/4000=\SI{3}{\milli\ampere}$.\\
A tomada mais alta tem, abaixo de si, $1+2=\SI{3}{\kilo\ohm}$:
\[
U_A=12\cdot\frac{3}{4}=\SI{9}{\volt}.
\]
A tomada inferior tem $\SI{2}{\kilo\ohm}$ abaixo:
\[
U_B=12\cdot\frac{2}{4}=\SI{6}{\volt}.
\]
\textbf{Cuidado:} cada saída conta apenas a resistência \emph{entre ela e o
terra}, não a do trecho acima. E as duas tomadas se \emph{influenciam}: carregar
uma altera a outra, pois compartilham a mesma cadeia.}
\end{questao}

\begin{questao}[2]{}
Uma fonte de $\SI{30}{\volt}$ alimenta dois resistores iguais em série, e o
ponto médio é adotado como referência (terra). Determine as tensões dos dois
extremos.
\resposta{$\SI{+15}{\volt}$ e $\SI{-15}{\volt}$}
\resolucao{Cada resistor cai $\SI{15}{\volt}$. Tomando o ponto médio como
$\SI{0}{\volt}$, o terminal superior fica em $\SI{+15}{\volt}$ e o inferior em
$\SI{-15}{\volt}$.\\
\textbf{Aplicação:} é assim que se obtém alimentação \emph{simétrica} para
amplificadores operacionais a partir de uma fonte simples. \\
\textbf{Limitação séria:} o "terra" assim criado é um ponto de divisor, com
resistência interna $R/2$. Qualquer corrente drenada de forma assimétrica pelos
dois ramos desloca a referência. Na prática, o divisor precisa ser reforçado por
um seguidor de tensão ou por capacitores de grande valor.}
\end{questao}

\begin{questao}[2]{}
Um divisor de $\SI{24}{\volt}$ tem $R_1=\SI{8}{\kilo\ohm}$ e
$R_2=\SI{4}{\kilo\ohm}$. Determine a saída, a corrente de repouso e a potência
total dissipada.
\resposta{$\SI{8}{\volt}$; $\SI{2}{\milli\ampere}$; $\SI{48}{\milli\watt}$}
\resolucao{
\[
U_o=24\cdot\frac{4}{12}=\SI{8}{\volt};
\qquad
I=\frac{24}{12000}=\SI{2}{\milli\ampere};
\]
\[
P=U_{in}I=24(\pot{2}{-3})=\SI{48}{\milli\watt}.
\]
\textbf{Distribuição:} $P_1=8000(\pot{2}{-3})^{2}=\SI{32}{\milli\watt}$ e
$P_2=\SI{16}{\milli\watt}$ --- proporcionais às resistências, como em qualquer
série.\\
Note que essa potência é consumida \emph{permanentemente}, mesmo sem carga
alguma conectada. Em equipamentos alimentados por bateria, divisores de repouso
são uma fonte silenciosa de descarga.}
\end{questao}

\begin{questao}[2]{}
Um divisor tem $R_1=\SI{2}{\kilo\ohm}$ fixo e $R_2$ ajustável de
$\SI{0}{\ohm}$ a $\SI{8}{\kilo\ohm}$, sob $\SI{10}{\volt}$. Determine a faixa
da tensão de saída.
\resposta{de $\SI{0}{\volt}$ a $\SI{8}{\volt}$}
\resolucao{
\[
R_2=0 \Rightarrow U_o=0;
\qquad
R_2=\SI{8}{\kilo\ohm} \Rightarrow U_o=10\cdot\frac{8}{10}=\SI{8}{\volt}.
\]
\textbf{Não linearidade do ajuste:} no meio do curso
($R_2=\SI{4}{\kilo\ohm}$), a saída é $10\cdot4/6=\SI{6.67}{\volt}$ --- e não
$\SI{4}{\volt}$, como uma escala linear sugeriria. A relação
$U_o(R_2)$ é uma hipérbole, não uma reta.\\
Para obter ajuste \emph{linear}, é preciso usar um potenciômetro (em que a soma
$R_1+R_2$ permanece constante), e não um reostato em série com resistor fixo.}
\end{questao}

\begin{questao}[2]{}
Um divisor deve entregar $\SI{4.5}{\volt}$ com $R_2=\SI{3}{\kilo\ohm}$.
Determine $U_{in}$ para $R_1=\SI{5}{\kilo\ohm}$.
\resposta{$U_{in}=\SI{12}{\volt}$}
\resolucao{
\[
4{,}5=U_{in}\cdot\frac{3}{5+3}
\;\Longrightarrow\;
U_{in}=4{,}5\cdot\frac{8}{3}=\SI{12}{\volt}.
\]
\textbf{Verificação:} $12\cdot3/8=\SI{4.5}{\volt}$ \checkmark, e a corrente é
$12/8000=\SI{1.5}{\milli\ampere}$.\\
A relação do divisor pode ser resolvida em qualquer das quatro variáveis ---
$U_{in}$, $U_o$, $R_1$ ou $R_2$ --- desde que as outras três sejam conhecidas.}
\end{questao}

\begin{questao}[2]{}
Um termistor NTC vale $\SI{10}{\kilo\ohm}$ a $\SI{25}{\celsius}$ e
$\SI{4}{\kilo\ohm}$ a $\SI{50}{\celsius}$. Ele ocupa a posição de $R_2$ em um
divisor com $R_1=\SI{10}{\kilo\ohm}$ sob $\SI{5}{\volt}$. Determine a saída nas
duas temperaturas.
\resposta{$\SI{2.5}{\volt}$ e $\SI{1.43}{\volt}$}
\resolucao{
\[
\SI{25}{\celsius}:\ U_o=5\cdot\frac{10}{20}=\SI{2.5}{\volt};
\qquad
\SI{50}{\celsius}:\ U_o=5\cdot\frac{4}{14}=\SI{1.43}{\volt}.
\]
Variação de $\SI{1.07}{\volt}$ em $\SI{25}{\celsius}$, ou cerca de
$\SI{43}{\milli\volt\per\celsius}$ na média --- sinal amplo e fácil de digitalizar.\\
\textbf{Ressalva:} a relação \emph{não} é linear, pois nem o termistor é linear
em $T$, nem o divisor é linear em $R_2$. Para converter tensão em temperatura é
preciso tabela, polinômio ou a equação de Steinhart--Hart. A escolha de
$R_1=\SI{10}{\kilo\ohm}$ não é casual --- ela maximiza a sensibilidade em torno
de $\SI{25}{\celsius}$, como se demonstra no bloco de desafio.}
\end{questao}

\begin{questao}[2]{}
Um LDR varia de $\SI{1}{\kilo\ohm}$ (claro) a $\SI{100}{\kilo\ohm}$ (escuro).
Ele é colocado como $R_1$ de um divisor com $R_2=\SI{10}{\kilo\ohm}$ sob
$\SI{5}{\volt}$. Determine a faixa de saída e o sentido da variação.
\resposta{de $\SI{4.55}{\volt}$ (claro) a $\SI{0.45}{\volt}$ (escuro)}
\resolucao{
\[
\text{claro}:\ U_o=5\cdot\frac{10}{11}=\SI{4.55}{\volt};
\qquad
\text{escuro}:\ U_o=5\cdot\frac{10}{110}=\SI{0.45}{\volt}.
\]
Com o LDR na posição \emph{superior}, mais luz significa \emph{maior} tensão de
saída.\\
\textbf{Inversão do comportamento:} trocando as posições (LDR embaixo), a saída
passaria de $\SI{0.45}{\volt}$ no claro a $\SI{4.55}{\volt}$ no escuro. A
escolha da posição define a polaridade do sensor --- decisão de projeto que
custa nada e evita inverter o sinal depois em software.}
\end{questao}

\begin{questao}[2]{}
Um divisor projetado com $R_1=\SI{8.7}{\kilo\ohm}$ e $R_2=\SI{3.3}{\kilo\ohm}$
para $12\to\SI{3.3}{\volt}$ deve usar valores E24. Adotando
$R_1=\SI{9.1}{\kilo\ohm}$, determine a saída real e o erro.
\resposta{$U_o=\SI{3.19}{\volt}$; erro de $-3{,}2\%$}
\resolucao{
\[
U_o=12\cdot\frac{3{,}3}{9{,}1+3{,}3}=12\cdot\frac{3{,}3}{12{,}4}=\SI{3.19}{\volt};
\qquad
\text{erro}=\frac{3{,}19-3{,}30}{3{,}30}=-3{,}2\%.
\]
\textbf{Alternativa melhor:} o par $\SI{10}{\kilo\ohm}/\SI{3.9}{\kilo\ohm}$ dá
$U_o=\SI{3.37}{\volt}$, erro de apenas $+2{,}0\%$. Vale sempre testar
\emph{vários} pares E24 antes de escolher --- a razão importa mais que os
valores individuais.\\
\textbf{Se $3\%$ for demais:} use a série E96 (tolerância $1\%$), que contém
$\SI{8.66}{\kilo\ohm}$ e resolve o problema diretamente.}
\end{questao}

\blocoaprofundamento

\begin{questao}[3]{}
Um divisor $R_1=\SI{6}{\kilo\ohm}$, $R_2=\SI{3}{\kilo\ohm}$ é alimentado por
uma fonte \textbf{real} de $\SI{12}{\volt}$ com $r=\SI{1}{\kilo\ohm}$.
\begin{itensq}
\item Determine a saída, considerando $r$.
\item Compare com a previsão que ignora $r$.
\item Determine a resistência de Thévenin vista na saída.
\end{itensq}
\resposta{(a) $\SI{3.6}{\volt}$; (b) previsão de $\SI{4}{\volt}$, erro de
$10\%$; (c) $R_{Th}=\SI{2.1}{\kilo\ohm}$}
\resolucao{(a) A resistência interna soma-se a $R_1$:
\[
U_o=12\cdot\frac{3}{1+6+3}=12\cdot\frac{3}{10}=\SI{3.6}{\volt}.
\]
(b) Ignorando $r$: $12\cdot3/9=\SI{4}{\volt}$. O erro é de $10\%$ --- nada
desprezível.\\
(c) Zerando a f.e.m. (mas mantendo $r$), a saída enxerga $R_2$ em paralelo com
$(r+R_1)$:
\[
R_{Th}=\frac{3\cdot7}{10}=\SI{2.1}{\kilo\ohm}.
\]
\textbf{Leitura:} a resistência interna da fonte se comporta como parte de
$R_1$ --- ela piora a razão \emph{e} aumenta a impedância de saída. Em divisores
de baixa impedância (resistores de dezenas de ohms), $r$ pode dominar
completamente o projeto.}
\end{questao}

\begin{questao}[3]{}
Dois divisores idênticos ($R_1=\SI{3}{\kilo\ohm}$, $R_2=\SI{1}{\kilo\ohm}$)
são alimentados pela mesma fonte de $\SI{8}{\volt}$, mas em um deles $R_2$ é
trocado por $\SI{1.1}{\kilo\ohm}$.
\begin{itensq}
\item Determine as duas saídas e a diferença entre elas.
\item Determine a sensibilidade da diferença a variações de $R_2$.
\end{itensq}
\resposta{(a) $\SI{2}{\volt}$ e $\SI{2.146}{\volt}$; diferença de
$\SI{146}{\milli\volt}$}
\resolucao{(a)
\[
U_A=8\cdot\frac{1}{4}=\SI{2}{\volt};
\qquad
U_B=8\cdot\frac{1{,}1}{4{,}1}=\SI{2.146}{\volt};
\qquad
\Delta U=\SI{146}{\milli\volt}.
\]
(b) Derivando $U=U_{in}R_2/(R_1+R_2)$:
\[
\frac{\mathrm{d}U}{\mathrm{d}R_2}
=\frac{U_{in}R_1}{(R_1+R_2)^{2}}
=\frac{8(3000)}{(4000)^{2}}=\SI{1.5}{\milli\volt\per\ohm}.
\]
Para $\Delta R_2=\SI{100}{\ohm}$, prevê-se $\SI{150}{\milli\volt}$ --- muito
próximo dos $\SI{146}{\milli\volt}$ exatos, pois a variação é pequena.\\
\textbf{Esta é a estrutura da ponte:} dois divisores comparados, em que o
\emph{desequilíbrio} carrega a informação. A grande vantagem é que variações
comuns aos dois ramos --- deriva da fonte, temperatura ambiente --- se cancelam
na diferença.}
\end{questao}

\begin{questao}[3]{}
Um divisor $R_1=R_2=\SI{100}{\kilo\ohm}$ sob $\SI{10}{\volt}$ é medido por um
voltímetro cuja resistência de entrada é $\SI{1}{\mega\ohm}$.
\begin{itensq}
\item Determine a leitura do instrumento.
\item Determine o erro percentual.
\item Determine a resistência mínima do voltímetro para erro inferior a
      $1\%$.
\end{itensq}
\resposta{(a) $\SI{4.76}{\volt}$; (b) $-4{,}8\%$; (c)
$\ge\SI{5}{\mega\ohm}$}
\resolucao{(a) O voltímetro entra em paralelo com $R_2$:
\[
R_2'=\frac{100\cdot1000}{1100}=\SI{90.9}{\kilo\ohm};
\qquad
U=10\cdot\frac{90{,}9}{100+90{,}9}=\SI{4.76}{\volt}.
\]
(b) Erro: $(4{,}76-5{,}00)/5{,}00=-4{,}8\%$.\\
(c) A resistência interna do divisor é
$R_{Th}=100\parallel100=\SI{50}{\kilo\ohm}$. O erro relativo é
aproximadamente $R_{Th}/R_V$, então
\[
\frac{50\,\mathrm{k}}{R_V}\le0{,}01
\;\Longrightarrow\;
R_V\ge\SI{5}{\mega\ohm}.
\]
\textbf{Lição de instrumentação:} o instrumento \emph{perturba} o circuito que
mede. Multímetros digitais comuns têm $\SI{10}{\mega\ohm}$ --- suficientes aqui,
mas insuficientes para divisores de megaohms. E note que a especificação
depende de $R_{Th}$ do circuito, não da tensão medida.}
\end{questao}

\begin{questao}[3]{}
Um divisor $R_1=\SI{4}{\kilo\ohm}$, $R_2=\SI{1}{\kilo\ohm}$ sob
$\SI{10}{\volt}$ fornece referência a um conversor A/D de $10$ bits com fundo
de escala de $\SI{2.5}{\volt}$.
\begin{itensq}
\item Determine a saída e a margem em relação ao fundo de escala.
\item Determine o degrau de quantização.
\item Compare o erro de quantização com o erro de resistores de $1\%$.
\end{itensq}
\resposta{(a) $\SI{2}{\volt}$, $80\%$ do fundo; (b) $\SI{2.44}{\milli\volt}$;
(c) resistores dominam ($\approx\SI{16}{\milli\volt}$)}
\resolucao{(a) $U_o=10\cdot1/5=\SI{2}{\volt}$, ou $80\%$ do fundo de escala ---
boa utilização da faixa, sem risco de saturação.\\
(b) Com $10$ bits, $2^{10}=1024$ degraus:
\[
\Delta=\frac{2{,}5}{1024}=\SI{2.44}{\milli\volt}.
\]
(c) Com resistores de $1\%$, o pior erro de razão vale
$\dfrac{R_1}{R_1+R_2}(t_1+t_2)=0{,}8(2\%)=1{,}6\%$, ou
$0{,}016\times2=\SI{32}{\milli\volt}$ --- \textbf{treze vezes} o degrau de
quantização.\\
\textbf{Conclusão de projeto:} não adianta usar um conversor de $12$ ou
$16$ bits com esse divisor. O gargalo é o divisor, e a solução é resistores de
$0{,}1\%$, calibração por software, ou uma referência de tensão dedicada.
Aumentar a resolução do conversor sem corrigir o divisor é gastar dinheiro em
dígitos sem significado.}
\end{questao}

\begin{questao}[3]{}
Um divisor de $\SI{48}{\volt}$ para $\SI{12}{\volt}$ deve alimentar uma carga
de $\SI{100}{\ohm}$.
\begin{itensq}
\item Determine $R_1$ e $R_2$ ignorando a carga, com corrente de repouso de
      $\SI{120}{\milli\ampere}$.
\item Determine a saída real com a carga conectada.
\item Avalie a viabilidade.
\end{itensq}
\resposta{(a) $R_1=\SI{300}{\ohm}$, $R_2=\SI{100}{\ohm}$;
(b) $\SI{6.86}{\volt}$; (c) inviável}
\resolucao{(a) $R_2=12/0{,}120=\SI{100}{\ohm}$;
$R_1=(48-12)/0{,}120=\SI{300}{\ohm}$.\\
(b) A carga de $\SI{100}{\ohm}$ fica em paralelo com $R_2$:
\[
R_2'=\frac{100\cdot100}{200}=\SI{50}{\ohm};
\qquad
U_o=48\cdot\frac{50}{350}=\SI{6.86}{\volt}.
\]
Queda de $43\%$ em relação ao projetado.\\
(c) \textbf{Inviável.} A carga consome $\SI{120}{\milli\ampere}$ a
$\SI{12}{\volt}$ --- exatamente a corrente de repouso, e não uma fração dela. A
regra de bolso exige repouso de $10$ a $20$ vezes a corrente da carga, o que
levaria $R_2$ a $\SI{10}{\ohm}$ e a dissipação total a
$48\times1{,}2=\SI{57.6}{\watt}$ para entregar $\SI{1.44}{\watt}$ --- rendimento
de $2{,}5\%$.\\
\textbf{Diagnóstico geral:} divisores resistivos servem para gerar
\emph{referências} de alta impedância, não para \emph{alimentar} cargas. Aqui o
caso pede regulador.}
\end{questao}

\begin{questao}[3]{}
Um divisor com $R_1=\SI{10}{\kilo\ohm}$ e $R_2=\SI{10}{\kilo\ohm}$ opera sob
$\SI{200}{\volt}$.
\begin{itensq}
\item Determine a potência de cada resistor e a tensão sobre cada um.
\item Especifique os resistores considerando potência \emph{e} tensão máxima de
      trabalho.
\end{itensq}
\resposta{$\SI{1}{\watt}$ e $\SI{100}{\volt}$ em cada; exige resistor de
potência com tensão de trabalho adequada}
\resolucao{(a) $I=200/20000=\SI{10}{\milli\ampere}$;
\[
U_1=U_2=\SI{100}{\volt};
\qquad
P_1=P_2=10^{4}(10^{-2})^{2}=\SI{1}{\watt}.
\]
(b) \textbf{Dois critérios independentes:}
\begin{itemize}
\item \textbf{Potência:} com fator $2$ de folga, resistores de
$\SI{2}{\watt}$.
\item \textbf{Tensão máxima de trabalho:} um resistor de filme comum de
$\frac14\,\si{\watt}$ suporta cerca de $\SI{250}{\volt}$; os de
$\SI{2}{\watt}$ chegam a $\SI{500}{\volt}$. Os $\SI{100}{\volt}$ exigidos são
atendidos com folga.
\end{itemize}
\textbf{Por que o segundo critério existe:} a tensão máxima é limitada pela
ruptura do revestimento e pelo arco entre as espiras do filme --- fenômeno que
independe da potência. Em divisores de \emph{alta} tensão (quilovolts), é comum
que a tensão, e não a potência, force o uso de vários resistores em série,
mesmo com dissipação irrisória.}
\end{questao}

\begin{questao}[3]{}
Um divisor deve fornecer $\SI{5}{\volt}$ a partir de $\SI{15}{\volt}$ para uma
entrada de microcontrolador com corrente de entrada de
$\SI{1}{\micro\ampere}$.
\begin{itensq}
\item Determine $R_1$ e $R_2$ para corrente de repouso de
      $\SI{10}{\micro\ampere}$.
\item Verifique o erro introduzido pela corrente de entrada.
\item Comente o compromisso com o consumo.
\end{itensq}
\resposta{(a) $R_1=\SI{1}{\mega\ohm}$, $R_2=\SI{500}{\kilo\ohm}$;
(b) erro de $\approx3\%$}
\resolucao{(a) $R_2=5/\pot{10}{-6}=\SI{500}{\kilo\ohm}$;
$R_1=10/\pot{10}{-6}=\SI{1}{\mega\ohm}$.\\
(b) A corrente de entrada de $\SI{1}{\micro\ampere}$ é $10\%$ da corrente de
repouso --- longe de desprezível. Modelando a entrada como resistência
equivalente $R_{in}=5/\pot{1}{-6}=\SI{5}{\mega\ohm}$:
\[
R_2'=\frac{500\cdot5000}{5500}=\SI{454.5}{\kilo\ohm};
\qquad
U_o=15\cdot\frac{454{,}5}{1454{,}5}=\SI{4.69}{\volt},
\]
erro de $-6{,}2\%$.\\
(c) \textbf{Compromisso.} Reduzir o erro exige aumentar a corrente de repouso ---
com $\SI{100}{\micro\ampere}$ ($R_1=\SI{100}{\kilo\ohm}$,
$R_2=\SI{50}{\kilo\ohm}$), o erro cai para cerca de $0{,}7\%$, mas o consumo em
repouso sobe dez vezes.\\
Em um dispositivo alimentado por bateria de $\SI{1000}{\milli\ampere\hour}$,
$\SI{100}{\micro\ampere}$ contínuos consomem a bateria em pouco mais de um ano
\emph{só com o divisor}. É o tipo de detalhe que decide a autonomia de um
produto.}
\end{questao}

\begin{questao}[3]{}
Determine, em um divisor, a expressão da resistência de Thévenin vista na saída
e avalie-a para $R_1=\SI{6}{\kilo\ohm}$ e $R_2=\SI{3}{\kilo\ohm}$.
\begin{itensq}
\item Deduza $R_{Th}$.
\item Mostre que $R_{Th}<\min(R_1,R_2)$.
\item Explique seu papel no carregamento.
\end{itensq}
\resposta{$R_{Th}=R_1\parallel R_2=\SI{2}{\kilo\ohm}$}
\resolucao{(a) Zerando a fonte de entrada (curto), a saída enxerga $R_1$ e
$R_2$ em paralelo:
\[
R_{Th}=\frac{R_1R_2}{R_1+R_2}=\frac{6\cdot3}{9}=\SI{2}{\kilo\ohm}.
\]
(b) O paralelo de dois resistores positivos é sempre menor que o menor deles ---
resultado geral já estabelecido para associações. Aqui,
$\SI{2}{\kilo\ohm}<\SI{3}{\kilo\ohm}$ \checkmark.\\
(c) \textbf{Papel no carregamento.} Conectada uma carga $R_L$, o divisor se
comporta como fonte de Thévenin:
\[
U_L=U_{Th}\cdot\frac{R_L}{R_L+R_{Th}}.
\]
O erro relativo é aproximadamente $R_{Th}/R_L$ para $R_L\gg R_{Th}$. Logo:
\begin{itemize}
\item $R_L=10R_{Th}$ $\Rightarrow$ erro $\approx9\%$;
\item $R_L=100R_{Th}$ $\Rightarrow$ erro $\approx1\%$.
\end{itemize}
\textbf{Consequência de projeto:} o único parâmetro que importa para o
carregamento é $R_{Th}$ --- não $R_1$, não $R_2$, não a razão. Reduzir os dois
resistores proporcionalmente reduz $R_{Th}$ e melhora a robustez, ao custo de
consumo.}
\end{questao}

\begin{questao}[3]{}
Um divisor alimenta uma carga que \textbf{varia} entre $\SI{10}{\kilo\ohm}$ e
$\SI{100}{\kilo\ohm}$. O divisor tem $R_{Th}=\SI{2}{\kilo\ohm}$ e
$U_{Th}=\SI{5}{\volt}$.
\begin{itensq}
\item Determine a faixa de tensão de saída.
\item Determine a regulação percentual.
\item Indique como melhorá-la.
\end{itensq}
\resposta{(a) de $\SI{4.17}{\volt}$ a $\SI{4.90}{\volt}$;
(b) $\approx15\%$}
\resolucao{(a)
\[
R_L=\SI{10}{\kilo\ohm}:\ U=5\cdot\frac{10}{12}=\SI{4.17}{\volt};
\qquad
R_L=\SI{100}{\kilo\ohm}:\ U=5\cdot\frac{100}{102}=\SI{4.90}{\volt}.
\]
(b) Regulação:
\[
\frac{4{,}90-4{,}17}{4{,}90}=15{,}0\%.
\]
(c) \textbf{Como melhorar.}
\begin{itemize}
\item \textbf{Reduzir $R_{Th}$.} Com $R_{Th}=\SI{200}{\ohm}$, a faixa passa a
$4{,}90$--$\SI{4.99}{\volt}$ e a regulação cai para $1{,}8\%$ --- ao custo de
dez vezes mais consumo de repouso.
\item \textbf{Inserir um seguidor de tensão} (amp-op em configuração de ganho
unitário). A impedância de saída cai para frações de ohm e a regulação
praticamente desaparece, sem aumentar o consumo do divisor.
\item \textbf{Usar regulador}, se a carga também exigir corrente apreciável.
\end{itemize}
\textbf{Observação:} a segunda opção é a mais elegante --- ela desacopla
\emph{razão} de \emph{impedância}, que no divisor puro estão inevitavelmente
amarradas.}
\end{questao}

\blocodesafio

\begin{questao}[4]{}
\textbf{(Condição de carregamento.)} Deduza a condição para que uma carga
$R_L$ altere a saída de um divisor em menos de um erro relativo $\varepsilon$.
\begin{itensq}
\item Obtenha a expressão exata do erro.
\item Deduza a condição aproximada para $\varepsilon$ pequeno.
\item Construa a tabela de $R_L/R_{Th}$ para
      $\varepsilon=10\%$, $1\%$ e $0{,}1\%$.
\end{itensq}
\resposta{$\varepsilon=\dfrac{R_{Th}}{R_L+R_{Th}}$; aproximadamente
$R_L\ge R_{Th}/\varepsilon$}
\resolucao{(a) Com carga,
$U_L=U_{Th}\dfrac{R_L}{R_L+R_{Th}}$, logo
\[
\varepsilon=\frac{U_{Th}-U_L}{U_{Th}}=1-\frac{R_L}{R_L+R_{Th}}
=\frac{R_{Th}}{R_L+R_{Th}}.
\]
(b) Para $R_L\gg R_{Th}$, o denominador tende a $R_L$:
\[
\varepsilon\approx\frac{R_{Th}}{R_L}
\;\Longrightarrow\;
R_L\ge\frac{R_{Th}}{\varepsilon}.
\]
(c)
\begin{center}
\begin{tabular}{ccc}
\hline
$\varepsilon$ & $R_L/R_{Th}$ (exato) & $R_L/R_{Th}$ (aprox.)\\
\hline
$10\%$ & $9$ & $10$\\
$1\%$ & $99$ & $100$\\
$0{,}1\%$ & $999$ & $1000$\\
\hline
\end{tabular}
\end{center}
A aproximação erra por exatamente $1$ unidade e é conservadora --- pode ser
usada sem receio.\\
\textbf{Regra de bolso resultante:} para $1\%$ de erro,
$R_L\ge100\,R_{Th}$. Note que $R_{Th}=R_1\parallel R_2$ é sempre \emph{menor}
que qualquer um dos resistores, o que torna a condição menos severa do que
comparar $R_L$ com $R_2$ diretamente --- erro comum que leva a
superdimensionar o divisor.}
\end{questao}

\begin{questao}[4]{}
\textbf{(Projeto sob restrição de potência.)} Projete um divisor que reduza
$\SI{100}{\volt}$ a $\SI{5}{\volt}$, dissipando no máximo
$\SI{0.5}{\watt}$.
\begin{itensq}
\item Determine $R_1$ e $R_2$.
\item Determine a potência e a tensão de cada resistor.
\item Determine a maior carga que pode ser atendida com erro de até $1\%$.
\end{itensq}
\resposta{(a) $R_1=\SI{19}{\kilo\ohm}$, $R_2=\SI{1}{\kilo\ohm}$;
(b) $0{,}475$ e $\SI{0.025}{\watt}$; (c) $R_L\ge\SI{95}{\kilo\ohm}$}
\resolucao{(a) A restrição de potência fixa a \emph{soma}:
\[
P=\frac{U_{in}^{2}}{R_1+R_2}\le0{,}5
\;\Longrightarrow\;
R_1+R_2\ge\frac{100^{2}}{0{,}5}=\SI{20}{\kilo\ohm}.
\]
A razão fixa a \emph{proporção}: $R_2/(R_1+R_2)=0{,}05$. Adotando o mínimo,
\[
R_2=\SI{1}{\kilo\ohm};\qquad R_1=\SI{19}{\kilo\ohm}.
\]
(b) $I=100/20000=\SI{5}{\milli\ampere}$;
\[
P_1=19000(\pot{5}{-3})^{2}=\SI{0.475}{\watt};
\qquad
P_2=\SI{0.025}{\watt};
\]
com $U_1=\SI{95}{\volt}$ e $U_2=\SI{5}{\volt}$.\\
\textbf{Especificação:} $R_1$ concentra $95\%$ da dissipação --- use
$\SI{1}{\watt}$ nele e $\frac18\,\si{\watt}$ em $R_2$. Alternativamente,
dividir $R_1$ em dois de $\SI{9.5}{\kilo\ohm}$ reparte o calor e a tensão.\\
(c) $R_{Th}=19\parallel1=\SI{950}{\ohm}$; pela condição anterior,
\[
R_L\ge100\,R_{Th}=\SI{95}{\kilo\ohm}.
\]
\textbf{Tensão de projeto:} o divisor entrega $\SI{5}{\volt}$ a uma carga que
consome no máximo $\SI{53}{\micro\ampere}$ --- serve para um conversor A/D, não
para alimentar nada.}
\end{questao}

\begin{questao}[4]{}
\textbf{(Potenciômetro carregado.)} Um potenciômetro de $\SI{10}{\kilo\ohm}$
é usado como divisor sob $\SI{10}{\volt}$, com carga
$R_L=\SI{10}{\kilo\ohm}$ na saída. Seja $x$ a fração do curso.
\begin{itensq}
\item Deduza $U_o(x)$ com a carga.
\item Avalie em $x=0{,}25$, $0{,}5$ e $0{,}75$ e compare com o caso ideal.
\item Determine onde o desvio é máximo e como reduzi-lo.
\end{itensq}
\resposta{Em $x=0{,}5$: $\SI{4.00}{\volt}$ contra $\SI{5}{\volt}$ ideal ---
desvio de $20\%$}
\resolucao{(a) A porção inferior vale $xR_p$ e fica em paralelo com $R_L$; a
superior vale $(1-x)R_p$:
\[
U_o=U_{in}\,\frac{xR_p\parallel R_L}{(1-x)R_p+\left(xR_p\parallel R_L\right)}.
\]
Com $R_p=R_L=\SI{10}{\kilo\ohm}$, tem-se
$xR_p\parallel R_L=\dfrac{10x}{1+x}\,\si{\kilo\ohm}$ e
\[
U_o=10\cdot\frac{\dfrac{10x}{1+x}}{10(1-x)+\dfrac{10x}{1+x}}
=10\cdot\frac{x}{(1-x)(1+x)+x}
=\frac{10x}{1+x-x^{2}}.
\]
(b)
\begin{center}
\begin{tabular}{cccc}
\hline
$x$ & ideal & com carga & desvio\\
\hline
$0{,}25$ & $\SI{2.50}{\volt}$ & $\SI{2.105}{\volt}$ & $-15{,}8\%$\\
$0{,}50$ & $\SI{5.00}{\volt}$ & $\SI{4.000}{\volt}$ & $-20{,}0\%$\\
$0{,}75$ & $\SI{7.50}{\volt}$ & $\SI{6.316}{\volt}$ & $-15{,}8\%$\\
\hline
\end{tabular}
\end{center}
(c) O desvio \emph{absoluto} é máximo perto do meio do curso, onde a
resistência de Thévenin do potenciômetro ($x(1-x)R_p$) atinge seu máximo,
$R_p/4=\SI{2.5}{\kilo\ohm}$. Nos extremos ($x\to0$ ou $x\to1$) a impedância
tende a zero e a carga deixa de importar.\\
\textbf{A curva deixa de ser reta:} o potenciômetro linear, carregado, comporta-se
de forma \textbf{não linear}, e o erro depende da posição --- o que é
inaceitável em um controle graduado.\\
\textbf{Correções:} \emph{(i)} usar $R_L\ge100R_p$, o que reduz o desvio máximo
a menos de $0{,}25\%$; \emph{(ii)} inserir um seguidor de tensão entre o cursor
e a carga, solução padrão em instrumentação; \emph{(iii)} usar potenciômetro de
lei "logarítmica reversa" pré-distorcida --- gambiarra que só funciona para uma
carga específica.}
\end{questao}

\begin{questao}[4]{}
\textbf{(Tolerância de razão.)} Um divisor usa resistores de $1\%$ com entrada
exata.
\begin{itensq}
\item Deduza a expressão do pior erro relativo da saída.
\item Avalie para as razões $12\to\SI{6}{\volt}$, $12\to\SI{5}{\volt}$ e
      $100\to\SI{5}{\volt}$.
\item Explique por que resistores \emph{casados} melhoram drasticamente o
      resultado.
\end{itensq}
\resposta{$\varepsilon=\dfrac{R_1}{R_1+R_2}(t_1+t_2)$; $1{,}00\%$,
$1{,}17\%$ e $1{,}90\%$}
\resolucao{(a) De $U_o=U_{in}\dfrac{R_2}{R_1+R_2}$, tomando logaritmos e
diferenciando:
\[
\frac{\Delta U_o}{U_o}
=\frac{R_1}{R_1+R_2}\left(\frac{\Delta R_2}{R_2}-\frac{\Delta R_1}{R_1}\right),
\]
cujo módulo máximo (desvios opostos) é
$\varepsilon=\dfrac{R_1}{R_1+R_2}(t_1+t_2)$.\\
(b) O fator $\dfrac{R_1}{R_1+R_2}$ vale $1-U_o/U_{in}$:
\begin{center}
\begin{tabular}{ccc}
\hline
Razão & $R_1/(R_1+R_2)$ & Pior erro\\
\hline
$12\to6$ & $0{,}500$ & $1{,}00\%$\\
$12\to5$ & $0{,}583$ & $1{,}17\%$\\
$100\to5$ & $0{,}950$ & $1{,}90\%$\\
\hline
\end{tabular}
\end{center}
\textbf{Padrão:} quanto \emph{maior} a razão de divisão, pior o erro --- ele
tende a $t_1+t_2=2\%$ quando $U_o\ll U_{in}$, e nunca o ultrapassa.\\
(c) \textbf{Resistores casados.} Em uma rede resistiva integrada, os dois
resistores são fabricados lado a lado, no mesmo processo e no mesmo substrato.
Seus desvios são \textbf{correlacionados}: se um sai $0{,}8\%$ acima do
nominal, o outro também sai. O termo entre parênteses,
$\dfrac{\Delta R_2}{R_2}-\dfrac{\Delta R_1}{R_1}$, tende a \emph{zero}.\\
Na prática, redes casadas oferecem razão com precisão de $0{,}05\%$ mesmo com
valores absolutos de $\pm20\%$. \textbf{O que importa no divisor é a razão, não
os valores} --- e o casamento é a única forma barata de obtê-la com precisão.
Esse é o princípio que viabiliza conversores A/D e amplificadores de
instrumentação.}
\end{questao}

\begin{questao}[4]{}
\textbf{(Sonda atenuadora.)} Uma sonda $10{:}1$ é formada por
$\SI{9}{\mega\ohm}$ na ponta, em série com a entrada de
$\SI{1}{\mega\ohm}$ do osciloscópio.
\begin{itensq}
\item Verifique a razão de atenuação e a resistência de entrada total.
\item Determine o erro ao medir uma fonte com $R_{Th}=\SI{100}{\kilo\ohm}$,
      com e sem a sonda.
\item Explique o que se perde ao usar a sonda.
\end{itensq}
\resposta{(a) $1/10$ e $\SI{10}{\mega\ohm}$; (b) $1{,}0\%$ com sonda,
$9{,}1\%$ sem}
\resolucao{(a) A entrada do osciloscópio \emph{é} o $R_2$ do divisor:
\[
\frac{U_o}{U_{in}}=\frac{1}{9+1}=\frac{1}{10}\ \checkmark;
\qquad
R_{entrada}=9+1=\SI{10}{\mega\ohm}.
\]
(b) \textbf{Com sonda:} o instrumento apresenta $\SI{10}{\mega\ohm}$ à fonte:
\[
\varepsilon=\frac{100\,\mathrm{k}}{10\,\mathrm{M}+100\,\mathrm{k}}=1{,}0\%.
\]
\textbf{Sem sonda} (cabo direto, $\SI{1}{\mega\ohm}$):
\[
\varepsilon=\frac{100\,\mathrm{k}}{1\,\mathrm{M}+100\,\mathrm{k}}=9{,}1\%.
\]
A sonda reduz o erro de carregamento por um fator $9$.\\
(c) \textbf{O que se perde: amplitude.} O sinal chega ao instrumento dez vezes
menor, o que reduz a relação sinal-ruído em $\SI{20}{\decibel}$ e torna
inviável medir sinais de poucos milivolts.\\
\textbf{O compromisso central da sonda} é exatamente esse: troca-se sensibilidade
por não perturbação. Por isso os osciloscópios trazem seletor $1\times/10\times$
--- sinais pequenos em fontes de baixa impedância pedem $1\times$; sinais grandes
em circuitos de alta impedância pedem $10\times$.\\
\emph{(Em sinais variáveis há ainda a capacitância parasita, que exige o
capacitor de compensação ajustável presente em toda sonda real --- assunto de
corrente alternada.)}}
\end{questao}

\begin{questao}[4]{}
\textbf{(Otimização de sensor.)} Um termistor de resistência $R_t$ ocupa a
posição inferior de um divisor com resistor fixo $R$ e alimentação $U$.
\begin{itensq}
\item Obtenha $\dfrac{\mathrm{d}U_o}{\mathrm{d}R_t}$.
\item Determine o valor de $R$ que maximiza a sensibilidade em um dado ponto de
      operação.
\item Avalie para $U=\SI{5}{\volt}$, $R_t=\SI{10}{\kilo\ohm}$ e coeficiente de
      $\SI{-400}{\ohm\per\celsius}$.
\end{itensq}
\resposta{(b) $R=R_t$; (c) $\SI{-50}{\milli\volt\per\celsius}$}
\resolucao{(a) De $U_o=U\dfrac{R_t}{R+R_t}$:
\[
\frac{\mathrm{d}U_o}{\mathrm{d}R_t}
=U\,\frac{(R+R_t)-R_t}{(R+R_t)^{2}}
=\frac{U\,R}{(R+R_t)^{2}}.
\]
(b) Maximizando em relação a $R$, com $R_t$ fixo:
\[
\frac{\mathrm{d}}{\mathrm{d}R}\left[\frac{R}{(R+R_t)^{2}}\right]
=\frac{(R+R_t)-2R}{(R+R_t)^{3}}
=\frac{R_t-R}{(R+R_t)^{3}}=0
\;\Longrightarrow\;
\boxed{R=R_t}
\]
--- e a saída no ponto ótimo é exatamente $U/2$.\\
\textbf{Mesma estrutura} do resultado de máxima potência: a derivada de
$R/(R+R_t)^{2}$ é a mesma que a de $P_R$ em uma série, e por isso o ótimo cai
em igualdade.\\
(c) Com $R=R_t=\SI{10}{\kilo\ohm}$:
\[
\frac{\mathrm{d}U_o}{\mathrm{d}R_t}=\frac{5(10^{4})}{(\pot{2}{4})^{2}}
=\SI{1.25e-4}{\volt\per\ohm},
\]
\[
\frac{\mathrm{d}U_o}{\mathrm{d}T}
=(\pot{1.25}{-4})(-400)=\SI{-50}{\milli\volt\per\celsius}.
\]
\textbf{Excelente sensibilidade:} um conversor A/D de $10$ bits com fundo de
$\SI{5}{\volt}$ tem degrau de $\SI{4.9}{\milli\volt}$, resolvendo cerca de
$\SI{0.1}{\celsius}$.\\
\textbf{Ressalva:} o ótimo vale para \emph{um} ponto de operação. Em faixa
ampla de temperatura, $R_t$ varia por ordens de grandeza e a sensibilidade cai
nas extremidades --- escolha $R$ igual a $R_t$ na temperatura de maior interesse,
não na média da faixa.}
\end{questao}

\begin{questao}[4]{}
\textbf{(Divisor contra regulador.)} Compare as duas soluções para obter
$\SI{5}{\volt}$ a partir de $\SI{12}{\volt}$, para uma carga que varia de
$\SI{1}{\milli\ampere}$ a $\SI{10}{\milli\ampere}$.
\begin{itensq}
\item Dimensione o divisor com repouso de dez vezes a carga máxima e determine
      a regulação.
\item Determine o rendimento de cada solução.
\item Estabeleça o critério de escolha.
\end{itensq}
\resposta{(a) $R_1=\SI{70}{\ohm}$, $R_2=\SI{50}{\ohm}$; regulação
$\approx5{,}3\%$; (b) $3{,}8\%$ contra $42\%$}
\resolucao{(a) Repouso de $\SI{100}{\milli\ampere}$:
$R_2=5/0{,}1=\SI{50}{\ohm}$ e
$R_1=7/0{,}1=\SI{70}{\ohm}$.
\[
R_{Th}=50\parallel70=\SI{29.2}{\ohm};
\qquad
\Delta U=R_{Th}\,\Delta I_L=29{,}2(\pot{9}{-3})=\SI{0.26}{\volt},
\]
ou seja, cerca de $5{,}3\%$ de variação.\\
(b) \textbf{Divisor:} consumo total
$\approx12\times0{,}11=\SI{1.32}{\watt}$ para entregar no máximo
$5\times0{,}01=\SI{50}{\milli\watt}$ --- rendimento de $3{,}8\%$.\\
\textbf{Regulador linear:} consome a mesma corrente da carga,
$12\times0{,}01=\SI{120}{\milli\watt}$ para entregar
$\SI{50}{\milli\watt}$ --- rendimento de $41{,}7\%$ (igual à razão
$U_o/U_{in}$), com regulação típica melhor que $0{,}1\%$.\\
(c) \textbf{Critério.}
\begin{itemize}
\item \textbf{Divisor} apenas quando a carga é de \emph{altíssima impedância} e
praticamente constante: entradas de conversores A/D, polarização de bases,
referências internas. Vantagens: dois componentes, custo desprezível, sem ruído
de comutação.
\item \textbf{Regulador} sempre que a carga consumir corrente apreciável ou
variar. O regulador linear resolve com rendimento $U_o/U_{in}$; o chaveado
ultrapassa $90\%$, ao custo de complexidade e ruído.
\end{itemize}
\textbf{Fronteira prática:} se a corrente da carga for maior que cerca de
$1\%$ da corrente de repouso que o divisor exigiria, o divisor já não é a
escolha certa.}
\end{questao}

\begin{questao}[4]{}
\textbf{(Divisor com proteção.)} Projete a interface entre um sinal de entrada
que varia de $\SI{0}{\volt}$ a $\SI{30}{\volt}$ e um conversor A/D de
$\SI{3.3}{\volt}$ e impedância de entrada muito alta.
\begin{itensq}
\item Determine a razão e escolha os resistores.
\item Verifique o comportamento em sobretensão de $\SI{50}{\volt}$ e proponha
      proteção.
\item Determine o consumo e discuta o compromisso com a impedância exigida pelo
      conversor.
\end{itensq}
\resposta{(a) razão $\approx1/10$: $R_1=\SI{91}{\kilo\ohm}$,
$R_2=\SI{10}{\kilo\ohm}$; (c) $\SI{297}{\micro\ampere}$ no fundo de escala}
\resolucao{(a) É preciso $30\to\SI{3.3}{\volt}$, razão $0{,}11$. Com
$R_2=\SI{10}{\kilo\ohm}$ e $R_1=\SI{91}{\kilo\ohm}$ (ambos E24):
\[
U_o=30\cdot\frac{10}{101}=\SI{2.97}{\volt},
\]
$10\%$ abaixo do fundo de escala --- margem prudente, que acomoda tolerância dos
resistores sem saturar o conversor.\\
(b) Com $\SI{50}{\volt}$ na entrada, a saída iria a
$50\times10/101=\SI{4.95}{\volt}$, acima do limite absoluto do conversor.\\
\textbf{Proteção:} um diodo Zener de $\SI{3.3}{\volt}$ (ou um diodo de
grampeamento para a alimentação) em paralelo com $R_2$. Em sobretensão, ele
conduz e limita a saída; a corrente é confinada por $R_1$ a
\[
I=\frac{50-3{,}3}{91\,\mathrm{k}}=\SI{0.51}{\milli\ampere},
\]
valor inofensivo tanto para o Zener quanto para $R_1$
($P=\SI{24}{\milli\watt}$). \textbf{É $R_1$ que faz a proteção funcionar} ---
sem ele, o diodo teria de absorver corrente ilimitada.\\
(c) No fundo de escala, $I=30/101\,\mathrm{k}=\SI{297}{\micro\ampere}$ e
$P\approx\SI{8.9}{\milli\watt}$.\\
\textbf{Compromisso.} $R_{Th}=91\parallel10=\SI{9}{\kilo\ohm}$. Conversores A/D
de aproximações sucessivas exigem impedância de fonte baixa (tipicamente
$<\SI{10}{\kilo\ohm}$) para carregar o capacitor de amostragem no tempo
disponível --- este projeto fica no limite.\\
\textbf{Soluções:} reduzir os resistores por um fator $10$ (custa dez vezes mais
consumo), acrescentar um capacitor de $\SI{100}{\nano\farad}$ na entrada do
conversor (fornece a carga instantânea de amostragem), ou inserir um seguidor
de tensão. A segunda opção é a mais comum, e a mais barata.}
\end{questao}

% END BANCO COMPLEMENTAR

\gabarito[4.1 Divisor de tensão]

% =====================================================================
\subsection{Divisor de corrente}
% =====================================================================

\blocofixacao

\begin{questao}[1]{}
Em um divisor de corrente de dois resistores em paralelo, a maior corrente
circula:
\begin{alternativas}
\item pelo maior resistor.
\item pelo menor resistor.
\item igualmente pelos dois, sempre.
\item pelo resistor mais próximo da fonte.
\item por nenhum: toda a corrente retorna à fonte.
\end{alternativas}
\resposta{(b)}
\resolucao{Ambos os ramos têm a mesma tensão; por $I = U/R$, o menor resistor
conduz a maior corrente. A corrente se reparte de forma \emph{inversamente}
proporcional às resistências.}
\end{questao}

\begin{questao}[1]{}
Uma corrente de $\SI{10}{\ampere}$ chega a um nó e se divide entre
$R_1=\SI{2}{\ohm}$ e $R_2=\SI{3}{\ohm}$ em paralelo. Calcule $i_1$ e $i_2$.

\circuitoq{%
\begin{circuitikz}[scale=1,transform shape,line width=0.8pt,american]
 \draw (0,0) to[isource,l={$\SI{10}{\ampere}$}] (0,2.6)
       to[short,-*,i={$I_t$}] (2,2.6)
       to[R,l={$R_1=\SI{2}{\ohm}$},i>_={$i_1$}] (2,0) -- (0,0);
 \draw (2,2.6) -- (4,2.6)
       to[R,l={$R_2=\SI{3}{\ohm}$},i>_={$i_2$}] (4,0) to[short,-*] (2,0);
\end{circuitikz}}

\resposta{$i_1 = \SI{6}{\ampere}$; $i_2 = \SI{4}{\ampere}$}
\resolucao{Pelo divisor de corrente (índices \emph{trocados}):
\[
i_1 = I_t\,\frac{R_2}{R_1+R_2}=10\cdot\frac{3}{5}=\SI{6}{\ampere},
\qquad
i_2 = I_t\,\frac{R_1}{R_1+R_2}=10\cdot\frac{2}{5}=\SI{4}{\ampere}.
\]
\textbf{Verificação:} $6+4 = \SI{10}{\ampere}$; e o menor resistor
(\SI{2}{\ohm}) ficou com a maior corrente.}
\end{questao}

\begin{questao}[1]{}
Uma corrente de $\SI{15}{\ampere}$ divide-se entre dois resistores em paralelo,
$\SI{3}{\ohm}$ e $\SI{6}{\ohm}$. As correntes valem, respectivamente:
\begin{alternativas}[2]
\item \SI{5}{\ampere} e \SI{10}{\ampere}
\item \SI{10}{\ampere} e \SI{5}{\ampere}
\item \SI{7.5}{\ampere} e \SI{7.5}{\ampere}
\item \SI{9}{\ampere} e \SI{6}{\ampere}
\item \SI{6}{\ampere} e \SI{9}{\ampere}
\end{alternativas}
\resposta{(b)}
\resolucao{$i_{3\Omega} = 15\cdot\dfrac{6}{9} = \SI{10}{\ampere}$ e
$i_{6\Omega} = 15\cdot\dfrac{3}{9} = \SI{5}{\ampere}$. O resistor de \SI{3}{\ohm},
por ser menor, conduz o dobro do de \SI{6}{\ohm}.}
\end{questao}

\blocoaplicacao

\begin{questao}[2]{}
No circuito, uma fonte de corrente de $\SI{30}{\ampere}$ alimenta três resistores
em paralelo. Calcule a corrente em cada um.

\circuitoq{%
\begin{circuitikz}[scale=1,transform shape,line width=0.8pt,american]
 \draw (0,0) to[isource,l={$\SI{30}{\ampere}$}] (0,2.6)
       to[short,-*,i={$I_t$}] (2,2.6)
       to[R,l={$\SI{2}{\ohm}$},i>_={$i_1$}] (2,0) -- (0,0);
 \draw (2,2.6) -- (4,2.6)
       to[R,l={$\SI{3}{\ohm}$},i>_={$i_2$}] (4,0) to[short,-*] (2,0);
 \draw (4,2.6) -- (6,2.6)
       to[R,l={$\SI{6}{\ohm}$},i>_={$i_3$}] (6,0) to[short,-*] (4,0);
\end{circuitikz}}

\resposta{$i_1=\SI{15}{\ampere}$, $i_2=\SI{10}{\ampere}$, $i_3=\SI{5}{\ampere}$}
\resolucao{Com três ou mais ramos, o mais prático é achar a tensão comum.
$\Req = 2\parallel3\parallel6$: como
$\dfrac{1}{2}+\dfrac{1}{3}+\dfrac{1}{6}=\dfrac{3+2+1}{6}=1$, temos
$\Req=\SI{1}{\ohm}$ e $U = 30\cdot1 = \SI{30}{\volt}$. Então:
\[
i_1=\frac{30}{2}=\SI{15}{\ampere},\quad
i_2=\frac{30}{3}=\SI{10}{\ampere},\quad
i_3=\frac{30}{6}=\SI{5}{\ampere}.
\]
\textbf{Verificação:} $15+10+5 = \SI{30}{\ampere}$.}
\end{questao}

\begin{questao}[2]{}
No divisor da figura, com três ramos, determine a corrente em cada resistor.

\circuitoq{%
\begin{circuitikz}[scale=1,transform shape,line width=0.8pt,american]
 \draw (0,0) to[isource,l={$\SI{25}{\ampere}$}] (0,2.6)
       to[short,-*,i={$I_t$}] (2,2.6)
       to[R,l={$\SI{12}{\ohm}$},i>_={$i_1$}] (2,0) -- (0,0);
 \draw (2,2.6) -- (4,2.6)
       to[R,l={$\SI{3}{\ohm}$},i>_={$i_2$}] (4,0) to[short,-*] (2,0);
 \draw (4,2.6) -- (6,2.6)
       to[R,l={$\SI{12}{\ohm}$},i>_={$i_3$}] (6,0) to[short,-*] (4,0);
\end{circuitikz}}

\resposta{$i_1=i_3\approx\SI{4.17}{\ampere}$; $i_2\approx\SI{16.7}{\ampere}$}
\resolucao{$\dfrac{1}{\Req}=\dfrac{1}{12}+\dfrac{1}{3}+\dfrac{1}{12}
=\dfrac{1+4+1}{12}=\dfrac{6}{12}$, logo $\Req=\SI{2}{\ohm}$ e
$U = 25\cdot2 = \SI{50}{\volt}$.
\[
i_1=i_3=\frac{50}{12}\approx\SI{4.17}{\ampere},
\qquad
i_2=\frac{50}{3}\approx\SI{16.7}{\ampere}.
\]
Os dois ramos de \SI{12}{\ohm} conduzem o mesmo (por simetria); o ramo de
\SI{3}{\ohm}, quatro vezes menor, conduz quatro vezes mais.}
\end{questao}

\blocoaprofundamento

\begin{questao}[3]{}
Uma corrente de $\SI{35}{\ampere}$ alimenta três resistores em paralelo,
$\SI{5}{\ohm}$, $\SI{10}{\ohm}$ e $\SI{20}{\ohm}$.
\begin{itensq}
\item Calcule a corrente em cada ramo.
\item Verifique que a soma das potências dissipadas nos ramos é igual à potência
      entregue pela fonte.
\end{itensq}
\resposta{(a) \SI{20}{\ampere}, \SI{10}{\ampere}, \SI{5}{\ampere};
(b) \SI{3500}{\watt} nos dois lados}
\resolucao{(a) $\dfrac{1}{\Req}=\dfrac{1}{5}+\dfrac{1}{10}+\dfrac{1}{20}
=\dfrac{4+2+1}{20}=\dfrac{7}{20}$, então $\Req = \dfrac{20}{7}\,\si{\ohm}$ e
$U = 35\cdot\dfrac{20}{7} = \SI{100}{\volt}$. As correntes:
\[
i_5=\frac{100}{5}=\SI{20}{\ampere},\quad
i_{10}=\frac{100}{10}=\SI{10}{\ampere},\quad
i_{20}=\frac{100}{20}=\SI{5}{\ampere}.
\]
(b) Potência em cada ramo ($P = U\,i$, com $U=\SI{100}{\volt}$):
\[
100\cdot20 + 100\cdot10 + 100\cdot5 = 2000+1000+500 = \SI{3500}{\watt}.
\]
Potência da fonte: $P = U\,I_t = 100\cdot35 = \SI{3500}{\watt}$ --- confere com a
conservação de energia.}
\end{questao}

\begin{questao}[3]{}
Um miliamperímetro tem fundo de escala de $\SI{1}{\milli\ampere}$ e resistência
interna $R_g = \SI{50}{\ohm}$. Para medir correntes de até $\SI{1}{\ampere}$,
liga-se em paralelo com ele um resistor \emph{shunt} $R_s$, que desvia o excesso
de corrente. Determine $R_s$.

\circuitoq{%
\begin{circuitikz}[scale=1,transform shape,line width=0.8pt,american]
 \draw (0,1.4) to[short,i={$I=\SI{1}{\ampere}$},-*] (2,1.4) coordinate(n1);
 \draw (n1) to[R,l={$R_g=\SI{50}{\ohm}$},i={$\SI{1}{\milli\ampere}$}] (5,1.4) coordinate(n2);
 \node[above,font=\footnotesize] at (3.5,1.75) {galvanômetro};
 \draw (n1) to[short] (2,0) to[R,l_={$R_s$},i_={$i_s$}] (5,0) to[short,-*] (n2);
 \draw (n2) to[short] (7,1.4);
\end{circuitikz}}

\resposta{$R_s \approx \SI{0.05}{\ohm}$}
\resolucao{O galvanômetro e o shunt estão em paralelo, logo têm a \emph{mesma
tensão}. No fundo de escala, o medidor conduz $\SI{1}{\milli\ampere}$ e o shunt
deve desviar o restante:
\[
i_s = 1 - 0{,}001 = \SI{0.999}{\ampere}.
\]
Igualando as tensões nos dois ramos ($R_g\,i_g = R_s\,i_s$):
\[
R_s = \frac{R_g\,i_g}{i_s} = \frac{50\cdot0{,}001}{0{,}999}\approx\SI{0.05}{\ohm}.
\]
\textbf{Ideia central:} o shunt é um divisor de corrente que amplia a faixa do
instrumento por um fator $\dfrac{I}{i_g} = \dfrac{1}{0{,}001} = 1000$. É assim que
um mesmo galvanômetro vira amperímetro de várias escalas.}
\end{questao}

\begin{questao}[3]{}
Um galvanômetro tem fundo de escala $i_g = \SI{1}{\milli\ampere}$ e resistência
interna $R_g = \SI{50}{\ohm}$. Deseja-se convertê-lo em um amperímetro de
\textbf{três escalas}: \SI{100}{\milli\ampere}, \SI{1}{\ampere} e
\SI{10}{\ampere}, usando shunts comutáveis.
\begin{itensq}
\item Deduza a expressão geral do shunt para uma escala $I_{\max}$.
\item Calcule os três shunts.
\item Explique por que a resistência do amperímetro \emph{diminui} conforme se
      aumenta a escala, e por que isso é desejável.
\end{itensq}
\resposta{(a) $R_s = \dfrac{R_g\,i_g}{I_{\max}-i_g}$;
(b) \SI{0.505}{\ohm}, \SI{50.05}{\milli\ohm}, \SI{5.0}{\milli\ohm}; (c) ver comentário}
\resolucao{\textbf{(a)} O shunt está em paralelo com o galvanômetro, logo ambos
têm a mesma tensão. No fundo de escala, o galvanômetro conduz $i_g$ e o shunt
desvia o restante, $I_{\max}-i_g$:
\[
R_g\,i_g = R_s\,(I_{\max}-i_g)
\;\Longrightarrow\;
\boxed{\,R_s=\frac{R_g\,i_g}{I_{\max}-i_g}\,}
\]
\textbf{(b)} Aplicando com $R_g i_g = 50\cdot\pot{1}{-3}=\SI{0.05}{\volt}$:
\begin{center}\small\sffamily
\begin{tabular}{@{}c c c@{}}
\toprule
\textbf{Escala} & \textbf{$R_s$} & \textbf{Fator de multiplicação}\\
\midrule
\SI{100}{\milli\ampere} & $\dfrac{0{,}05}{0{,}099}\approx\SI{0.505}{\ohm}$ & 100$\times$\\
\SI{1}{\ampere}   & $\dfrac{0{,}05}{0{,}999}\approx\SI{50.05}{\milli\ohm}$ & 1000$\times$\\
\SI{10}{\ampere}  & $\dfrac{0{,}05}{9{,}999}\approx\SI{5.0}{\milli\ohm}$ & 10\,000$\times$\\
\bottomrule
\end{tabular}
\end{center}
\textbf{(c)} A resistência total do instrumento é $R_g \parallel R_s$. Como $R_s$
diminui a cada escala maior, o paralelo também diminui --- na escala de
\SI{10}{\ampere}, o amperímetro apresenta apenas $\approx\SI{5}{\milli\ohm}$.\\
\textbf{Por que isso é desejável:} o amperímetro é ligado \emph{em série} com o
circuito, então sua resistência se soma à do circuito e \emph{altera} a corrente
que se pretende medir --- o chamado \textbf{erro de inserção}. Quanto menor a
resistência do instrumento, menor a perturbação. O amperímetro ideal teria
resistência nula (e o voltímetro ideal, resistência infinita).\\
\textbf{Estimativa do erro:} medindo \SI{10}{\ampere} em um circuito de
$\SI{1}{\ohm}$, o erro relativo seria $\dfrac{0{,}005}{1}=\SI{0.5}{\percent}$ ---
aceitável. Mas em um circuito de $\SI{0.05}{\ohm}$, o erro subiria para
\SI{10}{\percent}, tornando a medida inútil.}
\end{questao}

% BEGIN BANCO COMPLEMENTAR Divisor de corrente
% =====================================================================
%  Banco complementar --- Divisor de Corrente  (REVISADO)
%  Substitui a versao gerada por template (10 clones em N2, 9 em N3 e
%  8 questoes conceituais marcadas como N4).
%  RESTRICAO FORTE: o projeto de SHUNT ja aparece DUAS vezes no cap04
%  (miliamperimetro 1 mA/50 ohm e galvanometro 1 mA/50 ohm) e mais
%  duas vezes no banco de circuitos paralelo revisado (shunt 99 ohm e
%  shunt de Ayrton multiescala). Este banco NAO trata de shunt.
%  O cap04 tambem cobre: maior corrente (MC), 10 A entre 2/3, 15 A
%  entre 3/6, 30 A em tres paralelos, tres ramos com figura e 35 A
%  entre 5/10/20 -- de modo que a divisao basica de dois e tres ramos
%  tambem esta excluida.
%  Distribuicao mantida: 5 N1 + 10 N2 + 9 N3 + 8 N4 = 32
% =====================================================================

\subsubsection*{Banco complementar --- Divisor de corrente}

\begin{atencao}[Organização do banco]
A forma correta e geral do divisor de corrente é em \textbf{condutâncias}:
$I_k=I_T\dfrac{G_k}{\sum G_j}$. A versão com resistências,
$I_1=I_T\dfrac{R_2}{R_1+R_2}$, é apenas o caso particular de \emph{dois} ramos
--- e não se generaliza. O N3 trata de desequilíbrio entre condutores,
tolerância, falhas e o contraste entre alimentação por fonte de corrente e por
fonte de tensão; o N4 exige demonstração, projeto de balanceamento e análise de
sensibilidade.
\end{atencao}

\blocofixacao

\begin{questao}[1]{}
Três ramos com condutâncias $G_1=\SI{0.5}{\siemens}$,
$G_2=\SI{0.3}{\siemens}$ e $G_3=\SI{0.2}{\siemens}$ recebem
$I_T=\SI{10}{\ampere}$. Determine a corrente de cada ramo.
\resposta{$\SI{5}{\ampere}$, $\SI{3}{\ampere}$ e $\SI{2}{\ampere}$}
\resolucao{Com $\sum G=0{,}5+0{,}3+0{,}2=\SI{1}{\siemens}$:
\[
I_k=I_T\frac{G_k}{\sum G}
\;\Longrightarrow\;
5;\ 3;\ \SI{2}{\ampere}.
\]
Em condutâncias a conta é imediata: a corrente se reparte na \emph{mesma
proporção} das condutâncias. É por isso que essa é a forma natural do divisor de
corrente, e não a expressão com resistências.}
\end{questao}

\begin{questao}[1]{}
Dois ramos de $\SI{4}{\ohm}$ e $\SI{12}{\ohm}$ dividem uma corrente. Determine
a razão $I_1/I_2$.
\resposta{$I_1/I_2=3$}
\resolucao{Como a tensão é comum, $I=U/R$ e
\[
\frac{I_1}{I_2}=\frac{R_2}{R_1}=\frac{12}{4}=3.
\]
A razão das correntes é a \textbf{inversa} da razão das resistências --- o ramo
de menor resistência leva a maior corrente. Compare com a série, em que a razão
das \emph{tensões} é \emph{direta}.}
\end{questao}

\begin{questao}[1]{}
Um ramo de $\SI{20}{\ohm}$ de uma associação paralela conduz
$\SI{1.5}{\ampere}$; o outro ramo vale $\SI{30}{\ohm}$. Determine a corrente
total.
\resposta{$I_T=\SI{2.5}{\ampere}$}
\resolucao{A tensão comum sai do ramo conhecido:
\[
U=20(1{,}5)=\SI{30}{\volt}
\;\Longrightarrow\;
I_2=\frac{30}{30}=\SI{1}{\ampere},
\]
\[
I_T=1{,}5+1=\SI{2.5}{\ampere}.
\]
Um único ramo conhecido determina a tensão comum e, com ela, todo o resto ---
propriedade que torna qualquer resistor de um paralelo um sensor de corrente em
potencial.}
\end{questao}

\begin{questao}[1]{}
Um dos ramos de um divisor de corrente sofre \textbf{curto-circuito}. Determine
o que ocorre com as correntes.
\resposta{Toda a corrente passa pelo curto; os demais ramos ficam sem corrente}
\resolucao{O ramo em curto tem $R=0$, logo $G\to\infty$. Pelo divisor,
\[
I_k=I_T\frac{G_k}{\infty+\sum_{j\ne k}G_j}\to0
\]
para todos os ramos finitos, e $I_{curto}\to I_T$.\\
\textbf{Fisicamente:} o curto impõe $U=0$ à associação inteira, e sem tensão
nenhum resistor conduz. Os demais ramos continuam íntegros --- eles apenas
deixam de funcionar, sendo vítimas e não causa do defeito.}
\end{questao}

\begin{questao}[1]{}
Em um divisor de corrente alimentado por uma \textbf{fonte de corrente ideal},
um dos ramos é aberto. A corrente nos ramos restantes:
\begin{alternativas}[2]
\item permanece a mesma
\item aumenta
\item diminui
\item torna-se nula
\end{alternativas}
\resposta{(b)}
\resolucao{A fonte de corrente ideal mantém $I_T$ fixo. Com um caminho a menos,
a mesma corrente se distribui entre menos ramos --- cada um recebe \emph{mais}.\\
\textbf{Contraste essencial:} sob fonte de \emph{tensão} ideal, os ramos
restantes ficariam \emph{inalterados} (a tensão comum não muda) e apenas a
corrente total cairia. A mesma falha produz efeitos opostos conforme o tipo de
alimentação --- ponto explorado quantitativamente no bloco de
aprofundamento.}
\end{questao}

\blocoaplicacao

\begin{questao}[2]{}
Uma corrente de $\SI{12}{\ampere}$ divide-se entre quatro ramos de
$\SI{5}{\ohm}$, $\SI{10}{\ohm}$, $\SI{20}{\ohm}$ e $\SI{20}{\ohm}$.
\begin{itensq}
\item Determine a tensão da associação.
\item Determine a corrente de cada ramo.
\end{itensq}
\resposta{$U=\SI{30}{\volt}$; $6$, $3$, $1{,}5$ e $\SI{1.5}{\ampere}$}
\resolucao{
\[
\sum G=0{,}20+0{,}10+0{,}05+0{,}05=\SI{0.40}{\siemens}
\;\Longrightarrow\;
U=\frac{I_T}{\sum G}=\frac{12}{0{,}40}=\SI{30}{\volt}.
\]
\[
I_k=\frac{U}{R_k}\Rightarrow 6;\ 3;\ 1{,}5;\ \SI{1.5}{\ampere}.
\]
\textbf{Verificação:} $6+3+1{,}5+1{,}5=\SI{12}{\ampere}$ \checkmark.\\
\textbf{Método:} calcular a tensão comum primeiro é quase sempre mais rápido do
que aplicar a fórmula do divisor ramo a ramo --- e produz um valor intermediário
que serve de conferência.}
\end{questao}

\begin{questao}[2]{}
No circuito anterior, determine a fração percentual da corrente em cada ramo e
explique por que ela não depende de $I_T$.
\resposta{$50\%$, $25\%$, $12{,}5\%$ e $12{,}5\%$}
\resolucao{
\[
\frac{I_k}{I_T}=\frac{G_k}{\sum G}
\Rightarrow \frac{0{,}20}{0{,}40}=50\%;\quad25\%;\quad12{,}5\%;\quad12{,}5\%.
\]
A razão depende \emph{apenas} das condutâncias --- $I_T$ cancela. Dobrar a
corrente total dobra todas as parcelas, preservando as proporções.\\
\textbf{Consequência prática:} o perfil de divisão é uma propriedade fixa da
rede. Ele só muda se algum ramo mudar de valor, for aberto ou entrar em curto
--- o que faz da monitoração dessas frações uma técnica de diagnóstico.}
\end{questao}

\begin{questao}[2]{}
Uma corrente de $\SI{8}{\ampere}$ deve dividir-se de modo que
$\SI{25}{\percent}$ passe por um ramo $R$ e o restante por um ramo de
$\SI{15}{\ohm}$. Determine $R$.
\resposta{$R=\SI{45}{\ohm}$}
\resolucao{As correntes são $I_R=\SI{2}{\ampere}$ e
$I_{15}=\SI{6}{\ampere}$. A tensão é comum:
\[
U=15(6)=\SI{90}{\volt}
\;\Longrightarrow\;
R=\frac{90}{2}=\SI{45}{\ohm}.
\]
\textbf{Conferência pela razão:} $I_R/I_{15}=R_{15}/R=15/45=1/3$, e de fato
$2/6=1/3$ \checkmark.\\
Note que $R$ é o \emph{triplo} do outro ramo para conduzir um \emph{terço} da
corrente dele --- a proporcionalidade inversa em ação.}
\end{questao}

\begin{questao}[2]{}
Uma fonte de $\SI{60}{\volt}$ com resistência série $R_s=\SI{4}{\ohm}$ alimenta
dois ramos em paralelo de $\SI{10}{\ohm}$ e $\SI{15}{\ohm}$.
\begin{itensq}
\item Determine a corrente total.
\item Determine a corrente de cada ramo.
\end{itensq}
\resposta{$I_T=\SI{6}{\ampere}$; $\SI{3.6}{\ampere}$ e $\SI{2.4}{\ampere}$}
\resolucao{(a) $R_{par}=\dfrac{10\cdot15}{25}=\SI{6}{\ohm}$ e
\[
I_T=\frac{60}{4+6}=\SI{6}{\ampere}.
\]
(b) Pelo divisor de dois ramos:
\[
I_1=6\cdot\frac{15}{25}=\SI{3.6}{\ampere};
\qquad
I_2=6\cdot\frac{10}{25}=\SI{2.4}{\ampere}.
\]
\textbf{Ponto de atenção:} $R_s$ afeta a corrente \emph{total}, mas não a
\emph{proporção} da divisão --- ela é determinada só pelos dois ramos. Essa
separação entre "quanto chega" e "como se reparte" simplifica muito a análise de
circuitos reais.}
\end{questao}

\begin{questao}[2]{}
No circuito anterior, determine a potência dissipada em cada ramo e a tensão
sobre a associação.
\resposta{$U=\SI{36}{\volt}$; $\SI{129.6}{\watt}$ e $\SI{86.4}{\watt}$}
\resolucao{
\[
U=R_{par}I_T=6(6)=\SI{36}{\volt};
\qquad
P_1=\frac{36^{2}}{10}=\SI{129.6}{\watt};
\qquad
P_2=\frac{36^{2}}{15}=\SI{86.4}{\watt}.
\]
Total nos ramos: $\SI{216}{\watt}=36(6)$ \checkmark.\\
\textbf{Observe:} o ramo de \emph{menor} resistência dissipa mais --- regra do
paralelo ($P\propto1/R$ com $U$ comum). O resistor série dissipa ainda
$4(36)=\SI{144}{\watt}$, e a fonte entrega $60(6)=\SI{360}{\watt}$
\checkmark.}
\end{questao}

\begin{questao}[2]{}
Uma corrente de $\SI{10}{\ampere}$ divide-se entre $\SI{1}{\ohm}$ e
$\SI{1}{\kilo\ohm}$.
\begin{itensq}
\item Determine as duas correntes.
\item Avalie o erro de simplesmente desprezar o ramo de alta resistência.
\end{itensq}
\resposta{$\SI{9.99}{\ampere}$ e $\SI{9.99}{\milli\ampere}$; erro de
$0{,}1\%$}
\resolucao{(a)
\[
I_1=10\cdot\frac{1000}{1001}=\SI{9.990}{\ampere};
\qquad
I_2=10\cdot\frac{1}{1001}=\SI{9.99}{\milli\ampere}.
\]
(b) Desprezar o segundo ramo daria $I_1=\SI{10}{\ampere}$ --- erro de
$0{,}1\%$, muito abaixo da tolerância de qualquer resistor real.\\
\textbf{Regra prática:} um ramo mil vezes maior desvia um milésimo da corrente.
Se a razão for de $100{:}1$, o erro sobe para $1\%$; se for $10{:}1$, para
$9\%$ --- aí já não se pode desprezar.\\
\textbf{Mas atenção:} desprezar o ramo para calcular $I_1$ é legítimo;
\emph{esquecê-lo} pode não ser, pois aqueles $\SI{10}{\milli\ampere}$ podem ser
exatamente o sinal que interessa medir.}
\end{questao}

\begin{questao}[2]{}
Uma fonte de corrente de $\SI{12}{\ampere}$ alimenta três ramos de
$\SI{6}{\ohm}$, $\SI{12}{\ohm}$ e $\SI{12}{\ohm}$. O ramo de
$\SI{6}{\ohm}$ é removido. Determine as novas correntes.
\resposta{Antes: $6$, $3$, $\SI{3}{\ampere}$. Depois: $\SI{6}{\ampere}$ em cada}
\resolucao{\textbf{Antes:} $\sum G=\frac16+\frac1{12}+\frac1{12}
=\SI{0.3333}{\siemens}$, $U=12/0{,}3333=\SI{36}{\volt}$ e
$I=6;\ 3;\ \SI{3}{\ampere}$.\\
\textbf{Depois:} $\sum G=\SI{0.1667}{\siemens}$,
$U=12/0{,}1667=\SI{72}{\volt}$ e $I=\SI{6}{\ampere}$ em cada.\\
\textbf{A tensão dobrou} e a corrente de cada ramo remanescente também. Note que
os resistores não mudaram --- foi a \emph{fonte de corrente} que, obrigada a
manter $\SI{12}{\ampere}$, elevou a tensão para forçar a mesma corrente por
menos caminhos.\\
\textbf{Alerta de projeto:} em circuitos alimentados por fonte de corrente,
abrir um ramo \emph{sobrecarrega} os demais. É o oposto do que ocorre com fonte
de tensão.}
\end{questao}

\begin{questao}[2]{}
Um resistor de $\SI{10}{\ohm}$ está em paralelo com um \textbf{indutor} ideal,
e o conjunto recebe $\SI{4}{\ampere}$ em regime permanente CC. Determine a
corrente de cada ramo.
\resposta{$\SI{0}{\ampere}$ no resistor; $\SI{4}{\ampere}$ no indutor}
\resolucao{Em regime permanente CC, a corrente do indutor é constante e
$v=L\,\mathrm{d}i/\mathrm{d}t=0$: ele se comporta como um
\textbf{curto-circuito}.\\
Pelo divisor, $G_L\to\infty$ e toda a corrente vai para o indutor:
\[
I_L=\SI{4}{\ampere};
\qquad
I_R=0.
\]
\textbf{Consequência:} um indutor em paralelo com um resistor "apaga" o resistor
em CC. Isso é usado deliberadamente para curto-circuitar componentes em regime e
deixá-los ativos apenas durante transitórios.\\
\emph{(Um capacitor em paralelo faria o oposto: em CC ele é um aberto e toda a
corrente iria ao resistor.)}}
\end{questao}

\begin{questao}[2]{}
Um divisor de corrente tem três ramos e mediu-se $\SI{4}{\ampere}$,
$\SI{2}{\ampere}$ e $\SI{1.33}{\ampere}$, com tensão comum de
$\SI{20}{\volt}$. Verifique a consistência e determine as resistências.
\resposta{Consistente; $\SI{5}{\ohm}$, $\SI{10}{\ohm}$ e $\SI{15}{\ohm}$}
\resolucao{\textbf{Resistências:} $R_k=U/I_k$:
\[
\frac{20}{4}=5;\qquad \frac{20}{2}=10;\qquad \frac{20}{1{,}33}=\SI{15}{\ohm}.
\]
\textbf{Consistência:} $\sum G=0{,}2+0{,}1+0{,}0667=\SI{0.3667}{\siemens}$, logo
$I_T=20(0{,}3667)=\SI{7.33}{\ampere}$, que confere com
$4+2+1{,}33$ \checkmark.\\
\textbf{Dois testes independentes:} a tensão comum deve ser a mesma em todos os
ramos (senão as leituras são incompatíveis) e a LKC deve fechar. Se apenas o
segundo fechasse, o erro estaria na leitura de tensão.}
\end{questao}

\begin{questao}[2]{}
Explique por que a fórmula $I_1=I_T\dfrac{R_2}{R_1+R_2}$ \textbf{não} se
generaliza para três ramos, e mostre com um contraexemplo.
\resposta{A forma correta usa condutâncias; a versão com resistências vale só
para dois ramos}
\resolucao{\textbf{Contraexemplo.} Ramos de $\SI{6}{\ohm}$, $\SI{12}{\ohm}$ e
$\SI{12}{\ohm}$ com $I_T=\SI{12}{\ampere}$.\\
\emph{Forma correta:} $I_1=12\cdot\dfrac{1/6}{1/3}=\SI{6}{\ampere}$.\\
\emph{"Generalização" ingênua} $I_1=I_T\dfrac{R_2+R_3}{R_1+R_2+R_3}$:
$12\cdot\dfrac{24}{30}=\SI{9.6}{\ampere}$ --- errado, e nem sequer a soma das
três parcelas assim calculadas daria $\SI{12}{\ampere}$.\\
\textbf{Origem da confusão:} com \emph{dois} ramos,
$\dfrac{G_1}{G_1+G_2}=\dfrac{1/R_1}{1/R_1+1/R_2}=\dfrac{R_2}{R_1+R_2}$ --- a
simplificação funciona por acaso algébrico. Com três, a manipulação equivalente
produziria $\dfrac{R_2R_3}{R_2R_3+R_1R_3+R_1R_2}$, expressão que ninguém
memoriza.\\
\textbf{Recomendação:} use sempre condutâncias. A fórmula com resistências é
uma conveniência de dois ramos, não a regra geral.}
\end{questao}

\blocoaprofundamento

\begin{questao}[3]{}
Três condutores em paralelo, de resistências $\SI{0.10}{\ohm}$,
$\SI{0.11}{\ohm}$ e $\SI{0.12}{\ohm}$, conduzem $\SI{300}{\ampere}$.
\begin{itensq}
\item Determine a corrente de cada um.
\item Determine o desvio percentual em relação à divisão ideal.
\item Comente a implicação térmica.
\end{itensq}
\resposta{$109{,}4$; $99{,}5$ e $\SI{91.2}{\ampere}$; desvios de
$+9{,}4\%$, $-0{,}5\%$ e $-8{,}8\%$}
\resolucao{(a) $\sum G=10+9{,}091+8{,}333=\SI{27.42}{\siemens}$ e
$U=300/27{,}42=\SI{10.94}{\volt}$. As correntes valem
\[
109{,}4;\qquad 99{,}5;\qquad \SI{91.2}{\ampere}.
\]
(b) Em relação a $\SI{100}{\ampere}$ ideais: $+9{,}4\%$, $-0{,}5\%$ e
$-8{,}8\%$.\\
(c) \textbf{Implicação térmica.} A dissipação vai com $I^{2}$: o condutor mais
carregado dissipa $(109{,}4/100)^{2}=1{,}20$ vezes o previsto --- $20\%$ a
mais. Ele aquece mais que os outros dois.\\
\textbf{E aí surge a pergunta decisiva:} o cobre tem coeficiente térmico
\emph{positivo}. Aquecendo, sua resistência sobe, ele devolve corrente aos
vizinhos e o sistema se \textbf{estabiliza}. Se o material tivesse coeficiente
negativo, ocorreria o contrário --- fuga térmica.\\
\textbf{Prática de instalação:} condutores em paralelo devem ter mesmo
comprimento, seção, material e trajeto. Uma diferença de $\SI{20}{\percent}$ no
comprimento produz o desequilíbrio calculado aqui, e é por isso que a norma
exige lances idênticos.}
\end{questao}

\begin{questao}[3]{}
Uma corrente de $\SI{12}{\ampere}$ deve dividir-se na razão $1{:}2{:}3$.
\begin{itensq}
\item Determine as três correntes.
\item Determine as resistências, adotando a tensão comum de
      $\SI{12}{\volt}$.
\item Determine a potência de cada ramo.
\end{itensq}
\resposta{$2$, $4$ e $\SI{6}{\ampere}$; $\SI{6}{\ohm}$, $\SI{3}{\ohm}$ e
$\SI{2}{\ohm}$; $24$, $48$ e $\SI{72}{\watt}$}
\resolucao{(a) As frações são $\frac16$, $\frac26$ e $\frac36$ de
$\SI{12}{\ampere}$: $2$, $4$ e $\SI{6}{\ampere}$.\\
(b) Como $U$ é comum, $R_k=U/I_k$:
\[
\frac{12}{2}=6;\qquad \frac{12}{4}=3;\qquad \frac{12}{6}=\SI{2}{\ohm}.
\]
As resistências ficam na razão $6{:}3{:}2$ --- \emph{inversa} de
$1{:}2{:}3$, como esperado.\\
(c) $P_k=UI_k$: $24$, $48$ e $\SI{72}{\watt}$; total $\SI{144}{\watt}=12(12)$
\checkmark.\\
\textbf{Observação de projeto:} a tensão comum foi \emph{escolhida}. Adotando
$\SI{1.2}{\volt}$ em vez de $\SI{12}{\volt}$, as resistências cairiam por dez
($0{,}6$; $0{,}3$; $\SI{0.2}{\ohm}$) e a dissipação total também --- de
$144$ para $\SI{14.4}{\watt}$. A razão de divisão é a mesma; o que muda é a
perda. Esse grau de liberdade é explorado no bloco de desafio.}
\end{questao}

\begin{questao}[3]{}
Dois ramos nominalmente iguais de $\SI{10}{\ohm}$, com tolerância de
$\pm5\%$, dividem $\SI{10}{\ampere}$.
\begin{itensq}
\item Determine o pior desequilíbrio de corrente.
\item Compare o desvio percentual da corrente com o dos resistores.
\end{itensq}
\resposta{$5{,}25$ e $\SI{4.75}{\ampere}$; desvio de $\pm5\%$}
\resolucao{(a) O pior caso ocorre com desvios opostos: $R_1=\SI{9.5}{\ohm}$ e
$R_2=\SI{10.5}{\ohm}$.
\[
I_1=10\cdot\frac{10{,}5}{20}=\SI{5.25}{\ampere};
\qquad
I_2=\SI{4.75}{\ampere}.
\]
(b) Em relação aos $\SI{5}{\ampere}$ ideais, o desvio é de
$\pm5\%$ --- \emph{igual} à tolerância dos resistores.\\
\textbf{Por que dá exatamente isso:} a sensibilidade relativa da divisão é
$\dfrac{R_1}{R_1+R_2}=0{,}5$, e o pior caso combina as duas tolerâncias:
$0{,}5\times(5\%+5\%)=5\%$.\\
\textbf{Generalização:} com ramos \emph{desiguais}, o fator muda. Em um divisor
$1{:}9$, a sensibilidade do ramo pequeno é $0{,}9$ e o desequilíbrio chega a
$9\%$ --- o ramo minoritário é sempre o mais afetado por tolerâncias.}
\end{questao}

\begin{questao}[3]{}
Um ramo de $\SI{0.5}{\ohm}$ de um divisor adquire resistência de contato de
$\SI{0.05}{\ohm}$ em suas conexões. O outro ramo, também de
$\SI{0.5}{\ohm}$, permanece íntegro, e a corrente total é
$\SI{100}{\ampere}$.
\begin{itensq}
\item Determine as novas correntes.
\item Determine o desvio e a potência adicional no contato.
\end{itensq}
\resposta{$\SI{47.6}{\ampere}$ e $\SI{52.4}{\ampere}$; desvio de
$\mp4{,}8\%$; $\SI{113}{\watt}$ no contato}
\resolucao{(a) O ramo afetado passa a $\SI{0.55}{\ohm}$:
\[
I_1=100\cdot\frac{0{,}50}{1{,}05}=\SI{47.6}{\ampere};
\qquad
I_2=100\cdot\frac{0{,}55}{1{,}05}=\SI{52.4}{\ampere}.
\]
(b) Desvio de $-4{,}8\%$ e $+4{,}8\%$ em relação aos
$\SI{50}{\ampere}$ ideais.\\
No contato: $P=0{,}05(47{,}6)^{2}=\SI{113}{\watt}$ --- concentrados em uma área
minúscula.\\
\textbf{Por que é grave.} A resistência de contato \emph{cresce} com a corrosão
e com a temperatura. Mais resistência $\to$ mais aquecimento $\to$ mais
oxidação $\to$ mais resistência: realimentação positiva que termina em falha
térmica. Curiosamente, o desvio de corrente é \emph{estabilizador} (o ramo ruim
conduz menos), mas a concentração de potência no ponto de contato domina o
processo.\\
\textbf{Conclusão:} em divisores de alta corrente, a qualidade das conexões
importa mais que a tolerância dos resistores.}
\end{questao}

\begin{questao}[3]{}
Uma corrente de $\SI{10}{\ampere}$ divide-se entre $\SI{5}{\ohm}$ e
$\SI{20}{\ohm}$.
\begin{itensq}
\item Determine as correntes e as potências.
\item Determine qual ramo aquece mais e por qual fator.
\end{itensq}
\resposta{$8$ e $\SI{2}{\ampere}$; $320$ e $\SI{80}{\watt}$ --- o menor
resistor dissipa $4\times$ mais}
\resolucao{(a)
\[
I_1=10\cdot\frac{20}{25}=\SI{8}{\ampere};
\qquad
I_2=\SI{2}{\ampere};
\]
\[
P_1=5(8)^{2}=\SI{320}{\watt};
\qquad
P_2=20(2)^{2}=\SI{80}{\watt}.
\]
(b) O ramo de $\SI{5}{\ohm}$ dissipa quatro vezes mais.\\
\textbf{Razão algébrica:} com $U$ comum, $P=U^{2}/R$, logo
$P_1/P_2=R_2/R_1=4$. A razão das potências é a \emph{inversa} da razão das
resistências --- e igual à razão das correntes ($8/2=4$).\\
\textbf{Erro comum:} supor que o resistor maior aquece mais. Isso vale em
\emph{série} ($P\propto R$), não em paralelo. Em um divisor de corrente, o
componente crítico para dimensionamento térmico é sempre o de \emph{menor}
resistência.}
\end{questao}

\begin{questao}[3]{}
Três ramos de $\SI{10}{\ohm}$, $\SI{20}{\ohm}$ e $\SI{20}{\ohm}$ têm o
primeiro aberto. Compare o efeito sob dois tipos de alimentação:
\begin{itensq}
\item fonte de corrente ideal de $\SI{12}{\ampere}$;
\item fonte de tensão ideal de $\SI{60}{\volt}$.
\end{itensq}
\resposta{(a) tensão dobra ($60\to\SI{120}{\volt}$), ramos vão de $3$ a
$\SI{6}{\ampere}$; (b) tensão e ramos inalterados, total cai de $12$ a
$\SI{6}{\ampere}$}
\resolucao{\textbf{(a) Fonte de corrente.} Antes:
$\sum G=\SI{0.2}{\siemens}$, $U=\SI{60}{\volt}$,
$I=6;3;\SI{3}{\ampere}$.\\
Depois: $\sum G=\SI{0.1}{\siemens}$, e a fonte \emph{eleva a tensão} para
manter $\SI{12}{\ampere}$:
\[
U=\frac{12}{0{,}1}=\SI{120}{\volt};
\qquad
I=6;\ \SI{6}{\ampere}.
\]
Cada ramo remanescente \textbf{dobra} de corrente, e a potência quadruplica
(de $180$ para $\SI{720}{\watt}$ em cada).\\
\textbf{(b) Fonte de tensão.} $U=\SI{60}{\volt}$ permanece.
Antes: $6;3;\SI{3}{\ampere}$, total $\SI{12}{\ampere}$. Depois:
$3;\SI{3}{\ampere}$, total $\SI{6}{\ampere}$. \textbf{Os ramos remanescentes
não sentem nada.}\\
\textbf{Síntese.} A mesma falha, no mesmo circuito, produz sobrecarga severa em
um caso e nenhum efeito no outro. O que decide é a natureza da fonte:
\begin{itemize}
\item fonte de \emph{tensão} $\Rightarrow$ ramos desacoplados, falha localizada;
\item fonte de \emph{corrente} $\Rightarrow$ ramos acoplados, falha propagada.
\end{itemize}
É por isso que instalações elétricas são alimentadas por tensão, e por que
circuitos em série de LEDs acionados por corrente exigem proteção contra
abertura.}
\end{questao}

\begin{questao}[3]{}
Um ramo fixo de $\SI{10}{\ohm}$ está em paralelo com um reostato de
$\SI{0}{\ohm}$ a $\SI{40}{\ohm}$, sob corrente total de $\SI{5}{\ampere}$.
\begin{itensq}
\item Determine a faixa de corrente no ramo fixo.
\item Determine a posição do reostato para divisão igual.
\item Comente a linearidade do controle.
\end{itensq}
\resposta{(a) de $0$ a $\SI{4}{\ampere}$; (b) $R=\SI{10}{\ohm}$}
\resolucao{(a) Com o reostato em $\SI{0}{\ohm}$, ele curto-circuita o ramo
fixo: $I_{fixo}=0$. Com $\SI{40}{\ohm}$:
\[
I_{fixo}=5\cdot\frac{40}{50}=\SI{4}{\ampere}.
\]
A faixa é $0$ a $\SI{4}{\ampere}$ --- \emph{nunca} os $\SI{5}{\ampere}$
completos, pois o reostato sempre desvia alguma corrente.\\
(b) Divisão igual exige resistências iguais: $R=\SI{10}{\ohm}$, um quarto do
curso.\\
(c) \textbf{Fortemente não linear.} A relação
$I_{fixo}=5R/(R+10)$ é hiperbólica:
\begin{center}
\begin{tabular}{ccc}
\hline
$R$ & posição do curso & $I_{fixo}$\\
\hline
$\SI{10}{\ohm}$ & $25\%$ & $\SI{2.5}{\ampere}$\\
$\SI{20}{\ohm}$ & $50\%$ & $\SI{3.33}{\ampere}$\\
$\SI{40}{\ohm}$ & $100\%$ & $\SI{4.0}{\ampere}$\\
\hline
\end{tabular}
\end{center}
Metade da faixa útil de corrente é percorrida no primeiro quarto do curso --- o
ajuste é grosseiro no início e quase inerte no fim.}
\end{questao}

\begin{questao}[3]{}
Deduza a expressão da corrente em um ramo $R_k$ quando todos os demais são
substituídos por sua resistência equivalente $R_p$, e verifique com o divisor de
$6$, $12$ e $\SI{12}{\ohm}$ sob $\SI{12}{\ampere}$.
\resposta{$I_k=I_T\dfrac{R_p}{R_k+R_p}$; para $R_k=\SI{6}{\ohm}$ e
$R_p=\SI{6}{\ohm}$, $I_k=\SI{6}{\ampere}$}
\resolucao{\textbf{Dedução.} Agrupando todos os ramos exceto o $k$-ésimo em
$R_p$, sobram \emph{dois} ramos em paralelo --- e aí a fórmula de dois ramos é
válida:
\[
I_k=I_T\frac{R_p}{R_k+R_p},
\qquad\text{com }
\frac{1}{R_p}=\sum_{j\ne k}\frac{1}{R_j}.
\]
\textbf{Verificação.} Para $R_k=\SI{6}{\ohm}$, os demais dão
$R_p=12\parallel12=\SI{6}{\ohm}$:
\[
I_k=12\cdot\frac{6}{6+6}=\SI{6}{\ampere}.\ \checkmark
\]
\textbf{Utilidade:} é a maneira correta de \emph{recuperar} a fórmula de dois
ramos em uma rede com muitos ramos --- e mostra exatamente onde a
"generalização" ingênua erra: $R_p$ é o \emph{paralelo} dos demais, não a
\emph{soma}. No exemplo, a soma daria $\SI{24}{\ohm}$ e a corrente
$\SI{9.6}{\ampere}$, valor incorreto.}
\end{questao}

\begin{questao}[3]{}
Um divisor de corrente tem um ramo contendo uma \textbf{fonte de tensão} de
$\SI{10}{\volt}$ em série com $\SI{5}{\ohm}$; o outro ramo é um resistor de
$\SI{5}{\ohm}$. A corrente total injetada é $\SI{4}{\ampere}$.
\begin{itensq}
\item Explique por que o divisor simples não se aplica.
\item Determine as correntes.
\end{itensq}
\resposta{$I_1=\SI{1}{\ampere}$ e $I_2=\SI{3}{\ampere}$}
\resolucao{(a) O divisor de corrente pressupõe ramos \textbf{puramente
resistivos}, submetidos à mesma tensão e obedecendo $I=UG$. Um ramo com fonte
segue $I=(U-E)G$ --- relação \emph{afim}, não proporcional. A repartição deixa
de depender só das condutâncias.\\
(b) Chamando $U$ a tensão comum:
\[
\underbrace{\frac{U-10}{5}}_{I_1}+\underbrace{\frac{U}{5}}_{I_2}=4
\;\Longrightarrow\;
2U-10=20
\;\Longrightarrow\;
U=\SI{15}{\volt},
\]
\[
I_1=\frac{15-10}{5}=\SI{1}{\ampere};
\qquad
I_2=\frac{15}{5}=\SI{3}{\ampere}.
\]
\textbf{Compare:} o divisor simples, com dois ramos iguais, preveria
$\SI{2}{\ampere}$ em cada. A fonte de $\SI{10}{\volt}$ \emph{se opõe} à corrente
em seu ramo e desvia $\SI{1}{\ampere}$ para o outro.\\
\textbf{Método correto:} volte à LKC com a tensão comum como incógnita --- ou
seja, análise nodal. O divisor de corrente é um atalho, e atalhos têm
hipóteses.}
\end{questao}

\blocodesafio

\begin{questao}[4]{}
\textbf{(Demonstração geral.)} Deduza a fórmula do divisor de corrente para
$n$ ramos em paralelo.
\begin{itensq}
\item Obtenha $I_k$ em função das condutâncias.
\item Mostre que a forma com resistências vale apenas para $n=2$.
\item Verifique que $\sum I_k=I_T$ identicamente.
\end{itensq}
\resposta{$I_k=I_T\dfrac{G_k}{\sum_j G_j}$}
\resolucao{(a) Todos os ramos compartilham a tensão $U$. Pela LKC,
\[
I_T=\sum_j I_j=\sum_j UG_j=U\sum_j G_j
\;\Longrightarrow\;
U=\frac{I_T}{\sum_j G_j}.
\]
Logo
\[
I_k=UG_k=I_T\frac{G_k}{\sum_j G_j}.
\]
(b) Para $n=2$, substituindo $G=1/R$:
\[
\frac{G_1}{G_1+G_2}=\frac{1/R_1}{\dfrac{R_1+R_2}{R_1R_2}}
=\frac{R_2}{R_1+R_2}.
\]
A simplificação depende de o denominador ter \emph{exatamente} dois termos.
Para $n=3$, o resultado análogo seria
\[
I_1=I_T\frac{R_2R_3}{R_2R_3+R_1R_3+R_1R_2},
\]
que não tem a forma "resistência do outro sobre a soma" --- e para $n$ geral
envolve todos os produtos de $n-1$ resistências.\\
(c)
\[
\sum_k I_k=I_T\frac{\sum_k G_k}{\sum_j G_j}=I_T.
\]
A LKC é satisfeita \textbf{identicamente}, para quaisquer condutâncias --- ela
não é um resultado do cálculo, mas uma consequência da forma da expressão.
Verificar a soma, portanto, não detecta erro nas condutâncias; ela só detecta
erro de aritmética.}
\end{questao}

\begin{questao}[4]{}
\textbf{(Projeto de balanceamento.)} Deve-se dividir $\SI{60}{\ampere}$
igualmente entre três ramos, com perda total inferior a $\SI{10}{\watt}$.
\begin{itensq}
\item Determine a resistência de cada ramo e a queda de tensão.
\item Verifique a perda e determine a margem.
\item Discuta o compromisso entre precisão da divisão e perda.
\end{itensq}
\resposta{$R\le\SI{8.33}{\milli\ohm}$; queda de $\SI{0.167}{\volt}$;
$\SI{3.33}{\watt}$ por ramo no limite}
\resolucao{(a) Divisão igual exige três resistências iguais, com
$\SI{20}{\ampere}$ cada. A perda total é
\[
P=3RI_k^{2}=3R(400)=1200R\le10
\;\Longrightarrow\;
R\le\SI{8.33}{\milli\ohm}.
\]
Adotando $R=\SI{8.33}{\milli\ohm}$, a queda é
$U=8{,}33\,\mathrm{m}\times20=\SI{0.167}{\volt}$.\\
(b) No limite, cada ramo dissipa $\SI{3.33}{\watt}$ --- exigindo resistores de
pelo menos $\SI{10}{\watt}$ para operar com folga, e provavelmente ventilados.\\
(c) \textbf{O compromisso central.} O desequilíbrio relativo entre os ramos é
governado pela \emph{razão} entre a dispersão dos resistores de balanceamento e
as resistências \emph{parasitas} (contatos, trechos de barramento, cabos):
\begin{itemize}
\item $R$ \textbf{grande}: as parasitas tornam-se desprezíveis em comparação, e
a divisão fica precisa --- mas a perda cresce com $R$.
\item $R$ \textbf{pequeno}: pouca perda, mas as parasitas passam a dominar e o
balanceamento se degrada.
\end{itemize}
Com $R=\SI{8.33}{\milli\ohm}$ e contatos de $\SI{1}{\milli\ohm}$, a incerteza
de divisão já é da ordem de $12\%$ --- o resistor de balanceamento praticamente
não cumpre sua função.\\
\textbf{Conclusão:} não se pode ter simultaneamente divisão precisa e perda
desprezível com resistores passivos. Aplicações críticas usam balanceamento
\emph{ativo} (realimentação de corrente por ramo), que rompe o compromisso.}
\end{questao}

\begin{questao}[4]{}
\textbf{(Amplificação de desequilíbrio.)} Em $n$ condutores nominalmente
iguais, um deles tem resistência $R(1-\delta)$.
\begin{itensq}
\item Deduza a fração de corrente que ele conduz.
\item Obtenha a expressão aproximada para $\delta$ pequeno e identifique o fator
      de amplificação.
\item Avalie para $n=3$ e $\delta=5\%$, $10\%$ e $20\%$.
\end{itensq}
\resposta{Desvio relativo $\approx\dfrac{n-1}{n}\,\delta$; para $n=3$:
$3{,}4\%$, $7{,}1\%$ e $15{,}4\%$}
\resolucao{(a) As condutâncias são $\dfrac{1}{R(1-\delta)}$ e
$(n-1)$ vezes $\dfrac1R$:
\[
\frac{I_1}{I_T}=\frac{\dfrac{1}{1-\delta}}{\dfrac{1}{1-\delta}+(n-1)}
=\frac{1}{1+(n-1)(1-\delta)}.
\]
(b) Para $\delta\ll1$, expandindo em torno de $\delta=0$ (onde a fração vale
$1/n$):
\[
\frac{I_1}{I_T}\approx\frac1n\left(1+\frac{n-1}{n}\,\delta\right)
\;\Longrightarrow\;
\frac{\Delta I_1}{I_1}\approx\frac{n-1}{n}\,\delta.
\]
O \textbf{fator de amplificação} é $\dfrac{n-1}{n}$: sempre menor que $1$, mas
crescente com $n$ --- tende a $1$ para muitos condutores.\\
(c) Com $n=3$, o fator é $2/3$:
\begin{center}
\begin{tabular}{ccc}
\hline
$\delta$ & exato & aproximado\\
\hline
$5\%$ & $+3{,}45\%$ & $+3{,}33\%$\\
$10\%$ & $+7{,}14\%$ & $+6{,}67\%$\\
$20\%$ & $+15{,}4\%$ & $+13{,}3\%$\\
\hline
\end{tabular}
\end{center}
\textbf{Duas leituras.} A aproximação é boa até cerca de $10\%$ e subestima o
desvio real --- portanto \emph{não} é conservadora, e deve ser usada com
cuidado em verificação de projeto.\\
E o desequilíbrio de corrente é sempre \emph{menor} que o de resistência (fator
$<1$), o que é uma boa notícia. A má notícia é que a \textbf{potência} vai com
$I^{2}$: um desvio de $7\%$ na corrente significa $15\%$ a mais de
aquecimento no condutor mais carregado.}
\end{questao}

\begin{questao}[4]{}
\textbf{(Sensibilidade.)} Para um divisor de dois ramos:
\begin{itensq}
\item Obtenha $\dfrac{\partial I_1}{\partial R_1}$.
\item Obtenha a sensibilidade relativa e mostre que $|S_1|+|S_2|=1$.
\item Interprete os casos $R_1\ll R_2$ e $R_1\gg R_2$.
\end{itensq}
\resposta{$S_1=-\dfrac{R_1}{R_1+R_2}$, $S_2=+\dfrac{R_2}{R_1+R_2}$}
\resolucao{(a) De $I_1=I_T\dfrac{R_2}{R_1+R_2}$:
\[
\frac{\partial I_1}{\partial R_1}=-\frac{I_TR_2}{(R_1+R_2)^{2}}.
\]
(b) Sensibilidade relativa:
\[
S_1=\frac{\partial I_1}{\partial R_1}\cdot\frac{R_1}{I_1}
=-\frac{I_TR_2}{(R_1+R_2)^{2}}\cdot\frac{R_1(R_1+R_2)}{I_TR_2}
=-\frac{R_1}{R_1+R_2}.
\]
Analogamente, $S_2=+\dfrac{R_2}{R_1+R_2}$, e
\[
|S_1|+|S_2|=\frac{R_1+R_2}{R_1+R_2}=1.
\]
\textbf{Mesmo padrão} já encontrado em série, paralelo e divisor de tensão ---
consequência de $I_1/I_T$ ser função homogênea de grau zero nas resistências.\\
(c) \textbf{$R_1\ll R_2$:} $S_1\to0$. O ramo pequeno leva quase toda a corrente
e é \emph{insensível} ao próprio valor --- mas $S_2\to1$, ou seja, sua corrente
depende integralmente do ramo grande. Contraintuitivo e importante.\\
\textbf{$R_1\gg R_2$:} $S_1\to-1$. O ramo grande conduz pouco, e essa pouca
corrente é inteiramente determinada pelo seu próprio valor.\\
\textbf{Regra de projeto:} para dividir corrente com precisão, aperte a
tolerância do ramo que conduz \emph{menos} --- é ele que carrega a maior
sensibilidade relativa.}
\end{questao}

\begin{questao}[4]{}
\textbf{(Divisor de baixa perda para medição.)} Deseja-se medir
$\SI{100}{\ampere}$ desviando uma fração conhecida para um instrumento, com
perda mínima.
\begin{itensq}
\item Projete um divisor de razão $1000{:}1$ com ramo principal de
      $\SI{0.1}{\milli\ohm}$.
\item Determine a corrente de medição, a queda e a perda total.
\item Compare com um resistor \emph{shunt} direto de $\SI{1}{\milli\ohm}$.
\end{itensq}
\resposta{Ramo de medição de $\SI{100}{\milli\ohm}$; $\SI{99.9}{\milli\ampere}$;
perda de $\SI{1}{\watt}$ contra $\SI{10}{\watt}$}
\resolucao{(a) A razão de correntes é a razão inversa das resistências:
\[
\frac{I_{med}}{I_{princ}}=\frac{R_{princ}}{R_{med}}=\frac{1}{1000}
\;\Longrightarrow\;
R_{med}=1000\times0{,}1\,\mathrm{m}=\SI{100}{\milli\ohm}.
\]
(b)
\[
G_{tot}=10^{4}+10=\SI{10010}{\siemens};
\qquad
I_{med}=100\cdot\frac{10}{10010}=\SI{99.9}{\milli\ampere}.
\]
\[
R_{eq}=\SI{99.9}{\micro\ohm};
\qquad
U=100(\pot{99.9}{-6})=\SI{9.99}{\milli\volt};
\qquad
P=100(\pot{9.99}{-3})=\SI{0.999}{\watt}.
\]
(c) \textbf{Shunt direto de $\SI{1}{\milli\ohm}$:}
$U=\SI{100}{\milli\volt}$ e $P=\SI{10}{\watt}$ --- \textbf{dez vezes} mais
perda.\\
\textbf{O compromisso, explicitado.} A perda vale $P=I_T^{2}R_{eq}$, e
$R_{eq}$ é essencialmente o ramo principal. Reduzi-lo reduz a perda \emph{e} a
tensão de medição --- que já é de apenas $\SI{10}{\milli\volt}$, exigindo
amplificação de baixo ruído e cuidado com termopares parasitas nas junções.\\
\textbf{Alternativa moderna:} sensores de efeito Hall ou transformadores de
corrente medem sem inserir resistência alguma no caminho principal, eliminando a
perda por completo --- ao custo de exatidão e de resposta em CC.}
\end{questao}

\begin{questao}[4]{}
\textbf{(Estabilidade térmica da divisão.)} Dois ramos idênticos dividem uma
corrente elevada. Um deles aquece mais que o outro.
\begin{itensq}
\item Analise a evolução para material de coeficiente térmico \emph{positivo}.
\item Repita para coeficiente \emph{negativo}.
\item Determine a condição de estabilidade e cite exemplos.
\end{itensq}
\resposta{$\alpha>0$: estável (realimentação negativa); $\alpha<0$: instável
(fuga térmica)}
\resolucao{(a) \textbf{Coeficiente positivo} (cobre, manganina, a maioria dos
metais). O ramo mais quente tem sua resistência \emph{aumentada}; sua
condutância cai e ele passa a conduzir \emph{menos}. Menos corrente significa
menos aquecimento --- o desvio se corrige sozinho.\\
Trata-se de \textbf{realimentação negativa}: o sistema converge para um novo
equilíbrio, com desequilíbrio limitado.\\
(b) \textbf{Coeficiente negativo} (semicondutores, termistores NTC, carbono). O
ramo mais quente reduz sua resistência, atrai \emph{mais} corrente, aquece
ainda mais --- e o ciclo se acelera. \textbf{Realimentação positiva}, ou fuga
térmica: um ramo acaba conduzindo praticamente toda a corrente e se destrói.\\
(c) \textbf{Condição de estabilidade:} $\alpha>0$ nos ramos, ou --- se
$\alpha<0$ for inevitável --- resistências de balanceamento em série, de valor
suficiente para dominar a variação térmica.\\
\textbf{Exemplos práticos.}
\begin{itemize}
\item \textbf{Condutores em paralelo:} cobre tem $\alpha\approx+0{,}004/\si{\celsius}$
--- naturalmente estáveis.
\item \textbf{Transistores bipolares em paralelo:} o coeficiente da junção é
\emph{negativo}, e por isso cada emissor recebe um resistor de balanceamento.
Sem ele, um transistor assume toda a corrente e queima.
\item \textbf{Diodos e LEDs em paralelo:} mesmo problema --- daí a exigência de
resistor individual por ramo, e não um único resistor comum.
\item \textbf{Resistores de precisão para divisão:} usa-se manganina, que
combina $\alpha$ pequeno \emph{e} positivo.
\end{itemize}}
\end{questao}

\begin{questao}[4]{}
\textbf{(Divisor com elemento ativo.)} Um dos ramos de um divisor contém uma
fonte de corrente \textbf{dependente} que injeta $g\,U$, sendo $U$ a tensão
comum e $g=\SI{0.05}{\siemens}$. O outro ramo é um resistor de
$\SI{10}{\ohm}$, e a corrente total injetada é $\SI{6}{\ampere}$.
\begin{itensq}
\item Explique por que o divisor simples falha.
\item Determine a tensão comum e a corrente de cada ramo.
\item Determine o valor crítico de $g$.
\end{itensq}
\resposta{(b) $U=\SI{120}{\volt}$; $\SI{12}{\ampere}$ no resistor e
$\SI{6}{\ampere}$ injetados pelo ramo ativo; (c) $g_c=\SI{0.1}{\siemens}$}
\resolucao{(a) O divisor pressupõe que \emph{cada} ramo obedeça $I=UG$ com
$G>0$ fixo. A fonte dependente também é proporcional a $U$, mas com sinal que
pode \emph{retirar} corrente do nó --- ela funciona como condutância
\textbf{negativa} de valor $-g$. A soma $\sum G$ deixa de ser a soma de
condutâncias físicas.\\
(b) LKC no nó, com a fonte dependente \textbf{injetando} $gU$:
\[
\frac{U}{10}=6+gU
\;\Longrightarrow\;
U(0{,}1-0{,}05)=6
\;\Longrightarrow\;
U=\frac{6}{0{,}05}=\SI{120}{\volt},
\]
\[
I_R=\frac{120}{10}=\SI{12}{\ampere};
\qquad
I_{dep}=0{,}05(120)=\SI{6}{\ampere}.
\]
\textbf{Verificação da LKC:} entram $6$ (fonte externa) mais $6$ (fonte
dependente), totalizando $\SI{12}{\ampere}$ --- exatamente o que sai pelo
resistor \checkmark.\\
\textbf{Observe o absurdo aparente:} o resistor conduz $\SI{12}{\ampere}$, o
\textbf{dobro} da corrente injetada externamente. Nenhum divisor passivo
produziria isso, pois nele vale sempre $I_k\le I_T$.\\
(c) A condutância líquida é $\dfrac1{10}-g$, que se anula em
\[
g_c=\SI{0.1}{\siemens}.
\]
Nesse ponto o nó "flutua" e não há solução finita; acima dele, a resistência
equivalente torna-se negativa e o circuito é instável.\\
\textbf{Lição:} elementos ativos podem produzir correntes de ramo
\emph{maiores} que a corrente total injetada, ou de sinal invertido --- ambos
impossíveis em rede puramente passiva. O divisor de corrente é uma ferramenta
de redes passivas.}
\end{questao}

\begin{questao}[4]{}
\textbf{(Escolha de escala.)} Um divisor deve repartir $\SI{12}{\ampere}$ na
razão $1{:}2{:}3$. A razão fixa apenas as \emph{proporções}; resta escolher a
escala absoluta das resistências.
\begin{itensq}
\item Mostre que a perda total é proporcional à escala escolhida.
\item Determine a escala para perda de $\SI{14.4}{\watt}$ e para
      $\SI{144}{\watt}$.
\item Estabeleça o critério de escolha na prática.
\end{itensq}
\resposta{$P=I_T^{2}R_{eq}$, e $R_{eq}$ escala com o fator adotado; escalas de
$0{,}6/0{,}3/\SI{0.2}{\ohm}$ e $6/3/\SI{2}{\ohm}$}
\resolucao{(a) Multiplicando todas as resistências por $\lambda$, as
\emph{frações} $G_k/\sum G$ não mudam --- a divisão é preservada. Mas
\[
R_{eq}=\lambda R_{eq,0}
\;\Longrightarrow\;
P=I_T^{2}R_{eq}=\lambda\,I_T^{2}R_{eq,0}.
\]
A perda é \textbf{diretamente proporcional} à escala. (É a mesma propriedade de
homogeneidade de grau 1 já usada para tolerâncias em associações.)\\
(b) Com $R=6/3/\SI{2}{\ohm}$: $R_{eq}=\SI{1}{\ohm}$ e
$P=144(1)=\SI{144}{\watt}$.\\
Com $R=0{,}6/0{,}3/\SI{0.2}{\ohm}$: $R_{eq}=\SI{0.1}{\ohm}$ e
$P=\SI{14.4}{\watt}$ --- dez vezes menos, com \emph{exatamente} a mesma divisão
$2{:}4{:}6$ ampères.\\
(c) \textbf{Critério.} A escala é limitada por baixo e por cima:
\begin{itemize}
\item \textbf{Limite inferior:} resistências pequenas demais tornam-se
comparáveis às parasitas (contatos, trilhas, cabos), que passam a governar a
divisão. Regra prática: mantenha $R\ge20$ vezes a maior parasita estimada.
\item \textbf{Limite superior:} a perda $\lambda I_T^{2}R_{eq,0}$ e a queda de
tensão crescem linearmente, exigindo componentes maiores e comprometendo o
rendimento.
\end{itemize}
\textbf{Escolha:} a menor escala que ainda domine as parasitas com folga. Com
contatos de $\SI{1}{\milli\ohm}$, a escala de $0{,}6/0{,}3/\SI{0.2}{\ohm}$ é
segura (o menor ramo é $200\times$ a parasita); reduzir mais dez vezes já
deixaria a divisão à mercê da qualidade das conexões.}
\end{questao}

% END BANCO COMPLEMENTAR

\gabarito[4.2 Divisor de corrente]
              % completo

% Cap. 5 -------------------------------------------------------------
\setcounter{section}{4}
% =====================================================================
%  CAPITULO 5 - POTENCIA ELETRICA E MAXIMA TRANSFERENCIA
% =====================================================================
\section{Potência Elétrica e Máxima Transferência}

\begin{conceito}[Antes de começar]
\textbf{Potência dissipada em um resistor:}
$P = U\,I = R\,I^2 = \dfrac{U^2}{R}$.\par
\textbf{Fonte fornecendo ou absorvendo:} usando a convenção do gerador, se a
corrente sai pelo terminal positivo, a fonte \emph{fornece} potência
($P>0$); se entra pelo positivo, a fonte \emph{absorve} (está sendo carregada,
$P<0$).\par
\textbf{Máxima transferência de potência:} uma fonte real (Thévenin) de tensão
$V_{\text{Th}}$ e resistência interna $R_{\text{Th}}$ entrega a máxima potência a
uma carga $R_L$ quando
\[
\boxed{\,R_L = R_{\text{Th}}\,},\qquad
P_{\max}=\frac{V_{\text{Th}}^{\,2}}{4\,R_{\text{Th}}}.
\]
Nesse ponto, o rendimento é de apenas $\SI{50}{\percent}$ (metade da potência se
perde internamente).
\end{conceito}

\legendaniveis

\blocofixacao

\begin{questao}[1]{}
Um resistor de $\SI{7}{\ohm}$ é percorrido por uma corrente de $\SI{3}{\ampere}$.
A potência dissipada é:
\begin{alternativas}[3]
\item \SI{21}{\watt}
\item \SI{49}{\watt}
\item \SI{63}{\watt}
\item \SI{147}{\watt}
\item \SI{441}{\watt}
\end{alternativas}
\resposta{(c)}
\resolucao{$P = R\,I^2 = 7\cdot 3^2 = \SI{63}{\watt}$.}
\end{questao}

\begin{questao}[1]{}
Na condição de máxima transferência de potência de uma fonte real para uma carga,
a relação entre a resistência de carga $R_L$ e a resistência interna
$R_{\text{Th}}$ da fonte é:
\begin{alternativas}
\item $R_L$ deve ser a menor possível.
\item $R_L$ deve ser a maior possível.
\item $R_L = R_{\text{Th}}$.
\item $R_L = 2R_{\text{Th}}$.
\item $R_L = R_{\text{Th}}/2$.
\end{alternativas}
\resposta{(c)}
\resolucao{A potência na carga é máxima quando $R_L = R_{\text{Th}}$ (casamento de
impedância). Cargas muito pequenas recebem alta corrente mas baixa tensão; cargas
muito grandes, o contrário. O equilíbrio ótimo está na igualdade.}
\end{questao}

\begin{questao}[1]{}
Um dispositivo submetido a $\SI{12}{\volt}$ é percorrido por $\SI{3}{\ampere}$. A
potência que ele consome é:
\begin{alternativas}[3]
\item \SI{4}{\watt}
\item \SI{15}{\watt}
\item \SI{36}{\watt}
\item \SI{4}{\watt}
\item \SI{0.25}{\watt}
\end{alternativas}
\resposta{(c)}
\resolucao{$P = U\,I = 12\cdot3 = \SI{36}{\watt}$.}
\end{questao}

\blocoaplicacao

\begin{questao}[2]{}
Uma fonte de Thévenin tem $V_{\text{Th}} = \SI{12}{\volt}$ e
$R_{\text{Th}} = \SI{4}{\ohm}$.
\begin{itensq}
\item Qual o valor de $R_L$ para máxima transferência de potência?
\item Qual essa potência máxima?
\end{itensq}
\resposta{(a) \SI{4}{\ohm}; (b) \SI{9}{\watt}}
\resolucao{(a) $R_L = R_{\text{Th}} = \SI{4}{\ohm}$.\\
(b) $P_{\max} = \dfrac{V_{\text{Th}}^{\,2}}{4R_{\text{Th}}}
= \dfrac{12^2}{4\cdot4} = \dfrac{144}{16} = \SI{9}{\watt}$.\\
\textbf{Confira:} com $R_L = R_{\text{Th}}$, a corrente é
$I = \dfrac{12}{4+4} = \SI{1.5}{\ampere}$ e
$P_L = R_L I^2 = 4\cdot1{,}5^2 = \SI{9}{\watt}$.}
\end{questao}

\begin{questao}[2]{}
Uma fonte com resistência interna de $\SI{6}{\ohm}$ e fem de $\SI{18}{\volt}$
alimenta uma carga $R_L$. Determine $R_L$ para máxima transferência de potência e
o valor dessa potência.
\resposta{$R_L = \SI{6}{\ohm}$; $P_{\max} = \SI{13.5}{\watt}$}
\resolucao{$R_L = R_{\text{int}} = \SI{6}{\ohm}$, e
\[
P_{\max}=\frac{V^2}{4R_{\text{int}}}=\frac{18^2}{4\cdot6}=\frac{324}{24}
=\SI{13.5}{\watt}.
\]}
\end{questao}

\begin{questao}[2]{}
No circuito da figura, a fonte de $\SI{60}{\volt}$ está \textbf{absorvendo}
potência (sendo carregada). Sobre a corrente $I$ indicada (que entra pelo terminal
positivo da fonte de $\SI{60}{\volt}$), pode-se afirmar que:

\circuitoq{%
\begin{circuitikz}[scale=0.95,transform shape,line width=0.8pt]
 \draw (0,0) to[battery1,l_={$\SI{90}{\volt}$}] (0,3)
       to[R,l={$\SI{10}{\ohm}$}] (4,3)
       to[battery1,l={$\SI{60}{\volt}$},invert] (4,0) -- (0,0);
 \draw[->,Icol,line width=1pt] (0.8,3.4)--(2.0,3.4)
       node[midway,above,font=\footnotesize]{$I$};
\end{circuitikz}}

\begin{alternativas}
\item $I$ é positiva no sentido indicado.
\item $I$ é nula.
\item $I$ tem sentido tal que entra pelo $+$ da fonte de \SI{60}{\volt},
      confirmando a absorção.
\item a fonte de \SI{60}{\volt} nunca pode absorver potência.
\item nada se pode afirmar sem o valor do resistor.
\end{alternativas}
\resposta{(c)}
\resolucao{Para uma fonte \emph{absorver} potência, a corrente deve entrar pelo
seu terminal positivo (convenção do receptor). Aqui a fonte de \SI{90}{\volt}
impõe a corrente no sentido horário, que de fato entra pelo $+$ da fonte de
\SI{60}{\volt} --- por isso esta última é carregada.
A corrente vale $I = \dfrac{90-60}{10} = \SI{3}{\ampere}$, e a fonte de
\SI{60}{\volt} absorve $P = 60\cdot3 = \SI{180}{\watt}$. É exatamente o que ocorre
ao carregar uma bateria: uma fonte de maior tensão força corrente "para dentro"
dela.}
\end{questao}

\begin{questao}[2]{}
Uma bateria de fem $\SI{12}{\volt}$ e resistência interna $\SI{2}{\ohm}$ alimenta
uma carga $R_L = \SI{4}{\ohm}$.
\begin{itensq}
\item Qual a potência entregue à carga?
\item Compare com a potência máxima teórica (obtida quando $R_L = R_{\text{int}}$).
\end{itensq}
\resposta{(a) \SI{16}{\watt}; (b) máximo é \SI{18}{\watt} em $R_L=\SI{2}{\ohm}$}
\resolucao{(a) $I = \dfrac{12}{2+4} = \SI{2}{\ampere}$;
$P_L = R_L I^2 = 4\cdot4 = \SI{16}{\watt}$.\\
(b) O máximo ocorreria em $R_L = R_{\text{int}} = \SI{2}{\ohm}$:
$P_{\max} = \dfrac{12^2}{4\cdot2} = \SI{18}{\watt}$. A carga de \SI{4}{\ohm}
entrega \SI{16}{\watt}, um pouco abaixo do máximo --- coerente com a curva suave
de $P_L$ em torno do ponto ótimo.}
\end{questao}

\blocoaprofundamento

\begin{questao}[3]{}
Uma fonte real de $V_{\text{Th}} = \SI{20}{\volt}$ e $R_{\text{Th}} = \SI{10}{\ohm}$
alimenta uma carga variável $R_L$.
\begin{itensq}
\item Calcule a potência na carga para $R_L = \SI{5}{\ohm}$, $\SI{10}{\ohm}$ e
      $\SI{20}{\ohm}$.
\item Confirme que o máximo ocorre em $R_L = R_{\text{Th}}$.
\end{itensq}
\resposta{$P(5)\approx\SI{8.9}{\watt}$, $P(10)=\SI{10}{\watt}$,
$P(20)\approx\SI{8.9}{\watt}$; máximo em \SI{10}{\ohm}}
\resolucao{A corrente é $I = \dfrac{V_{\text{Th}}}{R_{\text{Th}}+R_L}$ e a potência
na carga, $P_L = R_L I^2 = \dfrac{V_{\text{Th}}^{\,2}\,R_L}{(R_{\text{Th}}+R_L)^2}$:
\[
P(5)=\frac{400\cdot5}{15^2}=\frac{2000}{225}\approx\SI{8.9}{\watt},\quad
P(10)=\frac{400\cdot10}{20^2}=\SI{10}{\watt},\quad
P(20)=\frac{400\cdot20}{30^2}=\frac{8000}{900}\approx\SI{8.9}{\watt}.
\]
O valor máximo é $P(10) = \SI{10}{\watt}$, confirmando
$R_L = R_{\text{Th}} = \SI{10}{\ohm}$. Note a \emph{simetria}: cargas de
\SI{5}{} e \SI{20}{\ohm} (uma metade, outra o dobro de $R_{\text{Th}}$) dão a mesma
potência. A curva $P_L \times R_L$ tem um máximo suave em $R_L = R_{\text{Th}}$.}
\end{questao}

\begin{questao}[3]{}
Explique por que, na condição de máxima transferência de potência
($R_L = R_{\text{Th}}$), o rendimento do circuito é de apenas
$\SI{50}{\percent}$, e por que, em sistemas de \emph{potência} (rede elétrica),
essa condição \textbf{não} é desejável.
\resposta{Ver comentário}
\resolucao{Com $R_L = R_{\text{Th}}$, a corrente se divide igualmente e a mesma
potência dissipada na carga é dissipada \emph{internamente} em $R_{\text{Th}}$.
O rendimento é
\[
\eta=\frac{P_{\text{carga}}}{P_{\text{total}}}
=\frac{R_L}{R_L+R_{\text{Th}}}=\frac{1}{2}=\SI{50}{\percent}.
\]
Metade da energia vira calor na fonte --- aceitável em sinais fracos (áudio,
antenas, onde importa extrair o \emph{máximo} sinal), mas inaceitável na rede
elétrica, onde se quer \emph{eficiência}. Por isso, em sistemas de potência,
opera-se com $R_L \gg R_{\text{Th}}$ (rendimento próximo de \SI{100}{\percent}),
sacrificando potência máxima em favor de baixo desperdício.}
\end{questao}

\begin{questao}[3]{}
Uma fonte de $V_{\text{Th}}=\SI{24}{\volt}$ e $R_{\text{Th}}=\SI{4}{\ohm}$
alimenta uma carga variável $R_L$. Complete mentalmente a análise abaixo e
responda.

\begin{table}[H]
\centering\small\sffamily
\setlength{\arraystretch}{1.25}
\begin{tabular}{@{}c c c c@{}}
\toprule
$R_L$ ($\Omega$) & $I$ (A) & $P_L$ (W) & $\eta$ (\%)\\
\midrule
1  & 4,8 & 23,0 & 20\\
2  & 4,0 & 32,0 & 33\\
\textbf{4}  & \textbf{3,0} & \textbf{36,0} & \textbf{50}\\
8  & 2,0 & 32,0 & 67\\
16 & 1,2 & 23,0 & 80\\
\bottomrule
\end{tabular}
\caption{Potência na carga e rendimento em função de $R_L$.}
\end{table}

\begin{itensq}
\item Verifique os valores da linha $R_L=\SI{8}{\ohm}$.
\item Explique a \emph{simetria} observada na coluna $P_L$.
\item Se o objetivo for \emph{eficiência energética}, qual $R_L$ escolher? E se
      for \emph{máxima potência}? Justifique a diferença.
\end{itensq}
\resposta{(a) $I=\SI{2}{\ampere}$, $P_L=\SI{32}{\watt}$, $\eta=\SI{67}{\percent}$;
(b) e (c) ver comentário}
\resolucao{\textbf{(a)} Para $R_L=\SI{8}{\ohm}$:
\[
I=\frac{24}{4+8}=\SI{2}{\ampere},\quad
P_L=R_LI^2=8\cdot4=\SI{32}{\watt},\quad
\eta=\frac{R_L}{R_{\text{Th}}+R_L}=\frac{8}{12}\approx\SI{67}{\percent}.\ \checkmark
\]
\textbf{(b) A simetria.} Observe que $R_L=\SI{2}{\ohm}$ e $R_L=\SI{8}{\ohm}$ dão a
\emph{mesma} potência (\SI{32}{\watt}); o mesmo ocorre com \SI{1}{} e
\SI{16}{\ohm} (\SI{23}{\watt}). Isso porque cargas que são \emph{recíprocas} em
relação a $R_{\text{Th}}$ --- isto é, $R_L$ e $R_{\text{Th}}^2/R_L$ --- produzem
potências idênticas:
\[
P_L=\frac{V^2R_L}{(R_{\text{Th}}+R_L)^2}
\quad\text{é invariante sob}\quad
R_L \leftrightarrow \frac{R_{\text{Th}}^{2}}{R_L}.
\]
Verifique: $2$ e $16/2=8$ \checkmark; $1$ e $16/1=16$ \checkmark. A curva é
simétrica em escala \emph{logarítmica}, com máximo em $R_L=R_{\text{Th}}$.

\textbf{(c) O conflito central de projeto.}
\begin{itemize}[leftmargin=1.6em,itemsep=1pt,topsep=2pt]
\item \textbf{Máxima potência} $\Rightarrow R_L = R_{\text{Th}} = \SI{4}{\ohm}$,
      mas com apenas \SI{50}{\percent} de rendimento --- metade da energia vira
      calor na fonte.
\item \textbf{Máxima eficiência} $\Rightarrow R_L$ o \emph{maior} possível
      ($\eta \to \SI{100}{\percent}$), mas a potência entregue tende a zero.
\end{itemize}
\textbf{Não existe ponto que otimize ambos} --- são objetivos conflitantes.
A escolha depende do que é escasso:
\begin{itemize}[leftmargin=1.6em,itemsep=1pt,topsep=2pt]
\item em \emph{telecomunicações} e áudio, o sinal é fraco e a energia é barata:
      casa-se a impedância para extrair o máximo sinal;
\item em \emph{sistemas de potência}, a energia é o produto: opera-se com
      $R_L \gg R_{\text{Th}}$, com rendimento acima de \SI{95}{\percent}.
\end{itemize}
\emph{Saber qual critério aplicar é mais importante que saber calcular ambos.}}
\end{questao}

\blocodesafio

\begin{questao}[4]{}
\textbf{Problema de síntese integrado.} Projete um circuito que satisfaça
simultaneamente:
\begin{itemize}[leftmargin=1.6em,itemsep=1pt,topsep=2pt]
\item alimentado por uma fonte de $\SI{24}{\volt}$ com resistência interna
      $\SI{2}{\ohm}$;
\item entregue \emph{máxima potência} a uma carga;
\item a corrente total drenada da fonte não ultrapasse $\SI{6}{\ampere}$.
\end{itemize}
Determine o valor da carga, a potência entregue e verifique a restrição de
corrente.
\resposta{$R_L = \SI{2}{\ohm}$; $P = \SI{72}{\watt}$; $I = \SI{6}{\ampere}$
(no limite)}
\resolucao{\textbf{Condição de máxima transferência:}
$R_L = R_{\text{int}} = \SI{2}{\ohm}$.\\
\textbf{Corrente resultante:}
\[
I = \frac{24}{R_{\text{int}}+R_L}=\frac{24}{4}=\SI{6}{\ampere}.
\]
A restrição é atendida \emph{exatamente no limite} --- não há margem.\\
\textbf{Potência na carga:}
\[
P_L = R_L I^2 = 2\cdot6^2 = \SI{72}{\watt}
\quad\text{(confere com } P_{\max}=\tfrac{24^2}{4\cdot2}=\SI{72}{\watt}).
\]
\textbf{Análise crítica do projeto:} a fonte dissipa internamente os mesmos
\SI{72}{\watt} (rendimento de \SI{50}{\percent}), totalizando \SI{144}{\watt}
consumidos. Como a corrente está no limite especificado, qualquer tolerância
para baixo em $R_L$ (por aquecimento, por exemplo) violaria a restrição.
\textbf{Recomendação de engenharia:} adotar $R_L$ ligeiramente \emph{acima} de
\SI{2}{\ohm} --- por exemplo \SI{2.2}{\ohm} --- sacrifica menos de
\SI{1}{\percent} da potência ($\SI{71.5}{\watt}$) e traz a corrente para
\SI{5.7}{\ampere}, criando margem de segurança. \emph{Projetar no limite exato é
um erro clássico de quem só sabe calcular; o profissional experiente sempre
reserva margem.}}
\end{questao}

% BEGIN BANCO COMPLEMENTAR Potência Elétrica e Máxima Transferência
% =====================================================================
%  Banco complementar --- Potencia Eletrica e Maxima Transferencia
%  REVISADO. Substitui a versao gerada por template (8 clones em N2,
%  9 em N3 e 7 questoes conceituais marcadas como N4).
%  O cap05 ja cobre: P=RI^2 e P=VI (MC), a condicao R_L=R_Th (MC),
%  Thevenin 12 V/4 ohm, fonte 18 V/6 ohm, fonte ABSORVENDO potencia,
%  bateria 12 V/2 ohm com R_L=4 ohm, DUAS tabelas de carga variavel
%  (20 V/10 ohm e 24 V/4 ohm), a explicacao de POR QUE o rendimento
%  vale 50% no ponto otimo, e um problema de sintese integrado.
%  Este banco evita todos esses casos e explora: a SIMETRIA dos pares
%  reciprocos de carga, a faixa de tolerancia em torno do otimo, o
%  otimo sob restricao de rendimento e de corrente, e o caso em que o
%  parametro livre e R_Th (onde nao ha otimo interior).
%  Distribuicao mantida: 5 N1 + 8 N2 + 9 N3 + 7 N4 = 29
% =====================================================================

\subsubsection*{Banco complementar --- Potência elétrica e máxima transferência}

\begin{atencao}[Organização do banco]
Duas expressões governam o tópico:
$P_L=\dfrac{V_{Th}^{2}R_L}{(R_L+R_{Th})^{2}}$ e
$P_{\max}=\dfrac{V_{Th}^{2}}{4R_{Th}}$ em $R_L=R_{Th}$. Convém trabalhar com a
variável adimensional $x=R_L/R_{Th}$, pois
$\dfrac{P_L}{P_{\max}}=\dfrac{4x}{(1+x)^{2}}$ e
$\eta=\dfrac{x}{1+x}$ --- duas curvas universais, independentes dos valores
absolutos. O N3 explora simetria, tolerância e casos de contorno; o N4 trata de
demonstração, projeto e otimização sob restrição.
\end{atencao}

\blocofixacao

\begin{questao}[1]{}
Uma fonte tem $V_{Th}=\SI{24}{\volt}$ e $R_{Th}=\SI{6}{\ohm}$. Determine
$R_L$ para máxima transferência e a potência correspondente.
\resposta{$R_L=\SI{6}{\ohm}$; $P_{\max}=\SI{24}{\watt}$}
\resolucao{A condição é $R_L=R_{Th}=\SI{6}{\ohm}$, e então
\[
P_{\max}=\frac{V_{Th}^{2}}{4R_{Th}}=\frac{576}{24}=\SI{24}{\watt}.
\]
\textbf{Guarde a fórmula do máximo:} ela dispensa calcular corrente e tensão, e
mostra que a potência máxima disponível de uma fonte depende \emph{apenas} de
$V_{Th}$ e $R_{Th}$ --- é uma característica da fonte, não da carga.}
\end{questao}

\begin{questao}[1]{}
Na condição anterior, determine a corrente e a tensão sobre a carga.
\resposta{$I=\SI{2}{\ampere}$; $U_L=\SI{12}{\volt}$}
\resolucao{Com $R_L=R_{Th}$, a resistência total é $2R_{Th}$:
\[
I=\frac{V_{Th}}{2R_{Th}}=\frac{24}{12}=\SI{2}{\ampere};
\qquad
U_L=R_LI=6(2)=\SI{12}{\volt}=\frac{V_{Th}}{2}.
\]
\textbf{Regra do ponto ótimo:} a tensão na carga é sempre \emph{metade} da
tensão em vazio, e a corrente é metade da corrente de curto-circuito. Esses dois
fatos identificam o ponto ótimo em qualquer medição de bancada.}
\end{questao}

\begin{questao}[1]{}
Uma fonte entrega no máximo $\SI{50}{\watt}$, com $R_{Th}=\SI{8}{\ohm}$.
Determine $V_{Th}$.
\resposta{$V_{Th}=\SI{40}{\volt}$}
\resolucao{De $P_{\max}=\dfrac{V_{Th}^{2}}{4R_{Th}}$:
\[
V_{Th}=\sqrt{4R_{Th}P_{\max}}=\sqrt{4(8)(50)}=\sqrt{1600}=\SI{40}{\volt}.
\]
\textbf{Conferência:} $\dfrac{40^{2}}{4(8)}=\dfrac{1600}{32}=\SI{50}{\watt}$
\checkmark.}
\end{questao}

\begin{questao}[1]{}
Uma fonte de $\SI{24}{\volt}$ entrega no máximo $\SI{24}{\watt}$. Determine
$R_{Th}$.
\resposta{$R_{Th}=\SI{6}{\ohm}$}
\resolucao{
\[
R_{Th}=\frac{V_{Th}^{2}}{4P_{\max}}=\frac{576}{96}=\SI{6}{\ohm}.
\]
\textbf{Método experimental:} medindo $V_{oc}$ e a potência máxima entregue a
uma carga ajustável, obtém-se $R_{Th}$ sem nunca curto-circuitar a fonte ---
alternativa útil quando o curto seria destrutivo.}
\end{questao}

\begin{questao}[1]{}
Se $R_L$ for ajustado para o \textbf{dobro} do valor ótimo, a potência
entregue:
\begin{alternativas}[2]
\item dobra
\item cai à metade
\item cai para $\SI{89}{\percent}$ do máximo
\item permanece máxima
\end{alternativas}
\resposta{(c)}
\resolucao{Com $x=R_L/R_{Th}=2$:
\[
\frac{P_L}{P_{\max}}=\frac{4x}{(1+x)^{2}}=\frac{8}{9}=88{,}9\%.
\]
\textbf{Fato importante:} o máximo é \textbf{achatado}. Errar a carga por um
fator $2$ custa apenas $11\%$ da potência --- o que torna o casamento muito
mais tolerante do que se costuma supor. A quantificação completa dessa
tolerância está no bloco de aprofundamento.}
\end{questao}

\blocoaplicacao

\begin{questao}[2]{}
Uma fonte com $V_{Th}=\SI{24}{\volt}$ e $R_{Th}=\SI{6}{\ohm}$ alimenta
sucessivamente $R_L=\SI{3}{\ohm}$ e $R_L=\SI{12}{\ohm}$. Determine a potência
em cada caso e comente.
\resposta{$\SI{21.33}{\watt}$ nos dois casos}
\resolucao{
\[
R_L=\SI{3}{\ohm}:\quad I=\frac{24}{9}=\SI{2.667}{\ampere},\quad
P=3(2{,}667)^{2}=\SI{21.33}{\watt};
\]
\[
R_L=\SI{12}{\ohm}:\quad I=\frac{24}{18}=\SI{1.333}{\ampere},\quad
P=12(1{,}333)^{2}=\SI{21.33}{\watt}.
\]
\textbf{Potências idênticas} --- e note que $3\times12=36=R_{Th}^{2}$. Cargas
que são \emph{recíprocas} em relação a $R_{Th}$ recebem exatamente a mesma
potência. A demonstração está no bloco de desafio.\\
\textbf{Mas os dois casos não são equivalentes:} com $\SI{3}{\ohm}$ a corrente
é o dobro e o rendimento é $33\%$; com $\SI{12}{\ohm}$, o rendimento é
$67\%$. Mesma potência entregue, consumo da fonte muito diferente.}
\end{questao}

\begin{questao}[2]{}
No ponto de máxima transferência da fonte anterior, determine a potência total
fornecida, a entregue à carga e a dissipada internamente.
\resposta{$\SI{48}{\watt}$, $\SI{24}{\watt}$ e $\SI{24}{\watt}$}
\resolucao{Com $I=\SI{2}{\ampere}$:
\[
P_{tot}=V_{Th}I=24(2)=\SI{48}{\watt};
\qquad
P_L=R_LI^{2}=\SI{24}{\watt};
\qquad
P_{int}=R_{Th}I^{2}=\SI{24}{\watt}.
\]
\textbf{Metade da energia é perdida dentro da fonte} --- consequência direta de
$R_L=R_{Th}$ e da corrente comum. É o preço do casamento, e a razão pela qual
ele nunca é usado em eletrotécnica de potência.}
\end{questao}

\begin{questao}[2]{}
Duas fontes têm a mesma potência máxima de $\SI{18}{\watt}$: a primeira com
$V_{Th}=\SI{12}{\volt}$ e a segunda com $V_{Th}=\SI{24}{\volt}$. Determine
$R_{Th}$ de cada uma.
\resposta{$\SI{2}{\ohm}$ e $\SI{8}{\ohm}$}
\resolucao{
\[
R_{Th}=\frac{V_{Th}^{2}}{4P_{\max}}
\Rightarrow
\frac{144}{72}=\SI{2}{\ohm};
\qquad
\frac{576}{72}=\SI{8}{\ohm}.
\]
\textbf{Leitura:} para a mesma potência disponível, dobrar a tensão permite
\emph{quadruplicar} a resistência interna. Fontes de tensão mais alta toleram
resistência interna maior --- princípio que fundamenta a transmissão de energia
em alta tensão.\\
\textbf{Diferença prática:} a primeira entrega $\SI{18}{\watt}$ a
$\SI{3}{\ampere}$; a segunda, a $\SI{1.5}{\ampere}$. Correntes menores
significam condutores mais finos.}
\end{questao}

\begin{questao}[2]{}
Uma carga fixa de $\SI{10}{\ohm}$ é alimentada, alternativamente, por duas
fontes de $\SI{20}{\volt}$: uma com $R_{Th}=\SI{2}{\ohm}$ e outra com
$R_{Th}=\SI{10}{\ohm}$. Compare a potência entregue.
\resposta{$\SI{27.8}{\watt}$ e $\SI{10}{\watt}$}
\resolucao{
\[
R_{Th}=\SI{2}{\ohm}:\quad I=\frac{20}{12}=\SI{1.667}{\ampere},\quad
P_L=10(1{,}667)^{2}=\SI{27.8}{\watt};
\]
\[
R_{Th}=\SI{10}{\ohm}:\quad I=\SI{1}{\ampere},\quad P_L=\SI{10}{\watt}.
\]
\textbf{Aparente paradoxo:} a segunda fonte está \emph{casada} com a carga
($R_L=R_{Th}$) e mesmo assim entrega bem menos que a primeira, descasada.\\
\textbf{Resolução:} o casamento maximiza a potência quando $R_L$ é o parâmetro
livre e a fonte é dada. Aqui a carga é fixa e a \emph{fonte} varia --- e nesse
caso quanto \emph{menor} $R_{Th}$, melhor. São problemas de otimização
diferentes, com respostas diferentes. O ponto é retomado no bloco de
aprofundamento.}
\end{questao}

\begin{questao}[2]{}
Uma fonte casada entrega $\SI{24}{\watt}$ durante $\SI{5}{\hour}$. Determine a
energia entregue à carga e a energia total consumida da fonte.
\resposta{$\SI{120}{\watt\hour}$ e $\SI{240}{\watt\hour}$}
\resolucao{
\[
E_L=P_Lt=24(5)=\SI{120}{\watt\hour}=\SI{0.12}{\kilo\watt\hour};
\]
\[
E_{tot}=48(5)=\SI{240}{\watt\hour}.
\]
Metade da energia --- $\SI{120}{\watt\hour}$ --- foi dissipada internamente e
perdida.\\
\textbf{Consequência para sistemas alimentados por bateria:} operar casado
\emph{reduz pela metade} a autonomia. Em um rádio portátil, isso significa
trocar tempo de operação por volume de som --- decisão que às vezes vale a pena,
mas que precisa ser consciente.}
\end{questao}

\begin{questao}[2]{}
Uma fonte de $\SI{24}{\volt}$ e $R_{Th}=\SI{6}{\ohm}$ alimenta
$R_L=\SI{18}{\ohm}$. Determine a potência na carga, o rendimento e compare com
o máximo disponível.
\resposta{$P_L=\SI{18}{\watt}$; $\eta=75\%$; $75\%$ do máximo}
\resolucao{
\[
I=\frac{24}{24}=\SI{1}{\ampere};
\qquad
P_L=18(1)^{2}=\SI{18}{\watt};
\qquad
\eta=\frac{18}{24}=75\%.
\]
Comparando com $P_{\max}=\SI{24}{\watt}$: entrega-se $75\%$ da potência
disponível, com rendimento de $75\%$ em vez de $50\%$.\\
\textbf{Compromisso quantificado:} abrir mão de $25\%$ da potência comprou
$25$ pontos percentuais de rendimento. Em quase toda aplicação prática esse é
um bom negócio --- e é por isso que sistemas reais operam com
$R_L\gg R_{Th}$.}
\end{questao}

\begin{questao}[2]{}
Uma fonte tem $V_{oc}=\SI{30}{\volt}$ e $I_{sc}=\SI{5}{\ampere}$. Determine
$R_{Th}$ e a potência máxima disponível.
\resposta{$R_{Th}=\SI{6}{\ohm}$; $P_{\max}=\SI{37.5}{\watt}$}
\resolucao{
\[
R_{Th}=\frac{V_{oc}}{I_{sc}}=\frac{30}{5}=\SI{6}{\ohm};
\qquad
P_{\max}=\frac{30^{2}}{4(6)}=\SI{37.5}{\watt}.
\]
\textbf{Forma alternativa e elegante:}
\[
P_{\max}=\frac{V_{oc}I_{sc}}{4}=\frac{30(5)}{4}=\SI{37.5}{\watt}.
\]
A potência máxima é um quarto do produto dos dois valores extremos --- e faz
sentido: no ponto ótimo, tensão e corrente valem, cada uma, metade do seu
extremo.}
\end{questao}

\begin{questao}[2]{}
Verifique o balanço de potência para a fonte de $\SI{24}{\volt}$ e
$R_{Th}=\SI{6}{\ohm}$ com $R_L=\SI{2}{\ohm}$.
\resposta{Fonte: $\SI{72}{\watt}$; carga: $\SI{18}{\watt}$; interna:
$\SI{54}{\watt}$}
\resolucao{
\[
I=\frac{24}{8}=\SI{3}{\ampere};
\qquad
P_{fonte}=24(3)=\SI{72}{\watt};
\]
\[
P_L=2(9)=\SI{18}{\watt};
\qquad
P_{int}=6(9)=\SI{54}{\watt};
\qquad 18+54=72\ \checkmark
\]
\textbf{Rendimento de apenas $25\%$} --- três quartos da energia aquecem a
fonte. Compare com a carga de $\SI{18}{\ohm}$ da questão anterior, que entrega
a \emph{mesma} $\SI{18}{\watt}$ com rendimento de $75\%$. Duas cargas
recíprocas, mesma potência, rendimentos opostos.}
\end{questao}

\blocoaprofundamento

\begin{questao}[3]{}
Para $V_{Th}=\SI{24}{\volt}$ e $R_{Th}=\SI{6}{\ohm}$, verifique a simetria dos
pares recíprocos de carga.
\begin{itensq}
\item Calcule $P_L$ para $R_L=\SI{2}{\ohm}$ e $R_L=\SI{18}{\ohm}$.
\item Repita para $\SI{1.5}{\ohm}$ e $\SI{24}{\ohm}$.
\item Identifique a relação entre os pares.
\end{itensq}
\resposta{(a) $\SI{18}{\watt}$ nos dois; (b) $\SI{15.36}{\watt}$ nos dois;
(c) $R_{L1}R_{L2}=R_{Th}^{2}$}
\resolucao{(a) $R_L=2$: $I=\SI{3}{\ampere}$, $P=\SI{18}{\watt}$.
$R_L=18$: $I=\SI{1}{\ampere}$, $P=\SI{18}{\watt}$.\\
(b) $R_L=1{,}5$: $I=24/7{,}5=\SI{3.2}{\ampere}$,
$P=1{,}5(10{,}24)=\SI{15.36}{\watt}$.
$R_L=24$: $I=\SI{0.8}{\ampere}$, $P=24(0{,}64)=\SI{15.36}{\watt}$.\\
(c) Em ambos os casos, $R_{L1}R_{L2}=R_{Th}^{2}=36$:
\[
2\times18=36;\qquad 1{,}5\times24=36.
\]
Em termos de $x=R_L/R_{Th}$, os pares são $x$ e $1/x$ --- \textbf{recíprocos}.\\
\textbf{Consequência para medição:} se uma varredura de carga encontra a mesma
potência em dois valores $R_{L1}$ e $R_{L2}$, então
$R_{Th}=\sqrt{R_{L1}R_{L2}}$ --- a \emph{média geométrica}. É um método de
determinar $R_{Th}$ sem localizar o máximo com precisão, o que é vantajoso
justamente porque o máximo é achatado.}
\end{questao}

\begin{questao}[3]{}
Determine a faixa de $R_L/R_{Th}$ que garante pelo menos $90\%$ da potência
máxima.
\begin{itensq}
\item Monte a equação e resolva.
\item Interprete a largura da faixa.
\end{itensq}
\resposta{$0{,}52\le x\le1{,}93$ --- razão de quase $4{:}1$}
\resolucao{(a) Impondo $\dfrac{4x}{(1+x)^{2}}=0{,}9$:
\[
4x=0{,}9(1+x)^{2}
\;\Longrightarrow\;
0{,}9x^{2}-2{,}2x+0{,}9=0,
\]
\[
x=\frac{2{,}2\pm\sqrt{4{,}84-3{,}24}}{1{,}8}
=\frac{2{,}2\pm1{,}265}{1{,}8}
\;\Longrightarrow\;
x=0{,}520\ \text{ou}\ 1{,}925.
\]
Note que $0{,}520\times1{,}925=1{,}000$ --- os limites são recíprocos, como a
simetria exige.\\
(b) \textbf{A faixa cobre uma razão de $3{,}7{:}1$.} Com
$R_{Th}=\SI{6}{\ohm}$, qualquer carga entre $\SI{3.1}{\ohm}$ e
$\SI{11.6}{\ohm}$ entrega ao menos $90\%$ do máximo.\\
\textbf{Implicação prática:} não vale a pena esforço para casar cargas com
precisão melhor que uns $20\%$. O máximo é achatado porque a função tem
derivada nula ali --- perto do topo, a curva é essencialmente uma parábola
suave. Esse é o motivo pelo qual "casamento aproximado" costuma bastar.}
\end{questao}

\begin{questao}[3]{}
Construa a tabela de rendimento e de fração de potência máxima para
$x=R_L/R_{Th}$ igual a $0{,}5$; $1$; $2$; $4$ e $9$, e interprete.
\resposta{$\eta$: $33$, $50$, $67$, $80$ e $90\%$; $P/P_{\max}$: $89$, $100$,
$89$, $64$ e $36\%$}
\resolucao{Com $\eta=\dfrac{x}{1+x}$ e
$\dfrac{P_L}{P_{\max}}=\dfrac{4x}{(1+x)^{2}}$:
\begin{center}
\begin{tabular}{ccc}
\hline
$x=R_L/R_{Th}$ & $\eta$ & $P_L/P_{\max}$\\
\hline
$0{,}5$ & $33{,}3\%$ & $88{,}9\%$\\
$1$ & $50{,}0\%$ & $100\%$\\
$2$ & $66{,}7\%$ & $88{,}9\%$\\
$4$ & $80{,}0\%$ & $64{,}0\%$\\
$9$ & $90{,}0\%$ & $36{,}0\%$\\
\hline
\end{tabular}
\end{center}
\textbf{Três leituras.}\\
\emph{(i)} $x=0{,}5$ e $x=2$ dão a mesma potência ($88{,}9\%$) mas rendimentos
opostos ($33\%$ e $67\%$). \textbf{Entre dois pares recíprocos, escolha sempre
o maior} --- mesma entrega, metade da perda.\\
\emph{(ii)} As duas curvas puxam em sentidos contrários: $\eta$ cresce
monotonicamente com $x$, enquanto $P_L/P_{\max}$ tem máximo em $x=1$ e decresce.
Não existe ponto que otimize ambos.\\
\emph{(iii)} A região $2\le x\le4$ é o compromisso usual: ainda se entrega de
$64\%$ a $89\%$ da potência disponível, com rendimento entre $67\%$ e
$80\%$.}
\end{questao}

\begin{questao}[3]{}
Uma carga fixa de $\SI{8}{\ohm}$ é alimentada por uma fonte de
$\SI{24}{\volt}$ cuja resistência interna pode ser reduzida por projeto.
\begin{itensq}
\item Calcule $P_L$ para $R_{Th}=16$, $8$, $4$, $2$ e $\SI{0}{\ohm}$.
\item Determine o valor ótimo de $R_{Th}$.
\item Explique por que a resposta não é $R_{Th}=R_L$.
\end{itensq}
\resposta{(b) $R_{Th}\to0$; a potência cresce monotonicamente conforme
$R_{Th}$ diminui}
\resolucao{(a)
\begin{center}
\begin{tabular}{cc}
\hline
$R_{Th}$ & $P_L$\\
\hline
$\SI{16}{\ohm}$ & $\SI{8}{\watt}$\\
$\SI{8}{\ohm}$ & $\SI{18}{\watt}$\\
$\SI{4}{\ohm}$ & $\SI{32}{\watt}$\\
$\SI{2}{\ohm}$ & $\SI{46.1}{\watt}$\\
$\SI{0}{\ohm}$ & $\SI{72}{\watt}$\\
\hline
\end{tabular}
\end{center}
(b) A potência cresce sem parar conforme $R_{Th}$ diminui: o ótimo é
$R_{Th}\to0$, no \emph{contorno} do domínio, não em um ponto interior.\\
(c) \textbf{A condição $R_L=R_{Th}$ não é simétrica.} Ela resolve o problema
"dada a fonte, escolha a carga". Aqui o problema é o inverso --- "dada a carga,
escolha a fonte" --- e a função
\[
P_L=\frac{V_{Th}^{2}R_L}{(R_L+R_{Th})^{2}}
\]
é \textbf{monotonicamente decrescente} em $R_{Th}$: sua derivada,
$-\dfrac{2V_{Th}^{2}R_L}{(R_L+R_{Th})^{3}}$, nunca se anula para
$R_{Th}\ge0$.\\
\textbf{Regra:} o casamento de impedâncias só faz sentido quando a fonte é
\emph{dada e imutável} --- uma antena, um transdutor, um sensor. Se a fonte pode
ser projetada, sempre convém $R_{Th}$ o menor possível.}
\end{questao}

\begin{questao}[3]{}
Uma bateria de $\EMF=\SI{12}{\volt}$ e $r=\SI{0.5}{\ohm}$ alimenta um motor de
partida.
\begin{itensq}
\item Determine a potência máxima disponível e a carga correspondente.
\item Determine a corrente nessa condição e avalie sua viabilidade.
\item Determine a potência com $R_L=\SI{2}{\ohm}$ e compare.
\end{itensq}
\resposta{(a) $\SI{72}{\watt}$ com $R_L=\SI{0.5}{\ohm}$; (b)
$\SI{12}{\ampere}$; (c) $\SI{46.1}{\watt}$ com $\eta=80\%$}
\resolucao{(a) $P_{\max}=\dfrac{144}{4(0{,}5)}=\SI{72}{\watt}$ com
$R_L=\SI{0.5}{\ohm}$.\\
(b) $I=\dfrac{12}{1}=\SI{12}{\ampere}$, com $\SI{72}{\watt}$ dissipados
\emph{dentro} da bateria. Uma bateria pequena não suporta isso por muito tempo:
o eletrólito aquece, $r$ cresce, e a potência disponível cai durante a própria
operação.\\
(c) $I=\dfrac{12}{2{,}5}=\SI{4.8}{\ampere}$;
$P_L=2(4{,}8)^{2}=\SI{46.1}{\watt}$; $\eta=2/2{,}5=80\%$.\\
\textbf{Comparação decisiva:} entrega-se $64\%$ da potência máxima, mas a perda
interna cai de $\SI{72}{\watt}$ para $\SI{11.5}{\watt}$ --- \emph{seis vezes}
menos aquecimento.\\
\textbf{Além disso:} o modelo de $r$ constante já é otimista a
$\SI{12}{\ampere}$. Baterias reais apresentam $r$ crescente com a corrente e
com o estado de descarga, o que torna a operação casada ainda menos atraente do
que o cálculo sugere.}
\end{questao}

\begin{questao}[3]{}
Uma fonte alimenta a carga através de uma rede resistiva intermediária. A fonte
tem $\SI{36}{\volt}$ e $\SI{4}{\ohm}$; a rede é um resistor de
$\SI{12}{\ohm}$ em série e um de $\SI{6}{\ohm}$ em paralelo com a saída.
\begin{itensq}
\item Determine o equivalente de Thévenin visto pela carga.
\item Determine $R_L$ ótimo e a potência máxima.
\item Compare com a potência máxima disponível na fonte original.
\end{itensq}
\resposta{(a) $V_{Th}=\SI{9.82}{\volt}$, $R_{Th}=\SI{4.36}{\ohm}$;
(b) $R_L=\SI{4.36}{\ohm}$ e $P_{\max}=\SI{5.53}{\watt}$; (c) contra
$\SI{81}{\watt}$ disponíveis na fonte}
\resolucao{(a) Com a carga removida, a rede é um divisor de
$4+12=\SI{16}{\ohm}$ e $\SI{6}{\ohm}$:
\[
V_{Th}=36\cdot\frac{6}{22}=\SI{9.82}{\volt};
\qquad
R_{Th}=\frac{16\cdot6}{22}=\SI{4.36}{\ohm}.
\]
(b) $R_L=\SI{4.36}{\ohm}$ e
\[
P_{\max}=\frac{(9{,}82)^{2}}{4(4{,}36)}=\SI{5.53}{\watt}.
\]
(c) A fonte original disponibilizaria
$\dfrac{36^{2}}{4(4)}=\SI{81}{\watt}$.\\
\textbf{A rede intermediária destruiu $93\%$ da potência disponível.} O
resistor série de $\SI{12}{\ohm}$ derruba a tensão e o de $\SI{6}{\ohm}$ desvia
corrente --- ambos dissipando.\\
\textbf{Conclusão:} redes resistivas \emph{não} casam impedâncias sem perda.
Elas alteram $R_{Th}$, sim, mas sempre ao custo de reduzir $V_{Th}$ e,
com ele, a potência disponível. Casamento sem perda exige elementos reativos
(transformador, rede LC) --- possíveis apenas em corrente alternada.}
\end{questao}

\begin{questao}[3]{}
Uma carga só pode ser escolhida entre $\SI{12}{\ohm}$ e $\SI{40}{\ohm}$
(valores disponíveis em estoque). A fonte tem $V_{Th}=\SI{24}{\volt}$ e
$R_{Th}=\SI{6}{\ohm}$.
\begin{itensq}
\item Determine qual escolher para máxima potência.
\item Determine qual escolher para máximo rendimento.
\item Discuta o critério de decisão.
\end{itensq}
\resposta{(a) $\SI{12}{\ohm}$ ($\SI{21.3}{\watt}$); (b) $\SI{40}{\ohm}$
($\eta=87\%$)}
\resolucao{(a)
\[
R_L=12:\ I=\frac{24}{18}=1{,}333;\ P_L=12(1{,}778)=\SI{21.3}{\watt};
\]
\[
R_L=40:\ I=\frac{24}{46}=0{,}522;\ P_L=40(0{,}272)=\SI{10.9}{\watt}.
\]
A de $\SI{12}{\ohm}$ é a mais próxima do ótimo ($x=2$) e entrega quase o dobro.\\
(b) $\eta=\dfrac{R_L}{R_L+6}$: $67\%$ contra $87\%$. A de $\SI{40}{\ohm}$
vence.\\
(c) \textbf{O critério depende do que é escasso.}
\begin{itemize}
\item Se a fonte é abundante (rede elétrica) e o custo está na energia, escolha
o \emph{rendimento}.
\item Se a fonte é limitada (bateria, sensor, antena) e o que importa é extrair
o máximo de um sinal fraco, escolha a \emph{potência}.
\item Se a limitação é térmica na própria fonte, escolha rendimento --- a de
$\SI{40}{\ohm}$ dissipa apenas $\SI{1.6}{\watt}$ internamente, contra
$\SI{10.7}{\watt}$ da outra.
\end{itemize}
Sem saber o que é escasso, a pergunta "qual é melhor" não tem resposta.}
\end{questao}

\begin{questao}[3]{}
Uma fonte de $\SI{24}{\volt}$ e $R_{Th}=\SI{6}{\ohm}$ tem sua corrente limitada
eletronicamente a $\SI{1.5}{\ampere}$.
\begin{itensq}
\item Determine se o ponto ótimo teórico é atingível.
\item Determine a carga que maximiza a potência sob essa restrição.
\item Compare com o máximo irrestrito.
\end{itensq}
\resposta{(a) não --- exigiria $\SI{2}{\ampere}$; (b)
$R_L=\SI{10}{\ohm}$; (c) $\SI{22.5}{\watt}$ contra $\SI{24}{\watt}$}
\resolucao{(a) O ótimo teórico exige $I=V_{Th}/(2R_{Th})=\SI{2}{\ampere}$,
acima do limite de $\SI{1.5}{\ampere}$. \textbf{O ponto ótimo é inatingível.}\\
(b) Como $P_L=R_LI^{2}$ cresce quando $R_L$ se aproxima de $R_{Th}$ por cima, o
melhor admissível é a \emph{menor} carga que ainda respeite o limite:
\[
I=\frac{24}{6+R_L}\le1{,}5
\;\Longrightarrow\;
6+R_L\ge16
\;\Longrightarrow\;
R_L\ge\SI{10}{\ohm}.
\]
Adotando $R_L=\SI{10}{\ohm}$: $I=\SI{1.5}{\ampere}$ e
$P_L=10(2{,}25)=\SI{22.5}{\watt}$.\\
(c) Isso é $93{,}8\%$ do máximo irrestrito --- perda de apenas $6\%$.\\
\textbf{Duas observações.} O ótimo restrito fica na \emph{fronteira} do domínio
admissível, não onde a derivada se anula: é a marca de um problema com
restrição ativa. E o custo da restrição é pequeno justamente porque o máximo é
achatado --- a mesma propriedade que torna o casamento tolerante torna as
restrições baratas.}
\end{questao}

\begin{questao}[3]{}
Determine o rendimento necessário para que uma fonte entregue $\SI{64}{\percent}$
da potência máxima, e verifique a coerência com a curva $\eta(x)$.
\resposta{$\eta=80\%$ (com $x=4$) ou $\eta=20\%$ (com $x=0{,}25$)}
\resolucao{Impondo $\dfrac{4x}{(1+x)^{2}}=0{,}64$:
\[
0{,}64x^{2}-2{,}72x+0{,}64=0
\;\Longrightarrow\;
x=4\ \text{ou}\ x=0{,}25,
\]
recíprocos, como esperado.\\
\textbf{Rendimentos:}
\[
x=4:\ \eta=\frac{4}{5}=80\%;
\qquad
x=0{,}25:\ \eta=\frac{0{,}25}{1{,}25}=20\%.
\]
\textbf{A mesma potência, dois rendimentos opostos} --- e a soma é
$100\%$, o que não é coincidência: se $x$ e $1/x$ são recíprocos, então
$\dfrac{x}{1+x}+\dfrac{1/x}{1+1/x}=\dfrac{x}{1+x}+\dfrac{1}{1+x}=1$.\\
\textbf{Decisão óbvia:} escolha sempre o ramo $x>1$. A escolha $x=0{,}25$
entregaria a mesma potência à carga dissipando \emph{quatro vezes} mais na
fonte.}
\end{questao}

\blocodesafio

\begin{questao}[4]{}
\textbf{(Demonstração completa.)} Prove a condição de máxima transferência de
potência.
\begin{itensq}
\item Derive $P_L(R_L)$ e obtenha o ponto crítico.
\item Confirme que é máximo pela segunda derivada ou pela análise dos limites.
\item Obtenha $P_{\max}$ e o rendimento correspondente.
\end{itensq}
\resposta{$R_L=R_{Th}$; $P_{\max}=\dfrac{V_{Th}^{2}}{4R_{Th}}$;
$\eta=50\%$}
\resolucao{(a) Com $I=\dfrac{V_{Th}}{R_{Th}+R_L}$:
\[
P_L=R_LI^{2}=\frac{V_{Th}^{2}R_L}{(R_{Th}+R_L)^{2}}.
\]
Derivando pela regra do quociente:
\[
\frac{\mathrm{d}P_L}{\mathrm{d}R_L}
=V_{Th}^{2}\,\frac{(R_{Th}+R_L)^{2}-R_L\cdot2(R_{Th}+R_L)}{(R_{Th}+R_L)^{4}}
=V_{Th}^{2}\,\frac{R_{Th}-R_L}{(R_{Th}+R_L)^{3}}.
\]
Anulando: $R_L=R_{Th}$.\\
(b) \textbf{Pelo sinal da derivada}, que é mais simples que a segunda derivada:
o denominador é sempre positivo, logo o sinal é o de $(R_{Th}-R_L)$ ---
\emph{positivo} para $R_L<R_{Th}$ e \emph{negativo} para $R_L>R_{Th}$. A função
cresce, atinge o topo e decresce: máximo confirmado.\\
\textbf{Confirmação pelos limites:} $P_L\to0$ quando $R_L\to0$ (tensão nula na
carga) e quando $R_L\to\infty$ (corrente nula). Entre dois zeros, uma função
positiva e suave tem máximo interior.\\
(c)
\[
P_{\max}=\frac{V_{Th}^{2}R_{Th}}{(2R_{Th})^{2}}=\frac{V_{Th}^{2}}{4R_{Th}};
\qquad
\eta=\frac{R_L}{R_L+R_{Th}}=\frac{1}{2}.
\]
\textbf{Observação estrutural:} a mesma derivada aparece na potência de um
resistor em série (máxima quando ele iguala o restante) e na sensibilidade do
divisor de tensão. É a mesma função $R/(R+a)^{2}$ em três contextos --- vale
reconhecê-la.}
\end{questao}

\begin{questao}[4]{}
\textbf{(Simetria dos pares recíprocos.)} Prove que $R_L$ e
$R_{Th}^{2}/R_L$ entregam a mesma potência.
\begin{itensq}
\item Demonstre algebricamente.
\item Mostre que os rendimentos correspondentes somam $100\%$.
\item Deduza um método para obter $R_{Th}$ a partir de duas medições.
\end{itensq}
\resposta{$P(x)=P(1/x)$; $\eta(x)+\eta(1/x)=1$;
$R_{Th}=\sqrt{R_{L1}R_{L2}}$}
\resolucao{(a) Em variável adimensional $x=R_L/R_{Th}$:
\[
\frac{P_L}{P_{\max}}=\frac{4x}{(1+x)^{2}}.
\]
Substituindo $x\to1/x$:
\[
\frac{4/x}{(1+1/x)^{2}}
=\frac{4/x}{\dfrac{(x+1)^{2}}{x^{2}}}
=\frac{4x}{(1+x)^{2}}.
\]
Idêntico --- a função é \textbf{invariante} sob $x\to1/x$, e o ponto fixo dessa
inversão é $x=1$, que é justamente o máximo.\\
(b)
\[
\eta(x)+\eta(1/x)=\frac{x}{1+x}+\frac{1/x}{1+1/x}
=\frac{x}{1+x}+\frac{1}{1+x}=1.
\]
Os rendimentos são complementares: $33\%$ e $67\%$; $20\%$ e $80\%$; e assim
por diante.\\
(c) \textbf{Método de medição.} Varie a carga e encontre \emph{dois} valores
$R_{L1}\ne R_{L2}$ que produzam a mesma potência. Como
$\dfrac{R_{L1}}{R_{Th}}=\dfrac{R_{Th}}{R_{L2}}$, segue
\[
R_{Th}=\sqrt{R_{L1}R_{L2}}\quad\text{(média geométrica)}.
\]
\textbf{Vantagem sobre buscar o máximo:} perto do topo a curva é chata, e
localizar o pico com precisão é difícil. Já os dois pontos de igual potência
podem ser tomados bem afastados do máximo, onde a curva é \emph{íngreme} --- e a
determinação fica muito mais precisa. É um caso em que fugir do ponto de
interesse melhora a medição.}
\end{questao}

\begin{questao}[4]{}
\textbf{(Tolerância em torno do ótimo.)} Deduza a faixa de $x=R_L/R_{Th}$ que
garante uma fração $f$ da potência máxima.
\begin{itensq}
\item Obtenha a equação e sua solução geral.
\item Construa a tabela para $f=0{,}90$, $0{,}95$ e $0{,}99$.
\item Interprete o comportamento da largura da faixa.
\end{itensq}
\resposta{$x=\dfrac{2-f\pm2\sqrt{1-f}}{f}$; faixas de $3{,}7{:}1$,
$2{,}5{:}1$ e $1{,}5{:}1$}
\resolucao{(a) De $\dfrac{4x}{(1+x)^{2}}=f$:
\[
fx^{2}+(2f-4)x+f=0
\;\Longrightarrow\;
x=\frac{(4-2f)\pm\sqrt{(2f-4)^{2}-4f^{2}}}{2f}
=\frac{2-f\pm2\sqrt{1-f}}{f}.
\]
Como o produto das raízes vale $f/f=1$, elas são sempre \textbf{recíprocas} ---
confirmação independente da simetria.\\
(b)
\begin{center}
\begin{tabular}{cccc}
\hline
$f$ & $x_{\min}$ & $x_{\max}$ & razão\\
\hline
$0{,}90$ & $0{,}520$ & $1{,}925$ & $3{,}70$\\
$0{,}95$ & $0{,}635$ & $1{,}576$ & $2{,}48$\\
$0{,}99$ & $0{,}818$ & $1{,}222$ & $1{,}49$\\
\hline
\end{tabular}
\end{center}
(c) \textbf{A largura encolhe como $\sqrt{1-f}$.} Passar de $90\%$ para
$99\%$ da potência exige apertar a faixa por um fator $2{,}5$; chegar a
$99{,}9\%$ exigiria outro fator $3$.\\
\textbf{Origem:} perto do máximo a derivada é nula, e a curva se comporta como
uma parábola --- a perda cresce com o \emph{quadrado} do desvio relativo.
Consequentemente, um desvio de $10\%$ na carga custa cerca de $0{,}25\%$ da
potência.\\
\textbf{Conclusão de engenharia:} casar impedâncias com precisão melhor que
$20\%$ raramente compensa. O esforço deve ir para reduzir $R_{Th}$, não para
refinar o casamento.}
\end{questao}

\begin{questao}[4]{}
\textbf{(Projeto de fonte.)} Projete uma fonte que entregue no máximo
$\SI{25}{\watt}$ a uma carga de $\SI{4}{\ohm}$.
\begin{itensq}
\item Determine $V_{Th}$ e $R_{Th}$.
\item Determine a corrente, a tensão de saída e a potência total.
\item Especifique a dissipação interna e discuta o dimensionamento.
\end{itensq}
\resposta{$V_{Th}=\SI{20}{\volt}$, $R_{Th}=\SI{4}{\ohm}$;
$I=\SI{2.5}{\ampere}$; $P_{tot}=\SI{50}{\watt}$}
\resolucao{(a) Para que $\SI{4}{\ohm}$ seja a carga ótima, é preciso
$R_{Th}=\SI{4}{\ohm}$. Então
\[
P_{\max}=\frac{V_{Th}^{2}}{4(4)}=25
\;\Longrightarrow\;
V_{Th}=\sqrt{400}=\SI{20}{\volt}.
\]
(b) $I=\dfrac{20}{8}=\SI{2.5}{\ampere}$;
$U_L=\SI{10}{\volt}$; $P_{tot}=20(2{,}5)=\SI{50}{\watt}$.\\
(c) \textbf{A fonte dissipa $\SI{25}{\watt}$ internamente} --- tanto quanto
entrega. Isso significa:
\begin{itemize}
\item O elemento que realiza $R_{Th}=\SI{4}{\ohm}$ precisa ser dimensionado para
$\SI{25}{\watt}$ contínuos (especifique $\SI{50}{\watt}$, com dissipador).
\item A fonte primária deve fornecer $\SI{50}{\watt}$ e
$\SI{2.5}{\ampere}$.
\item O rendimento global é $50\%$ por construção.
\end{itemize}
\textbf{Crítica ao enunciado do projeto.} Deliberadamente inserir
$R_{Th}=\SI{4}{\ohm}$ para "casar" é quase sempre um erro. Se o objetivo é
entregar $\SI{25}{\watt}$ a $\SI{4}{\ohm}$, basta uma fonte de
$\SI{10}{\volt}$ com $R_{Th}$ pequeno: ela entrega os mesmos
$\SI{25}{\watt}$ consumindo pouco mais que isso.\\
O casamento só se justifica quando $R_{Th}$ é \emph{imposto} pela física do
dispositivo --- resistência de saída de um transdutor, impedância característica
de uma linha --- e não quando é escolhido por projeto.}
\end{questao}

\begin{questao}[4]{}
\textbf{(Otimização com restrição de rendimento.)} Maximize a potência
entregue, com a restrição $\eta\ge80\%$, para
$V_{Th}=\SI{24}{\volt}$ e $R_{Th}=\SI{6}{\ohm}$.
\begin{itensq}
\item Traduza a restrição em uma condição sobre $x$.
\item Determine o ótimo restrito.
\item Compare com o ótimo irrestrito e interprete.
\end{itensq}
\resposta{$x\ge4$; $R_L=\SI{24}{\ohm}$ com $P_L=\SI{15.36}{\watt}$
($64\%$ do máximo)}
\resolucao{(a) $\eta=\dfrac{x}{1+x}\ge0{,}8
\Rightarrow x\ge0{,}8+0{,}8x
\Rightarrow 0{,}2x\ge0{,}8
\Rightarrow \boxed{x\ge4}$.\\
(b) No domínio admissível $x\ge4$, a função
$\dfrac{4x}{(1+x)^{2}}$ é \emph{decrescente} (pois já se passou do máximo em
$x=1$). O melhor valor está, portanto, na \textbf{fronteira}: $x=4$, isto é
$R_L=\SI{24}{\ohm}$.
\[
I=\frac{24}{30}=\SI{0.8}{\ampere};
\qquad
P_L=24(0{,}64)=\SI{15.36}{\watt};
\qquad
\eta=80\%\ \checkmark
\]
(c) O máximo irrestrito era $\SI{24}{\watt}$ com $\eta=50\%$. A restrição
custou $36\%$ da potência e comprou $30$ pontos de rendimento.\\
\textbf{Estrutura do problema.} Como as duas curvas são monótonas em sentidos
opostos na região $x>1$, \textbf{a restrição de rendimento é sempre ativa} --- o
ótimo cai exatamente sobre ela, nunca no interior. Isso simplifica muito o
projeto: basta resolver a igualdade $\eta=\eta_{\min}$.\\
\textbf{Generalização:} para rendimento mínimo $\eta_0$, o ótimo é
$x=\dfrac{\eta_0}{1-\eta_0}$, e a fração de potência obtida é
$4\eta_0(1-\eta_0)$. Para $\eta_0=0{,}5$ isso dá $100\%$ (o casamento); para
$\eta_0=0{,}9$, apenas $36\%$.}
\end{questao}

\begin{questao}[4]{}
\textbf{(Potência sob dois graus de liberdade.)} Analise a maximização de
$P_L$ quando \emph{ambos} $R_L$ e $R_{Th}$ podem variar, com $V_{Th}$ fixo.
\begin{itensq}
\item Mostre que não existe máximo interior.
\item Determine o comportamento no contorno.
\item Explique a relevância prática dessa análise.
\end{itensq}
\resposta{Não há máximo interior; $P_L\to\infty$ com
$R_{Th}\to0$ e $R_L\to0$}
\resolucao{(a) Com
$P_L=\dfrac{V_{Th}^{2}R_L}{(R_L+R_{Th})^{2}}$, as duas derivadas parciais são
\[
\frac{\partial P_L}{\partial R_L}=V_{Th}^{2}\frac{R_{Th}-R_L}{(R_L+R_{Th})^{3}};
\qquad
\frac{\partial P_L}{\partial R_{Th}}=-\frac{2V_{Th}^{2}R_L}{(R_L+R_{Th})^{3}}.
\]
A segunda é \textbf{estritamente negativa} para todo $R_L>0$ --- ela nunca se
anula. Não existe ponto em que ambas as derivadas sejam nulas, logo não há
máximo interior.\\
(b) O ótimo migra para o contorno $R_{Th}\to0$. Nesse limite,
\[
P_L\to\frac{V_{Th}^{2}}{R_L},
\]
que por sua vez cresce sem limite conforme $R_L\to0$. \textbf{A potência
diverge} --- o modelo perde sentido físico, pois pressupõe uma fonte de tensão
ideal capaz de fornecer corrente infinita.\\
(c) \textbf{Relevância.} A análise revela que "casamento de impedâncias" não é
um princípio universal de otimização, mas a solução de um problema \emph{muito
específico}: fonte fixa, carga livre.
\begin{itemize}
\item \emph{Fonte fixa, carga livre} $\Rightarrow$ $R_L=R_{Th}$ (casamento).
\item \emph{Carga fixa, fonte livre} $\Rightarrow$ $R_{Th}\to0$ (fonte rígida).
\item \emph{Ambos livres} $\Rightarrow$ problema mal posto, sem solução finita.
\end{itemize}
\textbf{Onde cada caso ocorre:} o primeiro em antenas, transdutores e sensores,
cuja impedância é imposta pela física. O segundo em fontes de alimentação, em
que o projeto persegue $R_{Th}$ o menor possível. Confundir os dois é o erro
conceitual mais comum do tópico --- e leva à ideia equivocada de que fontes de
alimentação deveriam ser "casadas" com suas cargas.}
\end{questao}

\begin{questao}[4]{}
\textbf{(Análise de um sistema real.)} Um painel fotovoltaico apresenta, sob
certa irradiância, $V_{oc}=\SI{22}{\volt}$ e $I_{sc}=\SI{5}{\ampere}$, com
ponto de máxima potência em $\SI{18}{\volt}$ e $\SI{4.5}{\ampere}$.
\begin{itensq}
\item Compare a potência máxima real com a prevista pelo modelo de Thévenin.
\item Explique a discrepância.
\item Discuta a implicação para o rastreamento do ponto de máxima potência.
\end{itensq}
\resposta{$\SI{81}{\watt}$ reais contra $\SI{27.5}{\watt}$ previstos --- o
painel não é uma fonte linear}
\resolucao{(a) \textbf{Modelo de Thévenin:}
$R_{Th}=22/5=\SI{4.4}{\ohm}$ e
\[
P_{\max}^{Th}=\frac{V_{oc}I_{sc}}{4}=\frac{22(5)}{4}=\SI{27.5}{\watt}.
\]
\textbf{Real:} $P=18(4{,}5)=\SI{81}{\watt}$ --- quase \textbf{três vezes}
mais.\\
(b) \textbf{O painel não é linear.} Sua curva $I\times V$ não é uma reta: ela é
quase horizontal em $I\approx I_{sc}$ até perto de $V_{oc}$, onde despenca
abruptamente. O modelo de Thévenin, que supõe a reta
$I=I_{sc}(1-V/V_{oc})$, subestima grosseiramente a área sob a curva.\\
A razão $\dfrac{P_{\max}}{V_{oc}I_{sc}}$ --- chamada \textbf{fator de forma} ---
vale $0{,}25$ para uma fonte linear e $0{,}74$ para este painel. Painéis bons
chegam a $0{,}80$.\\
(c) \textbf{Implicações para o rastreamento.}
\begin{itemize}
\item A carga ótima é $R_L=18/4{,}5=\SI{4}{\ohm}$ --- próxima de
$R_{Th}=\SI{4.4}{\ohm}$ por coincidência numérica, não por casamento.
\item O máximo é \textbf{muito mais agudo} que o de uma fonte linear, pois a
curva tem um "joelho" pronunciado. A tolerância generosa deduzida para fontes
lineares \emph{não} se aplica: errar a carga custa caro.
\item $V_{oc}$ e $I_{sc}$ variam com irradiância e temperatura, deslocando o
ponto ótimo continuamente ao longo do dia.
\end{itemize}
\textbf{Consequência:} é necessário um rastreador ativo (MPPT), que ajusta a
carga aparente em tempo real. E toda a teoria deste capítulo serve aqui apenas
como \emph{referência conceitual} --- o dimensionamento exige a curva real do
dispositivo.}
\end{questao}

% END BANCO COMPLEMENTAR

\gabarito[Capítulo 5 --- Potência e máxima transferência]
               % Potencia e maxima transferencia

% Cap. 6 -------------------------------------------------------------
\setcounter{section}{5}
% =====================================================================
%  CAPITULO 6 - TEOREMA DA SUPERPOSICAO
% =====================================================================
\section{Teorema da Superposição}

\begin{conceito}[Antes de começar]
Em um circuito \textbf{linear} com várias fontes independentes, a resposta
(corrente ou tensão) em qualquer elemento é a \textbf{soma algébrica} das
respostas produzidas por cada fonte agindo \emph{sozinha}.\par
\textbf{Como "desligar" uma fonte:}
\begin{itemize}[leftmargin=1.6em,itemsep=1pt,topsep=2pt]
\item fonte de \emph{tensão} ideal $\to$ substituída por um \textbf{curto-circuito}
      (fio);
\item fonte de \emph{corrente} ideal $\to$ substituída por um \textbf{circuito
      aberto}.
\end{itemize}
Resolve-se o circuito uma vez para cada fonte e somam-se os resultados, respeitando
os \emph{sinais} (sentidos adotados).
\end{conceito}

\begin{atencao}[Onde a superposição NÃO se aplica]
A superposição vale para tensões e correntes, que são lineares. \textbf{Não} vale
para \emph{potência}, pois $P = R\,I^2$ é quadrática:
$P \neq P_1 + P_2$. A potência deve ser calculada com a corrente \emph{total}.
\end{atencao}

\legendaniveis

\blocofixacao

\begin{questao}[1]{}
Ao aplicar o teorema da superposição, para "desligar" uma fonte de tensão ideal e
uma fonte de corrente ideal, deve-se, respectivamente:
\begin{alternativas}
\item abrir a de tensão e curto-circuitar a de corrente.
\item curto-circuitar a de tensão e abrir a de corrente.
\item abrir ambas.
\item curto-circuitar ambas.
\item remover ambas do circuito.
\end{alternativas}
\resposta{(b)}
\resolucao{Uma fonte de tensão ideal desligada mantém seus terminais no mesmo
potencial (tensão nula) --- equivale a um \emph{curto}. Uma fonte de corrente
ideal desligada não injeta corrente --- equivale a um \emph{circuito aberto}.}
\end{questao}

\begin{questao}[1]{}
O teorema da superposição pode ser aplicado diretamente para calcular:
\begin{alternativas}
\item a potência dissipada em um resistor.
\item a corrente e a tensão em um elemento.
\item apenas correntes, nunca tensões.
\item circuitos com componentes não lineares.
\item a energia armazenada em um capacitor.
\end{alternativas}
\resposta{(b)}
\resolucao{A superposição vale para grandezas \emph{lineares} --- correntes e
tensões --- em circuitos lineares. Potência (alternativa a) e energia
(alternativa e) são quadráticas e não podem ser superpostas. Componentes não
lineares (d) violam a hipótese do teorema.}
\end{questao}

\blocoaplicacao

\begin{questao}[2]{}
No circuito da figura, use a superposição para determinar a tensão $V$ no nó
central (sobre o resistor de $\SI{3}{\ohm}$).

\circuitoq{%
\begin{circuitikz}[scale=1,transform shape,line width=0.8pt]
 \draw (0,0) to[battery1,l_={$E_1=\SI{12}{\volt}$}] (0,3)
       to[R,l={$\SI{6}{\ohm}$}] (3,3) coordinate(n);
 \draw (n) to[R,l={$\SI{3}{\ohm}$}] (3,0) coordinate(b) -- (0,0);
 \draw (n) to[R,l={$\SI{6}{\ohm}$}] (6,3)
       to[battery1,l={$E_2=\SI{6}{\volt}$},invert] (6,0) -- (b);
 \node[above,Vcol] at (3,3.05) {$V$};
\end{circuitikz}}

\resposta{$V = \SI{4.5}{\volt}$}
\resolucao{\textbf{Só $E_1$} (curto em $E_2$): o nó vê $E_1$ através de
\SI{6}{\ohm} e, para o terra, \SI{3}{\ohm} em paralelo com \SI{6}{\ohm}
($3\parallel6 = \SI{2}{\ohm}$):
\[
V_1 = E_1\cdot\frac{2}{6+2}=12\cdot\frac{2}{8}=\SI{3}{\volt}.
\]
\textbf{Só $E_2$} (curto em $E_1$), por simetria do outro lado:
\[
V_2 = E_2\cdot\frac{3\parallel6}{6+(3\parallel6)}=6\cdot\frac{2}{8}=\SI{1.5}{\volt}.
\]
\textbf{Superpondo:} $V = V_1 + V_2 = 3 + 1{,}5 = \SI{4.5}{\volt}$.\\
\textbf{Verificação por nodal:}
$\dfrac{12-V}{6}+\dfrac{6-V}{6}=\dfrac{V}{3}\Rightarrow V=\SI{4.5}{\volt}$. Resultado consistente.}
\end{questao}

\begin{questao}[2]{}
No mesmo circuito da questão anterior, determine a corrente que atravessa o
resistor de $\SI{3}{\ohm}$.
\resposta{$I_{3\Omega} = \SI{1.5}{\ampere}$}
\resolucao{Conhecida a tensão do nó, $V = \SI{4.5}{\volt}$, e como o resistor de
\SI{3}{\ohm} liga o nó ao terra:
\[
I_{3\Omega}=\frac{V}{3}=\frac{4{,}5}{3}=\SI{1.5}{\ampere}.
\]
Poderíamos também superpor as correntes: cada fonte contribui com
$V_1/3 = \SI{1}{\ampere}$ e $V_2/3 = \SI{0.5}{\ampere}$, somando \SI{1.5}{\ampere}.}
\end{questao}

\begin{questao}[2]{}
No circuito da figura há uma fonte de \emph{tensão} e uma de \emph{corrente}. Use
a superposição para determinar o potencial $V_A$ do nó central.

\circuitoq{%
\begin{circuitikz}[scale=1,transform shape,line width=0.8pt,american]
 \draw (0,0) to[battery1,l_={$\SI{12}{\volt}$}] (0,3) -- (1.5,3)
       to[R,l={$R_1=\SI{5}{\ohm}$}] (3.5,3) coordinate(a);
 \node[above,Vcol] at (3.5,3.15){$V_A$};
 \draw (a) to[R,l={$R_2=\SI{5}{\ohm}$}] (3.5,0) coordinate(g) -- (0,0);
 \draw (a) -- (5.5,3) to[isource,l={$\SI{2}{\ampere}$}] (5.5,0) -- (g);
 \node[ground] at (3.5,0){};
\end{circuitikz}}

\resposta{$V_A = \SI{11}{\volt}$}
\resolucao{\textbf{Só a fonte de tensão} (abre a fonte de corrente): $R_1$ e $R_2$
formam um divisor,
\[
V_A' = 12\cdot\frac{R_2}{R_1+R_2}=12\cdot\frac{5}{10}=\SI{6}{\volt}.
\]
\textbf{Só a fonte de corrente} (curto-circuita a fonte de tensão): os
\SI{2}{\ampere} passam por $R_1\parallel R_2 = \SI{2.5}{\ohm}$,
\[
V_A'' = 2\cdot2{,}5 = \SI{5}{\volt}.
\]
\textbf{Superpondo:} $V_A = V_A' + V_A'' = 6 + 5 = \SI{11}{\volt}$.\\
\textbf{Verificação por nodal:}
$\dfrac{V_A-12}{5}+\dfrac{V_A}{5} = 2 \Rightarrow 2V_A - 12 = 10
\Rightarrow V_A = \SI{11}{\volt}$. Resultado consistente.}
\end{questao}

\blocoaprofundamento

\begin{questao}[3]{}
Ainda no circuito das questões anteriores, um estudante calculou a potência no
resistor de $\SI{3}{\ohm}$ somando as potências devidas a cada fonte
separadamente: $P_1 = \dfrac{V_1^2}{3}$ e $P_2 = \dfrac{V_2^2}{3}$. Mostre que
esse procedimento está \textbf{errado} e calcule a potência correta.
\resposta{Errado; a potência correta é $P = \SI{6.75}{\watt}$}
\resolucao{O cálculo do estudante daria
$P_1 = \dfrac{3^2}{3} = \SI{3}{\watt}$ e $P_2 = \dfrac{1{,}5^2}{3} = \SI{0.75}{\watt}$,
somando \SI{3.75}{\watt}. Mas a superposição \textbf{não vale para potência}: ela
é quadrática. O correto é usar a tensão \emph{total} $V = \SI{4.5}{\volt}$:
\[
P=\frac{V^2}{3}=\frac{4{,}5^2}{3}=\frac{20{,}25}{3}=\SI{6.75}{\watt}.
\]
A diferença ($6{,}75 \neq 3{,}75$) vem do termo cruzado $2V_1V_2/3$, que a soma
ingênua ignora. \textbf{Regra:} superponha tensões/correntes; só depois calcule a
potência com o valor total.}
\end{questao}

\begin{questao}[3]{}
No circuito da figura há \textbf{três} fontes independentes. Determine o
potencial $V$ do nó central usando o teorema da superposição, e em seguida
verifique o resultado por análise nodal.

\circuitoq{%
\begin{circuitikz}[scale=0.95,transform shape,line width=0.8pt]
 \draw (0,0) to[battery1,l_={$\SI{12}{\volt}$}] (0,3)
       to[R,l={$\SI{6}{\ohm}$}] (2.6,3) coordinate(n);
 \draw (n) to[R,l={$\SI{3}{\ohm}$}] (2.6,0) coordinate(b) -- (0,0);
 \draw (n) to[R,l={$\SI{6}{\ohm}$}] (5.2,3)
       to[battery1,l={$\SI{6}{\volt}$},invert] (5.2,0) -- (b);
 \draw (n) -- (2.6,4.4) -- (7.6,4.4) to[R,l={$\SI{6}{\ohm}$}] (7.6,1.6)
       to[battery1,l={$\SI{6}{\volt}$}] (7.6,0) -- (5.2,0);
 \node[above,Vcol] at (2.6,3.15) {$V$};
\end{circuitikz}}

\resposta{$V = \SI{4.8}{\volt}$}
\resolucao{\textbf{Superposição --- três passagens.} A cada vez, mantém-se uma
fonte e curto-circuitam-se as outras duas. Em todos os casos, o nó vê a fonte
ativa através de um resistor e, para o terra, o paralelo dos demais.

\smallskip
\textbf{(i) Só a fonte de \SI{12}{\volt}} (as outras em curto): o nó vê
$\SI{6}{\ohm}$ até a fonte e, para o terra, $3\parallel6\parallel6$. Como
$\dfrac13+\dfrac16+\dfrac16 = \dfrac23$, esse paralelo vale $\SI{1.5}{\ohm}$:
\[
V_1 = 12\cdot\frac{1{,}5}{6+1{,}5}=12\cdot\frac{1{,}5}{7{,}5}=\SI{2.4}{\volt}.
\]
\textbf{(ii) Só a primeira fonte de \SI{6}{\volt}:} por simetria de estrutura,
\[
V_2 = 6\cdot\frac{1{,}5}{7{,}5}=\SI{1.2}{\volt}.
\]
\textbf{(iii) Só a segunda fonte de \SI{6}{\volt}:} idem,
$V_3 = \SI{1.2}{\volt}$.

\smallskip
\textbf{Superpondo:} $V = 2{,}4+1{,}2+1{,}2 = \SI{4.8}{\volt}$.

\smallskip
\textbf{Verificação por análise nodal} (uma única equação):
\[
\frac{12-V}{6}+\frac{6-V}{6}+\frac{6-V}{6}=\frac{V}{3}.
\]
Multiplicando por 6: $(12-V)+(6-V)+(6-V) = 2V \Rightarrow 24-3V = 2V
\Rightarrow V = \SI{4.8}{\volt}$. \checkmark

\smallskip
\textbf{Comparação dos métodos.} A superposição exigiu \emph{três} análises
completas; a análise nodal exigiu apenas \emph{uma} equação. Isso ilustra que a
superposição, embora conceitualmente poderosa, raramente é o método mais
econômico quando há muitas fontes --- seu valor está em revelar a
\emph{contribuição individual} de cada fonte, informação que a nodal não fornece
diretamente.}
\end{questao}

% BEGIN BANCO COMPLEMENTAR Teorema da Superposição
% =====================================================================
%  Banco complementar --- Teorema da Superposicao  (REVISADO)
%  Substitui a versao gerada por template (9 clones em N2, 10 em N3 --
%  todos com o mesmo enunciado sobre potencia -- e 8 questoes
%  conceituais indevidamente marcadas como N4).
%  Sem sobreposicao com o cap06, que ja traz: como desativar fontes
%  ideais (MC), o que pode ser superposto (MC), V e I no resistor de
%  3 ohm com duas fontes, potencial com fonte de tensao + corrente,
%  o erro de somar potencias (analise do erro do estudante) e o
%  circuito com tres fontes independentes.
%  Distribuicao mantida: 6 N1 + 9 N2 + 10 N3 + 8 N4 = 33
% =====================================================================

\subsubsection*{Banco complementar --- Teorema da Superposição}

\begin{atencao}[Organização do banco]
A superposição decorre de uma única propriedade: em rede linear,
$\mathbf{v}=\mathbf{G}^{-1}\mathbf{b}$ é \emph{linear} no vetor de excitações.
Daí três consequências que organizam todo o tópico: tensões e correntes se
somam; \textbf{potências não}; e fontes \emph{dependentes} nunca são desativadas,
porque pertencem à rede, não à excitação. O N3 trata de fontes reais,
coeficientes de transferência e falha em circuitos não lineares; o N4 exige
demonstração, algoritmo e projeto de ensaio.
\end{atencao}

\blocofixacao

\begin{questao}[1]{}
Um nó $A$ é alimentado por uma fonte de $\SI{12}{\volt}$ através de
$\SI{6}{\ohm}$ e ligado ao terra por $\SI{3}{\ohm}$; há ainda uma fonte de
corrente ligada a $A$. Determine a contribuição \emph{apenas} da fonte de
tensão para $V_A$.
\resposta{$V_A'=\SI{4}{\volt}$}
\resolucao{Desativar a fonte de corrente significa substituí-la por um
\textbf{circuito aberto} --- o ramo simplesmente desaparece. Resta um divisor:
\[
V_A'=12\cdot\frac{3}{6+3}=\SI{4}{\volt}.
\]
Esta é apenas uma \emph{parcela}: o valor real de $V_A$ só sai depois de somada
a contribuição da outra fonte.}
\end{questao}

\begin{questao}[1]{}
No mesmo circuito, a fonte de corrente vale $\SI{2}{\ampere}$ e injeta em $A$.
Determine a contribuição \emph{apenas} dela.
\resposta{$V_A''=\SI{4}{\volt}$}
\resolucao{Desativar a fonte de tensão ideal significa substituí-la por um
\textbf{curto}. Os dois resistores ficam então em paralelo:
\[
R=\frac{6\cdot3}{9}=\SI{2}{\ohm}
\;\Longrightarrow\;
V_A''=2\cdot2=\SI{4}{\volt}.
\]
\textbf{Memorize a assimetria:} fonte de \emph{tensão} ideal $\to$ curto (pois
impõe $\SI{0}{\volt}$); fonte de \emph{corrente} ideal $\to$ aberto (pois impõe
$\SI{0}{\ampere}$).}
\end{questao}

\begin{questao}[1]{}
Some as duas contribuições anteriores e determine $V_A$.
\resposta{$V_A=\SI{8}{\volt}$}
\resolucao{
\[
V_A=V_A'+V_A''=4+4=\SI{8}{\volt}.
\]
\textbf{Verificação por análise nodal:}
\[
\frac{V_A-12}{6}+\frac{V_A}{3}=2
\;\Longrightarrow\;
V_A-12+2V_A=12
\;\Longrightarrow\;
V_A=\SI{8}{\volt}.\ \checkmark
\]
A concordância não é coincidência --- as duas técnicas são consequências da
mesma linearidade.}
\end{questao}

\begin{questao}[1]{}
Um circuito contém quatro fontes independentes. Quantas análises parciais a
superposição exige, e o que se faz em cada uma?
\resposta{Quatro análises; em cada uma, apenas uma fonte permanece ativa}
\resolucao{Uma análise por fonte independente: mantém-se \emph{uma} ativa e
desativam-se as outras três (tensão $\to$ curto, corrente $\to$ aberto). Ao
final, somam-se as quatro contribuições \emph{com sinal}.\\
\textbf{Observação de custo:} o esforço cresce \emph{linearmente} com o número
de fontes, mas cada análise parcial costuma ser bem mais simples que o circuito
completo. A vantagem se inverte quando há muitas fontes --- com quatro ou mais,
a análise nodal em geral é mais rápida.}
\end{questao}

\begin{questao}[1]{}
Duas fontes produzem, sobre o mesmo resistor, contribuições de
$\SI{+7}{\volt}$ e $\SI{-3}{\volt}$. Determine a tensão resultante e comente.
\resposta{$U=\SI{4}{\volt}$}
\resolucao{
\[
U=(+7)+(-3)=\SI{4}{\volt}.
\]
O sinal negativo indica que a segunda fonte produz, sozinha, uma tensão de
polaridade \emph{oposta} à adotada como referência. \textbf{Somar módulos seria
erro grave} --- daria $\SI{10}{\volt}$.\\
Fixe uma polaridade de referência \emph{antes} de começar e mantenha-a em todas
as análises parciais: é a única forma de os sinais saírem corretos
automaticamente.}
\end{questao}

\begin{questao}[1]{}
O teorema da superposição pode ser aplicado a um circuito que contenha um
resistor cuja resistência varia com a corrente?
\begin{alternativas}[2]
\item sim, sempre
\item sim, desde que haja apenas uma fonte
\item não, pois o circuito deixa de ser linear
\item não, pois a potência não se superpõe
\end{alternativas}
\resposta{(c)}
\resolucao{A superposição exige \textbf{linearidade}: a resposta deve ser
proporcional à excitação. Um resistor cuja resistência depende da corrente
quebra essa proporcionalidade, e a soma das respostas parciais deixa de valer.\\
A alternativa (d) enuncia um fato verdadeiro, mas irrelevante aqui --- a
potência não se superpõe nem mesmo em circuitos perfeitamente lineares, por uma
razão diferente (ela é \emph{quadrática}, não não linear).\\
E (b) é curiosa: com uma única fonte não há o que superpor, então a questão nem
se coloca.}
\end{questao}

\blocoaplicacao

\begin{questao}[2]{}
Um nó $A$ recebe $\SI{24}{\volt}$ através de $\SI{12}{\ohm}$,
$\SI{12}{\volt}$ através de $\SI{6}{\ohm}$, uma fonte de corrente de
$\SI{1}{\ampere}$, e liga-se ao terra por $\SI{4}{\ohm}$. Determine $V_A$ por
superposição.
\resposta{$V_A=\SI{10}{\volt}$}
\resolucao{A condutância total do nó é a mesma nas três análises:
\[
G=\frac1{12}+\frac16+\frac14=\SI{0.5}{\siemens}.
\]
\textbf{Só a fonte de $\SI{24}{\volt}$} (as outras: curto e aberto):
$V'=\dfrac{24/12}{0{,}5}=\SI{4}{\volt}$.\\
\textbf{Só a de $\SI{12}{\volt}$:} $V''=\dfrac{12/6}{0{,}5}=\SI{4}{\volt}$.\\
\textbf{Só a de corrente:} $V'''=\dfrac{1}{0{,}5}=\SI{2}{\volt}$.\\
\[
V_A=4+4+2=\SI{10}{\volt}.
\]
\textbf{Atalho revelado:} como todas as fontes de tensão são aterradas, $G$ não
muda entre as análises --- e a soma das três parcelas reproduz exatamente a
equação nodal única. Superposição e análise nodal coincidem passo a passo neste
caso.}
\end{questao}

\begin{questao}[2]{}
No circuito anterior, determine a contribuição \emph{percentual} de cada fonte
para $V_A$.
\resposta{$40\%$, $40\%$ e $20\%$}
\resolucao{
\[
\frac{4}{10}=40\%;\qquad \frac{4}{10}=40\%;\qquad \frac{2}{10}=20\%.
\]
\textbf{Utilidade prática:} a decomposição identifica qual fonte domina o
resultado. Se a fonte de corrente for a menos estável, seu erro afeta $V_A$
com peso de apenas $20\%$; se o objetivo for estabilizar $V_A$, o esforço deve
ir para as duas fontes de tensão.\\
Essa leitura --- impossível pela análise nodal direta, que entrega só o total ---
é a principal razão pela qual a superposição continua sendo usada mesmo quando
não é o método mais rápido.}
\end{questao}

\begin{questao}[2]{}
Duas fontes de corrente injetam $\SI{3}{\ampere}$ e retiram
$\SI{1}{\ampere}$ de um mesmo nó, ligado ao terra por
$\SI{10}{\ohm}\parallel\SI{15}{\ohm}$. Determine a tensão por superposição e
verifique.
\resposta{$V=\SI{12}{\volt}$}
\resolucao{A resistência do nó é $10\parallel15=\SI{6}{\ohm}$ em ambas as
análises (desativar uma fonte de corrente apenas abre seu ramo).
\[
V'=3(6)=\SI{18}{\volt};
\qquad
V''=-1(6)=\SI{-6}{\volt};
\qquad
V=18-6=\SI{12}{\volt}.
\]
\textbf{Verificação direta:} corrente líquida $3-1=\SI{2}{\ampere}$, logo
$V=2(6)=\SI{12}{\volt}$ \checkmark.\\
Com fontes de corrente no mesmo nó, a superposição é trivial --- elas se somam
algebricamente antes de qualquer conta. O método só ganha relevância quando as
fontes atuam em \emph{pontos diferentes} da rede.}
\end{questao}

\begin{questao}[2]{}
Duas fontes de tensão de $\SI{18}{\volt}$ e $\SI{6}{\volt}$, com polaridades
\textbf{opostas}, estão em série com $R=\SI{4}{\ohm}$. Determine a corrente por
superposição.
\resposta{$I=\SI{3}{\ampere}$}
\resolucao{Com a malha série, cada fonte age sozinha sobre o mesmo $R$ (a outra,
curto-circuitada, não altera a resistência):
\[
I'=\frac{18}{4}=\SI{4.5}{\ampere};
\qquad
I''=-\frac{6}{4}=\SI{-1.5}{\ampere};
\qquad
I=4{,}5-1{,}5=\SI{3}{\ampere}.
\]
\textbf{Verificação:} tensão líquida $18-6=\SI{12}{\volt}$, logo
$I=12/4=\SI{3}{\ampere}$ \checkmark.\\
O sinal negativo de $I''$ é essencial: a fonte de $\SI{6}{\volt}$ se opõe à de
$\SI{18}{\volt}$. Somar módulos daria $\SI{6}{\ampere}$ --- o dobro do
correto.}
\end{questao}

\begin{questao}[2]{}
Determine, por superposição, a \textbf{corrente} no resistor de
$\SI{4}{\ohm}$ do circuito de três fontes ($V_A=\SI{10}{\volt}$).
\resposta{$I=\SI{2.5}{\ampere}$}
\resolucao{As contribuições de tensão no nó já foram obtidas: $4$, $4$ e
$\SI{2}{\volt}$. Como o resistor liga $A$ ao terra,
\[
I'=\frac{4}{4}=1;\quad I''=\frac{4}{4}=1;\quad I'''=\frac{2}{4}=0{,}5
\;\Longrightarrow\;
I=\SI{2.5}{\ampere}.
\]
\textbf{Verificação:} $I=V_A/4=10/4=\SI{2.5}{\ampere}$ \checkmark.\\
\textbf{Observação:} correntes se superpõem tão livremente quanto tensões ---
ambas são \emph{lineares} nas excitações. É a mesma propriedade, aplicada a
outra grandeza.}
\end{questao}

\begin{questao}[2]{}
Em um circuito, a contribuição de $V_1$ para certa corrente é
$\SI{200}{\milli\ampere}$. Se $V_1$ for \textbf{triplicada}, mantidas as demais
fontes, o que ocorre com essa contribuição e com a corrente total?
\resposta{A contribuição triplica ($\SI{600}{\milli\ampere}$); a corrente total
aumenta em $\SI{400}{\milli\ampere}$}
\resolucao{Cada contribuição é \emph{proporcional} à sua fonte, portanto
$200\to\SI{600}{\milli\ampere}$.\\
As demais contribuições \textbf{não mudam}, pois nada foi alterado nas outras
fontes nem na rede. Logo a corrente total aumenta exatamente pelo acréscimo:
$600-200=\SI{400}{\milli\ampere}$.\\
\textbf{Poder do método:} para responder "o que acontece se eu mexer nesta
fonte", não é preciso resolver o circuito de novo --- basta escalar a parcela
correspondente. Essa é a base da análise de sensibilidade em redes lineares.}
\end{questao}

\begin{questao}[2]{}
Um circuito com duas fontes tem contribuições de $\SI{6}{\volt}$ e
$\SI{2}{\volt}$ sobre um resistor de $\SI{4}{\ohm}$. Calcule a potência de
duas formas: somando potências parciais e usando a tensão total. Compare.
\resposta{$\SI{10}{\watt}$ (errado) contra $\SI{16}{\watt}$ (correto)}
\resolucao{\textbf{Somando potências parciais (incorreto):}
\[
P'=\frac{6^{2}}{4}=9;\qquad P''=\frac{2^{2}}{4}=1;\qquad
P'+P''=\SI{10}{\watt}.
\]
\textbf{Pela tensão total (correto):}
\[
U=6+2=\SI{8}{\volt}
\;\Longrightarrow\;
P=\frac{8^{2}}{4}=\SI{16}{\watt}.
\]
A diferença é de $\SI{6}{\watt}$ --- erro de $37{,}5\%$ no resultado correto.\\
\textbf{Procedimento correto:} superponha \emph{tensões} (ou correntes), some, e
só então calcule a potência. Nunca superponha potências. A origem algébrica
dessa falha é explorada no bloco de desafio.}
\end{questao}

\begin{questao}[2]{}
Uma fonte \textbf{real} de $\SI{12}{\volt}$ com $r=\SI{2}{\ohm}$ e outra de
$\SI{6}{\volt}$ com $r=\SI{3}{\ohm}$ alimentam em paralelo uma carga de
$\SI{6}{\ohm}$. Determine a tensão na carga por superposição.
\resposta{$U=\SI{8}{\volt}$}
\resolucao{\textbf{Ponto crítico:} desativar uma fonte real significa zerar sua
\textbf{f.e.m.}, mantendo sua resistência interna no circuito.\\
Condutância total (a mesma nas duas análises):
$G=\frac12+\frac13+\frac16=\SI{1}{\siemens}$.
\[
U'=\frac{12/2}{1}=\SI{6}{\volt};
\qquad
U''=\frac{6/3}{1}=\SI{2}{\volt};
\qquad
U=\SI{8}{\volt}.
\]
\textbf{Erro comum:} remover a fonte de $\SI{6}{\volt}$ \emph{junto} com seus
$\SI{3}{\ohm}$. Isso daria $G=\frac12+\frac16=\SI{0.667}{\siemens}$ e
$U'=\SI{9}{\volt}$ --- resultado inconsistente, pois a soma não fecharia com a
análise nodal.}
\end{questao}

\begin{questao}[2]{}
Verifique o resultado da questão anterior por análise nodal e explique por que
os dois métodos devem concordar.
\resposta{$U=\SI{8}{\volt}$ pelos dois métodos}
\resolucao{\textbf{Nodal:}
\[
\frac{U-12}{2}+\frac{U-6}{3}+\frac{U}{6}=0.
\]
Multiplicando por $6$: $3U-36+2U-12+U=0\Rightarrow6U=48\Rightarrow
U=\SI{8}{\volt}$ \checkmark.\\
\textbf{Por que concordam:} a equação nodal tem a forma
$G\,U=\sum b_k$, com $b_k$ as contribuições de corrente de cada fonte. Como
$U=\left(\sum b_k\right)/G=\sum\left(b_k/G\right)$, cada parcela $b_k/G$ é
exatamente a contribuição individual da fonte $k$.\\
A superposição não é um método \emph{alternativo} à análise nodal --- ela é a
mesma equação, lida termo a termo. A demonstração formal está no bloco de
desafio.}
\end{questao}

\blocoaprofundamento

\begin{questao}[3]{}
Um circuito contém uma fonte independente de $\SI{12}{\volt}$ (através de
$\SI{2}{\kilo\ohm}$ até o nó $A$), $\SI{2}{\kilo\ohm}$ de $A$ ao terra, uma
fonte de corrente independente de $\SI{1}{\milli\ampere}$ injetando em $A$, e
uma fonte \textbf{dependente} que injeta $gV_A$ em $A$, com
$g=\SI{0.5}{\milli\siemens}$.
\begin{itensq}
\item Determine $V_A$ pela análise nodal.
\item Determine-o por superposição, aplicando o procedimento correto.
\item Mostre o resultado que se obteria desativando a fonte dependente.
\end{itensq}
\resposta{(a) e (b) $V_A=\SI{14}{\volt}$; (c) $\SI{7}{\volt}$ --- errado}
\resolucao{(a) LKC em $A$:
\[
\frac{V_A-12}{2000}+\frac{V_A}{2000}-\pot{0.5}{-3}V_A=\pot{1}{-3}.
\]
A condutância líquida é
$\pot{1}{-3}-\pot{0.5}{-3}=\pot{0.5}{-3}\,\si{\siemens}$, e o lado direito vale
$\pot{6}{-3}+\pot{1}{-3}=\pot{7}{-3}$. Logo
$V_A=\SI{14}{\volt}$.\\
(b) \textbf{A fonte dependente permanece ativa nas duas análises parciais} ---
ela não é excitação, é parte da rede. A condutância líquida é sempre
$\pot{0.5}{-3}\,\si{\siemens}$:
\[
V'=\frac{\pot{6}{-3}}{\pot{0.5}{-3}}=\SI{12}{\volt};
\qquad
V''=\frac{\pot{1}{-3}}{\pot{0.5}{-3}}=\SI{2}{\volt};
\qquad
V_A=\SI{14}{\volt}.\ \checkmark
\]
(c) Desativando indevidamente a dependente, a condutância voltaria a
$\pot{1}{-3}$:
\[
V'=6,\quad V''=1
\;\Longrightarrow\;
V_A=\SI{7}{\volt},
\]
exatamente \textbf{metade} do correto.\\
\textbf{Critério:} desative apenas fontes cujo valor é \emph{imposto de fora}.
Se o valor depende de uma variável do próprio circuito, a fonte fica.}
\end{questao}

\begin{questao}[3]{}
Escreva $V_A$ do circuito de três fontes na forma
$V_A=a\,V_1+b\,V_2+c\,I_3$ e determine os coeficientes.
\begin{itensq}
\item Obtenha $a$, $b$ e $c$ com suas unidades.
\item Verifique com os valores originais.
\item Explique o significado físico de cada coeficiente.
\end{itensq}
\resposta{$a=1/6$, $b=1/3$ (adimensionais), $c=\SI{2}{\ohm}$}
\resolucao{(a) Com $G=\SI{0.5}{\siemens}$, cada fonte de tensão $V_k$ contribui
com $(V_k/R_k)/G$ e a fonte de corrente com $I_3/G$:
\[
a=\frac{1/12}{0{,}5}=\frac16;
\qquad
b=\frac{1/6}{0{,}5}=\frac13;
\qquad
c=\frac{1}{0{,}5}=\SI{2}{\ohm}.
\]
(b) $\frac{24}{6}+\frac{12}{3}+2(1)=4+4+2=\SI{10}{\volt}$ \checkmark.\\
(c) $a$ e $b$ são \textbf{ganhos de tensão} (adimensionais): que fração de cada
fonte chega ao nó. $c$ é uma \textbf{resistência de transferência}
($\si{\volt\per\ampere}$): quantos volts aparecem por ampère injetado.\\
\textbf{Utilidade:} obtidos uma única vez, os coeficientes respondem a qualquer
combinação de valores de fonte sem refazer o circuito. É a forma compacta de
armazenar a solução de uma rede linear --- e é exatamente a linha correspondente
da matriz $\mathbf{G}^{-1}$.}
\end{questao}

\begin{questao}[3]{}
Compare o esforço de resolver por superposição e por análise nodal um circuito
com $2$ nós desconhecidos e $4$ fontes independentes.
\begin{itensq}
\item Quantifique cada abordagem.
\item Indique quando cada uma é preferível.
\end{itensq}
\resposta{Superposição: 4 sistemas $2\times2$; nodal: 1 sistema $2\times2$}
\resolucao{(a) \textbf{Nodal:} monta-se \emph{um} sistema $2\times2$ com todas
as fontes no vetor independente e resolve-se uma vez.\\
\textbf{Superposição:} quatro sistemas $2\times2$, um por fonte, e depois a
soma. Cerca de \textbf{quatro vezes} mais trabalho.\\
(b) \textbf{Nodal é preferível} quando se quer apenas o resultado numérico ---
quase sempre.\\
\textbf{Superposição é preferível} quando: \emph{(i)} se quer saber a
\emph{contribuição} de cada fonte (análise de sensibilidade, rastreamento de
ruído, identificação da fonte dominante); \emph{(ii)} as fontes são de naturezas
diferentes e se deseja tratá-las separadamente; \emph{(iii)} algumas análises
parciais são triviais, colapsando o circuito por curtos ou aberturas ---
situação frequente e que pode inverter completamente a comparação de esforço.\\
\textbf{Síntese:} a superposição raramente é o caminho mais curto para o
número, mas frequentemente é o único caminho para o \emph{entendimento}.}
\end{questao}

\begin{questao}[3]{}
Um circuito tem duas fontes cujas contribuições para a corrente em um resistor
de $\SI{5}{\ohm}$ são $\SI{2}{\ampere}$ e $\SI{-3}{\ampere}$.
\begin{itensq}
\item Determine a corrente e a potência reais.
\item Calcule o erro que se cometeria somando potências parciais.
\item Determine o termo cruzado e seu sinal.
\end{itensq}
\resposta{(a) $\SI{-1}{\ampere}$, $\SI{5}{\watt}$; (b) somaria
$\SI{65}{\watt}$; (c) termo cruzado $=\SI{-60}{\watt}$}
\resolucao{(a) $I=2-3=\SI{-1}{\ampere}$, logo
$P=5(1)^{2}=\SI{5}{\watt}$.\\
(b) $P'=5(4)=20$ e $P''=5(9)=\SI{45}{\watt}$; a soma daria
$\SI{65}{\watt}$ --- \textbf{treze vezes} o valor real.\\
(c)
\[
P=R(I'+I'')^{2}=R I'^{2}+R I''^{2}+2RI'I''
\;\Longrightarrow\;
P_{cruz}=2(5)(2)(-3)=\SI{-60}{\watt}.
\]
\textbf{Sinal negativo:} as contribuições têm sentidos opostos e se
\emph{cancelam} parcialmente. O termo cruzado subtrai.\\
\textbf{Caso extremo:} se as contribuições fossem $+2$ e $\SI{-2}{\ampere}$, a
corrente real seria nula e a potência também --- enquanto a soma de potências
parciais daria $\SI{40}{\watt}$ de dissipação inexistente. É o exemplo mais
claro de que potência não se superpõe.}
\end{questao}

\begin{questao}[3]{}
Duas fontes de tensão de $\SI{5}{\volt}$ e $\SI{-5}{\volt}$, cada uma através
de $\SI{1}{\kilo\ohm}$, alimentam um nó ligado ao terra por um \textbf{diodo}
(queda direta de $\SI{0.7}{\volt}$, bloqueio perfeito no reverso).
\begin{itensq}
\item Determine a tensão do nó com cada fonte agindo sozinha.
\item Some as contribuições.
\item Determine a tensão real e explique a discrepância.
\end{itensq}
\resposta{(a) $\SI{0.7}{\volt}$ e $\SI{-2.5}{\volt}$; (b) previsão
$\SI{-1.8}{\volt}$; (c) real $\SI{0}{\volt}$}
\resolucao{(a) \textbf{Só a fonte $+\SI{5}{\volt}$} (a outra substituída por
curto): o nó é alimentado por $\SI{5}{\volt}$ através de
$\SI{1}{\kilo\ohm}$ e ligado ao terra por outro $\SI{1}{\kilo\ohm}$ (o ramo
curto-circuitado). Sem o diodo, o divisor daria $\SI{2.5}{\volt}$ --- acima da
queda direta. O diodo conduz e grampeia o nó em
$V'=\SI{0.7}{\volt}$.\\
\textbf{Só a fonte $-\SI{5}{\volt}$:} pelo mesmo divisor, o nó tenderia a
$\SI{-2.5}{\volt}$. O diodo fica reversamente polarizado e não conduz, de modo
que $V''=\SI{-2.5}{\volt}$.\\
(b) Soma das parcelas: $0{,}7-2{,}5=\SI{-1.8}{\volt}$.\\
(c) \textbf{Situação real:} com as duas fontes ativas e simétricas, as correntes
pelos dois resistores iguais se cancelam e o nó assume $\SI{0}{\volt}$. Com
$\SI{0}{\volt}$ o diodo não conduz --- hipótese consistente. Logo
$V=\SI{0}{\volt}$.\\
A previsão da superposição ($\SI{-1.8}{\volt}$) está \textbf{errada}, e nem
sequer perto.\\
\textbf{Razão:} o diodo não é linear --- seu \emph{estado} (conduzindo ou
bloqueado) depende do circuito inteiro, e mudou entre as duas análises
parciais. Superpor respostas obtidas em \emph{configurações diferentes} não tem
significado. Este é o contraexemplo definitivo: a superposição não é uma
conveniência que se pode forçar, é uma consequência da linearidade.}
\end{questao}

\begin{questao}[3]{}
Explique por que, ao desativar uma fonte \textbf{real} de tensão, mantém-se sua
resistência interna, e quantifique o erro de removê-la no circuito de
$E_1=\SI{12}{\volt}/r_1=\SI{2}{\ohm}$, $E_2=\SI{6}{\volt}/r_2=\SI{3}{\ohm}$ e
$R_L=\SI{6}{\ohm}$.
\resposta{Contribuição correta de $E_1$: $\SI{6}{\volt}$; removendo $r_2$
obter-se-ia $\SI{9}{\volt}$ --- erro de $50\%$}
\resolucao{\textbf{Por quê:} a fonte real é um \emph{modelo} de dois elementos
--- uma f.e.m. ideal em série com $r$. Desativar a \emph{fonte} significa zerar
a f.e.m.; o resistor $r$ é um componente físico da rede e permanece, exatamente
como qualquer outro resistor.\\
\textbf{Correto:} com $E_2=0$ e $r_2$ presente,
$G=\frac12+\frac13+\frac16=\SI{1}{\siemens}$ e
\[
U'=\frac{12/2}{1}=\SI{6}{\volt}.
\]
\textbf{Incorreto:} removendo o ramo inteiro,
$G=\frac12+\frac16=\SI{0.667}{\siemens}$ e
$U'=\dfrac{6}{0{,}667}=\SI{9}{\volt}$ --- $50\%$ acima.\\
\textbf{Como o erro se denuncia:} a soma das parcelas incorretas não fecha com
a análise nodal ($\SI{8}{\volt}$). Conferir a soma total contra um método
independente é a salvaguarda padrão da superposição.}
\end{questao}

\begin{questao}[3]{}
Em um circuito com duas fontes, mediu-se: com ambas ativas,
$I=\SI{1.4}{\ampere}$; com apenas a primeira, $I'=\SI{2.0}{\ampere}$.
\begin{itensq}
\item Determine a contribuição da segunda fonte.
\item Determine a corrente se a segunda fonte for dobrada.
\item Determine o valor da segunda fonte que anularia $I$.
\end{itensq}
\resposta{(a) $\SI{-0.6}{\ampere}$; (b) $\SI{0.8}{\ampere}$; (c) o triplo do
valor original}
\resolucao{(a) $I''=I-I'=1{,}4-2{,}0=\SI{-0.6}{\ampere}$ --- ela se opõe à
primeira.\\
(b) Dobrando a segunda fonte, sua contribuição dobra e a da primeira não muda:
\[
I=2{,}0+2(-0{,}6)=\SI{0.8}{\ampere}.
\]
(c) Anular exige $I''=\SI{-2.0}{\ampere}$, ou seja, um fator
$2{,}0/0{,}6=3{,}33$ sobre o valor original.\\
\textbf{Poder do método:} duas medições permitem prever o comportamento para
\emph{qualquer} valor da segunda fonte, sem conhecer a topologia do circuito.
Basta a garantia de linearidade --- e é justamente por isso que a verificação
experimental da linearidade (bloco de desafio) tem valor prático.}
\end{questao}

\begin{questao}[3]{}
Uma ponte é alimentada por duas fontes independentes. Sabe-se que a
contribuição da primeira para a corrente do galvanômetro é
$\SI{+40}{\milli\ampere}$ e que a corrente total medida é
$\SI{0}{\milli\ampere}$.
\begin{itensq}
\item Determine a contribuição da segunda fonte.
\item Explique por que o equilíbrio observado \emph{não} significa que a ponte
      esteja balanceada no sentido usual.
\end{itensq}
\resposta{(a) $\SI{-40}{\milli\ampere}$; (b) trata-se de cancelamento entre
fontes, não de equilíbrio da rede}
\resolucao{(a) $I''=0-40=\SI{-40}{\milli\ampere}$.\\
(b) O \textbf{equilíbrio clássico} da ponte é uma propriedade da \emph{rede}
($R_1/R_2=R_3/R_4$): a corrente é nula qualquer que seja a alimentação, e o
valor de $R_g$ é irrelevante.\\
Aqui, ao contrário, a rede está desequilibrada --- cada fonte, sozinha, produz
corrente considerável. O zero resulta de um \textbf{cancelamento numérico}
entre as duas contribuições, e depende dos \emph{valores} das fontes.\\
\textbf{Consequência prática:} basta variar qualquer uma das fontes em $1\%$
para que o zero desapareça. Um ajuste desse tipo não tem a robustez do
equilíbrio verdadeiro, e por isso não serve como método de medição --- ele
mediria a razão entre as fontes, não a razão entre os resistores.}
\end{questao}

\begin{questao}[3]{}
Um sinal de $\SI{50}{\milli\volt}$ e uma interferência de rede de
$\SI{200}{\milli\volt}$ atuam simultaneamente sobre um circuito linear, com
ganhos de $8$ e $0{,}05$ respectivamente.
\begin{itensq}
\item Determine as contribuições na saída.
\item Determine a razão sinal-ruído resultante.
\item Explique como a superposição justifica o tratamento separado.
\end{itensq}
\resposta{(a) $\SI{400}{\milli\volt}$ e $\SI{10}{\milli\volt}$; (b) $40$, ou
$\SI{32}{\decibel}$}
\resolucao{(a) $V_{sinal}=8(50)=\SI{400}{\milli\volt}$;
$V_{ruído}=0{,}05(200)=\SI{10}{\milli\volt}$.\\
(b) Razão $=400/10=40$, ou
$20\log_{10}40\approx\SI{32}{\decibel}$.\\
Na entrada, a razão era $50/200=0{,}25$ --- o \emph{ruído era quatro vezes
maior que o sinal}. O circuito melhorou a relação em um fator $160$.\\
(c) \textbf{Justificativa:} em rede linear, sinal e interferência percorrem o
circuito \emph{independentemente}, cada um com seu próprio ganho. Isso permite
projetar o circuito para maximizar um ganho e minimizar o outro, tratando-os
como problemas separados.\\
\textbf{E é exatamente isso que se perde na não linearidade:} em um circuito não
linear surgem \emph{produtos de intermodulação} --- termos cruzados que geram
componentes em frequências novas, impossíveis de separar depois. É o análogo,
em sinais, do termo cruzado de potência.}
\end{questao}

\begin{questao}[3]{}
Um circuito linear tem três fontes. Sabe-se que $V_{saída}=\SI{15}{\volt}$ com
todas ativas, $\SI{9}{\volt}$ sem a terceira e $\SI{4}{\volt}$ apenas com a
primeira. Determine a contribuição de cada fonte.
\resposta{$\SI{4}{\volt}$, $\SI{5}{\volt}$ e $\SI{6}{\volt}$}
\resolucao{Sejam $a$, $b$ e $c$ as contribuições.
\begin{itemize}
\item Todas ativas: $a+b+c=15$.
\item Sem a terceira: $a+b=9$ $\Rightarrow$ $c=15-9=\SI{6}{\volt}$.
\item Apenas a primeira: $a=\SI{4}{\volt}$ $\Rightarrow$
      $b=9-4=\SI{5}{\volt}$.
\end{itemize}
\textbf{Verificação:} $4+5+6=\SI{15}{\volt}$ \checkmark.\\
\textbf{Observação de método:} bastaram \emph{três} medições para decompor
completamente um sistema de três fontes, e nenhuma delas exigiu isolar a segunda
ou a terceira sozinhas. Em geral, $m$ medições \textbf{independentes} bastam
para $m$ fontes --- resultado que decorre de o sistema ser linear e as
contribuições, aditivas.}
\end{questao}

\blocodesafio

\begin{questao}[4]{}
\textbf{(Demonstração formal.)} Seja uma rede resistiva descrita por
$\mathbf{G}\mathbf{v}=\mathbf{b}$, com
$\mathbf{b}=\mathbf{b}_1+\mathbf{b}_2$ as contribuições de duas fontes.
\begin{itensq}
\item Prove a superposição a partir da linearidade do sistema.
\item Generalize para $m$ fontes.
\item Identifique onde a demonstração falha se $\mathbf{G}$ depender de
      $\mathbf{v}$.
\end{itensq}
\resposta{$\mathbf{v}=\mathbf{G}^{-1}\mathbf{b}_1+\mathbf{G}^{-1}\mathbf{b}_2$}
\resolucao{(a) Como $\mathbf{G}$ é não singular em rede passiva conectada,
\[
\mathbf{v}=\mathbf{G}^{-1}\mathbf{b}
=\mathbf{G}^{-1}(\mathbf{b}_1+\mathbf{b}_2)
=\underbrace{\mathbf{G}^{-1}\mathbf{b}_1}_{\mathbf{v}_1}
+\underbrace{\mathbf{G}^{-1}\mathbf{b}_2}_{\mathbf{v}_2}.
\]
Ora, $\mathbf{v}_1$ é exatamente a solução com \emph{apenas} a primeira fonte
ativa (as demais zeradas anulam suas parcelas em $\mathbf{b}$). Logo
$\mathbf{v}=\mathbf{v}_1+\mathbf{v}_2$: o teorema.\\
(b) Por indução imediata,
$\mathbf{v}=\sum_{k=1}^{m}\mathbf{G}^{-1}\mathbf{b}_k$. A distributividade do
produto matricial sobre a soma é tudo o que se usa.\\
(c) \textbf{Se $\mathbf{G}=\mathbf{G}(\mathbf{v})$}, a matriz muda de uma
análise parcial para outra --- e $\mathbf{G}_1^{-1}\mathbf{b}_1
+\mathbf{G}_2^{-1}\mathbf{b}_2$ não é solução de nada. Foi precisamente isso que
ocorreu no exemplo do diodo: a matriz do circuito com o diodo conduzindo é
diferente da matriz com ele bloqueado.\\
\textbf{Três leituras do resultado.} \emph{(i)} A superposição é uma
propriedade da \emph{álgebra linear}, não da eletricidade --- vale para qualquer
sistema linear. \emph{(ii)} Fontes dependentes entram em $\mathbf{G}$, não em
$\mathbf{b}$ --- por isso não são desativadas. \emph{(iii)} A potência é
quadrática em $\mathbf{v}$, e nenhuma função quadrática de uma soma é a soma das
funções.}
\end{questao}

\begin{questao}[4]{}
\textbf{(O termo cruzado.)} Para um resistor $R$ sujeito a contribuições de
corrente $I_1$ e $I_2$:
\begin{itensq}
\item Expanda $P$ e identifique o termo cruzado.
\item Determine em que condições ele é nulo, positivo ou negativo.
\item Mostre que a soma de potências parciais pode ser maior \emph{ou} menor
      que a potência real, com exemplos numéricos.
\end{itensq}
\resposta{$P=RI_1^{2}+RI_2^{2}+2RI_1I_2$; o cruzado tem o sinal de
$I_1I_2$}
\resolucao{(a)
\[
P=R(I_1+I_2)^{2}=\underbrace{RI_1^{2}}_{P_1}+\underbrace{RI_2^{2}}_{P_2}
+\underbrace{2RI_1I_2}_{P_{cruz}}.
\]
(b) $P_{cruz}=0$ apenas se $I_1=0$ ou $I_2=0$ --- ou seja, praticamente nunca
com duas fontes atuantes. Ele é \textbf{positivo} quando as contribuições têm o
mesmo sentido e \textbf{negativo} quando se opõem.\\
\emph{(Nota: em corrente alternada, $P_{cruz}$ pode ter média nula se as
contribuições forem de frequências diferentes --- é a única situação em que
somar potências é legítimo, e ela não ocorre em CC.)}\\
(c) \textbf{Soma parcial menor que a real} ($R=\SI{4}{\ohm}$, contribuições de
tensão $6$ e $\SI{2}{\volt}$):
$P_1+P_2=9+1=\SI{10}{\watt}$ contra $P=\SI{16}{\watt}$. O cruzado soma
$\SI{6}{\watt}$.\\
\textbf{Soma parcial maior que a real} ($R=\SI{5}{\ohm}$, correntes
$+2$ e $\SI{-3}{\ampere}$): $P_1+P_2=20+45=\SI{65}{\watt}$ contra
$P=\SI{5}{\watt}$. O cruzado subtrai $\SI{60}{\watt}$.\\
\textbf{Conclusão:} o erro não tem sequer sinal previsível --- não existe
"correção" ou fator de segurança que salve o procedimento. A única saída é
superpor tensões ou correntes e calcular a potência ao final.}
\end{questao}

\begin{questao}[4]{}
\textbf{(Algoritmo.)} Formule um procedimento de superposição para $m$ fontes e
$n$ grandezas de interesse.
\begin{itensq}
\item Descreva o algoritmo e a estrutura de dados resultante.
\item Determine o custo computacional e compare com $m$ análises completas.
\item Indique como a matriz resultante é usada depois.
\end{itensq}
\resposta{Matriz de coeficientes $n\times m$; custo de uma fatoração mais $m$
substituições}
\resolucao{(a) \textbf{Algoritmo.}
\begin{enumerate}
\item Monte $\mathbf{G}$ uma única vez (ela não depende das fontes).
\item Fatore $\mathbf{G}=\mathbf{L}\mathbf{U}$ --- também uma única vez.
\item Para cada fonte $k=1,\dots,m$: monte $\mathbf{b}_k$ com apenas essa fonte
      ativa e resolva por substituição, obtendo $\mathbf{v}_k$.
\item Extraia as $n$ grandezas de interesse de cada $\mathbf{v}_k$ e armazene
      na coluna $k$ de uma matriz $\mathbf{H}$ de dimensão $n\times m$.
\end{enumerate}
(b) \textbf{Custo:} a fatoração custa $\mathcal{O}(N^{3}/3)$ para uma rede de
$N$ nós, e cada substituição custa $\mathcal{O}(N^{2})$. Total:
$\mathcal{O}(N^{3}/3+mN^{2})$.\\
Repetir a análise completa $m$ vezes custaria
$\mathcal{O}(mN^{3}/3)$. Para $N=20$ e $m=5$, a razão é de aproximadamente
$4{,}7$ contra $13{,}3$ (em milhares de operações) --- \textbf{quase três vezes
mais rápido}, e a vantagem cresce com $m$.\\
\textbf{Observação:} a superposição sai \emph{de graça} como subproduto da
fatoração. O custo extra em relação a uma única análise é apenas
$(m-1)N^{2}$.\\
(c) \textbf{Uso posterior:} qualquer combinação de valores de fonte
$\mathbf{s}$ produz as saídas por um simples produto
$\mathbf{y}=\mathbf{H}\mathbf{s}$ --- sem tocar no circuito. A matriz
$\mathbf{H}$ é o \emph{modelo linear} completo da rede, e é ela que permite
varredura de parâmetros, análise de pior caso e otimização.}
\end{questao}

\begin{questao}[4]{}
\textbf{(Projeto de ensaio.)} Elabore uma verificação experimental da
superposição em um circuito com duas fontes de tensão.
\begin{itensq}
\item Descreva o procedimento e o cuidado essencial na desativação.
\item Estabeleça o critério de aceitação, dadas incertezas de $\pm1\%$ nas
      medições.
\item Interprete uma discrepância de $8\%$.
\end{itensq}
\resposta{Critério: $|V-(V'+V'')|$ dentro de cerca de $\pm3\%$; $8\%$ indica
não linearidade ou erro de procedimento}
\resolucao{(a) \textbf{Procedimento.}
\begin{enumerate}
\item Com ambas as fontes ativas, meça $V$ no ponto escolhido.
\item Desative a fonte $2$ e meça $V'$.
\item Restaure a fonte $2$, desative a $1$ e meça $V''$.
\item Compare $V$ com $V'+V''$.
\end{enumerate}
\textbf{Cuidado essencial:} desativar significa \textbf{substituir por um
curto}, não simplesmente desligar a fonte nem retirá-la do circuito. Uma fonte
desligada apresenta impedância indefinida (alta, em geral) --- o que equivale a
\emph{abrir} o ramo e altera a topologia. Se a fonte tiver resistência interna
apreciável, ela deve ser mantida no lugar da f.e.m.\\
(b) \textbf{Critério.} As três medições são independentes, e as incertezas se
propagam. Como $V'$ e $V''$ entram somadas, a incerteza da soma é
$\sqrt{u'^{2}+u''^{2}}$; combinada com a de $V$, chega-se a algo em torno de
$\pm1{,}7\%$ para uma estimativa estatística, ou até $\pm3\%$ no pior caso.
\textbf{Adote $\pm3\%$} como faixa de aceitação.\\
(c) \textbf{Discrepância de $8\%$} está muito além da incerteza. Causas
prováveis, em ordem de probabilidade:
\begin{itemize}
\item \textbf{Erro de procedimento:} fonte desligada em vez de curto-circuitada
      --- de longe a causa mais comum.
\item \textbf{Resistência interna ignorada} na desativação.
\item \textbf{Não linearidade real:} aquecimento de resistores, diodo de
      proteção conduzindo, ou fonte entrando em limitação.
\end{itemize}
\textbf{Como distinguir:} repita o ensaio com tensões \emph{reduzidas à
metade}. Se a discrepância percentual permanecer, o problema é de procedimento;
se diminuir, é não linearidade dependente de nível.}
\end{questao}

\begin{questao}[4]{}
\textbf{(Comparação de métodos.)} Um circuito tem duas fontes de tensão reais e
uma fonte de corrente.
\begin{itensq}
\item Descreva a solução por transformação de fontes.
\item Compare com a superposição em número de etapas e em informação obtida.
\item Indique o critério de escolha.
\end{itensq}
\resposta{Transformação: mais rápida para o número; superposição: entrega a
decomposição}
\resolucao{(a) \textbf{Transformação de fontes.} Converta cada fonte real de
tensão $(E_k,r_k)$ em fonte de corrente $(E_k/r_k, r_k)$. Todas as fontes de
corrente ficam então em paralelo com todas as resistências, e podem ser
\emph{somadas algebricamente} em uma única fonte equivalente; as condutâncias
também se somam. Uma divisão final entrega a tensão do nó.\\
(b) \textbf{Etapas.}
\begin{center}
\begin{tabular}{lcc}
\hline
& Transformação & Superposição\\
\hline
Análises do circuito & 1 & 3\\
Resultado numérico & sim & sim\\
Contribuição por fonte & não & sim\\
\hline
\end{tabular}
\end{center}
A transformação é claramente mais econômica --- ela é, na prática, a
\emph{montagem} da equação nodal por outro caminho.\\
Mas note: as parcelas $E_k/r_k$ que a transformação soma são \emph{exatamente}
as contribuições que a superposição calcula separadamente. Os métodos não são
rivais; um é a forma agregada do outro.\\
(c) \textbf{Critério.} Use transformação de fontes (ou análise nodal) para
obter o número. Use superposição quando a pergunta for \emph{"quanto cada fonte
contribui"} --- rastreamento de ruído, análise de sensibilidade, diagnóstico de
qual fonte está causando um desvio. Se a resposta desejada for um número único,
a superposição é trabalho desperdiçado.}
\end{questao}

\begin{questao}[4]{}
\textbf{(Limites de aplicação.)} Determine se a superposição se aplica e
justifique.
\begin{itensq}
\item Circuito com lâmpada incandescente, cuja resistência cresce com a
      temperatura.
\item Circuito linear em que se deseja a \emph{energia} dissipada em um
      intervalo.
\item Circuito com amplificador operacional operando na região linear.
\end{itensq}
\resposta{(a) não; (b) não; (c) sim}
\resolucao{(a) \textbf{Não se aplica.} A resistência depende da potência
dissipada, que depende da corrente --- realimentação que torna o circuito não
linear. Cada análise parcial teria uma temperatura, e portanto uma
\emph{resistência}, diferente.\\
\emph{Ressalva:} para pequenas variações em torno de um ponto de operação, o
modelo incremental é linear e a superposição vale \emph{para as variações}.\\
(b) \textbf{Não se aplica.} Energia é a integral da potência, e potência é
quadrática --- herda o termo cruzado. Superpor energias comete o mesmo erro que
superpor potências, apenas acumulado no tempo.\\
\emph{Procedimento correto:} superponha correntes, obtenha $i(t)$ total,
calcule $p(t)=Ri^{2}$ e só então integre.\\
(c) \textbf{Aplica-se.} O amp-op na região linear é descrito por relações
lineares (curto virtual e corrente de entrada nula), e o ganho é uma constante.
Superposição é, aliás, a ferramenta padrão para analisar amplificadores
somadores e diferenciais: calcula-se o ganho de cada entrada com as demais
aterradas e somam-se as contribuições.\\
\emph{Fronteira:} se a saída satura, a linearidade se perde e a superposição
falha --- o mesmo padrão do diodo.}
\end{questao}

\begin{questao}[4]{}
\textbf{(Escolha de sinais.)} Estabeleça um procedimento sistemático para os
sinais das contribuições e aplique-o a um caso com resultado negativo.
\begin{itensq}
\item Formule a regra.
\item Aplique a duas fontes cujas contribuições de corrente são
      $\SI{2}{\ampere}$ (a favor) e $\SI{5}{\ampere}$ (contra) a referência.
\item Explique por que o resultado negativo não indica erro.
\end{itensq}
\resposta{$I=\SI{-3}{\ampere}$: mesma intensidade, sentido oposto ao adotado}
\resolucao{(a) \textbf{Regra em três passos.}
\begin{enumerate}
\item \textbf{Antes} de qualquer cálculo, desenhe uma seta de referência para a
      corrente (ou uma polaridade $+/-$ para a tensão) e \textbf{não a mude
      mais}.
\item Em cada análise parcial, determine o sentido real da contribuição.
Atribua sinal $+$ se coincidir com a referência e $-$ se for oposto.
\item Some algebricamente.
\end{enumerate}
O erro clássico é redesenhar a seta a cada análise parcial "para ficar
positivo" --- o que torna a soma final sem significado.\\
(b)
\[
I=(+2)+(-5)=\SI{-3}{\ampere}.
\]
(c) \textbf{O sinal negativo é informação, não erro.} Ele diz que a corrente
real tem módulo $\SI{3}{\ampere}$ e flui no sentido \emph{oposto} ao da seta
desenhada. Se a seta tivesse sido desenhada ao contrário desde o início, o
resultado seria $\SI{+3}{\ampere}$ --- e descreveria exatamente a mesma
realidade física.\\
\textbf{Verificação útil:} grandezas que não podem ser negativas --- como a
potência dissipada em um resistor --- servem de teste. Aqui,
$P=R(3)^{2}$ é positiva independentemente do sinal de $I$, como deve ser.}
\end{questao}

\begin{questao}[4]{}
\textbf{(Superposição e ruído.)} Um amplificador linear recebe um sinal
$V_s$ e sofre a injeção de uma interferência $V_n$ em um ponto interno, com
ganhos $A_s=40$ e $A_n=2$.
\begin{itensq}
\item Escreva a saída e a razão sinal-ruído em função das entradas.
\item Determine o ganho $A_n$ máximo admissível para
      $\mathrm{SNR}\ge\SI{40}{\decibel}$, com
      $V_s=\SI{10}{\milli\volt}$ e $V_n=\SI{100}{\milli\volt}$.
\item Explique por que essa análise seria impossível em circuito não linear.
\end{itensq}
\resposta{(a) $V_o=A_sV_s+A_nV_n$; (b) $A_n\le0{,}04$}
\resolucao{(a) Pela superposição, cada entrada percorre o circuito
independentemente:
\[
V_o=A_sV_s+A_nV_n;
\qquad
\mathrm{SNR}=\frac{A_sV_s}{A_nV_n}.
\]
(b) $\SI{40}{\decibel}$ corresponde a uma razão de $100$:
\[
\frac{40(10)}{A_n(100)}\ge100
\;\Longrightarrow\;
\frac{400}{100A_n}\ge100
\;\Longrightarrow\;
A_n\le0{,}04.
\]
O ganho de $2$ do enunciado é \textbf{cinquenta vezes} maior que o admissível
--- o projeto está muito longe da especificação, e a conta diz exatamente
quanto: são necessários $\SI{34}{\decibel}$ de atenuação adicional no caminho
da interferência.\\
(c) \textbf{Em circuito não linear, "o ganho do ruído" não existe como grandeza
independente.} A resposta à interferência dependeria do nível do sinal
\emph{simultâneo}, e vice-versa. Surgiriam produtos de intermodulação --- termos
proporcionais a $V_sV_n$ --- que criam componentes espectrais novas, dentro da
banda do sinal e portanto \textbf{impossíveis de filtrar}.\\
É o mesmo termo cruzado que impede somar potências, agora manifestado no
domínio dos sinais. A linearidade não é um detalhe matemático: é o que garante
que sinal e ruído permaneçam separáveis.}
\end{questao}

% END BANCO COMPLEMENTAR

\gabarito[Capítulo 6 --- Teorema da Superposição]
           % Teorema da Superposicao

% Cap. 7 -------------------------------------------------------------
\setcounter{section}{6}
% =====================================================================
%  CAPITULO 7 - ANALISE NODAL
% =====================================================================
\section{Análise Nodal}

\begin{conceito}[Antes de começar]
A análise nodal usa a \textbf{lei dos nós} de Kirchhoff para achar os potenciais
dos nós. Roteiro:
\begin{enumerate}[leftmargin=1.7em,itemsep=1pt,topsep=2pt]
\item Escolha um nó como \textbf{referência} (terra, $V=0$).
\item Nomeie os potenciais dos demais nós ($V_1, V_2, \dots$).
\item Em cada nó, escreva que a soma das correntes que \emph{saem} é nula. A
      corrente por um resistor entre os nós $a$ e $b$ é
      $\dfrac{V_a - V_b}{R}$.
\item Resolva o sistema.
\end{enumerate}
Uma fonte de corrente que injeta $I$ no nó entra como termo independente; uma
fonte de tensão entre dois nós cria uma relação direta entre eles.
\end{conceito}

\legendaniveis

\blocofixacao

\begin{questao}[1]{}
Na análise nodal, o nó de referência (terra) é escolhido para:
\begin{alternativas}
\item ter o maior potencial do circuito.
\item servir como ponto de potencial zero, ao qual os demais são comparados.
\item ser sempre o nó com a fonte de maior tensão.
\item eliminar as fontes de corrente.
\item garantir que todas as correntes sejam positivas.
\end{alternativas}
\resposta{(b)}
\resolucao{O terra é uma referência arbitrária de potencial nulo; os potenciais
dos outros nós são medidos em relação a ele. A escolha não altera as
\emph{diferenças} de potencial nem as correntes --- apenas simplifica as
equações.}
\end{questao}

\begin{questao}[1]{}
No circuito da figura, uma fonte de corrente de $\SI{9}{\ampere}$ alimenta dois
resistores em paralelo ligados ao terra. Determine o potencial do nó $V$.

\circuitoq{%
\begin{circuitikz}[scale=1,transform shape,line width=0.8pt,american]
 \draw (0,0) to[isource,l={$\SI{9}{\ampere}$}] (0,2.6)
       -- (1.5,2.6) coordinate(n);
 \node[above,Vcol] at (1.5,2.75) {$V$};
 \draw (n) to[R,l={$\SI{6}{\ohm}$}] (1.5,0) -- (0,0);
 \draw (n) -- (3.5,2.6) to[R,l={$\SI{3}{\ohm}$}] (3.5,0) -- (1.5,0);
 \node[ground] at (0,0){};
\end{circuitikz}}

\resposta{$V = \SI{18}{\volt}$}
\resolucao{A fonte injeta \SI{9}{\ampere} no nó, que escoam pelos dois resistores
em paralelo:
\[
9 = \frac{V}{6}+\frac{V}{3}=\frac{V+2V}{6}=\frac{3V}{6}=\frac{V}{2}
\ \Rightarrow\ V = \SI{18}{\volt}.
\]
Equivale a $V = I\cdot(6\parallel3) = 9\cdot2 = \SI{18}{\volt}$.}
\end{questao}

\blocoaplicacao

\begin{questao}[2]{}
Determine os potenciais nodais $V_1$ e $V_2$ do circuito da figura.

\circuitoq{%
\begin{circuitikz}[scale=1,transform shape,line width=0.8pt,american]
 \draw (0,0) to[isource,l={$\SI{6}{\ampere}$}] (0,2.8) -- (1.5,2.8) coordinate(a);
 \node[above,Vcol] at (1.5,2.95) {$V_1$};
 \draw (a) to[R,l={$\SI{2}{\ohm}$}] (1.5,0) coordinate(g);
 \draw (a) to[R,l={$\SI{3}{\ohm}$}] (4.5,2.8) coordinate(b);
 \node[above,Vcol] at (4.5,2.95) {$V_2$};
 \draw (b) to[R,l={$\SI{6}{\ohm}$}] (4.5,0) -- (g) -- (0,0);
 \node[ground] at (1.5,0){};
\end{circuitikz}}

\resposta{$V_1 \approx \SI{9.82}{\volt}$; $V_2 \approx \SI{6.55}{\volt}$}
\resolucao{Com o terra na base, escrevemos a lei dos nós em cada nó.\\
\textbf{Nó $V_1$} (a fonte injeta \SI{6}{\ampere}):
\[
6=\frac{V_1}{2}+\frac{V_1-V_2}{3}.
\]
\textbf{Nó $V_2$:}
\[
\frac{V_1-V_2}{3}=\frac{V_2}{6}.
\]
Da segunda equação, $2(V_1-V_2)=V_2 \Rightarrow V_1 = \dfrac{3V_2}{2}$.
Substituindo na primeira e resolvendo:
\[
V_1=\frac{108}{11}\approx\SI{9.82}{\volt},
\qquad
V_2=\frac{72}{11}\approx\SI{6.55}{\volt}.
\]
\textbf{Verificação:} corrente da fonte $= \dfrac{V_1}{2}+\dfrac{V_1-V_2}{3}
= \dfrac{108}{22}+\dfrac{36}{33}\approx 4{,}91+1{,}09 = \SI{6}{\ampere}$. Confere.}
\end{questao}

\begin{questao}[2]{}
No nó da figura, uma fonte injeta $\SI{3}{\ampere}$ e outra retira
$\SI{1}{\ampere}$; dois resistores de $\SI{6}{\ohm}$ e $\SI{3}{\ohm}$ ligam o nó ao
terra. Determine o potencial $V$.

\circuitoq{%
\begin{circuitikz}[scale=1,transform shape,line width=0.8pt,american]
 \draw (0,0) to[isource,l={$\SI{3}{\ampere}$}] (0,2.6) -- (1.5,2.6) coordinate(n);
 \node[above,Vcol] at (1.5,2.75){$V$};
 \draw (n) to[R,l={$\SI{6}{\ohm}$}] (1.5,0) -- (0,0);
 \draw (n) -- (3.5,2.6) to[R,l={$\SI{3}{\ohm}$}] (3.5,0) -- (1.5,0);
 \draw (n) -- (5,2.6) to[isource,l={$\SI{1}{\ampere}$},invert] (5,0) -- (3.5,0);
 \node[ground] at (1.5,0){};
\end{circuitikz}}

\resposta{$V = \SI{4}{\volt}$}
\resolucao{A corrente líquida injetada no nó é $3 - 1 = \SI{2}{\ampere}$, que
escoa pelos dois resistores em paralelo:
\[
V = I_{\text{líq}}\cdot(6\parallel3) = 2\cdot2 = \SI{4}{\volt}.
\]
Conferindo pela lei dos nós: $\dfrac{V}{6}+\dfrac{V}{3} = \dfrac{4}{6}+\dfrac{4}{3}
= 0{,}67 + 1{,}33 = \SI{2}{\ampere}$, igual à corrente líquida injetada.}
\end{questao}

\blocoaprofundamento

\begin{questao}[3]{}
No circuito da figura, há uma fonte de tensão e uma de corrente. Usando análise
nodal, determine o potencial $V_A$.

\circuitoq{%
\begin{circuitikz}[scale=1,transform shape,line width=0.8pt,american]
 \draw (0,0) to[battery1,l_={$\SI{12}{\volt}$}] (0,3) -- (1.5,3)
       to[R,l={$\SI{4}{\ohm}$}] (3.5,3) coordinate(a);
 \node[above,Vcol] at (3.5,3.15) {$V_A$};
 \draw (a) to[R,l={$\SI{6}{\ohm}$}] (3.5,0) coordinate(g) -- (0,0);
 \draw (a) -- (5.5,3);
 \draw (5.5,3) to[isource,l={$\SI{2}{\ampere}$}] (5.5,0) -- (g);
 \node[ground] at (3.5,0){};
\end{circuitikz}}

\resposta{$V_A = \SI{12}{\volt}$}
\resolucao{Chamando de $V_A$ o potencial do nó central e notando que o terminal
superior da fonte de tensão está fixo em \SI{12}{\volt}, a lei dos nós em $A$
(correntes saindo $=0$; a fonte de corrente injeta \SI{2}{\ampere} \emph{no} nó):
\[
\frac{V_A-12}{4}+\frac{V_A}{6}-2=0.
\]
Multiplicando por 12: $3(V_A-12)+2V_A-24=0 \Rightarrow 5V_A = 60
\Rightarrow V_A = \SI{12}{\volt}$.\\
\textbf{Interpretação:} nesse ponto, a corrente que a fonte de tensão empurraria
pelo resistor de \SI{4}{\ohm} é nula ($V_A = \SI{12}{\volt}$ iguala os dois lados),
e toda a corrente de \SI{2}{\ampere} da fonte de corrente escoa pelo resistor de
\SI{6}{\ohm}: $2 = 12/6$. Consistente.}
\end{questao}

\begin{questao}[3]{}
Determine os potenciais nodais $V_1$, $V_2$ e $V_3$ do circuito da figura,
montando e resolvendo o sistema de equações.

\circuitoq{%
\begin{circuitikz}[scale=1,transform shape,line width=0.8pt]
 \draw (0,0) to[\fonteV,l={$\SI{12}{\volt}$}] (0,3) -- (1.4,3)
       to[R,l={$\SI{2}{\ohm}$}] (3.2,3) coordinate(n1);
 \node[above,Vcol] at (3.2,3.2){$V_1$}; \fill (n1) circle (0.07);
 \draw (n1) to[R,l={$\SI{4}{\ohm}$}] (6,3) coordinate(n2);
 \node[above,Vcol] at (6,3.2){$V_2$}; \fill (n2) circle (0.07);
 \draw (n1) to[R,l_={$\SI{4}{\ohm}$}] (3.2,0) coordinate(n3);
 \node[below,Vcol] at (3.2,-0.25){$V_3$}; \fill (n3) circle (0.07);
 \draw (n2) to[R,l={$\SI{8}{\ohm}$}] (6,1.5) coordinate(m) -- (6,0) coordinate(b2);
 \draw (n3) to[R,l_={$\SI{4}{\ohm}$}] (0,0);
 \draw (n3) -- (b2);
 \draw (m) -- (7.6,1.5) to[R,l={$\SI{8}{\ohm}$}] (7.6,-1.2) -- (3.2,-1.2) -- (n3);
 \node[ground] at (3.2,-1.2){};
\end{circuitikz}}

\resposta{$V_1 = \dfrac{76}{9}\approx\SI{8.44}{\volt}$,
$V_2=\dfrac{16}{3}\approx\SI{5.33}{\volt}$,
$V_3=\dfrac{40}{9}\approx\SI{4.44}{\volt}$}
\resolucao{\textbf{Montagem.} Com o terra na base e aplicando a lei dos nós
(soma das correntes que \emph{saem} igual a zero) em cada nó:
\[
\text{Nó 1: }\ \frac{V_1-12}{2}+\frac{V_1-V_2}{4}+\frac{V_1-V_3}{4}=0,
\]
\[
\text{Nó 2: }\ \frac{V_2-V_1}{4}+\frac{V_2}{8}+\frac{V_2-V_3}{8}=0,
\]
\[
\text{Nó 3: }\ \frac{V_3-V_1}{4}+\frac{V_3-V_2}{8}+\frac{V_3}{4}=0.
\]
\textbf{Forma matricial} (multiplicando para eliminar denominadores e agrupando
as condutâncias):
\[
\begin{bmatrix}
1{,}0 & -0{,}25 & -0{,}25\\
-0{,}25 & 0{,}5 & -0{,}125\\
-0{,}25 & -0{,}125 & 0{,}625
\end{bmatrix}
\begin{bmatrix}V_1\\V_2\\V_3\end{bmatrix}
=
\begin{bmatrix}6\\0\\0\end{bmatrix}
\]
Note a \textbf{simetria} da matriz de condutâncias --- verificação de que a
montagem está correta.

\textbf{Solução:}
\[
V_1=\frac{76}{9}\approx\SI{8.44}{\volt},\quad
V_2=\frac{16}{3}\approx\SI{5.33}{\volt},\quad
V_3=\frac{40}{9}\approx\SI{4.44}{\volt}.
\]
\textbf{Verificação no nó 2:}
$\dfrac{5{,}33-8{,}44}{4}+\dfrac{5{,}33}{8}+\dfrac{5{,}33-4{,}44}{8}
= -0{,}778+0{,}667+0{,}111 = 0$. \checkmark\\
\textbf{Observação estrutural:} a matriz nodal é o \emph{dual} da matriz de
malhas --- diagonal com a soma das condutâncias do nó, fora da diagonal o negativo
da condutância compartilhada.}
\end{questao}

\begin{questao}[3]{}
No circuito da figura, uma fonte de $\SI{6}{\volt}$ está \emph{flutuante} entre os
nós 1 e 2 (nenhum deles é o terra). Uma fonte de corrente de $\SI{2}{\ampere}$
injeta corrente no nó 1. Determine $V_1$ e $V_2$ pelo método do \textbf{supernó}.

\circuitoq{%
\begin{circuitikz}[scale=1,transform shape,line width=0.8pt,american]
 % fonte de corrente
 \draw (0,0) to[isource,l={$\SI{2}{\ampere}$}] (0,3) -- (1.6,3) coordinate(n1);
 \node[above,Vcol] at (1.6,3.2){$V_1$};
 \fill (n1) circle (0.07);
 % R1 do no1 ao terra
 \draw (n1) to[R,l={$R_1=\SI{6}{\ohm}$}] (1.6,0) coordinate(g);
 % fonte flutuante entre no1 e no2
 \draw (n1) to[battery1,l={$\SI{6}{\volt}$}] (4.6,3) coordinate(n2);
 \node[above,Vcol] at (4.6,3.2){$V_2$};
 \fill (n2) circle (0.07);
 % R2 do no2 ao terra
 \draw (n2) to[R,l={$R_2=\SI{3}{\ohm}$}] (4.6,0) -- (g) -- (0,0);
 % marcacao do superno
 \draw[dashed,laranja,line width=1pt,rounded corners=8pt]
       (1.1,2.4) rectangle (5.1,3.75);
 \node[laranja,font=\footnotesize\sffamily,right] at (5.2,3.05) {supernó};
 \node[ground] at (1.6,0){};
\end{circuitikz}}

\resposta{$V_1 = \SI{8}{\volt}$; $V_2 = \SI{2}{\volt}$}
\resolucao{\textbf{Por que o método normal falha:} a corrente que atravessa a
fonte de \SI{6}{\volt} é desconhecida, então não se pode escrever a lei dos nós
isoladamente em 1 ou em 2.

\medskip
\textbf{Equação 1 --- lei dos nós no supernó.} Traçamos uma superfície fechada
englobando os nós 1 e 2 (tracejado laranja). Tudo o que \emph{entra} deve
\emph{sair}. Entra \SI{2}{\ampere} da fonte de corrente; saem as correntes por
$R_1$ e $R_2$:
\[
\frac{V_1}{R_1}+\frac{V_2}{R_2} = 2
\quad\Longrightarrow\quad
\frac{V_1}{6}+\frac{V_2}{3} = 2.
\]
Note que a corrente \emph{interna} à fonte não aparece --- ela entra e sai da
mesma superfície, cancelando-se. \textbf{Este é o truque do método.}

\medskip
\textbf{Equação 2 --- restrição da fonte.} A fonte impõe diretamente:
\[
V_1 - V_2 = 6.
\]

\medskip
\textbf{Resolvendo.} Da segunda, $V_2 = V_1 - 6$. Substituindo na primeira e
multiplicando por 6:
\[
V_1 + 2(V_1-6) = 12
\;\Rightarrow\;
3V_1 = 24
\;\Rightarrow\;
V_1 = \SI{8}{\volt},\quad V_2 = \SI{2}{\volt}.
\]
\textbf{Verificação:} $\dfrac{8}{6}+\dfrac{2}{3} = 1{,}33+0{,}67 =
\SI{2}{\ampere}$. Confere com a fonte de corrente.}
\end{questao}

% BEGIN BANCO COMPLEMENTAR Análise Nodal
% =====================================================================
% Banco complementar --- Analise Nodal (versao visual revisada)
% N1: fundamentos; N2/N3/N4: bancos separados, todos com circuitos.
% =====================================================================

\subsubsection*{Banco complementar --- Análise Nodal}

\begin{atencao}[Organização do banco revisado]
Nos níveis 2, 3 e 4, o enunciado parte de um \textbf{circuito desenhado}, não de
uma descrição textual da topologia. As questões avançam de redes de um e dois
nós para supernós, fontes dependentes, análise nodal modificada, projeto,
sensibilidade e determinação de equivalentes. O objetivo é treinar a leitura do
circuito e a montagem das equações, e não apenas manipular sistemas já prontos.
\end{atencao}

\blocofixacao

\begin{questao}[1]{}
Em um circuito com cinco nós, quantas tensões nodais independentes existem após a escolha do nó de referência?
\resposta{4}
\resolucao{Com um nó fixado em zero, restam $5-1=4$ potenciais desconhecidos.}
\end{questao}

\begin{questao}[1]{}
Um nó $A$ recebe $\SI{3}{\ampere}$ e está ligado ao terra por $\SI{10}{\ohm}$ e $\SI{15}{\ohm}$. Determine $V_A$.
\resposta{$V_A=\SI{18}{\volt}$}
\resolucao{$V_A(1/10+1/15)=3\Rightarrow V_A=\SI{18}{\volt}$.}
\end{questao}

\begin{questao}[1]{}
Escreva a corrente, no sentido $A\to B$, em um resistor $R$ entre os nós $A$ e $B$.
\resposta{$i_{AB}=(V_A-V_B)/R$}
\resolucao{É a lei de Ohm escrita diretamente em função das tensões nodais.}
\end{questao}

\begin{questao}[1]{}
Uma fonte de corrente ideal de $I$ sai de um nó. Na forma $\mathbf G\mathbf V=\mathbf I$, qual é o sinal de sua contribuição no vetor de correntes injetadas?
\resposta{Negativo}
\resolucao{O vetor $\mathbf I$ representa corrente líquida injetada. Uma fonte que retira corrente entra como $-I$.}
\end{questao}

\begin{questao}[1]{}
Converta $\SI{4}{\ohm}$, $\SI{5}{\ohm}$ e $\SI{20}{\ohm}$, todos ligados do mesmo nó ao terra, em uma única condutância equivalente.
\resposta{$G_{eq}=\SI{0.5}{\siemens}$}
\resolucao{$G_{eq}=1/4+1/5+1/20=0{,}5\,\siemens$.}
\end{questao}

\begin{questao}[1]{}
Uma fonte ideal de tensão de $\SI{12}{\volt}$ tem o terminal negativo aterrado. Qual é a tensão do terminal positivo?
\resposta{$\SI{12}{\volt}$}
\resolucao{O terminal positivo é um nó de potencial conhecido; não é necessário escrever uma equação nodal para determiná-lo.}
\end{questao}

\begin{questao}[1]{}
Quando uma fonte ideal de tensão liga dois nós desconhecidos e nenhum deles é o terra, qual técnica nodal é indicada?
\resposta{Supernó}
\resolucao{A corrente interna da fonte não é conhecida por uma relação resistiva; escreve-se a LKC no contorno do supernó e acrescenta-se a restrição de tensão.}
\end{questao}

\begin{questao}[1]{}
Após obter $V_A=\SI{18}{\volt}$, determine a corrente em um resistor de $\SI{15}{\ohm}$ entre $A$ e o terra.
\resposta{$\SI{1.2}{\ampere}$, de $A$ para o terra}
\resolucao{$i=18/15=\SI{1.2}{\ampere}$.}
\end{questao}

% =====================================================================
% NIVEL 2 --- ANALISE NODAL: APLICACAO COM CIRCUITOS
% 13 questoes, todas com figura
% =====================================================================
\blocoaplicacao

\begin{questao}[2]{}
Determine $V$ e as correntes nos dois resistores.
\circuitoq{\begin{circuitikz}[american,scale=.95,transform shape]
\draw (0,0) to[isource,l_={$\SI{5}{\ampere}$}] (0,2.7) -- (1.8,2.7) coordinate(n);
\draw (n) to[R,l={$\SI{10}{\ohm}$}] (1.8,0) -- (0,0);
\draw (n) -- (3.8,2.7) to[R,l={$\SI{20}{\ohm}$}] (3.8,0) -- (1.8,0);
\node[above,Vcol] at (n){$V$}; \node[ground] at (1.8,0){};
\end{circuitikz}}
\resposta{$V=\SI{33.33}{\volt}$; $i_{10}=\SI{3.33}{\ampere}$; $i_{20}=\SI{1.67}{\ampere}$}
\resolucao{$5=V/10+V/20\Rightarrow V=100/3\,\si{\volt}$. As correntes são $V/R$.}
\end{questao}

\begin{questao}[2]{}
Determine a tensão nodal $V$.
\circuitoq{\begin{circuitikz}[american,scale=.9,transform shape]
\draw (0,0) to[isource,l_={$\SI{6}{\ampere}$}] (0,2.6) -- (1.4,2.6) coordinate(n);
\draw (n) to[R,l={$\SI{8}{\ohm}$}] (1.4,0) -- (0,0);
\draw (n) -- (3.3,2.6) to[R,l={$\SI{8}{\ohm}$}] (3.3,0) -- (1.4,0);
\draw (n) -- (5,2.6) to[isource,l={$\SI{2}{\ampere}$},invert] (5,0) -- (3.3,0);
\node[above,Vcol] at (n){$V$}; \node[ground] at (1.4,0){};
\end{circuitikz}}
\resposta{$V=\SI{16}{\volt}$}
\resolucao{A corrente líquida injetada é $6-2=4$ A. Logo $V(1/8+1/8)=4\Rightarrow V=16$ V.}
\end{questao}

\begin{questao}[2]{}
Determine $V_A$ usando diretamente a LKC no nó.
\circuitoq{\begin{circuitikz}[american,scale=.9,transform shape]
\draw (0,0) to[battery1,l_={$\SI{12}{\volt}$}] (0,2.7) to[R,l={$\SI{4}{\ohm}$}] (2.1,2.7) coordinate(a);
\draw (a) to[R,l={$\SI{12}{\ohm}$}] (2.1,0) -- (0,0);
\draw (a) to[R,l={$\SI{6}{\ohm}$}] (4.4,2.7) -- (4.4,0) to[battery1,l={$\SI{6}{\volt}$},invert] (4.4,2.7);
\draw (4.4,0) -- (2.1,0); \node[ground] at (2.1,0){}; \node[above,Vcol] at (a){$V_A$};
\end{circuitikz}}
\resposta{$V_A=\SI{8}{\volt}$}
\resolucao{$\frac{V_A-12}{4}+\frac{V_A-6}{6}+\frac{V_A}{12}=0\Rightarrow V_A=8$ V.}
\end{questao}

\begin{questao}[2]{}
Determine $V_1$ e $V_2$.
\circuitoq{\begin{circuitikz}[american,scale=.95,transform shape]
\draw (0,0) to[isource,l_={$\SI{3}{\ampere}$}] (0,2.8) -- (1.5,2.8) coordinate(n1);
\draw (n1) to[R,l={$\SI{10}{\ohm}$}] (1.5,0);
\draw (n1) to[R,l={$\SI{10}{\ohm}$}] (4.2,2.8) coordinate(n2);
\draw (n2) to[R,l={$\SI{10}{\ohm}$}] (4.2,0) -- (1.5,0) -- (0,0);
\node[above,Vcol] at (n1){$V_1$}; \node[above,Vcol] at (n2){$V_2$}; \node[ground] at (1.5,0){};
\end{circuitikz}}
\resposta{$V_1=\SI{20}{\volt}$; $V_2=\SI{10}{\volt}$}
\resolucao{$V_1/10+(V_1-V_2)/10=3$ e $V_2/10+(V_2-V_1)/10=0$. Da segunda, $V_1=2V_2$; então $V_2=10$ V.}
\end{questao}

\begin{questao}[2]{}
Determine $V_A$. A fonte e o resistor de $\SI{10}{\ohm}$ formam uma fonte real de tensão.
\circuitoq{\begin{circuitikz}[american,scale=.95,transform shape]
\draw (0,0) to[battery1,l_={$\SI{30}{\volt}$}] (0,2.7) to[R,l={$\SI{10}{\ohm}$}] (2.7,2.7) coordinate(a);
\draw (a) to[R,l={$\SI{15}{\ohm}$}] (2.7,0) -- (0,0); \node[ground] at (2.7,0){}; \node[above,Vcol] at (a){$V_A$};
\end{circuitikz}}
\resposta{$V_A=\SI{18}{\volt}$}
\resolucao{$\frac{V_A-30}{10}+\frac{V_A}{15}=0\Rightarrow V_A=18$ V.}
\end{questao}

\begin{questao}[2]{}
Determine $V_A$ e a corrente no resistor de $\SI{5}{\ohm}$, positiva no sentido $A\to$ nó de $\SI{20}{\volt}$.
\circuitoq{\begin{circuitikz}[american,scale=.9,transform shape]
\draw (0,0) to[battery1,l_={$\SI{20}{\volt}$}] (0,2.7) to[R,l={$\SI{5}{\ohm}$}] (2.4,2.7) coordinate(a);
\draw (a) to[R,l={$\SI{20}{\ohm}$}] (2.4,0) -- (0,0);
\draw (4.4,0) to[isource,l_={$\SI{2}{\ampere}$}] (4.4,2.7) -- (a); \draw (4.4,0)--(2.4,0);
\node[ground] at (2.4,0){}; \node[above,Vcol] at (a){$V_A$};
\end{circuitikz}}
\resposta{$V_A=\SI{24}{\volt}$; $i=\SI{0.8}{\ampere}$}
\resolucao{$\frac{V_A-20}{5}+\frac{V_A}{20}=2\Rightarrow V_A=24$ V; $i=(24-20)/5=0{,}8$ A.}
\end{questao}

\begin{questao}[2]{}
Determine $V_1$ e $V_2$ na rede em escada.
\circuitoq{\begin{circuitikz}[american,scale=.95,transform shape]
\draw (0,0) to[isource,l_={$\SI{4}{\ampere}$}] (0,2.8) -- (1.4,2.8) coordinate(n1);
\draw (n1) to[R,l={$\SI{5}{\ohm}$}] (1.4,0);
\draw (n1) to[R,l={$\SI{10}{\ohm}$}] (4.1,2.8) coordinate(n2);
\draw (n2) to[R,l={$\SI{10}{\ohm}$}] (4.1,0) -- (1.4,0) -- (0,0);
\node[ground] at (1.4,0){}; \node[above,Vcol] at (n1){$V_1$}; \node[above,Vcol] at (n2){$V_2$};
\end{circuitikz}}
\resposta{$V_1=\SI{16}{\volt}$; $V_2=\SI{8}{\volt}$}
\resolucao{$V_1/5+(V_1-V_2)/10=4$ e $V_2/10+(V_2-V_1)/10=0$.}
\end{questao}

\begin{questao}[2]{}
A fonte de $\SI{1}{\ampere}$ transfere corrente do nó 1 para o nó 2. Determine $V_1$ e $V_2$.
\circuitoq{\begin{circuitikz}[american,scale=.9,transform shape]
\draw (0,0) to[isource,l_={$\SI{5}{\ampere}$}] (0,3) -- (1.5,3) coordinate(n1);
\draw (n1) to[R,l={$\SI{6}{\ohm}$}] (1.5,0);
\draw (n1) to[R,l={$\SI{6}{\ohm}$}] (4.1,3) coordinate(n2);
\draw (n2) to[R,l={$\SI{3}{\ohm}$}] (4.1,0) -- (1.5,0)--(0,0);
\draw (n1) to[isource,l={$\SI{1}{\ampere}$}] (n2);
\node[ground] at (1.5,0){}; \node[above,Vcol] at (n1){$V_1$}; \node[above,Vcol] at (n2){$V_2$};
\end{circuitikz}}
\resposta{$V_1=\SI{15.6}{\volt}$; $V_2=\SI{7.2}{\volt}$}
\resolucao{Nó 1: $V_1/6+(V_1-V_2)/6+1=5$. Nó 2: $V_2/3+(V_2-V_1)/6=1$. A solução é $V_1=78/5$ V e $V_2=36/5$ V.}
\end{questao}

\begin{questao}[2]{}
Determine os três potenciais da escada resistiva.
\circuitoq{\begin{circuitikz}[american,scale=.85,transform shape]
\draw (0,0) to[isource,l_={$\SI{4}{\ampere}$}] (0,2.8) -- (1.2,2.8) coordinate(n1);
\draw (n1) to[R,l={$\SI{10}{\ohm}$}] (1.2,0);
\draw (n1) to[R,l={$\SI{10}{\ohm}$}] (3.7,2.8) coordinate(n2);
\draw (n2) to[R,l={$\SI{10}{\ohm}$}] (3.7,0);
\draw (n2) to[R,l={$\SI{10}{\ohm}$}] (6.2,2.8) coordinate(n3);
\draw (n3) to[R,l={$\SI{10}{\ohm}$}] (6.2,0) -- (1.2,0) -- (0,0);
\node[ground] at (1.2,0){}; \node[above,Vcol] at(n1){$V_1$};\node[above,Vcol] at(n2){$V_2$};\node[above,Vcol] at(n3){$V_3$};
\end{circuitikz}}
\resposta{$V_1=\SI{25}{\volt}$; $V_2=\SI{10}{\volt}$; $V_3=\SI{5}{\volt}$}
\resolucao{As três LKC formam uma matriz tridiagonal. Da última, $V_2=2V_3$; da segunda, $V_1=2{,}5V_2$; na primeira obtém-se $V_2=10$ V.}
\end{questao}

\begin{questao}[2]{}
Duas fontes de corrente injetam corrente em nós diferentes. Determine $V_1$ e $V_2$.
\circuitoq{\begin{circuitikz}[american,scale=.9,transform shape]
\draw (0,0) to[isource,l_={$\SI{4}{\ampere}$}] (0,2.8) -- (1.5,2.8) coordinate(n1);
\draw (n1) to[R,l={$\SI{4}{\ohm}$}] (1.5,0);
\draw (n1) to[R,l={$\SI{8}{\ohm}$}] (4.5,2.8) coordinate(n2);
\draw (n2) to[R,l={$\SI{4}{\ohm}$}] (4.5,0);
\draw (6.2,0) to[isource,l_={$\SI{2}{\ampere}$}] (6.2,2.8) -- (n2); \draw (6.2,0)--(1.5,0)--(0,0);
\node[ground] at (1.5,0){};\node[above,Vcol] at(n1){$V_1$};\node[above,Vcol] at(n2){$V_2$};
\end{circuitikz}}
\resposta{$V_1=\SI{14}{\volt}$; $V_2=\SI{10}{\volt}$}
\resolucao{$V_1/4+(V_1-V_2)/8=4$ e $V_2/4+(V_2-V_1)/8=2$.}
\end{questao}

\begin{questao}[2]{}
Determine $V_1$ e $V_2$.
\circuitoq{\begin{circuitikz}[american,scale=.9,transform shape]
\draw (0,0) to[battery1,l_={$\SI{18}{\volt}$}] (0,2.8) to[R,l={$\SI{6}{\ohm}$}] (2.4,2.8) coordinate(n1);
\draw (n1) to[R,l={$\SI{12}{\ohm}$}] (2.4,0);
\draw (n1) to[R,l={$\SI{6}{\ohm}$}] (5.1,2.8) coordinate(n2);
\draw (n2) to[R,l={$\SI{6}{\ohm}$}] (5.1,0);
\draw (6.7,0) to[isource,l_={$\SI{1}{\ampere}$}] (6.7,2.8)--(n2);\draw (6.7,0)--(0,0);
\node[ground] at (2.4,0){};\node[above,Vcol] at(n1){$V_1$};\node[above,Vcol] at(n2){$V_2$};
\end{circuitikz}}
\resposta{$V_1=\SI{10.5}{\volt}$; $V_2=\SI{8.25}{\volt}$}
\resolucao{$\frac{V_1-18}{6}+\frac{V_1}{12}+\frac{V_1-V_2}{6}=0$ e $\frac{V_2}{6}+\frac{V_2-V_1}{6}=1$.}
\end{questao}

\begin{questao}[2]{}
A fonte horizontal força $\SI{2}{\ampere}$ do nó 1 para o nó 2. Determine as tensões nodais.
\circuitoq{\begin{circuitikz}[american,scale=.9,transform shape]
\draw (0,0) to[battery1,l_={$\SI{24}{\volt}$}] (0,2.8) to[R,l={$\SI{6}{\ohm}$}] (2.6,2.8) coordinate(n1);
\draw (n1) to[R,l={$\SI{12}{\ohm}$}] (2.6,0);
\draw (n1) to[isource,l={$\SI{2}{\ampere}$}] (5.2,2.8) coordinate(n2);
\draw (n2) to[R,l={$\SI{6}{\ohm}$}] (5.2,0)--(0,0);
\node[ground] at (2.6,0){};\node[above,Vcol] at(n1){$V_1$};\node[above,Vcol] at(n2){$V_2$};
\end{circuitikz}}
\resposta{$V_1=\SI{8}{\volt}$; $V_2=\SI{12}{\volt}$}
\resolucao{Nó 1: $(V_1-24)/6+V_1/12+2=0\Rightarrow V_1=8$ V. Nó 2: $V_2/6=2\Rightarrow V_2=12$ V.}
\end{questao}

\begin{questao}[2]{}
Determine $V_1$, $V_2$ e a corrente no resistor de $\SI{4}{\ohm}$ entre os nós.
\circuitoq{\begin{circuitikz}[american,scale=.9,transform shape]
\draw (0,0) to[battery1,l_={$\SI{10}{\volt}$}] (0,2.8) to[R,l={$\SI{2}{\ohm}$}] (2.4,2.8) coordinate(n1);
\draw (n1) to[R,l={$\SI{4}{\ohm}$}] (5,2.8) coordinate(n2);
\draw (n1) to[R,l={$\SI{4}{\ohm}$}] (2.4,0);
\draw (n2) to[R,l={$\SI{8}{\ohm}$}] (5,0)--(0,0);
\node[ground] at (2.4,0){};\node[above,Vcol] at(n1){$V_1$};\node[above,Vcol] at(n2){$V_2$};
\end{circuitikz}}
\resposta{$V_1=\SI{6}{\volt}$; $V_2=\SI{4}{\volt}$; $i_{12}=\SI{0.5}{\ampere}$}
\resolucao{$\frac{V_1-10}{2}+\frac{V_1}{4}+\frac{V_1-V_2}{4}=0$ e $\frac{V_2}{8}+\frac{V_2-V_1}{4}=0$. Logo $V_1=6$ V, $V_2=4$ V e $i_{12}=(6-4)/4=0{,}5$ A.}
\end{questao}

% =====================================================================
% NIVEL 3 --- ANALISE NODAL: APROFUNDAMENTO VISUAL
% 12 questoes; inclui 9 problemas com fontes dependentes
% =====================================================================
\blocoaprofundamento

\begin{questao}[3]{}
No circuito, $i_1$ e $i_2$ têm os sentidos indicados. Determine $i_1+i_2$.
\circuitoq{\begin{circuitikz}[american,scale=.92,transform shape]
\draw (0,0) to[isource,l_={$\SI{5}{\ampere}$}] (0,3) -- (1.8,3) coordinate(n);
\draw (n) to[R,l={$\SI{5}{\ohm}$},i>^={$i_2$}] (1.8,0) -- (0,0);
\draw (n) -- (3.9,3) to[cisource,l={$4i_2$}] (3.9,0) -- (1.8,0);
\draw (n) -- (6,3) to[R,l={$\SI{1}{\ohm}$},i>^={$i_1$}] (6,0) -- (3.9,0);
\node[ground] at (1.8,0){};
\end{circuitikz}}
\resposta{$i_1+i_2=\SI{15}{\ampere}$}
\resolucao{Como os resistores estão em paralelo, $i_1=V/1$ e $i_2=V/5$, logo $i_1=5i_2$. Pela LKC no nó superior, $5+4i_2=i_1+i_2$. Assim $5+4i_2=6i_2$, de onde $i_2=2{,}5$ A, $i_1=12{,}5$ A e $i_1+i_2=15$ A.}
\end{questao}

\begin{questao}[3]{}
A fonte dependente drena $0{,}1V_1$ ampère do nó 2. Determine $V_1$ e $V_2$.
\circuitoq{\begin{circuitikz}[american,scale=.9,transform shape]
\draw (0,0) to[isource,l_={$\SI{6}{\ampere}$}] (0,3) -- (1.5,3) coordinate(n1);
\draw (n1) to[R,l={$\SI{5}{\ohm}$}] (1.5,0);
\draw (n1) to[R,l={$\SI{5}{\ohm}$}] (4.4,3) coordinate(n2);
\draw (n2) to[R,l={$\SI{10}{\ohm}$}] (4.4,0);
\draw (n2) -- (6.3,3) to[cisource,l={$0{,}1V_1$},invert] (6.3,0) -- (0,0);
\node[ground] at (1.5,0){};\node[above,Vcol] at(n1){$V_1$};\node[above,Vcol] at(n2){$V_2$};
\end{circuitikz}}
\resposta{$V_1=\SI{18}{\volt}$; $V_2=\SI{6}{\volt}$}
\resolucao{$V_1/5+(V_1-V_2)/5=6$ e $V_2/10+(V_2-V_1)/5+0{,}1V_1=0$. A solução é $(18,6)$ V.}
\end{questao}

\begin{questao}[3]{}
A fonte de tensão controlada impõe $V_2=2V_1$. Determine $V_1$ e $V_2$.
\circuitoq{\begin{circuitikz}[american,scale=.9,transform shape]
\draw (0,0) to[isource,l_={$\SI{3}{\ampere}$}] (0,3) -- (1.6,3) coordinate(n1);
\draw (n1) to[R,l={$\SI{4}{\ohm}$}] (1.6,0);
\draw (n1) to[R,l={$\SI{8}{\ohm}$}] (4.6,3) coordinate(n2);
\draw (n2) to[cvsource,l={$2V_1$}] (4.6,0) -- (0,0);
\node[ground] at (1.6,0){};\node[above,Vcol] at(n1){$V_1$};\node[above,Vcol] at(n2){$V_2$};
\end{circuitikz}}
\resposta{$V_1=\SI{24}{\volt}$; $V_2=\SI{48}{\volt}$}
\resolucao{A LKC em 1 dá $V_1/4+(V_1-V_2)/8=3$. Com $V_2=2V_1$, resulta $V_1/8=3$, portanto $V_1=24$ V e $V_2=48$ V.}
\end{questao}

\begin{questao}[3]{}
A corrente de controle é $i_x=V_1/4$. A fonte dependente injeta $0{,}5i_x$ no nó 2. Determine $V_1$ e $V_2$.
\circuitoq{\begin{circuitikz}[american,scale=.9,transform shape]
\draw (0,0) to[isource,l_={$\SI{4}{\ampere}$}] (0,3) -- (1.5,3) coordinate(n1);
\draw (n1) to[R,l={$\SI{4}{\ohm}$},i>^={$i_x$}] (1.5,0);
\draw (n1) to[R,l={$\SI{8}{\ohm}$}] (4.5,3) coordinate(n2);
\draw (n2) to[R,l={$\SI{8}{\ohm}$}] (4.5,0);
\draw (6.3,0) to[cisource,l_={$0{,}5i_x$}] (6.3,3)--(n2);\draw (6.3,0)--(0,0);
\node[ground] at (1.5,0){};\node[above,Vcol] at(n1){$V_1$};\node[above,Vcol] at(n2){$V_2$};
\end{circuitikz}}
\resposta{$V_1=\SI{16}{\volt}$; $V_2=\SI{16}{\volt}$}
\resolucao{$V_1/4+(V_1-V_2)/8=4$ e $V_2/8+(V_2-V_1)/8=0{,}5(V_1/4)$. A solução é $V_1=V_2=16$ V; logo o resistor entre os nós não conduz corrente.}
\end{questao}

\begin{questao}[3]{}
A fonte de tensão controlada entre os nós satisfaz $V_1-V_2=5i_x$, com $i_x=V_2/5$. Determine $V_1$ e $V_2$ por supernó.
\circuitoq{\begin{circuitikz}[american,scale=.9,transform shape]
\draw (0,0) to[isource,l_={$\SI{5}{\ampere}$}] (0,3)--(1.5,3) coordinate(n1);
\draw (n1) to[R,l={$\SI{10}{\ohm}$}] (1.5,0);
\draw (n1) to[cvsource,l={$5i_x$}] (4.7,3) coordinate(n2);
\draw (n2) to[R,l={$\SI{5}{\ohm}$},i>^={$i_x$}] (4.7,0)--(0,0);
\node[ground] at (1.5,0){};\node[above,Vcol] at(n1){$V_1$};\node[above,Vcol] at(n2){$V_2$};
\end{circuitikz}}
\resposta{$V_1=\SI{25}{\volt}$; $V_2=\SI{12.5}{\volt}$}
\resolucao{Supernó: $V_1/10+V_2/5=5$. Restrição: $V_1-V_2=5(V_2/5)=V_2$, logo $V_1=2V_2$. Da LKC, $0{,}4V_2=5$, portanto $V_2=12{,}5$ V e $V_1=25$ V.}
\end{questao}

\begin{questao}[3]{}
Determine $V_1$, $V_2$ e a corrente na fonte de $\SI{20}{\volt}$, adotada positiva de 1 para 2.
\circuitoq{\begin{circuitikz}[american,scale=.9,transform shape]
\draw (0,0) to[isource,l_={$\SI{4}{\ampere}$}] (0,3)--(1.5,3) coordinate(n1);
\draw (n1) to[R,l={$\SI{10}{\ohm}$}] (1.5,0);
\draw (n1) to[battery1,l={$\SI{20}{\volt}$}] (4.7,3) coordinate(n2);
\draw (n2) to[R,l={$\SI{10}{\ohm}$}] (4.7,0)--(0,0);
\node[ground] at (1.5,0){};\node[above,Vcol] at(n1){$V_1$};\node[above,Vcol] at(n2){$V_2$};
\end{circuitikz}}
\resposta{$V_1=\SI{30}{\volt}$; $V_2=\SI{10}{\volt}$; $i_s=\SI{1}{\ampere}$ de 1 para 2}
\resolucao{Supernó: $V_1/10+V_2/10=4$ e $V_1-V_2=20$. Logo $(V_1,V_2)=(30,10)$ V. No nó 1, dos $4$ A injetados, $3$ A descem pelo resistor; o $1$ A restante atravessa a fonte de tensão de 1 para 2.}
\end{questao}

\begin{questao}[3]{}
Determine $V_1$, $V_2$, $V_3$ e a corrente $i_{23}$ no resistor de ponte de $\SI{6}{\ohm}$, positiva de 2 para 3.
\circuitoq{\begin{circuitikz}[american,scale=.82,transform shape]
\coordinate (g) at (0,-1.6);
\draw (g) to[isource,l_={$\SI{3}{\ampere}$}] (0,1.6) -- (1.3,1.6) coordinate(n1);
\draw (n1) to[R,l={$\SI{6}{\ohm}$}] (4.2,3.0) coordinate(n2);
\draw (n1) to[R,l_={$\SI{12}{\ohm}$}] (4.2,0.2) coordinate(n3);
\draw (n2) to[R,l={$\SI{6}{\ohm}$}] (n3);
\draw (n2) to[R,l={$\SI{12}{\ohm}$}] (4.2,-1.6);
\draw (n3) to[R,l_={$\SI{12}{\ohm}$}] (4.2,-1.6);
\draw (4.2,-1.6) -- (g);
\node[ground] at(g){};
\node[above,Vcol] at(n1){$V_1$};\node[above,Vcol] at(n2){$V_2$};\node[right,Vcol] at(n3){$V_3$};
\end{circuitikz}}
\resposta{$V_1=\dfrac{576}{19}\,\si{\volt}\approx\SI{30.32}{\volt}$; $V_2=\dfrac{360}{19}\,\si{\volt}\approx\SI{18.95}{\volt}$; $V_3=\dfrac{324}{19}\,\si{\volt}\approx\SI{17.05}{\volt}$; $i_{23}=\dfrac{6}{19}\,\si{\ampere}\approx\SI{0.316}{\ampere}$}
\resolucao{As LKC são $(V_1-V_2)/6+(V_1-V_3)/12=3$, $V_2/12+(V_2-V_1)/6+(V_2-V_3)/6=0$ e $V_3/12+(V_3-V_1)/12+(V_3-V_2)/6=0$. A solução é $(576,360,324)/19$ V. Assim $i_{23}=(V_2-V_3)/6=6/19$ A.}
\end{questao}

\begin{questao}[3]{}
A fonte controlada transfere corrente do nó 1 para o nó 2 no valor $0{,}05V_2$. Determine $V_1$ e $V_2$.
\circuitoq{\begin{circuitikz}[american,scale=.9,transform shape]
\draw (0,0) to[isource,l_={$\SI{5}{\ampere}$}] (0,3)--(1.5,3) coordinate(n1);
\draw (n1) to[R,l={$\SI{5}{\ohm}$}] (1.5,0);
\draw (n1) to[R,l={$\SI{10}{\ohm}$}] (4.5,3) coordinate(n2);
\draw (n2) to[R,l={$\SI{10}{\ohm}$}] (4.5,0);
\draw (n1) to[cisource,l={$0{,}05V_2$}] (n2);\draw (4.5,0)--(0,0);
\node[ground] at(1.5,0){};\node[above,Vcol] at(n1){$V_1$};\node[above,Vcol] at(n2){$V_2$};
\end{circuitikz}}
\resposta{$V_1=\SI{18.75}{\volt}$; $V_2=\SI{12.5}{\volt}$}
\resolucao{$V_1/5+(V_1-V_2)/10+0{,}05V_2=5$ e $V_2/10+(V_2-V_1)/10-0{,}05V_2=0$. Resulta $V_1=75/4$ V e $V_2=25/2$ V.}
\end{questao}

\begin{questao}[3]{}
A fonte flutuante é controlada e impõe $V_1-V_2=0{,}5V_2$. Determine as tensões nodais.
\circuitoq{\begin{circuitikz}[american,scale=.9,transform shape]
\draw (0,0) to[isource,l_={$\SI{3}{\ampere}$}] (0,3)--(1.5,3) coordinate(n1);
\draw (n1) to[R,l={$\SI{6}{\ohm}$}] (1.5,0);
\draw (n1) to[cvsource,l={$0{,}5V_2$}] (4.7,3) coordinate(n2);
\draw (n2) to[R,l={$\SI{3}{\ohm}$}] (4.7,0)--(0,0);
\node[ground] at(1.5,0){};\node[above,Vcol] at(n1){$V_1$};\node[above,Vcol] at(n2){$V_2$};
\end{circuitikz}}
\resposta{$V_1=\dfrac{54}{7}\,\si{\volt}\approx\SI{7.71}{\volt}$; $V_2=\dfrac{36}{7}\,\si{\volt}\approx\SI{5.14}{\volt}$}
\resolucao{Supernó: $V_1/6+V_2/3=3$. A restrição dá $V_1=1{,}5V_2$. Substituindo: $7V_2/12=3\Rightarrow V_2=36/7$ V.}
\end{questao}

\begin{questao}[3]{}
Determine a resistência vista pelo terminal usando a fonte de teste de $\SI{1}{\ampere}$. A fonte dependente injeta $0{,}05V$ no nó.
\circuitoq{\begin{circuitikz}[american,scale=.95,transform shape]
\draw (0,0) to[isource,l_={$I_t=\SI{1}{\ampere}$}] (0,3)--(2,3) coordinate(n);
\draw (n) to[R,l={$\SI{10}{\ohm}$}] (2,0)--(0,0);
\draw (4,0) to[cisource,l_={$0{,}05V$}] (4,3)--(n);\draw (4,0)--(2,0);
\node[ground] at(2,0){};\node[above,Vcol] at(n){$V$};
\end{circuitikz}}
\resposta{$R_{in}=\SI{20}{\ohm}$}
\resolucao{$V/10=1+0{,}05V\Rightarrow 0{,}05V=1\Rightarrow V=20$ V. Como $I_t=1$ A, $R_{in}=V/I_t=20\,\ohm$.}
\end{questao}

\begin{questao}[3]{}
No circuito, a fonte dependente injeta $0{,}1V$ ampère. Determine $V$ e verifique o balanço de potência.
\circuitoq{\begin{circuitikz}[american,scale=.9,transform shape]
\draw (0,0) to[isource,l_={$\SI{2}{\ampere}$}] (0,3)--(1.5,3) coordinate(n);
\draw (n) to[R,l={$\SI{4}{\ohm}$}] (1.5,0);
\draw (n)--(3.6,3) to[R,l={$\SI{8}{\ohm}$}] (3.6,0)--(1.5,0);
\draw (5.4,0) to[cisource,l_={$0{,}1V$}] (5.4,3)--(n);\draw (5.4,0)--(1.5,0);
\node[ground] at(1.5,0){};\node[above,Vcol] at(n){$V$};
\end{circuitikz}}
\resposta{$V=\dfrac{80}{11}\,\si{\volt}\approx\SI{7.27}{\volt}$; balanço fecha em $\SI{19.83}{\watt}$}
\resolucao{$V/4+V/8=2+0{,}1V\Rightarrow V=80/11$ V. Os resistores dissipam $V^2(1/4+1/8)=2400/121\approx19{,}83$ W. As fontes entregam $2V+0{,}1V^2=160/11+640/121=2400/121$ W.}
\end{questao}

\begin{questao}[3]{}
A fonte de $\SI{8}{\volt}$ é flutuante. Determine $V_1$, $V_2$ e $V_3$.
\circuitoq{\begin{circuitikz}[american,scale=.85,transform shape]
\draw (0,0) to[isource,l_={$\SI{5}{\ampere}$}] (0,3)--(1.3,3) coordinate(n1);
\draw (n1) to[R,l={$\SI{4}{\ohm}$}] (1.3,0);
\draw (n1) to[battery1,l={$\SI{8}{\volt}$}] (3.9,3) coordinate(n2);
\draw (n2) to[R,l={$\SI{8}{\ohm}$}] (3.9,0);
\draw (n2) to[R,l={$\SI{4}{\ohm}$}] (6.3,3) coordinate(n3);
\draw (n3) to[R,l={$\SI{4}{\ohm}$}] (6.3,0)--(0,0);
\node[ground] at(1.3,0){};\node[above,Vcol] at(n1){$V_1$};\node[above,Vcol] at(n2){$V_2$};\node[above,Vcol] at(n3){$V_3$};
\end{circuitikz}}
\resposta{$V_1=\SI{14}{\volt}$; $V_2=\SI{6}{\volt}$; $V_3=\SI{3}{\volt}$}
\resolucao{Supernó $1$--$2$: $V_1/4+V_2/8+(V_2-V_3)/4=5$; restrição $V_1-V_2=8$; nó 3: $V_3/4+(V_3-V_2)/4=0$. A solução é $(14,6,3)$ V.}
\end{questao}

% =====================================================================
% NIVEL 4 --- ANALISE NODAL: DESAFIOS COM CIRCUITOS
% 10 questoes, todas com figura
% =====================================================================
\blocodesafio

\begin{questao}[4]{}
A fonte dependente injeta $gV$ no nó. Obtenha $V(g)$, determine o valor crítico de $g$ e calcule $V$ para $g=\SI{0.04}{\siemens}$.
\circuitoq{\begin{circuitikz}[american,scale=.95,transform shape]
\draw (0,0) to[isource,l_={$\SI{1}{\ampere}$}] (0,3)--(2,3) coordinate(n);
\draw (n) to[R,l={$\SI{20}{\ohm}$}] (2,0)--(0,0);
\draw (4.2,0) to[cisource,l_={$gV$}] (4.2,3)--(n);\draw (4.2,0)--(2,0);
\node[ground] at(2,0){};\node[above,Vcol] at(n){$V$};
\end{circuitikz}}
\resposta{$V(g)=1/(0{,}05-g)$; $g_c=\SI{0.05}{\siemens}$; para $g=\SI{0.04}{\siemens}$, $V=\SI{100}{\volt}$}
\resolucao{LKC: $V/20=1+gV$. Assim $V=1/(0{,}05-g)$. Em $g=0{,}05$ S o coeficiente nodal se anula: a rede linear torna-se singular.}
\end{questao}

\begin{questao}[4]{}
A fonte controlada injeta $kV_1$ no nó 2. Determine $k$ para que $V_2=\SI{10}{\volt}$ e encontre $V_1$.
\circuitoq{\begin{circuitikz}[american,scale=.9,transform shape]
\draw (0,0) to[isource,l_={$\SI{2}{\ampere}$}] (0,3)--(1.5,3) coordinate(n1);
\draw (n1) to[R,l={$\SI{5}{\ohm}$}] (1.5,0);
\draw (n1) to[R,l={$\SI{5}{\ohm}$}] (4.5,3) coordinate(n2);
\draw (n2) to[R,l={$\SI{10}{\ohm}$}] (4.5,0);
\draw (6.3,0) to[cisource,l_={$kV_1$}] (6.3,3)--(n2);\draw (6.3,0)--(0,0);
\node[ground] at(1.5,0){};\node[above,Vcol] at(n1){$V_1$};\node[above,Vcol] at(n2){$V_2$};
\end{circuitikz}}
\resposta{$k=\SI{0.1}{\siemens}$; $V_1=\SI{10}{\volt}$}
\resolucao{Impondo $V_2=10$ V na LKC do nó 1: $V_1/5+(V_1-10)/5=2\Rightarrow V_1=10$ V. No nó 2: $10/10+(10-10)/5=k(10)$, logo $k=0{,}1$ S.}
\end{questao}

\begin{questao}[4]{}
A fonte dependente injeta $2i_x$ no nó 3, com $i_x=(V_1-V_2)/4$. Determine as três tensões e identifique um ramo que fica sem corrente.
\circuitoq{\begin{circuitikz}[american,scale=.82,transform shape]
\draw (0,0) to[isource,l_={$\SI{4}{\ampere}$}] (0,3)--(1.3,3) coordinate(n1);
\draw (n1) to[R,l={$\SI{4}{\ohm}$}] (1.3,0);
\draw (n1) to[R,l={$\SI{4}{\ohm}$},i>^={$i_x$}] (3.8,3) coordinate(n2);
\draw (n2) to[R,l={$\SI{8}{\ohm}$}] (3.8,0);
\draw (n2) to[R,l={$\SI{4}{\ohm}$}] (6.3,3) coordinate(n3);
\draw (n3) to[R,l={$\SI{4}{\ohm}$}] (6.3,0);
\draw (7.9,0) to[cisource,l_={$2i_x$}] (7.9,3)--(n3);\draw (7.9,0)--(0,0);
\node[ground] at(1.3,0){};\node[above,Vcol] at(n1){$V_1$};\node[above,Vcol] at(n2){$V_2$};\node[above,Vcol] at(n3){$V_3$};
\end{circuitikz}}
\resposta{$V_1=\SI{12}{\volt}$; $V_2=\SI{8}{\volt}$; $V_3=\SI{8}{\volt}$; o resistor entre 2 e 3 conduz corrente nula}
\resolucao{As LKC são $V_1/4+(V_1-V_2)/4=4$, $V_2/8+(V_2-V_1)/4+(V_2-V_3)/4=0$ e $V_3/4+(V_3-V_2)/4=2(V_1-V_2)/4$. A solução é $(12,8,8)$ V.}
\end{questao}

\begin{questao}[4]{}
A fonte de tensão controlada satisfaz $V_1-V_2=5i_x$, $i_x=V_2/10$, e a fonte de corrente controlada injeta $0{,}05V_1$ no nó 2. Determine $V_1$ e $V_2$.
\circuitoq{\begin{circuitikz}[american,scale=.88,transform shape]
\draw (0,0) to[isource,l_={$\SI{3}{\ampere}$}] (0,3)--(1.4,3) coordinate(n1);
\draw (n1) to[R,l={$\SI{10}{\ohm}$}] (1.4,0);
\draw (n1) to[cvsource,l={$5i_x$}] (4.5,3) coordinate(n2);
\draw (n2) to[R,l={$\SI{10}{\ohm}$},i>^={$i_x$}] (4.5,0);
\draw (6.3,0) to[cisource,l_={$0{,}05V_1$}] (6.3,3)--(n2);\draw (6.3,0)--(0,0);
\node[ground] at(1.4,0){};\node[above,Vcol] at(n1){$V_1$};\node[above,Vcol] at(n2){$V_2$};
\end{circuitikz}}
\resposta{$V_1=\dfrac{180}{7}\,\si{\volt}\approx\SI{25.71}{\volt}$; $V_2=\dfrac{120}{7}\,\si{\volt}\approx\SI{17.14}{\volt}$}
\resolucao{A restrição é $V_1-V_2=5(V_2/10)=0{,}5V_2$, logo $V_1=1{,}5V_2$. No supernó: $V_1/10+V_2/10=3+0{,}05V_1$. Substituindo, $0{,}175V_2=3$.}
\end{questao}

\begin{questao}[4]{}
Determine o equivalente de Thévenin visto no nó 2. A fonte dependente injeta $0{,}05V_1$ no nó 2.
\circuitoq{\begin{circuitikz}[american,scale=.9,transform shape]
\draw (0,0) to[isource,l_={$\SI{2}{\ampere}$}] (0,3)--(1.5,3) coordinate(n1);
\draw (n1) to[R,l={$\SI{10}{\ohm}$}] (1.5,0);
\draw (n1) to[R,l={$\SI{5}{\ohm}$}] (4.7,3) coordinate(n2);
\draw (6.5,0) to[cisource,l_={$0{,}05V_1$}] (6.5,3)--(n2);\draw (6.5,0)--(0,0);
\draw (n2)--(5.7,3);\draw (5.7,2.75)--(5.7,3.25);\node[right] at(5.8,3){porta};
\node[ground] at(1.5,0){};\node[above,Vcol] at(n1){$V_1$};\node[above,Vcol] at(n2){$V_2$};
\end{circuitikz}}
\resposta{$V_{Th}=\SI{50}{\volt}$; $R_{Th}=\SI{30}{\ohm}$}
\resolucao{Com a porta aberta: $(V_2-V_1)/5=0{,}05V_1\Rightarrow V_2=1{,}25V_1$. No nó 1, $V_1/10+(V_1-V_2)/5=2$, resultando $V_1=40$ V e $V_{Th}=50$ V. Para $R_{Th}$, abra a fonte de $2$ A e injete $I_t$ na porta. Da LKC do nó 1, $V_1=(2/3)V_2$; no nó 2 obtém-se $I_t=V_2/30$, logo $R_{Th}=30\,\ohm$.}
\end{questao}

\begin{questao}[4]{}
A fonte controlada impõe $V_2=kV_1$. Obtenha $V_1(k)$ e determine o valor de $k$ para o qual a solução se torna singular.
\circuitoq{\begin{circuitikz}[american,scale=.9,transform shape]
\draw (0,0) to[isource,l_={$\SI{3}{\ampere}$}] (0,3)--(1.5,3) coordinate(n1);
\draw (n1) to[R,l={$\SI{10}{\ohm}$}] (1.5,0);
\draw (n1) to[R,l={$\SI{5}{\ohm}$}] (4.5,3) coordinate(n2);
\draw (n2) to[cvsource,l={$kV_1$}] (4.5,0)--(0,0);
\node[ground] at(1.5,0){};\node[above,Vcol] at(n1){$V_1$};\node[above,Vcol] at(n2){$V_2$};
\end{circuitikz}}
\resposta{$V_1=3/(0{,}3-0{,}2k)$; singular em $k=1{,}5$}
\resolucao{LKC no nó 1: $V_1/10+(V_1-kV_1)/5=3$. O coeficiente é $0{,}1+0{,}2(1-k)=0{,}3-0{,}2k$. Ele zera em $k=1{,}5$.}
\end{questao}

\begin{questao}[4]{}
Escolha $R_L$ para obter $V=\SI{20}{\volt}$. Depois calcule a sensibilidade $dV/dR_L$ no ponto de projeto. A fonte dependente injeta $0{,}02V$ no nó.
\circuitoq{\begin{circuitikz}[american,scale=.88,transform shape]
\draw (0,0) to[isource,l_={$\SI{2}{\ampere}$}] (0,3)--(1.5,3) coordinate(n);
\draw (n) to[R,l={$\SI{20}{\ohm}$}] (1.5,0);
\draw (n)--(3.7,3) to[R,l={$R_L$}] (3.7,0)--(1.5,0);
\draw (5.5,0) to[cisource,l_={$0{,}02V$}] (5.5,3)--(n);\draw (5.5,0)--(1.5,0);
\node[ground] at(1.5,0){};\node[above,Vcol] at(n){$V$};
\end{circuitikz}}
\resposta{$R_L=\dfrac{100}{7}\,\ohm\approx\SI{14.29}{\ohm}$; $dV/dR_L\approx\SI{0.98}{\volt\per\ohm}$}
\resolucao{$V(1/20+1/R_L-0{,}02)=2$, portanto $V=2R_L/(1+0{,}03R_L)$. Impondo $V=20$ V, resulta $R_L=100/7\,\ohm$. A derivada é $dV/dR_L=2/(1+0{,}03R_L)^2\approx0{,}98$ V/$\ohm$.}
\end{questao}

\begin{questao}[4]{}
Resolva o circuito por análise nodal modificada, tomando como incógnitas $V_1$, $V_2$ e a corrente $i_s$ da fonte flutuante de $\SI{10}{\volt}$. A fonte controlada drena $0{,}05V_1$ do nó 2.
\circuitoq{\begin{circuitikz}[american,scale=.88,transform shape]
\draw (0,0) to[isource,l_={$\SI{3}{\ampere}$}] (0,3)--(1.5,3) coordinate(n1);
\draw (n1) to[R,l={$\SI{10}{\ohm}$}] (1.5,0);
\draw (n1) to[battery1,l={$\SI{10}{\volt}$}] (4.7,3) coordinate(n2);
\draw (n2) to[R,l={$\SI{20}{\ohm}$}] (4.7,0);
\draw (n2)--(6.5,3) to[cisource,l={$0{,}05V_1$}] (6.5,0)--(0,0);
\node[ground] at(1.5,0){};\node[above,Vcol] at(n1){$V_1$};\node[above,Vcol] at(n2){$V_2$};
\end{circuitikz}}
\resposta{$V_1=\SI{17.5}{\volt}$; $V_2=\SI{7.5}{\volt}$; $i_s=\SI{1.25}{\ampere}$ de 1 para 2}
\resolucao{Com $i_s$ positivo de 1 para 2: $V_1/10+i_s=3$; $V_2/20+0{,}05V_1-i_s=0$; $V_1-V_2=10$. A solução é $(17{,}5,7{,}5,1{,}25)$.}
\end{questao}

\begin{questao}[4]{}
Determine $R$ para que $V_2=\SI{10}{\volt}$. A fonte dependente drena $0{,}05V_1$ do nó 2.
\circuitoq{\begin{circuitikz}[american,scale=.9,transform shape]
\draw (0,0) to[isource,l_={$\SI{5}{\ampere}$}] (0,3)--(1.5,3) coordinate(n1);
\draw (n1) to[R,l={$\SI{10}{\ohm}$}] (1.5,0);
\draw (n1) to[R,l={$\SI{5}{\ohm}$}] (4.5,3) coordinate(n2);
\draw (n2) to[R,l={$R$}] (4.5,0);
\draw (n2)--(6.3,3) to[cisource,l={$0{,}05V_1$}] (6.3,0)--(0,0);
\node[ground] at(1.5,0){};\node[above,Vcol] at(n1){$V_1$};\node[above,Vcol] at(n2){$V_2$};
\end{circuitikz}}
\resposta{$V_1=\dfrac{70}{3}\,\si{\volt}$; $R=\dfrac{20}{3}\,\ohm\approx\SI{6.67}{\ohm}$}
\resolucao{Com $V_2=10$ V, $V_1/10+(V_1-10)/5=5\Rightarrow V_1=70/3$ V. No nó 2: $10/R+(10-V_1)/5+0{,}05V_1=0$, de onde $10/R=1{,}5$ e $R=20/3\,\ohm$.}
\end{questao}

\begin{questao}[4]{}
Determine $V_1$, $V_2$, $V_3$ e a corrente $i_{23}$ no resistor de ponte, positiva de 2 para 3. A fonte dependente injeta $0{,}125V_2$ no nó 3.
\circuitoq{\begin{circuitikz}[american,scale=.82,transform shape]
\coordinate (g) at (0,-1.7);
\draw (g) to[isource,l_={$\SI{6}{\ampere}$}] (0,1.6) -- (1.2,1.6) coordinate(n1);
\draw (n1) to[R,l={$\SI{4}{\ohm}$}] (4.0,3.0) coordinate(n2);
\draw (n1) to[R,l_={$\SI{4}{\ohm}$}] (4.0,0.2) coordinate(n3);
\draw (n2) to[R,l={$\SI{4}{\ohm}$}] (n3);
\draw (n2) to[R,l={$\SI{8}{\ohm}$}] (6.2,3.0) -- (6.2,-1.7) -- (g);
\draw (n3) to[R,l_={$\SI{8}{\ohm}$}] (4.0,-1.7) -- (g);
\draw (7.8,-1.7) to[cisource,l_={$0{,}125V_2$}] (7.8,0.2) -- (n3);
\draw (7.8,-1.7) -- (g);
\node[ground] at(g){};
\node[above,Vcol] at(n1){$V_1$};\node[above,Vcol] at(n2){$V_2$};\node[right,Vcol] at(n3){$V_3$};
\end{circuitikz}}
\resposta{$V_1=\SI{57}{\volt}$; $V_2=\SI{42}{\volt}$; $V_3=\SI{48}{\volt}$; $i_{23}=\SI{-1.5}{\ampere}$}
\resolucao{As LKC são $(V_1-V_2)/4+(V_1-V_3)/4=6$, $V_2/8+(V_2-V_1)/4+(V_2-V_3)/4=0$ e $V_3/8+(V_3-V_1)/4+(V_3-V_2)/4=0{,}125V_2$. A solução é $(57,42,48)$ V. Assim $i_{23}=(42-48)/4=-1{,}5$ A, isto é, $1{,}5$ A de 3 para 2.}
\end{questao}


% END BANCO COMPLEMENTAR

\gabarito[Capítulo 7 --- Análise Nodal]
          % Analise nodal

% Cap. 8 -------------------------------------------------------------
\setcounter{section}{7}
% =====================================================================
%  CAPITULO 8 - TEOREMAS DE THEVENIN E NORTON
% =====================================================================
\section{Teoremas de Thévenin e Norton}

\begin{conceito}[Antes de começar]
\textbf{Thévenin:} qualquer rede linear de dois terminais equivale a uma fonte de
tensão $V_{\text{Th}}$ em série com uma resistência $R_{\text{Th}}$.\par
\textbf{Norton:} a mesma rede equivale a uma fonte de corrente $I_{\text{N}}$ em
paralelo com $R_{\text{N}}$.\par
\textbf{Como obter:}
\begin{itemize}[leftmargin=1.6em,itemsep=1pt,topsep=2pt]
\item $V_{\text{Th}}$ = tensão de circuito \emph{aberto} nos terminais (retire a
      carga).
\item $I_{\text{N}}$ = corrente de \emph{curto-circuito} entre os terminais.
\item $R_{\text{Th}} = R_{\text{N}}$ = resistência vista dos terminais com as fontes
      \emph{desligadas} (tensão $\to$ curto, corrente $\to$ aberto).
\end{itemize}
\textbf{Relação entre eles:}
$V_{\text{Th}} = R_{\text{Th}}\,I_{\text{N}}$ e
$R_{\text{Th}} = R_{\text{N}}$.
\end{conceito}

\legendaniveis

\blocofixacao

\begin{questao}[1]{}
A tensão de Thévenin $V_{\text{Th}}$ de uma rede é obtida medindo-se, nos
terminais de saída:
\begin{alternativas}
\item a tensão de curto-circuito.
\item a tensão de circuito aberto (sem carga).
\item a corrente de curto-circuito.
\item a resistência interna.
\item a potência máxima.
\end{alternativas}
\resposta{(b)}
\resolucao{$V_{\text{Th}}$ é a tensão nos terminais \emph{sem carga} (circuito
aberto). A corrente de curto-circuito fornece $I_{\text{N}}$, e a razão entre elas
dá $R_{\text{Th}}$.}
\end{questao}

\begin{questao}[1]{}
A resistência de Thévenin de uma rede é calculada:
\begin{alternativas}
\item mantendo todas as fontes ativas.
\item desligando as fontes (tensão em curto, corrente em aberto) e olhando dos
      terminais.
\item somando todas as resistências do circuito.
\item medindo a tensão de saída em vazio.
\item apenas para circuitos com uma única fonte.
\end{alternativas}
\resposta{(b)}
\resolucao{Para achar $R_{\text{Th}}$, desligam-se as fontes independentes (tensão
$\to$ curto, corrente $\to$ aberto) e calcula-se a resistência equivalente vista
dos terminais de saída.}
\end{questao}

\blocoaplicacao

\begin{questao}[2]{}
Para o circuito da figura, determine o equivalente de Thévenin
($V_{\text{Th}}$ e $R_{\text{Th}}$) visto pelos terminais A--B.

\circuitoq{%
\begin{circuitikz}[scale=1,transform shape,line width=0.8pt]
 \draw (0,0) to[battery1,l_={$\SI{12}{\volt}$}] (0,3)
       to[R,l={$R_1=\SI{4}{\ohm}$}] (3,3) coordinate(a);
 \draw (a) to[R,l={$R_2=\SI{4}{\ohm}$}] (3,0) coordinate(b) -- (0,0);
 \draw (a) to[short,-o] (4.5,3) node[right,font=\bfseries]{A};
 \draw (b) to[short,-o] (4.5,0) node[right,font=\bfseries]{B};
\end{circuitikz}}

\resposta{$V_{\text{Th}} = \SI{6}{\volt}$; $R_{\text{Th}} = \SI{2}{\ohm}$}
\resolucao{\textbf{$V_{\text{Th}}$} (saída em aberto): sem carga, $R_1$ e $R_2$
formam um divisor de tensão:
\[
V_{\text{Th}}=12\cdot\frac{R_2}{R_1+R_2}=12\cdot\frac{4}{8}=\SI{6}{\volt}.
\]
\textbf{$R_{\text{Th}}$} (fonte de \SI{12}{\volt} em curto): $R_1$ e $R_2$ ficam
em paralelo, vistos de A--B:
\[
R_{\text{Th}}=R_1\parallel R_2=\frac{4\cdot4}{8}=\SI{2}{\ohm}.
\]}
\end{questao}

\begin{questao}[2]{}
Usando o equivalente de Thévenin obtido na questão anterior, calcule a tensão e a
corrente em uma carga $R_L = \SI{6}{\ohm}$ ligada entre A e B.
\resposta{$I_L = \SI{0.75}{\ampere}$; $U_L = \SI{4.5}{\volt}$}
\resolucao{Com o circuito reduzido a $V_{\text{Th}} = \SI{6}{\volt}$ em série com
$R_{\text{Th}} = \SI{2}{\ohm}$ e a carga $R_L = \SI{6}{\ohm}$:
\[
I_L=\frac{V_{\text{Th}}}{R_{\text{Th}}+R_L}=\frac{6}{2+6}=\SI{0.75}{\ampere},
\qquad
U_L=R_L\,I_L=6\cdot0{,}75=\SI{4.5}{\volt}.
\]
\textbf{Vantagem do método:} uma vez achado o equivalente, testar qualquer carga
é imediato --- não é preciso reanalisar o circuito inteiro a cada valor de $R_L$.}
\end{questao}

\begin{questao}[2]{}
Converta o equivalente de Thévenin ($V_{\text{Th}} = \SI{6}{\volt}$,
$R_{\text{Th}} = \SI{2}{\ohm}$) das questões anteriores em seu equivalente de
Norton.
\resposta{$I_{\text{N}} = \SI{3}{\ampere}$; $R_{\text{N}} = \SI{2}{\ohm}$}
\resolucao{A resistência é a mesma, $R_{\text{N}} = R_{\text{Th}} = \SI{2}{\ohm}$.
A corrente de Norton é a de curto-circuito:
\[
I_{\text{N}}=\frac{V_{\text{Th}}}{R_{\text{Th}}}=\frac{6}{2}=\SI{3}{\ampere}.
\]
Os dois equivalentes são intercambiáveis: uma fonte de \SI{6}{\volt} em série com
\SI{2}{\ohm} entrega exatamente a mesma corrente e tensão a qualquer carga que uma
fonte de \SI{3}{\ampere} em paralelo com \SI{2}{\ohm}.}
\end{questao}

\begin{questao}[2]{}
Determine o equivalente de Norton (corrente $I_{\text{N}}$ e resistência
$R_{\text{N}}$) visto pelos terminais A--B do circuito da figura.

\circuitoq{%
\begin{circuitikz}[scale=1,transform shape,line width=0.8pt]
 \draw (0,0) to[battery1,l_={$\SI{24}{\volt}$}] (0,3)
       to[R,l={$R_1=\SI{6}{\ohm}$}] (3,3) coordinate(a);
 \draw (a) to[R,l={$R_2=\SI{3}{\ohm}$}] (3,0) coordinate(b) -- (0,0);
 \draw (a) to[short,-o] (4.3,3) node[right,font=\bfseries]{A};
 \draw (b) to[short,-o] (4.3,0) node[right,font=\bfseries]{B};
\end{circuitikz}}

\resposta{$I_{\text{N}} = \SI{4}{\ampere}$; $R_{\text{N}} = \SI{2}{\ohm}$}
\resolucao{\textbf{$R_{\text{N}}$} (fonte em curto, $R_1\parallel R_2$):
$R_{\text{N}} = \dfrac{6\cdot3}{9} = \SI{2}{\ohm}$.\\
\textbf{$I_{\text{N}}$} (corrente de curto entre A e B): curto-circuitando a saída,
$R_2$ fica em paralelo com o curto (não conduz), e toda a corrente passa por $R_1$:
\[
I_{\text{N}} = \frac{24}{R_1} = \frac{24}{6} = \SI{4}{\ampere}.
\]
\textbf{Conferindo pela relação com Thévenin:}
$V_{\text{Th}} = 24\cdot\dfrac{3}{9} = \SI{8}{\volt}$ e
$I_{\text{N}} = \dfrac{V_{\text{Th}}}{R_{\text{Th}}} = \dfrac{8}{2}
= \SI{4}{\ampere}$. Confere.}
\end{questao}

\blocoaprofundamento

\begin{questao}[3]{}
No circuito da figura, determine o equivalente de Thévenin visto por $R_L$ e, em
seguida, o valor de $R_L$ que recebe a máxima potência, indicando essa potência.

\circuitoq{%
\begin{circuitikz}[scale=1,transform shape,line width=0.8pt]
 \draw (0,0) to[battery1,l_={$\SI{24}{\volt}$}] (0,3)
       to[R,l={$R_1=\SI{6}{\ohm}$}] (3,3) coordinate(a);
 \draw (a) to[R,l={$R_2=\SI{3}{\ohm}$}] (3,0) coordinate(b) -- (0,0);
 \draw (a) to[short,-o] (4.3,3) node[right,font=\bfseries]{A};
 \draw (b) to[short,-o] (4.3,0) node[right,font=\bfseries]{B};
 \node[right,font=\footnotesize,cinza] at (4.5,1.5) {$R_L$};
\end{circuitikz}}

\resposta{$V_{\text{Th}}=\SI{8}{\volt}$, $R_{\text{Th}}=\SI{2}{\ohm}$;
$R_L=\SI{2}{\ohm}$, $P_{\max}=\SI{8}{\watt}$}
\resolucao{\textbf{$V_{\text{Th}}$} (divisor sobre $R_2$, sem carga):
\[
V_{\text{Th}}=24\cdot\frac{3}{6+3}=24\cdot\frac13=\SI{8}{\volt}.
\]
\textbf{$R_{\text{Th}}$} (fonte em curto, $R_1\parallel R_2$):
\[
R_{\text{Th}}=\frac{6\cdot3}{9}=\SI{2}{\ohm}.
\]
\textbf{Máxima transferência:} $R_L = R_{\text{Th}} = \SI{2}{\ohm}$ e
\[
P_{\max}=\frac{V_{\text{Th}}^{\,2}}{4R_{\text{Th}}}=\frac{8^2}{4\cdot2}
=\frac{64}{8}=\SI{8}{\watt}.
\]
Este exercício reúne os capítulos 5 e 8: reduzir a rede a Thévenin é justamente o
que permite aplicar a condição $R_L = R_{\text{Th}}$ de forma direta.}
\end{questao}

\begin{questao}[3]{}
Um circuito foi reduzido ao equivalente de Thévenin $V_{\text{Th}}=\SI{18}{\volt}$,
$R_{\text{Th}}=\SI{3}{\ohm}$. Deseja-se saber, sem refazer as contas para cada
caso, a corrente na carga para $R_L = \SI{3}{\ohm}$, $\SI{6}{\ohm}$ e
$\SI{9}{\ohm}$, e em qual delas a potência entregue é máxima.
\resposta{$I(3)=\SI{3}{\ampere}$, $I(6)=\SI{2}{\ampere}$, $I(9)=\SI{1.5}{\ampere}$;
$P$ máxima em $R_L=\SI{3}{\ohm}$}
\resolucao{A vantagem do equivalente de Thévenin é justamente testar cargas
rapidamente com $I_L = \dfrac{V_{\text{Th}}}{R_{\text{Th}}+R_L}$:
\[
I(3)=\frac{18}{6}=\SI{3}{\ampere},\quad
I(6)=\frac{18}{9}=\SI{2}{\ampere},\quad
I(9)=\frac{18}{12}=\SI{1.5}{\ampere}.
\]
As potências: $P(3)=3^2\cdot3=\SI{27}{\watt}$,
$P(6)=2^2\cdot6=\SI{24}{\watt}$, $P(9)=1{,}5^2\cdot9=\SI{20.25}{\watt}$.
A máxima ocorre em $R_L = R_{\text{Th}} = \SI{3}{\ohm}$, com \SI{27}{\watt}
--- coerente com o teorema da máxima transferência.}
\end{questao}

\begin{questao}[3]{}
Determine o equivalente de Thévenin visto pelos terminais A--B do circuito da
figura, que possui \textbf{duas fontes}.

\circuitoq{%
\begin{circuitikz}[scale=1,transform shape,line width=0.8pt]
 \draw (0,0) to[\fonteV,l={$\SI{12}{\volt}$}] (0,3) -- (1,3)
       to[R,l={$\SI{4}{\ohm}$}] (3.4,3) coordinate(n);
 \fill (n) circle (0.07);
 \draw (n) to[R,l={$\SI{6}{\ohm}$}] (3.4,0) coordinate(g) -- (0,0);
 \draw (n) -- (5.6,3) to[R,l={$\SI{12}{\ohm}$}] (7.8,3)
       to[\fonteV,l={$\SI{6}{\volt}$},invert] (7.8,0) -- (g);
 \draw (n) -- (3.4,4.3) -- (9.4,4.3) to[short,-o] (10.2,4.3)
       node[right,font=\bfseries]{A};
 \draw (g) -- (3.4,-1.2) -- (9.4,-1.2) to[short,-o] (10.2,-1.2)
       node[right,font=\bfseries]{B};
\end{circuitikz}}

\resposta{$\Vth = \SI{7}{\volt}$; $\Rth = \SI{2}{\ohm}$}
\resolucao{\textbf{$\Vth$ --- tensão de circuito aberto.} Sem carga em A--B,
aplica-se a lei dos nós no nó central (potencial $V$):
\[
\frac{12-V}{4}+\frac{6-V}{12}=\frac{V}{6}.
\]
Multiplicando por 12: $3(12-V)+(6-V)=2V \Rightarrow 36-3V+6-V=2V
\Rightarrow 42 = 6V \Rightarrow V = \SI{7}{\volt}$.
Logo $\Vth = \SI{7}{\volt}$.

\textbf{$\Rth$ --- resistência com as fontes desligadas.} Curto-circuitando
\emph{ambas} as fontes de tensão, os três resistores ficam todos ligados entre o
nó central e o terra, ou seja, em \textbf{paralelo}:
\[
\frac{1}{\Rth}=\frac{1}{4}+\frac{1}{12}+\frac{1}{6}
=\frac{3+1+2}{12}=\frac{6}{12}
\;\Longrightarrow\;
\Rth=\SI{2}{\ohm}.
\]
\textbf{Equivalente de Norton:} $\IN = \dfrac{7}{2}=\SI{3.5}{\ampere}$ em paralelo
com \SI{2}{\ohm}.\\
\textbf{Máxima transferência:} com $R_L = \SI{2}{\ohm}$,
$P_{\max}=\dfrac{7^2}{4\cdot2}=\SI{6.125}{\watt}$.\\
\textbf{Atenção ao erro comum:} ao calcular $\Rth$, é preciso curto-circuitar
\emph{todas} as fontes de tensão --- não apenas uma. Aqui, esquecer a fonte de
\SI{6}{\volt} levaria a $4\parallel6=\SI{2.4}{\ohm}$, resultado errado.}
\end{questao}

\begin{questao}[3]{}
Uma fonte real desconhecida (caixa-preta) é caracterizada experimentalmente com
duas medições: em \emph{vazio}, o voltímetro indica $\SI{20}{\volt}$; com uma
carga de $\SI{8}{\ohm}$, a tensão nos terminais cai para $\SI{16}{\volt}$.

\circuitoq{%
\begin{circuitikz}[scale=0.95,transform shape,line width=0.8pt]
 \draw[dashed,cinza,line width=0.9pt,rounded corners=6pt] (-0.4,-0.5) rectangle (2.6,3.5);
 \node[cinza,font=\footnotesize\sffamily] at (1.1,3.75) {caixa-preta};
 \node[cinza,font=\Large] at (1.1,1.5) {?};
 \draw (2.6,3) to[short,-o] (4.0,3) node[right,font=\bfseries]{A};
 \draw (2.6,0) to[short,-o] (4.0,0) node[right,font=\bfseries]{B};
 \draw (4.3,3) to[R,l={$R_L$}] (4.3,0);
 \draw (4.0,3)--(4.3,3); \draw (4.0,0)--(4.3,0);
\end{circuitikz}}

\begin{itensq}
\item Determine o equivalente de Thévenin da caixa-preta.
\item Qual a corrente de curto-circuito nos terminais?
\item Qual carga extrairia a máxima potência, e qual seria essa potência?
\item Um técnico sugere medir $R_{\text{Th}}$ diretamente com um ohmímetro nos
      terminais A--B. Por que isso é um erro grave?
\end{itensq}
\resposta{(a) $\Vth=\SI{20}{\volt}$, $\Rth=\SI{2}{\ohm}$;
(b) $\SI{10}{\ampere}$; (c) $R_L=\SI{2}{\ohm}$, $P_{\max}=\SI{50}{\watt}$;
(d) ver comentário}
\resolucao{\textbf{(a)} Em vazio não há corrente, logo não há queda interna:
$\Vth = \SI{20}{\volt}$.\\
Com a carga, a corrente é $I = \dfrac{16}{8} = \SI{2}{\ampere}$. A queda interna
é $20-16 = \SI{4}{\volt}$, portanto
\[
\Rth=\frac{\Delta V}{I}=\frac{4}{2}=\SI{2}{\ohm}.
\]
\textbf{(b)} $I_{\text{cc}} = \dfrac{\Vth}{\Rth} = \dfrac{20}{2}
= \SI{10}{\ampere}$ (esta é também a corrente de Norton, $\IN$).\\
\textbf{(c)} Máxima transferência em $R_L = \Rth = \SI{2}{\ohm}$:
\[
P_{\max}=\frac{\Vth^{\,2}}{4\Rth}=\frac{400}{8}=\SI{50}{\watt}.
\]
\textbf{(d) Por que medir com ohmímetro é errado --- e perigoso.}
\begin{itemize}[leftmargin=1.6em,itemsep=1pt,topsep=2pt]
\item O ohmímetro pressupõe o circuito \textbf{despassivado}, isto é, sem fontes
      ativas. A definição de $\Rth$ exige \emph{desligar} as fontes internas
      (tensão $\to$ curto, corrente $\to$ aberto). Como não temos acesso ao
      interior da caixa, isso é impossível de garantir.
\item Com a fonte interna \emph{ativa}, o ohmímetro --- que injeta sua própria
      corrente de teste --- lerá um valor sem significado, podendo inclusive
      indicar valores negativos ou saturar.
\item Além disso, aplicar um ohmímetro sobre uma fonte ativa pode \textbf{danificar
      o instrumento}, e curto-circuitar terminais de fontes de baixa impedância
      pode causar arco elétrico.
\end{itemize}
\textbf{Método correto (não invasivo):} exatamente o usado neste problema ---
duas medições de \emph{tensão} sob condições de carga diferentes, e
$\Rth = \dfrac{\Delta V}{\Delta I}$. Alternativamente, mede-se a tensão em vazio
e a corrente de curto (esta última apenas se a fonte suportar), obtendo
$\Rth = \Vth/I_{\text{cc}}$.}
\end{questao}

% BEGIN BANCO COMPLEMENTAR Teoremas de Thévenin e Norton
% =====================================================================
%  Banco complementar --- Teoremas de Thevenin e Norton  (REVISADO)
%  Substitui a versao gerada por template (11 clones em N2, 11 em N3 --
%  estes ultimos com o mesmo enunciado repetido e uma "Extensao
%  orientada" colada ao final -- e 10 questoes conceituais
%  indevidamente marcadas como N4).
%  Sem sobreposicao com o cap08, que ja traz: definicoes de Vth e Rth
%  (MC), Thevenin e Norton de circuitos com figura, conversao entre os
%  modelos, Thevenin visto por RL com maxima transferencia,
%  dimensionamento para corrente desejada, rede com duas fontes e
%  caracterizacao de caixa-preta por medida em VAZIO + carga.
%  Maxima transferencia de potencia tem banco proprio: aqui ela so
%  aparece como contraponto a uma restricao de projeto.
%  Distribuicao mantida: 8 N1 + 11 N2 + 11 N3 + 10 N4 = 40
% =====================================================================

\subsubsection*{Banco complementar --- Teoremas de Thévenin e Norton}

\begin{atencao}[Organização do banco]
Duas grandezas resumem qualquer rede linear vista por uma porta:
$V_{Th}$ (tensão de circuito aberto) e $R_{Th}$ (resistência com as fontes
\emph{independentes} zeradas). Norton é a mesma informação em outra forma,
com $I_N=V_{Th}/R_{Th}$. O N3 trata de fontes dependentes, caracterização
experimental, múltiplas portas e tolerância; o N4 exige demonstração,
projeto de ensaio, otimização sob restrição e interpretação de
$R_{Th}<0$.
\end{atencao}

\blocofixacao

\begin{questao}[1]{}
Uma rede linear apresenta $V_{oc}=\SI{6}{\volt}$ e $I_{sc}=\SI{1.5}{\ampere}$.
Determine seu equivalente de Thévenin.
\resposta{$V_{Th}=\SI{6}{\volt}$; $R_{Th}=\SI{4}{\ohm}$}
\resolucao{Por definição, $V_{Th}=V_{oc}=\SI{6}{\volt}$. E como o curto impõe
$V=0$, toda a tensão interna cai sobre $R_{Th}$:
\[
R_{Th}=\frac{V_{oc}}{I_{sc}}=\frac{6}{1{,}5}=\SI{4}{\ohm}.
\]
Este é o método mais geral --- funciona inclusive com fontes dependentes, em
que zerar fontes não é possível.}
\end{questao}

\begin{questao}[1]{}
Converta o equivalente de Thévenin $V_{Th}=\SI{20}{\volt}$,
$R_{Th}=\SI{5}{\ohm}$ em equivalente de Norton.
\resposta{$I_N=\SI{4}{\ampere}$; $R_N=\SI{5}{\ohm}$}
\resolucao{
\[
I_N=\frac{V_{Th}}{R_{Th}}=\frac{20}{5}=\SI{4}{\ampere};
\qquad
R_N=R_{Th}=\SI{5}{\ohm}.
\]
A resistência é a \emph{mesma} nos dois modelos --- ela descreve a rede
passiva, que não muda. O que muda é a forma de representar a fonte: em
Thévenin, tensão \emph{em série}; em Norton, corrente \emph{em paralelo}.}
\end{questao}

\begin{questao}[1]{}
Converta o equivalente de Norton $I_N=\SI{250}{\milli\ampere}$,
$R_N=\SI{80}{\ohm}$ em equivalente de Thévenin.
\resposta{$V_{Th}=\SI{20}{\volt}$; $R_{Th}=\SI{80}{\ohm}$}
\resolucao{
\[
V_{Th}=I_NR_N=0{,}250\cdot80=\SI{20}{\volt};
\qquad
R_{Th}=\SI{80}{\ohm}.
\]
Note que $V_{Th}$ é exatamente a tensão que a fonte de corrente produz sobre
$R_N$ quando a porta está \emph{aberta} --- toda a corrente é obrigada a passar
por $R_N$.}
\end{questao}

\begin{questao}[1]{}
Ao calcular $R_{Th}$ pelo método de zerar as fontes, as fontes
\textbf{dependentes} devem ser:
\begin{alternativas}[2]
\item zeradas junto com as independentes
\item mantidas ativas
\item substituídas por resistores
\item removidas do circuito
\end{alternativas}
\resposta{(b)}
\resolucao{Fontes dependentes fazem parte da \emph{rede}, não da excitação:
seu valor é determinado por variáveis internas do circuito. Zerá-las
descaracterizaria a rede.\\
\textbf{Consequência prática:} com fontes dependentes, zerar as independentes
deixa um circuito que ainda pode responder --- e $R_{Th}$ precisa ser obtida
aplicando uma \textbf{fonte de teste} na porta e medindo $V/I$, ou pela razão
$V_{oc}/I_{sc}$.}
\end{questao}

\begin{questao}[1]{}
Um equivalente com $V_{Th}=\SI{18}{\volt}$ e $R_{Th}=\SI{6}{\ohm}$ alimenta
$R_L=\SI{12}{\ohm}$. Determine a corrente.
\resposta{$I=\SI{1}{\ampere}$}
\resolucao{
\[
I=\frac{V_{Th}}{R_{Th}+R_L}=\frac{18}{6+12}=\SI{1}{\ampere}.
\]
O equivalente reduz qualquer rede, por mais complexa, a um circuito série de
dois elementos. É essa economia que torna o teorema valioso.}
\end{questao}

\begin{questao}[1]{}
No mesmo equivalente, determine a tensão sobre a carga de $\SI{12}{\ohm}$.
\resposta{$U_L=\SI{12}{\volt}$}
\resolucao{Por divisor de tensão:
\[
U_L=V_{Th}\frac{R_L}{R_{Th}+R_L}=18\cdot\frac{12}{18}=\SI{12}{\volt}.
\]
\textbf{Conferência:} $U_L=R_LI=12(1)=\SI{12}{\volt}$ \checkmark. A queda
interna é $18-12=\SI{6}{\volt}$, coerente com $R_{Th}I=6(1)$.}
\end{questao}

\begin{questao}[1]{}
Sobre o equivalente de Thévenin, é correto afirmar que ele reproduz:
\begin{alternativas}[2]
\item o comportamento da rede apenas nos terminais da porta
\item todas as tensões e correntes internas da rede
\item a rede somente para uma carga específica
\item a rede somente em circuito aberto
\end{alternativas}
\resposta{(a)}
\resolucao{O equivalente é \textbf{externo}: ele reproduz a relação
$U\times I$ na porta para \emph{qualquer} carga --- o que descarta (c) e (d).
Mas ele descarta completamente o interior: a potência dissipada em $R_{Th}$
não corresponde à potência real dissipada dentro da rede, e nenhuma corrente
interna é recuperável a partir dele.\\
\textbf{Regra:} use o equivalente para analisar a carga; volte ao circuito
original se precisar de qualquer grandeza interna.}
\end{questao}

\begin{questao}[1]{}
Uma rede consiste em uma fonte ideal de $\SI{30}{\volt}$ em série com
$\SI{10}{\ohm}$, tendo $\SI{40}{\ohm}$ em paralelo com a porta. Determine
$R_{Th}$.
\resposta{$R_{Th}=\SI{8}{\ohm}$}
\resolucao{Zerando a fonte ideal de tensão (substituída por \textbf{curto}), a
porta enxerga $\SI{10}{\ohm}$ em paralelo com $\SI{40}{\ohm}$:
\[
R_{Th}=\frac{10\cdot40}{50}=\SI{8}{\ohm}.
\]
\textbf{Lembre das duas regras de zerar:} fonte de \emph{tensão} ideal vira
curto; fonte de \emph{corrente} ideal vira circuito aberto. Trocá-las é o erro
mais frequente do método.}
\end{questao}

\blocoaplicacao

\begin{questao}[2]{}
Uma fonte de $\SI{24}{\volt}$ alimenta $R_1=\SI{6}{\kilo\ohm}$ em série com
$R_2=\SI{3}{\kilo\ohm}$. Determine o equivalente de Thévenin visto sobre
$R_2$.
\resposta{$V_{Th}=\SI{8}{\volt}$; $R_{Th}=\SI{2}{\kilo\ohm}$}
\resolucao{\textbf{Tensão:} com a porta aberta, o divisor entrega
\[
V_{Th}=24\cdot\frac{3}{6+3}=\SI{8}{\volt}.
\]
\textbf{Resistência:} zerando a fonte (curto), $R_1$ e $R_2$ ficam em paralelo:
\[
R_{Th}=\frac{6\cdot3}{9}=\SI{2}{\kilo\ohm}.
\]
\textbf{Interpretação:} um divisor resistivo \emph{não} é fonte de tensão --- ele
tem resistência interna $R_1\parallel R_2$, e é ela que provoca a queda quando
uma carga é conectada.}
\end{questao}

\begin{questao}[2]{}
Converta o equivalente da questão anterior em Norton e verifique por
$I_{sc}$ calculada diretamente no circuito original.
\resposta{$I_N=\SI{4}{\milli\ampere}$; $R_N=\SI{2}{\kilo\ohm}$}
\resolucao{\textbf{Pela conversão:}
$I_N=\dfrac{8}{2000}=\SI{4}{\milli\ampere}$.\\
\textbf{Verificação direta:} curto-circuitando a porta, $R_2$ fica em paralelo
com o curto e não conduz; toda a corrente vem da fonte por $R_1$:
\[
I_{sc}=\frac{24}{6000}=\SI{4}{\milli\ampere}.\ \checkmark
\]
A concordância confirma o equivalente. Calcular $I_{sc}$ de forma independente
é o teste de consistência mais barato do método.}
\end{questao}

\begin{questao}[2]{}
Uma fonte de corrente de $\SI{3}{\milli\ampere}$ alimenta $R_1=\SI{4}{\kilo\ohm}$
ao terra, e a porta é acessada através de $R_2=\SI{2}{\kilo\ohm}$ em série.
Determine o equivalente de Thévenin.
\resposta{$V_{Th}=\SI{12}{\volt}$; $R_{Th}=\SI{6}{\kilo\ohm}$}
\resolucao{\textbf{Tensão:} com a porta aberta, nenhuma corrente passa por
$R_2$ --- logo não há queda nele, e $V_{Th}$ é a tensão sobre $R_1$:
\[
V_{Th}=(\pot{3}{-3})(4000)=\SI{12}{\volt}.
\]
\textbf{Resistência:} zerando a fonte de corrente (circuito \textbf{aberto}),
sobra $R_2$ em série com $R_1$:
\[
R_{Th}=2+4=\SI{6}{\kilo\ohm}.
\]
\textbf{Atenção:} zerar uma fonte de corrente \emph{abre} o ramo --- e por isso
$R_1$ e $R_2$ ficam em série, não em paralelo.}
\end{questao}

\begin{questao}[2]{}
Duas fontes aterradas alimentam a mesma porta: $\SI{20}{\volt}$ através de
$\SI{5}{\kilo\ohm}$ e $\SI{10}{\volt}$ através de $\SI{5}{\kilo\ohm}$.
Determine o equivalente de Thévenin.
\resposta{$V_{Th}=\SI{15}{\volt}$; $R_{Th}=\SI{2.5}{\kilo\ohm}$}
\resolucao{\textbf{Tensão (Millman):}
\[
V_{Th}=\frac{20/5+10/5}{1/5+1/5}=\frac{4+2}{0{,}4}=\SI{15}{\volt}.
\]
\textbf{Resistência:} zerando ambas as fontes, os dois resistores ficam em
paralelo:
\[
R_{Th}=\frac{5}{2}=\SI{2.5}{\kilo\ohm}.
\]
\textbf{Observação:} $V_{Th}$ é a média \emph{ponderada} das duas fontes, com
pesos dados pelas condutâncias. Como aqui os resistores são iguais, saiu a média
aritmética simples.}
\end{questao}

\begin{questao}[2]{}
Com $V_{Th}=\SI{8}{\volt}$ e $R_{Th}=\SI{2}{\kilo\ohm}$, complete a tabela de
tensão e corrente para $R_L=\SI{2}{\kilo\ohm}$, $\SI{6}{\kilo\ohm}$ e
$\SI{18}{\kilo\ohm}$.
\resposta{$(\SI{4}{\volt};\SI{2}{\milli\ampere})$,
$(\SI{6}{\volt};\SI{1}{\milli\ampere})$,
$(\SI{7.2}{\volt};\SI{0.4}{\milli\ampere})$}
\resolucao{Com $I=\dfrac{8}{2+R_L}$ (em $\si{\milli\ampere}$, com $R_L$ em
$\si{\kilo\ohm}$) e $U=R_LI$:
\begin{center}
\begin{tabular}{ccc}
\hline
$R_L$ & $I$ & $U$\\
\hline
$\SI{2}{\kilo\ohm}$ & $\SI{2}{\milli\ampere}$ & $\SI{4}{\volt}$\\
$\SI{6}{\kilo\ohm}$ & $\SI{1}{\milli\ampere}$ & $\SI{6}{\volt}$\\
$\SI{18}{\kilo\ohm}$ & $\SI{0.4}{\milli\ampere}$ & $\SI{7.2}{\volt}$\\
\hline
\end{tabular}
\end{center}
\textbf{Leitura:} conforme $R_L$ cresce, $U\to V_{Th}$ e $I\to0$. Os três
pontos são \emph{colineares} no plano $U\times I$ --- eles pertencem à reta
$U=8-2I$, que é a assinatura do equivalente. Essa reta é a "reta de carga da
fonte", retomada no bloco de desafio.}
\end{questao}

\begin{questao}[2]{}
Uma rede tem $V_{oc}=\SI{45}{\volt}$ e, sob carga de $\SI{20}{\ohm}$, entrega
$\SI{36}{\volt}$. Determine $R_{Th}$ e $I_{sc}$.
\resposta{$R_{Th}=\SI{5}{\ohm}$; $I_{sc}=\SI{9}{\ampere}$}
\resolucao{A corrente na carga é $I=36/20=\SI{1.8}{\ampere}$, e a queda interna
é $45-36=\SI{9}{\volt}$:
\[
R_{Th}=\frac{9}{1{,}8}=\SI{5}{\ohm};
\qquad
I_{sc}=\frac{45}{5}=\SI{9}{\ampere}.
\]
\textbf{Observação:} $I_{sc}$ foi obtida \emph{sem} curto-circuitar a rede ---
o que é essencial quando o curto seria destrutivo. O bloco de desafio
generaliza esse raciocínio.}
\end{questao}

\begin{questao}[2]{}
Determine o equivalente de Norton de uma fonte de $\SI{3}{\ampere}$ em paralelo
com $\SI{12}{\ohm}$, tendo ainda $\SI{6}{\ohm}$ em paralelo com a porta.
\resposta{$I_N=\SI{3}{\ampere}$; $R_N=\SI{4}{\ohm}$}
\resolucao{\textbf{Resistência:} zerando a fonte de corrente (aberto), sobram
os dois resistores em paralelo:
\[
R_N=\frac{12\cdot6}{18}=\SI{4}{\ohm}.
\]
\textbf{Corrente:} curto-circuitando a porta, os dois resistores ficam em
paralelo com o curto e não conduzem --- toda a corrente da fonte vai ao curto:
\[
I_N=\SI{3}{\ampere}.
\]
\textbf{Conferência:} $V_{Th}=I_NR_N=3(4)=\SI{12}{\volt}$, que também é a
tensão da associação $12\parallel6$ percorrida por $\SI{3}{\ampere}$
\checkmark.}
\end{questao}

\begin{questao}[2]{}
Um equivalente com $V_{Th}=\SI{12}{\volt}$ e $R_{Th}=\SI{4}{\ohm}$ alimenta
$R_L=\SI{8}{\ohm}$. Determine a potência na carga e a dissipada internamente.
\resposta{$P_L=\SI{8}{\watt}$; $P_{int}=\SI{4}{\watt}$}
\resolucao{$I=\dfrac{12}{12}=\SI{1}{\ampere}$, logo
\[
P_L=R_LI^{2}=8(1)=\SI{8}{\watt};
\qquad
P_{int}=R_{Th}I^{2}=4(1)=\SI{4}{\watt}.
\]
Total: $\SI{12}{\watt}=V_{Th}I$ \checkmark.\\
\textbf{Ressalva importante:} os $\SI{4}{\watt}$ \emph{não} são necessariamente
a potência realmente dissipada dentro da rede original --- $R_{Th}$ é um
artifício de modelagem externa. Para a potência interna verdadeira, é preciso
voltar ao circuito completo.}
\end{questao}

\begin{questao}[2]{}
Uma fonte de $\SI{12}{\volt}$ alimenta $R_1=\SI{6}{\kilo\ohm}$ até o nó $A$; de
$A$ saem $R_2=\SI{3}{\kilo\ohm}$ ao terra e $R_3=\SI{6}{\kilo\ohm}$ até o nó
$B$. Determine o equivalente de Thévenin visto de $B$.
\resposta{$V_{Th}=\SI{4}{\volt}$; $R_{Th}=\SI{8}{\kilo\ohm}$}
\resolucao{\textbf{Tensão:} com $B$ aberto, $R_3$ não conduz e não cai tensão
nele. Logo $V_B=V_A$, dado pelo divisor:
\[
V_{Th}=12\cdot\frac{3}{6+3}=\SI{4}{\volt}.
\]
\textbf{Resistência:} zerando a fonte, de $B$ enxerga-se $R_3$ em série com
$R_1\parallel R_2$:
\[
R_{Th}=6+\frac{6\cdot3}{9}=6+2=\SI{8}{\kilo\ohm}.
\]
Compare com a porta em $A$ ($R_{Th}=\SI{2}{\kilo\ohm}$): \textbf{mesma tensão,
resistência quatro vezes maior}. A porta escolhida altera o equivalente.}
\end{questao}

\begin{questao}[2]{}
Uma rede tem $I_N=\SI{5}{\ampere}$ e $R_N=\SI{3}{\ohm}$. Determine a tensão na
porta em aberto e a corrente em uma carga de $\SI{2}{\ohm}$.
\resposta{$V_{oc}=\SI{15}{\volt}$; $I_L=\SI{3}{\ampere}$}
\resolucao{\textbf{Em aberto:} toda a corrente passa por $R_N$:
$V_{oc}=5(3)=\SI{15}{\volt}$.\\
\textbf{Com carga:} pelo divisor de corrente,
\[
I_L=I_N\frac{R_N}{R_N+R_L}=5\cdot\frac{3}{5}=\SI{3}{\ampere}.
\]
\textbf{Conferência por Thévenin:} $V_{Th}=\SI{15}{\volt}$,
$I_L=15/(3+2)=\SI{3}{\ampere}$ \checkmark. Os dois modelos são intercambiáveis
--- escolha o que simplifica a conta.}
\end{questao}

\begin{questao}[2]{}
Explique por que o equivalente de Thévenin \textbf{não} pode ser usado para
calcular o rendimento interno de uma fonte real, e ilustre com o caso
$V_{Th}=\SI{10}{\volt}$, $R_{Th}=\SI{2}{\ohm}$, $R_L=\SI{8}{\ohm}$.
\resposta{$R_{Th}$ agrega efeitos que não correspondem a dissipação real
localizada}
\resolucao{No exemplo, $I=\SI{1}{\ampere}$,
$P_L=\SI{8}{\watt}$ e $P_{R_{Th}}=\SI{2}{\watt}$ --- aparentemente $80\%$ de
rendimento.\\
\textbf{Por que isso pode ser falso:} $R_{Th}$ resulta da \emph{combinação} de
todos os resistores da rede com as fontes zeradas. Se a rede original for, por
exemplo, um divisor, boa parte da potência é consumida por um caminho que
conduz corrente \emph{mesmo em circuito aberto} --- corrente que o modelo de
Thévenin simplesmente não representa.\\
\textbf{Caso extremo:} um divisor $6\,\mathrm{k}/3\,\mathrm{k}$ sob
$\SI{24}{\volt}$ consome $\SI{64}{\milli\watt}$ da fonte \emph{sem nenhuma
carga}, enquanto seu equivalente de Thévenin ($\SI{8}{\volt}$,
$\SI{2}{\kilo\ohm}$) prevê consumo nulo em aberto.\\
\textbf{Conclusão:} Thévenin é exato para tudo o que ocorre \emph{na porta} e
inútil para balanços energéticos internos.}
\end{questao}

\blocoaprofundamento

\begin{questao}[3]{}
Uma rede contém apenas uma fonte \textbf{dependente}: um nó $A$ ligado ao terra
por $R=\SI{1}{\kilo\ohm}$, com uma fonte de corrente controlada por tensão
injetando $g\,V_A$ no próprio nó, com $g=\SI{0.75}{\milli\siemens}$.
\begin{itensq}
\item Explique por que $V_{oc}$ e $I_{sc}$ são ambos nulos.
\item Determine $R_{Th}$ pelo método da fonte de teste.
\end{itensq}
\resposta{(a) sem fonte independente, não há excitação; (b)
$R_{Th}=\SI{4}{\kilo\ohm}$}
\resolucao{(a) Fontes dependentes não excitam nada por si só --- seu valor é
proporcional a uma variável do circuito. Sem fonte independente, a única
solução é a trivial: todas as tensões e correntes nulas. Logo
$V_{oc}=0$ e $I_{sc}=0$, e a razão $V_{oc}/I_{sc}$ é indeterminada
($0/0$).\\
(b) Aplique uma \textbf{fonte de teste} de $I_t=\SI{1}{\milli\ampere}$ na porta
e escreva a LKC no nó $A$:
\[
\frac{V_A}{1000}-\pot{0.75}{-3}V_A=\pot{1}{-3}
\;\Longrightarrow\;
V_A\left(\pot{1}{-3}-\pot{0.75}{-3}\right)=\pot{1}{-3},
\]
\[
V_A=\frac{\pot{1}{-3}}{\pot{0.25}{-3}}=\SI{4}{\volt}
\;\Longrightarrow\;
R_{Th}=\frac{V_A}{I_t}=\SI{4}{\kilo\ohm}.
\]
\textbf{Leitura:} a fonte dependente cancela $75\%$ da condutância do resistor,
quadruplicando a resistência vista. Redes com elementos ativos podem apresentar
$R_{Th}$ muito diferente de qualquer resistor físico presente nelas.}
\end{questao}

\begin{questao}[3]{}
Agora a rede tem \emph{também} uma fonte independente: $\SI{12}{\volt}$ através
de $R_1=\SI{2}{\kilo\ohm}$ até o nó $A$, com $R_2=\SI{2}{\kilo\ohm}$ de $A$ ao
terra e a mesma fonte dependente injetando $gV_A$ em $A$, com
$g=\SI{0.5}{\milli\siemens}$.
\begin{itensq}
\item Determine $V_{oc}$.
\item Determine $I_{sc}$.
\item Obtenha $R_{Th}$ e confirme pelo método da fonte de teste.
\end{itensq}
\resposta{(a) $\SI{12}{\volt}$; (b) $\SI{6}{\milli\ampere}$;
(c) $R_{Th}=\SI{2}{\kilo\ohm}$}
\resolucao{(a) Porta aberta, LKC em $A$:
\[
\frac{V_A-12}{2000}+\frac{V_A}{2000}-\pot{0.5}{-3}V_A=0
\;\Longrightarrow\;
V_A(\pot{1}{-3}-\pot{0.5}{-3})=\pot{6}{-3},
\]
\[
V_A=\frac{\pot{6}{-3}}{\pot{0.5}{-3}}=\SI{12}{\volt}=V_{oc}.
\]
(b) Com a porta em curto, $V_A=0$: a fonte dependente entrega
$gV_A=0$ e $R_2$ não conduz. Resta a corrente da fonte independente:
\[
I_{sc}=\frac{12}{2000}=\SI{6}{\milli\ampere}.
\]
(c) $R_{Th}=\dfrac{12}{\pot{6}{-3}}=\SI{2}{\kilo\ohm}$.\\
\textbf{Confirmação:} zerando só a fonte independente e aplicando teste,
$G=\pot{1}{-3}-\pot{0.5}{-3}=\pot{0.5}{-3}\,\si{\siemens}$, logo
$R_{Th}=\SI{2}{\kilo\ohm}$ \checkmark.\\
\textbf{Nota:} $V_{oc}=\SI{12}{\volt}$ é igual à fonte --- a dependente
compensou exatamente a queda em $R_1$. Coincidência numérica, mas ilustra que
elementos ativos podem "anular" resistências.}
\end{questao}

\begin{questao}[3]{}
Uma rede é caracterizada com \textbf{duas cargas} conhecidas:
$R_{L1}=\SI{2}{\kilo\ohm}$ dá $\SI{6}{\volt}$ e
$R_{L2}=\SI{6}{\kilo\ohm}$ dá $\SI{9}{\volt}$.
\begin{itensq}
\item Determine $V_{Th}$ e $R_{Th}$.
\item Explique a vantagem sobre a medição em vazio.
\end{itensq}
\resposta{(a) $V_{Th}=\SI{12}{\volt}$, $R_{Th}=\SI{2}{\kilo\ohm}$}
\resolucao{(a) De $U=V_{Th}\dfrac{R_L}{R_L+R_{Th}}$, escreva
$V_{Th}=U\dfrac{R_L+R_{Th}}{R_L}$ para as duas medidas e iguale:
\[
R_{Th}=\frac{R_{L1}R_{L2}(U_2-U_1)}{U_1R_{L2}-U_2R_{L1}}
=\frac{(2)(6)(3)}{6(6)-9(2)}=\frac{36}{18}=\SI{2}{\kilo\ohm},
\]
(valores em $\si{\kilo\ohm}$ e $\si{\volt}$), e então
\[
V_{Th}=6\cdot\frac{2+2}{2}=\SI{12}{\volt}.
\]
\textbf{Verificação com a segunda medida:}
$12\cdot6/8=\SI{9}{\volt}$ \checkmark.\\
(b) A medição em \emph{vazio} exige voltímetro de impedância muito maior que
$R_{Th}$; se a rede for de alta impedância, o próprio instrumento carrega o
circuito e a leitura não é $V_{Th}$. Além disso, algumas fontes não podem
operar sem carga mínima.\\
O método das duas cargas \textbf{extrapola} para $R_L\to\infty$ sem precisar
chegar lá --- e é o mesmo raciocínio usado para obter $\EMF$ e $r$ de uma
bateria por dois pontos de carga.}
\end{questao}

\begin{questao}[3]{}
Para a rede da questão anterior, avalie a sensibilidade do método: suponha que
a segunda leitura tenha sido $\SI{9.1}{\volt}$ em vez de $\SI{9}{\volt}$
($+1{,}1\%$).
\begin{itensq}
\item Recalcule $R_{Th}$ e $V_{Th}$.
\item Comente o resultado.
\end{itensq}
\resposta{$R_{Th}\approx\SI{2.09}{\kilo\ohm}$ ($+4{,}5\%$),
$V_{Th}\approx\SI{12.27}{\volt}$ ($+2{,}3\%$)}
\resolucao{(a) Com $U_2=\SI{9.1}{\volt}$, o denominador passa a
$U_1R_{L2}-U_2R_{L1}=6(6)-9{,}1(2)=36-18{,}2=17{,}8$, e
\[
R_{Th}=\frac{(2)(6)(3{,}1)}{17{,}8}=\frac{37{,}2}{17{,}8}
=\SI{2.09}{\kilo\ohm},
\qquad
V_{Th}=6\cdot\frac{2+2{,}09}{2}=\SI{12.27}{\volt}.
\]
(b) Um erro de $1{,}1\%$ na tensão produziu $4{,}5\%$ em $R_{Th}$ e
$2{,}3\%$ em $V_{Th}$ --- \textbf{amplificação de cerca de quatro vezes}.\\
\textbf{Origem do problema:} o denominador $U_1R_{L2}-U_2R_{L1}$ é uma
\emph{diferença de valores próximos} ($36$ contra $18$, mas que se aproximam
conforme as cargas se parecem), situação numericamente mal condicionada.\\
\textbf{Correção de método:} escolha cargas \textbf{bem separadas} (por
exemplo, $R_{L2}\ge10R_{L1}$), o que afasta os dois termos e reduz a
amplificação. E, sempre que possível, use três ou mais pontos com regressão
linear em vez de dois pontos isolados.}
\end{questao}

\begin{questao}[3]{}
Uma mesma rede é observada por duas portas: em $A$, obtém-se
$V_{Th}=\SI{4}{\volt}$ e $R_{Th}=\SI{2}{\kilo\ohm}$; em $B$,
$V_{Th}=\SI{4}{\volt}$ e $R_{Th}=\SI{8}{\kilo\ohm}$.
\begin{itensq}
\item Determine a corrente de curto em cada porta.
\item Determine, em cada porta, a tensão sobre uma carga de
      $\SI{4}{\kilo\ohm}$.
\item Interprete.
\end{itensq}
\resposta{(a) $\SI{2}{\milli\ampere}$ e $\SI{0.5}{\milli\ampere}$;
(b) $\SI{2.67}{\volt}$ e $\SI{1.33}{\volt}$}
\resolucao{(a) $I_{sc}=V_{Th}/R_{Th}$:
$4/2\mathrm{k}=\SI{2}{\milli\ampere}$ e
$4/8\mathrm{k}=\SI{0.5}{\milli\ampere}$.\\
(b) $U=4\cdot\dfrac{4}{4+R_{Th}}$:
\[
\text{porta }A:\ 4\cdot\frac{4}{6}=\SI{2.67}{\volt};
\qquad
\text{porta }B:\ 4\cdot\frac{4}{12}=\SI{1.33}{\volt}.
\]
(c) \textbf{Tensão em vazio idêntica, comportamento sob carga completamente
diferente.} A porta $A$ é "forte" (baixa impedância) e sustenta melhor a
tensão; a porta $B$ é "fraca" e colapsa.\\
\textbf{Consequência prática:} especificar apenas a tensão de saída de um
circuito é insuficiente --- é a resistência de saída que determina se ele serve
para a carga pretendida. Duas fontes de $\SI{5}{\volt}$ podem ser
completamente diferentes.}
\end{questao}

\begin{questao}[3]{}
No divisor $R_1=\SI{6}{\kilo\ohm}$ e $R_2=\SI{3}{\kilo\ohm}$ sob
$\SI{12}{\volt}$, ambos com tolerância de $\pm5\%$:
\begin{itensq}
\item Determine a faixa de $V_{Th}$.
\item Determine a faixa de $R_{Th}$.
\item Explique por que as duas se comportam de modo diferente.
\end{itensq}
\resposta{(a) $\SI{4}{\volt}$, de $3{,}74$ a $\SI{4.27}{\volt}$
($\approx\pm6{,}8\%$); (b) $\SI{2}{\kilo\ohm}$, de $1{,}90$ a
$\SI{2.10}{\kilo\ohm}$ ($\pm5\%$ exato)}
\resolucao{(a) $V_{Th}=12\dfrac{R_2}{R_1+R_2}$ é máxima com $R_2$ no máximo e
$R_1$ no mínimo:
\[
V_{Th}^{\max}=12\cdot\frac{3{,}15}{5{,}70+3{,}15}=\SI{4.271}{\volt};
\qquad
V_{Th}^{\min}=12\cdot\frac{2{,}85}{6{,}30+2{,}85}=\SI{3.738}{\volt}.
\]
Desvios de $+6{,}8\%$ e $-6{,}6\%$ --- \textbf{maiores que a tolerância dos
componentes}.\\
(b) $R_{Th}=R_1\parallel R_2$ é \textbf{homogênea de grau 1}: multiplicar ambos
por $1{,}05$ multiplica o resultado por $1{,}05$. Extremos:
$5{,}70\parallel2{,}85=\SI{1.90}{\kilo\ohm}$ e
$6{,}30\parallel3{,}15=\SI{2.10}{\kilo\ohm}$ --- exatamente $\pm5\%$.\\
(c) \textbf{A diferença está na estrutura.} $R_{Th}$ escala com os resistores;
$V_{Th}$ depende de uma \emph{razão}, e razões pioram quando numerador e
denominador desviam em sentidos opostos. Formalmente,
\[
\frac{\Delta V_{Th}}{V_{Th}}=\frac{R_1}{R_1+R_2}
\left(\frac{\Delta R_2}{R_2}-\frac{\Delta R_1}{R_1}\right),
\]
e com $R_1/(R_1+R_2)=2/3$ e desvios opostos de $5\%$, obtém-se
$\frac23(10\%)=6{,}7\%$ \checkmark.\\
\textbf{Regra:} para $V_{Th}$ preciso, use resistores \emph{casados} (mesmo
lote, mesma rede integrada) --- desvios correlacionados se cancelam na razão.}
\end{questao}

\begin{questao}[3]{}
Uma fonte de laboratório não pode ser curto-circuitada. Duas medições foram
feitas: com $R_{L1}=\SI{5.5}{\ohm}$ obteve-se $\SI{10}{\ampere}$; com
$R_{L2}=\SI{2.5}{\ohm}$, $\SI{20}{\ampere}$.
\begin{itensq}
\item Determine $V_{Th}$ e $R_{Th}$.
\item Determine o equivalente de Norton.
\item Comente a validade da extrapolação.
\end{itensq}
\resposta{(a) $\SI{60}{\volt}$ e $\SI{0.5}{\ohm}$;
(b) $I_N=\SI{120}{\ampere}$}
\resolucao{(a) As tensões medidas são $U_1=5{,}5(10)=\SI{55}{\volt}$ e
$U_2=2{,}5(20)=\SI{50}{\volt}$. Como $U=V_{Th}-R_{Th}I$:
\[
R_{Th}=-\frac{\Delta U}{\Delta I}=-\frac{50-55}{20-10}=\SI{0.5}{\ohm};
\qquad
V_{Th}=55+0{,}5(10)=\SI{60}{\volt}.
\]
(b) $I_N=V_{Th}/R_{Th}=60/0{,}5=\SI{120}{\ampere}$, com
$R_N=\SI{0.5}{\ohm}$.\\
(c) \textbf{Obtivemos $\SI{120}{\ampere}$ sem jamais aplicar um curto} --- o
que teria sido destrutivo. Mas a extrapolação pressupõe \textbf{linearidade até
o curto}, hipótese frágil: fontes reais entram em limitação de corrente,
aquecem (com $R_{Th}$ crescendo) e podem desarmar proteções.\\
\textbf{Uso correto do número:} $\SI{120}{\ampere}$ serve para
\emph{dimensionar proteção} e avaliar risco --- é um limite superior. Não deve
ser lido como previsão do que o instrumento realmente entregaria.}
\end{questao}

\begin{questao}[3]{}
Um equivalente $V_{Th}=\SI{5}{\volt}$, $R_{Th}=\SI{200}{\ohm}$ alimenta um LED
cuja queda direta é praticamente constante em $\SI{2}{\volt}$.
\begin{itensq}
\item Determine graficamente o ponto de operação (reta de carga).
\item Calcule a corrente.
\item Explique por que o método funciona com dispositivo não linear.
\end{itensq}
\resposta{$I=\SI{15}{\milli\ampere}$ em $\SI{2}{\volt}$}
\resolucao{(a) A rede impõe $U=5-200I$ --- uma \textbf{reta} no plano
$U\times I$, com intercepto $\SI{5}{\volt}$ e corrente de curto
$\SI{25}{\milli\ampere}$. O LED impõe sua própria curva, aqui aproximada pela
vertical $U=\SI{2}{\volt}$. O ponto de operação é a \textbf{interseção}.\\
(b)
\[
2=5-200I \;\Longrightarrow\; I=\frac{3}{200}=\SI{15}{\milli\ampere}.
\]
(c) \textbf{Thévenin exige linearidade apenas da rede, não da carga.} A rede
foi substituída por sua reta característica; a carga pode ter qualquer curva
$U\times I$. A solução é a interseção das duas características --- e existirá e
será única sempre que a curva da carga for monotônica.\\
\textbf{Poder do método:} sem ele, seria preciso resolver simultaneamente a
rede inteira e a equação não linear do dispositivo. Com ele, sobra uma
interseção de duas curvas --- técnica padrão para diodos, transistores e
células fotovoltaicas.}
\end{questao}

\begin{questao}[3]{}
Um equivalente tem $V_{Th}=\SI{24}{\volt}$ e $R_{Th}=\SI{8}{\ohm}$. A carga
deve receber \textbf{no mínimo} $\SI{20}{\volt}$.
\begin{itensq}
\item Determine a faixa admissível de $R_L$.
\item Determine a potência entregue no limite.
\item Compare com a potência máxima teórica.
\end{itensq}
\resposta{(a) $R_L\ge\SI{40}{\ohm}$; (b) $\SI{10}{\watt}$;
(c) $\SI{18}{\watt}$, porém a $\SI{12}{\volt}$}
\resolucao{(a)
\[
24\frac{R_L}{R_L+8}\ge20
\;\Longrightarrow\;
24R_L\ge20R_L+160
\;\Longrightarrow\;
R_L\ge\SI{40}{\ohm}.
\]
(b) Em $R_L=\SI{40}{\ohm}$: $U=\SI{20}{\volt}$,
$I=\SI{0.5}{\ampere}$, $P=\SI{10}{\watt}$.\\
(c) A potência máxima ocorreria em $R_L=R_{Th}=\SI{8}{\ohm}$, valendo
\[
P_{\max}=\frac{V_{Th}^{2}}{4R_{Th}}=\frac{576}{32}=\SI{18}{\watt},
\]
mas nessa condição $U=\SI{12}{\volt}$ --- \textbf{metade} de $V_{Th}$, violando
a restrição por larga margem.\\
\textbf{Conflito de objetivos:} maximizar potência e manter tensão são metas
incompatíveis em uma fonte de $R_{Th}$ fixa. A restrição de tensão vence
sempre que a carga for um circuito eletrônico, que simplesmente não funciona
abaixo de sua tensão mínima. A saída de engenharia é \emph{reduzir}
$R_{Th}$, não escolher outro $R_L$.}
\end{questao}

\begin{questao}[3]{}
Uma rede deve ser analisada com \textbf{oito} cargas diferentes.
\begin{itensq}
\item Compare o esforço de resolver a rede completa oito vezes com o de obter o
      equivalente uma vez.
\item Quantifique, supondo que a rede exija um sistema nodal $4\times4$.
\item Indique a limitação dessa estratégia.
\end{itensq}
\resposta{Duas soluções da rede contra oito; ganho cresce com o número de
cargas}
\resolucao{(a) \textbf{Sem Thévenin:} para cada carga, monta-se e resolve-se a
rede inteira --- oito sistemas.\\
\textbf{Com Thévenin:} resolve-se a rede \emph{duas} vezes (uma para
$V_{oc}$, outra para $R_{Th}$) e, em seguida, oito divisões triviais.\\
(b) Um sistema $4\times4$ por eliminação exige da ordem de $n^{3}/3\approx21$
operações de ponto flutuante estruturais. Oito resoluções custam
$\approx170$; duas custam $\approx43$, seguidas de oito contas de uma linha.
\textbf{Redução de cerca de $75\%$} --- e o ganho cresce com o número de cargas
e com o tamanho da rede.\\
(c) \textbf{Limitação:} o equivalente vale para \emph{uma} porta. Se as cargas
forem conectadas em portas diferentes, ou se houver duas cargas
simultaneamente, é preciso um equivalente por porta --- e, no caso de portas
que interagem, um modelo de dois acessos (matriz de resistências), não um
Thévenin isolado.\\
\textbf{E vale lembrar:} nenhuma grandeza interna é recuperável pelo
equivalente. Se o projeto exigir a corrente em um ramo interno para cada carga,
o ganho desaparece.}
\end{questao}

\begin{questao}[3]{}
Valide o equivalente do divisor ($\SI{24}{\volt}$, $\SI{6}{\kilo\ohm}$,
$\SI{3}{\kilo\ohm}$) com carga $R_L=\SI{2}{\kilo\ohm}$ por \textbf{três}
caminhos independentes: Thévenin, Norton e análise nodal.
\resposta{$U_L=\SI{4}{\volt}$ pelos três métodos}
\resolucao{\textbf{Thévenin} ($\SI{8}{\volt}$, $\SI{2}{\kilo\ohm}$):
\[
U_L=8\cdot\frac{2}{2+2}=\SI{4}{\volt}.
\]
\textbf{Norton} ($\SI{4}{\milli\ampere}$, $\SI{2}{\kilo\ohm}$): a corrente se
divide entre $R_N$ e $R_L$, iguais, logo $\SI{2}{\milli\ampere}$ em cada:
\[
U_L=(\pot{2}{-3})(2000)=\SI{4}{\volt}.
\]
\textbf{Nodal} (nó $A$, com $R_2$ e $R_L$ ao terra):
\[
\frac{V_A-24}{6}+\frac{V_A}{3}+\frac{V_A}{2}=0
\;\Longrightarrow\;
V_A(1+2+3)=24 \;\Longrightarrow\; V_A=\SI{4}{\volt}.
\]
\textbf{Valor da validação cruzada:} os três métodos usam caminhos algébricos
distintos; concordância tripla torna erro aritmético muito improvável.
Discordância localiza o problema --- se Thévenin e Norton concordam entre si mas
divergem do nodal, o erro está no equivalente; se divergem entre si, o erro está
na conversão.}
\end{questao}

\blocodesafio

\begin{questao}[4]{}
\textbf{(Demonstração do teorema.)} Prove o teorema de Thévenin usando o
princípio da superposição.
\begin{itensq}
\item Conecte uma fonte de corrente $I$ à porta e aplique superposição.
\item Identifique cada parcela e conclua a forma da relação $U\times I$.
\item Aponte exatamente onde a hipótese de linearidade é usada.
\end{itensq}
\resposta{$U=V_{oc}-R_{Th}I$, com $V_{oc}$ da parcela interna e $R_{Th}$ da
parcela externa}
\resolucao{(a) Seja $\mathcal{N}$ uma rede linear com uma porta. Conecte a ela
uma fonte de corrente ideal $I$ e chame de $U$ a tensão resultante na porta. O
circuito tem agora dois grupos de excitação: as fontes \emph{internas} de
$\mathcal{N}$ e a fonte de teste $I$. Por superposição,
\[
U=U'+U'',
\]
onde $U'$ é a resposta às fontes internas com $I=0$, e $U''$ é a resposta a $I$
com as fontes internas zeradas.\\
(b) \textbf{Parcela $U'$:} com $I=0$, a fonte de corrente é um circuito aberto
--- a porta está em vazio. Logo $U'=V_{oc}$, por definição.\\
\textbf{Parcela $U''$:} com as fontes internas zeradas, $\mathcal{N}$ reduz-se a
uma rede puramente passiva, que apresenta na porta uma resistência única; chame-a
$R_{Th}$. A fonte $I$ injeta corrente \emph{para dentro} da rede, produzindo
$U''=-R_{Th}I$.\\
Somando:
\[
\boxed{U=V_{oc}-R_{Th}\,I}
\]
--- que é exatamente a relação de uma fonte $V_{oc}$ em série com $R_{Th}$.
Como a relação vale para \emph{todo} $I$, os dois circuitos são
indistinguíveis do lado de fora.\\
(c) \textbf{A linearidade é usada em três pontos.} \emph{(i)} A superposição só
vale em redes lineares --- é o passo central. \emph{(ii)} A rede passiva
resultante só pode ser caracterizada por um \emph{único} número $R_{Th}$ porque
sua relação $U\times I$ é proporcional. \emph{(iii)} A soma $U'+U''$ pressupõe
que as parcelas não interagem.\\
\textbf{Consequência:} a demonstração explica por que fontes \emph{dependentes}
não invalidam o teorema (elas são lineares e permanecem na rede passiva,
afetando $R_{Th}$), enquanto um diodo o invalida (não é linear, e nenhuma das
três etapas se sustenta).}
\end{questao}

\begin{questao}[4]{}
\textbf{(Demonstração --- equivalência dos modelos.)} Prove que os equivalentes
de Thévenin e Norton descrevem exatamente a mesma porta.
\begin{itensq}
\item Escreva a relação $U\times I$ de cada modelo.
\item Mostre que elas coincidem se e somente se $I_N=V_{Th}/R_{Th}$ e
      $R_N=R_{Th}$.
\item Identifique o caso em que a conversão é impossível.
\end{itensq}
\resposta{Ambos produzem a reta $U=V_{Th}-R_{Th}I$; a conversão falha para
fontes ideais ($R_{Th}=0$ ou $R_N\to\infty$)}
\resolucao{(a) \textbf{Thévenin:} $U=V_{Th}-R_{Th}I$.\\
\textbf{Norton:} a corrente da fonte se reparte entre $R_N$ e a saída, ou seja
$I=I_N-U/R_N$, isto é
\[
U=R_N\left(I_N-I\right)=I_NR_N-R_NI.
\]
(b) Duas retas coincidem quando têm o mesmo intercepto e a mesma inclinação:
\[
I_NR_N=V_{Th}
\quad\text{e}\quad
R_N=R_{Th}
\;\Longrightarrow\;
I_N=\frac{V_{Th}}{R_{Th}}.
\]
Como o teorema de Thévenin garante que \emph{qualquer} rede linear vista por uma
porta produz uma relação $U\times I$ afim, e como toda reta afim de inclinação
negativa admite as duas representações, os modelos são intercambiáveis.\\
(c) \textbf{Fontes ideais.} Se $R_{Th}=0$ (fonte de tensão ideal), a reta é
vertical --- $I_N=V_{Th}/0$ diverge e não existe Norton equivalente. Se
$R_N\to\infty$ (fonte de corrente ideal), a reta é horizontal e não existe
Thévenin equivalente.\\
\textbf{Interpretação:} os dois modelos são equivalentes para toda rede
\emph{real} (com $0<R_{Th}<\infty$). A impossibilidade nos extremos não é falha
do teorema --- é a constatação de que uma fonte ideal de tensão simplesmente não
tem análogo em corrente.}
\end{questao}

\begin{questao}[4]{}
\textbf{(Método experimental generalizado.)} Deduza a caracterização de uma
rede linear a partir de $n\ge2$ pares $(I_k,U_k)$ obtidos com cargas
diferentes.
\begin{itensq}
\item Escreva o modelo e identifique a forma de reta.
\item Obtenha $V_{Th}$ e $R_{Th}$ por mínimos quadrados.
\item Explique como os resíduos revelam falha de linearidade.
\end{itensq}
\resposta{$U=V_{Th}-R_{Th}I$: intercepto e coeficiente angular de uma
regressão linear}
\resolucao{(a) O modelo é $U=V_{Th}-R_{Th}I$: uma reta com intercepto
$V_{Th}$ e coeficiente angular $-R_{Th}$. Cada medição sob carga fornece um
ponto.\\
(b) Por mínimos quadrados, com médias $\bar I$ e $\bar U$:
\[
-R_{Th}=\frac{\sum(I_k-\bar I)(U_k-\bar U)}{\sum(I_k-\bar I)^{2}},
\qquad
V_{Th}=\bar U+R_{Th}\bar I.
\]
Com dois pontos, a regressão degenera na fórmula de dois pontos --- que é o
caso particular, não o método geral.\\
(c) \textbf{Diagnóstico pelos resíduos.} Se a rede for linear, os resíduos
$U_k-\hat U_k$ se distribuem \emph{aleatoriamente} em torno de zero. Se
apresentarem \textbf{padrão sistemático} --- por exemplo, curvatura, com
resíduos negativos nas extremidades e positivos no centro --- a rede não é
linear na faixa ensaiada.\\
\textbf{Causas típicas:} aquecimento (resistência crescendo com a corrente),
entrada em limitação de corrente da fonte, ou saturação de um elemento
magnético.\\
\textbf{Vantagem decisiva sobre dois pontos:} com dois pontos \emph{sempre}
passa uma reta, e a não linearidade fica invisível. É preciso um terceiro ponto,
no mínimo, para que o modelo possa ser \emph{testado} --- e não apenas
ajustado.}
\end{questao}

\begin{questao}[4]{}
\textbf{(Projeto de ensaio seguro.)} Deve-se obter o equivalente de Norton de
uma fonte de $\SI{60}{\volt}$ cuja corrente de curto é estimada em
$\SI{120}{\ampere}$, usando instrumentos que suportam no máximo
$\SI{25}{\ampere}$.
\begin{itensq}
\item Determine as cargas de ensaio e justifique a escolha.
\item Determine a potência dissipada em cada carga.
\item Especifique um procedimento com proteção.
\end{itensq}
\resposta{(a) $\SI{5.5}{\ohm}$ e $\SI{2.5}{\ohm}$ ($10$ e
$\SI{20}{\ampere}$); (b) $\SI{550}{\watt}$ e $\SI{1000}{\watt}$}
\resolucao{(a) Exige-se $I\le\SI{25}{\ampere}$, ou seja
$R_{Th}+R_L\ge60/25=\SI{2.4}{\ohm}$. Com $R_{Th}\approx\SI{0.5}{\ohm}$
estimado, escolhem-se
\[
R_{L1}=\SI{5.5}{\ohm}\ (\SI{10}{\ampere})
\qquad\text{e}\qquad
R_{L2}=\SI{2.5}{\ohm}\ (\SI{20}{\ampere}).
\]
\textbf{Critério de separação:} as duas correntes devem diferir bastante ---
aqui, fator $2$ --- para que o denominador da regressão não seja uma diferença
de valores próximos, problema de mau condicionamento identificado no bloco
anterior.\\
(b) $P_1=5{,}5(10)^{2}=\SI{550}{\watt}$;
$P_2=2{,}5(20)^{2}=\SI{1000}{\watt}$. São cargas de potência --- resistores de
fio ventilados ou banco de água, jamais componentes de bancada comuns.\\
(c) \textbf{Procedimento.}
\begin{itemize}
\item Fusível ou disjuntor de $\SI{25}{\ampere}$ em série, sempre.
\item Começar pela carga \emph{maior} ($\SI{5.5}{\ohm}$), de menor corrente.
\item Medir $U$ e $I$ \emph{simultaneamente}, com o voltímetro nos terminais da
      carga (para não incluir a queda dos cabos de ensaio).
\item Leituras rápidas, antes que o aquecimento altere as resistências.
\item Registrar um terceiro ponto para verificar linearidade.
\end{itemize}
\textbf{Resultado:} $I_N=\SI{120}{\ampere}$ obtido sem que qualquer instrumento
veja mais que $\SI{20}{\ampere}$. A extrapolação substitui o ensaio
destrutivo --- este é o valor prático do teorema.}
\end{questao}

\begin{questao}[4]{}
\textbf{(Múltiplas portas.)} Uma rede com duas portas $A$ e $B$ tem, do lado
externo, comportamento descrito por
\[
\begin{bmatrix}U_A\\U_B\end{bmatrix}
=\begin{bmatrix}V_{A0}\\V_{B0}\end{bmatrix}
-\begin{bmatrix}R_{AA}&R_{AB}\\R_{BA}&R_{BB}\end{bmatrix}
\begin{bmatrix}I_A\\I_B\end{bmatrix}.
\]
\begin{itensq}
\item Explique por que um Thévenin isolado por porta é insuficiente.
\item Determine em que condição os dois equivalentes independentes bastam.
\item Relacione $R_{AB}$ e $R_{BA}$ com a reciprocidade.
\end{itensq}
\resposta{Basta um Thévenin por porta apenas se $R_{AB}=R_{BA}=0$ (portas
desacopladas)}
\resolucao{(a) Um equivalente de Thévenin isolado descreve
$U_A=V_{A0}-R_{AA}I_A$, ignorando o termo $R_{AB}I_B$. Se uma carga for
conectada em $B$, ela drena $I_B\ne0$ e \emph{altera} a tensão em $A$ --- efeito
que o modelo isolado não prevê. Caracterizar $A$ com $B$ aberto e depois usar
esse equivalente com $B$ carregado produz erro.\\
(b) Os equivalentes independentes bastam se, e somente se,
$R_{AB}=R_{BA}=0$: nenhuma corrente em uma porta afeta a outra. Isso ocorre
quando não há caminho resistivo comum --- por exemplo, portas alimentadas por
fontes separadas e sem terra compartilhado.\\
\textbf{Na prática}, um critério útil é $R_{AB}\ll R_{AA}$: o acoplamento existe,
mas é desprezível para a precisão desejada.\\
(c) Em rede \textbf{passiva}, a reciprocidade garante $R_{AB}=R_{BA}$ --- a
matriz é simétrica, como já visto na matriz de condutâncias nodal. Isso reduz de
quatro para três os parâmetros a determinar.\\
Com fontes \textbf{dependentes}, a simetria se perde e $R_{AB}\ne R_{BA}$: a
rede transmite preferencialmente em um sentido. Essa assimetria é exatamente o
que caracteriza um \emph{amplificador} --- e é medida como isolação reversa.}
\end{questao}

\begin{questao}[4]{}
\textbf{(Propagação de tolerância.)} Para o divisor $R_1$, $R_2$ sob $V$, com
tolerâncias relativas $t_1$ e $t_2$:
\begin{itensq}
\item Deduza a expressão da incerteza relativa de $V_{Th}$.
\item Deduza a de $R_{Th}$ e explique por que ela nunca excede
      $\max(t_1,t_2)$.
\item Avalie para $R_1=\SI{6}{\kilo\ohm}$, $R_2=\SI{3}{\kilo\ohm}$ e
      $t_1=t_2=5\%$, e proponha como reduzir a incerteza de $V_{Th}$.
\end{itensq}
\resposta{(a) $\dfrac{R_1}{R_1+R_2}(t_1+t_2)$ no pior caso;
(b) média ponderada das tolerâncias; (c) $6{,}7\%$ e $5\%$}
\resolucao{(a) De $V_{Th}=V\dfrac{R_2}{R_1+R_2}$, tomando o logaritmo e
diferenciando:
\[
\frac{\Delta V_{Th}}{V_{Th}}
=\frac{\Delta R_2}{R_2}-\frac{\Delta R_2+\Delta R_1}{R_1+R_2}
=\frac{R_1}{R_1+R_2}\left(\frac{\Delta R_2}{R_2}-\frac{\Delta R_1}{R_1}\right).
\]
No pior caso (desvios opostos), o módulo vale
$\dfrac{R_1}{R_1+R_2}(t_1+t_2)$.\\
(b) $R_{Th}=R_1\parallel R_2$ é homogênea de grau 1, e vale a combinação com
$\sum S_k=1$:
\[
t_{R_{Th}}=\frac{G_1t_1+G_2t_2}{G_1+G_2},
\]
uma \textbf{média ponderada} --- que, por construção, fica entre
$\min(t_1,t_2)$ e $\max(t_1,t_2)$ e nunca os ultrapassa.\\
(c) Com $R_1/(R_1+R_2)=2/3$:
$t_{V_{Th}}=\frac23(10\%)=6{,}7\%$, enquanto
$t_{R_{Th}}=5\%$.\\
\textbf{Como reduzir a incerteza de $V_{Th}$.}
\begin{itemize}
\item \textbf{Resistores casados:} em rede resistiva integrada, os desvios são
      \emph{correlacionados} ($\Delta R_1/R_1\approx\Delta R_2/R_2$) e o termo
      entre parênteses tende a zero --- a razão fica muito mais precisa que os
      valores absolutos. É o princípio que torna viáveis conversores
      analógico-digitais de alta resolução.
\item \textbf{Reduzir $R_1/(R_1+R_2)$}, isto é, trabalhar com razões de divisão
      próximas de $1$ --- nem sempre possível.
\item \textbf{Selecionar por medição} apenas o resistor crítico.
\end{itemize}
\textbf{Ponto central:} $V_{Th}$ depende de uma \emph{razão} e pode ser pior que
os componentes; $R_{Th}$ é homogênea e nunca é pior. Estruturas diferentes
propagam incerteza de modos diferentes.}
\end{questao}

\begin{questao}[4]{}
\textbf{(Resistência negativa.)} Uma rede ativa apresenta
$R_{Th}=\SI{-4}{\kilo\ohm}$.
\begin{itensq}
\item Interprete fisicamente o sinal.
\item Determine a condição sobre $R_L$ para que o circuito tenha solução
      estável.
\item Indique uma aplicação prática.
\end{itensq}
\resposta{(b) $R_L<\SI{4}{\kilo\ohm}$ em módulo, para que a resistência total
seja positiva}
\resolucao{(a) $R_{Th}<0$ significa que a tensão na porta \emph{aumenta} quando
a corrente drenada aumenta --- a rede \textbf{fornece} energia ao circuito
externo em resposta à perturbação, em vez de se opor a ela. Só é possível com
elementos ativos: a fonte dependente entrega mais corrente do que o resistor
consome, invertendo o sinal da condutância líquida.\\
(b) Com carga $R_L$, a resistência total da malha é $R_{Th}+R_L$. Para que a
solução seja estável, ela deve ser \textbf{positiva}:
\[
R_L+(-4\,\mathrm{k})>0
\;\Longrightarrow\;
R_L>\SI{4}{\kilo\ohm}.
\]
Com $R_L<\SI{4}{\kilo\ohm}$, qualquer perturbação cresce em vez de decair: o
circuito satura em um extremo ou oscila. E em $R_L=\SI{4}{\kilo\ohm}$
exatamente, a matriz do sistema torna-se singular --- não há solução finita.\\
\emph{(Observação: o critério depende da topologia; em configuração paralela o
sentido da desigualdade se inverte. O princípio é sempre o mesmo: a resistência
líquida vista pela malha deve ser positiva.)}\\
(c) \textbf{Aplicações.} A resistência negativa cancela perdas: em osciladores,
ela compensa exatamente a resistência do tanque LC e sustenta oscilação senoidal
permanente. Diodos túnel e diodos Gunn exibem trechos de resistência negativa em
suas curvas e são usados como osciladores de micro-ondas.\\
\textbf{Conexão com o diagnóstico:} um determinante que se anula em rede com
elementos ativos é o mesmo fenômeno visto na análise nodal --- não é problema
numérico, é aviso de instabilidade.}
\end{questao}

\begin{questao}[4]{}
\textbf{(Otimização sob duas restrições.)} Uma fonte com
$V_{Th}=\SI{24}{\volt}$ e $R_{Th}=\SI{8}{\ohm}$ deve alimentar uma carga que
exige, simultaneamente, $U\ge\SI{20}{\volt}$ e $P\ge\SI{12}{\watt}$.
\begin{itensq}
\item Mostre que as duas exigências são incompatíveis.
\item Determine o menor $R_{Th}$ que as tornaria compatíveis.
\item Proponha uma solução de engenharia.
\end{itensq}
\resposta{(a) $U\ge20$ exige $R_L\ge\SI{40}{\ohm}$, mas então
$P\le\SI{10}{\watt}$; (b) $R_{Th}\le\SI{6.67}{\ohm}$}
\resolucao{(a) A restrição de tensão dá $R_L\ge\SI{40}{\ohm}$. Nessa região, a
potência é
\[
P=\frac{U^{2}}{R_L}\le\frac{24^{2}R_L}{(R_L+8)^{2}},
\]
função \emph{decrescente} para $R_L>R_{Th}$. Seu valor máximo na região
admissível ocorre em $R_L=\SI{40}{\ohm}$, valendo
$P=20^{2}/40=\SI{10}{\watt}<\SI{12}{\watt}$. \textbf{Incompatível.}\\
(b) Impondo as duas condições no ponto crítico $U=\SI{20}{\volt}$ e
$P=\SI{12}{\watt}$: $I=P/U=\SI{0.6}{\ampere}$ e
$R_L=20/0{,}6=\SI{33.3}{\ohm}$. Da malha,
\[
R_{Th}=\frac{V_{Th}-U}{I}=\frac{24-20}{0{,}6}=\SI{6.67}{\ohm}.
\]
Qualquer $R_{Th}\le\SI{6.67}{\ohm}$ torna o par de requisitos atendível.\\
(c) \textbf{Soluções, em ordem de preferência.}
\begin{itemize}
\item \textbf{Reduzir $R_{Th}$} --- fio de maior seção, conexões melhores,
      fonte mais robusta. Ataca a causa.
\item \textbf{Elevar $V_{Th}$} --- com $\SI{28}{\volt}$, os
      $\SI{8}{\ohm}$ passam a ser aceitáveis.
\item \textbf{Regulador entre fonte e carga} --- desacopla completamente as
      duas restrições, ao custo de rendimento e complexidade.
\end{itemize}
\textbf{Lição de método:} quando duas especificações se mostram
incompatíveis, o cálculo não apenas revela o impasse --- ele quantifica
\emph{exatamente} o quanto um parâmetro precisa melhorar. Aqui, reduzir
$R_{Th}$ em $17\%$ resolve o problema.}
\end{questao}

\begin{questao}[4]{}
\textbf{(Reta de carga e ponto de operação.)} Uma fonte com
$V_{Th}=\SI{5}{\volt}$ e $R_{Th}=\SI{200}{\ohm}$ alimenta um dispositivo cuja
característica é $I=\SI{10}{\milli\ampere}\cdot\left(e^{U/0{,}5}-1\right)$, com
$U$ em volts.
\begin{itensq}
\item Descreva o método gráfico e obtenha o ponto de operação por iteração.
\item Determine a resistência \emph{dinâmica} do dispositivo nesse ponto.
\item Explique a diferença entre resistência estática e dinâmica.
\end{itensq}
\resposta{$U\approx\SI{0.90}{\volt}$, $I\approx\SI{20.5}{\milli\ampere}$;
$r_d\approx\SI{24}{\ohm}$}
\resolucao{(a) A reta de carga é $I=(5-U)/200$, indo de
$(0;\ \SI{25}{\milli\ampere})$ a $(\SI{5}{\volt};\ 0)$. O ponto de operação é a
interseção com a exponencial.\\
\emph{Iteração:} com $U=\SI{0.9}{\volt}$, o dispositivo pede
$10(e^{1{,}8}-1)=\SI{50.5}{\milli\ampere}$, enquanto a reta oferece
$\SI{20.5}{\milli\ampere}$ --- é preciso baixar $U$. Com
$U=\SI{0.80}{\volt}$: dispositivo $\SI{39.6}{\milli\ampere}$, reta
$\SI{21}{\milli\ampere}$. Com $U=\SI{0.70}{\volt}$: dispositivo
$\SI{30.2}{\milli\ampere}$, reta $\SI{21.5}{\milli\ampere}$. Com
$U=\SI{0.60}{\volt}$: dispositivo $\SI{23.2}{\milli\ampere}$, reta
$\SI{22}{\milli\ampere}$. Convergência próxima de
$U\approx\SI{0.59}{\volt}$, $I\approx\SI{22}{\milli\ampere}$.\\
(b) A resistência dinâmica é a inclinação local:
\[
r_d=\left(\frac{\mathrm{d}I}{\mathrm{d}U}\right)^{-1}
=\frac{0{,}5}{I+\SI{10}{\milli\ampere}}
\approx\frac{0{,}5}{\pot{32}{-3}}\approx\SI{16}{\ohm}.
\]
(c) \textbf{Estática:} $R=U/I\approx0{,}59/0{,}022\approx\SI{27}{\ohm}$ --- a
razão entre os valores no ponto.\\
\textbf{Dinâmica:} $r_d\approx\SI{16}{\ohm}$ --- a resposta a uma
\emph{pequena variação} em torno do ponto.\\
São grandezas distintas e, em dispositivos não lineares, sempre diferentes. A
estática governa o balanço de potência CC; a \textbf{dinâmica} governa a resposta
a sinais pequenos, e é ela que entra em qualquer modelo de pequenos sinais ---
razão pela qual o mesmo transistor tem "resistências" diferentes conforme a
pergunta.}
\end{questao}

\begin{questao}[4]{}
\textbf{(Limites do teorema.)} Analise se Thévenin se aplica em cada caso e
justifique.
\begin{itensq}
\item Rede com resistor dependente da temperatura, aquecendo com a corrente.
\item Rede contendo um diodo em série com a saída.
\item Rede linear observada em torno de um ponto de operação fixo.
\end{itensq}
\resposta{(a) não; (b) não; (c) sim, como modelo incremental}
\resolucao{(a) \textbf{Não se aplica.} A resistência depende da corrente, o que
torna a relação $U\times I$ não linear --- a característica deixa de ser uma
reta e passa a curvar-se. Dois pontos de medição dariam um "equivalente" que
falharia em qualquer terceiro ponto.\\
\emph{Ressalva útil:} se a variação térmica for pequena na faixa de interesse, o
modelo linear é uma boa aproximação local --- e os \emph{resíduos} da regressão
denunciariam o desvio.\\
(b) \textbf{Não se aplica à rede como um todo.} O diodo é não linear e, pior,
\emph{assimétrico}: a rede conduz em um sentido e bloqueia no outro. Nenhum par
$(V_{Th},R_{Th})$ reproduz esse comportamento.\\
\emph{Uso correto:} aplique Thévenin à parte \textbf{linear} da rede e trate o
diodo como carga externa, resolvendo por reta de carga. É exatamente o
procedimento das questões anteriores.\\
(c) \textbf{Aplica-se, na forma incremental.} Em torno de um ponto de operação,
qualquer rede suave se comporta linearmente para pequenas variações. Obtém-se
um equivalente de Thévenin \emph{de pequenos sinais}, com $R_{Th}$ construída a
partir das resistências \textbf{dinâmicas} --- não das estáticas.\\
\textbf{Síntese:} o teorema exige linearidade, mas linearidade
\emph{local} basta. É essa observação que sustenta toda a análise de pequenos
sinais em eletrônica: circuitos essencialmente não lineares são analisados com
as ferramentas lineares deste capítulo, desde que o sinal seja pequeno o
bastante.}
\end{questao}

% END BANCO COMPLEMENTAR

\gabarito[Capítulo 8 --- Thévenin e Norton]
        % Thevenin e Norton

% Cap. 9 -------------------------------------------------------------
\setcounter{section}{8}
% =====================================================================
%  CAPITULO 9 - METODO DE MAXWELL (ANALISE DE MALHAS)
% =====================================================================
\section{Método de Maxwell (Análise de Malhas)}

\begin{conceito}[Antes de começar]
O método de Maxwell (ou análise de malhas) usa a \textbf{lei das malhas} de
Kirchhoff, adotando uma \emph{corrente de malha} que circula em cada laço
fechado. Roteiro:
\begin{enumerate}[leftmargin=1.7em,itemsep=1pt,topsep=2pt]
\item Adote um sentido (geralmente horário) para a corrente de cada malha
      ($I_1, I_2, \dots$).
\item Em cada malha, some as quedas de tensão: um resistor \emph{compartilhado}
      por duas malhas é percorrido pela \emph{diferença} das correntes.
\item Iguale à soma das fem da malha (positivas se a corrente sai do $+$).
\item Resolva o sistema. A corrente real em um resistor compartilhado é a soma
      algébrica das correntes de malha que passam por ele.
\end{enumerate}
\end{conceito}

\legendaniveis

\blocofixacao

\begin{questao}[1]{}
No método das correntes de malha, um resistor compartilhado por duas malhas
adjacentes é percorrido por uma corrente igual a:
\begin{alternativas}
\item soma das duas correntes de malha, sempre.
\item diferença das duas correntes de malha.
\item média das duas correntes de malha.
\item apenas a corrente da primeira malha.
\item zero.
\end{alternativas}
\resposta{(b)}
\resolucao{As duas correntes de malha atravessam o resistor comum em sentidos, em
geral, opostos: a corrente real é a \emph{diferença} algébrica $I_1 - I_2$ (com o
sinal dependendo dos sentidos adotados).}
\end{questao}

\blocoaplicacao

\begin{questao}[2]{}
Aplique o método de Maxwell ao circuito da figura e determine as correntes de
malha $I_1$ e $I_2$.

\circuitoq{%
\begin{circuitikz}[scale=1,transform shape,line width=0.8pt]
 % malha 1 (esquerda)
 \draw (0,0) to[battery1,l_={$\SI{20}{\volt}$}] (0,3)
       to[R,l={$R_1=\SI{4}{\ohm}$}] (3,3) coordinate(top);
 \draw (top) to[R,l={$R_2=\SI{8}{\ohm}$}] (3,0) coordinate(bot) -- (0,0);
 % malha 2 (direita)
 \draw (top) to[R,l={$R_3=\SI{4}{\ohm}$}] (6,3) -- (6,0) -- (bot);
 \draw[->,color=loopcol,line width=1.4pt] (1.4,1.5) arc (0:310:0.5);
 \node[loopcol] at (1.4,1.5) {\footnotesize$I_1$};
 \draw[->,color=loopcol,line width=1.4pt] (4.6,1.5) arc (0:310:0.5);
 \node[loopcol] at (4.6,1.5) {\footnotesize$I_2$};
\end{circuitikz}}

\resposta{$I_1 = \SI{3}{\ampere}$; $I_2 = \SI{2}{\ampere}$}
\resolucao{Correntes de malha no sentido horário. O resistor $R_2$ (comum) é
percorrido por $I_1 - I_2$.\\
\textbf{Malha 1:} $20 = R_1 I_1 + R_2(I_1 - I_2) = 4I_1 + 8(I_1 - I_2)$, ou seja
$12I_1 - 8I_2 = 20$.\\
\textbf{Malha 2:} $0 = R_3 I_2 + R_2(I_2 - I_1) = 4I_2 + 8(I_2 - I_1)$, ou seja
$-8I_1 + 12I_2 = 0 \Rightarrow I_1 = \dfrac{3}{2}I_2$.\\
Substituindo: $12\cdot\dfrac{3}{2}I_2 - 8I_2 = 20 \Rightarrow 10I_2 = 20
\Rightarrow I_2 = \SI{2}{\ampere}$ e $I_1 = \SI{3}{\ampere}$.\\
\textbf{Corrente real em $R_2$:} $I_1 - I_2 = \SI{1}{\ampere}$ (de cima para
baixo).}
\end{questao}

\begin{questao}[2]{}
No circuito de duas malhas da figura, com duas fontes, determine $I_1$ e $I_2$.

\circuitoq{%
\begin{circuitikz}[scale=1,transform shape,line width=0.8pt]
 \draw (0,0) to[battery1,l_={$E_1=\SI{18}{\volt}$}] (0,3)
       to[R,l={$R_1=\SI{2}{\ohm}$}] (3,3) coordinate(top);
 \draw (top) to[R,l={$R_2=\SI{6}{\ohm}$}] (3,0) coordinate(bot) -- (0,0);
 \draw (top) to[R,l={$R_3=\SI{6}{\ohm}$}] (6,3)
       to[battery1,l={$E_2=\SI{6}{\volt}$},invert] (6,0) -- (bot);
 \draw[->,color=loopcol,line width=1.4pt] (1.4,1.5) arc (0:310:0.5);
 \node[loopcol] at (1.4,1.5) {\footnotesize$I_1$};
 \draw[->,color=loopcol,line width=1.4pt] (4.6,1.5) arc (0:310:0.5);
 \node[loopcol] at (4.6,1.5) {\footnotesize$I_2$};
\end{circuitikz}}

\resposta{$I_1 = \SI{3}{\ampere}$; $I_2 = \SI{1}{\ampere}$}
\resolucao{\textbf{Malha 1:} $E_1 = R_1 I_1 + R_2(I_1 - I_2)$:
\[
18 = 2I_1 + 6(I_1 - I_2) \Rightarrow 8I_1 - 6I_2 = 18.
\]
\textbf{Malha 2:} percorrendo com $E_2$ opondo-se a $I_2$:
$-E_2 = R_3 I_2 + R_2(I_2 - I_1)$, isto é
\[
-6 = 6I_2 + 6(I_2 - I_1) \Rightarrow -6I_1 + 12I_2 = -6.
\]
Resolvendo o sistema:
\[
\begin{cases}
8I_1 - 6I_2 = 18\\
-6I_1 + 12I_2 = -6
\end{cases}
\Rightarrow I_1 = \SI{3}{\ampere},\ I_2 = \SI{1}{\ampere}.
\]
A corrente no resistor central $R_2$ é $I_1 - I_2 = \SI{2}{\ampere}$.
\textbf{Verificação} (malha 1): $2\cdot3 + 6\cdot2 = 6+12 = \SI{18}{\volt} = E_1$.
Confere. (Este é o mesmo circuito resolvido por Kirchhoff no Capítulo 2 --- os
dois métodos dão o mesmo resultado.)}
\end{questao}

\begin{questao}[2]{}
No circuito da figura, com uma única fonte, aplique o método de Maxwell e
determine as correntes de malha $I_1$ e $I_2$ e a corrente real no resistor
central.

\circuitoq{%
\begin{circuitikz}[scale=1,transform shape,line width=0.8pt]
 \draw (0,0) to[battery1,l_={$\SI{24}{\volt}$}] (0,3)
       to[R,l={$R_1=\SI{2}{\ohm}$}] (3,3) coordinate(top);
 \draw (top) to[R,l={$R_2=\SI{6}{\ohm}$}] (3,0) coordinate(bot) -- (0,0);
 \draw (top) to[R,l={$R_3=\SI{3}{\ohm}$}] (6,3) -- (6,0) -- (bot);
 \draw[->,color=loopcol,line width=1.4pt] (1.4,1.5) arc (0:310:0.5);
 \node[loopcol] at (1.4,1.5){\footnotesize$I_1$};
 \draw[->,color=loopcol,line width=1.4pt] (4.6,1.5) arc (0:310:0.5);
 \node[loopcol] at (4.6,1.5){\footnotesize$I_2$};
\end{circuitikz}}

\resposta{$I_1 = \SI{6}{\ampere}$, $I_2 = \SI{4}{\ampere}$; $I_{R_2}=\SI{2}{\ampere}$}
\resolucao{Correntes de malha no sentido horário; $R_2$ (comum) conduz
$I_1 - I_2$.\\
\textbf{Malha 1:} $24 = R_1 I_1 + R_2(I_1-I_2) = 2I_1 + 6(I_1-I_2)$, ou
$8I_1 - 6I_2 = 24$.\\
\textbf{Malha 2:} $0 = R_3 I_2 + R_2(I_2-I_1) = 3I_2 + 6(I_2-I_1)$, ou
$-6I_1 + 9I_2 = 0 \Rightarrow I_1 = \tfrac{3}{2}I_2$.\\
Substituindo: $8\cdot\tfrac32 I_2 - 6I_2 = 24 \Rightarrow 6I_2 = 24
\Rightarrow I_2 = \SI{4}{\ampere}$ e $I_1 = \SI{6}{\ampere}$.\\
Corrente real em $R_2$: $I_1 - I_2 = \SI{2}{\ampere}$ (de cima para baixo).}
\end{questao}

\begin{questao}[2]{}
Sintetize, em uma tabela comparativa, quando usar cada técnica de análise
(nodal simples, supernó, malhas simples, supermalha), indicando o "sintoma" que
identifica cada caso.
\resposta{Ver tabela comentada}
\resolucao{
\begin{center}\small\sffamily
\setlength{\arraystretch}{1.35}
\begin{tabular}{@{}p{3.1cm} p{4.6cm} p{4.6cm}@{}}
\toprule
\textbf{Técnica} & \textbf{Sintoma que a indica} & \textbf{Equação extra}\\
\midrule
Nodal simples & Fontes de corrente; fontes de tensão com um terminal no terra
  & --- \\
\textbf{Supernó} & Fonte de tensão \emph{entre dois nós}, nenhum no terra
  & $V_a - V_b = E$ (restrição)\\
Malhas simples & Fontes de tensão; circuito planar
  & --- \\
\textbf{Supermalha} & Fonte de corrente \emph{compartilhada} por duas malhas
  & $I_b - I_a = I_{\text{fonte}}$ (restrição)\\
\bottomrule
\end{tabular}
\end{center}

\medskip
\textbf{O princípio unificador.} Em ambos os "super" métodos, a fonte
problemática é aquela cuja \emph{variável conjugada} é desconhecida: na fonte de
tensão não se sabe a \emph{corrente}; na fonte de corrente não se sabe a
\emph{tensão}. A estratégia é sempre a mesma:
\begin{enumerate}[leftmargin=1.7em,itemsep=1pt,topsep=2pt]
\item escreva a equação de conservação em um contorno que \emph{contorne} o
      elemento problemático (englobando-o, no supernó; evitando-o, na supermalha);
\item recupere a equação que falta a partir da \emph{própria definição} da fonte.
\end{enumerate}
Reconhecer esse padrão evita decorar dois procedimentos distintos --- é um só
princípio aplicado de forma dual.}
\end{questao}

\blocoaprofundamento

\begin{questao}[3]{}
Um estudante deve escolher entre análise nodal e análise de malhas para resolver
um circuito planar com $N$ nós e $M$ malhas. Explique o critério de escolha e
aplique-o ao circuito de duas malhas e três nós da questão anterior.
\resposta{Escolher o método com menor número de equações; aqui, malhas (2 eq.)}
\resolucao{Cada método gera um sistema com tantas equações quantas forem as
incógnitas: a análise \textbf{nodal} produz $(N-1)$ equações (nós menos o terra),
e a análise de \textbf{malhas}, $M$ equações (número de laços independentes).
Escolhe-se o que resulta em \emph{menos} equações.\\
No circuito anterior, há $N = 3$ nós (portanto $N-1 = 2$ equações nodais) e
$M = 2$ malhas (2 equações de malha) --- empate. Como havia \emph{fontes de
tensão}, as malhas costumam ser mais diretas (as fem entram sem manipulação),
então a análise de malhas foi a escolha natural. Se o circuito tivesse muitas
fontes de \emph{corrente} e poucos nós, a análise nodal seria preferível.
\textbf{Regra prática:} conte $N-1$ e $M$; havendo empate, deixe as fontes
decidirem --- fontes de tensão favorecem malhas; fontes de corrente favorecem
nós.}
\end{questao}

\begin{questao}[3]{}
Compare os resultados obtidos para um mesmo circuito de duas malhas resolvido por
três caminhos: (i) leis de Kirchhoff diretas, (ii) método de Maxwell (malhas) e
(iii) análise nodal. O que se pode concluir?
\resposta{Ver comentário}
\resolucao{Os três métodos são \emph{equivalentes} --- todos derivam das mesmas
leis de Kirchhoff (nós e malhas) e da lei de Ohm, e por isso levam
\textbf{exatamente ao mesmo resultado} para correntes e tensões. O que muda é a
\emph{organização} do trabalho:
\begin{itemize}[leftmargin=1.6em,itemsep=1pt,topsep=2pt]
\item \textbf{Kirchhoff direto}: intuitivo, mas gera muitas incógnitas
      (uma corrente por ramo) e equações.
\item \textbf{Maxwell (malhas)}: reduz o número de incógnitas ao número de malhas;
      ideal quando há fontes de \emph{tensão}.
\item \textbf{Nodal}: reduz ao número de nós menos um; ideal quando há fontes de
      \emph{corrente}.
\end{itemize}
A escolha é uma questão de \emph{eficiência}, não de correção. Dominar os três dá
ao estudante a flexibilidade de escolher o mais econômico para cada circuito.}
\end{questao}

\begin{questao}[3]{}
No circuito da figura, uma fonte de corrente de $\SI{2}{\ampere}$ é
\textbf{compartilhada} pelas malhas 1 e 2. Determine as correntes de malha $I_1$ e
$I_2$ pelo método da \textbf{supermalha}.

\circuitoq{%
\begin{circuitikz}[scale=1,transform shape,line width=0.8pt,american]
 % contorno externo
 \draw (0,0) to[battery1,l_={$\SI{24}{\volt}$}] (0,3)
       to[R,l={$R_1=\SI{4}{\ohm}$}] (3,3) coordinate(top);
 % fonte de corrente compartilhada (ramo central)
 \draw (top) to[isource,l={$\SI{2}{\ampere}$}] (3,0) coordinate(bot);
 \draw (bot) -- (0,0);
 % malha 2
 \draw (top) to[R,l={$R_2=\SI{4}{\ohm}$}] (6,3) -- (6,0) -- (bot);
 % setas de malha
 \draw[->,color=loopcol,line width=1.3pt] (1.4,1.5) arc (0:310:0.45);
 \node[loopcol] at (1.4,1.5){\footnotesize$I_1$};
 \draw[->,color=loopcol,line width=1.3pt] (4.7,1.5) arc (0:310:0.45);
 \node[loopcol] at (4.7,1.5){\footnotesize$I_2$};
 % contorno da supermalha
 \draw[dashed,laranja,line width=1pt,rounded corners=10pt]
       (-0.55,-0.5) rectangle (6.55,3.6);
 \node[laranja,font=\footnotesize\sffamily] at (3,-0.85) {supermalha (contorno externo)};
\end{circuitikz}}

\resposta{$I_1 = \SI{2}{\ampere}$; $I_2 = \SI{4}{\ampere}$}
\resolucao{\textbf{Por que o método normal falha:} a tensão sobre a fonte de
corrente é desconhecida, então não se pode escrever a lei das malhas
isoladamente na malha 1 nem na 2.

\medskip
\textbf{Equação 1 --- lei das malhas na supermalha.} Percorremos o
\emph{contorno externo} (tracejado laranja), que \textbf{evita} o ramo da fonte de
corrente. Nesse percurso encontramos a fonte de \SI{24}{\volt}, $R_1$ (percorrido
por $I_1$) e $R_2$ (percorrido por $I_2$):
\[
24 = R_1 I_1 + R_2 I_2 = 4I_1 + 4I_2.
\]

\medskip
\textbf{Equação 2 --- restrição da fonte de corrente.} A fonte no ramo
compartilhado impõe a diferença entre as correntes de malha:
\[
I_2 - I_1 = 2.
\]

\medskip
\textbf{Resolvendo.} Da segunda, $I_2 = I_1+2$. Substituindo:
\[
4I_1 + 4(I_1+2) = 24
\;\Rightarrow\;
8I_1 = 16
\;\Rightarrow\;
I_1 = \SI{2}{\ampere},\quad I_2 = \SI{4}{\ampere}.
\]
\textbf{Verificação:} $4\cdot2 + 4\cdot4 = 8+16 = \SI{24}{\volt}$. Confere.\\
\textbf{Consistência:} a fonte de corrente conduz $I_2 - I_1 = \SI{2}{\ampere}$,
exatamente seu valor nominal.}
\end{questao}

\begin{questao}[3]{}
Para o circuito de \textbf{três malhas} da figura, monte o sistema matricial por
inspeção e resolva-o pela \textbf{regra de Cramer}.

\circuitoq{%
\begin{circuitikz}[scale=0.95,transform shape,line width=0.8pt]
 % malha 1
 \draw (0,0) to[battery1,l_={$E=\SI{24}{\volt}$}] (0,3)
       to[R,l={$R_1=\SI{4}{\ohm}$}] (3,3) coordinate(t1);
 \draw (t1) to[R,l={$R_2=\SI{4}{\ohm}$}] (3,0) coordinate(b1) -- (0,0);
 % malha 2
 \draw (t1) to[R,l={$R_3=\SI{2}{\ohm}$}] (6,3) coordinate(t2);
 \draw (t2) to[R,l={$R_4=\SI{4}{\ohm}$}] (6,0) coordinate(b2) -- (b1);
 % malha 3
 \draw (t2) to[R,l={$R_5=\SI{4}{\ohm}$}] (9,3) -- (9,0) -- (b2);
 % setas
 \draw[->,color=loopcol,line width=1.3pt] (1.5,1.5) arc (0:310:0.42);
 \node[loopcol] at (1.5,1.5){\footnotesize$I_1$};
 \draw[->,color=loopcol,line width=1.3pt] (4.6,1.5) arc (0:310:0.42);
 \node[loopcol] at (4.6,1.5){\footnotesize$I_2$};
 \draw[->,color=loopcol,line width=1.3pt] (7.7,1.5) arc (0:310:0.42);
 \node[loopcol] at (7.7,1.5){\footnotesize$I_3$};
\end{circuitikz}}

\resposta{$I_1=\SI{4}{\ampere}$, $I_2=\SI{2}{\ampere}$, $I_3=\SI{1}{\ampere}$}
\resolucao{\textbf{Passo 1 --- montagem por inspeção.}
\begin{itemize}[leftmargin=1.6em,itemsep=1pt,topsep=2pt]
\item Malha 1: resistências $R_1+R_2 = 4+4 = 8$; compartilha $R_2=4$ com a malha 2;
      não toca a malha 3.
\item Malha 2: $R_2+R_3+R_4 = 4+2+4 = 10$; compartilha $R_2=4$ com a 1 e $R_4=4$
      com a 3.
\item Malha 3: $R_4+R_5 = 4+4 = 8$; compartilha $R_4=4$ com a malha 2.
\end{itemize}
\[
\begin{bmatrix}
8 & -4 & 0\\
-4 & 10 & -4\\
0 & -4 & 8
\end{bmatrix}
\begin{bmatrix}I_1\\I_2\\I_3\end{bmatrix}
=
\begin{bmatrix}24\\0\\0\end{bmatrix}
\]
(Repare na \emph{simetria} da matriz --- primeira verificação de que a montagem
está correta.)

\smallskip
\textbf{Passo 2 --- determinante principal.} Expandindo pela primeira linha:
\[
\Delta = 8\,(10\cdot8 - 16) - (-4)\,(-4\cdot8 - 0)
       = 8(64) - 4(32) = 512 - 128 = 384.
\]

\smallskip
\textbf{Passo 3 --- determinantes das colunas substituídas.}
\[
\Delta_1=\begin{vmatrix}24&-4&0\\0&10&-4\\0&-4&8\end{vmatrix}
=24(80-16)=1536,
\qquad
\Delta_2=\begin{vmatrix}8&24&0\\-4&0&-4\\0&0&8\end{vmatrix}=768,
\]
\[
\Delta_3=\begin{vmatrix}8&-4&24\\-4&10&0\\0&-4&0\end{vmatrix}=384.
\]

\smallskip
\textbf{Passo 4 --- aplicar Cramer.}
\[
I_1=\frac{1536}{384}=\SI{4}{\ampere},\quad
I_2=\frac{768}{384}=\SI{2}{\ampere},\quad
I_3=\frac{384}{384}=\SI{1}{\ampere}.
\]

\smallskip
\textbf{Verificação por balanço de potência} (o teste mais confiável):
\[
P_{\text{fonte}}=E\,I_1=24\cdot4=\SI{96}{\watt},
\]
\[
P_{\text{resistores}} = 4(4)^2 + 4(4-2)^2 + 2(2)^2 + 4(2-1)^2 + 4(1)^2
= 64+16+8+4+4 = \SI{96}{\watt}.
\]
Os valores coincidem: a solução está correta.}
\end{questao}

\begin{questao}[3]{}
No circuito da questão anterior, determine as correntes \emph{reais} nos
resistores compartilhados $R_2$ e $R_4$, e explique por que elas são menores que
as correntes de malha.
\resposta{$I_{R_2}=\SI{2}{\ampere}$; $I_{R_4}=\SI{1}{\ampere}$}
\resolucao{Os resistores compartilhados são percorridos pela \emph{diferença} das
correntes de malha que os atravessam:
\[
I_{R_2}=I_1-I_2 = 4-2 = \SI{2}{\ampere},
\qquad
I_{R_4}=I_2-I_3 = 2-1 = \SI{1}{\ampere}.
\]
Elas são menores porque as duas correntes de malha atravessam o resistor comum em
sentidos \emph{opostos}, cancelando-se parcialmente. As correntes de malha são um
\emph{artifício matemático} --- não são fisicamente mensuráveis; o que um
amperímetro mediria é a corrente de ramo, ou seja, essa diferença.\\
\textbf{Consequência prática:} ao dimensionar um resistor compartilhado, use a
corrente de \emph{ramo} (menor), não a de malha --- caso contrário o
superdimensionamento é desnecessário.}
\end{questao}

\begin{questao}[3]{}
Considere a matriz de malhas de um circuito de quatro malhas:
\[
[\mathbf{R}]=
\begin{bmatrix}
10 & -4 & 0 & -2\\
-4 & 12 & -3 & 0\\
0 & -3 & 9 & -2\\
-2 & 0 & -2 & 8
\end{bmatrix}
\]
\begin{itensq}
\item Sem resolver o sistema, extraia do arranjo a topologia: quais malhas são
      adjacentes e qual o valor de cada resistência compartilhada?
\item Qual a resistência total da malha 3?
\item Por que todos os elementos fora da diagonal são negativos ou nulos?
\end{itensq}
\resposta{Ver leitura estrutural}
\resolucao{(a) Os elementos \emph{fora} da diagonal revelam a topologia. Um valor
não-nulo em $R_{ij}$ indica que as malhas $i$ e $j$ compartilham um resistor de
valor $|R_{ij}|$:
\begin{itemize}[leftmargin=1.6em,itemsep=1pt,topsep=2pt]
\item malhas 1--2: compartilham $\SI{4}{\ohm}$;
\item malhas 1--4: compartilham $\SI{2}{\ohm}$;
\item malhas 2--3: compartilham $\SI{3}{\ohm}$;
\item malhas 3--4: compartilham $\SI{2}{\ohm}$;
\item malhas 1--3 e 2--4: \emph{não são adjacentes} (elementos nulos).
\end{itemize}
A topologia é, portanto, um \textbf{anel} de quatro malhas
($1\!-\!2\!-\!3\!-\!4\!-\!1$).\\
(b) A diagonal dá a soma de todas as resistências da malha:
$R_{33} = \SI{9}{\ohm}$, dos quais \SI{3}{\ohm} são compartilhados com a malha 2 e
\SI{2}{\ohm} com a malha 4; logo, \SI{4}{\ohm} são \emph{exclusivos} da malha 3.\\
(c) Porque as correntes de malha adjacentes percorrem o resistor comum em
sentidos \emph{opostos} (ambas adotadas no mesmo sentido de rotação). A queda de
tensão que a malha vizinha provoca é, assim, subtrativa --- daí o sinal negativo.
\textbf{Isso torna a matriz simétrica e diagonalmente dominante}, o que garante
$\Delta \neq 0$ e a existência de solução única.\\
\textbf{Habilidade desenvolvida:} \emph{ler} um circuito a partir da sua matriz é
tão importante quanto montá-la --- é assim que se depuram erros em simulações e
em softwares de análise (SPICE).}
\end{questao}

\begin{questao}[3]{}
O circuito da figura tem \textbf{três malhas} dispostas em "L" --- a terceira
malha fica \emph{abaixo} da segunda, e não alinhada com as demais. Monte o
sistema matricial e resolva por Cramer.

\circuitoq{%
\begin{circuitikz}[scale=0.95,transform shape,line width=0.8pt]
 % malha 1 (superior esquerda)
 \draw (0,2.6) to[\fonteV,l_={$E=\SI{12}{\volt}$}] (0,5.4)
       to[R,l={$\SI{2}{\ohm}$}] (3,5.4) coordinate(t1);
 \draw (t1) to[R,l={$R_a=\SI{2}{\ohm}$}] (3,2.6) coordinate(m1) -- (0,2.6);
 % malha 2 (superior direita)
 \draw (t1) to[R,l={$\SI{1}{\ohm}$}] (6,5.4) coordinate(t2);
 \draw (t2) -- (6,2.6) coordinate(m2) -- (m1);
 % malha 3 (INFERIOR, abaixo da malha 1)
 \draw (m1) to[R,l_={$R_b=\SI{2}{\ohm}$}] (3,0) coordinate(b1);
 \draw (b1) to[R,l_={$\SI{4}{\ohm}$}] (6,0) coordinate(b2) -- (m2);
 % setas de malha
 \draw[->,color=loopcol,line width=1.3pt] (1.5,4.0) arc (0:310:0.42);
 \node[loopcol] at (1.5,4.0){\footnotesize$I_1$};
 \draw[->,color=loopcol,line width=1.3pt] (4.5,4.0) arc (0:310:0.42);
 \node[loopcol] at (4.5,4.0){\footnotesize$I_2$};
 \draw[->,color=loopcol,line width=1.3pt] (4.5,1.3) arc (0:310:0.42);
 \node[loopcol] at (4.5,1.3){\footnotesize$I_3$};
\end{circuitikz}}

\resposta{$I_1=\SI{3}{\ampere}$, $I_2=\SI{2}{\ampere}$, $I_3=\SI{1}{\ampere}$}
\resolucao{\textbf{Passo 1 --- identificar os acoplamentos.} Apesar da disposição
em "L", a regra é a mesma: o que importa é \emph{quais malhas compartilham
resistores}.
\begin{itemize}[leftmargin=1.6em,itemsep=1pt,topsep=2pt]
\item Malhas 1 e 2 compartilham $R_a = \SI{2}{\ohm}$;
\item Malhas 1 e 3 compartilham $R_b = \SI{2}{\ohm}$;
\item Malhas 2 e 3 \textbf{não} são adjacentes (elemento nulo na matriz).
\end{itemize}
\textbf{Atenção:} a posição \emph{geométrica} das malhas no desenho é irrelevante
--- o que define a matriz é a \emph{topologia} (quem toca quem).

\textbf{Passo 2 --- montar a matriz.}
\[
\begin{bmatrix}
 6 & -2 & -2\\
-2 &  3 &  0\\
-2 &  0 &  6
\end{bmatrix}
\begin{bmatrix}I_1\\I_2\\I_3\end{bmatrix}
=
\begin{bmatrix}12\\0\\0\end{bmatrix}
\]
Diagonal: $R_{11}=2+2+2=6$; $R_{22}=2+1=3$; $R_{33}=2+4=6$.
Os zeros em $R_{23}=R_{32}=0$ confirmam a não-adjacência entre as malhas 2 e 3.

\textbf{Passo 3 --- determinantes e Cramer.}
\[
\Delta = 6(3\cdot6-0)-(-2)\big[(-2)(6)-0\big]+(-2)\big[0-(-2)(3)\big]
= 108-24-12 = 72.
\]
\[
\Delta_1 = 216,\qquad \Delta_2 = 144,\qquad \Delta_3 = 72,
\]
\[
I_1=\frac{216}{72}=\SI{3}{\ampere},\quad
I_2=\frac{144}{72}=\SI{2}{\ampere},\quad
I_3=\frac{72}{72}=\SI{1}{\ampere}.
\]
\textbf{Verificação rápida pelas linhas 2 e 3:}
$3I_2 = 2I_1 \Rightarrow I_2=\tfrac23\cdot3=2$ \checkmark;
$6I_3 = 2I_1 \Rightarrow I_3=\tfrac13\cdot3=1$ \checkmark.

\textbf{Passo 4 --- balanço de potência.}
\[
P_{\text{res}} = 2(3)^2+2(3-2)^2+2(3-1)^2+1(2)^2+4(1)^2
= 18+2+8+4+4 = \SI{36}{\watt},
\]
\[
P_{\text{fonte}} = 12\cdot3 = \SI{36}{\watt}.\ \checkmark
\]
\textbf{Lição estrutural.} Uma matriz de malhas com \emph{zeros fora da diagonal}
indica malhas não-adjacentes. Quanto mais zeros, mais "esparsa" a matriz e mais
fácil a solução --- e é justamente por isso que programas de simulação (SPICE)
usam algoritmos para matrizes esparsas: circuitos reais têm pouquíssimos
acoplamentos por malha.}
\end{questao}

\blocodesafio

\begin{questao}[4]{}
Um circuito planar tem $b$ ramos, $n$ nós e $m$ malhas independentes.
\begin{itensq}
\item Demonstre a relação $m = b - n + 1$ (teorema de Euler para grafos planares).
\item Um circuito tem 12 ramos e 7 nós. Quantas equações de malha e quantas
      equações nodais seriam necessárias? Qual método é mais econômico?
\end{itensq}
\resposta{(a) demonstração; (b) 6 de malha vs 6 nodais --- empate}
\resolucao{(a) Cada ramo carrega uma corrente desconhecida: são $b$ incógnitas.
A lei dos nós fornece $(n-1)$ equações independentes (a $n$-ésima é combinação
linear das demais). Restam, portanto,
\[
m = b - (n-1) = b - n + 1
\]
equações independentes a serem obtidas pela lei das malhas. Essa é exatamente a
relação de Euler aplicada a grafos planares conexos.\\
(b) Com $b=12$ e $n=7$:
\[
m = 12 - 7 + 1 = 6 \text{ equações de malha};
\qquad
n-1 = 6 \text{ equações nodais}.
\]
Há \textbf{empate} (6 equações em ambos). O critério de desempate passa a ser o
tipo de fonte: predominando fontes de \emph{tensão}, prefira malhas; predominando
fontes de \emph{corrente}, prefira nodal. Se houver muitas fontes de tensão
aterradas, a nodal ganha, pois cada uma delas \emph{elimina} uma incógnita
(o potencial daquele nó é conhecido).}
\end{questao}

\begin{questao}[4]{}
O circuito da figura possui \textbf{três malhas} e \textbf{duas fontes}. Monte o
sistema matricial, resolva-o pela regra de Cramer e interprete o sinal obtido
para $I_3$. Verifique o resultado pelo balanço de potência.

\circuitoq{%
\begin{circuitikz}[scale=0.95,transform shape,line width=0.8pt]
 \draw (0,0) to[\fonteV,l={$E_1=\SI{24}{\volt}$}] (0,3)
       to[R,l={$\SI{2}{\ohm}$}] (3,3) coordinate(t1);
 \draw (t1) to[R,l={$\SI{2}{\ohm}$}] (3,0) coordinate(b1) -- (0,0);
 \draw (t1) to[R,l={$\SI{2}{\ohm}$}] (6,3) coordinate(t2);
 \draw (t2) to[R,l={$\SI{2}{\ohm}$}] (6,0) coordinate(b2) -- (b1);
 \draw (t2) to[R,l={$\SI{6}{\ohm}$}] (9,3)
       to[\fonteV,l={$E_2=\SI{12}{\volt}$},invert] (9,0) -- (b2);
 \foreach \x/\n in {1.5/1,4.5/2,7.5/3}{
   \draw[->,color=loopcol,line width=1.3pt] (\x,1.5) arc (0:310:0.42);
   \node[loopcol] at (\x,1.5){\footnotesize$I_{\n}$};}
\end{circuitikz}}

\resposta{$I_1=\SI{7}{\ampere}$, $I_2=\SI{2}{\ampere}$, $I_3=\SI{-1}{\ampere}$
(sentido oposto ao adotado)}
\resolucao{\textbf{Passo 1 --- identificar os acoplamentos.} Apesar da disposição
em "L", a regra é a mesma: o que importa é \emph{quais malhas compartilham
resistores}.
\begin{itemize}[leftmargin=1.6em,itemsep=1pt,topsep=2pt]
\item Malhas 1 e 2 compartilham $R_a = \SI{4}{\ohm}$;
\item Malhas 1 e 3 compartilham $R_b = \SI{4}{\ohm}$;
\item Malhas 2 e 3 \textbf{não} são adjacentes (elemento nulo na matriz).
\end{itemize}
\textbf{Atenção:} a posição \emph{geométrica} das malhas no desenho é irrelevante
--- o que define a matriz é a \emph{topologia} (quem toca quem).

\textbf{Passo 2 --- montar a matriz.}
\[
\begin{bmatrix}
10 & -4 & -4\\
-4 & 6 & 0\\
-4 & 0 & 8
\end{bmatrix}
\begin{bmatrix}I_1\\I_2\\I_3\end{bmatrix}
=
\begin{bmatrix}24\\0\\0\end{bmatrix}
\]
Diagonal: $R_{11}=2+4+4=10$; $R_{22}=4+2=6$; $R_{33}=4+2+2=8$.
O zero em $R_{23}$ confirma a não-adjacência.

\textbf{Passo 3 --- Cramer.}
\[
\Delta = 10(48-0)-(-4)(-32-0)+(-4)(0+24)=480-128-96=\SI{256}{}.
\]
\[
I_1=\frac{\Delta_1}{\Delta},\quad
I_2=\frac{\Delta_2}{\Delta},\quad
I_3=\frac{\Delta_3}{\Delta}.
\]
Resolvendo o sistema, obtém-se
\[
I_1 = \SI{4}{\ampere},\qquad
I_2 = \SI{2.67}{\ampere},\qquad
I_3 = \SI{2}{\ampere}.
\]
\textbf{Verificação pela malha 2:} $6I_2 = 4I_1 \Rightarrow I_2 = \tfrac{2}{3}I_1$
\checkmark; \textbf{malha 3:} $8I_3 = 4I_1 \Rightarrow I_3 = \tfrac12 I_1$
\checkmark.

\textbf{Lição estrutural.} Uma matriz de malhas com \emph{zeros fora da diagonal}
indica malhas não-adjacentes. Quanto mais zeros, mais "esparsa" a matriz e mais
fácil a solução --- e é justamente por isso que programas de simulação (SPICE)
usam algoritmos para matrizes esparsas: circuitos reais têm pouquíssimos
acoplamentos por malha.}
\end{questao}

\begin{questao}[4]{}
O circuito da figura possui \textbf{quatro malhas} em arranjo $2\times2$.
\begin{itensq}
\item Monte a matriz de malhas por inspeção e verifique sua simetria.
\item Resolva o sistema pela regra de Cramer.
\item Determine as correntes reais em todos os resistores compartilhados.
\item Valide a solução pelo balanço de potência.
\end{itensq}

\circuitoq{%
\begin{circuitikz}[scale=0.95,transform shape,line width=0.8pt]
 % nos da grade 2x2
 % linha superior y=5.2 ; meio y=2.6 ; inferior y=0
 % colunas x=0, 3.4, 6.8
 \draw (0,2.6) to[\fonteV,l_={$E=\SI{48}{\volt}$}] (0,5.2)
       to[R,l={$\SI{4}{\ohm}$}] (3.4,5.2) coordinate(A);
 \draw (A) to[R,l={$\SI{2}{\ohm}$}] (6.8,5.2) coordinate(B);
 % ramo vertical central superior (compartilhado 1-2)
 \draw (A) to[R,l={$\SI{2}{\ohm}$}] (3.4,2.6) coordinate(C);
 % ramo vertical direito (malha 2)
 \draw (B) to[R,l={$\SI{2}{\ohm}$}] (6.8,2.6) coordinate(D);
 % linha do meio
 \draw (0,2.6) -- (C); \draw (C) -- (D);
 % ramo vertical esquerdo inferior (compartilhado 1-3)
 \draw (0,2.6) to[R,l_={$\SI{1}{\ohm}$}] (0,0) coordinate(E);
 % linha inferior
 \draw (E) to[R,l_={$\SI{2}{\ohm}$}] (3.4,0) coordinate(F);
 \draw (F) to[R,l_={$\SI{6}{\ohm}$}] (6.8,0) coordinate(G);
 % ramo vertical central inferior (compartilhado 3-4)
 \draw (C) to[R,l={$\SI{2}{\ohm}$}] (F);
 \draw (D) -- (G);
 % setas
 \draw[->,color=loopcol,line width=1.3pt] (1.7,3.9) arc (0:310:0.4);
 \node[loopcol] at (1.7,3.9){\footnotesize$I_1$};
 \draw[->,color=loopcol,line width=1.3pt] (5.1,3.9) arc (0:310:0.4);
 \node[loopcol] at (5.1,3.9){\footnotesize$I_2$};
 \draw[->,color=loopcol,line width=1.3pt] (1.7,1.3) arc (0:310:0.4);
 \node[loopcol] at (1.7,1.3){\footnotesize$I_3$};
 \draw[->,color=loopcol,line width=1.3pt] (5.1,1.3) arc (0:310:0.4);
 \node[loopcol] at (5.1,1.3){\footnotesize$I_4$};
\end{circuitikz}}

\resposta{$I_1=\SI{8}{\ampere}$, $I_2=\SI{3}{\ampere}$, $I_3=\SI{2}{\ampere}$,
$I_4=\SI{1}{\ampere}$}
\resolucao{\textbf{(a) Montagem por inspeção.} Adjacências: $1$--$2$
($\SI{2}{\ohm}$), $1$--$3$ ($\SI{1}{\ohm}$), $2$--$4$ ($\SI{2}{\ohm}$),
$3$--$4$ ($\SI{2}{\ohm}$); as malhas $1$--$4$ e $2$--$3$ são \emph{diagonais},
logo não-adjacentes (zeros).
\[
[\mathbf R]=
\begin{bmatrix}
 7 & -2 & -1 & 0\\
-2 &  6 &  0 & -2\\
-1 &  0 &  5 & -2\\
 0 & -2 & -2 & 10
\end{bmatrix},
\qquad
[\mathbf E]=
\begin{bmatrix}48\\0\\0\\0\end{bmatrix}
\]
A matriz é \textbf{simétrica} \checkmark\ --- primeira verificação de que a
montagem está correta.

\textbf{(b) Regra de Cramer.}
\[
\Delta=\det[\mathbf R]=1536,
\]
\[
\Delta_1=12288,\quad \Delta_2=4608,\quad \Delta_3=3072,\quad \Delta_4=1536.
\]
Portanto
\[
I_1=\frac{12288}{1536}=\SI{8}{\ampere},\quad
I_2=\frac{4608}{1536}=\SI{3}{\ampere},
\]
\[
I_3=\frac{3072}{1536}=\SI{2}{\ampere},\quad
I_4=\frac{1536}{1536}=\SI{1}{\ampere}.
\]

\textbf{(c) Correntes de ramo} (diferença entre as malhas adjacentes):
\begin{center}\small\sffamily
\begin{tabular}{@{}l c c@{}}
\toprule
\textbf{Resistor} & \textbf{Malhas} & \textbf{Corrente real}\\
\midrule
\SI{2}{\ohm} (central sup.) & 1--2 & $I_1-I_2=\SI{5}{\ampere}$\\
\SI{1}{\ohm} (esquerdo)     & 1--3 & $I_1-I_3=\SI{6}{\ampere}$\\
\SI{2}{\ohm} (direito)      & 2--4 & $I_2-I_4=\SI{2}{\ampere}$\\
\SI{2}{\ohm} (central inf.) & 3--4 & $I_3-I_4=\SI{1}{\ampere}$\\
\bottomrule
\end{tabular}
\end{center}

\textbf{(d) Balanço de potência.}
\[
P_{\text{res}} = 4(8)^2+2(5)^2+1(6)^2+2(3)^2+2(2)^2+2(2)^2+2(1)^2+6(1)^2
\]
\[
= 256+50+36+18+8+8+2+6 = \SI{384}{\watt}.
\]
\[
P_{\text{fonte}} = E\,I_1 = 48\cdot8 = \SI{384}{\watt}.\ \checkmark
\]
\textbf{Comentário sobre o método.} Com $n=4$ malhas, resolver "na mão" por
substituição seria trabalhoso e propenso a erro. A montagem matricial por
inspeção elimina a etapa de escrever equações uma a uma, e o balanço de potência
oferece uma verificação \emph{independente} do resultado --- se ele fecha, a
probabilidade de erro é mínima. \emph{Este é o fluxo de trabalho profissional:
montar sistematicamente, resolver com álgebra linear, validar por energia.}}
\end{questao}

% BEGIN BANCO COMPLEMENTAR Método de Maxwell (Análise de Malhas)
% =====================================================================
% BANCO COMPLEMENTAR --- ANALISE DE MALHAS / METODO DE MAXWELL
% Revisao 2026-08-15
% 7 N1 + 8 N2 + 5 N3 + 5 N4 = 25 questoes.
% N2/N3/N4: todas as questoes partem de CIRCUITOS diferentes e as setas
% de corrente de malha sao desenhadas explicitamente no sentido HORARIO.
% Fontes dependentes permanecem nos bancos especificos, evitando transformar
% toda a secao em repeticoes do mesmo tema.
% =====================================================================

\subsubsection*{Banco complementar --- Método de Maxwell (análise de malhas)}

\begin{atencao}[Convenção gráfica]
Em todos os circuitos deste banco, $I_1,I_2,\ldots$ são adotadas no
\textbf{sentido horário}. Uma corrente calculada negativa significa apenas que
a corrente real daquela malha circula no sentido oposto ao indicado. Os níveis
2, 3 e 4 usam topologias e objetivos diferentes: malhas triangulares, fontes em
ramos compartilhados, fontes de corrente periféricas, redes de três malhas,
ponte resistiva, síntese, máxima potência, identificação de parâmetros e rede
em roda de quatro malhas.
\end{atencao}

\providecommand{\fvMeshCW}[3]{%
  \draw[->,loopcol,line width=1.1pt] (#1)
    arc[start angle=0,end angle=-305,radius=#2cm];%
  \node[loopcol,fill=white,inner sep=1pt] at (#1){\footnotesize #3};}

% =====================================================================
% N1
% =====================================================================
\blocofixacao

\begin{questao}[1]{}
No desenho abaixo, a corrente de malha segue a convenção usada neste capítulo.
Identifique o sentido de $I_1$.
\circuitoq{%
\begin{circuitikz}[american,scale=.9,transform shape]
\draw (0,0) to[\fonteV,l^={$E$}] (0,2.8) to[R,l={$R$}] (4,2.8) -- (4,0)--(0,0);
\fvMeshCW{2.0,1.25}{.42}{$I_1$}
\end{circuitikz}}
\resposta{Horário}
\resolucao{A seta percorre o laço no sentido horário. Essa é apenas uma
convenção de referência; se o resultado algébrico de $I_1$ for negativo, a
corrente real circula no sentido anti-horário.}
\end{questao}

\begin{questao}[1]{}
Duas malhas horárias compartilham um resistor $R_c$. Qual é a corrente real no
resistor, tomando como positivo o sentido de $I_1$ no ramo comum?
\resposta{$i_c=I_1-I_2$}
\resolucao{As correntes horárias atravessam o ramo compartilhado em sentidos
opostos; por isso a corrente de ramo é a diferença algébrica.}
\end{questao}

\begin{questao}[1]{}
Um circuito planar conectado possui $b=11$ ramos e $n=7$ nós. Quantas correntes
de malha independentes são necessárias?
\resposta{$m=5$}
\resolucao{$m=b-n+1=11-7+1=5$.}
\end{questao}

\begin{questao}[1]{}
Uma fonte ideal de corrente está na periferia de uma única malha e sua seta
coincide com a corrente horária $I_2$. O que pode ser escrito imediatamente?
\resposta{$I_2=I_s$}
\resolucao{A fonte pertence somente àquela malha; portanto ela fixa diretamente
a corrente de malha e elimina uma incógnita do sistema.}
\end{questao}

\begin{questao}[1]{}
Uma fonte ideal de corrente está no ramo comum a duas malhas. Qual técnica deve
ser usada?
\resposta{Supermalha, mais a equação de restrição da fonte}
\resolucao{A tensão sobre uma fonte ideal de corrente é desconhecida; por isso
não se escreve LKT separadamente nas duas malhas. Contorna-se o ramo da fonte e
acrescenta-se a relação entre $I_1$ e $I_2$.}
\end{questao}

\begin{questao}[1]{}
Na matriz de uma rede resistiva com correntes de malha todas horárias, o que
representam $R_{kk}$ e $R_{jk}$, $j\ne k$?
\resposta{$R_{kk}$ é a soma das resistências da malha $k$; $R_{jk}$ é o negativo
da resistência compartilhada}
\resolucao{A diagonal reúne todos os resistores percorridos por $I_k$. Fora da
diagonal aparece $-R_c$ porque as correntes horárias percorrem o ramo comum em
sentidos opostos.}
\end{questao}

\begin{questao}[1]{}
Ao resolver uma malha adotada no sentido horário, obtém-se
$I_3=-\SI{0.8}{\ampere}$. Interprete fisicamente.
\resposta{A corrente real tem módulo $\SI{0.8}{\ampere}$ e circula no sentido anti-horário}
\resolucao{O sinal negativo não indica erro: ele informa que a escolha inicial
do sentido foi oposta à corrente física.}
\end{questao}

% =====================================================================
% N2 --- circuitos distintos
% =====================================================================
\blocoaplicacao

\begin{questao}[2]{}
Duas malhas triangulares compartilham o resistor diagonal de
$\SI{6}{\ohm}$. Determine $I_1$, $I_2$ e a corrente no ramo compartilhado.
\circuitoq{%
\begin{circuitikz}[american,scale=.88,transform shape,line width=.8pt]
\coordinate (A) at (0,0); \coordinate (B) at (3.4,4.2); \coordinate (C) at (6.8,0);
\draw (A) to[\fonteV,l^={$\SI{18}{\volt}$}] (0,2.0)
      to[R,l={$\SI{4}{\ohm}$}] (B);
\draw (B) to[R,l={$\SI{3}{\ohm}$}] (5.2,2.1)
      to[R,l={$\SI{3}{\ohm}$}] (C);
\draw (C)--(A);
\draw (B) to[R,l={$\SI{6}{\ohm}$}] (3.4,0)--(A);
\fvMeshCW{1.75,1.45}{.38}{$I_1$}
\fvMeshCW{5.05,1.45}{.38}{$I_2$}
\end{circuitikz}}
\resposta{$I_1=\dfrac{18}{7}\,\si{\ampere}$; $I_2=\dfrac{9}{7}\,\si{\ampere}$;
$i_c=\dfrac{9}{7}\,\si{\ampere}$}
\resolucao{As equações são
\[
10I_1-6I_2=18,\qquad -6I_1+12I_2=0.
\]
Da segunda, $I_1=2I_2$. Logo $14I_2=18$, resultando
$I_2=9/7$ A e $I_1=18/7$ A. No ramo comum,
$i_c=I_1-I_2=9/7$ A.}
\end{questao}

\begin{questao}[2]{}
A fonte de corrente de $\SI{1.5}{\ampere}$ está na periferia da malha 2 e sua
seta coincide com $I_2$. Determine $I_1$, $I_2$ e a corrente no resistor
central de $\SI{5}{\ohm}$.
\circuitoq{%
\begin{circuitikz}[american,scale=.92,transform shape,line width=.8pt]
\draw (0,0) to[\fonteV,l^={$\SI{30}{\volt}$}] (0,3.2)
      to[R,l={$\SI{5}{\ohm}$}] (3.2,3.2) coordinate(t);
\draw (t) to[R,l_={$\SI{5}{\ohm}$}] (3.2,0) coordinate(b)--(0,0);
\draw (t) to[R,l={$\SI{10}{\ohm}$}] (6.4,3.2)
      to[isource,l={$\SI{1.5}{\ampere}$}] (6.4,0)--(b);
\fvMeshCW{1.55,1.05}{.40}{$I_1$}
\fvMeshCW{4.85,1.05}{.40}{$I_2$}
\end{circuitikz}}
\resposta{$I_1=\SI{3.75}{\ampere}$; $I_2=\SI{1.5}{\ampere}$;
$i_c=\SI{2.25}{\ampere}$}
\resolucao{A fonte periférica fixa $I_2=1{,}5$ A. Para a malha 1,
\[
5I_1+5(I_1-I_2)=30,
\]
portanto $10I_1-7{,}5=30$ e $I_1=3{,}75$ A. Assim
$i_c=I_1-I_2=2{,}25$ A.}
\end{questao}

\begin{questao}[2]{}
O ramo comum contém um resistor de $\SI{4}{\ohm}$ em série com uma fonte de
$\SI{6}{\volt}$, com terminal positivo no topo. Determine $I_1$ e $I_2$.
\circuitoq{%
\begin{circuitikz}[american,scale=.9,transform shape,line width=.8pt]
\draw (0,0) to[\fonteV,l^={$\SI{24}{\volt}$}] (0,3.6)
      to[R,l={$\SI{4}{\ohm}$}] (3.4,3.6) coordinate(t);
\draw (t) to[R,l_={$\SI{4}{\ohm}$}] (3.4,2.0)
      to[\fonteV,l_={$\SI{6}{\volt}$}] (3.4,0) coordinate(b)--(0,0);
\draw (t) to[R,l={$\SI{8}{\ohm}$}] (6.8,3.6)--(6.8,0)--(b);
\fvMeshCW{1.65,1.15}{.40}{$I_1$}
\fvMeshCW{5.15,1.15}{.40}{$I_2$}
\end{circuitikz}}
\resposta{$I_1=\SI{3}{\ampere}$; $I_2=\SI{1.5}{\ampere}$}
\resolucao{Percorrendo as duas malhas no sentido horário:
\[
8I_1-4I_2=18,
\qquad
-4I_1+12I_2=6.
\]
A solução é $I_1=3$ A e $I_2=1{,}5$ A. A fonte compartilhada entra com sinais
opostos nas duas equações porque é atravessada em sentidos opostos.}
\end{questao}

\begin{questao}[2]{}
A rede possui três malhas triangulares em torno do nó central $O$. Monte a
matriz por inspeção e determine as três correntes.
\circuitoq{%
\begin{circuitikz}[american,scale=.78,transform shape,line width=.8pt]
\coordinate (A) at (0,0); \coordinate (B) at (4,6); \coordinate (C) at (8,0); \coordinate (O) at (4,2.1);
% borda AB: R=4 e E1=16
\draw (A) to[\fonteV,l^={$\SI{16}{\volt}$}] (1.15,1.75)
      to[R,l={$\SI{4}{\ohm}$}] (B);
% borda BC: R=5 e E2=11
\draw (B) to[R,l={$\SI{5}{\ohm}$}] (6.85,1.75)
      to[\fonteV,l={$\SI{11}{\volt}$}] (C);
% borda CA: R=6 e E3=1
\draw (C) to[R,l_={$\SI{6}{\ohm}$}] (4.7,0)
      to[\fonteV,l_={$\SI{1}{\volt}$}] (A);
% raios compartilhados
\draw (B) to[R,l={$\SI{2}{\ohm}$}] (O);
\draw (C) to[R,l={$\SI{3}{\ohm}$}] (O);
\draw (A) to[R,l_={$\SI{1}{\ohm}$}] (O);
\fvMeshCW{2.8,3.5}{.34}{$I_1$}
\fvMeshCW{5.4,3.5}{.34}{$I_2$}
\fvMeshCW{4.0,.95}{.34}{$I_3$}
\node[fill=white,inner sep=1pt] at(O){$O$};
\end{circuitikz}}
\resposta{$I_1=\SI{3}{\ampere}$; $I_2=\SI{2}{\ampere}$; $I_3=\SI{1}{\ampere}$}
\resolucao{Por inspeção,
\[
\begin{bmatrix}
7&-2&-1\\
-2&10&-3\\
-1&-3&10
\end{bmatrix}
\begin{bmatrix}I_1\\I_2\\I_3\end{bmatrix}
=
\begin{bmatrix}16\\11\\1\end{bmatrix}.
\]
Substituindo $(3,2,1)$ verifica-se cada linha; logo essa é a solução. O desenho
mostra que a matriz depende de quais raios são compartilhados, e não de uma
sequência linear de malhas.}
\end{questao}

\begin{questao}[2]{}
As duas fontes do circuito atuam em oposição. Determine as correntes horárias e
interprete o sinal de $I_2$.
\circuitoq{%
\begin{circuitikz}[american,scale=.92,transform shape,line width=.8pt]
\draw (0,0) to[\fonteV,l^={$\SI{20}{\volt}$}] (0,3.2)
      to[R,l={$\SI{6}{\ohm}$}] (3.2,3.2) coordinate(t);
\draw (t) to[R,l_={$\SI{4}{\ohm}$}] (3.2,0) coordinate(b)--(0,0);
\draw (t) to[R,l={$\SI{6}{\ohm}$}] (6.4,3.2)
      to[\fonteV,l_={$\SI{12}{\volt}$}] (6.4,0)--(b);
\fvMeshCW{1.55,1.05}{.40}{$I_1$}
\fvMeshCW{4.85,1.05}{.40}{$I_2$}
\end{circuitikz}}
\resposta{$I_1=\dfrac{38}{21}\,\si{\ampere}$;
$I_2=-\dfrac{10}{21}\,\si{\ampere}$}
\resolucao{O sistema é
\[
\begin{bmatrix}10&-4\\-4&10\end{bmatrix}
\begin{bmatrix}I_1\\I_2\end{bmatrix}
=
\begin{bmatrix}20\\-12\end{bmatrix}.
\]
Daí $I_1=38/21$ A e $I_2=-10/21$ A. Portanto a corrente física da segunda
malha tem módulo $10/21$ A e circula no sentido anti-horário.}
\end{questao}

\begin{questao}[2]{}
Escolha o valor e a polaridade da fonte $E$ para que a corrente no resistor
compartilhado de $\SI{4}{\ohm}$ seja nula.
\circuitoq{%
\begin{circuitikz}[american,scale=.92,transform shape,line width=.8pt]
\draw (0,0) to[\fonteV,l^={$\SI{24}{\volt}$}] (0,3.2)
      to[R,l={$\SI{4}{\ohm}$}] (3.2,3.2) coordinate(t);
\draw (t) to[R,l_={$\SI{4}{\ohm}$}] (3.2,0) coordinate(b)--(0,0);
\draw (t) to[R,l={$\SI{8}{\ohm}$}] (6.4,3.2)
      to[\fonteV,l_={$E$}] (6.4,0)--(b);
\fvMeshCW{1.55,1.05}{.40}{$I_1$}
\fvMeshCW{4.85,1.05}{.40}{$I_2$}
\end{circuitikz}}
\resposta{$E=\SI{48}{\volt}$, impulsionando $I_2$ no sentido horário}
\resolucao{Corrente nula no ramo comum exige $I_1=I_2=I$. A primeira malha dá
$(8-4)I=24$, ou $4I=24$, logo $I=6$ A. Na segunda,
$(12-4)I=E$, portanto $E=8\cdot6=48$ V, com polaridade que impulsione $I_2$.}
\end{questao}

\begin{questao}[2]{}
Determine as correntes de malha e faça o balanço de potência.
\circuitoq{%
\begin{circuitikz}[american,scale=.92,transform shape,line width=.8pt]
\draw (0,0) to[\fonteV,l^={$\SI{36}{\volt}$}] (0,3.2)
      to[R,l={$\SI{6}{\ohm}$}] (3.2,3.2) coordinate(t);
\draw (t) to[R,l_={$\SI{3}{\ohm}$}] (3.2,0) coordinate(b)--(0,0);
\draw (t) to[R,l={$\SI{6}{\ohm}$}] (6.4,3.2)
      to[\fonteV,l_={$\SI{12}{\volt}$},invert] (6.4,0)--(b);
\fvMeshCW{1.55,1.05}{.40}{$I_1$}
\fvMeshCW{4.85,1.05}{.40}{$I_2$}
\end{circuitikz}}
\resposta{$I_1=\SI{5}{\ampere}$; $I_2=\SI{3}{\ampere}$;
$P_{fontes}=P_R=\SI{216}{\watt}$}
\resolucao{O sistema
\[
\begin{bmatrix}9&-3\\-3&9\end{bmatrix}
\begin{bmatrix}I_1\\I_2\end{bmatrix}
=
\begin{bmatrix}36\\12\end{bmatrix}
\]
fornece $(I_1,I_2)=(5,3)$ A. A corrente no resistor comum é $2$ A. Assim
\[
P_R=6(5)^2+6(3)^2+3(2)^2=150+54+12=216\,\text{W}.
\]
As fontes fornecem $36(5)+12(3)=216$ W.}
\end{questao}

\begin{questao}[2]{}
A fonte de corrente da malha 2 fixa $I_2=\SI{2}{\ampere}$. Determine $I_1$ e a
tensão $v_s$ da fonte de corrente, definida positiva do terminal superior para
o inferior.
\circuitoq{%
\begin{circuitikz}[american,scale=.92,transform shape,line width=.8pt]
\draw (0,0) to[\fonteV,l^={$\SI{20}{\volt}$}] (0,3.2)
      to[R,l={$\SI{5}{\ohm}$}] (3.2,3.2) coordinate(t);
\draw (t) to[R,l_={$\SI{5}{\ohm}$}] (3.2,0) coordinate(b)--(0,0);
\draw (t) to[R,l={$\SI{10}{\ohm}$}] (6.4,3.2)
      to[isource,l={$\SI{2}{\ampere}$}] (6.4,0)--(b);
\node[right] at(6.75,1.6){$v_s$};
\fvMeshCW{1.55,1.05}{.40}{$I_1$}
\fvMeshCW{4.85,1.05}{.40}{$I_2$}
\end{circuitikz}}
\resposta{$I_1=\SI{3}{\ampere}$; $v_s=\SI{-15}{\volt}$}
\resolucao{Na malha 1,
$5I_1+5(I_1-2)=20$, então $I_1=3$ A. Na malha 2, percorrida no sentido horário,
\[
10(2)+v_s-5(3-2)=0,
\]
portanto $v_s=-15$ V. Isso significa que o terminal inferior da fonte está
$15$ V acima do terminal superior.}
\end{questao}

% =====================================================================
% N3 --- aprofundamento com novas topologias/objetivos
% =====================================================================
\blocoaprofundamento

\begin{questao}[3]{}
A fonte de corrente de $\SI{2}{\ampere}$ está entre as malhas 1 e 2. A fonte de
$\SI{8}{\volt}$ da malha 3 se opõe ao sentido horário indicado. Resolva usando
uma supermalha 1--2.
\circuitoq{%
\begin{circuitikz}[american,scale=.82,transform shape,line width=.8pt]
\draw (0,0) to[\fonteV,l^={$\SI{24}{\volt}$}] (0,3.4)
      to[R,l={$\SI{4}{\ohm}$}] (2.8,3.4) coordinate(t1);
\draw (t1) to[isource,l={$\SI{2}{\ampere}$},invert] (2.8,0) coordinate(b1)--(0,0);
\draw (t1) to[R,l={$\SI{2}{\ohm}$}] (5.6,3.4) coordinate(t2);
\draw (t2) to[R,l_={$\SI{4}{\ohm}$}] (5.6,0) coordinate(b2)--(b1);
\draw (t2) to[R,l={$\SI{4}{\ohm}$}] (8.4,3.4)
      to[\fonteV,l_={$\SI{8}{\volt}$}] (8.4,0)--(b2);
\fvMeshCW{1.25,1.05}{.34}{$I_1$}
\fvMeshCW{4.15,1.05}{.34}{$I_2$}
\fvMeshCW{7.05,1.05}{.34}{$I_3$}
\draw[dashed,laranja,line width=.9pt,rounded corners=7pt] (-.45,-.45) rectangle (6.05,3.85);
\node[laranja,fill=white,inner sep=1pt] at(2.8,-.72){\footnotesize supermalha 1--2};
\end{circuitikz}}
\resposta{$I_1=\SI{1.5}{\ampere}$; $I_2=\SI{3.5}{\ampere}$;
$I_3=\SI{0.75}{\ampere}$}
\resolucao{As equações são
\[
4I_1+6I_2-4I_3=24,
\qquad I_2-I_1=2,
\qquad -4I_2+8I_3=-8.
\]
Resolvendo, obtém-se $I_1=3/2$ A, $I_2=7/2$ A e $I_3=3/4$ A.}
\end{questao}

\begin{questao}[3]{}
O ramo comum é uma fonte ideal de $\SI{6}{\volt}$, positiva no topo. A fonte de
$\SI{18}{\volt}$ impulsiona $I_1$, enquanto a de $\SI{12}{\volt}$ se opõe a
$I_2$. Determine as correntes e a potência da fonte compartilhada.
\circuitoq{%
\begin{circuitikz}[american,scale=.92,transform shape,line width=.8pt]
\draw (0,0) to[\fonteV,l^={$\SI{18}{\volt}$}] (0,3.2)
      to[R,l={$\SI{6}{\ohm}$}] (3.2,3.2) coordinate(t);
\draw (t) to[\fonteV,l_={$\SI{6}{\volt}$}] (3.2,0) coordinate(b)--(0,0);
\draw (t) to[R,l={$\SI{4}{\ohm}$}] (6.4,3.2)
      to[\fonteV,l_={$\SI{12}{\volt}$}] (6.4,0)--(b);
\fvMeshCW{1.55,1.05}{.40}{$I_1$}
\fvMeshCW{4.85,1.05}{.40}{$I_2$}
\end{circuitikz}}
\resposta{$I_1=\SI{2}{\ampere}$; $I_2=\SI{-1.5}{\ampere}$;
$P_{6V}=\SI{21}{\watt}$ absorvidos}
\resolucao{Na malha 1, $6I_1+6=18$, logo $I_1=2$ A. Na malha 2,
$4I_2+12-6=0$, então $I_2=-1{,}5$ A. A corrente real na fonte comum, para
baixo, é $I_1-I_2=3{,}5$ A. Como ela entra no terminal positivo superior,
\[
P=6(3{,}5)=21\,\text{W},
\]
portanto a fonte de $6$ V absorve potência.}
\end{questao}

\begin{questao}[3]{}
As fontes de corrente periféricas fixam $I_1=\SI{2}{\ampere}$ e
$I_3=\SI{1}{\ampere}$. Determine $I_2$ e identifique qualquer ramo
compartilhado com corrente nula.
\circuitoq{%
\begin{circuitikz}[american,scale=.82,transform shape,line width=.8pt]
\draw (0,0) to[isource,l^={$\SI{2}{\ampere}$}] (0,3.4)
      --(2.8,3.4) coordinate(t1);
\draw (t1) to[R,l_={$\SI{2}{\ohm}$}] (2.8,0) coordinate(b1)--(0,0);
\draw (t1) to[R,l={$\SI{4}{\ohm}$}] (4.5,3.4)
      to[\fonteV,l={$\SI{11}{\volt}$}] (5.8,3.4) coordinate(t2);
\draw (t2) to[R,l_={$\SI{3}{\ohm}$}] (5.8,0) coordinate(b2)--(b1);
\draw (t2)--(8.6,3.4) to[isource,l={$\SI{1}{\ampere}$}] (8.6,0)--(b2);
\fvMeshCW{1.25,1.05}{.34}{$I_1$}
\fvMeshCW{4.30,1.05}{.34}{$I_2$}
\fvMeshCW{7.25,1.05}{.34}{$I_3$}
\end{circuitikz}}
\resposta{$I_2=\SI{2}{\ampere}$; o resistor de $\SI{2}{\ohm}$ conduz corrente nula}
\resolucao{A equação da malha central é
\[
(2+4+3)I_2-2I_1-3I_3=11.
\]
Com $I_1=2$ A e $I_3=1$ A:
$9I_2-4-3=11$, portanto $I_2=2$ A. Assim, no resistor compartilhado entre as
malhas 1 e 2, $I_1-I_2=0$.}
\end{questao}

\begin{questao}[3]{}
Analise a ponte resistiva pelo método das malhas e determine a corrente no
resistor de ponte de $\SI{5}{\ohm}$.
\circuitoq{%
\begin{circuitikz}[american,scale=.78,transform shape,line width=.8pt]
\coordinate (A) at (5,5); \coordinate (D) at (5,0); \coordinate (B) at (2.6,2.5); \coordinate (C) at (7.4,2.5);
\draw (A)--(0.5,5) to[\fonteV,l^={$\SI{18}{\volt}$}] (0.5,0)--(D);
\draw (A) to[R,l_={$\SI{2}{\ohm}$}] (B) to[R,l_={$\SI{4}{\ohm}$}] (D);
\draw (A) to[R,l={$\SI{3}{\ohm}$}] (C) to[R,l={$\SI{6}{\ohm}$}] (D);
\draw (B) to[R,l={$\SI{5}{\ohm}$}] (C);
\fvMeshCW{2.1,2.55}{.34}{$I_1$}
\fvMeshCW{5.0,3.55}{.32}{$I_2$}
\fvMeshCW{5.0,1.45}{.32}{$I_3$}
\end{circuitikz}}
\resposta{$I_1=\SI{5}{\ampere}$; $I_2=I_3=\SI{2}{\ampere}$;
$i_{ponte}=0$}
\resolucao{A matriz é
\[
\begin{bmatrix}
6&-2&-4\\
-2&10&-5\\
-4&-5&15
\end{bmatrix}
\begin{bmatrix}I_1\\I_2\\I_3\end{bmatrix}
=
\begin{bmatrix}18\\0\\0\end{bmatrix}.
\]
A solução é $(5,2,2)$ A. O resistor de ponte é comum às malhas 2 e 3, logo
$i_{ponte}=I_2-I_3=0$. O resultado expressa o equilíbrio da ponte.}
\end{questao}

\begin{questao}[3]{}
Use uma fonte de teste de $\SI{12}{\volt}$ e análise de malhas para determinar
a resistência equivalente vista pela fonte.
\circuitoq{%
\begin{circuitikz}[american,scale=.92,transform shape,line width=.8pt]
\draw (0,0) to[\fonteV,l^={$\SI{12}{\volt}$}] (0,3.2)
      to[R,l={$\SI{3}{\ohm}$}] (3.2,3.2) coordinate(t);
\draw (t) to[R,l_={$\SI{6}{\ohm}$}] (3.2,0) coordinate(b)--(0,0);
\draw (t) to[R,l={$\SI{6}{\ohm}$}] (6.4,3.2)--(6.4,0)--(b);
\fvMeshCW{1.55,1.05}{.40}{$I_1$}
\fvMeshCW{4.85,1.05}{.40}{$I_2$}
\end{circuitikz}}
\resposta{$I_1=\SI{2}{\ampere}$; $I_2=\SI{1}{\ampere}$;
$R_{eq}=\SI{6}{\ohm}$}
\resolucao{As equações são
\[
9I_1-6I_2=12,\qquad -6I_1+12I_2=0.
\]
Da segunda, $I_1=2I_2$; substituindo na primeira, $I_2=1$ A e $I_1=2$ A. A
fonte de teste fornece $2$ A, portanto $R_{eq}=12/2=6\,\ohm$.}
\end{questao}

% =====================================================================
% N4 --- sintese, otimizacao e problemas inversos
% =====================================================================
\blocodesafio

\begin{questao}[4]{}
A rede em ``roda'' possui quatro malhas triangulares em torno do nó $O$.
Monte a matriz e determine as quatro correntes.
\circuitoq{%
\begin{circuitikz}[american,scale=.72,transform shape,line width=.8pt]
\coordinate (A) at (-4,4); \coordinate (B) at (4,4); \coordinate (C) at (4,-4); \coordinate (D) at (-4,-4); \coordinate (O) at (0,0);
% lados externos: R=4 em serie com fontes
\draw (A) to[\fonteV,l^={$\SI{24}{\volt}$}] (-1.8,4) to[R,l={$\SI{4}{\ohm}$}] (B);
\draw (B) to[\fonteV,l={$\SI{12}{\volt}$}] (4,1.8) to[R,l={$\SI{4}{\ohm}$}] (C);
\draw (C) to[\fonteV,l_={$\SI{8}{\volt}$}] (1.8,-4) to[R,l_={$\SI{4}{\ohm}$}] (D);
\draw (D) to[\fonteV,l_={$\SI{4}{\volt}$},invert] (-4,1.8) to[R,l_={$\SI{4}{\ohm}$}] (A);
% raios
\draw (A) to[R,l={$\SI{2}{\ohm}$}] (O);
\draw (B) to[R,l_={$\SI{2}{\ohm}$}] (O);
\draw (C) to[R,l={$\SI{2}{\ohm}$}] (O);
\draw (D) to[R,l_={$\SI{2}{\ohm}$}] (O);
\fvMeshCW{0,2.45}{.34}{$I_1$}
\fvMeshCW{2.45,0}{.34}{$I_2$}
\fvMeshCW{0,-2.45}{.34}{$I_3$}
\fvMeshCW{-2.45,0}{.34}{$I_4$}
\node[fill=white,inner sep=1pt] at(O){$O$};
\end{circuitikz}}
\resposta{$I_1=\SI{4}{\ampere}$; $I_2=\SI{3}{\ampere}$;
$I_3=\SI{2}{\ampere}$; $I_4=\SI{1}{\ampere}$}
\resolucao{Cada malha contém $4\,\ohm$ exclusivo e dois raios de
$2\,\ohm$. A matriz é cíclica:
\[
\begin{bmatrix}
8&-2&0&-2\\
-2&8&-2&0\\
0&-2&8&-2\\
-2&0&-2&8
\end{bmatrix}
\begin{bmatrix}I_1\\I_2\\I_3\\I_4\end{bmatrix}
=
\begin{bmatrix}24\\12\\8\\-4\end{bmatrix}.
\]
A solução é $(4,3,2,1)$ A. O termo $-4$ representa a fonte da quarta malha
orientada contra a corrente horária de referência.}
\end{questao}

\begin{questao}[4]{}
Determine o resistor compartilhado $R$ para que $I_2=I_1/2$. Calcule também
as correntes resultantes.
\circuitoq{%
\begin{circuitikz}[american,scale=.92,transform shape,line width=.8pt]
\draw (0,0) to[\fonteV,l^={$\SI{20}{\volt}$}] (0,3.2)
      to[R,l={$\SI{4}{\ohm}$}] (3.2,3.2) coordinate(t);
\draw (t) to[R,l_={$R$}] (3.2,0) coordinate(b)--(0,0);
\draw (t) to[R,l={$\SI{6}{\ohm}$}] (6.4,3.2)--(6.4,0)--(b);
\fvMeshCW{1.55,1.05}{.40}{$I_1$}
\fvMeshCW{4.85,1.05}{.40}{$I_2$}
\end{circuitikz}}
\resposta{$R=\SI{6}{\ohm}$; $I_1=\dfrac{20}{7}\,\si{\ampere}$;
$I_2=\dfrac{10}{7}\,\si{\ampere}$}
\resolucao{Na malha 2,
\[
-RI_1+(6+R)I_2=0.
\]
Impondo $I_2=I_1/2$:
$-R+(6+R)/2=0$, logo $R=6\,\ohm$. Na malha 1,
$10I_1-6I_2=20$; com $I_2=I_1/2$, resulta $7I_1=20$.}
\end{questao}

\begin{questao}[4]{}
O resistor compartilhado é a carga $R_L$. Determine $R_L$ para que sua potência
seja máxima e calcule $P_{L,\max}$.
\circuitoq{%
\begin{circuitikz}[american,scale=.92,transform shape,line width=.8pt]
\draw (0,0) to[\fonteV,l^={$\SI{24}{\volt}$}] (0,3.2)
      to[R,l={$\SI{4}{\ohm}$}] (3.2,3.2) coordinate(t);
\draw (t) to[R,l_={$R_L$}] (3.2,0) coordinate(b)--(0,0);
\draw (t) to[R,l={$\SI{8}{\ohm}$}] (6.4,3.2)
      to[\fonteV,l_={$\SI{12}{\volt}$},invert] (6.4,0)--(b);
\fvMeshCW{1.55,1.05}{.40}{$I_1$}
\fvMeshCW{4.85,1.05}{.40}{$I_2$}
\end{circuitikz}}
\resposta{$R_L=\dfrac{8}{3}\,\ohm\approx\SI{2.67}{\ohm}$;
$P_{L,\max}=\SI{13.5}{\watt}$}
\resolucao{As equações são
\[
(4+R_L)I_1-R_LI_2=24,
\]
\[
-R_LI_1+(8+R_L)I_2=12.
\]
A corrente na carga é
\[
i_L=I_1-I_2=\frac{36}{3R_L+8}.
\]
Logo
\[
P_L=R_Li_L^2=\frac{1296R_L}{(3R_L+8)^2}.
\]
Derivando e igualando a zero, obtém-se $R_L=8/3\,\ohm$. Nesse ponto,
$i_L=36/16=2{,}25$ A e $P_{L,\max}=13{,}5$ W.}
\end{questao}

\begin{questao}[4]{}
A fonte $E$ deve ser escolhida para impor $I_2=0$. Determine $E$, sua polaridade
relativa ao sentido horário e as correntes $I_1$ e $I_3$.
\circuitoq{%
\begin{circuitikz}[american,scale=.82,transform shape,line width=.8pt]
\draw (0,0) to[\fonteV,l^={$\SI{24}{\volt}$}] (0,3.4)
      to[R,l={$\SI{6}{\ohm}$}] (2.8,3.4) coordinate(t1);
\draw (t1) to[R,l_={$\SI{2}{\ohm}$}] (2.8,0) coordinate(b1)--(0,0);
\draw (t1) to[R,l={$\SI{4}{\ohm}$}] (5.8,3.4) coordinate(t2);
\draw (t2) to[R,l_={$\SI{1}{\ohm}$}] (5.8,0) coordinate(b2)--(b1);
\draw (t2) to[R,l={$\SI{4}{\ohm}$}] (8.8,3.4)
      to[\fonteV,l_={$E$}] (8.8,0)--(b2);
\fvMeshCW{1.25,1.05}{.34}{$I_1$}
\fvMeshCW{4.30,1.05}{.34}{$I_2$}
\fvMeshCW{7.35,1.05}{.34}{$I_3$}
\end{circuitikz}}
\resposta{$I_1=\SI{3}{\ampere}$; $I_2=0$; $I_3=\SI{-6}{\ampere}$;
$E=\SI{30}{\volt}$ orientada contra $I_3$ horário}
\resolucao{O sistema é
\[
8I_1-2I_2=24,
\qquad
-2I_1+7I_2-I_3=0,
\qquad
-I_2+5I_3=E_s,
\]
onde $E_s$ é positivo se impulsiona $I_3$ horário. Impondo $I_2=0$:
$I_1=3$ A e, pela segunda equação, $I_3=-6$ A. Assim
$E_s=5(-6)=-30$ V. Portanto a fonte deve ter magnitude $30$ V e polaridade
oposta à referência horária.}
\end{questao}

\begin{questao}[4]{}
Em ensaio de uma rede de duas malhas, mediram-se $I_1=\SI{3}{\ampere}$ e
$I_2=\SI{2}{\ampere}$. Determine a resistência desconhecida $R$ do ramo comum
e a potência dissipada nela.
\circuitoq{%
\begin{circuitikz}[american,scale=.92,transform shape,line width=.8pt]
\draw (0,0) to[\fonteV,l^={$\SI{20}{\volt}$}] (0,3.2)
      to[R,l={$\SI{4}{\ohm}$}] (3.2,3.2) coordinate(t);
\draw (t) to[R,l_={$R$}] (3.2,0) coordinate(b)--(0,0);
\draw (t) to[R,l={$\SI{6}{\ohm}$}] (6.4,3.2)
      to[\fonteV,l_={$\SI{4}{\volt}$},invert] (6.4,0)--(b);
\fvMeshCW{1.55,1.05}{.40}{$I_1$}
\fvMeshCW{4.85,1.05}{.40}{$I_2$}
\end{circuitikz}}
\resposta{$R=\SI{8}{\ohm}$; $P_R=\SI{8}{\watt}$}
\resolucao{Pela malha 1,
\[
(4+R)3-R(2)=20 \Rightarrow 12+R=20 \Rightarrow R=8\,\ohm.
\]
A malha 2 confirma o mesmo valor:
\[
-3R+(6+R)2=4 \Rightarrow 12-R=4 \Rightarrow R=8\,\ohm.
\]
A corrente no ramo comum é $I_1-I_2=1$ A; portanto
$P_R=8(1)^2=8$ W.}
\end{questao}

% END BANCO COMPLEMENTAR

\gabarito[Capítulo 9 --- Método de Maxwell]
                % Metodo de Maxwell / malhas

% Cap. 10 ------------------------------------------------------------
\setcounter{section}{9}
% =====================================================================
%  CAPITULO 10 - CAPACITORES E INDUTORES
% =====================================================================
\section{Capacitores e Indutores}

\begin{conceito}[Antes de começar]
\textbf{Capacitor plano:} $C = \dfrac{\varepsilon_0\,A}{d}$, com
$\varepsilon_0 = \SI{8.85e-12}{\farad\per\meter}$.\par
\textbf{Carga e energia:} $Q = C\,U$ \quad e \quad
$E = \tfrac12 C\,U^2 = \tfrac12 Q\,U$.\par
\textbf{Associação (oposta à dos resistores):}
\begin{itemize}[leftmargin=1.6em,itemsep=1pt,topsep=2pt]
\item \emph{paralelo:} $\Ceq = C_1 + C_2 + \dots$ (as capacitâncias somam);
\item \emph{série:} $\dfrac{1}{\Ceq} = \dfrac{1}{C_1}+\dfrac{1}{C_2}+\dots$
\end{itemize}
\textbf{Circuito RC --- constante de tempo:} $\tau = R\,C$. Na \emph{carga},
$v_C(t) = U_0\big(1-e^{-t/\tau}\big)$; na \emph{descarga},
$v_C(t) = U_0\,e^{-t/\tau}$. Após $5\tau$, considera-se o regime permanente
atingido.\par
\textbf{Indutor:} $E = \tfrac12 L\,I^2$; a associação segue a \emph{mesma} regra
dos resistores (série soma, paralelo inverso). Constante de tempo RL:
$\tau = L/R$.
\end{conceito}

\legendaniveis

% =====================================================================
\subsection{Capacitores}
% =====================================================================

\blocofixacao

\begin{questao}[1]{}
A capacitância de um capacitor de placas paralelas:
\begin{alternativas}
\item aumenta com a distância entre as placas.
\item aumenta com a área das placas.
\item independe do meio entre as placas.
\item diminui ao inserir um dielétrico.
\item depende da tensão aplicada.
\end{alternativas}
\resposta{(b)}
\resolucao{De $C = \dfrac{\varepsilon_0 A}{d}$: a capacitância cresce com a área
$A$ e \emph{diminui} com a distância $d$ --- o que elimina (a). Ela não depende da
tensão (elimina e), e inserir um dielétrico \emph{aumenta} $C$ (elimina d).}
\end{questao}

\begin{questao}[1]{}
Calcule a capacitância de um capacitor plano de área $A = \SI{100}{\centi\meter\squared}$
e distância entre placas $d = \SI{1}{\milli\meter}$, no vácuo.
\dados{$\varepsilon_0 = \SI{8.85e-12}{\farad\per\meter}$.}
\resposta{$C \approx \SI{88.5}{\pico\farad}$}
\resolucao{Convertendo: $A = \SI{100}{\centi\meter\squared} = \pot{100}{-4}\,\si{\meter\squared}
= \pot{1}{-2}\,\si{\meter\squared}$ e $d = \pot{1}{-3}\,\si{\meter}$:
\[
C=\frac{\varepsilon_0 A}{d}
=\frac{\pot{8.85}{-12}\cdot\pot{1}{-2}}{\pot{1}{-3}}
=\pot{8.85}{-11}\,\si{\farad}\approx\SI{88.5}{\pico\farad}.
\]}
\end{questao}

\begin{questao}[1]{}
Um capacitor de $\SI{10}{\micro\farad}$ é submetido a uma tensão de
$\SI{12}{\volt}$. Determine a carga armazenada e a energia acumulada.
\resposta{$Q = \SI{120}{\micro\coulomb}$; $E = \SI{720}{\micro\joule}$}
\resolucao{$Q = C\,U = \pot{10}{-6}\cdot12 = \SI{120}{\micro\coulomb}$.\\
$E = \tfrac12 C\,U^2 = \tfrac12\cdot\pot{10}{-6}\cdot12^2
= \SI{720}{\micro\joule}$.}
\end{questao}

\begin{questao}[1]{}
Dois capacitores de $\SI{6}{\micro\farad}$ e $\SI{3}{\micro\farad}$ são ligados
em \textbf{série}. A capacitância equivalente é:
\begin{alternativas}[3]
\item \SI{2}{\micro\farad}
\item \SI{4.5}{\micro\farad}
\item \SI{9}{\micro\farad}
\item \SI{18}{\micro\farad}
\item \SI{0.5}{\micro\farad}
\end{alternativas}
\resposta{(a)}
\resolucao{Em série, $\Ceq = \dfrac{C_1 C_2}{C_1 + C_2} = \dfrac{6\cdot3}{9}
= \SI{2}{\micro\farad}$. Note que, no capacitor, a série se comporta como o
paralelo do resistor: o equivalente é \emph{menor} que o menor deles.}
\end{questao}

\begin{questao}[1]{}
Dois capacitores de $\SI{4}{\micro\farad}$ e $\SI{12}{\micro\farad}$ são ligados
em \textbf{paralelo}. A capacitância equivalente é:
\begin{alternativas}[3]
\item \SI{3}{\micro\farad}
\item \SI{8}{\micro\farad}
\item \SI{16}{\micro\farad}
\item \SI{48}{\micro\farad}
\item \SI{6}{\micro\farad}
\end{alternativas}
\resposta{(c)}
\resolucao{No paralelo, as capacitâncias \emph{somam}:
$\Ceq = 4 + 12 = \SI{16}{\micro\farad}$.}
\end{questao}

\begin{questao}[1]{}
Um capacitor de $\SI{100}{\micro\farad}$ é carregado a $\SI{50}{\volt}$. A energia
armazenada é:
\begin{alternativas}[3]
\item \SI{2.5}{\milli\joule}
\item \SI{5}{\milli\joule}
\item \SI{125}{\milli\joule}
\item \SI{250}{\milli\joule}
\item \SI{5}{\joule}
\end{alternativas}
\resposta{(c)}
\resolucao{$E = \tfrac12 C U^2 = \tfrac12\cdot\pot{100}{-6}\cdot50^2
= \tfrac12\cdot\pot{100}{-6}\cdot2500 = \SI{0.125}{\joule} = \SI{125}{\milli\joule}$.
Essa energia fica armazenada no campo elétrico entre as placas e pode ser liberada
rapidamente --- princípio do flash fotográfico e do desfibrilador.}
\end{questao}

\blocoaplicacao

\begin{questao}[2]{}
Uma associação mista tem $C_1 = \SI{6}{\micro\farad}$ e
$C_2 = \SI{3}{\micro\farad}$ em \emph{paralelo}, e esse conjunto em \emph{série}
com $C_3 = \SI{2}{\micro\farad}$. Calcule a capacitância total.

\circuitoq{%
\begin{circuitikz}[scale=1,transform shape,line width=0.8pt]
 \draw (0,0.75) to[short,o-] (1,0.75) coordinate(a);
 \draw (a) -- (1,1.5) to[C,l={$C_1=\SI{6}{\micro\farad}$}] (3,1.5) -- (3,0.75) coordinate(b);
 \draw (a) -- (1,0) to[C,l_={$C_2=\SI{3}{\micro\farad}$}] (3,0) -- (3,0.75);
 \draw (b) to[C,l={$C_3=\SI{2}{\micro\farad}$}] (5,0.75) to[short,-o] (6,0.75);
\end{circuitikz}}

\resposta{$\Ceq \approx \SI{1.64}{\micro\farad}$}
\resolucao{\textbf{Paralelo} $C_1$ e $C_2$: somam-se,
$C_{12} = 6 + 3 = \SI{9}{\micro\farad}$.\\
\textbf{Série} com $C_3$:
\[
\Ceq=\frac{C_{12}\,C_3}{C_{12}+C_3}=\frac{9\cdot2}{11}=\frac{18}{11}
\approx\SI{1.64}{\micro\farad}.
\]}
\end{questao}

\begin{questao}[2]{}
Uma associação é feita com dois capacitores de $\SI{8}{\micro\farad}$ em
\emph{série}, e esse conjunto em \emph{paralelo} com um capacitor de
$\SI{4}{\micro\farad}$. Calcule a capacitância equivalente.
\resposta{$\Ceq = \SI{8}{\micro\farad}$}
\resolucao{\textbf{Série} dos dois de \SI{8}{\micro\farad}:
$\dfrac{8}{2} = \SI{4}{\micro\farad}$.\\
\textbf{Paralelo} com o de \SI{4}{\micro\farad}:
$4 + 4 = \SI{8}{\micro\farad}$.}
\end{questao}

\begin{questao}[2]{}
Um circuito RC série tem $R = \SI{1}{\kilo\ohm}$ e $C = \SI{1}{\micro\farad}$,
ligado a uma fonte de $\SI{10}{\volt}$ em $t = 0$.
\begin{itensq}
\item Qual a constante de tempo?
\item Qual a tensão no capacitor após um intervalo igual a $\tau$?
\item Qual a tensão em regime permanente?
\end{itensq}
\resposta{(a) $\tau = \SI{1}{\milli\second}$; (b) $\approx\SI{6.32}{\volt}$;
(c) \SI{10}{\volt}}
\resolucao{(a) $\tau = RC = \pot{1}{3}\cdot\pot{1}{-6} = \SI{1}{\milli\second}$.\\
(b) Na carga, $v_C(\tau) = U_0(1 - e^{-1}) = 10\cdot0{,}632 \approx \SI{6.32}{\volt}$
--- o capacitor atinge cerca de $\SI{63}{\percent}$ da tensão final em um $\tau$.\\
(c) Em regime permanente ($t\to\infty$), o capacitor está totalmente carregado,
não circula corrente, e $v_C = U_0 = \SI{10}{\volt}$: o capacitor se comporta como
um circuito \emph{aberto}.}
\end{questao}

\begin{questao}[2]{}
No circuito da figura, determine a tensão no capacitor em regime permanente.

\circuitoq{%
\begin{circuitikz}[scale=1,transform shape,line width=0.8pt]
 \draw (0,0) to[battery1,l_={$\SI{12}{\volt}$}] (0,3)
       to[R,l={$R_1=\SI{4}{\ohm}$}] (3,3) coordinate(a);
 \draw (a) to[R,l={$R_2=\SI{2}{\ohm}$}] (3,0) coordinate(b) -- (0,0);
 \draw (a) -- (5,3) to[C,l={$C$}] (5,0) -- (b);
 \node[right,Vcol,font=\footnotesize] at (5.1,1.5) {$v_C$};
\end{circuitikz}}

\resposta{$v_C = \SI{4}{\volt}$}
\resolucao{Em regime permanente, o capacitor está carregado e \emph{não conduz}
(ramo aberto). Toda a corrente passa por $R_1$ e $R_2$ em série:
\[
I=\frac{12}{4+2}=\SI{2}{\ampere}.
\]
A tensão no capacitor iguala a tensão sobre $R_2$ (estão em paralelo):
\[
v_C = R_2\,I = 2\cdot2 = \SI{4}{\volt}.
\]
\textbf{Princípio-chave:} em regime permanente CC, capacitor = circuito aberto
(e indutor = curto-circuito). Isso transforma o problema em uma simples análise
resistiva.}
\end{questao}

\begin{questao}[2]{}
Três capacitores de $\SI{6}{\micro\farad}$ cada são associados. Determine a
capacitância equivalente se estiverem ligados (a) todos em série; (b) todos em
paralelo; (c) dois em paralelo, e o conjunto em série com o terceiro.
\resposta{(a) \SI{2}{\micro\farad}; (b) \SI{18}{\micro\farad}; (c) \SI{4}{\micro\farad}}
\resolucao{(a) Série de três iguais: $\Ceq = \dfrac{C}{3} = \SI{2}{\micro\farad}$.\\
(b) Paralelo de três iguais: $\Ceq = 3C = \SI{18}{\micro\farad}$.\\
(c) Dois em paralelo: $6+6 = \SI{12}{\micro\farad}$; em série com o terceiro:
$\dfrac{12\cdot6}{18} = \SI{4}{\micro\farad}$.\\
Repare como a topologia muda drasticamente o resultado, de \SI{2}{} a
\SI{18}{\micro\farad} com os mesmos três componentes.}
\end{questao}

\blocoaprofundamento

\begin{questao}[3]{}
Um capacitor de $\SI{2}{\micro\farad}$ está carregado a $\SI{6}{\volt}$ e é
descarregado através de um resistor de $\SI{500}{\ohm}$.
\begin{itensq}
\item Qual a constante de tempo da descarga?
\item Qual a corrente inicial de descarga?
\item Qual a carga restante no capacitor após um intervalo igual a $\tau$?
\end{itensq}
\resposta{(a) $\tau = \SI{1}{\milli\second}$; (b) $I_0 = \SI{12}{\milli\ampere}$;
(c) $\approx\SI{4.4}{\micro\coulomb}$}
\resolucao{(a) $\tau = RC = 500\cdot\pot{2}{-6} = \SI{1}{\milli\second}$.\\
(b) A corrente inicial é limitada apenas pelo resistor:
$I_0 = \dfrac{U_0}{R} = \dfrac{6}{500} = \SI{12}{\milli\ampere}$.\\
(c) A carga inicial é $Q_0 = C\,U_0 = \pot{2}{-6}\cdot6 = \SI{12}{\micro\coulomb}$.
Após um $\tau$, resta $Q(\tau) = Q_0\,e^{-1} = 12\cdot0{,}368
\approx\SI{4.4}{\micro\coulomb}$ --- cerca de $\SI{37}{\percent}$ do inicial.}
\end{questao}

\begin{questao}[3]{}
O gráfico mostra a tensão $v_C(t)$ sobre um capacitor durante a carga em um
circuito RC alimentado por $\SI{10}{\volt}$.

\circuitoqq{%
\begin{tikzpicture}
 \begin{axis}[width=9.5cm,height=5cm,xlabel={$t$ (ms)},ylabel={$v_C$ (V)},
   xmin=0,xmax=5,ymin=0,ymax=11,xtick={0,1,2,3,4,5},ytick={0,6.32,10},
   yticklabels={0,6.32,10},grid=both,axis lines=left,domain=0:5,samples=100]
   \addplot[Ccol,very thick] {10*(1-exp(-x))};
   \addplot[dashed,cinza] coordinates {(1,0)(1,6.32)};
   \addplot[dashed,cinza] coordinates {(0,6.32)(1,6.32)};
   \node[Ccol,anchor=west,font=\footnotesize] at (axis cs:1.1,6.0) {$(\tau,\,0{,}63\,U_0)$};
 \end{axis}
\end{tikzpicture}}{Carga de um capacitor em circuito RC.}

Sabendo que a tensão atinge $\SI{6.32}{\volt}$ em $t = \SI{1}{\milli\second}$,
determine a constante de tempo e a tensão final.
\resposta{$\tau = \SI{1}{\milli\second}$; $U_0 = \SI{10}{\volt}$}
\resolucao{O valor $\SI{6.32}{\volt}$ corresponde a $\SI{63}{\percent}$ da tensão
final --- por definição, esse é o instante $t = \tau$. Logo
$\tau = \SI{1}{\milli\second}$. A curva satisfaz $v_C = U_0(1 - e^{-t/\tau})$, e
sua assíntota horizontal é a tensão da fonte, $U_0 = \SI{10}{\volt}$. Lendo o
gráfico: a $\SI{63}{\percent}$ em $\SI{1}{\milli\second}$ e saturando em
\SI{10}{\volt}, confirma-se $\tau = \SI{1}{\milli\second}$.}
\end{questao}

\begin{questao}[3]{}
Um capacitor $C_1 = \SI{10}{\micro\farad}$ é carregado a $\SI{100}{\volt}$ e
depois conectado, por meio de uma chave, a um capacitor descarregado
$C_2 = \SI{5}{\micro\farad}$.

\circuitoq{%
\begin{circuitikz}[scale=1,transform shape,line width=0.8pt]
 \draw (0,0) to[C,l_={$C_1$}] (0,2.6) -- (1.6,2.6) coordinate(a);
 \draw (a) to[cspst,l={S}] (3.4,2.6) -- (4.6,2.6);
 \draw (4.6,2.6) to[C,l={$C_2$}] (4.6,0) -- (0,0);
 \node[font=\footnotesize] at (0,-0.45) {carregado a \SI{100}{\volt}};
 \node[font=\footnotesize] at (4.6,-0.45) {descarregado};
\end{circuitikz}}

Após fechar a chave:
\begin{itensq}
\item Determine a tensão final comum aos dois capacitores.
\item Calcule a energia armazenada antes e depois.
\item Explique o "paradoxo": para onde foi a energia que desapareceu?
\end{itensq}
\resposta{(a) $V_f \approx \SI{66.7}{\volt}$; (b) \SI{50}{\milli\joule} $\to$
\SI{33.3}{\milli\joule}; (c) ver comentário}
\resolucao{\textbf{(a)} A carga total se \emph{conserva} e se redistribui até que
as tensões se igualem:
\[
Q_0 = C_1V_0=(\pot{10}{-6})(100)=\SI{1}{\milli\coulomb},
\]
\[
V_f=\frac{Q_0}{C_1+C_2}=\frac{\pot{1}{-3}}{\pot{15}{-6}}
=\frac{200}{3}\approx\SI{66.7}{\volt}.
\]
\textbf{(b) Energias:}
\[
E_i=\tfrac12 C_1V_0^{\,2}=\tfrac12(\pot{10}{-6})(100)^2=\SI{50}{\milli\joule},
\]
\[
E_f=\tfrac12 (C_1+C_2)V_f^{\,2}
=\tfrac12(\pot{15}{-6})\left(\tfrac{200}{3}\right)^{2}
=\frac{100}{3}\approx\SI{33.3}{\milli\joule}.
\]
Perderam-se $\SI{16.7}{\milli\joule}$ --- exatamente $\dfrac{1}{3}$ da energia
inicial.

\textbf{(c) Resolução do paradoxo.} A carga se conserva, mas a energia
\textbf{não} precisa se conservar \emph{dentro do subsistema dos capacitores} ---
ela foi \emph{dissipada}. Os mecanismos:
\begin{itemize}[leftmargin=1.6em,itemsep=1pt,topsep=2pt]
\item \textbf{Resistência dos fios e da chave:} por menor que seja, a corrente
      transitória a atravessa e gera calor $\int Ri^2\,\mathrm{d}t$. O
      resultado notável é que \emph{a energia perdida independe do valor de $R$}
      --- só o \emph{tempo} do transitório muda.
\item \textbf{Radiação eletromagnética:} no limite $R\to0$, o circuito LC
      oscilaria e irradiaria a energia como ondas.
\end{itemize}
\textbf{Fórmula geral da perda:}
\[
\Delta E=\frac{1}{2}\,\frac{C_1C_2}{C_1+C_2}\,V_0^{\,2}
=\frac12\cdot\frac{10\cdot5}{15}\cdot10^{-6}\cdot 100^2
\approx\SI{16.7}{\milli\joule}.\ \checkmark
\]
\textbf{Relevância prática:} esse é o motivo pelo qual bancos de capacitores nunca
são conectados em paralelo sem resistores de pré-carga (\emph{soft-start}) ---
a corrente de equalização pode ser destrutiva.}
\end{questao}

\begin{questao}[3]{}
Um capacitor de placas paralelas com $C_0 = \SI{10}{\micro\farad}$ é carregado a
$U_0 = \SI{100}{\volt}$ e então \textbf{desconectado} da fonte. Em seguida,
introduz-se entre as placas um dielétrico de constante $\kappa = 4$.
\begin{itensq}
\item O que acontece com a carga, a capacitância, a tensão e a energia?
\item Calcule os valores finais de cada grandeza.
\item Para onde foi a energia "perdida"? Ela violou a conservação de energia?
\item E se o dielétrico fosse inserido com o capacitor \emph{ainda ligado} à
      fonte? Compare os dois casos.
\end{itensq}
\resposta{(b) $Q=\SI{1}{\milli\coulomb}$ (constante), $C=\SI{40}{\micro\farad}$,
$U=\SI{25}{\volt}$, $E$ cai de \SI{50}{\milli\joule} para \SI{12.5}{\milli\joule};
(c) e (d) ver comentário}
\resolucao{\textbf{(a) e (b) --- capacitor isolado ($Q$ constante).}
Desconectado, a carga não tem para onde ir:
\[
Q = C_0U_0 = \pot{10}{-6}\cdot100 = \SI{1}{\milli\coulomb}\ \text{(constante)}.
\]
O dielétrico multiplica a capacitância por $\kappa$:
\[
C = \kappa C_0 = \SI{40}{\micro\farad},
\qquad
U = \frac{Q}{C}=\frac{\pot{1}{-3}}{\pot{40}{-6}}=\SI{25}{\volt}
\quad\left(=\frac{U_0}{\kappa}\right).
\]
Energias:
\[
E_i=\tfrac12 C_0U_0^{\,2}=\SI{50}{\milli\joule},
\qquad
E_f=\tfrac12 CU^{2}=\tfrac12(\pot{40}{-6})(25)^2=\SI{12.5}{\milli\joule}
=\frac{E_i}{\kappa}.
\]
\textbf{(c) Para onde foi a energia?} \emph{Não houve violação alguma.} O
dielétrico é \textbf{atraído} para dentro do capacitor pelo campo de borda ---
existe uma força elétrica real puxando-o. Se ele for solto, essa força realiza
trabalho e o acelera; a energia "perdida" pelo campo converte-se em energia
cinética do dielétrico (e, ao final, em calor por atrito e vibração).
Para inseri-lo \emph{quase-estaticamente}, um agente externo precisa segurá-lo,
e é ele quem \emph{recebe} os \SI{37.5}{\milli\joule}.

\textbf{(d) Caso alternativo --- capacitor ligado à fonte ($U$ constante).}
Agora a tensão é fixada em \SI{100}{\volt}:
\[
C = \SI{40}{\micro\farad},\qquad
Q = CU = \SI{4}{\milli\coulomb}\ (\text{quadruplicou}),\qquad
E_f=\tfrac12 CU^{2}=\SI{200}{\milli\joule}\ (\text{quadruplicou}).
\]
A energia armazenada \emph{aumentou} \SI{150}{\milli\joule}. A fonte, porém,
forneceu $\Delta Q\cdot U = (\pot{3}{-3})(100)=\SI{300}{\milli\joule}$ --- o
dobro. A metade restante (\SI{150}{\milli\joule}) foi entregue como trabalho
mecânico ao dielétrico (que também é sugado para dentro).

\smallskip
\begin{center}\small\sffamily
\begin{tabular}{@{}l c c@{}}
\toprule
 & \textbf{Isolado ($Q$ const.)} & \textbf{Ligado ($U$ const.)}\\
\midrule
Carga & constante & $\times\kappa$\\
Tensão & $\div\kappa$ & constante\\
Energia & $\div\kappa$ (diminui) & $\times\kappa$ (aumenta)\\
\bottomrule
\end{tabular}
\end{center}
\textbf{Moral:} a mesma ação física (inserir o dielétrico) produz resultados
\emph{opostos} para a energia conforme a condição de contorno. Identificar
\emph{qual grandeza permanece constante} é o passo decisivo do problema.}
\end{questao}

% BEGIN BANCO COMPLEMENTAR Capacitores
% =====================================================================
%  Banco complementar --- Capacitores  (REVISADO)
%  Substitui a versao gerada por template (10 clones em N2, 11 em N3,
%  10 questoes conceituais indevidamente marcadas como N4).
%  Sem sobreposicao com o cap10, que ja traz: C de placas paralelas,
%  Q=CV com W, serie 6/3, paralelo 4/12, energia (MC), associacao mista
%  6||3, RC de 1 kohm/1 uF, tres de 6 uF, descarga de 2 uF/6 V, grafico
%  v_C(t), COMPARTILHAMENTO DE CARGA entre C1 e C2 e INSERCAO DE
%  DIELETRICO (a Q e a V constantes) --- estes dois ultimos motivaram a
%  remocao de duas N4 do banco antigo, que os duplicavam.
%  Distribuicao mantida: 4 N1 + 10 N2 + 11 N3 + 10 N4 = 35
% =====================================================================

\subsubsection*{Banco complementar --- Capacitores}

\begin{atencao}[Organização do banco]
Três relações organizam o tópico: $Q=CU$, $W=\frac12CU^{2}$ e
$i=C\dfrac{\mathrm{d}u}{\mathrm{d}t}$. Delas decorre o fato central dos
transitórios --- a \textbf{tensão} no capacitor não pode saltar (isso exigiria
corrente infinita), ao contrário do indutor, em que é a \emph{corrente} que é
contínua. O N3 trata de divisão de tensão em série, fuga, ESR e extração de
parâmetros; o N4 exige dedução da EDO, balanço energético, projeto de banco,
filtro e temporizador.
\end{atencao}

\blocofixacao

\begin{questao}[1]{}
Um capacitor armazena $\SI{300}{\micro\coulomb}$ sob $\SI{20}{\volt}$.
Determine sua capacitância.
\resposta{$C=\SI{15}{\micro\farad}$}
\resolucao{
\[
C=\frac{Q}{U}=\frac{\pot{300}{-6}}{20}=\pot{15}{-6}\,\si{\farad}
=\SI{15}{\micro\farad}.
\]
A capacitância é uma propriedade \emph{geométrica} do componente: ela não
depende de $Q$ nem de $U$, apenas da razão entre eles. Aplicar o dobro da
tensão dobra a carga e deixa $C$ inalterada.}
\end{questao}

\begin{questao}[1]{}
Um capacitor de $\SI{100}{\micro\farad}$ armazena $\SI{8}{\milli\joule}$.
Determine a tensão em seus terminais.
\resposta{$U\approx\SI{12.6}{\volt}$}
\resolucao{De $W=\frac12CU^{2}$,
\[
U=\sqrt{\frac{2W}{C}}=\sqrt{\frac{2(\pot{8}{-3})}{\pot{100}{-6}}}
=\sqrt{160}=\SI{12.65}{\volt}.
\]
A raiz quadrada indica que a energia é pouco sensível à tensão em valores
baixos e muito sensível em valores altos: dobrar a tensão \emph{quadruplica}
a energia armazenada.}
\end{questao}

\begin{questao}[1]{}
Sobre a tensão em um capacitor ideal, é correto afirmar que ela:
\begin{alternativas}[2]
\item pode variar instantaneamente
\item não pode variar instantaneamente
\item é sempre igual à da fonte
\item é sempre nula em regime permanente CC
\end{alternativas}
\resposta{(b)}
\resolucao{Um salto de tensão implicaria
$\mathrm{d}u/\mathrm{d}t\to\infty$ e, por
$i=C\,\mathrm{d}u/\mathrm{d}t$, corrente infinita --- impossível. Logo
$u(0^{+})=u(0^{-})$.\\
\textbf{Dualidade com o indutor:} no capacitor é a \emph{tensão} que é
contínua; no indutor, a \emph{corrente}. Em regime CC o capacitor tem corrente
nula (comporta-se como circuito aberto) e tensão em geral não nula --- o que
descarta a alternativa (d).}
\end{questao}

\begin{questao}[1]{}
Determine a constante de tempo de um circuito RC com
$R=\SI{4.7}{\kilo\ohm}$ e $C=\SI{10}{\micro\farad}$.
\resposta{$\tau=\SI{47}{\milli\second}$}
\resolucao{
\[
\tau=RC=4700\cdot\pot{10}{-6}=\pot{4.7}{-2}\,\si{\second}
=\SI{47}{\milli\second}.
\]
\textbf{Atenção à estrutura:} no RC, resistência \emph{maior} torna o
transitório \emph{mais lento} ($\tau=RC$); no RL ocorre o contrário
($\tau=L/R$). Verificar a dimensão ajuda:
$\si{\ohm}\times\si{\farad}=\si{\second}$.}
\end{questao}

\blocoaplicacao

\begin{questao}[2]{}
Dois capacitores de $\SI{10}{\micro\farad}$ em série são associados em paralelo
com um de $\SI{15}{\micro\farad}$. Determine $C_{eq}$.
\resposta{$C_{eq}=\SI{20}{\micro\farad}$}
\resolucao{\textbf{Série} (regra inversa, ao contrário dos resistores):
\[
C_s=\frac{10\cdot10}{20}=\SI{5}{\micro\farad}.
\]
\textbf{Paralelo} (soma direta):
\[
C_{eq}=5+15=\SI{20}{\micro\farad}.
\]
Memorize a inversão: capacitores em \emph{série} combinam como resistores em
paralelo, e vice-versa. A razão física é que a série reduz a capacitância
(equivale a afastar as placas), enquanto o paralelo a aumenta (equivale a somar
áreas).}
\end{questao}

\begin{questao}[2]{}
Capacitores de $\SI{1}{\micro\farad}$ e $\SI{2}{\micro\farad}$ estão em série
sob $\SI{30}{\volt}$. Determine $C_{eq}$, a carga e a tensão em cada um.
\resposta{$C_{eq}=\SI{0.667}{\micro\farad}$; $Q=\SI{20}{\micro\coulomb}$;
$U_1=\SI{20}{\volt}$, $U_2=\SI{10}{\volt}$}
\resolucao{
\[
C_{eq}=\frac{1\cdot2}{3}=\SI{0.667}{\micro\farad};
\qquad
Q=C_{eq}U=0{,}667\cdot30=\SI{20}{\micro\coulomb}.
\]
Em série, a \textbf{carga é a mesma} em todos (é a mesma corrente que os
carregou):
\[
U_1=\frac{20}{1}=\SI{20}{\volt},
\qquad
U_2=\frac{20}{2}=\SI{10}{\volt}.
\]
\textbf{Verificação:} $20+10=\SI{30}{\volt}$ \checkmark.\\
A tensão se reparte \emph{inversamente} à capacitância --- o capacitor
\textbf{menor} fica com a maior tensão. É o oposto da série de resistores, e é
a causa de um problema de projeto tratado adiante.}
\end{questao}

\begin{questao}[2]{}
A tensão sobre um capacitor de $\SI{10}{\micro\farad}$ cresce linearmente à
taxa de $\SI{500}{\volt\per\second}$. Determine a corrente.
\resposta{$i=\SI{5}{\milli\ampere}$}
\resolucao{
\[
i=C\frac{\mathrm{d}u}{\mathrm{d}t}=\pot{10}{-6}\cdot500
=\pot{5}{-3}\,\si{\ampere}=\SI{5}{\milli\ampere}.
\]
Enquanto a rampa durar, a corrente é \emph{constante}. Se a tensão parasse de
variar, a corrente cairia a zero instantaneamente --- o capacitor responde à
\textbf{variação} da tensão, não ao seu valor.}
\end{questao}

\begin{questao}[2]{}
Aplica-se a um capacitor de $\SI{2.2}{\micro\farad}$ uma tensão
\textbf{triangular} que sobe $\SI{10}{\volt}$ em $\SI{1}{\milli\second}$ e
desce igualmente no milissegundo seguinte. Descreva e quantifique a corrente.
\resposta{Onda \textbf{quadrada} de $\pm\SI{22}{\milli\ampere}$}
\resolucao{Na subida,
$\mathrm{d}u/\mathrm{d}t=10/10^{-3}=\SI{10000}{\volt\per\second}$:
\[
i=2{,}2\times10^{-6}\cdot10^{4}=\SI{22}{\milli\ampere}.
\]
Na descida a taxa inverte e $i=\SI{-22}{\milli\ampere}$.\\
\textbf{Resultado geral:} o capacitor \emph{deriva} a tensão --- triangular
entra, quadrada sai. O indutor faz o oposto (integra): tensão quadrada nele
produz corrente triangular. É a dualidade dos dois elementos em sua forma mais
visível no osciloscópio.}
\end{questao}

\begin{questao}[2]{}
Um circuito RC com $\tau=\SI{2}{\milli\second}$ é ligado a
$\SI{20}{\volt}$ com o capacitor inicialmente descarregado. Determine
$u_C$ em $t=\tau$, $2\tau$ e $3\tau$.
\resposta{$\SI{12.64}{\volt}$, $\SI{17.29}{\volt}$ e $\SI{19.00}{\volt}$}
\resolucao{Com $u_C=U\left(1-e^{-t/\tau}\right)$ e $U=\SI{20}{\volt}$:
\[
u(\tau)=20(0{,}632)=\SI{12.64}{\volt};
\quad
u(2\tau)=20(0{,}865)=\SI{17.29}{\volt};
\quad
u(3\tau)=20(0{,}950)=\SI{19.00}{\volt}.
\]
As frações $63{,}2\%$, $86{,}5\%$ e $95{,}0\%$ são as mesmas do circuito RL ---
elas dependem só da forma exponencial, não do tipo de elemento.}
\end{questao}

\begin{questao}[2]{}
No mesmo circuito, determine o instante em que a tensão atinge
$\SI{15}{\volt}$.
\resposta{$t\approx\SI{2.77}{\milli\second}$}
\resolucao{
\[
15=20\left(1-e^{-t/\tau}\right)
\;\Longrightarrow\;
e^{-t/\tau}=0{,}25
\;\Longrightarrow\;
t=\tau\ln4=2(1{,}386)=\SI{2.77}{\milli\second}.
\]
\textbf{Atalho útil:} atingir $75\%$ leva exatamente $\ln4\approx1{,}39\tau$;
atingir $50\%$ leva $\ln2\approx0{,}69\tau$. Esse último valor é a "meia-vida"
do transitório e vale a pena memorizar.}
\end{questao}

\begin{questao}[2]{}
Uma fonte de $\SI{24}{\volt}$ alimenta $R_1=\SI{8}{\kilo\ohm}$ em série com
$R_2=\SI{4}{\kilo\ohm}$. Um capacitor está ligado em paralelo com $R_2$.
Determine a tensão no capacitor em regime permanente e a corrente que o
atravessa.
\resposta{$U_C=\SI{8}{\volt}$; $i_C=0$}
\resolucao{Em regime CC o capacitor é um \textbf{circuito aberto}: não conduz.
Toda a corrente passa pelos dois resistores, e a tensão do capacitor é a de
$R_2$:
\[
U_C=24\cdot\frac{4}{8+4}=\SI{8}{\volt};
\qquad
i_C=0.
\]
\textbf{Método:} para determinar o estado final de qualquer circuito CC com
capacitores, basta \emph{removê-los} do esquema e resolver a rede resistiva
restante. Para indutores, o procedimento dual é substituí-los por curtos.}
\end{questao}

\begin{questao}[2]{}
Dois capacitores armazenam a mesma energia. O primeiro tem
$C_1=\SI{100}{\micro\farad}$ sob $\SI{10}{\volt}$. Se
$C_2=\SI{25}{\micro\farad}$, determine sua tensão.
\resposta{$U_2=\SI{20}{\volt}$}
\resolucao{
\[
\frac12C_1U_1^{2}=\frac12C_2U_2^{2}
\;\Longrightarrow\;
U_2=U_1\sqrt{\frac{C_1}{C_2}}=10\sqrt{4}=\SI{20}{\volt}.
\]
Ambos armazenam $\SI{5}{\milli\joule}$. Reduzir $C$ por um fator $4$ exige
\emph{dobrar} a tensão --- e a tensão é justamente o parâmetro caro e limitado
em capacitores reais, pois exige dielétrico mais espesso, o que reduz $C$.}
\end{questao}

\begin{questao}[2]{}
Um capacitor plano tem capacitância $C$. Determine a nova capacitância se:
\begin{itensq}
\item a área das placas for dobrada;
\item a distância entre placas for reduzida à metade;
\item ambas as alterações forem feitas.
\end{itensq}
\resposta{(a) $2C$; (b) $2C$; (c) $4C$}
\resolucao{De $C=\dfrac{\varepsilon A}{d}$: $C\propto A$ e
$C\propto1/d$.\\
(a) $A\to2A$ dá $2C$. (b) $d\to d/2$ dá $2C$. (c) os efeitos se multiplicam:
$4C$.\\
\textbf{Limite prático de (b):} reduzir $d$ aumenta o campo
$E=U/d$ para a mesma tensão. Chega-se rapidamente à rigidez dielétrica do
material, e o capacitor perfura. Há, portanto, um teto para essa estratégia ---
motivo pelo qual capacitores de alta capacitância usam grandes áreas
(folhas enroladas) em vez de dielétricos ultrafinos.}
\end{questao}

\begin{questao}[2]{}
Um capacitor carregado a $\SI{50}{\volt}$ é descarregado através de
$R=\SI{2.2}{\kilo\ohm}$. Determine a corrente inicial e explique sua evolução.
\resposta{$i_0=\SI{22.7}{\milli\ampere}$, decaindo exponencialmente}
\resolucao{No instante inicial a tensão ainda é $\SI{50}{\volt}$
(continuidade), e ela se aplica integralmente ao resistor:
\[
i_0=\frac{50}{2200}=\SI{22.7}{\milli\ampere}.
\]
À medida que o capacitor se descarrega, sua tensão cai e a corrente a acompanha:
$i(t)=i_0e^{-t/\tau}$.\\
\textbf{Observação:} a corrente é máxima no início e a tensão é máxima no
início --- diferentemente do circuito RL, em que a corrente parte de zero. Essa
diferença decorre justamente de qual grandeza é contínua em cada elemento.}
\end{questao}

\blocoaprofundamento

\begin{questao}[3]{}
$C_1=\SI{4}{\micro\farad}$ e $C_2=\SI{12}{\micro\farad}$, ambos com tensão
máxima de trabalho de $\SI{50}{\volt}$, são ligados em série.
\begin{itensq}
\item Determine $C_{eq}$ e a distribuição de tensão sob $\SI{100}{\volt}$.
\item Determine a máxima tensão que a associação suporta.
\item Explique por que o resultado é contraintuitivo.
\end{itensq}
\resposta{(a) $\SI{3}{\micro\farad}$; $\SI{75}{\volt}$ e $\SI{25}{\volt}$;
(b) $\SI{66.7}{\volt}$}
\resolucao{(a) $C_{eq}=\dfrac{4\cdot12}{16}=\SI{3}{\micro\farad}$. Como a carga
é comum, $U\propto1/C$:
\[
U_1=100\cdot\frac{12}{16}=\SI{75}{\volt},
\qquad
U_2=100\cdot\frac{4}{16}=\SI{25}{\volt}.
\]
(b) Sob $\SI{100}{\volt}$, $C_1$ já está a $\SI{75}{\volt}$ --- $50\%$ acima do
limite. Impondo $U_1\le\SI{50}{\volt}$:
\[
U_{\max}=50\cdot\frac{16}{12}=\SI{66.7}{\volt}.
\]
(c) \textbf{É contraintuitivo} porque associar dois componentes de
$\SI{50}{\volt}$ em série sugere $\SI{100}{\volt}$ de capacidade --- e isso
\emph{só} é verdade se eles forem \textbf{iguais}. Com valores diferentes, o
menor capacitor absorve a maior parte da tensão e limita todo o conjunto.\\
\textbf{Regra prática:} em série de capacitores, use sempre unidades idênticas
--- e ainda assim adicione resistores de equalização, como se detalha no bloco
de desafio.}
\end{questao}

\begin{questao}[3]{}
Capacitores de $\SI{2}{\micro\farad}$, $\SI{3}{\micro\farad}$ e
$\SI{5}{\micro\farad}$ estão em paralelo sob $\SI{20}{\volt}$. Determine a
carga de cada um, a carga total e $C_{eq}$, e compare com o caso série.
\resposta{$40$, $60$ e $\SI{100}{\micro\coulomb}$; total
$\SI{200}{\micro\coulomb}$; $C_{eq}=\SI{10}{\micro\farad}$}
\resolucao{Em paralelo a \textbf{tensão} é comum:
\[
Q_k=C_kU \Rightarrow 40,\ 60,\ \SI{100}{\micro\coulomb};
\qquad
Q_{tot}=\SI{200}{\micro\coulomb}.
\]
\[
C_{eq}=2+3+5=\SI{10}{\micro\farad}
\quad\text{(confere: }200/20=10\text{)}.
\]
\textbf{Comparação com a série:} lá, a \emph{carga} é comum e a tensão se
divide; aqui, a \emph{tensão} é comum e a carga se divide. O paralelo soma
capacitâncias e mantém a tensão de trabalho da menor unidade; a série reduz a
capacitância e (idealmente) soma as tensões. São estratégias opostas para
problemas opostos.}
\end{questao}

\begin{questao}[3]{}
Uma fonte de $\SI{12}{\volt}$ alimenta $R_1=\SI{6}{\kilo\ohm}$ em série com
$R_2=\SI{3}{\kilo\ohm}$, e um capacitor de $\SI{100}{\nano\farad}$ está em
paralelo com $R_2$.
\begin{itensq}
\item Determine a tensão final no capacitor.
\item Determine $\tau$ pelo equivalente de Thévenin.
\item Escreva $u_C(t)$.
\end{itensq}
\resposta{(a) $\SI{4}{\volt}$; (b) $\tau=\SI{0.2}{\milli\second}$;
(c) $u_C=4\left(1-e^{-t/0{,}2\,\mathrm{ms}}\right)$}
\resolucao{(a) Com o capacitor aberto em regime,
$U_C=12\cdot\dfrac{3}{9}=\SI{4}{\volt}$.\\
(b) Zerando a fonte, o capacitor vê $R_1\parallel R_2$:
\[
R_{Th}=\frac{6\cdot3}{9}=\SI{2}{\kilo\ohm}
\;\Longrightarrow\;
\tau=R_{Th}C=2000\cdot\pot{100}{-9}=\SI{0.2}{\milli\second}.
\]
(c) $u_C(t)=4\left(1-e^{-t/0{,}0002}\right)$ volts.\\
\textbf{Erro comum:} usar $\tau=R_1C$ ou $\tau=R_2C$. A constante de tempo é
governada pela resistência \emph{vista pelo capacitor} com as fontes zeradas
--- aqui, o paralelo das duas, que é menor que qualquer uma delas. O transitório
é, portanto, mais rápido do que qualquer estimativa ingênua sugeriria.}
\end{questao}

\begin{questao}[3]{}
Um capacitor já está a $\SI{2}{\volt}$ quando é ligado a um circuito que o
levaria a $\SI{10}{\volt}$, com $\tau=\SI{5}{\milli\second}$.
\begin{itensq}
\item Escreva $u_C(t)$.
\item Determine a tensão em $t=\SI{5}{\milli\second}$.
\item Determine a corrente inicial, sabendo que $R=\SI{5}{\kilo\ohm}$.
\end{itensq}
\resposta{(a) $u_C=10-8e^{-t/\tau}$; (b) $\SI{7.06}{\volt}$;
(c) $\SI{1.6}{\milli\ampere}$}
\resolucao{(a) A forma geral de qualquer transitório de primeira ordem é
\[
u(t)=U_f+\left(U_0-U_f\right)e^{-t/\tau}=10-8e^{-t/\tau}.
\]
(b) $u(5\,\mathrm{ms})=10-8e^{-1}=10-2{,}94=\SI{7.06}{\volt}$.\\
(c) Em $t=0^{+}$ a tensão do capacitor é $\SI{2}{\volt}$, e sobre o resistor
sobram $10-2=\SI{8}{\volt}$:
\[
i_0=\frac{8}{5000}=\SI{1.6}{\milli\ampere}.
\]
\textbf{Vantagem da forma geral:} ela cobre carga ($U_0=0$), descarga
($U_f=0$) e qualquer condição intermediária. Basta identificar os três
parâmetros $U_0$, $U_f$ e $\tau$ --- não há três fórmulas a memorizar.}
\end{questao}

\begin{questao}[3]{}
Um capacitor de $\SI{470}{\micro\farad}$ tem resistência de isolação de
$\SI{2}{\mega\ohm}$. Carregado a $\SI{100}{\volt}$ e deixado em circuito
aberto:
\begin{itensq}
\item Determine a constante de tempo de autodescarga.
\item Determine a tensão após $\SI{15}{\minute}$.
\item Comente a implicação de segurança.
\end{itensq}
\resposta{(a) $\tau=\SI{940}{\second}\approx\SI{15.7}{\minute}$;
(b) $\approx\SI{38}{\volt}$}
\resolucao{(a) A fuga é modelada por um resistor em paralelo:
\[
\tau=R_{iso}C=\pot{2}{6}\cdot\pot{470}{-6}=\SI{940}{\second}.
\]
(b) Com $t=\SI{900}{\second}$:
\[
u=100\,e^{-900/940}=100\,e^{-0{,}957}=\SI{38.4}{\volt}.
\]
(c) \textbf{Depois de quinze minutos desligado, o capacitor ainda guarda
$\SI{38}{\volt}$} --- e, em bancos de fonte chaveada com centenas de volts, o
resíduo é letal por horas. É por isso que equipamentos com barramento CC trazem
\emph{resistores de descarga} (bleeders) permanentes e o aviso de aguardar
antes de abrir o gabinete.\\
\textbf{Nota:} capacitores eletrolíticos ainda apresentam \emph{absorção
dielétrica} --- após descarga completa, a tensão "ressurge" parcialmente em
minutos. Curto-circuitar antes de manusear é procedimento obrigatório, não
precaução exagerada.}
\end{questao}

\begin{questao}[3]{}
Um capacitor eletrolítico tem ESR (resistência equivalente em série) de
$\SI{0.1}{\ohm}$ e conduz corrente de ondulação de
$\SI{1.5}{\ampere}$ eficazes.
\begin{itensq}
\item Determine a potência dissipada internamente.
\item Explique por que essa dissipação limita a vida do componente.
\item Indique duas formas de reduzi-la.
\end{itensq}
\resposta{(a) $P=\SI{0.225}{\watt}$}
\resolucao{(a)
\[
P=\mathrm{ESR}\cdot I_{rms}^{2}=0{,}1(1{,}5)^{2}=\SI{0.225}{\watt}.
\]
(b) Essa potência é dissipada \emph{dentro} do capacitor, elevando a
temperatura do eletrólito. A vida útil de eletrolíticos segue uma regra
empírica severa: ela \textbf{cai pela metade a cada $\SI{10}{\celsius}$} de
elevação. Um componente de $\SI{5000}{\hour}$ a $\SI{105}{\celsius}$ dura
$\SI{20000}{\hour}$ a $\SI{85}{\celsius}$ --- e o inverso também vale.
Pior: o eletrólito evapora, a ESR \emph{aumenta}, a dissipação cresce e o
processo se realimenta.\\
(c) \textbf{(i) Dividir a corrente:} usar $n$ capacitores em paralelo reduz a
corrente por unidade a $I/n$ e a dissipação individual a $P/n^{2}$.
Com dois, cada um dissipa $\SI{56}{\milli\watt}$ --- um quarto do original.\\
\textbf{(ii) Escolher tecnologia de baixa ESR:} capacitores de polímero sólido
têm ESR uma ordem de grandeza menor, ao custo de capacitância e tensão
menores.\\
\textbf{Consequência de diagnóstico:} em fontes chaveadas, o capacitor de saída
é quase sempre o primeiro componente a falhar --- e a ESR crescente, não a
perda de capacitância, é o sintoma que primeiro aparece.}
\end{questao}

\begin{questao}[3]{}
Um capacitor de $\SI{100}{\micro\farad}$, inicialmente descarregado, é
carregado a $\SI{20}{\volt}$ através de um resistor.
\begin{itensq}
\item Determine a energia fornecida pela fonte.
\item Determine a energia armazenada no capacitor.
\item Determine a energia dissipada no resistor e comente.
\end{itensq}
\resposta{(a) $\SI{40}{\milli\joule}$; (b) $\SI{20}{\milli\joule}$;
(c) $\SI{20}{\milli\joule}$ --- exatamente metade}
\resolucao{(a) A fonte mantém $\SI{20}{\volt}$ enquanto entrega toda a carga
final $Q=CU=\SI{2}{\milli\coulomb}$:
\[
W_{fonte}=QU=CU^{2}=\pot{100}{-6}(400)=\SI{40}{\milli\joule}.
\]
(b) $W_C=\frac12CU^{2}=\SI{20}{\milli\joule}$.\\
(c) Por conservação, $W_R=40-20=\SI{20}{\milli\joule}$.\\
\textbf{Resultado notável:} o rendimento do carregamento de um capacitor por
fonte de tensão através de resistor é \textbf{sempre $50\%$} --- não importa o
valor de $R$, nem de $C$, nem da tensão. Diminuir $R$ apenas acelera o
processo; a energia perdida é a mesma. A demonstração formal por integração
está no bloco de desafio.\\
\textbf{Consequência prática:} carregadores eficientes não usam resistor em série
--- usam conversores chaveados, que transferem energia em pacotes através de um
indutor e escapam desse limite de $50\%$.}
\end{questao}

\begin{questao}[3]{}
Um capacitor plano usa dielétrico com rigidez de
$\SI{20}{\kilo\volt\per\milli\meter}$ e espessura de
$\SI{0.1}{\milli\meter}$.
\begin{itensq}
\item Determine a tensão de ruptura.
\item Especifique a tensão nominal com margem de $50\%$.
\item Explique por que a rigidez limita o produto capacitância-tensão.
\end{itensq}
\resposta{(a) $\SI{2}{\kilo\volt}$; (b) $\SI{1}{\kilo\volt}$}
\resolucao{(a) $U_{rup}=E_{rig}\cdot d=20\cdot0{,}1=\SI{2}{\kilo\volt}$.\\
(b) Com margem de $50\%$, especifica-se $\SI{1}{\kilo\volt}$ --- e ainda assim
é prudente derrateá-la em alta temperatura e sob ondulação.\\
(c) Combinando $C=\dfrac{\varepsilon A}{d}$ com
$U_{\max}=E_{rig}d$:
\[
C\,U_{\max}=\varepsilon A E_{rig}
\quad\Longrightarrow\quad
W_{\max}=\frac12CU_{\max}^{2}=\frac12\varepsilon E_{rig}^{2}\,(A\,d).
\]
O termo $Ad$ é o \textbf{volume} de dielétrico. Ou seja, a energia máxima
armazenável depende apenas do volume e das propriedades do material
($\varepsilon$ e $E_{rig}$) --- \emph{não} da geometria escolhida.\\
\textbf{Consequência:} não existe truque geométrico para ganhar energia. Só
avança quem melhora o material. É por isso que o desenvolvimento de capacitores
é, essencialmente, ciência de materiais.}
\end{questao}

\begin{questao}[3]{}
Em um circuito RC de carga, mediu-se tensão final de $\SI{10}{\volt}$,
corrente inicial de $\SI{2}{\milli\ampere}$ e tensão de $\SI{6.32}{\volt}$ em
$\SI{3}{\milli\second}$. Determine $R$ e $C$.
\resposta{$R=\SI{5}{\kilo\ohm}$; $C=\SI{0.6}{\micro\farad}$}
\resolucao{\textbf{Resistência:} em $t=0^{+}$ o capacitor está descarregado e
toda a tensão cai em $R$:
\[
R=\frac{U_f}{i_0}=\frac{10}{\pot{2}{-3}}=\SI{5}{\kilo\ohm}.
\]
\textbf{Constante de tempo:} $6{,}32/10=0{,}632$ --- exatamente a fração de
$1\tau$. Logo $\tau=\SI{3}{\milli\second}$ e
\[
C=\frac{\tau}{R}=\frac{\pot{3}{-3}}{5000}=\pot{6}{-7}\,\si{\farad}
=\SI{0.6}{\micro\farad}.
\]
\textbf{Estrutura da medição:} a corrente inicial isola $R$; o transitório
isola $\tau$; e só a combinação entrega $C$. Nenhuma medida sozinha determina
o capacitor --- é o mesmo padrão do circuito RL, em que o regime permanente dá
$R$ e o transitório dá $L$.}
\end{questao}

\begin{questao}[3]{}
Compare o comportamento de capacitor e indutor preenchendo o raciocínio para os
instantes $t=0^{+}$ e $t\to\infty$, e justifique cada equivalência.
\resposta{Capacitor: curto em $t=0^{+}$ (se descarregado), aberto em regime;
indutor: aberto em $t=0^{+}$ (se sem corrente), curto em regime}
\resolucao{\textbf{Capacitor descarregado em $t=0^{+}$:} $u_C=0$, e um elemento
com tensão nula equivale a um \textbf{curto-circuito}. Em regime CC,
$\mathrm{d}u/\mathrm{d}t=0$ implica $i=0$: equivale a \textbf{circuito
aberto}.\\
\textbf{Indutor sem corrente em $t=0^{+}$:} $i_L=0$, o que equivale a
\textbf{circuito aberto}. Em regime, $\mathrm{d}i/\mathrm{d}t=0$ implica
$v=0$: equivale a \textbf{curto}.\\
\textbf{Atenção à condição inicial:} as equivalências em $t=0^{+}$ pressupõem
estado inicial nulo. Se o capacitor já estiver a $U_0$, ele equivale a uma
\emph{fonte de tensão} $U_0$; se o indutor já conduzir $I_0$, equivale a uma
\emph{fonte de corrente} $I_0$.\\
\textbf{Por que isso é útil:} esses quatro substitutos permitem obter os
valores inicial e final de qualquer transitório de primeira ordem resolvendo
apenas circuitos \emph{resistivos}. Junto com $\tau$, eles determinam a solução
completa sem integrar nada.}
\end{questao}

\begin{questao}[3]{}
Um filtro RC passa-baixas tem $R=\SI{10}{\kilo\ohm}$ e
$C=\SI{47}{\nano\farad}$.
\begin{itensq}
\item Determine a frequência de corte.
\item Relacione-a com a constante de tempo.
\item Descreva o comportamento muito abaixo e muito acima de $f_c$.
\end{itensq}
\resposta{(a) $f_c\approx\SI{339}{\hertz}$; (b) $f_c=1/(2\pi\tau)$}
\resolucao{(a)
\[
f_c=\frac{1}{2\pi RC}=\frac{1}{2\pi(10^{4})(\pot{47}{-9})}
=\SI{339}{\hertz}.
\]
(b) Como $\tau=RC=\SI{0.47}{\milli\second}$, vale
$f_c=\dfrac{1}{2\pi\tau}$: um circuito \emph{lento} (τ grande) filtra a partir
de frequências \emph{baixas}. As duas descrições --- temporal e frequencial ---
são a mesma informação.\\
(c) \textbf{Muito abaixo de $f_c$:} o capacitor tem reatância alta, quase não
carrega, e o sinal passa praticamente intacto (ganho $\approx1$).\\
\textbf{Muito acima:} o capacitor não tem tempo de carregar entre as variações;
a saída é fortemente atenuada, a $-\SI{20}{\decibel}$ por década.\\
\textbf{Exatamente em $f_c$:} o ganho vale $1/\sqrt2\approx0{,}707$
($-\SI{3}{\decibel}$), e a defasagem é de $\SI{45}{\degree}$. É a definição
convencional de frequência de corte.}
\end{questao}

\blocodesafio

\begin{questao}[4]{}
\textbf{(Dedução --- resposta do circuito RC.)} Um RC série com fonte $U$ tem a
chave fechada em $t=0$ e $u_C(0)=0$.
\begin{itensq}
\item Escreva a equação diferencial e resolva por separação de variáveis.
\item Obtenha $i(t)$ e verifique a LKT nos dois extremos.
\item Compare a estrutura da solução com a do circuito RL.
\end{itensq}
\resposta{$u_C=U\left(1-e^{-t/RC}\right)$;
$i=\dfrac{U}{R}e^{-t/RC}$}
\resolucao{(a) Pela LKT, com $i=C\,\mathrm{d}u_C/\mathrm{d}t$:
\[
U=Ri+u_C=RC\frac{\mathrm{d}u_C}{\mathrm{d}t}+u_C
\;\Longrightarrow\;
\frac{\mathrm{d}u_C}{U-u_C}=\frac{\mathrm{d}t}{RC}.
\]
Integrando de $0$ a $t$:
\[
-\ln\!\left(\frac{U-u_C}{U}\right)=\frac{t}{RC}
\;\Longrightarrow\;
u_C(t)=U\left(1-e^{-t/RC}\right).
\]
(b)
\[
i=C\frac{\mathrm{d}u_C}{\mathrm{d}t}
=CU\cdot\frac{1}{RC}e^{-t/RC}=\frac{U}{R}e^{-t/RC}.
\]
\emph{Em $t=0^{+}$:} $u_C=0$ e $Ri=U$; soma $=U$ \checkmark.\\
\emph{Em $t\to\infty$:} $u_C=U$ e $Ri=0$; soma $=U$ \checkmark.\\
(c) \textbf{Dualidade estrutural.} No RL, a \emph{corrente} sobe como
$1-e^{-t/\tau}$ e a tensão do indutor decai como $e^{-t/\tau}$. No RC, é a
\emph{tensão} que sobe e a corrente que decai. Trocando
$L\leftrightarrow C$, $i\leftrightarrow u$ e $R\leftrightarrow1/R$, uma solução
vira a outra --- razão pela qual basta dominar um dos dois casos.}
\end{questao}

\begin{questao}[4]{}
\textbf{(Demonstração --- rendimento do carregamento.)} Um capacitor
descarregado é carregado por uma fonte $U$ através de $R$.
\begin{itensq}
\item Integre $p_R=Ri^{2}$ para obter a energia dissipada.
\item Mostre que ela vale $\frac12CU^{2}$, independentemente de $R$.
\item Explique como conversores chaveados contornam esse limite.
\end{itensq}
\resposta{$W_R=\frac12CU^{2}$; rendimento sempre $50\%$}
\resolucao{(a) Com $i=\dfrac{U}{R}e^{-t/\tau}$ e $\tau=RC$:
\[
W_R=\int_0^{\infty}R\left(\frac{U}{R}\right)^{2}e^{-2t/\tau}\,\mathrm{d}t
=\frac{U^{2}}{R}\left[-\frac{\tau}{2}e^{-2t/\tau}\right]_0^{\infty}
=\frac{U^{2}\tau}{2R}.
\]
(b) Substituindo $\tau=RC$:
\[
W_R=\frac{U^{2}}{2R}\cdot RC=\frac12CU^{2}.
\]
O $R$ \textbf{cancela}: ele aparece na potência ($\propto1/R$) e desaparece na
duração ($\tau\propto R$). Como a fonte forneceu $CU^{2}$ e o capacitor reteve
$\frac12CU^{2}$, o rendimento é exatamente $50\%$ --- sempre.\\
(c) \textbf{O limite decorre de forçar a tensão do capacitor a saltar de $0$
para $U$ contra uma fonte rígida.} Conversores chaveados evitam isso: um
indutor recebe energia da fonte em regime de corrente, e depois a entrega ao
capacitor. Como o indutor é um elemento reativo (não dissipa), a transferência
pode, em princípio, ser feita sem perda --- e na prática rendimentos de $90\%$
a $95\%$ são rotineiros.\\
\textbf{Observação:} o mesmo resultado de $50\%$ aparece no compartilhamento de
carga entre dois capacitores e na descarga de um indutor. Em todos os casos, a
independência em relação a $R$ tem a mesma origem: potência e duração variam
inversamente.}
\end{questao}

\begin{questao}[4]{}
\textbf{(Projeto de banco.)} Deseja-se um banco de $\SI{100}{\micro\farad}$
capaz de operar em $\SI{100}{\volt}$, dispondo apenas de capacitores de
$\SI{100}{\micro\farad}$ e $\SI{50}{\volt}$.
\begin{itensq}
\item Determine a associação e o número de unidades.
\item Explique por que são necessários resistores de equalização.
\item Dimensione-os para uma corrente de fuga de até
      $\SI{10}{\micro\ampere}$.
\end{itensq}
\resposta{(a) duas séries de dois, em paralelo --- 4 unidades;
(c) $R\le\SI{50}{\kilo\ohm}$, dissipando $\SI{50}{\milli\watt}$ cada}
\resolucao{(a) \textbf{Série de dois:} suporta $\SI{100}{\volt}$, mas a
capacitância cai para $\SI{50}{\micro\farad}$.\\
\textbf{Duas dessas séries em paralelo:} $50+50=\SI{100}{\micro\farad}$,
mantendo $\SI{100}{\volt}$. Total: \textbf{4 capacitores}.\\
Note o custo: quadruplicou-se o número de componentes para dobrar a tensão. É o
preço inevitável --- a energia armazenada, essa sim, quadruplicou.\\
(b) \textbf{Sem equalização, a divisão de tensão em série é instável.} A tensão
de repouso é fixada pelas \emph{correntes de fuga}, não pelas capacitâncias: o
capacitor de menor fuga assume mais tensão. Como a fuga cresce com a tensão e
com a temperatura, o desequilíbrio se realimenta e pode levar um dos capacitores
a ultrapassar $\SI{50}{\volt}$ e perfurar --- o que joga toda a tensão no outro
e o destrói em seguida.\\
(c) Ligue um resistor em paralelo com cada capacitor, dimensionado para que sua
corrente domine a fuga. Adotando um fator $100$:
\[
I_R=100\cdot\pot{10}{-6}=\SI{1}{\milli\ampere}
\;\Longrightarrow\;
R\le\frac{50}{\pot{1}{-3}}=\SI{50}{\kilo\ohm}.
\]
Dissipação: $P=\dfrac{50^{2}}{\pot{50}{3}}=\SI{50}{\milli\watt}$ por resistor
--- especifique $\frac14\,\si{\watt}$.\\
\textbf{Bônus:} esses resistores também funcionam como \emph{bleeders},
descarregando o banco quando o equipamento é desligado. Com
$\tau=RC=50\,\mathrm{k}\times100\,\mu=\SI{5}{\second}$, o banco fica seguro em
menos de meio minuto.}
\end{questao}

\begin{questao}[4]{}
\textbf{(Projeto de filtro.)} Projete um passa-baixas RC de primeira ordem com
$f_c=\SI{100}{\hertz}$, usando $C=\SI{100}{\nano\farad}$.
\begin{itensq}
\item Determine $R$ e escolha o valor comercial E24.
\item Recalcule $f_c$ e avalie o erro.
\item Determine a atenuação em $\SI{1}{\kilo\hertz}$ e discuta a limitação da
      primeira ordem.
\end{itensq}
\resposta{(a) $R=\SI{15.9}{\kilo\ohm}\to\SI{16}{\kilo\ohm}$;
(b) $f_c=\SI{99.5}{\hertz}$, erro $0{,}5\%$;
(c) $\approx-\SI{20}{\decibel}$}
\resolucao{(a)
\[
R=\frac{1}{2\pi f_cC}=\frac{1}{2\pi(100)(\pot{100}{-9})}
=\SI{15915}{\ohm}.
\]
O valor E24 mais próximo é $\SI{16}{\kilo\ohm}$.\\
(b) $f_c=\dfrac{1}{2\pi(16000)(\pot{100}{-9})}=\SI{99.5}{\hertz}$ --- erro de
$0{,}5\%$, muito abaixo da tolerância do próprio capacitor (tipicamente
$\pm10\%$ ou pior). \textbf{O capacitor, e não o resistor, domina a incerteza
do projeto.}\\
(c) Uma década acima de $f_c$, o ganho é
\[
|H|=\frac{1}{\sqrt{1+(f/f_c)^{2}}}\approx\frac{1}{10}
\;\Longrightarrow\;
-\SI{20}{\decibel}.
\]
\textbf{Limitação:} $\SI{20}{\decibel}$ por década é pouco. Para atenuar
$\SI{60}{\decibel}$ seria preciso $f=\SI{100}{\kilo\hertz}$ --- três décadas
acima. Se a especificação exigir corte mais abrupto, é necessário aumentar a
\emph{ordem} do filtro (cascatear seções ou usar topologia ativa), não
apertar $R$ e $C$.\\
\textbf{Cuidado ao cascatear:} duas seções RC ligadas diretamente
\emph{carregam} uma à outra e o $f_c$ resultante não é o previsto. É preciso
desacoplá-las com um buffer, ou projetar a rede como um todo.}
\end{questao}

\begin{questao}[4]{}
\textbf{(Propagação de incerteza.)} Analise o efeito das tolerâncias de $R$ e
$C$ sobre a constante de tempo.
\begin{itensq}
\item Obtenha a expressão da incerteza relativa de $\tau$.
\item Avalie para $R$ com $\pm1\%$ e $C$ eletrolítico com $\pm20\%$, nos
      critérios de pior caso e estatístico.
\item Conclua sobre onde investir em precisão.
\end{itensq}
\resposta{(a) $\dfrac{\Delta\tau}{\tau}=\dfrac{\Delta R}{R}+\dfrac{\Delta C}{C}$;
(b) $21\%$ (pior caso) e $20{,}0\%$ (estatístico)}
\resolucao{(a) Como $\tau=RC$ é um \emph{produto}, tome o logaritmo e
diferencie:
\[
\ln\tau=\ln R+\ln C
\;\Longrightarrow\;
\frac{\Delta\tau}{\tau}=\frac{\Delta R}{R}+\frac{\Delta C}{C}.
\]
As incertezas \textbf{relativas} se somam --- diferentemente da soma, em que se
somam as absolutas.\\
(b) \emph{Pior caso:} $1\%+20\%=21\%$.\\
\emph{Estatístico:} $\sqrt{1^{2}+20^{2}}=\sqrt{401}=20{,}02\%$.\\
(c) \textbf{Os dois critérios praticamente coincidem} --- e isso é a
informação relevante. Quando uma parcela domina de forma tão desproporcional,
combinar em quadratura não ajuda: o resultado é essencialmente a maior
tolerância isolada.\\
\textbf{Conclusão de projeto:} trocar o resistor de $1\%$ por um de $0{,}1\%$
não muda nada ($20{,}00\%$ contra $20{,}02\%$). Todo o ganho está no capacitor.
Se a aplicação exigir temporização precisa, é preciso abandonar o eletrolítico
--- filme de poliéster oferece $\pm5\%$, e polipropileno chega a $\pm1\%$ com
excelente estabilidade térmica.\\
\textbf{Regra geral:} identifique a maior contribuição \emph{antes} de
especificar componentes. Precisão gasta em parcelas menores é dinheiro
jogado fora.}
\end{questao}

\begin{questao}[4]{}
\textbf{(Dedução --- força entre as placas.)} Considere um capacitor plano de
área $A$ e separação $d$.
\begin{itensq}
\item Obtenha a força entre as placas pelo método da energia, com carga
      constante.
\item Expresse o resultado em função de $U$ e discuta o caso a tensão
      constante.
\item Avalie para $A=\SI{100}{\centi\meter\squared}$,
      $d=\SI{1}{\milli\meter}$ e $U=\SI{1}{\kilo\volt}$.
\end{itensq}
\dados{$\varepsilon_0=\pot{8.85}{-12}\,\si{\farad\per\meter}$}
\resposta{$F=\dfrac{\varepsilon_0AU^{2}}{2d^{2}}\approx\SI{44}{\milli\newton}$,
atrativa}
\resolucao{(a) Com o capacitor \textbf{isolado}, $Q$ é constante e convém usar
\[
W=\frac{Q^{2}}{2C}=\frac{Q^{2}d}{2\varepsilon_0A}.
\]
A força que o sistema exerce é
\[
F=-\frac{\mathrm{d}W}{\mathrm{d}d}=-\frac{Q^{2}}{2\varepsilon_0A},
\]
negativa, isto é, \textbf{atrativa} --- as placas tendem a se aproximar,
reduzindo a energia.\\
(b) Substituindo $Q=CU=\dfrac{\varepsilon_0AU}{d}$:
\[
F=\frac{\varepsilon_0AU^{2}}{2d^{2}}\quad\text{(em módulo)}.
\]
\textbf{Sutileza importante:} a tensão constante, a energia \emph{aumenta} ao
aproximar as placas ($W=\frac12CU^{2}$ com $C$ crescendo), o que sugeriria
força repulsiva. O paradoxo se resolve incluindo a bateria: ela fornece
$U\,\mathrm{d}Q$, o dobro da variação de energia do campo. Feito o balanço
completo, a força é atrativa e de \emph{mesmo módulo} --- como tem de ser, pois
a força é uma grandeza física e não depende do que se escolhe manter
constante.\\
(c)
\[
F=\frac{(\pot{8.85}{-12})(\pot{1}{-2})(10^{6})}{2(10^{-6})}
=\SI{44.3}{\milli\newton}.
\]
\textbf{Ordem de grandeza:} cerca de $\SI{4.5}{\gram}$-força para
$\SI{1}{\kilo\volt}$ --- pequena, mas suficiente para fazer as placas vibrarem
audivelmente sob tensão alternada. É o "zumbido" de capacitores cerâmicos
classe 2, e o princípio de funcionamento de microfones e atuadores
eletrostáticos de MEMS.}
\end{questao}

\begin{questao}[4]{}
\textbf{(Projeto de temporizador.)} Um circuito deve atingir $70\%$ da tensão
final em exatamente $\SI{1}{\second}$.
\begin{itensq}
\item Determine $\tau$.
\item Escolha $C$ e determine $R$, com valores comerciais.
\item Avalie a repetibilidade real do temporizador.
\end{itensq}
\resposta{(a) $\tau=\SI{0.83}{\second}$; (b) $C=\SI{100}{\micro\farad}$ e
$R=\SI{8.2}{\kilo\ohm}$; (c) erro dominado pelo capacitor}
\resolucao{(a) De $0{,}70=1-e^{-t/\tau}$:
\[
e^{-t/\tau}=0{,}30
\;\Longrightarrow\;
\tau=\frac{t}{\ln(1/0{,}30)}=\frac{1}{1{,}204}=\SI{0.83}{\second}.
\]
(b) Com $C=\SI{100}{\micro\farad}$:
$R=\tau/C=0{,}83/\pot{100}{-6}=\SI{8306}{\ohm}$. O valor E24 mais próximo é
$\SI{8.2}{\kilo\ohm}$, que dá $\tau=\SI{0.82}{\second}$ e
\[
t_{70\%}=0{,}82(1{,}204)=\SI{0.987}{\second}\quad(\text{erro de }1{,}3\%).
\]
(c) \textbf{O erro de $1{,}3\%$ do arredondamento é irrelevante} diante da
tolerância do eletrolítico ($\pm20\%$), que se transfere integralmente para
$\tau$: o tempo real ficará entre $0{,}79$ e $\SI{1.19}{\second}$. Some-se a
deriva térmica (vários por cento entre $0$ e $\SI{60}{\celsius}$) e o
envelhecimento.\\
\textbf{Conclusão:} um temporizador RC com eletrolítico serve para "cerca de um
segundo", não para "um segundo". Para precisão, as alternativas são capacitor
de filme (levando a $R$ muito maior e $C$ menor) ou --- o caminho moderno ---
um contador digital com oscilador a cristal, cuja estabilidade é de partes por
milhão.\\
\textbf{Observação de método:} note que o requisito foi dado em $70\%$, e não
em $\tau$. Confundir os dois levaria a um erro de $20\%$ no projeto, pois
$t_{70\%}=1{,}204\,\tau$.}
\end{questao}

\begin{questao}[4]{}
\textbf{(Modelagem --- autodescarga.)} Um capacitor real é modelado por $C$
ideal em paralelo com uma resistência de fuga $R_f$.
\begin{itensq}
\item Deduza a lei de decaimento da tensão em circuito aberto.
\item Defina a constante de tempo de autodescarga e relacione-a com a
      especificação de "corrente de fuga" de catálogo.
\item Avalie para $C=\SI{1000}{\micro\farad}$ e
      $R_f=\SI{100}{\kilo\ohm}$, com tensão inicial de $\SI{300}{\volt}$.
\end{itensq}
\resposta{$u=U_0e^{-t/R_fC}$; $\tau=\SI{100}{\second}$; após
$\SI{5}{\minute}$ restam $\approx\SI{15}{\volt}$}
\resolucao{(a) Em circuito aberto, a única corrente é a de fuga:
\[
C\frac{\mathrm{d}u}{\mathrm{d}t}=-\frac{u}{R_f}
\;\Longrightarrow\;
u(t)=U_0e^{-t/R_fC}.
\]
(b) $\tau_f=R_fC$. Catálogos raramente informam $R_f$; informam a
\emph{corrente de fuga} $I_f$ a uma dada tensão. A conversão é imediata:
$R_f=U/I_f$, e portanto
\[
\tau_f=\frac{U\,C}{I_f}.
\]
Note que $\tau_f$ depende do \emph{produto} $R_fC$, e em eletrolíticos $R_f$
cai quando $C$ cresce --- razão pela qual é a autodescarga de capacitores grandes
não é proporcionalmente mais lenta.\\
(c) $\tau_f=\pot{100}{3}\cdot\pot{1000}{-6}=\SI{100}{\second}$. Após
$\SI{300}{\second}$:
\[
u=300\,e^{-3}=300(0{,}0498)=\SI{14.9}{\volt}.
\]
\textbf{Leitura de segurança:} cinco minutos depois de desligado, o barramento
ainda guarda $\SI{15}{\volt}$ --- e, o que é pior, $\SI{0.11}{\joule}$
concentrados. Aos $\SI{60}{\second}$ ainda havia $\SI{110}{\volt}$, plenamente
letais.\\
\textbf{Por isso o resistor de descarga é obrigatório}, e não opcional: com um
bleeder de $\SI{10}{\kilo\ohm}$, $\tau$ cai para $\SI{10}{\second}$ e o
equipamento fica seguro em menos de um minuto, ao custo de
$\SI{9}{\watt}$ apenas enquanto energizado.}
\end{questao}

\begin{questao}[4]{}
\textbf{(ESR e ondulação.)} Um capacitor de filtro de fonte chaveada tem
$\mathrm{ESR}=\SI{50}{\milli\ohm}$ e deve conduzir ondulação de
$\SI{2}{\ampere}$ eficazes.
\begin{itensq}
\item Determine a dissipação e a ondulação de tensão devida \emph{apenas} à
      ESR.
\item Compare com a ondulação devida à capacitância, para
      $C=\SI{470}{\micro\farad}$ a $\SI{100}{\kilo\hertz}$.
\item Conclua sobre o critério de seleção.
\end{itensq}
\resposta{(a) $\SI{0.2}{\watt}$ e $\approx\SI{100}{\milli\volt}$ de pico;
(b) $\approx\SI{5}{\milli\volt}$ --- a ESR domina}
\resolucao{(a) $P=\mathrm{ESR}\cdot I_{rms}^{2}=0{,}05(4)=\SI{0.2}{\watt}$.\\
A ondulação de tensão que a ESR produz acompanha a corrente instantaneamente:
\[
\Delta u_{ESR}\approx\mathrm{ESR}\cdot I_{pico}\approx0{,}05(2)=\SI{100}{\milli\volt}.
\]
(b) A parcela capacitiva depende da carga transferida por semiciclo. Tomando
$\Delta t\approx1/(2f)=\SI{5}{\micro\second}$:
\[
\Delta u_C\approx\frac{I\,\Delta t}{C}
=\frac{2(\pot{5}{-6})}{\pot{470}{-6}}\approx\SI{21}{\milli\volt}.
\]
(c) \textbf{A ESR domina} --- ela responde por cerca de $80\%$ da ondulação
total. Aumentar $C$ para $\SI{1000}{\micro\farad}$ reduziria $\Delta u_C$ pela
metade e deixaria a ondulação praticamente inalterada.\\
\textbf{Critério correto de seleção:} em fontes chaveadas de alta frequência,
escolhe-se o capacitor pela \textbf{ESR e pela corrente de ondulação
admissível}, não pela capacitância. É comum usar vários capacitores menores em
paralelo --- não para somar $C$, mas para dividir a ESR por $n$ e a dissipação
individual por $n^{2}$.\\
\textbf{Nota:} em frequências baixas ($\SI{120}{\hertz}$ de retificação de
rede), a relação se inverte e a capacitância volta a ser o critério dominante.
O parâmetro decisivo depende da faixa de operação.}
\end{questao}

\begin{questao}[4]{}
\textbf{(Comparação tecnológica.)} Um supercapacitor de $\SI{3000}{\farad}$ e
$\SI{2.7}{\volt}$ tem massa de $\SI{0.50}{\kilogram}$.
\begin{itensq}
\item Determine a energia armazenada e a densidade de energia.
\item Compare com uma bateria de íon-lítio ($\approx\SI{250}{\watt\hour\per\kilogram}$).
\item Aponte em que aspecto o supercapacitor é superior e por quê.
\end{itensq}
\resposta{(a) $\SI{10.9}{\kilo\joule}=\SI{3.04}{\watt\hour}$, ou
$\SI{6.1}{\watt\hour\per\kilogram}$; (b) cerca de $40$ vezes menos}
\resolucao{(a)
\[
W=\frac12CU^{2}=\frac12(3000)(2{,}7)^{2}=\SI{10935}{\joule}
=\frac{10935}{3600}=\SI{3.04}{\watt\hour}.
\]
Densidade: $3{,}04/0{,}50=\SI{6.1}{\watt\hour\per\kilogram}$.\\
(b) A bateria armazena cerca de \textbf{40 vezes mais} energia por quilograma.
Para igualar uma célula de íon-lítio de $\SI{20}{\gram}$, seriam necessários
$\SI{0.8}{\kilogram}$ de supercapacitores.\\
(c) \textbf{Potência específica e vida útil.} O supercapacitor pode ser
carregado e descarregado em \emph{segundos}, com correntes de centenas de
ampères, porque armazena energia \textbf{eletrostaticamente} --- sem reações
químicas, sem difusão iônica lenta, sem degradação de eletrodos. Suporta
centenas de milhares de ciclos, contra alguns milhares da bateria, e funciona
bem a baixas temperaturas.\\
\textbf{Consequência de aplicação:} baterias são \emph{reservatórios} de
energia; supercapacitores são \emph{amortecedores} de potência. Por isso
convivem em veículos elétricos --- o supercapacitor absorve o pico da frenagem
regenerativa e entrega o pico da aceleração, poupando a bateria justamente dos
regimes que mais a degradam.\\
\textbf{Ressalva:} a tensão do supercapacitor cai linearmente com a carga
(pois $Q=CU$), ao contrário do platô da bateria. Aproveitar mais de $75\%$ da
energia exige um conversor CC-CC de entrada ampla.}
\end{questao}

% END BANCO COMPLEMENTAR

\gabarito[10.1 Capacitores]

% =====================================================================
\subsection{Indutores}
% =====================================================================

\begin{conceito}[Antes de começar]
\textbf{Indutância de um solenoide:}
$L = \dfrac{\mu_0\,N^2\,A}{\ell}$, com
$\mu_0 = 4\pi\times10^{-7}\,\si{\henry\per\meter}$.\par
\textbf{Tensão induzida:} $v_L = L\,\dfrac{\mathrm{d}i}{\mathrm{d}t}$ --- o
indutor se opõe à \emph{variação} de corrente.\par
\textbf{Energia armazenada:} $E = \tfrac12 L\,I^2$.\par
\textbf{Associação (igual à dos resistores):} série soma, paralelo pelo inverso.\par
\textbf{Circuito RL:} $\tau = L/R$; na energização,
$i(t) = \dfrac{U_0}{R}\big(1 - e^{-t/\tau}\big)$.
\end{conceito}

\legendaniveis

\blocofixacao

\begin{questao}[1]{}
Um indutor ideal, em regime permanente CC (corrente constante), comporta-se como:
\begin{alternativas}
\item um circuito aberto.
\item um curto-circuito.
\item um resistor de valor $L$.
\item uma fonte de tensão.
\item um capacitor.
\end{alternativas}
\resposta{(b)}
\resolucao{Como $v_L = L\dfrac{\mathrm{d}i}{\mathrm{d}t}$ e, em regime permanente,
a corrente é constante ($\mathrm{d}i/\mathrm{d}t = 0$), a tensão no indutor é
\emph{nula}: ele se comporta como um \textbf{curto-circuito}. (É o oposto do
capacitor, que vira circuito aberto.)}
\end{questao}

\begin{questao}[1]{}
Calcule a indutância de um solenoide com $N = 500$ espiras, comprimento
$\ell = \SI{0.3}{\meter}$ e área $A = \SI{5}{\centi\meter\squared}$.
\dados{$\mu_0 = 4\pi\times10^{-7}\,\si{\henry\per\meter}$.}
\resposta{$L \approx \SI{0.52}{\milli\henry}$}
\resolucao{Com $A = \SI{5}{\centi\meter\squared} = \pot{5}{-4}\,\si{\meter\squared}$:
\[
L=\frac{\mu_0 N^2 A}{\ell}
=\frac{(\pot{4\pi}{-7})(500)^2(\pot{5}{-4})}{0{,}3}
\approx\pot{5.24}{-4}\,\si{\henry}=\SI{0.52}{\milli\henry}.
\]}
\end{questao}

\begin{questao}[1]{}
Um indutor de $\SI{0.1}{\henry}$ tem a corrente variando de $0$ a
$\SI{2}{\ampere}$ em $\SI{0.2}{\second}$. Calcule a tensão média induzida.
\resposta{$v_L = \SI{1}{\volt}$}
\resolucao{$v_L = L\dfrac{\Delta i}{\Delta t}
= 0{,}1\cdot\dfrac{2}{0{,}2} = \SI{1}{\volt}$.}
\end{questao}

\begin{questao}[1]{}
Dois indutores de $\SI{3}{\henry}$ e $\SI{9}{\henry}$ são ligados em paralelo. A
indutância equivalente é:
\begin{alternativas}[3]
\item \SI{12}{\henry}
\item \SI{6}{\henry}
\item \SI{2.25}{\henry}
\item \SI{3}{\henry}
\item \SI{27}{\henry}
\end{alternativas}
\resposta{(c)}
\resolucao{Indutores seguem a mesma regra dos resistores. Em paralelo:
\[
\Leq=\frac{L_1 L_2}{L_1+L_2}=\frac{3\cdot9}{12}=\frac{27}{12}=\SI{2.25}{\henry}.
\]}
\end{questao}

\blocoaplicacao

\begin{questao}[2]{}
Um indutor de $L = \SI{1}{\henry}$ é ligado em série com $R = \SI{100}{\ohm}$ a
uma fonte de $\SI{12}{\volt}$.
\begin{itensq}
\item Qual a constante de tempo?
\item Qual a corrente final (regime permanente)?
\item Qual a corrente após um intervalo igual a $\tau$?
\end{itensq}
\resposta{(a) $\tau = \SI{10}{\milli\second}$; (b) \SI{0.12}{\ampere};
(c) $\approx\SI{75.8}{\milli\ampere}$}
\resolucao{(a) $\tau = \dfrac{L}{R} = \dfrac{1}{100} = \SI{10}{\milli\second}$.\\
(b) Em regime permanente, o indutor vira curto e a corrente é limitada só por $R$:
$I_\infty = \dfrac{12}{100} = \SI{0.12}{\ampere}$.\\
(c) $i(\tau) = I_\infty(1 - e^{-1}) = 0{,}12\cdot0{,}632
\approx\SI{75.8}{\milli\ampere}$ --- cerca de $\SI{63}{\percent}$ do valor final,
o mesmo comportamento do capacitor na carga.}
\end{questao}

\begin{questao}[2]{}
Calcule a indutância equivalente para:
\begin{itensq}
\item $L_1 = \SI{2}{\milli\henry}$ e $L_2 = \SI{3}{\milli\henry}$ em série;
\item os mesmos dois indutores em paralelo.
\end{itensq}
\resposta{(a) \SI{5}{\milli\henry}; (b) \SI{1.2}{\milli\henry}}
\resolucao{(a) Série: $\Leq = L_1 + L_2 = 2 + 3 = \SI{5}{\milli\henry}$.\\
(b) Paralelo: $\Leq = \dfrac{2\cdot3}{5} = \SI{1.2}{\milli\henry}$.}
\end{questao}

\begin{questao}[2]{Apostila original revisada}
Dois indutores ideais de $\SI{5}{\henry}$ são ligados em paralelo. Essa associação
é conectada em série com um indutor de $\SI{2}{\henry}$. Determine a indutância
equivalente.
\resposta{$L_{\mathrm{eq}}=\SI{4.5}{\henry}$}
\resolucao{Os dois indutores iguais em paralelo resultam em
$L_p=5/2=\SI{2.5}{\henry}$. Em série com $\SI{2}{\henry}$:
$L_{\mathrm{eq}}=2{,}5+2=\SI{4.5}{\henry}$.}
\end{questao}

\blocoaprofundamento

\begin{questao}[3]{}
Compare, lado a lado, o comportamento de um capacitor e de um indutor em um
circuito de corrente contínua, preenchendo mentalmente a tabela abaixo e
justificando cada linha.

\begin{table}[H]
\centering
\small\sffamily
\setlength{\arraystretch}{1.3}
\begin{tabular}{@{}l c c@{}}
\toprule
 & \textbf{Capacitor} & \textbf{Indutor}\\
\midrule
Grandeza que não muda bruscamente & tensão $v_C$ & corrente $i_L$\\
Em regime permanente CC & circuito aberto & curto-circuito\\
Energia armazenada & $\tfrac12 C v^2$ (campo elétrico) & $\tfrac12 L i^2$ (campo magnético)\\
Constante de tempo & $\tau = RC$ & $\tau = L/R$\\
\bottomrule
\end{tabular}
\caption{Dualidade entre capacitor e indutor.}
\end{table}

Explique por que a tensão no capacitor e a corrente no indutor não podem variar
instantaneamente.
\resposta{Ver comentário}
\resolucao{A tensão do capacitor está ligada à sua carga por $v = Q/C$; mudar $v$
instantaneamente exigiria corrente infinita ($i = C\,\mathrm{d}v/\mathrm{d}t$),
o que é fisicamente impossível. De modo \emph{dual}, a corrente do indutor não
pode saltar, pois isso exigiria tensão infinita
($v = L\,\mathrm{d}i/\mathrm{d}t$). Por isso o capacitor "segura" a tensão e o
indutor "segura" a corrente --- é essa continuidade que dá origem às respostas
exponenciais com constante de tempo $\tau$. A tabela evidencia a bela
\textbf{dualidade} entre os dois componentes: trocando tensão$\leftrightarrow$corrente
e $C\leftrightarrow L$, uma descrição vira a outra.}
\end{questao}

\begin{questao}[3]{}
Um indutor de $L = \SI{0.5}{\henry}$ conduz corrente estabilizada de
$\SI{4}{\ampere}$, alimentado por uma fonte através de um resistor. Em
$t=0$, a chave que liga a fonte é \textbf{aberta bruscamente}.
\begin{itensq}
\item Qual a energia armazenada no indutor imediatamente antes da abertura?
\item Por que surge um arco elétrico (faísca) nos contatos da chave?
\item Estime a tensão induzida se a corrente for interrompida em
      $\SI{1}{\milli\second}$.
\item Qual solução prática se usa para proteger o circuito? Explique seu
      funcionamento.
\end{itensq}
\resposta{(a) $E=\SI{4}{\joule}$; (c) $\approx\SI{2000}{\volt}$; (b) e (d) ver comentário}
\resolucao{\textbf{(a)} $E = \tfrac12 L I^2 = \tfrac12(0{,}5)(4)^2
= \SI{4}{\joule}$ --- armazenados no \emph{campo magnético}.

\textbf{(b) Origem do arco.} A corrente em um indutor \textbf{não pode variar
instantaneamente}, pois isso exigiria tensão infinita:
$v_L = L\,\dfrac{\mathrm{d}i}{\mathrm{d}t}$. Ao abrir a chave, tenta-se impor
$\mathrm{d}i/\mathrm{d}t \to \infty$; o indutor reage gerando uma tensão
altíssima, de polaridade \emph{invertida}, que tenta manter a corrente
circulando. Essa tensão ioniza o ar entre os contatos e a corrente continua
fluindo pelo \textbf{arco} --- que é o caminho que o circuito "encontra" para
dissipar os \SI{4}{\joule} armazenados.

\textbf{(c) Estimativa da sobretensão.} Com $\Delta i = \SI{4}{\ampere}$ em
$\Delta t = \SI{1}{\milli\second}$:
\[
|v_L| = L\,\frac{\Delta i}{\Delta t}
= 0{,}5\cdot\frac{4}{\pot{1}{-3}} = \SI{2000}{\volt}.
\]
Duzentas vezes maior que uma alimentação típica de \SI{12}{\volt}! Se a abertura
fosse dez vezes mais rápida, a tensão seria dez vezes maior.

\textbf{(d) Proteção --- o diodo de roda livre (\emph{freewheeling}).}
Liga-se um diodo em \textbf{antiparalelo} com o indutor (catodo no positivo).
Em operação normal ele fica reversamente polarizado e não conduz. No instante da
abertura, a tensão do indutor inverte, o diodo entra em condução e oferece um
caminho fechado para a corrente, que decai suavemente com constante de tempo
$\tau = L/R$, dissipando a energia no próprio circuito em vez de no arco.\\
Alternativas: \emph{snubber} RC, varistor (MOV) ou diodo Zener --- este último
quando se deseja desmagnetização \emph{rápida}, aceitando uma sobretensão
controlada.\\
\textbf{Onde isso aparece na prática:} todo relé, contator, solenoide ou motor
CC acionado por transistor precisa desse diodo. Sem ele, o transistor de
acionamento é destruído na primeira comutação --- um dos erros mais comuns de
quem começa em eletrônica de potência.}
\end{questao}

\begin{questao}[3]{Apostila original revisada}
Um indutor de $L=\SI{2}{\milli\henry}$ permanece por muito tempo em série com um
resistor de $\SI{6}{\kilo\ohm}$ e uma fonte de $\SI{12}{\volt}$. Em $t=0$, a fonte
e o resistor de $\SI{6}{\kilo\ohm}$ são desconectados, e o indutor passa a descarregar
por um resistor de $\SI{2}{\kilo\ohm}$.
\begin{itensq}
\item Determine $i_L(0^-)$ e $i_L(0^+)$.
\item Calcule a constante de tempo da descarga.
\item Escreva $i_L(t)$ para $t\ge 0$ e determine o módulo da tensão inicial no resistor.
\end{itensq}
\resposta{$i_L(0^-)=i_L(0^+)=\SI{2}{\milli\ampere}$; $\tau=\SI{1}{\micro\second}$;
$i_L(t)=\SI{2}{\milli\ampere}e^{-t/(\SI{1}{\micro\second})}$; $|v_R(0^+)|=\SI{4}{\volt}$}
\resolucao{Em regime permanente, o indutor é curto-circuito: $i=12/6000=\SI{2}{\milli\ampere}$.
A corrente do indutor é contínua na comutação. Na descarga, $\tau=L/R=2\,\mathrm{mH}/2\,\mathrm{k\Omega}=\SI{1}{\micro\second}$.
A tensão inicial no resistor vale $Ri=2000\times0{,}002=\SI{4}{\volt}$; sua polaridade é tal que mantém a corrente.}
\end{questao}

% BEGIN BANCO COMPLEMENTAR Indutores
% =====================================================================
%  Banco complementar --- Indutores  (REVISADO)
%  Substitui a versao gerada por template (12 clones em N2, 12 em N3,
%  10 questoes conceituais indevidamente marcadas como N4).
%  Sem sobreposicao com o cap10, que ja traz: indutor ideal em CC (MC),
%  indutancia de solenoide N=500, rampa de corrente basica, paralelo
%  3H/9H, RL com L=1H e R=100 ohm, Leq serie/paralelo e a tabela
%  comparativa capacitor x indutor.
%  Distribuicao mantida: 6 N1 + 12 N2 + 12 N3 + 10 N4 = 40
% =====================================================================

\subsubsection*{Banco complementar --- Indutores}

\begin{atencao}[Organização do banco]
Três relações sustentam tudo: $v=L\dfrac{\mathrm{d}i}{\mathrm{d}t}$,
$W=\frac12LI^{2}$ e $\lambda=LI$. Delas seguem os dois fatos que governam o
comportamento transitório --- a corrente \textbf{não pode saltar} (isso exigiria
tensão infinita) e o indutor \textbf{se opõe à variação}, não à corrente em si.
O N3 trata de sobretensão, acoplamento, saturação e extração de parâmetros; o
N4 exige dedução da EDO, integração de energia e projeto com restrições
físicas.
\end{atencao}

\blocofixacao

\begin{questao}[1]{}
A tensão sobre um indutor vale $\SI{50}{\volt}$ quando a corrente varia à taxa
de $\SI{250}{\ampere\per\second}$. Determine a indutância.
\resposta{$L=\SI{0.2}{\henry}$}
\resolucao{Isolando $L$ na relação fundamental:
\[
v=L\frac{\mathrm{d}i}{\mathrm{d}t}
\;\Longrightarrow\;
L=\frac{v}{\mathrm{d}i/\mathrm{d}t}=\frac{50}{250}=\SI{0.2}{\henry}.
\]
O henry é, portanto, volt-segundo por ampère: $\SI{1}{\henry}$ é a indutância
que produz $\SI{1}{\volt}$ para uma variação de $\SI{1}{\ampere\per\second}$.}
\end{questao}

\begin{questao}[1]{}
Um indutor de $\SI{0.4}{\henry}$ conduz $\SI{5}{\ampere}$ constantes. Determine
a energia armazenada no campo magnético.
\resposta{$W=\SI{5}{\joule}$}
\resolucao{
\[
W=\frac12LI^{2}=\frac12(0{,}4)(5)^{2}=\SI{5}{\joule}.
\]
A dependência é \emph{quadrática} na corrente: dobrar $I$ quadruplica a
energia. Note que, embora a tensão sobre o indutor ideal seja nula em regime CC
(corrente constante), a energia armazenada não é --- ela ficou guardada no
campo durante a fase de crescimento da corrente.}
\end{questao}

\begin{questao}[1]{}
Sobre a corrente em um indutor ideal, é correto afirmar que ela:
\begin{alternativas}[2]
\item pode variar instantaneamente
\item não pode variar instantaneamente
\item é sempre constante
\item é sempre nula em corrente contínua
\end{alternativas}
\resposta{(b)}
\resolucao{Um salto de corrente significaria
$\mathrm{d}i/\mathrm{d}t\to\infty$ e, por $v=L\,\mathrm{d}i/\mathrm{d}t$, uma
tensão infinita --- fisicamente impossível. Logo a corrente no indutor é uma
função \textbf{contínua} do tempo: $i(0^{+})=i(0^{-})$.\\
Cuidado com (c) e (d): a corrente pode variar (e varia, nos transitórios), e em
CC ela é constante e em geral \emph{não nula}. O que é nula em CC é a
\emph{tensão}.}
\end{questao}

\begin{questao}[1]{}
Um indutor de $\SI{200}{\milli\henry}$ está em série com $\SI{50}{\ohm}$.
Determine a constante de tempo do circuito.
\resposta{$\tau=\SI{4}{\milli\second}$}
\resolucao{
\[
\tau=\frac{L}{R}=\frac{0{,}200}{50}=\SI{4e-3}{\second}=\SI{4}{\milli\second}.
\]
\textbf{Atenção à estrutura:} no circuito RL, $\tau=L/R$ --- resistência
\emph{maior} torna o transitório \emph{mais rápido}. No RC é o oposto
($\tau=RC$). Trocar as duas fórmulas é o erro mais comum entre os dois
elementos armazenadores.}
\end{questao}

\begin{questao}[1]{}
Determine o fluxo concatenado em um indutor de $\SI{0.25}{\henry}$ percorrido
por $\SI{8}{\ampere}$.
\resposta{$\lambda=\SI{2}{\weber}$}
\resolucao{
\[
\lambda=LI=0{,}25\cdot8=\SI{2}{\weber}.
\]
O fluxo concatenado $\lambda=N\Phi$ é o produto do fluxo por espira pelo número
de espiras. A indutância é justamente a constante de proporcionalidade entre
$\lambda$ e $I$ --- é essa a definição primária de $L$, da qual
$v=L\,\mathrm{d}i/\mathrm{d}t$ decorre pela lei de Faraday.}
\end{questao}

\begin{questao}[1]{}
Aplicam-se $\SI{30}{\volt}$ a um indutor de $\SI{0.15}{\henry}$ inicialmente
sem corrente. Determine a taxa inicial de variação da corrente.
\resposta{$\mathrm{d}i/\mathrm{d}t=\SI{200}{\ampere\per\second}$}
\resolucao{
\[
\frac{\mathrm{d}i}{\mathrm{d}t}=\frac{v}{L}=\frac{30}{0{,}15}
=\SI{200}{\ampere\per\second}.
\]
No instante inicial a corrente é nula, mas sua \emph{taxa} é máxima. Se a
tensão se mantivesse constante e não houvesse resistência, a corrente cresceria
indefinidamente em rampa --- é a resistência do circuito que estabelece o
patamar final.}
\end{questao}

\blocoaplicacao

\begin{questao}[2]{}
Um indutor de $\SI{0.5}{\henry}$ conduz $\SI{3}{\ampere}$.
\begin{itensq}
\item Determine a energia armazenada.
\item Determine a tensão necessária para levar a corrente a zero
      \emph{linearmente} em $\SI{50}{\milli\second}$.
\end{itensq}
\resposta{(a) $\SI{2.25}{\joule}$; (b) $\SI{30}{\volt}$}
\resolucao{(a) $W=\frac12(0{,}5)(9)=\SI{2.25}{\joule}$.\\
(b) A taxa é $\mathrm{d}i/\mathrm{d}t=(0-3)/0{,}050
=\SI{-60}{\ampere\per\second}$, e
\[
|v|=L\left|\frac{\mathrm{d}i}{\mathrm{d}t}\right|=0{,}5\cdot60=\SI{30}{\volt}.
\]
O sinal negativo indica \textbf{inversão de polaridade}: o indutor passa a
funcionar como fonte, devolvendo ao circuito os $\SI{2.25}{\joule}$
armazenados. É esse comportamento que produz as sobretensões de abertura.}
\end{questao}

\begin{questao}[2]{}
A corrente em um indutor de $\SI{100}{\milli\henry}$ cai linearmente de
$\SI{4}{\ampere}$ a zero em $\SI{20}{\milli\second}$. Determine a tensão
induzida e sua polaridade.
\resposta{$\SI{20}{\volt}$, com polaridade \emph{invertida} em relação à fase
de crescimento}
\resolucao{
\[
|v|=L\left|\frac{\Delta i}{\Delta t}\right|
=0{,}100\cdot\frac{4}{0{,}020}=\SI{20}{\volt}.
\]
Durante o \emph{crescimento} da corrente, o indutor absorve energia e sua
tensão se opõe ao aumento. Durante o \emph{decrescimento}, ele devolve energia
e a tensão inverte, tentando sustentar a corrente.\\
Em ambos os casos vale a lei de Lenz: o indutor se opõe à \emph{variação}, não
ao sentido da corrente.}
\end{questao}

\begin{questao}[2]{}
Três indutores não acoplados de $\SI{6}{\milli\henry}$,
$\SI{12}{\milli\henry}$ e $\SI{12}{\milli\henry}$ são associados. Determine
$L_{eq}$ em série e em paralelo.
\resposta{Série: $\SI{30}{\milli\henry}$; paralelo: $\SI{3}{\milli\henry}$}
\resolucao{\textbf{Série:} $L_{eq}=6+12+12=\SI{30}{\milli\henry}$.\\
\textbf{Paralelo:}
\[
\frac{1}{L_{eq}}=\frac16+\frac1{12}+\frac1{12}=\frac{2+1+1}{12}=\frac13
\;\Longrightarrow\;
L_{eq}=\SI{3}{\milli\henry}.
\]
As regras são \emph{idênticas} às dos resistores --- e opostas às dos
capacitores. A ressalva é a hipótese de \textbf{não acoplamento}: se houver
fluxo mútuo, as fórmulas mudam (ver bloco de aprofundamento).}
\end{questao}

\begin{questao}[2]{}
Um circuito RL série tem $L=\SI{0.6}{\henry}$, $R=\SI{30}{\ohm}$ e fonte de
$\SI{24}{\volt}$, com a chave fechada em $t=0$.
\begin{itensq}
\item Determine $\tau$ e a corrente final.
\item Determine a tensão sobre o indutor em $t=0^{+}$.
\end{itensq}
\resposta{(a) $\tau=\SI{20}{\milli\second}$, $I_f=\SI{0.8}{\ampere}$;
(b) $\SI{24}{\volt}$}
\resolucao{(a) $\tau=L/R=0{,}6/30=\SI{20}{\milli\second}$;
$I_f=V/R=24/30=\SI{0.8}{\ampere}$.\\
(b) Em $t=0^{+}$ a corrente ainda é nula (continuidade), logo o resistor não
tem queda e \textbf{toda} a tensão da fonte aparece sobre o indutor:
$v_L(0^{+})=\SI{24}{\volt}$.\\
\textbf{Os dois extremos:} em $t=0^{+}$ o indutor se comporta como circuito
\emph{aberto} (corrente nula); em $t\to\infty$, como \emph{curto}
(tensão nula). Reconhecer esses dois limites resolve a maioria dos problemas
sem integrar nada.}
\end{questao}

\begin{questao}[2]{}
No circuito anterior, determine a corrente em $t=\tau$, $t=2\tau$ e
$t=3\tau$.
\resposta{$\SI{0.506}{\ampere}$, $\SI{0.692}{\ampere}$ e
$\SI{0.760}{\ampere}$}
\resolucao{Com $i(t)=I_f\left(1-e^{-t/\tau}\right)$ e
$I_f=\SI{0.8}{\ampere}$:
\[
i(\tau)=0{,}8(1-e^{-1})=0{,}8(0{,}632)=\SI{0.506}{\ampere},
\]
\[
i(2\tau)=0{,}8(0{,}865)=\SI{0.692}{\ampere},
\qquad
i(3\tau)=0{,}8(0{,}950)=\SI{0.760}{\ampere}.
\]
As frações $63{,}2\%$, $86{,}5\%$ e $95{,}0\%$ são universais --- valem para
qualquer RL ou RC. Após $5\tau$ atinge-se $99{,}3\%$, e por convenção o
transitório é dado por encerrado.}
\end{questao}

\begin{questao}[2]{}
Aplica-se a um indutor de $\SI{5}{\milli\henry}$ uma onda quadrada de tensão
que alterna entre $+\SI{10}{\volt}$ e $-\SI{10}{\volt}$, com meio período de
$\SI{1}{\milli\second}$. Determine a variação de corrente pico a pico em
regime.
\resposta{$\Delta i=\SI{2}{\ampere}$}
\resolucao{Durante o semiciclo positivo, a tensão constante produz rampa de
corrente:
\[
\Delta i=\frac{v\,\Delta t}{L}=\frac{10\cdot0{,}001}{0{,}005}=\SI{2}{\ampere}.
\]
No semiciclo negativo a corrente desce pela mesma quantidade, e o regime é
simétrico.\\
\textbf{Resultado geral:} tensão quadrada no indutor produz corrente
\textbf{triangular} --- o indutor \emph{integra} a tensão. O capacitor faz o
oposto: corrente quadrada nele produz tensão triangular.}
\end{questao}

\begin{questao}[2]{}
Um indutor \emph{real} de $\SI{0.5}{\henry}$ possui resistência de enrolamento
$r=\SI{8}{\ohm}$ e é ligado a uma fonte CC de $\SI{24}{\volt}$.
\begin{itensq}
\item Determine a corrente em regime permanente.
\item Determine a energia armazenada e a potência dissipada.
\end{itensq}
\resposta{(a) $\SI{3}{\ampere}$; (b) $W=\SI{2.25}{\joule}$ e
$P=\SI{72}{\watt}$}
\resolucao{(a) Em regime, $\mathrm{d}i/\mathrm{d}t=0$ e o indutor ideal não cai
tensão: só $r$ limita a corrente.
\[
I=\frac{24}{8}=\SI{3}{\ampere}.
\]
(b) $W=\frac12(0{,}5)(9)=\SI{2.25}{\joule}$;
$P=rI^{2}=8(9)=\SI{72}{\watt}$.\\
\textbf{Contraste essencial:} a energia é \emph{armazenada} e recuperável; a
potência é \emph{dissipada} continuamente e perdida. Manter esses
$\SI{2.25}{\joule}$ no campo custa $\SI{72}{\watt}$ permanentes --- razão pela
qual eletroímãs de grande porte usam supercondutores, em que $r=0$ e a
manutenção do campo é gratuita.}
\end{questao}

\begin{questao}[2]{}
Dois indutores armazenam a mesma energia: o primeiro tem
$L_1=\SI{2}{\henry}$ com $I_1=\SI{1}{\ampere}$. Se
$L_2=\SI{0.5}{\henry}$, determine $I_2$.
\resposta{$I_2=\SI{2}{\ampere}$}
\resolucao{
\[
\frac12L_1I_1^{2}=\frac12L_2I_2^{2}
\;\Longrightarrow\;
I_2=I_1\sqrt{\frac{L_1}{L_2}}=1\cdot\sqrt{4}=\SI{2}{\ampere}.
\]
Ambos armazenam $\SI{1}{\joule}$. Reduzir $L$ por um fator $4$ exige
\emph{dobrar} a corrente --- e, como as perdas vão com $I^{2}$, o indutor menor
custa quatro vezes mais em dissipação para armazenar o mesmo tanto. É o
compromisso central no projeto de indutores de potência.}
\end{questao}

\begin{questao}[2]{}
Um solenoide de $N=200$ espiras tem indutância $L$. Determine a nova
indutância se o número de espiras for dobrado, mantidos o comprimento e a
seção.
\resposta{$4L$}
\resolucao{Para um solenoide longo,
\[
L=\frac{\mu N^{2}A}{\ell},
\]
de modo que $L\propto N^{2}$. Dobrar $N$ quadruplica $L$.\\
\textbf{Por que o quadrado:} dobrar as espiras dobra o campo produzido por
uma dada corrente \emph{e} dobra o número de espiras que concatenam esse campo.
Os dois efeitos se multiplicam.\\
\textbf{Ressalva prática:} com o mesmo espaço de bobina, dobrar $N$ exige fio
de seção menor, o que aumenta $r$ e as perdas. Na prática, o ganho de $L$ vem
acompanhado de penalidade em corrente admissível.}
\end{questao}

\begin{questao}[2]{}
Um indutor com núcleo de ar tem $L_0=\SI{2}{\milli\henry}$. Introduz-se um
núcleo ferromagnético de permeabilidade relativa $\mu_r=500$. Determine a nova
indutância e comente as limitações.
\resposta{$L=\SI{1}{\henry}$}
\resolucao{Como $L\propto\mu=\mu_r\mu_0$:
\[
L=\mu_rL_0=500\cdot2=\SI{1000}{\milli\henry}=\SI{1}{\henry}.
\]
\textbf{Três limitações do núcleo.} \emph{(i)} \textbf{Saturação:} acima de
certa corrente, $\mu_r$ despenca e $L$ colapsa. \emph{(ii)} \textbf{Perdas:}
histerese e correntes de Foucault dissipam energia proporcionalmente à
frequência. \emph{(iii)} \textbf{Não linearidade:} $L$ deixa de ser constante,
e as fórmulas lineares passam a ser aproximações válidas apenas em torno de um
ponto de operação.\\
Núcleos de ar têm $L$ minúsculo, mas são perfeitamente lineares e sem perdas
magnéticas --- por isso dominam aplicações de radiofrequência.}
\end{questao}

\begin{questao}[2]{}
No circuito RL com $\tau=\SI{20}{\milli\second}$ e $I_f=\SI{0.8}{\ampere}$,
determine o instante em que a corrente atinge $\SI{0.5}{\ampere}$.
\resposta{$t\approx\SI{19.6}{\milli\second}$}
\resolucao{De $i=I_f(1-e^{-t/\tau})$, isole a exponencial:
\[
e^{-t/\tau}=1-\frac{0{,}5}{0{,}8}=0{,}375
\;\Longrightarrow\;
t=\tau\ln\!\left(\frac{1}{0{,}375}\right)=0{,}020\cdot0{,}981
=\SI{19.6}{\milli\second}.
\]
Ou seja, praticamente $1\tau$ --- coerente com o fato de que
$\SI{0.5}{\ampere}$ corresponde a $62{,}5\%$ de $I_f$, muito próximo dos
$63{,}2\%$ atingidos em exatamente $\tau$.}
\end{questao}

\begin{questao}[2]{}
Um indutor de $\SI{0.4}{\henry}$ conduz $\SI{2}{\ampere}$ quando a fonte é
substituída por um curto através de $R=\SI{20}{\ohm}$. Determine $\tau$ e a
corrente após $\SI{40}{\milli\second}$.
\resposta{$\tau=\SI{20}{\milli\second}$; $i=\SI{0.27}{\ampere}$}
\resolucao{$\tau=L/R=0{,}4/20=\SI{20}{\milli\second}$.\\
Na descarga, a corrente decai exponencialmente a partir do valor inicial:
\[
i(t)=I_0e^{-t/\tau}=2\,e^{-40/20}=2e^{-2}=2(0{,}135)=\SI{0.27}{\ampere}.
\]
\textbf{Continuidade:} no instante da comutação a corrente permanece
$\SI{2}{\ampere}$ --- ela não pode saltar. O que muda instantaneamente é a
\emph{tensão}, que assume $-RI_0=\SI{-40}{\volt}$ para manter essa corrente
circulando pelo resistor.}
\end{questao}

\blocoaprofundamento

\begin{questao}[3]{}
Um indutor de $\SI{0.2}{\henry}$ conduz $\SI{5}{\ampere}$ e é descarregado
através de um resistor.
\begin{itensq}
\item Determine a energia total dissipada no resistor.
\item Explique por que ela não depende do valor de $R$.
\item Determine como $R$ afeta o processo.
\end{itensq}
\resposta{(a) $\SI{2.5}{\joule}$; (b) conservação de energia; (c) $R$ altera a
\emph{duração} e a \emph{potência}, não a energia total}
\resolucao{(a) Toda a energia armazenada acaba no resistor:
\[
W=\frac12LI_0^{2}=\frac12(0{,}2)(25)=\SI{2.5}{\joule}.
\]
(b) O indutor ideal não dissipa e o circuito é fechado: a energia magnética não
tem outro destino. A independência em relação a $R$ é consequência direta da
conservação --- e é confirmada por integração no bloco de desafio.\\
(c) $R$ governa $\tau=L/R$ e, portanto, o \emph{ritmo}:
\begin{itemize}
\item $R$ pequeno $\Rightarrow$ $\tau$ grande: descarga lenta, potência baixa e
prolongada.
\item $R$ grande $\Rightarrow$ $\tau$ pequeno: descarga rápida, com pico de
potência $P_0=RI_0^{2}$ muito elevado.
\end{itemize}
\textbf{Exemplo:} com $R=\SI{10}{\ohm}$, $P_0=\SI{250}{\watt}$ por
$\SI{20}{\milli\second}$; com $R=\SI{1000}{\ohm}$, $P_0=\SI{25}{\kilo\watt}$
por $\SI{0.2}{\milli\second}$. Mesma energia, potências que diferem em cem
vezes.}
\end{questao}

\begin{questao}[3]{}
Um indutor de $\SI{0.1}{\henry}$ conduz $\SI{2}{\ampere}$ quando uma chave
mecânica é aberta, interrompendo a corrente em cerca de
$\SI{1}{\micro\second}$.
\begin{itensq}
\item Estime a tensão induzida.
\item Explique por que esse valor não se observa na prática.
\item Indique a consequência real.
\end{itensq}
\resposta{(a) $\SI{200}{\kilo\volt}$ (estimativa idealizada);
(b) o ar rompe muito antes; (c) formação de arco elétrico nos contatos}
\resolucao{(a)
\[
|v|=L\frac{\Delta i}{\Delta t}=0{,}1\cdot\frac{2}{10^{-6}}
=\pot{2}{5}\,\si{\volt}=\SI{200}{\kilo\volt}.
\]
(b) A rigidez dielétrica do ar é da ordem de
$\SI{3}{\kilo\volt\per\milli\meter}$. Com contatos separados por décimos de
milímetro, o ar ioniza a alguns \emph{quilovolts} --- muito antes dos
$\SI{200}{\kilo\volt}$. O cálculo é uma \textbf{estimativa por absurdo}: ele
mostra que a hipótese "a corrente cessa em $\SI{1}{\micro\second}$" é
insustentável.\\
(c) Forma-se um \textbf{arco elétrico}, que mantém a corrente circulando
enquanto o campo se esvazia. O arco limita a tensão a algumas centenas de
volts, mas corrói os contatos, gera interferência eletromagnética e reduz
drasticamente a vida da chave.\\
\textbf{Leitura correta:} não é o indutor que "gera" $\SI{200}{\kilo\volt}$ ---
é a impossibilidade de interromper a corrente que \emph{força} o circuito a
encontrar um caminho, e ele o encontra no ar.}
\end{questao}

\begin{questao}[3]{}
Um diodo de roda livre é colocado em antiparalelo com uma bobina de
$\SI{0.1}{\henry}$ e $r=\SI{5}{\ohm}$, que conduz $\SI{2}{\ampere}$.
\begin{itensq}
\item Explique o funcionamento no instante da abertura.
\item Determine a tensão máxima sobre a chave.
\item Determine a constante de tempo de extinção.
\end{itensq}
\resposta{(b) $\approx\SI{0.7}{\volt}$ acima da alimentação;
(c) $\tau=\SI{20}{\milli\second}$}
\resolucao{(a) Com a chave aberta, a corrente da bobina precisa de um caminho.
O diodo, que estava reversamente polarizado, entra em condução e fecha uma
malha bobina--diodo. A corrente continua circulando e decai exponencialmente,
dissipando-se em $r$ e no diodo --- sem arco.\\
(b) A tensão sobre a bobina fica grampeada em $-V_D\approx\SI{-0.7}{\volt}$, de
modo que a chave suporta apenas a tensão de alimentação acrescida de
$\SI{0.7}{\volt}$. Compare com os quilovolts do caso anterior.\\
(c) $\tau=L/r=0{,}1/5=\SI{20}{\milli\second}$.\\
\textbf{Custo do método:} a extinção fica \emph{lenta}. Em um relé, isso
prolonga o tempo de desarme; em acionamentos rápidos, prefere-se um diodo em
série com um resistor (ou um Zener), que eleva a tensão de grampeamento e
acelera a queda --- trocando proteção máxima por velocidade.}
\end{questao}

\begin{questao}[3]{}
Dois indutores acoplados, $L_1=\SI{4}{\milli\henry}$ e
$L_2=\SI{9}{\milli\henry}$, têm indutância mútua $M=\SI{3}{\milli\henry}$.
\begin{itensq}
\item Determine $L_{eq}$ em série nas duas orientações possíveis.
\item Determine o coeficiente de acoplamento $k$.
\item Explique como identificar experimentalmente a orientação.
\end{itensq}
\resposta{(a) $\SI{19}{\milli\henry}$ e $\SI{7}{\milli\henry}$; (b) $k=0{,}5$}
\resolucao{(a) Em série,
\[
L_{eq}=L_1+L_2\pm2M
\;\Longrightarrow\;
4+9+6=\SI{19}{\milli\henry}
\ \text{ou}\ 4+9-6=\SI{7}{\milli\henry}.
\]
O sinal $+$ vale para \textbf{acoplamento aditivo} (os fluxos se reforçam; a
corrente entra pelos dois pontos marcados); o sinal $-$, para
\textbf{subtrativo}.\\
(b)
\[
k=\frac{M}{\sqrt{L_1L_2}}=\frac{3}{\sqrt{36}}=0{,}5.
\]
Como $0\le k\le1$, este é um acoplamento médio --- típico de bobinas
próximas com núcleo de ar. Transformadores de potência atingem $k>0{,}99$.\\
(c) \textbf{Método experimental:} meça $L_{eq}$ nas duas ligações possíveis. A
\emph{maior} corresponde ao acoplamento aditivo e identifica os pontos.
Além disso, $M$ sai de graça:
\[
M=\frac{L_{aditivo}-L_{subtrativo}}{4}=\frac{19-7}{4}=\SI{3}{\milli\henry}.
\]
É assim que se determina $M$ sem acesso à geometria interna das bobinas.}
\end{questao}

\begin{questao}[3]{}
Uma fonte de $\SI{12}{\volt}$ alimenta, através de $R_1=\SI{20}{\ohm}$, um nó
do qual saem $R_3=\SI{60}{\ohm}$ ao terra e um ramo com
$R_2=\SI{30}{\ohm}$ em série com $L=\SI{0.9}{\henry}$.
\begin{itensq}
\item Determine a resistência de Thévenin vista pelo indutor.
\item Determine $\tau$ e a corrente final no indutor.
\end{itensq}
\resposta{(a) $R_{Th}=\SI{45}{\ohm}$; (b) $\tau=\SI{20}{\milli\second}$ e
$I_f=\SI{0.2}{\ampere}$}
\resolucao{(a) Zerando a fonte (curto), o indutor vê $R_2$ em série com
$R_1\parallel R_3$:
\[
R_{Th}=30+\frac{20\cdot60}{80}=30+15=\SI{45}{\ohm}.
\]
(b) $\tau=L/R_{Th}=0{,}9/45=\SI{20}{\milli\second}$.\\
A tensão de Thévenin sai com o indutor removido (sem corrente em $R_2$, logo
sem queda nele):
\[
V_{Th}=12\cdot\frac{60}{20+60}=\SI{9}{\volt}
\;\Longrightarrow\;
I_f=\frac{9}{45}=\SI{0.2}{\ampere}.
\]
\textbf{Método geral:} em qualquer circuito com \emph{um} elemento armazenador,
substitua todo o restante pelo equivalente de Thévenin visto por ele. O
problema se reduz a um RL série, com $\tau=L/R_{Th}$ e
$I_f=V_{Th}/R_{Th}$ --- e não é preciso resolver a rede inteira em regime
transitório.}
\end{questao}

\begin{questao}[3]{}
Um indutor de $\SI{0.5}{\henry}$ já conduz $\SI{1}{\ampere}$ quando é ligado a
uma fonte que estabeleceria $\SI{3}{\ampere}$ em regime, com
$R=\SI{25}{\ohm}$.
\begin{itensq}
\item Escreva $i(t)$.
\item Determine a corrente após $\SI{20}{\milli\second}$.
\item Determine a tensão inicial sobre o indutor.
\end{itensq}
\resposta{(a) $i(t)=3-2e^{-t/\tau}$ com $\tau=\SI{20}{\milli\second}$;
(b) $\SI{2.26}{\ampere}$; (c) $\SI{50}{\volt}$}
\resolucao{(a) A forma geral do transitório de primeira ordem é
\[
i(t)=I_f+\left(I_0-I_f\right)e^{-t/\tau},
\]
com $I_0=\SI{1}{\ampere}$, $I_f=\SI{3}{\ampere}$ e
$\tau=0{,}5/25=\SI{20}{\milli\second}$:
\[
i(t)=3-2e^{-t/0{,}020}.
\]
(b) $i(0{,}020)=3-2e^{-1}=3-0{,}736=\SI{2.26}{\ampere}$.\\
(c) Em $t=0^{+}$, $i=\SI{1}{\ampere}$ e o resistor cai
$25(1)=\SI{25}{\volt}$; a fonte vale $RI_f=25(3)=\SI{75}{\volt}$. Logo
\[
v_L(0^{+})=75-25=\SI{50}{\volt}.
\]
\textbf{Generalidade:} a fórmula de (a) cobre carga, descarga e condições
iniciais quaisquer --- basta identificar $I_0$, $I_f$ e $\tau$. Não é preciso
memorizar três expressões diferentes.}
\end{questao}

\begin{questao}[3]{}
Em um circuito RL, mediu-se corrente final de $\SI{2}{\ampere}$ com fonte de
$\SI{50}{\volt}$, e verificou-se que a corrente atingiu $\SI{1.264}{\ampere}$
em $\SI{8}{\milli\second}$. Determine $R$ e $L$.
\resposta{$R=\SI{25}{\ohm}$; $L=\SI{0.2}{\henry}$}
\resolucao{\textbf{Resistência:} do regime permanente,
$R=V/I_f=50/2=\SI{25}{\ohm}$.\\
\textbf{Constante de tempo:} note que
$1{,}264/2=0{,}632$ --- exatamente a fração característica de $1\tau$. Logo
$\tau=\SI{8}{\milli\second}$, e
\[
L=R\tau=25(0{,}008)=\SI{0.2}{\henry}.
\]
\textbf{Se a fração não fosse reconhecível}, o caminho geral seria
$\tau=-t/\ln(1-i/I_f)$.\\
\textbf{Método:} o regime permanente entrega $R$; o transitório entrega $\tau$
e, por consequência, $L$. Duas medidas de naturezas diferentes são necessárias,
porque nenhuma delas isola $L$ sozinha.}
\end{questao}

\begin{questao}[3]{}
Analise o sinal da potência instantânea $p=vi$ em um indutor durante um ciclo
completo de carga e descarga, e relacione-o com a energia armazenada.
\resposta{$p>0$ na carga (absorve), $p<0$ na descarga (devolve); a integral em
um ciclo completo é nula}
\resolucao{\textbf{Carga:} $i>0$ e crescente $\Rightarrow$
$v=L\,\mathrm{d}i/\mathrm{d}t>0$ $\Rightarrow$ $p>0$: o indutor \emph{absorve}
energia da fonte e a armazena no campo.\\
\textbf{Descarga:} $i>0$ mas decrescente $\Rightarrow$ $v<0$ $\Rightarrow$
$p<0$: o indutor \emph{fornece} energia ao circuito.\\
\textbf{Integral:} como
\[
\int p\,\mathrm{d}t=\int Li\,\frac{\mathrm{d}i}{\mathrm{d}t}\,\mathrm{d}t
=\int Li\,\mathrm{d}i=\frac12Li^{2},
\]
a energia depende \emph{apenas} da corrente instantânea, não do caminho
percorrido. Se o ciclo retorna à mesma corrente, a integral é nula: o indutor
ideal não consome energia líquida.\\
\textbf{Consequência:} indutor e capacitor são elementos \textbf{reativos} ---
trocam energia com o circuito sem consumi-la. Só a resistência dissipa. Essa
distinção é a base de todo o tratamento de potência ativa e reativa em corrente
alternada.}
\end{questao}

\begin{questao}[3]{}
Determine a densidade de energia magnética em um núcleo onde
$B=\SI{1.2}{\tesla}$, supondo $\mu\approx\mu_0$ no entreferro.
\begin{itensq}
\item Calcule $u$.
\item Compare com a densidade de energia elétrica em um campo de
      $\SI{3}{\mega\volt\per\meter}$.
\end{itensq}
\dados{$\mu_0=4\pi\times10^{-7}\,\si{\henry\per\meter}$;
$\varepsilon_0=\pot{8.85}{-12}\,\si{\farad\per\meter}$}
\resposta{(a) $u\approx\SI{573}{\kilo\joule\per\meter\cubed}$;
(b) $\approx\SI{40}{\joule\per\meter\cubed}$}
\resolucao{(a)
\[
u_B=\frac{B^{2}}{2\mu_0}=\frac{(1{,}2)^{2}}{2(4\pi\times10^{-7})}
=\frac{1{,}44}{\pot{2.513}{-6}}
\approx\pot{5.73}{5}\,\si{\joule\per\meter\cubed}.
\]
(b)
\[
u_E=\frac12\varepsilon_0E^{2}
=\frac12(\pot{8.85}{-12})(3\times10^{6})^{2}
\approx\SI{40}{\joule\per\meter\cubed}.
\]
\textbf{Razão: cerca de $14\,000$ vezes.} O campo elétrico usado é o limite de
ruptura do ar, e o magnético é típico de aço-silício --- ou seja, comparam-se
dois valores \emph{praticáveis}.\\
\textbf{Consequência tecnológica:} armazenar energia magneticamente é muito
mais denso do que eletrostaticamente. É por isso que transformadores e motores
operam com campos magnéticos, e por que capacitores só competem em energia
quando se usa dielétrico de altíssima constante e espessura nanométrica ---
como nos supercapacitores.}
\end{questao}

\begin{questao}[3]{}
Duas bobinas têm $L_1=\SI{50}{\milli\henry}$ e $L_2=\SI{200}{\milli\henry}$.
\begin{itensq}
\item Determine o valor máximo possível de $M$.
\item Se $M=\SI{80}{\milli\henry}$, determine $k$.
\item Interprete fisicamente $k=1$ e $k=0$.
\end{itensq}
\resposta{(a) $M_{\max}=\SI{100}{\milli\henry}$; (b) $k=0{,}8$}
\resolucao{(a) Como $k\le1$,
\[
M_{\max}=\sqrt{L_1L_2}=\sqrt{50\cdot200}=\sqrt{10\,000}
=\SI{100}{\milli\henry}.
\]
(b) $k=80/100=0{,}8$ --- acoplamento forte, típico de bobinas sobre um mesmo
núcleo ferromagnético com pequeno entreferro.\\
(c) $k=1$ significa \textbf{acoplamento perfeito}: \emph{todo} o fluxo produzido
por uma bobina concatena a outra, sem dispersão. É o transformador ideal ---
inatingível, mas aproximado em núcleos toroidais fechados de alta
permeabilidade.\\
$k=0$ significa \textbf{desacoplamento total}: nenhuma influência mútua. Obtém-se
afastando as bobinas ou dispondo seus eixos perpendicularmente, de modo que o
fluxo de uma atravesse a outra tangencialmente. É a razão de indutores vizinhos
em placas serem montados em orientações ortogonais.}
\end{questao}

\begin{questao}[3]{}
Um indutor com núcleo apresenta $L=\SI{100}{\micro\henry}$ até
$\SI{3}{\ampere}$; acima disso, a saturação reduz $L$ a
$\SI{20}{\micro\henry}$. Ele opera com $\SI{12}{\volt}$ aplicados durante
$\SI{10}{\micro\second}$ a partir de $\SI{2}{\ampere}$.
\begin{itensq}
\item Determine a corrente ao final do pulso, ignorando a saturação.
\item Refaça considerando a saturação.
\item Comente a consequência prática.
\end{itensq}
\resposta{(a) $\SI{3.2}{\ampere}$; (b) $\SI{5}{\ampere}$; (c) risco de
destruição do interruptor}
\resolucao{(a) Com $L$ constante,
\[
\Delta i=\frac{v\,\Delta t}{L}=\frac{12\cdot10^{-5}}{10^{-4}}
=\SI{1.2}{\ampere}
\;\Longrightarrow\;
i_{final}=2+1{,}2=\SI{3.2}{\ampere}.
\]
(b) A corrente sobe de $2$ a $\SI{3}{\ampere}$ com
$L=\SI{100}{\micro\henry}$, o que consome
\[
\Delta t_1=\frac{L\,\Delta i}{v}=\frac{10^{-4}\cdot1}{12}
=\SI{8.33}{\micro\second}.
\]
Restam $\SI{1.67}{\micro\second}$, agora com $L=\SI{20}{\micro\henry}$:
\[
\Delta i_2=\frac{12\cdot\pot{1.67}{-6}}{\pot{2}{-5}}=\SI{1}{\ampere}
\;\Longrightarrow\;
i_{final}=\SI{5}{\ampere}.
\]
(c) A corrente real é $56\%$ maior que a prevista pelo modelo linear, e a
subida final é \emph{cinco vezes} mais íngreme. Em um conversor chaveado, isso
significa que o transistor pode ver o dobro da corrente esperada em uma fração
do tempo --- destruição por sobrecorrente antes que qualquer proteção lenta
atue.\\
\textbf{Regra de projeto:} dimensione o indutor para que a corrente de pico
fique confortavelmente abaixo do joelho de saturação, e verifique
$L$ \emph{na corrente de operação}, nunca no valor de catálogo medido a
corrente nula.}
\end{questao}

\begin{questao}[3]{}
Uma fonte de \textbf{corrente} de $\SI{2}{\ampere}$ alimenta um indutor de
$\SI{0.5}{\henry}$ em paralelo com $R=\SI{10}{\ohm}$, em regime permanente.
\begin{itensq}
\item Determine a corrente no indutor e no resistor.
\item A fonte é subitamente desligada (aberta). Determine a tensão inicial
      sobre o resistor.
\item Determine o tempo de extinção.
\end{itensq}
\resposta{(a) $\SI{2}{\ampere}$ no indutor, $\SI{0}{\ampere}$ no resistor;
(b) $\SI{-20}{\volt}$; (c) $\tau=\SI{50}{\milli\second}$}
\resolucao{(a) Em regime CC o indutor ideal é um \textbf{curto}: ele desvia
toda a corrente, e o resistor em paralelo fica com tensão nula --- logo
corrente nula. Todos os $\SI{2}{\ampere}$ passam por $L$.\\
(b) Ao abrir a fonte, a corrente do indutor \emph{não pode saltar}: continua
$\SI{2}{\ampere}$, e o único caminho disponível é o resistor. A corrente agora
o percorre em sentido oposto ao anterior:
\[
v_R=-RI_0=-10(2)=\SI{-20}{\volt}.
\]
(c) $\tau=L/R=0{,}5/10=\SI{50}{\milli\second}$; a extinção é praticamente
completa em $5\tau=\SI{250}{\milli\second}$.\\
\textbf{Observação:} o resistor, que parecia inútil em regime (corrente nula),
é justamente o que evita a sobretensão destrutiva na abertura. Ele é um
\emph{caminho de escape} projetado --- exatamente a função do diodo de roda
livre, aqui desempenhada por um resistor.}
\end{questao}

\blocodesafio

\begin{questao}[4]{}
\textbf{(Dedução --- resposta do circuito RL.)} Um circuito RL série com fonte
$V$ tem a chave fechada em $t=0$, com $i(0)=0$.
\begin{itensq}
\item Escreva a equação diferencial e resolva-a por separação de variáveis.
\item Identifique $\tau$ e o valor final.
\item Obtenha $v_L(t)$ e verifique a LKT em $t=0^{+}$ e $t\to\infty$.
\end{itensq}
\resposta{$i(t)=\dfrac{V}{R}\left(1-e^{-Rt/L}\right)$;
$\tau=L/R$; $v_L=Ve^{-Rt/L}$}
\resolucao{(a) Pela LKT,
\[
V=Ri+L\frac{\mathrm{d}i}{\mathrm{d}t}
\;\Longrightarrow\;
\frac{\mathrm{d}i}{V/R-i}=\frac{R}{L}\,\mathrm{d}t.
\]
Integrando de $0$ a $t$ (com $i(0)=0$):
\[
-\ln\!\left(\frac{V/R-i}{V/R}\right)=\frac{R}{L}t
\;\Longrightarrow\;
i(t)=\frac{V}{R}\left(1-e^{-Rt/L}\right).
\]
(b) O expoente deve ser adimensional, o que identifica
$\tau=L/R$; e $i(\infty)=V/R$, coerente com o indutor tornando-se um curto.\\
(c)
\[
v_L=L\frac{\mathrm{d}i}{\mathrm{d}t}
=L\cdot\frac{V}{R}\cdot\frac{R}{L}e^{-Rt/L}=V\,e^{-Rt/L}.
\]
\emph{Em $t=0^{+}$:} $v_L=V$ e $Ri=0$; soma $=V$ \checkmark.\\
\emph{Em $t\to\infty$:} $v_L=0$ e $Ri=V$; soma $=V$ \checkmark.\\
\textbf{Estrutura geral:} a solução é (regime permanente) $+$ (transitório
exponencial). Essa decomposição vale para \emph{qualquer} circuito de primeira
ordem e dispensa refazer a integração a cada problema.}
\end{questao}

\begin{questao}[4]{}
\textbf{(Integração --- energia na descarga.)} Um indutor com corrente inicial
$I_0$ é descarregado através de $R$.
\begin{itensq}
\item Escreva $i(t)$ e a potência instantânea no resistor.
\item Integre para obter a energia total dissipada.
\item Mostre que o resultado independe de $R$ e interprete.
\end{itensq}
\resposta{$W=\frac12LI_0^{2}$, independente de $R$}
\resolucao{(a) Sem fonte, $Ri+L\,\mathrm{d}i/\mathrm{d}t=0$, cuja solução é
$i(t)=I_0e^{-t/\tau}$ com $\tau=L/R$. Então
\[
p_R(t)=Ri^{2}=RI_0^{2}e^{-2t/\tau}.
\]
(b)
\[
W=\int_0^{\infty}RI_0^{2}e^{-2t/\tau}\,\mathrm{d}t
=RI_0^{2}\left[-\frac{\tau}{2}e^{-2t/\tau}\right]_0^{\infty}
=\frac{RI_0^{2}\tau}{2}.
\]
Substituindo $\tau=L/R$:
\[
W=\frac{RI_0^{2}}{2}\cdot\frac{L}{R}=\frac12LI_0^{2}.
\]
(c) O $R$ \textbf{cancela} exatamente: ele aparece na potência
($\propto R$) e desaparece na duração ($\tau\propto1/R$). Aumentar $R$ dez
vezes decuplica a potência inicial e divide o tempo por dez --- a área sob a
curva não muda.\\
\textbf{Significado:} a energia dissipada é determinada pelo \emph{estado
inicial} do indutor, não pelo caminho de descarga. É a mesma lógica que faz um
capacitor descarregado entregar sempre $\frac12CV_0^{2}$, e confirma por
integração o argumento de conservação usado no bloco anterior.}
\end{questao}

\begin{questao}[4]{}
\textbf{(Projeto de proteção.)} Uma bobina de relé com $L=\SI{0.2}{\henry}$ e
$r=\SI{100}{\ohm}$ opera em $\SI{24}{\volt}$. O transistor de comando suporta
no máximo $\SI{60}{\volt}$.
\begin{itensq}
\item Determine a corrente de operação e a energia armazenada.
\item Dimensione um grampeador resistivo (diodo em série com $R_g$) que
      respeite o limite do transistor.
\item Compare o tempo de extinção com o do diodo simples.
\end{itensq}
\resposta{(a) $\SI{0.24}{\ampere}$, $\SI{5.76}{\milli\joule}$;
(b) $R_g\le\SI{150}{\ohm}$; (c) $\SI{0.8}{\milli\second}$ contra
$\SI{2}{\milli\second}$}
\resolucao{(a) $I=24/100=\SI{0.24}{\ampere}$;
$W=\frac12(0{,}2)(0{,}24)^{2}=\SI{5.76}{\milli\joule}$.\\
(b) Na abertura, a corrente $\SI{0.24}{\ampere}$ é forçada pelo diodo e por
$R_g$. A tensão no coletor sobe para $V+R_gI$ (desprezando o diodo). Exigindo
$\le\SI{60}{\volt}$:
\[
24+R_g(0{,}24)\le60
\;\Longrightarrow\;
R_g\le\frac{36}{0{,}24}=\SI{150}{\ohm}.
\]
Adote $R_g=\SI{150}{\ohm}$ e verifique a dissipação: o pulso entrega
$\SI{5.76}{\milli\joule}$ e, a $\SI{10}{\hertz}$ de comutação, isso são
$\SI{57.6}{\milli\watt}$ médios --- um resistor de
$\frac14\,\si{\watt}$ basta, desde que suporte o pico de
$0{,}24^{2}(150)=\SI{8.6}{\watt}$ instantâneos.\\
(c) \emph{Com grampeador:} $\tau=L/(r+R_g)=0{,}2/250=\SI{0.8}{\milli\second}$.\\
\emph{Com diodo simples:} $\tau=L/r=0{,}2/100=\SI{2}{\milli\second}$.\\
\textbf{Compromisso quantificado:} o grampeador reduz o tempo de desarme em
$60\%$, ao custo de submeter o transistor a $\SI{60}{\volt}$ em vez de
$\SI{24.7}{\volt}$. Quanto maior $R_g$, mais rápido o relé desarma e maior a
tensão no interruptor --- o projeto consiste em usar toda a margem do
transistor, e nada além dela.}
\end{questao}

\begin{questao}[4]{}
\textbf{(Projeto de armazenamento.)} Deseja-se um indutor que armazene
$\SI{5}{\joule}$ conduzindo $\SI{10}{\ampere}$.
\begin{itensq}
\item Determine $L$.
\item Para um núcleo toroidal com $A=\SI{10}{\centi\meter\squared}$,
      $\ell=\SI{0.25}{\meter}$ e $\mu_r=2000$, determine o número de espiras.
\item Verifique a densidade de fluxo e avalie a viabilidade.
\end{itensq}
\resposta{(a) $L=\SI{0.1}{\henry}$; (b) $N\approx35$ espiras;
(c) $B\approx\SI{3.5}{\tesla}$ --- \textbf{inviável}, o núcleo satura}
\resolucao{(a) $W=\frac12LI^{2}\Rightarrow
L=\dfrac{2W}{I^{2}}=\dfrac{10}{100}=\SI{0.1}{\henry}$.\\
(b) De $L=\dfrac{\mu_r\mu_0N^{2}A}{\ell}$:
\[
N=\sqrt{\frac{L\ell}{\mu_r\mu_0A}}
=\sqrt{\frac{0{,}1\cdot0{,}25}{2000(4\pi\times10^{-7})(10^{-3})}}
=\sqrt{9{,}95}\approx35\ \text{espiras}.
\]
(c) O fluxo é
\[
B=\frac{\mu_r\mu_0NI}{\ell}
=\frac{2000(4\pi\times10^{-7})(35)(10)}{0{,}25}\approx\SI{3.5}{\tesla}.
\]
Materiais ferromagnéticos comuns saturam entre $1{,}2$ e
$\SI{1.8}{\tesla}$. O projeto é \textbf{inviável}: muito antes de atingir
$\SI{10}{\ampere}$, o núcleo satura, $\mu_r$ despenca e $L$ some.\\
\textbf{Diagnóstico e correção.} A energia armazenada em material
ferromagnético é limitada justamente porque $\mu_r$ alto significa $B$ alto para
pouca corrente. A solução clássica é introduzir um \textbf{entreferro}: ele
reduz a permeabilidade efetiva (exigindo mais espiras), mas permite corrente
muito maior antes da saturação --- e a maior parte da energia passa a residir
no entreferro, onde $u=B^{2}/2\mu_0$ é enorme.\\
\textbf{Lição:} calcular $L$ e $N$ não basta. Sem verificar $B$, obtém-se um
projeto aritmeticamente correto e fisicamente impossível.}
\end{questao}

\begin{questao}[4]{}
\textbf{(Projeto de constante de tempo.)} Deseja-se $\tau=\SI{20}{\milli\second}$
com $R=\SI{50}{\ohm}$.
\begin{itensq}
\item Determine $L$.
\item Avalie a viabilidade física dessa indutância.
\item Proponha duas alternativas para obter o mesmo $\tau$ de forma prática.
\end{itensq}
\resposta{(a) $L=\SI{1}{\henry}$; (c) aumentar $R$, ou usar RC equivalente}
\resolucao{(a) $L=\tau R=0{,}020(50)=\SI{1}{\henry}$.\\
(b) \textbf{Um henry é um componente grande.} Exige núcleo ferromagnético e
centenas de espiras; a resistência do próprio enrolamento facilmente atinge
dezenas de ohms --- comparável aos $\SI{50}{\ohm}$ do projeto, o que
\emph{altera} o $\tau$ pretendido. Some-se a isso massa, custo, sensibilidade a
campos externos e saturação.\\
(c) \textbf{Alternativa 1 --- elevar $R$.} Com $R=\SI{500}{\ohm}$,
$L=\SI{10}{\henry}$... pior. O caminho correto é o inverso: fixar um $L$
razoável e ajustar $R$. Com $L=\SI{10}{\milli\henry}$, obter
$\tau=\SI{20}{\milli\second}$ exigiria $R=\SI{0.5}{\ohm}$ --- baixo demais para
ser prático em circuitos de sinal.\\
\textbf{Alternativa 2 --- usar RC.} Um capacitor de
$\SI{20}{\micro\farad}$ com $\SI{1}{\kilo\ohm}$ dá $\tau=\SI{20}{\milli\second}$
em um componente barato, pequeno e estável.\\
\textbf{Conclusão de engenharia:} para constantes de tempo da ordem de
milissegundos ou mais, \textbf{RC vence RL} quase sempre. Indutores são
preferíveis quando se precisa de armazenamento de energia com corrente elevada
ou de comportamento em frequência --- não para gerar atrasos.}
\end{questao}

\begin{questao}[4]{}
\textbf{(Demonstração --- acoplamento em série.)} Duas bobinas acopladas em
série conduzem a mesma corrente $i$.
\begin{itensq}
\item Deduza $L_{eq}=L_1+L_2\pm2M$ a partir das tensões induzidas.
\item Demonstre que $M\le\sqrt{L_1L_2}$, usando o fato de a energia ser não
      negativa.
\item Explique o significado da convenção dos pontos.
\end{itensq}
\resposta{(b) da condição $W\ge0$ para toda corrente segue
$M^{2}\le L_1L_2$}
\resolucao{(a) Cada bobina sofre autoindução \emph{e} indução mútua:
\[
v_1=L_1\frac{\mathrm{d}i}{\mathrm{d}t}\pm M\frac{\mathrm{d}i}{\mathrm{d}t},
\qquad
v_2=L_2\frac{\mathrm{d}i}{\mathrm{d}t}\pm M\frac{\mathrm{d}i}{\mathrm{d}t}.
\]
Somando e comparando com
$v=L_{eq}\,\mathrm{d}i/\mathrm{d}t$:
\[
L_{eq}=L_1+L_2\pm2M.
\]
(b) Com correntes independentes $i_1$ e $i_2$, a energia do par é
\[
W=\frac12L_1i_1^{2}+\frac12L_2i_2^{2}+Mi_1i_2.
\]
Como um sistema passivo não pode armazenar energia negativa, $W\ge0$ para
\emph{quaisquer} $i_1,i_2$. Tomando $i_1=x$ e $i_2=-1$:
\[
\frac12L_1x^{2}-Mx+\frac12L_2\ge0\quad\forall x,
\]
o que exige discriminante não positivo:
\[
M^{2}-L_1L_2\le0
\;\Longrightarrow\;
M\le\sqrt{L_1L_2}
\;\Longrightarrow\;
k=\frac{M}{\sqrt{L_1L_2}}\le1.
\]
(c) Os pontos marcam terminais \textbf{magneticamente correspondentes}:
correntes que entram pelos terminais pontuados produzem fluxos que se
\emph{reforçam}. Se ambas as correntes entram (ou ambas saem) pelos pontos, o
termo mútuo é positivo; caso contrário, negativo.\\
\textbf{Por que a convenção é necessária:} o esquema elétrico não informa o
sentido de enrolamento das bobinas, que é uma propriedade \emph{física} da
construção. Sem os pontos, $L_{eq}$ seria ambíguo entre
$\SI{19}{\milli\henry}$ e $\SI{7}{\milli\henry}$ --- e, em um transformador, a
polaridade da saída ficaria indeterminada.}
\end{questao}

\begin{questao}[4]{}
\textbf{(Indutor real em regime.)} Um indutor com $L=\SI{0.5}{\henry}$ e
$r=\SI{2}{\ohm}$ é alimentado por $\SI{10}{\volt}$ CC.
\begin{itensq}
\item Determine a corrente, a energia armazenada e a potência dissipada em
      regime.
\item Determine o tempo necessário para atingir $99\%$ da corrente final.
\item Discuta o "custo de manutenção" da energia armazenada.
\end{itensq}
\resposta{(a) $\SI{5}{\ampere}$, $\SI{6.25}{\joule}$, $\SI{50}{\watt}$;
(b) $\approx\SI{1.15}{\second}$}
\resolucao{(a) $I=10/2=\SI{5}{\ampere}$;
$W=\frac12(0{,}5)(25)=\SI{6.25}{\joule}$;
$P=rI^{2}=2(25)=\SI{50}{\watt}$.\\
(b) $\tau=L/r=0{,}5/2=\SI{0.25}{\second}$, e $99\%$ correspondem a
$t=\tau\ln100=0{,}25(4{,}605)=\SI{1.15}{\second}$.\\
(c) \textbf{O contraste é gritante.} Carregar o campo custou
$\SI{6.25}{\joule}$; mantê-lo custa $\SI{50}{\watt}$ \emph{continuamente}.
Em apenas $\SI{0.125}{\second}$ de operação, a dissipação já iguala toda a
energia armazenada; em um minuto, dissipam-se $\SI{3000}{\joule}$ --- $480$
vezes o conteúdo do campo.\\
\textbf{Figura de mérito:} a razão $W/P=L/(2r)=\tau/2$ tem dimensão de tempo e
mede a "qualidade" do indutor em CC. Aqui vale $\SI{0.125}{\second}$.\\
\textbf{Consequência:} indutores são péssimos \emph{reservatórios} de energia
para armazenamento prolongado --- ao contrário de baterias ou capacitores, que
não exigem corrente para manter o estado. Eles são excelentes para
\emph{transferir} energia em ciclos curtos, que é exatamente o que fazem em
conversores chaveados, onde o ciclo dura microssegundos.}
\end{questao}

\begin{questao}[4]{}
\textbf{(Conversor chaveado.)} Um conversor opera a $\SI{50}{\kilo\hertz}$ com
razão cíclica $D=0{,}5$, aplicando $\SI{12}{\volt}$ sobre um indutor de
$\SI{100}{\micro\henry}$ durante a condução.
\begin{itensq}
\item Determine a ondulação de corrente pico a pico.
\item Determine a corrente de pico se a componente contínua for
      $\SI{3}{\ampere}$.
\item O indutor satura em $\SI{3.5}{\ampere}$. Proponha e quantifique duas
      correções.
\end{itensq}
\resposta{(a) $\Delta i=\SI{1.2}{\ampere}$; (b) $I_{pico}=\SI{3.6}{\ampere}$;
(c) satura --- elevar $L$ ou a frequência}
\resolucao{(a) Durante o tempo de condução
$t_{on}=D/f=0{,}5/\pot{5}{4}=\SI{10}{\micro\second}$:
\[
\Delta i=\frac{V\,t_{on}}{L}=\frac{12(10^{-5})}{10^{-4}}=\SI{1.2}{\ampere}.
\]
(b) A ondulação se distribui simetricamente em torno da média:
\[
I_{pico}=3+\frac{1{,}2}{2}=\SI{3.6}{\ampere}\ >\ \SI{3.5}{\ampere}.
\]
\textbf{O indutor satura} --- e note que a corrente \emph{média}
($\SI{3}{\ampere}$) estava confortavelmente abaixo do limite. Dimensionar pela
média é o erro clássico.\\
(c) Como $\Delta i=\dfrac{VD}{fL}$, há dois caminhos:\\
\textbf{(i) Aumentar $L$.} Para $I_{pico}\le3{,}5$ é preciso
$\Delta i\le\SI{1}{\ampere}$, ou seja
\[
L\ge\frac{12(10^{-5})}{1}=\SI{120}{\micro\henry}.
\]
Custo: indutor maior e mais caro.\\
\textbf{(ii) Aumentar a frequência.} Mantendo
$L=\SI{100}{\micro\henry}$, exige-se
\[
f\ge\frac{12(0{,}5)}{10^{-4}(1)}=\SI{60}{\kilo\hertz}.
\]
Custo: mais perdas de comutação no transistor e maior interferência
eletromagnética.\\
\textbf{Compromisso:} $\Delta i\propto1/(fL)$ --- o produto $fL$ é o que
importa. Elevar a frequência permite indutores menores, e é exatamente essa a
tendência que tornou as fontes chaveadas modernas tão compactas.}
\end{questao}

\begin{questao}[4]{}
\textbf{(Demonstração --- continuidade da corrente.)} Prove que a corrente em
um indutor ideal não pode sofrer descontinuidade, usando um argumento
energético.
\begin{itensq}
\item Suponha um salto de $I_1$ para $I_2$ em intervalo $\Delta t\to0$ e
      analise a tensão.
\item Analise a potência e a energia envolvidas.
\item Explique por que capacitores em paralelo apresentam o problema dual.
\end{itensq}
\resposta{Um salto exigiria tensão e potência infinitas; a energia é finita,
mas a taxa de transferência não pode ser}
\resolucao{(a) Um salto implica
\[
v=L\lim_{\Delta t\to0}\frac{I_2-I_1}{\Delta t}\to\infty.
\]
Nenhuma fonte real fornece tensão infinita.\\
(b) A potência $p=vi$ também diverge. A \emph{energia} necessária é finita
---
\[
\Delta W=\frac12L\left(I_2^{2}-I_1^{2}\right),
\]
--- mas ela teria de ser transferida em tempo nulo, exigindo potência
infinita. \textbf{É a taxa, não a quantidade, que torna o salto impossível.}\\
(c) \textbf{Problema dual:} dois capacitores com tensões diferentes ligados em
paralelo exigiriam salto de \emph{tensão}, o que demandaria corrente infinita
($i=C\,\mathrm{d}v/\mathrm{d}t$). Na prática, a resistência inevitável dos
condutores limita a corrente e dissipa a diferença de energia --- e demonstra-se
que a perda independe do valor dessa resistência, resultado análogo ao da
descarga do indutor.\\
\textbf{Regra prática:} nunca ligue em série indutores com correntes distintas,
nem em paralelo capacitores com tensões distintas. Em ambos os casos, o modelo
ideal falha e o circuito real encontra uma saída --- arco elétrico no primeiro
caso, pico de corrente no segundo.}
\end{questao}

\begin{questao}[4]{}
\textbf{(Modelagem --- sensor indutivo.)} Um sensor de proximidade usa uma
bobina cuja indutância varia com a distância $x$ a um alvo ferromagnético,
segundo $L(x)=L_0\left(1+\dfrac{a}{x}\right)$, com
$L_0=\SI{10}{\milli\henry}$ e $a=\SI{1}{\milli\meter}$.
\begin{itensq}
\item Determine $L$ para $x=\SI{1}{\milli\meter}$ e $x=\SI{5}{\milli\meter}$.
\item Obtenha a sensibilidade $\mathrm{d}L/\mathrm{d}x$ e avalie nos dois
      pontos.
\item Discuta a consequência para o projeto do sensor.
\end{itensq}
\resposta{(a) $\SI{20}{\milli\henry}$ e $\SI{12}{\milli\henry}$;
(b) $-\SI{10}{\milli\henry\per\milli\meter}$ e
$-\SI{0.4}{\milli\henry\per\milli\meter}$}
\resolucao{(a)
\[
L(1)=10(1+1)=\SI{20}{\milli\henry};
\qquad
L(5)=10(1+0{,}2)=\SI{12}{\milli\henry}.
\]
(b)
\[
\frac{\mathrm{d}L}{\mathrm{d}x}=-\frac{L_0a}{x^{2}}.
\]
Em $x=\SI{1}{\milli\meter}$:
$-10(1)/1=-\SI{10}{\milli\henry\per\milli\meter}$.\\
Em $x=\SI{5}{\milli\meter}$:
$-10(1)/25=-\SI{0.4}{\milli\henry\per\milli\meter}$.\\
A sensibilidade cai com $1/x^{2}$: é \textbf{25 vezes menor} a
$\SI{5}{\milli\meter}$.\\
(c) \textbf{Três consequências.} \emph{(i)} O sensor é fortemente
\textbf{não linear} --- a leitura bruta de $L$ não é proporcional a $x$, e a
conversão exige linearização ou tabela. \emph{(ii)} A \textbf{faixa útil} é
curta: além de poucos milímetros, a variação se confunde com a deriva térmica
da própria bobina. \emph{(iii)} A \textbf{resolução} depende da posição: perto
do alvo, décimos de micrômetro são detectáveis; longe, nem décimos de
milímetro.\\
\textbf{Solução usual:} operar sempre na região próxima e usar o sensor como
detector de \emph{limiar} (presença/ausência), em vez de medidor de distância
--- que é exatamente como os sensores indutivos industriais são
especificados.}
\end{questao}

% END BANCO COMPLEMENTAR

\gabarito[10.2 Indutores]
 % Capacitores e indutores

% Cap. 11 ------------------------------------------------------------
\setcounter{section}{10}
% =====================================================================
%  CAPITULO 11 - ELETROMAGNETISMO
% =====================================================================
\section{Eletromagnetismo}

\begin{conceito}[Antes de começar]
\textbf{Campo de um fio reto longo:}
$B = \dfrac{\mu_0\,I}{2\pi\,r}$.\par
\textbf{Campo no centro de uma espira circular:}
$B = \dfrac{\mu_0\,I}{2\,R}$.\par
\textbf{Campo no interior de um solenoide:}
$B = \mu_0\,n\,I$, onde $n = N/\ell$ é o número de espiras por metro.\par
\textbf{Fluxo magnético:} $\Phi = B\,A\cos\theta$, medido em webers (Wb).\par
\textbf{Força sobre um condutor:} $F = B\,I\,\ell\sin\theta$.\par
\textbf{Força sobre uma carga em movimento:} $F = q\,v\,B\sin\theta$.\par
\textbf{Lei de Faraday--Lenz:} $\varepsilon = -N\,\dfrac{\mathrm{d}\Phi}{\mathrm{d}t}$;
o sinal negativo (Lenz) indica que a corrente induzida se opõe à variação de fluxo
que a criou. Constante: $\mu_0 = 4\pi\times10^{-7}\,\si{\tesla\meter\per\ampere}$.
\end{conceito}

\legendaniveis

% =====================================================================
\subsection{Conceitos iniciais e materiais magnéticos}
% =====================================================================

\blocofixacao

\begin{questao}[1]{}
O princípio da inseparabilidade dos polos magnéticos garante que:
\begin{alternativas}
\item existem monopolos magnéticos isolados.
\item ao dividir um ímã, obtêm-se dois novos ímãs, cada um com dois polos.
\item o polo norte pode ser isolado do polo sul.
\item as linhas de campo magnético têm início e fim.
\item um ímã pode ter apenas um polo.
\end{alternativas}
\resposta{(b)}
\resolucao{Diferentemente das cargas elétricas, não existem monopolos magnéticos:
ao partir um ímã, cada pedaço volta a ter polos norte e sul. Por isso as linhas de
campo magnético são sempre \emph{fechadas} (sem começo nem fim), o que elimina
(d).}
\end{questao}

\begin{questao}[1]{}
Quanto à permeabilidade magnética relativa $\mu_r$, um material
\textbf{ferromagnético} caracteriza-se por:
\begin{alternativas}
\item $\mu_r \approx 1$.
\item $\mu_r$ ligeiramente menor que 1.
\item $\mu_r$ muito maior que 1 (centenas a milhares).
\item $\mu_r = 0$.
\item $\mu_r$ negativa.
\end{alternativas}
\resposta{(c)}
\resolucao{Ferromagnéticos (ferro, níquel, cobalto) têm $\mu_r$ da ordem de
centenas ou milhares --- por isso concentram intensamente as linhas de campo e são
usados em núcleos de transformadores e motores. Paramagnéticos têm
$\mu_r\gtrsim1$; diamagnéticos, $\mu_r\lesssim1$.}
\end{questao}

\begin{questao}[1]{}
Dois materiais são submetidos ao mesmo campo externo $H$. Um apresenta
$B_1 = \SI{0.002}{\tesla}$ e o outro, $B_2 = \SI{1.8}{\tesla}$. Qual é o
ferromagnético?
\resposta{O de $B_2 = \SI{1.8}{\tesla}$}
\resolucao{Como $B = \mu\,H$, para o mesmo $H$ o material com maior $B$ tem maior
permeabilidade. O valor $B_2 = \SI{1.8}{\tesla}$ é muito superior, indicando
$\mu$ elevado: material \textbf{ferromagnético}. O de $B_1 = \SI{0.002}{\tesla}$ é
paramagnético ou diamagnético (resposta muito fraca ao campo).}
\end{questao}

\begin{questao}[1]{Apostila original revisada}
Dois materiais apresentam permeabilidades relativas
$\mu_{rA}=1{,}0005$ e $\mu_{rB}=0{,}99998$. Classifique cada material como
paramagnético ou diamagnético e descreva, qualitativamente, como reage a um campo
magnético externo.
\resposta{$A$: paramagnético; $B$: diamagnético}
\resolucao{Materiais com $\mu_r$ ligeiramente maior que 1 são paramagnéticos e são
fracamente atraídos pelo campo. Materiais com $\mu_r$ ligeiramente menor que 1 são
diamagnéticos e são fracamente repelidos.}
\end{questao}

\blocoaplicacao

\begin{questao}[2]{}
Explique por que os materiais ferromagnéticos são preferidos na construção de
núcleos de máquinas elétricas (transformadores, motores, geradores).
\resposta{Ver comentário}
\resolucao{Por terem $\mu_r$ muito elevado, os ferromagnéticos \emph{concentram} e
intensificam o campo magnético para uma mesma corrente de excitação, o que aumenta
o fluxo e, portanto, a eficiência do acoplamento magnético (mais tensão induzida,
mais torque). Além disso, oferecem um "caminho" de baixa relutância para o fluxo,
canalizando-o pelo núcleo em vez de dispersá-lo pelo ar. A desvantagem são as
perdas por histerese e correntes parasitas, minimizadas usando núcleos laminados
e ligas específicas (aço-silício).}
\end{questao}

\blocoaprofundamento

\begin{questao}[3]{}
Um mesmo enrolamento é usado com dois núcleos diferentes: um de ar
($\mu_r \approx 1$) e outro de ferro-silício ($\mu_r = 4000$). Mantendo a mesma
corrente:
\begin{itensq}
\item Qual é a razão entre as induções magnéticas $B$ obtidas nos dois casos?
\item Se o núcleo de ferro \emph{saturar} a partir de certo ponto, o que acontece
      com essa razão ao se aumentar muito a corrente?
\item Por que a saturação impõe um limite de projeto às máquinas elétricas?
\end{itensq}
\resposta{(a) $B_{\text{Fe}}/B_{\text{ar}} = 4000$;
(b) a razão \emph{cai}; (c) ver comentário}
\resolucao{(a) Como $B = \mu H = \mu_r\mu_0 H$, com $H$ fixo (mesma corrente e
geometria):
\[
\frac{B_{\text{Fe}}}{B_{\text{ar}}}=\frac{\mu_{r,\text{Fe}}}{\mu_{r,\text{ar}}}
=4000.
\]
O núcleo ferromagnético multiplica o campo por quatro mil --- é o que torna
viáveis transformadores e motores compactos.\\
(b) Na \textbf{saturação}, praticamente todos os domínios magnéticos já estão
alinhados; o material não consegue mais contribuir. A partir daí, aumentos de
corrente produzem acréscimos de $B$ apenas na proporção do ar
($\mu_r \to 1$ incremental). A razão $B_{\text{Fe}}/B_{\text{ar}}$, portanto,
\emph{diminui} progressivamente --- a curva $B\times H$ "achata".\\
(c) Porque o torque de um motor e a tensão de um transformador dependem do fluxo,
e o fluxo não pode crescer indefinidamente. Ultrapassado o joelho da curva de
saturação:
\begin{itemize}[leftmargin=1.6em,itemsep=1pt,topsep=2pt]
\item a corrente de magnetização aumenta acentuadamente (aquecendo o enrolamento) sem ganho
      proporcional de fluxo;
\item a forma de onda distorce, gerando harmônicos;
\item as perdas aumentam desproporcionalmente.
\end{itemize}
Por isso o projetista dimensiona a \emph{seção do núcleo} para operar
\textbf{abaixo} do joelho --- tipicamente em torno de \SI{1.5}{\tesla} para
aço-silício. \emph{Saturação é o limite físico que define o tamanho mínimo de
uma máquina elétrica.}}
\end{questao}

\begin{questao}[3]{}
Um material ferromagnético submetido a um campo alternado percorre um
\textbf{ciclo de histerese}: a curva $B \times H$ não é uma reta e não retorna
pelo mesmo caminho.

\circuitoqq{%
\begin{tikzpicture}
 \begin{axis}[width=8.2cm,height=5.6cm,xlabel={$H$ (A/m)},ylabel={$B$ (T)},
   xmin=-3,xmax=3,ymin=-1.6,ymax=1.6,xtick=\empty,ytick=\empty,
   axis lines=middle,axis line style={-Stealth}]
   \addplot[Icol,very thick,smooth] coordinates
     {(-2.6,-1.35)(-1.6,-1.25)(-0.6,-1.0)(0,-0.75)(0.5,-0.2)(1.0,0.75)(1.7,1.2)(2.6,1.35)};
   \addplot[Vcol,very thick,smooth] coordinates
     {(2.6,1.35)(1.6,1.25)(0.6,1.0)(0,0.75)(-0.5,0.2)(-1.0,-0.75)(-1.7,-1.2)(-2.6,-1.35)};
   \node[font=\footnotesize,Vcol] at (axis cs:0.42,0.95) {$B_r$};
   \node[font=\footnotesize,Icol] at (axis cs:-1.0,-0.35) {$H_c$};
 \end{axis}
\end{tikzpicture}}{Ciclo de histerese de um material ferromagnético.}

\begin{itensq}
\item Identifique o que representam $B_r$ (remanência) e $H_c$ (coercividade).
\item A \emph{área} interna do ciclo tem significado físico. Qual?
\item Compare os requisitos de um material para (i) núcleo de transformador e
      (ii) ímã permanente. Qual deve ter ciclo "largo" e qual "estreito"?
\end{itensq}
\resposta{Ver análise comentada}
\resolucao{\textbf{(a)} $B_r$ é a \textbf{remanência}: a indução que \emph{permanece}
no material quando o campo externo é removido ($H = 0$). É o que faz um prego
continuar magnetizado depois de encostar num ímã.\\
$H_c$ é a \textbf{coercividade}: o campo externo, de sentido oposto, necessário
para \emph{anular} a magnetização ($B = 0$). Quantifica a resistência do material à desmagnetização.

\textbf{(b)} A área interna do ciclo representa a \textbf{energia dissipada por
unidade de volume e por ciclo} --- as chamadas \emph{perdas por histerese}. Cada
inversão do campo exige reorientar os domínios magnéticos, e esse processo é
irreversível, gerando calor. A potência perdida é
\[
P_h = k\,f\,V\,(\text{área do ciclo}),
\]
proporcional à \emph{frequência} $f$: por isso máquinas de alta frequência exigem
materiais especiais.

\textbf{(c) Requisitos opostos.}
\begin{center}\small\sffamily
\begin{tabular}{@{}l l l@{}}
\toprule
 & \textbf{Núcleo de transformador} & \textbf{Ímã permanente}\\
\midrule
Ciclo desejado & \textbf{estreito} (fino) & \textbf{largo} (amplo)\\
$H_c$ & baixa & alta\\
$B_r$ & baixa & alta\\
Objetivo & minimizar perdas & reter magnetização\\
Material típico & ferro-silício, ferrite & alnico, neodímio (NdFeB)\\
Denominação & magneticamente \emph{mole} & magneticamente \emph{duro}\\
\bottomrule
\end{tabular}
\end{center}
\textbf{Raciocínio central:} o transformador inverte o campo 50 ou 60 vezes por
segundo --- cada inversão custa a área do ciclo em calor, logo o ciclo deve ser o
mais estreito possível. O ímã permanente, ao contrário, precisa \emph{resistir} à
desmagnetização: quanto maior a área do ciclo, maior é sua resistência à desmagnetização. \emph{O mesmo
gráfico, lido com objetivos opostos, leva a escolhas opostas de material.}}
\end{questao}

% BEGIN BANCO COMPLEMENTAR Conceitos iniciais e materiais magnéticos
% =====================================================================
%  Banco complementar --- Conceitos Iniciais e Materiais Magneticos
%  REVISADO. Substitui a versao gerada por template (10 clones em N2,
%  9 em N3 e 8 questoes conceituais de uma linha marcadas como N4).
%  O cap11 ja cobre: inseparabilidade dos polos (MC), definicao de
%  ferromagnetico por mu_r (MC), comparacao de dois materiais por B sob
%  mesmo H, classificacao para/diamagnetica a partir de mu_r, por que
%  ferromagneticos sao usados em nucleos, comparacao ar x ferro-silicio
%  no mesmo enrolamento e a descricao qualitativa do ciclo de histerese.
%  A LAMINACAO ja foi tratada quantitativamente no banco de Lenz, e o
%  ENTREFERRO no banco de fluxo -- por isso as N4 antigas sobre esses
%  dois temas foram removidas daqui.
%  Este banco explora: H, B e M como grandezas distintas, curva de
%  magnetizacao inicial e mu diferencial, temperatura de Curie, energia
%  do ciclo, materiais moles x duros com numeros, ponto de operacao em
%  nucleo nao linear, circuito magnetico com dois materiais e ensaio
%  experimental do ciclo.
%  Distribuicao mantida: 4 N1 + 11 N2 + 10 N3 + 8 N4 = 33
% =====================================================================

\subsubsection*{Banco complementar --- Conceitos iniciais e materiais magnéticos}

\begin{atencao}[Organização do banco]
Três grandezas convivem e são frequentemente confundidas. $\vec H$
($\si{\ampere\per\meter}$) é o campo \emph{aplicado}, determinado apenas pelas
correntes de condução; $\vec M$ é a \emph{magnetização} do material; e
$\vec B=\mu_0(\vec H+\vec M)$ é o que efetivamente age sobre cargas e produz
fluxo. Em materiais lineares, $\vec B=\mu_r\mu_0\vec H$ --- mas ferromagnéticos
\textbf{não} são lineares, e neles $\mu_r$ é uma função, não uma constante. O N4
trata de energia do ciclo, ponto de operação e ensaio.
\end{atencao}

\blocofixacao

\begin{questao}[1]{}
Um material linear tem $\mu_r=500$ e é submetido a
$H=\SI{100}{\ampere\per\meter}$. Determine $B$.
\dados{$\mu_0=4\pi\times10^{-7}\,\si{\tesla\meter\per\ampere}$}
\resposta{$B=\SI{62.8}{\milli\tesla}$}
\resolucao{
\[
B=\mu_r\mu_0H=(500)(4\pi\times10^{-7})(100)=\pot{6.28}{-2}\,\si{\tesla}.
\]
\textbf{Sem o material} ($\mu_r=1$), o mesmo $H$ produziria apenas
$\SI{126}{\micro\tesla}$. O material multiplicou $B$ por $500$ --- e é
exatamente esse o papel de um núcleo magnético.\\
\textbf{Atenção:} $H$ é fixado pelas correntes; $B$ é o que o material
\emph{responde}. Trocar o núcleo muda $B$, mas não muda $H$.}
\end{questao}

\begin{questao}[1]{}
Determine a magnetização $M$ do material da questão anterior.
\resposta{$M=\SI{49.9}{\kilo\ampere\per\meter}$}
\resolucao{De $B=\mu_0(H+M)$:
\[
M=\frac{B}{\mu_0}-H=\mu_rH-H=(\mu_r-1)H
=(499)(100)=\pot{4.99}{4}\,\si{\ampere\per\meter}.
\]
\textbf{Interpretação:} a magnetização do material equivale a uma corrente
superficial de $\SI{49.9}{\kilo\ampere\per\meter}$ --- quinhentas vezes maior
que a corrente de condução que a provocou.\\
\textbf{É por isso que o núcleo "amplifica":} ele não cria energia, apenas
alinha os momentos magnéticos já existentes no material, cujo efeito somado
supera em muito o do enrolamento.}
\end{questao}

\begin{questao}[1]{}
Define-se a suscetibilidade magnética como $\chi=\mu_r-1$. Classifique os
materiais com $\chi=+\pot{5}{-4}$ e $\chi=-\pot{2}{-5}$.
\resposta{Paramagnético e diamagnético, respectivamente}
\resolucao{\textbf{Paramagnético} ($\chi>0$, pequeno): os momentos atômicos
alinham-se \emph{fracamente} com o campo, reforçando-o de forma quase
imperceptível. Exemplos: alumínio, platina, oxigênio líquido.\\
\textbf{Diamagnético} ($\chi<0$): o campo induz momentos \emph{opostos}, e o
material é levemente repelido. Exemplos: cobre, água, bismuto, grafite. Todo
material tem componente diamagnética; ela só é observável quando não há
efeito maior a mascará-la.\\
\textbf{Ordens de grandeza:} $|\chi|\sim10^{-5}$ a $10^{-4}$ nesses dois casos,
contra $\chi\sim10^{3}$ em ferromagnéticos --- diferença de sete ordens de
grandeza.}
\end{questao}

\begin{questao}[1]{}
Um enrolamento de $500$ espiras percorrido por $\SI{0.4}{\ampere}$ envolve um
núcleo de caminho médio $\SI{25}{\centi\meter}$. Determine $H$.
\resposta{$H=\SI{800}{\ampere\per\meter}$}
\resolucao{Da lei de Ampère aplicada ao caminho médio,
\[
H=\frac{NI}{\ell}=\frac{500(0{,}4)}{0{,}25}=\SI{800}{\ampere\per\meter}.
\]
\textbf{Observação central:} $H$ \textbf{não depende do material}. Ele é
determinado exclusivamente pelas correntes e pela geometria.\\
Trocar o núcleo de ar por ferro não altera $H$ --- altera $B$, e portanto o
fluxo. É essa separação que torna útil distinguir as duas grandezas.}
\end{questao}

\blocoaplicacao

\begin{questao}[2]{}
Um núcleo de $\SI{10}{\centi\meter}$ de caminho médio, seção
$\SI{1}{\centi\meter\squared}$ e $\mu_r=1000$ recebe $\SI{200}{}$
ampère-espiras.
\begin{itensq}
\item Determine a relutância.
\item Determine o fluxo e a densidade de fluxo.
\item Avalie o resultado.
\end{itensq}
\resposta{$\mathcal{R}=\pot{7.96}{5}\,\si{\per\henry}$;
$\Phi=\SI{0.251}{\milli\weber}$; $B=\SI{2.51}{\tesla}$ --- \textbf{saturado}}
\resolucao{(a)
\[
\mathcal{R}=\frac{\ell}{\mu_r\mu_0A}
=\frac{0{,}10}{(1000)(4\pi\times10^{-7})(\pot{1}{-4})}
=\pot{7.96}{5}\,\si{\per\henry}.
\]
(b)
\[
\Phi=\frac{NI}{\mathcal{R}}=\frac{200}{\pot{7.96}{5}}
=\pot{2.51}{-4}\,\si{\weber};
\qquad
B=\frac{\Phi}{A}=\SI{2.51}{\tesla}.
\]
(c) \textbf{Resultado impossível.} Nenhum aço comum sustenta
$\SI{2.51}{\tesla}$ --- todos saturam entre $1{,}5$ e $\SI{2}{\tesla}$.\\
\textbf{O que realmente ocorreria:} conforme $B$ se aproxima da saturação,
$\mu_r$ desaba de $1000$ para valores próximos de $1$. A relutância dispara e o
fluxo estaciona em torno de $B_{sat}A=\SI{0.15}{\milli\weber}$.\\
\textbf{Regra de método:} calcular com $\mu_r$ constante e \emph{depois}
verificar $B$ contra $B_{sat}$. Se o valor exceder, o modelo linear não se
aplica e é preciso recorrer à curva $B\times H$ real.}
\end{questao}

\begin{questao}[2]{}
Um mesmo enrolamento produz $H=\SI{500}{\ampere\per\meter}$. Determine $B$ para
núcleo de ar, de alumínio ($\mu_r=1{,}00002$) e de ferro-silício
($\mu_r=4000$).
\resposta{$\SI{628}{\micro\tesla}$; $\SI{628}{\micro\tesla}$;
$\SI{2.51}{\tesla}$}
\resolucao{
\[
\text{ar}:\ B=\mu_0(500)=\SI{628}{\micro\tesla};
\]
\[
\text{alumínio}:\ B=(1{,}00002)\mu_0(500)=\SI{628}{\micro\tesla};
\]
\[
\text{ferro-silício}:\ B=(4000)\mu_0(500)=\SI{2.51}{\tesla}
\ \text{(saturaria antes)}.
\]
\textbf{Duas leituras.}\\
\emph{(i)} O alumínio é \emph{indistinguível} do ar para fins práticos --- sua
suscetibilidade de $\pot{2}{-5}$ altera $B$ na quinta casa decimal. É por isso
que se pode enrolar bobinas em carretéis de alumínio sem consequência magnética.\\
\emph{(ii)} O ferro-silício produziria $\SI{2.51}{\tesla}$ \emph{se fosse
linear} --- mas ele satura antes. Com $H=\SI{500}{\ampere\per\meter}$, um
ferro-silício real já está bem dentro do joelho da curva, com $B$ real em torno
de $1{,}5$ a $\SI{1.7}{\tesla}$.}
\end{questao}

\begin{questao}[2]{}
Um material apresenta $B=\SI{1.2}{\tesla}$ para $H=\SI{400}{\ampere\per\meter}$
e $B=\SI{1.5}{\tesla}$ para $H=\SI{1200}{\ampere\per\meter}$.
\begin{itensq}
\item Determine $\mu_r$ em cada ponto.
\item Interprete a variação.
\end{itensq}
\resposta{$\mu_r=2390$ e $995$}
\resolucao{(a)
\[
\mu_r=\frac{B}{\mu_0H}:\quad
\frac{1{,}2}{(4\pi\times10^{-7})(400)}=2390;
\qquad
\frac{1{,}5}{(4\pi\times10^{-7})(1200)}=995.
\]
(b) \textbf{A permeabilidade caiu para menos da metade} enquanto $H$
triplicou.\\
\textbf{Interpretação:} o material está entrando na região de \emph{joelho} da
curva de magnetização. Os domínios já estão em boa parte alinhados, e cada
acréscimo de $H$ produz cada vez menos acréscimo de $B$.\\
\textbf{Consequência prática:} triplicar a corrente aumentou o fluxo em apenas
$25\%$. Operar nessa região é ineficiente --- o cobre gasta muito para
render pouco.\\
\textbf{Lição de fundo:} $\mu_r$ de catálogo é sempre um valor \emph{para uma
condição específica}. Usá-lo fora dessa condição é um erro comum e silencioso.}
\end{questao}

\begin{questao}[2]{}
Explique a diferença entre a permeabilidade \textbf{normal}
($\mu=B/H$) e a permeabilidade \textbf{diferencial}
($\mu_d=\mathrm{d}B/\mathrm{d}H$), e indique quando cada uma se aplica.
\resposta{Normal para o ponto de operação; diferencial para pequenos sinais
sobrepostos}
\resolucao{\textbf{Permeabilidade normal:} razão entre $B$ e $H$ no ponto de
operação. Ela governa o \emph{fluxo total} e, portanto, a relutância e o
dimensionamento do núcleo.\\
\textbf{Permeabilidade diferencial:} inclinação da curva naquele ponto. Ela
governa a resposta a \emph{pequenas variações} sobrepostas --- e portanto a
indutância incremental.\\
\textbf{Onde diferem drasticamente:} perto da saturação, $B/H$ ainda é grande
(centenas), mas $\mathrm{d}B/\mathrm{d}H$ já caiu para quase $\mu_0$. Um
indutor com corrente contínua elevada pode ter indutância incremental dezenas
de vezes menor que a nominal.\\
\textbf{Consequência de projeto:} em filtros com componente contínua --- muito
comuns em fontes chaveadas --- é a permeabilidade \emph{diferencial} que
determina o desempenho, e é ela que precisa ser consultada no catálogo.}
\end{questao}

\begin{questao}[2]{}
Compare a temperatura de Curie do ferro ($\SI{770}{\celsius}$), do níquel
($\SI{358}{\celsius}$) e de uma ferrite típica
($\approx\SI{250}{\celsius}$), e explique o que ocorre acima dela.
\resposta{Acima de $T_C$ o material torna-se paramagnético e perde
praticamente toda a permeabilidade}
\resolucao{\textbf{O que ocorre acima de $T_C$.} A agitação térmica supera a
interação de troca que mantém os momentos atômicos alinhados dentro dos
domínios. A estrutura de domínios se desfaz, e o material passa de
\textbf{ferromagnético} a \textbf{paramagnético}:
\[
\mu_r:\ 10^{3}\text{--}10^{4}\ \longrightarrow\ \approx1.
\]
\textbf{A transição é abrupta} --- trata-se de uma transição de fase de segunda
ordem, e $\mu_r$ despenca em poucos graus em torno de $T_C$.\\
\textbf{Consequências práticas.}
\begin{itemize}
\item \textbf{Ferrites são as mais vulneráveis.} Com $T_C$ próxima de
$\SI{250}{\celsius}$, e temperatura de operação já em torno de
$\SI{100}{\celsius}$, a margem é pequena. Um transformador de fonte chaveada
que superaqueça perde indutância, a corrente dispara e ele se destrói --- fuga
térmica.
\item \textbf{Aplicação deliberada:} panelas de indução com fundo cuja $T_C$ é
próxima da temperatura desejada param de aquecer sozinhas ao atingi-la. É um
termostato sem sensor.
\item \textbf{Desmagnetização por aquecimento:} aquecer um ímã acima de
$T_C$ e resfriá-lo sem campo o deixa completamente desmagnetizado.
\end{itemize}}
\end{questao}

\begin{questao}[2]{}
Um ímã permanente de ferrite tem coeficiente de temperatura de
$\SI{-0.2}{\percent\per\celsius}$ para a remanência. Determine a variação de
$B_r$ entre $\SI{20}{\celsius}$ e $\SI{120}{\celsius}$.
\resposta{Redução de $20\%$}
\resolucao{
\[
\frac{\Delta B_r}{B_r}=(-0{,}002)(100)=-0{,}20=-20\%.
\]
\textbf{Reversível ou não?} Até certo limite, a perda é \textbf{reversível} --- o
ímã recupera a magnetização ao esfriar. Acima de uma temperatura crítica
(bem abaixo de $T_C$), parte da perda torna-se \textbf{irreversível}, porque
domínios se reorientam permanentemente.\\
\textbf{Comparação entre materiais:}
\begin{center}
\begin{tabular}{lcc}
\hline
Material & Coef. de $B_r$ & $T_{\max}$ típica\\
\hline
Ferrite & $\SI{-0.2}{\percent\per\celsius}$ & $\SI{250}{\celsius}$\\
Neodímio (NdFeB) & $\SI{-0.12}{\percent\per\celsius}$ & $\SI{80}{}$ a
 $\SI{200}{\celsius}$\\
Samário-cobalto & $\SI{-0.03}{\percent\per\celsius}$ & $\SI{300}{\celsius}$\\
\hline
\end{tabular}
\end{center}
\textbf{Compromisso revelador:} o neodímio tem a maior densidade de energia mas
péssima estabilidade térmica; o samário-cobalto é bem mais estável e mais caro.
Em motores que aquecem, a escolha raramente é o material mais forte.}
\end{questao}

\begin{questao}[2]{}
Um núcleo de aço-silício dissipa perdas por histerese proporcionais à área do
ciclo. Com área de $\SI{200}{\joule\per\meter\cubed}$ por ciclo e volume de
$\SI{1000}{\centi\meter\cubed}$, determine a potência a $\SI{60}{\hertz}$ e a
$\SI{400}{\hertz}$.
\resposta{$\SI{12}{\watt}$ e $\SI{80}{\watt}$}
\resolucao{
\[
P_h=f\times(\text{área})\times V:
\quad
60(200)(\pot{1}{-3})=\SI{12}{\watt};
\qquad
400(200)(\pot{1}{-3})=\SI{80}{\watt}.
\]
\textbf{Proporcionalidade linear com $f$} --- cada ciclo custa a mesma energia,
e mais ciclos por segundo significam mais potência.\\
\textbf{Contraste com as correntes parasitas}, cujas perdas vão com
$f^{2}$. Em baixa frequência a histerese domina; em alta, as parasitas.\\
\textbf{Consequência:} o mesmo núcleo que funciona bem em $\SI{60}{\hertz}$
pode ser inviável em $\SI{400}{\hertz}$ --- e é por isso que equipamentos
aeronáuticos, que operam nessa frequência, usam ligas especiais de lâminas
finíssimas.}
\end{questao}

\begin{questao}[2]{}
Um material magnético \textbf{mole} tem $H_c=\SI{50}{\ampere\per\meter}$ e um
material \textbf{duro} tem $H_c=\SI{500}{\kilo\ampere\per\meter}$. Compare-os
e indique aplicações.
\resposta{Razão de $10^{4}$ na coercividade; moles para núcleos, duros para ímãs
permanentes}
\resolucao{\textbf{Coercividade $H_c$} é o campo necessário para \emph{anular} a
magnetização remanente. Ela mede o quanto o material "resiste" a mudar de
estado.
\begin{center}
\begin{tabular}{lccl}
\hline
Tipo & $H_c$ & Área do ciclo & Aplicação\\
\hline
Mole & baixa & pequena & núcleos de transformador e motor\\
Duro & alta & grande & ímãs permanentes\\
\hline
\end{tabular}
\end{center}
\textbf{Materiais moles} (ferro-silício, permalloy, ferrites moles):
magnetizam e desmagnetizam facilmente, com pouca perda por ciclo. É o que se
quer em algo que inverte a magnetização $120$ vezes por segundo.\\
\textbf{Materiais duros} (alnico, ferrite dura, NdFeB): resistem à
desmagnetização, mantendo o campo sem corrente alguma.\\
\textbf{O critério é a área do ciclo:} pequena para quem vai percorrê-lo muitas
vezes, grande para quem precisa permanecer magnetizado.}
\end{questao}

\begin{questao}[2]{}
Explique por que a curva de magnetização \textbf{inicial} de um material
virgem difere do ciclo de histerese subsequente.
\resposta{A curva inicial parte da origem; o ciclo nunca mais passa por ela}
\resolucao{\textbf{Material virgem:} os domínios estão orientados
aleatoriamente, e a magnetização líquida é nula. Aplicando $H$ crescente, a
curva parte da \textbf{origem} e sobe até a saturação --- é a \emph{curva de
magnetização inicial}.\\
\textbf{Ao reduzir $H$ a zero}, a magnetização \emph{não} retorna a zero: resta
a \textbf{remanência} $B_r$. Os domínios que se reorientaram não voltam
espontaneamente.\\
\textbf{A partir daí}, ciclando $H$ entre $\pm H_{\max}$, o material percorre
um laço fechado que \textbf{nunca mais passa pela origem}.\\
\textbf{Como voltar ao estado virgem:} aplicar campo alternado de amplitude
decrescente até zero (desmagnetização por ciclos decrescentes), ou aquecer acima
de $T_C$ e resfriar sem campo.\\
\textbf{Importância prática:} ao levantar experimentalmente a curva
$B\times H$, é preciso desmagnetizar o material antes --- caso contrário, os
primeiros pontos refletem a história anterior, não a propriedade do material.}
\end{questao}

\begin{questao}[2]{}
Um transformador é desligado no pico da corrente de magnetização. Explique o
que ocorre na religação e por que a corrente de \emph{inrush} pode ser elevada.
\resposta{O fluxo remanente soma-se ao fluxo imposto, podendo levar o núcleo a
saturação profunda}
\resolucao{\textbf{Ao desligar}, o núcleo retém fluxo remanente
$\Phi_r$ --- consequência direta da histerese.\\
\textbf{Ao religar}, o fluxo evolui a partir desse valor. No pior caso, a tensão
é aplicada no instante em que exigiria fluxo de mesma polaridade, e o fluxo
total pode chegar a
\[
\Phi\approx\Phi_r+2\Phi_{\max},
\]
muito além da saturação.\\
\textbf{Consequência:} saturado, o núcleo perde permeabilidade, a indutância
desaba e a corrente é limitada apenas pela resistência do enrolamento. A
corrente de \emph{inrush} pode atingir dez a vinte vezes a nominal, por alguns
ciclos.\\
\textbf{Soluções empregadas:}
\begin{itemize}
\item \textbf{Chaveamento sincronizado:} energizar no instante correto da onda
de tensão.
\item \textbf{Resistor ou NTC de pré-carga}, curto-circuitado após alguns
ciclos.
\item \textbf{Proteção com curva temporizada}, que tolera o transitório sem
atuar.
\end{itemize}
\textbf{Note a conexão:} este fenômeno só existe \emph{porque} há histerese ---
um núcleo hipotético sem remanência não teria \emph{inrush} por essa causa.}
\end{questao}

\begin{questao}[2]{}
Um circuito magnético tem dois trechos em série: ferro
($\ell_1=\SI{30}{\centi\meter}$, $\mu_r=2000$) e ar
($\ell_2=\SI{2}{\milli\meter}$), ambos com seção
$\SI{5}{\centi\meter\squared}$. Determine a relutância total e a fração de cada
trecho.
\resposta{$\mathcal{R}_{fe}=\pot{2.39}{5}$, $\mathcal{R}_{ar}=\pot{3.18}{6}$;
o ar responde por $93\%$}
\resolucao{
\[
\mathcal{R}_{fe}=\frac{0{,}30}{(2000)\mu_0(\pot{5}{-4})}
=\pot{2.39}{5}\,\si{\per\henry};
\]
\[
\mathcal{R}_{ar}=\frac{\pot{2}{-3}}{\mu_0(\pot{5}{-4})}
=\pot{3.18}{6}\,\si{\per\henry}.
\]
\[
\mathcal{R}_{tot}=\pot{3.42}{6};
\qquad
\frac{\mathcal{R}_{ar}}{\mathcal{R}_{tot}}=93\%.
\]
\textbf{Dois milímetros de ar dominam trinta centímetros de ferro} --- porque a
razão de permeabilidades ($2000$) supera de longe a razão de comprimentos
($150$).\\
\textbf{As relutâncias somam-se em série}, exatamente como resistências. E o
fluxo é o mesmo nos dois trechos, o que garante $B$ igual --- já que a seção
também é a mesma.\\
\textbf{Consequência:} o comportamento do circuito é governado quase
inteiramente pelo entreferro, e a não linearidade do ferro fica mascarada. É
esse efeito que \emph{lineariza} indutores com entreferro.}
\end{questao}

\blocoaprofundamento

\begin{questao}[3]{}
Distinga rigorosamente $\vec H$, $\vec B$ e $\vec M$.
\begin{itensq}
\item Escreva a relação entre elas e as unidades.
\item Determine qual delas muda ao inserir um núcleo em um solenoide, mantida a
      corrente.
\item Explique qual é "mais fundamental" e por quê.
\end{itensq}
\resposta{$\vec B=\mu_0(\vec H+\vec M)$; ao inserir o núcleo, $H$ não muda e
$B$ aumenta}
\resolucao{(a)
\[
\vec B=\mu_0\left(\vec H+\vec M\right),
\]
com $[H]=[M]=\si{\ampere\per\meter}$ e $[B]=\si{\tesla}$.\\
\emph{($H$ é às vezes chamado "intensidade de campo magnético" e $B$
"densidade de fluxo" --- nomenclatura infeliz, porque sugere que $H$ é o campo
principal.)}\\
(b) \textbf{Ao inserir o núcleo com corrente constante:}
\begin{itemize}
\item $H=NI/\ell$ \textbf{não muda} --- ele depende só das correntes de
condução.
\item $M$ \textbf{passa de zero a $(\mu_r-1)H$} --- o material se magnetiza.
\item $B$ \textbf{aumenta por} $\mu_r$ --- consequência das duas anteriores.
\end{itemize}
(c) \textbf{Qual é mais fundamental: $\vec B$.} Argumentos:
\begin{itemize}
\item É $\vec B$ que aparece na força de Lorentz
($\vec F=q\vec v\times\vec B$) --- é ele que age sobre cargas.
\item É $\vec B$ que aparece na lei de Faraday --- é seu fluxo que induz f.e.m.
\item $\nabla\cdot\vec B=0$ é lei universal; $\nabla\cdot\vec H=-\nabla\cdot
\vec M$ não é nula em geral.
\end{itemize}
\textbf{Para que serve $\vec H$, então:} ele é conveniente porque sua circulação
depende \emph{apenas} das correntes de condução ($\oint\vec H\cdot
\mathrm{d}\vec\ell=NI$), sem precisar contabilizar as correntes de
magnetização do material. É uma grandeza \emph{auxiliar} de enorme utilidade
prática --- mas auxiliar.}
\end{questao}

\begin{questao}[3]{}
A curva $B\times H$ de um material apresenta três regiões distintas.
\begin{itensq}
\item Descreva cada uma em termos de domínios.
\item Identifique onde $\mu_r$ é máxima.
\item Indique a região de operação recomendada.
\end{itensq}
\resposta{Deslocamento reversível, deslocamento irreversível e rotação;
$\mu_r$ máxima no trecho intermediário}
\resolucao{(a) \textbf{Três regiões.}
\begin{enumerate}
\item \textbf{Campos baixos --- deslocamento reversível de paredes.} As paredes
entre domínios deslocam-se ligeiramente, favorecendo os domínios alinhados com
$H$. Removido o campo, elas retornam. A curva é pouco inclinada.
\item \textbf{Campos médios --- deslocamento irreversível (saltos de
Barkhausen).} As paredes superam obstáculos (impurezas, contornos de grão) em
saltos abruptos. É aqui que a curva é mais \textbf{íngreme} e onde nasce a
histerese, pois os saltos não se desfazem ao reduzir $H$.
\item \textbf{Campos altos --- rotação dos domínios.} Já não há paredes a
mover; os momentos giram progressivamente até se alinharem com $H$. A curva
achata-se: é o \emph{joelho} e depois a saturação.
\end{enumerate}
(b) \textbf{$\mu_r$ é máxima na região intermediária}, onde
$\mathrm{d}B/\mathrm{d}H$ é maior. Nem no início (paredes presas) nem no fim
(domínios já alinhados).\\
\textbf{Consequência contraintuitiva:} a permeabilidade em campos muito baixos é
\emph{menor} que a máxima --- daí catálogos especificarem $\mu_i$ (inicial) e
$\mu_{\max}$ separadamente.\\
(c) \textbf{Região recomendada:} abaixo do joelho, tipicamente com
$B_{\max}$ entre $60$ e $80\%$ de $B_{sat}$. Isso aproveita a alta
permeabilidade sem risco de saturação em transitórios.}
\end{questao}

\begin{questao}[3]{}
Um núcleo tem curva $B\times H$ não linear, e é preciso determinar o ponto de
operação com entreferro.
\begin{itensq}
\item Formule o problema.
\item Descreva o método gráfico da reta de entreferro.
\item Explique por que o entreferro estabiliza o ponto de operação.
\end{itensq}
\resposta{Interseção da curva $B\times H$ do material com a reta imposta pelo
entreferro}
\resolucao{(a) \textbf{Formulação.} A lei de Ampère no circuito magnético dá
\[
NI=H_{fe}\,\ell_{fe}+H_g\,g,
\]
e no entreferro $H_g=B/\mu_0$ (com $B$ igual ao do ferro, por conservação do
fluxo e seção constante). Então
\[
NI=H_{fe}\ell_{fe}+\frac{B}{\mu_0}g.
\]
Duas incógnitas ($B$ e $H_{fe}$) e uma equação --- a segunda relação é a
\emph{curva do material}, que não tem forma algébrica simples.\\
(b) \textbf{Método gráfico.} Reescrevendo,
\[
B=\frac{\mu_0\,NI}{g}-\frac{\mu_0\ell_{fe}}{g}\,H_{fe},
\]
que é uma \textbf{reta} no plano $B\times H$ --- a \emph{reta de entreferro} ou
\emph{reta de carga magnética}.\\
Traçando-a sobre a curva $B\times H$ do material, o \textbf{ponto de operação} é
a interseção.\\
\emph{(É exatamente o mesmo procedimento da reta de carga usada com diodos: uma
característica não linear do dispositivo e uma reta imposta pelo circuito.)}\\
(c) \textbf{Por que o entreferro estabiliza.} A inclinação da reta é
$-\mu_0\ell_{fe}/g$: quanto \emph{maior} o entreferro, mais \textbf{deitada} a
reta.\\
Uma reta deitada cruza a curva do material em uma região de inclinação
comparável --- e ali a interseção é \emph{pouco sensível} a variações de
$\mu_r$ do material. Trocar o lote de ferro, ou aquecê-lo, quase não desloca o
ponto.\\
\textbf{Sem entreferro}, a reta é quase vertical e a interseção depende
inteiramente da curva do material --- que varia com temperatura, lote e
história magnética. \textbf{O entreferro compra previsibilidade ao custo de mais
ampère-espiras.}}
\end{questao}

\begin{questao}[3]{}
Compare aço-silício e ferrite para aplicação em transformadores.
\begin{itensq}
\item Compare $B_{sat}$, resistividade e faixa de frequência.
\item Determine em qual faixa cada um é preferível.
\item Justifique fisicamente.
\end{itensq}
\resposta{Aço-silício: alto $B_{sat}$, baixa resistividade, até
$\si{\kilo\hertz}$; ferrite: baixo $B_{sat}$, altíssima resistividade,
até $\si{\mega\hertz}$}
\resolucao{(a)
\begin{center}
\begin{tabular}{lcc}
\hline
Propriedade & Aço-silício & Ferrite (MnZn)\\
\hline
$B_{sat}$ & $1{,}5$--$\SI{2}{\tesla}$ & $0{,}3$--$\SI{0.5}{\tesla}$\\
Resistividade & $\approx\pot{5}{-7}\,\si{\ohm\meter}$ &
 $10^{-1}$ a $\SI{10}{\ohm\meter}$\\
$\mu_r$ típica & $2000$--$8000$ & $1000$--$5000$\\
Faixa útil & até $\approx\SI{1}{\kilo\hertz}$ & $\SI{20}{\kilo\hertz}$ a
 $\si{\mega\hertz}$\\
\hline
\end{tabular}
\end{center}
\textbf{Note a resistividade:} a ferrite é um \emph{cerâmico}, com
resistividade até $10^{7}$ vezes maior que a do aço.\\
(b) \textbf{Baixa frequência} ($\SI{50}{}$--$\SI{400}{\hertz}$): aço-silício.
Seu $B_{sat}$ elevado permite núcleos menores para a mesma potência, e as perdas
parasitas são controladas por laminação.\\
\textbf{Alta frequência} ($>\SI{20}{\kilo\hertz}$): ferrite. Nenhuma laminação
seria fina o bastante para conter as correntes parasitas no aço.\\
(c) \textbf{Justificativa física.} As perdas por correntes parasitas vão com
\[
P_F\propto\frac{B^{2}f^{2}t^{2}}{\rho}.
\]
Em alta frequência, o termo $f^{2}$ explode. Há duas formas de compensar:
reduzir $t$ (laminação) ou aumentar $\rho$.\\
A ferrite ataca $\rho$ --- e ganha por muitas ordens de grandeza, o que nenhuma
laminação alcança.\\
\textbf{O preço:} $B_{sat}$ baixo obriga a usar seção maior ou aceitar menos
fluxo. \textbf{Mas em alta frequência isso não importa}, porque
$V=4{,}44Nf\Phi_{\max}$ --- com $f$ mil vezes maior, o fluxo necessário é mil
vezes menor. \textbf{É a alta frequência que torna a ferrite viável, e a ferrite
que torna a alta frequência viável.}}
\end{questao}

\begin{questao}[3]{}
Um ímã permanente é retirado de seu circuito magnético e deixado com os polos
livres ao ar.
\begin{itensq}
\item Explique o que ocorre com seu ponto de operação.
\item Determine o efeito de aproximar uma peça de ferro.
\item Justifique a prática de armazenar ímãs com "quebra-fluxo".
\end{itensq}
\resposta{Ao ar, o ponto de operação desce na curva de desmagnetização; o ferro
o eleva}
\resolucao{(a) \textbf{Ao ar livre.} O ímã precisa "empurrar" o fluxo por um
caminho de alta relutância --- o próprio ar em torno dele. Isso equivale a um
entreferro enorme, e a reta de carga fica muito deitada.\\
O ponto de operação desce na \emph{curva de desmagnetização} (segundo quadrante
do laço), afastando-se de $B_r$. O ímã fica submetido a campo desmagnetizante
produzido por ele mesmo.\\
(b) \textbf{Aproximando ferro.} O ferro oferece caminho de baixa relutância, e a
reta de carga se levanta. O ponto de operação sobe, aproximando-se de
$B_r$ --- o ímã fica "mais forte".\\
É por isso que um ímã em ferradura fechado por uma armadura de ferro sustenta
muito mais peso que aberto.\\
(c) \textbf{Quebra-fluxo (\emph{keeper}).} É uma peça de ferro doce colocada
entre os polos durante o armazenamento. Ela:
\begin{itemize}
\item mantém o ponto de operação alto, evitando desmagnetização progressiva;
\item protege contra campos externos adversos;
\item reduz a atração de limalha e detritos.
\end{itemize}
\textbf{Ressalva histórica:} a prática era essencial com alnico, cuja curva de
desmagnetização é frágil. Com neodímio moderno --- cuja curva é praticamente
reta no segundo quadrante --- o efeito é bem menor, embora o quebra-fluxo
continue útil por outras razões.}
\end{questao}

\begin{questao}[3]{}
Um material tem $B_r=\SI{1.2}{\tesla}$ e $H_c=\SI{900}{\kilo\ampere\per\meter}$.
\begin{itensq}
\item Estime o produto de energia máximo.
\item Compare com ferrite ($B_r=\SI{0.4}{\tesla}$,
      $H_c=\SI{250}{\kilo\ampere\per\meter}$).
\item Explique o significado do produto de energia.
\end{itensq}
\resposta{$(BH)_{\max}\approx\SI{270}{\kilo\joule\per\meter\cubed}$ contra
$\SI{25}{\kilo\joule\per\meter\cubed}$}
\resolucao{(a) Para curva de desmagnetização aproximadamente \emph{reta} (caso
do neodímio), o produto máximo ocorre no ponto médio:
\[
(BH)_{\max}\approx\frac{B_rH_c}{4}
=\frac{(1{,}2)(\pot{9}{5})}{4}
=\SI{270}{\kilo\joule\per\meter\cubed}.
\]
(b) Para a ferrite:
\[
(BH)_{\max}\approx\frac{(0{,}4)(\pot{2.5}{5})}{4}
=\SI{25}{\kilo\joule\per\meter\cubed},
\]
cerca de \textbf{onze vezes menor}.\\
(c) \textbf{Significado.} O produto de energia mede a energia magnética
disponível \emph{por unidade de volume} do ímã. Ele é a figura de mérito que
determina quão pequeno o ímã pode ser para produzir dado campo em dado
entreferro.\\
\textbf{Consequência prática:} para o mesmo desempenho, um ímã de neodímio ocupa
cerca de um décimo do volume de um de ferrite. Isso é o que viabilizou
servomotores compactos, fones de ouvido pequenos e discos rígidos.\\
\textbf{Contrapartida:} o neodímio é caro, oxida (precisa de revestimento),
desmagnetiza acima de $\SI{80}{}$--$\SI{200}{\celsius}$ e depende de
terras-raras. A ferrite é barata, quimicamente estável e térmicamente mais
tolerante --- e por isso continua dominando em volume de produção.}
\end{questao}

\begin{questao}[3]{}
Explique a diferença entre desmagnetização por \textbf{aquecimento} e por
\textbf{campo alternado decrescente}.
\begin{itensq}
\item Descreva cada processo.
\item Compare eficácia e aplicabilidade.
\end{itensq}
\resposta{Aquecimento acima de $T_C$ desordena os domínios; campo alternado os
distribui simetricamente}
\resolucao{(a) \textbf{Aquecimento acima de $T_C$.} A agitação térmica destrói o
alinhamento dentro dos domínios. Resfriando \emph{sem campo aplicado}, os
domínios se reformam em orientações aleatórias, e a magnetização líquida é
nula.\\
\textbf{Campo alternado decrescente.} Aplica-se campo alternado de amplitude
inicialmente alta (o bastante para saturar) e reduz-se lentamente a zero. O
material percorre laços de histerese cada vez menores, espiralando até a origem.
Ao final, tantos domínios apontam em um sentido quanto no outro.\\
(b) \textbf{Comparação.}
\begin{center}
\begin{tabular}{lcc}
\hline
& Aquecimento & Campo alternado\\
\hline
Eficácia & completa & muito boa\\
Danifica a peça? & possivelmente & não\\
Aplicável a peça montada? & raramente & sim\\
Rapidez & lento & segundos\\
\hline
\end{tabular}
\end{center}
\textbf{Na prática industrial, usa-se campo alternado} --- em desmagnetizadores
de ferramentas, de peças usinadas e de fitas magnéticas. O aquecimento é
inviável para peças que não toleram $\SI{770}{\celsius}$.\\
\textbf{Cuidado essencial:} a redução do campo deve ser \emph{lenta} em relação
ao período. Desligar bruscamente o desmagnetizador pode deixar a peça
magnetizada no valor correspondente à última meia-onda --- pior que antes.\\
\textbf{Alternativa equivalente:} afastar lentamente a peça de um campo
alternado fixo, o que produz o mesmo decaimento gradual.}
\end{questao}

\begin{questao}[3]{}
Um núcleo opera com $B_{\max}=\SI{1.2}{\tesla}$ a $\SI{60}{\hertz}$. Analise o
efeito de elevar $B_{\max}$ para $\SI{1.5}{\tesla}$.
\begin{itensq}
\item Determine o efeito sobre o número de espiras.
\item Determine o efeito sobre as perdas.
\item Avalie o compromisso.
\end{itensq}
\resposta{$N$ cai $20\%$; perdas por histerese sobem cerca de $60\%$ e
parasitas $56\%$}
\resolucao{(a) De $V=4{,}44Nf\Phi_{\max}=4{,}44NfB_{\max}A$, com $V$, $f$ e
$A$ fixos:
\[
N\propto\frac{1}{B_{\max}}
\;\Longrightarrow\;
\frac{N'}{N}=\frac{1{,}2}{1{,}5}=0{,}80.
\]
\textbf{Vinte por cento menos espiras} --- menos cobre, menos resistência, menos
perdas no enrolamento.\\
(b) \textbf{Perdas no núcleo aumentam.}
\begin{itemize}
\item \textbf{Parasitas:} $P_F\propto B_{\max}^{2}$, logo
$(1{,}5/1{,}2)^{2}=1{,}56$ --- aumento de $56\%$.
\item \textbf{Histerese:} $P_h\propto B_{\max}^{n}$ com $n\approx1{,}6$ a
$2$ (expoente de Steinmetz), dando aumento de $60$ a $56\%$.
\end{itemize}
(c) \textbf{O compromisso.}
\begin{center}
\begin{tabular}{lc}
\hline
Elevar $B_{\max}$ & Efeito\\
\hline
Espiras & $-20\%$\\
Perdas no cobre & diminuem\\
Perdas no núcleo & $+\approx58\%$\\
Volume e custo & diminuem\\
Margem para saturação & \textbf{diminui}\\
\hline
\end{tabular}
\end{center}
\textbf{Existe um ótimo econômico}, em que a soma das perdas é mínima --- e ele
depende do custo relativo do cobre e do aço e do regime de carga.\\
\textbf{Mas há um limite duro:} $\SI{1.5}{\tesla}$ já está próximo do joelho. A
margem para sobretensões e para o transitório de energização encolhe, e o risco
de \emph{inrush} severo cresce. \textbf{Projetar perto da saturação economiza
material e compra problemas operacionais.}}
\end{questao}

\begin{questao}[3]{}
Explique por que a permeabilidade de um material ferromagnético não é uma
constante, e quais parâmetros a afetam.
\resposta{$\mu$ depende de $H$, da temperatura, da frequência e da história
magnética}
\resolucao{\textbf{Quatro dependências, todas relevantes.}
\begin{enumerate}
\item \textbf{Da intensidade $H$.} É a mais óbvia: $\mu_r$ parte de um valor
inicial, cresce até um máximo e despenca na saturação. Variação de uma ordem de
grandeza ou mais.
\item \textbf{Da temperatura.} $\mu_r$ tende a \emph{crescer} ao se aproximar de
$T_C$ (efeito Hopkinson) e depois colapsa abruptamente. Isso torna o
comportamento térmico não monotônico e difícil de prever.
\item \textbf{Da frequência.} Em alta frequência, as paredes de domínio não
conseguem acompanhar o campo --- há \emph{relaxação}. A permeabilidade efetiva
cai, e aparece uma componente de perda. Catálogos de ferrite trazem
$\mu$ complexo, com parte real e imaginária.
\item \textbf{Da história magnética.} Por causa da histerese, o mesmo $H$ pode
corresponder a diferentes $B$, conforme o material esteja subindo ou descendo o
laço.
\end{enumerate}
\textbf{Consequência para o projetista.} Um valor único de $\mu_r$ em catálogo é
sempre uma \emph{simplificação} referida a condições específicas --- geralmente
campo baixo, temperatura ambiente e baixa frequência.\\
\textbf{Prática correta:} usar $\mu_r$ para estimativas iniciais, e recorrer às
\emph{curvas} do fabricante (de $B\times H$, de perdas por frequência, de
$\mu$ por temperatura) para o dimensionamento final. Projetos de fontes
chaveadas que ignoram isso costumam funcionar na bancada e falhar em campo.}
\end{questao}

\begin{questao}[3]{}
Um sensor de proximidade indutivo detecta metais pela variação de indutância.
\begin{itensq}
\item Explique o princípio para metais ferromagnéticos.
\item Explique para metais não ferromagnéticos (alumínio, cobre).
\item Determine por que o alcance difere entre eles.
\end{itensq}
\resposta{Ferromagnéticos aumentam $L$ por permeabilidade; não ferromagnéticos
reduzem $L$ e o fator de qualidade por correntes parasitas}
\resolucao{(a) \textbf{Ferromagnéticos.} A aproximação do material oferece
caminho de menor relutância, aumentando a indutância da bobina sensora. A
variação de $L$ é detectada pelo circuito.\\
Simultaneamente, correntes parasitas no aço amortecem o oscilador.\\
(b) \textbf{Não ferromagnéticos.} Alumínio e cobre têm $\mu_r\approx1$ --- não
alteram a relutância. Mas são \emph{bons condutores}, e o campo alternado induz
neles correntes parasitas intensas.\\
Pela lei de Lenz, essas correntes criam campo oposto, o que \textbf{reduz} a
indutância efetiva e, sobretudo, \textbf{drena energia} do oscilador, reduzindo
seu fator de qualidade.\\
(c) \textbf{Por que o alcance difere.} Sensores indutivos comuns são
especificados para aço, e apresentam \emph{fatores de correção}:
\begin{center}
\begin{tabular}{lc}
\hline
Material & Fator de alcance\\
\hline
Aço & $1{,}00$\\
Aço inoxidável & $\approx0{,}70$\\
Latão & $\approx0{,}50$\\
Alumínio & $\approx0{,}40$\\
Cobre & $\approx0{,}30$\\
\hline
\end{tabular}
\end{center}
\textbf{Razão:} no aço, os dois efeitos (permeabilidade e parasitas) somam-se.
Nos não ferromagnéticos, resta apenas o segundo --- e ele é tanto mais fraco
quanto \emph{melhor} o condutor, porque as correntes se concentram numa
película superficial cada vez mais fina.\\
\textbf{Consequência de aplicação:} um sensor ajustado para detectar peças de
aço a $\SI{10}{\milli\meter}$ só detectará cobre a $\SI{3}{\milli\meter}$ ---
erro de projeto frequente em linhas de montagem com peças mistas.}
\end{questao}

\blocodesafio

\begin{questao}[4]{}
\textbf{(Energia do ciclo de histerese.)} Demonstre que a área do laço
$B\times H$ representa energia dissipada por unidade de volume e por ciclo.
\begin{itensq}
\item Deduza a expressão.
\item Determine a potência dissipada em função da frequência.
\item Explique a origem microscópica da perda.
\end{itensq}
\resposta{$w=\oint H\,\mathrm{d}B$ por ciclo; $P=fVw$}
\resolucao{(a) \textbf{Dedução.} A potência instantânea fornecida ao núcleo pela
fonte é $p=\varepsilon i$. Com $\varepsilon=NA\,\mathrm{d}B/\mathrm{d}t$ e
$i=H\ell/N$:
\[
p=NA\frac{\mathrm{d}B}{\mathrm{d}t}\cdot\frac{H\ell}{N}
=(A\ell)\,H\,\frac{\mathrm{d}B}{\mathrm{d}t}
=V\,H\,\frac{\mathrm{d}B}{\mathrm{d}t}.
\]
A energia por ciclo é
\[
W=\int p\,\mathrm{d}t=V\oint H\,\mathrm{d}B,
\]
e $\oint H\,\mathrm{d}B$ é precisamente a \textbf{área do laço} no plano
$B\times H$, com unidade $\si{\joule\per\meter\cubed}$.\\
\textbf{Se o material fosse sem histerese}, o laço colapsaria em uma curva
única, a integral de ida cancelaria a de volta e a energia por ciclo seria
\textbf{zero} --- toda a energia magnética seria devolvida.\\
(b)
\[
P_h=f\,V\oint H\,\mathrm{d}B,
\]
\textbf{linear em $f$}. (As perdas por correntes parasitas, distintas destas,
vão com $f^{2}$.)\\
(c) \textbf{Origem microscópica.} A perda vem dos \textbf{saltos de Barkhausen}:
as paredes de domínio ficam presas em impurezas, defeitos e contornos de grão, e
se libertam abruptamente quando o campo cresce o suficiente.\\
Cada salto é um processo \textbf{irreversível} --- a energia elástica acumulada
converte-se em vibração da rede, ou seja, em \textbf{calor}.\\
\textbf{Consequência de engenharia:} materiais mais puros e com grãos maiores e
melhor orientados têm menos obstáculos, laços mais estreitos e menores perdas.
É exatamente isso que o processo de \emph{grão orientado} do aço-silício
persegue.}
\end{questao}

\begin{questao}[4]{}
\textbf{(Circuito magnético com dois materiais.)} Um circuito tem trecho de
ferro ($\ell_1$, $\mu_1$) em série com trecho de ferrite ($\ell_2$,
$\mu_2$), mesma seção.
\begin{itensq}
\item Formule as equações.
\item Determine $B$ e $H$ em cada trecho.
\item Explique qual grandeza é contínua na interface e por quê.
\end{itensq}
\resposta{$B$ é contínuo; $H$ é descontínuo, na razão inversa das
permeabilidades}
\resolucao{(a) \textbf{Duas equações.}\\
\emph{Conservação do fluxo} (equivalente à LKC):
\[
\Phi_1=\Phi_2
\;\Longrightarrow\;
B_1A=B_2A
\;\Longrightarrow\;
B_1=B_2=B.
\]
\emph{Lei de Ampère} (equivalente à LKT):
\[
NI=H_1\ell_1+H_2\ell_2.
\]
(b) Com $H_i=B/\mu_i$:
\[
NI=B\left(\frac{\ell_1}{\mu_1}+\frac{\ell_2}{\mu_2}\right)
\;\Longrightarrow\;
B=\frac{NI}{\dfrac{\ell_1}{\mu_1}+\dfrac{\ell_2}{\mu_2}}.
\]
\[
H_1=\frac{B}{\mu_1};
\qquad
H_2=\frac{B}{\mu_2};
\qquad
\frac{H_1}{H_2}=\frac{\mu_2}{\mu_1}.
\]
(c) \textbf{Qual grandeza é contínua --- e por quê.}
\begin{itemize}
\item \textbf{$B$ é contínuo} através da interface (na componente normal),
porque $\nabla\cdot\vec B=0$: o fluxo que chega precisa continuar. Não há
"acúmulo de fluxo" possível.
\item \textbf{$H$ é descontínuo}, na razão inversa das permeabilidades. O
material de menor $\mu$ exige \emph{muito mais} $H$ para sustentar o mesmo
$B$.
\end{itemize}
\textbf{Exemplo numérico revelador.} Com $\mu_{r1}=2000$ (ferro) e
$\mu_{r2}=2$ (um material quase não magnético), $H_2$ seria \emph{mil vezes}
$H_1$ --- e praticamente toda a força magnetomotriz seria consumida naquele
trecho, ainda que ele fosse curtíssimo.\\
\textbf{É exatamente o que acontece no entreferro} ($\mu_r=1$), e por isso ele
domina o circuito. Um entreferro não é um caso especial: é apenas o caso extremo
deste problema de dois materiais.}
\end{questao}

\begin{questao}[4]{}
\textbf{(Ensaio do ciclo de histerese.)} Projete um experimento para levantar o
laço $B\times H$ de um material em osciloscópio.
\begin{itensq}
\item Descreva o circuito e o que cada canal mede.
\item Determine as constantes de conversão.
\item Aponte os cuidados experimentais.
\end{itensq}
\resposta{Canal horizontal $\propto H$ (corrente por shunt); vertical
$\propto B$ (f.e.m. integrada)}
\resolucao{(a) \textbf{Circuito.} Um núcleo toroidal com dois enrolamentos:
primário $N_1$ excitado por fonte alternada e secundário $N_2$ para medição.
\begin{itemize}
\item \textbf{Canal horizontal ($X$):} tensão sobre um resistor \emph{shunt}
$R_s$ em série com o primário. Como $H=N_1i_1/\ell$, essa tensão é proporcional
a $H$.
\item \textbf{Canal vertical ($Y$):} a f.e.m. do secundário passa por um
\textbf{integrador} $RC$. Como
$\varepsilon_2=N_2A\,\mathrm{d}B/\mathrm{d}t$, sua integral é proporcional a
$B$.
\end{itemize}
Em modo $XY$, o osciloscópio traça diretamente o laço.\\
(b) \textbf{Constantes de conversão.}
\[
H=\frac{N_1}{\ell R_s}\,v_X;
\qquad
B=\frac{R_iC_i}{N_2A}\,v_Y,
\]
com $R_iC_i$ os componentes do integrador.\\
\textbf{Condição do integrador:} $R_iC_i\gg1/f$, para que ele integre de fato em
vez de apenas atenuar. Com $f=\SI{60}{\hertz}$, adote
$R_iC_i\ge\SI{0.1}{\second}$.\\
(c) \textbf{Cuidados experimentais.}
\begin{itemize}
\item \textbf{Desmagnetizar antes} --- caso contrário os primeiros ciclos
refletem a história anterior. Comece com amplitude alta e reduza.
\item \textbf{$R_s$ pequeno} (fração de ohm), para não distorcer a corrente nem
dissipar demais. Use quatro fios se possível.
\item \textbf{Núcleo toroidal}, que garante $H$ uniforme e caminho médio bem
definido --- em núcleos em E o caminho é ambíguo.
\item \textbf{Deriva do integrador:} qualquer \emph{offset} faz o laço "subir"
lentamente na tela. Use acoplamento CA ou integrador com realimentação CC.
\item \textbf{Varrer a amplitude} para obter a família de laços e localizar a
curva de magnetização normal, que é o lugar geométrico dos vértices.
\end{itemize}
\textbf{O que se extrai:} $B_r$ (interseção com o eixo vertical), $H_c$
(interseção com o horizontal), $B_{sat}$ e a área --- e daí as perdas por
ciclo.}
\end{questao}

\begin{questao}[4]{}
\textbf{(Domínios magnéticos.)} Explique a estrutura de domínios e sua relação
com as propriedades macroscópicas.
\begin{itensq}
\item Explique por que os domínios existem.
\item Relacione a estrutura com $B_r$, $H_c$ e $\mu_r$.
\item Explique a diferença entre materiais moles e duros nesse nível.
\end{itensq}
\resposta{Domínios minimizam a energia magnetostática; a mobilidade das paredes
determina $\mu_r$ e $H_c$}
\resolucao{(a) \textbf{Por que existem.} Um cristal ferromagnético
uniformemente magnetizado criaria campo externo intenso, com alta energia
magnetostática armazenada no espaço.\\
Dividir-se em \textbf{domínios} de orientações opostas faz esse campo externo
quase desaparecer --- as linhas se fecham internamente. O custo é a energia das
\emph{paredes de domínio} (paredes de Bloch), onde os momentos giram
gradualmente.\\
\textbf{O equilíbrio} entre esses dois termos fixa o tamanho típico dos
domínios, tipicamente de micrômetros a dezenas de micrômetros.\\
(b) \textbf{Relação com as propriedades macroscópicas.}
\begin{itemize}
\item \textbf{$\mu_r$ alta} exige paredes \emph{móveis} --- material puro, sem
tensões internas, com grãos grandes e bem orientados.
\item \textbf{$H_c$ baixa} tem a mesma origem: paredes que deslizam facilmente
não resistem à inversão.
\item \textbf{$B_r$ alta} exige que os domínios permaneçam alinhados após a
retirada do campo --- o que requer \emph{anisotropia} cristalina forte,
"ancorando" a magnetização em direções preferenciais.
\end{itemize}
(c) \textbf{Moles contra duros, no nível microscópico.}
\begin{center}
\begin{tabular}{lll}
\hline
& Moles & Duros\\
\hline
Paredes & muito móveis & fixadas ou inexistentes\\
Impurezas & minimizadas & introduzidas de propósito\\
Anisotropia & baixa & alta\\
Grãos & grandes, orientados & finos, monodomínio\\
\hline
\end{tabular}
\end{center}
\textbf{O contraste é elegante:} ímãs modernos de neodímio são feitos de
partículas tão pequenas que cada uma é um \textbf{monodomínio} --- não há
paredes para mover, e inverter a magnetização exige girar todos os momentos
contra a anisotropia. Daí a coercividade enorme.\\
\textbf{Já o aço-silício de grão orientado} persegue o oposto: grãos grandes,
alinhados por laminação a frio, com paredes livres para deslizar.\\
\textbf{Mesma física, objetivos opostos --- e processos de fabricação
inteiramente diferentes.}}
\end{questao}

\begin{questao}[4]{}
\textbf{(Projeto sob restrição de perdas.)} Um transformador de
$\SI{60}{\hertz}$ deve ter perdas no núcleo inferiores a $\SI{2}{\watt}$, com
material cuja perda específica é
$\SI{1.5}{\watt\per\kilogram}$ a $\SI{1.5}{\tesla}$ e
$\SI{0.9}{\watt\per\kilogram}$ a $\SI{1.2}{\tesla}$.
\begin{itensq}
\item Determine a massa admissível em cada condição.
\item Determine o efeito sobre o número de espiras.
\item Estabeleça o critério de escolha.
\end{itensq}
\resposta{$\SI{1.33}{\kilogram}$ a $\SI{1.5}{\tesla}$;
$\SI{2.22}{\kilogram}$ a $\SI{1.2}{\tesla}$}
\resolucao{(a)
\[
m\le\frac{2}{1{,}5}=\SI{1.33}{\kilogram}\ \text{a } \SI{1.5}{\tesla};
\qquad
m\le\frac{2}{0{,}9}=\SI{2.22}{\kilogram}\ \text{a } \SI{1.2}{\tesla}.
\]
(b) De $V=4{,}44NfB_{\max}A$, com tensão e frequência fixas,
\[
N\,B_{\max}A=\text{const}.
\]
Reduzir $B_{\max}$ de $1{,}5$ para $\SI{1.2}{\tesla}$ exige aumentar
$NA$ em $25\%$ --- mais espiras, ou seção maior, ou ambos.\\
(c) \textbf{Critério de escolha.} As duas opções levam a projetos diferentes:
\begin{center}
\begin{tabular}{lcc}
\hline
& $B=\SI{1.5}{\tesla}$ & $B=\SI{1.2}{\tesla}$\\
\hline
Massa de núcleo & menor & maior ($+67\%$)\\
Espiras & menos & mais ($+25\%$)\\
Perdas no cobre & menores & maiores\\
Margem para saturação & pequena & confortável\\
Corrente de \emph{inrush} & severa & moderada\\
\hline
\end{tabular}
\end{center}
\textbf{Note o conflito:} operar em $B$ alto economiza \emph{ferro} mas gasta
mais \emph{cobre} por unidade de perda; operar em $B$ baixo faz o inverso.\\
\textbf{Decisão:}
\begin{itemize}
\item \textbf{Se o custo de material dominar} (transformadores de baixo custo,
uso intermitente), escolha $B$ alto.
\item \textbf{Se o custo de energia dominar} (operação contínua por décadas),
escolha $B$ baixo --- as perdas em vazio existem $\SI{8760}{\hour}$ por ano.
\item \textbf{Se houver risco de sobretensão}, escolha $B$ baixo por segurança:
um transformador operando a $\SI{1.5}{\tesla}$ satura com apenas $10\%$ de
sobretensão.
\end{itemize}}
\end{questao}

\begin{questao}[4]{}
\textbf{(Anisotropia e grão orientado.)} Explique o que é o aço-silício de grão
orientado e por que ele é usado em transformadores mas não em motores.
\begin{itensq}
\item Explique o processo e o efeito.
\item Compare as perdas nas duas direções.
\item Justifique a diferença de aplicação.
\end{itensq}
\resposta{Grãos alinhados numa direção preferencial; excelente ao longo dela,
ruim transversalmente}
\resolucao{(a) \textbf{Processo e efeito.} O aço-silício é laminado a frio e
recozido de modo que os grãos cristalinos se alinhem com o eixo de fácil
magnetização do ferro (a direção $\langle100\rangle$ da célula cúbica) paralelo
à direção de laminação.\\
\textbf{Resultado:} nessa direção, magnetizar é muito mais fácil --- as paredes
de domínio se movem livremente, $\mu_r$ é altíssima e as perdas caem
drasticamente.\\
(b) \textbf{Comparação por direção.}
\begin{center}
\begin{tabular}{lcc}
\hline
& Direção de laminação & Transversal\\
\hline
Perdas específicas & $\approx\SI{1}{\watt\per\kilogram}$ &
 $2$ a $3\times$ maiores\\
$\mu_r$ & muito alta & bem menor\\
$B$ para dado $H$ & alto & reduzido\\
\hline
\end{tabular}
\end{center}
\textbf{A anisotropia é grande} --- o material é excelente em uma direção e
medíocre nas outras.\\
(c) \textbf{Por que serve a transformadores e não a motores.}
\begin{itemize}
\item \textbf{Transformador:} o fluxo percorre um caminho \emph{fixo e
conhecido} --- as colunas e as culatras do núcleo. É possível cortar as lâminas
de modo que o fluxo siga sempre a direção favorável. \emph{(Nos cantos, onde o
fluxo precisa dobrar, usam-se juntas em $\ang{45}$ para minimizar a
transição.)}
\item \textbf{Motor:} o fluxo no estator \textbf{gira}, percorrendo todas as
direções ao longo de cada revolução. Não existe direção preferencial a
privilegiar --- metade do tempo o fluxo estaria na direção ruim.
\end{itemize}
\textbf{Por isso motores usam aço de grão \emph{não} orientado}, deliberadamente
\textbf{isotrópico}: pior que o grão orientado em sua melhor direção, mas
uniforme em todas --- e portanto melhor \emph{na média} sobre uma rotação
completa.\\
\textbf{Princípio geral:} otimizar para o caso médio quando a solicitação é
variável, e para o caso específico quando ela é fixa.}
\end{questao}

\begin{questao}[4]{}
\textbf{(Blindagem magnética com material de alta permeabilidade.)} Analise
quantitativamente a blindagem de campo estático por invólucro de mu-metal.
\begin{itensq}
\item Explique o mecanismo por relutância.
\item Estime o fator de blindagem para invólucro cilíndrico.
\item Aponte as limitações práticas.
\end{itensq}
\resposta{Fator $\approx\mu_r t/D$; o mecanismo é desvio de fluxo, não
cancelamento}
\resolucao{(a) \textbf{Mecanismo.} Ao contrário da blindagem elétrica --- que
\emph{cancela} o campo por redistribuição de cargas --- a blindagem magnética
\textbf{desvia} o fluxo.\\
O invólucro de alta $\mu_r$ oferece caminho de relutância muito menor que o ar
interno. O fluxo "prefere" circular pela parede e contorna a região protegida.\\
\textbf{Analogia:} é um divisor de fluxo, em que o material é o ramo de baixa
relutância e a cavidade, o de alta.\\
(b) \textbf{Fator de blindagem.} Para um cilindro de diâmetro $D$ e parede
$t$, com $t\ll D$, a estimativa clássica é
\[
S\approx\frac{\mu_r\,t}{D}.
\]
Com $\mu_r=20000$ (mu-metal recozido), $t=\SI{1}{\milli\meter}$ e
$D=\SI{100}{\milli\meter}$:
\[
S\approx\frac{20000(1)}{100}=200,
\]
ou seja, atenuação de cerca de $\SI{46}{\decibel}$.\\
\textbf{Compare com a blindagem elétrica}, que atinge muitas ordens de grandeza
com uma folha fina de qualquer metal. A magnética é \emph{muito} mais modesta.\\
(c) \textbf{Limitações práticas.}
\begin{itemize}
\item \textbf{Saturação.} Se o campo externo for intenso, a parede satura,
$\mu_r$ desaba e a blindagem falha justamente quando mais se precisa dela.
\emph{Solução:} camadas múltiplas, com a externa de material de alto
$B_{sat}$ (aço comum) e a interna de mu-metal.
\item \textbf{Sensibilidade mecânica.} O mu-metal perde permeabilidade se for
dobrado, cortado ou submetido a choque. Exige \textbf{recozimento após a
conformação}, em atmosfera de hidrogênio --- processo caro.
\item \textbf{Aberturas.} Furos e fendas quebram o caminho de baixa relutância e
degradam a blindagem muito mais que no caso elétrico.
\item \textbf{Custo e peso}, ambos elevados.
\end{itemize}
\textbf{Conclusão:} blindar campo magnético estático é caro, frágil e de eficácia
limitada. Sempre que possível, é preferível \emph{afastar} a fonte, reduzir a
área das espiras ou compensar ativamente o campo com bobinas.}
\end{questao}

\begin{questao}[4]{}
\textbf{(Ponto de operação de ímã permanente.)} Determine o ponto de operação
de um ímã em um circuito magnético com entreferro.
\begin{itensq}
\item Formule a equação da reta de carga.
\item Determine a condição de máximo aproveitamento.
\item Explique o efeito de aumentar o entreferro.
\end{itensq}
\resposta{A reta de carga tem inclinação
$-\dfrac{\ell_m A_g}{g A_m}$; o aproveitamento é máximo em
$(BH)_{\max}$}
\resolucao{(a) \textbf{Reta de carga.} Duas condições:
\emph{conservação do fluxo} ($B_mA_m=B_gA_g$) e \emph{lei de Ampère}
($H_m\ell_m+H_gg=0$, pois não há corrente).\\
Da segunda, $H_m\ell_m=-H_gg=-\dfrac{B_g}{\mu_0}g$. Combinando com a primeira:
\[
B_m=-\mu_0\,\frac{\ell_m A_g}{g\,A_m}\,H_m.
\]
É uma reta pela \textbf{origem} no segundo quadrante ($H_m<0$, $B_m>0$) ---
a chamada \emph{linha de carga} do ímã.\\
\textbf{Note:} o ímã opera com $H$ \emph{negativo} --- ele está sujeito ao
próprio campo desmagnetizante. Sem entreferro ($g\to0$), a reta seria vertical e
o ponto de operação estaria em $B_r$.\\
(b) \textbf{Máximo aproveitamento.} A energia disponível no entreferro é
proporcional ao produto $B_mH_m$. O ponto ótimo é aquele em que
$|B_mH_m|$ é \textbf{máximo} --- o já mencionado $(BH)_{\max}$.\\
Dimensiona-se o circuito \emph{para que a reta de carga passe por esse ponto},
escolhendo a razão $\dfrac{\ell_m A_g}{g A_m}$ adequada. Isso minimiza o volume
de ímã necessário.\\
(c) \textbf{Efeito de aumentar o entreferro.} A inclinação
$-\mu_0\ell_mA_g/(gA_m)$ diminui em módulo --- a reta \textbf{deita}. O ponto de
operação desce na curva de desmagnetização:
\begin{itemize}
\item $B_m$ diminui (menos fluxo);
\item $|H_m|$ aumenta (mais campo desmagnetizante sobre o próprio ímã);
\item se o entreferro for grande demais, o ponto pode ultrapassar o
\textbf{joelho} da curva de desmagnetização --- e aí ocorre perda
\textbf{irreversível} de magnetização.
\end{itemize}
\textbf{Consequência prática decisiva:} desmontar um circuito magnético com ímã
de alnico pode danificá-lo permanentemente, pois o ponto de operação despenca.
Ímãs de neodímio, cuja curva de desmagnetização é praticamente reta até o
segundo quadrante inteiro, toleram isso muito melhor --- e é uma das razões de
sua popularidade, além do produto de energia.}
\end{questao}

% END BANCO COMPLEMENTAR

\gabarito[12.1 Conceitos iniciais e materiais magnéticos]

% =====================================================================
\subsection{Campo magnético}
% =====================================================================

\blocofixacao

\begin{questao}[1]{}
O campo magnético a uma distância $r$ de um fio retilíneo longo percorrido por
corrente $I$:
\begin{alternativas}
\item é proporcional a $r$.
\item é inversamente proporcional a $r$.
\item é inversamente proporcional a $r^2$.
\item independe de $r$.
\item é proporcional a $r^2$.
\end{alternativas}
\resposta{(b)}
\resolucao{Para um fio reto longo, $B = \dfrac{\mu_0 I}{2\pi r}$: o campo é
\emph{inversamente} proporcional à distância $r$ (não ao quadrado, como no campo
elétrico de uma carga puntiforme).}
\end{questao}

\begin{questao}[1]{}
Um fio reto longo conduz $I = \SI{10}{\ampere}$. Determine o módulo do campo
magnético a $\SI{5}{\centi\meter}$ do fio.
\dados{$\mu_0 = 4\pi\times10^{-7}\,\si{\tesla\meter\per\ampere}$.}
\resposta{$B = \pot{4}{-5}\,\si{\tesla} = \SI{40}{\micro\tesla}$}
\resolucao{$B = \dfrac{\mu_0 I}{2\pi r}
= \dfrac{(\pot{4\pi}{-7})(10)}{2\pi(0{,}05)}
= \dfrac{\pot{2}{-6}}{0{,}05} = \pot{4}{-5}\,\si{\tesla}$.\\
Um atalho útil: $\dfrac{\mu_0}{2\pi} = \pot{2}{-7}$, então
$B = \pot{2}{-7}\cdot\dfrac{I}{r}$.}
\end{questao}

\begin{questao}[1]{}
Um fio retilíneo é percorrido por corrente para cima (saindo do chão em direção ao
teto). Usando a regra da mão direita, as linhas de campo magnético ao redor do fio:
\begin{alternativas}
\item apontam radialmente para fora do fio.
\item apontam radialmente para dentro do fio.
\item formam círculos concêntricos em torno do fio.
\item são paralelas ao fio.
\item não existem (fio reto não gera campo).
\end{alternativas}
\resposta{(c)}
\resolucao{O campo de um fio reto forma \emph{círculos concêntricos} em planos
perpendiculares ao fio. Pela regra da mão direita (polegar no sentido da corrente,
dedos envolvendo o fio), determina-se o sentido dessas linhas. O campo nunca é
radial (isso é característico do campo \emph{elétrico} de uma carga).}
\end{questao}

\blocoaplicacao

\begin{questao}[2]{}
Um solenoide de $\SI{1000}{}$ espiras por metro é percorrido por
$I = \SI{2}{\ampere}$. Determine o campo magnético no seu interior.
\dados{$\mu_0 = 4\pi\times10^{-7}\,\si{\tesla\meter\per\ampere}$.}
\resposta{$B \approx \SI{2.51}{\milli\tesla}$}
\resolucao{$B = \mu_0\,n\,I = (\pot{4\pi}{-7})(1000)(2)
= \pot{8\pi}{-4} \approx\SI{2.51}{\milli\tesla}$. O campo no interior de um
solenoide é \emph{uniforme} e independe do raio.}
\end{questao}

\begin{questao}[2]{}
Uma espira circular de raio $R = \SI{10}{\centi\meter}$ conduz
$I = \SI{5}{\ampere}$. Determine o campo magnético no seu centro.
\dados{$\mu_0 = 4\pi\times10^{-7}\,\si{\tesla\meter\per\ampere}$.}
\resposta{$B \approx \SI{31.4}{\micro\tesla}$}
\resolucao{$B = \dfrac{\mu_0 I}{2R}
= \dfrac{(\pot{4\pi}{-7})(5)}{2(0{,}1)}
= \dfrac{\pot{2\pi}{-6}}{0{,}2}=\pot{\pi}{-5}
\approx\SI{31.4}{\micro\tesla}$.}
\end{questao}

\begin{questao}[2]{}
Um fio reto conduz $I = \SI{4}{\ampere}$. Determine o campo magnético a
$\SI{2}{\centi\meter}$ do fio.
\dados{$\mu_0 = 4\pi\times10^{-7}\,\si{\tesla\meter\per\ampere}$.}
\resposta{$B = \SI{40}{\micro\tesla}$}
\resolucao{$B = \dfrac{\mu_0 I}{2\pi r} = \pot{2}{-7}\cdot\dfrac{4}{0{,}02}
= \pot{4}{-5}\,\si{\tesla} = \SI{40}{\micro\tesla}$
(usando o atalho $\mu_0/2\pi = \pot{2}{-7}$).}
\end{questao}

\blocoaprofundamento

\begin{questao}[3]{}
Dois fios paralelos, distantes $\SI{10}{\centi\meter}$, conduzem correntes de
$\SI{3}{\ampere}$ e $\SI{6}{\ampere}$ no \textbf{mesmo sentido}. Determine o campo
magnético resultante no ponto médio entre eles e diga se os fios se atraem ou se
repelem.
\dados{$\mu_0 = 4\pi\times10^{-7}\,\si{\tesla\meter\per\ampere}$.}
\resposta{$B \approx \SI{12}{\micro\tesla}$; os fios se atraem}
\resolucao{No ponto médio ($r = \SI{5}{\centi\meter}$ de cada fio), os campos dos
dois fios têm \emph{sentidos opostos} (correntes no mesmo sentido), então os
módulos se subtraem:
\[
B = \frac{\mu_0}{2\pi r}(I_2 - I_1)
= \pot{2}{-7}\cdot\frac{6-3}{0{,}05}
= \pot{2}{-7}\cdot60 = \pot{1.2}{-5}\,\si{\tesla}\approx\SI{12}{\micro\tesla}.
\]
Quanto à força \emph{entre} os fios: correntes de mesmo sentido se
\textbf{atraem} (regra fundamental do eletromagnetismo, base da definição do
ampère). Note a distinção: o \emph{campo} no ponto médio se subtrai, mas a
\emph{força} entre os fios é atrativa --- são perguntas diferentes.}
\end{questao}

\begin{questao}[3]{}
Dois fios paralelos e longos, separados por $\SI{5}{\centi\meter}$, conduzem
$\SI{10}{\ampere}$ cada, em sentidos \textbf{opostos}. Determine a força por
unidade de comprimento entre eles e diga se é de atração ou repulsão.
\dados{$\mu_0 = 4\pi\times10^{-7}\,\si{\tesla\meter\per\ampere}$.}
\resposta{$F/\ell = \SI{4e-4}{\newton\per\meter}$, repulsão}
\resolucao{A força por comprimento entre dois fios paralelos é
\[
\frac{F}{\ell}=\frac{\mu_0 I_1 I_2}{2\pi d}
=\pot{2}{-7}\cdot\frac{10\cdot10}{0{,}05}
=\pot{4}{-4}\,\si{\newton\per\meter}.
\]
Correntes em sentidos \emph{opostos} se \textbf{repelem} (e as de mesmo sentido se
atraem). Esta força é a base da definição do ampère no SI.}
\end{questao}

\begin{questao}[3]{}
Um \textbf{toroide} de raio médio $r = \SI{10}{\centi\meter}$ possui $N = 200$
espiras percorridas por $I = \SI{3}{\ampere}$.
\begin{itensq}
\item Aplicando a lei de Ampère a uma circunferência de raio $r$ interna ao
      toroide, deduza a expressão de $B$ e calcule seu valor.
\item Mostre que o campo \emph{fora} do toroide é nulo.
\item Compare com um solenoide reto de mesmo $N$ e comprimento igual ao
      perímetro médio do toroide. O que se conclui?
\end{itensq}
\dados{$\mu_0 = 4\pi\times10^{-7}\,\si{\tesla\meter\per\ampere}$.}
\resposta{(a) $B = \dfrac{\mu_0 N I}{2\pi r} = \SI{1.2}{\milli\tesla}$;
(b) e (c) ver comentário}
\resolucao{\textbf{(a)} A lei de Ampère estabelece
$\oint \vec B\cdot\mathrm{d}\vec\ell = \mu_0 I_{\text{env}}$. Escolhendo uma
circunferência de raio $r$ concêntrica ao toroide, por simetria $B$ é constante e
tangente ao longo dela:
\[
B\,(2\pi r) = \mu_0\,N I
\;\Longrightarrow\;
B=\frac{\mu_0 N I}{2\pi r}
=\frac{(\pot{4\pi}{-7})(200)(3)}{2\pi(0{,}1)}
=\pot{1.2}{-3}\,\si{\tesla}=\SI{1.2}{\milli\tesla}.
\]
\textbf{(b) Campo externo nulo.} Para um percurso \emph{fora} do toroide há dois
casos: (i) circunferência de raio maior que o externo --- a corrente entra e sai
do enrolamento, de modo que a corrente líquida envolvida é \emph{zero};
(ii) circunferência de raio menor que o interno --- não envolve corrente alguma.
Em ambos, $\oint \vec B\cdot\mathrm{d}\vec\ell = 0$ e, pela simetria,
$B = 0$.\\
\textbf{Consequência prática:} o toroide \textbf{confina} o campo em seu interior,
praticamente sem dispersão. Por isso transformadores toroidais produzem muito
menos interferência eletromagnética que os de núcleo E-I, e são preferidos em
áudio e instrumentação de precisão.

\textbf{(c) Comparação com o solenoide.} O perímetro médio é
$\ell = 2\pi r = \SI{0.628}{\meter}$. Um solenoide reto com o mesmo $N$ e esse
comprimento teria
\[
B = \mu_0\,\frac{N}{\ell}\,I = \frac{\mu_0 N I}{2\pi r},
\]
\emph{exatamente a mesma expressão}. Ou seja: o toroide é um solenoide
"fechado sobre si mesmo". A diferença essencial é que o solenoide reto tem
\textbf{extremidades}, onde as linhas de campo escapam (efeito de borda) e o
campo é menos uniforme; o toroide não tem bordas, sendo o dispositivo ideal para
estudar o comportamento magnético puro de um material --- razão pela qual as
curvas de histerese são levantadas em amostras toroidais.}
\end{questao}

\begin{questao}[3]{Apostila original revisada}
Duas espiras circulares, concêntricas e coplanares, têm raios
$R_1=10\pi\,\si{\centi\metre}$ e $R_2=2{,}5\pi\,\si{\centi\metre}$. A primeira é
percorrida por $I_1=\SI{2}{\ampere}$. Determine o valor e o sentido de $I_2$ para
que o campo magnético resultante no centro seja nulo.
\resposta{$I_2=\SI{0.5}{\ampere}$, em sentido oposto a $I_1$}
\resolucao{No centro de uma espira, $B=\mu_0 I/(2R)$. Para cancelamento,
$I_1/R_1=I_2/R_2$. Assim,
$I_2=I_1(R_2/R_1)=2(2{,}5/10)=\SI{0.5}{\ampere}$. Os sentidos devem ser opostos.}
\end{questao}

% BEGIN BANCO COMPLEMENTAR Campo magnético
% =====================================================================
%  Banco complementar --- Campo Magnetico  (REVISADO)
%  Substitui a versao gerada por template (5 clones em N1, 9 em N2,
%  8 em N3 e 8 questoes conceituais marcadas como N4).
%  O cap11 ja cobre: B proporcional a 1/r (MC), fio de 10 A a 5 cm,
%  regra da mao direita, solenoide de 1000 esp/m, ESPIRA CIRCULAR de
%  R=10 cm no centro, fio de 4 A a 2 cm, DOIS FIOS PARALELOS (mesmo
%  sentido e sentidos opostos), TOROIDE e duas espiras concentricas.
%  Este banco evita esses casos e explora: fio de secao finita por
%  lei de Ampere, cabo coaxial, campo no eixo de espira, dipolo
%  magnetico a grande distancia, projeto de solenoide, mapeamento com
%  sensor Hall, blindagem magnetica e os limites praticos de eletroimas.
%  Distribuicao mantida: 5 N1 + 9 N2 + 8 N3 + 8 N4 = 30
% =====================================================================

\subsubsection*{Banco complementar --- Campo magnético}

\begin{atencao}[Organização do banco]
Três resultados cobrem quase tudo: fio retilíneo longo
($B=\mu_0I/2\pi r$), centro de espira ($B=\mu_0I/2R$) e solenoide longo
($B=\mu_0nI$, \emph{independente} do raio). Note que apenas o primeiro decai com
a distância --- dentro do solenoide o campo é uniforme. O N3 trata de condutores
de seção finita, cabo coaxial e campo fora do eixo; o N4 exige dedução por
Biot--Savart, projeto, blindagem e limites físicos.
\end{atencao}

\blocofixacao

\begin{questao}[1]{}
Um fio retilíneo longo conduz $\SI{2}{\ampere}$. Determine o campo magnético a
$\SI{1}{\centi\meter}$ dele.
\dados{$\mu_0=4\pi\times10^{-7}\,\si{\tesla\meter\per\ampere}$}
\resposta{$B=\SI{40}{\micro\tesla}$}
\resolucao{
\[
B=\frac{\mu_0 I}{2\pi r}=\frac{(4\pi\times10^{-7})(2)}{2\pi(0{,}01)}
=\frac{\pot{2}{-7}\times2}{0{,}01}=\pot{4}{-5}\,\si{\tesla}.
\]
\textbf{Atalho útil:} $\dfrac{\mu_0}{2\pi}=\pot{2}{-7}$ exatamente, de modo que
$B=\pot{2}{-7}\dfrac{I}{r}$ --- evita carregar $\pi$ pela conta toda.\\
\textbf{Ordem de grandeza:} $\SI{40}{\micro\tesla}$ é comparável ao campo
magnético da Terra ($\approx\SI{50}{\micro\tesla}$). Uma bússola perto desse fio
já seria perturbada.}
\end{questao}

\begin{questao}[1]{}
Um solenoide tem $\SI{2000}{}$ espiras por metro e conduz
$\SI{3}{\ampere}$. Determine o campo em seu interior.
\resposta{$B=\SI{7.54}{\milli\tesla}$}
\resolucao{
\[
B=\mu_0 nI=(4\pi\times10^{-7})(2000)(3)=\pot{7.54}{-3}\,\si{\tesla}.
\]
\textbf{Note o que \emph{não} aparece na fórmula:} o raio do solenoide e o
número total de espiras. Só importa a \textbf{densidade linear}
$n=N/L$.\\
Dois solenoides, um fino e um grosso, com a mesma densidade de espiras e a mesma
corrente, produzem o \emph{mesmo} campo interno.}
\end{questao}

\begin{questao}[1]{}
Determine a corrente necessária para que um fio retilíneo produza
$B=\SI{100}{\micro\tesla}$ a $\SI{2}{\centi\meter}$.
\resposta{$I=\SI{10}{\ampere}$}
\resolucao{
\[
I=\frac{2\pi rB}{\mu_0}=\frac{rB}{\pot{2}{-7}}
=\frac{(0{,}02)(\pot{1}{-4})}{\pot{2}{-7}}=\SI{10}{\ampere}.
\]
\textbf{Leitura prática:} são necessários $\SI{10}{\ampere}$ para produzir, a
dois centímetros, apenas o dobro do campo terrestre. Campos magnéticos
intensos exigem correntes altas --- ou muitas espiras.}
\end{questao}

\begin{questao}[1]{}
Sobre o campo magnético de um solenoide longo, é correto afirmar:
\begin{alternativas}[2]
\item é uniforme no interior e praticamente nulo no exterior
\item é maior no exterior
\item aumenta com a distância ao eixo, no interior
\item é nulo no interior
\end{alternativas}
\resposta{(a)}
\resolucao{Aplicando a lei de Ampère a um retângulo com um lado dentro e outro
fora, obtém-se $B_{int}=\mu_0nI$ e $B_{ext}\approx0$ para um solenoide
\emph{ideal} (infinitamente longo).\\
\textbf{Por que o campo externo quase se anula:} as linhas que saem de uma
extremidade retornam pela outra, espalhando-se por um volume enorme --- a
densidade de linhas fora é desprezível diante da interna.\\
\textbf{Em solenoides reais}, o campo externo não é exatamente nulo, e nas
extremidades o campo interno cai para cerca de \emph{metade} do valor central.}
\end{questao}

\begin{questao}[1]{}
O campo magnético da Terra vale aproximadamente $\SI{50}{\micro\tesla}$.
Expresse-o em gauss.
\dados{$\SI{1}{\tesla}=10^{4}\,\si{\gauss}$}
\resposta{$\SI{0.5}{\gauss}$}
\resolucao{
\[
B=\pot{5}{-5}\,\si{\tesla}\times10^{4}=\SI{0.5}{\gauss}.
\]
\textbf{Escala de referência para o tópico:}
\begin{center}
\begin{tabular}{lcc}
\hline
Situação & Tesla & Gauss\\
\hline
Campo terrestre & $\pot{5}{-5}$ & $0{,}5$\\
Ímã de geladeira & $\pot{5}{-3}$ & $50$\\
Ímã de neodímio (superfície) & $0{,}5$ & $5000$\\
Ressonância magnética & $1{,}5$ a $3$ & $\pot{1.5}{4}$ a $\pot{3}{4}$\\
\hline
\end{tabular}
\end{center}
O tesla é uma unidade \emph{grande}: campos de laboratório raramente passam de
alguns tesla, e o gauss continua em uso justamente por ser mais próximo dos
valores cotidianos.}
\end{questao}

\blocoaplicacao

\begin{questao}[2]{}
Um solenoide de $100$ espiras e $\SI{10}{\centi\meter}$ de comprimento conduz
$\SI{1}{\ampere}$, com núcleo de ar.
\begin{itensq}
\item Determine a densidade de espiras.
\item Determine o campo interno.
\end{itensq}
\resposta{$n=\SI{1000}{\per\meter}$; $B=\SI{1.26}{\milli\tesla}$}
\resolucao{(a) $n=\dfrac{100}{0{,}10}=\SI{1000}{\per\meter}$.\\
(b) $B=\mu_0nI=(4\pi\times10^{-7})(1000)(1)
=\pot{1.257}{-3}\,\si{\tesla}$.\\
\textbf{Comparação instrutiva:} isso é apenas $25$ vezes o campo terrestre ---
um solenoide modesto produz campo modesto. Para chegar a
$\SI{0.1}{\tesla}$ seriam necessários $\SI{80}{\ampere}$ com a mesma
geometria, o que exigiria refrigeração.\\
\textbf{Caminho alternativo:} inserir núcleo ferromagnético, que multiplica
$B$ por $\mu_r$ --- centenas ou milhares de vezes.}
\end{questao}

\begin{questao}[2]{}
Um fio conduz corrente constante. Compare o campo a $\SI{4}{\centi\meter}$ com
o campo a $\SI{12}{\centi\meter}$.
\resposta{$B_1/B_2=3$}
\resolucao{Como $B\propto1/r$:
\[
\frac{B_1}{B_2}=\frac{r_2}{r_1}=\frac{12}{4}=3.
\]
\textbf{Contraste importante com o campo elétrico:} o campo de uma carga
puntiforme cai com $1/r^{2}$, mas o de um fio \emph{infinito} cai com
$1/r$.\\
A razão é a mesma da eletrostática: o fio é uma fonte \emph{estendida} em uma
dimensão, e a superfície amperiana (um cilindro) tem área proporcional a $r$,
não a $r^{2}$.}
\end{questao}

\begin{questao}[2]{}
Uma espira circular de raio $\SI{5}{\centi\meter}$ deve produzir
$\SI{100}{\micro\tesla}$ em seu centro. Determine a corrente.
\resposta{$I\approx\SI{7.96}{\ampere}$}
\resolucao{Do campo no centro de uma espira, $B=\dfrac{\mu_0I}{2R}$:
\[
I=\frac{2RB}{\mu_0}=\frac{2(0{,}05)(\pot{1}{-4})}{4\pi\times10^{-7}}
=\SI{7.96}{\ampere}.
\]
\textbf{Comparação com o fio reto:} um fio retilíneo a
$\SI{5}{\centi\meter}$ precisaria de $\SI{25}{\ampere}$ para o mesmo campo. A
espira é mais eficiente porque \emph{todos} os seus elementos contribuem no
mesmo sentido no centro --- a razão entre as duas fórmulas é exatamente
$\pi$.}
\end{questao}

\begin{questao}[2]{}
Um solenoide tem $N$ espiras e comprimento $L$. Determine o que ocorre com o
campo interno se $N$ e $L$ forem \textbf{ambos dobrados}.
\resposta{O campo não se altera}
\resolucao{Como $B=\mu_0\dfrac{N}{L}I$, dobrar $N$ e $L$ mantém a razão
$N/L$ inalterada:
\[
B'=\mu_0\frac{2N}{2L}I=B.
\]
\textbf{Interpretação:} o solenoide dobrado equivale a dois solenoides
originais colocados em sequência. Cada trecho "enxerga" a mesma vizinhança de
espiras, e o campo local não muda.\\
\textbf{O que muda:} a região de campo uniforme fica mais longa, e a resistência
do enrolamento dobra --- exigindo o dobro da tensão para a mesma corrente.}
\end{questao}

\begin{questao}[2]{}
Dois fios longos e paralelos, distantes $\SI{8}{\centi\meter}$, conduzem
$\SI{4}{\ampere}$ cada, no \textbf{mesmo} sentido. Determine o campo no ponto
médio entre eles.
\resposta{$B=0$}
\resolucao{Cada fio produz, no ponto médio,
\[
B=\pot{2}{-7}\frac{4}{0{,}04}=\pot{2}{-5}\,\si{\tesla}.
\]
Pela regra da mão direita, com as correntes no \emph{mesmo} sentido, os dois
campos têm sentidos \textbf{opostos} no ponto médio --- e módulos iguais.
Cancelam-se:
\[
B_{total}=0.
\]
\textbf{Cuidado com a intuição:} correntes de mesmo sentido \emph{atraem-se},
mas seus campos se \emph{cancelam} entre elas. Força e campo são coisas
diferentes.\\
\textbf{Fora do par}, ao contrário, os campos se somam --- e é por isso que dois
fios juntos conduzindo no mesmo sentido produzem, de longe, o campo de um único
fio de corrente $2I$.}
\end{questao}

\begin{questao}[2]{}
Um fio reto é imerso em um meio de permeabilidade relativa $\mu_r=200$.
Determine o efeito sobre o campo, comparado ao ar.
\resposta{O campo fica $200$ vezes maior}
\resolucao{Em meio linear,
\[
B=\frac{\mu_r\mu_0I}{2\pi r},
\]
de modo que $B$ é multiplicado por $\mu_r=200$.\\
\textbf{Mecanismo:} o material se \emph{magnetiza}, e seus dipolos alinhados
somam sua própria contribuição ao campo aplicado. Diferentemente do dielétrico
--- que \emph{reduz} o campo elétrico --- o material ferromagnético o
\textbf{amplifica}.\\
\textbf{Ressalva séria:} $\mu_r$ não é constante. Ele cai drasticamente quando o
material satura, e a relação $B\times H$ deixa de ser linear. A fórmula acima
vale apenas em campos baixos.}
\end{questao}

\begin{questao}[2]{}
Determine a distância a que um fio de $\SI{50}{\ampere}$ produz campo igual ao
terrestre ($\SI{50}{\micro\tesla}$).
\resposta{$r=\SI{20}{\centi\meter}$}
\resolucao{
\[
r=\frac{\pot{2}{-7}I}{B}=\frac{(\pot{2}{-7})(50)}{\pot{5}{-5}}
=\SI{0.20}{\meter}.
\]
\textbf{Consequência prática:} a menos de $\SI{20}{\centi\meter}$ de um
condutor de $\SI{50}{\ampere}$, uma bússola aponta mais para o fio do que para
o norte.\\
É por isso que instrumentos magnéticos sensíveis --- e monitores de tubo, na
época em que existiam --- precisam ficar afastados de cabos de potência.}
\end{questao}

\begin{questao}[2]{}
Um toroide e um solenoide têm a mesma densidade de espiras e a mesma corrente.
Compare seus campos internos e externos.
\resposta{Campo interno praticamente igual; o toroide tem campo externo
\emph{rigorosamente} nulo}
\resolucao{\textbf{Campo interno:} ambos valem $\approx\mu_0nI$. No toroide,
$n=N/(2\pi r)$ varia ligeiramente com o raio, de modo que o campo não é
perfeitamente uniforme na seção --- mas para toroides finos a diferença é
pequena.\\
\textbf{Campo externo:} aqui está a diferença essencial.
\begin{itemize}
\item \textbf{Solenoide:} as linhas precisam retornar por fora, e o campo
externo, embora fraco, \emph{existe}.
\item \textbf{Toroide:} as linhas fecham-se \emph{inteiramente} dentro do
núcleo. Aplicando a lei de Ampère a um círculo externo, a corrente líquida
envolvida é \textbf{zero} --- e o campo também.
\end{itemize}
\textbf{Consequência de projeto:} indutores toroidais praticamente não irradiam
nem captam interferência, e podem ser montados lado a lado sem acoplamento
mútuo. É por isso que dominam aplicações de filtragem em fontes chaveadas.}
\end{questao}

\begin{questao}[2]{}
Dois fios perpendiculares entre si, ambos no mesmo plano da folha, conduzem
$\SI{6}{\ampere}$ e $\SI{8}{\ampere}$. Determine o campo resultante em um
ponto a $\SI{3}{\centi\meter}$ do primeiro e $\SI{4}{\centi\meter}$ do
segundo.
\resposta{$B=\SI{56.6}{\micro\tesla}$}
\resolucao{Cada fio produz campo perpendicular ao plano definido por ele e pelo
ponto. Como os fios são perpendiculares entre si e coplanares, ambos os campos
saem (ou entram) \emph{perpendicularmente à folha} --- portanto são
\textbf{paralelos} e se somam algebricamente:
\[
B_1=\pot{2}{-7}\frac{6}{0{,}03}=\SI{40}{\micro\tesla};
\qquad
B_2=\pot{2}{-7}\frac{8}{0{,}04}=\SI{40}{\micro\tesla}.
\]
Se os sentidos coincidirem: $B=\SI{80}{\micro\tesla}$; se forem opostos,
$B=0$.\\
\textbf{Atenção ao caso geral:} se os fios \emph{não} fossem coplanares, os
campos teriam direções diferentes e a soma seria vetorial --- com
$\SI{40}{\micro\tesla}$ perpendiculares, o resultado seria
$40\sqrt2=\SI{56.6}{\micro\tesla}$. \textbf{A geometria decide se a soma é
algébrica ou vetorial}, e é ela que precisa ser estabelecida antes de qualquer
conta.}
\end{questao}

\blocoaprofundamento

\begin{questao}[3]{}
Um condutor cilíndrico maciço de raio $a=\SI{2}{\milli\meter}$ conduz
$\SI{100}{\ampere}$ uniformemente distribuídos na seção.
\begin{itensq}
\item Determine o campo a $\SI{1}{\milli\meter}$, na superfície e a
      $\SI{5}{\milli\meter}$ do eixo.
\item Identifique onde o campo é máximo.
\end{itensq}
\resposta{$5$, $10$ e $\SI{4}{\milli\tesla}$; máximo na superfície}
\resolucao{\textbf{Dentro ($r<a$):} a superfície amperiana envolve apenas a
fração $(r/a)^{2}$ da corrente:
\[
B=\frac{\mu_0 I r}{2\pi a^{2}}
=\pot{2}{-7}\frac{100\,r}{(0{,}002)^{2}}.
\]
Em $r=\SI{1}{\milli\meter}$: $B=\SI{5}{\milli\tesla}$.\\
\textbf{Fora ($r>a$):} toda a corrente é envolvida, e vale a fórmula do fio
fino:
\[
B=\pot{2}{-7}\frac{100}{r}.
\]
Em $r=\SI{5}{\milli\meter}$: $B=\SI{4}{\milli\tesla}$.\\
\textbf{Na superfície ($r=a$):} as duas expressões coincidem em
$B=\SI{10}{\milli\tesla}$ --- o campo é contínuo.\\
(b) \textbf{O máximo está na superfície.} Dentro, $B$ cresce linearmente com
$r$; fora, decai com $1/r$. O pico ocorre exatamente em $r=a$.\\
\textbf{Analogia com a esfera isolante carregada:} lá o campo elétrico também
cresce linearmente dentro e cai com $1/r^{2}$ fora, com máximo na superfície. A
estrutura matemática é a mesma --- muda apenas a lei de Gauss ou de Ampère
aplicada.}
\end{questao}

\begin{questao}[3]{}
Um cabo coaxial tem condutor central de raio $a$ conduzindo $I$ e blindagem
externa conduzindo $I$ em sentido \textbf{oposto}, com $I=\SI{5}{\ampere}$.
\begin{itensq}
\item Determine o campo entre os condutores, a $\SI{2}{\milli\meter}$ do eixo.
\item Determine o campo fora do cabo.
\item Explique a consequência prática.
\end{itensq}
\resposta{$B=\SI{0.5}{\milli\tesla}$ entre os condutores; $B=0$ fora}
\resolucao{(a) Entre os condutores, a superfície amperiana envolve apenas a
corrente central:
\[
B=\pot{2}{-7}\frac{5}{0{,}002}=\pot{5}{-4}\,\si{\tesla}=\SI{0.5}{\milli\tesla}.
\]
(b) \textbf{Fora do cabo}, a superfície amperiana envolve $+I$ e $-I$ --- a
corrente líquida é \textbf{zero}:
\[
\oint\vec B\cdot\mathrm{d}\vec\ell=\mu_0(0)=0
\;\Longrightarrow\;
B=0.
\]
(c) \textbf{Consequência prática: o cabo coaxial não irradia nem capta campo
magnético.} Todo o campo fica confinado entre os condutores.\\
Isso o torna a escolha padrão para transmitir sinais em ambientes com
interferência --- e explica por que um cabo paralelo comum, cujos condutores
estão afastados, é muito pior: ali as correntes opostas não se cancelam
perfeitamente e o par forma uma espira que irradia.\\
\textbf{Mesma ideia, outra forma:} o par trançado obtém efeito semelhante ao
inverter a orientação da espira a cada torção, de modo que as contribuições
sucessivas se anulem.}
\end{questao}

\begin{questao}[3]{}
Uma espira de raio $R=\SI{10}{\centi\meter}$ conduz $\SI{5}{\ampere}$.
\begin{itensq}
\item Determine o campo no centro.
\item Determine o campo no eixo, a $\SI{10}{\centi\meter}$ e a
      $\SI{20}{\centi\meter}$ do centro.
\item Comente a taxa de queda.
\end{itensq}
\resposta{$31{,}4$; $11{,}1$ e $\SI{2.81}{\micro\tesla}$}
\resolucao{(a) $B_0=\dfrac{\mu_0I}{2R}=\dfrac{(4\pi\times10^{-7})(5)}{0{,}20}
=\pot{3.14}{-5}\,\si{\tesla}$.\\
(b) No eixo,
\[
B=\frac{\mu_0IR^{2}}{2\left(R^{2}+x^{2}\right)^{3/2}}.
\]
Em $x=R=\SI{0.10}{\meter}$:
$B=\pot{1.11}{-5}\,\si{\tesla}$, ou seja
$B_0/2^{3/2}=0{,}354\,B_0$.\\
Em $x=2R=\SI{0.20}{\meter}$: $B=\pot{2.81}{-6}\,\si{\tesla}$, cerca de
$0{,}089\,B_0$.\\
(c) \textbf{A queda é rápida.} A apenas um raio de distância, o campo já caiu
para $35\%$; a dois raios, para $9\%$.\\
Para $x\gg R$, a expressão tende a
$B\to\dfrac{\mu_0IR^{2}}{2x^{3}}$ --- decaimento com $1/x^{3}$, típico de
\textbf{dipolo}. É o mesmo expoente do dipolo elétrico, e pela mesma razão: a
espira não tem "monopolo magnético" para produzir termo $1/x^{2}$.}
\end{questao}

\begin{questao}[3]{}
Compare o campo no centro de um solenoide \emph{finito} com o do solenoide
ideal, e determine o campo em sua extremidade.
\begin{itensq}
\item Indique a condição para a aproximação ideal.
\item Determine o campo na extremidade.
\item Discuta a consequência para projeto.
\end{itensq}
\resposta{Na extremidade, $B\approx\mu_0nI/2$ --- metade do valor central}
\resolucao{(a) \textbf{Condição:} $L\gg D$, isto é, comprimento muito maior que
o diâmetro. Na prática, $L/D>10$ garante erro inferior a $1\%$ no centro.\\
(b) \textbf{Na extremidade}, apenas "metade" do solenoide contribui --- as
espiras estão todas de um lado. O resultado exato para solenoide longo é
\[
B_{extremidade}=\frac{\mu_0nI}{2},
\]
exatamente \textbf{metade} do valor central.\\
(c) \textbf{Consequências de projeto.}
\begin{itemize}
\item A região de campo verdadeiramente uniforme é \emph{menor} que o
solenoide: ela se restringe à parte central, tipicamente à metade do
comprimento.
\item Para uniformidade melhor que $1\%$ em uma região de comprimento
$\ell$, o solenoide precisa ter comprimento bem maior que $\ell$ --- ou usar
enrolamento com densidade \emph{corrigida} nas extremidades.
\item A alternativa clássica é o par de \textbf{bobinas de Helmholtz}: duas
bobinas iguais separadas por uma distância igual ao raio, arranjo que anula a
segunda derivada do campo no centro e produz região uniforme com pouco
material.
\end{itemize}}
\end{questao}

\begin{questao}[3]{}
Três fios longos e paralelos, coplanares e igualmente espaçados de
$\SI{5}{\centi\meter}$, conduzem $\SI{10}{\ampere}$ cada, todos no mesmo
sentido. Determine o campo no fio central.
\resposta{$B=0$ no fio central}
\resolucao{O fio central não produz campo sobre si mesmo. Restam as
contribuições dos dois vizinhos, ambos a $\SI{5}{\centi\meter}$:
\[
B_1=B_2=\pot{2}{-7}\frac{10}{0{,}05}=\SI{40}{\micro\tesla}.
\]
Pela regra da mão direita, com os dois vizinhos conduzindo no mesmo sentido e
estando em \emph{lados opostos} do fio central, seus campos ali têm sentidos
\textbf{contrários}:
\[
B=40-40=0.
\]
\textbf{Mas a força não é nula --- ela também é.} Como $\vec F=I\vec\ell\times
\vec B$ e $\vec B=\vec 0$ no fio central, a força sobre ele é zero: as atrações
dos dois vizinhos se equilibram.\\
\textbf{Nos fios externos, porém}, a situação é diferente: cada um sofre atração
resultante para o centro. O conjunto tende a se \emph{comprimir} --- efeito real
em barramentos de alta corrente, que precisam de espaçadores mecânicos.}
\end{questao}

\begin{questao}[3]{}
Um eletroímã deve produzir $\SI{0.1}{\tesla}$ com núcleo de ar, usando
solenoide de $\SI{20}{\centi\meter}$.
\begin{itensq}
\item Determine o produto $NI$ necessário.
\item Avalie a viabilidade com $N=1000$ espiras.
\item Proponha alternativa.
\end{itensq}
\resposta{$NI=\SI{15.9}{\kilo\ampere}$-espira; exigiria $\SI{15.9}{\ampere}$
com $1000$ espiras}
\resolucao{(a) De $B=\mu_0NI/L$:
\[
NI=\frac{BL}{\mu_0}=\frac{(0{,}1)(0{,}20)}{4\pi\times10^{-7}}
=\pot{1.59}{4}\ \text{A-espira}.
\]
(b) Com $N=1000$: $I=\SI{15.9}{\ampere}$.\\
\textbf{Viável, mas problemático.} Mil espiras conduzindo
$\SI{16}{\ampere}$ exigem fio grosso (a densidade de corrente admissível limita
a seção), o que aumenta o volume do enrolamento. E a dissipação é considerável:
com resistência de apenas $\SI{1}{\ohm}$, seriam
$\SI{256}{\watt}$ --- exigindo refrigeração.\\
(c) \textbf{Alternativa: núcleo ferromagnético.} Com $\mu_r=1000$, a mesma
corrente produziria $\SI{100}{\tesla}$ --- valor impossível, porque o material
\textbf{satura} em torno de $1{,}5$ a $\SI{2}{\tesla}$.\\
\textbf{O que o núcleo realmente entrega:} os $\SI{0.1}{\tesla}$ pedidos com
corrente cerca de mil vezes menor ($\SI{16}{\milli\ampere}$), enquanto o
material permanecer na região linear.\\
\textbf{Limite físico:} acima da saturação, o núcleo deixa de ajudar e o
eletroímã volta a se comportar como se fosse de ar. É por isso que campos acima
de $\SI{2}{\tesla}$ exigem bobinas supercondutoras, sem núcleo.}
\end{questao}

\begin{questao}[3]{}
Um cabo de par paralelo tem condutores separados por $\SI{2}{\milli\meter}$,
conduzindo $\SI{1}{\ampere}$ em sentidos opostos.
\begin{itensq}
\item Determine o campo no ponto médio entre eles.
\item Determine o campo a $\SI{10}{\centi\meter}$ do par.
\item Compare com um fio único.
\end{itensq}
\resposta{$\SI{400}{\micro\tesla}$ no meio; $\approx\SI{0.04}{\micro\tesla}$ a
$\SI{10}{\centi\meter}$}
\resolucao{(a) No ponto médio, cada fio dista $\SI{1}{\milli\meter}$, e com
correntes \emph{opostas} os campos se \textbf{somam}:
\[
B=2\times\pot{2}{-7}\frac{1}{0{,}001}=\SI{400}{\micro\tesla}.
\]
(b) A $\SI{10}{\centi\meter}$, os dois fios estão praticamente à mesma
distância e seus campos quase se cancelam. O resultado é aproximadamente o de um
\textbf{dipolo}:
\[
B\approx\frac{\mu_0 I d}{2\pi r^{2}}
=\pot{2}{-7}\frac{(1)(0{,}002)}{(0{,}10)^{2}}
=\pot{4}{-8}\,\si{\tesla}.
\]
(c) \textbf{Um fio único} de $\SI{1}{\ampere}$ produziria, a
$\SI{10}{\centi\meter}$:
$B=\pot{2}{-6}\,\si{\tesla}$ --- \textbf{cinquenta vezes mais}.\\
\textbf{Conclusão:} manter os condutores de ida e volta \emph{juntos} reduz
drasticamente o campo irradiado, porque o par se comporta como dipolo
($1/r^{2}$) em vez de fio isolado ($1/r$). Quanto menor a separação $d$, melhor
--- e é exatamente por isso que se recomenda não separar os condutores de um
mesmo circuito.}
\end{questao}

\begin{questao}[3]{}
Um sensor Hall indica $\SI{2.5}{\milli\volt}$ quando exposto a
$\SI{100}{\milli\tesla}$, com resposta linear e tensão de repouso nula.
\begin{itensq}
\item Determine a sensibilidade.
\item Determine o campo correspondente a uma leitura de
      $\SI{0.6}{\milli\volt}$.
\item Discuta as fontes de erro.
\end{itensq}
\resposta{$\SI{25}{\micro\volt\per\milli\tesla}$; $B=\SI{24}{\milli\tesla}$}
\resolucao{(a)
\[
S=\frac{2{,}5\,\si{\milli\volt}}{100\,\si{\milli\tesla}}
=\SI{25}{\micro\volt\per\milli\tesla}.
\]
(b) $B=\dfrac{0{,}6}{0{,}025}=\SI{24}{\milli\tesla}$.\\
(c) \textbf{Fontes de erro, em ordem de importância.}
\begin{itemize}
\item \textbf{Tensão de \emph{offset}:} sensores reais indicam algo mesmo com
$B=0$, por assimetria construtiva. Valores de dezenas de microvolts são comuns
--- comparáveis ao sinal de alguns militesla. \emph{Correção:} medir com o
sensor blindado e subtrair.
\item \textbf{Deriva térmica:} tanto o \emph{offset} quanto a sensibilidade
variam com a temperatura.
\item \textbf{Orientação:} o sensor responde apenas à componente
\emph{perpendicular} à sua face. Um erro angular $\theta$ introduz fator
$\cos\theta$ --- desprezível para ângulos pequenos, mas $13\%$ a
$\ang{30}$.
\item \textbf{Campo terrestre} somado à medida, se não for compensado.
\end{itemize}
\textbf{Procedimento robusto:} medir com o sensor em uma orientação, depois
girá-lo $\ang{180}$ e medir de novo. A média das duas leituras cancela o
\emph{offset}, e a semidiferença dá o campo.}
\end{questao}

\blocodesafio

\begin{questao}[4]{}
\textbf{(Dedução por Biot--Savart.)} Deduza o campo magnético no centro de uma
espira circular de raio $R$ percorrida por corrente $I$.
\begin{itensq}
\item Aplique a lei de Biot--Savart a um elemento de corrente.
\item Integre e obtenha o resultado.
\item Explique por que a integração é simples neste caso.
\end{itensq}
\resposta{$B=\dfrac{\mu_0 I}{2R}$}
\resolucao{(a) A lei de Biot--Savart estabelece que um elemento
$I\,\mathrm{d}\vec\ell$ produz, a uma distância $\vec r$,
\[
\mathrm{d}\vec B=\frac{\mu_0}{4\pi}\,
\frac{I\,\mathrm{d}\vec\ell\times\hat r}{r^{2}}.
\]
Para o \textbf{centro} da espira, todo elemento está à mesma distância $R$, e
$\mathrm{d}\vec\ell$ é sempre \emph{perpendicular} a $\hat r$ --- logo
$|\mathrm{d}\vec\ell\times\hat r|=\mathrm{d}\ell$.\\
(b) Além disso, todos os $\mathrm{d}\vec B$ apontam na \emph{mesma} direção
(perpendicular ao plano da espira), de modo que a integral vetorial vira
escalar:
\[
B=\frac{\mu_0I}{4\pi R^{2}}\oint\mathrm{d}\ell
=\frac{\mu_0I}{4\pi R^{2}}(2\pi R)
=\frac{\mu_0 I}{2R}.
\]
(c) \textbf{Três simplificações coincidiram:} distância constante, produto
vetorial de módulo máximo (ângulo reto) e contribuições todas paralelas. É a
combinação mais favorável possível.\\
\textbf{Fora do centro isso se perde:} no eixo, a distância deixa de ser $R$ e
as contribuições ganham componentes radiais que se cancelam --- exigindo
projeção. Fora do eixo, a integral não tem forma fechada elementar e recorre-se a
métodos numéricos ou a funções elípticas.}
\end{questao}

\begin{questao}[4]{}
\textbf{(Ampère contra Biot--Savart.)} Compare os dois métodos de cálculo de
campo magnético.
\begin{itensq}
\item Enuncie cada lei e sua condição de uso.
\item Determine quando a lei de Ampère é praticável.
\item Classifique quatro configurações quanto ao método adequado.
\end{itensq}
\resposta{Ampère exige simetria suficiente para tirar $B$ da integral;
Biot--Savart é geral mas trabalhoso}
\resolucao{(a) \textbf{Biot--Savart:}
$\mathrm{d}\vec B=\dfrac{\mu_0}{4\pi}
\dfrac{I\,\mathrm{d}\vec\ell\times\hat r}{r^{2}}$ --- vale \emph{sempre}, para
qualquer distribuição de corrente. É a lei fundamental.\\
\textbf{Ampère:} $\oint\vec B\cdot\mathrm{d}\vec\ell=\mu_0I_{env}$ --- também
sempre verdadeira, mas só \emph{útil} para calcular $B$ em casos especiais.\\
(b) \textbf{Ampère é praticável quando} existe uma curva fechada ao longo da
qual $|\vec B|$ é constante e $\vec B$ é paralelo (ou perpendicular) ao
caminho. Só assim $B$ pode ser retirado da integral:
\[
\oint\vec B\cdot\mathrm{d}\vec\ell=B\oint\mathrm{d}\ell=B\,\ell.
\]
Isso exige \textbf{simetria} --- cilíndrica, plana ou toroidal --- que precisa
ser \emph{conhecida de antemão}, por argumento físico.\\
(c)
\begin{center}
\begin{tabular}{lll}
\hline
Configuração & Método & Motivo\\
\hline
Fio retilíneo infinito & Ampère & simetria cilíndrica\\
Solenoide/toroide ideal & Ampère & simetria evidente\\
Espira circular (centro) & Biot--Savart & sem curva de $B$ constante\\
Segmento finito de fio & Biot--Savart & sem simetria alguma\\
\hline
\end{tabular}
\end{center}
\textbf{Erro comum:} tentar aplicar Ampère a uma espira circular usando a
própria espira como caminho. O campo \emph{não} é constante nem tangente ao
longo dela --- a integral existe, mas não permite isolar $B$.\\
\textbf{Analogia com a eletrostática:} Gauss está para Coulomb assim como
Ampère está para Biot--Savart. Em ambos os casos, a lei integral é elegante e
poderosa \emph{apenas} quando a simetria coopera.}
\end{questao}

\begin{questao}[4]{}
\textbf{(Projeto de solenoide.)} Projete um solenoide de $500$ espiras e
$\SI{25}{\centi\meter}$ para produzir $\SI{1}{\milli\tesla}$.
\begin{itensq}
\item Determine a corrente necessária.
\item Determine a potência dissipada, para fio de resistência total
      $\SI{2}{\ohm}$.
\item Avalie a uniformidade do campo.
\end{itensq}
\resposta{$I=\SI{0.398}{\ampere}$; $P=\SI{0.32}{\watt}$}
\resolucao{(a) $n=\dfrac{500}{0{,}25}=\SI{2000}{\per\meter}$, e
\[
I=\frac{B}{\mu_0 n}=\frac{\pot{1}{-3}}{(4\pi\times10^{-7})(2000)}
=\SI{0.398}{\ampere}.
\]
(b) $P=RI^{2}=2(0{,}398)^{2}=\SI{0.32}{\watt}$ --- desprezível, sem necessidade
de refrigeração.\\
(c) \textbf{Uniformidade.} Supondo diâmetro de $\SI{4}{\centi\meter}$, a razão
$L/D=6{,}25$ --- razoável, mas não excelente. Consequências:
\begin{itemize}
\item No \textbf{centro}, o campo fica cerca de $1$ a $2\%$ abaixo do valor
ideal $\mu_0nI$.
\item Nas \textbf{extremidades}, cai para aproximadamente metade.
\item A região com uniformidade melhor que $1\%$ limita-se ao terço central,
cerca de $\SI{8}{\centi\meter}$.
\end{itemize}
\textbf{Se for necessária melhor uniformidade:} alongar o solenoide (mantendo
$n$, o que exige mais espiras e mais tensão), ou adotar bobinas de Helmholtz.\\
\textbf{Verificação de coerência:} $\SI{1}{\milli\tesla}$ é vinte vezes o campo
terrestre --- o experimento precisa ser orientado ou compensado, sob pena de o
campo ambiente introduzir erro de $5\%$.}
\end{questao}

\begin{questao}[4]{}
\textbf{(Dipolo magnético.)} Uma espira de raio $R$ e corrente $I$ é observada
no eixo, a distância $x\gg R$.
\begin{itensq}
\item Obtenha a forma assintótica do campo.
\item Defina o momento de dipolo magnético e reescreva o resultado.
\item Compare com o dipolo elétrico e comente a ausência de monopolos.
\end{itensq}
\resposta{$B\to\dfrac{\mu_0 m}{2\pi x^{3}}$, com $m=I\pi R^{2}$}
\resolucao{(a) Partindo de
$B=\dfrac{\mu_0IR^{2}}{2(R^{2}+x^{2})^{3/2}}$ e fazendo $x\gg R$:
\[
B\to\frac{\mu_0IR^{2}}{2x^{3}}.
\]
(b) Definindo o \textbf{momento de dipolo magnético}
$m=I\,A=I\pi R^{2}$:
\[
B=\frac{\mu_0 m}{2\pi x^{3}}.
\]
\textbf{Verificação numérica:} para $R=\SI{10}{\centi\meter}$ e
$I=\SI{5}{\ampere}$, tem-se $m=\SI{0.157}{\ampere\meter\squared}$ e, a
$\SI{1}{\meter}$, $B=\pot{3.14}{-8}\,\si{\tesla}$. O valor exato é
$\pot{3.10}{-8}$ --- a aproximação erra $1{,}5\%$ a dez raios de distância.\\
(c) \textbf{Comparação com o dipolo elétrico.} Ambos decaem com $1/r^{3}$, e as
expressões são formalmente análogas ($p=2qa$ lá, $m=IA$ aqui).\\
\textbf{Mas há uma diferença fundamental.} O dipolo elétrico é feito de duas
cargas que \emph{podem} ser separadas --- e então cada uma produz campo
$1/r^{2}$. Já o dipolo magnético \textbf{não pode} ser separado: não existem
monopolos magnéticos.\\
\textbf{Consequência:} o dipolo é o termo \emph{dominante} de qualquer fonte
magnética estacionária. Não há termo $1/r^{2}$ em magnetostática --- e é por isso
que $\nabla\cdot\vec B=0$ figura entre as equações de Maxwell, sem análogo
elétrico.}
\end{questao}

\begin{questao}[4]{}
\textbf{(Blindagem magnética.)} Compare a blindagem de campos magnéticos com a
de campos elétricos.
\begin{itensq}
\item Explique por que uma gaiola metálica não blinda campo magnético estático.
\item Descreva os dois mecanismos disponíveis.
\item Discuta a dependência com a frequência.
\end{itensq}
\resposta{Não há "cargas magnéticas" a redistribuir; usa-se desvio por material
de alto $\mu_r$ ou correntes induzidas}
\resolucao{(a) \textbf{Por que a gaiola falha.} A blindagem elétrica funciona
porque o condutor possui \emph{cargas livres} que se redistribuem até anular o
campo interno. Não existe equivalente magnético: \textbf{não há monopolos
magnéticos livres} para migrar.\\
Um campo magnético \emph{estático} atravessa o cobre e o alumínio praticamente
sem atenuação.\\
(b) \textbf{Dois mecanismos disponíveis.}
\begin{itemize}
\item \textbf{Desvio de fluxo} (campos estáticos ou lentos). Um invólucro de
material de alto $\mu_r$ --- \emph{mu-metal}, com $\mu_r$ de dezenas de
milhares --- oferece caminho de baixa relutância. O fluxo prefere circular pelo
material e "contorna" a região protegida. Não há cancelamento: há \emph{desvio}.
\item \textbf{Correntes induzidas} (campos variáveis). Pela lei de Lenz, um
campo alternado induz correntes em um invólucro condutor, e essas correntes
geram campo oposto. A atenuação cresce com a frequência e com a condutividade.
\end{itemize}
(c) \textbf{Dependência com a frequência.}
\begin{center}
\begin{tabular}{lll}
\hline
Faixa & Mecanismo eficaz & Material\\
\hline
Estático (CC) & desvio de fluxo & mu-metal\\
Baixa (dezenas de Hz) & desvio, pouca indução & mu-metal\\
Alta (kHz e acima) & correntes induzidas & cobre, alumínio\\
\hline
\end{tabular}
\end{center}
\textbf{Consequência prática:} blindar campo elétrico é fácil e barato --- basta
qualquer metal aterrado. Blindar campo magnético estático é caro, exige material
especial, e a atenuação típica fica em uma ou duas ordens de grandeza, contra
muitas ordens no caso elétrico.\\
É por isso que salas de ressonância magnética têm blindagem magnética de
toneladas, enquanto a blindagem elétrica de um cabo é uma malha fina.}
\end{questao}

\begin{questao}[4]{}
\textbf{(Mapeamento experimental.)} Elabore um procedimento para mapear o campo
magnético de uma bobina usando sensor Hall.
\begin{itensq}
\item Descreva o procedimento e o tratamento do \emph{offset}.
\item Indique como obter as três componentes do campo.
\item Estabeleça o critério de validação dos dados.
\end{itensq}
\resposta{Medir com inversão de orientação para cancelar \emph{offset};
três orientações ortogonais para as componentes}
\resolucao{(a) \textbf{Procedimento.}
\begin{enumerate}
\item \textbf{Caracterizar o zero:} com a bobina \emph{desligada}, registrar a
leitura em cada ponto. Isso captura tanto o \emph{offset} do sensor quanto o
campo ambiente (terrestre e de equipamentos próximos).
\item \textbf{Medir com a bobina ligada} nos mesmos pontos.
\item \textbf{Subtrair} ponto a ponto. O resultado é o campo da bobina,
livre de \emph{offset} e de campo ambiente.
\end{enumerate}
\textbf{Alternativa mais robusta:} inverter o sentido da corrente na bobina e
tomar a \emph{semidiferença} das duas leituras. Isso cancela tudo o que não
depende da corrente.\\
(b) \textbf{Três componentes.} O sensor Hall responde apenas à componente
perpendicular à sua face. Para obter $\vec B$ completo, meça em três orientações
ortogonais no mesmo ponto:
\[
|\vec B|=\sqrt{B_x^{2}+B_y^{2}+B_z^{2}}.
\]
Um sensor triaxial faz isso simultaneamente e evita erros de reposicionamento.\\
(c) \textbf{Critérios de validação.}
\begin{itemize}
\item \textbf{Linearidade com a corrente:} $B$ deve ser proporcional a $I$.
Meça em três correntes e verifique --- desvio indica saturação de núcleo ou
não linearidade do sensor.
\item \textbf{Simetria:} pontos simétricos em relação ao eixo devem dar valores
iguais. Discrepância revela desalinhamento do sensor ou da bobina.
\item \textbf{Valor central:} deve concordar com $\mu_0nI$ dentro da incerteza,
descontada a correção de comprimento finito.
\item \textbf{Reprodutibilidade:} repita o primeiro ponto ao final. Diferença
indica deriva térmica do sensor.
\end{itemize}
\textbf{Registre sempre a temperatura} --- a sensibilidade Hall varia
tipicamente algumas centenas de ppm por grau.}
\end{questao}

\begin{questao}[4]{}
\textbf{(Limites de eletroímãs.)} Analise o que limita, na prática, o campo
magnético alcançável.
\begin{itensq}
\item Determine o limite imposto pela saturação do núcleo.
\item Determine o limite térmico de bobinas de cobre.
\item Descreva as soluções empregadas para superá-los.
\end{itensq}
\resposta{Núcleos saturam entre $1{,}5$ e $\SI{2}{\tesla}$; acima disso é
preciso bobina sem núcleo, com corrente muito alta}
\resolucao{(a) \textbf{Saturação.} Materiais ferromagnéticos alinham seus
domínios completamente em torno de $1{,}5$ a $\SI{2}{\tesla}$ (ferro-silício,
$\approx\SI{2}{\tesla}$; ferrites, bem menos). Acima disso, $\mu_r$ cai para
próximo de $1$ e o núcleo deixa de contribuir --- \textbf{acrescentar corrente
passa a render tão pouco quanto no ar}.\\
Este é um limite \emph{de material}, não de projeto: nenhuma geometria o
contorna.\\
(b) \textbf{Limite térmico.} Sem núcleo, $B=\mu_0nI$ exige corrente enorme. A
potência dissipada vai com $I^{2}$, e a densidade de corrente admissível em
cobre refrigerado a água fica em torno de
$\SI{20}{\ampere\per\milli\meter\squared}$.\\
Eletroímãs resistivos de laboratório atingem $1$ a $\SI{2}{\tesla}$ consumindo
\emph{megawatts} --- e boa parte do custo operacional é a refrigeração.\\
(c) \textbf{Soluções empregadas.}
\begin{itemize}
\item \textbf{Supercondutores:} resistência nula elimina o limite térmico.
Bobinas de NbTi e $\mathrm{Nb_3Sn}$ chegam a $\SI{20}{\tesla}$ e são a base de
equipamentos de ressonância magnética e de aceleradores. O custo passa a ser a
criogenia.
\item \textbf{Campos pulsados:} correntes altíssimas por milissegundos, antes
que o calor se acumule. Alcançam $\SI{100}{\tesla}$, com a bobina
frequentemente destruída no processo.
\item \textbf{Ímãs híbridos:} bobina supercondutora externa somada a uma
resistiva interna --- combinação que detém os recordes de campo contínuo, acima
de $\SI{45}{\tesla}$.
\end{itemize}
\textbf{Limite último:} acima de algumas dezenas de tesla, a \emph{pressão
magnética} $B^{2}/2\mu_0$ supera a resistência mecânica de qualquer material. A
$\SI{100}{\tesla}$ ela vale cerca de $\SI{4}{\giga\pascal}$ --- a bobina
literalmente se rompe.}
\end{questao}

\begin{questao}[4]{}
\textbf{(Superposição e simetria.)} Um condutor cilíndrico de raio $a$ possui
uma cavidade cilíndrica de raio $b$, paralela ao eixo e deslocada de $d$.
A corrente $I$ distribui-se uniformemente na seção restante.
\begin{itensq}
\item Formule a estratégia de solução por superposição.
\item Determine o campo dentro da cavidade.
\item Interprete o resultado.
\end{itensq}
\resposta{O campo na cavidade é \textbf{uniforme}, de módulo
$\mu_0Jd/2$}
\resolucao{(a) \textbf{Estratégia.} O condutor furado equivale à
\emph{superposição} de:
\begin{itemize}
\item um cilindro \textbf{maciço} de raio $a$ com densidade de corrente
$+J$; \emph{mais}
\item um cilindro de raio $b$, na posição da cavidade, com densidade
$-J$.
\end{itemize}
A soma reproduz corrente $J$ na região sólida e zero na cavidade --- exatamente
a situação real.\\
(b) Dentro de um cilindro maciço de densidade $J$, a lei de Ampère dá
$B=\dfrac{\mu_0 J r}{2}$, ou vetorialmente
$\vec B=\dfrac{\mu_0}{2}\,\vec J\times\vec r$, com $\vec r$ medido a partir do
\emph{eixo daquele} cilindro.\\
Somando as duas contribuições em um ponto da cavidade, com $\vec r_1$ e
$\vec r_2$ medidos a partir de cada eixo:
\[
\vec B=\frac{\mu_0}{2}\vec J\times\vec r_1
-\frac{\mu_0}{2}\vec J\times\vec r_2
=\frac{\mu_0}{2}\vec J\times(\vec r_1-\vec r_2)
=\frac{\mu_0}{2}\vec J\times\vec d.
\]
\textbf{O vetor $\vec d$ é constante} --- ele liga os dois eixos e não depende do
ponto.\\
(c) \textbf{Interpretação: o campo na cavidade é rigorosamente uniforme}, de
módulo $\mu_0Jd/2$ e direção perpendicular a $\vec d$.\\
Isso é notável: uma geometria sem simetria evidente produz, na região vazia, o
campo mais simples possível.\\
\textbf{Aplicação:} o mesmo resultado, com $\vec J$ substituído por densidade de
carga, dá campo elétrico uniforme na cavidade de uma esfera carregada --- e é
usado no projeto de aceleradores e de ímãs de deflexão, em que se busca campo
uniforme em uma região de acesso.}
\end{questao}

% END BANCO COMPLEMENTAR

\gabarito[12.2 Campo magnético]

% =====================================================================
\subsection{Fluxo magnético}
% =====================================================================

\blocofixacao

\begin{questao}[1]{}
O fluxo magnético através de uma espira é máximo quando o plano da espira está:
\begin{alternativas}
\item paralelo ao campo.
\item perpendicular ao campo (campo atravessando a espira frontalmente).
\item inclinado a $\SI{45}{\degree}$ com o campo.
\item em qualquer orientação (o fluxo não depende do ângulo).
\item girando rapidamente.
\end{alternativas}
\resposta{(b)}
\resolucao{Como $\Phi = B\,A\cos\theta$, sendo $\theta$ o ângulo entre o campo e a
\emph{normal} à espira, o fluxo é máximo quando $\theta = 0$ --- ou seja, quando o
campo atravessa a espira perpendicularmente ao seu plano. Se o plano fica paralelo
ao campo, $\theta = \SI{90}{\degree}$ e o fluxo é nulo.}
\end{questao}

\begin{questao}[1]{}
Uma espira de área $\SI{20}{\centi\meter\squared}$ é atravessada
perpendicularmente por um campo uniforme de $\SI{0.5}{\tesla}$. Calcule o fluxo
magnético.
\resposta{$\Phi = \SI{1}{\milli\weber}$}
\resolucao{Com $\theta = 0$ ($\cos\theta = 1$) e
$A = \SI{20}{\centi\meter\squared} = \pot{20}{-4}\,\si{\meter\squared}$:
\[
\Phi = B\,A = 0{,}5\cdot\pot{20}{-4} = \pot{1}{-3}\,\si{\weber}=\SI{1}{\milli\weber}.
\]}
\end{questao}

\blocoaplicacao

\begin{questao}[2]{}
Uma espira circular de raio $\SI{8}{\centi\meter}$ está imersa em um campo
uniforme de $\SI{0.3}{\tesla}$, perpendicular ao seu plano.
\begin{itensq}
\item Calcule o fluxo magnético através da espira.
\item Se o campo variar linearmente de $\SI{0.3}{}$ para $\SI{0.1}{\tesla}$ em
      $\SI{0.5}{\second}$, qual é a fem induzida?
\end{itensq}
\resposta{(a) $\Phi \approx \SI{6.03}{\milli\weber}$; (b) $\varepsilon \approx \SI{8.04}{\milli\volt}$}
\resolucao{(a) A área é $A = \pi R^2 = \pi(0{,}08)^2 \approx \SI{0.0201}{\meter\squared}$:
\[
\Phi = B\,A = 0{,}3\cdot0{,}0201 \approx \SI{6.03}{\milli\weber}.
\]
(b) A fem é o módulo da taxa de variação do fluxo. Como só $B$ muda:
\[
\varepsilon = \left|\frac{\Delta\Phi}{\Delta t}\right|
= A\,\frac{|\Delta B|}{\Delta t}
= 0{,}0201\cdot\frac{0{,}2}{0{,}5}\approx\SI{8.04}{\milli\volt}.
\]}
\end{questao}

\begin{questao}[2]{Apostila original revisada}
Um campo magnético uniforme de módulo $B=\SI{1e-4}{\tesla}$ atravessa uma espira de
área $A=\SI{1e-5}{\metre\squared}$. O campo forma $\SI{45}{\degree}$ com a normal
ao plano da espira. Calcule o fluxo magnético.
\resposta{$\Phi=\SI{7.07e-10}{\weber}$}
\resolucao{$\Phi=BA\cos\theta=(10^{-4})(10^{-5})\cos45^\circ
=7{,}07\times10^{-10}\,\si{\weber}$.}
\end{questao}

\blocoaprofundamento

\begin{questao}[3]{}
O gráfico mostra a variação do fluxo magnético $\Phi(t)$ através de uma espira.
Determine a fem induzida em cada trecho e esboce mentalmente a forma de onda da
fem.

\circuitoqq{%
\begin{tikzpicture}
 \begin{axis}[width=10cm,height=5cm,xlabel={$t$ (s)},ylabel={$\Phi$ (mWb)},
   xmin=0,xmax=6,ymin=0,ymax=5,xtick={0,1,2,3,4,5,6},ytick={0,2,4},
   grid=both,axis lines=left]
   \addplot[Icol,very thick] coordinates {(0,0)(2,4)(4,4)(6,0)};
 \end{axis}
\end{tikzpicture}}{Fluxo magnético em função do tempo.}

\resposta{$0$--$2\,$s: $\varepsilon=\SI{2}{\milli\volt}$; $2$--$4\,$s:
$\varepsilon=0$; $4$--$6\,$s: $\varepsilon=\SI{-2}{\milli\volt}$}
\resolucao{A fem é a \emph{inclinação} do gráfico $\Phi \times t$ (com sinal
trocado, por Lenz):
\begin{itemize}[leftmargin=1.6em,itemsep=1pt,topsep=2pt]
\item \textbf{$0$ a $\SI{2}{\second}$:} $\Phi$ sobe de $0$ a
      \SI{4}{\milli\weber}, inclinação $\dfrac{4}{2} = \SI{2}{\milli\weber\per\second}$,
      logo $|\varepsilon| = \SI{2}{\milli\volt}$.
\item \textbf{$2$ a $\SI{4}{\second}$:} $\Phi$ constante, inclinação nula,
      $\varepsilon = 0$.
\item \textbf{$4$ a $\SI{6}{\second}$:} $\Phi$ desce de \SI{4}{} a $0$, inclinação
      $-\SI{2}{\milli\weber\per\second}$, $\varepsilon = \SI{-2}{\milli\volt}$
      (sentido oposto).
\end{itemize}
A fem só existe quando há \emph{variação} de fluxo, e seu sinal se inverte entre a
subida e a descida --- exatamente o princípio de um gerador.}
\end{questao}

\blocodesafio

\begin{questao}[4]{}
Uma espira retangular de área $A = \SI{200}{\centi\meter\squared}$ gira com
velocidade angular constante $\omega$ em um campo magnético uniforme
$B = \SI{0.4}{\tesla}$, completando $\SI{1800}{}$ rotações por minuto.
\begin{itensq}
\item Escreva a expressão do fluxo $\Phi(t)$ e calcule seu valor máximo.
\item Obtenha a fem induzida $\varepsilon(t)$ e seu valor de pico.
\item Determine a frequência da tensão gerada e explique por que a saída é
      \emph{senoidal}.
\item Se a espira tivesse $N = 100$ voltas, o que mudaria?
\end{itensq}
\resposta{(a) $\Phi_{\max}=\SI{8}{\milli\weber}$;
(b) $\varepsilon_{\text{pico}}\approx\SI{1.51}{\volt}$; (c) \SI{30}{\hertz};
(d) tudo $\times100$}
\resolucao{\textbf{(a)} Com a espira girando, o ângulo entre a normal e o campo
varia: $\theta = \omega t$. Logo
\[
\Phi(t)=B\,A\cos(\omega t),
\qquad
\Phi_{\max}=B\,A=0{,}4\cdot\pot{200}{-4}=\SI{8}{\milli\weber}.
\]
\textbf{(b)} A velocidade angular é
\[
\omega = \frac{2\pi\cdot1800}{60}=60\pi\approx\SI{188.5}{\radian\per\second}.
\]
Pela lei de Faraday,
\[
\varepsilon(t)=-\frac{\mathrm{d}\Phi}{\mathrm{d}t}
= B A\,\omega\sin(\omega t),
\qquad
\varepsilon_{\text{pico}}=BA\omega
=(\pot{8}{-3})(188{,}5)\approx\SI{1.51}{\volt}.
\]
\textbf{(c)} A frequência é
$f = \dfrac{1800}{60} = \SI{30}{\hertz}$ (uma volta $=$ um ciclo completo).\\
\textbf{Por que senoidal?} Porque o \emph{fluxo} varia cossenoidalmente com o
ângulo (projeção da área sobre a direção do campo), e a fem é sua
\emph{derivada}. A derivada do cosseno é o seno --- daí a defasagem de
$\SI{90}{\degree}$ entre fluxo e tensão. Este é literalmente o princípio de
funcionamento do \textbf{alternador}: a forma de onda senoidal da rede elétrica
nasce da geometria circular do movimento.

\textbf{(d)} Com $N$ espiras, todas concatenam o mesmo fluxo e suas fem se somam
em série:
\[
\varepsilon_{\text{pico}} = N B A \omega \approx 100\cdot1{,}51 = \SI{151}{\volt}.
\]
O fluxo \emph{por espira} não muda; o que muda é o \emph{fluxo concatenado}
$N\Phi$. É assim que geradores reais atingem centenas de volts com campos
modestos: aumentando $N$, $B$, $A$ ou $\omega$ --- os quatro parâmetros de
projeto de uma máquina elétrica.}
\end{questao}

% BEGIN BANCO COMPLEMENTAR Fluxo magnético
% =====================================================================
%  Banco complementar --- Fluxo Magnetico  (REVISADO)
%  Substitui a versao gerada por template (6 clones em N1, 10 em N2,
%  11 em N3 e 7 questoes conceituais marcadas como N4).
%  O cap11 ja cobre: condicao de fluxo maximo (MC), espira de 20 cm2
%  perpendicular a 0,5 T, espira circular de R=8 cm em 0,3 T, campo de
%  1e-4 T em area de 1e-5 m2, grafico Phi(t) com fem por trecho e --
%  como N4 -- a ESPIRA RETANGULAR GIRANDO com velocidade angular
%  constante. Esta ultima motivou a remocao da N4 antiga sobre fluxo em
%  espira girante, que a duplicava.
%  Este banco explora: campo nao uniforme por integracao, fluxo de um
%  fio reto atraves de espira retangular, CIRCUITO MAGNETICO com
%  relutancia e entreferro, dispersao e coeficiente de acoplamento,
%  independencia da superficie e nucleo de secao variavel.
%  A INDUCAO propriamente dita tem banco proprio (Lenz); aqui a fem so
%  aparece onde o fluxo e o objeto de interesse.
%  Distribuicao mantida: 6 N1 + 10 N2 + 11 N3 + 7 N4 = 34
% =====================================================================

\subsubsection*{Banco complementar --- Fluxo magnético}

\begin{atencao}[Organização do banco]
O fluxo é $\Phi=\int\vec B\cdot\mathrm{d}\vec A$, que se reduz a
$BA\cos\theta$ apenas quando o campo é \textbf{uniforme} e a superfície
\textbf{plana} --- $\theta$ sendo o ângulo entre $\vec B$ e a \emph{normal}, e
não entre $\vec B$ e o plano. Duas propriedades estruturam o tópico: o fluxo por
qualquer superfície \emph{fechada} é nulo, e o fluxo por superfícies de mesma
\emph{borda} é o mesmo. O N3 trata de campos não uniformes, circuitos magnéticos
e dispersão; o N4 exige demonstrações e projeto.
\end{atencao}

\blocofixacao

\begin{questao}[1]{}
Uma superfície plana de $\SI{0.01}{\meter\squared}$ está em campo uniforme de
$\SI{0.1}{\tesla}$, com a normal formando $\ang{60}$ com o campo. Determine o
fluxo.
\resposta{$\Phi=\SI{0.5}{\milli\weber}$}
\resolucao{
\[
\Phi=BA\cos\theta=(0{,}1)(0{,}01)\cos\ang{60}
=(0{,}1)(0{,}01)(0{,}5)=\pot{5}{-4}\,\si{\weber}.
\]
\textbf{Atenção ao ângulo:} $\theta$ é medido entre $\vec B$ e a \textbf{normal}
à superfície. Se o enunciado der o ângulo com o \emph{plano}, use o complemento
--- é a troca de $\cos$ por $\sin$ que mais gera erro neste tópico.}
\end{questao}

\begin{questao}[1]{}
Determine o fluxo quando o campo é \textbf{paralelo} ao plano da superfície.
\resposta{$\Phi=0$}
\resolucao{Campo paralelo ao plano significa campo \emph{perpendicular} à
normal, logo $\theta=\ang{90}$ e $\cos\theta=0$:
\[
\Phi=0.
\]
\textbf{Interpretação geométrica:} nenhuma linha de campo \emph{atravessa} a
superfície --- elas passam rentes a ela.\\
\textbf{Casos-limite úteis:} $\Phi$ é máximo com $\vec B$ perpendicular ao plano
(paralelo à normal) e nulo com $\vec B$ contido no plano.}
\end{questao}

\begin{questao}[1]{}
Um fluxo de $\SI{2}{\milli\weber}$ atravessa perpendicularmente uma superfície
em campo de $\SI{0.4}{\tesla}$. Determine a área.
\resposta{$A=\SI{50}{\centi\meter\squared}$}
\resolucao{
\[
A=\frac{\Phi}{B}=\frac{\pot{2}{-3}}{0{,}4}=\pot{5}{-3}\,\si{\meter\squared}
=\SI{50}{\centi\meter\squared}.
\]
\textbf{Lembre da conversão:}
$\SI{1}{\meter\squared}=\SI{e4}{\centi\meter\squared}$, pois a área é o
\emph{quadrado} de um comprimento.}
\end{questao}

\begin{questao}[1]{}
Uma espira de $\SI{30}{\centi\meter\squared}$ é atravessada
perpendicularmente por fluxo de $\SI{1.5}{\milli\weber}$. Determine a densidade
de fluxo.
\resposta{$B=\SI{0.5}{\tesla}$}
\resolucao{
\[
B=\frac{\Phi}{A}=\frac{\pot{1.5}{-3}}{\pot{3}{-3}}=\SI{0.5}{\tesla}.
\]
\textbf{Nomenclatura:} $B$ é chamado \emph{densidade de fluxo} justamente por
ser fluxo por unidade de área --- e $\SI{1}{\tesla}=\SI{1}{\weber\per\meter\squared}$.\\
As duas grandezas descrevem a mesma realidade em escalas diferentes: $B$ é
local, $\Phi$ é integral.}
\end{questao}

\begin{questao}[1]{}
O fluxo magnético líquido através de qualquer superfície \textbf{fechada} é:
\begin{alternativas}[2]
\item proporcional à corrente envolvida
\item sempre nulo
\item proporcional ao volume
\item máximo quando a superfície é esférica
\end{alternativas}
\resposta{(b)}
\resolucao{As linhas de campo magnético são \textbf{fechadas} --- não têm início
nem fim. Toda linha que entra em uma superfície fechada necessariamente sai dela,
e as contribuições se cancelam:
\[
\oint\vec B\cdot\mathrm{d}\vec A=0.
\]
\textbf{Contraste com o campo elétrico:} lá, a lei de Gauss dá
$\oint\vec E\cdot\mathrm{d}\vec A=Q_{env}/\varepsilon_0$ --- diferente de zero
sempre que houver carga envolvida.\\
\textbf{A diferença tem uma causa única:} existem cargas elétricas isoladas
(monopolos elétricos), mas \textbf{não existem monopolos magnéticos}. Esta é uma
das equações de Maxwell, e sua forma diferencial é $\nabla\cdot\vec B=0$.}
\end{questao}

\begin{questao}[1]{}
Converta $\SI{5}{\milli\weber}$ para maxwells.
\dados{$\SI{1}{\weber}=10^{8}\,\si{\maxwell}$}
\resposta{$\pot{5}{5}\,\si{\maxwell}$}
\resolucao{
\[
\Phi=\pot{5}{-3}\times10^{8}=\pot{5}{5}\,\si{\maxwell}.
\]
\textbf{Origem da unidade:} o maxwell pertence ao sistema CGS, e vale
$\SI{1}{\gauss\centi\meter\squared}$. A conversão decorre diretamente de
$\SI{1}{\tesla}=\SI{e4}{\gauss}$ e
$\SI{1}{\meter\squared}=\SI{e4}{\centi\meter\squared}$:
\[
10^{4}\times10^{4}=10^{8}.
\]
Embora obsoleto no SI, o maxwell ainda aparece em catálogos de ímãs e em
literatura mais antiga.}
\end{questao}

\blocoaplicacao

\begin{questao}[2]{}
Uma bobina de $50$ espiras é atravessada por fluxo de
$\SI{1}{\milli\weber}$ por espira.
\begin{itensq}
\item Determine o fluxo concatenado.
\item Determine a indutância, se a corrente for $\SI{2}{\ampere}$.
\end{itensq}
\resposta{$\lambda=\SI{50}{\milli\weber}$-espira; $L=\SI{25}{\milli\henry}$}
\resolucao{(a)
\[
\lambda=N\Phi=50(\pot{1}{-3})=\pot{5}{-2}\,\si{\weber}\text{-espira}.
\]
(b)
\[
L=\frac{\lambda}{I}=\frac{0{,}05}{2}=\SI{25}{\milli\henry}.
\]
\textbf{Distinção essencial.} $\Phi$ é o fluxo através de \emph{uma} espira;
$\lambda=N\Phi$ é o \emph{fluxo concatenado}, que conta cada espira
separadamente. É $\lambda$, e não $\Phi$, que entra na definição de indutância e
na lei de Faraday para bobinas ($\varepsilon=-\mathrm{d}\lambda/\mathrm{d}t$).}
\end{questao}

\begin{questao}[2]{}
Uma espira de $\SI{200}{\centi\meter\squared}$ está em campo de
$\SI{0.25}{\tesla}$, com a normal a $\ang{60}$ do campo. Determine o fluxo e o
valor que ele teria na posição de fluxo máximo.
\resposta{$\Phi=\SI{2.5}{\milli\weber}$; máximo de $\SI{5}{\milli\weber}$}
\resolucao{
\[
\Phi=BA\cos\ang{60}=(0{,}25)(\pot{2}{-2})(0{,}5)=\pot{2.5}{-3}\,\si{\weber};
\]
\[
\Phi_{\max}=BA=\pot{5}{-3}\,\si{\weber}\quad(\theta=0).
\]
\textbf{Aplicação:} girar a espira entre essas duas posições varia o fluxo de
$\SI{2.5}{\milli\weber}$. Feito em $\SI{10}{\milli\second}$, isso induziria
$\varepsilon=\SI{0.25}{\volt}$ por espira --- princípio do gerador.}
\end{questao}

\begin{questao}[2]{}
Um solenoide de $\SI{1000}{}$ espiras por metro conduz $\SI{2}{\ampere}$ e tem
seção de $\SI{10}{\centi\meter\squared}$. Determine o fluxo por espira.
\resposta{$\Phi=\SI{2.51}{\micro\weber}$}
\resolucao{
\[
B=\mu_0nI=(4\pi\times10^{-7})(1000)(2)=\pot{2.51}{-3}\,\si{\tesla};
\]
\[
\Phi=BA=(\pot{2.51}{-3})(\pot{1}{-3})=\pot{2.51}{-6}\,\si{\weber}.
\]
\textbf{Ordem de grandeza:} apenas $\SI{2.5}{\micro\weber}$ --- fluxos em
núcleo de ar são pequenos. Com núcleo ferromagnético de
$\mu_r=1000$, o mesmo solenoide produziria $\SI{2.5}{\milli\weber}$.\\
\textbf{É por isso que máquinas elétricas usam núcleo:} sem ele, o fluxo (e
portanto a f.e.m. induzida e o torque) seria mil vezes menor.}
\end{questao}

\begin{questao}[2]{}
O fluxo através de uma bobina de $100$ espiras varia linearmente de zero a
$\SI{4}{\milli\weber}$ em $\SI{20}{\milli\second}$. Determine a f.e.m.
induzida.
\resposta{$\varepsilon=\SI{20}{\volt}$}
\resolucao{
\[
|\varepsilon|=N\frac{\Delta\Phi}{\Delta t}
=100\,\frac{\pot{4}{-3}}{\pot{2}{-2}}=\SI{20}{\volt}.
\]
\textbf{Como a variação é linear}, a f.e.m. é \emph{constante} durante todo o
intervalo. Se o fluxo variasse de outra forma, seria preciso derivar ponto a
ponto.\\
\textbf{Note o papel de $N$:} cem espiras multiplicam a f.e.m. por cem, embora o
fluxo \emph{por espira} seja o mesmo. É o fluxo concatenado que importa.}
\end{questao}

\begin{questao}[2]{}
Duas superfícies diferentes apoiam-se na \textbf{mesma} espira: uma plana e
outra em forma de calota. Compare o fluxo através delas.
\resposta{São iguais}
\resolucao{O fluxo através de qualquer superfície depende apenas de sua
\textbf{borda}.\\
\textbf{Demonstração:} as duas superfícies, juntas, formam uma superfície
\emph{fechada}. Como
$\oint\vec B\cdot\mathrm{d}\vec A=0$, o fluxo que entra por uma sai pela outra
--- e portanto os fluxos através delas são iguais (atentando para a orientação
das normais).\\
\textbf{Consequência prática:} na lei de Faraday, pode-se escolher
\emph{qualquer} superfície apoiada no circuito. Escolha a que tornar a
integral mais simples --- geralmente a plana, mas às vezes uma superfície
curva evita regiões de campo complicado.\\
\textbf{Contraste:} isso \emph{não} vale para o fluxo elétrico, que depende da
carga envolvida e portanto da superfície escolhida.}
\end{questao}

\begin{questao}[2]{}
Uma espira quadrada de lado $\SI{20}{\centi\meter}$ é girada de
$\ang{0}$ a $\ang{90}$ (normal em relação ao campo) em campo de
$\SI{0.6}{\tesla}$. Determine a variação de fluxo.
\resposta{$\Delta\Phi=\SI{-24}{\milli\weber}$}
\resolucao{
\[
\Phi_i=BA=(0{,}6)(0{,}04)=\SI{24}{\milli\weber};
\qquad
\Phi_f=BA\cos\ang{90}=0.
\]
\[
\Delta\Phi=0-0{,}024=\SI{-24}{\milli\weber}.
\]
\textbf{Duas formas de variar o fluxo aparecem no tópico:} mudar $B$, mudar
$A$, ou mudar a \emph{orientação}. Esta última é a base do gerador rotativo, em
que $\theta=\omega t$ e o fluxo varia senoidalmente.\\
\textbf{Sinal:} o fluxo diminuiu. Pela lei de Lenz, a corrente induzida
circularia no sentido de \emph{manter} o fluxo original.}
\end{questao}

\begin{questao}[2]{}
Um núcleo toroidal tem seção de $\SI{4}{\centi\meter\squared}$ e conduz fluxo
de $\SI{0.6}{\milli\weber}$. Determine a densidade de fluxo e avalie a
proximidade da saturação.
\resposta{$B=\SI{1.5}{\tesla}$ --- próximo da saturação típica}
\resolucao{
\[
B=\frac{\Phi}{A}=\frac{\pot{6}{-4}}{\pot{4}{-4}}=\SI{1.5}{\tesla}.
\]
\textbf{Avaliação:} aços-silício saturam entre $1{,}5$ e $\SI{2}{\tesla}$;
ferrites, bem antes (tipicamente $0{,}3$ a $\SI{0.5}{\tesla}$).\\
Portanto, para aço-silício este núcleo opera \textbf{no limite}. Qualquer
aumento de corrente levaria à saturação, com queda abrupta de $\mu_r$ e
distorção severa da forma de onda.\\
\textbf{Prática de projeto:} dimensione para $B_{\max}$ entre $60$ e
$80\%$ da saturação, deixando margem para transitórios e para a variação com a
temperatura.}
\end{questao}

\begin{questao}[2]{}
Uma bobina de $200$ espiras e área $\SI{15}{\centi\meter\squared}$ está em
campo de $\SI{0.4}{\tesla}$ perpendicular. Determine o fluxo concatenado.
\resposta{$\lambda=\SI{120}{\milli\weber}$-espira}
\resolucao{
\[
\Phi=BA=(0{,}4)(\pot{1.5}{-3})=\pot{6}{-4}\,\si{\weber};
\]
\[
\lambda=N\Phi=200(\pot{6}{-4})=\pot{0.12}{}\,\si{\weber}\text{-espira}.
\]
\textbf{Se essa bobina fosse retirada do campo em $\SI{50}{\milli\second}$:}
\[
|\varepsilon|=\frac{\Delta\lambda}{\Delta t}=\frac{0{,}12}{0{,}05}
=\SI{2.4}{\volt}.
\]
É esse o princípio do \emph{fluxímetro}: integrando a f.e.m. medida, recupera-se
$\Delta\lambda$ e, com $N$ e $A$ conhecidos, o campo.}
\end{questao}

\begin{questao}[2]{}
Uma espira está totalmente fora de uma região de campo uniforme e é empurrada
até ficar totalmente dentro. Descreva a variação do fluxo.
\resposta{Cresce de zero a $BA$, com taxa constante enquanto a espira atravessa
a fronteira}
\resolucao{\textbf{Três fases:}
\begin{itemize}
\item \textbf{Totalmente fora:} $\Phi=0$ e $\mathrm{d}\Phi/\mathrm{d}t=0$ ---
nenhuma f.e.m.
\item \textbf{Atravessando a fronteira:} a área \emph{dentro} do campo cresce
linearmente com o deslocamento. Se a velocidade for constante,
$\Phi$ cresce linearmente e a f.e.m. é \textbf{constante}.
\item \textbf{Totalmente dentro:} $\Phi=BA$ e volta a ser constante ---
nenhuma f.e.m., mesmo que a espira continue se movendo.
\end{itemize}
\textbf{Ponto conceitual importante:} não é o \emph{movimento} que induz f.e.m.,
e sim a \textbf{variação de fluxo}. Uma espira movendo-se rapidamente dentro de
um campo uniforme não gera nada.\\
Com largura $\ell$ e velocidade $v$, a f.e.m. na fase intermediária vale
$\varepsilon=B\ell v$.}
\end{questao}

\begin{questao}[2]{}
Explique a diferença entre fluxo magnético e densidade de fluxo, e indique
quando cada um é a grandeza adequada.
\resposta{$\Phi$ é integral (weber); $B$ é local (tesla). $\Phi$ governa a
indução; $B$ governa forças e saturação}
\resolucao{\textbf{Densidade de fluxo $\vec B$:} grandeza \emph{local} e
vetorial, medida em tesla. Ela é o que determina:
\begin{itemize}
\item a força sobre cargas e condutores ($\vec F=q\vec v\times\vec B$);
\item a saturação do material ($B$ próximo de $B_{sat}$);
\item o que um sensor Hall mede.
\end{itemize}
\textbf{Fluxo $\Phi$:} grandeza \emph{integral} e escalar, medida em weber.
Ela é o que determina:
\begin{itemize}
\item a f.e.m. induzida ($\varepsilon=-\mathrm{d}\Phi/\mathrm{d}t$);
\item a indutância ($L=\lambda/I$);
\item o acoplamento entre bobinas.
\end{itemize}
\textbf{Relação:} $\Phi=\int\vec B\cdot\mathrm{d}\vec A$, e no caso uniforme
$B=\Phi/A$.\\
\textbf{Erro comum:} falar em "fluxo" quando se quer dizer "campo". Um ímã
pequeno e um grande podem ter o mesmo $B$ na superfície e fluxos totais muito
diferentes --- e é o fluxo que determina o desempenho em um circuito magnético.}
\end{questao}

\blocoaprofundamento

\begin{questao}[3]{}
Um campo não uniforme $B(x)=B_0\dfrac{x}{L}$, perpendicular ao plano, atravessa
uma superfície retangular de largura $L$ (na direção $x$) e altura $h$, com
$B_0=\SI{0.5}{\tesla}$, $L=\SI{20}{\centi\meter}$ e
$h=\SI{10}{\centi\meter}$.
\begin{itensq}
\item Determine o fluxo por integração.
\item Compare com as estimativas usando $B_0$ e usando o valor médio.
\end{itensq}
\resposta{$\Phi=\SI{5}{\milli\weber}$ --- igual à estimativa pelo valor médio}
\resolucao{(a) Dividindo em faixas verticais de largura $\mathrm{d}x$ e área
$h\,\mathrm{d}x$:
\[
\Phi=\int_0^{L}B_0\frac{x}{L}\,h\,\mathrm{d}x
=\frac{B_0h}{L}\left[\frac{x^{2}}{2}\right]_0^{L}
=\frac{B_0hL}{2}.
\]
\[
\Phi=\frac{(0{,}5)(0{,}10)(0{,}20)}{2}=\pot{5}{-3}\,\si{\weber}.
\]
(b)
\begin{center}
\begin{tabular}{lc}
\hline
Estimativa & $\Phi$\\
\hline
Usando $B_0$ (valor máximo) & $\SI{10}{\milli\weber}$\\
\textbf{Integração exata} & $\SI{5}{\milli\weber}$\\
Usando $\bar B=B_0/2$ & $\SI{5}{\milli\weber}$\\
\hline
\end{tabular}
\end{center}
\textbf{A média coincide com o exato} --- mas apenas porque $B$ varia
\emph{linearmente}. Para $B\propto x^{2}$, a integração daria $B_0hL/3$,
enquanto a média aritmética dos extremos daria $B_0hL/2$ --- erro de
$50\%$.\\
\textbf{Regra:} substituir o campo pela média só é legítimo em variação linear.
Fora disso, integre.}
\end{questao}

\begin{questao}[3]{}
Um fio retilíneo longo conduz $\SI{10}{\ampere}$. Uma espira retangular de
$\SI{20}{\centi\meter}$ de comprimento (paralelo ao fio) está no mesmo plano,
com os lados a $\SI{5}{\centi\meter}$ e $\SI{20}{\centi\meter}$ do fio.
\begin{itensq}
\item Deduza a expressão do fluxo.
\item Calcule seu valor.
\end{itensq}
\resposta{$\Phi=\dfrac{\mu_0Ia}{2\pi}\ln\dfrac{d+b}{d}=\SI{0.55}{\micro\weber}$}
\resolucao{(a) O campo do fio \textbf{não é uniforme} na região da espira:
$B(r)=\dfrac{\mu_0I}{2\pi r}$. Dividindo a espira em faixas de largura
$\mathrm{d}r$ e comprimento $a$:
\[
\Phi=\int_d^{d+b}\frac{\mu_0I}{2\pi r}\,a\,\mathrm{d}r
=\frac{\mu_0Ia}{2\pi}\ln\!\left(\frac{d+b}{d}\right).
\]
(b) Com $a=\SI{0.20}{\meter}$, $d=\SI{0.05}{\meter}$ e
$b=\SI{0.15}{\meter}$ (borda externa em $\SI{0.20}{\meter}$):
\[
\Phi=\pot{2}{-7}(10)(0{,}20)\ln 4
=(\pot{4}{-7})(1{,}386)=\pot{5.55}{-7}\,\si{\weber}.
\]
\textbf{Observe o logaritmo:} afastar a espira reduz o fluxo muito devagar.
Dobrar $d$ e $b$ juntos (mantendo a razão) \emph{não muda nada} --- o fluxo
depende apenas da razão $(d+b)/d$ e do comprimento $a$.\\
\textbf{Consequência:} a indutância mútua entre um fio e uma espira próxima é
notavelmente insensível à distância absoluta. É por isso que interferência
magnética de cabos de potência é difícil de eliminar apenas afastando-os.}
\end{questao}

\begin{questao}[3]{}
Um núcleo magnético de ferro tem comprimento médio $\SI{50}{\centi\meter}$,
seção de $\SI{4}{\centi\meter\squared}$ e $\mu_r=2000$, com enrolamento de
$\SI{1000}{}$ ampère-espiras.
\begin{itensq}
\item Determine a relutância do núcleo.
\item Determine o fluxo e a densidade de fluxo.
\item Avalie o resultado.
\end{itensq}
\resposta{$\mathcal{R}=\pot{4.97}{5}\,\si{\per\henry}$;
$\Phi=\SI{2}{\milli\weber}$; $B=\SI{5}{\tesla}$ --- \textbf{impossível}}
\resolucao{(a) Em analogia ao circuito elétrico, a \textbf{relutância} vale
\[
\mathcal{R}=\frac{\ell}{\mu_r\mu_0A}
=\frac{0{,}50}{(2000)(4\pi\times10^{-7})(\pot{4}{-4})}
=\pot{4.97}{5}\,\si{\per\henry}.
\]
(b) Com a força magnetomotriz $\mathcal{F}=NI=\SI{1000}{}$ ampère-espiras:
\[
\Phi=\frac{\mathcal{F}}{\mathcal{R}}=\frac{1000}{\pot{4.97}{5}}
=\pot{2.01}{-3}\,\si{\weber};
\qquad
B=\frac{\Phi}{A}=\SI{5.03}{\tesla}.
\]
(c) \textbf{O resultado é impossível.} Nenhum material ferromagnético comum
sustenta $\SI{5}{\tesla}$ --- todos saturam entre $1{,}5$ e
$\SI{2}{\tesla}$.\\
\textbf{O que realmente aconteceria:} conforme $B$ se aproxima de
$\SI{2}{\tesla}$, $\mu_r$ despenca de $2000$ para valores próximos de $1$. A
relutância dispara, e o fluxo estaciona em torno de
$B_{sat}A=\SI{0.8}{\milli\weber}$.\\
\textbf{Lição de método:} o modelo linear com $\mu_r$ constante é válido
\emph{apenas} abaixo da saturação. Calcular e depois \textbf{verificar $B$}
contra $B_{sat}$ é etapa obrigatória --- exatamente como se verifica a rigidez
dielétrica em problemas de campo elétrico.}
\end{questao}

\begin{questao}[3]{}
O mesmo núcleo da questão anterior recebe um \textbf{entreferro} de
$\SI{1}{\milli\meter}$.
\begin{itensq}
\item Determine a relutância do entreferro e a total.
\item Determine o novo fluxo e $B$.
\item Explique por que o entreferro domina.
\end{itensq}
\resposta{$\mathcal{R}_g=\pot{1.99}{6}$, quatro vezes a do ferro;
$B=\SI{1.0}{\tesla}$}
\resolucao{(a) No entreferro, $\mu_r=1$:
\[
\mathcal{R}_g=\frac{g}{\mu_0A}
=\frac{\pot{1}{-3}}{(4\pi\times10^{-7})(\pot{4}{-4})}
=\pot{1.99}{6}\,\si{\per\henry}.
\]
Em série com o ferro:
$\mathcal{R}_{tot}=\pot{4.97}{5}+\pot{1.99}{6}=\pot{2.49}{6}$.\\
(b)
\[
\Phi=\frac{1000}{\pot{2.49}{6}}=\pot{4.02}{-4}\,\si{\weber};
\qquad
B=\SI{1.005}{\tesla}.
\]
\textbf{Agora sim, valor fisicamente sustentável} --- e com margem para
saturação.\\
(c) \textbf{Por que o entreferro domina.} Ele tem apenas
$\SI{1}{\milli\meter}$ contra $\SI{500}{\milli\meter}$ de ferro --- é
$500$ vezes mais curto. Mas sua permeabilidade é $2000$ vezes menor. O saldo:
\[
\frac{\mathcal{R}_g}{\mathcal{R}_{fe}}=\frac{2000}{500}=4.
\]
\textbf{Um milímetro de ar vale quatro vezes meio metro de ferro.}\\
\textbf{Consequências de projeto.}
\begin{itemize}
\item O entreferro \emph{lineariza} o circuito magnético: como ele domina, a
relação $\Phi\times I$ torna-se quase reta, mesmo com o ferro não linear.
\item Ele permite armazenar muito mais energia (que se concentra no ar) --- daí
seu uso em indutores de conversores chaveados.
\item Em contrapartida, reduz a indutância e exige mais ampère-espiras.
\item \textbf{Entreferros parasitas} (juntas mal ajustadas de núcleos em E ou U)
degradam severamente o desempenho, mesmo com dezenas de micrômetros.
\end{itemize}}
\end{questao}

\begin{questao}[3]{}
Duas bobinas compartilham parte do fluxo. A bobina $1$ produz
$\SI{5}{\milli\weber}$, dos quais $\SI{4}{\milli\weber}$ atravessam a bobina
$2$.
\begin{itensq}
\item Determine o fluxo de dispersão.
\item Determine o coeficiente de acoplamento aproximado.
\item Discuta como aumentá-lo.
\end{itensq}
\resposta{$\SI{1}{\milli\weber}$ de dispersão; $k\approx0{,}8$}
\resolucao{(a) \textbf{Dispersão:} a parcela que não atravessa a outra bobina:
\[
\Phi_{disp}=5-4=\SI{1}{\milli\weber}\quad(20\%).
\]
(b) O coeficiente de acoplamento é a fração de fluxo compartilhada:
\[
k\approx\frac{\Phi_{12}}{\Phi_1}=\frac{4}{5}=0{,}8.
\]
\emph{(A definição rigorosa é $k=M/\sqrt{L_1L_2}$, que coincide com esta razão
quando as bobinas são simétricas.)}\\
(c) \textbf{Como aumentar $k$.}
\begin{itemize}
\item \textbf{Núcleo fechado} de alta permeabilidade: o fluxo prefere o caminho
de baixa relutância e "não escapa". Toroides atingem $k>0{,}99$.
\item \textbf{Enrolamentos concêntricos ou intercalados}, em vez de lado a lado
--- reduz drasticamente o caminho de fuga.
\item \textbf{Eliminar entreferros} desnecessários, que forçam o fluxo a
divergir.
\end{itemize}
\textbf{Compromisso:} $k$ alto é desejável em transformadores (transferência
eficiente), mas dispersão \emph{controlada} é útil em alguns casos --- ela
funciona como indutância série natural, limitando corrente de curto em
transformadores de solda e em reatores.}
\end{questao}

\begin{questao}[3]{}
Um núcleo tem seção variável: $\SI{6}{\centi\meter\squared}$ em um trecho e
$\SI{2}{\centi\meter\squared}$ em outro, sem dispersão.
\begin{itensq}
\item Compare o fluxo nos dois trechos.
\item Compare a densidade de fluxo.
\item Identifique onde ocorre a saturação primeiro.
\end{itensq}
\resposta{Mesmo fluxo; $B$ três vezes maior na seção estreita, que satura
primeiro}
\resolucao{(a) \textbf{O fluxo é o mesmo.} Sem dispersão, todas as linhas que
passam por um trecho passam pelo outro --- é a consequência direta de
$\oint\vec B\cdot\mathrm{d}\vec A=0$ aplicada a uma superfície que envolve a
junção.\\
\textbf{Analogia elétrica:} é o equivalente magnético da LKC. O fluxo se
conserva ao longo de um circuito magnético sem ramificação, como a corrente em
um circuito série.\\
(b) Como $B=\Phi/A$:
\[
\frac{B_2}{B_1}=\frac{A_1}{A_2}=\frac{6}{2}=3.
\]
(c) \textbf{A seção estreita satura primeiro}, por ter $B$ três vezes maior.\\
\textbf{Consequência de projeto:} o dimensionamento de um núcleo é governado
pela sua \textbf{menor} seção. Aumentar a seção larga não adianta nada ---
é preciso alargar o gargalo.\\
\textbf{Onde isso aparece na prática:} cantos de núcleos em E, junções mal
projetadas e regiões com furos de fixação. São esses pontos que saturam
primeiro e produzem aquecimento e distorção localizados.}
\end{questao}

\begin{questao}[3]{}
Uma espira de $\SI{25}{\centi\meter\squared}$ está parcialmente dentro de uma
região de campo de $\SI{0.8}{\tesla}$, com $60\%$ de sua área na região.
\begin{itensq}
\item Determine o fluxo.
\item Determine a f.e.m. se a espira for retirada completamente em
      $\SI{40}{\milli\second}$.
\end{itensq}
\resposta{$\Phi=\SI{1.2}{\milli\weber}$; $\varepsilon=\SI{30}{\milli\volt}$}
\resolucao{(a) Apenas a área \emph{dentro} do campo contribui:
\[
\Phi=B\,A_{efetiva}=(0{,}8)(0{,}60)(\pot{2.5}{-3})
=\pot{1.2}{-3}\,\si{\weber}.
\]
(b)
\[
|\varepsilon|=\frac{\Delta\Phi}{\Delta t}
=\frac{\pot{1.2}{-3}}{\pot{4}{-2}}=\pot{3}{-2}\,\si{\volt}.
\]
\textbf{Observação importante:} o fluxo depende da área \emph{efetivamente}
imersa no campo, não da área geométrica total da espira. Fora da região,
$B=0$ e a contribuição é nula.\\
\textbf{Se a espira tivesse $N$ espiras}, a f.e.m. seria $N$ vezes maior --- e é
assim que se aumenta a sensibilidade de sensores de fluxo sem aumentar a área.}
\end{questao}

\begin{questao}[3]{}
Um transformador tem núcleo de seção $\SI{10}{\centi\meter\squared}$ e opera
com $B_{\max}=\SI{1.2}{\tesla}$.
\begin{itensq}
\item Determine o fluxo máximo.
\item Determine o número de espiras do primário para
      $\SI{220}{\volt}$ a $\SI{60}{\hertz}$.
\end{itensq}
\resposta{$\Phi_{\max}=\SI{1.2}{\milli\weber}$; $N\approx688$ espiras}
\resolucao{(a) $\Phi_{\max}=B_{\max}A=(1{,}2)(\pot{1}{-3})
=\pot{1.2}{-3}\,\si{\weber}$.\\
(b) Para fluxo senoidal, a f.e.m. eficaz vale
\[
V_{ef}=4{,}44\,N\,f\,\Phi_{\max},
\]
\emph{(o fator $4{,}44=2\pi/\sqrt2$ decorre de derivar uma senoide e converter
para valor eficaz)}. Logo
\[
N=\frac{220}{4{,}44(60)(\pot{1.2}{-3})}=\frac{220}{0{,}3197}\approx688.
\]
\textbf{Compromisso central do projeto:} menos espiras significam $B$ maior (e
risco de saturação); mais espiras significam mais cobre, mais resistência e mais
perdas.\\
\textbf{Por que a frequência importa:} $N\propto1/f$. Em $\SI{400}{\hertz}$
(aviação) ou em dezenas de quilohertz (fontes chaveadas), o mesmo transformador
precisa de muito menos espiras --- e o núcleo pode ser muito menor. É essa a
razão física de fontes chaveadas serem tão mais compactas que as lineares.}
\end{questao}

\begin{questao}[3]{}
Uma superfície fechada envolve completamente um ímã permanente. Determine o
fluxo magnético líquido através dela e justifique.
\resposta{Nulo, independentemente da posição e da intensidade do ímã}
\resolucao{\textbf{Resultado:} $\oint\vec B\cdot\mathrm{d}\vec A=0$, sempre.\\
\textbf{Justificativa:} as linhas de campo de um ímã \emph{não começam nem
terminam} nele --- elas saem do polo norte, percorrem o espaço externo, entram
pelo sul e \textbf{continuam dentro do material}, fechando-se sobre si mesmas.\\
Toda linha que atravessa a superfície para fora atravessa também para dentro,
em algum outro ponto. O balanço é rigorosamente nulo.\\
\textbf{Contraste com a carga elétrica.} Se a superfície envolvesse uma carga
$+Q$, o fluxo elétrico seria $Q/\varepsilon_0\ne0$ --- porque as linhas
elétricas \emph{nascem} na carga.\\
\textbf{O que isso significa fisicamente:} não existe "carga magnética". Cortar
um ímã ao meio não separa os polos --- produz dois ímãs completos, cada um com
seus dois polos. Não há como isolar um norte de um sul.\\
\textbf{Consequência para a teoria:} $\nabla\cdot\vec B=0$ é uma equação de
Maxwell, e é ela que garante a existência do potencial vetor $\vec A$ com
$\vec B=\nabla\times\vec A$.}
\end{questao}

\begin{questao}[3]{}
Uma bobina de $500$ espiras e área $\SI{8}{\centi\meter\squared}$ é usada como
sensor de campo. Ela é retirada rapidamente do campo, e um integrador registra
$\SI{6}{\milli\volt\second}$.
\begin{itensq}
\item Determine o fluxo concatenado inicial.
\item Determine o campo.
\item Explique a vantagem do método.
\end{itensq}
\resposta{$\lambda=\SI{6}{\milli\weber}$-espira; $B=\SI{15}{\milli\tesla}$}
\resolucao{(a) Integrando a lei de Faraday:
\[
\int\varepsilon\,\mathrm{d}t=\Delta\lambda
\;\Longrightarrow\;
\lambda_i=\pot{6}{-3}\,\si{\weber}\text{-espira}
\]
(pois o fluxo final é zero).\\
(b)
\[
\Phi=\frac{\lambda}{N}=\frac{\pot{6}{-3}}{500}=\pot{1.2}{-5}\,\si{\weber};
\qquad
B=\frac{\Phi}{A}=\frac{\pot{1.2}{-5}}{\pot{8}{-4}}=\pot{1.5}{-2}\,\si{\tesla}.
\]
(c) \textbf{Vantagem do método integral.} O resultado depende \emph{apenas} da
variação total de fluxo --- \textbf{não} de \emph{como} a bobina foi retirada.
Rápido ou devagar, com solavancos ou suavemente, a integral é a mesma.\\
Isso torna a medição notavelmente robusta: não é preciso controlar a velocidade
nem a trajetória.\\
\textbf{Limitações:} o integrador acumula qualquer \emph{offset} do
amplificador, o que limita o tempo de medição (problema já discutido na
integração de sinais ruidosos). E o método mede apenas \emph{variações} --- para
campos estáticos permanentes, é preciso mover a bobina ou usar sensor Hall.}
\end{questao}

\begin{questao}[3]{}
Avalie o erro de assumir campo uniforme ao calcular o fluxo através de uma
espira de $\SI{10}{\centi\meter}$ de lado, situada entre $\SI{5}{}$ e
$\SI{15}{\centi\meter}$ de um fio retilíneo.
\begin{itensq}
\item Determine o fluxo exato e o estimado pelo campo no centro.
\item Determine o erro.
\end{itensq}
\resposta{Erro de cerca de $+2{,}5\%$ na estimativa pelo centro}
\resolucao{\textbf{Exato} (com $I$ genérico e $a=\SI{0.10}{\meter}$):
\[
\Phi_{ex}=\frac{\mu_0Ia}{2\pi}\ln\frac{0{,}15}{0{,}05}
=\pot{2}{-8}I\ln3=\pot{2.197}{-8}\,I.
\]
\textbf{Estimativa pelo centro} ($r=\SI{0.10}{\meter}$):
\[
\Phi_{est}=\frac{\mu_0I}{2\pi(0{,}10)}\,a^{2}
=\pot{2}{-7}\frac{I}{0{,}10}(0{,}01)=\pot{2}{-8}\,I.
\]
\[
\text{erro}=\frac{2{,}00-2{,}197}{2{,}197}=-9{,}0\%.
\]
\textbf{A estimativa pelo centro \emph{subestima}} o fluxo --- porque
$B\propto1/r$ é uma função \textbf{convexa}, e a média de uma função convexa é
maior que a função da média (desigualdade de Jensen).\\
\textbf{Quando a aproximação é aceitável:} quando a variação de $B$ ao longo da
espira for pequena. Aqui, $B$ varia por um fator $3$ --- muito. Se a espira
estivesse entre $95$ e $\SI{105}{\centi\meter}$, a variação seria de apenas
$10\%$ e o erro cairia para frações de por cento.\\
\textbf{Critério prático:} use campo uniforme se
$B_{\max}/B_{\min}<1{,}1$; caso contrário, integre.}
\end{questao}

\blocodesafio

\begin{questao}[4]{}
\textbf{(Demonstração --- fluxo por superfície fechada.)} Prove que
$\oint\vec B\cdot\mathrm{d}\vec A=0$ e explore suas consequências.
\begin{itensq}
\item Demonstre a partir da inexistência de monopolos.
\item Deduza a conservação do fluxo em circuitos magnéticos.
\item Compare com a lei de Gauss elétrica.
\end{itensq}
\resposta{Linhas fechadas $\Rightarrow$ fluxo líquido nulo $\Rightarrow$ o fluxo
se conserva ao longo de um circuito magnético}
\resolucao{(a) \textbf{Demonstração.} As linhas de $\vec B$ são curvas
\emph{fechadas} --- fato experimental, equivalente à inexistência de monopolos
magnéticos. Uma curva fechada que penetre uma superfície fechada precisa
atravessá-la um número \emph{par} de vezes, alternando entrada e saída.\\
Cada entrada contribui com fluxo negativo e cada saída com positivo, de módulos
iguais. A soma se cancela:
\[
\oint\vec B\cdot\mathrm{d}\vec A=0.
\]
Pelo teorema da divergência, isso equivale a $\nabla\cdot\vec B=0$ em todo
ponto.\\
(b) \textbf{Conservação do fluxo.} Considere uma junção em um núcleo magnético,
onde um trecho de seção $A_1$ encontra outro de seção $A_2$. Envolvendo a junção
com uma superfície fechada:
\[
-\Phi_1+\Phi_2=0
\;\Longrightarrow\;
\Phi_1=\Phi_2.
\]
\textbf{O fluxo se conserva} --- exatamente como a corrente em um nó. Se o
núcleo se ramificar em dois caminhos:
\[
\Phi_{total}=\Phi_a+\Phi_b,
\]
que é a \textbf{lei dos nós magnética}, análoga à LKC.\\
(c) \textbf{Comparação com Gauss elétrica.}
\begin{center}
\begin{tabular}{ll}
\hline
Elétrico & Magnético\\
\hline
$\oint\vec E\cdot\mathrm{d}\vec A=Q/\varepsilon_0$
 & $\oint\vec B\cdot\mathrm{d}\vec A=0$\\
Linhas nascem e morrem em cargas & Linhas são fechadas\\
Existem monopolos (cargas) & Não existem monopolos\\
$\nabla\cdot\vec E=\rho/\varepsilon_0$ & $\nabla\cdot\vec B=0$\\
\hline
\end{tabular}
\end{center}
\textbf{Consequência profunda:} a ausência de monopolos permite escrever
$\vec B=\nabla\times\vec A$, base de toda a formulação moderna do
eletromagnetismo. Se monopolos existissem, essa construção seria impossível.}
\end{questao}

\begin{questao}[4]{}
\textbf{(Independência da superfície.)} Prove que o fluxo através de qualquer
superfície apoiada em uma dada curva fechada é o mesmo.
\begin{itensq}
\item Demonstre.
\item Explique a importância para a lei de Faraday.
\item Indique o que ocorre se as normais forem orientadas incoerentemente.
\end{itensq}
\resposta{Duas superfícies de mesma borda formam superfície fechada, cujo fluxo
é nulo}
\resolucao{(a) \textbf{Demonstração.} Sejam $S_1$ e $S_2$ duas superfícies
apoiadas na mesma curva $C$. Juntas, elas delimitam um \textbf{volume fechado}.
Aplicando o resultado anterior:
\[
\oint\vec B\cdot\mathrm{d}\vec A
=\Phi_{S_2}^{(saindo)}-\Phi_{S_1}^{(saindo)}=0
\;\Longrightarrow\;
\Phi_{S_1}=\Phi_{S_2}.
\]
(b) \textbf{Importância para Faraday.} A lei
$\varepsilon=-\dfrac{\mathrm{d}\Phi}{\mathrm{d}t}$ só faz sentido se
$\Phi$ estiver \emph{bem definido} --- ou seja, se não depender da superfície
escolhida.\\
Se dependesse, a f.e.m. de um mesmo circuito teria valores diferentes conforme a
superfície imaginada --- absurdo, pois a f.e.m. é medível.\\
\textbf{Liberdade prática:} escolha a superfície que simplificar a integral. Ao
calcular o fluxo através de uma espira que envolve um solenoide, por exemplo,
convém uma superfície que corte o solenoide perpendicularmente.\\
(c) \textbf{Orientação das normais.} A igualdade exige que as normais sejam
orientadas \emph{coerentemente} com o sentido de percurso da borda, pela regra
da mão direita.\\
Se uma delas for invertida, aparece um sinal negativo --- e a conclusão errônea
de que os fluxos são opostos. \textbf{Fixe o sentido de percurso da curva
primeiro}, e derive as normais dele.}
\end{questao}

\begin{questao}[4]{}
\textbf{(Circuito magnético.)} Estabeleça a analogia formal entre circuitos
magnéticos e elétricos, e discuta seus limites.
\begin{itensq}
\item Construa a tabela de correspondências.
\item Deduza a expressão da relutância.
\item Aponte três diferenças que quebram a analogia.
\end{itensq}
\resposta{$\mathcal{F}=\Phi\mathcal{R}$, com
$\mathcal{R}=\ell/\mu A$; a analogia falha por não linearidade, dispersão e
ausência de dissipação}
\resolucao{(a)
\begin{center}
\begin{tabular}{ll}
\hline
Circuito elétrico & Circuito magnético\\
\hline
F.e.m. $\varepsilon$ & Força magnetomotriz $\mathcal{F}=NI$\\
Corrente $I$ & Fluxo $\Phi$\\
Resistência $R=\ell/\sigma A$ & Relutância
 $\mathcal{R}=\ell/\mu A$\\
$\varepsilon=RI$ & $\mathcal{F}=\mathcal{R}\Phi$\\
Condutividade $\sigma$ & Permeabilidade $\mu$\\
LKC nos nós & Conservação do fluxo\\
Resistências em série somam & Relutâncias em série somam\\
\hline
\end{tabular}
\end{center}
(b) \textbf{Dedução da relutância.} Da lei de Ampère aplicada ao caminho médio
do núcleo, $NI=H\ell$. Com $B=\mu H$ e $\Phi=BA$:
\[
NI=\frac{B\ell}{\mu}=\frac{\Phi\,\ell}{\mu A}
\;\Longrightarrow\;
\mathcal{R}=\frac{\ell}{\mu A}.
\]
(c) \textbf{Três diferenças que quebram a analogia.}
\begin{enumerate}
\item \textbf{Não linearidade.} $\sigma$ é constante em metais, mas
$\mu$ varia enormemente com $B$ --- e desaba na saturação. A "resistência
magnética" depende da própria corrente que a atravessa.
\item \textbf{Dispersão.} A corrente elétrica fica confinada no condutor porque
a razão de condutividades metal/ar é de $10^{20}$. Já a razão de
permeabilidades ferro/ar é de apenas $10^{3}$ --- e por isso parte do fluxo
\emph{escapa} do núcleo. Não existe "isolante magnético".
\item \textbf{Ausência de dissipação.} A corrente dissipa $RI^{2}$ em calor.
O fluxo \emph{não dissipa} nada --- manter $\Phi$ constante não custa energia
(as perdas por histerese e Foucault ocorrem apenas quando o fluxo \emph{varia},
e por outros mecanismos).
\end{enumerate}
\textbf{Conclusão:} a analogia é excelente como \emph{ferramenta de cálculo}
abaixo da saturação, e enganosa como \emph{modelo físico}. Use-a para
dimensionar; não a use para raciocinar sobre energia.}
\end{questao}

\begin{questao}[4]{}
\textbf{(Projeto de entreferro.)} Um indutor de $\SI{100}{\micro\henry}$ deve
conduzir $\SI{5}{\ampere}$ de pico, com núcleo de ferrite de seção
$\SI{1}{\centi\meter\squared}$, caminho médio $\SI{6}{\centi\meter}$,
$\mu_r=2000$ e $B_{sat}=\SI{0.3}{\tesla}$.
\begin{itensq}
\item Determine a energia armazenada e o número mínimo de espiras.
\item Determine o entreferro necessário.
\item Determine a fração de energia armazenada no entreferro.
\end{itensq}
\resposta{$W=\SI{1.25}{\milli\joule}$; $N=17$ espiras;
$g\approx\SI{0.33}{\milli\meter}$; $92\%$ da energia no entreferro}
\resolucao{(a)
\[
W=\tfrac12LI^{2}=\tfrac12(\pot{1}{-4})(25)=\SI{1.25}{\milli\joule}.
\]
O fluxo concatenado é $\lambda=LI=\pot{5}{-4}\,\si{\weber}$-espira, e o fluxo
máximo por espira é limitado pela saturação:
\[
\Phi_{\max}=B_{sat}A=(0{,}3)(\pot{1}{-4})=\SI{30}{\micro\weber}.
\]
\[
N\ge\frac{\lambda}{\Phi_{\max}}=\frac{\pot{5}{-4}}{\pot{3}{-5}}=16{,}7
\;\Longrightarrow\;
N=17\ \text{espiras}.
\]
\emph{(Verificação: com $N=17$, $B=\lambda/(NA)=\SI{0.294}{\tesla}$ ---
logo abaixo do limite.)}\\
(b) De $L=N^{2}/\mathcal{R}$:
\[
\mathcal{R}_{tot}=\frac{17^{2}}{\pot{1}{-4}}=\pot{2.89}{6}\,\si{\per\henry}.
\]
A parcela do ferrite é
$\mathcal{R}_{fe}=\dfrac{0{,}06}{(2000)\mu_0(\pot{1}{-4})}
=\pot{2.39}{5}$, restando para o entreferro
$\mathcal{R}_g=\pot{2.65}{6}$:
\[
g=\mathcal{R}_g\,\mu_0A=(\pot{2.65}{6})(4\pi\times10^{-7})(\pot{1}{-4})
=\SI{0.33}{\milli\meter}.
\]
\textbf{Note que o entreferro responde por $92\%$ da relutância total}, embora
seja $180$ vezes mais curto que o caminho de ferrite.\\
(c) Como $B$ é o mesmo nos dois meios e $u=B^{2}/2\mu$:
\[
\frac{W_g}{W_{fe}}=\frac{g/\mu_0}{\ell_{fe}/(\mu_r\mu_0)}
=\frac{g\,\mu_r}{\ell_{fe}}
=\frac{(\pot{3.33}{-4})(2000)}{0{,}06}=11{,}1,
\]
ou seja, $\SI{91.7}{\percent}$ da energia está no \textbf{entreferro}.\\
\textbf{Por que isso importa.} A densidade de energia
$u=B^{2}/2\mu$ é \emph{inversamente} proporcional a $\mu$ --- para o mesmo
$B$, o ar armazena duas mil vezes mais energia por volume que o ferrite.\\
\textbf{O entreferro é o reservatório; o núcleo apenas conduz o fluxo até
ele.} Um indutor sem entreferro satura com corrente mínima e armazena quase
nada --- ao contrário de um transformador, que \emph{não deve} armazenar
energia e por isso dispensa entreferro.}
\end{questao}

\begin{questao}[4]{}
\textbf{(Dispersão e acoplamento.)} Analise o efeito da dispersão de fluxo no
desempenho de um transformador.
\begin{itensq}
\item Relacione $k$ com as indutâncias própria e mútua.
\item Determine o efeito sobre a tensão de saída.
\item Discuta quando a dispersão é desejável.
\end{itensq}
\resposta{$k=M/\sqrt{L_1L_2}$; a dispersão aparece como indutância série e
causa queda de tensão com a carga}
\resolucao{(a) Por definição,
\[
k=\frac{M}{\sqrt{L_1L_2}},
\qquad 0\le k\le1,
\]
resultado já demonstrado a partir da não negatividade da energia do par de
bobinas.\\
Decompondo cada indutância em parcela \emph{magnetizante} (fluxo compartilhado)
e parcela de \emph{dispersão}:
\[
L_1=L_{m1}+L_{d1},
\qquad
k^{2}=\frac{L_{m1}L_{m2}}{L_1L_2}.
\]
(b) \textbf{Efeito na saída.} A indutância de dispersão comporta-se como uma
\emph{impedância em série} com o enrolamento. Sob carga, ela produz queda de
tensão proporcional à corrente:
\begin{itemize}
\item \textbf{Em vazio}, a relação de tensões é praticamente $N_2/N_1$ --- a
dispersão não se manifesta.
\item \textbf{Com carga}, a tensão de saída \emph{cai}, e a regulação piora
proporcionalmente à dispersão.
\end{itemize}
Um transformador com $k=0{,}98$ tem cerca de $2\%$ de dispersão e regulação
tipicamente na faixa de $2$ a $5\%$.\\
(c) \textbf{Quando a dispersão é desejável.}
\begin{itemize}
\item \textbf{Transformadores de solda:} a dispersão elevada (às vezes
$k\approx0{,}8$) limita naturalmente a corrente de arco, dispensando reator
externo. A característica "mole" é justamente o que se quer.
\item \textbf{Transformadores de forno e de descarga:} mesma lógica --- a carga
é essencialmente um curto controlado.
\item \textbf{Limitação de curto-circuito:} em transformadores de distribuição,
uma dispersão mínima é \emph{exigida} por norma, para limitar a corrente de
falta a valores que a proteção consiga interromper.
\end{itemize}
\textbf{Conclusão:} $k\to1$ não é universalmente melhor. A dispersão é um
parâmetro de projeto, e não apenas um defeito a minimizar.}
\end{questao}

\begin{questao}[4]{}
\textbf{(Fluxo e energia.)} Deduza a energia armazenada em um circuito
magnético e mostre onde ela reside.
\begin{itensq}
\item Deduza $W=\frac12\mathcal{R}\Phi^{2}$.
\item Determine a densidade de energia em função de $B$ e $\mu$.
\item Determine a fração de energia no entreferro do exemplo anterior.
\end{itensq}
\resposta{$u=B^{2}/2\mu$; praticamente toda a energia fica no entreferro}
\resolucao{(a) Com $L=N^{2}/\mathcal{R}$ e $\lambda=N\Phi=LI$:
\[
W=\frac12LI^{2}=\frac{\lambda^{2}}{2L}
=\frac{(N\Phi)^{2}}{2N^{2}/\mathcal{R}}
=\frac{1}{2}\mathcal{R}\,\Phi^{2}.
\]
\textbf{Analogia:} é a forma magnética de $W=\frac12 CU^{2}$ escrita como
$Q^{2}/2C$ --- o fluxo faz o papel da carga e a relutância, o do inverso da
capacitância.\\
(b) Escrevendo $\mathcal{R}=\ell/\mu A$ e $\Phi=BA$:
\[
W=\frac{1}{2}\frac{\ell}{\mu A}(BA)^{2}=\frac{B^{2}}{2\mu}(A\ell)
\;\Longrightarrow\;
u=\frac{B^{2}}{2\mu}.
\]
\textbf{Ponto decisivo:} a densidade de energia é \emph{inversamente}
proporcional a $\mu$. Para o mesmo $B$, o ar armazena $\mu_r$ vezes \emph{mais}
energia por volume que o ferro.\\
(c) \textbf{Fração no entreferro.} Como $B$ é o mesmo nos dois meios (fluxo
conservado e mesma seção), a razão das energias é a razão dos produtos
$\ell/\mu$:
\[
\frac{W_g}{W_{fe}}
=\frac{g/\mu_0}{\ell_{fe}/(\mu_r\mu_0)}
=\frac{g\,\mu_r}{\ell_{fe}}
=\frac{(\pot{0.56}{-3})(2000)}{0{,}50}=2{,}24.
\]
Ou seja, cerca de $69\%$ da energia está no entreferro, contra $31\%$ no
ferro --- e a proporção cresce rapidamente com $g$.\\
\textbf{Por isso o entreferro é o elemento de projeto de um indutor de
potência:} ele é o \emph{reservatório}. O ferro apenas conduz o fluxo até ele.}
\end{questao}

\begin{questao}[4]{}
\textbf{(Núcleo de seção variável.)} Formule o cálculo do fluxo em um núcleo
cuja seção varia ao longo do caminho, sem dispersão.
\begin{itensq}
\item Escreva a relutância total.
\item Determine o critério de dimensionamento.
\item Analise o caso de um núcleo com estreitamento localizado.
\end{itensq}
\resposta{$\mathcal{R}=\displaystyle\int\frac{\mathrm{d}\ell}{\mu A(\ell)}$;
dimensiona-se pela \emph{menor} seção}
\resolucao{(a) Cada trecho infinitesimal contribui como uma relutância em
série:
\[
\mathcal{R}=\int_0^{L}\frac{\mathrm{d}\ell}{\mu\,A(\ell)}.
\]
Para trechos de seção constante, isso vira a soma
$\sum\ell_i/(\mu A_i)$.\\
(b) \textbf{Critério de dimensionamento.} Sem dispersão, o fluxo é o
\emph{mesmo} em toda a extensão. Logo
\[
B(\ell)=\frac{\Phi}{A(\ell)},
\]
e $B$ é máximo onde $A$ é \textbf{mínimo}. A condição de não saturação é
\[
\Phi\le B_{sat}\,A_{\min}.
\]
\textbf{Toda a capacidade do núcleo é ditada pela sua menor seção} --- alargar
o resto não ajuda em nada.\\
(c) \textbf{Estreitamento localizado.} Suponha um núcleo de
$\SI{6}{\centi\meter\squared}$ com um trecho curto de
$\SI{2}{\centi\meter\squared}$:
\begin{itemize}
\item \textbf{Efeito na relutância:} se o trecho estreito for curto, sua
contribuição a $\mathcal{R}$ é pequena --- o fluxo total quase não muda.
\item \textbf{Efeito na saturação:} ali $B$ é \emph{três vezes} maior, e o
trecho satura com um terço do fluxo que o resto suportaria.
\item \textbf{Efeito local:} saturado, aquele trecho aquece por perdas
concentradas e sua permeabilidade cai --- criando um entreferro
\emph{efetivo} não projetado.
\end{itemize}
\textbf{Onde isso ocorre na prática:} furos de fixação, cantos internos de
núcleos em E, e emendas de lâminas. São defeitos que passam despercebidos no
cálculo de relutância (pequeno efeito) mas dominam o comportamento na saturação
(efeito grande).\\
\textbf{Regra:} calcule $\mathcal{R}$ pelo caminho todo, mas verifique
$B$ na \emph{seção crítica}.}
\end{questao}


 
 

% END BANCO COMPLEMENTAR

\gabarito[12.3 Fluxo magnético]

% =====================================================================
\subsection{Força magnética}
% =====================================================================

\blocofixacao

\begin{questao}[1]{}
A força magnética sobre uma carga em movimento é \textbf{nula} quando:
\begin{alternativas}
\item a velocidade é perpendicular ao campo.
\item a velocidade é paralela (ou antiparalela) ao campo.
\item a carga é positiva.
\item o campo é muito intenso.
\item a velocidade é máxima.
\end{alternativas}
\resposta{(b)}
\resolucao{Como $F = qvB\sin\theta$, a força se anula quando $\sin\theta = 0$, ou
seja, quando $\vec{v}$ é paralela (ou antiparalela) a $\vec{B}$. É máxima quando
$\vec{v}\perp\vec{B}$ ($\theta = \SI{90}{\degree}$).}
\end{questao}

\begin{questao}[1]{}
Um condutor retilíneo de $\SI{30}{\centi\meter}$ conduz $\SI{2}{\ampere}$
perpendicularmente a um campo de $\SI{0.5}{\tesla}$. Determine a força sobre ele.
\resposta{$F = \SI{0.3}{\newton}$}
\resolucao{$F = B\,I\,\ell = 0{,}5\cdot2\cdot0{,}3 = \SI{0.3}{\newton}$
(com $\sin\theta = 1$, pois são perpendiculares).}
\end{questao}

\begin{questao}[1]{}
Uma carga positiva move-se horizontalmente para a direita, em uma região onde o
campo magnético aponta para dentro do plano da página. A força magnética sobre a
carga aponta:
\begin{alternativas}[2]
\item para cima.
\item para baixo.
\item para a direita.
\item para a esquerda.
\item para dentro da página.
\end{alternativas}
\resposta{(a)}
\resolucao{Pela regra da mão direita para $\vec F = q\vec v\times\vec B$: dedos
apontando no sentido de $\vec v$ (direita), curvando-se para $\vec B$ (para dentro
da página), o polegar aponta para \emph{cima}. Como a carga é positiva, a força
tem esse mesmo sentido. (Se fosse negativa, seria para baixo.)}
\end{questao}

\blocoaplicacao

\begin{questao}[2]{}
Uma carga de $\SI{1.6e-19}{\coulomb}$ move-se a
$\SI{e6}{\meter\per\second}$, perpendicularmente a um campo de $\SI{0.2}{\tesla}$.
Determine o módulo da força magnética e descreva a trajetória resultante.
\resposta{$F = \SI{3.2e-14}{\newton}$; trajetória circular}
\resolucao{$F = q\,v\,B = \pot{1.6}{-19}\cdot\pot{1}{6}\cdot0{,}2
= \pot{3.2}{-14}\,\si{\newton}$.\\
Como a força é sempre \emph{perpendicular} à velocidade, ela não altera o módulo
de $v$ (não realiza trabalho), apenas curva a trajetória: o resultado é um
\textbf{movimento circular uniforme}. É o princípio dos cíclotrons e dos
espectrômetros de massa.}
\end{questao}

\begin{questao}[2]{}
Uma espira retangular percorrida por corrente é colocada em um campo magnético
uniforme. Explique por que ela tende a \textbf{girar}, e não a se deslocar, e
relacione isso com o funcionamento de um motor elétrico.
\resposta{Ver comentário}
\resolucao{Em um campo uniforme, os dois lados opostos da espira conduzem corrente
em sentidos \emph{contrários}, sofrendo forças magnéticas iguais e opostas
($\vec F = B I \ell$). Essas forças formam um \textbf{binário} (conjugado): a
resultante é nula (não há translação), mas o \emph{torque} não é --- a espira
gira. Esse é exatamente o princípio do \textbf{motor de corrente contínua}: o
torque $\tau = N B I A \sin\theta$ faz o rotor girar, e o comutador inverte a
corrente a cada meia-volta para manter a rotação no mesmo sentido.}
\end{questao}

\blocoaprofundamento

\begin{questao}[3]{}
Um próton ($q = \SI{1.6e-19}{\coulomb}$, $m = \SI{1.67e-27}{\kilogram}$) entra
perpendicularmente em um campo uniforme de $\SI{0.5}{\tesla}$ com velocidade
$\SI{2e6}{\meter\per\second}$. Determine o raio da trajetória circular descrita.
\resposta{$r \approx \SI{4.2}{\centi\meter}$}
\resolucao{A força magnética atua como resultante centrípeta:
$qvB = \dfrac{mv^2}{r}$. Isolando $r$:
\[
r=\frac{mv}{qB}
=\frac{(\pot{1.67}{-27})(\pot{2}{6})}{(\pot{1.6}{-19})(0{,}5)}
\approx\pot{4.2}{-2}\,\si{\meter}=\SI{4.2}{\centi\meter}.
\]
O raio cresce com a quantidade de movimento ($mv$) e diminui com a carga e o
campo. Medindo $r$ conhece-se a razão carga/massa da partícula --- é assim que
funciona o espectrômetro de massa.}
\end{questao}

\begin{questao}[3]{}
Um seletor de velocidades é formado por campos elétrico e magnético
\textbf{cruzados} (perpendiculares entre si e à velocidade). Partículas
carregadas entram com velocidades variadas e apenas as de uma velocidade
específica atravessam sem desviar.

\circuitoq{%
\begin{tikzpicture}[scale=1,line width=0.8pt,>=Stealth]
 \draw[cinza!50,fill=cinza!8] (0,-1.2) rectangle (5.4,1.2);
 % campo E (setas para baixo)
 \foreach \x in {0.7,1.6,2.5,3.4,4.3}
   \draw[->,Vcol!70] (\x,1.05)--(\x,-1.05);
 \node[Vcol,font=\footnotesize] at (5.05,0.9) {$\vec E$};
 % campo B (para dentro da pagina)
 \foreach \x in {1.15,2.05,2.95,3.85}
   \foreach \y in {-0.6,0,0.6}
     {\node[Icol,font=\scriptsize] at (\x,\y) {$\otimes$};}
 \node[Icol,font=\footnotesize] at (0.35,-0.9) {$\vec B$};
 % particula
 \draw[->,laranja,line width=1.2pt] (-1.3,0)--(0,0)
       node[midway,above,font=\footnotesize]{$v$};
 \draw[->,laranja,line width=1.2pt] (5.4,0)--(6.5,0);
 \fill[laranja] (-1.3,0) circle (0.1);
\end{tikzpicture}}{Seletor de velocidades com campos cruzados.}

\begin{itensq}
\item Deduza a condição para que a partícula atravesse sem desvio.
\item Mostre que essa condição \emph{não depende} da carga nem da massa da
      partícula.
\item Com $E = \SI{3e5}{\volt\per\meter}$ e $B = \SI{0.15}{\tesla}$, qual a
      velocidade selecionada?
\item Após o seletor, coloca-se um campo magnético $B'$ que curva a trajetória.
      Explique como essa combinação permite medir a razão $q/m$.
\end{itensq}
\resposta{(a) $qE = qvB \Rightarrow v = E/B$; (c) $v = \SI{2e6}{\meter\per\second}$;
(b) e (d) ver comentário}
\resolucao{\textbf{(a)} Sobre a partícula atuam duas forças perpendiculares à
trajetória e \emph{opostas} entre si:
\[
F_e = qE \quad(\text{para baixo, se } q>0),
\qquad
F_m = qvB \quad(\text{para cima, pela regra da mão direita}).
\]
Para não haver desvio, a resultante deve ser nula:
\[
qE = qvB
\;\Longrightarrow\;
\boxed{\,v=\frac{E}{B}\,}
\]
\textbf{(b)} A carga $q$ \textbf{cancela-se} nos dois lados da igualdade, e a
massa sequer aparece --- não há aceleração, logo a segunda lei de Newton não é
invocada. Consequência notável: o seletor filtra por \emph{velocidade}, e apenas
por ela. Elétrons, prótons e íons pesados com a mesma $v$ atravessam igualmente.
Partículas mais rápidas são desviadas pela força magnética (que cresce com $v$);
mais lentas, pela força elétrica.

\textbf{(c)} $v = \dfrac{E}{B} = \dfrac{\pot{3}{5}}{0{,}15}
= \pot{2}{6}\,\si{\meter\per\second}$.

\textbf{(d) O espectrômetro de massa.} Com a velocidade agora \emph{conhecida e
uniforme}, a partícula entra em uma região de campo $B'$ puro, onde a força
magnética atua como centrípeta:
\[
qvB' = \frac{mv^{2}}{r}
\;\Longrightarrow\;
\frac{q}{m}=\frac{v}{B'r}=\frac{E}{B\,B'\,r}.
\]
Medindo o raio $r$ da trajetória (pela marca deixada no detector) e conhecendo
$E$, $B$ e $B'$, obtém-se $q/m$ diretamente. Íons de massas diferentes descrevem
raios diferentes e se separam espacialmente --- daí o nome
\emph{espectrômetro}.\\
\textbf{Impacto histórico:} foi com esse arranjo que Thomson mediu $q/m$ do
elétron (1897) e que Aston descobriu os \textbf{isótopos} (1919), rendendo-lhe o
Nobel. Hoje é ferramenta padrão em química analítica, datação e antidoping.}
\end{questao}

% BEGIN BANCO COMPLEMENTAR Força magnética
% =====================================================================
%  Banco complementar --- Forca Magnetica  (REVISADO)
%  Substitui a versao gerada por template (5 clones em N1, 10 em N2,
%  10 em N3 e 8 questoes conceituais marcadas como N4).
%  O cap11 ja cobre: condicao de forca nula (MC), condutor de 30 cm com
%  2 A em 0,5 T, regra da mao para carga movendo-se com B entrando no
%  plano, carga elementar a 1e6 m/s em 0,2 T, espira retangular girando
%  (qualitativo), proton em campo uniforme (raio) e o SELETOR DE
%  VELOCIDADES com campos cruzados -- motivo pelo qual é a N4 antiga sobre
%  projeto de seletor foi removida deste banco. Ela e substituida pelo
%  ESPECTROMETRO DE MASSA, que usa o seletor como etapa e vai alem.
%  Este banco explora: forca entre fios e a definicao do ampere,
%  trajetoria helicoidal, torque e momento magnetico, efeito Hall
%  quantitativo, campo nao uniforme, projeto de atuador e a demonstracao
%  de que a forca magnetica nao realiza trabalho.
%  ACRESCIMO POSTERIOR: tres questoes sobre a FORCA DE LORENTZ COMPLETA,
%  F = q(E + v x B). As duas parcelas ja apareciam separadamente e os
%  campos cruzados surgiam no seletor de velocidades e na deriva ExB,
%  mas a expressao unificada nao estava escrita em lugar nenhum.
%  Distribuicao: 5 N1 + 11 N2 + 11 N3 + 9 N4 = 36
% =====================================================================

\subsubsection*{Banco complementar --- Força magnética}

\begin{atencao}[Organização do banco]
Duas expressões governam tudo: $\vec F=q\,\vec v\times\vec B$ para cargas e
$\vec F=I\,\vec\ell\times\vec B$ para condutores. Ambas são \textbf{produtos
vetoriais} --- a força é \emph{perpendicular} tanto à velocidade quanto ao
campo, e se anula quando eles são paralelos. Daí decorre o fato central: a força
magnética \textbf{nunca realiza trabalho}. O N3 trata de trajetórias,
espectrometria e efeito Hall; o N4 exige demonstrações e projeto.
\end{atencao}

\blocofixacao

\begin{questao}[1]{}
Uma carga $q=+\SI{2}{\micro\coulomb}$ move-se a $\SI{500}{\meter\per\second}$
perpendicularmente a um campo de $\SI{0.2}{\tesla}$. Determine a força.
\resposta{$F=\SI{0.2}{\milli\newton}$}
\resolucao{
\[
F=qvB\sin\theta=(\pot{2}{-6})(500)(0{,}2)(1)=\pot{2}{-4}\,\si{\newton}.
\]
\textbf{Direção:} perpendicular \emph{simultaneamente} a $\vec v$ e a
$\vec B$, dada pela regra da mão direita aplicada a
$\vec v\times\vec B$ (e invertida, se a carga for negativa).}
\end{questao}

\begin{questao}[1]{}
Um condutor retilíneo de $\SI{40}{\centi\meter}$ conduz $\SI{3}{\ampere}$
perpendicularmente a $B=\SI{0.25}{\tesla}$. Determine a força.
\resposta{$F=\SI{0.3}{\newton}$}
\resolucao{
\[
F=BIL\sin\theta=(0{,}25)(3)(0{,}40)(1)=\SI{0.3}{\newton}.
\]
\textbf{Origem da fórmula:} ela é a soma das forças sobre todos os portadores do
condutor. Escrevendo $I=nqv_dA$ e $L$ o comprimento, o produto
$nqv_dA\cdot L$ reproduz $qv_d$ somado sobre $nAL$ portadores.\\
A força é transmitida ao \emph{fio} pelas colisões dos portadores com a rede
cristalina --- não são os elétrons que "empurram" o fio diretamente.}
\end{questao}

\begin{questao}[1]{}
A mesma carga da primeira questão move-se agora formando $\ang{30}$ com o
campo. Determine a força.
\resposta{$F=\SI{0.1}{\milli\newton}$}
\resolucao{
\[
F=qvB\sin\ang{30}=(\pot{2}{-4})(0{,}5)=\pot{1}{-4}\,\si{\newton}.
\]
\textbf{Metade} do valor perpendicular.\\
\textbf{Casos-limite:} $\theta=\ang{90}$ dá força máxima; $\theta=0$ ou
$\ang{180}$ (velocidade \emph{paralela} ao campo) dá força \textbf{nula} --- a
carga segue em linha reta, sem sentir o campo.}
\end{questao}

\begin{questao}[1]{}
Uma carga positiva move-se para a \textbf{direita} em uma região onde
$\vec B$ aponta para \textbf{cima}. Determine o sentido da força.
\begin{alternativas}[2]
\item para dentro do plano
\item para fora do plano
\item para cima
\item para a direita
\end{alternativas}
\resposta{(a)}
\resolucao{Pela regra da mão direita para $\vec v\times\vec B$: dedos apontando
no sentido de $\vec v$ (direita), curvando-se para $\vec B$ (cima) --- o polegar
aponta para \textbf{dentro} do plano.\\
\textbf{Verificação:} a força deve ser perpendicular a ambos. "Direita" e
"cima" definem o plano da folha, logo a força só pode ser perpendicular a ele
--- o que elimina (c) e (d) imediatamente.\\
\textbf{Se a carga fosse negativa}, a força seria para fora do plano.}
\end{questao}

\begin{questao}[1]{}
A partir de $F=qvB$, mostre que $\SI{1}{\tesla}$ equivale a
$\SI{1}{\newton\second\per\coulomb\per\meter}$.
\resposta{$\si{\tesla}=\dfrac{\si{\newton}}{\si{\coulomb}\cdot
\si{\meter\per\second}}$}
\resolucao{
\[
B=\frac{F}{qv}
\;\Longrightarrow\;
[\si{\tesla}]=\frac{\si{\newton}}{\si{\coulomb}\cdot\si{\meter\per\second}}
=\frac{\si{\newton}\cdot\si{\second}}{\si{\coulomb}\cdot\si{\meter}}
=\frac{\si{\newton}}{\si{\ampere}\cdot\si{\meter}}.
\]
A última forma, $\si{\newton\per\ampere\per\meter}$, decorre diretamente de
$F=BIL$ --- e é a mais prática para conferir contas com condutores.\\
\textbf{O tesla é uma unidade grande:} $\SI{1}{\tesla}$ exerce
$\SI{1}{\newton}$ sobre um metro de fio conduzindo $\SI{1}{\ampere}$.}
\end{questao}

\blocoaplicacao

\begin{questao}[2]{}
Um fio de $\SI{10}{\centi\meter}$ conduz $\SI{2}{\ampere}$ em um campo de
$\SI{0.2}{\tesla}$, formando $\ang{30}$ com ele. Determine a força.
\resposta{$F=\SI{20}{\milli\newton}$}
\resolucao{
\[
F=BIL\sin\ang{30}=(0{,}2)(2)(0{,}10)(0{,}5)=\pot{2}{-2}\,\si{\newton}.
\]
\textbf{Interpretação do seno:} apenas a componente de $\vec\ell$
\emph{perpendicular} a $\vec B$ contribui. Equivalentemente, o comprimento
\emph{efetivo} é $L\sin\theta=\SI{5}{\centi\meter}$.\\
\textbf{Direção da força:} perpendicular ao plano definido pelo fio e pelo
campo --- e \emph{não} perpendicular apenas ao fio.}
\end{questao}

\begin{questao}[2]{}
Um próton entra perpendicularmente em um campo de $\SI{0.5}{\tesla}$ com
$v=\pot{2}{6}\,\si{\meter\per\second}$.
\begin{itensq}
\item Determine o raio da trajetória.
\item Determine o período.
\end{itensq}
\dados{$e=\pot{1.6}{-19}\,\si{\coulomb}$; $m_p=\pot{1.67}{-27}\,\si{\kilogram}$}
\resposta{$r=\SI{4.18}{\centi\meter}$; $T=\SI{131}{\nano\second}$}
\resolucao{(a) A força magnética é centrípeta:
\[
qvB=\frac{mv^{2}}{r}
\;\Longrightarrow\;
r=\frac{mv}{qB}=\frac{(\pot{1.67}{-27})(\pot{2}{6})}{(\pot{1.6}{-19})(0{,}5)}
=\SI{0.0418}{\meter}.
\]
(b)
\[
T=\frac{2\pi r}{v}=\frac{2\pi m}{qB}
=\frac{2\pi(\pot{1.67}{-27})}{(\pot{1.6}{-19})(0{,}5)}
=\pot{1.31}{-7}\,\si{\second}.
\]
\textbf{Fato notável:} o período \textbf{não depende da velocidade} nem do raio.
Partículas rápidas descrevem círculos maiores, mas levam o mesmo tempo para
completá-los.\\
É esse fato que torna possível o \textbf{cíclotron}: a frequência de aceleração
pode ser fixa, embora a partícula ganhe energia a cada volta.}
\end{questao}

\begin{questao}[2]{}
Um elétron e um próton, com a mesma velocidade, entram no mesmo campo
magnético perpendicularmente. Compare seus raios e períodos.
\resposta{$r_p/r_e=1836$; $T_p/T_e=1836$}
\resolucao{Como $r=mv/(qB)$ e $T=2\pi m/(qB)$, ambos são
\emph{proporcionais à massa} (as cargas têm o mesmo módulo):
\[
\frac{r_p}{r_e}=\frac{T_p}{T_e}=\frac{m_p}{m_e}
=\frac{\pot{1.67}{-27}}{\pot{9.11}{-31}}=1836.
\]
\textbf{Sentidos opostos:} tendo cargas de sinais contrários, as duas partículas
curvam-se para lados opostos --- é assim que se identifica o sinal da carga em
câmaras de detecção.\\
\textbf{Consequência experimental:} para conter elétrons e prótons no mesmo
aparelho, o raio do próton é quase dois mil vezes maior. Aceleradores de prótons
são, por isso, muito maiores que os de elétrons para a mesma energia.}
\end{questao}

\begin{questao}[2]{}
Dois fios paralelos longos, separados por $\SI{5}{\centi\meter}$, conduzem
$\SI{10}{\ampere}$ cada. Determine a força por unidade de comprimento.
\dados{$\mu_0=4\pi\times10^{-7}\,\si{\tesla\meter\per\ampere}$}
\resposta{$F/L=\SI{0.4}{\milli\newton\per\meter}$}
\resolucao{Um fio cria campo $B=\dfrac{\mu_0I_1}{2\pi d}$ na posição do outro,
que sofre $F=BI_2L$:
\[
\frac{F}{L}=\frac{\mu_0I_1I_2}{2\pi d}
=\pot{2}{-7}\frac{(10)(10)}{0{,}05}=\pot{4}{-4}\,\si{\newton\per\meter}.
\]
\textbf{Sentido:} correntes de \emph{mesmo} sentido \textbf{atraem-se};
de sentidos opostos, \textbf{repelem-se}.\\
\textbf{Cuidado com a intuição elétrica:} cargas de mesmo sinal se repelem, mas
correntes de mesmo sentido se atraem. As duas situações não são análogas.}
\end{questao}

\begin{questao}[2]{}
Uma espira retangular de área $\SI{100}{\centi\meter\squared}$ com $50$ voltas
conduz $\SI{2}{\ampere}$ em um campo de $\SI{0.3}{\tesla}$.
\begin{itensq}
\item Determine o momento de dipolo magnético.
\item Determine o torque máximo.
\end{itensq}
\resposta{$m=\SI{1}{\ampere\meter\squared}$; $\tau=\SI{0.3}{\newton\meter}$}
\resolucao{(a) $A=\SI{100}{\centi\meter\squared}=\pot{1}{-2}\,\si{\meter\squared}$:
\[
m=NIA=50(2)(\pot{1}{-2})=\SI{1}{\ampere\meter\squared}.
\]
(b)
\[
\tau=mB\sin\theta
\;\Longrightarrow\;
\tau_{\max}=mB=(1)(0{,}3)=\SI{0.3}{\newton\meter}.
\]
\textbf{Quando ocorre o máximo:} com o plano da espira \emph{paralelo} ao campo,
isto é, $\vec m$ perpendicular a $\vec B$. O torque é \textbf{nulo} quando
$\vec m$ está alinhado com $\vec B$ --- posição de equilíbrio.\\
\textbf{Analogia útil:} a espira comporta-se como uma agulha de bússola, e
$\vec m$ faz o papel do "norte" dela.}
\end{questao}

\begin{questao}[2]{}
Explique por que a força magnética não altera o módulo da velocidade de uma
partícula carregada.
\resposta{Ela é sempre perpendicular a $\vec v$, logo não realiza trabalho}
\resolucao{A potência entregue por uma força é $P=\vec F\cdot\vec v$. Como
\[
\vec F=q\,\vec v\times\vec B
\;\Longrightarrow\;
\vec F\perp\vec v
\;\Longrightarrow\;
\vec F\cdot\vec v=0.
\]
Sem potência, não há variação de energia cinética --- e portanto o \emph{módulo}
da velocidade permanece constante.\\
\textbf{O que a força magnética faz:} muda a \emph{direção} do movimento, e só
isso. Ela curva a trajetória sem acelerar nem frear.\\
\textbf{Consequência prática:} um campo magnético sozinho jamais acelera
partículas. Aceleradores usam campos \emph{elétricos} para ganhar energia e
campos magnéticos apenas para \emph{guiar} o feixe.}
\end{questao}

\begin{questao}[2]{}
Um elétron é abandonado \textbf{em repouso} em uma região onde coexistem
$E=\SI{500}{\volt\per\meter}$ e $B=\SI{0.02}{\tesla}$, perpendiculares entre
si.
\begin{itensq}
\item Determine a força e a aceleração no instante inicial.
\item Descreva qualitativamente o movimento subsequente.
\end{itensq}
\dados{$e=\pot{1.6}{-19}\,\si{\coulomb}$; $m_e=\pot{9.11}{-31}\,\si{\kilogram}$}
\resposta{$F=\SI{8e-17}{\newton}$ (só elétrica);
$a=\pot{8.8}{13}\,\si{\meter\per\second\squared}$}
\resolucao{(a) A força total é a \textbf{força de Lorentz}:
\[
\vec F=q\left(\vec E+\vec v\times\vec B\right).
\]
Com $\vec v=\vec 0$, a parcela magnética \textbf{se anula} --- resta apenas a
elétrica:
\[
F=eE=(\pot{1.6}{-19})(500)=\pot{8}{-17}\,\si{\newton};
\qquad
a=\frac{F}{m_e}=\pot{8.8}{13}\,\si{\meter\per\second\squared}.
\]
(b) \textbf{Movimento subsequente.} Assim que o elétron adquire velocidade, a
parcela magnética \emph{deixa de ser nula} e passa a encurvar a trajetória. A
força elétrica continua acelerando; a magnética continua desviando.\\
O resultado é uma trajetória \textbf{cicloidal}, com deriva média
\[
v_d=\frac{E}{B}=\frac{500}{0{,}02}=\pot{2.5}{4}\,\si{\meter\per\second},
\]
perpendicular a ambos os campos.\\
\textbf{Ponto essencial:} a força magnética \emph{depende da velocidade} --- ela
não age sobre cargas paradas. Já a elétrica age sempre. É essa assimetria que
faz o movimento começar de um jeito e evoluir para outro.}
\end{questao}

\begin{questao}[2]{}
Um fio em forma de semicircunferência de raio $R$ conduz corrente $I$ em um
campo uniforme $\vec B$ perpendicular ao plano do semicírculo. Determine a
força resultante.
\resposta{$F=2RIB$, como se o fio fosse retilíneo de comprimento $2R$}
\resolucao{Para um fio de forma qualquer em campo \textbf{uniforme},
\[
\vec F=I\left(\int\mathrm{d}\vec\ell\right)\times\vec B
=I\,\vec L_{ef}\times\vec B,
\]
onde $\vec L_{ef}$ é o vetor que liga o \textbf{início ao fim} do fio --- a
\emph{corda}, não o comprimento percorrido.\\
Para a semicircunferência, a corda vale $2R$:
\[
F=2RIB.
\]
\textbf{Poder do resultado:} a forma do fio é irrelevante em campo uniforme.
Um fio retorcido, ondulado ou em zigue-zague sofre a mesma força que o segmento
reto entre suas extremidades.\\
\textbf{Corolário importante:} em uma \emph{espira fechada}, a corda é nula ---
logo a força resultante sobre qualquer espira em campo uniforme é
\textbf{zero}. Só há torque.}
\end{questao}

\begin{questao}[2]{}
Uma balança de corrente tem um fio horizontal de $\SI{20}{\centi\meter}$ em um
campo vertical de $\SI{0.4}{\tesla}$. Determine a corrente necessária para
equilibrar uma massa de $\SI{2}{\gram}$.
\dados{$g=\SI{9.8}{\meter\per\second\squared}$}
\resposta{$I=\SI{0.245}{\ampere}$}
\resolucao{Equilíbrio entre peso e força magnética:
\[
BIL=mg
\;\Longrightarrow\;
I=\frac{mg}{BL}=\frac{(\pot{2}{-3})(9{,}8)}{(0{,}4)(0{,}20)}
=\SI{0.245}{\ampere}.
\]
\textbf{Uso do arranjo:} conhecendo $m$, $L$ e $I$, obtém-se $B$ com boa
exatidão --- é o método clássico de calibração absoluta de campo magnético, sem
depender de sensores.\\
\textbf{Inversamente}, conhecendo $B$, mede-se corrente. Foi assim que o ampère
foi definido durante décadas: pela força entre condutores.}
\end{questao}

\begin{questao}[2]{}
Determine a velocidade de uma partícula que atravessa sem desvio uma região com
$E=\SI{2e4}{\volt\per\meter}$ e $B=\SI{0.1}{\tesla}$ perpendiculares entre si
e à velocidade.
\resposta{$v=\pot{2}{5}\,\si{\meter\per\second}$}
\resolucao{Para não haver desvio, as forças elétrica e magnética devem se
cancelar:
\[
qE=qvB
\;\Longrightarrow\;
v=\frac{E}{B}=\frac{\pot{2}{4}}{0{,}1}=\pot{2}{5}\,\si{\meter\per\second}.
\]
\textbf{Duas propriedades notáveis:}
\begin{itemize}
\item A condição \textbf{não depende da carga} --- $q$ cancela. Íons e elétrons
com a mesma velocidade passam igualmente.
\item Também \textbf{não depende da massa}. É pura seleção de \emph{velocidade}.
\end{itemize}
\textbf{Consequência:} partículas mais rápidas são desviadas pela força
magnética (maior), e mais lentas pela elétrica. Apenas a velocidade $E/B$
atravessa em linha reta.}
\end{questao}

\begin{questao}[2]{}
Uma espira quadrada de lado $\SI{10}{\centi\meter}$ conduz $\SI{5}{\ampere}$
em um campo uniforme de $\SI{0.2}{\tesla}$ paralelo ao seu plano.
\begin{itensq}
\item Determine a força resultante.
\item Determine o torque.
\end{itensq}
\resposta{$\vec F=\vec 0$; $\tau=\SI{10}{\milli\newton\meter}$}
\resolucao{(a) \textbf{Força resultante nula.} Em campo uniforme, os quatro
lados sofrem forças que se cancelam duas a duas --- lados opostos conduzem em
sentidos contrários.\\
(b) \textbf{Mas o torque não é nulo.} As forças nos dois lados perpendiculares
ao campo formam um \emph{binário}:
\[
\tau=mB\sin\ang{90}=IAB=(5)(\pot{1}{-2})(0{,}2)
=\pot{1}{-2}\,\si{\newton\meter}.
\]
\textbf{Situação de torque máximo:} o campo está no plano da espira, ou seja,
$\vec m$ (perpendicular ao plano) está a $\ang{90}$ de $\vec B$.\\
\textbf{Este é o princípio do motor elétrico} --- e também sua limitação: o
torque cai a zero quando a espira gira até o alinhamento, o que exige comutador
para inverter a corrente e manter a rotação.}
\end{questao}

\blocoaprofundamento

\begin{questao}[3]{}
Uma partícula de carga $q$ e massa $m$ entra perpendicularmente em um campo
uniforme $B$ com velocidade $v$.
\begin{itensq}
\item Deduza o raio e o período da trajetória.
\item Determine a frequência angular (frequência de cíclotron).
\item Explique por que o período independe de $v$.
\end{itensq}
\resposta{$r=\dfrac{mv}{qB}$; $T=\dfrac{2\pi m}{qB}$;
$\omega=\dfrac{qB}{m}$}
\resolucao{(a) A força magnética é perpendicular a $\vec v$ e de módulo
constante --- exatamente as condições de um movimento \textbf{circular
uniforme}. Igualando à resultante centrípeta:
\[
qvB=\frac{mv^{2}}{r}
\;\Longrightarrow\;
r=\frac{mv}{qB}.
\]
\[
T=\frac{2\pi r}{v}=\frac{2\pi m}{qB}.
\]
(b)
\[
\omega=\frac{2\pi}{T}=\frac{qB}{m},
\]
a \textbf{frequência angular de cíclotron}.\\
(c) \textbf{Por que $T$ independe de $v$.} Dobrar a velocidade dobra a força
magnética \emph{e} dobra o raio necessário --- os dois efeitos se compensam
exatamente. A partícula percorre um círculo duas vezes maior a velocidade duas
vezes maior, no mesmo tempo.\\
Algebricamente, $v$ aparece em $r$ e é cancelado por $v$ em $T=2\pi r/v$.\\
\textbf{Consequência:} o cíclotron pode operar com frequência \emph{fixa}.
\emph{(Isso vale enquanto $v\ll c$; em regime relativístico a massa efetiva
cresce, $T$ aumenta e é preciso modular a frequência --- o síncrotron.)}}
\end{questao}

\begin{questao}[3]{}
Uma partícula entra em um campo uniforme formando $\ang{30}$ com $\vec B$, com
$v=\pot{1}{6}\,\si{\meter\per\second}$ e $B=\SI{0.1}{\tesla}$ (próton).
\begin{itensq}
\item Decomponha a velocidade e descreva a trajetória.
\item Determine o raio e o passo da hélice.
\end{itensq}
\resposta{Hélice com $r=\SI{5.2}{\centi\meter}$ e passo de
$\SI{0.57}{\meter}$}
\resolucao{(a) \textbf{Decomposição:}
\[
v_\perp=v\sin\ang{30}=\pot{5}{5}\,\si{\meter\per\second};
\qquad
v_\parallel=v\cos\ang{30}=\pot{8.66}{5}\,\si{\meter\per\second}.
\]
A componente \textbf{perpendicular} sofre força e produz movimento circular. A
componente \textbf{paralela} não sofre força alguma ($\vec v\times\vec B=0$
para ela) e produz movimento retilíneo uniforme.\\
A composição é uma \textbf{hélice} de eixo paralelo a $\vec B$.\\
(b)
\[
r=\frac{mv_\perp}{qB}
=\frac{(\pot{1.67}{-27})(\pot{5}{5})}{(\pot{1.6}{-19})(0{,}1)}
=\SI{0.052}{\meter};
\]
\[
T=\frac{2\pi m}{qB}=\pot{6.56}{-7}\,\si{\second};
\qquad
\text{passo}=v_\parallel T=(\pot{8.66}{5})(\pot{6.56}{-7})=\SI{0.568}{\meter}.
\]
\textbf{Aplicação:} é essa trajetória que confina partículas nos cinturões de
Van Allen e em espelhos magnéticos --- as partículas espiralam ao longo das
linhas de campo, e o passo diminui onde o campo se intensifica, até a reversão
do movimento.}
\end{questao}

\begin{questao}[3]{}
Um espectrômetro de massa usa um seletor de velocidades ($v=\pot{1}{5}\,
\si{\meter\per\second}$) seguido de um campo de $\SI{0.1}{\tesla}$. Íons de
neônio de massas $20\,u$ e $22\,u$, todos com carga $+e$, entram no segundo
campo.
\begin{itensq}
\item Determine o raio de cada trajetória.
\item Determine a separação entre os pontos de chegada no detector.
\end{itensq}
\dados{$u=\pot{1.66}{-27}\,\si{\kilogram}$}
\resposta{$r_{20}=\SI{20.8}{\centi\meter}$, $r_{22}=\SI{22.8}{\centi\meter}$;
separação de $\SI{4.2}{\centi\meter}$}
\resolucao{(a) Com $r=mv/(qB)$ e velocidade já selecionada:
\[
r_{20}=\frac{20(\pot{1.66}{-27})(\pot{1}{5})}{(\pot{1.6}{-19})(0{,}1)}
=\SI{0.2075}{\meter};
\qquad
r_{22}=\SI{0.2283}{\meter}.
\]
(b) Após meia volta, cada íon atinge o detector a uma distância $2r$ do ponto
de entrada:
\[
\Delta=2\left(r_{22}-r_{20}\right)=2(0{,}0208)=\SI{4.15}{\centi\meter}.
\]
\textbf{Por que o seletor é indispensável:} sem ele, íons de massas diferentes
mas velocidades diferentes poderiam produzir o \emph{mesmo} raio --- pois
$r\propto mv$. O seletor fixa $v$, tornando $r$ função apenas de $m/q$.\\
\textbf{Resolução do instrumento:} a separação relativa é
$\Delta r/r=\Delta m/m=10\%$ para esses isótopos --- larga. Separar isótopos
com $\Delta m/m$ de $0{,}1\%$ exigiria feixe muito bem colimado e detector de
alta resolução espacial.}
\end{questao}

\begin{questao}[3]{}
A definição do ampère baseava-se na força entre condutores paralelos.
\begin{itensq}
\item Determine a força por metro entre dois fios distantes
      $\SI{1}{\meter}$, conduzindo $\SI{1}{\ampere}$ cada.
\item Explique o papel dessa relação na definição da unidade.
\item Comente a mudança recente.
\end{itensq}
\resposta{$F/L=\pot{2}{-7}\,\si{\newton\per\meter}$}
\resolucao{(a)
\[
\frac{F}{L}=\frac{\mu_0I^{2}}{2\pi d}
=\frac{(4\pi\times10^{-7})(1)}{2\pi(1)}=\pot{2}{-7}\,\si{\newton\per\meter}.
\]
(b) \textbf{Definição antiga (até 2019):} o ampère era \emph{definido} como a
corrente que, em dois condutores paralelos infinitos separados por
$\SI{1}{\meter}$ no vácuo, produzisse exatamente
$\pot{2}{-7}\,\si{\newton}$ por metro.\\
Essa definição \textbf{fixava $\mu_0$} em $4\pi\times10^{-7}$ por construção ---
não era um valor medido, mas uma consequência da escolha.\\
(c) \textbf{Desde 2019}, o ampère é definido fixando-se a \emph{carga
elementar} em $e=\pot{1.602176634}{-19}\,\si{\coulomb}$ exatamente.\\
\textbf{Consequência da mudança:} $\mu_0$ deixou de ser exato e passou a ser uma
grandeza \textbf{medida}, com incerteza experimental --- embora ainda
igual a $4\pi\times10^{-7}$ dentro de partes por bilhão.\\
A troca reflete uma mudança de filosofia: definir unidades por
\emph{constantes fundamentais} em vez de por experimentos difíceis de
reproduzir.}
\end{questao}

\begin{questao}[3]{}
Um sensor Hall de espessura $\SI{0.1}{\milli\meter}$ conduz
$\SI{1}{\ampere}$ em campo de $\SI{0.5}{\tesla}$.
\begin{itensq}
\item Determine a tensão Hall para cobre ($n=\pot{8.5}{28}\,\si{\per\meter\cubed}$).
\item Repita para um semicondutor com $n=\pot{1}{22}\,\si{\per\meter\cubed}$.
\item Explique a diferença.
\end{itensq}
\resposta{$\SI{0.37}{\micro\volt}$ no cobre; $\SI{3.1}{\volt}$ no
semicondutor}
\resolucao{(a) A tensão Hall vale
\[
V_H=\frac{IB}{nqt}
=\frac{(1)(0{,}5)}{(\pot{8.5}{28})(\pot{1.6}{-19})(\pot{1}{-4})}
=\pot{3.68}{-7}\,\si{\volt}.
\]
(b) Com $n=\pot{1}{22}$:
\[
V_H=\frac{0{,}5}{(\pot{1}{22})(\pot{1.6}{-19})(\pot{1}{-4})}=\SI{3.1}{\volt}.
\]
\textbf{Diferença de sete ordens de grandeza.}\\
(c) \textbf{Explicação.} A tensão Hall é inversamente proporcional a $n$. Metais
têm um portador por átomo --- densidade altíssima --- e cada portador precisa se
mover \emph{devagar} para produzir a mesma corrente. Velocidade de deriva baixa
significa força magnética pequena e, portanto, separação de cargas mínima.\\
Semicondutores têm densidade de portadores milhões de vezes menor; eles se movem
muito mais rápido, e o efeito Hall se torna macroscópico.\\
\textbf{Consequência tecnológica:} \emph{todos} os sensores Hall comerciais são
de semicondutor. Um sensor de cobre exigiria amplificação de nanovolts,
inviável na presença de ruído térmico e de tensões termoelétricas.}
\end{questao}

\begin{questao}[3]{}
Um motor elementar tem espira de $200$ voltas, área
$\SI{50}{\centi\meter\squared}$, corrente de $\SI{3}{\ampere}$ e campo de
$\SI{0.4}{\tesla}$.
\begin{itensq}
\item Determine o torque máximo.
\item Determine a potência mecânica a $\SI{1200}{rpm}$, supondo torque médio
      igual a $2/\pi$ do máximo.
\end{itensq}
\resposta{$\tau_{\max}=\SI{1.2}{\newton\meter}$;
$P\approx\SI{96}{\watt}$}
\resolucao{(a)
\[
m=NIA=200(3)(\pot{5}{-3})=\SI{3}{\ampere\meter\squared};
\qquad
\tau_{\max}=mB=3(0{,}4)=\SI{1.2}{\newton\meter}.
\]
\emph{(Atenção à conversão: $\SI{50}{\centi\meter\squared}
=\pot{5}{-3}\,\si{\meter\squared}$, e não $\pot{5}{-2}$.)}\\
(b) O torque varia como $\sin\theta$ ao longo da volta, e seu valor médio sobre
meio ciclo é $\dfrac{2}{\pi}\tau_{\max}=\SI{0.764}{\newton\meter}$.
\[
\omega=1200\ \text{rpm}=\frac{1200(2\pi)}{60}=\SI{125.7}{\radian\per\second};
\]
\[
P=\tau_{\text{méd}}\,\omega=0{,}764(125{,}7)=\SI{96}{\watt}.
\]
\textbf{Por que o comutador é necessário:} sem inverter a corrente a cada meia
volta, o torque mudaria de sinal e a espira apenas oscilaria em torno da posição
de alinhamento.\\
\textbf{Por que motores reais usam muitas bobinas:} com uma só, o torque
\emph{pulsa} entre zero e o máximo. Várias bobinas defasadas produzem torque
quase constante --- e é isso que se busca em qualquer máquina prática.}
\end{questao}

\begin{questao}[3]{}
Uma espira de corrente é colocada em um campo magnético \textbf{não uniforme},
mais intenso de um lado.
\begin{itensq}
\item Determine se há força resultante.
\item Determine o sentido dessa força.
\item Relacione com o comportamento de ímãs.
\end{itensq}
\resposta{Há força resultante, dirigida para a região de campo mais intenso (se
$\vec m$ estiver alinhado com $\vec B$)}
\resolucao{(a) \textbf{Sim.} Em campo \emph{uniforme}, as forças nos lados
opostos da espira se cancelam exatamente e resta apenas torque. Em campo
\emph{não uniforme}, os lados estão em campos de módulos diferentes --- e o
cancelamento é incompleto.\\
(b) \textbf{Sentido.} A força resultante é dada pelo gradiente:
\[
\vec F=\nabla\left(\vec m\cdot\vec B\right).
\]
Se $\vec m$ estiver \emph{alinhado} com $\vec B$, a espira é atraída para a
região de campo mais intenso. Se estiver \emph{antialinhado}, é repelida.\\
(c) \textbf{Relação com ímãs.} Um ímã permanente é, essencialmente, um conjunto
de espiras microscópicas (correntes de magnetização) com momentos alinhados.
Daí:
\begin{itemize}
\item \textbf{Dois ímãs se atraem ou repelem} conforme a orientação relativa dos
momentos --- exatamente o comportamento da espira.
\item \textbf{Um ímã atrai ferro não magnetizado} porque primeiro o magnetiza
(alinhando seus domínios com o campo aplicado) e depois o puxa para a região de
campo mais forte. A atração é sempre no sentido do gradiente.
\item \textbf{Em campo uniforme, um ímã não é atraído} --- ele apenas gira até
se alinhar. É por isso que a agulha de uma bússola aponta, mas não é puxada em
direção ao norte.
\end{itemize}}
\end{questao}

\begin{questao}[3]{}
Compare a força entre dois fios paralelos conduzindo correntes iguais no mesmo
sentido e em sentidos opostos.
\begin{itensq}
\item Determine o sentido em cada caso.
\item Determine o módulo.
\item Explique a diferença em termos de campo.
\end{itensq}
\resposta{Mesmo sentido: atração; opostos: repulsão. Módulos iguais}
\resolucao{(a) e (b) O \textbf{módulo} é o mesmo nos dois casos:
\[
\frac{F}{L}=\frac{\mu_0I_1I_2}{2\pi d}.
\]
O que muda é o \emph{sentido}: correntes paralelas \textbf{atraem-se},
antiparalelas \textbf{repelem-se}.\\
(c) \textbf{Explicação pelo campo.} Considere o fio 2 imerso no campo do fio 1.
\begin{itemize}
\item \textbf{Mesmo sentido:} na posição do fio 2, o campo de 1 aponta em certo
sentido; aplicando $\vec F=I\vec\ell\times\vec B$, a força aponta \emph{para} o
fio 1.
\item \textbf{Sentidos opostos:} $\vec\ell$ inverte, e a força inverte junto.
\end{itemize}
\textbf{Leitura pelas linhas de campo.} Com correntes paralelas, o campo se
\emph{cancela} entre os fios e se \emph{reforça} fora --- as linhas envolvem o
par, e a "tensão" ao longo delas puxa os fios um contra o outro. Com correntes
opostas, ocorre o inverso: campo intenso entre os fios, que os afasta.\\
\textbf{Consequência prática:} em curtos-circuitos, barramentos paralelos
sofrem forças enormes. Com $\SI{20}{\kilo\ampere}$ a $\SI{10}{\centi\meter}$,
a força chega a $\SI{800}{\newton\per\meter}$ --- razão de existirem
espaçadores mecânicos robustos em painéis de potência.}
\end{questao}

\begin{questao}[3]{}
Um próton atravessa uma região com $E=\SI{2e4}{\volt\per\meter}$ e
$B=\SI{0.1}{\tesla}$, perpendiculares entre si e à velocidade, dispostos de
modo que as forças se oponham.
\begin{itensq}
\item Determine as duas parcelas da força de Lorentz para
      $v=\pot{1}{5}$, $\pot{2}{5}$ e $\pot{4}{5}\,\si{\meter\per\second}$.
\item Identifique qual parcela domina em cada caso.
\item Interprete o resultado.
\end{itensq}
\resposta{Domina a elétrica em $\pot{1}{5}$, equilíbrio em $\pot{2}{5}$ e a
magnética em $\pot{4}{5}\,\si{\meter\per\second}$}
\resolucao{(a) A parcela elétrica \textbf{não depende de $v$}:
\[
F_E=qE=(\pot{1.6}{-19})(\pot{2}{4})=\pot{3.2}{-15}\,\si{\newton}.
\]
A magnética é \textbf{proporcional a $v$}: $F_B=qvB$.
\begin{center}
\begin{tabular}{cccl}
\hline
$v$ ($\si{\meter\per\second}$) & $F_E$ ($\si{\newton}$)
 & $F_B$ ($\si{\newton}$) & Resultante\\
\hline
$\pot{1}{5}$ & $\pot{3.2}{-15}$ & $\pot{1.6}{-15}$ & elétrica domina\\
$\pot{2}{5}$ & $\pot{3.2}{-15}$ & $\pot{3.2}{-15}$ & \textbf{nula}\\
$\pot{4}{5}$ & $\pot{3.2}{-15}$ & $\pot{6.4}{-15}$ & magnética domina\\
\hline
\end{tabular}
\end{center}
(b) e (c) \textbf{Interpretação.} O equilíbrio ocorre em
\[
v=\frac{E}{B}=\frac{\pot{2}{4}}{0{,}1}=\pot{2}{5}\,\si{\meter\per\second},
\]
e apenas nessa velocidade a partícula segue reta.\\
\textbf{As três situações da tabela descrevem o mesmo dispositivo}: partículas
lentas são empurradas para um lado pela força elétrica, rápidas para o outro
pela magnética, e só as de velocidade $E/B$ atravessam sem desvio. É a
\emph{seleção de velocidades}, vista agora como competição entre as duas
parcelas de $\vec F=q(\vec E+\vec v\times\vec B)$.\\
\textbf{Observe a estrutura:} uma parcela constante e outra linear em $v$. A
igualdade só pode ocorrer em \emph{um} valor de velocidade --- e é isso que dá
seletividade ao instrumento.}
\end{questao}

\begin{questao}[3]{}
Uma partícula carregada move-se em uma região com campos $\vec E$ e $\vec B$
perpendiculares entre si, mas com velocidade \emph{diferente} de $E/B$.
\begin{itensq}
\item Descreva qualitativamente o movimento.
\item Determine a velocidade média de deriva.
\item Indique uma aplicação.
\end{itensq}
\resposta{Movimento cicloidal, com deriva média
$\vec v_d=\dfrac{\vec E\times\vec B}{B^{2}}$}
\resolucao{(a) \textbf{Movimento.} A partícula descreve uma trajetória
\textbf{cicloidal}: um movimento circular (devido a $\vec B$) superposto a um
deslocamento uniforme. Se ela parte do repouso, a curva é uma cicloide clássica
--- como um ponto na borda de uma roda que rola.\\
(b) \textbf{Velocidade de deriva.} O centro do círculo desloca-se com
\[
\vec v_d=\frac{\vec E\times\vec B}{B^{2}},
\qquad
|v_d|=\frac{E}{B},
\]
perpendicular a \emph{ambos} os campos.\\
\textbf{Duas propriedades notáveis:} a deriva \textbf{não depende da carga}
(íons e elétrons derivam no mesmo sentido) nem da \textbf{massa}.\\
(c) \textbf{Aplicações.}
\begin{itemize}
\item \textbf{Magnetrons} (fornos de micro-ondas): elétrons em campos cruzados
descrevem trajetórias cicloidais que excitam cavidades ressonantes.
\item \textbf{Confinamento de plasma:} a deriva $\vec E\times\vec B$ é um dos
mecanismos que governam o transporte em tokamaks --- e como ela não separa
cargas, não gera corrente própria.
\item \textbf{Magnetosfera terrestre:} explica parte do movimento de partículas
aprisionadas.
\end{itemize}}
\end{questao}

\begin{questao}[3]{}
Um condutor de $\SI{50}{\centi\meter}$ e $\SI{100}{\gram}$ está apoiado sobre
dois trilhos horizontais em campo vertical de $\SI{0.8}{\tesla}$. O
coeficiente de atrito estático é $0{,}25$.
\begin{itensq}
\item Determine a corrente mínima para iniciar o movimento.
\item Determine a aceleração inicial se a corrente for o dobro, com atrito
      cinético de $0{,}20$.
\end{itensq}
\resposta{$I=\SI{0.613}{\ampere}$; $a\approx\SI{2.9}{\meter\per\second\squared}$}
\resolucao{(a) A força magnética deve vencer o atrito estático:
\[
BIL=\mu_e mg
\;\Longrightarrow\;
I=\frac{0{,}25(0{,}100)(9{,}8)}{(0{,}8)(0{,}50)}=\SI{0.613}{\ampere}.
\]
(b) Com $I=\SI{1.226}{\ampere}$:
\[
F=BIL=(0{,}8)(1{,}226)(0{,}50)=\SI{0.490}{\newton};
\]
\[
f_c=\mu_c mg=0{,}20(0{,}100)(9{,}8)=\SI{0.196}{\newton};
\]
\[
a=\frac{0{,}490-0{,}196}{0{,}100}=\SI{2.94}{\meter\per\second\squared}.
\]
\textbf{Este é o princípio do canhão eletromagnético (\emph{railgun})} --- e
também do motor linear. \\
\textbf{Ressalva importante:} conforme o condutor acelera, ele passa a induzir
f.e.m. contrária ($\varepsilon=BLv$), que reduz a corrente. A aceleração
calculada vale apenas no \emph{instante inicial}; o movimento tende a uma
velocidade-limite em que a f.e.m. induzida equilibra a fonte.}
\end{questao}

\blocodesafio

\begin{questao}[4]{}
\textbf{(Demonstração --- trabalho nulo.)} Prove que a força magnética nunca
realiza trabalho sobre uma carga puntiforme.
\begin{itensq}
\item Demonstre formalmente.
\item Concilie com o fato de que motores elétricos realizam trabalho.
\item Explique o que ocorre com a energia em um motor.
\end{itensq}
\resposta{$\vec F\cdot\vec v=0$ sempre; o trabalho no motor vem da fonte, não do
campo}
\resolucao{(a) A potência é
\[
P=\vec F\cdot\vec v=q\left(\vec v\times\vec B\right)\cdot\vec v.
\]
O produto misto $(\vec v\times\vec B)\cdot\vec v$ é \textbf{identicamente
nulo}, pois $\vec v\times\vec B$ é perpendicular a $\vec v$ por construção. Logo
$P=0$ e $\Delta E_c=0$: o \emph{módulo} da velocidade nunca muda.\\
(b) \textbf{O aparente paradoxo do motor.} Um motor claramente realiza trabalho
--- ele levanta cargas e move veículos. Como conciliar?\\
\textbf{Resolução:} a força magnética atua sobre os \emph{portadores} do
condutor, e é perpendicular à velocidade \emph{deles}. Mas o fio se move, e a
velocidade total de um portador é a soma da deriva ao longo do fio com a
velocidade do fio.\\
A força magnética \emph{redistribui} energia: ela empurra o fio (realizando
trabalho positivo sobre ele) e, simultaneamente, freia a deriva dos portadores
(trabalho negativo de mesmo módulo). A soma continua sendo zero.\\
(c) \textbf{De onde vem a energia.} O freio à deriva manifesta-se como
\textbf{f.e.m. contraeletromotriz}, e é a \emph{fonte} que precisa vencê-la:
\[
P_{fonte}=\varepsilon_{fem}\,I=P_{\text{mecânica}}+P_{Joule}.
\]
\textbf{O campo magnético é um intermediário, não uma fonte.} Ele transfere
energia da fonte elétrica para o eixo, sem fornecer nem armazenar nada
líquido.\\
\emph{(É o mesmo papel de uma polia em mecânica: ela redireciona força sem criar
energia.)}}
\end{questao}

\begin{questao}[4]{}
\textbf{(Dedução do torque.)} Deduza a expressão do torque sobre uma espira
retangular em campo uniforme.
\begin{itensq}
\item Analise as forças em cada lado.
\item Obtenha $\tau=NIAB\sin\theta$.
\item Generalize para espira de forma qualquer.
\end{itensq}
\resposta{$\vec\tau=\vec m\times\vec B$, com $\vec m=NI\vec A$}
\resolucao{(a) Seja a espira de lados $a$ e $b$, com $\vec B$ no plano
horizontal e o eixo de rotação vertical. Chamando $\theta$ o ângulo entre
$\vec B$ e a \emph{normal} à espira:
\begin{itemize}
\item Os lados de comprimento $a$ \emph{paralelos} ao eixo sofrem forças
$F=BIa$, perpendiculares tanto ao fio quanto ao campo, em sentidos opostos.
Elas formam um \textbf{binário}.
\item Os lados de comprimento $b$ sofrem forças que se cancelam e cuja linha de
ação passa pelo eixo --- não produzem torque.
\end{itemize}
(b) O braço de alavanca de cada força do binário é
$\dfrac{b}{2}\sin\theta$:
\[
\tau=2\left(BIa\right)\frac{b}{2}\sin\theta=BI(ab)\sin\theta=BIA\sin\theta.
\]
Com $N$ espiras, $\tau=NIAB\sin\theta$.\\
(c) \textbf{Generalização.} Definindo o momento magnético
$\vec m=NI\vec A$ (com $\vec A$ normal à espira, sentido dado pela regra da mão
direita):
\[
\boxed{\vec\tau=\vec m\times\vec B}
\]
Essa forma vale para espira de \textbf{qualquer} formato --- basta decompô-la em
retângulos infinitesimais, cujas correntes internas se cancelam aos pares.\\
\textbf{Analogia formal com o dipolo elétrico:} $\vec\tau=\vec p\times\vec E$.
As duas expressões são idênticas em estrutura, e por isso a espira é
literalmente um \emph{dipolo magnético}.}
\end{questao}

\begin{questao}[4]{}
\textbf{(Projeto de espectrômetro.)} Projete um espectrômetro capaz de separar
isótopos com $\Delta m/m=1\%$.
\begin{itensq}
\item Determine a separação espacial em função dos parâmetros.
\item Estabeleça o requisito de colimação do feixe.
\item Discuta as limitações práticas.
\end{itensq}
\resposta{$\Delta r/r=\Delta m/m$; a resolução exige feixe com dispersão angular
e de velocidade menores que $1\%$}
\resolucao{(a) De $r=\dfrac{mv}{qB}$ com $v$ e $B$ fixos:
\[
\frac{\Delta r}{r}=\frac{\Delta m}{m}=1\%.
\]
Com $r=\SI{20}{\centi\meter}$, a separação de raios é
$\SI{2}{\milli\meter}$, e no detector (após meia volta) a separação é
$2\Delta r=\SI{4}{\milli\meter}$.\\
(b) \textbf{Requisito de colimação.} Para \emph{resolver} duas linhas separadas
por $\SI{4}{\milli\meter}$, a largura de cada linha deve ser bem menor. Duas
fontes de alargamento:
\begin{itemize}
\item \textbf{Dispersão de velocidade:} como $r\propto v$, uma dispersão
$\Delta v/v$ produz alargamento $\Delta r/r$ de mesma magnitude. É preciso
$\Delta v/v\ll1\%$ --- daí a necessidade do seletor de velocidades.
\item \textbf{Dispersão angular:} íons entrando com pequenos ângulos diferentes
descrevem círculos deslocados. \emph{Felizmente}, a geometria de meia volta é
\textbf{autofocalizante}: íons que divergem de um ponto reconvergem após
$\ang{180}$, em primeira ordem. O alargamento residual é de segunda ordem no
ângulo.
\end{itemize}
(c) \textbf{Limitações práticas.}
\begin{itemize}
\item \textbf{Vácuo:} colisões com gás residual espalham o feixe. É preciso
pressão abaixo de $\pot{1}{-5}\,\si{\pascal}$.
\item \textbf{Uniformidade do campo:} variações de $B$ ao longo da trajetória
introduzem erro proporcional. Exige-se uniformidade melhor que a resolução
desejada.
\item \textbf{Carga espacial:} feixes intensos se repelem eletrostaticamente e
se alargam --- limitando a corrente utilizável.
\item \textbf{Íons de mesma razão $m/q$:} não são separáveis. $\mathrm{N_2^+}$ e
$\mathrm{CO^+}$ (ambos $28\,u$) exigem instrumentos de altíssima resolução para
distinguir.
\end{itemize}}
\end{questao}

\begin{questao}[4]{}
\textbf{(Projeto de atuador linear.)} Projete um atuador que produza
$\SI{5}{\newton}$ com curso de $\SI{2}{\centi\meter}$.
\begin{itensq}
\item Determine as combinações de $B$, $I$ e $L$ possíveis.
\item Avalie o compromisso entre corrente e número de espiras.
\item Determine a potência dissipada e o trabalho realizado.
\end{itensq}
\resposta{Com $B=\SI{0.5}{\tesla}$ e $L=\SI{0.2}{\meter}$: $I=\SI{50}{\ampere}$
com uma espira, ou $\SI{2.5}{\ampere}$ com vinte}
\resolucao{(a) De $F=NBIL$:
\[
NIL=\frac{F}{B}=\frac{5}{0{,}5}=\SI{10}{\ampere\meter}.
\]
Qualquer combinação que satisfaça esse produto serve --- há dois graus de
liberdade.\\
(b) \textbf{Compromisso.} Com $L=\SI{0.2}{\meter}$ fixo:
\begin{center}
\begin{tabular}{ccc}
\hline
$N$ & $I$ & Observação\\
\hline
$1$ & $\SI{50}{\ampere}$ & fio muito grosso, fonte de alta corrente\\
$20$ & $\SI{2.5}{\ampere}$ & equilíbrio razoável\\
$200$ & $\SI{0.25}{\ampere}$ & bobina volumosa, alta indutância\\
\hline
\end{tabular}
\end{center}
\textbf{Mais espiras} reduzem a corrente (fio mais fino, fonte mais simples),
mas aumentam volume, massa, resistência e \emph{indutância} --- o que torna o
atuador mais lento a responder.\\
\textbf{Menos espiras} dão resposta rápida, ao custo de corrente elevada.\\
(c) \textbf{Dissipação.} Com $N=20$, $I=\SI{2.5}{\ampere}$ e resistência total
estimada em $\SI{2}{\ohm}$:
\[
P=RI^{2}=2(6{,}25)=\SI{12.5}{\watt}.
\]
\textbf{Trabalho útil:}
\[
W=F\,d=5(0{,}02)=\SI{0.1}{\joule}.
\]
\textbf{Comparação reveladora:} se o curso levar $\SI{0.1}{\second}$, a potência
mecânica é $\SI{1}{\watt}$ contra $\SI{12.5}{\watt}$ dissipados --- rendimento
de $7\%$. Atuadores eletromagnéticos de curso curto são intrinsecamente
ineficientes, e é por isso que costumam operar em regime \emph{pulsado}, e não
contínuo.}
\end{questao}

\begin{questao}[4]{}
\textbf{(Efeito Hall.)} Deduza a expressão da tensão Hall e analise suas
aplicações.
\begin{itensq}
\item Deduza $V_H=\dfrac{IB}{nqt}$.
\item Explique como o sinal revela o tipo de portador.
\item Cite três aplicações e o que cada uma mede.
\end{itensq}
\resposta{Equilíbrio entre força magnética e força elétrica transversal}
\resolucao{(a) \textbf{Dedução.} Os portadores, com velocidade de deriva
$v_d$, sofrem força magnética $qv_dB$ transversal e acumulam-se em uma face.
Isso cria campo elétrico transversal $E_H$, que cresce até equilibrar:
\[
qE_H=qv_dB
\;\Longrightarrow\;
E_H=v_dB.
\]
Com largura $w$, a tensão é $V_H=E_Hw=v_dBw$. Usando
$I=nqv_d(wt)$, elimina-se $v_d$:
\[
v_d=\frac{I}{nqwt}
\;\Longrightarrow\;
V_H=\frac{IB}{nqt}.
\]
\textbf{Note:} a largura $w$ cancela --- só a \emph{espessura} $t$ importa.\\
(b) \textbf{Sinal e tipo de portador.} Elétrons (negativos) e lacunas
(positivas) movem-se em sentidos opostos para produzir a mesma corrente. A força
$q\vec v\times\vec B$ os desvia para o \textbf{mesmo} lado --- mas com cargas de
sinais contrários, a polaridade de $V_H$ \textbf{se inverte}.\\
\textbf{É o único método simples de determinar experimentalmente o sinal dos
portadores} --- e foi assim que se estabeleceu que semicondutores tipo P
conduzem por lacunas.\\
(c) \textbf{Três aplicações.}
\begin{itemize}
\item \textbf{Medição de campo magnético} (gaussímetro): com $I$, $n$ e
$t$ conhecidos, $V_H\propto B$.
\item \textbf{Medição de corrente sem contato:} coloca-se o sensor no entreferro
de um núcleo que envolve o condutor. Mede $B$, e daí $I$ --- com isolação
galvânica completa.
\item \textbf{Caracterização de materiais:} com $B$ e $I$ conhecidos, obtém-se
$n$ (densidade de portadores) e o \emph{sinal} deles. Combinado com a
condutividade, fornece a mobilidade.
\end{itemize}
\textbf{Vantagem sobre alternativas:} sensores Hall respondem a campo
\emph{estático}, ao contrário de bobinas, que só detectam variação.}
\end{questao}

\begin{questao}[4]{}
\textbf{(Estabilidade de espira.)} Analise o equilíbrio de uma espira de
corrente em campo magnético uniforme.
\begin{itensq}
\item Determine as posições de equilíbrio.
\item Classifique cada uma quanto à estabilidade.
\item Deduza o período de pequenas oscilações.
\end{itensq}
\resposta{$\theta=0$ é estável, $\theta=\pi$ é instável;
$T=2\pi\sqrt{J/(mB)}$}
\resolucao{(a) O torque é $\tau=-mB\sin\theta$, com $\theta$ o ângulo entre
$\vec m$ e $\vec B$. Ele se anula em
\[
\theta=0 \quad\text{e}\quad \theta=\pi.
\]
(b) \textbf{Classificação pela energia.} A energia potencial de um dipolo
magnético é
\[
U=-\vec m\cdot\vec B=-mB\cos\theta.
\]
\begin{itemize}
\item $\theta=0$ ($\vec m$ paralelo a $\vec B$): $U=-mB$, \textbf{mínimo}
$\Rightarrow$ equilíbrio \textbf{estável}. Deslocada, a espira sofre torque
restaurador.
\item $\theta=\pi$ (antiparalelo): $U=+mB$, \textbf{máximo} $\Rightarrow$
equilíbrio \textbf{instável}. Qualquer perturbação a faz girar $\ang{180}$.
\end{itemize}
(c) \textbf{Pequenas oscilações.} Para $\theta\ll1$, $\sin\theta\approx\theta$ e
\[
J\ddot\theta=-mB\theta
\;\Longrightarrow\;
\ddot\theta+\frac{mB}{J}\theta=0,
\]
equação do MHS com
\[
\omega=\sqrt{\frac{mB}{J}};
\qquad
T=2\pi\sqrt{\frac{J}{mB}},
\]
onde $J$ é o momento de inércia da espira.\\
\textbf{Aplicação prática:} é assim que funcionam o \textbf{galvanômetro de
bobina móvel} --- em que uma mola acrescenta torque restaurador proporcional a
$\theta$, tornando a deflexão proporcional à corrente --- e a \textbf{bússola},
cujo período de oscilação permite medir o campo terrestre se $J$ e $m$ forem
conhecidos.\\
\textbf{Diferença de $U$ entre os dois equilíbrios:} $2mB$ --- energia
necessária para inverter a espira, e a mesma que um ímã precisa receber para ter
sua polarização revertida.}
\end{questao}

\begin{questao}[4]{}
\textbf{(Força de Lorentz --- estrutura e consequências.)} Analise a expressão
completa $\vec F=q\left(\vec E+\vec v\times\vec B\right)$.
\begin{itensq}
\item Compare as duas parcelas quanto à dependência da velocidade, à direção e
      ao trabalho realizado.
\item Avalie um próton com $v=\pot{3}{4}\,\si{\meter\per\second}$ em
      $E=\SI{500}{\volt\per\meter}$ e $B=\SI{0.02}{\tesla}$
      \emph{perpendiculares entre si}, com $\vec v$ paralelo a $\vec E$.
\item Explique por que aceleradores usam campo elétrico para acelerar e
      magnético para guiar.
\end{itensq}
\resposta{$F_E=\SI{8e-17}{\newton}$, $F_B=\SI{9.6e-17}{\newton}$,
perpendiculares; $|F|=\SI{1.25e-16}{\newton}$ a $\ang{50.2}$ de $\vec E$}
\resolucao{(a) \textbf{Comparação das duas parcelas.}
\begin{center}
\begin{tabular}{lll}
\hline
& Parcela elétrica & Parcela magnética\\
\hline
Expressão & $q\vec E$ & $q\,\vec v\times\vec B$\\
Depende de $v$? & \textbf{não} & \textbf{sim} (linear)\\
Age em repouso? & sim & não\\
Direção & paralela a $\vec E$ & $\perp$ a $\vec v$ \emph{e} a $\vec B$\\
Realiza trabalho? & \textbf{sim} & \textbf{nunca}\\
Altera $|v|$? & sim & não\\
\hline
\end{tabular}
\end{center}
\textbf{A linha decisiva é a do trabalho.} Como
$\left(\vec v\times\vec B\right)\cdot\vec v=0$ identicamente, a parcela
magnética tem potência nula --- ela desvia sem acelerar. A elétrica, ao
contrário, é a única capaz de variar a energia cinética.\\
(b) \textbf{Avaliação numérica.}
\[
F_E=qE=(\pot{1.6}{-19})(500)=\pot{8}{-17}\,\si{\newton};
\]
\[
F_B=qvB=(\pot{1.6}{-19})(\pot{3}{4})(0{,}02)=\pot{9.6}{-17}\,\si{\newton}.
\]
Com $\vec v$ paralelo a $\vec E$ e $\vec B$ perpendicular a ambos, a força
magnética é \textbf{perpendicular} à elétrica:
\[
|\vec F|=\sqrt{F_E^{2}+F_B^{2}}=\pot{1.25}{-16}\,\si{\newton},
\qquad
\theta=\arctan\frac{9{,}6}{8{,}0}=\ang{50.2}
\]
em relação a $\vec E$.\\
\textbf{Razão entre as parcelas:}
$\dfrac{F_B}{F_E}=\dfrac{vB}{E}=1{,}20$ --- adimensional, e igual a $1$
exatamente na velocidade $E/B$ do seletor.\\
\textbf{Trabalho em um deslocamento de $\SI{1}{\centi\meter}$ ao longo de
$\vec E$:} apenas a parcela elétrica contribui,
\[
W=F_E d=(\pot{8}{-17})(0{,}01)=\pot{8}{-19}\,\si{\joule}=\SI{5}{\electronvolt},
\]
que é simplesmente $qEd$ --- ou seja, $q$ vezes a tensão percorrida.\\
(c) \textbf{Divisão de papéis em aceleradores.}
\begin{itemize}
\item \textbf{Campo elétrico acelera.} É o único que realiza trabalho, e o ganho
de energia por passagem é $qU$ --- daí as cavidades de radiofrequência.
\item \textbf{Campo magnético guia.} Ele curva a trajetória sem custo
energético, permitindo fechar o percurso em círculo e reaproveitar as mesmas
cavidades milhares de vezes.
\item \textbf{Campo magnético também focaliza}, com quadrupolos que comprimem o
feixe transversalmente --- de novo, sem alterar a energia.
\end{itemize}
\textbf{Se fosse possível acelerar com campo magnético}, aceleradores lineares
de quilômetros seriam desnecessários. A impossibilidade não é técnica: é
consequência direta de $\left(\vec v\times\vec B\right)\perp\vec v$.\\
\textbf{Nota histórica:} a expressão completa reúne dois fenômenos que foram
descobertos e estudados separadamente. Sua unificação --- e a constatação de que
o que é campo elétrico em um referencial pode ser magnético em outro --- é um
dos pontos de partida da relatividade restrita.}
\end{questao}

\begin{questao}[4]{}
\textbf{(Projeto de cíclotron.)} Um cíclotron acelera prótons até
$\SI{10}{\mega\electronvolt}$ usando campo de $\SI{1.5}{\tesla}$.
\begin{itensq}
\item Determine a frequência de operação.
\item Determine o raio máximo e o número de voltas, com
      $\SI{50}{\kilo\volt}$ de ganho por passagem.
\item Determine se a correção relativística é necessária.
\end{itensq}
\dados{$e=\pot{1.6}{-19}\,\si{\coulomb}$; $m_p=\pot{1.67}{-27}\,\si{\kilogram}$;
$\SI{1}{\mega\electronvolt}=\pot{1.6}{-13}\,\si{\joule}$}
\resposta{$f=\SI{22.9}{\mega\hertz}$; $r=\SI{30.5}{\centi\meter}$;
$100$ voltas; correção de $1{,}1\%$}
\resolucao{(a) A frequência de cíclotron independe da energia:
\[
f=\frac{qB}{2\pi m}=\frac{(\pot{1.6}{-19})(1{,}5)}{2\pi(\pot{1.67}{-27})}
=\SI{22.9}{\mega\hertz}.
\]
\textbf{É isso que torna o aparelho viável:} o oscilador que alimenta os "dês"
opera em frequência \emph{fixa}, embora a partícula ganhe energia
continuamente.\\
(b) Da energia final,
\[
v=\sqrt{\frac{2E}{m}}
=\sqrt{\frac{2(10)(\pot{1.6}{-13})}{\pot{1.67}{-27}}}
=\pot{4.38}{7}\,\si{\meter\per\second};
\]
\[
r=\frac{mv}{qB}=\SI{0.305}{\meter}
\;\Longrightarrow\;
\text{diâmetro}\approx\SI{61}{\centi\meter}.
\]
Com dois cruzamentos do intervalo por volta, cada volta rende
$\SI{100}{\kilo\electronvolt}$:
\[
N=\frac{10\,\si{\mega\electronvolt}}{0{,}1\,\si{\mega\electronvolt}}
=100\ \text{voltas}.
\]
\textbf{Tempo total:} $100/f=\SI{4.4}{\micro\second}$.\\
(c) \textbf{Correção relativística.} Com
$v/c=0{,}146$:
\[
\gamma=\frac{1}{\sqrt{1-(v/c)^{2}}}=1{,}0108,
\]
ou seja, a massa efetiva é $1{,}1\%$ maior no fim da aceleração.\\
\textbf{Consequência:} a frequência de cíclotron cai na mesma proporção, e a
partícula começa a chegar \emph{atrasada} ao intervalo acelerador. Após muitas
voltas o defasamento acumula e a aceleração cessa.\\
\textbf{Aqui, $1{,}1\%$ é tolerável} --- o cíclotron clássico funciona. Mas para
$\SI{100}{\mega\electronvolt}$ o desvio passaria de $10\%$ e seria proibitivo.\\
\textbf{Soluções:} o \textbf{sincrociclotron} modula a frequência ao longo do
ciclo; o \textbf{síncrotron} mantém o raio fixo e varia o campo. Ambos abandonam
a simplicidade do cíclotron justamente porque o período deixa de ser constante
--- e a raiz disso está na expressão $T=2\pi m/qB$, válida apenas enquanto
$m$ o for.}
\end{questao}

\begin{questao}[4]{}
\textbf{(Força entre correntes e o limite de barramentos.)} Um barramento de
subestação tem dois condutores paralelos separados por $\SI{15}{\centi\meter}$,
que devem suportar um curto-circuito de $\SI{40}{\kilo\ampere}$.
\begin{itensq}
\item Determine a força por metro durante o curto.
\item Determine o espaçamento máximo entre suportes, para deflexão admissível.
\item Discuta o caráter dessa solicitação.
\end{itensq}
\resposta{$F/L\approx\SI{2.1}{\kilo\newton\per\meter}$ --- comparável ao peso
de uma pessoa por metro}
\resolucao{(a)
\[
\frac{F}{L}=\frac{\mu_0I^{2}}{2\pi d}
=\pot{2}{-7}\frac{(\pot{4}{4})^{2}}{0{,}15}
=\pot{2}{-7}\frac{\pot{1.6}{9}}{0{,}15}=\SI{2133}{\newton\per\meter}.
\]
Mais de duas toneladas-força por metro de barramento --- e \textbf{repulsiva},
se as correntes forem opostas (que é o caso em um curto entre fases).\\
(b) \textbf{Espaçamento entre suportes.} Tratando o barramento como viga
biapoiada com carga distribuída $w=F/L$, a flecha é
\[
\delta=\frac{5wL_s^{4}}{384EJ}.
\]
Para deflexão admissível $\delta$, o vão máximo cresce apenas com
$L_s\propto\delta^{1/4}$ --- \textbf{dependência muito fraca}. Dobrar a flecha
admissível permite aumentar o vão em apenas $19\%$.\\
Na prática, os suportes ficam a cada $\SI{0.5}{}$ a $\SI{1}{\meter}$.\\
(c) \textbf{Caráter da solicitação.}
\begin{itemize}
\item \textbf{É transitória} --- dura o tempo de atuação da proteção, tipicamente
$\SI{100}{\milli\second}$. Mas a força \emph{de pico} pode ser o dobro da
calculada, devido à componente assimétrica da corrente de curto.
\item \textbf{Vai com $I^{2}$} --- dobrar a corrente de curto quadruplica a
força. É por isso que a corrente de curto-circuito presumida é um dado de
projeto tão crítico.
\item \textbf{É dinâmica:} se a frequência natural da estrutura estiver próxima
de $\SI{120}{\hertz}$ (dobro da rede, pois $F\propto i^{2}$), pode haver
ressonância mecânica.
\end{itemize}
\textbf{Conclusão:} em regime normal ($\SI{1}{\kilo\ampere}$, digamos), a força
seria $\SI{1.3}{\newton\per\meter}$ --- desprezível. O dimensionamento mecânico
de barramentos é governado inteiramente pela condição de \emph{falta}.}
\end{questao}

% END BANCO COMPLEMENTAR

\gabarito[12.4 Força magnética]

% =====================================================================
\subsection{Lei de Lenz e indução eletromagnética}
% =====================================================================

\blocofixacao

\begin{questao}[1]{}
Segundo a lei de Lenz, a corrente induzida em uma espira tem sentido tal que:
\begin{alternativas}
\item reforça a variação do fluxo que a originou.
\item se opõe à variação do fluxo que a originou.
\item é sempre horária.
\item independe do sentido da variação do fluxo.
\item é proporcional ao fluxo, e não à sua variação.
\end{alternativas}
\resposta{(b)}
\resolucao{A lei de Lenz é uma manifestação da conservação de energia: a corrente
induzida cria um campo que \emph{se opõe} à variação de fluxo que a gerou. Se
reforçasse a variação (alternativa a), teríamos energia criada do nada.}
\end{questao}

\begin{questao}[1]{}
Um campo magnético variável atravessa perpendicularmente uma espira, com fluxo
$\Phi(t) = 0{,}02\,t$ (em Wb, com $t$ em segundos). Determine a fem induzida.
\resposta{$\varepsilon = \SI{20}{\milli\volt}$}
\resolucao{A fem é a taxa de variação do fluxo:
\[
\varepsilon=\left|\frac{\mathrm{d}\Phi}{\mathrm{d}t}\right|
=\SI{0.02}{\weber\per\second}=\SI{20}{\milli\volt}.
\]
Como o fluxo cresce linearmente, a fem é \emph{constante}.}
\end{questao}

\blocoaplicacao

\begin{questao}[2]{}
O campo magnético no interior de uma bobina aumenta linearmente de $0$ a
$\SI{0.5}{\tesla}$ em $\SI{2}{\second}$. A bobina tem $\SI{50}{}$ espiras e área de
$\SI{20}{\centi\meter\squared}$. Calcule a fem induzida.
\resposta{$\varepsilon = \SI{25}{\milli\volt}$}
\resolucao{Com $N = 50$ espiras e
$A = \pot{20}{-4}\,\si{\meter\squared}$:
\[
\varepsilon = N\,A\,\frac{\Delta B}{\Delta t}
= 50\cdot\pot{20}{-4}\cdot\frac{0{,}5}{2}
= 50\cdot\pot{20}{-4}\cdot0{,}25 = \SI{25}{\milli\volt}.
\]}
\end{questao}

\begin{questao}[2]{}
Uma bobina de $\SI{100}{}$ espiras e área $\SI{30}{\centi\meter\squared}$ está em
um campo que varia segundo $B(t) = 0{,}5 - 0{,}1\,t$ (em T). Calcule a fem induzida
e diga se a corrente induzida cresce, decresce ou é constante.
\resposta{$\varepsilon = \SI{30}{\milli\volt}$, constante}
\resolucao{A taxa de variação do campo é
$\dfrac{\mathrm{d}B}{\mathrm{d}t} = \SI{-0.1}{\tesla\per\second}$ (constante):
\[
|\varepsilon| = N\,A\,\left|\frac{\mathrm{d}B}{\mathrm{d}t}\right|
= 100\cdot\pot{30}{-4}\cdot0{,}1 = \SI{30}{\milli\volt}.
\]
Como a taxa de variação é constante, a fem é constante --- e, num circuito de
resistência fixa, a corrente induzida também é \emph{constante}. O fluxo está
diminuindo, então a corrente induzida atua no sentido de \emph{reforçá-lo}
(Lenz).}
\end{questao}

\begin{questao}[2]{}
Um ímã é aproximado, com o polo norte à frente, de uma espira condutora, como na
figura. Determine o sentido da corrente induzida e a natureza da força entre o
ímã e a espira.

\circuitoq{%
\begin{tikzpicture}[scale=0.9,line width=0.8pt,>=Stealth]
 % ima
 \draw[fill=Icol!20,draw=Icol] (-4,-0.35) rectangle (-2.4,0.35);
 \node at (-3.5,0){\footnotesize S}; \node at (-2.7,0){\footnotesize N};
 \draw[->,Vcol,thick] (-2.2,0)--(-1.2,0); \node[Vcol,font=\footnotesize] at (-1.7,0.3){$v$};
 % espira
 \draw[azulprof,very thick] (0,0) ellipse (0.35 and 1.1);
 \node[azulprof,font=\footnotesize] at (0,-1.4){espira};
\end{tikzpicture}}{Ímã aproximando-se de uma espira.}

\begin{alternativas}
\item A corrente induzida cria um polo norte voltado para o ímã, repelindo-o.
\item A corrente induzida cria um polo sul voltado para o ímã, atraindo-o.
\item Não há corrente induzida, pois o ímã não toca a espira.
\item A corrente induzida não gera força sobre o ímã.
\item A espira é atraída, acelerando o ímã.
\end{alternativas}
\resposta{(a)}
\resolucao{Ao aproximar o polo norte, o fluxo através da espira \emph{aumenta}.
Pela lei de Lenz, a corrente induzida se opõe a esse aumento, criando um campo que
\emph{repele} o ímã --- ou seja, a face da espira voltada para o ímã se comporta
como um polo \textbf{norte}. Consequentemente, surge uma força de \emph{repulsão}
que dificulta a aproximação. É a manifestação da conservação de energia: para
aproximar o ímã é preciso realizar trabalho, que se converte na energia elétrica
da corrente induzida.}
\end{questao}

\begin{questao}[2]{}
Um transformador ideal tem no primário uma bobina onde o fluxo magnético varia à
taxa de $\SI{0.01}{\weber\per\second}$. Se o enrolamento tem $\SI{200}{}$ espiras,
qual é a fem induzida?
\resposta{$\varepsilon = \SI{2}{\volt}$}
\resolucao{$\varepsilon = N\dfrac{\mathrm{d}\Phi}{\mathrm{d}t}
= 200\cdot0{,}01 = \SI{2}{\volt}$. É assim que o transformador transfere energia:
o fluxo variável criado pelo primário induz tensão no secundário,
proporcionalmente ao número de espiras de cada lado.}
\end{questao}

\blocoaprofundamento

\begin{questao}[3]{}
O fluxo magnético em uma espira varia senoidalmente:
$\Phi(t) = \Phi_0\cos(\omega t)$, com $\Phi_0 = \SI{0.01}{\weber}$ e
$\omega = \SI{100}{\radian\per\second}$.
\begin{itensq}
\item Obtenha a expressão da fem induzida.
\item Qual é o valor máximo da fem?
\end{itensq}
\resposta{(a) $\varepsilon(t) = \Phi_0\omega\sin(\omega t)$;
(b) $\varepsilon_{\max} = \SI{1}{\volt}$}
\resolucao{(a) A fem é a derivada do fluxo (com sinal trocado):
\[
\varepsilon = -\frac{\mathrm{d}\Phi}{\mathrm{d}t}
= -\Phi_0\cdot(-\omega\sin(\omega t)) = \Phi_0\,\omega\sin(\omega t).
\]
(b) O valor máximo ocorre quando $\sin(\omega t) = 1$:
\[
\varepsilon_{\max} = \Phi_0\,\omega = 0{,}01\cdot100 = \SI{1}{\volt}.
\]
\textbf{Observação:} um fluxo cossenoidal gera uma fem \emph{senoidal} --- há uma
defasagem de $\SI{90}{\degree}$ entre eles. É exatamente o que ocorre em um
gerador de corrente alternada: a espira girando em campo uniforme produz fluxo
senoidal e, portanto, tensão senoidal.}
\end{questao}

\begin{questao}[3]{Apostila original revisada}
A corrente induzida em um circuito puramente resistivo é
$i(t)=I_0\sin(\omega t)$. Determine a expressão da fem induzida e indique em quais
instantes a corrente muda de sentido durante um período.
\resposta{$\mathcal{E}(t)=RI_0\sin(\omega t)$; inversões em $t=0$, $\pi/\omega$ e $2\pi/\omega$}
\resolucao{Como o circuito é resistivo, $\mathcal{E}=Ri$. O sinal da senoide muda a
cada múltiplo de $\pi$, portanto o sentido da corrente se inverte a cada meio período.}
\end{questao}

\blocodesafio

\begin{questao}[4]{}
Uma barra condutora de comprimento $\ell = \SI{40}{\centi\meter}$ desliza sem
atrito, com velocidade constante $v = \SI{3}{\meter\per\second}$, sobre dois
trilhos paralelos fechados por um resistor $R = \SI{2}{\ohm}$. O conjunto está
imerso em campo magnético uniforme $B = \SI{0.5}{\tesla}$, perpendicular ao plano
dos trilhos.

\circuitoq{%
\begin{tikzpicture}[scale=1,line width=0.8pt,>=Stealth]
 % trilhos
 \draw[cinza,line width=1.4pt] (0,0)--(6,0);
 \draw[cinza,line width=1.4pt] (0,2.2)--(6,2.2);
 % resistor
 \draw (0,0) to[R,l_={$R$}] (0,2.2);
 % barra
 \draw[Vcol,line width=2.2pt] (3.6,-0.15)--(3.6,2.35);
 \node[Vcol,font=\footnotesize,above] at (3.6,2.4) {barra};
 \draw[->,laranja,line width=1.3pt] (3.9,1.1)--(5.1,1.1)
       node[midway,above,font=\footnotesize]{$v$};
 % campo B saindo
 \foreach \x in {1.0,1.8,2.6}
   \foreach \y in {0.55,1.1,1.65}
     {\node[Icol,font=\scriptsize] at (\x,\y) {$\odot$};}
 \node[Icol,font=\footnotesize] at (1.8,0.15) {$\vec B$ (saindo)};
 \draw[<->,cinza] (6.3,0)--(6.3,2.2) node[midway,right,font=\footnotesize]{$\ell$};
\end{tikzpicture}}{Barra deslizante sobre trilhos --- fem de movimento.}

\begin{itensq}
\item Calcule a fem induzida e a corrente no circuito.
\item Determine a força magnética que atua sobre a barra e seu sentido.
\item Calcule a potência dissipada no resistor e a potência mecânica necessária
      para manter a barra em movimento uniforme. Compare-as.
\item Interprete o resultado à luz da lei de Lenz e da conservação de energia.
\end{itensq}
\resposta{(a) $\varepsilon=\SI{0.6}{\volt}$, $i=\SI{0.3}{\ampere}$;
(b) $F=\SI{0.06}{\newton}$, opondo-se ao movimento;
(c) ambas $=\SI{0.18}{\watt}$; (d) ver comentário}
\resolucao{\textbf{(a) Fem de movimento.} À medida que a barra avança, a área do
circuito aumenta a uma taxa $\dfrac{\mathrm{d}A}{\mathrm{d}t} = \ell v$. Logo
\[
\varepsilon=\left|\frac{\mathrm{d}\Phi}{\mathrm{d}t}\right|
= B\,\ell\,v = 0{,}5\cdot0{,}4\cdot3=\SI{0.6}{\volt},
\qquad
i=\frac{\varepsilon}{R}=\frac{0{,}6}{2}=\SI{0.3}{\ampere}.
\]
\textbf{(b) Força sobre a barra.} A barra, agora percorrida por corrente,
encontra-se no campo $B$ e sofre
\[
F = B\,i\,\ell = 0{,}5\cdot0{,}3\cdot0{,}4=\SI{0.06}{\newton}.
\]
Pela lei de Lenz, essa força \textbf{opõe-se ao movimento} --- ela freia a barra.
(Se favorecesse o movimento, teríamos aceleração espontânea e energia criada do
nada.)

\textbf{(c) Balanço de potência.}
\[
P_{\text{dissipada}} = R\,i^{2} = 2(0{,}3)^{2}=\SI{0.18}{\watt},
\]
\[
P_{\text{mecânica}} = F\,v = 0{,}06\cdot3 = \SI{0.18}{\watt}.
\]
\textbf{São exatamente iguais.}

\textbf{(d) Interpretação.} O resultado não é coincidência: é a
\textbf{conservação de energia} em ação. Para manter a barra em movimento
\emph{uniforme}, um agente externo deve aplicar uma força igual e oposta à força
magnética de frenagem, realizando trabalho a uma taxa $Fv$. Toda essa potência
mecânica é convertida em potência elétrica e, finalmente, dissipada como calor no
resistor. Nada se perde, nada se cria:
\[
\underbrace{P_{\text{mec}}}_{\text{trabalho externo}}
\;\longrightarrow\;
\underbrace{\varepsilon\,i}_{\text{potência elétrica}}
\;\longrightarrow\;
\underbrace{R i^{2}}_{\text{calor}}.
\]
Verificando pela via elétrica: $\varepsilon\,i = 0{,}6\cdot0{,}3
= \SI{0.18}{\watt}$ \checkmark\\
\textbf{Este é o princípio do gerador elétrico} --- e, invertendo o raciocínio, o
do \emph{freio eletromagnético} usado em trens e em esteiras de academia: quanto
maior a velocidade, maior a corrente induzida e maior a frenagem, sem contato
mecânico nem desgaste.}
\end{questao}

% BEGIN BANCO COMPLEMENTAR Lei de Lenz e indução eletromagnética
% =====================================================================
%  Banco complementar --- Lei de Lenz e Inducao Eletromagnetica
%  REVISADO. Substitui a versao gerada por template (8 clones em N1,
%  11 em N2, 13 em N3 e 9 questoes conceituais marcadas como N4).
%  O cap11 ja cobre: enunciado da lei de Lenz (MC), Phi(t)=0,02t,
%  campo crescendo linearmente de 0 a 0,5 T, bobina de 100 espiras com
%  B(t), IMA APROXIMADO COM O POLO NORTE (N2), transformador ideal com
%  dPhi/dt dado, fluxo senoidal, corrente induzida senoidal e -- como
%  N4 -- a BARRA DESLIZANDO SOBRE TRILHOS. As duas ultimas motivaram a
%  remocao de duas N4 do banco antigo, que as duplicavam.
%  Este banco explora: CARGA induzida (que independe do tempo), freio
%  por correntes parasitas, laminacao de nucleos, velocidade-limite,
%  inducao mutua com convencao dos pontos, projeto de bobina e de
%  experimento, e a razao pela qual transformadores nao operam em CC.
%  Distribuicao mantida: 8 N1 + 11 N2 + 13 N3 + 9 N4 = 41
% =====================================================================

\subsubsection*{Banco complementar --- Lei de Lenz e indução eletromagnética}

\begin{atencao}[Organização do banco]
A lei de Faraday é $\varepsilon=-N\dfrac{\mathrm{d}\Phi}{\mathrm{d}t}$: o que
induz f.e.m. é a \textbf{variação} do fluxo, não seu valor. O sinal negativo é a
lei de Lenz, e ela não é um detalhe --- é a conservação da energia em forma
elétrica. O N3 trata de velocidade-limite, carga induzida, correntes parasitas e
indução mútua; o N4 exige demonstrações, projeto e a análise de por que
transformadores exigem corrente alternada.
\end{atencao}

\blocofixacao

\begin{questao}[1]{}
Uma bobina de $200$ espiras sofre variação de fluxo de
$\SI{2}{\milli\weber}$ por espira em $\SI{5}{\milli\second}$. Determine a
f.e.m. média induzida.
\resposta{$\varepsilon=\SI{80}{\volt}$}
\resolucao{
\[
|\varepsilon|=N\frac{\Delta\Phi}{\Delta t}
=200\,\frac{\pot{2}{-3}}{\pot{5}{-3}}=\SI{80}{\volt}.
\]
\textbf{Note o papel de $N$:} cada espira contribui com a mesma f.e.m., e elas
estão em série. Duzentas espiras multiplicam por duzentos --- é assim que se
obtêm tensões altas a partir de variações de fluxo modestas.}
\end{questao}

\begin{questao}[1]{}
Uma barra de $\SI{30}{\centi\meter}$ move-se a $\SI{4}{\meter\per\second}$
perpendicularmente a um campo de $\SI{0.5}{\tesla}$. Determine a f.e.m.
induzida.
\resposta{$\varepsilon=\SI{0.6}{\volt}$}
\resolucao{
\[
\varepsilon=B\ell v=(0{,}5)(0{,}30)(4)=\SI{0.6}{\volt}.
\]
\textbf{Origem da fórmula:} os portadores da barra movem-se com ela e sofrem
força magnética $qvB$ ao longo do comprimento. Eles se acumulam nas
extremidades até que o campo elétrico resultante equilibre a força magnética ---
e a diferença de potencial daí resultante é $B\ell v$.\\
\textbf{Equivalência com Faraday:} se a barra fizer parte de um circuito, a área
varia à taxa $\ell v$ e $\mathrm{d}\Phi/\mathrm{d}t=B\ell v$ --- o mesmo
resultado.}
\end{questao}

\begin{questao}[1]{}
Determine o número de espiras necessário para obter $\SI{12}{\volt}$ com
variação de fluxo de $\SI{0.04}{\weber\per\second}$.
\resposta{$N=300$}
\resolucao{
\[
N=\frac{\varepsilon}{\mathrm{d}\Phi/\mathrm{d}t}=\frac{12}{0{,}04}=300.
\]
\textbf{Compromisso de projeto:} mais espiras dão mais tensão, mas aumentam a
resistência do enrolamento (mais perdas) e o volume. Aumentar
$\mathrm{d}\Phi/\mathrm{d}t$ --- por campo mais intenso ou variação mais rápida
--- costuma ser preferível quando possível.}
\end{questao}

\begin{questao}[1]{}
Uma bobina de $100$ espiras deve gerar $\SI{50}{\volt}$ com variação de fluxo
de $\SI{5}{\milli\weber}$. Determine o tempo em que essa variação deve ocorrer.
\resposta{$\Delta t=\SI{10}{\milli\second}$}
\resolucao{
\[
\Delta t=\frac{N\,\Delta\Phi}{\varepsilon}
=\frac{100(\pot{5}{-3})}{50}=\pot{1}{-2}\,\si{\second}.
\]
\textbf{Leitura:} a f.e.m. é inversamente proporcional ao tempo. A \emph{mesma}
variação de fluxo, feita dez vezes mais rápido, produz dez vezes mais tensão ---
princípio de funcionamento das bobinas de ignição e dos geradores de alta
tensão por interrupção de corrente.}
\end{questao}

\begin{questao}[1]{}
O sinal negativo na lei de Faraday expressa:
\begin{alternativas}[2]
\item a lei de Lenz e a conservação da energia
\item que a f.e.m. é sempre negativa
\item um erro de convenção
\item que a corrente é sempre horária
\end{alternativas}
\resposta{(a)}
\resolucao{O sinal indica que a corrente induzida circula no sentido de
\textbf{opor-se} à variação de fluxo que a produziu.\\
\textbf{Por que precisa ser assim:} se a corrente induzida \emph{reforçasse} a
variação, ela aumentaria o fluxo, que induziria mais corrente, e assim por
diante --- energia surgiria do nada. A lei de Lenz é a garantia de que isso não
ocorre.\\
\textbf{Consequência prática:} qualquer variação de fluxo encontra
\emph{resistência}. Aproximar um ímã de uma espira fechada exige força; girar
um gerador sob carga exige torque. É de onde vem a energia elétrica.}
\end{questao}

\begin{questao}[1]{}
Uma espira está imersa em campo magnético uniforme e \textbf{constante}, e é
deslocada rapidamente dentro da região. Determine a f.e.m. induzida.
\resposta{$\varepsilon=0$}
\resolucao{O fluxo é $\Phi=BA$, com $B$ e $A$ constantes --- ele \textbf{não
varia}, e portanto
\[
\varepsilon=-\frac{\mathrm{d}\Phi}{\mathrm{d}t}=0.
\]
\textbf{Ponto conceitual central:} \emph{movimento não induz f.e.m.}; variação
de fluxo é que induz. Uma espira pode mover-se a grande velocidade sem gerar
nada, desde que o fluxo através dela permaneça o mesmo.\\
\textbf{Onde há f.e.m.:} apenas ao \emph{atravessar a fronteira} da região de
campo, quando a área imersa muda.}
\end{questao}

\begin{questao}[1]{}
Mostre que $\si{\weber\per\second}$ equivale a volt.
\resposta{$\si{\weber}=\si{\volt\second}$}
\resolucao{Da lei de Faraday, $\varepsilon=\mathrm{d}\Phi/\mathrm{d}t$, logo
\[
[\si{\volt}]=\frac{[\si{\weber}]}{[\si{\second}]}
\;\Longrightarrow\;
\si{\weber}=\si{\volt}\cdot\si{\second}.
\]
\textbf{Leitura útil:} o weber é um "volt-segundo". Isso torna imediato o
resultado de que a \emph{integral} da f.e.m. no tempo é a variação de fluxo ---
base do fluxímetro e do método de medição por integração.}
\end{questao}

\begin{questao}[1]{}
Segundo a lei de Lenz, quando o fluxo através de uma espira \textbf{aumenta},
a corrente induzida cria um campo próprio que:
\begin{alternativas}[2]
\item reforça o campo externo
\item opõe-se ao campo externo
\item é perpendicular ao campo externo
\item é nulo
\end{alternativas}
\resposta{(b)}
\resolucao{A corrente induzida produz campo que se \emph{opõe ao aumento} --- ou
seja, no sentido contrário ao campo externo dentro da espira.\\
\textbf{Se o fluxo diminuísse}, a corrente inverteria, criando campo no
\emph{mesmo} sentido do externo, para tentar sustentá-lo.\\
\textbf{Regra prática:} a espira "não gosta de mudanças". Ela sempre reage no
sentido de manter o fluxo como estava --- nunca no de acelerar a mudança.}
\end{questao}

\blocoaplicacao

\begin{questao}[2]{}
Uma barra de $\SI{50}{\centi\meter}$ desliza a $\SI{3}{\meter\per\second}$
sobre trilhos em campo de $\SI{0.4}{\tesla}$, com resistência total de
$\SI{2}{\ohm}$.
\begin{itensq}
\item Determine a f.e.m. e a corrente.
\item Determine a força necessária para manter a velocidade.
\item Verifique o balanço de potência.
\end{itensq}
\resposta{$\varepsilon=\SI{0.6}{\volt}$; $I=\SI{0.3}{\ampere}$;
$F=\SI{60}{\milli\newton}$; $P=\SI{0.18}{\watt}$ pelos dois caminhos}
\resolucao{(a)
\[
\varepsilon=B\ell v=(0{,}4)(0{,}5)(3)=\SI{0.6}{\volt};
\qquad
I=\frac{0{,}6}{2}=\SI{0.3}{\ampere}.
\]
(b) A corrente induzida na barra, imersa no campo, sofre força magnética
\emph{oposta} ao movimento (lei de Lenz):
\[
F=BI\ell=(0{,}4)(0{,}3)(0{,}5)=\SI{0.06}{\newton}.
\]
Para manter a velocidade constante, o agente externo deve aplicar força igual e
contrária.\\
(c) \textbf{Balanço.}
\[
P_{\text{mecânica}}=Fv=(0{,}06)(3)=\SI{0.18}{\watt};
\qquad
P_{\text{elétrica}}=\frac{\varepsilon^{2}}{R}=\frac{0{,}36}{2}=\SI{0.18}{\watt}.
\]
\textbf{Coincidem exatamente} --- toda a energia mecânica fornecida vira calor no
resistor. É a lei de Lenz em forma quantitativa: a força de oposição é
precisamente a necessária para que a conta feche.}
\end{questao}

\begin{questao}[2]{}
Uma bobina de $100$ espiras e área $\SI{200}{\centi\meter\squared}$ gira a
$\SI{377}{\radian\per\second}$ em campo de $\SI{0.5}{\tesla}$.
\begin{itensq}
\item Determine a f.e.m. máxima.
\item Determine a f.e.m. eficaz.
\end{itensq}
\resposta{$\varepsilon_{\max}=\SI{377}{\volt}$;
$\varepsilon_{ef}=\SI{267}{\volt}$}
\resolucao{(a) Com $\Phi=BA\cos\omega t$:
\[
\varepsilon=-N\frac{\mathrm{d}\Phi}{\mathrm{d}t}=NBA\omega\sin\omega t
\;\Longrightarrow\;
\varepsilon_{\max}=NBA\omega.
\]
\[
\varepsilon_{\max}=100(0{,}5)(\pot{2}{-2})(377)=\SI{377}{\volt}.
\]
(b) Para forma senoidal,
\[
\varepsilon_{ef}=\frac{\varepsilon_{\max}}{\sqrt2}=\SI{267}{\volt}.
\]
\textbf{Note:} $\omega=\SI{377}{\radian\per\second}$ corresponde a
$\SI{60}{\hertz}$ --- é o valor $2\pi(60)$, e aparece constantemente em
eletrotécnica.\\
\textbf{Compromisso do gerador:} $\varepsilon\propto\omega$, mas a frequência é
fixada pela rede. Resta aumentar $N$, $B$ ou $A$ --- e é por isso que geradores
de grande porte têm núcleos volumosos e campos próximos da saturação.}
\end{questao}

\begin{questao}[2]{}
Duas bobinas têm indutância mútua $M=\SI{50}{\milli\henry}$. A corrente na
primeira varia à taxa de $\SI{200}{\ampere\per\second}$. Determine a f.e.m.
induzida na segunda.
\resposta{$\varepsilon_2=\SI{10}{\volt}$}
\resolucao{
\[
\varepsilon_2=M\frac{\mathrm{d}i_1}{\mathrm{d}t}=(0{,}05)(200)=\SI{10}{\volt}.
\]
\textbf{Simetria da indução mútua:} $M_{12}=M_{21}$ sempre, mesmo que as bobinas
sejam muito diferentes. Uma corrente variando na bobina pequena induz na grande
exatamente a mesma f.e.m. que o inverso produziria --- resultado de
reciprocidade, análogo ao que vale em redes resistivas.\\
\textbf{Aplicação:} é este o princípio do transformador, e também da
interferência indesejada entre circuitos vizinhos.}
\end{questao}

\begin{questao}[2]{}
Um transformador ideal tem $1000$ espiras no primário e $100$ no secundário,
com $\SI{220}{\volt}$ na entrada.
\begin{itensq}
\item Determine a tensão de saída.
\item Determine a corrente de entrada para carga que exige
      $\SI{5}{\ampere}$.
\end{itensq}
\resposta{$V_2=\SI{22}{\volt}$; $I_1=\SI{0.5}{\ampere}$}
\resolucao{(a) Como o \emph{mesmo} fluxo atravessa os dois enrolamentos,
\[
\frac{V_2}{V_1}=\frac{N_2}{N_1}
\;\Longrightarrow\;
V_2=220\left(\frac{100}{1000}\right)=\SI{22}{\volt}.
\]
(b) No transformador ideal não há perdas, e a potência se conserva:
\[
V_1I_1=V_2I_2
\;\Longrightarrow\;
I_1=\frac{22(5)}{220}=\SI{0.5}{\ampere}.
\]
\textbf{Relação inversa:} a corrente transforma-se na razão \emph{oposta} à da
tensão. Abaixar tensão dez vezes eleva a corrente dez vezes.\\
\textbf{Consequência para os enrolamentos:} o secundário precisa de fio dez
vezes mais grosso, embora tenha dez vezes menos espiras. O volume de cobre acaba
sendo comparável nos dois lados.}
\end{questao}

\begin{questao}[2]{}
O fluxo em uma espira varia segundo $\Phi(t)=0{,}05\,t^{2}$ (weber, com $t$ em
segundos). Determine a f.e.m. em $t=\SI{2}{\second}$ e a f.e.m. média nos dois
primeiros segundos.
\resposta{$\varepsilon(2)=\SI{0.2}{\volt}$; média de $\SI{0.1}{\volt}$}
\resolucao{\textbf{Instantânea:}
\[
\varepsilon=-\frac{\mathrm{d}\Phi}{\mathrm{d}t}=-0{,}10\,t
\;\Longrightarrow\;
|\varepsilon(2)|=\SI{0.2}{\volt}.
\]
\textbf{Média:}
\[
|\bar\varepsilon|=\frac{\Delta\Phi}{\Delta t}
=\frac{0{,}05(4)-0}{2}=\SI{0.1}{\volt}.
\]
\textbf{As duas diferem por um fator $2$} --- e ambas estão corretas, para
perguntas diferentes.\\
\textbf{Quando coincidem:} apenas se o fluxo variar \emph{linearmente}. Aqui a
variação é quadrática, e a f.e.m. cresce ao longo do intervalo, partindo de zero
e chegando a $\SI{0.2}{\volt}$ --- a média é o valor intermediário.}
\end{questao}

\begin{questao}[2]{}
Uma espira de resistência $\SI{10}{\ohm}$ e $50$ espiras sofre variação de
fluxo de $\SI{4}{\milli\weber}$. Determine a \textbf{carga} total que circula.
\resposta{$q=\SI{20}{\milli\coulomb}$}
\resolucao{Integrando a corrente induzida:
\[
q=\int i\,\mathrm{d}t=\int\frac{|\varepsilon|}{R}\,\mathrm{d}t
=\frac{N}{R}\int\left|\frac{\mathrm{d}\Phi}{\mathrm{d}t}\right|\mathrm{d}t
=\frac{N\,\Delta\Phi}{R}.
\]
\[
q=\frac{50(\pot{4}{-3})}{10}=\pot{2}{-2}\,\si{\coulomb}.
\]
\textbf{Resultado notável: a carga \emph{não depende do tempo}.} Rápida ou
lenta, a mesma variação de fluxo faz circular a mesma carga --- porque
f.e.m. e duração variam inversamente e o produto se conserva.\\
\textbf{É o princípio do fluxímetro balístico:} integrando a corrente com um
galvanômetro de bobina móvel, mede-se $\Delta\Phi$ sem controlar a velocidade do
movimento.}
\end{questao}

\begin{questao}[2]{}
Uma espira quadrada de lado $\SI{10}{\centi\meter}$ entra com velocidade
constante de $\SI{2}{\meter\per\second}$ em uma região de campo de
$\SI{0.6}{\tesla}$.
\begin{itensq}
\item Determine a f.e.m. durante a entrada.
\item Determine a f.e.m. quando totalmente dentro.
\end{itensq}
\resposta{$\SI{0.12}{\volt}$ durante a entrada; zero depois}
\resolucao{(a) Enquanto atravessa a fronteira, apenas o lado dianteiro está no
campo e funciona como barra móvel:
\[
\varepsilon=B\ell v=(0{,}6)(0{,}10)(2)=\SI{0.12}{\volt}.
\]
\emph{(Equivalentemente: a área imersa cresce a $\ell v$, e
$\mathrm{d}\Phi/\mathrm{d}t=B\ell v$.)}\\
(b) \textbf{Totalmente dentro}, os lados dianteiro e traseiro estão ambos no
campo e suas f.e.m. se \emph{cancelam}. O fluxo é constante e
\[
\varepsilon=0.
\]
\textbf{Perfil completo:} pulso retangular de $\SI{0.12}{\volt}$ durante a
entrada, zero no percurso interno, e pulso de \emph{polaridade invertida}
durante a saída. Cada pulso dura $\ell/v=\SI{50}{\milli\second}$.}
\end{questao}

\begin{questao}[2]{}
Um disco metálico gira entre os polos de um ímã. Explique o efeito observado
sobre seu movimento.
\resposta{Ele freia --- correntes de Foucault opõem-se ao movimento}
\resolucao{O disco, ao girar, tem diferentes regiões atravessadas por fluxos
variáveis. Isso induz \textbf{correntes de Foucault} (ou parasitas) --- redemoinhos
de corrente no próprio material.\\
Pela lei de Lenz, essas correntes produzem forças que se \textbf{opõem ao
movimento}: o disco desacelera, e a energia cinética vira calor.\\
\textbf{Duas leituras.}
\begin{itemize}
\item \textbf{Como recurso:} é o \emph{freio magnético}, usado em trens, em
esteiras ergométricas e no disco do antigo medidor de energia. Não há contato,
portanto não há desgaste.
\item \textbf{Como problema:} em núcleos de transformadores e motores, as
mesmas correntes representam perda pura. Combate-se \textbf{laminando} o núcleo
--- assunto retomado no bloco de desafio.
\end{itemize}
\textbf{Propriedade importante:} a força de frenagem é proporcional à
\emph{velocidade}. O freio é forte em alta velocidade e desaparece no repouso ---
por isso ele nunca segura um veículo parado.}
\end{questao}

\begin{questao}[2]{}
Uma bobina com núcleo de ferro tem $\SI{500}{}$ espiras. O fluxo varia de
$\SI{0}{}$ a $\SI{1.2}{\milli\weber}$ em $\SI{8}{\milli\second}$. Determine a
f.e.m. média.
\resposta{$\varepsilon=\SI{75}{\volt}$}
\resolucao{
\[
|\varepsilon|=N\frac{\Delta\Phi}{\Delta t}
=500\,\frac{\pot{1.2}{-3}}{\pot{8}{-3}}=\SI{75}{\volt}.
\]
\textbf{Papel do núcleo:} sem ele, o mesmo enrolamento e a mesma corrente
produziriam fluxo centenas de vezes menor --- e f.e.m. proporcionalmente menor.
O ferro é o que torna a indução tecnologicamente útil.\\
\textbf{Ressalva:} se o fluxo variar a partir de um valor já próximo da
saturação, a relação deixa de ser linear e a f.e.m. real fica abaixo da
prevista.}
\end{questao}

\begin{questao}[2]{}
Uma espira é conectada a um resistor e um ímã é afastado dela. Determine o
sentido da força que a espira exerce sobre o ímã.
\resposta{Atrativa --- ela tenta impedir o afastamento}
\resolucao{Ao afastar o ímã, o fluxo através da espira \textbf{diminui}. Pela lei
de Lenz, a corrente induzida circula no sentido de \emph{sustentar} o fluxo ---
criando campo no mesmo sentido do original.\\
A espira comporta-se então como um ímã com o polo oposto voltado ao ímã que se
afasta, e portanto o \textbf{atrai}.\\
\textbf{Pela terceira lei de Newton}, o ímã atrai a espira com força igual e
contrária.\\
\textbf{Conclusão geral:} a espira sempre \emph{resiste} ao que se está fazendo.
Aproximar exige empurrar contra repulsão; afastar exige puxar contra atração.
Em ambos os casos, o trabalho realizado é o que aparece como energia elétrica
--- e o resistor aquece.}
\end{questao}

\begin{questao}[2]{}
Uma bobina de $200$ espiras, área $\SI{100}{\centi\meter\squared}$ e
resistência $\SI{40}{\ohm}$ está em campo de $\SI{0.8}{\tesla}$
perpendicular. O campo é desligado em $\SI{20}{\milli\second}$.
\begin{itensq}
\item Determine a f.e.m. e a corrente induzidas.
\item Determine o sentido da corrente em relação ao campo original.
\item Determine a carga total.
\end{itensq}
\resposta{$\varepsilon=\SI{80}{\volt}$; $I=\SI{2}{\ampere}$;
$q=\SI{40}{\milli\coulomb}$}
\resolucao{(a)
\[
|\varepsilon|=N\frac{\Delta\Phi}{\Delta t}
=200\,\frac{(0{,}8)(\pot{1}{-2})}{\pot{2}{-2}}=\SI{80}{\volt};
\qquad
I=\frac{80}{40}=\SI{2}{\ampere}.
\]
(b) O fluxo \textbf{diminui}, logo a corrente induzida circula no sentido de
\emph{sustentá-lo} --- criando campo próprio no \textbf{mesmo} sentido do campo
original.\\
\emph{(Se o campo estivesse sendo ligado, a corrente circularia ao contrário.)}\\
(c)
\[
q=\frac{N\,\Delta\Phi}{R}=\frac{200(\pot{8}{-3})}{40}
=\pot{4}{-2}\,\si{\coulomb}.
\]
\textbf{Observação de segurança:} $\SI{80}{\volt}$ surgem em uma bobina que
antes não tinha tensão alguma. Desligar bruscamente eletroímãs de grande porte
produz tensões muito superiores --- razão pela qual eles exigem circuitos de
desmagnetização controlada.}
\end{questao}

\blocoaprofundamento

\begin{questao}[3]{}
Uma barra de $\SI{100}{\gram}$ e $\SI{40}{\centi\meter}$ desliza sem atrito em
trilhos \textbf{verticais}, em campo horizontal de $\SI{0.5}{\tesla}$, com
resistência total de $\SI{0.5}{\ohm}$. Ela é solta do repouso.
\begin{itensq}
\item Escreva a equação do movimento.
\item Determine a velocidade-limite.
\item Determine a potência dissipada nesse regime.
\end{itensq}
\dados{$g=\SI{9.8}{\meter\per\second\squared}$}
\resposta{$v_{lim}=\SI{12.25}{\meter\per\second}$; $P=\SI{12}{\watt}$}
\resolucao{(a) A barra sofre o peso e a força magnética de oposição:
\[
m\frac{\mathrm{d}v}{\mathrm{d}t}=mg-\frac{B^{2}\ell^{2}v}{R},
\]
pois $F=BI\ell$ com $I=B\ell v/R$.\\
\textbf{Estrutura da equação:} é idêntica à de um corpo em queda com arrasto
\emph{linear} --- a força de oposição é proporcional à velocidade.\\
(b) Na velocidade-limite, a aceleração se anula:
\[
v_{lim}=\frac{mgR}{B^{2}\ell^{2}}
=\frac{(0{,}1)(9{,}8)(0{,}5)}{(0{,}25)(0{,}16)}
=\frac{0{,}49}{0{,}04}=\SI{12.25}{\meter\per\second}.
\]
(c)
\[
P=mgv_{lim}=(0{,}1)(9{,}8)(12{,}25)=\SI{12}{\watt},
\]
que é exatamente $\dfrac{(B\ell v_{lim})^{2}}{R}$ \checkmark.\\
\textbf{Interpretação:} em regime, toda a energia potencial gravitacional
perdida vira calor no resistor. A barra não ganha mais energia cinética.\\
\textbf{Efeito de $R$:} reduzir $R$ \emph{diminui} a velocidade-limite --- o
freio fica mais forte. Com $R\to0$ (trilhos supercondutores), a barra mal se
moveria.}
\end{questao}

\begin{questao}[3]{}
Uma bobina de $N$ espiras e resistência $R$ é retirada de um campo. Demonstre
que a carga induzida independe da velocidade do movimento.
\begin{itensq}
\item Deduza a expressão.
\item Explique fisicamente.
\item Indique a aplicação em instrumentação.
\end{itensq}
\resposta{$q=\dfrac{N\Delta\Phi}{R}$ --- independe de $\Delta t$}
\resolucao{(a)
\[
q=\int i\,\mathrm{d}t
=\int\frac{N}{R}\frac{\mathrm{d}\Phi}{\mathrm{d}t}\,\mathrm{d}t
=\frac{N}{R}\int\mathrm{d}\Phi
=\frac{N\,\Delta\Phi}{R}.
\]
O tempo \textbf{cancela} na integração.\\
(b) \textbf{Explicação física.} Retirar a bobina depressa produz f.e.m. alta
durante pouco tempo; devagar, f.e.m. baixa durante muito tempo. O
\emph{produto} --- que é a carga --- permanece o mesmo.\\
Algebricamente: $\varepsilon\propto1/\Delta t$ e a duração é
$\Delta t$; o produto não depende dela.\\
(c) \textbf{Aplicação: fluxímetro balístico.} Um galvanômetro balístico responde
à \emph{carga} total, não à corrente instantânea. Conectado à bobina, sua
deflexão máxima é proporcional a $\Delta\Phi$ --- e o operador não precisa
controlar a velocidade com que retira a bobina.\\
\textbf{Robustez notável:} nem a velocidade, nem a suavidade, nem a trajetória
influenciam o resultado. É uma das medições mais tolerantes a erro de
procedimento em todo o eletromagnetismo.}
\end{questao}

\begin{questao}[3]{}
Um núcleo de transformador maciço de $\SI{10}{\milli\meter}$ de espessura é
substituído por $20$ lâminas isoladas de $\SI{0.5}{\milli\meter}$.
\begin{itensq}
\item Determine a redução das perdas por correntes parasitas.
\item Explique o mecanismo.
\item Indique a limitação prática.
\end{itensq}
\resposta{Redução de $400$ vezes}
\resolucao{(a) As perdas por correntes de Foucault em lâminas obedecem a
\[
P\propto\frac{B_{\max}^{2}f^{2}t^{2}}{\rho},
\]
com $t$ a espessura da lâmina. Reduzindo $t$ de $10$ para
$\SI{0.5}{\milli\meter}$:
\[
\frac{P'}{P}=\left(\frac{0{,}5}{10}\right)^{2}=\frac{1}{400}.
\]
(b) \textbf{Mecanismo.} As correntes parasitas circulam em laços
\emph{perpendiculares} ao fluxo. Laços maiores encerram mais fluxo (mais f.e.m.)
e encontram menos resistência (caminho mais largo) --- e por isso a perda cresce
com o \emph{quadrado} da dimensão.\\
Isolar as lâminas quebra os laços grandes em muitos laços pequenos, cada um com
f.e.m. menor e resistência maior.\\
\textbf{Importante:} as lâminas devem ser paralelas ao fluxo e isoladas
\emph{entre si} --- um verniz de poucos micrômetros basta. Se houver rebarba
metálica unindo-as, o efeito se perde.\\
(c) \textbf{Limitações.}
\begin{itemize}
\item \textbf{Fator de empilhamento:} o isolamento ocupa espaço, e a seção
magnética efetiva cai para $95$ a $97\%$ da geométrica.
\item \textbf{Custo:} lâminas mais finas são mais caras de produzir e montar.
\item \textbf{Frequência:} como $P\propto f^{2}t^{2}$, em alta frequência nem
lâminas finíssimas bastam --- daí o uso de \textbf{ferrites}, que são cerâmicas
de altíssima resistividade e praticamente não conduzem correntes parasitas.
\end{itemize}}
\end{questao}

\begin{questao}[3]{}
Um solenoide longo de $n_1$ espiras por metro e seção $A$ tem, em torno de si,
uma bobina de $N_2$ espiras.
\begin{itensq}
\item Deduza a indutância mútua.
\item Avalie para $n_1=\SI{2000}{\per\meter}$, $A=\SI{5}{\centi\meter\squared}$
      e $N_2=100$.
\item Explique por que $M$ não depende do raio da bobina externa.
\end{itensq}
\resposta{$M=\mu_0n_1N_2A$; $M=\SI{1.26}{\milli\henry}$}
\resolucao{(a) A corrente $i_1$ no solenoide cria $B=\mu_0n_1i_1$ em seu
interior, e fluxo $\Phi=BA$ por espira da bobina externa. O fluxo concatenado
nela é
\[
\lambda_2=N_2\mu_0n_1i_1A
\;\Longrightarrow\;
M=\frac{\lambda_2}{i_1}=\mu_0n_1N_2A.
\]
(b)
\[
M=(4\pi\times10^{-7})(2000)(100)(\pot{5}{-4})
=\pot{1.26}{-4}\,\si{\henry}=\SI{0.126}{\milli\henry}.
\]
(c) \textbf{Por que o raio externo não importa.} O campo do solenoide longo é
praticamente nulo do lado de fora. Toda a contribuição ao fluxo vem da região
\emph{interna} ao solenoide, cuja área é $A$ --- não a área da bobina
externa.\\
Enrolar a bobina justa ou folgada dá o \emph{mesmo} $M$.\\
\textbf{Consequência de projeto:} para aumentar o acoplamento, aumente a seção
do solenoide, não o diâmetro do enrolamento secundário. E note que
$M=M_{21}=M_{12}$ --- calcular pelo caminho fácil (do solenoide para a bobina) dá
o mesmo resultado que o caminho difícil.}
\end{questao}

\begin{questao}[3]{}
Uma espira retangular está próxima a um fio retilíneo cuja corrente varia
segundo $i(t)=I_0\sin\omega t$.
\begin{itensq}
\item Escreva o fluxo através da espira.
\item Determine a f.e.m. induzida.
\item Explique por que essa interferência é difícil de eliminar.
\end{itensq}
\resposta{$\varepsilon=-M\omega I_0\cos\omega t$, com
$M=\dfrac{\mu_0a}{2\pi}\ln\dfrac{d+b}{d}$}
\resolucao{(a) Do resultado já obtido para o fluxo de um fio através de espira
retangular:
\[
\Phi(t)=\frac{\mu_0\,i(t)\,a}{2\pi}\ln\!\left(\frac{d+b}{d}\right)
=M\,i(t),
\]
com $M$ a indutância mútua entre fio e espira.\\
(b)
\[
\varepsilon=-\frac{\mathrm{d}\Phi}{\mathrm{d}t}
=-M\,\frac{\mathrm{d}i}{\mathrm{d}t}
=-M\omega I_0\cos\omega t.
\]
A amplitude é $M\omega I_0$ --- \textbf{proporcional à frequência}.\\
(c) \textbf{Por que é difícil eliminar.}
\begin{itemize}
\item $M$ depende do \textbf{logaritmo} da razão de distâncias. Afastar a espira
reduz a interferência muito devagar --- dobrar $d$ e $b$ juntos não muda
\emph{nada}.
\item A amplitude cresce com $\omega$: transitórios rápidos e harmônicos de alta
ordem acoplam-se muito mais que a fundamental.
\item Blindagem magnética é cara e pouco eficaz em baixa frequência, ao
contrário da blindagem elétrica.
\end{itemize}
\textbf{O que realmente funciona:} reduzir a \emph{área} da espira receptora
(cabos trançados, roteamento junto ao retorno) e reduzir a área da fonte (par
trançado no lado do fio). Atacar a geometria, não a distância.}
\end{questao}

\begin{questao}[3]{}
Uma bobina de $\SI{200}{\milli\henry}$ conduz corrente que varia de
$\SI{2}{\ampere}$ a zero em $\SI{5}{\milli\second}$.
\begin{itensq}
\item Determine a f.e.m. autoinduzida.
\item Determine a energia envolvida.
\item Explique o risco associado.
\end{itensq}
\resposta{$\varepsilon=\SI{80}{\volt}$; $W=\SI{0.4}{\joule}$}
\resolucao{(a)
\[
|\varepsilon|=L\left|\frac{\Delta i}{\Delta t}\right|
=0{,}200\,\frac{2}{\pot{5}{-3}}=\SI{80}{\volt}.
\]
(b) $W=\frac12Li^{2}=\frac12(0{,}2)(4)=\SI{0.4}{\joule}$.\\
(c) \textbf{Risco.} A f.e.m. autoinduzida opõe-se à \emph{redução} da corrente
--- ela tenta mantê-la circulando. Se a interrupção for mais rápida, a tensão
sobe proporcionalmente: interromper em $\SI{5}{\micro\second}$ produziria
$\SI{80}{\kilo\volt}$.\\
Na prática, o ar rompe muito antes e forma-se \textbf{arco elétrico} nos
contatos, que os corrói e gera interferência.\\
\textbf{Proteção:} diodo de roda livre em antiparalelo, que oferece caminho à
corrente e limita a tensão. Os $\SI{0.4}{\joule}$ então se dissipam na
resistência do próprio circuito, sem arco.\\
\textbf{Ponto conceitual:} a autoindução é indução mútua da bobina
\emph{consigo mesma} --- a mesma lei de Faraday, aplicada ao próprio fluxo.}
\end{questao}

\begin{questao}[3]{}
Um gerador fornece $\SI{500}{\watt}$ a uma carga. Analise o torque necessário
para acioná-lo e sua relação com a lei de Lenz.
\begin{itensq}
\item Determine o torque a $\SI{1800}{rpm}$, supondo rendimento de
      $\SI{90}{\percent}$.
\item Explique o que acontece se a carga for desligada.
\end{itensq}
\resposta{$\tau=\SI{2.95}{\newton\meter}$; sem carga, o torque cai quase a
zero}
\resolucao{(a) A potência mecânica de entrada é
\[
P_{mec}=\frac{500}{0{,}90}=\SI{556}{\watt};
\qquad
\omega=\frac{1800(2\pi)}{60}=\SI{188.5}{\radian\per\second};
\]
\[
\tau=\frac{P_{mec}}{\omega}=\frac{556}{188{,}5}=\SI{2.95}{\newton\meter}.
\]
(b) \textbf{Sem carga}, não circula corrente no enrolamento. Sem corrente, não
há força magnética de oposição --- e o torque necessário cai a praticamente
zero (resta apenas o atrito e as perdas em vazio).\\
\textbf{É a lei de Lenz em ação.} O torque resistente \emph{surge} porque a
corrente induzida cria campo que se opõe ao movimento. Mais carga significa mais
corrente, mais oposição, mais torque exigido.\\
\textbf{Experiência reveladora:} girar um gerador manualmente com os terminais
abertos é fácil; curto-circuitá-los torna a rotação quase impossível. A energia
elétrica não vem "do campo" --- vem do braço de quem gira.}
\end{questao}

\begin{questao}[3]{}
Um anel condutor é colocado sobre o núcleo de um solenoide alimentado por
corrente alternada. Observa-se que ele é \textbf{arremessado} para cima.
\begin{itensq}
\item Explique o fenômeno.
\item Determine o que muda se o anel for cortado.
\item Determine o que muda com corrente contínua em regime.
\end{itensq}
\resposta{Repulsão por corrente induzida; anel cortado não salta; em CC de
regime, nada acontece}
\resolucao{(a) \textbf{Explicação.} A corrente alternada no solenoide produz
fluxo variável através do anel, induzindo nele corrente. Pela lei de Lenz, essa
corrente cria campo que se \emph{opõe} ao do solenoide --- e dois campos opostos
se \textbf{repelem}.\\
Como a repulsão persiste em ambos os semiciclos (a corrente induzida inverte
junto com a indutora), a força média é sempre para cima.\\
\emph{(É o experimento clássico do "anel saltante" de Elihu Thomson.)}\\
(b) \textbf{Anel cortado:} não há caminho fechado, logo não circula corrente. Sem
corrente, não há campo induzido nem força. \textbf{O anel não salta} --- embora a
f.e.m. continue sendo induzida em seus terminais.\\
Este é o teste decisivo: mostra que a força depende da \emph{corrente}, não da
f.e.m.\\
(c) \textbf{Corrente contínua em regime:} o fluxo é constante, não há indução,
não há força. \textbf{Nada acontece.}\\
\emph{Mas no instante em que a corrente CC é \textbf{ligada}, há um transitório
de fluxo --- e o anel dá um salto único.} Isso reforça o ponto: é a
\emph{variação} que induz, não o campo.}
\end{questao}

\begin{questao}[3]{}
Uma espira de área $\SI{50}{\centi\meter\squared}$ está em campo que varia
segundo $B(t)=0{,}4\sin(100\pi t)$ tesla, e \emph{simultaneamente} sua área
diminui a $\SI{10}{\centi\meter\squared\per\second}$.
\begin{itensq}
\item Escreva a expressão geral da f.e.m.
\item Identifique as duas contribuições.
\end{itensq}
\resposta{$\varepsilon=-\left(A\dfrac{\mathrm{d}B}{\mathrm{d}t}
+B\dfrac{\mathrm{d}A}{\mathrm{d}t}\right)$}
\resolucao{(a) Como $\Phi=BA$ com \emph{ambos} variáveis, aplica-se a regra do
produto:
\[
\varepsilon=-\frac{\mathrm{d}}{\mathrm{d}t}(BA)
=-\left(A\frac{\mathrm{d}B}{\mathrm{d}t}+B\frac{\mathrm{d}A}{\mathrm{d}t}\right).
\]
(b) \textbf{Duas contribuições de naturezas diferentes.}
\begin{itemize}
\item \textbf{Termo $A\,\mathrm{d}B/\mathrm{d}t$:} f.e.m. induzida por
\emph{campo variável} --- o mecanismo do transformador. Aqui,
$\mathrm{d}B/\mathrm{d}t=40\pi\cos(100\pi t)$, com amplitude
$\SI{125.7}{\tesla\per\second}$, dando até
$\SI{0.63}{\volt}$.
\item \textbf{Termo $B\,\mathrm{d}A/\mathrm{d}t$:} f.e.m. por \emph{movimento}
--- o mecanismo do gerador. Com
$\mathrm{d}A/\mathrm{d}t=\SI{-1e-3}{\meter\squared\per\second}$, contribui com
até $\SI{0.4}{\milli\volt}$.
\end{itemize}
\textbf{Neste caso o primeiro termo domina} por três ordens de grandeza.\\
\textbf{Ponto conceitual:} os dois mecanismos são fisicamente distintos --- um
envolve campo elétrico induzido, o outro força magnética sobre portadores --- mas
a lei de Faraday os unifica em uma única expressão. Essa unificação é uma das
observações que levaram Einstein à relatividade restrita.}
\end{questao}

\begin{questao}[3]{}
Uma bobina de $500$ espiras, área $\SI{20}{\centi\meter\squared}$ e resistência
$\SI{25}{\ohm}$ é girada de $\ang{180}$ em campo de $\SI{0.3}{\tesla}$.
\begin{itensq}
\item Determine a variação de fluxo por espira.
\item Determine a carga total induzida.
\end{itensq}
\resposta{$\Delta\Phi=\SI{1.2}{\milli\weber}$; $q=\SI{24}{\milli\coulomb}$}
\resolucao{(a) Girar $\ang{180}$ inverte o sinal do fluxo:
\[
\Phi_i=BA=(0{,}3)(\pot{2}{-3})=\SI{0.6}{\milli\weber};
\qquad
\Phi_f=-\SI{0.6}{\milli\weber};
\]
\[
|\Delta\Phi|=\SI{1.2}{\milli\weber}\quad\text{--- o \textbf{dobro} de }
\Phi_i.
\]
(b)
\[
q=\frac{N|\Delta\Phi|}{R}=\frac{500(\pot{1.2}{-3})}{25}
=\pot{2.4}{-2}\,\si{\coulomb}.
\]
\textbf{Erro comum:} usar $\Delta\Phi=BA$ em vez de $2BA$. A inversão completa
faz o fluxo passar de $+\Phi$ a $-\Phi$, e a variação é o dobro do valor
inicial.\\
\textbf{Vantagem prática do giro de $\ang{180}$:} ele dobra o sinal em relação a
simplesmente retirar a bobina do campo --- técnica padrão em medições com
fluxímetro.}
\end{questao}

\begin{questao}[3]{}
Um transformador com $k=0{,}95$ e $L_1=\SI{100}{\milli\henry}$,
$L_2=\SI{25}{\milli\henry}$.
\begin{itensq}
\item Determine a indutância mútua.
\item Determine a f.e.m. no secundário para
      $\mathrm{d}i_1/\mathrm{d}t=\SI{50}{\ampere\per\second}$.
\item Determine a fração de fluxo que não se acopla.
\end{itensq}
\resposta{$M=\SI{47.5}{\milli\henry}$; $\varepsilon_2=\SI{2.38}{\volt}$;
$5\%$ de dispersão}
\resolucao{(a)
\[
M=k\sqrt{L_1L_2}=0{,}95\sqrt{(0{,}100)(0{,}025)}
=0{,}95(0{,}05)=\SI{47.5}{\milli\henry}.
\]
(b)
\[
\varepsilon_2=M\frac{\mathrm{d}i_1}{\mathrm{d}t}=(0{,}0475)(50)
=\SI{2.38}{\volt}.
\]
(c) A fração acoplada é $k=95\%$; portanto \textbf{$5\%$ do fluxo se
dispersa}, não atravessando o secundário.\\
\textbf{Efeito da dispersão:} ela aparece como indutância em série
(\emph{indutância de dispersão}) e provoca queda de tensão sob carga. Um
transformador com $k=0{,}95$ tem regulação bem pior que um com
$k=0{,}999$.\\
\textbf{Comparação de topologias:} núcleos toroidais fechados atingem
$k>0{,}99$; núcleos em E com entreferro parasita ficam em torno de $0{,}95$;
transformadores com enrolamentos lado a lado, bem menos.}
\end{questao}

\begin{questao}[3]{}
Analise o perfil de f.e.m. produzido por um fluxo que cresce linearmente por
$\SI{2}{\second}$, permanece constante por $\SI{3}{\second}$ e decresce
linearmente por $\SI{1}{\second}$, com amplitude de $\SI{6}{\milli\weber}$ e
$N=100$.
\begin{itensq}
\item Determine a f.e.m. em cada trecho.
\item Verifique a área total sob a curva de f.e.m.
\end{itensq}
\resposta{$+0{,}3$; $0$; $\SI{-0.6}{\volt}$; área líquida nula}
\resolucao{(a)
\begin{center}
\begin{tabular}{lccc}
\hline
Trecho & $\mathrm{d}\Phi/\mathrm{d}t$ & $\varepsilon$ & Duração\\
\hline
Subida & $\SI{3}{\milli\weber\per\second}$ & $\SI{-0.3}{\volt}$
 & $\SI{2}{\second}$\\
Platô & $0$ & $0$ & $\SI{3}{\second}$\\
Descida & $\SI{-6}{\milli\weber\per\second}$ & $\SI{+0.6}{\volt}$
 & $\SI{1}{\second}$\\
\hline
\end{tabular}
\end{center}
(Os sinais decorrem de $\varepsilon=-N\,\mathrm{d}\Phi/\mathrm{d}t$.)\\
\textbf{Observe:} a descida é duas vezes mais rápida, e a f.e.m. é o
\emph{dobro} em módulo.\\
(b) \textbf{Área sob a curva:}
\[
(-0{,}3)(2)+(0)(3)+(0{,}6)(1)=-0{,}6+0{,}6=0.
\]
\textbf{Isso não é coincidência.} A área é $\int\varepsilon\,\mathrm{d}t
=-N\Delta\Phi_{total}$, e o fluxo terminou onde começou. \textbf{Sempre que o
fluxo retorna ao valor inicial, a integral da f.e.m. é nula.}\\
\textbf{Consequência prática:} em regime periódico, a tensão média sobre um
indutor é sempre zero --- princípio do \emph{balanço volt-segundo}, base do
projeto de conversores chaveados.}
\end{questao}

\begin{questao}[3]{}
Um freio magnético produz força proporcional à velocidade, $F=bv$, com
$b=\SI{20}{\newton\second\per\meter}$, atuando sobre um corpo de
$\SI{5}{\kilogram}$ inicialmente a $\SI{10}{\meter\per\second}$.
\begin{itensq}
\item Escreva a equação do movimento e resolva.
\item Determine a constante de tempo e a distância percorrida.
\item Explique por que o freio não para o corpo.
\end{itensq}
\resposta{$v=v_0e^{-t/\tau}$ com $\tau=\SI{0.25}{\second}$;
$d=\SI{2.5}{\meter}$}
\resolucao{(a)
\[
m\frac{\mathrm{d}v}{\mathrm{d}t}=-bv
\;\Longrightarrow\;
v(t)=v_0e^{-bt/m}.
\]
(b) $\tau=m/b=5/20=\SI{0.25}{\second}$.\\
A distância total é
\[
d=\int_0^{\infty}v_0e^{-t/\tau}\mathrm{d}t=v_0\tau=10(0{,}25)=\SI{2.5}{\meter}.
\]
\textbf{Mesma estrutura} da descarga exponencial de um capacitor --- e a mesma
propriedade: a "área" (distância) equivale à velocidade inicial mantida por
$\tau$.\\
(c) \textbf{Por que não para.} A força é proporcional a $v$: quando a velocidade
tende a zero, a força também tende a zero. Matematicamente, o corpo leva tempo
\emph{infinito} para parar --- ele apenas se aproxima assintoticamente do
repouso.\\
\textbf{Consequência prática decisiva:} freios magnéticos são excelentes para
\emph{desacelerar} e péssimos para \emph{segurar}. Trens de alta velocidade os
usam na faixa alta e completam a parada com freios de atrito --- que funcionam
justamente onde o magnético falha.}
\end{questao}

\blocodesafio

\begin{questao}[4]{}
\textbf{(Lenz e conservação da energia.)} Demonstre que a lei de Lenz é uma
consequência necessária da conservação da energia.
\begin{itensq}
\item Suponha o contrário e derive a contradição.
\item Quantifique com o exemplo da barra em trilhos.
\item Discuta o que a lei de Lenz \emph{não} determina.
\end{itensq}
\resposta{Sinal oposto levaria a crescimento ilimitado de energia}
\resolucao{(a) \textbf{Demonstração por absurdo.} Suponha que a corrente
induzida \emph{reforçasse} a variação de fluxo. Ao aproximar um ímã de uma
espira, ela o \textbf{atrairia} --- acelerando a aproximação, o que aumentaria
$\mathrm{d}\Phi/\mathrm{d}t$, que aumentaria a corrente, que aumentaria a
atração.\\
O sistema entraria em fuga: energia cinética \emph{e} energia elétrica cresceriam
simultaneamente, sem fonte. \textbf{Contradição} com a conservação da energia.\\
Portanto a corrente deve \emph{opor-se} --- e o sinal negativo não é convenção,
é necessidade.\\
(b) \textbf{Quantificação.} Na barra em trilhos:
\[
P_{mec}=Fv=\left(\frac{B^{2}\ell^{2}v}{R}\right)v
=\frac{B^{2}\ell^{2}v^{2}}{R};
\qquad
P_{el}=\frac{(B\ell v)^{2}}{R}=\frac{B^{2}\ell^{2}v^{2}}{R}.
\]
\textbf{Idênticas, termo a termo.} A força de oposição tem exatamente o módulo
necessário para que a potência mecânica iguale a elétrica --- nem mais, nem
menos.\\
\textbf{Se a força fosse maior}, sobraria energia mecânica sem destino; se fosse
menor, apareceria energia elétrica do nada.\\
(c) \textbf{O que a lei de Lenz não determina.} Ela fixa apenas o \emph{sentido}
da corrente. O \textbf{módulo} vem da lei de Faraday combinada com a
resistência do circuito.\\
E ela não diz nada sobre \emph{quanta} energia é convertida --- isso depende da
carga. Uma espira aberta induz f.e.m. mas não corrente, e não há oposição
alguma: aproximar um ímã de uma espira cortada não exige força.}
\end{questao}

\begin{questao}[4]{}
\textbf{(Projeto de bobina sensora.)} Projete uma bobina que gere
$\SI{5}{\volt}$ médios quando o fluxo variar $\SI{0.2}{\milli\weber}$ em
$\SI{4}{\milli\second}$.
\begin{itensq}
\item Determine o número de espiras.
\item Determine a resistência admissível para corrente de saída inferior a
      $\SI{1}{\milli\ampere}$ em carga de $\SI{10}{\kilo\ohm}$.
\item Avalie a resposta em frequência.
\end{itensq}
\resposta{$N=100$ espiras}
\resolucao{(a)
\[
N=\frac{\varepsilon\,\Delta t}{\Delta\Phi}
=\frac{5(\pot{4}{-3})}{\pot{2}{-4}}=100\ \text{espiras}.
\]
(b) Com carga de $\SI{10}{\kilo\ohm}$ e $\SI{5}{\volt}$, a corrente é
$\SI{0.5}{\milli\ampere}$ --- já abaixo do limite. A queda na resistência
interna deve ser pequena:
\[
\frac{R_{bobina}}{R_{bobina}+10\,\mathrm{k}}\le1\%
\;\Longrightarrow\;
R_{bobina}\le\SI{100}{\ohm}.
\]
Com $100$ espiras, isso é facilmente atingível --- fio de $\SI{0.2}{\milli\meter}$
e alguns metros de comprimento dão poucos ohms.\\
(c) \textbf{Resposta em frequência.} A bobina é um derivador: sua saída é
proporcional a $\mathrm{d}\Phi/\mathrm{d}t$, logo a sensibilidade
\textbf{cresce com a frequência}.
\begin{itemize}
\item \textbf{Vantagem:} excelente para detectar transitórios rápidos e
descargas.
\item \textbf{Limitação inferior:} ela \emph{não} detecta campo estático ---
para isso é preciso sensor Hall ou fluxgate.
\item \textbf{Limitação superior:} a indutância própria da bobina, associada à
capacitância entre espiras, forma um ressonador. Acima da frequência de
autorressonância a resposta deixa de ser proporcional a
$\mathrm{d}\Phi/\mathrm{d}t$.
\end{itemize}
\textbf{Para recuperar $\Phi$ (e não sua derivada):} acrescente um integrador na
saída --- é assim que funciona a bobina de Rogowski para medição de corrente.}
\end{questao}

\begin{questao}[4]{}
\textbf{(Freio por correntes parasitas.)} Modele quantitativamente um freio
magnético de disco.
\begin{itensq}
\item Estabeleça a dependência da força com os parâmetros.
\item Determine por que a força é proporcional à velocidade.
\item Analise o comportamento em velocidades muito altas.
\end{itensq}
\resposta{$F\propto\sigma B^{2}tv$; proporcionalidade a $v$ decorre de
$\varepsilon\propto v$ e $i\propto\varepsilon$}
\resolucao{(a) \textbf{Dependência.} A f.e.m. induzida localmente é
proporcional a $Bv$; a corrente parasita, a $\sigma Bv$ (com $\sigma$ a
condutividade); e a força, ao produto da corrente pelo campo:
\[
F\propto\sigma\,B^{2}\,t\,v\,A,
\]
com $t$ a espessura e $A$ a área sob o polo.\\
\textbf{Note a dependência quadrática em $B$} --- dobrar o campo quadruplica a
frenagem, o que torna ímãs de neodímio muito mais eficazes que ferrites.\\
(b) \textbf{Por que $F\propto v$.} Dois fatores lineares se combinam:
\[
\varepsilon\propto v
\;\Longrightarrow\;
i\propto v
\;\Longrightarrow\;
F=Bi\ell\propto v.
\]
É o mesmo mecanismo do arrasto viscoso --- e leva à mesma equação diferencial,
com decaimento exponencial da velocidade.\\
(c) \textbf{Velocidades muito altas.} A proporcionalidade \emph{deixa de valer}
por dois motivos:
\begin{itemize}
\item \textbf{Reação de armadura:} as correntes parasitas tornam-se intensas o
bastante para criar campo próprio comparável ao aplicado, distorcendo-o. A força
satura e chega a \emph{diminuir}.
\item \textbf{Efeito pelicular:} em alta frequência efetiva, as correntes se
concentram na superfície do disco, reduzindo a espessura útil $t$.
\end{itemize}
\textbf{Existe, portanto, uma velocidade de força máxima}, acima da qual é o freio
perde eficácia --- comportamento oposto ao do freio de atrito e que precisa ser
levado em conta no dimensionamento.}
\end{questao}

\begin{questao}[4]{}
\textbf{(Convenção dos pontos na indução mútua.)} Explique o significado físico
da convenção dos pontos e sua determinação experimental.
\begin{itensq}
\item Enuncie a convenção.
\item Determine o sinal de $M$ nas equações de malha.
\item Descreva o procedimento experimental.
\end{itensq}
\resposta{Pontos marcam terminais magneticamente correspondentes; o sinal de
$M$ depende de as correntes entrarem ou não pelos pontos}
\resolucao{(a) \textbf{Convenção.} Correntes que entram pelos terminais
\emph{pontuados} produzem fluxos que se \textbf{reforçam} no núcleo comum.\\
Isso é informação \emph{física} (sentido de enrolamento), que o desenho
esquemático não transmite --- daí a necessidade da marcação.\\
(b) \textbf{Sinal nas equações.} Para duas bobinas acopladas:
\[
v_1=L_1\frac{\mathrm{d}i_1}{\mathrm{d}t}
\pm M\frac{\mathrm{d}i_2}{\mathrm{d}t};
\qquad
v_2=L_2\frac{\mathrm{d}i_2}{\mathrm{d}t}
\pm M\frac{\mathrm{d}i_1}{\mathrm{d}t}.
\]
Usa-se $+M$ se \emph{ambas} as correntes entram pelos pontos (ou ambas saem), e
$-M$ caso contrário.\\
\textbf{Em série}, isso dá $L_{eq}=L_1+L_2\pm2M$ --- resultado já obtido no banco
de indutores.\\
(c) \textbf{Procedimento experimental.} Meça $L_{eq}$ com as bobinas ligadas em
série nas \emph{duas} formas possíveis:
\[
M=\frac{L_{aditivo}-L_{subtrativo}}{4}.
\]
A ligação de \emph{maior} indutância identifica o acoplamento aditivo --- e
portanto os terminais correspondentes.\\
\textbf{Método alternativo, mais direto:} aplique um pulso de tensão a uma
bobina e observe a polaridade do pulso induzido na outra com osciloscópio. Se
ambas as pontas positivas estiverem nos terminais correspondentes, os sinais têm
a mesma polaridade.\\
\textbf{Por que isso importa na prática:} inverter a marcação de um
transformador em um conversor inverte a fase da saída --- e pode transformar
realimentação negativa em positiva, destruindo o circuito.}
\end{questao}

\begin{questao}[4]{}
\textbf{(Verificação experimental da lei de Faraday.)} Projete um experimento
para verificar $\varepsilon\propto N$ e $\varepsilon\propto v$.
\begin{itensq}
\item Descreva o arranjo e o procedimento.
\item Estabeleça o critério de validação.
\item Aponte as principais fontes de erro.
\end{itensq}
\resposta{Ímã caindo por tubo com bobinas de $N$ variável; verificar
linearidade em $N$ e em $v$}
\resolucao{(a) \textbf{Arranjo.} Um ímã permanente cai por dentro de um tubo
isolante vertical, atravessando bobinas montadas ao longo dele. Um osciloscópio
registra a f.e.m. de cada bobina.
\begin{enumerate}
\item \textbf{Variar $N$:} bobinas de $50$, $100$, $200$ e $400$ espiras, todas
com a mesma área e posição relativa. Compare os picos de f.e.m.
\item \textbf{Variar $v$:} use \emph{a mesma} bobina em alturas diferentes. Como
$v=\sqrt{2gh}$, a velocidade é conhecida a partir da altura de queda.
\end{enumerate}
(b) \textbf{Critério de validação.}
\begin{itemize}
\item Gráfico de $\varepsilon_{pico}$ contra $N$: deve ser uma \textbf{reta pela
origem}. Ajuste linear com $R^{2}>0{,}99$ e intercepto compatível com zero.
\item Gráfico de $\varepsilon_{pico}$ contra $\sqrt{h}$: também reta pela
origem.
\item \textbf{Teste mais forte:} a \emph{área} sob cada pulso de f.e.m. deve ser
$N\Delta\Phi$ --- \emph{independente} da velocidade. Verificar isso confirma
simultaneamente a lei e a independência da carga induzida.
\end{itemize}
(c) \textbf{Fontes de erro.}
\begin{itemize}
\item \textbf{Frenagem do próprio ímã} pelas correntes induzidas: se as bobinas
estiverem em curto ou a carga for baixa, o ímã desacelera e $v$ deixa de ser
$\sqrt{2gh}$. \emph{Solução:} medir em circuito aberto (osciloscópio de alta
impedância).
\item \textbf{Tubo condutor:} um tubo de alumínio ou cobre freia
drasticamente o ímã. Use acrílico ou vidro.
\item \textbf{Posição de disparo do osciloscópio} e largura de banda
insuficiente, que arredondam o pico.
\item \textbf{Orientação do ímã} variando durante a queda --- use guia.
\end{itemize}}
\end{questao}

\begin{questao}[4]{}
\textbf{(Transformador em corrente contínua.)} Explique por que um
transformador não transfere potência em regime de corrente contínua.
\begin{itensq}
\item Analise o regime permanente.
\item Analise o transitório de ligação.
\item Determine o que ocorre com o enrolamento primário.
\end{itensq}
\resposta{Fluxo constante $\Rightarrow$ $\mathrm{d}\Phi/\mathrm{d}t=0
\Rightarrow$ nenhuma f.e.m. no secundário}
\resolucao{(a) \textbf{Regime permanente.} Com corrente contínua constante, o
fluxo no núcleo é constante:
\[
\frac{\mathrm{d}\Phi}{\mathrm{d}t}=0
\;\Longrightarrow\;
\varepsilon_2=0.
\]
\textbf{Nenhuma tensão no secundário}, por maior que seja a corrente no
primário. A indução exige \emph{variação}.\\
(b) \textbf{Transitório de ligação.} No instante em que a fonte CC é ligada, o
fluxo sobe de zero ao valor final --- e \emph{há} f.e.m. no secundário durante
esse intervalo. Aparece um pulso, cuja integral é $N_2\Delta\Phi$.\\
O mesmo ocorre, com polaridade invertida, ao desligar.\\
\textbf{É esse pulso que se explora na bobina de ignição} de motores: uma
corrente contínua é interrompida abruptamente, e o pulso resultante atinge
dezenas de quilovolts.\\
(c) \textbf{O que ocorre com o primário --- e é o ponto crítico.} Em regime CC, a
única coisa que limita a corrente é a \textbf{resistência ôhmica} do
enrolamento, que é propositalmente pequena (poucos ohms).\\
Um primário projetado para $\SI{220}{\volt}$ em $\SI{60}{\hertz}$ tem sua
corrente limitada pela \emph{reatância} $\omega L$, tipicamente centenas de
ohms. Em CC, $\omega=0$ e essa limitação \textbf{desaparece}:
\[
I_{CC}=\frac{220}{R_{enrolamento}}\gg I_{CA}.
\]
\textbf{Resultado:} corrente dezenas de vezes maior que a nominal, saturação
imediata do núcleo e queima do enrolamento em segundos.\\
\textbf{Conclusão:} não é apenas que o transformador "não funciona" em CC --- ele
se \emph{destrói}.}
\end{questao}

\begin{questao}[4]{}
\textbf{(Espira entrando em campo.)} Deduza a f.e.m. durante a entrada de uma
espira retangular em uma região de campo uniforme, com velocidade constante.
\begin{itensq}
\item Deduza pela variação de área e pela força sobre portadores.
\item Determine a força necessária para manter a velocidade.
\item Analise o balanço de energia.
\end{itensq}
\resposta{$\varepsilon=B\ell v$ pelos dois caminhos;
$F=\dfrac{B^{2}\ell^{2}v}{R}$}
\resolucao{(a) \textbf{Pela variação de área.} Com o lado dianteiro a uma
profundidade $x$ dentro do campo, $A_{imersa}=\ell x$ e
\[
\Phi=B\ell x
\;\Longrightarrow\;
\varepsilon=-\frac{\mathrm{d}\Phi}{\mathrm{d}t}=-B\ell\frac{\mathrm{d}x}{\mathrm{d}t}
=-B\ell v.
\]
\textbf{Pela força sobre portadores.} Os portadores do lado dianteiro movem-se
com velocidade $v$ e sofrem força $qvB$ ao longo do lado. O trabalho por unidade
de carga ao percorrer o comprimento $\ell$ é
\[
\varepsilon=\frac{W}{q}=vB\ell.
\]
\textbf{Os dois caminhos coincidem} --- e essa coincidência não é trivial: um
usa a lei de Faraday (variação de fluxo), o outro a força de Lorentz. Sua
equivalência é uma das simetrias mais notáveis do eletromagnetismo.\\
(b) \textbf{Força necessária.} A corrente $i=B\ell v/R$ percorre o lado
dianteiro, imerso no campo, sofrendo
\[
F=Bi\ell=\frac{B^{2}\ell^{2}v}{R},
\]
oposta ao movimento. É essa força que o agente externo deve equilibrar.\\
\textbf{Note:} os lados laterais também estão parcialmente no campo, mas suas
forças são opostas entre si e se cancelam.\\
(c) \textbf{Balanço.}
\[
P_{mec}=Fv=\frac{B^{2}\ell^{2}v^{2}}{R}
=\frac{(B\ell v)^{2}}{R}=\frac{\varepsilon^{2}}{R}=P_{el}.
\]
\textbf{Fecha exatamente.} E note o que acontece quando a espira fica
\emph{totalmente dentro}: os dois lados no campo geram f.e.m. opostas, a
corrente cessa, a força desaparece --- e a espira desliza livremente. A oposição
existe apenas enquanto o fluxo varia.}
\end{questao}

\begin{questao}[4]{}
\textbf{(Balanço volt-segundo.)} Demonstre que, em regime periódico, a tensão
média sobre um indutor é nula, e explore as consequências.
\begin{itensq}
\item Demonstre.
\item Aplique a um conversor com dois patamares de tensão.
\item Explique o que ocorre se o balanço não for respeitado.
\end{itensq}
\resposta{$\bar v_L=0$; daí $V_1t_1=V_2t_2$ em regime}
\resolucao{(a) \textbf{Demonstração.} Em regime periódico de período $T$, a
corrente retorna ao mesmo valor a cada ciclo: $i(T)=i(0)$. Como
$v_L=L\,\mathrm{d}i/\mathrm{d}t$,
\[
\int_0^{T}v_L\,\mathrm{d}t=L\int_{i(0)}^{i(T)}\mathrm{d}i=L\left[i(T)-i(0)\right]=0.
\]
Logo $\bar v_L=0$. Equivalentemente, o \textbf{produto volt-segundo} aplicado
em um sentido deve igualar o aplicado no outro.\\
(b) \textbf{Aplicação a conversor.} Se o indutor recebe $+V_1$ durante
$t_1$ e $-V_2$ durante $t_2$:
\[
V_1t_1=V_2t_2.
\]
\textbf{Exemplo (conversor abaixador):} com entrada $\SI{12}{\volt}$, saída
$\SI{5}{\volt}$ e razão cíclica $D$, o indutor vê $(12-5)$ durante $DT$ e
$-5$ durante $(1-D)T$:
\[
7D=5(1-D)
\;\Longrightarrow\;
D=\frac{5}{12}=0{,}417.
\]
\textbf{O balanço volt-segundo produziu a relação de conversão} $V_o=DV_i$ sem
nenhuma análise de circuito --- é o método padrão de projeto.\\
(c) \textbf{Se o balanço não for respeitado.} A corrente do indutor não retorna
ao valor inicial: ela cresce (ou decresce) a cada ciclo, acumulando
indefinidamente.\\
Em um transformador, isso significa \textbf{saturação progressiva} do núcleo até
o colapso da indutância e a destruição do interruptor --- o fenômeno conhecido
como \emph{staircase saturation}.\\
\textbf{É por isso que conversores em ponte exigem simetria rigorosa} entre os
semiciclos, e por que se acrescenta capacitor de bloqueio em série com o
primário: ele impede qualquer componente contínua de tensão, garantindo o
balanço por construção.}
\end{questao}

\begin{questao}[4]{}
\textbf{(Correntes parasitas em condutores.)} Analise por que condutores
maciços de grande seção apresentam perdas superiores às previstas em corrente
alternada.
\begin{itensq}
\item Explique o efeito pelicular.
\item Determine a profundidade de penetração e avalie para cobre a
      $\SI{60}{\hertz}$.
\item Indique as soluções empregadas.
\end{itensq}
\resposta{$\delta=\sqrt{\dfrac{2\rho}{\omega\mu}}\approx\SI{8.5}{\milli\meter}$
para cobre a $\SI{60}{\hertz}$}
\resolucao{(a) \textbf{Efeito pelicular.} A corrente alternada em um condutor
cria fluxo variável \emph{dentro dele}. Esse fluxo induz correntes parasitas
que, pela lei de Lenz, se opõem à corrente no centro e a reforçam na periferia.\\
Resultado: a corrente concentra-se junto à superfície, a seção efetiva diminui e
a \textbf{resistência aumenta}.\\
(b) \textbf{Profundidade de penetração.}
\[
\delta=\sqrt{\frac{2\rho}{\omega\mu}}.
\]
Para cobre ($\rho=\pot{1.7}{-8}$, $\mu\approx\mu_0$) a
$\SI{60}{\hertz}$:
\[
\delta=\sqrt{\frac{2(\pot{1.7}{-8})}{2\pi(60)(4\pi\times10^{-7})}}
\approx\SI{8.5}{\milli\meter}.
\]
\textbf{Interpretação:} condutores com raio muito menor que
$\SI{8.5}{\milli\meter}$ praticamente não sofrem o efeito em
$\SI{60}{\hertz}$. Já um barramento de $\SI{50}{\milli\meter}$ tem seu centro
essencialmente inútil.\\
\textbf{Dependência com a frequência:} $\delta\propto1/\sqrt{f}$. A
$\SI{1}{\mega\hertz}$, $\delta$ cai para cerca de
$\SI{66}{\micro\meter}$.\\
(c) \textbf{Soluções.}
\begin{itemize}
\item \textbf{Condutores tubulares} ou perfis chatos, que aproveitam melhor a
periferia --- comuns em barramentos de subestação.
\item \textbf{Fio Litz:} muitos fios finos isolados e \emph{transpostos}, de modo
que cada um ocupe alternadamente centro e periferia. Usado em enrolamentos de
alta frequência.
\item \textbf{Vários condutores em paralelo}, transpostos ao longo do percurso.
\end{itemize}
\textbf{Observação:} apenas colocar fios em paralelo \emph{sem} transposição não
resolve --- os externos conduziriam mais que os internos, pelo mesmo mecanismo.
É a transposição que equaliza.}
\end{questao}



\begin{questao}[4]{}
 Um feixe de prótons atravessa inicialmente uma região onde atuam simultaneamente um campo elétrico uniforme de intensidade
\[
E = 70\,\mathrm{kV/m}
\]
e um campo magnético uniforme de intensidade
\[
B = 50\,\mathrm{mT},
\]
perpendiculares entre si e também à direção de movimento das partículas. Os campos são ajustados de modo que apenas os prótons que atravessam essa região sem sofrer desvio alcancem uma segunda região, na qual existe apenas o campo magnético de mesma intensidade, conforme ilustrado na figura.

Desprezando os efeitos gravitacionais e considerando
\[
m_p = 1{,}67\times10^{-27}\,\mathrm{kg},
\qquad
q_p = 1{,}60\times10^{-19}\,\mathrm{C},
\]
determine:

\begin{itensq}
    \item a velocidade dos prótons que atravessam o seletor sem sofrer desvio;
    \item o raio \(r\) da trajetória circular descrita pelos prótons na região onde atua apenas o campo magnético;
    \item a distância \(2r\) entre o ponto de entrada na região magnética e o ponto em que os prótons atingem a placa detectora.
\end{itensq}



 \begin{tikzpicture}[
    >=Latex,
    line cap=round,
    line join=round,
    font=\sffamily
]

% ==================================================
% PARÂMETROS
% ==================================================
\def\r{2.35}
\def\xleft{0.0}
\def\xfilter{3.4}
\def\xslit{6.45}
\pgfmathsetmacro{\xfield}{\xslit-0.35}
\def\ybeam{0.0}
\def\ybottom{-4.7}

% Centro do arco semicircular
\coordinate (C) at (\xfield,\ybottom/2);
\coordinate (TopArc) at ($(C)+(90:\r)$);
\coordinate (BotArc) at ($(C)+(-90:\r)$);

% ==================================================
% PRÓTON INICIAL
% ==================================================
\draw[thick] (-1.05,0) circle (0.23);
\node at (-1.05,0) {\large $+$};

% Feixe horizontal
\draw[thick,dashed] (-0.75,0) -- (\xfield,0);

% ==================================================
% FENDA DE ENTRADA
% ==================================================
\draw[line width=3.6pt] (0.12,0.32) -- (0.12,0.82);
\draw[line width=3.6pt] (0.12,-0.32) -- (0.12,-0.82);

% ==================================================
% SELETOR DE VELOCIDADES
% ==================================================
% Placas
\draw[line width=2.2pt] (2.15,0.78) -- (4.65,0.78);
\draw[line width=2.2pt] (2.15,-0.78) -- (4.65,-0.78);

% Terminais das placas
\draw[line width=1.8pt] (3.40,0.78) -- (3.40,1.62);
\draw[line width=1.8pt] (3.40,-0.78) -- (3.40,-1.62);

% Polaridade
\node[font=\Large] at (2.75,1.15) {$-$};
\node[font=\Large] at (2.75,-1.15) {$+$};

% Campo elétrico (setas para cima)
\foreach \x in {2.45,2.95,3.45,3.95,4.45}{
    \draw[myE,line width=1.5pt,-{Latex[length=3mm,width=2mm]}]
        (\x,-0.62) -- (\x,0.62);
}

% Pontos de B no seletor
\foreach \x in {2.55,3.25,3.95,4.35}{
    \foreach \y in {-0.42,0,0.42}{
        \fill[myB] (\x,\y) circle (2.3pt);
    }
}

% Rótulos vetoriais
\node[myE,above] at (4.95,0.25) {$\vec E$};
\node[myB,right] at (4.55,0.0) {$\vec B$};

% Velocidade
\draw[thick,-{Latex[length=3mm,width=2mm]}] (1.05,0.28) -- (1.95,0.28);
\node[above] at (1.5,0.28) {$\vec v$};

% ==================================================
% FENDA DE SAÍDA
% ==================================================
\draw[line width=3.6pt] (\xslit-0.35,0.32) -- (\xslit-0.35,0.82);
\draw[line width=3.6pt] (\xslit-0.35,-0.32) -- (\xslit-0.35,-0.82);

% ==================================================
% REGIÃO COM CAMPO MAGNÉTICO
% ==================================================
\foreach \x in {6.7,7.4,8.1,8.8,9.5}{
    \foreach \y in {0.45,-0.15,-0.85,-1.55,-2.25,-2.95,-3.65,-4.35, -5}{
        \fill[myB] (\x,\y) circle (2.5pt);
    }
}
\foreach \x/\y in {
    7.05/-0.50,7.75/-1.20,8.45/-0.50,9.15/-1.20,
    7.05/-1.90,7.75/-2.60,8.45/-1.90,9.15/-2.60,
    7.05/-3.30,7.75/-4.00,8.45/-3.30,9.15/-4.00
}{
    \fill[myB] (\x,\y) circle (2.5pt);
}

\node[myB] at (9.75,-0.25) {$\vec B$};

% ==================================================
% TRAJETÓRIA SEMICIRCULAR
% ==================================================
\draw[thick,dashed,-{Latex[length=3mm,width=2mm]}]
    (TopArc) arc[start angle=90,end angle=-90,radius=\r];

% ==================================================
% MEDIDA 2r
% ==================================================
\draw[thick,{Latex[length=2.8mm]}-{Latex[length=2.8mm]}]
    (5.65,-0.02) -- (5.65,\ybottom+0.02);

\draw[thin] (5.42,0) -- (5.90,0);
\draw[thin] (5.42,\ybottom) -- (5.90,\ybottom);

\node[fill=white,inner sep=1.5pt,font=\Large] at (5.65,-2.35) {$2r$};

% ==================================================
% PLACA DETECTORA
% ==================================================
\draw[line width=4pt] (\xslit-0.35,\ybottom+0.55) -- (\xslit-0.35,\ybottom-0.55);
\draw[line width=1.7pt] (\xslit-0.35,\ybottom) -- (\xslit+0.15,\ybottom);

\node[align=center,font=\Large] at (3.95,-4.55) {Detector\\plate};

% ==================================================
% TEXTOS DOS CAMPOS
% ==================================================
\node[anchor=west,text=myE,font=\bfseries\Large] at (-0.55,-2.05)
    {$E = 70\ \mathrm{kV/m}$};

\node[anchor=west,text=myB,font=\bfseries\Large] at (-0.55,-2.85)
    {$B = 50\ \mathrm{mT}$};

\end{tikzpicture}

\resposta{$v=\dfrac{E}{B}=\SI{1.4e6}{\meter\per\second}$,
\quad
$r=\dfrac{m_pv}{q_pB}\approx\SI{0.292}{\meter}$,
\quad
$2r\approx\SI{0.585}{\meter}$}

\resolucao{%
(a) \textbf{Velocidade dos prótons no seletor.}
Para que o próton atravesse a região de campos cruzados sem sofrer desvio,
a força elétrica e a força magnética devem possuir o mesmo módulo e sentidos
opostos:
\[
F_E=F_B.
\]
Como
\[
F_E=qE
\qquad\text{e}\qquad
F_B=qvB,
\]
tem-se
\[
qE=qvB.
\]
Cancelando a carga $q$:
\[
v=\frac{E}{B}.
\]
Substituindo
\[
E=\SI{70}{\kilo\volt\per\meter}
\qquad\text{e}\qquad
B=\SI{50}{\milli\tesla}=\SI{0.050}{\tesla},
\]
resulta
\[
v=
\frac{70\times10^{3}}{0.050}
=
\boxed{\SI{1.4e6}{\meter\per\second}}.
\]

\textbf{Interpretação:} somente os prótons com essa velocidade atravessam o
seletor em linha reta. Prótons mais rápidos ou mais lentos sofrem uma força
resultante e são desviados.\\

(b) \textbf{Raio da trajetória na região magnética.}
Após deixar o seletor, o próton entra em uma região onde atua apenas o campo
magnético. Como a força magnética é sempre perpendicular à velocidade, ela não
altera o módulo da velocidade, mas muda continuamente sua direção, atuando como
força centrípeta:
\[
F_B=F_c.
\]
Portanto,
\[
q_pvB=\frac{m_pv^2}{r}.
\]
Isolando o raio:
\[
r=\frac{m_pv}{q_pB}.
\]
Para
\[
m_p=\SI{1.67e-27}{\kilogram},
\qquad
q_p=\SI{1.60e-19}{\coulomb},
\]
tem-se
\[
r=
\frac{
(\pot{1.67}{-27})(\pot{1.4}{6})
}{
(\pot{1.60}{-19})(0.050)
}
\approx
\boxed{\SI{0.292}{\meter}}.
\]
Assim,
\[
r\approx\SI{29.2}{\centi\meter}.
\]

(c) \textbf{Distância até a placa detectora.}
A trajetória mostrada corresponde a uma semicircunferência. Dessa forma, a
distância vertical entre o ponto de entrada na região magnética e o ponto de
impacto na placa detectora corresponde ao diâmetro da trajetória:
\[
d=2r.
\]
Logo,
\[
d=2(0.292)
\approx
\boxed{\SI{0.585}{\meter}}
\]
ou, equivalentemente,
\[
\boxed{d\approx\SI{58.5}{\centi\meter}}.
\]

\textbf{Observação:} o seletor de velocidades utiliza a compensação entre as
forças elétrica e magnética para selecionar partículas com uma velocidade
específica. Na segunda região, a presença apenas do campo magnético permite
determinar a razão entre momento linear e carga da partícula a partir da
curvatura de sua trajetória.
}
\end{questao}



% ============================================================
% QUESTÃO 5 — PARTÍCULA ACELERADA E ANALISADOR MAGNÉTICO
% ============================================================

\begin{questao}[4]{}

Um elétron parte praticamente do repouso e é acelerado por uma diferença de
potencial de
\[
\Delta V=\SI{1.8}{\kilo\volt}.
\]
Após atravessar uma pequena abertura, ele penetra perpendicularmente em uma
região onde existe um campo magnético uniforme de intensidade
\[
B=\SI{25}{\milli\tesla},
\]
perpendicular ao plano da figura.

Considere
\[
m_e=\SI{9.11e-31}{\kilogram},
\qquad
|q_e|=\SI{1.60e-19}{\coulomb},
\]
e despreze efeitos relativísticos.

Determine:

\begin{itensq}
    \item a velocidade do elétron ao deixar a região de aceleração;
    \item o raio da trajetória circular no campo magnético;
    \item o período do movimento circular;
    \item o tempo necessário para o elétron percorrer um quarto de circunferência.
\end{itensq}

\begin{center}
\begin{tikzpicture}[
    >=Latex,
    line cap=round,
    line join=round,
    scale=0.95,
    font=\sffamily
]

% --------------------------------------------------
% Região de aceleração
% --------------------------------------------------

\node at (-4.6,0) {\Large $e^-$};

\draw[thick,-{Latex[length=3mm]}]
    (-4.25,0) -- (-3.65,0)
    node[midway,above] {$\vec v$};

% Placas aceleradoras
\draw[line width=2.5pt] (-3.3,-1.1) -- (-3.3,1.1);
\draw[line width=2.5pt] (-1.3,-1.1) -- (-1.3,1.1);

\node[font=\Large] at (-3.65,0.65) {$-$};
\node[font=\Large] at (-0.95,0.65) {$+$};

% Campo elétrico
\foreach \y in {-0.7,-0.35,0,0.35,0.7}{
    \draw[myE,thick,-{Latex[length=2.5mm]}]
        (-1.55,\y) -- (-3.05,\y);
}

\node[myE] at (-2.3,1.45) {$\Delta V=\SI{1.8}{\kilo\volt}$};

% Fenda
\draw[line width=3pt] (-0.65,0.28) -- (-0.65,1.0);
\draw[line width=3pt] (-0.65,-0.28) -- (-0.65,-1.0);

% Feixe após aceleração
\draw[thick,dashed,-{Latex[length=3mm]}]
    (-0.6,0) -- (0.6,0);

% --------------------------------------------------
% Região magnética
% --------------------------------------------------

\foreach \x in {0.8,1.5,2.2,2.9,3.6}{
    \foreach \y in {-2.1,-1.4,-0.7,0,0.7,1.4,2.1}{
        \draw[myB,thick] (\x-0.08,\y-0.08) -- (\x+0.08,\y+0.08);
        \draw[myB,thick] (\x-0.08,\y+0.08) -- (\x+0.08,\y-0.08);
    }
}

\node[myB] at (3.65,2.55)
    {$\vec B$ entrando no plano};

% Trajetória
\draw[very thick,dashed,-{Latex[length=3mm]}]
    (0.6,0)
    arc[start angle=90,end angle=0,radius=2.0];

% Centro
\fill (2.6,-2.0) circle (1.5pt);

\draw[thin,dashed]
    (2.6,-2.0) -- (0.6,0);

\node at (1.55,-1.05) {$r$};

\end{tikzpicture}
\end{center}

\resposta{$v\approx\SI{2.52e7}{\meter\per\second}$,
\quad
$r\approx\SI{5.72}{\milli\meter}$,
\quad
$T\approx\SI{1.43}{\nano\second}$,
\quad
$t_{1/4}\approx\SI{0.357}{\nano\second}$}

\resolucao{%
(a) \textbf{Velocidade após a aceleração.}
A energia potencial elétrica adquirida pelo elétron transforma-se em energia
cinética:
\[
|q_e|\Delta V=\frac{1}{2}m_ev^2.
\]
Portanto,
\[
v=\sqrt{\frac{2|q_e|\Delta V}{m_e}}.
\]
Substituindo os valores:
\[
v=
\sqrt{
\frac{
2(1.60\times10^{-19})(1.8\times10^3)
}{
9.11\times10^{-31}
}}
\approx
\boxed{\SI{2.52e7}{\meter\per\second}}.
\]

(b) \textbf{Raio da trajetória.}
Como $\vec v\perp\vec B$, a força magnética atua como força centrípeta:
\[
|q_e|vB=\frac{m_ev^2}{r}.
\]
Assim,
\[
r=\frac{m_ev}{|q_e|B}.
\]
Logo,
\[
r=
\frac{(9.11\times10^{-31})(2.52\times10^7)}
{(1.60\times10^{-19})(25\times10^{-3})}
\approx
\boxed{\SI{5.72}{\milli\meter}}.
\]

(c) \textbf{Período do movimento.}
\[
T=\frac{2\pi m_e}{|q_e|B}.
\]
Portanto,
\[
T\approx
\boxed{\SI{1.43e-9}{\second}}
=
\boxed{\SI{1.43}{\nano\second}}.
\]

(d) Para um quarto de circunferência:
\[
t_{1/4}=\frac{T}{4}.
\]
Assim,
\[
t_{1/4}\approx
\boxed{\SI{3.57e-10}{\second}}
=
\boxed{\SI{0.357}{\nano\second}}.
\]

\textbf{Interpretação:} a diferença de potencial determina a energia cinética
da partícula, enquanto o campo magnético modifica apenas a direção da
velocidade, produzindo uma trajetória circular.
}

\end{questao}


% ============================================================
% QUESTÃO 6 — ESPECTRÔMETRO DE MASSAS COM DOIS ISÓTOPOS
% ============================================================

\begin{questao}[4]{}

Um espectrômetro de massas é utilizado para separar dois isótopos de neônio,
$^{20}\mathrm{Ne}^{+}$ e $^{22}\mathrm{Ne}^{+}$. Inicialmente, os íons
atravessam um seletor de velocidades constituído por campos elétrico e
magnético perpendiculares, de intensidades
\[
E=\SI{30}{\kilo\volt\per\meter},
\qquad
B_s=\SI{20}{\milli\tesla}.
\]

Os íons selecionados entram em seguida em uma região contendo apenas um campo
magnético uniforme de intensidade
\[
B_a=\SI{0.35}{\tesla}.
\]

Considere
\[
u=\SI{1.66e-27}{\kilogram},
\qquad
q=+e=\SI{1.60e-19}{\coulomb}.
\]

Determine:

\begin{itensq}
    \item a velocidade dos íons que atravessam o seletor sem desvio;
    \item os raios das trajetórias dos dois isótopos;
    \item a distância entre os pontos de impacto na placa detectora;
    \item qual isótopo atinge a região mais distante do ponto de entrada.
\end{itensq}

\begin{center}
\begin{tikzpicture}[
    >=Latex,
    line cap=round,
    line join=round,
    scale=0.88,
    font=\sffamily
]

% Fonte de íons
\draw[thick] (-4.7,0) circle (0.3);
\node at (-4.7,0) {$+$};

\draw[thick,-{Latex[length=3mm]}]
    (-4.35,0) -- (-3.65,0);

% Seletor
\draw[line width=2pt] (-3.2,0.8) -- (-1.0,0.8);
\draw[line width=2pt] (-3.2,-0.8) -- (-1.0,-0.8);

\node at (-2.85,1.12) {$-$};
\node at (-2.85,-1.12) {$+$};

\foreach \x in {-2.9,-2.4,-1.9,-1.4}{
    \draw[myE,thick,-{Latex[length=2.5mm]}]
        (\x,-0.6) -- (\x,0.6);
}

\foreach \x in {-2.7,-2.0,-1.3}{
    \foreach \y in {-0.4,0,0.4}{
        \fill[myB] (\x,\y) circle (2pt);
    }
}

\node at (-2.1,1.55) {Seletor de velocidades};

% Fenda
\draw[line width=3pt] (-0.45,0.3) -- (-0.45,1);
\draw[line width=3pt] (-0.45,-0.3) -- (-0.45,-1);

% Região analisadora
\foreach \x in {0.2,0.9,1.6,2.3,3.0,3.7}{
    \foreach \y in {-4.1,-3.4,-2.7,-2.0,-1.3,-0.6,0.1}{
        \fill[myB] (\x,\y) circle (2.3pt);
    }
}

\node[myB] at (3.5,0.7) {$\vec B_a$};

% Trajetória 20
\draw[very thick,dashed]
    (-0.35,0)
    arc[start angle=90,end angle=-90,radius=1.8];

\node at (1.25,-1.1) {$^{20}\mathrm{Ne}^{+}$};

% Trajetória 22
\draw[very thick,dash dot]
    (-0.35,0)
    arc[start angle=90,end angle=-90,radius=2.15];

\node at (2.65,-2.2) {$^{22}\mathrm{Ne}^{+}$};

% Detector
\draw[line width=4pt] (-0.45,-3.2) -- (-0.45,-4.8);

\node[rotate=90] at (-0.85,-4.0) {Detector};

% Marcas de impacto
\fill (-0.35,-3.6) circle (2.2pt);
\fill (-0.35,-4.3) circle (2.2pt);

\draw[<->,thick]
    (-1.25,-3.6) -- (-1.25,-4.3)
    node[midway,left] {$\Delta d$};

\end{tikzpicture}
\end{center}

\resposta{$v=\SI{1.50e6}{\meter\per\second}$,
\quad
$r_{20}\approx\SI{0.889}{\meter}$,
\quad
$r_{22}\approx\SI{0.977}{\meter}$,
\quad
$\Delta d\approx\SI{0.178}{\meter}$}

\resolucao{%
(a) \textbf{Velocidade selecionada.}
No seletor:
\[
qE=qvB_s.
\]
Logo,
\[
v=\frac{E}{B_s}.
\]
Assim,
\[
v=
\frac{30\times10^3}{20\times10^{-3}}
=
\boxed{\SI{1.50e6}{\meter\per\second}}.
\]

(b) \textbf{Raios das trajetórias.}
Para um íon de massa $m$:
\[
r=\frac{mv}{qB_a}.
\]

Para o $^{20}\mathrm{Ne}^{+}$:
\[
m_{20}=20u=20(1.66\times10^{-27}),
\]
resultando em
\[
r_{20}\approx
\boxed{\SI{0.889}{\meter}}.
\]

Para o $^{22}\mathrm{Ne}^{+}$:
\[
m_{22}=22u,
\]
e
\[
r_{22}\approx
\boxed{\SI{0.977}{\meter}}.
\]

(c) Como cada partícula descreve uma semicircunferência, a posição de impacto
é determinada pelo diâmetro:
\[
d=2r.
\]
A separação entre os impactos é
\[
\Delta d=2(r_{22}-r_{20}).
\]
Logo,
\[
\Delta d
=
2(0.977-0.889)
\approx
\boxed{\SI{0.178}{\meter}}.
\]

(d) Como
\[
r=\frac{mv}{qB},
\]
mantidos $v$, $q$ e $B$, o raio é diretamente proporcional à massa. Portanto,
o isótopo
\[
\boxed{^{22}\mathrm{Ne}^{+}}
\]
atinge a região mais distante.

\textbf{Interpretação:} o espectrômetro converte diferenças de massa em
diferenças geométricas mensuráveis na posição de impacto das partículas.
}

\end{questao}


% ============================================================
% QUESTÃO 7 — TRAJETÓRIA HELICOIDAL
% ============================================================

\begin{questao}[4]{}

Uma partícula alfa entra em uma região onde existe um campo magnético uniforme
de intensidade
\[
B=\SI{0.18}{\tesla}.
\]
A velocidade inicial da partícula possui módulo
\[
v=\SI{7.0e5}{\meter\per\second}
\]
e forma um ângulo de
\[
\theta=40^\circ
\]
com a direção do campo magnético.

Considere
\[
m_\alpha=\SI{6.64e-27}{\kilogram},
\qquad
q_\alpha=\SI{3.20e-19}{\coulomb}.
\]

Determine:

\begin{itensq}
    \item as componentes $v_{\parallel}$ e $v_{\perp}$ da velocidade;
    \item o raio da trajetória helicoidal;
    \item o período de uma volta;
    \item o passo $p$ da hélice;
    \item o que aconteceria com a trajetória se $\theta=90^\circ$.
\end{itensq}

\begin{center}
\begin{tikzpicture}[
    >=Latex,
    line cap=round,
    line join=round,
    scale=1.0,
    font=\sffamily
]

% Eixo do campo
\draw[myB,line width=1.3pt,-{Latex[length=3mm]}]
    (-4,0) -- (4.7,0)
    node[right] {$\vec B$};

% Linhas paralelas de campo
\foreach \y in {-1.4,-0.7,0.7,1.4}{
    \draw[myB!70,thin,-{Latex[length=2mm]}]
        (-3.8,\y) -- (4.3,\y);
}

% Hélice estilizada
\draw[very thick]
plot[
    domain=0:720,
    samples=150,
    smooth
]
({-3.3+0.010*\x},{0.60*sin(\x)});

% Partícula
\fill (-3.3,0) circle (2.8pt);
\node[above left] at (-3.3,0) {$\alpha$};

% Vetor velocidade
\draw[thick,-{Latex[length=3mm]}]
    (-3.3,0) -- (-2.35,0.8)
    node[above] {$\vec v$};

% Componente paralela
\draw[dashed,thick,-{Latex[length=2.5mm]}]
    (-3.3,-0.15) -- (-2.0,-0.15)
    node[midway,below] {$v_{\parallel}$};

% Componente perpendicular
\draw[dashed,thick,-{Latex[length=2.5mm]}]
    (-3.3,0) -- (-3.3,0.85)
    node[left] {$v_{\perp}$};

% Ângulo
\draw (-2.75,0)
    arc[start angle=0,end angle=40,radius=0.55];

\node at (-2.55,0.25) {$\theta$};

% Passo
\draw[<->,thick]
    (0.3,-1.15) -- (3.9,-1.15)
    node[midway,fill=white,inner sep=2pt] {$p$};

\node at (0,-2.0)
    {Trajetória helicoidal};

\end{tikzpicture}
\end{center}

\resposta{$v_{\parallel}\approx\SI{5.36e5}{\meter\per\second}$,
\quad
$v_{\perp}\approx\SI{4.50e5}{\meter\per\second}$,
\quad
$r\approx\SI{5.18}{\centi\meter}$,
\quad
$T\approx\SI{0.724}{\micro\second}$,
\quad
$p\approx\SI{0.388}{\meter}$}

\resolucao{%
(a) \textbf{Decomposição da velocidade.}
A componente paralela ao campo é
\[
v_{\parallel}=v\cos\theta,
\]
e a componente perpendicular:
\[
v_{\perp}=v\sin\theta.
\]

Logo,
\[
v_{\parallel}
=
(7.0\times10^5)\cos40^\circ
\approx
\boxed{\SI{5.36e5}{\meter\per\second}},
\]
e
\[
v_{\perp}
=
(7.0\times10^5)\sin40^\circ
\approx
\boxed{\SI{4.50e5}{\meter\per\second}}.
\]

(b) \textbf{Raio da hélice.}
Somente $v_{\perp}$ produz movimento circular:
\[
q_\alpha v_{\perp}B
=
\frac{m_\alpha v_{\perp}^2}{r}.
\]
Portanto,
\[
r=
\frac{m_\alpha v_{\perp}}{q_\alpha B}.
\]
Substituindo:
\[
r\approx
\boxed{\SI{5.18}{\centi\meter}}.
\]

(c) \textbf{Período.}
\[
T=
\frac{2\pi m_\alpha}{q_\alpha B}.
\]
Assim,
\[
T\approx
\boxed{\SI{7.24e-7}{\second}}
=
\boxed{\SI{0.724}{\micro\second}}.
\]

(d) \textbf{Passo da hélice.}
Durante uma volta, o deslocamento paralelo ao campo é
\[
p=v_{\parallel}T.
\]
Logo,
\[
p=
(5.36\times10^5)(7.24\times10^{-7})
\approx
\boxed{\SI{0.388}{\meter}}.
\]

(e) Para $\theta=90^\circ$:
\[
v_{\parallel}=0.
\]
Consequentemente, desaparece o avanço longitudinal da hélice e a trajetória
passa a ser
\[
\boxed{\text{circular}}.
\]

\textbf{Interpretação:} o movimento helicoidal resulta da superposição de um
movimento retilíneo uniforme paralelo ao campo com um movimento circular
uniforme perpendicular ao campo.
}

\end{questao}


% ============================================================
% QUESTÃO 8 — CÍCLOTRON
% ============================================================

\begin{questao}[4]{}

Um cíclotron é utilizado para acelerar prótons em uma região de campo magnético
uniforme de intensidade
\[
B=\SI{0.80}{\tesla}.
\]
Os prótons descrevem órbitas de raio progressivamente maior até atingirem o
raio máximo
\[
R=\SI{25}{\centi\meter}.
\]

Considere
\[
m_p=\SI{1.67e-27}{\kilogram},
\qquad
q_p=\SI{1.60e-19}{\coulomb},
\]
e despreze correções relativísticas.

Determine:

\begin{itensq}
    \item a frequência que deve ser aplicada entre os eletrodos do cíclotron;
    \item o período do movimento dos prótons;
    \item a velocidade do próton ao atingir o raio máximo;
    \item a energia cinética máxima do próton, em joules e em elétron-volts.
\end{itensq}

\begin{center}
\begin{tikzpicture}[
    >=Latex,
    line cap=round,
    line join=round,
    scale=0.95,
    font=\sffamily
]

% Região magnética
\foreach \x in {-3,-2.25,-1.5,-0.75,0.75,1.5,2.25,3}{
    \foreach \y in {-2.2,-1.5,-0.8,0,0.8,1.5,2.2}{
        \fill[myB!75] (\x,\y) circle (1.8pt);
    }
}

% Dees
\draw[line width=2pt]
    (0.18,2.55)
    arc[start angle=90,end angle=270,radius=2.55];

\draw[line width=2pt]
    (-0.18,-2.55)
    arc[start angle=-90,end angle=90,radius=2.55];

% Fenda central
\draw[line width=2pt] (-0.18,-2.55) -- (-0.18,2.55);
\draw[line width=2pt] (0.18,-2.55) -- (0.18,2.55);

\node at (-1.45,0) {\Large D};
\node at (1.45,0) {\Large D};

% Campo elétrico alternado na fenda
\draw[myE,very thick,<->]
    (-0.13,0.75) -- (0.13,0.75);

\node[myE,right] at (0.25,0.75) {$\vec E(t)$};

% Espiral
\draw[very thick,-{Latex[length=3mm]}]
plot[
    domain=0:900,
    samples=250,
    smooth
]
({0.0027*\x*cos(\x)},
 {0.0027*\x*sin(\x)});

% Centro
\fill (0,0) circle (2.5pt);

% Raio final
\draw[dashed,thick]
    (0,0) -- (2.43,0);

\node[below] at (1.25,0) {$R$};

% Campo
\node[myB] at (3.7,2.0)
    {$\vec B$};

% Saída
\draw[thick,-{Latex[length=3mm]}]
    (2.45,-0.35) -- (3.65,-0.7);

\node[right] at (3.65,-0.7)
    {feixe};

\end{tikzpicture}
\end{center}

\resposta{$f\approx\SI{12.2}{\mega\hertz}$,
\quad
$T\approx\SI{82.0}{\nano\second}$,
\quad
$v_{\max}\approx\SI{1.92e7}{\meter\per\second}$,
\quad
$K_{\max}\approx\SI{3.07e-13}{\joule}
\approx\SI{1.92}{\mega\electronvolt}$}

\resolucao{%
(a) \textbf{Frequência ciclotrônica.}
A força magnética fornece a força centrípeta:
\[
q_pvB=\frac{m_pv^2}{r}.
\]
Como
\[
v=\omega r,
\]
obtém-se
\[
\omega=\frac{q_pB}{m_p}.
\]
Portanto,
\[
f=\frac{\omega}{2\pi}
=
\frac{q_pB}{2\pi m_p}.
\]
Substituindo:
\[
f=
\frac{
(1.60\times10^{-19})(0.80)
}{
2\pi(1.67\times10^{-27})
}
\approx
\boxed{\SI{12.2}{\mega\hertz}}.
\]

(b) \textbf{Período.}
\[
T=\frac{1}{f}.
\]
Logo,
\[
T\approx
\boxed{\SI{82.0}{\nano\second}}.
\]

(c) \textbf{Velocidade máxima.}
Da condição de movimento circular:
\[
r=\frac{m_pv}{q_pB}.
\]
Para $r=R$:
\[
v_{\max}=
\frac{q_pBR}{m_p}.
\]
Assim,
\[
v_{\max}
=
\frac{
(1.60\times10^{-19})(0.80)(0.25)
}{
1.67\times10^{-27}
}
\approx
\boxed{\SI{1.92e7}{\meter\per\second}}.
\]

(d) \textbf{Energia cinética máxima.}
\[
K_{\max}
=
\frac{1}{2}m_pv_{\max}^2.
\]
Substituindo:
\[
K_{\max}
\approx
\boxed{\SI{3.07e-13}{\joule}}.
\]

Como
\[
\SI{1}{\electronvolt}
=
\SI{1.60e-19}{\joule},
\]
tem-se
\[
K_{\max}
=
\frac{3.07\times10^{-13}}
{1.60\times10^{-19}}
\approx
1.92\times10^6\,\mathrm{eV}.
\]
Portanto,
\[
\boxed{K_{\max}\approx\SI{1.92}{\mega\electronvolt}}.
\]

\textbf{Interpretação:} a frequência ciclotrônica permanece aproximadamente
constante enquanto os efeitos relativísticos forem desprezíveis. A cada
passagem pela região entre os eletrodos, o campo elétrico fornece energia à
partícula, aumentando sua velocidade e, consequentemente, o raio da órbita.
}

\end{questao}

% END BANCO COMPLEMENTAR

\gabarito[12.5 Lei de Lenz e indução eletromagnética]
      % Eletromagnetismo


% Cap. 12 ------------------------------------------------------------
% REMOVIDO: "Medicoes Eletricas e Propagacao de Incertezas".
% O capitulo e o banco complementar correspondente foram retirados da obra.
% O tratamento de incertezas NAO se perdeu: ele permanece distribuido nos
% blocos de Nivel 4 dos capitulos e bancos que o utilizam de fato --
% tolerancia em associacoes (serie e paralelo), tolerancia de razao em
% divisores e em Thevenin, propagacao em RC, incerteza na verificacao
% experimental da LKC e no metodo das duas cargas.

% Nota: o conteudo avancado (redes infinitas, superno, supermalha,
% solucao matricial e sintese) foi DISTRIBUIDO nas secoes correspondentes,
% como blocos de "Nivel 4 - Desafio".

\end{document}
